\PassOptionsToPackage{unicode}{hyperref}
\PassOptionsToPackage{hyphens}{url}
\PassOptionsToPackage{dvipsnames,svgnames,x11names}{xcolor}

\documentclass[
  11pt,twoside,letterpaper]{book}

\usepackage{xcolor}
\usepackage{amsmath,amssymb}
\usepackage{iftex}
\usepackage{etoolbox}
\usepackage{pdflscape} 
% Automatically rotate wide tables onto landscape pages
% This wraps every table and longtable environment in a landscape page.
% If you prefer conditional rotation only for very-wide tables, tell me
% and I can implement a measurement-based approach instead.
\AtBeginEnvironment{table}{\begin{landscape}}
\AtEndEnvironment{table}{\end{landscape}}
\usepackage{calc}
\usepackage{longtable}
\usepackage{booktabs}
\usepackage{array}
\usepackage{url}
\usepackage{hyperref}
\usepackage{listings}
\usepackage{upquote}
\usepackage{showframe}

% ============================================================
% Admonition Boxes (NOTE, TIP, IMPORTANT, WARNING, CAUTION)
% Used by the admonitions.lua Pandoc filter to render
% GitHub/mdbook-style >  [!NOTE] block-quotes in LaTeX.
% ============================================================
\usepackage[most]{tcolorbox}

% Shared base style — override colors per environment below
\tcbset{
  admonitionbase/.style={
    enhanced,
    breakable,
    left=6pt, right=6pt, top=4pt, bottom=4pt,
    fonttitle=\bfseries\small,
    attach boxed title to top left={yshift=-2mm, xshift=4mm},
    boxed title style={
      sharp corners,
      interior style={fill=white},
      frame hidden,
      top=1pt, bottom=1pt, left=2pt, right=2pt,
    },
    sharp corners,
    boxrule=0.6pt,
    parbox=false,
  }
}

% ---- NOTE (blue) ----
\definecolor{noteframe}{HTML}{0055AA}
\definecolor{notebg}{HTML}{EEF4FF}
\newtcolorbox{noteenv}{
  admonitionbase,
  colframe=noteframe,
  colback=notebg,
  title={\textcolor{noteframe}{ℹ\ NOTE}},
  boxed title style={colback=notebg},
}

% ---- TIP (green) ----
\definecolor{tipframe}{HTML}{2E7D32}
\definecolor{tipbg}{HTML}{F0FAF0}
\newtcolorbox{tipenv}{
  admonitionbase,
  colframe=tipframe,
  colback=tipbg,
  title={\textcolor{tipframe}{✦\ TIP}},
  boxed title style={colback=tipbg},
}

% ---- IMPORTANT (purple) ----
\definecolor{importantframe}{HTML}{6A0DAD}
\definecolor{importantbg}{HTML}{F5EEFF}
\newtcolorbox{importantenv}{
  admonitionbase,
  colframe=importantframe,
  colback=importantbg,
  title={\textcolor{importantframe}{★\ IMPORTANT}},
  boxed title style={colback=importantbg},
}

% ---- WARNING (amber) ----
\definecolor{warningframe}{HTML}{B45309}
\definecolor{warningbg}{HTML}{FFFBEB}
\newtcolorbox{warningenv}{
  admonitionbase,
  colframe=warningframe,
  colback=warningbg,
  title={\textcolor{warningframe}{⚠\ WARNING}},
  boxed title style={colback=warningbg},
}

% ---- CAUTION (red) ----
\definecolor{cautionframe}{HTML}{B22222}
\definecolor{cautionbg}{HTML}{FFF0F0}
\newtcolorbox{cautionenv}{
  admonitionbase,
  colframe=cautionframe,
  colback=cautionbg,
  title={\textcolor{cautionframe}{✗\ CAUTION}},
  boxed title style={colback=cautionbg},
}
% ============================================================

% Hyperlink Color Configuration
\definecolor{linkblue}{HTML}{0055AA} 
\definecolor{exturl}{HTML}{008000}   

\hypersetup{
    colorlinks=true,
    linkcolor=linkblue,
    citecolor=linkblue,
    filecolor=linkblue,
    urlcolor=exturl,
    pdfnewwindow=true,
    breaklinks=true
}


\ifLuaTeX
  \usepackage{fontspec}
  \usepackage[sfdefault]{atkinson}
  \setmonofont{AtkinsonHyperlegibleMono}[
    Extension=.otf,
    UprightFont=*-Regular,
    BoldFont=*-Bold,
    ItalicFont=*-Italic,
    BoldItalicFont=*-BoldItalic
  ]
  \directlua{
    luaotfload.add_fallback("symbolsfallback", {
      "DejaVuSans:mode=harf;",
      "NotoSans:mode=harf;"
    })
  }
  \setmainfont{AtkinsonHyperlegibleNext}[
    Extension=.otf,
    UprightFont=*-Regular,
    BoldFont=*-Bold,
    ItalicFont=*-Italic,
    BoldItalicFont=*-BoldItalic,
    BoldFeatures={RawFeature={fallback=symbolsfallback}},
    UprightFeatures={RawFeature={fallback=symbolsfallback}},
    ItalicFeatures={RawFeature={fallback=symbolsfallback}},
    BoldItalicFeatures={RawFeature={fallback=symbolsfallback}}
  ]
\fi

\newfontfamily\listingfont{AtkinsonHyperlegibleMono}[
  Extension=.otf,
  UprightFont=*-Regular,
  BoldFont=*-Bold,
  ItalicFont=*-Italic,
  BoldItalicFont=*-BoldItalic
]

% Unicode replacements for glyphs missing in Atkinson fonts
% Use a dedicated symbol font (DejaVu Sans / Noto Sans) for these glyphs
\usepackage{newunicodechar}
\usepackage{textcomp}
\usepackage{pifont}
% Symbol fallback font (uses system DejaVu Sans if available)
\newfontfamily\symbolfont{DejaVu Sans}[Scale=MatchLowercase]
% Fallback mappings: render specific unicode glyphs using \symbolfont
\newunicodechar{↓}{{\symbolfont↓}}
\newunicodechar{←}{{\symbolfont←}}
\newunicodechar{─}{{\symbolfont─}}
\newunicodechar{⚠}{{\symbolfont⚠}}
\newunicodechar{✗}{{\symbolfont✗}}
\newunicodechar{□}{{\symbolfont□}}
\newunicodechar{✓}{\checkmark}

% Footnotes to Endnotes Configuration
\usepackage{endnotes}
\let\footnote=\endnote
\renewcommand{\notesname}{Notes}

% Automatically print endnotes at the end of each chapter
\preto\chapter{%
  \ifnum\value{endnote}>0
    \begingroup
    \parindent 0pt
    \parskip 2ex
    \def\enotesize{\normalsize}
    \theendnotes
    \endgroup
  \fi
}

% Code Highlighting Colors (Alabaster)
\definecolor{alabaster-bg}{HTML}{FAFAFA}
\definecolor{alabaster-keyword}{HTML}{C92C2C}
\definecolor{alabaster-string}{HTML}{008000}
\definecolor{alabaster-comment}{HTML}{0078D7}
\definecolor{alabaster-number}{HTML}{333333}
\definecolor{alabaster-identifier}{HTML}{000000}
\definecolor{BrickRed}{HTML}{B22222}

\newlength{\AlabasterTabWidth}
\settowidth{\AlabasterTabWidth}{{\listingfont\footnotesize\ \ \ \ }}
\newcommand*\AlabasterTab{\makebox[\AlabasterTabWidth][c]{\color{BrickRed}\scriptsize\ensuremath{\rightarrow}}}

\lstset{
  numbers=left,
  numberstyle=\tiny\color{alabaster-number},
  breaklines=true,
  prebreak=\space\mbox{\textcolor{green}{\ensuremath{\rightarrow}}},
  postbreak=\mbox{\textcolor{green}{\ensuremath{\hookrightarrow}\space}},
  breakatwhitespace=false,
  tabsize=4,
  showtabs=true,
  tab=\AlabasterTab,
  showstringspaces=false,
  showspaces=false,
  firstnumber=1,
  numberfirstline=true,
  stepnumber=2,
  basicstyle=\footnotesize\listingfont\color{alabaster-identifier},
  columns=fullflexible,
  extendedchars=true,
  inputencoding=utf8,
  aboveskip=0pt,
  belowskip=0pt
}

\lstdefinestyle{Alabaster}{
  frame=lines,
  backgroundcolor=\color{alabaster-bg},
  basicstyle=\footnotesize\listingfont\color{alabaster-identifier},
  keywordstyle=\bfseries\color{alabaster-keyword},
  identifierstyle=\color{alabaster-identifier},
  stringstyle=\color{alabaster-string},
  commentstyle=\color{alabaster-comment}\itshape,
  showtabs=true,
  belowcaptionskip=0.25em,
  abovecaptionskip=0.25em,
  aboveskip=1em,
  belowskip=0.25em,
  captionpos=b,
  rulesep=0.5em,
  framesep=3pt,
  framerule=0.4pt,
  xleftmargin=3.4pt,
  xrightmargin=3.4pt
}

\lstdefinelanguage{openscad}{
  sensitive=true,
  morekeywords={module,function,if,else,for,intersection,difference,union,translate,rotate,scale,mirror,resize,render,include,use,assign,let,return,echo},
  morecomment=[l]{//},
  morecomment=[s]{/*}{*/},
  morestring=[b]{"},
  morestring=[b]{'},
  keywordstyle=\color{alabaster-keyword}\bfseries,
  commentstyle=\color{alabaster-comment}\itshape,
  stringstyle=\color{alabaster-string},
}

\lstdefinelanguage{powershell}{
  sensitive=true,
  morekeywords={function,param,begin,process,end,if,else,foreach,while,return,break,continue,try,catch,finally,throw,exit},
  morecomment=[l]{\#},
  morestring=[b]{"},
  morestring=[b]{'},
  keywordstyle=\color{alabaster-keyword}\bfseries,
  commentstyle=\color{alabaster-comment}\itshape,
  stringstyle=\color{alabaster-string},
}

\lstdefinelanguage{cmd}{
  sensitive=false,
  morekeywords={if,else,for,call,exit,echo,set,cd,copy,del,move,ren,start,dir,type,shift,md,rd,cls},
  morecomment=[l]{REM},
  morecomment=[l]{::},
  morestring=[b]{"},
  alsoletter={@,\%},
  keywordstyle=\color{alabaster-keyword}\bfseries,
  commentstyle=\color{alabaster-comment}\itshape,
  stringstyle=\color{alabaster-string},
}

\lstdefinelanguage{json}{
  sensitive=true,
  morestring=[b]{"},
  morecomment=[l]{//},
  morecomment=[s]{/*}{*/},
  literate=
    *{0}{{{\color{alabaster-keyword}0}}}{1}
     {1}{{{\color{alabaster-keyword}1}}}{1}
     {2}{{{\color{alabaster-keyword}2}}}{1}
     {3}{{{\color{alabaster-keyword}3}}}{1}
     {4}{{{\color{alabaster-keyword}4}}}{1}
     {5}{{{\color{alabaster-keyword}5}}}{1}
     {6}{{{\color{alabaster-keyword}6}}}{1}
     {7}{{{\color{alabaster-keyword}7}}}{1}
     {8}{{{\color{alabaster-keyword}8}}}{1}
     {9}{{{\color{alabaster-keyword}9}}}{1}
     {:}{{{\color{alabaster-comment}{:}}}}{1}
     {,}{{{\color{alabaster-comment}{,}}}}{1}
     {\{}{{{\color{alabaster-identifier}{\{}}}}{1}
     {\}}{{{\color{alabaster-identifier}{\}}}}}{1}
     {[}{{{\color{alabaster-identifier}{[}}}}{1}
     {]}{{{\color{alabaster-identifier}{]}}}}{1},
  keywordstyle=\color{alabaster-keyword}\bfseries,
  stringstyle=\color{alabaster-string},
  commentstyle=\color{alabaster-comment}\itshape,
}

\lstset{style=Alabaster}
\lstalias{openscad}{openscad}
\lstalias{batch}{cmd}
\lstalias{powershell}{powershell}
\lstalias{json}{json}

\setcounter{secnumdepth}{-\maxdimen}

\usepackage{iftex}
\ifPDFTeX
  \usepackage[T1]{fontenc}
  \usepackage[utf8]{inputenc}
  \usepackage{textcomp} % provide euro and other symbols
\else % if luatex or xetex
  \usepackage{unicode-math} % this also loads fontspec
  \defaultfontfeatures{Scale=MatchLowercase}
  \defaultfontfeatures[\rmfamily]{Ligatures=TeX,Scale=1}
\fi
\usepackage{lmodern}
\ifPDFTeX\else
  % xetex/luatex font selection
\fi
% Use upquote if available, for straight quotes in verbatim environments
\IfFileExists{upquote.sty}{\usepackage{upquote}}{}
\IfFileExists{microtype.sty}{% use microtype if available
  \usepackage[]{microtype}
  \UseMicrotypeSet[protrusion]{basicmath} % disable protrusion for tt fonts
}{}
\makeatletter
\@ifundefined{KOMAClassName}{% if non-KOMA class
  \IfFileExists{parskip.sty}{%
    \usepackage{parskip}
  }{% else
    \setlength{\parindent}{0pt}
    \setlength{\parskip}{6pt plus 2pt minus 1pt}}
}{% if KOMA class
  \KOMAoptions{parskip=half}}
\makeatother
\usepackage{longtable,booktabs,array}
\newcounter{none} % for unnumbered tables
\usepackage{calc} % for calculating minipage widths
% Correct order of tables after \paragraph or \subparagraph
\usepackage{etoolbox}
\makeatletter
\patchcmd\longtable{\par}{\if@noskipsec\mbox{}\fi\par}{}{}
\makeatother
% Allow footnotes in longtable head/foot
\IfFileExists{footnotehyper.sty}{\usepackage{footnotehyper}}{\usepackage{footnote}}
\makesavenoteenv{longtable}
\ifLuaTeX
\usepackage[bidi=basic,shorthands=off]{babel}
\else
\usepackage[bidi=default,shorthands=off]{babel}
\fi
\ifLuaTeX
  \usepackage{selnolig} % disable illegal ligatures
\fi
\setlength{\emergencystretch}{3em} % prevent overfull lines
\providecommand{\tightlist}{%
  \setlength{\itemsep}{0pt}\setlength{\parskip}{0pt}}

\IfFileExists{fvextra.sty}{% use fvextra if available to break long lines in code blocks
  \usepackage{fvextra}
  \fvset{breaklines}
}{}

\usepackage{bookmark}
\IfFileExists{xurl.sty}{\usepackage{xurl}}{} % add URL line breaks if available
\urlstyle{same}
\hypersetup{
  pdftitle={3D Printing Course for Blind and Visually Impaired Learners},
  pdfauthor={Michael Ryan Hunsaker, M.Ed., Ph.D.},
  pdflang={en},
  colorlinks=true,
  linkcolor={Maroon},
  filecolor={Maroon},
  citecolor={Blue},
  urlcolor={Blue},
  pdfcreator={LaTeX via pandoc}}


\title{3D Printing Course for Blind and Visually Impaired Learners}
\author{Michael Ryan Hunsaker, M.Ed., Ph.D.}
\date{2026-02-24}

% Configure longtable for better PDF output
% Don't automatically wrap in landscape - causes nesting errors with document environment
% Tables will resize columns as needed with LuaTeX's text width handling
\setlength\LTpost{0pt}
\setlength\LTpre{0pt}

% Define \tightlist macro for Pandoc-generated lists
% Provides proper spacing for bullet and numbered lists
\providecommand{\tightlist}{%
  \setlength{\itemsep}{0pt}%
  \setlength{\parskip}{0pt}%
}

\begin{document}

% --- HARDCODED COVER SECTION ---
\begin{titlepage}
    \centering
    \vspace*{2cm}
    {\Huge\bfseries 3D Printing Course for Blind and Visually Impaired Learners\par}
    \vspace{1.5cm}
    {\Large Michael Ryan Hunsaker, M.Ed., Ph.D.}
    \vfill
    {\large Updated: \today}
\end{titlepage}

% Copyright Page (Verso)
\newpage
\thispagestyle{empty}
\vspace*{\fill}
{\small \textcopyright\ 2026 Michael Ryan Hunsaker, M.Ed., Ph.D. 
All rights reserved.}

\vspace{1em}
{\small Creative Commons CC BY-NC-ND license. To view a copy of this license, visit \url{https://creativecommons.org/licenses/by-nc-nd/4.0/}\par}

{\small Book URL: \url{https://mrhunsaker.github.io/VI_3DMake_OpenSCAD_Curriculum/}}
\newpage
% --- END COVER SECTION ---

\frontmatter


\cleardoublepage 
{
\hypersetup{linkcolor=linkblue}
\setcounter{tocdepth}{2}
\tableofcontents
}



\mainmatter

\chapter{Introduction}\label{docs__pandoc__latex__src__readme.md__readme}

This course teaches 3D design and digital fabrication using a fully accessible, command-line-driven toolchain centered on OpenSCAD (text-based CAD), 3DMake (non-visual build automation), and accessible editors (VS Code, Notepad++, command-line editors) with screen reader support.

The curriculum is explicitly designed for blind and visually impaired learners who use screen readers (NVDA, JAWS, Orca, VoiceOver). It eliminates GUI navigation and visual feedback in favor of keyboard-driven, text-based workflows that screen readers can fully access. Accessibility is not an add-on. It is the foundation of every tool, workflow, and lesson in this curriculum.

\section*{Curriculum Structure}\label{docs__pandoc__latex__src__readme.md__curriculum-structure}

\subsection*{Screen Reader Setup \& Accessibility Fundamentals}\label{docs__pandoc__latex__src__readme.md__screen-reader-setup--accessibility-fundamentals}

Before choosing a command-line pathway, students learn to optimize their screen reader setup for terminal work and understand accessibility options.

What\textquotesingle s Included:

\begin{itemize}
\tightlist
\item
  \hyperref[docs__pandoc__latex__src__setup_accessibility_complete.md__accessibility-guide]{Screen Reader Accessibility Guide} - Comprehensive NVDA \& JAWS reference with terminal tips for PowerShell, CMD, and Git Bash
\item
  \hyperref[docs__pandoc__latex__src__setup_accessibility_complete.md__screen-reader-choice]{Screen Reader Choice: Windows CLI} - Help choosing between NVDA and JAWS
\item
  \hyperref[docs__pandoc__latex__src__setup_accessibility_complete.md__braille-displays]{Braille Display \& Terminal Mode} - Advanced accessibility setup with refreshable braille displays
\item
  \hyperref[docs__pandoc__latex__src__setup_accessibility_complete.md__editor-setup]{Editor Selection and Setup} - Choosing and configuring an accessible text editor
\end{itemize}

Why start here? Screen reader optimization is the foundation for all command-line work. Students learn keyboard shortcuts, navigation techniques, and terminal-specific accessibility before diving into any CLI pathway.

\subsection*{Command-Line Foundation (Choose Your Path)}\label{docs__pandoc__latex__src__readme.md__command-line-foundation-choose-your-path}

Students master terminal/command-line fundamentals before learning 3D design. You choose between three equivalent pathways based on your operating system and learning preferences.

\subsubsection*{Compare All Three Pathways}\label{docs__pandoc__latex__src__readme.md__compare-all-three-pathways}

\hyperref[docs__pandoc__latex__src__command_line_interface_selection__command_line_interface_selection.md__command-line-fundamentals---choose-your-path]{Command Line Interface Selection Guide}

This guide provides a comprehensive comparison including:

\begin{itemize}
\tightlist
\item
  Feature matrix comparing all three shells
\item
  Command comparison table (navigation, file operations, scripting)
\item
  Learner profiles to help you choose
\item
  Goal-based recommendations
\item
  FAQ addressing common questions
\end{itemize}

\subsubsection*{PowerShell Foundation (Recommended for Windows Users)}\label{docs__pandoc__latex__src__readme.md__powershell-foundation-recommended-for-windows-users}

Where to Start: \href{docs\%5Cpandoc\%5Clatex\%5Csrc\%5CPS_Complete_Curriculum.md\#part-1-powershell}{PowerShell Curriculum Overview}

Lessons Included:

\begin{itemize}
\tightlist
\item
  Powershell Lesson Pre: Your First Terminal - Opening PowerShell, first commands, screen reader tricks
\item
  Powershell Lesson 0: Getting Started - Paths, shortcuts, tab completion
\item
  Powershell Lesson 1: Navigation - Moving around the file system confidently
\item
  Powershell Lesson 2: File and Folder Manipulation - Creating, editing, moving files
\item
  Powershell Lesson 3: Input, Output, and Piping - Redirecting output, piping commands
\item
  Powershell Lesson 4: Environment Variables and Aliases - Setting variables, creating shortcuts
\item
  Powershell Lesson 5: Filling in the Gaps - Control flow, profiles, useful tricks
\item
  Powershell Lesson 6: Advanced Terminal Techniques - Scripts, functions, professional workflows
\item
  PowerShell Unit Test \& Practice - Comprehensive self-assessment
\end{itemize}

Time commitment: 30-45 hours (for screen reader users)
Best for: Users who want modern Windows features and advanced automation
Skills: Terminal navigation, piping, advanced scripting, professional automation

\subsubsection*{Windows Command Prompt Foundation (Simplified Alternative)}\label{docs__pandoc__latex__src__readme.md__windows-command-prompt-foundation-simplified-alternative}

Where to Start: \href{docs\%5Cpandoc\%5Clatex\%5Csrc\%5CCMD_Complete_Curriculum.md\#part-2-cmd}{CMD Curriculum Overview}

Lessons Included:

\begin{itemize}
\tightlist
\item
  Command Line Lesson Pre: Your First Terminal - Opening CMD, first commands, screen reader tricks
\item
  Command Line Lesson 0: Getting Started - Paths, shortcuts, command basics
\item
  Command Line Lesson 1: Navigation - Moving around the file system confidently
\item
  Command Line Lesson 2: File and Folder Manipulation - Creating, editing, moving files
\item
  Command Line Lesson 3: Input, Output \& Redirection - Redirecting output, piping commands
\item
  Command Line Lesson 4: Environment Variables \& Shortcuts - Setting variables, shortcuts
\item
  Command Line Lesson 5: Filling in the Gaps - Batch files, advanced techniques
\item
  Command Line Lesson 6: Advanced Terminal Techniques - Scripts, automation, professional workflows
\item
  CMD Unit Test \& Practice - Comprehensive self-assessment
\end{itemize}

Time commitment: 30-45 hours (for screen reader users)
Best for: Absolute beginners or users who prefer simplicity
Skills: Terminal navigation, file operations, batch scripting, basic automation

\subsubsection*{Git Bash Foundation (Cross-Platform Skills)}\label{docs__pandoc__latex__src__readme.md__git-bash-foundation-cross-platform-skills}

Where to Start: \href{docs\%5Cpandoc\%5Clatex\%5Csrc\%5CGitBash_Complete_Curriculum.md\#part-3-git-bash}{Git Bash Curriculum Overview}

Introductory \& Reference Materials:

Lessons Included:

\begin{itemize}
\tightlist
\item
  Bash Lesson Pre: Your First Terminal - Opening Git Bash, first commands, screen reader tricks
\item
  Bash Lesson 0: Getting Started - Paths, shortcuts, Unix path notation
\item
  Bash Lesson 1: Navigation - Moving around the file system confidently
\item
  Bash Lesson 2: File and Folder Manipulation - Creating, editing, moving files
\item
  Bash Lesson 3: Input, Output \& Piping - Redirecting output, piping commands
\item
  Bash Lesson 4: Environment Variables \& Aliases - Setting variables, creating shortcuts
\item
  Bash Lesson 5: Filling in the Gaps - Shell profiles, history, debugging
\item
  Bash Lesson 6: Advanced Terminal Techniques - Scripts, functions, professional workflows
\item
  GitBash Unit Test \& Practice - Comprehensive self-assessment
\end{itemize}

Time commitment: 30-45 hours (for screen reader users)
Best for: Users who want cross-platform skills (Windows, macOS, Linux)
Skills: Unix/bash commands, shell scripting, cross-platform automation

Important

Choose ONE pathway and complete it fully. All three teach identical fundamental concepts using different tools. Each integrates fully with 3D design workflows.

\subsection*{3dMake Foundation \& Design (11 Lessons + 6 Appendices)}\label{docs__pandoc__latex__src__readme.md__3dmake-foundation--design-11-lessons--6-appendices}

Where to Start: \href{docs\%5Cpandoc\%5Clatex\%5Csrc\%5C3dMake_Foundation\%5C3dMake_Intro\%5C3dMake_Intro.md}{3dMake Introduction}

Students build parametric 3D designs using OpenSCAD and automate the workflow with 3DMake.

\subsubsection*{Main Curriculum: 11 Progressive Lessons}\label{docs__pandoc__latex__src__readme.md__main-curriculum-11-progressive-lessons}

{\def\LTcaptype{none} % do not increment counter
\begin{longtable}[]{@{}
  >{\raggedright\arraybackslash}p{(\linewidth - 6\tabcolsep) * \real{0.1885}}
  >{\raggedright\arraybackslash}p{(\linewidth - 6\tabcolsep) * \real{0.0738}}
  >{\raggedright\arraybackslash}p{(\linewidth - 6\tabcolsep) * \real{0.6475}}
  >{\raggedright\arraybackslash}p{(\linewidth - 6\tabcolsep) * \real{0.0902}}@{}}
\toprule\noalign{}
\begin{minipage}[b]{\linewidth}\raggedright
Part
\end{minipage} & \begin{minipage}[b]{\linewidth}\raggedright
Lessons
\end{minipage} & \begin{minipage}[b]{\linewidth}\raggedright
Focus
\end{minipage} & \begin{minipage}[b]{\linewidth}\raggedright
Duration
\end{minipage} \\
\midrule\noalign{}
\endhead
\bottomrule\noalign{}
\endlastfoot
Foundations & 1-3 &
Environment setup, primitives, parametric design + advanced language features
& 4-5 hours \\
Verification \& Safety & 4-5 &
AI verification, safety protocols, materials & 2-3 hours \\
Applied Projects & 6-8 &
Practical commands, text functions, transforms, assembly patterns &
5-6 hours \\
Advanced Topics & 9-10 &
Automation with file I/O, troubleshooting, mastery & 3-4 hours \\
Leadership & 11 & Stakeholder-centric design & 2-3 hours \\
\end{longtable}
}

Total: 45-55 hours instruction + practice projects

\subsubsection*{Reference Appendices}\label{docs__pandoc__latex__src__readme.md__reference-appendices}

\textbf{Command Line Tools:}

\begin{itemize}
\tightlist
\item
  \hyperref[docs__pandoc__latex__src__cli_appendices_complete.md__appendix-a]{Appendix A: Command Line (CMD/Batch) Integration for SCAD Workflows} - Batch processing and automation for CMD users
\item
  \hyperref[docs__pandoc__latex__src__cli_appendices_complete.md__appendix-b]{Appendix B: PowerShell Integration for SCAD Workflows} - Batch processing, automation scripts, and advanced workflow integration for PowerShell users
\item
  \hyperref[docs__pandoc__latex__src__cli_appendices_complete.md__appendix-c]{Appendix C: Git Bash / POSIX Shell Integration for SCAD Workflows} - Shell scripting and automation for Git Bash / Unix shell users
\end{itemize}

\textbf{3dMake:}

\begin{itemize}
\tightlist
\item
  \hyperref[docs__pandoc__latex__src__3dmake_appendices_complete.md__appendix-a]{3DMake CLI Cheat Sheet} - Quick reference for 3dm command-line tool
\item
  \hyperref[docs__pandoc__latex__src__3dmake_appendices_complete.md__appendix-b]{Appendix A: Comprehensive Slicing Guide} - Complete reference for PrusaSlicer, Bambu Studio, Cura, and OrcaSlicer configuration
\item
  \hyperref[docs__pandoc__latex__src__3dmake_appendices_complete.md__appendix-c]{Appendix B: Material Properties \& Selection Guide} - Detailed material reference including shrinkage data, print settings, and properties
\item
  \hyperref[docs__pandoc__latex__src__3dmake_appendices_complete.md__appendix-d]{Appendix C: Tolerance Testing \& Quality Assurance Matrix} - Comprehensive QA procedures and tolerance validation methods
\item
  \hyperref[docs__pandoc__latex__src__3dmake_appendices_complete.md__appendix-e]{Appendix D: Advanced OpenSCAD Concepts} - Specialized topics including gears, batch processing, performance, print orientation, and recursion
\item
  \href{docs\%5Cpandoc\%5Clatex\%5Csrc\%5Cassets\%5C3dMake_Foundation\%5CREADME.md}{3dMake Code Examples \& Assets} - OpenSCAD code examples and reference designs
\end{itemize}

\section*{The Accessible Toolchain}\label{docs__pandoc__latex__src__readme.md__the-accessible-toolchain}

\subsection*{Screen Reader Compatibility Throughout}\label{docs__pandoc__latex__src__readme.md__screen-reader-compatibility-throughout}

This course uses tools designed for screen reader access:

\begin{itemize}
\tightlist
\item
  Terminal/Command line - Text-based, fully accessible to NVDA, JAWS, Orca, VoiceOver
\item
  OpenSCAD - Free, open-source text-based CAD (no visual-only GUI dependency)
\item
  3DMake - Command-line build tool eliminating GUI navigation
\item
  Accessible text editors - Notepad, VS Code, Notepad++, Nano, Vim, etc. (all keyboard-driven, screen reader friendly)
\end{itemize}

See \hyperref[docs__pandoc__latex__src__3dmake_appendices_complete.md__appendix-a]{Appendix A: 3DMake Cheat Sheet} for detailed keyboard shortcuts and configuration.

\subsection*{3DMake: Non-Visual Build Automation}\label{docs__pandoc__latex__src__readme.md__3dmake-non-visual-build-automation}

3DMake makes the entire design-to-print pipeline accessible:

\begin{lstlisting}[style=Alabaster, language=bash]
3dm build        -> Compiles main.scad to main.stl
3dm info         -> Validates geometry and runs diagnostics
3dm slice        -> Prepares model for printing

\end{lstlisting}

\begin{itemize}
\tightlist
\item
  No GUI navigation needed
\item
  Automation eliminates repetitive manual steps
\item
  Configuration files store parameters as human-readable text
\item
  Error reporting is text-based (screen reader accessible)
\item
  Works with command-line slicers
\end{itemize}

\subsection*{Iterative, Non-Visual Design}\label{docs__pandoc__latex__src__readme.md__iterative-non-visual-design}

Students learn to design through code and testing, not visual previews:

\begin{itemize}
\tightlist
\item
  Write parametric OpenSCAD code in accessible editors
\item
  Run \texttt{3dm\ build} to compile to printable file
\item
  Use measurement-based verification (calipers, scales, functional testing)
\item
  Iterate by editing parameters and rebuilding
\item
  No reliance on 3D preview or visual feedback
\end{itemize}

\section*{Project-Based Learning}\label{docs__pandoc__latex__src__readme.md__project-based-learning}

3dMake lessons include hands-on projects:

\begin{itemize}
\tightlist
\item
  Lesson 6: Keychain with embossed text (3dm commands)
\item
  Lesson 7: Phone stand (parametric transforms)
\item
  Lesson 8: Stackable bins (interlocking features, tolerances)
\item
  Lesson 9: Build automation (PowerShell batch processing)
\item
  Lesson 10: QA testing + accessibility audit (measurement, troubleshooting)
\item
  Lesson 11: Beaded jewelry holder (stakeholder-driven design)
\end{itemize}

Each project requires:

\begin{itemize}
\tightlist
\item
  Parametric OpenSCAD code (clean, well-commented)
\item
  Functional prototype (tested, iterated)
\item
  Complete documentation (reflections, measurements, decisions)
\end{itemize}

\section*{Learning Support}\label{docs__pandoc__latex__src__readme.md__learning-support}

\subsection*{Primary Navigation \& Overview}\label{docs__pandoc__latex__src__readme.md__primary-navigation--overview}

\href{docs\%5Cpandoc\%5Clatex\%5Csrc\%5C3dMake_Foundation\%5C3dMake_Foundation_Curriculum_Guide.md}{Curriculum Guide} - Detailed overview of all lessons and appendices\\
\hyperref[docs__pandoc__latex__src__3dmake_appendices_complete.md__appendix-a]{Quick Reference} - At-a-glance command and syntax reference\\
\hyperref[docs__pandoc__latex__src__3dmake_all_lessons.md__3dmake-foundation--design--complete-course-compendium]{mdBook Navigation Guide} - How to navigate this curriculum

\subsection*{Setup \& Configuration Guides}\label{docs__pandoc__latex__src__readme.md__setup--configuration-guides}

\hyperref[docs__pandoc__latex__src__3dmake_all_lessons.md__3dmake-foundation--design--complete-course-compendium]{3dMake Setup Guide} - Installation walkthrough\\
\hyperref[docs__pandoc__latex__src__3dmake_all_lessons.md__3dmake-foundation--design--complete-course-compendium]{VSCode Setup Guide} - Accessibility configuration\\
\hyperref[docs__pandoc__latex__src__3dmake_all_lessons.md__3dmake-foundation--design--complete-course-compendium]{Screen Reader Coding Tips (NVDA \& JAWS)} - Keyboard shortcuts and coding configuration

\subsection*{Quick Reference Materials}\label{docs__pandoc__latex__src__readme.md__quick-reference-materials}

\hyperref[docs__pandoc__latex__src__3dmake_appendices_complete.md__appendix-a]{OpenSCAD Cheat Sheet} - Syntax quick-reference\\
\hyperref[docs__pandoc__latex__src__3dmake_appendices_complete.md__appendix-a]{3DMake CLI Cheat Sheet} - Complete 3dm command reference\\
\hyperref[docs__pandoc__latex__src__3dmake_appendices_complete.md__appendix-c]{Filament Comparison Table} - Material selection guide\\
\hyperref[docs__pandoc__latex__src__3dmake_appendices_complete.md__appendix-d]{Master Rubric} - Project assessment criteria

\subsection*{Safety \& Maintenance}\label{docs__pandoc__latex__src__readme.md__safety--maintenance}

\hyperref[docs__pandoc__latex__src__3dmake_foundation__lessons_3dmake_5__lessons_3dmake_5.md__safety-protocols-and-the-physical-fabrication-interface]{Safety Checklist} - Pre-print safety procedures\\
\hyperref[docs__pandoc__latex__src__3dmake_foundation__lessons_3dmake_5__lessons_3dmake_5.md__safety-protocols-and-the-physical-fabrication-interface]{Maintenance Log} - Printer maintenance tracking

\subsection*{Troubleshooting \& Diagnostics}\label{docs__pandoc__latex__src__readme.md__troubleshooting--diagnostics}

\hyperref[docs__pandoc__latex__src__backmatter_complete.md__troubleshooting]{Common Issues and Solutions} - Comprehensive troubleshooting guide\\
\hyperref[docs__pandoc__latex__src__backmatter_complete.md__troubleshooting]{Diagnostic Checklist} - Systematic problem diagnosis\\
\hyperref[docs__pandoc__latex__src__backmatter_complete.md__troubleshooting]{Technical Reference} - Advanced systems architecture and professional reference

\subsection*{Glossaries \& Terminology}\label{docs__pandoc__latex__src__readme.md__glossaries--terminology}

\hyperref[docs__pandoc__latex__src__backmatter_complete.md__student-glossary]{Student Glossary} - Complete terminology reference for OpenSCAD, 3dMake, 3D Printing, and CLI\\
\hyperref[docs__pandoc__latex__src__backmatter_complete.md__teacher-glossary]{Teacher Glossary} - Pedagogical framework and instructional design terminology

\subsection*{Command Line Integration Appendices}\label{docs__pandoc__latex__src__readme.md__command-line-integration-appendices}

\href{docs\%5Cpandoc\%5Clatex\%5Csrc\%5CAppendices\%5CCLI_Appendices.md}{CLI Appendices Overview} - All command line integration guides\\
\hyperref[docs__pandoc__latex__src__cli_appendices_complete.md__appendix-a]{Appendix A: CMD/Batch Integration} - Windows CMD automation\\
\hyperref[docs__pandoc__latex__src__cli_appendices_complete.md__appendix-b]{Appendix B: PowerShell Integration} - PowerShell workflow automation\\
\hyperref[docs__pandoc__latex__src__cli_appendices_complete.md__appendix-c]{Appendix C: Git Bash Integration} - Unix shell integration

\subsection*{3dMake \& OpenSCAD Appendices}\label{docs__pandoc__latex__src__readme.md__3dmake--openscad-appendices}

\href{docs\%5Cpandoc\%5Clatex\%5Csrc\%5CAppendices\%5C3DMake_Appendices.md}{3dMake Appendices Overview} - All 3dMake reference materials\\
\hyperref[docs__pandoc__latex__src__3dmake_appendices_complete.md__appendix-b]{Appendix A: Comprehensive Slicing Guide} - PrusaSlicer, Bambu Studio, Cura, OrcaSlicer\\
\hyperref[docs__pandoc__latex__src__3dmake_appendices_complete.md__appendix-c]{Appendix B: Material Properties \& Selection Guide} - Detailed material specifications\\
\hyperref[docs__pandoc__latex__src__3dmake_appendices_complete.md__appendix-d]{Appendix C: Tolerance Testing \& QA Matrix} - Quality assurance procedures\\
\hyperref[docs__pandoc__latex__src__3dmake_appendices_complete.md__appendix-e]{Appendix D: Advanced OpenSCAD Concepts} - Gears, recursion, performance optimization

\subsection*{Code Examples \& Assets}\label{docs__pandoc__latex__src__readme.md__code-examples--assets}

\href{docs\%5Cpandoc\%5Clatex\%5Csrc\%5Cassets\%5C3dMake_Foundation\%5CREADME.md}{3dMake Code Examples \& Assets} - OpenSCAD code examples and reference designs

\subsection*{Further Reading \& References}\label{docs__pandoc__latex__src__readme.md__further-reading--references}

\hyperref[docs__pandoc__latex__src__backmatter_complete.md__further-reading]{References \& Further Reading} - Peer-reviewed research and scholarly works

\section*{Supplemental Resources \& Textbooks}\label{docs__pandoc__latex__src__readme.md__supplemental-resources--textbooks}

\subsection*{Textbook Options}\label{docs__pandoc__latex__src__readme.md__textbook-options}

\textbf{Beginner-Friendly Textbooks:}

\begin{itemize}
\tightlist
\item
  \href{https://nostarch.com/programmingopenscad}{Programming with OpenSCAD: A Beginner\textquotesingle s Guide to Coding 3D-Printable Objects} by Justin Gohde and Marius Kintel - Comprehensive reference covering OpenSCAD syntax, geometry concepts, and design patterns. Ideal for deep dives into specific topics and as a reference guide throughout the course. Available in EPUB format (screen reader accessible).
\item
  \href{https://www.packtpub.com/en-us/product/simplifying-3d-printing-with-openscad-9781801813174}{Simplifying 3D Printing with OpenSCAD} by Alicia Noors - Focused on practical workflows, optimization, and real-world printing scenarios. Available in EPUB/PDF formats.
\item
  \href{https://www.apress.com/gp/book/9781484213292}{OpenSCAD for 3D Printing} by Al Williams - Hands-on introduction to OpenSCAD with practical 3D printing projects. Covers basics through intermediate topics with step-by-step tutorials.
\item
  \href{https://www.amazon.com/Mastering-OpenSCAD-Jochen-Kerdels/dp/1803243538}{Mastering OpenSCAD} by Jochen Kerdels - Project-based learning approach with emphasis on parametric design and reusable code modules.
\end{itemize}

\textbf{Free Online Resources:}

\begin{itemize}
\tightlist
\item
  \href{https://en.wikibooks.org/wiki/OpenSCAD_User_Manual}{Official OpenSCAD Documentation} - Complete language reference and user manual on Wikibooks. Fully accessible and searchable.
\item
  \href{https://openscad.org/documentation.html}{OpenSCAD Tutorial} - Official beginner tutorial covering core concepts and syntax.
\item
  \href{https://openscad.org/cheatsheet/}{OpenSCAD Cheat Sheet (Official)} - Quick reference for all OpenSCAD commands and functions.
\end{itemize}

\subsection*{Companion Teaching Resources}\label{docs__pandoc__latex__src__readme.md__companion-teaching-resources}

\textbf{Curriculum-Specific Materials:}

\href{https://programmingwithopenscad.github.io/learning.html}{Practice Worksheets} - Printable worksheets for visualization practice, decomposition exercises, vocabulary building, and assessment.\\
\href{https://programmingwithopenscad.github.io/quick-reference.html}{Visual Quick Reference} - Command syntax guides and geometry reference.\\
\href{https://github.com/ProgrammingWithOpenSCAD/CodeSolutions}{Code Solutions Repository} - Working OpenSCAD examples organized by topic (3D shapes, transformations, loops, modules, if-statements, advanced techniques).

\textbf{Community Resources:}

\begin{itemize}
\tightlist
\item
  \href{https://www.thingiverse.com/search?q=openscad&type=things}{Thingiverse OpenSCAD Designs} - Community-created parametric designs with source code available for study and modification.
\item
  \href{https://forum.openscad.org/}{OpenSCAD Forums} - Active community support for troubleshooting and learning.
\item
  \href{https://www.youtube.com/results?search_query=openscad+tutorial+beginner}{OpenSCAD Tutorials (YouTube)} - Video tutorials (note: supplemental to text-based curriculum; transcripts recommended for accessibility).
\end{itemize}

\section*{Getting Started}\label{docs__pandoc__latex__src__readme.md__getting-started}

\subsection*{For Students}\label{docs__pandoc__latex__src__readme.md__for-students}

\textbf{Step 1: Screen Reader Setup \& Accessibility}

Before beginning any technical work, optimize your screen reader configuration for command-line and coding work:

\begin{itemize}
\tightlist
\item
  Start with \hyperref[docs__pandoc__latex__src__setup_accessibility_complete.md__accessibility-guide]{Screen Reader Accessibility Guide} - Configure NVDA/JAWS for terminal work
\item
  Review \hyperref[docs__pandoc__latex__src__setup_accessibility_complete.md__screen-reader-choice]{Screen Reader Choice: Windows CLI} - Help choosing between NVDA and JAWS
\item
  Optional: \hyperref[docs__pandoc__latex__src__setup_accessibility_complete.md__braille-displays]{Braille Display \& Terminal Mode} - Refreshable braille display configuration
\item
  Complete \hyperref[docs__pandoc__latex__src__setup_accessibility_complete.md__editor-setup]{Editor Selection and Setup} - Choose and configure an accessible text editor
\end{itemize}

\textbf{Step 2: Choose Your Command Line Pathway}

Read \hyperref[docs__pandoc__latex__src__command_line_interface_selection__command_line_interface_selection.md__command-line-fundamentals---choose-your-path]{Command Line Interface Selection Guide} to compare all three CLI options, then select ONE pathway:

\begin{itemize}
\tightlist
\item
  \href{docs\%5Cpandoc\%5Clatex\%5Csrc\%5CPS_Complete_Curriculum.md\#part-1-powershell}{PowerShell Curriculum Overview} - Recommended for Windows users
\item
  \href{docs\%5Cpandoc\%5Clatex\%5Csrc\%5CCMD_Complete_Curriculum.md\#part-2-cmd}{CMD Curriculum Overview} - Simplified alternative for beginners
\item
  \href{docs\%5Cpandoc\%5Clatex\%5Csrc\%5CGitBash_Complete_Curriculum.md\#part-3-git-bash}{Git Bash Curriculum Overview} - Cross-platform Unix skills
\end{itemize}

Complete ALL lessons in your chosen pathway (Pre through Lesson 6 + Unit Test) before moving to Step 3.

\textbf{Step 3: Begin 3dMake Foundation}

Once you\textquotesingle ve completed your CLI foundation:

\begin{enumerate}
\tightlist
\item
  Read \href{docs\%5Cpandoc\%5Clatex\%5Csrc\%5C3dMake_Foundation\%5C3dMake_Intro\%5C3dMake_Intro.md}{3dMake Introduction} - Overview of OpenSCAD and 3dMake
\item
  Start \hyperref[docs__pandoc__latex__src__3dmake_all_lessons.md__3dmake-foundation--design--complete-course-compendium]{Lesson 1: Environmental Configuration} - Install OpenSCAD and 3dMake
\item
  Progress through Lessons 2-11 sequentially
\item
  Complete all project templates as you encounter them
\item
  Review \href{docs\%5Cpandoc\%5Clatex\%5Csrc\%5C3dMake_Foundation\%5C3dMake_Foundation_Curriculum_Guide.md}{3dMake Foundation Curriculum Guide} for detailed lesson descriptions
\end{enumerate}

\textbf{Step 4: Projects \& Documentation}

Use the \href{docs\%5Cpandoc\%5Clatex\%5Csrc\%5C3dMake_Foundation\%5CTemplates\%5CStudent}{Student Templates} for all project documentation:

\begin{itemize}
\tightlist
\item
  Your First Print (Lesson 1)
\item
  Your Second Print \& Bonus Print (Lesson 2)
\item
  Dice Dice Dice (Lesson 4)
\item
  Parametric Keychain (Lesson 6)
\item
  Miniature Assembly (Lesson 7)
\item
  Snap Fit Clip (Lesson 8)
\item
  Accessibility Audit (Lesson 10)
\item
  Projects 0, 1, and 3 (Lesson 11)
\end{itemize}

\textbf{Step 5: Final Assessment}

Complete \hyperref[docs__pandoc__latex__src__3dmake_all_lessons.md__3dmake-foundation--design--complete-course-compendium]{3dMake Foundation Final Exam} to demonstrate mastery of OpenSCAD, 3DMake, and parametric design principles.

\begin{center}\rule{0.5\linewidth}{0.5pt}\end{center}

\subsection*{For Instructors}\label{docs__pandoc__latex__src__readme.md__for-instructors}

\textbf{Course Setup \& Planning:}

\begin{enumerate}
\tightlist
\item
  Review \hyperref[docs__pandoc__latex__src__syllabus.md__syllabus]{Syllabus} - Complete course structure, philosophy, learning objectives, and assessment framework
\item
  Read \href{docs\%5Cpandoc\%5Clatex\%5Csrc\%5C3dMake_Foundation\%5C3dMake_Foundation_Curriculum_Guide.md}{3dMake Foundation Curriculum Guide} - Detailed lesson-by-lesson breakdown with timing and objectives
\item
  Study \hyperref[docs__pandoc__latex__src__3dmake_appendices_complete.md__appendix-d]{Master Rubric} - Comprehensive project assessment criteria
\item
  Explore \hyperref[docs__pandoc__latex__src__backmatter_complete.md__teacher-glossary]{Teacher Glossary} - Pedagogical framework and instructional terminology
\end{enumerate}

\textbf{Assessment Materials:}

Use the \href{docs\%5Cpandoc\%5Clatex\%5Csrc\%5C3dMake_Foundation\%5CTemplates\%5CTeacher}{Teacher Templates} for grading and feedback:

\begin{itemize}
\tightlist
\item
  your\_first\_print\_teacher\_template.md
\item
  your\_second\_print\_teacher\_template.md
\item
  bonus\_print\_teacher\_template.md
\item
  dice\_dice\_dice\_teacher\_template.md
\item
  parametric\_keychain\_teacher\_template.md
\item
  miniature\_assembly\_teacher\_template.md
\item
  snap\_fit\_clip\_teacher\_template.md
\item
  accessibility\_audit\_teacher\_template.md
\item
  project\_0\_teacher\_template.md
\item
  project\_1\_teacher\_template.md
\item
  project\_3\_teacher\_template.md
\end{itemize}

\textbf{Final Exam:}

Administer \hyperref[docs__pandoc__latex__src__3dmake_all_lessons.md__3dmake-foundation--design--complete-course-compendium]{3dMake Foundation Final Exam} - 25-question comprehensive assessment covering error detection, code analysis, design challenges, and real-world application (100 points total).

\textbf{Accessibility Support:}

\begin{itemize}
\tightlist
\item
  Ensure all students complete \hyperref[docs__pandoc__latex__src__setup_accessibility_complete.md__accessibility-guide]{Screen Reader Setup} before beginning coursework
\item
  Reference \hyperref[docs__pandoc__latex__src__3dmake_all_lessons.md__3dmake-foundation--design--complete-course-compendium]{Screen Reader Coding Tips (NVDA \& JAWS)} for troubleshooting student screen reader issues
\item
  Provide access to \hyperref[docs__pandoc__latex__src__backmatter_complete.md__student-glossary]{Student Glossary} for terminology reference
\end{itemize}

\textbf{Additional Resources:}

\begin{itemize}
\tightlist
\item
  \hyperref[docs__pandoc__latex__src__backmatter_complete.md__troubleshooting]{Technical Reference} - Advanced troubleshooting and systems architecture
\item
  \hyperref[docs__pandoc__latex__src__backmatter_complete.md__further-reading]{Further Reading} - Peer-reviewed research supporting curriculum design
\item
  \hyperref[docs__pandoc__latex__src__3dmake_appendices_complete.md__appendix-a]{Quick Reference} - At-a-glance lesson and command reference for student handouts
\end{itemize}

\section*{Screen Readers}\label{docs__pandoc__latex__src__readme.md__screen-readers}

We know that users of this curriculum will primarily be JAWS and NVDA screenreader users, or else users of Orca if on a Linux-based system. Dolphin SuperNova (commercial) and Windows Narrator (built-in) are also supported; the workflows and recommendations in this document apply to them. See \url{https://yourdolphin.com/ScreenReader-Training} and \url{https://support.microsoft.com/en-us/windows/complete-guide-to-narrator-e4397a0d-ef4f-b386-d8ae-c172f109bdb1} for vendor documentation.

\section{3D Design \& Printing Curriculum - Non-Visual Toolchain Edition}\label{docs__pandoc__latex__src__syllabus.md__syllabus}

Author: Michael Ryan Hunsaker, M.Ed., Ph.D.\\
Last Updated: 2026-03-04

\subsection*{Overview}\label{docs__pandoc__latex__src__syllabus.md__overview}

This curriculum teaches 3D design and digital fabrication using a fully accessible, command-line-driven toolchain centered on OpenSCAD (text-based CAD), 3DMake (non-visual build automation), and accessible editors (VS Code, Notepad++, command-line editors) with screen reader support. Students progress from foundational command-line skills through guided projects to real-world, stakeholder-driven design challenges.

\subsubsection*{Who This Course Is For}\label{docs__pandoc__latex__src__syllabus.md__who-this-course-is-for}

This course is explicitly designed for blind and visually impaired learners who use screen readers (NVDA, JAWS, Orca, VoiceOver, Windows Narrator, Dolphin SuperNova). It eliminates GUI navigation and visual feedback in favor of keyboard-driven, text-based workflows that screen readers can fully access.

Accessibility is not an add-on. It is the foundation of every tool, workflow, and lesson in this curriculum.

\subsubsection*{Core Philosophy}\label{docs__pandoc__latex__src__syllabus.md__core-philosophy}

\begin{enumerate}
\item
  Text-First Design: All core work happens in text editors and command-line interfaces - no graphical CAD previews, no mouse-dependent menu navigation.
\item
  Parametric Thinking: Students learn to express geometry as code using OpenSCAD, enabling precise, reproducible, and iterable designs without visual feedback.
\item
  Automation and Independence: 3DMake automates the journey from code to printed object, handling compilation, slicing orchestration, and metadata management through simple command-line commands and text configuration files.
\item
  Screen Reader Mastery: Students develop fluency with accessibility technologies (NVDA, JAWS, VoiceOver) and accessible editors, building skills that apply to careers in software, engineering, and digital fabrication.
\item
  Real-World Impact: Projects culminate in designing assistive-technology solutions for real stakeholders, combining technical skill with human-centered design and documentation.
\end{enumerate}

\subsection*{Curriculum Structure \& Scope/Sequence}\label{docs__pandoc__latex__src__syllabus.md__curriculum-structure--scopesequence}

\subsubsection*{Setup \& Accessibility Fundamentals (Prerequisite - 2-3 hours)}\label{docs__pandoc__latex__src__syllabus.md__setup--accessibility-fundamentals-prerequisite---2-3-hours}

Start here: \hyperref[docs__pandoc__latex__src__setup_accessibility_complete.md__accessibility-guide]{Screen Reader Accessibility Guide}

Before choosing a command-line pathway, students optimize their screen reader setup for terminal work.

{\def\LTcaptype{none} % do not increment counter
\begin{longtable}[]{@{}
  >{\raggedright\arraybackslash}p{(\linewidth - 4\tabcolsep) * \real{0.5665}}
  >{\raggedright\arraybackslash}p{(\linewidth - 4\tabcolsep) * \real{0.0640}}
  >{\raggedright\arraybackslash}p{(\linewidth - 4\tabcolsep) * \real{0.3695}}@{}}
\toprule\noalign{}
\begin{minipage}[b]{\linewidth}\raggedright
Component
\end{minipage} & \begin{minipage}[b]{\linewidth}\raggedright
Duration
\end{minipage} & \begin{minipage}[b]{\linewidth}\raggedright
Content
\end{minipage} \\
\midrule\noalign{}
\endhead
\bottomrule\noalign{}
\endlastfoot
\hyperref[docs__pandoc__latex__src__setup_accessibility_complete.md__accessibility-guide]{Screen Reader Accessibility Guide}
& 1-1.5 hours &
NVDA/JAWS reference for PowerShell, CMD, and Git Bash; keyboard shortcuts \\
\hyperref[docs__pandoc__latex__src__setup_accessibility_complete.md__screen-reader-choice]{Screen Reader Choice: Windows CLI}
& 30 min &
Comparing NVDA, JAWS, Narrator, and Dolphin; choosing the right tool \\
\hyperref[docs__pandoc__latex__src__setup_accessibility_complete.md__braille-displays]{Braille Display \& Terminal Mode}
& 30 min &
Optional: configuring refreshable braille displays for terminal work \\
\hyperref[docs__pandoc__latex__src__setup_accessibility_complete.md__editor-setup]{Editor Selection and Setup}
& 30 min &
Choosing Notepad, Notepad++, or VS Code; configuring indent announcement \\
\end{longtable}
}

\subsubsection*{Command-Line Foundation (Choose Your Path)}\label{docs__pandoc__latex__src__syllabus.md__command-line-foundation-choose-your-path}

Start here: \hyperref[docs__pandoc__latex__src__command_line_interface_selection__command_line_interface_selection.md__command-line-fundamentals---choose-your-path]{Command Line Interface Selection Guide}

Students master terminal/command-line fundamentals before learning 3D design. Choose one of three equivalent pathways based on your operating system and preferences. All three pathways teach the same concepts and prepare you equally well for 3dMake work.

Important

Choose ONE pathway and complete it fully. All three teach identical fundamental concepts using different tools. Each integrates fully with 3D design workflows.

\paragraph*{Pathway A: PowerShell Foundation (Recommended for Windows)}\label{docs__pandoc__latex__src__syllabus.md__pathway-a-powershell-foundation-recommended-for-windows}

Total Duration: 30-45 hours
Start here: \href{docs\%5Cpandoc\%5Clatex\%5Csrc\%5CPS_Complete_Curriculum.md\#part-1-powershell}{PowerShell Curriculum Overview}

{\def\LTcaptype{none} % do not increment counter
\begin{longtable}[]{@{}
  >{\raggedright\arraybackslash}p{(\linewidth - 4\tabcolsep) * \real{0.6218}}
  >{\raggedright\arraybackslash}p{(\linewidth - 4\tabcolsep) * \real{0.0546}}
  >{\raggedright\arraybackslash}p{(\linewidth - 4\tabcolsep) * \real{0.3235}}@{}}
\toprule\noalign{}
\begin{minipage}[b]{\linewidth}\raggedright
Component
\end{minipage} & \begin{minipage}[b]{\linewidth}\raggedright
Duration
\end{minipage} & \begin{minipage}[b]{\linewidth}\raggedright
Content
\end{minipage} \\
\midrule\noalign{}
\endhead
\bottomrule\noalign{}
\endlastfoot
\href{docs\%5Cpandoc\%5Clatex\%5Csrc\%5CPS_Complete_Curriculum.md\#part-1-powershell}{PowerShell Introduction}
& 20-30 min &
Screen-reader-friendly quick-start (JAWS/NVDA); essential commands overview \\
\href{docs\%5Cpandoc\%5Clatex\%5Csrc\%5CPS_Complete_Curriculum.md\#part-1-powershell}{PowerShell Tutorial}
& 30-45 min &
Hands-on tutorial: paths, navigation, wildcards, running scripts \\
\href{docs\%5Cpandoc\%5Clatex\%5Csrc\%5CPS_Complete_Curriculum.md\#ps-lesson-pre}{Powershell Lesson Pre: Your First Terminal}
& 2-2.5 hours &
Opening PowerShell, first commands, basic navigation, screen reader tricks \\
\href{docs\%5Cpandoc\%5Clatex\%5Csrc\%5CPS_Complete_Curriculum.md\#ps-lesson-0}{Powershell Lesson 0: Getting Started}
& 1.5 hours & Paths, shortcuts, tab completion \\
\href{docs\%5Cpandoc\%5Clatex\%5Csrc\%5CPS_Complete_Curriculum.md\#ps-lesson-1}{Powershell Lesson 1: Navigation}
& 2-2.5 hours & Moving around the file system confidently \\
\href{docs\%5Cpandoc\%5Clatex\%5Csrc\%5CPS_Complete_Curriculum.md\#ps-lesson-2}{Powershell Lesson 2: File \& Folder Manipulation}
& 2.5-3 hours &
Creating, editing, moving, copying, deleting files and directories \\
\href{docs\%5Cpandoc\%5Clatex\%5Csrc\%5CPS_Complete_Curriculum.md\#ps-lesson-3}{Powershell Lesson 3: Input, Output \& Piping}
& 2.5-3 hours &
Redirecting output, piping commands, understanding data flow \\
\href{docs\%5Cpandoc\%5Clatex\%5Csrc\%5CPS_Complete_Curriculum.md\#ps-lesson-4}{Powershell Lesson 4: Environment Variables \& Aliases}
& 2-2.5 hours &
Setting variables, creating shortcuts, persistent configurations \\
\href{docs\%5Cpandoc\%5Clatex\%5Csrc\%5CPS_Complete_Curriculum.md\#ps-lesson-5}{Powershell Lesson 5: Filling in the Gaps}
& 2.5-3 hours &
Control flow, profiles, useful tricks, scripting fundamentals \\
\href{docs\%5Cpandoc\%5Clatex\%5Csrc\%5CPS_Complete_Curriculum.md\#ps-lesson-6}{Powershell Lesson 6: Advanced Terminal Techniques}
& 4-4.5 hours &
Scripts, functions, loops, professional workflows, automation patterns \\
\href{docs\%5Cpandoc\%5Clatex\%5Csrc\%5CPS_Complete_Curriculum.md\#ps-unit-exam}{PowerShell Unit Test \& Practice}
& 2-3 hours & Practice exercises, assessment, reinforcement \\
\end{longtable}
}

Outcomes: Terminal fluency, file system mastery, basic scripting, screen reader optimization, automation readiness

\paragraph*{Pathway B: Windows Command Prompt (CMD) (Simpler alternative)}\label{docs__pandoc__latex__src__syllabus.md__pathway-b-windows-command-prompt-cmd-simpler-alternative}

Total Duration: 30-45 hours
Start here: \href{docs\%5Cpandoc\%5Clatex\%5Csrc\%5CCMD_Complete_Curriculum.md\#part-2-cmd}{CMD Curriculum Overview}

{\def\LTcaptype{none} % do not increment counter
\begin{longtable}[]{@{}
  >{\raggedright\arraybackslash}p{(\linewidth - 4\tabcolsep) * \real{0.6372}}
  >{\raggedright\arraybackslash}p{(\linewidth - 4\tabcolsep) * \real{0.0575}}
  >{\raggedright\arraybackslash}p{(\linewidth - 4\tabcolsep) * \real{0.3053}}@{}}
\toprule\noalign{}
\begin{minipage}[b]{\linewidth}\raggedright
Component
\end{minipage} & \begin{minipage}[b]{\linewidth}\raggedright
Duration
\end{minipage} & \begin{minipage}[b]{\linewidth}\raggedright
Content
\end{minipage} \\
\midrule\noalign{}
\endhead
\bottomrule\noalign{}
\endlastfoot
\href{docs\%5Cpandoc\%5Clatex\%5Csrc\%5CCMD_Complete_Curriculum.md\#cmd-lesson-pre}{Command Line Lesson Pre: Your First Terminal}
& 2-2.5 hours &
Opening CMD, first commands, basic navigation, screen reader tricks \\
\href{docs\%5Cpandoc\%5Clatex\%5Csrc\%5CCMD_Complete_Curriculum.md\#cmd-lesson-0}{Command Line Lesson 0: Getting Started}
& 1.5 hours & Paths, shortcuts, command basics \\
\href{docs\%5Cpandoc\%5Clatex\%5Csrc\%5CCMD_Complete_Curriculum.md\#cmd-lesson-1}{Command Line Lesson 1: Navigation}
& 2-2.5 hours & Moving around the file system confidently \\
\href{docs\%5Cpandoc\%5Clatex\%5Csrc\%5CCMD_Complete_Curriculum.md\#cmd-lesson-2}{Command Line Lesson 2: File \& Folder Manipulation}
& 2.5-3 hours &
Creating, editing, moving, copying, deleting files and directories \\
\href{docs\%5Cpandoc\%5Clatex\%5Csrc\%5CCMD_Complete_Curriculum.md\#cmd-lesson-3}{Command Line Lesson 3: Input, Output \& Redirection}
& 2-2.5 hours &
Redirecting output, piping commands, understanding data flow \\
\href{docs\%5Cpandoc\%5Clatex\%5Csrc\%5CCMD_Complete_Curriculum.md\#cmd-lesson-4}{Command Line Lesson 4: Environment Variables \& Shortcuts}
& 2-2.5 hours &
Setting variables, creating shortcuts, persistent configurations \\
\href{docs\%5Cpandoc\%5Clatex\%5Csrc\%5CCMD_Complete_Curriculum.md\#cmd-lesson-5}{Command Line Lesson 5: Filling in the Gaps}
& 2.5-3 hours &
Batch files, advanced techniques, scripting fundamentals \\
\href{docs\%5Cpandoc\%5Clatex\%5Csrc\%5CCMD_Complete_Curriculum.md\#cmd-lesson-6}{Command Line Lesson 6: Advanced Terminal Techniques}
& 3-3.5 hours & Scripts, automation patterns, professional workflows \\
\href{docs\%5Cpandoc\%5Clatex\%5Csrc\%5CCMD_Complete_Curriculum.md\#cmd-unit-exam}{CMD Unit Test \& Practice}
& 2-3 hours & Practice exercises, assessment, reinforcement \\
\end{longtable}
}

Outcomes: Terminal fluency, file system mastery, batch scripting, screen reader optimization, automation readiness

\paragraph*{Pathway C: Git Bash (Best for macOS/Linux or cross-platform development)}\label{docs__pandoc__latex__src__syllabus.md__pathway-c-git-bash-best-for-macoslinux-or-cross-platform-development}

Total Duration: 20-25 hours
Start here: \href{docs\%5Cpandoc\%5Clatex\%5Csrc\%5CGitBash_Complete_Curriculum.md\#part-3-git-bash}{Git Bash Curriculum Overview}

{\def\LTcaptype{none} % do not increment counter
\begin{longtable}[]{@{}
  >{\raggedright\arraybackslash}p{(\linewidth - 4\tabcolsep) * \real{0.6454}}
  >{\raggedright\arraybackslash}p{(\linewidth - 4\tabcolsep) * \real{0.0598}}
  >{\raggedright\arraybackslash}p{(\linewidth - 4\tabcolsep) * \real{0.2948}}@{}}
\toprule\noalign{}
\begin{minipage}[b]{\linewidth}\raggedright
Component
\end{minipage} & \begin{minipage}[b]{\linewidth}\raggedright
Duration
\end{minipage} & \begin{minipage}[b]{\linewidth}\raggedright
Content
\end{minipage} \\
\midrule\noalign{}
\endhead
\bottomrule\noalign{}
\endlastfoot
\href{docs\%5Cpandoc\%5Clatex\%5Csrc\%5CGitBash_Complete_Curriculum.md\#part-3-git-bash}{Git Bash Introduction}
& 20-30 min &
Screen-reader-friendly quick-start (JAWS/NVDA); essential commands \\
\href{docs\%5Cpandoc\%5Clatex\%5Csrc\%5CGitBash_Complete_Curriculum.md\#gb-tutorial}{Git Bash Tutorial}
& 30-45 min &
Hands-on tutorial: paths, navigation, wildcards, running scripts \\
\href{docs\%5Cpandoc\%5Clatex\%5Csrc\%5CGitBash_Complete_Curriculum.md\#gb-accessibility-guide}{Screen Reader Accessibility Guide for Git Bash}
& 30-45 min & NVDA and JAWS configuration specific to Git Bash \\
\href{docs\%5Cpandoc\%5Clatex\%5Csrc\%5CGitBash_Complete_Curriculum.md\#gb-lesson-pre}{Bash Lesson Pre: Your First Terminal}
& 2-2.5 hours &
Opening Git Bash, first commands, basic navigation, screen reader tricks \\
\href{docs\%5Cpandoc\%5Clatex\%5Csrc\%5CGitBash_Complete_Curriculum.md\#gb-lesson-0}{Bash Lesson 0: Getting Started}
& 1.5 hours &
Unix-style paths, shortcuts, command basics, Windows path conversion \\
\href{docs\%5Cpandoc\%5Clatex\%5Csrc\%5CGitBash_Complete_Curriculum.md\#gb-lesson-1}{Bash Lesson 1: Navigation}
& 2-2.5 hours & Moving around the file system confidently \\
\href{docs\%5Cpandoc\%5Clatex\%5Csrc\%5CGitBash_Complete_Curriculum.md\#gb-lesson-2}{Bash Lesson 2: File and Folder Manipulation}
& 2-2.5 hours &
Creating, editing, moving, copying, deleting files and directories \\
\href{docs\%5Cpandoc\%5Clatex\%5Csrc\%5CGitBash_Complete_Curriculum.md\#gb-lesson-3}{Bash Lesson 3: Input, Output \& Piping}
& 2-2.5 hours &
Redirecting output, piping with grep/sort/wc, understanding data flow \\
\href{docs\%5Cpandoc\%5Clatex\%5Csrc\%5CGitBash_Complete_Curriculum.md\#gb-lesson-4}{Bash Lesson 4: Environment Variables \& Aliases}
& 1.5-2 hours & Setting variables, creating aliases, editing .bashrc \\
\href{docs\%5Cpandoc\%5Clatex\%5Csrc\%5CGitBash_Complete_Curriculum.md\#gb-lesson-5}{Bash Lesson 5: Filling in the Gaps}
& 2-2.5 hours & Shell profiles, command history, debugging \\
\href{docs\%5Cpandoc\%5Clatex\%5Csrc\%5CGitBash_Complete_Curriculum.md\#gb-lesson-6}{Bash Lesson 6: Advanced Terminal Techniques}
& 2.5-3.5 hours &
Shell scripts, functions, loops, professional workflows \\
\href{docs\%5Cpandoc\%5Clatex\%5Csrc\%5CGitBash_Complete_Curriculum.md\#gb-unit-exam}{GitBash Unit Test \& Practice}
& 2-2.5 hours & Practice exercises, assessment, reinforcement \\
\end{longtable}
}

Outcomes: Terminal fluency, file system mastery, bash scripting, version control basics, automation readiness

\subsubsection*{Common Outcomes (All Pathways)}\label{docs__pandoc__latex__src__syllabus.md__common-outcomes-all-pathways}

\begin{itemize}
\tightlist
\item
  Comfort with terminal/command-line interface
\item
  File system navigation and manipulation
\item
  Basic scripting and automation
\item
  Screen reader optimization for terminal work
\item
  Foundation for 3DMake automation tasks
\end{itemize}

\subsubsection*{3dMake Foundation (Main Curriculum - 16-20 hours)}\label{docs__pandoc__latex__src__syllabus.md__3dmake-foundation-main-curriculum---16-20-hours}

Start here: \href{docs\%5Cpandoc\%5Clatex\%5Csrc\%5C3dMake_Foundation\%5C3dMake_Intro\%5C3dMake_Intro.md}{3dMake Introduction}

11 progressive lessons building from foundational concepts to leadership-level design thinking, organized in 5 parts. Version 2.1 adds comprehensive advanced programming and design topics throughout.

\paragraph*{\texorpdfstring{Foundations (Lessons 1-3 \textbar{} \textasciitilde{}4-5 hours)}{Foundations (Lessons 1-3 \textbar{} \textasciitilde4-5 hours)}}\label{docs__pandoc__latex__src__syllabus.md__foundations-lessons-1-3--4-5-hours}

{\def\LTcaptype{none} % do not increment counter
\begin{longtable}[]{@{}
  >{\raggedright\arraybackslash}p{(\linewidth - 6\tabcolsep) * \real{0.4416}}
  >{\raggedright\arraybackslash}p{(\linewidth - 6\tabcolsep) * \real{0.4221}}
  >{\raggedright\arraybackslash}p{(\linewidth - 6\tabcolsep) * \real{0.0779}}
  >{\raggedright\arraybackslash}p{(\linewidth - 6\tabcolsep) * \real{0.0584}}@{}}
\toprule\noalign{}
\begin{minipage}[b]{\linewidth}\raggedright
Lesson
\end{minipage} & \begin{minipage}[b]{\linewidth}\raggedright
Focus
\end{minipage} & \begin{minipage}[b]{\linewidth}\raggedright
Duration
\end{minipage} & \begin{minipage}[b]{\linewidth}\raggedright
Project
\end{minipage} \\
\midrule\noalign{}
\endhead
\bottomrule\noalign{}
\endlastfoot
\hyperref[docs__pandoc__latex__src__3dmake_all_lessons.md__3dmake-foundation--design--complete-course-compendium]{Lesson 1}
& Environmental Configuration + Code Documentation Standards & 60-90 min
& None \\
\hyperref[docs__pandoc__latex__src__3dmake_all_lessons.md__3dmake-foundation--design--complete-course-compendium]{Lesson 2}
& Primitives \& Boolean Operations + Modifier Characters Debugging &
75-90 min & None \\
\hyperref[docs__pandoc__latex__src__3dmake_all_lessons.md__3dmake-foundation--design--complete-course-compendium]{Lesson 3}
& Parametric Architecture + Advanced Programming Concepts & 90-120 min &
None \\
\end{longtable}
}

\paragraph*{\texorpdfstring{Verification \& Safety (Lessons 4-5 \textbar{} \textasciitilde{}2 hours)}{Verification \& Safety (Lessons 4-5 \textbar{} \textasciitilde2 hours)}}\label{docs__pandoc__latex__src__syllabus.md__verification--safety-lessons-4-5--2-hours}

{\def\LTcaptype{none} % do not increment counter
\begin{longtable}[]{@{}
  >{\raggedright\arraybackslash}p{(\linewidth - 6\tabcolsep) * \real{0.5231}}
  >{\raggedright\arraybackslash}p{(\linewidth - 6\tabcolsep) * \real{0.3231}}
  >{\raggedright\arraybackslash}p{(\linewidth - 6\tabcolsep) * \real{0.0846}}
  >{\raggedright\arraybackslash}p{(\linewidth - 6\tabcolsep) * \real{0.0692}}@{}}
\toprule\noalign{}
\begin{minipage}[b]{\linewidth}\raggedright
Lesson
\end{minipage} & \begin{minipage}[b]{\linewidth}\raggedright
Focus
\end{minipage} & \begin{minipage}[b]{\linewidth}\raggedright
Duration
\end{minipage} & \begin{minipage}[b]{\linewidth}\raggedright
Project
\end{minipage} \\
\midrule\noalign{}
\endhead
\bottomrule\noalign{}
\endlastfoot
\hyperref[docs__pandoc__latex__src__3dmake_all_lessons.md__3dmake-foundation--design--complete-course-compendium]{Lesson 4}
& AI-Enhanced Verification \& Feedback & 45-60 min & None \\
\hyperref[docs__pandoc__latex__src__3dmake_foundation__lessons_3dmake_5__lessons_3dmake_5.md__safety-protocols-and-the-physical-fabrication-interface]{Lesson 5}
& Safety Protocols \& Material Introduction & 60-90 min & None \\
\end{longtable}
}

\paragraph*{\texorpdfstring{Applied Projects (Lessons 6-8 \textbar{} \textasciitilde{}5-6 hours)}{Applied Projects (Lessons 6-8 \textbar{} \textasciitilde5-6 hours)}}\label{docs__pandoc__latex__src__syllabus.md__applied-projects-lessons-6-8--5-6-hours}

{\def\LTcaptype{none} % do not increment counter
\begin{longtable}[]{@{}
  >{\raggedright\arraybackslash}p{(\linewidth - 6\tabcolsep) * \real{0.4690}}
  >{\raggedright\arraybackslash}p{(\linewidth - 6\tabcolsep) * \real{0.2966}}
  >{\raggedright\arraybackslash}p{(\linewidth - 6\tabcolsep) * \real{0.0897}}
  >{\raggedright\arraybackslash}p{(\linewidth - 6\tabcolsep) * \real{0.1448}}@{}}
\toprule\noalign{}
\begin{minipage}[b]{\linewidth}\raggedright
Lesson
\end{minipage} & \begin{minipage}[b]{\linewidth}\raggedright
Focus
\end{minipage} & \begin{minipage}[b]{\linewidth}\raggedright
Duration
\end{minipage} & \begin{minipage}[b]{\linewidth}\raggedright
Project
\end{minipage} \\
\midrule\noalign{}
\endhead
\bottomrule\noalign{}
\endlastfoot
\hyperref[docs__pandoc__latex__src__3dmake_all_lessons.md__3dmake-foundation--design--complete-course-compendium]{Lesson 6}
& Practical 3dm Commands + String Functions & 75-105 min &
Customizable Keycap \\
\hyperref[docs__pandoc__latex__src__3dmake_all_lessons.md__3dmake-foundation--design--complete-course-compendium]{Lesson 7}
& Parametric Transforms + Math Functions & 90-120 min & Phone Stand \\
\hyperref[docs__pandoc__latex__src__3dmake_all_lessons.md__3dmake-foundation--design--complete-course-compendium]{Lesson 8}
& Advanced Design + Assembly Best Practices & 105-150 min &
Stackable Bins \\
\end{longtable}
}

\paragraph*{\texorpdfstring{Advanced Topics (Lessons 9-10 \textbar{} \textasciitilde{}3-4 hours)}{Advanced Topics (Lessons 9-10 \textbar{} \textasciitilde3-4 hours)}}\label{docs__pandoc__latex__src__syllabus.md__advanced-topics-lessons-9-10--3-4-hours}

{\def\LTcaptype{none} % do not increment counter
\begin{longtable}[]{@{}
  >{\raggedright\arraybackslash}p{(\linewidth - 6\tabcolsep) * \real{0.4011}}
  >{\raggedright\arraybackslash}p{(\linewidth - 6\tabcolsep) * \real{0.3333}}
  >{\raggedright\arraybackslash}p{(\linewidth - 6\tabcolsep) * \real{0.0734}}
  >{\raggedright\arraybackslash}p{(\linewidth - 6\tabcolsep) * \real{0.1921}}@{}}
\toprule\noalign{}
\begin{minipage}[b]{\linewidth}\raggedright
Lesson
\end{minipage} & \begin{minipage}[b]{\linewidth}\raggedright
Focus
\end{minipage} & \begin{minipage}[b]{\linewidth}\raggedright
Duration
\end{minipage} & \begin{minipage}[b]{\linewidth}\raggedright
Project
\end{minipage} \\
\midrule\noalign{}
\endhead
\bottomrule\noalign{}
\endlastfoot
\hyperref[docs__pandoc__latex__src__3dmake_all_lessons.md__3dmake-foundation--design--complete-course-compendium]{Lesson 9}
& Automation + File Import/Export (requires CLI Foundation) & 75-105 min
& Batch Processing Automation \\
\hyperref[docs__pandoc__latex__src__3dmake_all_lessons.md__3dmake-foundation--design--complete-course-compendium]{Lesson 10}
& Troubleshooting \& Mastery with Measurement & 120-150 min &
QA Testing + Accessibility Audit \\
\end{longtable}
}

\paragraph*{\texorpdfstring{Leadership (Lesson 11 \textbar{} \textasciitilde{}2 hours)}{Leadership (Lesson 11 \textbar{} \textasciitilde2 hours)}}\label{docs__pandoc__latex__src__syllabus.md__leadership-lesson-11--2-hours}

{\def\LTcaptype{none} % do not increment counter
\begin{longtable}[]{@{}
  >{\raggedright\arraybackslash}p{(\linewidth - 6\tabcolsep) * \real{0.4671}}
  >{\raggedright\arraybackslash}p{(\linewidth - 6\tabcolsep) * \real{0.3026}}
  >{\raggedright\arraybackslash}p{(\linewidth - 6\tabcolsep) * \real{0.0789}}
  >{\raggedright\arraybackslash}p{(\linewidth - 6\tabcolsep) * \real{0.1513}}@{}}
\toprule\noalign{}
\begin{minipage}[b]{\linewidth}\raggedright
Lesson
\end{minipage} & \begin{minipage}[b]{\linewidth}\raggedright
Focus
\end{minipage} & \begin{minipage}[b]{\linewidth}\raggedright
Duration
\end{minipage} & \begin{minipage}[b]{\linewidth}\raggedright
Project
\end{minipage} \\
\midrule\noalign{}
\endhead
\bottomrule\noalign{}
\endlastfoot
\hyperref[docs__pandoc__latex__src__3dmake_all_lessons.md__3dmake-foundation--design--complete-course-compendium]{Lesson 11}
& Stakeholder-Centric Design \& Design Thinking & 90-120 min &
Beaded Jewelry Holder \\
\end{longtable}
}

Total: 16-20 hours of instruction + projects

\subsubsection*{Reference Appendices}\label{docs__pandoc__latex__src__syllabus.md__reference-appendices}

\paragraph*{Overview Documents}\label{docs__pandoc__latex__src__syllabus.md__overview-documents}

\begin{itemize}
\tightlist
\item
  \textbf{\href{docs\%5Cpandoc\%5Clatex\%5Csrc\%5CAppendices\%5CCLI_Appendices.md}{CLI Appendices Overview}} - Summary of all command line integration guides
\item
  \textbf{\href{docs\%5Cpandoc\%5Clatex\%5Csrc\%5CAppendices\%5C3DMake_Appendices.md}{3dMake Appendices Overview}} - Summary of all 3dMake reference materials
\end{itemize}

\paragraph*{Command Line Integration Appendices}\label{docs__pandoc__latex__src__syllabus.md__command-line-integration-appendices}

Located in \texttt{Command\_Line\_Interface\_Selection/}:

{\def\LTcaptype{none} % do not increment counter
\begin{longtable}[]{@{}
  >{\raggedright\arraybackslash}p{(\linewidth - 4\tabcolsep) * \real{0.5216}}
  >{\raggedright\arraybackslash}p{(\linewidth - 4\tabcolsep) * \real{0.2112}}
  >{\raggedright\arraybackslash}p{(\linewidth - 4\tabcolsep) * \real{0.2672}}@{}}
\toprule\noalign{}
\begin{minipage}[b]{\linewidth}\raggedright
Appendix
\end{minipage} & \begin{minipage}[b]{\linewidth}\raggedright
Title
\end{minipage} & \begin{minipage}[b]{\linewidth}\raggedright
Use When
\end{minipage} \\
\midrule\noalign{}
\endhead
\bottomrule\noalign{}
\endlastfoot
\hyperref[docs__pandoc__latex__src__cli_appendices_complete.md__appendix-a]{Appendix A}
& CMD/Batch Integration for SCAD Workflows &
Automating with Windows Command Prompt and batch files \\
\hyperref[docs__pandoc__latex__src__cli_appendices_complete.md__appendix-b]{Appendix B}
& PowerShell Integration for SCAD Workflows &
Automating with PowerShell scripts and advanced workflows \\
\hyperref[docs__pandoc__latex__src__cli_appendices_complete.md__appendix-c]{Appendix C}
& Git Bash / POSIX Shell Integration for SCAD &
Automating with Unix/bash shell scripts for cross-platform \\
\end{longtable}
}

\paragraph*{3dMake \& OpenSCAD Appendices}\label{docs__pandoc__latex__src__syllabus.md__3dmake--openscad-appendices}

Located in \texttt{3dMake\_Foundation/}:

{\def\LTcaptype{none} % do not increment counter
\begin{longtable}[]{@{}
  >{\raggedright\arraybackslash}p{(\linewidth - 4\tabcolsep) * \real{0.4611}}
  >{\raggedright\arraybackslash}p{(\linewidth - 4\tabcolsep) * \real{0.2383}}
  >{\raggedright\arraybackslash}p{(\linewidth - 4\tabcolsep) * \real{0.3005}}@{}}
\toprule\noalign{}
\begin{minipage}[b]{\linewidth}\raggedright
Appendix
\end{minipage} & \begin{minipage}[b]{\linewidth}\raggedright
Title
\end{minipage} & \begin{minipage}[b]{\linewidth}\raggedright
Use When
\end{minipage} \\
\midrule\noalign{}
\endhead
\bottomrule\noalign{}
\endlastfoot
\hyperref[docs__pandoc__latex__src__3dmake_appendices_complete.md__appendix-a]{3DMake Cheat Sheet}
& Complete 3dm CLI Reference &
Quick command lookup, troubleshooting 3dm commands \\
\hyperref[docs__pandoc__latex__src__3dmake_appendices_complete.md__appendix-b]{Appendix A}
& Comprehensive Slicing Guide &
Slicing questions, switching slicers, quality issues \\
\hyperref[docs__pandoc__latex__src__3dmake_appendices_complete.md__appendix-c]{Appendix B}
& Material Properties \& Selection Guide &
Choosing material, troubleshooting prints, cost analysis \\
\hyperref[docs__pandoc__latex__src__3dmake_appendices_complete.md__appendix-d]{Appendix C}
& Tolerance Testing \& Quality Assurance Matrix &
Starting a project, verifying dimensions, quality issues \\
\hyperref[docs__pandoc__latex__src__3dmake_appendices_complete.md__appendix-e]{Appendix D}
& Advanced OpenSCAD Concepts &
Building mechanical systems, optimizing complex models \\
\href{docs\%5Cpandoc\%5Clatex\%5Csrc\%5Cassets\%5C3dMake_Foundation\%5CREADME.md}{Code Examples \& Assets}
& 3dMake Code Examples \& Assets &
Reference designs, working code examples \\
\end{longtable}
}

\subsubsection*{Project Templates \& Assessment Materials}\label{docs__pandoc__latex__src__syllabus.md__project-templates--assessment-materials}

\paragraph*{Student Project Templates}\label{docs__pandoc__latex__src__syllabus.md__student-project-templates}

Located in \texttt{3dMake\_Foundation/Templates/Student/}, these templates guide students through structured documentation for each project:

\begin{itemize}
\tightlist
\item
  \textbf{your\_first\_print\_student\_template.md} - Lesson 1 extension project
\item
  \textbf{your\_second\_print\_student\_template.md} - Lesson 2 basic parametric design
\item
  \textbf{bonus\_print\_student\_template.md} - Lesson 2 advanced practice
\item
  \textbf{dice\_dice\_dice\_student\_template.md} - Lesson 4 verification project
\item
  \textbf{parametric\_keychain\_student\_template.md} - Lesson 6 practical 3dm commands
\item
  \textbf{miniature\_assembly\_student\_template.md} - Lesson 7 transforms project
\item
  \textbf{snap\_fit\_clip\_student\_template.md} - Lesson 8 assembly design
\item
  \textbf{accessibility\_audit\_student\_template.md} - Lesson 10 QA testing
\item
  \textbf{project\_0\_student\_template.md} - Lesson 11 preliminary design
\item
  \textbf{project\_1\_student\_template.md} - Lesson 11 prototype iteration
\item
  \textbf{project\_3\_student\_template.md} - Lesson 11 final stakeholder project
\end{itemize}

Each template includes sections for problem statement, design decisions, code documentation, testing results, measurements, reflections, and iteration notes.

\paragraph*{Teacher Assessment Templates}\label{docs__pandoc__latex__src__syllabus.md__teacher-assessment-templates}

Located in \texttt{3dMake\_Foundation/Templates/Teacher/}, corresponding templates provide grading rubrics and feedback frameworks:

\begin{itemize}
\tightlist
\item
  \textbf{your\_first\_print\_teacher\_template.md}
\item
  \textbf{your\_second\_print\_teacher\_template.md}
\item
  \textbf{bonus\_print\_teacher\_template.md}
\item
  \textbf{dice\_dice\_dice\_teacher\_template.md}
\item
  \textbf{parametric\_keychain\_teacher\_template.md}
\item
  \textbf{miniature\_assembly\_teacher\_template.md}
\item
  \textbf{snap\_fit\_clip\_teacher\_template.md}
\item
  \textbf{accessibility\_audit\_teacher\_template.md}
\item
  \textbf{project\_0\_teacher\_template.md}
\item
  \textbf{project\_1\_teacher\_template.md}
\item
  \textbf{project\_3\_teacher\_template.md}
\end{itemize}

Each teacher template includes the 0-9 rubric scoring guide, common issues to watch for, and focus areas for feedback.

\paragraph*{Final Assessment}\label{docs__pandoc__latex__src__syllabus.md__final-assessment}

\textbf{\hyperref[docs__pandoc__latex__src__3dmake_all_lessons.md__3dmake-foundation--design--complete-course-compendium]{3dMake Foundation Final Exam}} - Comprehensive 25-question assessment (100 points total) covering:

\begin{itemize}
\tightlist
\item
  Section 1: Error Detection \& Code Analysis (Problems 1-10)
\item
  Section 2: Design Challenges \& Real-World Application (Problems 11-25)
\end{itemize}

Questions include error identification, code explanation, behavioral analysis, and design problem-solving.

\paragraph*{Master Rubric}\label{docs__pandoc__latex__src__syllabus.md__master-rubric}

\textbf{\hyperref[docs__pandoc__latex__src__3dmake_appendices_complete.md__appendix-d]{Master Rubric}} - Unified assessment framework used across all projects, providing detailed criteria for the three scoring categories (Problem \& Solution, Design \& Code Quality, Documentation).

\subsubsection*{Reference Materials for Students \& Teachers}\label{docs__pandoc__latex__src__syllabus.md__reference-materials-for-students--teachers}

\paragraph*{Glossaries}\label{docs__pandoc__latex__src__syllabus.md__glossaries}

\begin{itemize}
\tightlist
\item
  \textbf{\hyperref[docs__pandoc__latex__src__backmatter_complete.md__student-glossary]{Student Glossary}} - Complete terminology reference for OpenSCAD, 3dMake, 3D Printing, and CLI
\item
  \textbf{\hyperref[docs__pandoc__latex__src__backmatter_complete.md__teacher-glossary]{Teacher Glossary}} - Pedagogical framework, instructional design strategies, assessment concepts, and accessibility terminology
\end{itemize}

\paragraph*{Technical References}\label{docs__pandoc__latex__src__syllabus.md__technical-references}

\begin{itemize}
\tightlist
\item
  \textbf{\hyperref[docs__pandoc__latex__src__backmatter_complete.md__troubleshooting]{Technical Reference \& Troubleshooting}} - Advanced systems architecture, OpenSCAD professional modeling reference, 3DMake professional usage, CLI references (CMD, PowerShell, Git Bash), and comprehensive troubleshooting compendium
\item
  \textbf{\hyperref[docs__pandoc__latex__src__backmatter_complete.md__further-reading]{Further Reading \& References}} - Peer-reviewed research and scholarly works supporting curriculum design
\end{itemize}

\paragraph*{Quick References}\label{docs__pandoc__latex__src__syllabus.md__quick-references}

\begin{itemize}
\tightlist
\item
  \textbf{\href{docs\%5Cpandoc\%5Clatex\%5Csrc\%5C3dMake_Foundation\%5C3dMake_Foundation_Curriculum_Guide.md}{3dMake Foundation Curriculum Guide}} - Detailed lesson-by-lesson breakdown with objectives and timing
\item
  \textbf{\hyperref[docs__pandoc__latex__src__3dmake_appendices_complete.md__appendix-a]{3dMake Quick Reference}} - At-a-glance command and syntax reference
\end{itemize}

\subsection*{Learning Progression: Student Roles}\label{docs__pandoc__latex__src__syllabus.md__learning-progression-student-roles}

Students move through roles across the curriculum:

{\def\LTcaptype{none} % do not increment counter
\begin{longtable}[]{@{}
  >{\raggedright\arraybackslash}p{(\linewidth - 6\tabcolsep) * \real{0.1736}}
  >{\raggedright\arraybackslash}p{(\linewidth - 6\tabcolsep) * \real{0.1488}}
  >{\raggedright\arraybackslash}p{(\linewidth - 6\tabcolsep) * \real{0.2810}}
  >{\raggedright\arraybackslash}p{(\linewidth - 6\tabcolsep) * \real{0.3967}}@{}}
\toprule\noalign{}
\begin{minipage}[b]{\linewidth}\raggedright
Phase
\end{minipage} & \begin{minipage}[b]{\linewidth}\raggedright
Role
\end{minipage} & \begin{minipage}[b]{\linewidth}\raggedright
Core Tools
\end{minipage} & \begin{minipage}[b]{\linewidth}\raggedright
Focus
\end{minipage} \\
\midrule\noalign{}
\endhead
\bottomrule\noalign{}
\endlastfoot
CLI Foundation & Observer/Learner & Terminal, command line &
CLI fundamentals and keyboard navigation \\
3dMake Lessons 1-5 & Observer/Learner & OpenSCAD, 3DMake, editor &
Using CLI tools, safety, concepts, measurement \\
3dMake Lessons 6-8 & Operator & Editor, OpenSCAD, 3DMake, slicer &
Hands-on practice with structured projects \\
3dMake Lessons 9-10 & Designer & Full toolchain &
Parametric design, automation, troubleshooting \\
3dMake Lesson 11 & Problem-Solver & Full toolchain + documentation &
Stakeholder design, real-world impact \\
\end{longtable}
}

\subsection*{The Accessible Toolchain: How It Works}\label{docs__pandoc__latex__src__syllabus.md__the-accessible-toolchain-how-it-works}

\subsubsection*{OpenSCAD - Text-Based 3D Design}\label{docs__pandoc__latex__src__syllabus.md__openscad---text-based-3d-design}

OpenSCAD is a free, open-source CAD tool that uses a programming language to describe 3D geometry. Students write code that defines shapes, transforms them, and combines them using Boolean operations.

Why OpenSCAD?

\begin{itemize}
\tightlist
\item
  Screen reader friendly: All work happens in a text editor; no visual-only 3D preview.
\item
  Repeatable: Code is version-controlled, documented, and shareable.
\item
  Parametric: Variables allow students to design once and generate variations by changing numbers.
\item
  No visual dependency: Students reason about geometry through code structure and testing.
\end{itemize}

\subsubsection*{3DMake - The Non-Visual Build Bridge}\label{docs__pandoc__latex__src__syllabus.md__3dmake---the-non-visual-build-bridge}

3DMake is a command-line tool that automates the journey from OpenSCAD code to a printable file:

\begin{lstlisting}[style=Alabaster, language=bash]
3dm build
3dm info
3dm slice

\end{lstlisting}

Why 3DMake?

\begin{itemize}
\tightlist
\item
  No GUI navigation: All interaction is keyboard-driven and text-based.
\item
  Automation: Eliminates repetitive manual steps.
\item
  Metadata tracking: Configuration files store parameters as human-readable text.
\item
  Error reporting: Diagnostic output is text that screen readers can read aloud.
\end{itemize}

\subsubsection*{Accessible Editors}\label{docs__pandoc__latex__src__syllabus.md__accessible-editors}

Students write OpenSCAD code using screen reader-accessible editors:

\begin{itemize}
\tightlist
\item
  VS Code (Windows, macOS, Linux): Industry-standard with built-in screen reader support
\item
  Notepad++ (Windows): Lightweight, keyboard-driven, excellent screen reader support
\item
  Command-line editors (Nano, Vim, Emacs): Full keyboard control, no mouse needed
\end{itemize}

\subsection*{Prerequisites by Section}\label{docs__pandoc__latex__src__syllabus.md__prerequisites-by-section}

{\def\LTcaptype{none} % do not increment counter
\begin{longtable}[]{@{}
  >{\raggedright\arraybackslash}p{(\linewidth - 4\tabcolsep) * \real{0.1591}}
  >{\raggedright\arraybackslash}p{(\linewidth - 4\tabcolsep) * \real{0.2955}}
  >{\raggedright\arraybackslash}p{(\linewidth - 4\tabcolsep) * \real{0.5455}}@{}}
\toprule\noalign{}
\begin{minipage}[b]{\linewidth}\raggedright
Section
\end{minipage} & \begin{minipage}[b]{\linewidth}\raggedright
Prerequisites
\end{minipage} & \begin{minipage}[b]{\linewidth}\raggedright
What You\textquotesingle ll Learn
\end{minipage} \\
\midrule\noalign{}
\endhead
\bottomrule\noalign{}
\endlastfoot
Setup & None - start here &
Screen reader optimization, editor selection, accessibility setup \\
CLI Foundation & Setup &
Terminal basics, keyboard navigation, file operations, basic scripting \\
3dMake Lessons 1-5 & Setup (CLI Foundation recommended) &
3D printing concepts, safety, measurement, OpenSCAD basics, debugging \\
3dMake Lessons 6-8 & Lessons 1-5 &
Building projects, parametric design, transforms, tolerances \\
3dMake Lessons 9-10 & Lessons 6-8 + CLI Foundation required &
Automation, troubleshooting, advanced measurement and QA \\
3dMake Lesson 11 & Lessons 9-10 &
Stakeholder design, real-world prototyping, leadership \\
\end{longtable}
}

\subsection*{Grading Rubric}\label{docs__pandoc__latex__src__syllabus.md__grading-rubric}

All projects are scored on a 0-9 scale across three equally weighted categories (3 points each):

{\def\LTcaptype{none} % do not increment counter
\begin{longtable}[]{@{}
  >{\raggedright\arraybackslash}p{(\linewidth - 4\tabcolsep) * \real{0.1484}}
  >{\raggedright\arraybackslash}p{(\linewidth - 4\tabcolsep) * \real{0.0516}}
  >{\raggedright\arraybackslash}p{(\linewidth - 4\tabcolsep) * \real{0.8000}}@{}}
\toprule\noalign{}
\begin{minipage}[b]{\linewidth}\raggedright
Category
\end{minipage} & \begin{minipage}[b]{\linewidth}\raggedright
Points
\end{minipage} & \begin{minipage}[b]{\linewidth}\raggedright
What We Measure
\end{minipage} \\
\midrule\noalign{}
\endhead
\bottomrule\noalign{}
\endlastfoot
Problem \& Solution & 0-3 &
Does the design solve the stated problem? Are all functional requirements met? \\
Design \& Code Quality & 0-3 &
Is the OpenSCAD code clean, well-commented, and well-structured? Does the print work well? Is there evidence of iteration? \\
Documentation & 0-3 &
Are all sections complete? Are reflections thoughtful and specific? Are measurements recorded? \\
\end{longtable}
}

\subsubsection*{Category 1: Problem \& Solution (0-3 points)}\label{docs__pandoc__latex__src__syllabus.md__category-1-problem--solution-0-3-points}

{\def\LTcaptype{none} % do not increment counter
\begin{longtable}[]{@{}
  >{\raggedright\arraybackslash}p{(\linewidth - 2\tabcolsep) * \real{0.0400}}
  >{\raggedright\arraybackslash}p{(\linewidth - 2\tabcolsep) * \real{0.9600}}@{}}
\toprule\noalign{}
\begin{minipage}[b]{\linewidth}\raggedright
Score
\end{minipage} & \begin{minipage}[b]{\linewidth}\raggedright
Description
\end{minipage} \\
\midrule\noalign{}
\endhead
\bottomrule\noalign{}
\endlastfoot
3 &
The prototype clearly and effectively solves the stated problem. All functional requirements are met. The solution shows evidence of testing against the requirements. \\
2 &
The prototype mostly meets the problem. Most functional requirements are met. Minor gaps between the design and the requirements. \\
1 &
The prototype partially addresses the problem. Several functional requirements are not met or were not clearly tested. \\
0 &
The prototype does not address the stated problem, or no functional requirements were established. \\
\end{longtable}
}

\subsubsection*{Category 2: Design \& Code Quality (0-3 points)}\label{docs__pandoc__latex__src__syllabus.md__category-2-design--code-quality-0-3-points}

OpenSCAD code is central to this course. We evaluate the clarity, structure, and documentation of your code as much as the print quality.

{\def\LTcaptype{none} % do not increment counter
\begin{longtable}[]{@{}
  >{\raggedright\arraybackslash}p{(\linewidth - 2\tabcolsep) * \real{0.0341}}
  >{\raggedright\arraybackslash}p{(\linewidth - 2\tabcolsep) * \real{0.9659}}@{}}
\toprule\noalign{}
\begin{minipage}[b]{\linewidth}\raggedright
Score
\end{minipage} & \begin{minipage}[b]{\linewidth}\raggedright
Description
\end{minipage} \\
\midrule\noalign{}
\endhead
\bottomrule\noalign{}
\endlastfoot
3 &
Code is clean, well-organized, and thoroughly commented. Variables/modules are used appropriately. Print quality is excellent. Design shows original thinking and at least one meaningful iteration. \\
2 &
Code works but lacks clear structure or comments. Variables are used but could be better named. Print quality is acceptable. Some iteration evident. \\
1 &
Code is functional but poorly organized. Comments are minimal or missing. Print quality has defects. Little or no iteration. \\
0 &
Code does not work, is not submitted, or shows no original thinking. Print is not functional. \\
\end{longtable}
}

\subsubsection*{Category 3: Documentation (0-3 points)}\label{docs__pandoc__latex__src__syllabus.md__category-3-documentation-0-3-points}

{\def\LTcaptype{none} % do not increment counter
\begin{longtable}[]{@{}
  >{\raggedright\arraybackslash}p{(\linewidth - 2\tabcolsep) * \real{0.0354}}
  >{\raggedright\arraybackslash}p{(\linewidth - 2\tabcolsep) * \real{0.9646}}@{}}
\toprule\noalign{}
\begin{minipage}[b]{\linewidth}\raggedright
Score
\end{minipage} & \begin{minipage}[b]{\linewidth}\raggedright
Description
\end{minipage} \\
\midrule\noalign{}
\endhead
\bottomrule\noalign{}
\endlastfoot
3 &
All required sections are present, complete, and specific. Reflections are thoughtful and reference specific decisions, problems encountered, and learning. Photos/measurements are included. \\
2 &
Most required sections are present. Some sections are vague or missing detail. Reflections show some thought but are brief or generic. \\
1 &
Documentation is incomplete. Major sections are missing or consist of one-line responses. Reflections are minimal. \\
0 & Documentation is not submitted or is essentially empty. \\
\end{longtable}
}

\subsubsection*{Score Interpretation}\label{docs__pandoc__latex__src__syllabus.md__score-interpretation}

{\def\LTcaptype{none} % do not increment counter
\begin{longtable}[]{@{}
  >{\raggedright\arraybackslash}p{(\linewidth - 4\tabcolsep) * \real{0.1226}}
  >{\raggedright\arraybackslash}p{(\linewidth - 4\tabcolsep) * \real{0.3491}}
  >{\raggedright\arraybackslash}p{(\linewidth - 4\tabcolsep) * \real{0.5283}}@{}}
\toprule\noalign{}
\begin{minipage}[b]{\linewidth}\raggedright
Total Score
\end{minipage} & \begin{minipage}[b]{\linewidth}\raggedright
Interpretation
\end{minipage} & \begin{minipage}[b]{\linewidth}\raggedright
Next Step
\end{minipage} \\
\midrule\noalign{}
\endhead
\bottomrule\noalign{}
\endlastfoot
8-9 & Excellent work & Move on to next project \\
6-7 & Good work with room for improvement &
Move on; instructor may suggest revisiting one element \\
4-5 & Meets basic expectations &
Resubmission of specific weak areas recommended \\
2-3 & Does not meet expectations & Resubmission required \\
0-1 & Missing major deliverables &
Meet with instructor; create a completion plan \\
\end{longtable}
}

\subsubsection*{Resubmission Policy}\label{docs__pandoc__latex__src__syllabus.md__resubmission-policy}

Students may resubmit any project as many times as they need to improve their score. Resubmissions must include a one-paragraph explanation of what was changed and why. The resubmission score replaces the original score.

\subsection*{Quick Links to Essential Tools \& Setup}\label{docs__pandoc__latex__src__syllabus.md__quick-links-to-essential-tools--setup}

\subsubsection*{Core Design Toolchain}\label{docs__pandoc__latex__src__syllabus.md__core-design-toolchain}

OpenSCAD

\begin{itemize}
\tightlist
\item
  \href{https://openscad.org/downloads.html}{OpenSCAD Download} - Free, cross-platform CAD (all major OS)
\item
  \href{https://openscad.org/documentation.html}{OpenSCAD Documentation} - Official reference
\item
  \hyperref[docs__pandoc__latex__src__3dmake_appendices_complete.md__appendix-a]{OpenSCAD Cheat Sheet} - Quick syntax reference
\end{itemize}

3DMake

\begin{itemize}
\tightlist
\item
  \href{https://github.com/tdeck/3dmake}{3DMake Documentation \& Installation} - Command-line build tool for OpenSCAD
\item
  \hyperref[docs__pandoc__latex__src__3dmake_appendices_complete.md__appendix-a]{3DMake CLI Cheat Sheet} - Complete command reference
\item
  \hyperref[docs__pandoc__latex__src__3dmake_appendices_complete.md__appendix-a]{3dMake Quick Reference} - Command and workflow reference
\end{itemize}

Editors

\begin{itemize}
\tightlist
\item
  \href{https://code.visualstudio.com/}{VS Code Download} - Free, screen-reader-accessible code editor
\item
  \href{https://notepad-plus-plus.org/}{Notepad++ Download} - Free, lightweight Windows editor
\item
  \hyperref[docs__pandoc__latex__src__setup_accessibility_complete.md__editor-setup]{Editor Selection and Setup Guide} - Accessibility-focused setup guide
\end{itemize}

\subsubsection*{Screen Reader \& Accessibility}\label{docs__pandoc__latex__src__syllabus.md__screen-reader--accessibility}

Screen Readers

\begin{itemize}
\tightlist
\item
  \href{https://www.nvaccess.org/}{NVDA Download} - Free, open-source screen reader (Windows)
\item
  \href{https://www.freedomscientific.com/products/software/jaws/}{JAWS Screen Reader} - Commercial screen reader (Windows, macOS)
\item
  \href{https://help.gnome.org/users/orca/stable/}{Orca Screen Reader} - Built-in screen reader for Linux/GNOME desktop
\item
  \href{https://support.microsoft.com/en-us/windows/complete-guide-to-narrator-e4397a0d-ef4f-b386-d8ae-c172f109bdb1}{Windows Narrator} - Built-in Windows screen reader
\item
  \href{https://yourdolphin.com/ScreenReader-Training}{Dolphin SuperNova} - Commercial screen reader (Windows)
\end{itemize}

Note: This curriculum is designed to work with all major screen readers. The workflows are tested primarily with NVDA and JAWS but apply to all screen readers listed above.

Accessibility Configuration

\begin{itemize}
\tightlist
\item
  \hyperref[docs__pandoc__latex__src__setup_accessibility_complete.md__accessibility-guide]{Screen Reader Accessibility Guide} - Comprehensive terminal accessibility guide
\item
  \hyperref[docs__pandoc__latex__src__3dmake_all_lessons.md__3dmake-foundation--design--complete-course-compendium]{Screen Reader Coding Tips (NVDA \& JAWS)} - Keyboard shortcuts and configuration
\item
  \hyperref[docs__pandoc__latex__src__3dmake_all_lessons.md__3dmake-foundation--design--complete-course-compendium]{VSCode Setup Guide} - Accessibility-focused editor configuration
\item
  \href{docs\%5Cpandoc\%5Clatex\%5Csrc\%5CGitBash_Complete_Curriculum.md\#gb-accessibility-guide}{Git Bash Screen Reader Guide} - NVDA and JAWS configuration for Git Bash
\end{itemize}

\subsubsection*{Slicing Software}\label{docs__pandoc__latex__src__syllabus.md__slicing-software}

\begin{itemize}
\tightlist
\item
  \href{https://www.prusa3d.com/page/prusaslicer424/}{PrusaSlicer}
\item
  \href{https://bambulab.com/en/download/studio}{Bambu Studio}
\item
  \href{https://ultimaker.com/software/ultimaker-cura}{Cura}
\item
  \href{https://github.com/SoftFever/OrcaSlicer/releases}{OrcaSlicer}
\item
  \hyperref[docs__pandoc__latex__src__3dmake_appendices_complete.md__appendix-b]{Appendix A: Comprehensive Slicing Guide} - Detailed setup guides for all major slicers
\end{itemize}

\subsubsection*{Supplemental Textbooks \& Learning Resources}\label{docs__pandoc__latex__src__syllabus.md__supplemental-textbooks--learning-resources}

\textbf{Beginner-Friendly Textbooks:}

\begin{itemize}
\tightlist
\item
  \href{https://nostarch.com/programmingopenscad}{Programming with OpenSCAD: A Beginner\textquotesingle s Guide to Coding 3D-Printable Objects} by Justin Gohde and Marius Kintel - Comprehensive reference covering OpenSCAD syntax, geometry concepts, and design patterns. Available in EPUB format (screen reader accessible).
\item
  \href{https://www.packtpub.com/en-us/product/simplifying-3d-printing-with-openscad-9781801813174}{Simplifying 3D Printing with OpenSCAD} by Alicia Noors - Focused on practical workflows, optimization, and real-world printing scenarios. Available in EPUB/PDF formats.
\item
  \href{https://www.apress.com/gp/book/9781484213292}{OpenSCAD for 3D Printing} by Al Williams - Hands-on introduction to OpenSCAD with practical 3D printing projects. Covers basics through intermediate topics.
\item
  \href{https://www.amazon.com/Mastering-OpenSCAD-Jochen-Kerdels/dp/1803243538}{Mastering OpenSCAD} by Jochen Kerdels - Project-based learning approach with emphasis on parametric design and reusable code modules.
\end{itemize}

\textbf{Free Online Resources:}

\begin{itemize}
\tightlist
\item
  \href{https://en.wikibooks.org/wiki/OpenSCAD_User_Manual}{Official OpenSCAD Documentation} - Complete language reference and user manual on Wikibooks. Fully accessible and searchable.
\item
  \href{https://openscad.org/documentation.html}{OpenSCAD Tutorial} - Official beginner tutorial covering core concepts and syntax.
\item
  \href{https://openscad.org/cheatsheet/}{OpenSCAD Cheat Sheet (Official)} - Quick reference for all OpenSCAD commands and functions.
\end{itemize}

\textbf{Curriculum-Specific Resources:}

\begin{itemize}
\tightlist
\item
  \href{https://programmingwithopenscad.github.io/learning.html}{Programming with OpenSCAD Companion Resources} - Practice worksheets for visualization, decomposition exercises, vocabulary building, and assessment.
\item
  \href{https://programmingwithopenscad.github.io/quick-reference.html}{Visual Quick Reference} - Command syntax guides and geometry reference.
\item
  \href{https://github.com/ProgrammingWithOpenSCAD/CodeSolutions}{Code Solutions Repository} - Working OpenSCAD examples organized by topic (3D shapes, transformations, loops, modules, if-statements, advanced techniques).
\end{itemize}

\textbf{Community Resources:}

\begin{itemize}
\tightlist
\item
  \href{https://www.thingiverse.com/search?q=openscad&type=things}{Thingiverse OpenSCAD Designs} - Community-created parametric designs with source code for study and modification.
\item
  \href{https://forum.openscad.org/}{OpenSCAD Forums} - Active community support for troubleshooting and learning.
\item
  \href{https://www.youtube.com/results?search_query=openscad+tutorial+beginner}{OpenSCAD Tutorials (YouTube)} - Video tutorials (note: supplemental to text-based curriculum; transcripts recommended for accessibility).
\end{itemize}

\subsection*{Local Resources: Utah Makerspaces \& Community Printing}\label{docs__pandoc__latex__src__syllabus.md__local-resources-utah-makerspaces--community-printing}

\subsubsection*{Public Library Make Spaces}\label{docs__pandoc__latex__src__syllabus.md__public-library-make-spaces}

Salt Lake City Public Library

\begin{itemize}
\tightlist
\item
  \href{https://services.slcpl.org/creativelab}{SLC Public Creative Lab} - Main Library (Level 1)

  \begin{itemize}
  \tightlist
  \item
    Hardware: Prusa i3 MK3, LulzBot Taz 5, Elegoo Mars 2 (resin)
  \item
    Pricing: Free for prints under 6 hours; \$0.50/hr + material cost otherwise
  \end{itemize}
\end{itemize}

Salt Lake County Library System

\begin{itemize}
\tightlist
\item
  \href{https://www.slcolibrary.org/what-we-have/create}{County Library "Create" Spaces}

  \begin{itemize}
  \tightlist
  \item
    Hardware: Flashforge Adventurer 5M Pro, LulzBot Workhorse, laser cutters
  \item
    Pricing: \$0.06 per gram of filament used
  \end{itemize}
\end{itemize}

\subsubsection*{Makerspaces \& Community Centers}\label{docs__pandoc__latex__src__syllabus.md__makerspaces--community-centers}

Make Salt Lake

\begin{itemize}
\tightlist
\item
  Location: 663 W 100 S, Salt Lake City, UT 84101
\item
  Website: \url{https://makesaltlake.org/}
\item
  Equipment: Full metal shop, CNC machines, large-scale FDM and resin printing
\end{itemize}

University of Utah Maker Spaces

\begin{itemize}
\tightlist
\item
  \href{https://www.theblackbookofinnovation.com/lassonde-studios}{Lassonde Studios}
\item
  \href{https://lib.utah.edu/protospace.php}{Marriott Library ProtoSpace}
\end{itemize}

\subsection*{Troubleshooting \& Getting Help}\label{docs__pandoc__latex__src__syllabus.md__troubleshooting--getting-help}

\subsubsection*{Quick Reference Materials}\label{docs__pandoc__latex__src__syllabus.md__quick-reference-materials}

\textbf{Setup \& Configuration:}

\begin{itemize}
\tightlist
\item
  \hyperref[docs__pandoc__latex__src__3dmake_all_lessons.md__3dmake-foundation--design--complete-course-compendium]{3dMake Setup Guide} - Installation walkthrough
\item
  \hyperref[docs__pandoc__latex__src__3dmake_all_lessons.md__3dmake-foundation--design--complete-course-compendium]{VSCode Setup Guide} - Accessibility-focused editor configuration
\item
  \hyperref[docs__pandoc__latex__src__3dmake_all_lessons.md__3dmake-foundation--design--complete-course-compendium]{mdBook Navigation Guide} - How to navigate this curriculum
\end{itemize}

\textbf{Syntax \& Commands:}

\begin{itemize}
\tightlist
\item
  \hyperref[docs__pandoc__latex__src__3dmake_appendices_complete.md__appendix-a]{OpenSCAD Cheat Sheet} - Complete syntax quick-reference
\item
  \hyperref[docs__pandoc__latex__src__3dmake_appendices_complete.md__appendix-a]{3DMake CLI Cheat Sheet} - All 3dm commands and options
\end{itemize}

\textbf{Materials \& Safety:}

\begin{itemize}
\tightlist
\item
  \hyperref[docs__pandoc__latex__src__3dmake_appendices_complete.md__appendix-c]{Filament Comparison Table} - Material selection guide
\item
  \hyperref[docs__pandoc__latex__src__3dmake_foundation__lessons_3dmake_5__lessons_3dmake_5.md__safety-protocols-and-the-physical-fabrication-interface]{Safety Checklist} - Pre-print safety procedures
\item
  \hyperref[docs__pandoc__latex__src__3dmake_foundation__lessons_3dmake_5__lessons_3dmake_5.md__safety-protocols-and-the-physical-fabrication-interface]{Maintenance Log} - Printer maintenance tracking
\end{itemize}

\subsubsection*{When You\textquotesingle re Stuck}\label{docs__pandoc__latex__src__syllabus.md__when-youre-stuck}

\textbf{For Technical Issues:}

\begin{enumerate}
\tightlist
\item
  Check \hyperref[docs__pandoc__latex__src__backmatter_complete.md__troubleshooting]{Common Issues and Solutions} - Comprehensive troubleshooting guide
\item
  Review \hyperref[docs__pandoc__latex__src__backmatter_complete.md__troubleshooting]{Diagnostic Checklist} - Systematic problem diagnosis
\item
  Consult \hyperref[docs__pandoc__latex__src__backmatter_complete.md__troubleshooting]{Technical Reference} - Advanced systems architecture and professional reference
\item
  Post in \href{https://discord.gg/F2Nx2VxTB7}{OpenSCAD Discord} or \href{https://www.reddit.com/r/openscad/}{Reddit r/openscad}
\item
  Visit your local makerspace for hands-on support
\end{enumerate}

\textbf{For Accessibility Support:}

\begin{itemize}
\tightlist
\item
  Contact your NVDA/JAWS vendor directly for technical assistance
\item
  Refer to the \hyperref[docs__pandoc__latex__src__setup_accessibility_complete.md__accessibility-guide]{Screen Reader Accessibility Guide}
\item
  Review \hyperref[docs__pandoc__latex__src__3dmake_all_lessons.md__3dmake-foundation--design--complete-course-compendium]{Screen Reader Coding Tips (NVDA \& JAWS)}
\item
  Check the \href{docs\%5Cpandoc\%5Clatex\%5Csrc\%5CGitBash_Complete_Curriculum.md\#gb-accessibility-guide}{Git Bash Screen Reader Guide} if using Git Bash
\end{itemize}

\textbf{For Design \& Concept Questions:}

\begin{itemize}
\tightlist
\item
  Review \hyperref[docs__pandoc__latex__src__backmatter_complete.md__student-glossary]{Student Glossary} for terminology clarification
\item
  Check project-specific guides in lesson folders
\item
  Consult the \hyperref[docs__pandoc__latex__src__3dmake_appendices_complete.md__appendix-d]{Master Rubric} for assessment expectations
\end{itemize}

\section{Setup and Accessibility Fundamentals --- Complete Guide}\label{docs__pandoc__latex__src__setup_accessibility_complete.md__setup-and-accessibility-fundamentals--complete-guide}

Everything you need to configure an accessible working environment before starting the command-line or 3dMake curriculum. Read these sections in order the first time through; return to individual sections as reference when needed.

\begin{center}\rule{0.5\linewidth}{0.5pt}\end{center}

\subsection*{Contents}\label{docs__pandoc__latex__src__setup_accessibility_complete.md__contents}

{\def\LTcaptype{none} % do not increment counter
\begin{longtable}[]{@{}
  >{\raggedright\arraybackslash}p{(\linewidth - 2\tabcolsep) * \real{0.4452}}
  >{\raggedright\arraybackslash}p{(\linewidth - 2\tabcolsep) * \real{0.5548}}@{}}
\toprule\noalign{}
\begin{minipage}[b]{\linewidth}\raggedright
Section
\end{minipage} & \begin{minipage}[b]{\linewidth}\raggedright
Title
\end{minipage} \\
\midrule\noalign{}
\endhead
\bottomrule\noalign{}
\endlastfoot
\hyperref[docs__pandoc__latex__src__setup_accessibility_complete.md__screen-reader-choice]{Part 1: Choosing a Screen Reader}
&
Compare NVDA, JAWS, Windows Narrator, and Dolphin SuperNova for command-line work \\
\hyperref[docs__pandoc__latex__src__setup_accessibility_complete.md__accessibility-guide]{Part 2: Screen Reader Accessibility Guide}
& Key commands and terminal strategies for all four screen readers \\
\hyperref[docs__pandoc__latex__src__setup_accessibility_complete.md__braille-displays]{Part 3: Braille Display Setup}
& Connecting and configuring a braille display with any screen reader \\
\hyperref[docs__pandoc__latex__src__setup_accessibility_complete.md__editor-setup]{Part 4: Editor Selection and Setup}
&
Choosing and configuring VS Code or other editors for accessible OpenSCAD work \\
\end{longtable}
}

\textbf{Recommended reading order:} Part 1 → Part 2 → Part 3 (if applicable) → Part 4. Parts 3 and 4 can be deferred until you need them.

\begin{center}\rule{0.5\linewidth}{0.5pt}\end{center}

\protect\phantomsection\label{docs__pandoc__latex__src__setup_accessibility_complete.md__screen-reader-choice}{}

\section*{Part 1: Choosing a Screen Reader for Windows Command-Line Work}\label{docs__pandoc__latex__src__setup_accessibility_complete.md__part-1-choosing-a-screen-reader-for-windows-command-line-work}

\begin{center}\rule{0.5\linewidth}{0.5pt}\end{center}

Screen readers provide speech output and, where supported, braille output that make text-based interfaces accessible to people who are blind or have low vision. In command-line environments they convert terminal text, prompts, and keyboard navigation into audible or tactile feedback so users can work independently in shells such as Windows Terminal, Command Prompt, PowerShell, and Git Bash.

Choosing a screen reader is a personal decision. There is no single reader that is universally better or worse for every user. Preferences depend on workflow, training, budget, and comfort with customization. Each product has genuine strengths and trade-offs. If possible, try the available options and pick what feels most natural for you.

\begin{center}\rule{0.5\linewidth}{0.5pt}\end{center}

\subsection*{The Four Screen Readers Covered in This Curriculum}\label{docs__pandoc__latex__src__setup_accessibility_complete.md__the-four-screen-readers-covered-in-this-curriculum}

\subsubsection*{NVDA (NonVisual Desktop Access)}\label{docs__pandoc__latex__src__setup_accessibility_complete.md__nvda-nonvisual-desktop-access}

NVDA is a free, open-source screen reader developed by NV Access. It is one of the most widely used screen readers in the world.

\textbf{Advantages:}

\begin{itemize}
\tightlist
\item
  Free to download and use; supported by donations.
\item
  Actively developed with frequent updates.
\item
  Excellent support for Windows terminals (Command Prompt, PowerShell, Windows Terminal, Git Bash).
\item
  Extensible through a rich add-on ecosystem.
\item
  Strong community support, forums, and troubleshooting resources.
\item
  Good braille display support via built-in braille drivers and the BRLTTY/LibLouis library.
\end{itemize}

\textbf{Disadvantages:}

\begin{itemize}
\tightlist
\item
  Advanced paid support is community-driven rather than available from a commercial vendor helpdesk.
\item
  Some edge-case compatibility differences compared to commercial readers.
\end{itemize}

\textbf{Cost:} Free (donations welcomed at \url{https://www.nvaccess.org/}).

\textbf{Trial:} Not applicable --- NVDA is always free.

\textbf{Download:} \url{https://www.nvaccess.org/download/}

\begin{center}\rule{0.5\linewidth}{0.5pt}\end{center}

\subsubsection*{JAWS (Job Access With Speech)}\label{docs__pandoc__latex__src__setup_accessibility_complete.md__jaws-job-access-with-speech}

JAWS is a commercial screen reader developed by Freedom Scientific (part of Vispero). It is one of the most feature-rich and widely deployed screen readers in enterprise and educational settings.

\textbf{Advantages:}

\begin{itemize}
\tightlist
\item
  Mature product with decades of development and broad application compatibility.
\item
  Extensive scripting system allows deep customization for specific applications.
\item
  Vendor-provided paid support, training, and documentation.
\item
  Advanced braille display support and configuration.
\item
  Widely used in workplaces and schools, so shared knowledge base is large.
\end{itemize}

\textbf{Disadvantages:}

\begin{itemize}
\tightlist
\item
  Commercial product with significant licensing cost.
\item
  JAWS demo mode limits session length to approximately 40 minutes before requiring a computer restart or session reset to continue. This can be disruptive during extended practice sessions.
\item
  Some advanced customization (JAWS scripts) requires a learning curve.
\end{itemize}

\textbf{Cost:} Commercial. Typical U.S. pricing is several hundred dollars per year for a software maintenance agreement (SMA), or a higher one-time perpetual license fee. Check \url{https://www.freedomscientific.com/products/software/jaws/} for current pricing.

\textbf{Trial:} JAWS offers a time-limited demonstration mode (approximately 40 minutes per session). After the demo period expires, you must restart the computer to continue using JAWS without a license. This applies to the full JAWS installer --- there is no separate shorter trial download.

\textbf{Download:} \url{https://www.freedomscientific.com/downloads/jaws/}

\begin{center}\rule{0.5\linewidth}{0.5pt}\end{center}

\subsubsection*{Windows Narrator}\label{docs__pandoc__latex__src__setup_accessibility_complete.md__windows-narrator}

Windows Narrator is Microsoft\textquotesingle s built-in screen reader, included with Windows 10 and Windows 11 at no extra cost. It has improved significantly with recent Windows releases.

\textbf{Advantages:}

\begin{itemize}
\tightlist
\item
  Already installed on every Windows 10 and Windows 11 computer --- no download or installation needed.
\item
  No cost.
\item
  Simple to start: press \textbf{Windows+Ctrl+Enter} to toggle Narrator on and off.
\item
  Reasonable support for Windows Terminal, Command Prompt, and PowerShell.
\item
  Integrates tightly with Windows accessibility features.
\end{itemize}

\textbf{Disadvantages:}

\begin{itemize}
\tightlist
\item
  Fewer customization options than JAWS or NVDA for complex command-line workflows.
\item
  Braille display support exists but is more limited than dedicated screen readers.
\item
  Less powerful for advanced scripting, code editing, or rapid navigation in large terminal outputs.
\item
  Less community documentation for terminal-specific workflows compared to NVDA and JAWS.
\end{itemize}

\textbf{Cost:} Free (included with Windows).

\textbf{Trial:} Not applicable --- Narrator is always available on Windows.

\textbf{Enable:} Settings \textgreater{} Accessibility \textgreater{} Narrator, or press \textbf{Windows+Ctrl+Enter}.

\begin{center}\rule{0.5\linewidth}{0.5pt}\end{center}

\subsubsection*{Dolphin SuperNova}\label{docs__pandoc__latex__src__setup_accessibility_complete.md__dolphin-supernova}

Dolphin SuperNova (sometimes referred to as "Dolphin Screen Reader" or simply "SuperNova") is a commercial product from Dolphin Computer Access. It is available in several editions: Magnifier, Magnifier \& Speech, and Screen Reader. The Screen Reader and Magnifier \& Speech editions provide full screen reader functionality.

\textbf{Advantages:}

\begin{itemize}
\tightlist
\item
  Commercial product with vendor support, training options, and documentation.
\item
  Offers combined magnification and speech in a single product --- useful for users with low vision who benefit from both.
\item
  Good braille display support.
\item
  Works with Windows terminals (Command Prompt, PowerShell, Windows Terminal).
\item
  Full-featured 30-day trial license available for evaluation before purchase.
\end{itemize}

\textbf{Disadvantages:}

\begin{itemize}
\tightlist
\item
  Commercial licensing cost.
\item
  Less widely used than NVDA or JAWS, so community resources and third-party documentation are more limited.
\item
  Configuration for some terminal workflows may require additional setup.
\end{itemize}

\textbf{Cost:} Commercial. Pricing varies by edition (Magnifier only, Magnifier \& Speech, or Screen Reader). Check \url{https://yourdolphin.com/supernova/} for current pricing. Expect costs broadly similar to other commercial screen readers.

\textbf{Trial:} Dolphin provides a full-featured 30-day trial license. This is more generous than JAWS\textquotesingle s per-session demo and allows uninterrupted evaluation over a full month.

\textbf{Download:} \url{https://yourdolphin.com/supernova/}

\begin{center}\rule{0.5\linewidth}{0.5pt}\end{center}

\subsection*{Choosing Between Them for Command-Line Work}\label{docs__pandoc__latex__src__setup_accessibility_complete.md__choosing-between-them-for-command-line-work}

Use this decision guide to narrow your choice:

\textbf{If budget is the primary constraint:}\\
Start with NVDA or Windows Narrator. Both are free. NVDA is generally more capable and better documented for terminal work. Windows Narrator requires no installation and is a good quick-start option.

\textbf{If you already own or have access to JAWS or Dolphin through school or work:}\\
Use what you have. Both work well for this curriculum.

\textbf{If you have low vision (partial sight) in addition to using a screen reader:}\\
Consider Dolphin SuperNova, which combines magnification and speech in one product, or use Windows Magnifier alongside NVDA.

\textbf{If your organization requires vendor-supported software:}\\
JAWS and Dolphin SuperNova both offer commercial support contracts.

\textbf{If you are unsure:}\\
Install NVDA (free, no risk) and try a few lessons. You can always switch later, and the terminal commands you learn work identically regardless of which screen reader you use.

\begin{center}\rule{0.5\linewidth}{0.5pt}\end{center}

\subsection*{Terminal Command Reference by Screen Reader}\label{docs__pandoc__latex__src__setup_accessibility_complete.md__terminal-command-reference-by-screen-reader}

All four screen readers can perform the same tasks in the terminal. The table below shows the key actions and how to perform them with each reader. Exact key names can vary by keyboard layout and screen reader configuration.

{\def\LTcaptype{none} % do not increment counter
\begin{longtable}[]{@{}
  >{\raggedright\arraybackslash}p{(\linewidth - 8\tabcolsep) * \real{0.1980}}
  >{\raggedright\arraybackslash}p{(\linewidth - 8\tabcolsep) * \real{0.1980}}
  >{\raggedright\arraybackslash}p{(\linewidth - 8\tabcolsep) * \real{0.2030}}
  >{\raggedright\arraybackslash}p{(\linewidth - 8\tabcolsep) * \real{0.1827}}
  >{\raggedright\arraybackslash}p{(\linewidth - 8\tabcolsep) * \real{0.2183}}@{}}
\toprule\noalign{}
\begin{minipage}[b]{\linewidth}\raggedright
Task
\end{minipage} & \begin{minipage}[b]{\linewidth}\raggedright
NVDA
\end{minipage} & \begin{minipage}[b]{\linewidth}\raggedright
JAWS
\end{minipage} & \begin{minipage}[b]{\linewidth}\raggedright
Windows Narrator
\end{minipage} & \begin{minipage}[b]{\linewidth}\raggedright
Dolphin SuperNova
\end{minipage} \\
\midrule\noalign{}
\endhead
\bottomrule\noalign{}
\endlastfoot
\textbf{Read current line or prompt} & NVDA+Home (reads current line) &
Insert+Up Arrow (re-reads current line) &
Narrator+D (reads current line) & CapsLock+L (reads current line) \\
\textbf{Read all output from here} &
NVDA+Down Arrow (read to end of screen) &
Insert+Ctrl+Down (read to end of screen) & Narrator+R (read from cursor)
& CapsLock+Numpad Plus (say all) \\
\textbf{Move to previous line of output} & Up Arrow in review mode &
Insert+Up Arrow (line by line) & Narrator+Up Arrow (scan mode) &
CapsLock+Numpad 8 \\
\textbf{Move to next line of output} & Down Arrow in review mode &
Insert+Down Arrow (line by line) & Narrator+Down Arrow (scan mode) &
CapsLock+Numpad 2 \\
\textbf{Jump to end of screen/buffer} & NVDA+End & Insert+End &
Narrator+End & CapsLock+Page Down \\
\textbf{Jump to start of screen/buffer} & NVDA+Home & Insert+Home &
Narrator+Home & CapsLock+Page Up \\
\textbf{Recall last command (shell history)} &
Up Arrow (shell native; SR reads it) &
Up Arrow (shell native; SR reads it) &
Up Arrow (shell native; SR reads it) &
Up Arrow (shell native; SR reads it) \\
\textbf{Stop reading} & Ctrl & Ctrl & Ctrl & Ctrl \\
\textbf{Open screen reader settings} & NVDA+N \textgreater{} Preferences
& Insert+F3 \textgreater{} Options/Settings &
Windows+Ctrl+N (Narrator settings) &
CapsLock+SpaceBar (SuperNova control panel) \\
\end{longtable}
}

\begin{quote}
\textbf{Notes on modifier keys:}

\begin{itemize}
\tightlist
\item
  \textbf{NVDA key} is Insert by default; can be changed to CapsLock in NVDA Preferences \textgreater{} Keyboard.
\item
  \textbf{JAWS key} is Insert by default; can be changed to CapsLock in JAWS Settings Center.
\item
  \textbf{Narrator key} is CapsLock or Insert; configurable in Settings \textgreater{} Accessibility \textgreater{} Narrator \textgreater{} Keyboard shortcuts.
\item
  \textbf{Dolphin key} is CapsLock by default; configurable in SuperNova Control Panel \textgreater{} Keyboard.
\end{itemize}
\end{quote}

\begin{center}\rule{0.5\linewidth}{0.5pt}\end{center}

\subsection*{Important Notes for This Curriculum}\label{docs__pandoc__latex__src__setup_accessibility_complete.md__important-notes-for-this-curriculum}

\begin{enumerate}
\item
  \textbf{All four screen readers work with all three shells} (CMD, PowerShell, Git Bash) covered in this curriculum. You do not need to change screen readers based on which shell pathway you choose.
\item
  \textbf{The best practice for all four screen readers} when dealing with long terminal output is to redirect it to a file and open it in Notepad:

  \begin{lstlisting}[style=Alabaster]
  command > output.txt
  notepad output.txt

  \end{lstlisting}

  This works in CMD, PowerShell, and Git Bash, and gives you a stable document that all four screen readers can read comfortably.
\item
  \textbf{Braille display users} can use a braille display alongside any of these four screen readers. See the Braille Display Setup guide for detailed instructions.
\item
  \textbf{The screen reader tips in each lesson} of this curriculum are written for NVDA, JAWS, Windows Narrator, and Dolphin SuperNova. If the specific key commands differ from what you hear when practicing, check your screen reader\textquotesingle s documentation for the equivalent command --- the underlying concept is always the same.
\end{enumerate}

\begin{center}\rule{0.5\linewidth}{0.5pt}\end{center}

\subsection*{Additional Resources}\label{docs__pandoc__latex__src__setup_accessibility_complete.md__additional-resources}

{\def\LTcaptype{none} % do not increment counter
\begin{longtable}[]{@{}
  >{\raggedright\arraybackslash}p{(\linewidth - 4\tabcolsep) * \real{0.0717}}
  >{\raggedright\arraybackslash}p{(\linewidth - 4\tabcolsep) * \real{0.4979}}
  >{\raggedright\arraybackslash}p{(\linewidth - 4\tabcolsep) * \real{0.4304}}@{}}
\toprule\noalign{}
\begin{minipage}[b]{\linewidth}\raggedright
Screen Reader
\end{minipage} & \begin{minipage}[b]{\linewidth}\raggedright
Official Documentation
\end{minipage} & \begin{minipage}[b]{\linewidth}\raggedright
Support / Help
\end{minipage} \\
\midrule\noalign{}
\endhead
\bottomrule\noalign{}
\endlastfoot
NVDA & \url{https://www.nvaccess.org/} &
\url{https://www.nvaccess.org/download/} \\
JAWS & \url{https://www.freedomscientific.com/products/software/jaws/} &
\url{https://www.freedomscientific.com/downloads/jaws/} \\
Windows Narrator & \url{https://support.microsoft.com/narrator} &
\url{https://support.microsoft.com/narrator} \\
Dolphin SuperNova & \url{https://yourdolphin.com/supernova/} &
\url{https://yourdolphin.com/supernova/} \\
\end{longtable}
}

\begin{center}\rule{0.5\linewidth}{0.5pt}\end{center}

\protect\phantomsection\label{docs__pandoc__latex__src__setup_accessibility_complete.md__accessibility-guide}{}

\section*{Part 2: Screen Reader Accessibility Guide for Command-Line Terminals}\label{docs__pandoc__latex__src__setup_accessibility_complete.md__part-2-screen-reader-accessibility-guide-for-command-line-terminals}

\begin{center}\rule{0.5\linewidth}{0.5pt}\end{center}

\textbf{Target Users:} NVDA, JAWS, Windows Narrator, and Dolphin SuperNova users using PowerShell, CMD, or Git Bash\\
\textbf{Last Updated:} February 2026

This guide is used throughout the Command-Line Foundation curriculum (PowerShell, CMD, and Git Bash pathways) to help screen reader users navigate and work efficiently with the terminal.

\begin{center}\rule{0.5\linewidth}{0.5pt}\end{center}

\subsection*{Table of Contents}\label{docs__pandoc__latex__src__setup_accessibility_complete.md__table-of-contents}

\begin{enumerate}
\tightlist
\item
  \hyperref[docs__pandoc__latex__src__setup_accessibility_complete.md__which-terminal]{Which Terminal Should I Use?}
\item
  \hyperref[docs__pandoc__latex__src__setup_accessibility_complete.md__getting-started]{Getting Started with Screen Readers}
\item
  \hyperref[docs__pandoc__latex__src__setup_accessibility_complete.md__nvda-tips]{NVDA-Specific Tips}
\item
  \hyperref[docs__pandoc__latex__src__setup_accessibility_complete.md__jaws-tips]{JAWS-Specific Tips}
\item
  \hyperref[docs__pandoc__latex__src__setup_accessibility_complete.md__narrator-tips]{Windows Narrator-Specific Tips}
\item
  \hyperref[docs__pandoc__latex__src__setup_accessibility_complete.md__dolphin-tips]{Dolphin SuperNova-Specific Tips}
\item
  \hyperref[docs__pandoc__latex__src__setup_accessibility_complete.md__general-terminal]{General Terminal Accessibility}
\item
  \hyperref[docs__pandoc__latex__src__setup_accessibility_complete.md__long-output]{Working with Long Output}
\item
  \hyperref[docs__pandoc__latex__src__setup_accessibility_complete.md__shortcuts]{Keyboard Shortcuts Reference}
\item
  \hyperref[docs__pandoc__latex__src__setup_accessibility_complete.md__shell-differences]{Shell-Specific Differences}
\item
  \hyperref[docs__pandoc__latex__src__setup_accessibility_complete.md__troubleshooting]{Troubleshooting}
\end{enumerate}

\begin{center}\rule{0.5\linewidth}{0.5pt}\end{center}

\protect\phantomsection\label{docs__pandoc__latex__src__setup_accessibility_complete.md__which-terminal}{}

\subsection*{Which Terminal Should I Use?}\label{docs__pandoc__latex__src__setup_accessibility_complete.md__which-terminal-should-i-use}

All three command-line shells (PowerShell, CMD, Git Bash) work well with screen readers. Choose the one that best fits your setup:

{\def\LTcaptype{none} % do not increment counter
\begin{longtable}[]{@{}
  >{\raggedright\arraybackslash}p{(\linewidth - 6\tabcolsep) * \real{0.2632}}
  >{\raggedright\arraybackslash}p{(\linewidth - 6\tabcolsep) * \real{0.2105}}
  >{\raggedright\arraybackslash}p{(\linewidth - 6\tabcolsep) * \real{0.2316}}
  >{\raggedright\arraybackslash}p{(\linewidth - 6\tabcolsep) * \real{0.2947}}@{}}
\toprule\noalign{}
\begin{minipage}[b]{\linewidth}\raggedright
Feature
\end{minipage} & \begin{minipage}[b]{\linewidth}\raggedright
PowerShell
\end{minipage} & \begin{minipage}[b]{\linewidth}\raggedright
CMD
\end{minipage} & \begin{minipage}[b]{\linewidth}\raggedright
Git Bash
\end{minipage} \\
\midrule\noalign{}
\endhead
\bottomrule\noalign{}
\endlastfoot
\textbf{Platform} & Windows only & Windows only &
Windows, macOS, Linux \\
\textbf{Screen Reader Support} & Excellent & Excellent & Excellent \\
\textbf{Learning Curve} & Moderate & Gentle & Moderate-Steep \\
\textbf{Use Case} & Modern automation & Familiar simplicity &
Cross-platform/Git workflows \\
\textbf{Recommended For} & 3D design automation & Getting started easily
& Software developers \\
\end{longtable}
}

\textbf{Key Point:} Choose one pathway and complete it fully. All three teach the same skills and prepare you equally well for 3dMake work. Don\textquotesingle t switch mid-curriculum.

\begin{center}\rule{0.5\linewidth}{0.5pt}\end{center}

\protect\phantomsection\label{docs__pandoc__latex__src__setup_accessibility_complete.md__getting-started}{}

\subsection*{Getting Started with Screen Readers}\label{docs__pandoc__latex__src__setup_accessibility_complete.md__getting-started-with-screen-readers}

\subsubsection*{Which Screen Reader Should I Use?}\label{docs__pandoc__latex__src__setup_accessibility_complete.md__which-screen-reader-should-i-use}

NVDA, JAWS, Windows Narrator, and Dolphin SuperNova all work with the three shells covered in this curriculum. Here is a brief comparison:

{\def\LTcaptype{none} % do not increment counter
\begin{longtable}[]{@{}
  >{\raggedright\arraybackslash}p{(\linewidth - 8\tabcolsep) * \real{0.1869}}
  >{\raggedright\arraybackslash}p{(\linewidth - 8\tabcolsep) * \real{0.2056}}
  >{\raggedright\arraybackslash}p{(\linewidth - 8\tabcolsep) * \real{0.2056}}
  >{\raggedright\arraybackslash}p{(\linewidth - 8\tabcolsep) * \real{0.2336}}
  >{\raggedright\arraybackslash}p{(\linewidth - 8\tabcolsep) * \real{0.1682}}@{}}
\toprule\noalign{}
\begin{minipage}[b]{\linewidth}\raggedright
Feature
\end{minipage} & \begin{minipage}[b]{\linewidth}\raggedright
NVDA
\end{minipage} & \begin{minipage}[b]{\linewidth}\raggedright
JAWS
\end{minipage} & \begin{minipage}[b]{\linewidth}\raggedright
Windows Narrator
\end{minipage} & \begin{minipage}[b]{\linewidth}\raggedright
Dolphin SuperNova
\end{minipage} \\
\midrule\noalign{}
\endhead
\bottomrule\noalign{}
\endlastfoot
\textbf{Cost} & Free (open-source) & Commercial (paid) &
Free (built into Windows) & Commercial (paid) \\
\textbf{Installation} & Simple download & Complex but thorough &
Already installed & Standard installer \\
\textbf{Terminal Support} & Excellent (all shells) &
Excellent (all shells) & Good (all shells) & Good (all shells) \\
\textbf{Learning Curve} & Gentle & Steeper & Minimal & Moderate \\
\textbf{Customization} & Good (add-ons) & Extensive (scripts) & Limited
& Good \\
\end{longtable}
}

\textbf{Recommendation:} If you are new to screen readers, start with NVDA (free, well-documented, good community support) or Windows Narrator (already on your computer, no installation needed). JAWS and Dolphin SuperNova are excellent choices if you already own a license or your organization provides one.

\subsubsection*{Before You Start}\label{docs__pandoc__latex__src__setup_accessibility_complete.md__before-you-start}

\begin{enumerate}
\tightlist
\item
  Make sure your screen reader is running \textbf{before} opening your terminal.
\item
  Open your terminal (PowerShell, CMD, or Git Bash) and let your screen reader read the title and prompt.
\item
  If you do not hear anything, press \textbf{Alt+Tab} to cycle windows and find your terminal.
\item
  Use your screen reader\textquotesingle s review/virtual cursor to understand the layout.
\end{enumerate}

\begin{center}\rule{0.5\linewidth}{0.5pt}\end{center}

\protect\phantomsection\label{docs__pandoc__latex__src__setup_accessibility_complete.md__nvda-tips}{}

\subsection*{NVDA-Specific Tips}\label{docs__pandoc__latex__src__setup_accessibility_complete.md__nvda-specific-tips}

NVDA is free and available from \url{https://www.nvaccess.org/}

These tips work across all shells: PowerShell, CMD, and Git Bash.

\subsubsection*{Key Commands for Terminals}\label{docs__pandoc__latex__src__setup_accessibility_complete.md__key-commands-for-terminals}

{\def\LTcaptype{none} % do not increment counter
\begin{longtable}[]{@{}
  >{\raggedright\arraybackslash}p{(\linewidth - 2\tabcolsep) * \real{0.3611}}
  >{\raggedright\arraybackslash}p{(\linewidth - 2\tabcolsep) * \real{0.6389}}@{}}
\toprule\noalign{}
\begin{minipage}[b]{\linewidth}\raggedright
Command
\end{minipage} & \begin{minipage}[b]{\linewidth}\raggedright
What It Does
\end{minipage} \\
\midrule\noalign{}
\endhead
\bottomrule\noalign{}
\endlastfoot
\textbf{NVDA+Home} & Read the current line (your command or output) \\
\textbf{NVDA+Down Arrow} & Read from cursor position to end of screen \\
\textbf{NVDA+Up Arrow} & Read from top of screen to cursor \\
\textbf{NVDA+Page Down} & Read next page \\
\textbf{NVDA+Page Up} & Read previous page \\
\textbf{NVDA+F7} & Open the Elements List / Review Mode viewer \\
\textbf{NVDA+Shift+Right Arrow} & Read next word \\
\textbf{NVDA+Shift+Down Arrow} & Read entire visible screen \\
\textbf{NVDA+End} & Jump to end of line \\
\textbf{NVDA+Home} & Jump to start of line \\
\textbf{NVDA+Numpad 5} &
Announce the character under the review cursor \\
\textbf{NVDA+Shift+Numpad 5} &
Announce the word under the review cursor \\
\end{longtable}
}

\begin{quote}
\textbf{Note on NVDA key:} On desktop keyboards the NVDA modifier is typically the \textbf{Insert} key. On laptops without a numpad the NVDA key is often \textbf{CapsLock}. You can change this in NVDA Preferences \textgreater{} Keyboard.
\end{quote}

\subsubsection*{Example: Reading Long Output}\label{docs__pandoc__latex__src__setup_accessibility_complete.md__example-reading-long-output}

\textbf{Scenario:} You ran a command and it listed 50 files. You cannot hear them all at once.

\textbf{Solution with NVDA:}

\begin{enumerate}
\tightlist
\item
  After the command finishes, press \textbf{NVDA+Home} to read the current line.
\item
  Press \textbf{NVDA+Down Arrow} repeatedly to read all output line by line.
\item
  Or redirect to a file first: \texttt{command\ \textgreater{}\ output.txt}, then open with \texttt{notepad\ output.txt} for more comfortable reading.
\end{enumerate}

\subsubsection*{NVDA Settings for Terminals}\label{docs__pandoc__latex__src__setup_accessibility_complete.md__nvda-settings-for-terminals}

\begin{enumerate}
\tightlist
\item
  Press \textbf{NVDA+N} to open the NVDA menu.
\item
  Go to \textbf{Preferences \textgreater{} Settings \textgreater{} Speech}.
\item
  Increase or decrease verbosity to taste.
\item
  Under \textbf{Document Formatting}, enable \textbf{"Report indentation"} (important for reading code).
\item
  Under \textbf{Terminal}, enable \textbf{"Speak typed characters"} so you hear each key as you type.
\end{enumerate}

\begin{center}\rule{0.5\linewidth}{0.5pt}\end{center}

\protect\phantomsection\label{docs__pandoc__latex__src__setup_accessibility_complete.md__jaws-tips}{}

\subsection*{JAWS-Specific Tips}\label{docs__pandoc__latex__src__setup_accessibility_complete.md__jaws-specific-tips}

JAWS is a commercial screen reader available from \url{https://www.freedomscientific.com/products/software/jaws/}

These tips work across all shells: PowerShell, CMD, and Git Bash.

\subsubsection*{Key Commands for Terminals}\label{docs__pandoc__latex__src__setup_accessibility_complete.md__key-commands-for-terminals-1}

{\def\LTcaptype{none} % do not increment counter
\begin{longtable}[]{@{}ll@{}}
\toprule\noalign{}
Command & What It Does \\
\midrule\noalign{}
\endhead
\bottomrule\noalign{}
\endlastfoot
\textbf{Insert+Down Arrow} & Read all / say all from cursor downward \\
\textbf{Insert+Up Arrow} & Re-read current line \\
\textbf{Insert+Page Down} & Read next page of terminal output \\
\textbf{Insert+Page Up} & Read previous page of terminal output \\
\textbf{Insert+End} & Jump to end of text on screen \\
\textbf{Insert+Home} & Jump to start of text on screen \\
\textbf{Insert+Ctrl+Down} & Read to end of screen \\
\textbf{Insert+Ctrl+Up} & Read to beginning of screen \\
\textbf{Insert+Shift+Page Down} & Select and read next page \\
\textbf{Insert+F3} & Open JAWS menu \\
\textbf{Insert+Z} & Toggle JAWS virtual cursor on/off \\
\textbf{Insert+F2} & Open JAWS Manager (Settings Center) \\
\end{longtable}
}

\begin{quote}
\textbf{Note on JAWS key:} The JAWS modifier is typically \textbf{Insert} on a standard keyboard. On laptops, JAWS can be configured to use \textbf{CapsLock} as the modifier instead.
\end{quote}

\subsubsection*{Example: Reading Long Output}\label{docs__pandoc__latex__src__setup_accessibility_complete.md__example-reading-long-output-1}

\textbf{Scenario:} You ran a command and saved output to a file.

\textbf{Solution with JAWS:}

\begin{enumerate}
\tightlist
\item
  Open the file in Notepad: \texttt{notepad\ file.txt}
\item
  In Notepad, press \textbf{Insert+Ctrl+Down} to hear all content from the top.
\item
  Use \textbf{Insert+Down Arrow} to read line by line at your own pace.
\item
  Use \textbf{Insert+F} (JAWS Find) to search for specific text within the file.
\end{enumerate}

\subsubsection*{JAWS Settings for Terminals}\label{docs__pandoc__latex__src__setup_accessibility_complete.md__jaws-settings-for-terminals}

\begin{enumerate}
\tightlist
\item
  Press \textbf{Insert+F3} to open the JAWS menu.
\item
  Go to \textbf{Options \textgreater{} Settings Center} (or \textbf{Utilities \textgreater{} Settings Manager} in older versions).
\item
  Search for \textbf{"terminal"} or \textbf{"console"}.
\item
  Enable \textbf{"Announce output"} and verify \textbf{"Speak when program speaks"} is on.
\item
  For indent reporting, search for \textbf{"Indent"} and set to \textbf{"Tones"} or \textbf{"Tones and Speech"}.
\end{enumerate}

\begin{center}\rule{0.5\linewidth}{0.5pt}\end{center}

\protect\phantomsection\label{docs__pandoc__latex__src__setup_accessibility_complete.md__narrator-tips}{}

\subsection*{Windows Narrator-Specific Tips}\label{docs__pandoc__latex__src__setup_accessibility_complete.md__windows-narrator-specific-tips}

Windows Narrator is built into Windows 10 and Windows 11 at no extra cost. Enable it via \textbf{Settings \textgreater{} Accessibility \textgreater{} Narrator} or press \textbf{Windows+Ctrl+Enter}.

These tips work across all shells: PowerShell, CMD, and Git Bash.

\subsubsection*{Key Commands for Terminals}\label{docs__pandoc__latex__src__setup_accessibility_complete.md__key-commands-for-terminals-2}

{\def\LTcaptype{none} % do not increment counter
\begin{longtable}[]{@{}
  >{\raggedright\arraybackslash}p{(\linewidth - 2\tabcolsep) * \real{0.3766}}
  >{\raggedright\arraybackslash}p{(\linewidth - 2\tabcolsep) * \real{0.6234}}@{}}
\toprule\noalign{}
\begin{minipage}[b]{\linewidth}\raggedright
Command
\end{minipage} & \begin{minipage}[b]{\linewidth}\raggedright
What It Does
\end{minipage} \\
\midrule\noalign{}
\endhead
\bottomrule\noalign{}
\endlastfoot
\textbf{Narrator+D} & Read current line \\
\textbf{Narrator+M} & Read next line \\
\textbf{Narrator+I} & Read current item or focused element \\
\textbf{Narrator+R} & Read from cursor / read all from here \\
\textbf{Narrator+Ctrl+R} & Read page / read all \\
\textbf{Narrator+Left/Right Arrow} & Move to previous/next item \\
\textbf{Narrator+Up/Down Arrow} &
Move to previous/next line in scan mode \\
\textbf{Narrator+Enter} & Activate/interact with current item \\
\textbf{Narrator+Space} &
Toggle Narrator scan mode (browse vs. type mode) \\
\textbf{Narrator+F1} & Open Narrator help \\
\textbf{Windows+Ctrl+Enter} & Start or stop Narrator \\
\end{longtable}
}

\begin{quote}
\textbf{Note on Narrator key:} The Narrator modifier is \textbf{CapsLock} or \textbf{Insert} by default. You can change this in \textbf{Settings \textgreater{} Accessibility \textgreater{} Narrator \textgreater{} Keyboard shortcuts}.
\end{quote}

\subsubsection*{Working with Terminals Using Narrator}\label{docs__pandoc__latex__src__setup_accessibility_complete.md__working-with-terminals-using-narrator}

Narrator works well in \textbf{focus mode} (the default when a terminal is active). In focus mode, arrow keys move through command history and output rather than triggering Narrator navigation.

\textbf{Tips:}

\begin{itemize}
\tightlist
\item
  After running a command, press \textbf{Narrator+D} to read the current line.
\item
  Use \textbf{Narrator+R} to read all output from the current position.
\item
  For long outputs, always redirect to a file: \texttt{command\ \textgreater{}\ output.txt}, then open in Notepad where Narrator\textquotesingle s reading experience is more comfortable.
\item
  In Windows Terminal, enable \textbf{"Accessible terminal"} mode in Windows Terminal settings for better Narrator integration.
\end{itemize}

\subsubsection*{Narrator Settings for Terminals}\label{docs__pandoc__latex__src__setup_accessibility_complete.md__narrator-settings-for-terminals}

\begin{enumerate}
\tightlist
\item
  Open \textbf{Settings \textgreater{} Accessibility \textgreater{} Narrator}.
\item
  Under \textbf{"Change what you hear when typing"}, enable \textbf{"Hear characters as you type"}.
\item
  Under \textbf{"Change what Narrator reads"}, enable \textbf{"Read hints for controls and buttons"}.
\item
  Set verbosity level to \textbf{3 or 4} for a good balance of detail.
\item
  Under \textbf{"Choose a Narrator voice"}, select a voice you find comfortable for extended use.
\end{enumerate}

\textbf{Known limitation:} Narrator has fewer customization options than JAWS or NVDA for advanced terminal work. If you find Narrator insufficient for complex workflows, consider switching to NVDA (free) or JAWS.

\begin{center}\rule{0.5\linewidth}{0.5pt}\end{center}

\protect\phantomsection\label{docs__pandoc__latex__src__setup_accessibility_complete.md__dolphin-tips}{}

\subsection*{Dolphin SuperNova-Specific Tips}\label{docs__pandoc__latex__src__setup_accessibility_complete.md__dolphin-supernova-specific-tips}

Dolphin SuperNova is a commercial screen reader (with optional magnification) available from \url{https://yourdolphin.com/supernova/}

These tips work across all shells: PowerShell, CMD, and Git Bash.

\subsubsection*{Key Commands for Terminals}\label{docs__pandoc__latex__src__setup_accessibility_complete.md__key-commands-for-terminals-3}

{\def\LTcaptype{none} % do not increment counter
\begin{longtable}[]{@{}
  >{\raggedright\arraybackslash}p{(\linewidth - 2\tabcolsep) * \real{0.4872}}
  >{\raggedright\arraybackslash}p{(\linewidth - 2\tabcolsep) * \real{0.5128}}@{}}
\toprule\noalign{}
\begin{minipage}[b]{\linewidth}\raggedright
Command
\end{minipage} & \begin{minipage}[b]{\linewidth}\raggedright
What It Does
\end{minipage} \\
\midrule\noalign{}
\endhead
\bottomrule\noalign{}
\endlastfoot
\textbf{Caps Lock+Down Arrow} (or NumPad 2) & Read current line \\
\textbf{Caps Lock+Numpad 5} & Read current word \\
\textbf{Caps Lock+Numpad 6} & Read next word \\
\textbf{Caps Lock+Numpad 4} & Read previous word \\
\textbf{Caps Lock+Numpad 8} & Read from cursor to top \\
\textbf{Caps Lock+Numpad 2} & Read from cursor to bottom \\
\textbf{Caps Lock+Numpad Plus} &
Say all / read all from current position \\
\textbf{Caps Lock+L} & Re-read current line \\
\textbf{Caps Lock+Right Arrow} & Read next character \\
\textbf{Caps Lock+Left Arrow} & Read previous character \\
\textbf{Caps Lock+Page Down} & Jump to end of document or screen \\
\textbf{Caps Lock+Page Up} & Jump to start of document or screen \\
\end{longtable}
}

\begin{quote}
\textbf{Note on Dolphin key:} Dolphin SuperNova typically uses \textbf{Caps Lock} as its modifier key (referred to as the "Dolphin key"). This can be changed in \textbf{SuperNova Control Panel \textgreater{} Keyboard}.
\end{quote}

\subsubsection*{Working with Terminals Using Dolphin SuperNova}\label{docs__pandoc__latex__src__setup_accessibility_complete.md__working-with-terminals-using-dolphin-supernova}

\begin{enumerate}
\tightlist
\item
  Open SuperNova \textbf{before} opening your terminal application.
\item
  Dolphin SuperNova automatically tracks focus as you move between the editor, terminal, and file explorer.
\item
  In the terminal, use the \textbf{Dolphin key+Numpad Plus} (say all) to have the screen read after running a command.
\item
  For complex output, redirect to a file and open in Notepad: \texttt{command\ \textgreater{}\ output.txt}, then \texttt{notepad\ output.txt}.
\end{enumerate}

\subsubsection*{Dolphin SuperNova Settings for Terminals}\label{docs__pandoc__latex__src__setup_accessibility_complete.md__dolphin-supernova-settings-for-terminals}

\begin{enumerate}
\tightlist
\item
  Open \textbf{SuperNova Control Panel} (press Caps Lock+SpaceBar, or click the SuperNova icon).
\item
  Go to \textbf{Speech \textgreater{} Verbosity} and set a comfortable level (3 or 4 recommended for terminal work).
\item
  Go to \textbf{Speech \textgreater{} Text Processing \textgreater{} Indentation} and enable indentation announcement (important for code).
\item
  Go to \textbf{Braille \textgreater{} Settings} if you are using a braille display alongside speech.
\item
  In \textbf{General \textgreater{} Application Settings}, you can set terminal-specific verbosity so SuperNova behaves differently in your terminal versus other programs.
\end{enumerate}

\textbf{Note:} Dolphin SuperNova\textquotesingle s magnification features can be helpful if you have some remaining vision. Terminal text can be magnified while still receiving speech output.

\begin{center}\rule{0.5\linewidth}{0.5pt}\end{center}

\protect\phantomsection\label{docs__pandoc__latex__src__setup_accessibility_complete.md__general-terminal}{}

\subsection*{General Terminal Accessibility}\label{docs__pandoc__latex__src__setup_accessibility_complete.md__general-terminal-accessibility}

\subsubsection*{Understanding the Terminal Layout}\label{docs__pandoc__latex__src__setup_accessibility_complete.md__understanding-the-terminal-layout}

All terminals (PowerShell, CMD, Git Bash) follow the same basic layout:

\begin{enumerate}
\tightlist
\item
  \textbf{Title bar:} Window name (e.g., "Windows PowerShell", "Command Prompt", "Git Bash")
\item
  \textbf{Content area:} Command history and output
\item
  \textbf{Prompt:} The area where you type (e.g., \texttt{PS\textgreater{}}, \texttt{C:\textbackslash{}\textgreater{}}, \texttt{\$})
\end{enumerate}

\textbf{Your screen reader reads from top to bottom, but focus is at the prompt (bottom).}

\subsubsection*{Navigation Sequence}\label{docs__pandoc__latex__src__setup_accessibility_complete.md__navigation-sequence}

\textbf{When you open any terminal:}

\begin{enumerate}
\tightlist
\item
  Your screen reader announces the window title.
\item
  Then it announces the prompt line.
\item
  Anything before the prompt is previous output.
\item
  Anything after the prompt is where new output will appear.
\end{enumerate}

\subsubsection*{Reading Output Effectively}\label{docs__pandoc__latex__src__setup_accessibility_complete.md__reading-output-effectively}

\textbf{Strategy 1: Immediate Output (Small Amount)}

\begin{itemize}
\tightlist
\item
  Run a command.
\item
  Your screen reader announces output immediately.
\item
  This works well for short outputs (a few lines).
\end{itemize}

\textbf{Strategy 2: Large Output (Many Lines)}

\begin{itemize}
\tightlist
\item
  Redirect to a file: \texttt{command\ \textgreater{}\ output.txt}
\item
  Open the file: \texttt{notepad\ output.txt} (works in all shells on Windows)
\item
  Read in Notepad --- easier and more controllable for all screen readers.
\end{itemize}

\textbf{Strategy 3: Filtering Output}

\begin{itemize}
\tightlist
\item
  Use filtering commands to reduce output.
\item
  PowerShell: \texttt{command\ \textbar{}\ Select-String\ "pattern"}
\item
  CMD: \texttt{command\ \textbar{}\ find\ "pattern"}
\item
  Git Bash: \texttt{command\ \textbar{}\ grep\ "pattern"}
\end{itemize}

\begin{center}\rule{0.5\linewidth}{0.5pt}\end{center}

\protect\phantomsection\label{docs__pandoc__latex__src__setup_accessibility_complete.md__long-output}{}

\subsection*{Working with Long Output}\label{docs__pandoc__latex__src__setup_accessibility_complete.md__working-with-long-output}

This is one of the most common challenges for screen reader users. Here are proven solutions:

\subsubsection*{Solution 1: Redirect to a File (Recommended)}\label{docs__pandoc__latex__src__setup_accessibility_complete.md__solution-1-redirect-to-a-file-recommended}

\begin{lstlisting}[style=Alabaster]
command > output.txt
notepad output.txt

\end{lstlisting}

\textbf{Advantages:} Easy to navigate at your own pace, works with all screen readers, output does not scroll away, you can save it for later reference.

\subsubsection*{Solution 2: Use Pagination}\label{docs__pandoc__latex__src__setup_accessibility_complete.md__solution-2-use-pagination}

\begin{lstlisting}[style=Alabaster]
command | more

\end{lstlisting}

Press \textbf{Space} to see the next page, \textbf{Q} to quit. Note: some screen readers struggle with \texttt{more}; Solution 1 is generally preferred.

\subsubsection*{Solution 3: Filter Output}\label{docs__pandoc__latex__src__setup_accessibility_complete.md__solution-3-filter-output}

\begin{lstlisting}[style=Alabaster]
# PowerShell
ls | Select-String "\.scad"

# CMD
dir | find ".scad"

# Git Bash
ls -la | grep ".scad"

\end{lstlisting}

\subsubsection*{Solution 4: Count Before Displaying}\label{docs__pandoc__latex__src__setup_accessibility_complete.md__solution-4-count-before-displaying}

\begin{lstlisting}[style=Alabaster, language=powershell]
# PowerShell
(ls).Count

\end{lstlisting}

This tells you how many items there are before deciding whether to redirect to a file.

\begin{center}\rule{0.5\linewidth}{0.5pt}\end{center}

\protect\phantomsection\label{docs__pandoc__latex__src__setup_accessibility_complete.md__shortcuts}{}

\subsection*{Keyboard Shortcuts Reference}\label{docs__pandoc__latex__src__setup_accessibility_complete.md__keyboard-shortcuts-reference}

\subsubsection*{Shell Navigation (Works Regardless of Screen Reader)}\label{docs__pandoc__latex__src__setup_accessibility_complete.md__shell-navigation-works-regardless-of-screen-reader}

{\def\LTcaptype{none} % do not increment counter
\begin{longtable}[]{@{}
  >{\raggedright\arraybackslash}p{(\linewidth - 4\tabcolsep) * \real{0.1750}}
  >{\raggedright\arraybackslash}p{(\linewidth - 4\tabcolsep) * \real{0.4250}}
  >{\raggedright\arraybackslash}p{(\linewidth - 4\tabcolsep) * \real{0.4000}}@{}}
\toprule\noalign{}
\begin{minipage}[b]{\linewidth}\raggedright
Key
\end{minipage} & \begin{minipage}[b]{\linewidth}\raggedright
Action
\end{minipage} & \begin{minipage}[b]{\linewidth}\raggedright
Why It Matters
\end{minipage} \\
\midrule\noalign{}
\endhead
\bottomrule\noalign{}
\endlastfoot
\textbf{Up Arrow} & Show previous command &
Repeat commands without retyping \\
\textbf{Down Arrow} & Show next command & Navigate through history \\
\textbf{Tab} & Auto-complete path or command name &
Faster and more accurate \\
\textbf{Shift+Tab} & Cycle backward through completions &
If Tab went too far \\
\textbf{Home} & Jump to start of command line &
Edit beginning of a command \\
\textbf{End} & Jump to end of command line & Edit end of a command \\
\textbf{Ctrl+A} & Select all text on line & Copy entire command \\
\textbf{Ctrl+C} & Stop / interrupt running command &
Abort long-running tasks \\
\textbf{Ctrl+L} & Clear screen & Start fresh visually \\
\textbf{Enter} & Run command & Execute what you typed \\
\end{longtable}
}

\subsubsection*{Screen Reader Navigation Quick Reference}\label{docs__pandoc__latex__src__setup_accessibility_complete.md__screen-reader-navigation-quick-reference}

{\def\LTcaptype{none} % do not increment counter
\begin{longtable}[]{@{}
  >{\raggedright\arraybackslash}p{(\linewidth - 8\tabcolsep) * \real{0.2396}}
  >{\raggedright\arraybackslash}p{(\linewidth - 8\tabcolsep) * \real{0.1979}}
  >{\raggedright\arraybackslash}p{(\linewidth - 8\tabcolsep) * \real{0.1771}}
  >{\raggedright\arraybackslash}p{(\linewidth - 8\tabcolsep) * \real{0.1771}}
  >{\raggedright\arraybackslash}p{(\linewidth - 8\tabcolsep) * \real{0.2083}}@{}}
\toprule\noalign{}
\begin{minipage}[b]{\linewidth}\raggedright
Task
\end{minipage} & \begin{minipage}[b]{\linewidth}\raggedright
NVDA
\end{minipage} & \begin{minipage}[b]{\linewidth}\raggedright
JAWS
\end{minipage} & \begin{minipage}[b]{\linewidth}\raggedright
Windows Narrator
\end{minipage} & \begin{minipage}[b]{\linewidth}\raggedright
Dolphin SuperNova
\end{minipage} \\
\midrule\noalign{}
\endhead
\bottomrule\noalign{}
\endlastfoot
Read current line & NVDA+Home & Insert+Up Arrow & Narrator+D &
CapsLock+L \\
Read all from here & NVDA+Down Arrow & Insert+Ctrl+Down & Narrator+R &
CapsLock+Numpad Plus \\
Read next line & Down Arrow (review) & Insert+Down Arrow & Narrator+M &
CapsLock+Numpad 2 \\
Read previous line & Up Arrow (review) & Insert+Up Arrow &
Narrator+Up Arrow & CapsLock+Numpad 8 \\
Jump to end of output & NVDA+End & Insert+End & Narrator+End &
CapsLock+Page Down \\
Jump to start of output & NVDA+Home & Insert+Home & Narrator+Home &
CapsLock+Page Up \\
Stop reading & Ctrl & Ctrl & Ctrl & Ctrl \\
\end{longtable}
}

\begin{center}\rule{0.5\linewidth}{0.5pt}\end{center}

\protect\phantomsection\label{docs__pandoc__latex__src__setup_accessibility_complete.md__shell-differences}{}

\subsection*{Shell-Specific Differences}\label{docs__pandoc__latex__src__setup_accessibility_complete.md__shell-specific-differences}

While all three shells work equally well with screen readers, there are differences in commands and syntax:

\subsubsection*{File Listing and Navigation}\label{docs__pandoc__latex__src__setup_accessibility_complete.md__file-listing-and-navigation}

{\def\LTcaptype{none} % do not increment counter
\begin{longtable}[]{@{}
  >{\raggedright\arraybackslash}p{(\linewidth - 6\tabcolsep) * \real{0.2532}}
  >{\raggedright\arraybackslash}p{(\linewidth - 6\tabcolsep) * \real{0.2911}}
  >{\raggedright\arraybackslash}p{(\linewidth - 6\tabcolsep) * \real{0.2278}}
  >{\raggedright\arraybackslash}p{(\linewidth - 6\tabcolsep) * \real{0.2278}}@{}}
\toprule\noalign{}
\begin{minipage}[b]{\linewidth}\raggedright
Task
\end{minipage} & \begin{minipage}[b]{\linewidth}\raggedright
PowerShell
\end{minipage} & \begin{minipage}[b]{\linewidth}\raggedright
CMD
\end{minipage} & \begin{minipage}[b]{\linewidth}\raggedright
Git Bash
\end{minipage} \\
\midrule\noalign{}
\endhead
\bottomrule\noalign{}
\endlastfoot
\textbf{List files} & \texttt{ls} or \texttt{Get-ChildItem} &
\texttt{dir} & \texttt{ls\ -la} \\
\textbf{List names only} & \texttt{ls\ -Name} & \texttt{dir\ /B} &
\texttt{ls\ -1} \\
\textbf{Change directory} & \texttt{cd\ path} & \texttt{cd\ path} &
\texttt{cd\ path} \\
\textbf{Current location} & \texttt{pwd} & \texttt{cd} (no args) &
\texttt{pwd} \\
\textbf{Create folder} & \texttt{mkdir\ foldername} &
\texttt{mkdir\ foldername} & \texttt{mkdir\ foldername} \\
\textbf{Delete file} & \texttt{Remove-Item\ file.txt} &
\texttt{del\ file.txt} & \texttt{rm\ file.txt} \\
\textbf{Home directory} & \texttt{\$HOME} or \texttt{\textasciitilde{}}
& \texttt{\%USERPROFILE\%} & \texttt{\textasciitilde{}} \\
\end{longtable}
}

\subsubsection*{Output Redirection}\label{docs__pandoc__latex__src__setup_accessibility_complete.md__output-redirection}

{\def\LTcaptype{none} % do not increment counter
\begin{longtable}[]{@{}
  >{\raggedright\arraybackslash}p{(\linewidth - 6\tabcolsep) * \real{0.1748}}
  >{\raggedright\arraybackslash}p{(\linewidth - 6\tabcolsep) * \real{0.3204}}
  >{\raggedright\arraybackslash}p{(\linewidth - 6\tabcolsep) * \real{0.2330}}
  >{\raggedright\arraybackslash}p{(\linewidth - 6\tabcolsep) * \real{0.2718}}@{}}
\toprule\noalign{}
\begin{minipage}[b]{\linewidth}\raggedright
Task
\end{minipage} & \begin{minipage}[b]{\linewidth}\raggedright
PowerShell
\end{minipage} & \begin{minipage}[b]{\linewidth}\raggedright
CMD
\end{minipage} & \begin{minipage}[b]{\linewidth}\raggedright
Git Bash
\end{minipage} \\
\midrule\noalign{}
\endhead
\bottomrule\noalign{}
\endlastfoot
\textbf{Save to file} & \texttt{command\ \textgreater{}\ file.txt} &
\texttt{command\ \textgreater{}\ file.txt} &
\texttt{command\ \textgreater{}\ file.txt} \\
\textbf{Append to file} &
\texttt{command\ \textgreater{}\textgreater{}\ file.txt} &
\texttt{command\ \textgreater{}\textgreater{}\ file.txt} &
\texttt{command\ \textgreater{}\textgreater{}\ file.txt} \\
\textbf{Filter output} &
\texttt{command\ \textbar{}\ Select-String\ "text"} &
\texttt{command\ \textbar{}\ find\ "text"} &
\texttt{command\ \textbar{}\ grep\ "text"} \\
\textbf{Open file} & \texttt{notepad\ file.txt} &
\texttt{notepad\ file.txt} & \texttt{notepad\ file.txt} (Windows) \\
\end{longtable}
}

\subsubsection*{Scripting and Automation}\label{docs__pandoc__latex__src__setup_accessibility_complete.md__scripting-and-automation}

{\def\LTcaptype{none} % do not increment counter
\begin{longtable}[]{@{}
  >{\raggedright\arraybackslash}p{(\linewidth - 6\tabcolsep) * \real{0.2769}}
  >{\raggedright\arraybackslash}p{(\linewidth - 6\tabcolsep) * \real{0.2923}}
  >{\raggedright\arraybackslash}p{(\linewidth - 6\tabcolsep) * \real{0.2308}}
  >{\raggedright\arraybackslash}p{(\linewidth - 6\tabcolsep) * \real{0.2000}}@{}}
\toprule\noalign{}
\begin{minipage}[b]{\linewidth}\raggedright
Feature
\end{minipage} & \begin{minipage}[b]{\linewidth}\raggedright
PowerShell
\end{minipage} & \begin{minipage}[b]{\linewidth}\raggedright
CMD
\end{minipage} & \begin{minipage}[b]{\linewidth}\raggedright
Git Bash
\end{minipage} \\
\midrule\noalign{}
\endhead
\bottomrule\noalign{}
\endlastfoot
\textbf{File extension} & \texttt{.ps1} & \texttt{.bat} &
\texttt{.sh} \\
\textbf{Comment} & \texttt{\#} & \texttt{REM} or \texttt{::} &
\texttt{\#} \\
\textbf{Variable} & \texttt{\$var\ =\ "value"} & \texttt{set\ var=value}
& \texttt{var=value} \\
\textbf{Echo text} & \texttt{Write-Host\ "text"} & \texttt{echo\ text} &
\texttt{echo\ text} \\
\textbf{Get help} & \texttt{Get-Help\ command} & \texttt{help\ command}
& \texttt{man\ command} \\
\end{longtable}
}

\textbf{Choose ONE pathway and stick with it.} Each curriculum teaches you the concepts using that specific shell\textquotesingle s syntax. The accessibility experience is virtually identical across all four screen readers in all three shells --- only the command syntax differs.

\begin{center}\rule{0.5\linewidth}{0.5pt}\end{center}

\subsection*{Troubleshooting}\label{docs__pandoc__latex__src__setup_accessibility_complete.md__troubleshooting}

\subsubsection*{Problem 1: "I Can\textquotesingle t Hear the Output After Running a Command"}\label{docs__pandoc__latex__src__setup_accessibility_complete.md__problem-1-i-cant-hear-the-output-after-running-a-command}

\textbf{Causes and Solutions:}

\begin{enumerate}
\tightlist
\item
  \textbf{Cursor is not at the prompt} --- Press \textbf{End} or \textbf{Ctrl+End} to go to the end of text, then use screen reader commands to review.
\item
  \textbf{Output scrolled off-screen} --- Redirect to file: \texttt{command\ \textgreater{}\ output.txt}, then \texttt{notepad\ output.txt}.
\item
  \textbf{Screen reader focus is on window title, not content} --- Press \textbf{Tab} or arrow keys to move into the content area.
\item
  \textbf{Large output overwhelming screen reader} --- Use filtering: \texttt{command\ \textbar{}\ grep\ "pattern"} (Git Bash), \texttt{command\ \textbar{}\ find\ "text"} (CMD), or \texttt{command\ \textbar{}\ Select-String\ "text"} (PowerShell).
\end{enumerate}

\subsubsection*{Problem 2: "Tab Completion Isn\textquotesingle t Working"}\label{docs__pandoc__latex__src__setup_accessibility_complete.md__problem-2-tab-completion-isnt-working}

\begin{enumerate}
\tightlist
\item
  You must type at least one character before pressing Tab.
\item
  The folder or file must exist --- run \texttt{dir} (CMD), \texttt{ls} (PowerShell/Git Bash) to check.
\item
  If multiple matches exist, press Tab again to cycle through them.
\end{enumerate}

\subsubsection*{Problem 3: "Access Denied or Permission Denied"}\label{docs__pandoc__latex__src__setup_accessibility_complete.md__problem-3-access-denied-or-permission-denied}

\begin{enumerate}
\tightlist
\item
  Close your terminal, right-click it, and choose \textbf{Run as administrator}. Confirm the UAC prompt.
\item
  If the file is in use by another program, close that program and try again.
\item
  If the path contains spaces, use quotes: \texttt{cd\ "Program\ Files"}.
\end{enumerate}

\subsubsection*{Problem 4: "A Command Runs Forever and Won\textquotesingle t Stop"}\label{docs__pandoc__latex__src__setup_accessibility_complete.md__problem-4-a-command-runs-forever-and-wont-stop}

Press \textbf{Ctrl+C} to interrupt any running command.

\subsubsection*{Problem 5: "I Need to Edit My Last Command"}\label{docs__pandoc__latex__src__setup_accessibility_complete.md__problem-5-i-need-to-edit-my-last-command}

\begin{enumerate}
\tightlist
\item
  Press \textbf{Up Arrow} to recall the previous command.
\item
  Use \textbf{Left/Right Arrow} keys to move through it.
\item
  Edit as needed and press \textbf{Enter} to run the modified version.
\end{enumerate}

\subsubsection*{Problem 6: Screen Reader is Not Announcing Indentation}\label{docs__pandoc__latex__src__setup_accessibility_complete.md__problem-6-screen-reader-is-not-announcing-indentation}

Proper indent announcement is critical for reading code. Enable it for your screen reader:

\begin{itemize}
\tightlist
\item
  \textbf{NVDA:} NVDA Menu \textgreater{} Preferences \textgreater{} Settings \textgreater{} Document Formatting \textgreater{} enable \textbf{"Report indentation"}, set to \textbf{"Tones and speech"}.
\item
  \textbf{JAWS:} Settings Center \textgreater{} search \textbf{"Indent"} \textgreater{} enable \textbf{"Announce Indentation"}, set mode to \textbf{"Tones"} or \textbf{"Tones and Speech"}.
\item
  \textbf{Windows Narrator:} Settings \textgreater{} Accessibility \textgreater{} Narrator \textgreater{} Advanced Options \textgreater{} enable \textbf{"Report indentation"} (limited options compared to NVDA/JAWS).
\item
  \textbf{Dolphin SuperNova:} SuperNova Control Panel \textgreater{} Speech \textgreater{} Text Processing \textgreater{} Indentation \textgreater{} enable and set preferred announcement style.
\end{itemize}

\begin{center}\rule{0.5\linewidth}{0.5pt}\end{center}

\subsection*{Pro Tips for Efficiency}\label{docs__pandoc__latex__src__setup_accessibility_complete.md__pro-tips-for-efficiency}

\subsubsection*{Create Aliases for Frequently Used Commands}\label{docs__pandoc__latex__src__setup_accessibility_complete.md__create-aliases-for-frequently-used-commands}

\textbf{PowerShell:} \texttt{Set-Alias\ -Name\ ll\ -Value\ "Get-ChildItem\ -Name"}\\
\textbf{CMD:} Create a \texttt{.bat} file in a folder on your PATH.\\
\textbf{Git Bash:} \texttt{alias\ ll=\textquotesingle{}ls\ -la\textquotesingle{}} (add to \texttt{\textasciitilde{}/.bashrc} to make permanent)

\subsubsection*{Use Command History}\label{docs__pandoc__latex__src__setup_accessibility_complete.md__use-command-history}

All shells support \textbf{Up/Down Arrow} to navigate history. In PowerShell, \texttt{history} shows a numbered list and \texttt{Invoke-History\ 5} reruns command number 5. In Git Bash, \texttt{history} and \texttt{!5} work similarly.

\subsubsection*{Redirect Everything to Files for Accessibility}\label{docs__pandoc__latex__src__setup_accessibility_complete.md__redirect-everything-to-files-for-accessibility}

If a command produces output, save it:

\begin{lstlisting}[style=Alabaster]
command > results.txt
notepad results.txt

\end{lstlisting}

This is always more accessible than reading from the terminal directly.

\begin{center}\rule{0.5\linewidth}{0.5pt}\end{center}

\subsection*{Quick Reference Card}\label{docs__pandoc__latex__src__setup_accessibility_complete.md__quick-reference-card}

\begin{lstlisting}[style=Alabaster]
EVERY SESSION STARTS WITH:
1. pwd / cd            (where am I?)
2. dir / ls            (what is here?)
3. cd path             (go there)

LONG OUTPUT?
-> command > file.txt
-> notepad file.txt

STUCK?
-> Ctrl+C

WANT TO REPEAT?
-> Up Arrow

NEED HELP?
-> PowerShell:  `Get-Help command-name`
-> CMD:         `help command-name`
-> Git Bash:    `man command-name`

\end{lstlisting}

\begin{center}\rule{0.5\linewidth}{0.5pt}\end{center}

\subsection*{Additional Resources}\label{docs__pandoc__latex__src__setup_accessibility_complete.md__additional-resources-1}

\begin{itemize}
\tightlist
\item
  \textbf{NVDA:} \url{https://www.nvaccess.org/}
\item
  \textbf{JAWS:} \url{https://www.freedomscientific.com/products/software/jaws/}
\item
  \textbf{Windows Narrator:} \url{https://support.microsoft.com/narrator}
\item
  \textbf{Dolphin SuperNova:} \url{https://yourdolphin.com/supernova/}
\item
  \textbf{PowerShell Docs:} \url{https://example.com}
\item
  \textbf{CMD Documentation:} \url{https://example.com}
\item
  \textbf{Git Bash / Git for Windows:} \url{https://example.com}
\item
  \textbf{Accessibility Best Practices:} \url{https://example.com}
\end{itemize}

\begin{center}\rule{0.5\linewidth}{0.5pt}\end{center}

\protect\phantomsection\label{docs__pandoc__latex__src__setup_accessibility_complete.md__braille-displays}{}

\section*{Part 3: Using a Braille Display with the Curriculum}\label{docs__pandoc__latex__src__setup_accessibility_complete.md__part-3-using-a-braille-display-with-the-curriculum}

\begin{center}\rule{0.5\linewidth}{0.5pt}\end{center}

Many blind and low-vision programmers prefer tactile output alongside or instead of speech. A refreshable braille display presents terminal text as a row (or multiple rows) of braille cells and allows you to read at your own pace, re-read lines without interrupting speech, and work quietly. This page covers how to connect and configure a braille display with each screen reader used in this curriculum.

\begin{center}\rule{0.5\linewidth}{0.5pt}\end{center}

\subsection*{How a Braille Display Works with a Terminal}\label{docs__pandoc__latex__src__setup_accessibility_complete.md__how-a-braille-display-works-with-a-terminal}

A refreshable braille display uses small pins that rise and fall to form braille characters. When connected to a computer and paired with a screen reader, the display shows the line of text where the screen reader\textquotesingle s focus or review cursor is located. You can pan left and right across a long line or press routing buttons to move the screen reader cursor to a specific position.

In a terminal environment:

\begin{itemize}
\tightlist
\item
  The display shows the current command line (what you are typing).
\item
  After running a command, you can use your screen reader\textquotesingle s review cursor and the display\textquotesingle s pan keys to read through the output line by line.
\item
  Routing keys (small buttons above or below each braille cell) let you move the cursor directly to that point in the text --- useful for re-running or editing commands from history.
\end{itemize}

\begin{center}\rule{0.5\linewidth}{0.5pt}\end{center}

\subsection*{Connecting a Braille Display}\label{docs__pandoc__latex__src__setup_accessibility_complete.md__connecting-a-braille-display}

\subsubsection*{USB Connection}\label{docs__pandoc__latex__src__setup_accessibility_complete.md__usb-connection}

\begin{enumerate}
\tightlist
\item
  Plug the display\textquotesingle s USB cable into your computer.
\item
  Windows will attempt to install drivers automatically. If the display is not recognized, install the vendor\textquotesingle s driver software first (see vendor links at the end of this page).
\item
  Open your screen reader\textquotesingle s braille settings and select the connected display.
\end{enumerate}

\subsubsection*{Bluetooth Connection}\label{docs__pandoc__latex__src__setup_accessibility_complete.md__bluetooth-connection}

\begin{enumerate}
\tightlist
\item
  Put the braille display into pairing mode (refer to the device\textquotesingle s quick-start guide).
\item
  On Windows, go to \textbf{Settings \textgreater{} Bluetooth \& devices \textgreater{} Add a device} and pair the display.
\item
  Once paired, open your screen reader\textquotesingle s braille settings and select the display.
\end{enumerate}

\subsubsection*{Before Opening the Terminal}\label{docs__pandoc__latex__src__setup_accessibility_complete.md__before-opening-the-terminal}

Always confirm the display is connected and showing output before opening your terminal. A quick test: open Notepad, type a few words, and verify they appear on the display. This confirms the screen reader is routing output correctly before you add the complexity of a terminal session.

\begin{center}\rule{0.5\linewidth}{0.5pt}\end{center}

\subsection*{Configuring Each Screen Reader for Braille}\label{docs__pandoc__latex__src__setup_accessibility_complete.md__configuring-each-screen-reader-for-braille}

\subsubsection*{NVDA}\label{docs__pandoc__latex__src__setup_accessibility_complete.md__nvda}

NVDA includes built-in braille support for a wide range of displays without requiring separate vendor software for many devices.

\begin{enumerate}
\tightlist
\item
  Connect your display via USB or Bluetooth.
\item
  Open the NVDA menu: press \textbf{NVDA+N} (Insert+N or CapsLock+N depending on your modifier key setting).
\item
  Go to \textbf{Preferences \textgreater{} Settings \textgreater{} Braille} (or press \textbf{NVDA+Ctrl+B} as a shortcut on some versions).
\item
  In the \textbf{"Braille display"} dropdown, select your display from the list. NVDA supports automatic detection for many displays --- you can also select \textbf{"Automatic"} and let NVDA detect it.
\item
  Set your preferred \textbf{braille output table} (e.g., Unified English Braille Grade 1 or Grade 2, or your country\textquotesingle s national table).
\item
  Set \textbf{"Braille input table"} if you want to type commands using the display\textquotesingle s braille keyboard (if equipped).
\item
  Enable \textbf{"Show cursor"} so you can see the caret position on the display.
\item
  Click \textbf{OK} and test by moving the cursor in Notepad.
\end{enumerate}

\textbf{Recommended NVDA braille settings for terminal work:}

\begin{itemize}
\tightlist
\item
  Output table: Unified English Braille Grade 1 (uncontracted) is recommended for reading code --- contractions in Grade 2 can make code harder to read.
\item
  Cursor shape: Dots 7 and 8 (underline) is a common choice.
\item
  Enable \textbf{"Word wrap"} to avoid cutting words at the display\textquotesingle s edge.
\end{itemize}

\textbf{NVDA braille keyboard shortcut reference (while braille display is active):}

{\def\LTcaptype{none} % do not increment counter
\begin{longtable}[]{@{}
  >{\raggedright\arraybackslash}p{(\linewidth - 2\tabcolsep) * \real{0.5493}}
  >{\raggedright\arraybackslash}p{(\linewidth - 2\tabcolsep) * \real{0.4507}}@{}}
\toprule\noalign{}
\begin{minipage}[b]{\linewidth}\raggedright
Action
\end{minipage} & \begin{minipage}[b]{\linewidth}\raggedright
NVDA Command
\end{minipage} \\
\midrule\noalign{}
\endhead
\bottomrule\noalign{}
\endlastfoot
Pan braille display right & Display\textquotesingle s right pan key \\
Pan braille display left & Display\textquotesingle s left pan key \\
Toggle braille tethered to focus/review & NVDA+Ctrl+T \\
Move to previous braille line & NVDA+Shift+Up (or display key) \\
Move to next braille line & NVDA+Shift+Down (or display key) \\
Open braille settings & NVDA+Ctrl+B (some versions) \\
\end{longtable}
}

\begin{center}\rule{0.5\linewidth}{0.5pt}\end{center}

\subsubsection*{JAWS}\label{docs__pandoc__latex__src__setup_accessibility_complete.md__jaws}

JAWS supports a wide range of braille displays through its built-in braille manager and additional vendor drivers.

\begin{enumerate}
\tightlist
\item
  Connect your display via USB or Bluetooth.
\item
  Install the vendor\textquotesingle s display driver if required (check the vendor\textquotesingle s website --- some displays work without a driver in JAWS, others need one).
\item
  Open the JAWS menu: press \textbf{Insert+F3} (or your JAWS key+F3).
\item
  Go to \textbf{Options \textgreater{} Basics \textgreater{} Braille} or use the \textbf{JAWS Settings Center} and search for \textbf{"Braille"}.
\item
  In Braille Settings, select your display model from the list.
\item
  Choose your \textbf{braille translation table} (e.g., English Unified Grade 1 for uncontracted, Grade 2 for contracted).
\item
  Configure \textbf{cursor routing} so that pressing a routing button on the display moves the JAWS cursor to that position.
\item
  Test in Notepad before moving to a terminal.
\end{enumerate}

\textbf{JAWS braille settings for terminal work:}

\begin{itemize}
\tightlist
\item
  Use \textbf{Grade 1 (uncontracted) braille} for reading code and terminal output. Grade 2 contractions can make commands and file paths difficult to parse.
\item
  Enable \textbf{"Show active cursor"} in JAWS braille settings.
\item
  In the JAWS Settings Center, look for \textbf{"Braille for Applications"} --- you can set different braille behavior for specific applications (e.g., a terminal vs. a word processor).
\end{itemize}

\textbf{Accessing braille settings:}

\begin{itemize}
\tightlist
\item
  \textbf{Insert+F3} \textgreater{} Options \textgreater{} Basics \textgreater{} Braille
\item
  Or: \textbf{Insert+F2} (JAWS Manager) \textgreater{} Settings Center \textgreater{} search "braille"
\end{itemize}

\begin{center}\rule{0.5\linewidth}{0.5pt}\end{center}

\subsubsection*{Windows Narrator}\label{docs__pandoc__latex__src__setup_accessibility_complete.md__windows-narrator-1}

Windows Narrator supports braille output through the \textbf{Windows braille service} (Windows 10 version 1903 and later, and Windows 11).

\begin{enumerate}
\tightlist
\item
  Connect your display via USB or Bluetooth.
\item
  Install the vendor\textquotesingle s display driver if required.
\item
  Open \textbf{Settings \textgreater{} Accessibility \textgreater{} Narrator \textgreater{} Braille} (Windows 11) or \textbf{Settings \textgreater{} Ease of Access \textgreater{} Narrator \textgreater{} Use Braille} (Windows 10).
\item
  Turn on \textbf{"Install braille"} (first-time setup may require downloading the braille translation component --- approximately 70 MB).
\item
  Select your \textbf{display brand and model} from the list.
\item
  Choose a \textbf{braille table} (Grade 1 recommended for terminal work).
\item
  Test in Notepad before using in a terminal.
\end{enumerate}

\textbf{Important Narrator braille notes:}

\begin{itemize}
\tightlist
\item
  Narrator\textquotesingle s braille support uses the \textbf{BRLTTY} back-end on Windows, which supports most mainstream displays.
\item
  If your display is not listed, check for a Windows Update or visit the display vendor\textquotesingle s website for a Windows Narrator-compatible driver.
\item
  Narrator\textquotesingle s braille customization is less extensive than NVDA\textquotesingle s or JAWS\textquotesingle s. If you need more control, consider using NVDA (free) for braille-heavy work.
\item
  The \textbf{braille display showing Narrator output} does not change the fact that Narrator\textquotesingle s keyboard commands remain the same --- your display shows what Narrator is currently reading.
\end{itemize}

\begin{center}\rule{0.5\linewidth}{0.5pt}\end{center}

\subsubsection*{Dolphin SuperNova}\label{docs__pandoc__latex__src__setup_accessibility_complete.md__dolphin-supernova-1}

Dolphin SuperNova includes integrated braille support as part of its screen reader editions (SuperNova Magnifier \& Speech and SuperNova Screen Reader).

\begin{enumerate}
\tightlist
\item
  Connect your display via USB or Bluetooth.
\item
  Install the vendor\textquotesingle s driver if required.
\item
  Open the \textbf{SuperNova Control Panel}: press \textbf{CapsLock+SpaceBar} or click the SuperNova icon in the taskbar.
\item
  Go to \textbf{Braille \textgreater{} Display} and select your display model from the list.
\item
  Go to \textbf{Braille \textgreater{} Translation} and select your braille table (Grade 1 / uncontracted recommended for code).
\item
  Go to \textbf{Braille \textgreater{} Settings} to configure cursor style, word wrap, and other options.
\item
  Click \textbf{Apply} and test in Notepad.
\end{enumerate}

\textbf{Dolphin SuperNova braille tips for terminal work:}

\begin{itemize}
\tightlist
\item
  SuperNova supports a wide range of displays from Freedom Scientific (Focus), HumanWare (Brailliant, BrailleNote as terminal), HIMS, Optelec, and others.
\item
  Use \textbf{Grade 1 (uncontracted)} braille for terminal and code work.
\item
  SuperNova\textquotesingle s \textbf{"Braille cursor tracking"} setting determines whether the braille display follows the text cursor or the screen review cursor --- set this to \textbf{"Cursor"} when working in a terminal.
\item
  Dolphin provides a \textbf{braille viewer} on screen (a visual representation of what is on the braille display) --- useful when setting up or troubleshooting.
\end{itemize}

\begin{center}\rule{0.5\linewidth}{0.5pt}\end{center}

\subsection*{Using the Braille Display in a Terminal}\label{docs__pandoc__latex__src__setup_accessibility_complete.md__using-the-braille-display-in-a-terminal}

Once your display is configured, here is how to work in a terminal:

\begin{enumerate}
\tightlist
\item
  \textbf{Reading the prompt:} The display shows your current prompt (e.g., \texttt{C:\textbackslash{}Users\textbackslash{}YourName\textgreater{}} or \texttt{PS\ C:\textbackslash{}\textgreater{}}). Pan right if the prompt is longer than your display width.
\item
  \textbf{Reading output:} After running a command, use your screen reader\textquotesingle s review cursor commands (see the Screen Reader Accessibility Guide) to move through the output. The display shows one line at a time and updates as you move.
\item
  \textbf{Redirecting long output:} For long terminal output, redirect to a file and open in Notepad:

  \begin{lstlisting}[style=Alabaster]
  command > output.txt
  notepad output.txt

  \end{lstlisting}

  In Notepad, you can pan through the braille display comfortably without the output scrolling away.
\item
  \textbf{Using routing keys:} Press the routing key above a word or character to move your cursor to that position. This is especially useful for editing commands recalled from shell history.
\end{enumerate}

\begin{center}\rule{0.5\linewidth}{0.5pt}\end{center}

\subsection*{Braille Grade and Code}\label{docs__pandoc__latex__src__setup_accessibility_complete.md__braille-grade-and-code}

For all screen readers, use \textbf{Grade 1 (uncontracted) braille} when working in terminals and with code. Contracted braille (Grade 2) uses abbreviations designed for reading prose, and those contractions will make command names, file paths, and code syntax difficult or impossible to read correctly.

Example: In Grade 2, the letters "cd" might be represented as a contraction meaning something other than the Change Directory command. In Grade 1, \texttt{cd} is always shown as the letters c and d.

\begin{center}\rule{0.5\linewidth}{0.5pt}\end{center}

\subsection*{Notetakers and BrailleNote / BrailleSense Devices}\label{docs__pandoc__latex__src__setup_accessibility_complete.md__notetakers-and-braillenote--braillesense-devices}

BrailleNote Touch+ (HumanWare) and BrailleSense (HIMS) devices run Android as their primary operating system. \textbf{They cannot run \texttt{3dMake} or Windows-based terminal workflows natively.} These devices are not supported as standalone authoring environments for this curriculum.

\textbf{Workaround --- using them as braille terminals:} Both device families can function as a braille display when connected to a Windows computer via USB or Bluetooth. In that configuration:

\begin{itemize}
\tightlist
\item
  The device acts as a standard refreshable braille display.
\item
  Your Windows screen reader (NVDA, JAWS, Narrator, or SuperNova) drives the output.
\item
  Commands are typed on the Windows keyboard; the BrailleNote/BrailleSense shows the output as braille.
\end{itemize}

To set this up, refer to your device\textquotesingle s documentation for "Terminal mode" or "PC connection mode" and then follow the screen reader braille configuration steps above for your chosen screen reader.

\begin{center}\rule{0.5\linewidth}{0.5pt}\end{center}

\subsection*{Multiline Braille Displays}\label{docs__pandoc__latex__src__setup_accessibility_complete.md__multiline-braille-displays}

Standard refreshable braille displays show a single line of braille (typically 32, 40, or 80 cells). Multiline displays show several lines simultaneously, which can significantly improve the experience of reading code or terminal output because you can see context above and below the current line.

\subsubsection*{Examples}\label{docs__pandoc__latex__src__setup_accessibility_complete.md__examples}

\textbf{Monarch (APH / HumanWare partnership):}\\
A multiline braille device developed by the American Printing House for the Blind (APH) in partnership with HumanWare. It presents multiple lines of braille and supports graphical braille output, making it useful for reviewing blocks of code, diagrams, and structured text. Confirm current driver and screen reader compatibility before purchase.

\textbf{DotPad (Dot Inc.):}\\
A multiline braille and tactile graphics display. Its multiline capability allows programmers to scan several lines of code or terminal output at once. Useful for reviewing structured output without panning repeatedly. Check current Windows compatibility and screen reader support status.

\textbf{Graphiti and Graphiti Plus (Orbit Research):}\\
Multiline tactile displays from Orbit Research designed for reading text and tactile graphics. Their expanded line count helps with reading code structure and terminal output. Check current driver support for NVDA, JAWS, Narrator, and SuperNova.

\subsubsection*{Considerations for Multiline Displays}\label{docs__pandoc__latex__src__setup_accessibility_complete.md__considerations-for-multiline-displays}

\begin{itemize}
\tightlist
\item
  \textbf{Driver and screen reader support:} Multiline displays are newer technology and support varies. Confirm that your chosen screen reader (NVDA, JAWS, Narrator, or SuperNova) has a driver for the specific device before purchase.
\item
  \textbf{Size and portability:} Multiline displays are larger and heavier than single-line models. Consider your workspace and whether you need portability.
\item
  \textbf{Cost:} Multiline displays are significantly more expensive than single-line displays. Evaluate whether the workflow benefit justifies the cost for your situation.
\item
  \textbf{Learning curve:} Using a multiline display effectively --- navigating across lines, understanding the spatial layout --- requires some practice beyond what is needed for single-line displays.
\end{itemize}

\begin{center}\rule{0.5\linewidth}{0.5pt}\end{center}

\subsection*{Troubleshooting}\label{docs__pandoc__latex__src__setup_accessibility_complete.md__troubleshooting-1}

\textbf{Display shows nothing / no braille output:}

\begin{itemize}
\tightlist
\item
  Verify the USB or Bluetooth connection and check that the display is powered on.
\item
  In your screen reader\textquotesingle s braille settings, confirm the correct display model is selected.
\item
  Try disconnecting and reconnecting. Restart the screen reader if needed.
\item
  Install or reinstall the vendor\textquotesingle s display driver.
\end{itemize}

\textbf{Display shows output but navigation is out of sync:}

\begin{itemize}
\tightlist
\item
  Toggle your screen reader\textquotesingle s braille mode off and on (usually in the braille settings).
\item
  Restart the screen reader (not the whole computer) and reopen the terminal.
\item
  Try a different terminal emulator (Windows Terminal, Command Prompt, PowerShell) to see if the issue is terminal-specific.
\end{itemize}

\textbf{Display driver not found:}

\begin{itemize}
\tightlist
\item
  Go to the vendor\textquotesingle s website (links below) and download the latest driver for your operating system version.
\item
  Some displays require firmware updates to work with newer Windows versions.
\end{itemize}

\textbf{Routing keys not moving the cursor:}

\begin{itemize}
\tightlist
\item
  In your screen reader\textquotesingle s braille settings, confirm that cursor routing is enabled.
\item
  In NVDA, this is under Preferences \textgreater{} Settings \textgreater{} Braille \textgreater{} enable "Move system cursor when routing review cursor".
\item
  In JAWS, check the routing settings in the JAWS braille configuration.
\end{itemize}

\begin{center}\rule{0.5\linewidth}{0.5pt}\end{center}

\subsection*{Vendor Resources}\label{docs__pandoc__latex__src__setup_accessibility_complete.md__vendor-resources}

{\def\LTcaptype{none} % do not increment counter
\begin{longtable}[]{@{}
  >{\raggedright\arraybackslash}p{(\linewidth - 4\tabcolsep) * \real{0.1488}}
  >{\raggedright\arraybackslash}p{(\linewidth - 4\tabcolsep) * \real{0.3306}}
  >{\raggedright\arraybackslash}p{(\linewidth - 4\tabcolsep) * \real{0.5207}}@{}}
\toprule\noalign{}
\begin{minipage}[b]{\linewidth}\raggedright
Vendor
\end{minipage} & \begin{minipage}[b]{\linewidth}\raggedright
Products
\end{minipage} & \begin{minipage}[b]{\linewidth}\raggedright
Support / Drivers
\end{minipage} \\
\midrule\noalign{}
\endhead
\bottomrule\noalign{}
\endlastfoot
HumanWare & Brailliant, BrailleNote Touch+ &
https://www.humanware.com/ \\
Freedom Scientific & Focus Blue series &
https://www.freedomscientific.com/Products/Blindness/FocusBlue/ \\
HIMS & BrailleSense, Braille EDGE, Smart Beetle &
https://www.hims-inc.com/ \\
Optelec & Braille STAR, Easy Link & https://www.optelec.com/ \\
Orbit Research & Graphiti, Orbit 20/40 &
https://www.orbitresearch.com/ \\
APH & Monarch & https://www.aph.org/ \\
Dot Inc. & DotPad & https://dotincorp.com/ \\
Dolphin & GuideConnect, braille accessories &
https://yourdolphin.com/ \\
\end{longtable}
}

\begin{center}\rule{0.5\linewidth}{0.5pt}\end{center}

\subsection*{Connecting a Braille Display: Summary Table}\label{docs__pandoc__latex__src__setup_accessibility_complete.md__connecting-a-braille-display-summary-table}

{\def\LTcaptype{none} % do not increment counter
\begin{longtable}[]{@{}
  >{\raggedright\arraybackslash}p{(\linewidth - 4\tabcolsep) * \real{0.1491}}
  >{\raggedright\arraybackslash}p{(\linewidth - 4\tabcolsep) * \real{0.5175}}
  >{\raggedright\arraybackslash}p{(\linewidth - 4\tabcolsep) * \real{0.3333}}@{}}
\toprule\noalign{}
\begin{minipage}[b]{\linewidth}\raggedright
Screen Reader
\end{minipage} & \begin{minipage}[b]{\linewidth}\raggedright
Where to Find Braille Settings
\end{minipage} & \begin{minipage}[b]{\linewidth}\raggedright
Recommended Braille Table for Code
\end{minipage} \\
\midrule\noalign{}
\endhead
\bottomrule\noalign{}
\endlastfoot
NVDA &
NVDA Menu \textgreater{} Preferences \textgreater{} Settings \textgreater{} Braille
& Unified English Braille Grade 1 \\
JAWS &
Insert+F3 \textgreater{} Options \textgreater{} Basics \textgreater{} Braille (or Settings Center)
& English Braille Grade 1 (uncontracted) \\
Windows Narrator &
Settings \textgreater{} Accessibility \textgreater{} Narrator \textgreater{} Braille
& English Grade 1 (uncontracted) \\
Dolphin SuperNova &
CapsLock+SpaceBar \textgreater{} Braille \textgreater{} Display and Translation
& Grade 1 (uncontracted) \\
\end{longtable}
}

This page gives a practical overview. Always consult your braille display vendor\textquotesingle s documentation and your screen reader\textquotesingle s braille guide for device-specific setup steps and advanced configuration options.

\begin{center}\rule{0.5\linewidth}{0.5pt}\end{center}

\protect\phantomsection\label{docs__pandoc__latex__src__setup_accessibility_complete.md__editor-setup}{}

\section*{Part 4: Editor Selection and Accessibility Setup Guide}\label{docs__pandoc__latex__src__setup_accessibility_complete.md__part-4-editor-selection-and-accessibility-setup-guide}

\begin{center}\rule{0.5\linewidth}{0.5pt}\end{center}

\subsection*{Table of Contents}\label{docs__pandoc__latex__src__setup_accessibility_complete.md__table-of-contents-1}

\begin{enumerate}
\tightlist
\item
  \hyperref[docs__pandoc__latex__src__setup_accessibility_complete.md__editor-comparison]{Editor Comparison}
\item
  \hyperref[docs__pandoc__latex__src__setup_accessibility_complete.md__editor-setup-for-3dmake]{Editor Setup for 3dMake}
\item
  \hyperref[docs__pandoc__latex__src__setup_accessibility_complete.md__screen-reader-indent-announcement-configuration]{Screen Reader Indent Announcement Configuration}
\item
  \hyperref[docs__pandoc__latex__src__setup_accessibility_complete.md__curriculum-specific-editor-recommendations]{Curriculum-Specific Editor Recommendations}
\end{enumerate}

\subsection*{Editor Comparison}\label{docs__pandoc__latex__src__setup_accessibility_complete.md__editor-comparison}

\subsubsection*{Overview Table}\label{docs__pandoc__latex__src__setup_accessibility_complete.md__overview-table}

{\def\LTcaptype{none} % do not increment counter
\begin{longtable}[]{@{}
  >{\raggedright\arraybackslash}p{(\linewidth - 6\tabcolsep) * \real{0.2396}}
  >{\raggedright\arraybackslash}p{(\linewidth - 6\tabcolsep) * \real{0.1562}}
  >{\raggedright\arraybackslash}p{(\linewidth - 6\tabcolsep) * \real{0.2500}}
  >{\raggedright\arraybackslash}p{(\linewidth - 6\tabcolsep) * \real{0.3542}}@{}}
\toprule\noalign{}
\begin{minipage}[b]{\linewidth}\raggedright
Feature
\end{minipage} & \begin{minipage}[b]{\linewidth}\raggedright
Notepad
\end{minipage} & \begin{minipage}[b]{\linewidth}\raggedright
Notepad++
\end{minipage} & \begin{minipage}[b]{\linewidth}\raggedright
Visual Studio Code
\end{minipage} \\
\midrule\noalign{}
\endhead
\bottomrule\noalign{}
\endlastfoot
Cost & Free (built-in) & Free & Free \\
Learning Curve & Minimal & Low & Moderate \\
Screen Reader Support & Good (basic) & Good (syntax features) &
Excellent (built-in accessibility) \\
Extension/Plugin System & None & Limited & Extensive \\
Keyboard Navigation & Good & Good & Excellent \\
Customization & None & Moderate & Very high \\
Performance & Excellent & Very good & Good \\
Syntax Highlighting & None & Yes (OpenSCAD available) &
Yes (OpenSCAD available) \\
Terminal Integration & None & None & Built-in \\
Real-time Feedback & None & None & Yes (with extensions) \\
Hot Key Customization & Limited & Good & Excellent \\
File Size Handling & Good & Good & Excellent \\
Built-in Debugging & None & None & Limited \\
\end{longtable}
}

\subsubsection*{Detailed Comparison}\label{docs__pandoc__latex__src__setup_accessibility_complete.md__detailed-comparison}

\paragraph*{Notepad}\label{docs__pandoc__latex__src__setup_accessibility_complete.md__notepad}

Advantages:

\begin{itemize}
\tightlist
\item
  Minimal interface with no distractions-excellent for absolute beginners
\item
  Very predictable behavior for screen reader users
\item
  Extremely fast file operations
\item
  No configuration required; works immediately
\item
  Pure text editing with no formatting surprises
\end{itemize}

Disadvantages:

\begin{itemize}
\tightlist
\item
  No syntax highlighting (OpenSCAD code appears as plain text)
\item
  No keyboard shortcuts for common editing tasks
\item
  Limited undo/redo capabilities compared to modern editors
\item
  No integrated terminal (requires separate command prompt window)
\item
  No way to run 3dm commands directly
\end{itemize}

Best For: Users who prefer absolute simplicity and want minimal cognitive load during learning phase

Screen Reader Experience: Windows Narrator reads all content clearly; JAWS and NVDA work well with standard keyboard navigation

\paragraph*{Notepad++}\label{docs__pandoc__latex__src__setup_accessibility_complete.md__notepad-1}

Advantages:

\begin{itemize}
\tightlist
\item
  Lightweight and fast
\item
  Good syntax highlighting for OpenSCAD code
\item
  Customizable interface with configurable keyboard shortcuts
\item
  Tab management for multiple files
\item
  Better organization than Notepad for managing projects
\item
  Good screen reader support for basic operations
\end{itemize}

Disadvantages:

\begin{itemize}
\tightlist
\item
  No built-in terminal (requires external command prompt)
\item
  Plugin system is limited compared to VSCode
\item
  Screen reader experience with syntax highlighting features can be inconsistent
\item
  Less extensive keyboard customization than VSCode
\item
  No integrated development environment features
\end{itemize}

Best For: Users who want a lightweight editor with syntax highlighting but prefer not to use a full IDE

Screen Reader Experience: JAWS and NVDA handle navigation well; Windows Narrator works but may struggle with complex UI elements

\paragraph*{Visual Studio Code}\label{docs__pandoc__latex__src__setup_accessibility_complete.md__visual-studio-code}

Advantages:

\begin{itemize}
\tightlist
\item
  Extensive built-in accessibility features (accessibility inspector, keyboard navigation shortcuts)
\item
  Excellent OpenSCAD extension available (scad-preview)
\item
  Integrated terminal allows running 3dm commands without context switching
\item
  Powerful keyboard shortcut customization
\item
  Rich extension ecosystem for workflow enhancement
\item
  Remote development capabilities
\item
  Native support for multiple projects and workspaces
\item
  Excellent search and find/replace functionality
\end{itemize}

Disadvantages:

\begin{itemize}
\tightlist
\item
  Steeper learning curve than Notepad or Notepad++
\item
  Requires initial configuration for accessibility
\item
  More resource-intensive than lighter editors
\item
  Built-in terminal can be distracting for some users (alternative: Alt-Tab to standalone terminal)
\end{itemize}

Best For: Users building comprehensive accessibility workflow and wanting integrated development environment

Screen Reader Experience: Excellent; built specifically with accessibility in mind. NVDA, JAWS, Windows Narrator, and Dolphin all receive high-quality support.

\subsection*{Editor Setup for 3dMake}\label{docs__pandoc__latex__src__setup_accessibility_complete.md__editor-setup-for-3dmake}

\subsubsection*{Setting Default Editor in 3dMake}\label{docs__pandoc__latex__src__setup_accessibility_complete.md__setting-default-editor-in-3dmake}

The 3dMake tool uses your system default text editor. However it can be changed with an edit to the global configuration file by typing this into any terminal.

\begin{lstlisting}[style=Alabaster]
3dm edit-global-config

\end{lstlisting}

You get a file like this in your default text edit program, typically notepad on Windows.

\begin{lstlisting}[style=Alabaster]
view = "3sil"
model_name = "main"
auto_Start_prints = true
printer_profile = "bambu_labs_X1_carbon"

\end{lstlisting}

Add one of the following lines to the end of the file:

\texttt{editor\ =\ "code"} for VSCode
\texttt{editor\ =\ "notepad"} for notepad (already the default)
\texttt{editor\ =\ \textquotesingle{}\textquotesingle{}\textquotesingle{}C:\textbackslash{}Program\ Files\ (x86)\textbackslash{}Notepad++\textbackslash{}notepad++.exe\textquotesingle{}\textquotesingle{}\textquotesingle{}} for Notepad++ (the triple apostrophes mean you do not have to escape the spaces)

\subsubsection*{Notepad Setup}\label{docs__pandoc__latex__src__setup_accessibility_complete.md__notepad-setup}

Configuration Steps:

\begin{enumerate}
\item
  Open 3dMake-generated .scad files directly:

  \begin{itemize}
  \tightlist
  \item
    Navigate to your project folder in File Explorer
  \item
    Right-click the \texttt{.scad} file -\textgreater{} "Open with" -\textgreater{} "Notepad"
  \item
    File opens immediately for editing
  \end{itemize}
\item
  Edit and Save:

  \begin{itemize}
  \tightlist
  \item
    Make changes to your code
  \item
    Press \texttt{Ctrl+S} to save
  \item
    File updates automatically if 3dMake renders in background
  \end{itemize}
\item
  Run 3dMake Commands:

  \begin{itemize}
  \tightlist
  \item
    Save your edits (Ctrl+S)
  \item
    Alt-Tab to Command Prompt or PowerShell window
  \item
    Navigate to project directory: \texttt{cd\ C:\textbackslash{}path\textbackslash{}to\textbackslash{}project}
  \item
    Run 3dMake command: \texttt{3dMake\ render\ project.scad}
  \end{itemize}
\end{enumerate}

Keyboard Shortcuts:

\begin{itemize}
\tightlist
\item
  \texttt{Ctrl+H}: Find and Replace
\item
  \texttt{Ctrl+G}: Go to Line (newer versions)
\item
  \texttt{Ctrl+Z}: Undo
\item
  \texttt{Ctrl+Y}: Redo
\end{itemize}

\subsubsection*{Notepad++ Setup}\label{docs__pandoc__latex__src__setup_accessibility_complete.md__notepad-setup-1}

Installation and Configuration:

\begin{enumerate}
\item
  Download from: \url{https://notepad-plus-plus.org/downloads/}
\item
  Choose: Standard Installer (for automatic system integration)
\item
  Configure OpenSCAD Language Support:

  \begin{itemize}
  \tightlist
  \item
    Open Notepad++
  \item
    Language -\textgreater{} User Defined Language -\textgreater{} Import... (if OpenSCAD UDL available)
  \item
    Or manually set syntax highlighting:

    \begin{itemize}
    \tightlist
    \item
      Language -\textgreater{} OpenSCAD (if available in language menu)
    \item
      Otherwise Language -\textgreater{} C++ (provides similar highlighting)
    \end{itemize}
  \end{itemize}
\item
  Customize for Accessibility:

  \begin{itemize}
  \tightlist
  \item
    Settings -\textgreater{} Preferences -\textgreater{} General
  \item
    Check: "Minimize to system tray" (optional)
  \item
    Settings -\textgreater{} Preferences -\textgreater{} MISC.
  \item
    Ensure Word Wrap is set to preference
  \item
    Settings -\textgreater{} Preferences -\textgreater{} Backup
  \item
    Enable regular backups of your work
  \end{itemize}
\item
  Set as Default Editor (see registry method above)
\end{enumerate}

Recommended Keyboard Shortcuts (User-Defined):

Create custom shortcuts by:

\begin{itemize}
\tightlist
\item
  Settings -\textgreater{} Shortcut Mapper
\item
  Add shortcuts for frequently used actions:

  \begin{itemize}
  \tightlist
  \item
    Save and Switch to Terminal: Alt+T
  \item
    Copy File Path: Ctrl+Shift+C
  \end{itemize}
\end{itemize}

Tab Management:

\begin{itemize}
\tightlist
\item
  Open multiple files: Each opens in a separate tab
\item
  Switch between tabs: Ctrl+Tab (next) / Ctrl+Shift+Tab (previous)
\item
  Close tab: Ctrl+W
\end{itemize}

Running 3dMake Commands:

\begin{itemize}
\tightlist
\item
  Save file with Ctrl+S
\item
  Alt-Tab to standalone terminal or command prompt
\item
  Run: \texttt{3dMake\ render\ filename.scad}
\end{itemize}

\subsubsection*{Visual Studio Code Setup}\label{docs__pandoc__latex__src__setup_accessibility_complete.md__visual-studio-code-setup}

Installation:

\begin{enumerate}
\tightlist
\item
  Download from: \url{https://code.visualstudio.com/} or type \texttt{winget\ search\ VSCode} in a terminal and then \texttt{winget\ install\ } whichever option you prefer
\item
  Install: Run installer and follow prompts
\item
  Launch: Open VSCode
\end{enumerate}

Enable OpenSCAD Support:

\begin{enumerate}
\tightlist
\item
  Open Extensions (Ctrl+Shift+X)
\item
  Search: "OpenSCAD"
\item
  Install: "scad-preview" by Antyos
\item
  Install: "OpenSCAD Syntax Highlighter" (optional, for better syntax highlighting)
\end{enumerate}

Initial Accessibility Configuration:

\begin{enumerate}
\item
  Open Settings: Ctrl+,
\item
  Search: "accessibility"
\item
  Enable Key Settings:

  \begin{itemize}
  \tightlist
  \item
    Accessible View: Toggle ON
  \item
    Screen Reader: Select your screen reader (NVDA, JAWS, Narrator, Dolphin)
  \item
    Keyboard Navigation: Ensure enabled
  \item
    Bracket Pair Guides: Can help with code structure understanding
  \end{itemize}
\item
  Configure Editor Font:

  \begin{itemize}
  \tightlist
  \item
    Search: "editor.fontSize"
  \item
    Set to comfortable size (recommend 14-16 for better readability)
  \item
    Search: "editor.fontFamily"
  \item
    Select monospace font (e.g., "Consolas" or "Courier New")
  \end{itemize}
\end{enumerate}

Set as Default Editor (see registry method above)

VSCode Terminal Options:

\paragraph*{Option 1: Using Built-in Terminal (Less Accessible)}\label{docs__pandoc__latex__src__setup_accessibility_complete.md__option-1-using-built-in-terminal-less-accessible}

\begin{enumerate}
\tightlist
\item
  View -\textgreater{} Terminal (or Ctrl+`)
\item
  Terminal opens at bottom of VSCode window
\item
  Run commands: \texttt{3dMake\ render\ filename.scad}
\item
  Note: Switching focus between editor and terminal requires Tab navigation, which can be cumbersome for screen reader users
\end{enumerate}

Keyboard Navigation:

\begin{itemize}
\tightlist
\item
  Ctrl+` : Toggle terminal visibility
\item
  Ctrl+Shift+` : Create new terminal
\item
  Alt+\^{}/v : Switch between terminals
\end{itemize}

\paragraph*{Option 2: Alt-Tab to Standalone Terminal (RECOMMENDED for Accessibility)}\label{docs__pandoc__latex__src__setup_accessibility_complete.md__option-2-alt-tab-to-standalone-terminal-recommended-for-accessibility}

\begin{enumerate}
\item
  Keep Command Prompt/PowerShell open:

  \begin{itemize}
  \tightlist
  \item
    Open Command Prompt (Win+R, type "cmd", Enter)
  \item
    Position window or minimize to taskbar
  \item
    Navigate to project directory: \texttt{cd\ C:\textbackslash{}path\textbackslash{}to\textbackslash{}project}
  \end{itemize}
\item
  From VSCode, switch terminals:

  \begin{itemize}
  \tightlist
  \item
    Alt+Tab to Command Prompt
  \item
    Run command: \texttt{3dMake\ render\ filename.scad}
  \item
    Alt+Tab back to VSCode
  \item
    Continue editing
  \end{itemize}
\end{enumerate}

Why This Is More Accessible:

\begin{itemize}
\tightlist
\item
  Screen reader focus switches clearly between two applications
\item
  Terminal output is read without VSCode context interference
\item
  Cleaner context switching for command-line workflows
\item
  Easier to diagnose issues when editor and terminal are separate
\end{itemize}

Keyboard Shortcuts for Common Tasks:

{\def\LTcaptype{none} % do not increment counter
\begin{longtable}[]{@{}ll@{}}
\toprule\noalign{}
Action & Shortcut \\
\midrule\noalign{}
\endhead
\bottomrule\noalign{}
\endlastfoot
Save & Ctrl+S \\
Find & Ctrl+F \\
Find and Replace & Ctrl+H \\
Go to Line & Ctrl+G \\
Open File & Ctrl+O \\
Open Folder & Ctrl+K, Ctrl+O \\
Open Terminal & Ctrl+` \\
Alt-Tab to Another Window & Alt+Tab \\
Command Palette & Ctrl+Shift+P \\
\end{longtable}
}

Project Organization in VSCode:

\begin{enumerate}
\item
  Open Project Folder:

  \begin{itemize}
  \tightlist
  \item
    File -\textgreater{} Open Folder (Ctrl+K, Ctrl+O)
  \item
    Select your project directory
  \item
    All project files appear in Explorer sidebar
  \end{itemize}
\item
  File Navigation:

  \begin{itemize}
  \tightlist
  \item
    Press Ctrl+P for Quick Open
  \item
    Type filename to search and jump to file
  \item
    Press Enter to open
  \end{itemize}
\item
  Quick Switch Between Files:

  \begin{itemize}
  \tightlist
  \item
    Ctrl+Tab : Open recent files list
  \item
    Arrow keys to select
  \item
    Enter to open
  \end{itemize}
\end{enumerate}

\subsection*{Screen Reader Indent Announcement Configuration}\label{docs__pandoc__latex__src__setup_accessibility_complete.md__screen-reader-indent-announcement-configuration}

Proper indent announcement is critical for OpenSCAD development, as indentation indicates code nesting and structure.

\subsubsection*{NVDA (NonVisual Desktop Access)}\label{docs__pandoc__latex__src__setup_accessibility_complete.md__nvda-nonvisual-desktop-access-1}

Enable Indent Announcement:

\begin{enumerate}
\tightlist
\item
  Open NVDA Menu: Alt+N or right-click NVDA icon
\item
  Preferences -\textgreater{} Settings (Ctrl+Comma)
\item
  Document Formatting: Tab to it
\item
  Check: "Report indentation"
\item
  In the "Indentation reporting" dropdown: Select "Tones and speech"
\item
  Tone Description: NVDA will announce indent level as progressively higher tones (or speaking indent amount)
\item
  Apply: Click OK
\end{enumerate}

Additional Tab Stop Configuration:

\begin{enumerate}
\tightlist
\item
  Preferences -\textgreater{} Settings -\textgreater{} Document Formatting
\item
  Check: "Report line indentation"
\item
  This will announce: "Indent level 4" or similar as you navigate code
\end{enumerate}

Testing:

\begin{itemize}
\tightlist
\item
  Open a \texttt{.scad} file with nested code (e.g., difference() \{ cube(); sphere(); \})
\item
  Press Down Arrow to move line by line
\item
  NVDA announces indentation level on each new indented line
\end{itemize}

Keyboard Control:

\begin{itemize}
\tightlist
\item
  NVDA+3 on Numpad: Cycles indent/outline level reporting
\end{itemize}

\subsubsection*{JAWS (Freedom Scientific)}\label{docs__pandoc__latex__src__setup_accessibility_complete.md__jaws-freedom-scientific}

Enable Indent Announcement:

\begin{enumerate}
\tightlist
\item
  Open JAWS Manager: Press JAWSKey+F2 (or right-click JAWS icon)
\item
  Utilities -\textgreater{} Settings Manager
\item
  Search: "Indent"
\item
  Look for setting: "Announce Indentation" or "Report Indentation"
\item
  Enable: Set to "Tones" or "Tones and Speech"
\item
  Speech Indent Announcement: Speak indent level
\item
  Tone Indent Announcement: Pitch increases with indent level
\end{enumerate}

Advanced Configuration (Custom Scripts):

If built-in settings don\textquotesingle t work:

\begin{enumerate}
\tightlist
\item
  JAWSKey+F2 -\textgreater{} Utilities -\textgreater{} Settings Manager
\item
  Search: "Line Breaks" or "Formatting"
\item
  Ensure: "Report line indentation" is enabled
\item
  Set tone adjustment: Higher pitch for deeper indents
\end{enumerate}

Testing:

\begin{itemize}
\tightlist
\item
  Open \texttt{.scad} file with nested code
\item
  Navigate with arrow keys
\item
  JAWS announces indent changes with tones or speech
\end{itemize}

Keyboard Shortcuts:

\begin{itemize}
\tightlist
\item
  JAWSKey+Alt+I: Toggle indent reporting
\item
  JAWSKey+Alt+Shift+I: Cycle between indent reporting modes (speech/tones/off)
\end{itemize}

\subsubsection*{Windows Narrator}\label{docs__pandoc__latex__src__setup_accessibility_complete.md__windows-narrator-2}

Enable Indent Announcement:

\begin{enumerate}
\tightlist
\item
  Open Settings: Win+I
\item
  Ease of Access -\textgreater{} Narrator
\item
  Advanced Options: Scroll down
\item
  Check: "Report indentation"
\item
  Indentation Reporting: Select "Tones" (less intrusive) or "Speech" (explicit)
\item
  Apply settings
\end{enumerate}

Narrator Keyboard Shortcuts:

\begin{itemize}
\tightlist
\item
  Narrator+Page Down: Read from current position to end of window
\item
  Narrator+Alt+Arrow Keys: Navigate text
\item
  Narrator+V, I: Customize indentation reporting (in Narrator settings)
\end{itemize}

Testing:

\begin{itemize}
\tightlist
\item
  Open \texttt{.scad} file
\item
  Use Narrator+Page Down to read through code
\item
  Listen for indent tone changes or announcements
\end{itemize}

Note: Windows Narrator has fewer customization options than JAWS/NVDA; consider NVDA or JAWS for deeper indent control

\subsubsection*{Dolphin EasyConverter (Dolphin Screen Reader)}\label{docs__pandoc__latex__src__setup_accessibility_complete.md__dolphin-easyconverter-dolphin-screen-reader}

Enable Indent Announcement:

\begin{enumerate}
\tightlist
\item
  Open Dolphin Central: Right-click Dolphin icon or click Dolphin icon in taskbar
\item
  Utilities -\textgreater{} Settings -\textgreater{} Text Processing
\item
  Look for: "Indentation" section
\item
  Enable: "Announce indentation"
\item
  Mode: Select "Tones", "Speech", or "Tones and Speech"
\item
  Tone Pitch: Configure pitch increase for deeper indents
\item
  Apply
\end{enumerate}

ECO (Ease of Cursor Operation) Customization:

\begin{enumerate}
\tightlist
\item
  Dolphin Central -\textgreater{} Utilities -\textgreater{} ECO Settings
\item
  Text Options -\textgreater{} Indentation Reporting
\item
  Set preferred announcement style
\end{enumerate}

Testing:

\begin{itemize}
\tightlist
\item
  Open \texttt{.scad} file with indented code
\item
  Navigate with arrows
\item
  Dolphin announces indent changes
\end{itemize}

Keyboard Shortcuts:

\begin{itemize}
\tightlist
\item
  Ctrl+Dolphin+I: Toggle indent reporting on/off
\item
  Ctrl+Dolphin+Shift+I: Cycle indent reporting mode
\end{itemize}

\subsection*{Curriculum-Specific Editor Recommendations}\label{docs__pandoc__latex__src__setup_accessibility_complete.md__curriculum-specific-editor-recommendations}

\subsubsection*{For Absolute Beginners (Lesson 1-2)}\label{docs__pandoc__latex__src__setup_accessibility_complete.md__for-absolute-beginners-lesson-1-2}

Recommended: Notepad or Notepad++

Rationale:

\begin{itemize}
\tightlist
\item
  Minimal interface reduces cognitive load
\item
  Focus stays on learning OpenSCAD syntax, not editor features
\item
  Keyboard navigation is straightforward
\item
  Screen reader experience is predictable
\end{itemize}

Setup:

\begin{enumerate}
\tightlist
\item
  Use Notepad or Notepad++ as default editor
\item
  Configure screen reader indent announcement
\item
  Keep separate Command Prompt window open for 3dMake commands
\item
  Alt-Tab workflow between editor and terminal
\end{enumerate}

Workflow Example:

\begin{lstlisting}[style=Alabaster]
1. Open Command Prompt -> Navigate to project folder
2. Run: 3dMake new myproject
3. Alt+Tab to file explorer, open myproject.scad
4. Notepad++ opens file
5. Edit code
6. Ctrl+S to save
7. Alt+Tab to Command Prompt
8. Run: 3dMake render myproject.scad
9. Check output, return to Notepad++ to refine code

\end{lstlisting}

\subsubsection*{For Intermediate Users (Lesson 3-6)}\label{docs__pandoc__latex__src__setup_accessibility_complete.md__for-intermediate-users-lesson-3-6}

Recommended: Notepad++ or VSCode

Notepad++:

\begin{itemize}
\tightlist
\item
  Adds project organization without overwhelming complexity
\item
  Tab support for managing multiple files
\item
  Syntax highlighting improves code understanding
\item
  Still lightweight and predictable
\end{itemize}

VSCode:

\begin{itemize}
\tightlist
\item
  Opens doors to more sophisticated workflows
\item
  Extension system enables advanced features (OpenSCAD preview)
\item
  Keyboard customization becomes valuable
\item
  Terminal integration useful but use Alt-Tab method
\end{itemize}

Setup Decision Tree:

\begin{itemize}
\tightlist
\item
  Choose Notepad++ if: You prefer simplicity and want to focus on code logic
\item
  Choose VSCode if: You\textquotesingle re ready to invest time learning editor features for long-term benefit
\end{itemize}

Workflow Example with VSCode:

\begin{lstlisting}[style=Alabaster]
1. Open VSCode (folder view of project)
2. Press Ctrl+P to open file search
3. Type filename and press Enter to open
4. Edit code with autocomplete
5. Ctrl+H for find/replace across project
6. Ctrl+S to save
7. Alt+Tab to Command Prompt
8. Run: 3dMake render filename.scad
9. Alt+Tab back to VSCode to refine code

\end{lstlisting}

\subsubsection*{For Advanced Users (Lesson 7-11)}\label{docs__pandoc__latex__src__setup_accessibility_complete.md__for-advanced-users-lesson-7-11}

Recommended: Visual Studio Code

Rationale:

\begin{itemize}
\tightlist
\item
  Powerful search/replace across large projects
\item
  Extension system enables specialized workflows
\item
  Keyboard customization reaches full potential
\item
  Workspace management for complex projects
\item
  Debugging capabilities aid troubleshooting
\end{itemize}

Advanced Setup:

\begin{enumerate}
\item
  Create custom keyboard shortcuts:

  \begin{itemize}
  \tightlist
  \item
    Ctrl+Alt+R: Save and render current file
  \item
    Ctrl+Alt+P: Preview (if using scad-preview extension)
  \end{itemize}
\item
  Install Additional Extensions:

  \begin{itemize}
  \tightlist
  \item
    "scad-preview": Real-time 3D preview
  \item
    "Better Comments": Categorize comments with colors/tones
  \item
    "Bracket Pair Colorizer": Visual/tonal bracket matching
  \item
    "GitLens": Track code changes over time
  \end{itemize}
\item
  Create Task Runner for Common Commands:

  \begin{itemize}
  \tightlist
  \item
    Ctrl+Shift+B: Configure build task to run 3dMake
  \item
    Create separate tasks for render, export, etc.
  \end{itemize}
\item
  Use Workspaces:

  \begin{itemize}
  \tightlist
  \item
    File -\textgreater{} Save Workspace As...
  \item
    Save project-specific workspace with all settings
  \item
    Reopen same workspace configuration automatically
  \end{itemize}
\end{enumerate}

Advanced Workflow Example:

\begin{lstlisting}[style=Alabaster]
1. Open VSCode with project workspace
2. Ctrl+Shift+P -> Run Task -> "3dMake Render Current"
3. Renders file and shows output
4. Use scad-preview extension for real-time 3D view
5. Edit code with advanced search/replace
6. Ctrl+Alt+R saves and renders automatically
7. Use version control (Git) for tracking changes

\end{lstlisting}

\subsection*{Quick Reference: Editor Comparison for Curriculum}\label{docs__pandoc__latex__src__setup_accessibility_complete.md__quick-reference-editor-comparison-for-curriculum}

\subsubsection*{Foundations (Lessons 1-2)}\label{docs__pandoc__latex__src__setup_accessibility_complete.md__foundations-lessons-1-2}

\begin{itemize}
\tightlist
\item
  Primary: Notepad or Notepad++
\item
  Focus: Learn syntax and basic concepts
\item
  Terminal: Standalone Command Prompt (Alt-Tab)
\end{itemize}

\subsubsection*{Core Skills (Lessons 3-6)}\label{docs__pandoc__latex__src__setup_accessibility_complete.md__core-skills-lessons-3-6}

\begin{itemize}
\tightlist
\item
  Primary: Notepad++ or VSCode
\item
  New Features: Begin using editor syntax highlighting
\item
  Terminal: Standalone (continue Alt-Tab method)
\item
  Skills: File organization, search/replace basics
\end{itemize}

\subsubsection*{Advanced Projects (Lessons 7-11)}\label{docs__pandoc__latex__src__setup_accessibility_complete.md__advanced-projects-lessons-7-11}

\begin{itemize}
\tightlist
\item
  Primary: VSCode (strongly recommended)
\item
  Advanced: Use extensions, real-time preview, complex project management
\item
  Terminal: Choose Alt-Tab or built-in based on preference
\item
  Skills: Workspaces, task automation, version control
\end{itemize}

\subsection*{Troubleshooting Common Issues}\label{docs__pandoc__latex__src__setup_accessibility_complete.md__troubleshooting-common-issues}

\subsubsection*{Problem: Screen Reader Not Announcing Indent}\label{docs__pandoc__latex__src__setup_accessibility_complete.md__problem-screen-reader-not-announcing-indent}

Solution:

\begin{enumerate}
\tightlist
\item
  Verify indent announcement is enabled in screen reader settings (see above)
\item
  Test with existing \texttt{.scad} file with clear indentation
\item
  Try different announcement modes (tones vs. speech)
\item
  Restart screen reader: Alt+Ctrl+N (NVDA) or app restart (JAWS)
\end{enumerate}

\subsubsection*{Problem: 3dMake Commands Not Running from VSCode Terminal}\label{docs__pandoc__latex__src__setup_accessibility_complete.md__problem-3dmake-commands-not-running-from-vscode-terminal}

Solution:

\begin{enumerate}
\tightlist
\item
  Ensure 3dMake is in your system PATH
\item
  Use standalone terminal instead (Alt-Tab method) - more reliable
\item
  In VSCode terminal, manually navigate to correct directory first
\item
  Verify command syntax: \texttt{3dMake\ render\ filename.scad}
\end{enumerate}

\subsubsection*{Problem: File Not Saving in Editor}\label{docs__pandoc__latex__src__setup_accessibility_complete.md__problem-file-not-saving-in-editor}

Solution:

\begin{enumerate}
\tightlist
\item
  Verify you pressed Ctrl+S
\item
  Check file permissions on project folder
\item
  Try "Save As" instead
\item
  Ensure filename includes \texttt{.scad} extension
\end{enumerate}

\subsubsection*{Problem: Syntax Highlighting Not Working}\label{docs__pandoc__latex__src__setup_accessibility_complete.md__problem-syntax-highlighting-not-working}

Solution:

\begin{enumerate}
\tightlist
\item
  Verify file has \texttt{.scad} extension
\item
  In Notepad++: Language menu -\textgreater{} select OpenSCAD or C++
\item
  In VSCode: Install OpenSCAD extension (search Extensions)
\item
  Restart editor
\end{enumerate}

\subsubsection*{Problem: Alt-Tab Not Switching Between Windows}\label{docs__pandoc__latex__src__setup_accessibility_complete.md__problem-alt-tab-not-switching-between-windows}

Solution:

\begin{enumerate}
\tightlist
\item
  Ensure Command Prompt is open and minimized (not closed)
\item
  Press Alt+Tab and hold briefly to see window switcher
\item
  Use Alt+Tab multiple times if more than 2 windows open
\item
  Alternatively, click taskbar directly (Alt+Tab usually more accessible)
\end{enumerate}

\subsection*{Next Steps}\label{docs__pandoc__latex__src__setup_accessibility_complete.md__next-steps}

After completing this setup guide:

\begin{enumerate}
\tightlist
\item
  Choose your editor based on the recommendations for your skill level
\item
  Configure screen reader indent announcement immediately (critical for code structure understanding)
\item
  Set editor as default for \texttt{.scad} files
\item
  Test with a simple file: Create a test project and edit it
\item
  Practice Alt-Tab workflow before moving to Lesson 1
\item
  Document your setup in a personal note for reference
\end{enumerate}

You are now ready to begin Lesson 1: Environmental Configuration and Developer Workflow

\subsection*{Additional Resources}\label{docs__pandoc__latex__src__setup_accessibility_complete.md__additional-resources-2}

\begin{itemize}
\tightlist
\item
  NVDA Documentation: \url{https://www.nvaccess.org/documentation/}
\item
  JAWS Documentation: \url{https://www.freedomscientific.com/products/software/jaws/}
\item
  VSCode Accessibility: \url{https://code.visualstudio.com/docs/editor/accessibility}
\item
  Windows Narrator Guide: \url{https://support.microsoft.com/en-us/help/22798/windows-11-narrator-get-started}
\item
  Notepad++ Documentation: \url{https://notepad-plus-plus.org/online-help/}
\end{itemize}

\section{Safety Protocols and the Physical Fabrication Interface}\label{docs__pandoc__latex__src__3dmake_foundation__lessons_3dmake_5__lessons_3dmake_5.md__safety-protocols-and-the-physical-fabrication-interface}

\textbf{Estimated Time:} 90--120 minutes\\
\textbf{Course:} 3dMake Certification --- High School Distance Learning Track

\begin{center}\rule{0.5\linewidth}{0.5pt}\end{center}

\subsection*{Before You Start}\label{docs__pandoc__latex__src__3dmake_foundation__lessons_3dmake_5__lessons_3dmake_5.md__before-you-start}

Every lesson so far has been entirely digital: code on a screen, commands in a terminal, virtual geometry. Now we\textquotesingle re moving into the physical world. 3D printers are machines that get very hot, draw substantial electrical current, and produce airborne particles and vapors. This lesson exists because these machines are genuinely safe when you know what you\textquotesingle re doing, and genuinely dangerous when you don\textquotesingle t.

The good news is that 3D printing safety is not complicated. The hazards are well-understood, the mitigations are straightforward, and the incidents that do happen are almost always caused by ignoring simple rules. The goal of this lesson is to make sure you know the rules and understand \emph{why} they exist --- not just as things to memorize for a quiz, but as habits you can apply in any shop or makerspace you ever work in.

\begin{center}\rule{0.5\linewidth}{0.5pt}\end{center}

\subsection*{Learning Objectives}\label{docs__pandoc__latex__src__3dmake_foundation__lessons_3dmake_5__lessons_3dmake_5.md__learning-objectives}

By the end of this lesson you will be able to:

\begin{itemize}
\tightlist
\item
  Identify and apply the Hierarchy of Controls for 3D printing hazards
\item
  Apply fire, electrical, and fume safety protocols for FDM printing
\item
  Understand filament-specific hazards and safe material handling
\item
  Follow a safe printer startup, monitoring, and shutdown sequence
\item
  Select appropriate filament materials for given projects
\item
  Use pre-flight build scripts to estimate cost before committing to a print
\item
  Design parts that minimize safety risks during printing
\end{itemize}

\begin{center}\rule{0.5\linewidth}{0.5pt}\end{center}

\subsection*{Concept 1: The Hierarchy of Controls}\label{docs__pandoc__latex__src__3dmake_foundation__lessons_3dmake_5__lessons_3dmake_5.md__concept-1-the-hierarchy-of-controls}

\subsubsection*{Why We Don\textquotesingle t Just "Be Careful"}\label{docs__pandoc__latex__src__3dmake_foundation__lessons_3dmake_5__lessons_3dmake_5.md__why-we-dont-just-be-careful}

When it comes to workplace and lab safety, the common advice "just be careful" is the least effective approach available. It relies entirely on individual attention, and attention is inconsistent --- everyone gets tired, distracted, or hurried sometimes. Professional safety systems are designed to work even when individual attention lapses.

The \textbf{Hierarchy of Controls} is a framework developed by OSHA (the U.S. Occupational Safety and Health Administration) and adopted worldwide. It ranks safety strategies from most effective to least effective, and it says: apply the most effective strategy you can afford, and layer additional strategies below it.

The five levels, from most to least effective:

\textbf{1. Elimination} --- Remove the hazard entirely. If a filament requires a sealed enclosure with active filtration and you\textquotesingle re in an open classroom, the safest choice is simply not to use that filament. If a print requires supports made from difficult material, redesign the print to not need supports. Elimination always wins.

\textbf{2. Substitution} --- Replace the hazard with a less dangerous one. Use PLA instead of ABS. Use a PEI spring steel bed instead of a glass bed that can shatter. Substitution keeps the functionality while reducing the risk.

\textbf{3. Engineering Controls} --- Build protection into the environment. A fume enclosure with a HEPA + activated carbon filter. Thermal runaway protection in the printer firmware. A smoke detector mounted above the printer. These controls work automatically without relying on human behavior.

\textbf{4. Administrative Controls} --- Rules, procedures, and training. The startup checklist, the monitoring schedule, the training program you\textquotesingle re in right now. Administrative controls are less reliable than engineering ones because they depend on people following rules consistently.

\textbf{5. Personal Protective Equipment (PPE)} --- Gloves, goggles, respirators. PPE is the last resort --- it protects one person, it can fail or be used incorrectly, and it doesn\textquotesingle t actually remove the hazard. Always apply higher-level controls first and use PPE as a supplement.

\textbf{The key insight:} many beginners think that putting on gloves and goggles makes them safe. In the hierarchy, PPE is the weakest protection. Real safety comes from eliminating or substituting hazards, and then engineering the remaining risks away. PPE is what you use when those options aren\textquotesingle t available.

\begin{center}\rule{0.5\linewidth}{0.5pt}\end{center}

\subsection*{Step 1 --- Fire and Electrical Safety}\label{docs__pandoc__latex__src__3dmake_foundation__lessons_3dmake_5__lessons_3dmake_5.md__step-1--fire-and-electrical-safety}

\subsubsection*{The Golden Rule}\label{docs__pandoc__latex__src__3dmake_foundation__lessons_3dmake_5__lessons_3dmake_5.md__the-golden-rule}

\textbf{Never leave an FDM printer completely unattended during the first layer.} Print failures most often happen at the start --- the first layer either bonds to the bed or it doesn\textquotesingle t. If it doesn\textquotesingle t bond and nobody is watching, the printer may continue for hours, dragging filament across the bed in a mess called "spaghetti" that can tangle around the print head, cause jams, or --- in rare worst cases --- create a fire hazard.

Be present and watching for at least the first 15 minutes of every print. For long prints, check in every 30 minutes or set up remote monitoring via OctoPrint or a camera.

\subsubsection*{Fire Hazards}\label{docs__pandoc__latex__src__3dmake_foundation__lessons_3dmake_5__lessons_3dmake_5.md__fire-hazards}

Fire risk in FDM printing comes from two main sources. First, \textbf{electrical failure}: wiring that cracks from repeated movement, loose connections at terminal blocks, or a faulty temperature sensor that allows the heater to run without limit. Second, \textbf{filament accumulation}: plastic that melts and collects on the hot nozzle can sometimes ignite if the build-up becomes large enough.

Fire safety protocols:

\begin{itemize}
\tightlist
\item
  Keep a \textbf{Class ABC fire extinguisher} within arm\textquotesingle s reach of every printer
\item
  Never run the printer while away from the building
\item
  Keep the print area clear of flammable materials: paper, fabric, solvent containers
\item
  Mount a \textbf{smoke detector} directly above the printer so rising smoke reaches it quickly
\item
  Know where the power switch is \emph{before} you start --- cutting power is the first response to any emergency
\item
  If smoke appears: cut power immediately, do not try to continue the print, and report the incident
\end{itemize}

\subsubsection*{Electrical Safety}\label{docs__pandoc__latex__src__3dmake_foundation__lessons_3dmake_5__lessons_3dmake_5.md__electrical-safety}

\begin{itemize}
\tightlist
\item
  Use a \textbf{surge protector} to protect the printer\textquotesingle s control board from voltage spikes
\item
  Inspect all wiring periodically, especially the flexible cable bundle that runs to the heated bed --- this bundle flexes with every print and can crack the insulation over time
\item
  \textbf{Thermal runaway protection must be enabled.} This firmware feature monitors the temperature sensor and shuts the heater off if the temperature climbs unexpectedly or the sensor fails. A printer without thermal runaway protection can overheat catastrophically. Never disable this protection. If your printer doesn\textquotesingle t have it, do not use it.
\item
  Never open the power supply enclosure or work on any electrical component while the printer is plugged in. The power supply contains mains voltage (120V or 240V) that can kill.
\end{itemize}

\begin{center}\rule{0.5\linewidth}{0.5pt}\end{center}

\subsection*{Step 2 --- Fume Safety}\label{docs__pandoc__latex__src__3dmake_foundation__lessons_3dmake_5__lessons_3dmake_5.md__step-2--fume-safety}

\subsubsection*{Why Fumes Matter}\label{docs__pandoc__latex__src__3dmake_foundation__lessons_3dmake_5__lessons_3dmake_5.md__why-fumes-matter}

All FDM printing produces some airborne emissions. These fall into two categories: \textbf{particulates} (tiny plastic particles, especially ultrafine particles below 0.1 microns) and \textbf{VOCs} (volatile organic compounds --- chemical vapors that off-gas from the hot plastic). Both can affect respiratory health with repeated exposure, and the risk varies significantly by material.

{\def\LTcaptype{none} % do not increment counter
\begin{longtable}[]{@{}
  >{\raggedright\arraybackslash}p{(\linewidth - 4\tabcolsep) * \real{0.1280}}
  >{\raggedright\arraybackslash}p{(\linewidth - 4\tabcolsep) * \real{0.0800}}
  >{\raggedright\arraybackslash}p{(\linewidth - 4\tabcolsep) * \real{0.7920}}@{}}
\toprule\noalign{}
\begin{minipage}[b]{\linewidth}\raggedright
Material
\end{minipage} & \begin{minipage}[b]{\linewidth}\raggedright
Fume Risk
\end{minipage} & \begin{minipage}[b]{\linewidth}\raggedright
Required Mitigation
\end{minipage} \\
\midrule\noalign{}
\endhead
\bottomrule\noalign{}
\endlastfoot
PLA & Low & Good room ventilation (open window or HVAC on) \\
PETG & Low--Medium &
Active ventilation; avoid enclosed spaces without air exchange \\
ABS & High &
Dedicated fume enclosure with HEPA + activated carbon filter required; NOT safe for open classrooms \\
ASA & High & Same as ABS \\
TPU & Medium &
Active ventilation; keep nozzle temperature at the minimum needed \\
Resin (SLA/MSLA) & Very High &
Full enclosure; nitrile gloves; UV-blocking goggles; no open-workspace use \\
\end{longtable}
}

\subsubsection*{The Safe Choice: PLA}\label{docs__pandoc__latex__src__3dmake_foundation__lessons_3dmake_5__lessons_3dmake_5.md__the-safe-choice-pla}

\textbf{PLA (Polylactic Acid) is the safest classroom filament available.} It is made from bio-derived sources (corn starch or sugarcane), prints at relatively low temperatures (190--220°C), produces the lowest emissions of any common FDM material, and is forgiving of settings variations. For any classroom project where the mechanical requirements don\textquotesingle t specifically demand a different material, default to PLA.

\subsubsection*{What "Good Ventilation" Actually Means}\label{docs__pandoc__latex__src__3dmake_foundation__lessons_3dmake_5__lessons_3dmake_5.md__what-good-ventilation-actually-means}

"Good ventilation" means that the air in the room is actively exchanging with outside air. This is not the same as "a window that\textquotesingle s crackable." Active ventilation means:

\begin{itemize}
\tightlist
\item
  A window open with a fan pulling air through the room, or
\item
  An HVAC system that is running and moving air, not just recirculating internal air
\end{itemize}

If you\textquotesingle re printing ABS or ASA, room ventilation is not sufficient. You need an enclosure with a filtered exhaust system --- a box around the printer with an internal fan that pulls fumes through a HEPA filter (for particles) and an activated carbon filter (for VOCs) before exhausting outside or into a large ventilated space.

\begin{center}\rule{0.5\linewidth}{0.5pt}\end{center}

\subsection*{Step 3 --- Material Selection}\label{docs__pandoc__latex__src__3dmake_foundation__lessons_3dmake_5__lessons_3dmake_5.md__step-3--material-selection}

\subsubsection*{Choosing the Right Filament}\label{docs__pandoc__latex__src__3dmake_foundation__lessons_3dmake_5__lessons_3dmake_5.md__choosing-the-right-filament}

The right material for a project depends on what the part needs to do and where it will be used. This table gives you a practical guide:

{\def\LTcaptype{none} % do not increment counter
\begin{longtable}[]{@{}
  >{\raggedright\arraybackslash}p{(\linewidth - 10\tabcolsep) * \real{0.0530}}
  >{\raggedright\arraybackslash}p{(\linewidth - 10\tabcolsep) * \real{0.0728}}
  >{\raggedright\arraybackslash}p{(\linewidth - 10\tabcolsep) * \real{0.1126}}
  >{\raggedright\arraybackslash}p{(\linewidth - 10\tabcolsep) * \real{0.1589}}
  >{\raggedright\arraybackslash}p{(\linewidth - 10\tabcolsep) * \real{0.1060}}
  >{\raggedright\arraybackslash}p{(\linewidth - 10\tabcolsep) * \real{0.4967}}@{}}
\toprule\noalign{}
\begin{minipage}[b]{\linewidth}\raggedright
Filament
\end{minipage} & \begin{minipage}[b]{\linewidth}\raggedright
Strength
\end{minipage} & \begin{minipage}[b]{\linewidth}\raggedright
Flexibility
\end{minipage} & \begin{minipage}[b]{\linewidth}\raggedright
Heat Resistance
\end{minipage} & \begin{minipage}[b]{\linewidth}\raggedright
Print Difficulty
\end{minipage} & \begin{minipage}[b]{\linewidth}\raggedright
Best Uses
\end{minipage} \\
\midrule\noalign{}
\endhead
\bottomrule\noalign{}
\endlastfoot
PLA & Medium & Rigid & Low (softens at 50--60°C) & Easy &
Prototypes, display models, artwork, educational parts \\
PETG & Medium-High & Slightly flexible & Medium (75--80°C) &
Easy--Medium &
Functional parts, food-adjacent applications, water contact \\
ABS & High & Semi-rigid & High (85--100°C) & Hard &
Hot environments, automotive interiors, post-processing (acetone smoothing) \\
TPU & Medium & Very flexible & Medium & Medium &
Gaskets, grips, phone cases, flexible joints \\
Nylon & Very High & Flexible & High & Hard &
Mechanical parts, gears, load-bearing functional parts \\
ASA & High & Semi-rigid & Very High (90°C+) & Hard &
Outdoor use, UV-exposed parts \\
\end{longtable}
}

\subsubsection*{What "Hard" to Print Actually Means}\label{docs__pandoc__latex__src__3dmake_foundation__lessons_3dmake_5__lessons_3dmake_5.md__what-hard-to-print-actually-means}

When a material is listed as "hard," that means it requires one or more of the following: a higher nozzle temperature, a heated enclosure to prevent warping, special bed adhesion (glue stick, hairspray, specialized plates), slower print speeds, or more careful tuning of retraction settings. Starting with a difficult material before you have PLA dialed in is a recipe for frustration and wasted filament.

\textbf{The practical rule:} use PLA until you have a specific reason not to. When you need better heat resistance, try PETG next. ABS and specialty materials come later, with appropriate safety infrastructure.

\subsubsection*{PLA and Sustainability --- A Realistic Picture}\label{docs__pandoc__latex__src__3dmake_foundation__lessons_3dmake_5__lessons_3dmake_5.md__pla-and-sustainability--a-realistic-picture}

PLA is often marketed as "biodegradable" and "eco-friendly." These claims deserve some scrutiny. PLA \emph{is} bio-derived --- made from fermented plant starch rather than petroleum --- which reduces its carbon footprint somewhat compared to petroleum plastics. But PLA only biodegrades under very specific industrial composting conditions (temperatures above 58°C sustained for weeks). It does not break down in a home compost bin or a landfill.

The most sustainable approach to FDM printing is not any particular material --- it\textquotesingle s designing durable parts that don\textquotesingle t need to be reprinted, using the minimum infill that meets your mechanical requirements, and avoiding test prints you don\textquotesingle t need (which is exactly why this course emphasizes verification before printing).

\begin{center}\rule{0.5\linewidth}{0.5pt}\end{center}

\subsection*{Step 4 --- Estimating Print Cost Before You Print}\label{docs__pandoc__latex__src__3dmake_foundation__lessons_3dmake_5__lessons_3dmake_5.md__step-4--estimating-print-cost-before-you-print}

\subsubsection*{Always Estimate First}\label{docs__pandoc__latex__src__3dmake_foundation__lessons_3dmake_5__lessons_3dmake_5.md__always-estimate-first}

Sending a model to print without checking the cost is like placing a restaurant order without looking at the price. For short prints of small parts, the filament cost is trivial (often under \$0.50). But for large, high-infill parts, it can be several dollars --- and more importantly, it represents several hours of printer time that you can\textquotesingle t get back if the part turns out to need revision.

The estimate formula from Lesson 1 applies here:

\begin{lstlisting}[style=Alabaster]
mass (grams)  = (volume_mm³ ÷ 1000) × filament_density_g_per_cm³
cost ($)      = mass × (spool_price_$ ÷ spool_weight_g)

\end{lstlisting}

The preflight scripts below automate this calculation and require you to confirm before proceeding.

\subsubsection*{Pre-Flight Cost Check Scripts}\label{docs__pandoc__latex__src__3dmake_foundation__lessons_3dmake_5__lessons_3dmake_5.md__pre-flight-cost-check-scripts}

\textbf{Linux / macOS (bash):}

\begin{lstlisting}[style=Alabaster, language=bash]
#!/bin/bash
set -e

3dm build
INFO=$(3dm info)
echo "$INFO"

# Extract volume from 3dm info output
VOLUME=$(echo "$INFO" | grep -oP '(?<=Volume:\s)\d+\.?\d*')
if [ -z "$VOLUME" ]; then
  echo "ERROR: Could not parse volume from 3dm info output"
  exit 1
fi

DENSITY=1.24    # g/cm³ — PLA density (change for other materials)
COST_PER_G=0.02 # $/g  — adjust for your spool price

MASS=$(echo "$VOLUME $DENSITY" | awk '{printf "%.1f", ($1/1000)*$2}')
COST=$(echo "$MASS $COST_PER_G" | awk '{printf "%.2f", $1*$2}')

echo ""
echo "=== PRE-FLIGHT SUMMARY ==="
echo "Volume:   ${VOLUME} mm³"
echo "Mass:     ${MASS} g  (PLA @ ${DENSITY} g/cm³)"
echo "Filament: \$${COST}  (@ \$${COST_PER_G}/g)"
echo ""
echo "Proceed? [y/N]"
read CONFIRM
if [ "$CONFIRM" != "y" ]; then
  echo "Cancelled."
  exit 0
fi
echo "Approved. Send to slicer."

\end{lstlisting}

\textbf{Windows (PowerShell):}

\begin{lstlisting}[style=Alabaster, language=powershell]
3dm build
$info = 3dm info
Write-Host $info

$volumeMatch = $info | Select-String -Pattern "Volume:\s+([\d.]+)"
if (-not $volumeMatch) { Write-Error "Could not parse volume"; exit 1 }
$volume = [double]$volumeMatch.Matches.Groups[1].Value

$density  = 1.24    # PLA
$costPerG = 0.02
$mass     = ($volume / 1000) * $density
$cost     = $mass * $costPerG

Write-Host ""
Write-Host "=== PRE-FLIGHT SUMMARY ==="
Write-Host ("Volume:   {0:F0} mm3" -f $volume)
Write-Host ("Mass:     {0:F1} g  (PLA)" -f $mass)
Write-Host ("`$Filament: {0:F2}" -f $cost)
Write-Host ""
$confirm = Read-Host "Proceed? [y/N]"
if ($confirm -ne "y") { Write-Host "Cancelled."; exit 0 }
Write-Host "Approved. Send to slicer."

\end{lstlisting}

\begin{center}\rule{0.5\linewidth}{0.5pt}\end{center}

\subsection*{Step 5 --- The Printer Startup Checklist}\label{docs__pandoc__latex__src__3dmake_foundation__lessons_3dmake_5__lessons_3dmake_5.md__step-5--the-printer-startup-checklist}

\subsubsection*{Before Every Print Session}\label{docs__pandoc__latex__src__3dmake_foundation__lessons_3dmake_5__lessons_3dmake_5.md__before-every-print-session}

Use this checklist before starting any print. It takes about 5 minutes and prevents the most common causes of print failure.

\begin{lstlisting}[style=Alabaster]
PRINTER STARTUP CHECKLIST
□ Bed level verified (run auto-level, or manually check at all four corners)
□ Build plate clean (wipe with IPA — remove finger oils and dust)
□ Nozzle cleared (cold pull if returning from more than a day idle)
□ Filament loaded and extruding cleanly (run 50mm manual extrude purge)
□ First layer height set correctly (0.2mm gap typical for 0.4mm nozzle)
□ Print fan operational (watch it spin up at the start)
□ Workspace clear of flammable materials
□ You will remain present for the first 15 minutes of the print
□ Power switch is accessible and you know where it is

\end{lstlisting}

\subsubsection*{Why Each Item Matters}\label{docs__pandoc__latex__src__3dmake_foundation__lessons_3dmake_5__lessons_3dmake_5.md__why-each-item-matters}

\textbf{Bed leveling} is the single most critical setting for successful prints. The first layer needs to be squished gently against the bed --- close enough to bond, but not so close that filament can\textquotesingle t flow. A bed that\textquotesingle s 0.1 mm too high or too low can cause every print to fail. Most modern printers have automatic bed leveling (ABL), but you should still understand how it works and check that it\textquotesingle s running correctly.

\textbf{Clean build plate} prevents adhesion failures. Skin oils from touching the plate are invisible but dramatically reduce how well plastic sticks. An isopropyl alcohol (IPA) wipe before every print takes 30 seconds and prevents a common failure mode.

\textbf{Cold pull (nozzle clearing)} removes degraded filament from inside the nozzle. Heat the nozzle to print temperature, push some filament through, then let it cool to about 90°C and pull the filament out firmly. The end of the pulled filament should bring out any debris or degraded material from inside the hot end.

\textbf{Filament loaded correctly} means the extruder is feeding without clicking or skipping, the filament is straight in the Bowden tube (if present), and the purge line at the start of the print flows smoothly with no bubbles or gaps.

\begin{center}\rule{0.5\linewidth}{0.5pt}\end{center}

\subsection*{Step 6 --- Monitoring During Prints}\label{docs__pandoc__latex__src__3dmake_foundation__lessons_3dmake_5__lessons_3dmake_5.md__step-6--monitoring-during-prints}

\subsubsection*{What to Watch For}\label{docs__pandoc__latex__src__3dmake_foundation__lessons_3dmake_5__lessons_3dmake_5.md__what-to-watch-for}

Checking in periodically during a print takes only a few seconds each time and lets you catch problems early. A problem caught at layer 10 is a minor inconvenience. The same problem discovered at layer 100 means you\textquotesingle ve wasted 90 layers of material and print time.

\textbf{Spaghetti} --- when the part lifts off the bed and the nozzle starts dragging filament through the air instead of depositing it in layers. The result looks like a pile of plastic spaghetti. This is the most common failure mode and the most important to catch early. If you see it, pause or cancel the print immediately.

\textbf{Layer shifts} --- a horizontal displacement somewhere in the print, where part of the model is offset from the rest. This usually indicates a mechanical problem (belt tension, loose pulley, or a motor missing steps). Cancel the print and diagnose the cause before trying again.

\textbf{Nozzle clog} --- the extruder starts making a clicking or ticking sound as the gear slips because material isn\textquotesingle t flowing. Filament may start grinding. Pause the print and clear the clog before restarting.

\textbf{Stringing} --- thin wisps of filament connecting different parts of the print. Usually a retraction or temperature setting issue rather than an emergency, but it may require post-processing cleanup.

\textbf{Smoke --- immediate emergency.} Cut power. Do not attempt to continue or diagnose while the printer is running. Notify your instructor. Do not restart the printer until the cause has been fully identified and corrected.

\begin{center}\rule{0.5\linewidth}{0.5pt}\end{center}

\subsection*{Step 7 --- Safe Shutdown and Part Removal}\label{docs__pandoc__latex__src__3dmake_foundation__lessons_3dmake_5__lessons_3dmake_5.md__step-7--safe-shutdown-and-part-removal}

\subsubsection*{The Right Way to End a Print}\label{docs__pandoc__latex__src__3dmake_foundation__lessons_3dmake_5__lessons_3dmake_5.md__the-right-way-to-end-a-print}

A print ending doesn\textquotesingle t mean the work is done. Follow these steps after every print:

\begin{enumerate}
\item
  \textbf{Wait for the nozzle to cool below 50°C} before reaching anywhere near the hot end. The nozzle can reach 200°C+ during printing and retains heat for several minutes after shutdown.
\item
  \textbf{Wait for the bed to cool below 30°C} before removing the part. Parts printed on a hot bed can warp slightly if removed while the bed is still warm. Many PEI-coated beds release parts naturally as they cool --- you may not even need tools.
\item
  \textbf{Use a spatula or palette knife} to gently separate the part from the build plate. Never force parts off with your fingers --- the edge of a part against a spring steel bed can cut skin. Apply the spatula at a low angle and use a gentle rocking motion.
\item
  \textbf{Store filament in a sealed bag with desiccant.} Filament absorbs moisture from the air (this is called hygroscopy), and moisture in filament causes bubbling, popping, and reduced mechanical properties during extrusion. This is especially important for Nylon and PETG. A small silica gel desiccant packet inside a ziplock bag is sufficient for short-term storage.
\end{enumerate}

\begin{center}\rule{0.5\linewidth}{0.5pt}\end{center}

\subsection*{Concept 2: Design for Printability and Safety}\label{docs__pandoc__latex__src__3dmake_foundation__lessons_3dmake_5__lessons_3dmake_5.md__concept-2-design-for-printability-and-safety}

\subsubsection*{Your Design Choices Affect Safety}\label{docs__pandoc__latex__src__3dmake_foundation__lessons_3dmake_5__lessons_3dmake_5.md__your-design-choices-affect-safety}

The decisions you make in OpenSCAD directly affect how safe and reliable the printing process will be. Good design doesn\textquotesingle t just mean a part that looks right --- it means a part that prints reliably without requiring dangerous materials, excessive supports, or heroic printer tuning.

\textbf{Minimize overhangs.} Any geometry that extends more than about 45° from vertical needs supports --- temporary scaffolding that the printer adds automatically. Supports increase print time, use more material, often leave rough marks where they\textquotesingle re removed, and increase the risk of print failure due to tangled support material. Good design reorients parts or modifies geometry to eliminate or reduce overhangs.

\textbf{Maintain minimum wall thickness.} For a 0.4 mm nozzle, walls below 0.8 mm (two nozzle widths) may not print reliably. Walls below 1.2 mm (three nozzle widths) are marginal for any kind of structural use. Walls below about 0.8 mm may not appear in the sliced output at all --- the slicer simply omits them as too small.

\textbf{Use the safest material that meets requirements.} If PLA is strong enough, use PLA. Don\textquotesingle t introduce PETG, ABS, or specialty materials unless you have a specific reason and the appropriate safety infrastructure.

\textbf{Encode printability constraints in your code:}

\begin{lstlisting}[style=Alabaster, language=openscad]
// Printability constants — these drive safe design decisions
wall_min     = 1.2;    // mm — minimum wall for 0.4mm nozzle (3× nozzle width)
hole_min_r   = 1.5;    // mm — minimum hole radius for reliable circular holes
clearance    = 0.2;    // mm — fitting clearance between mating faces
overhang_max = 45;     // degrees — maximum overhang without supports

module printable_box(w, d, h, wall=2) {
  // Validate against printability limits
  assert(wall >= wall_min, str("wall (", wall, "mm) must be >= ", wall_min, "mm"));
  assert(w > 2 * wall, "width must be greater than two wall thicknesses");
  assert(d > 2 * wall, "depth must be greater than two wall thicknesses");

  // Open-top box: no overhangs on the inside
  difference() {
    cube([w, d, h]);
    translate([wall, wall, wall])
      cube([w - 2*wall, d - 2*wall, h]);  // open at top — no overhang required
  }
}

printable_box(50, 40, 30);

\end{lstlisting}

\begin{center}\rule{0.5\linewidth}{0.5pt}\end{center}

\subsection*{Exercises}\label{docs__pandoc__latex__src__3dmake_foundation__lessons_3dmake_5__lessons_3dmake_5.md__exercises}

\textbf{Exercise 5.1:} Run \texttt{3dm\ info} on your model from Lesson 3. Calculate the material cost in PLA and in PETG (density 1.27 g/cm³). If PLA spools cost \$20/kg and PETG spools cost \$25/kg, what is the cost difference for your part?

\textbf{Exercise 5.2:} Write a pre-flight script (bash or PowerShell) that builds your model, reports the estimated cost, and requires a \texttt{y} confirmation before proceeding.

\textbf{Exercise 5.3 (Advanced):} Modify the pre-flight script to accept a \texttt{-\/-material} argument (options: \texttt{pla}, \texttt{petg}, \texttt{abs}, \texttt{tpu}) and calculate cost using the correct density for the selected material.

\begin{center}\rule{0.5\linewidth}{0.5pt}\end{center}

\subsection*{Quiz --- Lesson 5 (15 questions)}\label{docs__pandoc__latex__src__3dmake_foundation__lessons_3dmake_5__lessons_3dmake_5.md__quiz--lesson-5-15-questions}

\begin{enumerate}
\tightlist
\item
  What are the five levels of the Hierarchy of Controls, from most to least effective?
\item
  Why is Elimination considered the most effective safety control?
\item
  Which filament requires a dedicated fume enclosure with HEPA + activated carbon filtration and is NOT recommended for open classrooms?
\item
  At what temperature should you wait before removing a print from the bed?
\item
  What should you do immediately if you see smoke coming from your 3D printer?
\item
  Give three reasons why PLA is considered the safest classroom filament.
\item
  What is the minimum reliable wall thickness for a 0.4 mm nozzle, and why does this minimum exist?
\item
  What is the purpose of storing filament with desiccant?
\item
  True or False: It is safe to leave a 3D printer completely unattended for hours once the first layer has printed successfully.
\item
  What is "spaghetti" in the context of a 3D printing failure, and what causes it?
\item
  What is thermal runaway protection, and why is it critical that it remain enabled?
\item
  What are the two types of airborne emissions from FDM printing, and which filaments produce the most hazardous levels of each?
\item
  Describe one engineering control and one administrative control you could implement in a classroom 3D printing lab.
\item
  Explain why minimizing overhangs in your design is both a safety consideration and a quality consideration.
\item
  A student wants to print a car dashboard mount in ABS because it will be in a hot car. The classroom has no fume enclosure. Using the Hierarchy of Controls, what recommendations would you make?
\end{enumerate}

\begin{center}\rule{0.5\linewidth}{0.5pt}\end{center}

\subsection*{Extension Problems (15)}\label{docs__pandoc__latex__src__3dmake_foundation__lessons_3dmake_5__lessons_3dmake_5.md__extension-problems-15}

\begin{enumerate}
\tightlist
\item
  Create a material selection flowchart that guides a user from mechanical requirements to the safest filament that meets them.
\item
  Conduct a ventilation assessment of your classroom or makerspace. Document air exchange methods, window locations, and any existing filtration.
\item
  Write a Standard Operating Procedure (SOP) for ABS printing that includes all required safety equipment and ventilation infrastructure.
\item
  Design and describe a "first layer test tile" --- a 100 × 100 × 0.4 mm flat tile --- and explain how you\textquotesingle d use it to verify bed leveling.
\item
  Build a parametric filament storage clip in OpenSCAD: a clip that holds the free end of a filament spool. Parameters: filament diameter, clip width.
\item
  Research OSHA\textquotesingle s published guidance on VOC exposure in 3D printing environments. Summarize the key recommendations.
\item
  Compare the Safety Data Sheets (SDS) for PLA and ABS filament from two manufacturers. Document differences in recommended exposure limits and PPE.
\item
  Design a simple printer enclosure in OpenSCAD: four walls, a front door opening, and a top panel with a circular vent hole. Document the parametric variables.
\item
  Create a safety poster covering the five most important 3D printing safety rules, organized using the Hierarchy of Controls framework.
\item
  Write a one-page risk assessment for a new FDM printer being added to a classroom.
\item
  Design and describe a parametric filament moisture indicator holder: a small box that holds a humidity indicator card inside a resealable filament storage bag.
\item
  Build a "print monitoring log" template: a paper form with columns for time, nozzle temperature, bed temperature, layer number, visual observations, and action taken.
\item
  Research the difference between particle emissions and VOC emissions from FDM printing. Which is considered more hazardous at typical classroom distances from the printer?
\item
  Design a parametric build plate corner protector: a small clip that attaches to the edge of a glass or spring steel build plate to protect it from spatula damage.
\item
  Write a hypothetical "near-miss incident report" for a fictional 3D printing incident using a standard workplace incident report format (date, location, description, contributing factors, corrective actions).
\end{enumerate}

\begin{center}\rule{0.5\linewidth}{0.5pt}\end{center}

\subsection*{References and Helpful Resources}\label{docs__pandoc__latex__src__3dmake_foundation__lessons_3dmake_5__lessons_3dmake_5.md__references-and-helpful-resources}

\begin{itemize}
\tightlist
\item
  OSHA Hierarchy of Controls --- https://www.osha.gov/hierarchy-of-controls / NIOSH --- https://www.cdc.gov/niosh/topics/hierarchy/default.html
\item
  UL Research Institutes --- 3D Printing Emissions Study --- https://www.ul.com/news/ul-research-institutes-releases-3d-printing-emissions-study
\item
  NIOSH Science Blog --- Health and Safety Considerations for 3D Printing --- https://blogs.cdc.gov/niosh-science-blog/2020/05/14/3d-printing/
\item
  PrusaSlicer Documentation --- https://docs.prusa3d.com/en/
\item
  All3DP Filament Types --- https://all3dp.com/1/3d-printer-filament-types-3d-printing-3d-filament/
\item
  MatterHackers Filament Compare --- https://www.matterhackers.com/3d-printer-filament-compare
\item
  Prusa Research Materials Guide --- https://help.prusa3d.com/materials
\end{itemize}

\subsubsection*{Supplemental Resources}\label{docs__pandoc__latex__src__3dmake_foundation__lessons_3dmake_5__lessons_3dmake_5.md__supplemental-resources}

\begin{itemize}
\tightlist
\item
  3DMake GitHub Repository --- https://github.com/tdeck/3dmake
\item
  OpenSCAD User Manual --- https://en.wikibooks.org/wiki/OpenSCAD\_User\_Manual
\item
  All other reference documents are embedded as appendices below --- no separate files needed.
\end{itemize}

\begin{center}\rule{0.5\linewidth}{0.5pt}\end{center}

\section*{Part 2: Reference Appendices}\label{docs__pandoc__latex__src__3dmake_foundation__lessons_3dmake_5__lessons_3dmake_5.md__part-2-reference-appendices}

\emph{All supporting documents for Lesson 5 are included below. You do not need to open any additional files.}

\begin{center}\rule{0.5\linewidth}{0.5pt}\end{center}

\subsection*{Appendix A: Safety Checklist for 3D Printing}\label{docs__pandoc__latex__src__3dmake_foundation__lessons_3dmake_5__lessons_3dmake_5.md__appendix-a-safety-checklist-for-3d-printing}

\emph{Complete this checklist before and after each printing session. Print it out and keep a copy at your printer station.}

\begin{center}\rule{0.5\linewidth}{0.5pt}\end{center}

\subsubsection*{Pre-Print Setup}\label{docs__pandoc__latex__src__3dmake_foundation__lessons_3dmake_5__lessons_3dmake_5.md__pre-print-setup}

\textbf{Work Area}

\begin{itemize}
\tightlist
\item[$\square$]
  Desk or table cleared of unnecessary items
\item[$\square$]
  Tripping hazards removed from around the printer
\item[$\square$]
  Adequate ventilation confirmed (window open, fan running, or fume enclosure active)
\item[$\square$]
  Class ABC fire extinguisher within arm\textquotesingle s reach
\item[$\square$]
  Smoke detector mounted above printer is functional
\end{itemize}

\textbf{Printer Inspection}

\begin{itemize}
\tightlist
\item[$\square$]
  Nozzle is clean --- no old filament residue baked on
\item[$\square$]
  Build plate is level (re-level if the printer has been moved or bumped)
\item[$\square$]
  Heated bed temperature sensor is attached and secure
\item[$\square$]
  All cables are secure and none are frayed or cracked --- especially the flexible bed cable bundle
\item[$\square$]
  Thermal runaway protection is enabled (check firmware settings if unsure)
\end{itemize}

\textbf{Filament Preparation}

\begin{itemize}
\tightlist
\item[$\square$]
  Filament spool rotates freely on the holder --- no binding or tangling
\item[$\square$]
  Filament path is clear of obstructions from spool to extruder
\item[$\square$]
  Extruder drive gear is clean --- not clogged with ground plastic
\item[$\square$]
  Filament is correctly loaded and primed (a few mm of filament extruded before starting)
\end{itemize}

\textbf{Environmental Conditions}

\begin{itemize}
\tightlist
\item[$\square$]
  Room temperature is adequate (18--25°C ideal --- extreme cold slows bed heating; extreme heat can soften PLA)
\item[$\square$]
  No direct drafts from windows or air conditioning blowing on the printer (drafts cause warping and layer splitting)
\item[$\square$]
  Adequate lighting to clearly see the first layer as it prints
\end{itemize}

\begin{center}\rule{0.5\linewidth}{0.5pt}\end{center}

\subsubsection*{During Print}\label{docs__pandoc__latex__src__3dmake_foundation__lessons_3dmake_5__lessons_3dmake_5.md__during-print}

\textbf{First Layer Monitoring (Do Not Leave)}

\begin{itemize}
\tightlist
\item[$\square$]
  Watch the first 2--3 layers without interruption
\item[$\square$]
  Bed adhesion is correct --- filament sticks flat, not lifting at corners or curling
\item[$\square$]
  Nozzle temperature is stable at target value
\item[$\square$]
  No unusual sounds (grinding, clicking, popping)
\end{itemize}

\textbf{Regular Checks (Every 15--30 Minutes)}

\begin{itemize}
\tightlist
\item[$\square$]
  Print is building correctly --- no layer shifting, no spaghetti
\item[$\square$]
  Filament is feeding smoothly from the spool
\item[$\square$]
  No grinding or skipping sounds from the extruder
\item[$\square$]
  No warping visible at part edges
\end{itemize}

\textbf{Temperature Stability}

\begin{itemize}
\tightlist
\item[$\square$]
  Heated bed maintains consistent temperature throughout
\item[$\square$]
  Nozzle temperature does not fluctuate more than ±5°C
\item[$\square$]
  No thermal runaway warnings on the printer display
\end{itemize}

\begin{center}\rule{0.5\linewidth}{0.5pt}\end{center}

\subsubsection*{Stop the Print Immediately If}\label{docs__pandoc__latex__src__3dmake_foundation__lessons_3dmake_5__lessons_3dmake_5.md__stop-the-print-immediately-if}

\begin{itemize}
\tightlist
\item[$\square$]
  The nozzle jams or extrudes unevenly
\item[$\square$]
  Filament stops feeding entirely
\item[$\square$]
  Any burning smell is detected
\item[$\square$]
  Visible layer shifting occurs
\item[$\square$]
  Grinding or skipping sounds from the extruder
\item[$\square$]
  Any unusual chemical or burning odor
\item[$\square$]
  Smoke appears --- \textbf{cut power immediately}
\end{itemize}

\begin{center}\rule{0.5\linewidth}{0.5pt}\end{center}

\subsubsection*{Post-Print}\label{docs__pandoc__latex__src__3dmake_foundation__lessons_3dmake_5__lessons_3dmake_5.md__post-print}

\textbf{Cool Down}

\begin{itemize}
\tightlist
\item[$\square$]
  Allow the heated bed to cool naturally for 10--15 minutes before attempting print removal
\item[$\square$]
  Keep hands clear of the nozzle until it reads below 50°C
\item[$\square$]
  Confirm the printer is idle and all temperatures are dropping before touching
\end{itemize}

\textbf{Print Removal}

\begin{itemize}
\tightlist
\item[$\square$]
  Use a proper spatula or scraper --- not fingers or improvised tools
\item[$\square$]
  Remove the print only when the bed has fully cooled (PLA releases easily from PEI beds when cool)
\item[$\square$]
  Inspect the print for sharp edges, support stubs, or layer delamination
\item[$\square$]
  Sand or file sharp edges if the part will be handled
\end{itemize}

\textbf{Equipment Cleaning}

\begin{itemize}
\tightlist
\item[$\square$]
  Wipe the nozzle with a brass wire brush when cooled to 50°C (never cold --- residue must be slightly soft)
\item[$\square$]
  Remove any plastic debris and failed print material from the build plate
\item[$\square$]
  Check the extruder for embedded filament residue
\item[$\square$]
  Clear any visible dust from the electronics area
\end{itemize}

\textbf{Workspace Cleanup}

\begin{itemize}
\tightlist
\item[$\square$]
  Return all tools to their proper storage location
\item[$\square$]
  Dispose of support material and failed prints in a bin (not on the floor)
\item[$\square$]
  Store filament in a sealed bag with desiccant
\item[$\square$]
  Confirm the workspace is clear and safe for the next user
\end{itemize}

\begin{center}\rule{0.5\linewidth}{0.5pt}\end{center}

\subsubsection*{Hazard Reference Summary}\label{docs__pandoc__latex__src__3dmake_foundation__lessons_3dmake_5__lessons_3dmake_5.md__hazard-reference-summary}

{\def\LTcaptype{none} % do not increment counter
\begin{longtable}[]{@{}
  >{\raggedright\arraybackslash}p{(\linewidth - 4\tabcolsep) * \real{0.2841}}
  >{\raggedright\arraybackslash}p{(\linewidth - 4\tabcolsep) * \real{0.2614}}
  >{\raggedright\arraybackslash}p{(\linewidth - 4\tabcolsep) * \real{0.4545}}@{}}
\toprule\noalign{}
\begin{minipage}[b]{\linewidth}\raggedright
Hazard
\end{minipage} & \begin{minipage}[b]{\linewidth}\raggedright
Temperature / Condition
\end{minipage} & \begin{minipage}[b]{\linewidth}\raggedright
Response
\end{minipage} \\
\midrule\noalign{}
\endhead
\bottomrule\noalign{}
\endlastfoot
Hot nozzle & 190--260°C & Do not touch; wait for \textless{} 50°C \\
Hot bed & 50--110°C & Do not touch; allow full cool-down \\
Fumes (ABS/PETG) & At operating temp &
Ventilate; use enclosure with filtration \\
Electrical (power supply) & 120V / 240V mains &
Never open while plugged in \\
Tangled filament & Any & Clear before printing; never force-pull \\
\end{longtable}
}

\subsubsection*{Emergency Actions}\label{docs__pandoc__latex__src__3dmake_foundation__lessons_3dmake_5__lessons_3dmake_5.md__emergency-actions}

\begin{itemize}
\tightlist
\item
  \textbf{Fire:} Unplug immediately. Use a dry chemical (Class ABC) extinguisher --- NOT water on electrical fires. Evacuate and call emergency services.
\item
  \textbf{Burns:} Rinse with cool (not cold) running water for 15+ minutes. Seek medical attention for any burn larger than a coin.
\item
  \textbf{Fume inhalation:} Move to fresh air immediately. Seek medical attention if symptoms (headache, dizziness, nausea) persist.
\end{itemize}

\emph{Last Reviewed:} \_\_\_\_\_\_\_\_\_\_ \emph{Reviewed By:} \_\_\_\_\_\_\_\_\_\_

\begin{center}\rule{0.5\linewidth}{0.5pt}\end{center}

\subsection*{Appendix B: Filament Comparison Table}\label{docs__pandoc__latex__src__3dmake_foundation__lessons_3dmake_5__lessons_3dmake_5.md__appendix-b-filament-comparison-table}

\emph{Quick reference for choosing the right filament. All temperature values are typical ranges --- always check the label on your specific spool. The spool always wins over this table.}

\begin{center}\rule{0.5\linewidth}{0.5pt}\end{center}

\subsubsection*{Side-by-Side Comparison}\label{docs__pandoc__latex__src__3dmake_foundation__lessons_3dmake_5__lessons_3dmake_5.md__side-by-side-comparison}

{\def\LTcaptype{none} % do not increment counter
\begin{longtable}[]{@{}
  >{\raggedright\arraybackslash}p{(\linewidth - 8\tabcolsep) * \real{0.2830}}
  >{\raggedright\arraybackslash}p{(\linewidth - 8\tabcolsep) * \real{0.1226}}
  >{\raggedright\arraybackslash}p{(\linewidth - 8\tabcolsep) * \real{0.1981}}
  >{\raggedright\arraybackslash}p{(\linewidth - 8\tabcolsep) * \real{0.1604}}
  >{\raggedright\arraybackslash}p{(\linewidth - 8\tabcolsep) * \real{0.2358}}@{}}
\toprule\noalign{}
\begin{minipage}[b]{\linewidth}\raggedright
Property
\end{minipage} & \begin{minipage}[b]{\linewidth}\raggedright
PLA
\end{minipage} & \begin{minipage}[b]{\linewidth}\raggedright
PETG
\end{minipage} & \begin{minipage}[b]{\linewidth}\raggedright
TPU (Flex)
\end{minipage} & \begin{minipage}[b]{\linewidth}\raggedright
ABS
\end{minipage} \\
\midrule\noalign{}
\endhead
\bottomrule\noalign{}
\endlastfoot
\textbf{Difficulty} & Easiest & Moderate & Hard & Very Hard \\
\textbf{Nozzle temp} & 190--220°C & 230--250°C & 220--240°C &
230--250°C \\
\textbf{Bed temp} & 50--60°C & 70--85°C & 30--60°C & 90--110°C \\
\textbf{Enclosure needed?} & No & No & No & Yes \\
\textbf{Rigidity} & High & Medium & None --- flexible & High \\
\textbf{Impact resistance} & Low (brittle) & Medium & Very high &
Medium \\
\textbf{Heat resistance} & Low (\textasciitilde{}60°C) &
Medium (\textasciitilde{}80°C) & Medium &
High (\textasciitilde{}100°C) \\
\textbf{Flexibility} & None & Slight & Rubber-like & None \\
\textbf{Moisture sensitivity} & Low & Medium & High & Medium \\
\textbf{Fume concern} & Lowest & Low & Moderate & High \\
\textbf{Bed adhesion} & Easy & Use glue stick on PEI & Easy &
Requires enclosure + glue \\
\textbf{Print speed} & Normal & Normal & Slow (20--30 mm/s) & Normal \\
\textbf{Stringing tendency} & Low & Medium--High & High & Low \\
\textbf{Approximate cost} & \$15--25/kg & \$15--30/kg & \$20--40/kg &
\$15--25/kg \\
\textbf{Recommended for beginners?} & ✅ Yes & After PLA &
After experience & ❌ No \\
\end{longtable}
}

\begin{center}\rule{0.5\linewidth}{0.5pt}\end{center}

\subsubsection*{When to Use Each Material}\label{docs__pandoc__latex__src__3dmake_foundation__lessons_3dmake_5__lessons_3dmake_5.md__when-to-use-each-material}

\textbf{PLA --- Use for:}

\begin{itemize}
\tightlist
\item
  Prototypes and test prints
\item
  Classroom projects
\item
  Decorative objects and tactile models
\item
  Anything that won\textquotesingle t be exposed to heat, moisture, or heavy mechanical stress
\end{itemize}

\textbf{PLA --- Avoid for:}

\begin{itemize}
\tightlist
\item
  Objects left in a hot car or direct sun (softens above \textasciitilde{}60°C)
\item
  Parts that need to flex or bend without cracking
\item
  High-impact applications
\end{itemize}

\textbf{PETG --- Use for:}

\begin{itemize}
\tightlist
\item
  Functional parts that need to be tougher than PLA
\item
  Parts exposed to mild heat or moisture
\item
  Mechanical components: brackets, clips, mounts
\item
  Food-contact applications (check your specific brand for food-safe certification)
\end{itemize}

\textbf{PETG --- Avoid for:}

\begin{itemize}
\tightlist
\item
  Very fine surface detail (it strings more than PLA)
\item
  Projects where you need the easiest possible first-time print
\end{itemize}

\textbf{TPU / Flexible --- Use for:}

\begin{itemize}
\tightlist
\item
  Wearable objects (wristbands, phone cases, watch straps)
\item
  Bumpers and shock absorbers
\item
  Grips and handles
\item
  Objects that must deform under pressure and return to shape
\end{itemize}

\textbf{TPU --- Avoid for:}

\begin{itemize}
\tightlist
\item
  Fine surface detail
\item
  Printers with Bowden extruders (the flexible filament buckles in the tube --- direct drive only)
\item
  Your first few prints while still learning
\end{itemize}

\textbf{ABS --- Use for:}

\begin{itemize}
\tightlist
\item
  High-heat environments (car interiors, near motors)
\item
  Parts requiring post-processing (ABS sands, bonds with acetone vapor, and paints easily)
\item
  Professional or industrial contexts with proper dedicated ventilation
\end{itemize}

\textbf{ABS --- Avoid for:}

\begin{itemize}
\tightlist
\item
  Any classroom without a dedicated sealed fume enclosure and active filtration
\item
  Beginners
\item
  Any print where warping at the edges would be a problem and you have no enclosure
\end{itemize}

\begin{center}\rule{0.5\linewidth}{0.5pt}\end{center}

\subsubsection*{Slicer Profile Quick Reference}\label{docs__pandoc__latex__src__3dmake_foundation__lessons_3dmake_5__lessons_3dmake_5.md__slicer-profile-quick-reference}

{\def\LTcaptype{none} % do not increment counter
\begin{longtable}[]{@{}
  >{\raggedright\arraybackslash}p{(\linewidth - 4\tabcolsep) * \real{0.1026}}
  >{\raggedright\arraybackslash}p{(\linewidth - 4\tabcolsep) * \real{0.2436}}
  >{\raggedright\arraybackslash}p{(\linewidth - 4\tabcolsep) * \real{0.6538}}@{}}
\toprule\noalign{}
\begin{minipage}[b]{\linewidth}\raggedright
Filament
\end{minipage} & \begin{minipage}[b]{\linewidth}\raggedright
PrusaSlicer Profile
\end{minipage} & \begin{minipage}[b]{\linewidth}\raggedright
Key Changes from PLA Default
\end{minipage} \\
\midrule\noalign{}
\endhead
\bottomrule\noalign{}
\endlastfoot
PLA & Generic PLA & Baseline --- no changes needed \\
PETG & Generic PETG &
Higher temps; slower cooling; glue stick on PEI bed \\
TPU & Generic Flex & Reduce speed to 20--30 mm/s; reduce retraction \\
ABS & Generic ABS & Higher temps; disable cooling fan; use enclosure \\
\end{longtable}
}

\begin{center}\rule{0.5\linewidth}{0.5pt}\end{center}

\subsubsection*{Filament Storage}\label{docs__pandoc__latex__src__3dmake_foundation__lessons_3dmake_5__lessons_3dmake_5.md__filament-storage}

All filament absorbs moisture from the air over time. Wet filament causes: popping or crackling sounds during printing, bubbles in extruded plastic, excessive stringing, and weak or brittle finished parts.

{\def\LTcaptype{none} % do not increment counter
\begin{longtable}[]{@{}
  >{\raggedright\arraybackslash}p{(\linewidth - 4\tabcolsep) * \real{0.0777}}
  >{\raggedright\arraybackslash}p{(\linewidth - 4\tabcolsep) * \real{0.1942}}
  >{\raggedright\arraybackslash}p{(\linewidth - 4\tabcolsep) * \real{0.7282}}@{}}
\toprule\noalign{}
\begin{minipage}[b]{\linewidth}\raggedright
Filament
\end{minipage} & \begin{minipage}[b]{\linewidth}\raggedright
Moisture Sensitivity
\end{minipage} & \begin{minipage}[b]{\linewidth}\raggedright
Storage Recommendation
\end{minipage} \\
\midrule\noalign{}
\endhead
\bottomrule\noalign{}
\endlastfoot
PLA & Low & Sealed bag with desiccant when not in use \\
PETG & Medium &
Sealed bag with desiccant; dry before use if stored open for weeks \\
TPU & High &
Always store sealed; dry in oven at 65°C for 4--6 hours if moisture-absorbed \\
ABS & Medium & Sealed bag with desiccant \\
\end{longtable}
}

\emph{Sources: All3DP Filament Types Guide; MatterHackers Filament Compare; Prusa Research Materials Guide}

\begin{center}\rule{0.5\linewidth}{0.5pt}\end{center}

\subsection*{Appendix C: Material Properties Quick Reference}\label{docs__pandoc__latex__src__3dmake_foundation__lessons_3dmake_5__lessons_3dmake_5.md__appendix-c-material-properties-quick-reference}

\emph{Essential material data for FDM/FFF printing. Use alongside the Filament Comparison Table for material selection decisions.}

\begin{center}\rule{0.5\linewidth}{0.5pt}\end{center}

\subsubsection*{Properties at a Glance}\label{docs__pandoc__latex__src__3dmake_foundation__lessons_3dmake_5__lessons_3dmake_5.md__properties-at-a-glance}

{\def\LTcaptype{none} % do not increment counter
\begin{longtable}[]{@{}
  >{\raggedright\arraybackslash}p{(\linewidth - 8\tabcolsep) * \real{0.2879}}
  >{\raggedright\arraybackslash}p{(\linewidth - 8\tabcolsep) * \real{0.1515}}
  >{\raggedright\arraybackslash}p{(\linewidth - 8\tabcolsep) * \real{0.1515}}
  >{\raggedright\arraybackslash}p{(\linewidth - 8\tabcolsep) * \real{0.1970}}
  >{\raggedright\arraybackslash}p{(\linewidth - 8\tabcolsep) * \real{0.2121}}@{}}
\toprule\noalign{}
\begin{minipage}[b]{\linewidth}\raggedright
Property
\end{minipage} & \begin{minipage}[b]{\linewidth}\raggedright
PLA
\end{minipage} & \begin{minipage}[b]{\linewidth}\raggedright
PETG
\end{minipage} & \begin{minipage}[b]{\linewidth}\raggedright
ABS
\end{minipage} & \begin{minipage}[b]{\linewidth}\raggedright
TPU
\end{minipage} \\
\midrule\noalign{}
\endhead
\bottomrule\noalign{}
\endlastfoot
Nozzle temp (°C) & 200--220 & 235--250 & 240--260 & 220--240 \\
Bed temp (°C) & 20--60 & 70--90 & 100--110 & 20--60 \\
Strength & Moderate & High & Very High & Low (flexible) \\
Flexibility & None & Slight & Slight & High \\
Ease of printing & Very Easy & Easy & Hard & Moderate \\
Cost (\$/kg) & \$15--25 & \$20--30 & \$18--28 & \$30--50 \\
Print speed & 40--60 mm/s & 30--50 mm/s & 20--40 mm/s & 10--30 mm/s \\
Bed adhesion & Easy & Moderate & Requires prep & Easy \\
UV resistance & Low & Moderate & High & Low \\
Chemical resistance & No & Moderate & Yes & No \\
\end{longtable}
}

\begin{center}\rule{0.5\linewidth}{0.5pt}\end{center}

\subsubsection*{Heat Resistance After Printing}\label{docs__pandoc__latex__src__3dmake_foundation__lessons_3dmake_5__lessons_3dmake_5.md__heat-resistance-after-printing}

This is one of the most common reasons to choose a material other than PLA.

{\def\LTcaptype{none} % do not increment counter
\begin{longtable}[]{@{}
  >{\raggedright\arraybackslash}p{(\linewidth - 6\tabcolsep) * \real{0.0870}}
  >{\centering\arraybackslash}p{(\linewidth - 6\tabcolsep) * \real{0.1848}}
  >{\centering\arraybackslash}p{(\linewidth - 6\tabcolsep) * \real{0.1087}}
  >{\raggedright\arraybackslash}p{(\linewidth - 6\tabcolsep) * \real{0.6196}}@{}}
\toprule\noalign{}
\begin{minipage}[b]{\linewidth}\raggedright
Material
\end{minipage} & \begin{minipage}[b]{\linewidth}\centering
Safe to use up to
\end{minipage} & \begin{minipage}[b]{\linewidth}\centering
Softens at
\end{minipage} & \begin{minipage}[b]{\linewidth}\raggedright
Notes
\end{minipage} \\
\midrule\noalign{}
\endhead
\bottomrule\noalign{}
\endlastfoot
PLA & \textasciitilde{}50--60°C & \textasciitilde{}60--70°C &
Not suitable for car interiors or items near heat sources \\
PETG & \textasciitilde{}80--100°C & \textasciitilde{}100--110°C &
Handles warm environments well \\
ABS & \textasciitilde{}90--110°C & \textasciitilde{}110--120°C &
Good heat resistance; most demanding to print \\
TPU & \textasciitilde{}60--80°C & \textasciitilde{}80--100°C &
Limited heat resistance despite flex durability \\
\end{longtable}
}

\begin{center}\rule{0.5\linewidth}{0.5pt}\end{center}

\subsubsection*{Chemical Resistance}\label{docs__pandoc__latex__src__3dmake_foundation__lessons_3dmake_5__lessons_3dmake_5.md__chemical-resistance}

{\def\LTcaptype{none} % do not increment counter
\begin{longtable}[]{@{}lccccc@{}}
\toprule\noalign{}
Material & Water & Alcohol & Acetone & Oils & Acids \\
\midrule\noalign{}
\endhead
\bottomrule\noalign{}
\endlastfoot
PLA & Low & Low & Dissolves & Low & Low \\
PETG & High & Moderate & Moderate & High & Moderate \\
ABS & High & Moderate & Dissolves & High & Moderate \\
TPU & Moderate & Low & Low & Moderate & Low \\
\end{longtable}
}

\textbf{Acetone note:} Can be used to smooth ABS surfaces (vapor or liquid application). Also dissolves PLA --- keep acetone away from PLA prints and filament spools.

\begin{center}\rule{0.5\linewidth}{0.5pt}\end{center}

\subsubsection*{Strength After Printing --- Curing Time}\label{docs__pandoc__latex__src__3dmake_foundation__lessons_3dmake_5__lessons_3dmake_5.md__strength-after-printing--curing-time}

Freshly printed parts are not at full strength.

{\def\LTcaptype{none} % do not increment counter
\begin{longtable}[]{@{}ll@{}}
\toprule\noalign{}
Time After Printing & Approximate Strength \\
\midrule\noalign{}
\endhead
\bottomrule\noalign{}
\endlastfoot
0--24 hours & 70--80\% of final strength \\
24--48 hours & 90--95\% of final strength \\
48--72 hours & \textasciitilde{}99\% of final strength \\
7+ days & Maximum strength \\
\end{longtable}
}

Do not stress-test functional parts immediately after printing. Allow at least 24 hours before load testing.

\begin{center}\rule{0.5\linewidth}{0.5pt}\end{center}

\subsubsection*{Storage Conditions}\label{docs__pandoc__latex__src__3dmake_foundation__lessons_3dmake_5__lessons_3dmake_5.md__storage-conditions}

{\def\LTcaptype{none} % do not increment counter
\begin{longtable}[]{@{}ll@{}}
\toprule\noalign{}
Factor & Recommended Range \\
\midrule\noalign{}
\endhead
\bottomrule\noalign{}
\endlastfoot
Storage temperature & 15--25°C \\
Relative humidity & 30--50\% \\
Light exposure & Away from direct sunlight \\
Container & Sealed with desiccant packets \\
\end{longtable}
}

\begin{center}\rule{0.5\linewidth}{0.5pt}\end{center}

\subsubsection*{Quick Material Selection Guide}\label{docs__pandoc__latex__src__3dmake_foundation__lessons_3dmake_5__lessons_3dmake_5.md__quick-material-selection-guide}

{\def\LTcaptype{none} % do not increment counter
\begin{longtable}[]{@{}ll@{}}
\toprule\noalign{}
Question & Answer \\
\midrule\noalign{}
\endhead
\bottomrule\noalign{}
\endlastfoot
First time printing? & PLA \\
Need durability? & PETG or ABS \\
Need flexibility? & TPU \\
Need high strength? & ABS or PETG \\
Quick test / prototype? & PLA \\
Outdoor use? & PETG or ABS \\
Food contact? & Food-grade certified PETG (verify brand) \\
Heat resistant? & ABS \\
Easiest possible print? & PLA \\
\end{longtable}
}

\begin{center}\rule{0.5\linewidth}{0.5pt}\end{center}

\subsection*{Appendix D: Printer Maintenance Log}\label{docs__pandoc__latex__src__3dmake_foundation__lessons_3dmake_5__lessons_3dmake_5.md__appendix-d-printer-maintenance-log}

\emph{Copy this log and keep a printed or digital copy with your printer. Fill in the Printer Information section first, then log every maintenance event and issue as it occurs. A complete log is your best diagnostic tool when problems arise.}

\begin{center}\rule{0.5\linewidth}{0.5pt}\end{center}

\subsubsection*{Printer Information}\label{docs__pandoc__latex__src__3dmake_foundation__lessons_3dmake_5__lessons_3dmake_5.md__printer-information}

{\def\LTcaptype{none} % do not increment counter
\begin{longtable}[]{@{}ll@{}}
\toprule\noalign{}
Field & Value \\
\midrule\noalign{}
\endhead
\bottomrule\noalign{}
\endlastfoot
Printer model & \\
Serial number & \\
Purchase date & \\
Firmware version & \\
Last service date & \\
Warranty expiry & \\
Manufacturer support contact & \\
\end{longtable}
}

\begin{center}\rule{0.5\linewidth}{0.5pt}\end{center}

\subsubsection*{Maintenance Schedule}\label{docs__pandoc__latex__src__3dmake_foundation__lessons_3dmake_5__lessons_3dmake_5.md__maintenance-schedule}

\textbf{Before each use (Daily)}

\begin{itemize}
\tightlist
\item[$\square$]
  Visual inspection --- no visible damage, cables intact
\item[$\square$]
  Clean nozzle tip if residue is present
\item[$\square$]
  Verify build plate is level (re-level if printer was moved)
\item[$\square$]
  Confirm filament spool rotates freely with no tangles
\end{itemize}

\textbf{Weekly}

\begin{itemize}
\tightlist
\item[$\square$]
  Clean extruder drive gear (compressed air or brush out ground plastic)
\item[$\square$]
  Inspect and wipe Z-axis rails and lead screw
\item[$\square$]
  Check all cable connections at the control board
\item[$\square$]
  Test emergency stop / power cut
\end{itemize}

\textbf{Monthly}

\begin{itemize}
\tightlist
\item[$\square$]
  Full build plate leveling procedure (manual or automatic)
\item[$\square$]
  Clean interior of print chamber
\item[$\square$]
  Inspect heating elements and thermistor connections
\item[$\square$]
  Run a temperature calibration test print
\end{itemize}

\textbf{Quarterly (every 3 months)}

\begin{itemize}
\tightlist
\item[$\square$]
  Replace nozzle if inner bore shows wear or repeat clogging
\item[$\square$]
  Full mechanical inspection --- belt tension, eccentric nuts, frame screws
\item[$\square$]
  Check firmware for updates
\item[$\square$]
  Run a full calibration test print and document results
\end{itemize}

\begin{center}\rule{0.5\linewidth}{0.5pt}\end{center}

\subsubsection*{Maintenance Log}\label{docs__pandoc__latex__src__3dmake_foundation__lessons_3dmake_5__lessons_3dmake_5.md__maintenance-log}

{\def\LTcaptype{none} % do not increment counter
\begin{longtable}[]{@{}
  >{\raggedright\arraybackslash}p{(\linewidth - 12\tabcolsep) * \real{0.0741}}
  >{\raggedright\arraybackslash}p{(\linewidth - 12\tabcolsep) * \real{0.0741}}
  >{\raggedright\arraybackslash}p{(\linewidth - 12\tabcolsep) * \real{0.2037}}
  >{\raggedright\arraybackslash}p{(\linewidth - 12\tabcolsep) * \real{0.1852}}
  >{\raggedright\arraybackslash}p{(\linewidth - 12\tabcolsep) * \real{0.2222}}
  >{\raggedright\arraybackslash}p{(\linewidth - 12\tabcolsep) * \real{0.1852}}
  >{\raggedright\arraybackslash}p{(\linewidth - 12\tabcolsep) * \real{0.0556}}@{}}
\toprule\noalign{}
\begin{minipage}[b]{\linewidth}\raggedright
Date
\end{minipage} & \begin{minipage}[b]{\linewidth}\raggedright
Type
\end{minipage} & \begin{minipage}[b]{\linewidth}\raggedright
Description
\end{minipage} & \begin{minipage}[b]{\linewidth}\raggedright
Time Spent
\end{minipage} & \begin{minipage}[b]{\linewidth}\raggedright
Issues Found
\end{minipage} & \begin{minipage}[b]{\linewidth}\raggedright
Resolution
\end{minipage} & \begin{minipage}[b]{\linewidth}\raggedright
By
\end{minipage} \\
\midrule\noalign{}
\endhead
\bottomrule\noalign{}
\endlastfoot
& & & & & & \\
& & & & & & \\
& & & & & & \\
& & & & & & \\
& & & & & & \\
& & & & & & \\
& & & & & & \\
& & & & & & \\
\end{longtable}
}

\emph{Types: Routine / Repair / Calibration / Cleaning / Part replacement}

\begin{center}\rule{0.5\linewidth}{0.5pt}\end{center}

\subsubsection*{Issue Tracking Log}\label{docs__pandoc__latex__src__3dmake_foundation__lessons_3dmake_5__lessons_3dmake_5.md__issue-tracking-log}

{\def\LTcaptype{none} % do not increment counter
\begin{longtable}[]{@{}llllll@{}}
\toprule\noalign{}
Date & Symptom & Diagnosis & Action Taken & Status & Notes \\
\midrule\noalign{}
\endhead
\bottomrule\noalign{}
\endlastfoot
& & & & & \\
& & & & & \\
& & & & & \\
& & & & & \\
& & & & & \\
\end{longtable}
}

\emph{Statuses: Open / In progress / Resolved / Monitoring}

\begin{center}\rule{0.5\linewidth}{0.5pt}\end{center}

\subsubsection*{Parts Replacement Record}\label{docs__pandoc__latex__src__3dmake_foundation__lessons_3dmake_5__lessons_3dmake_5.md__parts-replacement-record}

{\def\LTcaptype{none} % do not increment counter
\begin{longtable}[]{@{}llllll@{}}
\toprule\noalign{}
Date & Part Replaced & Reason & Supplier & Cost & Notes \\
\midrule\noalign{}
\endhead
\bottomrule\noalign{}
\endlastfoot
& & & & & \\
& & & & & \\
& & & & & \\
\end{longtable}
}

\begin{center}\rule{0.5\linewidth}{0.5pt}\end{center}

\subsubsection*{Filament Compatibility Notes}\label{docs__pandoc__latex__src__3dmake_foundation__lessons_3dmake_5__lessons_3dmake_5.md__filament-compatibility-notes}

Record your tested settings for each material on this specific printer. Published temperature ranges are starting points --- your printer may run hotter or cooler than its display indicates.

{\def\LTcaptype{none} % do not increment counter
\begin{longtable}[]{@{}
  >{\raggedright\arraybackslash}p{(\linewidth - 12\tabcolsep) * \real{0.1481}}
  >{\raggedright\arraybackslash}p{(\linewidth - 12\tabcolsep) * \real{0.0926}}
  >{\centering\arraybackslash}p{(\linewidth - 12\tabcolsep) * \real{0.2037}}
  >{\centering\arraybackslash}p{(\linewidth - 12\tabcolsep) * \real{0.1481}}
  >{\centering\arraybackslash}p{(\linewidth - 12\tabcolsep) * \real{0.0926}}
  >{\centering\arraybackslash}p{(\linewidth - 12\tabcolsep) * \real{0.2222}}
  >{\raggedright\arraybackslash}p{(\linewidth - 12\tabcolsep) * \real{0.0926}}@{}}
\toprule\noalign{}
\begin{minipage}[b]{\linewidth}\raggedright
Filament
\end{minipage} & \begin{minipage}[b]{\linewidth}\raggedright
Brand
\end{minipage} & \begin{minipage}[b]{\linewidth}\centering
Nozzle Temp
\end{minipage} & \begin{minipage}[b]{\linewidth}\centering
Bed Temp
\end{minipage} & \begin{minipage}[b]{\linewidth}\centering
Speed
\end{minipage} & \begin{minipage}[b]{\linewidth}\centering
Success Rate
\end{minipage} & \begin{minipage}[b]{\linewidth}\raggedright
Notes
\end{minipage} \\
\midrule\noalign{}
\endhead
\bottomrule\noalign{}
\endlastfoot
PLA & & & & & & \\
PETG & & & & & & \\
ABS & & & & & & \\
TPU & & & & & & \\
Other & & & & & & \\
\end{longtable}
}

\begin{center}\rule{0.5\linewidth}{0.5pt}\end{center}

\subsubsection*{Troubleshooting Quick Reference}\label{docs__pandoc__latex__src__3dmake_foundation__lessons_3dmake_5__lessons_3dmake_5.md__troubleshooting-quick-reference}

\textbf{Nozzle Clog}

\begin{itemize}
\tightlist
\item
  First response: Cold pull (heat to printing temp, push filament through manually, cool to 90°C, pull sharply)
\item
  Second response: Acupuncture needle or nozzle cleaning filament
\item
  Last resort: Replace nozzle
\end{itemize}

\textbf{Bed Adhesion Problems}

\begin{itemize}
\tightlist
\item
  Check: Bed level, nozzle height from bed, bed surface condition (clean? worn?)
\item
  Fix by material: PLA → clean PEI with IPA; PETG → apply glue stick; ABS → glue stick + enclosure
\end{itemize}

\textbf{Layer Shifting}

\begin{itemize}
\tightlist
\item
  Common causes: Belt too loose, motor current too low, print speed too high, X/Y axis obstruction
\item
  Check belt tension first --- it\textquotesingle s the most common cause
\end{itemize}

\textbf{Extrusion Issues (Under-extrusion)}

\begin{itemize}
\tightlist
\item
  Common causes: Partial clog, worn drive gear, incorrect extruder steps/mm, wet filament
\item
  Run a flow rate test: extrude 100 mm and measure actual output
\end{itemize}

\begin{center}\rule{0.5\linewidth}{0.5pt}\end{center}

\subsubsection*{Monthly Performance Summary}\label{docs__pandoc__latex__src__3dmake_foundation__lessons_3dmake_5__lessons_3dmake_5.md__monthly-performance-summary}

{\def\LTcaptype{none} % do not increment counter
\begin{longtable}[]{@{}lcll@{}}
\toprule\noalign{}
Month & Print Success Rate & Common Issues & Overall Assessment \\
\midrule\noalign{}
\endhead
\bottomrule\noalign{}
\endlastfoot
& & & \\
& & & \\
& & & \\
& & & \\
\end{longtable}
}

\emph{Last log entry:} \_\_\_\_\_\_\_\_\_\_ \emph{Logged by:} \_\_\_\_\_\_\_\_\_\_

\newpage

\chapter{Command-Line Fundamentals - Choose Your Path}\label{docs__pandoc__latex__src__command_line_interface_selection__command_line_interface_selection.md__command-line-fundamentals---choose-your-path}

Welcome! Before diving into 3D design with OpenSCAD, you\textquotesingle ll master command-line fundamentals. This page will help you understand what command-line interfaces are and choose the best path for you.

\section*{What is a Command-Line Interface (CLI)?}\label{docs__pandoc__latex__src__command_line_interface_selection__command_line_interface_selection.md__what-is-a-command-line-interface-cli}

A command-line interface is a text-based way to control your computer by typing commands instead of clicking buttons. It\textquotesingle s like sending written instructions to your computer.

\subsection*{Why learn it?}\label{docs__pandoc__latex__src__command_line_interface_selection__command_line_interface_selection.md__why-learn-it}

\begin{itemize}
\tightlist
\item
  Speed: Text commands are often faster than clicking through menus
\item
  Precision: Exact control over what your computer does
\item
  Accessibility: Perfect for screen readers - text is naturally readable
\item
  Automation: Repeat tasks automatically
\item
  3D Printing: Essential for batch processing models and integrating tools
\end{itemize}

\subsection*{Real-world example}\label{docs__pandoc__latex__src__command_line_interface_selection__command_line_interface_selection.md__real-world-example}

Instead of:

\begin{enumerate}
\tightlist
\item
  Opening File Explorer (click)
\item
  Navigating folders (click, click, click)
\item
  Right-clicking a file (click)
\item
  Selecting "Copy" (click)
\item
  Navigating to destination (click, click)
\item
  Right-clicking (click)
\item
  Selecting "Paste" (click)
\end{enumerate}

You type: \texttt{cp\ myfile.txt\ backup/} and press Enter. Done.

\section*{Three Command-Line Options on Windows}\label{docs__pandoc__latex__src__command_line_interface_selection__command_line_interface_selection.md__three-command-line-options-on-windows}

Windows offers three ways to use the command line. All are accessible with screen readers. Here\textquotesingle s how they compare:

\subsection*{Option 1: Windows Command Prompt (CMD)}\label{docs__pandoc__latex__src__command_line_interface_selection__command_line_interface_selection.md__option-1-windows-command-prompt-cmd}

What it is: The original Windows command-line (1981-present)

Best for: Absolute beginners, maximum simplicity

Pros:

\begin{itemize}
\tightlist
\item
  Simple commands and syntax
\item
  Minimal learning curve
\item
  Easy to understand error messages
\item
  Great for basic file operations
\item
  Perfect entry point to command-line world
\end{itemize}

Cons:

\begin{itemize}
\tightlist
\item
  Limited advanced features
\item
  Less powerful than alternatives
\item
  No built-in piping (but available)
\item
  Smaller ecosystem
\end{itemize}

Typical command:

\begin{lstlisting}[style=Alabaster, language=cmd]
copy myfile.txt backup\

\end{lstlisting}

\subsection*{Option 2: PowerShell}\label{docs__pandoc__latex__src__command_line_interface_selection__command_line_interface_selection.md__option-2-powershell}

What it is: Microsoft\textquotesingle s modern, powerful shell (2006-present)

Best for: Intermediate users, advanced automation

Pros:

\begin{itemize}
\tightlist
\item
  Very powerful for scripting
\item
  Modern syntax and features
\item
  Excellent for 3D printing automation
\item
  Professional workflows
\item
  Large community
\end{itemize}

Cons:

\begin{itemize}
\tightlist
\item
  Steeper learning curve than CMD
\item
  More complex syntax
\item
  More "wordy" commands
\item
  Overkill for simple tasks
\end{itemize}

Typical command:

\begin{lstlisting}[style=Alabaster, language=powershell]
Copy-Item -Path myfile.txt -Destination backup/

\end{lstlisting}

\subsection*{Option 3: Git Bash}\label{docs__pandoc__latex__src__command_line_interface_selection__command_line_interface_selection.md__option-3-git-bash}

What it is: A Unix/Linux shell on Windows (runs bash inside Git for Windows)

Best for: Programmers, users familiar with Linux, advanced users

Pros:

\begin{itemize}
\tightlist
\item
  Familiar if you know Linux/Unix
\item
  Powerful piping and text processing
\item
  Consistent with other platforms (macOS, Linux)
\item
  Excellent for advanced workflows
\item
  Industry-standard for developers
\end{itemize}

Cons:

\begin{itemize}
\tightlist
\item
  Requires Git installation
\item
  Steeper learning curve
\item
  Path syntax is different from native Windows
\item
  Less integrated with Windows system tools
\item
  May be "too much" for beginners
\end{itemize}

Typical command:

\begin{lstlisting}[style=Alabaster, language=bash]
cp myfile.txt backup/

\end{lstlisting}

\section*{Command Comparison Table}\label{docs__pandoc__latex__src__command_line_interface_selection__command_line_interface_selection.md__command-comparison-table}

Here\textquotesingle s how common tasks compare across the three options:

{\def\LTcaptype{none} % do not increment counter
\begin{longtable}[]{@{}
  >{\raggedright\arraybackslash}p{(\linewidth - 6\tabcolsep) * \real{0.2110}}
  >{\raggedright\arraybackslash}p{(\linewidth - 6\tabcolsep) * \real{0.2477}}
  >{\raggedright\arraybackslash}p{(\linewidth - 6\tabcolsep) * \real{0.3028}}
  >{\raggedright\arraybackslash}p{(\linewidth - 6\tabcolsep) * \real{0.2385}}@{}}
\toprule\noalign{}
\begin{minipage}[b]{\linewidth}\raggedright
Task
\end{minipage} & \begin{minipage}[b]{\linewidth}\raggedright
Command Prompt
\end{minipage} & \begin{minipage}[b]{\linewidth}\raggedright
PowerShell
\end{minipage} & \begin{minipage}[b]{\linewidth}\raggedright
Git Bash
\end{minipage} \\
\midrule\noalign{}
\endhead
\bottomrule\noalign{}
\endlastfoot
Show current location & \texttt{cd} & \texttt{pwd} & \texttt{pwd} \\
List files & \texttt{dir\ /B} & \texttt{ls\ -n} & \texttt{ls} \\
Go to folder & \texttt{cd\ Documents} & \texttt{cd\ Documents} &
\texttt{cd\ Documents} \\
Go up one level & \texttt{cd\ ..} & \texttt{cd\ ..} & \texttt{cd\ ..} \\
Go home & \texttt{cd\ \%USERPROFILE\%} & \texttt{cd\ \textasciitilde{}}
& \texttt{cd\ \textasciitilde{}} \\
Create folder & \texttt{mkdir\ Projects} & \texttt{mkdir\ Projects} &
\texttt{mkdir\ Projects} \\
Create file & \texttt{echo\ text\ \textgreater{}\ file.txt} &
\texttt{echo\ "text"\ \textgreater{}\ file.txt} &
\texttt{echo\ "text"\ \textgreater{}\ file.txt} \\
Copy file & \texttt{copy\ old.txt\ new.txt} &
\texttt{Copy-Item\ old.txt\ new.txt} & \texttt{cp\ old.txt\ new.txt} \\
Move file & \texttt{move\ old.txt\ folder/} &
\texttt{Move-Item\ old.txt\ folder/} & \texttt{mv\ old.txt\ folder/} \\
Delete file & \texttt{del\ file.txt} & \texttt{Remove-Item\ file.txt} &
\texttt{rm\ file.txt} \\
List with filter & \texttt{dir\ /B\ *.txt} & \texttt{ls\ *.txt} &
\texttt{ls\ *.txt} \\
Save output to file & \texttt{dir\ \textgreater{}\ list.txt} &
\texttt{ls\ \textgreater{}\ list.txt} &
\texttt{ls\ \textgreater{}\ list.txt} \\
Page through output & \texttt{dir\ \textbar{}\ more} &
\texttt{ls\ \textbar{}\ more} & \texttt{ls\ \textbar{}\ less} \\
Search in files & \texttt{findstr\ "text"\ file.txt} &
\texttt{Select-String\ "text"\ file.txt} &
\texttt{grep\ "text"\ file.txt} \\
Show file contents & \texttt{type\ file.txt} &
\texttt{cat\ file.txt} or \texttt{Get-Content} &
\texttt{cat\ file.txt} \\
Create script & \texttt{.bat} files & \texttt{.ps1} files &
\texttt{.sh} files \\
Run script & \texttt{script.bat} & \texttt{.\textbackslash{}script.ps1}
& \texttt{./script.sh} \\
\end{longtable}
}

\section*{Feature Comparison Table}\label{docs__pandoc__latex__src__command_line_interface_selection__command_line_interface_selection.md__feature-comparison-table}

{\def\LTcaptype{none} % do not increment counter
\begin{longtable}[]{@{}
  >{\raggedright\arraybackslash}p{(\linewidth - 6\tabcolsep) * \real{0.3824}}
  >{\raggedright\arraybackslash}p{(\linewidth - 6\tabcolsep) * \real{0.2059}}
  >{\raggedright\arraybackslash}p{(\linewidth - 6\tabcolsep) * \real{0.2059}}
  >{\raggedright\arraybackslash}p{(\linewidth - 6\tabcolsep) * \real{0.2059}}@{}}
\toprule\noalign{}
\begin{minipage}[b]{\linewidth}\raggedright
Feature
\end{minipage} & \begin{minipage}[b]{\linewidth}\raggedright
CMD
\end{minipage} & \begin{minipage}[b]{\linewidth}\raggedright
PowerShell
\end{minipage} & \begin{minipage}[b]{\linewidth}\raggedright
Git Bash
\end{minipage} \\
\midrule\noalign{}
\endhead
\bottomrule\noalign{}
\endlastfoot
Simplicity & Easiest & Moderate & Hardest \\
Beginner-Friendly & Best & Good & Challenging \\
Power/Capability & Basic & Excellent & Excellent \\
Screen Reader Compatible & Perfect & Perfect & Perfect \\
Linux/macOS Skills & Windows-only & Some overlap & Full overlap \\
3D Printing Automation & Adequate & Excellent & Adequate \\
Learning Curve & Gentle & Moderate & Steep \\
Community Support & Moderate & Excellent & Excellent \\
Windows Integration & Perfect & Perfect & Good \\
Installation Difficulty & Built-in & Built-in & Requires Git \\
\end{longtable}
}

\section*{Quick Learner Profile Test}\label{docs__pandoc__latex__src__command_line_interface_selection__command_line_interface_selection.md__quick-learner-profile-test}

Answer these questions to find your best match:

\subsection*{Question 1: Experience Level}\label{docs__pandoc__latex__src__command_line_interface_selection__command_line_interface_selection.md__question-1-experience-level}

\begin{itemize}
\tightlist
\item
  A: I\textquotesingle ve never used a command line Easier paths better (CMD or PowerShell)
\item
  B: I\textquotesingle ve used terminals before Any path works
\item
  C: I use macOS or Linux Git Bash most natural
\end{itemize}

\subsection*{Question 2: What matters most?}\label{docs__pandoc__latex__src__command_line_interface_selection__command_line_interface_selection.md__question-2-what-matters-most}

\begin{itemize}
\tightlist
\item
  A: Simplicity and quick learning Choose CMD
\item
  B: Power and advanced features Choose PowerShell
\item
  C: Consistency across Windows/Mac/Linux Choose Git Bash
\end{itemize}

\subsection*{Question 3: Future goals}\label{docs__pandoc__latex__src__command_line_interface_selection__command_line_interface_selection.md__question-3-future-goals}

\begin{itemize}
\tightlist
\item
  A: Just need to manage files for 3D printing CMD is fine
\item
  B: Advanced automation and scripting PowerShell recommended
\item
  C: Professional development workflows Git Bash best
\end{itemize}

\subsection*{Question 4: Your main concern}\label{docs__pandoc__latex__src__command_line_interface_selection__command_line_interface_selection.md__question-4-your-main-concern}

\begin{itemize}
\tightlist
\item
  A: Don\textquotesingle t want steep learning curve CMD
\item
  B: Want industry-standard skills Git Bash
\item
  C: Want Microsoft\textquotesingle s modern tool PowerShell
\end{itemize}

\section*{Recommendation by Goal}\label{docs__pandoc__latex__src__command_line_interface_selection__command_line_interface_selection.md__recommendation-by-goal}

\subsection*{Goal: "I want to learn the basics and get to 3D printing quickly"}\label{docs__pandoc__latex__src__command_line_interface_selection__command_line_interface_selection.md__goal-i-want-to-learn-the-basics-and-get-to-3d-printing-quickly}

\subsubsection*{Start with CMD (Command Prompt)}\label{docs__pandoc__latex__src__command_line_interface_selection__command_line_interface_selection.md__start-with-cmd-command-prompt}

\begin{itemize}
\tightlist
\item
  Simplest syntax
\item
  Fastest to get productive
\item
  All core concepts transfer to others
\item
  Can switch later if needed
\end{itemize}

\href{docs\%5Cpandoc\%5Clatex\%5Csrc\%5CCMD_Complete_Curriculum.md\#part-2-cmd}{Start CMD Foundation}

\subsection*{Goal: "I want power and professional automation"}\label{docs__pandoc__latex__src__command_line_interface_selection__command_line_interface_selection.md__goal-i-want-power-and-professional-automation}

\subsubsection*{Start with PowerShell}\label{docs__pandoc__latex__src__command_line_interface_selection__command_line_interface_selection.md__start-with-powershell}

\begin{itemize}
\tightlist
\item
  Microsoft\textquotesingle s modern, recommended tool
\item
  Professional-grade capabilities
\item
  Better for complex 3D printing workflows
\item
  Skills are in-demand
\end{itemize}

\href{docs\%5Cpandoc\%5Clatex\%5Csrc\%5CPS_Complete_Curriculum.md\#part-1-powershell}{Start PowerShell Foundation}

\subsection*{Goal: "I want skills that work on Windows, Mac, and Linux"}\label{docs__pandoc__latex__src__command_line_interface_selection__command_line_interface_selection.md__goal-i-want-skills-that-work-on-windows-mac-and-linux}

\subsubsection*{Start with Git Bash}\label{docs__pandoc__latex__src__command_line_interface_selection__command_line_interface_selection.md__start-with-git-bash}

\begin{itemize}
\tightlist
\item
  Unix/bash skills transfer everywhere
\item
  Great preparation for professional development
\item
  Consistent across all platforms
\item
  Growing standard in 3D printing tools
\end{itemize}

\href{docs\%5Cpandoc\%5Clatex\%5Csrc\%5CGitBash_Complete_Curriculum.md\#part-3-git-bash}{Start Git Bash Foundation}

\section*{Can I Switch Paths Later?}\label{docs__pandoc__latex__src__command_line_interface_selection__command_line_interface_selection.md__can-i-switch-paths-later}

Yes, absolutely! All three teach the same fundamental concepts:

\begin{itemize}
\tightlist
\item
  File navigation and organization
\item
  Creating and managing files/folders
\item
  Combining commands for powerful workflows
\item
  Scripting and automation basics
\end{itemize}

Once you learn one, switching to another is quick. The concepts are identical; only the syntax changes.

Example: If you learn CMD first, then later want PowerShell\textquotesingle s power, you\textquotesingle ll find it easy. The command \texttt{cd\ Documents} works the same way in all three.

\section*{Important: All Are Equally Accessible}\label{docs__pandoc__latex__src__command_line_interface_selection__command_line_interface_selection.md__important-all-are-equally-accessible}

\subsection*{Screen readers work perfectly with all three}\label{docs__pandoc__latex__src__command_line_interface_selection__command_line_interface_selection.md__screen-readers-work-perfectly-with-all-three}

\begin{itemize}
\tightlist
\item
  Text-based by nature (perfect for NVDA, JAWS)
\item
  No mouse required
\item
  Output is naturally readable
\item
  Keyboard-only workflows
\end{itemize}

Don\textquotesingle t let accessibility concerns influence your choice. All are fully accessible.

\section*{Getting Started: Your Decision}\label{docs__pandoc__latex__src__command_line_interface_selection__command_line_interface_selection.md__getting-started-your-decision}

Take a moment and choose:

\subsection*{1. I want the simplest path}\label{docs__pandoc__latex__src__command_line_interface_selection__command_line_interface_selection.md__1-i-want-the-simplest-path}

\href{docs\%5Cpandoc\%5Clatex\%5Csrc\%5CCMD_Complete_Curriculum.md\#part-2-cmd}{Command Prompt Foundation}

\begin{itemize}
\tightlist
\item
  Time to first success: \textasciitilde{}30 minutes
\item
  Learning curve: Gentlest
\item
  When to upgrade: Once you\textquotesingle re comfortable and want power
\end{itemize}

\subsection*{2. I want modern, powerful Windows tools}\label{docs__pandoc__latex__src__command_line_interface_selection__command_line_interface_selection.md__2-i-want-modern-powerful-windows-tools}

\href{docs\%5Cpandoc\%5Clatex\%5Csrc\%5CPS_Complete_Curriculum.md\#part-1-powershell}{PowerShell Foundation}

\begin{itemize}
\tightlist
\item
  Time to first success: \textasciitilde{}45 minutes
\item
  Learning curve: Moderate
\item
  Best for: Professional automation, 3D printing workflows
\end{itemize}

\subsection*{3. I want Unix/Linux skills that work everywhere}\label{docs__pandoc__latex__src__command_line_interface_selection__command_line_interface_selection.md__3-i-want-unixlinux-skills-that-work-everywhere}

\href{docs\%5Cpandoc\%5Clatex\%5Csrc\%5CGitBash_Complete_Curriculum.md\#part-3-git-bash}{Git Bash Foundation}

\begin{itemize}
\tightlist
\item
  Time to first success: \textasciitilde{}1 hour
\item
  Learning curve: Steeper but rewarding
\item
  Best for: Professional development, cross-platform work
\end{itemize}

\section*{Not Sure? Here\textquotesingle s What Most People Do}\label{docs__pandoc__latex__src__command_line_interface_selection__command_line_interface_selection.md__not-sure-heres-what-most-people-do}

If you\textquotesingle re reading this and unsure:

\begin{quote}
Start with Command Prompt (CMD). It\textquotesingle s the gentlest introduction. You\textquotesingle ll be productive quickly and can always switch to PowerShell or Git Bash later. The skills transfer completely.
\end{quote}

After completing CMD:

\begin{itemize}
\tightlist
\item
  Want more power? PowerShell is next
\item
  Want Linux skills? Git Bash is next
\item
  Want to stick with CMD? You have all the skills you need
\end{itemize}

\section*{FAQ}\label{docs__pandoc__latex__src__command_line_interface_selection__command_line_interface_selection.md__faq}

Q: Do I need to pick now and stick with it forever?
A: No. Start with one, try another, switch between them. They\textquotesingle re tools. Use what works.

Q: Will my 3D printing skills work in all three?
A: Yes. Once you understand the concepts (file organization, automation, piping), they apply everywhere.

Q: If I pick CMD, can I learn PowerShell later?
A: Absolutely. Many learners do exactly this. CMD gets you productive; PowerShell adds power.

Q: Is Git Bash harder?
A: Slightly, due to path syntax and Unix conventions. But not dramatically. If you take time with it, you\textquotesingle ll learn it.

Q: Which do professional 3D printing developers use?
A: Mix of all three, but Git Bash/Linux is most common in cross-platform teams.

\section*{Ready to Begin?}\label{docs__pandoc__latex__src__command_line_interface_selection__command_line_interface_selection.md__ready-to-begin}

Choose your path above and click to start. Remember:

\begin{itemize}
\tightlist
\item
  Each lesson includes practice exercises
\item
  You can\textquotesingle t break anything
\item
  Mistakes are learning opportunities
\item
  Ask for help if stuck
\end{itemize}

Let\textquotesingle s get you comfortable with the command line!

Other Screen Readers

Dolphin SuperNova (commercial) and Windows Narrator (built-in) are also supported; the workflows and recommendations in this document apply to them. See \url{https://yourdolphin.com/supernova/} and \url{https://support.microsoft.com/narrator} for vendor documentation.

\section{Foundation of Command Line Tools --- Complete Curriculum}\label{docs__pandoc__latex__src__cli_complete_curriculum.md__foundation-of-command-line-tools--complete-curriculum}

\subsection*{Command-Line Fundamentals --- Choose Your Path}\label{docs__pandoc__latex__src__cli_complete_curriculum.md__command-line-fundamentals--choose-your-path}

Welcome! Before diving into 3D design with OpenSCAD, you\textquotesingle ll master command-line fundamentals. This page will help you understand what command-line interfaces are and choose the best path for you.

\subsection*{What is a Command-Line Interface (CLI)?}\label{docs__pandoc__latex__src__cli_complete_curriculum.md__what-is-a-command-line-interface-cli}

A command-line interface is a text-based way to control your computer by typing commands instead of clicking buttons. It\textquotesingle s like sending written instructions to your computer.

\subsubsection*{Why learn it?}\label{docs__pandoc__latex__src__cli_complete_curriculum.md__why-learn-it}

\begin{itemize}
\tightlist
\item
  Speed: Text commands are often faster than clicking through menus
\item
  Precision: Exact control over what your computer does
\item
  Accessibility: Perfect for screen readers - text is naturally readable
\item
  Automation: Repeat tasks automatically
\item
  3D Printing: Essential for batch processing models and integrating tools
\end{itemize}

\subsubsection*{Real-world example}\label{docs__pandoc__latex__src__cli_complete_curriculum.md__real-world-example}

Instead of:

\begin{enumerate}
\tightlist
\item
  Opening File Explorer (click)
\item
  Navigating folders (click, click, click)
\item
  Right-clicking a file (click)
\item
  Selecting "Copy" (click)
\item
  Navigating to destination (click, click)
\item
  Right-clicking (click)
\item
  Selecting "Paste" (click)
\end{enumerate}

You type: \texttt{cp\ myfile.txt\ backup/} and press Enter. Done.

\subsection*{Three Command-Line Options on Windows}\label{docs__pandoc__latex__src__cli_complete_curriculum.md__three-command-line-options-on-windows}

Windows offers three ways to use the command line. All are accessible with screen readers. Here\textquotesingle s how they compare:

\subsubsection*{Option 1: Windows Command Prompt (CMD)}\label{docs__pandoc__latex__src__cli_complete_curriculum.md__option-1-windows-command-prompt-cmd}

What it is: The original Windows command-line (1981-present)

Best for: Absolute beginners, maximum simplicity

Pros:

\begin{itemize}
\tightlist
\item
  Simple commands and syntax
\item
  Minimal learning curve
\item
  Easy to understand error messages
\item
  Great for basic file operations
\item
  Perfect entry point to command-line world
\end{itemize}

Cons:

\begin{itemize}
\tightlist
\item
  Limited advanced features
\item
  Less powerful than alternatives
\item
  No built-in piping (but available)
\item
  Smaller ecosystem
\end{itemize}

Typical command:

\begin{lstlisting}[style=Alabaster, language=cmd]
copy myfile.txt backup\

\end{lstlisting}

\subsubsection*{Option 2: PowerShell}\label{docs__pandoc__latex__src__cli_complete_curriculum.md__option-2-powershell}

What it is: Microsoft\textquotesingle s modern, powerful shell (2006-present)

Best for: Intermediate users, advanced automation

Pros:

\begin{itemize}
\tightlist
\item
  Very powerful for scripting
\item
  Modern syntax and features
\item
  Excellent for 3D printing automation
\item
  Professional workflows
\item
  Large community
\end{itemize}

Cons:

\begin{itemize}
\tightlist
\item
  Steeper learning curve than CMD
\item
  More complex syntax
\item
  More "wordy" commands
\item
  Overkill for simple tasks
\end{itemize}

Typical command:

\begin{lstlisting}[style=Alabaster, language=powershell]
Copy-Item -Path myfile.txt -Destination backup/

\end{lstlisting}

\subsubsection*{Option 3: Git Bash}\label{docs__pandoc__latex__src__cli_complete_curriculum.md__option-3-git-bash}

What it is: A Unix/Linux shell on Windows (runs bash inside Git for Windows)

Best for: Programmers, users familiar with Linux, advanced users

Pros:

\begin{itemize}
\tightlist
\item
  Familiar if you know Linux/Unix
\item
  Powerful piping and text processing
\item
  Consistent with other platforms (macOS, Linux)
\item
  Excellent for advanced workflows
\item
  Industry-standard for developers
\end{itemize}

Cons:

\begin{itemize}
\tightlist
\item
  Requires Git installation
\item
  Steeper learning curve
\item
  Path syntax is different from native Windows
\item
  Less integrated with Windows system tools
\item
  May be "too much" for beginners
\end{itemize}

Typical command:

\begin{lstlisting}[style=Alabaster, language=bash]
cp myfile.txt backup/

\end{lstlisting}

\subsection*{Command Comparison Table}\label{docs__pandoc__latex__src__cli_complete_curriculum.md__command-comparison-table}

Here\textquotesingle s how common tasks compare across the three options:

{\def\LTcaptype{none} % do not increment counter
\begin{longtable}[]{@{}
  >{\raggedright\arraybackslash}p{(\linewidth - 6\tabcolsep) * \real{0.2110}}
  >{\raggedright\arraybackslash}p{(\linewidth - 6\tabcolsep) * \real{0.2477}}
  >{\raggedright\arraybackslash}p{(\linewidth - 6\tabcolsep) * \real{0.3028}}
  >{\raggedright\arraybackslash}p{(\linewidth - 6\tabcolsep) * \real{0.2385}}@{}}
\toprule\noalign{}
\begin{minipage}[b]{\linewidth}\raggedright
Task
\end{minipage} & \begin{minipage}[b]{\linewidth}\raggedright
Command Prompt
\end{minipage} & \begin{minipage}[b]{\linewidth}\raggedright
PowerShell
\end{minipage} & \begin{minipage}[b]{\linewidth}\raggedright
Git Bash
\end{minipage} \\
\midrule\noalign{}
\endhead
\bottomrule\noalign{}
\endlastfoot
Show current location & \texttt{cd} & \texttt{pwd} & \texttt{pwd} \\
List files & \texttt{dir\ /B} & \texttt{ls\ -n} & \texttt{ls} \\
Go to folder & \texttt{cd\ Documents} & \texttt{cd\ Documents} &
\texttt{cd\ Documents} \\
Go up one level & \texttt{cd\ ..} & \texttt{cd\ ..} & \texttt{cd\ ..} \\
Go home & \texttt{cd\ \%USERPROFILE\%} & \texttt{cd\ \textasciitilde{}}
& \texttt{cd\ \textasciitilde{}} \\
Create folder & \texttt{mkdir\ Projects} & \texttt{mkdir\ Projects} &
\texttt{mkdir\ Projects} \\
Create file & \texttt{echo\ text\ \textgreater{}\ file.txt} &
\texttt{echo\ "text"\ \textgreater{}\ file.txt} &
\texttt{echo\ "text"\ \textgreater{}\ file.txt} \\
Copy file & \texttt{copy\ old.txt\ new.txt} &
\texttt{Copy-Item\ old.txt\ new.txt} & \texttt{cp\ old.txt\ new.txt} \\
Move file & \texttt{move\ old.txt\ folder/} &
\texttt{Move-Item\ old.txt\ folder/} & \texttt{mv\ old.txt\ folder/} \\
Delete file & \texttt{del\ file.txt} & \texttt{Remove-Item\ file.txt} &
\texttt{rm\ file.txt} \\
List with filter & \texttt{dir\ /B\ *.txt} & \texttt{ls\ *.txt} &
\texttt{ls\ *.txt} \\
Save output to file & \texttt{dir\ \textgreater{}\ list.txt} &
\texttt{ls\ \textgreater{}\ list.txt} &
\texttt{ls\ \textgreater{}\ list.txt} \\
Page through output & \texttt{dir\ \textbar{}\ more} &
\texttt{ls\ \textbar{}\ more} & \texttt{ls\ \textbar{}\ less} \\
Search in files & \texttt{findstr\ "text"\ file.txt} &
\texttt{Select-String\ "text"\ file.txt} &
\texttt{grep\ "text"\ file.txt} \\
Show file contents & \texttt{type\ file.txt} &
\texttt{cat\ file.txt} or \texttt{Get-Content} &
\texttt{cat\ file.txt} \\
Create script & \texttt{.bat} files & \texttt{.ps1} files &
\texttt{.sh} files \\
Run script & \texttt{script.bat} & \texttt{.\textbackslash{}script.ps1}
& \texttt{./script.sh} \\
\end{longtable}
}

\subsection*{Feature Comparison Table}\label{docs__pandoc__latex__src__cli_complete_curriculum.md__feature-comparison-table}

{\def\LTcaptype{none} % do not increment counter
\begin{longtable}[]{@{}
  >{\raggedright\arraybackslash}p{(\linewidth - 6\tabcolsep) * \real{0.3824}}
  >{\raggedright\arraybackslash}p{(\linewidth - 6\tabcolsep) * \real{0.2059}}
  >{\raggedright\arraybackslash}p{(\linewidth - 6\tabcolsep) * \real{0.2059}}
  >{\raggedright\arraybackslash}p{(\linewidth - 6\tabcolsep) * \real{0.2059}}@{}}
\toprule\noalign{}
\begin{minipage}[b]{\linewidth}\raggedright
Feature
\end{minipage} & \begin{minipage}[b]{\linewidth}\raggedright
CMD
\end{minipage} & \begin{minipage}[b]{\linewidth}\raggedright
PowerShell
\end{minipage} & \begin{minipage}[b]{\linewidth}\raggedright
Git Bash
\end{minipage} \\
\midrule\noalign{}
\endhead
\bottomrule\noalign{}
\endlastfoot
Simplicity & Easiest & Moderate & Hardest \\
Beginner-Friendly & Best & Good & Challenging \\
Power/Capability & Basic & Excellent & Excellent \\
Screen Reader Compatible & Perfect & Perfect & Perfect \\
Linux/macOS Skills & Windows-only & Some overlap & Full overlap \\
3D Printing Automation & Adequate & Excellent & Adequate \\
Learning Curve & Gentle & Moderate & Steep \\
Community Support & Moderate & Excellent & Excellent \\
Windows Integration & Perfect & Perfect & Good \\
Installation Difficulty & Built-in & Built-in & Requires Git \\
\end{longtable}
}

\subsection*{Quick Learner Profile Test}\label{docs__pandoc__latex__src__cli_complete_curriculum.md__quick-learner-profile-test}

Answer these questions to find your best match:

\subsubsection*{Question 1: Experience Level}\label{docs__pandoc__latex__src__cli_complete_curriculum.md__question-1-experience-level}

\begin{itemize}
\tightlist
\item
  A: I\textquotesingle ve never used a command line Easier paths better (CMD or PowerShell)
\item
  B: I\textquotesingle ve used terminals before Any path works
\item
  C: I use macOS or Linux Git Bash most natural
\end{itemize}

\subsubsection*{Question 2: What matters most?}\label{docs__pandoc__latex__src__cli_complete_curriculum.md__question-2-what-matters-most}

\begin{itemize}
\tightlist
\item
  A: Simplicity and quick learning Choose CMD
\item
  B: Power and advanced features Choose PowerShell
\item
  C: Consistency across Windows/Mac/Linux Choose Git Bash
\end{itemize}

\subsubsection*{Question 3: Future goals}\label{docs__pandoc__latex__src__cli_complete_curriculum.md__question-3-future-goals}

\begin{itemize}
\tightlist
\item
  A: Just need to manage files for 3D printing CMD is fine
\item
  B: Advanced automation and scripting PowerShell recommended
\item
  C: Professional development workflows Git Bash best
\end{itemize}

\subsubsection*{Question 4: Your main concern}\label{docs__pandoc__latex__src__cli_complete_curriculum.md__question-4-your-main-concern}

\begin{itemize}
\tightlist
\item
  A: Don\textquotesingle t want steep learning curve CMD
\item
  B: Want industry-standard skills Git Bash
\item
  C: Want Microsoft\textquotesingle s modern tool PowerShell
\end{itemize}

\subsection*{Recommendation by Goal}\label{docs__pandoc__latex__src__cli_complete_curriculum.md__recommendation-by-goal}

\subsubsection*{Goal: "I want to learn the basics and get to 3D printing quickly"}\label{docs__pandoc__latex__src__cli_complete_curriculum.md__goal-i-want-to-learn-the-basics-and-get-to-3d-printing-quickly}

\paragraph*{Start with CMD (Command Prompt)}\label{docs__pandoc__latex__src__cli_complete_curriculum.md__start-with-cmd-command-prompt}

\begin{itemize}
\tightlist
\item
  Simplest syntax
\item
  Fastest to get productive
\item
  All core concepts transfer to others
\item
  Can switch later if needed
\end{itemize}

\hyperref[docs__pandoc__latex__src__cli_complete_curriculum.md__part-2-cmd]{Start CMD Foundation --- jump to Part 2 ↓}

\subsubsection*{Goal: "I want power and professional automation"}\label{docs__pandoc__latex__src__cli_complete_curriculum.md__goal-i-want-power-and-professional-automation}

\paragraph*{Start with PowerShell}\label{docs__pandoc__latex__src__cli_complete_curriculum.md__start-with-powershell}

\begin{itemize}
\tightlist
\item
  Microsoft\textquotesingle s modern, recommended tool
\item
  Professional-grade capabilities
\item
  Better for complex 3D printing workflows
\item
  Skills are in-demand
\end{itemize}

\hyperref[docs__pandoc__latex__src__cli_complete_curriculum.md__part-1-powershell]{Start PowerShell Foundation --- jump to Part 1 ↓}

\subsubsection*{Goal: "I want skills that work on Windows, Mac, and Linux"}\label{docs__pandoc__latex__src__cli_complete_curriculum.md__goal-i-want-skills-that-work-on-windows-mac-and-linux}

\paragraph*{Start with Git Bash}\label{docs__pandoc__latex__src__cli_complete_curriculum.md__start-with-git-bash}

\begin{itemize}
\tightlist
\item
  Unix/bash skills transfer everywhere
\item
  Great preparation for professional development
\item
  Consistent across all platforms
\item
  Growing standard in 3D printing tools
\end{itemize}

\hyperref[docs__pandoc__latex__src__cli_complete_curriculum.md__part-3-gitbash]{Start Git Bash Foundation --- jump to Part 3 ↓}

\subsection*{Can I Switch Paths Later?}\label{docs__pandoc__latex__src__cli_complete_curriculum.md__can-i-switch-paths-later}

Yes, absolutely! All three teach the same fundamental concepts:

\begin{itemize}
\tightlist
\item
  File navigation and organization
\item
  Creating and managing files/folders
\item
  Combining commands for powerful workflows
\item
  Scripting and automation basics
\end{itemize}

Once you learn one, switching to another is quick. The concepts are identical; only the syntax changes.

Example: If you learn CMD first, then later want PowerShell\textquotesingle s power, you\textquotesingle ll find it easy. The command \texttt{cd\ Documents} works the same way in all three.

\subsection*{Important: All Are Equally Accessible}\label{docs__pandoc__latex__src__cli_complete_curriculum.md__important-all-are-equally-accessible}

\subsubsection*{Screen readers work perfectly with all three}\label{docs__pandoc__latex__src__cli_complete_curriculum.md__screen-readers-work-perfectly-with-all-three}

\begin{itemize}
\tightlist
\item
  Text-based by nature (perfect for NVDA, JAWS)
\item
  No mouse required
\item
  Output is naturally readable
\item
  Keyboard-only workflows
\end{itemize}

Note

All three shells --- CMD, PowerShell, and Git Bash --- work perfectly with NVDA, JAWS, Dolphin SuperNova, and Windows Narrator. Don\textquotesingle t let accessibility concerns influence your choice.

Don\textquotesingle t let accessibility concerns influence your choice. All are fully accessible.

\subsection*{Getting Started: Your Decision}\label{docs__pandoc__latex__src__cli_complete_curriculum.md__getting-started-your-decision}

Take a moment and choose:

\subsubsection*{1. I want the simplest path}\label{docs__pandoc__latex__src__cli_complete_curriculum.md__1-i-want-the-simplest-path}

\hyperref[docs__pandoc__latex__src__cli_complete_curriculum.md__part-2-cmd]{Command Prompt Foundation --- jump to Part 2 ↓}

\begin{itemize}
\tightlist
\item
  Time to first success: \textasciitilde{}30 minutes
\item
  Learning curve: Gentlest
\item
  When to upgrade: Once you\textquotesingle re comfortable and want power
\end{itemize}

\subsubsection*{2. I want modern, powerful Windows tools}\label{docs__pandoc__latex__src__cli_complete_curriculum.md__2-i-want-modern-powerful-windows-tools}

\hyperref[docs__pandoc__latex__src__cli_complete_curriculum.md__part-1-powershell]{PowerShell Foundation --- jump to Part 1 ↓}

\begin{itemize}
\tightlist
\item
  Time to first success: \textasciitilde{}45 minutes
\item
  Learning curve: Moderate
\item
  Best for: Professional automation, 3D printing workflows
\end{itemize}

\subsubsection*{3. I want Unix/Linux skills that work everywhere}\label{docs__pandoc__latex__src__cli_complete_curriculum.md__3-i-want-unixlinux-skills-that-work-everywhere}

\hyperref[docs__pandoc__latex__src__cli_complete_curriculum.md__part-3-gitbash]{Git Bash Foundation --- jump to Part 3 ↓}

\begin{itemize}
\tightlist
\item
  Time to first success: \textasciitilde{}1 hour
\item
  Learning curve: Steeper but rewarding
\item
  Best for: Professional development, cross-platform work
\end{itemize}

\subsection*{Not Sure? Here\textquotesingle s What Most People Do}\label{docs__pandoc__latex__src__cli_complete_curriculum.md__not-sure-heres-what-most-people-do}

If you\textquotesingle re reading this and unsure:

Tip

Start with Command Prompt (CMD). It\textquotesingle s the gentlest introduction, and you\textquotesingle ll be productive quickly. You can always switch to PowerShell or Git Bash later --- the concepts transfer completely.

After completing CMD:

\begin{itemize}
\tightlist
\item
  Want more power? PowerShell is next
\item
  Want Linux skills? Git Bash is next
\item
  Want to stick with CMD? You have all the skills you need
\end{itemize}

\subsection*{FAQ}\label{docs__pandoc__latex__src__cli_complete_curriculum.md__faq}

Q: Do I need to pick now and stick with it forever?
A: No. Start with one, try another, switch between them. They\textquotesingle re tools. Use what works.

Q: Will my 3D printing skills work in all three?
A: Yes. Once you understand the concepts (file organization, automation, piping), they apply everywhere.

Q: If I pick CMD, can I learn PowerShell later?
A: Absolutely. Many learners do exactly this. CMD gets you productive; PowerShell adds power.

Q: Is Git Bash harder?
A: Slightly, due to path syntax and Unix conventions. But not dramatically. If you take time with it, you\textquotesingle ll learn it.

Q: Which do professional 3D printing developers use?
A: Mix of all three, but Git Bash/Linux is most common in cross-platform teams.

\subsection*{Ready to Begin?}\label{docs__pandoc__latex__src__cli_complete_curriculum.md__ready-to-begin}

Choose your path above and click to start. Remember:

\begin{itemize}
\tightlist
\item
  Each lesson includes practice exercises
\end{itemize}

Tip

You can\textquotesingle t break your computer by typing a wrong command. The worst that happens is an error message. Mistakes are learning opportunities --- just read the message and try again.

\begin{itemize}
\tightlist
\item
  You can\textquotesingle t break anything
\item
  Mistakes are learning opportunities
\item
  Ask for help if stuck
\end{itemize}

Let\textquotesingle s get you comfortable with the command line!

Other Screen Readers

\subsection*{\texorpdfstring{Dolphin SuperNova (commercial) and Windows Narrator (built-in) are also supported; the workflows and recommendations in this document apply to them. See \url{https://yourdolphin.com/supernova/} and \url{https://support.microsoft.com/narrator} for vendor documentation.}{Dolphin SuperNova (commercial) and Windows Narrator (built-in) are also supported; the workflows and recommendations in this document apply to them. See https://yourdolphin.com/supernova/ and https://support.microsoft.com/narrator for vendor documentation.}}\label{docs__pandoc__latex__src__cli_complete_curriculum.md__dolphin-supernova-commercial-and-windows-narrator-built-in-are-also-supported-the-workflows-and-recommendations-in-this-document-apply-to-them-see-httpsyourdolphincomsupernova-and-httpssupportmicrosoftcomnarrator-for-vendor-documentation}

\textbf{Last Updated:} February 2026
\textbf{Total Duration:} \textasciitilde{}30--45 hours per track (screen reader users)
\textbf{Target Users:} Anyone using a screen reader --- NVDA, JAWS, Dolphin SuperNova, Windows Narrator

This combined curriculum covers three command-line environments taught in parallel. Each section follows the same lesson structure (Pre → 0 → 1 → 2 → 3 → 4 → 5 → 6 → Unit Exam) with shell-specific syntax and examples.

\begin{center}\rule{0.5\linewidth}{0.5pt}\end{center}

\subsection*{Master Table of Contents}\label{docs__pandoc__latex__src__cli_complete_curriculum.md__master-table-of-contents}

\subsubsection*{Part 1 --- PowerShell}\label{docs__pandoc__latex__src__cli_complete_curriculum.md__part-1--powershell}

{\def\LTcaptype{none} % do not increment counter
\begin{longtable}[]{@{}
  >{\raggedright\arraybackslash}p{(\linewidth - 4\tabcolsep) * \real{0.3333}}
  >{\raggedright\arraybackslash}p{(\linewidth - 4\tabcolsep) * \real{0.3333}}
  >{\raggedright\arraybackslash}p{(\linewidth - 4\tabcolsep) * \real{0.3333}}@{}}
\toprule\noalign{}
\begin{minipage}[b]{\linewidth}\raggedright
Section
\end{minipage} & \begin{minipage}[b]{\linewidth}\raggedright
Title
\end{minipage} & \begin{minipage}[b]{\linewidth}\raggedright
Time
\end{minipage} \\
\midrule\noalign{}
\endhead
\bottomrule\noalign{}
\endlastfoot
\hyperref[docs__pandoc__latex__src__cli_complete_curriculum.md__ps-accessibility-guide]{Screen Reader Accessibility Guide}
& NVDA, JAWS, and terminal accessibility reference & 1.5 hrs \\
\hyperref[docs__pandoc__latex__src__cli_complete_curriculum.md__ps-lesson-pre]{Lesson Pre: Your First Terminal}
& Opening PowerShell, first commands, basic navigation & 2--2.5 hrs \\
\hyperref[docs__pandoc__latex__src__cli_complete_curriculum.md__ps-lesson-0]{Lesson 0: Getting Started}
& Paths, shortcuts, tab completion & 1.5 hrs \\
\hyperref[docs__pandoc__latex__src__cli_complete_curriculum.md__ps-lesson-1]{Lesson 1: Navigation}
& Moving around the file system confidently & 2--2.5 hrs \\
\hyperref[docs__pandoc__latex__src__cli_complete_curriculum.md__ps-lesson-2]{Lesson 2: File and Folder Manipulation}
& Create, copy, move, delete files and folders & 2.5--3 hrs \\
\hyperref[docs__pandoc__latex__src__cli_complete_curriculum.md__ps-lesson-3]{Lesson 3: Input, Output, and Piping}
& Chain commands together, redirect output & 2.5--3 hrs \\
\hyperref[docs__pandoc__latex__src__cli_complete_curriculum.md__ps-lesson-4]{Lesson 4: Environment Variables and Aliases}
& Automate settings, create shortcuts & 2--2.5 hrs \\
\hyperref[docs__pandoc__latex__src__cli_complete_curriculum.md__ps-lesson-5]{Lesson 5: Filling in the Gaps}
& Profiles, history, debugging & 2--2.5 hrs \\
\hyperref[docs__pandoc__latex__src__cli_complete_curriculum.md__ps-lesson-6]{Lesson 6: Advanced Techniques}
& Scripts, functions, loops, professional workflows & 4--4.5 hrs \\
\hyperref[docs__pandoc__latex__src__cli_complete_curriculum.md__ps-unit-exam]{Unit Exam}
& Comprehensive assessment --- Lessons 0 through 6 & 60--90 min \\
\end{longtable}
}

\subsubsection*{Part 2 --- Windows Command Prompt (CMD)}\label{docs__pandoc__latex__src__cli_complete_curriculum.md__part-2--windows-command-prompt-cmd}

{\def\LTcaptype{none} % do not increment counter
\begin{longtable}[]{@{}
  >{\raggedright\arraybackslash}p{(\linewidth - 4\tabcolsep) * \real{0.3333}}
  >{\raggedright\arraybackslash}p{(\linewidth - 4\tabcolsep) * \real{0.3333}}
  >{\raggedright\arraybackslash}p{(\linewidth - 4\tabcolsep) * \real{0.3333}}@{}}
\toprule\noalign{}
\begin{minipage}[b]{\linewidth}\raggedright
Section
\end{minipage} & \begin{minipage}[b]{\linewidth}\raggedright
Title
\end{minipage} & \begin{minipage}[b]{\linewidth}\raggedright
Time
\end{minipage} \\
\midrule\noalign{}
\endhead
\bottomrule\noalign{}
\endlastfoot
\hyperref[docs__pandoc__latex__src__cli_complete_curriculum.md__cmd-accessibility-guide]{Screen Reader Accessibility Guide}
& NVDA, JAWS, Narrator, and SuperNova reference for CMD & 1.5 hrs \\
\hyperref[docs__pandoc__latex__src__cli_complete_curriculum.md__cmd-lesson-pre]{Lesson Pre: Your First Terminal}
& Opening CMD, first commands, basic navigation & 1.5--2 hrs \\
\hyperref[docs__pandoc__latex__src__cli_complete_curriculum.md__cmd-lesson-0]{Lesson 0: Getting Started}
& Paths, shortcuts, tab completion & 1.5 hrs \\
\hyperref[docs__pandoc__latex__src__cli_complete_curriculum.md__cmd-lesson-1]{Lesson 1: Navigation}
& Moving around the file system confidently & 2--2.5 hrs \\
\hyperref[docs__pandoc__latex__src__cli_complete_curriculum.md__cmd-lesson-2]{Lesson 2: File and Folder Manipulation}
& Create, copy, move, delete files and folders & 2.5--3 hrs \\
\hyperref[docs__pandoc__latex__src__cli_complete_curriculum.md__cmd-lesson-3]{Lesson 3: Input, Output, and Piping}
& Redirect output, chain commands with pipes & 2.5--3 hrs \\
\hyperref[docs__pandoc__latex__src__cli_complete_curriculum.md__cmd-lesson-4]{Lesson 4: Environment Variables and Aliases}
& Automate settings, create doskey shortcuts & 2--2.5 hrs \\
\hyperref[docs__pandoc__latex__src__cli_complete_curriculum.md__cmd-lesson-5]{Lesson 5: Filling in the Gaps}
& Autorun scripts, history, debugging & 2--2.5 hrs \\
\hyperref[docs__pandoc__latex__src__cli_complete_curriculum.md__cmd-lesson-6]{Lesson 6: Advanced Techniques}
& Batch scripts, subroutines, loops, professional workflows &
4--4.5 hrs \\
\hyperref[docs__pandoc__latex__src__cli_complete_curriculum.md__cmd-unit-exam]{Unit Exam}
& Comprehensive assessment --- Lessons 0 through 6 & 60--90 min \\
\end{longtable}
}

\subsubsection*{Part 3 --- Git Bash}\label{docs__pandoc__latex__src__cli_complete_curriculum.md__part-3--git-bash}

{\def\LTcaptype{none} % do not increment counter
\begin{longtable}[]{@{}
  >{\raggedright\arraybackslash}p{(\linewidth - 4\tabcolsep) * \real{0.3333}}
  >{\raggedright\arraybackslash}p{(\linewidth - 4\tabcolsep) * \real{0.3333}}
  >{\raggedright\arraybackslash}p{(\linewidth - 4\tabcolsep) * \real{0.3333}}@{}}
\toprule\noalign{}
\begin{minipage}[b]{\linewidth}\raggedright
Section
\end{minipage} & \begin{minipage}[b]{\linewidth}\raggedright
Title
\end{minipage} & \begin{minipage}[b]{\linewidth}\raggedright
Time
\end{minipage} \\
\midrule\noalign{}
\endhead
\bottomrule\noalign{}
\endlastfoot
\hyperref[docs__pandoc__latex__src__cli_complete_curriculum.md__gb-accessibility-guide]{Screen Reader Accessibility Guide}
& NVDA, JAWS, Narrator, and SuperNova reference for Git Bash &
1--1.5 hrs \\
\hyperref[docs__pandoc__latex__src__cli_complete_curriculum.md__gb-tutorial]{Tutorial: Git Bash Quick Start}
& Prompt orientation, core commands, path concepts & 30--45 min \\
\hyperref[docs__pandoc__latex__src__cli_complete_curriculum.md__gb-lesson-pre]{Lesson Pre: Your First Terminal}
& Installing Git Bash, opening it, first commands & 1.5--2 hrs \\
\hyperref[docs__pandoc__latex__src__cli_complete_curriculum.md__gb-lesson-0]{Lesson 0: Getting Started}
&
Paths, \texttt{\textasciitilde{}}, \texttt{./}, \texttt{../}, tab completion
& 20--30 min \\
\hyperref[docs__pandoc__latex__src__cli_complete_curriculum.md__gb-lesson-1]{Lesson 1: Navigation}
& \texttt{pwd}, \texttt{cd}, \texttt{ls} and flags, wildcards &
1 class period \\
\hyperref[docs__pandoc__latex__src__cli_complete_curriculum.md__gb-lesson-2]{Lesson 2: File and Folder Manipulation}
&
\texttt{touch}, \texttt{mkdir}, \texttt{cp}, \texttt{mv}, \texttt{rm}, \texttt{rmdir}
& 30--45 min \\
\hyperref[docs__pandoc__latex__src__cli_complete_curriculum.md__gb-lesson-3]{Lesson 3: Input, Output, and Piping}
&
\texttt{echo}, \texttt{cat}, \texttt{\textgreater{}}, \texttt{\textgreater{}\textgreater{}}, \texttt{\textbar{}}, \texttt{clip}, \texttt{grep}
& 1 class period \\
\hyperref[docs__pandoc__latex__src__cli_complete_curriculum.md__gb-lesson-4]{Lesson 4: Environment Variables and Aliases}
&
\texttt{\$VARNAME}, \texttt{PATH}, \texttt{which}, \texttt{alias}, \texttt{.bashrc}
& 30--45 min \\
\hyperref[docs__pandoc__latex__src__cli_complete_curriculum.md__gb-lesson-5]{Lesson 5: Filling in the Gaps}
& \texttt{history}, \texttt{Ctrl+R}, \texttt{.bashrc}, \texttt{./} &
30--45 min \\
\hyperref[docs__pandoc__latex__src__cli_complete_curriculum.md__gb-lesson-6]{Lesson 6: Advanced Techniques}
& Shell scripts, functions, loops, professional workflows &
4--4.5 hrs \\
\hyperref[docs__pandoc__latex__src__cli_complete_curriculum.md__gb-unit-exam]{Unit Exam}
& Comprehensive assessment --- Lessons 0 through 6 & 60--90 min \\
\end{longtable}
}

\begin{center}\rule{0.5\linewidth}{0.5pt}\end{center}

\subsection*{Quick Jump --- By Lesson Across All Three Shells}\label{docs__pandoc__latex__src__cli_complete_curriculum.md__quick-jump--by-lesson-across-all-three-shells}

{\def\LTcaptype{none} % do not increment counter
\begin{longtable}[]{@{}
  >{\raggedright\arraybackslash}p{(\linewidth - 6\tabcolsep) * \real{0.2353}}
  >{\raggedright\arraybackslash}p{(\linewidth - 6\tabcolsep) * \real{0.3235}}
  >{\raggedright\arraybackslash}p{(\linewidth - 6\tabcolsep) * \real{0.1471}}
  >{\raggedright\arraybackslash}p{(\linewidth - 6\tabcolsep) * \real{0.2941}}@{}}
\toprule\noalign{}
\begin{minipage}[b]{\linewidth}\raggedright
Lesson
\end{minipage} & \begin{minipage}[b]{\linewidth}\raggedright
PowerShell
\end{minipage} & \begin{minipage}[b]{\linewidth}\raggedright
CMD
\end{minipage} & \begin{minipage}[b]{\linewidth}\raggedright
Git Bash
\end{minipage} \\
\midrule\noalign{}
\endhead
\bottomrule\noalign{}
\endlastfoot
Accessibility Guide &
\hyperref[docs__pandoc__latex__src__cli_complete_curriculum.md__ps-accessibility-guide]{PS}
&
\hyperref[docs__pandoc__latex__src__cli_complete_curriculum.md__cmd-accessibility-guide]{CMD}
&
\hyperref[docs__pandoc__latex__src__cli_complete_curriculum.md__gb-accessibility-guide]{Git Bash} \\
Lesson Pre &
\hyperref[docs__pandoc__latex__src__cli_complete_curriculum.md__ps-lesson-pre]{PS}
&
\hyperref[docs__pandoc__latex__src__cli_complete_curriculum.md__cmd-lesson-pre]{CMD}
&
\hyperref[docs__pandoc__latex__src__cli_complete_curriculum.md__gb-lesson-pre]{Git Bash} \\
Lesson 0: Getting Started &
\hyperref[docs__pandoc__latex__src__cli_complete_curriculum.md__ps-lesson-0]{PS}
&
\hyperref[docs__pandoc__latex__src__cli_complete_curriculum.md__cmd-lesson-0]{CMD}
&
\hyperref[docs__pandoc__latex__src__cli_complete_curriculum.md__gb-lesson-0]{Git Bash} \\
Lesson 1: Navigation &
\hyperref[docs__pandoc__latex__src__cli_complete_curriculum.md__ps-lesson-1]{PS}
&
\hyperref[docs__pandoc__latex__src__cli_complete_curriculum.md__cmd-lesson-1]{CMD}
&
\hyperref[docs__pandoc__latex__src__cli_complete_curriculum.md__gb-lesson-1]{Git Bash} \\
Lesson 2: File \& Folder Manipulation &
\hyperref[docs__pandoc__latex__src__cli_complete_curriculum.md__ps-lesson-2]{PS}
&
\hyperref[docs__pandoc__latex__src__cli_complete_curriculum.md__cmd-lesson-2]{CMD}
&
\hyperref[docs__pandoc__latex__src__cli_complete_curriculum.md__gb-lesson-2]{Git Bash} \\
Lesson 3: Input, Output \& Piping &
\hyperref[docs__pandoc__latex__src__cli_complete_curriculum.md__ps-lesson-3]{PS}
&
\hyperref[docs__pandoc__latex__src__cli_complete_curriculum.md__cmd-lesson-3]{CMD}
&
\hyperref[docs__pandoc__latex__src__cli_complete_curriculum.md__gb-lesson-3]{Git Bash} \\
Lesson 4: Environment Variables &
\hyperref[docs__pandoc__latex__src__cli_complete_curriculum.md__ps-lesson-4]{PS}
&
\hyperref[docs__pandoc__latex__src__cli_complete_curriculum.md__cmd-lesson-4]{CMD}
&
\hyperref[docs__pandoc__latex__src__cli_complete_curriculum.md__gb-lesson-4]{Git Bash} \\
Lesson 5: Filling in the Gaps &
\hyperref[docs__pandoc__latex__src__cli_complete_curriculum.md__ps-lesson-5]{PS}
&
\hyperref[docs__pandoc__latex__src__cli_complete_curriculum.md__cmd-lesson-5]{CMD}
&
\hyperref[docs__pandoc__latex__src__cli_complete_curriculum.md__gb-lesson-5]{Git Bash} \\
Lesson 6: Advanced Techniques &
\hyperref[docs__pandoc__latex__src__cli_complete_curriculum.md__ps-lesson-6]{PS}
&
\hyperref[docs__pandoc__latex__src__cli_complete_curriculum.md__cmd-lesson-6]{CMD}
&
\hyperref[docs__pandoc__latex__src__cli_complete_curriculum.md__gb-lesson-6]{Git Bash} \\
Unit Exam &
\hyperref[docs__pandoc__latex__src__cli_complete_curriculum.md__ps-unit-exam]{PS}
&
\hyperref[docs__pandoc__latex__src__cli_complete_curriculum.md__cmd-unit-exam]{CMD}
&
\hyperref[docs__pandoc__latex__src__cli_complete_curriculum.md__gb-unit-exam]{Git Bash} \\
\end{longtable}
}

\begin{center}\rule{0.5\linewidth}{0.5pt}\end{center}

\protect\phantomsection\label{docs__pandoc__latex__src__cli_complete_curriculum.md__part-1-powershell}{}

\section*{Part 1 --- PowerShell}\label{docs__pandoc__latex__src__cli_complete_curriculum.md__part-1--powershell-1}

\begin{center}\rule{0.5\linewidth}{0.5pt}\end{center}

\textbf{Last Updated:} February 2026\\
\textbf{Total Duration:} 30--45 hours of instruction and practice (for screen reader users)\\
\textbf{Target Users:} Anyone using a screen reader (NVDA, JAWS, Dolphin SuperNova, Windows Narrator, or other)

\emph{This single file contains the complete PowerShell Foundation curriculum in reading order. All lessons, quizzes, extension problems, and the unit exam are included. No additional files are needed.}

\begin{center}\rule{0.5\linewidth}{0.5pt}\end{center}

\subsection*{Table of Contents}\label{docs__pandoc__latex__src__cli_complete_curriculum.md__table-of-contents}

{\def\LTcaptype{none} % do not increment counter
\begin{longtable}[]{@{}
  >{\raggedright\arraybackslash}p{(\linewidth - 4\tabcolsep) * \real{0.3333}}
  >{\raggedright\arraybackslash}p{(\linewidth - 4\tabcolsep) * \real{0.3333}}
  >{\raggedright\arraybackslash}p{(\linewidth - 4\tabcolsep) * \real{0.3333}}@{}}
\toprule\noalign{}
\begin{minipage}[b]{\linewidth}\raggedright
Section
\end{minipage} & \begin{minipage}[b]{\linewidth}\raggedright
Title
\end{minipage} & \begin{minipage}[b]{\linewidth}\raggedright
Estimated Time
\end{minipage} \\
\midrule\noalign{}
\endhead
\bottomrule\noalign{}
\endlastfoot
\hyperref[docs__pandoc__latex__src__cli_complete_curriculum.md__ps-accessibility-guide]{Screen Reader Accessibility Guide}
& NVDA, JAWS, and terminal accessibility reference & 1.5 hours \\
\hyperref[docs__pandoc__latex__src__cli_complete_curriculum.md__ps-lesson-pre]{Lesson Pre: Your First Terminal}
& Opening PowerShell, first commands, basic navigation & 2--2.5 hours \\
\hyperref[docs__pandoc__latex__src__cli_complete_curriculum.md__ps-lesson-0]{Lesson 0: Getting Started}
& Paths, shortcuts, tab completion & 1.5 hours \\
\hyperref[docs__pandoc__latex__src__cli_complete_curriculum.md__ps-lesson-1]{Lesson 1: Navigation}
& Moving around the file system confidently & 2--2.5 hours \\
\hyperref[docs__pandoc__latex__src__cli_complete_curriculum.md__ps-lesson-2]{Lesson 2: File and Folder Manipulation}
& Create, copy, move, delete files and folders & 2.5--3 hours \\
\hyperref[docs__pandoc__latex__src__cli_complete_curriculum.md__ps-lesson-3]{Lesson 3: Input, Output, and Piping}
& Chain commands together, redirect output & 2.5--3 hours \\
\hyperref[docs__pandoc__latex__src__cli_complete_curriculum.md__ps-lesson-4]{Lesson 4: Environment Variables and Aliases}
& Automate settings, create shortcuts & 2--2.5 hours \\
\hyperref[docs__pandoc__latex__src__cli_complete_curriculum.md__ps-lesson-5]{Lesson 5: Filling in the Gaps}
& Profiles, history, debugging & 2--2.5 hours \\
\hyperref[docs__pandoc__latex__src__cli_complete_curriculum.md__ps-lesson-6]{Lesson 6: Advanced Techniques}
& Scripts, functions, loops, professional workflows & 4--4.5 hours \\
\hyperref[docs__pandoc__latex__src__cli_complete_curriculum.md__ps-unit-exam]{Unit Exam}
& Comprehensive assessment --- Lessons 0 through 6 & 60--90 minutes \\
\end{longtable}
}

\textbf{Total:} approximately 30--45 hours including practice exercises

\begin{center}\rule{0.5\linewidth}{0.5pt}\end{center}

\subsection*{How to Use This File}\label{docs__pandoc__latex__src__cli_complete_curriculum.md__how-to-use-this-file}

\textbf{If you have never used a terminal before:} read from the top in order. Do not skip lessons --- each builds on the one before it.

\textbf{If you have used a terminal but not with a screen reader:} start with the Accessibility Guide, skim Lesson Pre, then continue from Lesson 0.

\textbf{If you are experienced with both:} jump directly to the lesson you need. Use the quiz at the end of each lesson to self-test before moving on.

\textbf{Reference while working:} the Accessibility Guide and Lesson 0 contain material you will return to throughout the curriculum. Keep a copy open or bookmark those sections.

\begin{center}\rule{0.5\linewidth}{0.5pt}\end{center}

\subsection*{Curriculum Rules (Read These First)}\label{docs__pandoc__latex__src__cli_complete_curriculum.md__curriculum-rules-read-these-first}

\textbf{Rule 1 --- Always know where you are.}\\
Every session, first thing: \texttt{pwd}. If you don\textquotesingle t know your path, you\textquotesingle ll get lost.

\textbf{Rule 2 --- Check before you delete.}\\
Before any \texttt{rm} command: \texttt{ls\ -n} to confirm what you\textquotesingle re about to remove.

\textbf{Rule 3 --- Use \texttt{-n} with \texttt{ls}.}\\
Always \texttt{ls\ -n} --- never bare \texttt{ls}. The \texttt{-n} flag gives one name per line, which screen readers can follow.

\textbf{Rule 4 --- When lost, redirect to a file.}\\
If output is confusing: \texttt{command\ \textgreater{}\ output.txt} then \texttt{notepad.exe\ output.txt}. This always works.

\textbf{Rule 5 --- Save everything you create.}\\
Keep all practice files in a dedicated learning folder: \texttt{mkdir\ \textasciitilde{}/PS\_Learning}.

\begin{center}\rule{0.5\linewidth}{0.5pt}\end{center}

\begin{center}\rule{0.5\linewidth}{0.5pt}\end{center}

\protect\phantomsection\label{docs__pandoc__latex__src__cli_complete_curriculum.md__ps-accessibility-guide}{}

\section*{Screen Reader Accessibility Guide}\label{docs__pandoc__latex__src__cli_complete_curriculum.md__screen-reader-accessibility-guide}

\begin{center}\rule{0.5\linewidth}{0.5pt}\end{center}

Target Users: NVDA, JAWS, and other screen reader users\\
Last Updated: 2026

This guide supports the PowerShell Foundation curriculum and helps screen reader users navigate and work efficiently with PowerShell on Windows.

\subsection*{Table of Contents}\label{docs__pandoc__latex__src__cli_complete_curriculum.md__table-of-contents-1}

\begin{enumerate}
\tightlist
\item
  \hyperref[docs__pandoc__latex__src__cli_complete_curriculum.md__ps-getting-started]{Getting Started with Screen Readers}
\item
  \hyperref[docs__pandoc__latex__src__cli_complete_curriculum.md__ps-nvda-tips]{NVDA-Specific Tips}
\item
  \hyperref[docs__pandoc__latex__src__cli_complete_curriculum.md__ps-jaws-tips]{JAWS-Specific Tips}
\item
  \hyperref[docs__pandoc__latex__src__cli_complete_curriculum.md__ps-general-terminal]{General Terminal Accessibility}
\item
  \hyperref[docs__pandoc__latex__src__cli_complete_curriculum.md__ps-long-output]{Working with Long Output}
\item
  \hyperref[docs__pandoc__latex__src__cli_complete_curriculum.md__ps-shortcuts]{Keyboard Shortcuts Reference}
\item
  \hyperref[docs__pandoc__latex__src__cli_complete_curriculum.md__ps-troubleshooting]{Troubleshooting}
\end{enumerate}

\protect\phantomsection\label{docs__pandoc__latex__src__cli_complete_curriculum.md__ps-getting-started}{}

\subsection*{Getting Started with Screen Readers}\label{docs__pandoc__latex__src__cli_complete_curriculum.md__getting-started-with-screen-readers}

\subsubsection*{Which Screen Reader Should I Use?}\label{docs__pandoc__latex__src__cli_complete_curriculum.md__which-screen-reader-should-i-use}

NVDA is free and often recommended for new users; JAWS is a powerful commercial option. Also consider Dolphin SuperNova and Windows Narrator:

\begin{itemize}
\tightlist
\item
  Dolphin SuperNova: commercial speech, braille and magnification (check vendor docs for keyboard mappings).
\item
  Windows Narrator: built into Windows and useful for quick access without installing third-party software.
\end{itemize}

\subsubsection*{Before You Start}\label{docs__pandoc__latex__src__cli_complete_curriculum.md__before-you-start}

\begin{enumerate}
\tightlist
\item
  Start your screen reader before opening PowerShell.
\item
  Open PowerShell and listen for the window title and prompt.
\item
  If silent, press Alt+Tab to find the window.
\end{enumerate}

\subsubsection*{What is PowerShell?}\label{docs__pandoc__latex__src__cli_complete_curriculum.md__what-is-powershell}

PowerShell is the modern Windows shell and scripting environment. Common commands include \texttt{Get-ChildItem}, \texttt{Get-Content}, and \texttt{Out-File}. PowerShell provides richer objects and piping than CMD.

\protect\phantomsection\label{docs__pandoc__latex__src__cli_complete_curriculum.md__ps-nvda-tips}{}

\subsection*{NVDA-Specific Tips}\label{docs__pandoc__latex__src__cli_complete_curriculum.md__nvda-specific-tips}

NVDA is available from \url{https://www.nvaccess.org/}

\subsubsection*{Dolphin SuperNova}\label{docs__pandoc__latex__src__cli_complete_curriculum.md__dolphin-supernova}

Dolphin SuperNova: \url{https://yourdolphin.com/supernova/} --- commercial option providing speech and magnification; consult vendor guides for features and commands.

\subsubsection*{Windows Narrator}\label{docs__pandoc__latex__src__cli_complete_curriculum.md__windows-narrator}

Windows Narrator: \url{https://support.microsoft.com/narrator} --- built-in Narrator has a different set of commands; it can be enabled via Windows Settings \textgreater{} Accessibility.

\subsubsection*{Key Commands for PowerShell}\label{docs__pandoc__latex__src__cli_complete_curriculum.md__key-commands-for-powershell}

{\def\LTcaptype{none} % do not increment counter
\begin{longtable}[]{@{}
  >{\raggedright\arraybackslash}p{(\linewidth - 2\tabcolsep) * \real{0.2361}}
  >{\raggedright\arraybackslash}p{(\linewidth - 2\tabcolsep) * \real{0.7639}}@{}}
\toprule\noalign{}
\begin{minipage}[b]{\linewidth}\raggedright
Command
\end{minipage} & \begin{minipage}[b]{\linewidth}\raggedright
What It Does
\end{minipage} \\
\midrule\noalign{}
\endhead
\bottomrule\noalign{}
\endlastfoot
NVDA+Home & Read the current line (your command or output) \\
NVDA+Down Arrow & Read from cursor to end of screen \\
NVDA+Up Arrow & Read from top to cursor \\
NVDA+Page Down & Read next page \\
NVDA+Page Up & Read previous page \\
NVDA+F7 & Open the Review Mode viewer (can scroll through text) \\
\end{longtable}
}

\subsubsection*{Example: Reading Long Output}\label{docs__pandoc__latex__src__cli_complete_curriculum.md__example-reading-long-output}

If \texttt{Get-ChildItem} produces many lines, redirect to a file and open it in Notepad:

\begin{lstlisting}[style=Alabaster, language=powershell]
Get-ChildItem -Name > list.txt
notepad list.txt

\end{lstlisting}

\texttt{-Name} prints one item per line (screen reader friendly).

\protect\phantomsection\label{docs__pandoc__latex__src__cli_complete_curriculum.md__ps-jaws-tips}{}

\subsection*{JAWS-Specific Tips}\label{docs__pandoc__latex__src__cli_complete_curriculum.md__jaws-specific-tips}

JAWS is available from \url{https://www.freedomscientific.com/}

\subsubsection*{Key Commands for PowerShell}\label{docs__pandoc__latex__src__cli_complete_curriculum.md__key-commands-for-powershell-1}

{\def\LTcaptype{none} % do not increment counter
\begin{longtable}[]{@{}ll@{}}
\toprule\noalign{}
Command & What It Does \\
\midrule\noalign{}
\endhead
\bottomrule\noalign{}
\endlastfoot
Insert+Down Arrow & Read line by line downward \\
Insert+Up Arrow & Read line by line upward \\
Insert+Page Down & Read next page of text \\
Insert+Page Up & Read previous page of text \\
\end{longtable}
}

\subsubsection*{Example: Reading Long Output}\label{docs__pandoc__latex__src__cli_complete_curriculum.md__example-reading-long-output-1}

\begin{enumerate}
\tightlist
\item
  Redirect: \texttt{Get-ChildItem\ -Name\ \textgreater{}\ list.txt}
\item
  Open Notepad: \texttt{notepad\ list.txt}
\item
  Use Insert+Ctrl+Down to read full contents.
\end{enumerate}

\protect\phantomsection\label{docs__pandoc__latex__src__cli_complete_curriculum.md__ps-general-terminal}{}

\subsection*{General Terminal Accessibility}\label{docs__pandoc__latex__src__cli_complete_curriculum.md__general-terminal-accessibility}

\subsubsection*{Understanding the PowerShell Layout}\label{docs__pandoc__latex__src__cli_complete_curriculum.md__understanding-the-powershell-layout}

PowerShell shows a title bar, content area, and a prompt that looks like:

\begin{lstlisting}[style=Alabaster, language=powershell]
PS C:\Users\YourName>

\end{lstlisting}

\subsubsection*{Navigation Sequence}\label{docs__pandoc__latex__src__cli_complete_curriculum.md__navigation-sequence}

\begin{enumerate}
\tightlist
\item
  Screen reader announces the title
\item
  Then it announces the prompt line
\item
  Anything above prompt is prior output
\end{enumerate}

\protect\phantomsection\label{docs__pandoc__latex__src__cli_complete_curriculum.md__ps-long-output}{}

\subsection*{Working with Long Output}\label{docs__pandoc__latex__src__cli_complete_curriculum.md__working-with-long-output}

\subsubsection*{Solution 1: Redirect to a File}\label{docs__pandoc__latex__src__cli_complete_curriculum.md__solution-1-redirect-to-a-file}

\begin{lstlisting}[style=Alabaster, language=powershell]
Get-ChildItem -Name > list.txt
notepad list.txt

\end{lstlisting}

\subsubsection*{Solution 2: Use Pagination}\label{docs__pandoc__latex__src__cli_complete_curriculum.md__solution-2-use-pagination}

\begin{lstlisting}[style=Alabaster, language=powershell]
Get-Content largefile.txt | more

\end{lstlisting}

\subsubsection*{Solution 3: Filter Output}\label{docs__pandoc__latex__src__cli_complete_curriculum.md__solution-3-filter-output}

\begin{lstlisting}[style=Alabaster, language=powershell]
Get-ChildItem -Name | Where-Object { $_ -like "*.scad" }

\end{lstlisting}

\subsubsection*{Solution 4: Count Before Displaying}\label{docs__pandoc__latex__src__cli_complete_curriculum.md__solution-4-count-before-displaying}

\begin{lstlisting}[style=Alabaster, language=powershell]
(Get-ChildItem -Name).Count

\end{lstlisting}

\protect\phantomsection\label{docs__pandoc__latex__src__cli_complete_curriculum.md__ps-shortcuts}{}

\subsection*{Keyboard Shortcuts Reference}\label{docs__pandoc__latex__src__cli_complete_curriculum.md__keyboard-shortcuts-reference}

{\def\LTcaptype{none} % do not increment counter
\begin{longtable}[]{@{}ll@{}}
\toprule\noalign{}
Key & Action \\
\midrule\noalign{}
\endhead
\bottomrule\noalign{}
\endlastfoot
Up Arrow & Show previous command \\
Down Arrow & Show next command \\
Tab & Auto-complete file/folder names \\
Home & Jump to start of line \\
End & Jump to end of line \\
Ctrl+C & Stop command \\
Enter & Run command \\
\end{longtable}
}

\protect\phantomsection\label{docs__pandoc__latex__src__cli_complete_curriculum.md__ps-troubleshooting}{}

\subsection*{Troubleshooting}\label{docs__pandoc__latex__src__cli_complete_curriculum.md__troubleshooting}

\subsubsection*{Problem: "I Can\textquotesingle t Hear the Output"}\label{docs__pandoc__latex__src__cli_complete_curriculum.md__problem-i-cant-hear-the-output}

\begin{enumerate}
\tightlist
\item
  Redirect to file and open in Notepad.
\item
  Use End to jump to the end of text.
\end{enumerate}

\subsubsection*{Problem: "Tab Completion Isn\textquotesingle t Working"}\label{docs__pandoc__latex__src__cli_complete_curriculum.md__problem-tab-completion-isnt-working}

\begin{enumerate}
\tightlist
\item
  Type some characters before Tab.
\end{enumerate}

\subsubsection*{Problem: "Command Not Found"}\label{docs__pandoc__latex__src__cli_complete_curriculum.md__problem-command-not-found}

\begin{enumerate}
\tightlist
\item
  Use \texttt{Get-Command\ programname} to check availability.
\end{enumerate}

\subsection*{Pro Tips}\label{docs__pandoc__latex__src__cli_complete_curriculum.md__pro-tips}

\begin{enumerate}
\tightlist
\item
  Use \texttt{Get-ChildItem\ -Name} for one-per-line listings.
\item
  Use \texttt{Out-File\ -FilePath\ list.txt} to capture output with encoding options.
\end{enumerate}

\subsection*{Recommended Workflow}\label{docs__pandoc__latex__src__cli_complete_curriculum.md__recommended-workflow}

\begin{enumerate}
\tightlist
\item
  \texttt{Set-Location} (or \texttt{cd}) to the project folder
\item
  \texttt{Get-ChildItem\ -Name} to list files
\item
  Redirect large output to files and open in Notepad
\end{enumerate}

\subsection*{Additional Resources}\label{docs__pandoc__latex__src__cli_complete_curriculum.md__additional-resources}

\begin{itemize}
\tightlist
\item
  NVDA Documentation: \url{https://www.nvaccess.org/documentation/}
\item
  JAWS Documentation: \url{https://www.freedomscientific.com/support/}
\item
  PowerShell Documentation: \url{https://docs.microsoft.com/powershell}
\item
  NVDA Documentation: \url{https://www.nvaccess.org/documentation/}
\item
  JAWS Documentation: \url{https://www.freedomscientific.com/support/}
\item
  Dolphin SuperNova: \url{https://yourdolphin.com/supernova/}
\item
  Windows Narrator: \url{https://support.microsoft.com/narrator}
\end{itemize}

\begin{center}\rule{0.5\linewidth}{0.5pt}\end{center}

\protect\phantomsection\label{docs__pandoc__latex__src__cli_complete_curriculum.md__ps-lesson-pre}{}

\section*{Lesson Pre: Your First Terminal}\label{docs__pandoc__latex__src__cli_complete_curriculum.md__lesson-pre-your-first-terminal}

\begin{center}\rule{0.5\linewidth}{0.5pt}\end{center}

Duration: 1.5-2 hours (for screen reader users)\\
Prerequisites: None - this is the starting point\\
Accessibility Note: This lesson is designed specifically for screen reader users (NVDA, JAWS)

\subsection*{What is a Terminal?}\label{docs__pandoc__latex__src__cli_complete_curriculum.md__what-is-a-terminal}

A terminal (also called a command line or shell) is a text-based interface where you type commands instead of clicking buttons. Think of it like sending written instructions to your computer instead of pointing and clicking.

Why learn this?

\begin{itemize}
\tightlist
\item
  Faster and more precise work (especially for 3D printing scripts and automation)
\item
  Essential for programming and using tools like OpenSCAD
\item
  Accessibility: Many command line tools work perfectly with screen readers
\item
  Scripting: Automate repetitive tasks
\end{itemize}

\subsection*{Opening PowerShell for the First Time}\label{docs__pandoc__latex__src__cli_complete_curriculum.md__opening-powershell-for-the-first-time}

\subsubsection*{3DMake: Your First Real-World Goal}\label{docs__pandoc__latex__src__cli_complete_curriculum.md__3dmake-your-first-real-world-goal}

\textbf{3DMake} (\href{https://github.com/tdeck/3dmake}{github.com/tdeck/3dmake}) is a free, accessibility-first command-line tool for designing and 3D-printing models entirely from the terminal --- no mouse required. Every PowerShell skill in this curriculum directly supports a real 3DMake workflow.

Here is the end goal --- one command that builds your OpenSCAD model, slices it, and sends it to the printer:

\begin{lstlisting}[style=Alabaster, language=powershell]
# Navigate to your project folder
cd ~/Documents/3d_projects/my_cube

# Build, slice, and print in a single command
3dm build slice print

\end{lstlisting}

Keep this goal in mind as you work through each lesson.

\subsection*{Opening PowerShell for the First Time}\label{docs__pandoc__latex__src__cli_complete_curriculum.md__opening-powershell-for-the-first-time-1}

\subsubsection*{On Windows}\label{docs__pandoc__latex__src__cli_complete_curriculum.md__on-windows}

Method 1: Search (Easiest)

\begin{enumerate}
\tightlist
\item
  Press the Windows key alone
\item
  You should hear "Search"
\item
  Type: \texttt{PowerShell}
\item
  You\textquotesingle ll hear search results appear
\item
  Press Enter to open the first result (Windows PowerShell)
\item
  PowerShell will open in a new window
\end{enumerate}

Method 2: Using the Start Menu

\begin{enumerate}
\tightlist
\item
  Press Windows key + X (opens the Quick Link menu)
\item
  Look for "Windows PowerShell" or "Terminal"
\item
  Press Enter
\end{enumerate}

Method 3: From File Explorer

\begin{enumerate}
\tightlist
\item
  Open File Explorer
\item
  Navigate to any folder
\item
  In the menu bar, select "File" -\textgreater{} "Open Windows PowerShell here"
\item
  PowerShell opens in that folder location
\end{enumerate}

\subsubsection*{First Connection: Understanding the Prompt}\label{docs__pandoc__latex__src__cli_complete_curriculum.md__first-connection-understanding-the-prompt}

When PowerShell opens, your screen reader will announce the window title and then the prompt. The prompt is where you type commands.

What you\textquotesingle ll hear:

\begin{lstlisting}[style=Alabaster]
PS C:\Users\YourName>

\end{lstlisting}

What this means:

\begin{itemize}
\tightlist
\item
  \texttt{PS} = "PowerShell" indicator
\item
  \texttt{C:\textbackslash{}Users\textbackslash{}YourName} = Your current location (the "path")
\item
  \texttt{\textgreater{}} = The prompt is ready for your input
\end{itemize}

Important

Your cursor is blinking right after the \texttt{\textgreater{}}. This is where you type.

\subsection*{Your First Commands (Screen Reader Edition)}\label{docs__pandoc__latex__src__cli_complete_curriculum.md__your-first-commands-screen-reader-edition}

\subsubsection*{\texorpdfstring{Command 1: "Where Am I?" - \texttt{pwd}}{Command 1: "Where Am I?" - pwd}}\label{docs__pandoc__latex__src__cli_complete_curriculum.md__command-1-where-am-i---pwd}

What it does: Tells you your current location

Type this:

\begin{lstlisting}[style=Alabaster, language=powershell]
pwd

\end{lstlisting}

Press Enter

What you\textquotesingle ll hear:
Your screen reader will announce the current path, something like:

\begin{lstlisting}[style=Alabaster]
C:\Users\YourName

\end{lstlisting}

Understanding paths:

\begin{itemize}
\tightlist
\item
  Paths show your location in the file system (like a mailing address)
\item
  Windows paths use backslashes: \texttt{C:\textbackslash{}Users\textbackslash{}YourName\textbackslash{}Documents}
\item
  Think of it like folders inside folders: \texttt{C:\textbackslash{}} (main drive) -\textgreater{} \texttt{Users} -\textgreater{} \texttt{YourName} -\textgreater{} \texttt{Documents}
\end{itemize}

\subsubsection*{\texorpdfstring{Command 2: "What\textquotesingle s Here?" - \texttt{ls\ -n}}{Command 2: "What\textquotesingle s Here?" - ls -n}}\label{docs__pandoc__latex__src__cli_complete_curriculum.md__command-2-whats-here---ls--n}

What it does: Lists all files and folders in your current location. The \texttt{-n} flag makes it screen-reader friendly (names only, one per line)

Type this:

\begin{lstlisting}[style=Alabaster, language=powershell]
ls -n

\end{lstlisting}

Press Enter

What you\textquotesingle ll hear:
Your screen reader will announce each file and folder name, one per line:

\begin{lstlisting}[style=Alabaster]
Desktop
Documents
Downloads
Music
Pictures
...

\end{lstlisting}

Why \texttt{-n}?

\begin{itemize}
\tightlist
\item
  Without \texttt{-n}, PowerShell shows files in columns (hard to read with a screen reader)
\item
  With \texttt{-n}, each file/folder is on its own line (perfect for screen readers)
\end{itemize}

\subsubsection*{\texorpdfstring{Command 3: "Go There" - \texttt{cd\ Documents}}{Command 3: "Go There" - cd Documents}}\label{docs__pandoc__latex__src__cli_complete_curriculum.md__command-3-go-there---cd-documents}

What it does: Changes your location (navigates to a folder)

Type this:

\begin{lstlisting}[style=Alabaster, language=powershell]
cd Documents

\end{lstlisting}

Press Enter

What you\textquotesingle ll hear:
The prompt changes to show your new location. You might hear something like:

\begin{lstlisting}[style=Alabaster]
PS C:\Users\YourName\Documents>

\end{lstlisting}

Practice navigation:

\begin{enumerate}
\tightlist
\item
  Run \texttt{pwd} to confirm you\textquotesingle re in Documents
\item
  Run \texttt{ls\ -n} to see what files are in Documents
\item
  Try going back: \texttt{cd\ ..} (the \texttt{..} means "go up one level")
\item
  Run \texttt{pwd} again to confirm
\item
  Go back to Documents: \texttt{cd\ Documents}
\end{enumerate}

\subsection*{Reading Screen Reader Output (Critical Skills)}\label{docs__pandoc__latex__src__cli_complete_curriculum.md__reading-screen-reader-output-critical-skills}

\subsubsection*{Dealing with Long Lists}\label{docs__pandoc__latex__src__cli_complete_curriculum.md__dealing-with-long-lists}

When you run \texttt{ls\ -n} in a folder with many files, the list might be very long. Your screen reader might announce 50+ items rapidly.

Solution 1: Save to a File

\begin{lstlisting}[style=Alabaster, language=powershell]
ls -n > list.txt
notepad.exe list.txt

\end{lstlisting}

This saves the list to a file and opens it in Notepad where you can read it more slowly.

Solution 2: Search Within the Output

\begin{lstlisting}[style=Alabaster, language=powershell]
ls -n | findstr "search-term"

\end{lstlisting}

Example: If you\textquotesingle re looking for files containing "scad", type:

\begin{lstlisting}[style=Alabaster, language=powershell]
ls -n | findstr "scad"

\end{lstlisting}

\subsubsection*{Navigating Tab Completion}\label{docs__pandoc__latex__src__cli_complete_curriculum.md__navigating-tab-completion}

One of the most powerful screen reader tricks is Tab completion:

How it works:

\begin{enumerate}
\tightlist
\item
  Type the first few letters of a folder or file name
\item
  Press Tab
\item
  PowerShell automatically completes the rest
\end{enumerate}

Example:

\begin{enumerate}
\tightlist
\item
  You\textquotesingle re in \texttt{C:\textbackslash{}Users\textbackslash{}YourName\textgreater{}}
\item
  Type: \texttt{cd\ Doc}
\item
  Press Tab
\item
  PowerShell auto-completes it to: \texttt{cd\ Documents}
\item
  Press Enter to go there
\end{enumerate}

With a screen reader:

\begin{enumerate}
\tightlist
\item
  As you type \texttt{Doc}, your screen reader announces each letter
\item
  When you press Tab, PowerShell types the rest and your screen reader announces the full command
\item
  This is much faster than typing the whole thing
\end{enumerate}

\subsection*{Creating and Editing Files}\label{docs__pandoc__latex__src__cli_complete_curriculum.md__creating-and-editing-files}

\subsubsection*{Create a Simple File}\label{docs__pandoc__latex__src__cli_complete_curriculum.md__create-a-simple-file}

Type this:

\begin{lstlisting}[style=Alabaster, language=powershell]
echo "Hello, PowerShell!" > hello.txt

\end{lstlisting}

What this does:

\begin{itemize}
\tightlist
\item
  \texttt{echo} sends text to the screen (or file)
\item
  \texttt{"Hello,\ PowerShell!"} is the text
\item
  \texttt{\textgreater{}} redirects it to a file called \texttt{hello.txt}
\end{itemize}

\subsubsection*{Read the File Back}\label{docs__pandoc__latex__src__cli_complete_curriculum.md__read-the-file-back}

Type this:

\begin{lstlisting}[style=Alabaster, language=powershell]
cat hello.txt

\end{lstlisting}

What you\textquotesingle ll hear: Your screen reader announces:

\begin{lstlisting}[style=Alabaster]
Hello, PowerShell!

\end{lstlisting}

\subsubsection*{Open and Edit the File}\label{docs__pandoc__latex__src__cli_complete_curriculum.md__open-and-edit-the-file}

Type this:

\begin{lstlisting}[style=Alabaster, language=powershell]
notepad.exe hello.txt

\end{lstlisting}

This opens the file in Notepad where you can edit it with your screen reader.

\subsection*{Essential Keyboard Shortcuts}\label{docs__pandoc__latex__src__cli_complete_curriculum.md__essential-keyboard-shortcuts}

These work in PowerShell and are crucial for screen reader users:

{\def\LTcaptype{none} % do not increment counter
\begin{longtable}[]{@{}
  >{\raggedright\arraybackslash}p{(\linewidth - 2\tabcolsep) * \real{0.2152}}
  >{\raggedright\arraybackslash}p{(\linewidth - 2\tabcolsep) * \real{0.7848}}@{}}
\toprule\noalign{}
\begin{minipage}[b]{\linewidth}\raggedright
Key Combination
\end{minipage} & \begin{minipage}[b]{\linewidth}\raggedright
What It Does
\end{minipage} \\
\midrule\noalign{}
\endhead
\bottomrule\noalign{}
\endlastfoot
Up Arrow &
Shows your previous command (press again to go further back) \\
Down Arrow & Shows your next command (if you went back) \\
Tab & Auto-completes folder/file names \\
\end{longtable}
}

Important

\textbf{Ctrl+C} stops any running command immediately. If a command seems frozen or is producing endless output, press Ctrl+C to regain control of the prompt.

\textbar{} Ctrl+C          \textbar{} Stops a running command                                      \textbar{}
\textbar{} Ctrl+L          \textbar{} Clears the screen                                            \textbar{}
\textbar{} Enter           \textbar{} Runs the command                                             \textbar{}

Screen reader tip: These all work perfectly with your screen reader. Try them!

\subsection*{Screen Reader-Specific Tips}\label{docs__pandoc__latex__src__cli_complete_curriculum.md__screen-reader-specific-tips}

\subsubsection*{NVDA Users}\label{docs__pandoc__latex__src__cli_complete_curriculum.md__nvda-users}

\begin{enumerate}
\item
  Reading Command Output:

  \begin{itemize}
  \tightlist
  \item
    Use NVDA+Home to read the current line
  \item
    Use NVDA+Down Arrow to read to the end of the screen
  \item
    Use NVDA+Page Down to read the next page
  \end{itemize}
\item
  Reviewing Text:

  \begin{itemize}
  \tightlist
  \item
    Use NVDA+Shift+Page Up to review text above
  \end{itemize}
\end{enumerate}

\subsubsection*{JAWS Users}\label{docs__pandoc__latex__src__cli_complete_curriculum.md__jaws-users}

\begin{enumerate}
\item
  Reading Output:

  \begin{itemize}
  \tightlist
  \item
    Use Insert+Down Arrow to read line-by-line
  \item
    Use Insert+Page Down to read by page
  \item
    Use Insert+End to jump to the end of text
  \end{itemize}
\item
  Reading All Text:

  \begin{itemize}
  \tightlist
  \item
    Use Insert+Down Arrow repeatedly
  \item
    Or use Insert+Ctrl+Down to read to the end
  \end{itemize}
\end{enumerate}

\subsubsection*{Common Issue: "I Can\textquotesingle t Hear the Output"}\label{docs__pandoc__latex__src__cli_complete_curriculum.md__common-issue-i-cant-hear-the-output}

Problem: You run a command but don\textquotesingle t hear the output

Solutions:

\begin{enumerate}
\tightlist
\item
  Make sure your cursor is at the prompt (try pressing End or Ctrl+End)
\item
  Use Up Arrow to go back to your previous command and review it
\item
  Try redirecting to a file: \texttt{command\ \textgreater{}\ output.txt} then open the file
\item
  In NVDA: Try pressing NVDA+F7 to open the Review Mode viewer
\end{enumerate}

\subsection*{Practice Exercises}\label{docs__pandoc__latex__src__cli_complete_curriculum.md__practice-exercises}

Complete these in order. Take your time with each one:

\subsubsection*{Exercise 1: Basic Navigation}\label{docs__pandoc__latex__src__cli_complete_curriculum.md__exercise-1-basic-navigation}

\begin{enumerate}
\tightlist
\item
  Open PowerShell
\item
  Run \texttt{pwd} and note your location
\item
  Run \texttt{ls\ -n} and listen to what\textquotesingle s there
\item
  Try \texttt{cd\ Documents} or another folder
\item
  Run \texttt{pwd} to confirm your new location
\item
  Run \texttt{ls\ -n} in this new location
\end{enumerate}

Goal: You should be comfortable knowing where you are and what\textquotesingle s around you

\subsubsection*{Exercise 2: Using Tab Completion}\label{docs__pandoc__latex__src__cli_complete_curriculum.md__exercise-2-using-tab-completion}

\begin{enumerate}
\tightlist
\item
  In your home directory, type \texttt{cd\ D} (just the letter D)
\item
  Press Tab
\item
  PowerShell should auto-complete to a folder starting with D
\item
  Repeat with other folder names
\item
  Try typing a longer name: \texttt{cd\ Down} and Tab to \texttt{Downloads}
\end{enumerate}

Goal: Tab completion should feel natural

\subsubsection*{Exercise 3: Creating and Viewing Files}\label{docs__pandoc__latex__src__cli_complete_curriculum.md__exercise-3-creating-and-viewing-files}

\begin{enumerate}
\tightlist
\item
  Create a file: \texttt{echo\ "Test\ content"\ \textgreater{}\ test.txt}
\item
  View it: \texttt{cat\ test.txt}
\item
  Create another: \texttt{echo\ "Line\ 2"\ \textgreater{}\ another.txt}
\item
  List both: \texttt{ls\ -n}
\end{enumerate}

Goal: You understand create, view, and list operations

\subsubsection*{Exercise 4: Going Up Levels}\label{docs__pandoc__latex__src__cli_complete_curriculum.md__exercise-4-going-up-levels}

\begin{enumerate}
\tightlist
\item
  Navigate into several folders: \texttt{cd\ Documents}, then \texttt{cd\ folder1}, etc.
\item
  From deep inside, use \texttt{cd\ ..} multiple times to go back up
\item
  After each \texttt{cd\ ..}, run \texttt{pwd} to confirm your location
\end{enumerate}

Goal: You understand relative navigation with \texttt{..}

\subsubsection*{Exercise 5: Redirecting Output}\label{docs__pandoc__latex__src__cli_complete_curriculum.md__exercise-5-redirecting-output}

\begin{enumerate}
\tightlist
\item
  Create a list: \texttt{ls\ -n\ \textgreater{}\ directory\_list.txt}
\item
  Open it: \texttt{notepad.exe\ directory\_list.txt}
\item
  Read it with your screen reader
\item
  Close Notepad
\item
  Verify the file exists: \texttt{ls\ -n\ \textbar{}\ findstr\ "directory"}
\end{enumerate}

Goal: You can save long outputs to files for easier reading

\subsection*{Checkpoint Questions}\label{docs__pandoc__latex__src__cli_complete_curriculum.md__checkpoint-questions}

After completing this lesson, you should be able to answer:

\begin{enumerate}
\tightlist
\item
  What does \texttt{pwd} do?
\item
  What does \texttt{ls\ -n} do?
\item
  Why do we use \texttt{-n} with \texttt{ls}?
\item
  What path are you in right now?
\item
  How do you navigate to a new folder?
\item
  How do you go up one level?
\item
  What\textquotesingle s the Tab key for?
\item
  What does \texttt{echo\ "text"\ \textgreater{}\ file.txt} do?
\item
  How do you read a file back?
\item
  How do you stop a command that\textquotesingle s running?
\end{enumerate}

You should be able to answer all 10 with confidence before moving to Lesson 0.

\subsection*{Common Questions}\label{docs__pandoc__latex__src__cli_complete_curriculum.md__common-questions}

Q: Do I need to use PowerShell? Can I use Command Prompt (cmd.exe)?
A: PowerShell is more powerful and works better with modern tools. We recommend PowerShell, but Command Prompt basics are similar.

Tip

If your screen reader isn\textquotesingle t announcing terminal output, use \texttt{command\ \textgreater{}\ file.txt} to save output to a file, then open it with \texttt{notepad.exe\ file.txt} for reliable reading. This trick works with all four screen readers throughout the curriculum.

Q: Why is my screen reader not reading the output?
A: This is common. Use \texttt{command\ \textgreater{}\ file.txt} to save output to a file, then open it with Notepad for reliable reading.

Q: What if I type something wrong?
A: Just press Enter and you\textquotesingle ll see an error message. Type the correct command on the next line. No harm done!

Q: How do I get help with a command?
A: Type \texttt{Get-Help\ command-name} (we\textquotesingle ll cover this in Lesson 0)

Q: Can I make PowerShell more accessible?
A: Yes! We\textquotesingle ll cover customization in Lesson 5.

\subsection*{Next Steps}\label{docs__pandoc__latex__src__cli_complete_curriculum.md__next-steps}

Once you\textquotesingle re comfortable with these basics:

\begin{itemize}
\tightlist
\item
  Move to Lesson 0: Getting Started for deeper path understanding
\item
  Then continue through Lessons 1--5 for full terminal mastery
\end{itemize}

\subsection*{Resources}\label{docs__pandoc__latex__src__cli_complete_curriculum.md__resources}

\begin{itemize}
\tightlist
\item
  Microsoft PowerShell Docs: \url{https://docs.microsoft.com/powershell/}
\item
  NVDA Screen Reader: \url{https://www.nvaccess.org/}
\item
  JAWS Screen Reader: \url{https://www.freedomscientific.com/products/software/jaws/}
\item
  Windows Terminal Accessibility: \url{https://docs.microsoft.com/windows/terminal/}
\end{itemize}

\protect\phantomsection\label{docs__pandoc__latex__src__cli_complete_curriculum.md__ps-troubleshooting}{}

\subsection*{Troubleshooting}\label{docs__pandoc__latex__src__cli_complete_curriculum.md__troubleshooting-1}

{\def\LTcaptype{none} % do not increment counter
\begin{longtable}[]{@{}
  >{\raggedright\arraybackslash}p{(\linewidth - 2\tabcolsep) * \real{0.2545}}
  >{\raggedright\arraybackslash}p{(\linewidth - 2\tabcolsep) * \real{0.7455}}@{}}
\toprule\noalign{}
\begin{minipage}[b]{\linewidth}\raggedright
Issue
\end{minipage} & \begin{minipage}[b]{\linewidth}\raggedright
Solution
\end{minipage} \\
\midrule\noalign{}
\endhead
\bottomrule\noalign{}
\endlastfoot
PowerShell won\textquotesingle t open &
Try searching Windows, or right-click a folder and select "Open PowerShell here" \\
Can\textquotesingle t hear the output &
Try redirecting to a file: \texttt{command\ \textgreater{}\ output.txt} \\
Tab completion not working &
Make sure you typed at least one character before pressing Tab \\
Command not found &
Make sure you spelled it correctly; try \texttt{Get-Command} to see available commands \\
Stuck in a command & Press Ctrl+C to stop it \\
\end{longtable}
}

Still stuck? The checkpoint questions and exercises are your best teacher. Work through them multiple times until comfortable.

Other Screen Readers

Dolphin SuperNova (commercial) and Windows Narrator (built-in) are also supported; the workflows and recommendations in this document apply to them. See \url{https://yourdolphin.com/supernova/} and \url{https://support.microsoft.com/narrator} for vendor documentation.

\begin{center}\rule{0.5\linewidth}{0.5pt}\end{center}

\protect\phantomsection\label{docs__pandoc__latex__src__cli_complete_curriculum.md__ps-lesson-0}{}

\section*{Lesson 0: Getting Started --- Layout, Paths, and the Shell}\label{docs__pandoc__latex__src__cli_complete_curriculum.md__lesson-0-getting-started--layout-paths-and-the-shell}

\begin{center}\rule{0.5\linewidth}{0.5pt}\end{center}

Estimated time: 20-30 minutes

\subsection*{Learning Objectives}\label{docs__pandoc__latex__src__cli_complete_curriculum.md__learning-objectives}

\begin{itemize}
\tightlist
\item
  Launch PowerShell and locate the prompt
\item
  Understand path notation and shortcuts (\texttt{\textasciitilde{}}, \texttt{./}, \texttt{../})
\item
  Use tab completion to navigate quickly
\end{itemize}

\subsection*{Materials}\label{docs__pandoc__latex__src__cli_complete_curriculum.md__materials}

\begin{itemize}
\tightlist
\item
  Computer with PowerShell
\item
  Editor (Notepad/VS Code)
\end{itemize}

\subsection*{Step-by-step Tasks}\label{docs__pandoc__latex__src__cli_complete_curriculum.md__step-by-step-tasks}

\begin{enumerate}
\tightlist
\item
  Open PowerShell and note the prompt (it includes the current path).
\item
  Run \texttt{pwd} and say or note the printed path.
\item
  Use \texttt{ls\ -n} to list names in your home directory.
\item
  Practice \texttt{cd\ Documents}, \texttt{cd\ ../} and \texttt{cd\ \textasciitilde{}} until comfortable.
\item
  Try tab-completion: type \texttt{cd\ \textasciitilde{}/D} and press Tab.
\end{enumerate}

\subsection*{Checkpoints}\label{docs__pandoc__latex__src__cli_complete_curriculum.md__checkpoints}

\begin{itemize}
\tightlist
\item
  Confirm you can state your current path and move to \texttt{Documents}.
\end{itemize}

\subsection*{Quiz - Lesson PS.0}\label{docs__pandoc__latex__src__cli_complete_curriculum.md__quiz---lesson-ps0}

\subsection*{3DMake Quick-Start: Your First Real Project}\label{docs__pandoc__latex__src__cli_complete_curriculum.md__3dmake-quick-start-your-first-real-project}

You now know enough PowerShell to set up a real 3DMake project. Do this now --- it is the fastest path to a printed object.

\textbf{Step 1 --- Download and run the 3DMake setup wizard:}

\begin{lstlisting}[style=Alabaster, language=powershell]
# Go to your Downloads folder (where you extracted 3dmake_windows.zip)
cd ~/Downloads/3dmake

# Run the guided setup (answer the questions about your printer)
./3dm setup

\end{lstlisting}

\textbf{Step 2 --- Create your first project:}

\begin{lstlisting}[style=Alabaster, language=powershell]
# Create a folder for all your 3D projects
mkdir ~/Documents/3d_projects

# Create your first project folder and initialize it
mkdir ~/Documents/3d_projects/my_cube
cd ~/Documents/3d_projects/my_cube
3dm new

\end{lstlisting}

\textbf{Step 3 --- Confirm the project structure:}

\begin{lstlisting}[style=Alabaster, language=powershell]
# You should see three items: src, build, and 3dmake.toml
ls -n

\end{lstlisting}

You just used \texttt{mkdir}, \texttt{cd}, and \texttt{ls\ -n} on a real project. You are ready to print.

\subsection*{Checkpoints}\label{docs__pandoc__latex__src__cli_complete_curriculum.md__checkpoints-1}

\begin{itemize}
\tightlist
\item
  Confirm you can state your current path and move to \texttt{Documents}.
\end{itemize}

\subsection*{Quiz - Lesson PS.0}\label{docs__pandoc__latex__src__cli_complete_curriculum.md__quiz---lesson-ps0-1}

\begin{enumerate}
\tightlist
\item
  What is a path?
\item
  What does \texttt{\textasciitilde{}} mean?
\item
  How do you autocomplete a path?
\item
  How do you go up one directory?
\item
  What command lists only names (\texttt{ls} flag)?
\item
  True or False: On Windows, PowerShell uses backslashes (\texttt{\textbackslash{}}) in paths, but forward slashes (/) are also accepted.
\item
  Explain the difference between an absolute path and a relative path.
\item
  If you are in \texttt{C:\textbackslash{}Users\textbackslash{}YourName\textbackslash{}Documents} and you type \texttt{cd\ ../}, where do you end up?
\item
  What happens when you press Tab while typing a folder name in PowerShell?
\item
  Describe a practical reason why understanding paths is important for a 3D printing workflow.
\item
  What does \texttt{./} mean in a path, and when would you use it?
\item
  If a folder path contains spaces (e.g., \texttt{Program\ Files}), how do you navigate to it with \texttt{cd}?
\item
  Explain what the prompt \texttt{PS\ C:\textbackslash{}Users\textbackslash{}YourName\textgreater{}} tells you about your current state.
\item
  How would you navigate to your home directory from any location using a single command?
\item
  What is the advantage of using relative paths (like \texttt{../}) versus absolute paths in automation scripts?
\end{enumerate}

\subsection*{Extension Problems}\label{docs__pandoc__latex__src__cli_complete_curriculum.md__extension-problems}

\begin{enumerate}
\tightlist
\item
  Create a nested folder and practice \texttt{cd} into it by typing partial names and using Tab.
\item
  Use \texttt{ls\ -n\ -af} to list only files in a folder.
\item
  Save \texttt{pwd} output to a file and open it in Notepad.
\item
  Try \texttt{cd} into a folder whose name contains spaces; observe how quotes are handled.
\item
  Create a short note file and open it from PowerShell.
\item
  Build a folder structure that mirrors your project organization; navigate to each level and document the path.
\item
  Create a script that prints your current path and the total number of files in it; run it from different locations.
\item
  Investigate the special paths (e.g., \texttt{\$HOME}, \texttt{\$PSScriptRoot}); write down what each contains and when you\textquotesingle d use them.
\item
  Compare absolute vs. relative paths by navigating to the same folder using each method; explain which is easier for automation.
\item
  Create a PowerShell function that changes to a frequently-used folder and lists its contents in one command; test it from different starting locations.
\item
  Navigate to three different locations and at each one note the prompt, the path from \texttt{pwd}, and verify you understand what each shows.
\item
  Create a complex folder tree (at least 5 levels deep) and navigate it using only relative paths; verify your location at each step.
\item
  Document all shortcuts you know (\texttt{\textasciitilde{}}, \texttt{./}, \texttt{../}, \texttt{\$HOME}) and demonstrate each one works as expected.
\item
  Write a guide for a peer on how to understand the PowerShell prompt and path notation without using GUI file explorer.
\item
  Create a troubleshooting flowchart: if someone says "I don\textquotesingle t know where I am," what commands do you give them to find out?
\end{enumerate}

\subsection*{References}\label{docs__pandoc__latex__src__cli_complete_curriculum.md__references}

\begin{itemize}
\tightlist
\item
  Microsoft. (2024). \emph{PowerShell scripting overview and documentation}. \url{https://learn.microsoft.com/powershell/scripting/overview}
\item
  Microsoft. (2024). \emph{Filesystem navigation in PowerShell}. \url{https://learn.microsoft.com/powershell/scripting/learn/shell/navigate-the-filesystem}
\item
  Microsoft. (2024). \emph{Accessibility features in PowerShell ISE}. \url{https://learn.microsoft.com/powershell/scripting/windows-powershell/ise/accessibility-in-windows-powershell-ise}
\end{itemize}

\subsection*{Helpful Resources}\label{docs__pandoc__latex__src__cli_complete_curriculum.md__helpful-resources}

\begin{itemize}
\tightlist
\item
  \href{https://learn.microsoft.com/powershell/scripting/overview}{PowerShell Basics - Microsoft Learn}
\item
  \href{https://learn.microsoft.com/powershell/scripting/learn/shell/navigate-the-filesystem}{Filesystem Navigation Guide}
\item
  \href{https://poshcode.gitbook.io/powershell-faq/src/getting-started/filesystem-navigation}{Understanding Path Notation}
\item
  \href{https://learn.microsoft.com/powershell/module/psreadline/about/about_psreadline_functions}{Tab Completion Reference}
\item
  \href{https://learn.microsoft.com/powershell/scripting/windows-powershell/ise/accessibility-in-windows-powershell-ise}{Accessibility in PowerShell ISE}
\end{itemize}

\begin{center}\rule{0.5\linewidth}{0.5pt}\end{center}

\protect\phantomsection\label{docs__pandoc__latex__src__cli_complete_curriculum.md__ps-lesson-1}{}

\section*{Lesson 1: Navigation --- Moving Around Your File System}\label{docs__pandoc__latex__src__cli_complete_curriculum.md__lesson-1-navigation--moving-around-your-file-system}

\begin{center}\rule{0.5\linewidth}{0.5pt}\end{center}

Duration: 1 class period\\
Prerequisite: Lesson 0 (Getting Started)

Learning Objectives By the end of this lesson, you will be able to:

\begin{itemize}
\tightlist
\item
  Use \texttt{pwd} to print your current location
\item
  Use \texttt{cd} to move between directories
\item
  Use \texttt{ls} (and its flags) to list files and folders
\item
  Use wildcards \texttt{*} and \texttt{?} to filter listings
\item
  Navigate relative vs. absolute paths
\item
  Search for files by name and extension
\end{itemize}

Materials

\begin{itemize}
\tightlist
\item
  PowerShell
\item
  Text editor (Notepad or VS Code)
\end{itemize}

\subsection*{Commands Covered in This Lesson}\label{docs__pandoc__latex__src__cli_complete_curriculum.md__commands-covered-in-this-lesson}

{\def\LTcaptype{none} % do not increment counter
\begin{longtable}[]{@{}
  >{\raggedright\arraybackslash}p{(\linewidth - 2\tabcolsep) * \real{0.2571}}
  >{\raggedright\arraybackslash}p{(\linewidth - 2\tabcolsep) * \real{0.7429}}@{}}
\toprule\noalign{}
\begin{minipage}[b]{\linewidth}\raggedright
Command
\end{minipage} & \begin{minipage}[b]{\linewidth}\raggedright
What It Does
\end{minipage} \\
\midrule\noalign{}
\endhead
\bottomrule\noalign{}
\endlastfoot
\texttt{pwd} & Print Working Directory - shows where you are \\
\texttt{cd\ path} & Change Directory - move to a new location \\
\texttt{ls} & List - shows files and folders in current location \\
\texttt{ls\ -n} & List names only (screen reader friendly) \\
\texttt{ls\ -n\ -af} & List names of files only \\
\texttt{ls\ -n\ -ad} & List names of directories only \\
\texttt{ls\ *.extension} & List files matching a pattern \\
\end{longtable}
}

\subsection*{\texorpdfstring{\texttt{pwd} - Where Am I?}{pwd - Where Am I?}}\label{docs__pandoc__latex__src__cli_complete_curriculum.md__pwd---where-am-i}

Type \texttt{pwd} and press \texttt{Enter}. PowerShell prints the full path to your current location.

\begin{lstlisting}[style=Alabaster, language=powershell]
pwd
# Output: C:\Users\YourName

\end{lstlisting}

When to use: Always run this if you\textquotesingle re unsure of your current location.

\subsection*{\texorpdfstring{\texttt{cd} - Changing Directories}{cd - Changing Directories}}\label{docs__pandoc__latex__src__cli_complete_curriculum.md__cd---changing-directories}

\texttt{cd} stands for "change directory."

\begin{lstlisting}[style=Alabaster, language=powershell]
# Go to Documents
cd Documents
# Go up one level to parent directory
cd ..
# Go to home directory
cd ~
# Go to a specific path
cd C:\Users\YourName\Documents\3D_Projects

\end{lstlisting}

\subsection*{\texorpdfstring{\texttt{ls} - Listing Files and Folders}{ls - Listing Files and Folders}}\label{docs__pandoc__latex__src__cli_complete_curriculum.md__ls---listing-files-and-folders}

Use \texttt{ls\ -n} for screen reader compatibility.

\begin{lstlisting}[style=Alabaster, language=powershell]
# List all files and folders (names only)
ls -n
# List only files (no folders)
ls -n -af
# List only folders (no files)
ls -n -ad

\end{lstlisting}

\subsection*{Wildcards - Finding Files by Pattern}\label{docs__pandoc__latex__src__cli_complete_curriculum.md__wildcards---finding-files-by-pattern}

Wildcards help you find files without typing the full name.

\texttt{*} (asterisk) matches any number of characters:

\begin{lstlisting}[style=Alabaster, language=powershell]
# List all .scad files
ls -n *.scad
# List all files starting with "part"
ls -n part*
# List all files ending with "_final"
ls -n *_final*

\end{lstlisting}

\texttt{?} (question mark) matches exactly one character:

\begin{lstlisting}[style=Alabaster, language=powershell]
# Find files like model1.scad, model2.scad (but not model12.scad)
ls -n model?.scad

\end{lstlisting}

\subsection*{Step-by-step Practice}\label{docs__pandoc__latex__src__cli_complete_curriculum.md__step-by-step-practice}

\subsection*{3DMake in Practice: Navigating a Project}\label{docs__pandoc__latex__src__cli_complete_curriculum.md__3dmake-in-practice-navigating-a-project}

Navigation commands are exactly what you need to explore and manage 3DMake projects:

\begin{lstlisting}[style=Alabaster, language=powershell]
# Go to your project root
cd ~/Documents/3d_projects/my_cube

# See the project layout: src/, build/, 3dmake.toml
ls -n

# Enter the source folder where your .scad model files live
cd src
ls -n

# Open the main model in your text editor
3dm edit-model

# Go back to project root, then check the build output
cd ..
cd build
ls -n *.stl        # list any STL files already built

# Return to project root
cd ..

# Use wildcards to find all projects with "cube" in the name
ls -n ~/Documents/3d_projects/*cube*

\end{lstlisting}

Tab completion makes navigation fast: type \texttt{cd\ \textasciitilde{}/Doc} and press Tab.

\subsection*{Step-by-step Practice}\label{docs__pandoc__latex__src__cli_complete_curriculum.md__step-by-step-practice-1}

\begin{enumerate}
\tightlist
\item
  Run \texttt{pwd} and confirm your location
\item
  Move to \texttt{Documents}: \texttt{cd\ Documents}
\item
  Confirm you moved: \texttt{pwd}
\item
  List files and folders: \texttt{ls\ -n}
\item
  List only files: \texttt{ls\ -n\ -af}
\item
  Go back up: \texttt{cd\ ..}
\item
  Search for files: \texttt{ls\ -n\ *.txt}
\end{enumerate}

\subsection*{Checkpoints}\label{docs__pandoc__latex__src__cli_complete_curriculum.md__checkpoints-2}

After this lesson, you should be able to:

\begin{itemize}
\tightlist
\item[$\square$]
  Navigate to any folder using \texttt{cd}
\item[$\square$]
  Confirm your location with \texttt{pwd}
\item[$\square$]
  List files and folders with \texttt{ls\ -n}
\item[$\square$]
  Use wildcards to find files by pattern
\item[$\square$]
  Move between absolute and relative paths confidently
\end{itemize}

\subsection*{Quiz - Lesson PS.1}\label{docs__pandoc__latex__src__cli_complete_curriculum.md__quiz---lesson-ps1}

\begin{enumerate}
\tightlist
\item
  What does \texttt{pwd} show?
\item
  How do you list directories only with \texttt{ls}?
\item
  What wildcard matches any number of characters?
\item
  How do you list files with the \texttt{.scad} extension?
\item
  Give an example of an absolute path and a relative path.
\item
  True or False: The \texttt{*} wildcard matches exactly one character.
\item
  Explain the difference between \texttt{ls\ -n} and \texttt{ls\ -n\ -ad}.
\item
  Write a command that would list all \texttt{.txt} files in your Documents folder using a wildcard.
\item
  How would you search for files containing "part" in their name across multiple files?
\item
  Describe a practical scenario where using wildcards saves time in a 3D printing workflow.
\item
  What happens when you use \texttt{ls\ -n\ part?.scad} versus \texttt{ls\ -n\ part*.scad}?
\item
  How would you navigate to a folder whose name contains both spaces and special characters?
\item
  If you\textquotesingle re in \texttt{/Documents/Projects/3D} and you want to go to \texttt{/Documents/Resources}, what command would you use?
\item
  Write a command sequence that navigates to the Downloads folder, lists only files, then returns to home.
\item
  Explain the purpose of using \texttt{ls\ -n\ -af} specifically in a screen reader context.
\end{enumerate}

\subsection*{Extension Problems}\label{docs__pandoc__latex__src__cli_complete_curriculum.md__extension-problems-1}

\begin{enumerate}
\tightlist
\item
  Write a one-line script that lists \texttt{.scad} files and saves to \texttt{scad\_list.txt}.
\item
  Use \texttt{ls\ -n\ \textasciitilde{}/Documents\ \textbar{}\ more} to page through long listings.
\item
  Combine \texttt{ls} with \texttt{Select-String} to search for a filename pattern.
\item
  Create a shortcut alias in the session for a long path and test it.
\item
  Practice tab-completion in a directory with many similarly named files.
\item
  Build a PowerShell script that recursively lists all \texttt{.scad} and \texttt{.stl} files in a directory tree; save the results to a file.
\item
  Compare the output of \texttt{ls}, \texttt{Get-ChildItem}, and \texttt{gci} to understand PowerShell aliasing; document what each command does.
\item
  Create a filtering command that displays only files modified in the last 7 days; test it on your documents folder.
\item
  Write a non-visual guide to PowerShell navigation; include descriptions of common patterns and how to verify directory contents audibly.
\item
  Develop a navigation workflow for a typical 3D printing project: move between CAD, slicing, and print-log folders efficiently; document the commands.
\item
  Create a complex wildcard search: find all files in a folder and subfolders that match multiple patterns (e.g., \texttt{*\_v1.*\ OR\ *\_final.*}).
\item
  Build a script that navigates through a folder tree, counts files at each level, and reports the structure.
\item
  Document the output differences between \texttt{ls\ -n}, \texttt{ls\ -n\ -af}, \texttt{ls\ -n\ -ad}, and \texttt{Get-ChildItem}; explain when to use each.
\item
  Create a navigation "cheat sheet" as a PowerShell script that prints common paths and how to navigate to them.
\item
  Design a project folder structure on your computer, document each path, then create a script that validates all folders exist.
\end{enumerate}

\subsection*{References}\label{docs__pandoc__latex__src__cli_complete_curriculum.md__references-1}

\begin{itemize}
\tightlist
\item
  Microsoft. (2024). \emph{Get-ChildItem cmdlet reference}. \url{https://learn.microsoft.com/powershell/module/microsoft.powershell.management/get-childitem}
\item
  Microsoft. (2024). \emph{PowerShell wildcards and filtering}. \url{https://learn.microsoft.com/powershell/scripting/learn/shell/using-wildcards}
\item
  Microsoft. (2024). \emph{Navigation best practices in PowerShell}. \url{https://learn.microsoft.com/powershell/scripting/learn/shell/navigate-the-filesystem}
\end{itemize}

\subsection*{Helpful Resources}\label{docs__pandoc__latex__src__cli_complete_curriculum.md__helpful-resources-1}

\begin{itemize}
\tightlist
\item
  \href{https://learn.microsoft.com/powershell/module/microsoft.powershell.management/get-childitem}{Get-ChildItem Cmdlet Reference}
\item
  \href{https://learn.microsoft.com/powershell/scripting/learn/shell/using-wildcards}{PowerShell Wildcards and Filtering}
\item
  \href{https://learn.microsoft.com/powershell/scripting/learn/shell/navigate-the-filesystem}{Navigation Best Practices}
\item
  \href{https://poshcode.gitbook.io/powershell-faq/src/getting-started/filesystem-navigation}{Relative and Absolute Paths}
\item
  \href{https://learn.microsoft.com/powershell/scripting/windows-powershell/ise/accessibility-in-windows-powershell-ise}{Screen Reader Tips for PowerShell}
\end{itemize}

\begin{center}\rule{0.5\linewidth}{0.5pt}\end{center}

\protect\phantomsection\label{docs__pandoc__latex__src__cli_complete_curriculum.md__ps-lesson-2}{}

\section*{Lesson 2: File and Folder Manipulation}\label{docs__pandoc__latex__src__cli_complete_curriculum.md__lesson-2-file-and-folder-manipulation}

\begin{center}\rule{0.5\linewidth}{0.5pt}\end{center}

Estimated time: 30-45 minutes

Learning Objectives

\begin{itemize}
\tightlist
\item
  Create, copy, move, and delete files and folders from PowerShell
\item
  Use \texttt{ni}, \texttt{mkdir}, \texttt{cp}, \texttt{mv}, \texttt{rm}, and \texttt{rmdir} safely
\item
  Understand when operations are permanent and how to confirm results
\end{itemize}

Materials

\begin{itemize}
\tightlist
\item
  PowerShell
\item
  Small practice folder for exercises
\end{itemize}

Step-by-step Tasks

\begin{enumerate}
\tightlist
\item
  Create a practice directory: \texttt{mkdir\ \textasciitilde{}/Documents/PS\_Practice} and \texttt{cd} into it.
\item
  Create two files: \texttt{ni\ file1.txt} and \texttt{ni\ file2.txt}.
\item
  Copy \texttt{file1.txt} to \texttt{file1\_backup.txt} with \texttt{cp} and confirm with \texttt{ls\ -n}.
\item
  Rename \texttt{file2.txt} to \texttt{notes.txt} using \texttt{mv} and confirm.
\item
  Delete \texttt{file1.txt} with \texttt{rm} and verify the backup remains.
\end{enumerate}

Checkpoints

\begin{itemize}
\tightlist
\item
  After step 3 you should see both the original and the backup file.
\end{itemize}

\subsection*{3DMake in Practice: Organizing Your Projects}\label{docs__pandoc__latex__src__cli_complete_curriculum.md__3dmake-in-practice-organizing-your-projects}

File management commands are essential for 3DMake project organization:

\begin{lstlisting}[style=Alabaster, language=powershell]
# Create a new project from scratch
mkdir ~/Documents/3d_projects/bottle_cap
cd ~/Documents/3d_projects/bottle_cap
3dm new

# Back up a working project before making big changes
cd ~/Documents/3d_projects
cp -r bottle_cap bottle_cap_backup

# Rename the main model file to something more specific
cd bottle_cap/src
mv main.scad bottle_cap.scad

# List only .scad files to confirm the rename
ls -n *.scad

# Delete a draft STL you no longer need
rm ~/Documents/3d_projects/bottle_cap/build/draft.stl

# Now build the renamed model (using the -m flag to specify the name)
cd ~/Documents/3d_projects/bottle_cap
3dm build -m bottle_cap

# See the fresh output in the build folder
ls -n build/

\end{lstlisting}

Checkpoints

\begin{itemize}
\tightlist
\item
  After step 3 you should see both the original and the backup file.
\end{itemize}

\subsection*{Quiz - Lesson PS.2}\label{docs__pandoc__latex__src__cli_complete_curriculum.md__quiz---lesson-ps2}

\begin{enumerate}
\tightlist
\item
  How do you create an empty file from PowerShell?
\item
  What command copies a file?
\item
  How do you rename a file?
\item
  What does \texttt{rm\ -r} do?
\item
  Why is \texttt{rm} potentially dangerous?
\item
  True or False: \texttt{cp} requires the \texttt{-r} flag to copy both files and folders.
\item
  Explain the difference between \texttt{rm} and \texttt{rmdir}.
\item
  If you delete a file with \texttt{rm}, can you recover it from PowerShell?
\item
  Write a command that would copy an entire folder and all its contents to a new location.
\item
  Describe a practical safety check you would perform before running \texttt{rm\ -r} on a folder.
\end{enumerate}

Caution

By default, \texttt{cp} \textbf{silently overwrites} an existing file at the destination with no warning. Use \texttt{cp\ -i} (interactive) to get a confirmation prompt before overwriting, especially when working with important files.

\begin{enumerate}
\setcounter{enumi}{10}
\tightlist
\item
  What happens if you \texttt{cp} a file to a destination where a file with the same name already exists? How would you handle this safely?
\item
  Compare \texttt{mv\ old\_name.txt\ new\_name.txt} vs \texttt{mv\ old\_name.txt\ \textasciitilde{}/Documents/new\_name.txt}. What is the key difference?
\item
  Design a workflow to safely delete 50 files matching the pattern \texttt{*.bak} from a folder containing 500 files. What commands and verifications would you use?
\item
  Explain how you could back up all \texttt{.scad} files from a project folder into a timestamped backup folder in one command.
\item
  When organizing a 3D printing project, you need to move completed designs to an archive folder and delete failed prototypes. How would you structure this as a safe, auditable process?
\end{enumerate}

\subsection*{Extension Problems}\label{docs__pandoc__latex__src__cli_complete_curriculum.md__extension-problems-2}

\begin{enumerate}
\tightlist
\item
  Create a folder tree and copy it to a new location with \texttt{cp\ -r}.
\item
  Write a one-line command that creates three files named \texttt{a\ b\ c} and lists them.
\item
  Move a file into a new folder and confirm the move.
\item
  Use wildcards to delete files matching a pattern in a safe test folder.
\item
  Export a listing of the practice folder to \texttt{practice\_listing.txt}.
\item
  Create a backup script that copies all \texttt{.scad} files from your project folder to a backup folder with timestamp naming.
\item
  Build a safe deletion workflow: list files matching a pattern, verify count, then delete with confirmation; document the steps.
\item
  Write a PowerShell script that organizes files by extension into subfolders; test it on a sample folder tree.
\item
  Create a file operation audit trail: log all copy, move, and delete operations to a text file for review.
\item
  Develop a project template generator: a script that creates a standard folder structure for a new 3D printing project with essential subfolders.
\item
  Implement a file conflict handler: write a script that handles cases where \texttt{cp} would overwrite an existing file by renaming the existing file with a timestamp before copying.
\item
  Create a batch rename operation: use a script to rename all files in a folder from \texttt{old\_prefix\_*} to \texttt{new\_prefix\_*}; test with actual files and verify the results.
\item
  Build a folder comparison tool: list all files in two folders and identify which files exist in one but not the other; output to a report.
\item
  Write a destructive operation validator: before executing \texttt{rm\ -r}, create a script that lists exactly what will be deleted, shows file counts by type, and requires explicit user confirmation to proceed.
\item
  Design a complete project lifecycle workflow: create folders for active projects, completed designs, and archive; include move operations between folders, backup steps, and verification that all files arrive intact.
\end{enumerate}

\subsection*{References}\label{docs__pandoc__latex__src__cli_complete_curriculum.md__references-2}

\begin{itemize}
\tightlist
\item
  Microsoft. (2024). \emph{New-Item cmdlet reference}. \url{https://learn.microsoft.com/powershell/module/microsoft.powershell.management/new-item}
\item
  Microsoft. (2024). \emph{Copy-Item and Move-Item cmdlets}. \url{https://learn.microsoft.com/powershell/module/microsoft.powershell.management/copy-item}
\item
  Microsoft. (2024). \emph{File system operations guide}. \url{https://learn.microsoft.com/powershell/scripting/learn/shell/manipulating-items}
\end{itemize}

\subsection*{Helpful Resources}\label{docs__pandoc__latex__src__cli_complete_curriculum.md__helpful-resources-2}

\begin{itemize}
\tightlist
\item
  \href{https://learn.microsoft.com/powershell/module/microsoft.powershell.management/new-item}{New-Item Cmdlet Reference}
\item
  \href{https://learn.microsoft.com/powershell/module/microsoft.powershell.management/copy-item}{Copy-Item and Move-Item}
\item
  \href{https://learn.microsoft.com/powershell/module/microsoft.powershell.management/remove-item}{Remove-Item Cmdlet Reference}
\item
  \href{https://learn.microsoft.com/powershell/scripting/learn/shell/manipulating-items}{File System Operations Guide}
\item
  \href{https://poshcode.gitbook.io/powershell-faq/src/getting-started/file-management}{Safe Deletion Practices}
\end{itemize}

\begin{center}\rule{0.5\linewidth}{0.5pt}\end{center}

\protect\phantomsection\label{docs__pandoc__latex__src__cli_complete_curriculum.md__ps-lesson-3}{}

\section*{Lesson 3: Input, Output, and Piping}\label{docs__pandoc__latex__src__cli_complete_curriculum.md__lesson-3-input-output-and-piping}

\begin{center}\rule{0.5\linewidth}{0.5pt}\end{center}

Duration: 1 class period
Prerequisite: Lesson 2 (File and Folder Manipulation)

\subsection*{Learning Objectives}\label{docs__pandoc__latex__src__cli_complete_curriculum.md__learning-objectives-1}

By the end of this lesson, you will be able to:

\begin{itemize}
\tightlist
\item
  Use \texttt{echo} to print text to the screen
\item
  Use \texttt{cat} to read file contents
\item
  Use \texttt{\textgreater{}} to redirect output into a file
\item
  Use \texttt{\textbar{}} (pipe) to send one command\textquotesingle s output to another
\item
  Copy output to the clipboard with \texttt{clip}
\item
  Open files with a text editor from the command line
\end{itemize}

\subsection*{Commands Covered}\label{docs__pandoc__latex__src__cli_complete_curriculum.md__commands-covered}

{\def\LTcaptype{none} % do not increment counter
\begin{longtable}[]{@{}
  >{\raggedright\arraybackslash}p{(\linewidth - 2\tabcolsep) * \real{0.3288}}
  >{\raggedright\arraybackslash}p{(\linewidth - 2\tabcolsep) * \real{0.6712}}@{}}
\toprule\noalign{}
\begin{minipage}[b]{\linewidth}\raggedright
Command
\end{minipage} & \begin{minipage}[b]{\linewidth}\raggedright
What It Does
\end{minipage} \\
\midrule\noalign{}
\endhead
\bottomrule\noalign{}
\endlastfoot
\texttt{echo\ "text"} & Print text to the screen \\
\texttt{cat\ filename} & Print the contents of a file \\
\texttt{\textgreater{}\ filename} &
Redirect output into a file (overwrites) \\
\texttt{\textgreater{}\textgreater{}\ filename} &
Append output to a file (adds to end) \\
\texttt{\textbar{}} & Pipe - send output from one command to the next \\
\texttt{clip} & Copy piped input to the Windows clipboard \\
\texttt{notepad.exe\ filename} & Open a file in Notepad \\
\end{longtable}
}

\subsection*{\texorpdfstring{\texttt{echo} - Printing Text}{echo - Printing Text}}\label{docs__pandoc__latex__src__cli_complete_curriculum.md__echo---printing-text}

\texttt{echo} prints text to the screen. It is useful for testing, for writing text into files, and for understanding how piping works.

\begin{lstlisting}[style=Alabaster, language=powershell]
echo "Hello, World"
echo "This is a test"

\end{lstlisting}

\subsection*{\texorpdfstring{\texttt{cat} - Reading Files}{cat - Reading Files}}\label{docs__pandoc__latex__src__cli_complete_curriculum.md__cat---reading-files}

\texttt{cat} prints the contents of a file to the screen.

\begin{lstlisting}[style=Alabaster, language=powershell]
# Read a text file
cat ~/Documents/notes.txt
# Read an OpenSCAD file
cat ~/Documents/OpenSCAD_Projects/project0.scad

\end{lstlisting}

With a long file, use \texttt{cat\ filename\ \textbar{}\ more} to read it page by page (press \texttt{Space} to advance, \texttt{Q} to quit).

\subsection*{\texorpdfstring{\texttt{\textgreater{}} - Redirecting Output to a File}{\textgreater{} - Redirecting Output to a File}}\label{docs__pandoc__latex__src__cli_complete_curriculum.md__---redirecting-output-to-a-file}

The \texttt{\textgreater{}} symbol redirects output from the screen into a file instead.

\begin{lstlisting}[style=Alabaster, language=powershell]
# Create a file with a single line
echo "Author: Your Name" > header.txt
# Confirm the file was created and has content
cat header.txt

\end{lstlisting}

Warning

\texttt{\textgreater{}} overwrites the file if it already exists. Use \texttt{\textgreater{}\textgreater{}} to append instead:

\begin{lstlisting}[style=Alabaster, language=powershell]
echo "Date: 2025" >> header.txt
echo "Project: Floor Marker" >> header.txt
cat header.txt

\end{lstlisting}

\subsection*{\texorpdfstring{\texttt{\textbar{}} - Piping}{\textbar{} - Piping}}\label{docs__pandoc__latex__src__cli_complete_curriculum.md__---piping}

The pipe symbol \texttt{\textbar{}} sends the output of one command to the input of the next. This lets you chain commands together.

\begin{lstlisting}[style=Alabaster, language=powershell]
# List files and send the list to clip (copies to clipboard)
ls -n | clip
## 3DMake in Practice: Capturing Build Output

Piping and redirection are essential for working with 3DMake's output accessibly:

```powershell
# Build and save the full log to a file so you can read it in Notepad
3dm build > build_log.txt
notepad.exe build_log.txt

# Build and slice, keeping only lines with errors or warnings
3dm build slice 2>&1 | Select-String -Pattern "error|warning" -CaseSensitive:$false

# List all STL files produced so far
ls -n build/ | findstr ".stl"

# Get a list of available printer profiles and save for reference
3dm list-profiles > profiles.txt
notepad.exe profiles.txt

# Get a list of material overlays (PLA, PETG, supports, etc.)
3dm list-overlays > overlays.txt
notepad.exe overlays.txt

# Get an AI description of your model and save it
3dm info > model_description.txt
notepad.exe model_description.txt

\end{lstlisting}

\begin{lstlisting}[style=Alabaster, language=powershell]
# List files and send the list to clip (copies to clipboard)
ls -n | clip
# Now paste with Ctrl + V anywhere

\end{lstlisting}

\begin{lstlisting}[style=Alabaster, language=powershell]
# Search within a file's contents using Select-String (like grep)
cat project0.scad | Select-String "cube"

\end{lstlisting}

\subsection*{\texorpdfstring{\texttt{clip} - Copying to Clipboard}{clip - Copying to Clipboard}}\label{docs__pandoc__latex__src__cli_complete_curriculum.md__clip---copying-to-clipboard}

\texttt{clip} takes whatever is piped to it and puts it on the Windows clipboard.

\begin{lstlisting}[style=Alabaster, language=powershell]
# Copy your current directory path to the clipboard
pwd | clip
# Copy a file listing to clipboard
ls -n | clip
# Copy the contents of a file to clipboard
cat notes.txt | clip

\end{lstlisting}

After any of these, press \texttt{Ctrl\ +\ V} in any application to paste.

\subsection*{Opening Files in Notepad}\label{docs__pandoc__latex__src__cli_complete_curriculum.md__opening-files-in-notepad}

\begin{lstlisting}[style=Alabaster, language=powershell]
# Open a file in Notepad
notepad.exe ~/Documents/notes.txt
# Open a .scad file
notepad.exe ~/Documents/OpenSCAD_Projects/project0.scad
# Create a new file and open it
notepad new_notes.txt

\end{lstlisting}

\subsection*{Step-by-step Tasks}\label{docs__pandoc__latex__src__cli_complete_curriculum.md__step-by-step-tasks-1}

\begin{enumerate}
\tightlist
\item
  Create \texttt{practice.txt} with three lines using \texttt{echo} and \texttt{\textgreater{}}/\texttt{\textgreater{}\textgreater{}}.
\item
  Read the file with \texttt{cat\ practice.txt}.
\item
  Pipe the file into \texttt{Select-String} to search for a word.
\item
  Copy the file contents to clipboard with \texttt{cat\ practice.txt\ \textbar{}\ clip}.
\item
  Redirect \texttt{ls\ -n} into \texttt{list.txt} and open it in Notepad.
  Checkpoints
\end{enumerate}

\begin{itemize}
\tightlist
\item
  After step 3 you should be able to find a keyword using piping.
\end{itemize}

\subsection*{Quiz - Lesson PS.3}\label{docs__pandoc__latex__src__cli_complete_curriculum.md__quiz---lesson-ps3}

\begin{enumerate}
\tightlist
\item
  What is the difference between \texttt{\textgreater{}} and \texttt{\textgreater{}\textgreater{}}?
\item
  What does the pipe \texttt{\textbar{}} do?
\item
  How do you copy output to the clipboard?
\item
  How would you page through long output?
\item
  How do you suppress output to nowhere?
\item
  True or False: The pipe operator \texttt{\textbar{}} connects the output of one command to the input of another.
\item
  Explain why redirecting output to a file is useful for screen reader users.
\item
  Write a command that would search for the word "sphere" in all \texttt{.scad} files in a directory.
\item
  How would you count the number of lines in a file using PowerShell piping?
\item
  Describe a practical scenario in 3D printing where you would pipe or redirect command output.
\item
  What would be the difference in output between \texttt{echo\ "test"\ \textgreater{}\ file.txt} (run twice) vs \texttt{echo\ "test"\ \textgreater{}\textgreater{}\ file.txt} (run twice)? Show the expected file contents.
\item
  Design a three-step piping chain: read a file, filter for specific content, and save the results; explain what each pipe does.
\item
  You have a 500-line \texttt{.scad} file and need to find all instances of \texttt{sphere()} and count them. Write the command.
\item
  Explain how \texttt{clip} is particularly valuable for screen reader users when working with file paths or long output strings.
\item
  Describe how you would use pipes and redirection to create a timestamped backup report of all \texttt{.stl} files in a 3D printing project folder.
\end{enumerate}

\subsection*{Extension Problems}\label{docs__pandoc__latex__src__cli_complete_curriculum.md__extension-problems-3}

\begin{enumerate}
\tightlist
\item
  Use piping to count lines in a file (hint: \texttt{Select-String\ -Pattern\ \textquotesingle{}.\textquotesingle{}\ \textbar{}\ Measure-Object}).
\item
  Save a long \texttt{ls\ -n} output and search it with \texttt{Select-String}.
\item
  Chain multiple pipes to filter and then save results.
\item
  Practice copying different command outputs to clipboard and pasting.
\item
  Create a small script that generates a report (counts of files by extension).
\item
  Build a data processing pipeline: read a CSV file, filter rows, and export results; document each step.
\item
  Write a script that pipes directory listing to Count occurrences of each file extension; create a summary report.
\item
  Create a log analysis command: read a log file, filter for errors, and save matching lines to a separate error log.
\item
  Design a piping workflow for 3D printing file management: find \texttt{.stl} files, extract their sizes, and generate a report.
\item
  Develop a reusable piping function library: create functions for common filtering, sorting, and reporting patterns; test each function with different inputs.
\item
  Build a complex filter pipeline: read a \texttt{.scad} file, extract lines containing specific geometry commands, count each type, and output a summary to both screen and file.
\item
  Create an interactive piping tool: build a script that accepts user input for a search term, pipes through multiple filters, and displays paginated results.
\item
  Develop a performance analysis tool: use piping to combine file listing, metadata extraction, and statistical reporting; export results to a dated report file.
\item
  Implement a comprehensive error-handling pipeline: read output, catch errors, log them separately, and generate a summary of successes vs failures.
\item
  Design and execute a real-world project backup workflow: use piping to verify file integrity, count files by type, generate a backup manifest, and create audit logs-all in one integrated command pipeline.
\end{enumerate}

\subsection*{References}\label{docs__pandoc__latex__src__cli_complete_curriculum.md__references-3}

\begin{itemize}
\tightlist
\item
  Microsoft. (2024). \emph{Out-File cmdlet for redirection}. \url{https://learn.microsoft.com/en-us/powershell/module/microsoft.powershell.utility/out-file}
\item
  Microsoft. (2024). \emph{Select-String piping reference}. \url{https://learn.microsoft.com/en-us/powershell/module/microsoft.powershell.utility/select-string}
\item
  Microsoft. (2024). \emph{PowerShell pipeline concepts}. \url{https://learn.microsoft.com/en-us/powershell/scripting/learn/deep-dives/everything-about-pipeline}
\end{itemize}

\subsection*{Helpful Resources}\label{docs__pandoc__latex__src__cli_complete_curriculum.md__helpful-resources-3}

\begin{itemize}
\tightlist
\item
  \href{https://learn.microsoft.com/powershell/module/microsoft.powershell.utility/out-file}{Using Out-File for Redirection}
\item
  \href{https://learn.microsoft.com/powershell/module/microsoft.powershell.utility/select-string}{Piping and Select-String}
\item
  \href{https://learn.microsoft.com/powershell/module/microsoft.powershell.management/get-content}{Get-Content Cmdlet Reference}
\item
  \href{https://learn.microsoft.com/powershell/module/microsoft.powershell.utility/measure-object}{Measure-Object for Counting}
\item
  \href{https://learn.microsoft.com/powershell/scripting/learn/shell/using-the-pipeline}{PowerShell Pipeline Concept}
\end{itemize}

\begin{center}\rule{0.5\linewidth}{0.5pt}\end{center}

\protect\phantomsection\label{docs__pandoc__latex__src__cli_complete_curriculum.md__ps-lesson-4}{}

\section*{Lesson 4: Environment Variables, PATH, and Aliases}\label{docs__pandoc__latex__src__cli_complete_curriculum.md__lesson-4-environment-variables-path-and-aliases}

\begin{center}\rule{0.5\linewidth}{0.5pt}\end{center}

Estimated time: 30-45 minutes

Learning Objectives

\begin{itemize}
\tightlist
\item
  Read environment variables with \texttt{\$env:VARNAME}
\item
  Inspect and verify programs in the \texttt{PATH}
\item
  Create temporary aliases and understand making them persistent via the profile
\end{itemize}

Materials

\begin{itemize}
\tightlist
\item
  PowerShell (with rights to open profile if desired)
\end{itemize}

Step-by-step Tasks

\begin{enumerate}
\tightlist
\item
  Show your username and home path with \texttt{echo\ \$env:USERNAME} and \texttt{echo\ \$env:USERPROFILE}.
\item
  Inspect \texttt{echo\ \$env:PATH} and identify whether \texttt{openscad} or \texttt{code} would be found.
\item
  Run \texttt{Get-Command\ openscad} and note the result.
\item
  Create a temporary alias: \texttt{Set-Alias\ -Name\ preview\ -Value\ openscad} and test \texttt{preview\ myfile.scad}.
\item
  Open your profile (\texttt{notepad.exe\ \$PROFILE}) and add the alias line to make it persistent (advanced).
\end{enumerate}

Checkpoints

\begin{itemize}
\tightlist
\item
  After step 3 you can determine whether a program will be found by \texttt{PATH}.
\end{itemize}

\subsection*{Quiz - Lesson PS.4}\label{docs__pandoc__latex__src__cli_complete_curriculum.md__quiz---lesson-ps4}

\subsection*{3DMake in Practice: Aliases and Profile Setup}\label{docs__pandoc__latex__src__cli_complete_curriculum.md__3dmake-in-practice-aliases-and-profile-setup}

PowerShell aliases and profile scripts make your 3DMake workflow much faster:

\begin{lstlisting}[style=Alabaster, language=powershell]
# Quick function to jump to your 3D projects folder
function 3dp { Set-Location ~/Documents/3d_projects }

# Quick function to run the full build-slice-print pipeline
function 3dprint { 3dm build slice print }

# Quick function to build and get an AI model description
function 3dinfo { 3dm build; 3dm info }

# Add these functions to your PowerShell profile so they load every session
notepad.exe $PROFILE
# (Paste the three functions above into the file, then save)

# Reload your profile without restarting PowerShell
. $PROFILE

# Now use them:
3dp            # jump to your 3D projects folder
3dprint        # run the full pipeline
3dinfo         # build and describe the model

\end{lstlisting}

Checkpoints

\begin{itemize}
\tightlist
\item
  After step 3 you can determine whether a program will be found by \texttt{PATH}.
\end{itemize}

\subsection*{Quiz - Lesson PS.4}\label{docs__pandoc__latex__src__cli_complete_curriculum.md__quiz---lesson-ps4-1}

\begin{enumerate}
\tightlist
\item
  How do you print an environment variable?
\item
  What is the purpose of \texttt{PATH}?
\item
  How do you check whether \texttt{openscad} is available?
\item
  How do you create a temporary alias?
\item
  Where would you make an alias permanent?
\item
  True or False: Environment variables are case-sensitive on all platforms.
\item
  Explain why having a program in your PATH is useful compared to always using its full file path.
\item
  Write a command that would create an alias called \texttt{slicer} for the OpenSCAD executable.
\item
  What file would you edit to make an alias persist across PowerShell sessions?
\item
  Describe a practical benefit of using the \texttt{\$env:TEMP} directory for temporary files in a 3D printing workflow.
\item
  You have a custom script at \texttt{C:\textbackslash{}Scripts\textbackslash{}backup\_models.ps1} that you want to run from anywhere as \texttt{backup-now}. What steps would you take to make this work?
\item
  Explain the difference between setting an environment variable in the current session vs. adding it to your profile for permanence.
\item
  Design a profile strategy for managing multiple 3D printing projects, each with different tool paths and directories; show how to structure environment variables for each.
\item
  If a program is not found by \texttt{Get-Command}, what are the possible reasons, and how would you troubleshoot?
\item
  Describe how you would verify that your PowerShell profile is loading correctly and how to debug issues if aliases or environment variables don\textquotesingle t appear after restarting PowerShell.
\end{enumerate}

\subsection*{Extension Problems}\label{docs__pandoc__latex__src__cli_complete_curriculum.md__extension-problems-4}

\begin{enumerate}
\tightlist
\item
  Add a folder to PATH for a test program (describe steps; do not change system PATH without admin).
\item
  Create a short profile snippet that sets two aliases and test re-opening PowerShell.
\item
  Use \texttt{Get-Command} to list the path for several common programs.
\item
  Explore \texttt{\$env:TEMP} and create a file there.
\item
  Save a copy of your current PATH to a text file and examine it in your editor.
\item
  Create a PowerShell profile script that loads custom aliases and environment variables for your 3D printing workflow; test it in a new session.
\item
  Build a "project profile" that sets environment variables for CAD, slicing, and print directories; switch between profiles for different projects.
\item
  Write a script that audits your current environment variables and creates a summary report of what\textquotesingle s set and why.
\item
  Design a custom alias system for common 3D printing commands; document the aliases and their purposes.
\item
  Create a profile migration guide: document how to export and import your PowerShell profile across machines for consistent workflows.
\item
  Implement a safe PATH modification script: create a utility that allows you to add/remove directories from PATH for the current session only; show how to make it permanent in your profile.
\item
  Build a comprehensive profile framework with modules: create separate .ps1 files for aliases, environment variables, and functions; have your main profile load all of them dynamically.
\item
  Develop an environment validation tool: write a script that checks whether all required programs (OpenSCAD, slicers, etc.) are accessible via PATH; report findings and suggest fixes.
\item
  Create a project-switching alias system: design a function that changes all environment variables and aliases based on the current project; test switching between multiple projects.
\item
  Build a profile troubleshooting guide: create a script that exports your current environment state (variables, aliases, PATH) to a timestamped file, allowing you to compare states before and after changes and identify what broke.
\end{enumerate}

\subsection*{References}\label{docs__pandoc__latex__src__cli_complete_curriculum.md__references-4}

\begin{itemize}
\tightlist
\item
  Microsoft. (2024). \emph{Environment variables in PowerShell}. \url{https://learn.microsoft.com/powershell/scripting/learn/shell/using-environment-variables}
\item
  Microsoft. (2024). \emph{Set-Alias cmdlet reference}. \url{https://learn.microsoft.com/powershell/module/microsoft.powershell.utility/set-alias}
\item
  Microsoft. (2024). \emph{Creating and using PowerShell profiles}. \url{https://learn.microsoft.com/powershell/module/microsoft.powershell.core/about/about_profiles}
\end{itemize}

\subsection*{Helpful Resources}\label{docs__pandoc__latex__src__cli_complete_curriculum.md__helpful-resources-4}

\begin{itemize}
\tightlist
\item
  \href{https://learn.microsoft.com/powershell/scripting/learn/shell/using-environment-variables}{Environment Variables in PowerShell}
\item
  \href{https://poshcode.gitbook.io/powershell-faq/src/getting-started/environment-variables}{Understanding the PATH Variable}
\item
  \href{https://learn.microsoft.com/powershell/module/microsoft.powershell.utility/set-alias}{Set-Alias Cmdlet Reference}
\item
  \href{https://learn.microsoft.com/powershell/module/microsoft.powershell.core/about/about_profiles}{Creating a PowerShell Profile}
\item
  \href{https://learn.microsoft.com/powershell/module/microsoft.powershell.core/get-command}{Get-Command for Locating Programs}
\end{itemize}

\begin{center}\rule{0.5\linewidth}{0.5pt}\end{center}

\protect\phantomsection\label{docs__pandoc__latex__src__cli_complete_curriculum.md__ps-lesson-5}{}

\section*{Lesson 5: Filling in the Gaps --- Profiles and History}\label{docs__pandoc__latex__src__cli_complete_curriculum.md__lesson-5-filling-in-the-gaps--profiles-and-history}

\begin{center}\rule{0.5\linewidth}{0.5pt}\end{center}

Estimated time: 30-45 minutes

Learning Objectives

\begin{itemize}
\tightlist
\item
  Use history and abort commands (\texttt{history}, \texttt{Ctrl+C})
\item
  Inspect and edit your PowerShell profile for persistent settings
\item
  Run programs by full path using the \texttt{\&} operator
\item
  Handle common screen reader edge cases when using the terminal
\end{itemize}

Materials

\begin{itemize}
\tightlist
\item
  PowerShell and an editor (Notepad/ VS Code)
\end{itemize}

Step-by-step Tasks

\begin{enumerate}
\tightlist
\item
  Run several simple commands (e.g., \texttt{pwd}, \texttt{ls\ -n}, \texttt{echo\ hi}) then run \texttt{history} to view them.
\item
  Use \texttt{Invoke-History\ \textless{}n\textgreater{}} to re-run a previous command (replace \texttt{\textless{}n\textgreater{}} with a history number).
\item
  Practice aborting a long-running command with \texttt{Ctrl\ +\ C} (for example, \texttt{ping\ 8.8.8.8}).
\item
  Open your profile: \texttt{notepad.exe\ \$PROFILE}; if it doesn\textquotesingle t exist, create it: \texttt{ni\ \$PROFILE\ -Force}.
\item
  Add a persistent alias line to your profile (example: \texttt{Set-Alias\ -Name\ preview\ -Value\ openscad}), save, and reopen PowerShell to verify.
\end{enumerate}

Checkpoints

\begin{itemize}
\tightlist
\item
  After step 2 you can re-run a recent command by history number.
\item
  After step 5 your alias should persist across sessions.
\end{itemize}

\subsection*{Quiz - Lesson PS.5}\label{docs__pandoc__latex__src__cli_complete_curriculum.md__quiz---lesson-ps5}

\begin{enumerate}
\tightlist
\item
  How do you view the command history?
\item
  Which key combination aborts a running command?
\item
  What does \texttt{echo\ \$PROFILE} show?
\item
  How does the \texttt{\&} operator help run executables?
\item
  What is one strategy if terminal output stops being announced by your screen reader?
\item
  True or False: Using \texttt{Ctrl+C} permanently deletes any files created by the command you abort.
\item
  Explain the difference between \texttt{history} and \texttt{Get-History} in PowerShell.
\item
  If you place code in your profile but it doesn\textquotesingle t take effect after opening a new PowerShell window, what should you verify?
\item
  Write a command that would run a program at the path \texttt{C:\textbackslash{}Program\ Files\textbackslash{}OpenSCAD\textbackslash{}openscad.exe} directly.
\item
  Describe a practical workflow scenario where having keyboard shortcuts (aliases) in your profile would save time.
\item
  Explain how to re-run the 5th command from your history, and what would happen if that command had file operations (creates/deletes).
\item
  Design a profile initialization strategy that separates utilities for different projects; explain how you would switch between them.
\item
  Walk through a troubleshooting workflow: your screen reader stops announcing output after running a long command. What steps would you take to diagnose and resolve the issue?
\end{enumerate}

Caution

Before any destructive operation --- mass delete, bulk overwrite, or recursive removal --- use history review (\texttt{Get-History}) and dry-run techniques (e.g., \texttt{What-If}) to verify the command is correct. A single mistaken path in a bulk delete can be unrecoverable.

\begin{enumerate}
\setcounter{enumi}{13}
\tightlist
\item
  Create a safety checkpoint system: before any destructive operation (mass delete, overwrite), how would you use profile functions and history to verify the command is correct?
\item
  Develop a comprehensive capstone scenario: integrate everything from Lessons 0--5 (navigation, file operations, piping, environment setup, history) to design an automated 3D printing project workflow with error handling and logging.
\end{enumerate}

\subsection*{Extension Problems}\label{docs__pandoc__latex__src__cli_complete_curriculum.md__extension-problems-5}

\begin{enumerate}
\tightlist
\item
  Add an alias and an environment variable change to your profile and document the behavior after reopening PowerShell.
\item
  Create a short script that automates creating a project folder and an initial .scad file.
\item
  Experiment with running OpenSCAD by full path using \texttt{\&} and by placing it in PATH; compare results.
\item
  Practice redirecting \texttt{Get-Help} output to a file and reading it in Notepad for screen reader clarity.
\item
  Document three screen reader troubleshooting steps you used and when they helped.
\item
  Build a comprehensive PowerShell profile that includes aliases, environment variables, and helper functions for your 3D printing workflow.
\item
  Create a script that troubleshoots common PowerShell issues (module loading, permission errors, command not found); test at least three scenarios.
\item
  Write a PowerShell function that coordinates multiple tasks: creates a project folder, starts OpenSCAD, and opens slicing software.
\item
  Design a screen-reader accessibility guide for PowerShell: document commands, outputs, and accessible navigation patterns.
\item
  Develop an advanced PowerShell workflow: implement error handling, logging, and confirmation prompts for risky operations.
\item
  Implement a "undo" system using history: create a function that logs destructive commands (rm, mv, cp -Force) and allows you to review/rollback the last operation.
\item
  Build a profile debugger: create a script that compares two PowerShell sessions\textquotesingle{} environment states (variables, aliases, functions) to identify what loaded/failed to load.
\item
  Develop a multi-project profile manager: design a system where you can switch entire environments (paths, aliases, variables) for different 3D printing projects by running a single command.
\item
  Create a comprehensive accessibility analyzer: write a script that tests whether key PowerShell commands produce screen-reader-friendly output; document workarounds for commands that don\textquotesingle t.
\item
  Design a complete capstone project: build an integrated automation suite that manages a 3D printing workflow (project setup, file organization, CAD/slicing tool automation, output logging, error recovery, and audit trails) with full error handling and documentation.
\end{enumerate}

\subsection*{References}\label{docs__pandoc__latex__src__cli_complete_curriculum.md__references-5}

\begin{itemize}
\tightlist
\item
  Microsoft. (2024). \emph{PowerShell history and recall functionality}. \url{https://learn.microsoft.com/powershell/scripting/learn/shell/using-history}
\item
  Microsoft. (2024). \emph{Understanding and creating PowerShell profiles}. \url{https://learn.microsoft.com/powershell/module/microsoft.powershell.core/about/about_profiles}
\item
  Microsoft. (2024). \emph{The call operator (\&) for running executables}. \url{https://learn.microsoft.com/powershell/module/microsoft.powershell.core/about/about_operators\#call-operator-}
\end{itemize}

\subsection*{Helpful Resources}\label{docs__pandoc__latex__src__cli_complete_curriculum.md__helpful-resources-5}

\begin{itemize}
\tightlist
\item
  \href{https://learn.microsoft.com/powershell/scripting/learn/shell/using-history}{PowerShell History and Recall}
\item
  \href{https://learn.microsoft.com/powershell/module/microsoft.powershell.core/about/about_profiles}{Understanding Profiles}
\item
  \href{https://learn.microsoft.com/powershell/module/microsoft.powershell.core/invoke-history}{Invoke-History Cmdlet Reference}
\item
  \href{https://learn.microsoft.com/powershell/module/microsoft.powershell.core/about/about_operators\#call-operator-}{The Call Operator (\&)}
\item
  \href{https://learn.microsoft.com/powershell/scripting/windows-powershell/ise/accessibility-in-windows-powershell-ise}{Screen Reader Tips and Tricks}
\end{itemize}

\begin{center}\rule{0.5\linewidth}{0.5pt}\end{center}

\protect\phantomsection\label{docs__pandoc__latex__src__cli_complete_curriculum.md__ps-lesson-6}{}

\section*{Lesson 6: Advanced Techniques --- Scripts, Functions, and Professional Workflows}\label{docs__pandoc__latex__src__cli_complete_curriculum.md__lesson-6-advanced-techniques--scripts-functions-and-professional-workflows}

\begin{center}\rule{0.5\linewidth}{0.5pt}\end{center}

Duration: 4-4.5 hours (for screen reader users)\\
Prerequisites: Complete Lessons 0--5\\
Skill Level: Advanced intermediate

\subsection*{3DMake in Practice: Automating Your Print Workflow}\label{docs__pandoc__latex__src__cli_complete_curriculum.md__3dmake-in-practice-automating-your-print-workflow}

Use PowerShell scripting to fully automate a 3DMake workflow across multiple projects:

\begin{lstlisting}[style=Alabaster, language=powershell]
# batch_print.ps1 - build and print every 3DMake project in a folder
param(
    [string]$ProjectsDir = "$HOME\Documents\3d_projects"
)

Get-ChildItem -Path $ProjectsDir -Directory | ForEach-Object {
    $project = $_.FullName
    $name = $_.Name
    Write-Host "Processing: $name"
    Set-Location $project

    # Build and slice; capture all output
    $output = & 3dm build slice 2>&1
    if ($LASTEXITCODE -eq 0) {
        Write-Host "  OK - sending to printer"
        & 3dm print
    } else {
        Write-Host "  FAILED - log saved to build_error.txt"
        $output | Out-File "build_error.txt"
        notepad.exe "build_error.txt"
    }
}

Write-Host "Batch complete."

\end{lstlisting}

Save as \texttt{batch\_print.ps1} and run with \texttt{.\textbackslash{}batch\_print.ps1}. This script uses loops, conditionals, and piping --- all from this lesson.

This lesson extends PowerShell skills to professional-level workflows. You\textquotesingle ll learn to automate complex tasks, write reusable code, and integrate tools for 3D printing workflows.

\subsection*{Learning Objectives}\label{docs__pandoc__latex__src__cli_complete_curriculum.md__learning-objectives-2}

By the end of this lesson, you will be able to:

\begin{itemize}
\tightlist
\item
  Create and run PowerShell scripts (.ps1 files)
\item
  Write functions that accept parameters
\item
  Use loops to repeat tasks automatically
\item
  Automate batch processing of 3D models
\item
  Debug scripts when something goes wrong
\item
  Create professional workflows combining multiple tools
\end{itemize}

\subsection*{PowerShell Scripts Basics}\label{docs__pandoc__latex__src__cli_complete_curriculum.md__powershell-scripts-basics}

\subsubsection*{What\textquotesingle s a Script?}\label{docs__pandoc__latex__src__cli_complete_curriculum.md__whats-a-script}

A script is a file containing PowerShell commands that run in sequence. Instead of typing commands one by one, you put them in a file and run them all at once.

Why use scripts?

\begin{itemize}
\tightlist
\item
  Repeatability: Run the same task 100 times identically
\item
  Documentation: Commands are written down for reference
\item
  Complexity: Combine many commands logically
\item
  Automation: Schedule scripts to run automatically
\end{itemize}

\subsubsection*{Creating Your First Script}\label{docs__pandoc__latex__src__cli_complete_curriculum.md__creating-your-first-script}

Step 1: Open a text editor

\begin{lstlisting}[style=Alabaster, language=powershell]
notepad.exe my-first-script.ps1

\end{lstlisting}

Step 2: Type this script

\begin{lstlisting}[style=Alabaster, language=powershell]
# This is a comment - screen readers will read it
Write-Output "Script is running!"
pwd
ls -n
Write-Output "Script is done!"

\end{lstlisting}

Step 3: Save the file

\begin{itemize}
\tightlist
\item
  In Notepad: Ctrl+S
\item
  Make sure filename ends in \texttt{.ps1}
\item
  Save in an easy-to-find location (like Documents)
\end{itemize}

Step 4: Run the script

\begin{lstlisting}[style=Alabaster, language=powershell]
.\my-first-script.ps1

\end{lstlisting}

What happens: PowerShell runs each command in sequence and shows output.

\subsubsection*{Important: Script Execution Policy}\label{docs__pandoc__latex__src__cli_complete_curriculum.md__important-script-execution-policy}

Caution

If you get an "execution policy" error when running a \texttt{.ps1} script, PowerShell\textquotesingle s security is blocking unsigned scripts. Use \texttt{-Scope\ Process} to bypass it only for the current session --- this is safer than changing the policy machine-wide.

On some Windows systems, you might get an error about "execution policy". This is a security feature.

To fix it temporarily:

\begin{lstlisting}[style=Alabaster, language=powershell]
Set-ExecutionPolicy -ExecutionPolicy Bypass -Scope Process

\end{lstlisting}

Then try your script again:

\begin{lstlisting}[style=Alabaster, language=powershell]
.\my-first-script.ps1

\end{lstlisting}

Note for screen readers: Your screen reader will announce the error if there is one. Use \texttt{Get-Help\ Get-ExecutionPolicy} for more information.

\subsection*{Variables and Parameters}\label{docs__pandoc__latex__src__cli_complete_curriculum.md__variables-and-parameters}

\subsubsection*{Using Variables}\label{docs__pandoc__latex__src__cli_complete_curriculum.md__using-variables}

Variables store values you want to use later.

Example script:

\begin{lstlisting}[style=Alabaster, language=powershell]
$mypath = "C:\Users\YourName\Documents"
cd $mypath
Write-Output "I am now in:"
pwd
ls -n

\end{lstlisting}

Breaking it down:

\begin{itemize}
\tightlist
\item
  \texttt{\$mypath} = variable name (always starts with \texttt{\$})
\item
  \texttt{=} = assign the value after this
\item
  \texttt{"C:\textbackslash{}Users..."} = the value (a path)
\item
  \texttt{cd\ \$mypath} = use the variable (replace \texttt{\$mypath} with its value)
\end{itemize}

\subsubsection*{Functions with Parameters}\label{docs__pandoc__latex__src__cli_complete_curriculum.md__functions-with-parameters}

A function is reusable code that you can run with different inputs.

Example: A function that lists files in a folder

\begin{lstlisting}[style=Alabaster, language=powershell]
function ListFolder {
    param(
        [string]$path
    )
    Write-Output "Contents of: $path"
    cd $path
    ls -n
}
# Use the function:
ListFolder -path "C:\Users\YourName\Documents"
ListFolder -path "C:\Users\YourName\Downloads"

\end{lstlisting}

What\textquotesingle s happening:

\begin{itemize}
\tightlist
\item
  \texttt{function\ ListFolder} = name of the function
\item
  \texttt{param({[}string{]}\$path)} = the function accepts a parameter called \texttt{\$path}
\item
  Inside the function, use \texttt{\$path} like any variable
\item
  Call the function with \texttt{-path\ "value"}
\end{itemize}

Screen reader tip: When you call a function, PowerShell will announce the results just like any command.

\subsection*{Loops - Repeating Tasks}\label{docs__pandoc__latex__src__cli_complete_curriculum.md__loops---repeating-tasks}

\subsubsection*{Loop Over Files}\label{docs__pandoc__latex__src__cli_complete_curriculum.md__loop-over-files}

Imagine you have 10 SCAD files and want to print their contents. You could do it 10 times manually, or use a loop.

Example: Print every .scad file in a folder

\begin{lstlisting}[style=Alabaster, language=powershell]
$scadFiles = ls -n *.scad
foreach ($file in $scadFiles) {
    Write-Output "=== File: $file ==="
    cat $file
    Write-Output ""
}

\end{lstlisting}

What\textquotesingle s happening:

\begin{itemize}
\tightlist
\item
  \texttt{\$scadFiles\ =\ ls\ -n\ *.scad} = find all .scad files and store in variable
\item
  \texttt{foreach\ (\$file\ in\ \$scadFiles)} = for each file, do this:

  \begin{itemize}
  \tightlist
  \item
    \texttt{Write-Output\ "===\ File:\ \$file\ ==="} = announce which file
  \item
    \texttt{cat\ \$file} = show contents
  \item
    \texttt{Write-Output\ ""} = blank line between files
  \end{itemize}
\end{itemize}

Result: All files printed one after another, organized and readable.

\subsubsection*{Loop with a Counter}\label{docs__pandoc__latex__src__cli_complete_curriculum.md__loop-with-a-counter}

Example: Do something 5 times

\begin{lstlisting}[style=Alabaster, language=powershell]
for ($i = 1; $i -le 5; $i++) {
    Write-Output "This is iteration number $i"
    # Do something here
}

\end{lstlisting}

What\textquotesingle s happening:

\begin{itemize}
\tightlist
\item
  \texttt{for\ (\$i\ =\ 1;\ \$i\ -le\ 5;\ \$i++)} = loop from 1 to 5
\item
  \texttt{\$i} = counter variable (starts at 1, increases each loop)
\item
  \texttt{-le} = "less than or equal to" (stop when \$i \textgreater{} 5)
\item
  \texttt{\$i++} = add 1 to \$i each time through
\end{itemize}

\subsection*{Real-World Example - Batch Processing SCAD Files}\label{docs__pandoc__latex__src__cli_complete_curriculum.md__real-world-example---batch-processing-scad-files}

\subsubsection*{Scenario}\label{docs__pandoc__latex__src__cli_complete_curriculum.md__scenario}

You have 10 OpenSCAD (.scad) files in a folder. You want to:

\begin{enumerate}
\tightlist
\item
  List them all
\item
  Check how many there are
\item
  For each one, verify it exists
\end{enumerate}

\subsubsection*{The Script}\label{docs__pandoc__latex__src__cli_complete_curriculum.md__the-script}

\begin{lstlisting}[style=Alabaster, language=powershell]
# Batch Processing Script for SCAD Files
$scadFolder = "C:\Users\YourName\Documents\3D_Projects"
$scadFiles = ls $scadFolder -Filter *.scad -Name
Write-Output "Processing SCAD files in: $scadFolder"
Write-Output "Found $($scadFiles.Count) files"
Write-Output ""
foreach ($file in $scadFiles) {
    $fullPath = Join-Path -Path $scadFolder -ChildPath $file
    if (Test-Path -Path $fullPath) {
        Write-Output " Found: $file"
    } else {
        Write-Output " Missing: $file"
    }
}
Write-Output ""
Write-Output "Batch processing complete!"

\end{lstlisting}

Breaking it down:

\begin{itemize}
\tightlist
\item
  \texttt{\$scadFolder} = where to look
\item
  \texttt{ls\ \$scadFolder\ -Filter\ *.scad\ -Name} = find .scad files, show names only
\item
  \texttt{foreach} = process each file
\item
  \texttt{Join-Path} = combine folder and filename into full path
\item
  \texttt{Test-Path} = check if file exists
\item
  \texttt{if} = do different things based on condition
\end{itemize}

\subsubsection*{Running the Script}\label{docs__pandoc__latex__src__cli_complete_curriculum.md__running-the-script}

\begin{enumerate}
\tightlist
\item
  Save as \texttt{batch-process.ps1}
\item
  Edit \texttt{\$scadFolder} to match your real folder
\item
  Run it:

  \begin{lstlisting}[style=Alabaster, language=powershell]
  .\batch-process.ps1

  \end{lstlisting}
\end{enumerate}

Screen reader output:

\begin{lstlisting}[style=Alabaster]
Processing SCAD files in: C:\Users\YourName\Documents\3D_Projects
Found 10 files

 Found: model1.scad
 Found: model2.scad
 Found: model3.scad
[... more files ...]
Batch processing complete!

\end{lstlisting}

\subsection*{Error Handling}\label{docs__pandoc__latex__src__cli_complete_curriculum.md__error-handling}

\subsubsection*{Try-Catch Blocks}\label{docs__pandoc__latex__src__cli_complete_curriculum.md__try-catch-blocks}

What if something goes wrong? Use try-catch:

Example:

\begin{lstlisting}[style=Alabaster, language=powershell]
try {
    $file = "C:\nonexistent\path\file.txt"
    $content = cat $file
    Write-Output $content
} catch {
    Write-Output "Error: Could not read file"
    Write-Output "Details: $_"
}

\end{lstlisting}

What\textquotesingle s happening:

\begin{itemize}
\tightlist
\item
  \texttt{try} = run these commands
\item
  If an error happens, PowerShell jumps to \texttt{catch}
\item
  \texttt{catch} = handle the error gracefully
\item
  \texttt{\$\_} = the error message
\end{itemize}

Screen reader advantage: Errors are announced clearly instead of crashing silently.

\subsubsection*{Validating Input}\label{docs__pandoc__latex__src__cli_complete_curriculum.md__validating-input}

Example: Make sure a folder exists before processing

\begin{lstlisting}[style=Alabaster, language=powershell]
function ProcessFolder {
    param([string]$folderPath)
    if (-not (Test-Path -Path $folderPath)) {
        Write-Output "Error: Folder does not exist: $folderPath"
        return
    }
    Write-Output "Processing folder: $folderPath"
    ls -n $folderPath
}
ProcessFolder -folderPath "C:\Users\YourName\Documents"

\end{lstlisting}

What\textquotesingle s happening:

\begin{itemize}
\tightlist
\item
  \texttt{Test-Path} = check if folder exists
\item
  \texttt{-not} = if NOT true
\item
  \texttt{return} = exit the function early if error
\end{itemize}

\subsection*{Debugging Scripts}\label{docs__pandoc__latex__src__cli_complete_curriculum.md__debugging-scripts}

\subsubsection*{Common Errors and Solutions}\label{docs__pandoc__latex__src__cli_complete_curriculum.md__common-errors-and-solutions}

\paragraph*{Error 1: "Command not found"}\label{docs__pandoc__latex__src__cli_complete_curriculum.md__error-1-command-not-found}

Cause: Typo in command name

Fix: Check spelling

\begin{lstlisting}[style=Alabaster, language=powershell]
# Wrong:
writ-output "hello"
# Correct:
Write-Output "hello"

\end{lstlisting}

\paragraph*{Error 2: "Variable is null"}\label{docs__pandoc__latex__src__cli_complete_curriculum.md__error-2-variable-is-null}

Cause: Variable was never assigned

Fix: Make sure variable is set before using it

\begin{lstlisting}[style=Alabaster, language=powershell]
$myvar = "hello"  # Set first
Write-Output $myvar  # Then use

\end{lstlisting}

\paragraph*{Error 3: "Cannot find path"}\label{docs__pandoc__latex__src__cli_complete_curriculum.md__error-3-cannot-find-path}

Cause: Wrong folder path

Fix: Verify path exists

\begin{lstlisting}[style=Alabaster, language=powershell]
# Check if path exists:
Test-Path -Path "C:\Users\YourName\Documents"
# If false, the path is wrong

\end{lstlisting}

Warning

\textbf{"Access denied"} means PowerShell doesn\textquotesingle t have permission to perform the operation. Always understand \emph{why} you need administrator privileges before elevating --- running everything as admin creates unnecessary security risk.

\paragraph*{Error 4: "Access denied"}\label{docs__pandoc__latex__src__cli_complete_curriculum.md__error-4-access-denied}

Cause: Don\textquotesingle t have permission

Fix: Run PowerShell as administrator

\begin{itemize}
\tightlist
\item
  Right-click PowerShell -\textgreater{} Run as administrator
\end{itemize}

\subsubsection*{Debugging Technique: Trace Output}\label{docs__pandoc__latex__src__cli_complete_curriculum.md__debugging-technique-trace-output}

Add \texttt{Write-Output} statements to track what\textquotesingle s happening:

\begin{lstlisting}[style=Alabaster, language=powershell]
$path = "C:\Users\YourName\Documents"
Write-Output "Starting script. Path is: $path"
$files = ls -n $path
Write-Output "Found $($files.Count) files"
foreach ($file in $files) {
    Write-Output "Processing: $file"
    # Do something with $file
    Write-Output "Done with: $file"
}
Write-Output "Script complete"

\end{lstlisting}

Your screen reader will announce each step, so you know where errors happen.

\subsection*{Creating Professional Workflows}\label{docs__pandoc__latex__src__cli_complete_curriculum.md__creating-professional-workflows}

\subsubsection*{Example 1: Automated Project Setup}\label{docs__pandoc__latex__src__cli_complete_curriculum.md__example-1-automated-project-setup}

Scenario: You start a new 3D printing project regularly. Instead of creating folders manually:

\begin{lstlisting}[style=Alabaster, language=powershell]
function SetupNewProject {
    param([string]$projectName)
    $baseFolder = "C:\Users\YourName\Documents\3D_Projects"
    $projectFolder = Join-Path -Path $baseFolder -ChildPath $projectName
    # Create folder structure
    mkdir $projectFolder -Force
    mkdir "$projectFolder\designs" -Force
    mkdir "$projectFolder\output" -Force
    mkdir "$projectFolder\notes" -Force
    # Create a README
    $readmeContent = @"
# $projectName
Created: $(Get-Date)
## Designs
All .scad files go here.
## Output
STL and other exports go here.
## Notes
Project notes and observations.
"@
    $readmeContent | Out-File -FilePath "$projectFolder\README.txt" -Encoding utf8
    Write-Output "Project setup complete: $projectFolder"
}
# Use it:
SetupNewProject -projectName "MyKeychain"
SetupNewProject -projectName "PhoneStand"

\end{lstlisting}

What it does:

\begin{itemize}
\tightlist
\item
  Creates folder structure for a new project
\item
  Sets up subfolders for designs, output, notes
\item
  Creates a README file automatically
\end{itemize}

\subsubsection*{Example 2: Batch File Verification}\label{docs__pandoc__latex__src__cli_complete_curriculum.md__example-2-batch-file-verification}

Scenario: Before processing, verify all required files exist:

\begin{lstlisting}[style=Alabaster, language=powershell]
function VerifyProjectFiles {
    param([string]$projectFolder)
    $requiredFiles = @(
        "README.txt",
        "designs",
        "output",
        "notes"
    )
    $allGood = $true
    foreach ($item in $requiredFiles) {
        $path = Join-Path -Path $projectFolder -ChildPath $item
        if (Test-Path -Path $path) {
            Write-Output " Found: $item"
        } else {
            Write-Output " Missing: $item"
            $allGood = $false
        }
    }
    if ($allGood) {
        Write-Output "All checks passed!"
        return $true
    } else {
        Write-Output "Some files are missing!"
        return $false
    }
}
# Use it:
$verified = VerifyProjectFiles -projectFolder "C:\Users\YourName\Documents\3D_Projects\MyKeychain"
if ($verified) {
    Write-Output "Safe to proceed with processing"
}

\end{lstlisting}

\subsection*{Screen Reader Tips for Scripts}\label{docs__pandoc__latex__src__cli_complete_curriculum.md__screen-reader-tips-for-scripts}

\subsubsection*{Making Script Output Readable}\label{docs__pandoc__latex__src__cli_complete_curriculum.md__making-script-output-readable}

Problem: Script runs but output scrolls too fast or is hard to follow

Solution 1: Save to file

\begin{lstlisting}[style=Alabaster, language=powershell]
.\my-script.ps1 > output.txt
notepad.exe output.txt

\end{lstlisting}

Solution 2: Use Write-Output with clear sections

\begin{lstlisting}[style=Alabaster, language=powershell]
Write-Output "========== STARTING =========="
Write-Output ""
# ... script ...
Write-Output ""
Write-Output "========== COMPLETE =========="

\end{lstlisting}

Solution 3: Pause between major sections

\begin{lstlisting}[style=Alabaster, language=powershell]
Write-Output "Pausing... Press Enter to continue"
Read-Host

\end{lstlisting}

Your screen reader will announce the pause, give you time to read output.

\subsubsection*{Announcing Progress}\label{docs__pandoc__latex__src__cli_complete_curriculum.md__announcing-progress}

For long-running scripts:

\begin{lstlisting}[style=Alabaster, language=powershell]
$files = ls -n *.scad
$count = 0
foreach ($file in $files) {
    $count++
    Write-Output "Processing $count of $($files.Count): $file"
    # Do something with $file
}
Write-Output "All $count files processed!"

\end{lstlisting}

\subsection*{Practice Exercises}\label{docs__pandoc__latex__src__cli_complete_curriculum.md__practice-exercises-1}

\subsubsection*{Exercise 1: Your First Script}\label{docs__pandoc__latex__src__cli_complete_curriculum.md__exercise-1-your-first-script}

Goal: Create and run a simple script

Steps:

\begin{enumerate}
\tightlist
\item
  Create file: \texttt{notepad.exe\ hello.ps1}
\item
  Type:

  \begin{lstlisting}[style=Alabaster, language=powershell]
  Write-Output "Hello from my first PowerShell script!"
  pwd
  ls -n

  \end{lstlisting}
\item
  Save and run: \texttt{.\textbackslash{}hello.ps1}
\end{enumerate}

Checkpoint: You should see output for each command.

\subsubsection*{Exercise 2: Script with a Variable}\label{docs__pandoc__latex__src__cli_complete_curriculum.md__exercise-2-script-with-a-variable}

Goal: Use a variable to make the script flexible

Steps:

\begin{enumerate}
\tightlist
\item
  Create file: \texttt{notepad.exe\ smart-listing.ps1}
\item
  Type:

  \begin{lstlisting}[style=Alabaster, language=powershell]
  $targetFolder = "C:\Users\YourName\Documents"
  Write-Output "Listing contents of: $targetFolder"
  ls -n $targetFolder

  \end{lstlisting}
\item
  Edit \texttt{\$targetFolder} to a real folder on your computer
\item
  Run: \texttt{.\textbackslash{}smart-listing.ps1}
\end{enumerate}

Checkpoint: You should see listing of that specific folder.

\subsubsection*{Exercise 3: Function}\label{docs__pandoc__latex__src__cli_complete_curriculum.md__exercise-3-function}

Goal: Create a reusable function

Steps:

\begin{enumerate}
\tightlist
\item
  Create file: \texttt{notepad.exe\ navigate.ps1}
\item
  Type:

  \begin{lstlisting}[style=Alabaster, language=powershell]
  function GoTo {
      param([string]$path)
      if (Test-Path -Path $path) {
          cd $path
          Write-Output "Now in: $(pwd)"
          Write-Output "Contents:"
          ls -n
      } else {
          Write-Output "Path does not exist: $path"
      }
  }
  # Test the function:
  GoTo -path "C:\Users\YourName\Documents"
  GoTo -path "C:\Users\YourName\Downloads"

  \end{lstlisting}
\item
  Run: \texttt{.\textbackslash{}navigate.ps1}
\end{enumerate}

Checkpoint: Both commands should work, showing contents of each folder.

\subsubsection*{Exercise 4: Loop}\label{docs__pandoc__latex__src__cli_complete_curriculum.md__exercise-4-loop}

Goal: Use a loop to repeat an action

Steps:

\begin{enumerate}
\tightlist
\item
  Create file: \texttt{notepad.exe\ repeat.ps1}
\item
  Type:

  \begin{lstlisting}[style=Alabaster, language=powershell]
  Write-Output "Demonstrating a loop:"
  for ($i = 1; $i -le 5; $i++) {
      Write-Output "Iteration $i: Hello!"
  }
  Write-Output "Loop complete!"

  \end{lstlisting}
\item
  Run: \texttt{.\textbackslash{}repeat.ps1}
\end{enumerate}

Checkpoint: Should print "Iteration 1" through "Iteration 5".

\subsubsection*{Exercise 5: Real-World Script}\label{docs__pandoc__latex__src__cli_complete_curriculum.md__exercise-5-real-world-script}

Goal: Create a useful script for a real task

Steps:

\begin{enumerate}
\tightlist
\item
  Create a folder: \texttt{mkdir\ C:\textbackslash{}Users\textbackslash{}YourName\textbackslash{}Documents\textbackslash{}TestFiles}
\item
  Create some test files:

  \begin{lstlisting}[style=Alabaster, language=powershell]
  echo "test" > C:\Users\YourName\Documents\TestFiles\file1.txt
  echo "test" > C:\Users\YourName\Documents\TestFiles\file2.txt
  echo "test" > C:\Users\YourName\Documents\TestFiles\file3.txt

  \end{lstlisting}
\item
  Create script: \texttt{notepad.exe\ report.ps1}
\item
  Type:

  \begin{lstlisting}[style=Alabaster, language=powershell]
  $folder = "C:\Users\YourName\Documents\TestFiles"
  $files = ls -n $folder
  Write-Output "=== FILE REPORT ==="
  Write-Output "Folder: $folder"
  Write-Output "Total files: $($files.Count)"
  Write-Output ""
  Write-Output "Files:"
  foreach ($file in $files) {
      Write-Output "  - $file"
  }
  Write-Output ""
  Write-Output "=== END REPORT ==="

  \end{lstlisting}
\item
  Run: \texttt{.\textbackslash{}report.ps1}
\end{enumerate}

Checkpoint: Should show report of all files in the test folder.

\subsection*{Quiz --- Lesson 6}\label{docs__pandoc__latex__src__cli_complete_curriculum.md__quiz--lesson-6}

\begin{enumerate}
\tightlist
\item
  What is a PowerShell script?
\item
  What file extension do PowerShell scripts use?
\item
  What is a variable and how do you create one?
\item
  What is a function and why would you use one?
\item
  How do you run a script?
\item
  What is a loop and what does \texttt{foreach} do?
\item
  What does \texttt{Test-Path} do?
\item
  How do you handle errors in a script?
\item
  When would you use \texttt{Try-Catch}?
\item
  What technique makes script output readable for screen readers?
\end{enumerate}

\subsection*{Extension Problems}\label{docs__pandoc__latex__src__cli_complete_curriculum.md__extension-problems-6}

\begin{enumerate}
\tightlist
\item
  Auto-Backup Script: Create a script that copies all files from one folder to another, announcing progress
\item
  File Counter: Write a function that counts files by extension (.txt, .scad, .stl, etc.)
\item
  Folder Cleaner: Script that deletes files older than 30 days (with user confirmation)
\item
  Project Template: Function that creates a complete project folder structure with all needed files
\item
  Batch Rename: Script that renames all files in a folder according to a pattern
\item
  Log Generator: Create a script that records what it does to a log file for later review
\item
  Scheduled Task: Set up a script to run automatically every day at a specific time
\item
  File Verifier: Check that all SCAD files in a folder have corresponding STL exports
\item
  Report Generator: Create a summary report of all projects in a folder
\item
  Error Tracker: Script that lists all commands that had errors in your recent history
\end{enumerate}

\subsection*{Important Notes}\label{docs__pandoc__latex__src__cli_complete_curriculum.md__important-notes}

\begin{itemize}
\tightlist
\item
  Always test scripts on small sets of files first before running them on important data
\item
  Save your work regularly - use version control if possible
\item
  Test error handling - make sure errors don\textquotesingle t crash silently
\item
  Document your scripts - use comments so you remember what each part does
\item
  Backup before batch operations - if something goes wrong, you have the original
\end{itemize}

\subsection*{References}\label{docs__pandoc__latex__src__cli_complete_curriculum.md__references-6}

\begin{itemize}
\tightlist
\item
  Microsoft PowerShell Scripting Guide: \url{https://docs.microsoft.com/powershell/scripting/}
\item
  Function Documentation: \url{https://docs.microsoft.com/powershell/module/microsoft.powershell.core/about/about_functions_advanced}
\item
  Error Handling: \url{https://docs.microsoft.com/powershell/module/microsoft.powershell.core/about/about_error_handling}
\item
  Loops: \url{https://docs.microsoft.com/powershell/module/microsoft.powershell.core/about/about_foreach}
\end{itemize}

Next Steps: After mastering this lesson, explore PowerShell modules, remoting, and 3D printing integration in the main curriculum.

\begin{center}\rule{0.5\linewidth}{0.5pt}\end{center}

\protect\phantomsection\label{docs__pandoc__latex__src__cli_complete_curriculum.md__ps-unit-exam}{}

\section*{Unit Exam --- Comprehensive Assessment}\label{docs__pandoc__latex__src__cli_complete_curriculum.md__unit-exam--comprehensive-assessment}

\begin{center}\rule{0.5\linewidth}{0.5pt}\end{center}

Estimated time: 60-90 minutes

\subsection*{Key Learning Outcomes Assessed}\label{docs__pandoc__latex__src__cli_complete_curriculum.md__key-learning-outcomes-assessed}

By completing this unit test, you will demonstrate:

\begin{enumerate}
\tightlist
\item
  Understanding of file system navigation and path concepts
\item
  Proficiency with file and folder manipulation commands
\item
  Ability to redirect and pipe command output
\item
  Knowledge of environment variables and aliases
\item
  Screen-reader accessibility best practices in terminal environments
\item
  Problem-solving and command chaining skills
\end{enumerate}

\subsection*{Target Audience:}\label{docs__pandoc__latex__src__cli_complete_curriculum.md__target-audience}

Users who have completed Lessons 0--5 and need to demonstrate mastery of PowerShell fundamentals.

\subsection*{Instructions:}\label{docs__pandoc__latex__src__cli_complete_curriculum.md__instructions}

Complete all sections below. For multiple choice, select the best answer. For short answers, write one to two sentences. For hands-on tasks, capture evidence (screenshots or output files) and submit alongside your answers.

\subsection*{Part A: Multiple Choice Questions (20 questions)}\label{docs__pandoc__latex__src__cli_complete_curriculum.md__part-a-multiple-choice-questions-20-questions}

Select the best answer for each question. Each question is worth 1 point.

\begin{enumerate}
\item
  What is the primary purpose of the \texttt{PATH} environment variable?

  \begin{itemize}
  \tightlist
  \item
    A) Store your home directory location
  \item
    B) Tell the shell where to find executable programs
  \item
    C) Configure the visual appearance of the terminal
  \item
    D) Store the current working directory name
  \end{itemize}
\item
  Which command prints your current working directory?

  \begin{itemize}
  \tightlist
  \item
    A) \texttt{ls\ -n}
  \item
    B) \texttt{cd}
  \item
    C) \texttt{pwd}
  \item
    D) \texttt{whoami}
  \end{itemize}
\item
  What does the \texttt{\textasciitilde{}} symbol represent in PowerShell paths?

  \begin{itemize}
  \tightlist
  \item
    A) The root directory
  \item
    B) The current directory
  \item
    C) The parent directory
  \item
    D) The home directory
  \end{itemize}
\item
  How do you list only file names (not full details) in a way that is screen-reader friendly?

  \begin{itemize}
  \tightlist
  \item
    A) \texttt{ls}
  \item
    B) \texttt{ls\ -n}
  \item
    C) \texttt{ls\ -l}
  \item
    D) \texttt{cat\ -n}
  \end{itemize}
\item
  Which command creates a new empty file?

  \begin{itemize}
  \tightlist
  \item
    A) \texttt{mkdir\ filename}
  \item
    B) \texttt{ni\ filename}
  \item
    C) \texttt{touch\ filename}
  \item
    D) \texttt{echo\ filename}
  \end{itemize}
\end{enumerate}

Important

\texttt{\textgreater{}} \textbf{overwrites} a file completely --- all previous content is lost. \texttt{\textgreater{}\textgreater{}} \textbf{appends} to the end of a file, preserving existing content. Using \texttt{\textgreater{}} when you meant \texttt{\textgreater{}\textgreater{}} will permanently destroy the file\textquotesingle s previous contents.

\begin{enumerate}
\setcounter{enumi}{5}
\item
  What is the difference between \texttt{\textgreater{}} and \texttt{\textgreater{}\textgreater{}}?

  \begin{itemize}
  \tightlist
  \item
    A) \texttt{\textgreater{}} redirects to file, \texttt{\textgreater{}\textgreater{}} displays on screen
  \item
    B) \texttt{\textgreater{}} overwrites a file, \texttt{\textgreater{}\textgreater{}} appends to a file
  \item
    C) They do the same thing
  \item
    D) \texttt{\textgreater{}} is for text, \texttt{\textgreater{}\textgreater{}} is for binary
  \end{itemize}
\item
  What does the pipe operator \texttt{\textbar{}} do?

  \begin{itemize}
  \tightlist
  \item
    A) Creates a folder
  \item
    B) Sends the output of one command to the input of another
  \item
    C) Deletes files matching a pattern
  \item
    D) Lists all processes
  \end{itemize}
\item
  Which command copies a file?

  \begin{itemize}
  \tightlist
  \item
    A) \texttt{mv}
  \item
    B) \texttt{rm}
  \item
    C) \texttt{cp}
  \item
    D) \texttt{cd}
  \end{itemize}
\item
  How do you rename a file from \texttt{oldname.txt} to \texttt{newname.txt}?

  \begin{itemize}
  \tightlist
  \item
    A) \texttt{cp\ oldname.txt\ newname.txt}
  \item
    B) \texttt{mv\ oldname.txt\ newname.txt}
  \item
    C) \texttt{rename\ oldname.txt\ newname.txt}
  \item
    D) \texttt{rn\ oldname.txt\ newname.txt}
  \end{itemize}
\item
  What is the purpose of \texttt{Select-String}?

  \begin{itemize}
  \tightlist
  \item
    A) Select files in a directory
  \item
    B) Search for text patterns within a file
  \item
    C) Select a string to copy to clipboard
  \item
    D) Select which shell to use
  \end{itemize}
\item
  Which key combination allows you to autocomplete a path in PowerShell?

  \begin{itemize}
  \tightlist
  \item
    A) \texttt{Ctrl\ +\ A}
  \item
    B) \texttt{Ctrl\ +\ E}
  \item
    C) \texttt{Tab}
  \item
    D) \texttt{Space}
  \end{itemize}
\item
  How do you copy text to the Windows clipboard from PowerShell?

  \begin{itemize}
  \tightlist
  \item
    A) \texttt{cat\ filename\ \textgreater{}\ clipboard}
  \item
    B) \texttt{cat\ filename\ \textbar{}\ clip}
  \item
    C) \texttt{copy\ filename}
  \item
    D) \texttt{cat\ filename\ \textbar{}\ paste}
  \end{itemize}
\item
  What does \texttt{Get-Command\ openscad} do?

  \begin{itemize}
  \tightlist
  \item
    A) Opens the OpenSCAD application
  \item
    B) Gets help about the OpenSCAD command
  \item
    C) Locates the full path of the openscad executable
  \item
    D) Lists all available commands
  \end{itemize}
\item
  Which wildcard matches any single character?

  \begin{itemize}
  \tightlist
  \item
    A) \texttt{*}
  \item
    B) \texttt{?}
  \item
    C) \texttt{\%}
  \item
    D) \texttt{\#}
  \end{itemize}
\item
  What is the purpose of the \texttt{\&} call operator?

  \begin{itemize}
  \tightlist
  \item
    A) Run a script or executable by full path
  \item
    B) Execute all commands in parallel
  \item
    C) Combine multiple commands
  \item
    D) Create an alias
  \end{itemize}
\item
  How do you create a temporary alias for a command?

  \begin{itemize}
  \tightlist
  \item
    A) \texttt{alias\ preview=\textquotesingle{}openscad\textquotesingle{}}
  \item
    B) \texttt{Set-Alias\ -Name\ preview\ -Value\ openscad}
  \item
    C) \texttt{New-Alias\ preview\ openscad}
  \item
    D) \texttt{Alias\ preview\ =\ openscad}
  \end{itemize}
\item
  Where is your PowerShell profile typically stored?

  \begin{itemize}
  \tightlist
  \item
    A) C:\textbackslash Program Files\textbackslash PowerShell\textbackslash profile.ps1
  \item
    B) The location returned by \texttt{echo\ \$PROFILE}
  \item
    C) \textasciitilde{}/PowerShell/profile.ps1
  \item
    D) \textasciitilde{}/.bashrc
  \end{itemize}
\item
  How do you abort a long-running command in PowerShell?

  \begin{itemize}
  \tightlist
  \item
    A) Press \texttt{Escape}
  \item
    B) Press \texttt{Ctrl\ +\ X}
  \item
    C) Press \texttt{Ctrl\ +\ C}
  \item
    D) Press \texttt{Alt\ +\ F4}
  \end{itemize}
\item
  What command shows the history of previously run commands?

  \begin{itemize}
  \tightlist
  \item
    A) \texttt{history}
  \item
    B) \texttt{Get-History}
  \item
    C) \texttt{Show-History}
  \item
    D) Both A and B
  \end{itemize}
\item
  How do you permanently set an alias so it persists across PowerShell sessions?

  \begin{itemize}
  \tightlist
  \item
    A) Use \texttt{Set-Alias} in the terminal every time
  \item
    B) Add the \texttt{Set-Alias} line to your PowerShell profile
  \item
    C) Use the Windows Control Panel
  \item
    D) Aliases cannot be made permanent
  \end{itemize}
\end{enumerate}

\subsection*{Part B: Short Answer Questions (10 questions)}\label{docs__pandoc__latex__src__cli_complete_curriculum.md__part-b-short-answer-questions-10-questions}

Answer each question in one to two sentences. Each question is worth 2 points.

\begin{enumerate}
\item
  Explain the difference between absolute and relative paths. Give one example of each.
\item
  Why is \texttt{ls\ -n} preferred over \texttt{ls} for screen reader users? Describe what flag you would use to list only files.
\item
  What is the purpose of redirecting output to a file, and give an example of when you would use \texttt{\textgreater{}} instead of \texttt{\textgreater{}\textgreater{}}?
\end{enumerate}

Warning

\texttt{rm\ -r\ \textasciitilde{}/Documents/my\_folder} permanently deletes the folder and \textbf{everything inside it} --- all files, all subdirectories --- with no confirmation and no undo. Always verify the path with \texttt{ls\ -1} first and consider making a backup with \texttt{cp\ -r} before running recursive deletions.

\begin{enumerate}
\setcounter{enumi}{3}
\item
  Describe what would happen if you ran \texttt{rm\ -r\ \textasciitilde{}/Documents/my\_folder} and why this command should be used carefully.
\item
  How would you search for all files with a \texttt{.scad} extension in your current directory? Write the command.
\item
  Explain what happens when you pipe the output of \texttt{ls\ -n} into \texttt{clip}. What would you do next?
\item
  What is an environment variable, and give one example of how you might use it in PowerShell.
\item
  If a program is not in your \texttt{PATH}, what two methods could you use to run it from PowerShell?
\item
  Describe how you would open a file in Notepad and also add a line to it from PowerShell.
\item
  What is one strategy you would use if your screen reader stops announcing terminal output while using PowerShell?
\end{enumerate}

\subsection*{Part C: Hands-On Tasks (10 tasks)}\label{docs__pandoc__latex__src__cli_complete_curriculum.md__part-c-hands-on-tasks-10-tasks}

Complete each task and capture evidence (screenshots, output files, or command transcripts). Each task is worth 3 points.

\subsubsection*{Tasks 1-5: File System and Navigation}\label{docs__pandoc__latex__src__cli_complete_curriculum.md__tasks-1-5-file-system-and-navigation}

\begin{enumerate}
\item
  Create a folder structure \texttt{\textasciitilde{}/Documents/PowerShell\_Assessment/Projects} using a single command. Capture the \texttt{ls\ -n} output showing the creation.
\item
  Create five files named \texttt{project\_1.scad}, \texttt{project\_2.scad}, \texttt{project\_3.txt}, \texttt{notes\_1.txt}, and \texttt{notes\_2.txt} inside the \texttt{Projects} folder. Use wildcards to list only \texttt{.scad} files, then capture the output.
\item
  Copy the entire \texttt{Projects} folder to \texttt{Projects\_Backup} using \texttt{cp\ -r}. Capture the \texttt{ls\ -n} output showing both folders exist.
\item
  Move (rename) \texttt{project\_1.scad} to \texttt{project\_1\_final.scad}. Capture the \texttt{ls\ -n} output showing the renamed file.
\item
  Delete \texttt{notes\_1.txt} and \texttt{notes\_2.txt} using a single \texttt{rm} command with wildcards. Capture the final \texttt{ls\ -n} output.
\end{enumerate}

\subsubsection*{Tasks 6-10: Advanced Operations and Scripting}\label{docs__pandoc__latex__src__cli_complete_curriculum.md__tasks-6-10-advanced-operations-and-scripting}

\begin{enumerate}
\setcounter{enumi}{5}
\item
  Create a file called \texttt{my\_data.txt} with at least four lines using \texttt{echo} and \texttt{\textgreater{}\textgreater{}}. Then read it with \texttt{cat\ my\_data.txt} and capture the output.
\item
  Use \texttt{Select-String} to search for a keyword (e.g., "project") in \texttt{my\_data.txt} and pipe the results to \texttt{clip}. Paste the results into Notepad and capture a screenshot.
\item
  List all files in the \texttt{Projects} folder and redirect the output to \texttt{projects\_list.txt}. Open it in Notepad and capture a screenshot of the file.
\item
  Create a temporary alias called \texttt{myls} that runs \texttt{ls\ -n}, test it, and capture the output. Then explain what would be required to make it permanent.
\item
  Run \texttt{Get-Help\ Get-ChildItem} and redirect the output to \texttt{help\_output.txt}. Open the file in Notepad and capture a screenshot showing at least the first page of help content.
\end{enumerate}

\subsection*{Grading Rubric}\label{docs__pandoc__latex__src__cli_complete_curriculum.md__grading-rubric}

{\def\LTcaptype{none} % do not increment counter
\begin{longtable}[]{@{}llll@{}}
\toprule\noalign{}
Section & Questions & Points Each & Total \\
\midrule\noalign{}
\endhead
\bottomrule\noalign{}
\endlastfoot
Multiple Choice & 20 & 1 & 20 \\
Short Answer & 10 & 2 & 20 \\
Hands-On Tasks & 10 & 3 & 30 \\
Total & 40 & - & 70 \\
\end{longtable}
}

Passing Score: 49 points (70\%)

\subsection*{Helpful Resources for Review}\label{docs__pandoc__latex__src__cli_complete_curriculum.md__helpful-resources-for-review}

\begin{itemize}
\tightlist
\item
  \href{https://learn.microsoft.com/powershell/scripting/overview}{PowerShell Command Reference}
\item
  \href{https://learn.microsoft.com/powershell/scripting/learn/shell/navigate-the-filesystem}{Navigation and File System}
\item
  \href{https://learn.microsoft.com/powershell/scripting/learn/shell/using-the-pipeline}{Using Pipes and Filtering}
\item
  \href{https://learn.microsoft.com/powershell/module/microsoft.powershell.core/about/about_profiles}{Profile and Aliases}
\item
  \href{https://learn.microsoft.com/powershell/scripting/windows-powershell/ise/accessibility-in-windows-powershell-ise}{Screen Reader Accessibility Tips}
\end{itemize}

\subsection*{Submission Checklist}\label{docs__pandoc__latex__src__cli_complete_curriculum.md__submission-checklist}

\begin{itemize}
\tightlist
\item[$\square$]
  All 20 multiple choice questions answered
\item[$\square$]
  All 10 short answer questions answered (1-2 sentences each)
\item[$\square$]
  All 10 hands-on tasks completed with evidence captured
\item[$\square$]
  Files/screenshots organized and labeled clearly
\item[$\square$]
  Submission includes this checklist
\end{itemize}

\begin{center}\rule{0.5\linewidth}{0.5pt}\end{center}

\protect\phantomsection\label{docs__pandoc__latex__src__cli_complete_curriculum.md__part-2-cmd}{}

\section*{Part 2 --- Windows Command Prompt (CMD)}\label{docs__pandoc__latex__src__cli_complete_curriculum.md__part-2--windows-command-prompt-cmd-1}

\begin{center}\rule{0.5\linewidth}{0.5pt}\end{center}

\textbf{Last Updated:} February 2026\\
\textbf{Total Duration:} 30--45 hours of instruction and practice (for screen reader users)\\
\textbf{Target Users:} Anyone using a screen reader --- NVDA, JAWS, Windows Narrator, and Dolphin SuperNova are all covered throughout this curriculum.\\
\textbf{Alternate to:} PowerShell curriculum (same concepts, CMD syntax)

\emph{This single file contains the complete CMD Foundation curriculum in reading order. All lessons, quizzes, extension problems, and the unit exam are included. No additional files are needed.}

\begin{center}\rule{0.5\linewidth}{0.5pt}\end{center}

\subsection*{Table of Contents}\label{docs__pandoc__latex__src__cli_complete_curriculum.md__table-of-contents-2}

{\def\LTcaptype{none} % do not increment counter
\begin{longtable}[]{@{}
  >{\raggedright\arraybackslash}p{(\linewidth - 4\tabcolsep) * \real{0.3333}}
  >{\raggedright\arraybackslash}p{(\linewidth - 4\tabcolsep) * \real{0.3333}}
  >{\raggedright\arraybackslash}p{(\linewidth - 4\tabcolsep) * \real{0.3333}}@{}}
\toprule\noalign{}
\begin{minipage}[b]{\linewidth}\raggedright
Section
\end{minipage} & \begin{minipage}[b]{\linewidth}\raggedright
Title
\end{minipage} & \begin{minipage}[b]{\linewidth}\raggedright
Estimated Time
\end{minipage} \\
\midrule\noalign{}
\endhead
\bottomrule\noalign{}
\endlastfoot
\hyperref[docs__pandoc__latex__src__cli_complete_curriculum.md__cmd-accessibility-guide]{Screen Reader Accessibility Guide}
& NVDA, JAWS, Narrator, and SuperNova reference for CMD & 1.5 hours \\
\hyperref[docs__pandoc__latex__src__cli_complete_curriculum.md__cmd-lesson-pre]{Lesson Pre: Your First Terminal}
& Opening CMD, first commands, basic navigation & 1.5--2 hours \\
\hyperref[docs__pandoc__latex__src__cli_complete_curriculum.md__cmd-lesson-0]{Lesson 0: Getting Started}
& Paths, shortcuts, tab completion & 1.5 hours \\
\hyperref[docs__pandoc__latex__src__cli_complete_curriculum.md__cmd-lesson-1]{Lesson 1: Navigation}
& Moving around the file system confidently & 2--2.5 hours \\
\hyperref[docs__pandoc__latex__src__cli_complete_curriculum.md__cmd-lesson-2]{Lesson 2: File and Folder Manipulation}
& Create, copy, move, delete files and folders & 2.5--3 hours \\
\hyperref[docs__pandoc__latex__src__cli_complete_curriculum.md__cmd-lesson-3]{Lesson 3: Input, Output, and Piping}
& Redirect output, chain commands with pipes & 2.5--3 hours \\
\hyperref[docs__pandoc__latex__src__cli_complete_curriculum.md__cmd-lesson-4]{Lesson 4: Environment Variables and Aliases}
& Automate settings, create doskey shortcuts & 2--2.5 hours \\
\hyperref[docs__pandoc__latex__src__cli_complete_curriculum.md__cmd-lesson-5]{Lesson 5: Filling in the Gaps}
& Autorun scripts, history, debugging & 2--2.5 hours \\
\hyperref[docs__pandoc__latex__src__cli_complete_curriculum.md__cmd-lesson-6]{Lesson 6: Advanced Techniques}
& Batch scripts, subroutines, loops, professional workflows &
4--4.5 hours \\
\hyperref[docs__pandoc__latex__src__cli_complete_curriculum.md__cmd-unit-exam]{Unit Exam}
& Comprehensive assessment --- Lessons 0 through 6 & 60--90 minutes \\
\end{longtable}
}

\textbf{Total:} approximately 30--45 hours including practice exercises

\begin{center}\rule{0.5\linewidth}{0.5pt}\end{center}

\subsection*{How to Use This File}\label{docs__pandoc__latex__src__cli_complete_curriculum.md__how-to-use-this-file-1}

\textbf{If you have never used a terminal before:} read from the top in order. Do not skip lessons --- each builds on the one before it.

\textbf{If you have used a terminal but not with a screen reader:} start with the Accessibility Guide, skim Lesson Pre, then continue from Lesson 0.

\textbf{If you are experienced with both:} jump directly to the lesson you need. Use the quiz at the end of each lesson to self-test before moving on.

\textbf{Reference while working:} the Accessibility Guide contains material you will return to throughout the curriculum. Keep a copy open or bookmark that section.

\textbf{PowerShell users:} this curriculum covers the same concepts as the PowerShell curriculum but uses CMD syntax. The two curricula are interchangeable as learning paths --- choose one and stay consistent.

\begin{center}\rule{0.5\linewidth}{0.5pt}\end{center}

\subsection*{Curriculum Rules (Read These First)}\label{docs__pandoc__latex__src__cli_complete_curriculum.md__curriculum-rules-read-these-first-1}

Important

\textbf{Rule 1 --- Always know where you are.} Start every session with \texttt{cd} (no arguments) to print your current location. If you don\textquotesingle t know your path, you will get lost and may operate on the wrong files.

\textbf{Rule 1 --- Always know where you are.}\\
Every session, first thing: \texttt{cd} (with no arguments). If you do not know your path, you will get lost.

Warning

\textbf{Files deleted from CMD do not go to the Recycle Bin.} They are permanently gone immediately. Always run \texttt{dir\ /B} to confirm exactly what you are about to delete before running \texttt{del} or \texttt{rmdir\ /S}.

\textbf{Rule 2 --- Check before you delete.}\\
Before any \texttt{del} or \texttt{rmdir} command: \texttt{dir\ /B} to confirm exactly what you are about to remove. Deleted files do not go to the Recycle Bin from CMD.

Tip

Always use \texttt{dir\ /B} --- never bare \texttt{dir}. The \texttt{/B} (bare) flag outputs one name per line, which screen readers can follow reliably. Bare \texttt{dir} outputs multi-column details that can be difficult to navigate audibly.

\textbf{Rule 3 --- Use \texttt{dir\ /B} for listings.}\\
Always \texttt{dir\ /B} --- never bare \texttt{dir}. The \texttt{/B} flag gives one name per line, which all screen readers can follow reliably.

Tip

When terminal output is confusing or overwhelming, redirect it to a file: \texttt{command\ \textgreater{}\ output.txt}, then open it with \texttt{notepad.exe\ output.txt}. This is clearer than reading raw terminal output and works with all four screen readers.

\textbf{Rule 4 --- When output is confusing, redirect to a file.}\\
\texttt{command\ \textgreater{}\ output.txt} then \texttt{notepad.exe\ output.txt}. This works with all four screen readers.

\textbf{Rule 5 --- Save everything you create.}\\
Keep all practice files in a dedicated learning folder: \texttt{mkdir\ \%USERPROFILE\%\textbackslash{}CMD\_Learning}

\begin{center}\rule{0.5\linewidth}{0.5pt}\end{center}

\begin{center}\rule{0.5\linewidth}{0.5pt}\end{center}

\protect\phantomsection\label{docs__pandoc__latex__src__cli_complete_curriculum.md__cmd-accessibility-guide}{}

\section*{Screen Reader Accessibility Guide}\label{docs__pandoc__latex__src__cli_complete_curriculum.md__screen-reader-accessibility-guide-1}

\begin{center}\rule{0.5\linewidth}{0.5pt}\end{center}

Target Users: NVDA, JAWS, and other screen reader users\\
Last Updated: 2026

This guide supports the CMD Foundation curriculum and helps screen reader users navigate and work efficiently with the Windows Command Prompt.

\subsection*{Table of Contents}\label{docs__pandoc__latex__src__cli_complete_curriculum.md__table-of-contents-3}

\begin{enumerate}
\tightlist
\item
  \hyperref[docs__pandoc__latex__src__cli_complete_curriculum.md__cmd-getting-started]{Getting Started with Screen Readers}
\item
  \hyperref[docs__pandoc__latex__src__cli_complete_curriculum.md__cmd-nvda-tips]{NVDA-Specific Tips}
\item
  \hyperref[docs__pandoc__latex__src__cli_complete_curriculum.md__cmd-jaws-tips]{JAWS-Specific Tips}
\item
  \hyperref[docs__pandoc__latex__src__cli_complete_curriculum.md__cmd-general-terminal]{General Terminal Accessibility}
\item
  \hyperref[docs__pandoc__latex__src__cli_complete_curriculum.md__cmd-long-output]{Working with Long Output}
\item
  \hyperref[docs__pandoc__latex__src__cli_complete_curriculum.md__cmd-shortcuts]{Keyboard Shortcuts Reference}
\item
  \hyperref[docs__pandoc__latex__src__cli_complete_curriculum.md__cmd-troubleshooting]{Troubleshooting}
\end{enumerate}

\protect\phantomsection\label{docs__pandoc__latex__src__cli_complete_curriculum.md__cmd-getting-started}{}

\subsection*{Getting Started with Screen Readers}\label{docs__pandoc__latex__src__cli_complete_curriculum.md__getting-started-with-screen-readers-1}

\subsubsection*{Which Screen Reader Should I Use?}\label{docs__pandoc__latex__src__cli_complete_curriculum.md__which-screen-reader-should-i-use-1}

Both NVDA and JAWS work with CMD; NVDA is free and often easiest to start with. JAWS provides advanced features for power users. Dolphin SuperNova and Windows Narrator are also options:

\begin{itemize}
\tightlist
\item
  Dolphin SuperNova: commercial speech, braille, and magnification (use vendor docs for key mappings).
\item
  Windows Narrator: built into Windows for quick, no-install access.
\end{itemize}

\subsubsection*{Before You Start}\label{docs__pandoc__latex__src__cli_complete_curriculum.md__before-you-start-1}

\begin{enumerate}
\tightlist
\item
  Start your screen reader before opening CMD.
\item
  Open Command Prompt and listen for the window title and prompt.
\item
  If silent, press Alt+Tab to find the window.
\end{enumerate}

\subsubsection*{What is CMD?}\label{docs__pandoc__latex__src__cli_complete_curriculum.md__what-is-cmd}

CMD (Command Prompt) is the classic Windows shell. Common commands include \texttt{dir}, \texttt{type}, and \texttt{more}. Paths use backslashes (e.g., \texttt{C:\textbackslash{}Users\textbackslash{}You}).

\protect\phantomsection\label{docs__pandoc__latex__src__cli_complete_curriculum.md__cmd-nvda-tips}{}

\subsection*{NVDA-Specific Tips}\label{docs__pandoc__latex__src__cli_complete_curriculum.md__nvda-specific-tips-1}

NVDA is available from \url{https://www.nvaccess.org/}

\subsubsection*{Dolphin SuperNova}\label{docs__pandoc__latex__src__cli_complete_curriculum.md__dolphin-supernova-1}

Dolphin SuperNova: \url{https://yourdolphin.com/supernova/} --- commercial screen reader and magnifier; check institutional licensing.

\subsubsection*{Windows Narrator}\label{docs__pandoc__latex__src__cli_complete_curriculum.md__windows-narrator-1}

Windows Narrator: \url{https://support.microsoft.com/narrator} --- built-in, simpler command set; useful when third-party readers are unavailable.

\subsubsection*{Key Commands for CMD}\label{docs__pandoc__latex__src__cli_complete_curriculum.md__key-commands-for-cmd}

{\def\LTcaptype{none} % do not increment counter
\begin{longtable}[]{@{}
  >{\raggedright\arraybackslash}p{(\linewidth - 2\tabcolsep) * \real{0.2361}}
  >{\raggedright\arraybackslash}p{(\linewidth - 2\tabcolsep) * \real{0.7639}}@{}}
\toprule\noalign{}
\begin{minipage}[b]{\linewidth}\raggedright
Command
\end{minipage} & \begin{minipage}[b]{\linewidth}\raggedright
What It Does
\end{minipage} \\
\midrule\noalign{}
\endhead
\bottomrule\noalign{}
\endlastfoot
NVDA+Home & Read the current line (your command or output) \\
NVDA+Down Arrow & Read from cursor to end of screen \\
NVDA+Up Arrow & Read from top to cursor \\
NVDA+Page Down & Read next page \\
NVDA+Page Up & Read previous page \\
NVDA+F7 & Open the Review Mode viewer (can scroll through text) \\
\end{longtable}
}

\subsubsection*{Example: Reading Long Output}\label{docs__pandoc__latex__src__cli_complete_curriculum.md__example-reading-long-output-2}

If \texttt{dir} produces many lines, redirect to a file and open it in Notepad:

\begin{lstlisting}[style=Alabaster, language=cmd]
dir /b > list.txt
notepad list.txt

\end{lstlisting}

\texttt{dir\ /b} shows one item per line (screen reader friendly).

\protect\phantomsection\label{docs__pandoc__latex__src__cli_complete_curriculum.md__cmd-jaws-tips}{}

\subsection*{JAWS-Specific Tips}\label{docs__pandoc__latex__src__cli_complete_curriculum.md__jaws-specific-tips-1}

JAWS is available from \url{https://www.freedomscientific.com/}

\subsubsection*{Key Commands for CMD}\label{docs__pandoc__latex__src__cli_complete_curriculum.md__key-commands-for-cmd-1}

{\def\LTcaptype{none} % do not increment counter
\begin{longtable}[]{@{}ll@{}}
\toprule\noalign{}
Command & What It Does \\
\midrule\noalign{}
\endhead
\bottomrule\noalign{}
\endlastfoot
Insert+Down Arrow & Read line by line downward \\
Insert+Up Arrow & Read line by line upward \\
Insert+Page Down & Read next page of text \\
Insert+Page Up & Read previous page of text \\
\end{longtable}
}

\subsubsection*{Example: Reading Long Output}\label{docs__pandoc__latex__src__cli_complete_curriculum.md__example-reading-long-output-3}

\begin{enumerate}
\tightlist
\item
  Redirect: \texttt{dir\ /b\ \textgreater{}\ list.txt}
\item
  Open Notepad: \texttt{notepad\ list.txt}
\item
  Use Insert+Ctrl+Down to read full contents.
\end{enumerate}

\protect\phantomsection\label{docs__pandoc__latex__src__cli_complete_curriculum.md__cmd-general-terminal}{}

\subsection*{General Terminal Accessibility}\label{docs__pandoc__latex__src__cli_complete_curriculum.md__general-terminal-accessibility-1}

\subsubsection*{Understanding the CMD Layout}\label{docs__pandoc__latex__src__cli_complete_curriculum.md__understanding-the-cmd-layout}

The Command Prompt window shows a title bar, a content area, and the prompt (e.g., \texttt{C:\textbackslash{}Users\textbackslash{}YourName\textgreater{}}).

\subsubsection*{The CMD Prompt}\label{docs__pandoc__latex__src__cli_complete_curriculum.md__the-cmd-prompt}

Example:

\begin{lstlisting}[style=Alabaster]
C:\Users\YourName>

\end{lstlisting}

\subsubsection*{Navigation Sequence}\label{docs__pandoc__latex__src__cli_complete_curriculum.md__navigation-sequence-1}

\begin{enumerate}
\tightlist
\item
  Screen reader announces the title
\item
  Then it announces the prompt line
\item
  Anything above prompt is prior output
\end{enumerate}

\protect\phantomsection\label{docs__pandoc__latex__src__cli_complete_curriculum.md__cmd-long-output}{}

\subsection*{Working with Long Output}\label{docs__pandoc__latex__src__cli_complete_curriculum.md__working-with-long-output-1}

\subsubsection*{Solution 1: Redirect to a File}\label{docs__pandoc__latex__src__cli_complete_curriculum.md__solution-1-redirect-to-a-file-1}

\begin{lstlisting}[style=Alabaster, language=cmd]
dir /b > list.txt
notepad list.txt

\end{lstlisting}

\subsubsection*{Solution 2: Use Pagination}\label{docs__pandoc__latex__src__cli_complete_curriculum.md__solution-2-use-pagination-1}

\begin{lstlisting}[style=Alabaster, language=cmd]
type largefile.txt | more

\end{lstlisting}

Use Space for next page and Q to quit.

\subsubsection*{Solution 3: Filter Output}\label{docs__pandoc__latex__src__cli_complete_curriculum.md__solution-3-filter-output-1}

\begin{lstlisting}[style=Alabaster, language=cmd]
dir /b | findstr /R "\.scad$"

\end{lstlisting}

\subsubsection*{Solution 4: Count Before Displaying}\label{docs__pandoc__latex__src__cli_complete_curriculum.md__solution-4-count-before-displaying-1}

\begin{lstlisting}[style=Alabaster, language=cmd]
dir /b | find /v "" /c

\end{lstlisting}

\protect\phantomsection\label{docs__pandoc__latex__src__cli_complete_curriculum.md__cmd-shortcuts}{}

\subsection*{Keyboard Shortcuts Reference}\label{docs__pandoc__latex__src__cli_complete_curriculum.md__keyboard-shortcuts-reference-1}

{\def\LTcaptype{none} % do not increment counter
\begin{longtable}[]{@{}ll@{}}
\toprule\noalign{}
Key & Action \\
\midrule\noalign{}
\endhead
\bottomrule\noalign{}
\endlastfoot
Up Arrow & Show previous command \\
Down Arrow & Show next command \\
Tab & Auto-complete file/folder names \\
Home & Jump to start of line \\
End & Jump to end of line \\
Ctrl+C & Stop command \\
Enter & Run command \\
\end{longtable}
}

\protect\phantomsection\label{docs__pandoc__latex__src__cli_complete_curriculum.md__cmd-troubleshooting}{}

\subsection*{Troubleshooting}\label{docs__pandoc__latex__src__cli_complete_curriculum.md__troubleshooting-2}

\subsubsection*{Problem: "I Can\textquotesingle t Hear the Output"}\label{docs__pandoc__latex__src__cli_complete_curriculum.md__problem-i-cant-hear-the-output-1}

\begin{enumerate}
\tightlist
\item
  Redirect to file and open in Notepad.
\item
  Use End to jump to the end of text.
\end{enumerate}

\subsubsection*{Problem: "Tab Completion Isn\textquotesingle t Working"}\label{docs__pandoc__latex__src__cli_complete_curriculum.md__problem-tab-completion-isnt-working-1}

\begin{enumerate}
\tightlist
\item
  Type at least one character before Tab.
\end{enumerate}

\subsubsection*{Problem: "Command Not Found"}\label{docs__pandoc__latex__src__cli_complete_curriculum.md__problem-command-not-found-1}

\begin{enumerate}
\tightlist
\item
  Use \texttt{where\ programname} to find installed programs.
\end{enumerate}

\subsection*{Pro Tips}\label{docs__pandoc__latex__src__cli_complete_curriculum.md__pro-tips-1}

\begin{enumerate}
\tightlist
\item
  Use \texttt{dir\ /b} for one-per-line listings.
\item
  Create a personal notes file and open it in Notepad for quick reference.
\end{enumerate}

\subsection*{Recommended Workflow}\label{docs__pandoc__latex__src__cli_complete_curriculum.md__recommended-workflow-1}

\begin{enumerate}
\tightlist
\item
  \texttt{cd} to the project folder
\item
  \texttt{dir\ /b} to list files
\item
  Redirect large output to files and open in Notepad
\end{enumerate}

\subsection*{Additional Resources}\label{docs__pandoc__latex__src__cli_complete_curriculum.md__additional-resources-1}

\begin{itemize}
\tightlist
\item
  NVDA Documentation: \url{https://www.nvaccess.org/documentation/}
\item
  JAWS Documentation: \url{https://www.freedomscientific.com/support/}
\item
  Windows CMD Reference: \url{https://docs.microsoft.com/windows-server/administration/windows-commands/windows-commands}
\item
  NVDA Documentation: \url{https://www.nvaccess.org/documentation/}
\item
  JAWS Documentation: \url{https://www.freedomscientific.com/support/}
\item
  Dolphin SuperNova: \url{https://yourdolphin.com/supernova/}
\item
  Windows Narrator: \url{https://support.microsoft.com/narrator}
\end{itemize}

\begin{center}\rule{0.5\linewidth}{0.5pt}\end{center}

\protect\phantomsection\label{docs__pandoc__latex__src__cli_complete_curriculum.md__cmd-lesson-pre}{}

\section*{Lesson Pre: Your First Terminal --- Screen Reader Navigation Fundamentals}\label{docs__pandoc__latex__src__cli_complete_curriculum.md__lesson-pre-your-first-terminal--screen-reader-navigation-fundamentals}

\begin{center}\rule{0.5\linewidth}{0.5pt}\end{center}

\textbf{Duration:} 1.5-2 hours (for screen reader users)\\
\textbf{Prerequisites:} None --- this is the starting point\\
\textbf{Accessibility Note:} This lesson is designed specifically for screen reader users. Tips are provided for NVDA, JAWS, Windows Narrator, and Dolphin SuperNova.

\begin{center}\rule{0.5\linewidth}{0.5pt}\end{center}

\subsection*{What is a Terminal?}\label{docs__pandoc__latex__src__cli_complete_curriculum.md__what-is-a-terminal-1}

A terminal (also called command line or Command Prompt) is a text-based interface where you type commands instead of clicking buttons. Think of it like sending written instructions to your computer instead of pointing and clicking.

\textbf{Why learn this?}

\begin{itemize}
\tightlist
\item
  Faster and more precise work, especially for 3D printing scripts and automation.
\item
  Essential for using tools like OpenSCAD.
\item
  Accessibility: command-line tools work reliably with screen readers --- output is plain text.
\item
  Simple automation without complex graphical interfaces.
\end{itemize}

\begin{center}\rule{0.5\linewidth}{0.5pt}\end{center}

\subsection*{Opening Command Prompt for the First Time}\label{docs__pandoc__latex__src__cli_complete_curriculum.md__opening-command-prompt-for-the-first-time}

\subsubsection*{3DMake: Your First Real-World Goal}\label{docs__pandoc__latex__src__cli_complete_curriculum.md__3dmake-your-first-real-world-goal-1}

\textbf{3DMake} (\href{https://github.com/tdeck/3dmake}{github.com/tdeck/3dmake}) is a free, accessibility-first command-line tool for designing and 3D-printing models entirely from the terminal --- no mouse required. Every CMD skill in this curriculum directly supports a real 3DMake workflow.

Here is the end goal --- one command that builds your OpenSCAD model, slices it, and sends it to the printer:

\begin{lstlisting}[style=Alabaster, language=cmd]
REM Navigate to your project folder
cd %USERPROFILE%\Documents\3d_projects\my_cube

REM Build, slice, and print in a single command
3dm build slice print

\end{lstlisting}

Keep this goal in mind as you work through each lesson.

\begin{center}\rule{0.5\linewidth}{0.5pt}\end{center}

\subsection*{Opening Command Prompt for the First Time}\label{docs__pandoc__latex__src__cli_complete_curriculum.md__opening-command-prompt-for-the-first-time-1}

\subsubsection*{Method 1: Search (Easiest)}\label{docs__pandoc__latex__src__cli_complete_curriculum.md__method-1-search-easiest}

\begin{enumerate}
\tightlist
\item
  Press the \textbf{Windows key} alone.
\item
  Type: \texttt{cmd}
\item
  You will hear search results appear.
\item
  Press \textbf{Enter} to open the first result (Command Prompt).
\end{enumerate}

\subsubsection*{Method 2: Using the Run Dialog}\label{docs__pandoc__latex__src__cli_complete_curriculum.md__method-2-using-the-run-dialog}

\begin{enumerate}
\tightlist
\item
  Press \textbf{Windows key + R} (opens the Run dialog).
\item
  Type: \texttt{cmd}
\item
  Press \textbf{Enter}.
\end{enumerate}

\subsubsection*{Method 3: From the Start Menu}\label{docs__pandoc__latex__src__cli_complete_curriculum.md__method-3-from-the-start-menu}

\begin{enumerate}
\tightlist
\item
  Press the \textbf{Windows key}.
\item
  Navigate to \textbf{Windows System} or \textbf{Windows Tools} (depending on your Windows version).
\item
  Open \textbf{Command Prompt}.
\end{enumerate}

\subsubsection*{First Connection: Understanding the Prompt}\label{docs__pandoc__latex__src__cli_complete_curriculum.md__first-connection-understanding-the-prompt-1}

When Command Prompt opens, your screen reader will announce the window title and then the \textbf{prompt}. The prompt is where you type commands.

\textbf{What you will hear:}

\begin{lstlisting}[style=Alabaster]
C:\Users\YourName>

\end{lstlisting}

\textbf{What this means:}

\begin{itemize}
\tightlist
\item
  \texttt{C:\textbackslash{}Users\textbackslash{}YourName} = Your current location (the path).
\item
  \texttt{\textgreater{}} = The prompt is ready for input.
\end{itemize}

Your cursor is right after the \texttt{\textgreater{}}. This is where you type.

\begin{center}\rule{0.5\linewidth}{0.5pt}\end{center}

\subsection*{Your First Commands (Screen Reader Edition)}\label{docs__pandoc__latex__src__cli_complete_curriculum.md__your-first-commands-screen-reader-edition-1}

\subsubsection*{\texorpdfstring{Command 1: "Where Am I?" --- \texttt{cd}}{Command 1: "Where Am I?" --- cd}}\label{docs__pandoc__latex__src__cli_complete_curriculum.md__command-1-where-am-i--cd}

Running \texttt{cd} with no arguments shows your current directory.

\textbf{Type this and press Enter:}

\begin{lstlisting}[style=Alabaster, language=cmd]
cd

\end{lstlisting}

\textbf{What you will hear:}
Your screen reader announces the current path, for example:

\begin{lstlisting}[style=Alabaster]
C:\Users\YourName

\end{lstlisting}

\textbf{Understanding paths:}

\begin{itemize}
\tightlist
\item
  A path shows your location in the file system, like a mailing address.
\item
  Windows paths use backslashes: \texttt{C:\textbackslash{}Users\textbackslash{}YourName\textbackslash{}Documents}
\item
  Think of it as nested folders: \texttt{C:\textbackslash{}} (the main drive) \textgreater{} \texttt{Users} \textgreater{} \texttt{YourName} \textgreater{} \texttt{Documents}
\end{itemize}

\subsubsection*{\texorpdfstring{Command 2: "What\textquotesingle s Here?" --- \texttt{dir}}{Command 2: "What\textquotesingle s Here?" --- dir}}\label{docs__pandoc__latex__src__cli_complete_curriculum.md__command-2-whats-here--dir}

\textbf{Type this and press Enter:}

\begin{lstlisting}[style=Alabaster, language=cmd]
dir /B

\end{lstlisting}

\textbf{What you will hear:}
Your screen reader announces folder and file names, one per line. The \texttt{/B} flag gives a clean "bare" listing --- names only, no dates or sizes --- which is much easier to follow with a screen reader.

If you want more detail (sizes, dates), run \texttt{dir} without \texttt{/B}. But for navigation, \texttt{dir\ /B} is preferred.

\subsubsection*{\texorpdfstring{Command 3: "Go There" --- \texttt{cd\ Documents}}{Command 3: "Go There" --- cd Documents}}\label{docs__pandoc__latex__src__cli_complete_curriculum.md__command-3-go-there--cd-documents}

\textbf{Type this and press Enter:}

\begin{lstlisting}[style=Alabaster, language=cmd]
cd Documents

\end{lstlisting}

\textbf{What you will hear:}
The prompt changes to show your new location:

\begin{lstlisting}[style=Alabaster]
C:\Users\YourName\Documents>

\end{lstlisting}

\textbf{Practice navigation:}

\begin{enumerate}
\tightlist
\item
  Run \texttt{cd} to confirm you are in Documents.
\item
  Run \texttt{dir\ /B} to see what is there.
\item
  Go back up: \texttt{cd\ ..} (the \texttt{..} means "go up one level").
\item
  Run \texttt{cd} again to confirm.
\item
  Return to Documents: \texttt{cd\ Documents}
\end{enumerate}

\begin{center}\rule{0.5\linewidth}{0.5pt}\end{center}

\subsection*{Reading Screen Reader Output (Critical Skills)}\label{docs__pandoc__latex__src__cli_complete_curriculum.md__reading-screen-reader-output-critical-skills-1}

\subsubsection*{Dealing with Long Lists}\label{docs__pandoc__latex__src__cli_complete_curriculum.md__dealing-with-long-lists-1}

When you run \texttt{dir} in a folder with many files, the list may be very long. Here is how to manage it:

\textbf{Solution 1: Save to a File}

\begin{lstlisting}[style=Alabaster, language=cmd]
dir /B > list.txt
notepad.exe list.txt

\end{lstlisting}

This saves the listing to a file and opens it in Notepad, where you can read it calmly with your screen reader.

\textbf{Solution 2: Use Pause}

\begin{lstlisting}[style=Alabaster, language=cmd]
dir /B | more

\end{lstlisting}

Shows output one page at a time. Press \textbf{Space} to go to the next page, \textbf{Q} to quit. Note: some screen readers read \texttt{more} output less reliably than a saved file --- if this is difficult, use Solution 1.

\subsubsection*{Tab Completion}\label{docs__pandoc__latex__src__cli_complete_curriculum.md__tab-completion}

Tab completion is one of the most powerful screen reader techniques in the terminal.

\textbf{How it works:}

\begin{enumerate}
\tightlist
\item
  Type the first few letters of a folder or file name.
\item
  Press \textbf{Tab}.
\item
  Command Prompt automatically completes the rest.
\end{enumerate}

\textbf{Example:}

\begin{enumerate}
\tightlist
\item
  You are at \texttt{C:\textbackslash{}Users\textbackslash{}YourName\textgreater{}}
\item
  Type: \texttt{cd\ Doc}
\item
  Press \textbf{Tab}
\item
  Command Prompt completes it to: \texttt{cd\ Documents}
\item
  Press \textbf{Enter} to navigate there.
\end{enumerate}

With a screen reader, when you press Tab the screen reader announces the completed command --- much faster and more accurate than typing the whole thing.

\begin{center}\rule{0.5\linewidth}{0.5pt}\end{center}

\subsection*{Creating and Viewing Files}\label{docs__pandoc__latex__src__cli_complete_curriculum.md__creating-and-viewing-files}

\subsubsection*{Create a Simple File}\label{docs__pandoc__latex__src__cli_complete_curriculum.md__create-a-simple-file-1}

\begin{lstlisting}[style=Alabaster, language=cmd]
echo Hello, Command Prompt! > hello.txt

\end{lstlisting}

\begin{itemize}
\tightlist
\item
  \texttt{echo} outputs text.
\item
  \texttt{\textgreater{}} redirects it to a file called \texttt{hello.txt}.
\end{itemize}

\subsubsection*{Read the File Back}\label{docs__pandoc__latex__src__cli_complete_curriculum.md__read-the-file-back-1}

\begin{lstlisting}[style=Alabaster, language=cmd]
type hello.txt

\end{lstlisting}

Your screen reader announces the file contents.

\subsubsection*{Open and Edit the File}\label{docs__pandoc__latex__src__cli_complete_curriculum.md__open-and-edit-the-file-1}

\begin{lstlisting}[style=Alabaster, language=cmd]
notepad.exe hello.txt

\end{lstlisting}

Opens the file in Notepad for comfortable editing with your screen reader.

\begin{center}\rule{0.5\linewidth}{0.5pt}\end{center}

\subsection*{Essential Keyboard Shortcuts}\label{docs__pandoc__latex__src__cli_complete_curriculum.md__essential-keyboard-shortcuts-1}

{\def\LTcaptype{none} % do not increment counter
\begin{longtable}[]{@{}
  >{\raggedright\arraybackslash}p{(\linewidth - 2\tabcolsep) * \real{0.2099}}
  >{\raggedright\arraybackslash}p{(\linewidth - 2\tabcolsep) * \real{0.7901}}@{}}
\toprule\noalign{}
\begin{minipage}[b]{\linewidth}\raggedright
Key Combination
\end{minipage} & \begin{minipage}[b]{\linewidth}\raggedright
What It Does
\end{minipage} \\
\midrule\noalign{}
\endhead
\bottomrule\noalign{}
\endlastfoot
\textbf{Up Arrow} &
Shows your previous command (press again for earlier commands) \\
\textbf{Down Arrow} & Shows a more recent command in history \\
\textbf{Tab} & Auto-completes folder and file names \\
\textbf{Ctrl+C} & Stops a running command \\
\textbf{cls} & Clears the screen (type this and press Enter) \\
\textbf{Enter} & Runs the command you have typed \\
\textbf{Home} & Moves cursor to start of command line \\
\textbf{End} & Moves cursor to end of command line \\
\end{longtable}
}

\begin{center}\rule{0.5\linewidth}{0.5pt}\end{center}

\subsection*{Screen Reader-Specific Tips}\label{docs__pandoc__latex__src__cli_complete_curriculum.md__screen-reader-specific-tips-1}

\subsubsection*{NVDA Users}\label{docs__pandoc__latex__src__cli_complete_curriculum.md__nvda-users-1}

NVDA is available free from \url{https://www.nvaccess.org/}

\begin{enumerate}
\tightlist
\item
  \textbf{After running a command}, press \textbf{NVDA+Home} (Insert+Home or CapsLock+Home) to read the current line.
\item
  \textbf{To read all output}, press \textbf{NVDA+Down Arrow} to read from the cursor to the end of the screen.
\item
  \textbf{To browse previous output}, use the NVDA review cursor: press \textbf{Numpad 8/2} (desktop) or the equivalent laptop keys to move up and down through lines.
\item
  \textbf{If output scrolled past}, redirect to a file: \texttt{command\ \textgreater{}\ output.txt}, then \texttt{notepad.exe\ output.txt}.
\item
  \textbf{To open NVDA settings}: press \textbf{NVDA+N} (opens NVDA menu), then go to Preferences.
\end{enumerate}

\begin{quote}
\textbf{NVDA key note:} The NVDA modifier key is \textbf{Insert} by default on desktops. On laptops it is often \textbf{CapsLock}. Check your NVDA settings if these shortcuts do not work as described.
\end{quote}

\subsubsection*{JAWS Users}\label{docs__pandoc__latex__src__cli_complete_curriculum.md__jaws-users-1}

JAWS is available from \url{https://www.freedomscientific.com/products/software/jaws/}

\begin{enumerate}
\tightlist
\item
  \textbf{After running a command}, press \textbf{Insert+Up Arrow} to re-read the current line.
\item
  \textbf{To read all output}, press \textbf{Insert+Ctrl+Down} to read to the end of the screen.
\item
  \textbf{To move through output line by line}, press \textbf{Insert+Down Arrow} repeatedly.
\item
  \textbf{If output is too long}, redirect to a file: \texttt{command\ \textgreater{}\ output.txt}, then \texttt{notepad.exe\ output.txt}.
\item
  \textbf{JAWS demo mode reminder:} If you are using JAWS without a license, sessions are limited to approximately 40 minutes. Save your work and restart the session if JAWS stops responding.
\end{enumerate}

\begin{quote}
\textbf{JAWS key note:} The JAWS modifier key is \textbf{Insert} by default. On some laptops it can be set to \textbf{CapsLock} in JAWS Settings Center.
\end{quote}

\subsubsection*{Windows Narrator Users}\label{docs__pandoc__latex__src__cli_complete_curriculum.md__windows-narrator-users}

Windows Narrator is built into Windows 10 and Windows 11. Enable it with \textbf{Windows+Ctrl+Enter}.

\begin{enumerate}
\tightlist
\item
  \textbf{After running a command}, press \textbf{Narrator+D} (CapsLock+D or Insert+D) to read the current line.
\item
  \textbf{To read all output from here}, press \textbf{Narrator+R} to read from the cursor position.
\item
  \textbf{To move through output}, use \textbf{Narrator+Up/Down Arrow} in scan mode (press \textbf{Narrator+Space} to toggle scan mode on).
\item
  \textbf{For long outputs}, always redirect to a file: \texttt{command\ \textgreater{}\ output.txt}, then \texttt{notepad.exe\ output.txt}. Notepad is more comfortable for Narrator users than reading directly from the terminal.
\item
  \textbf{Narrator settings}: Windows key \textgreater{} Settings \textgreater{} Accessibility \textgreater{} Narrator.
\end{enumerate}

\begin{quote}
\textbf{Narrator key note:} The Narrator modifier key is \textbf{CapsLock} or \textbf{Insert}. You can change it in Settings \textgreater{} Accessibility \textgreater{} Narrator \textgreater{} Keyboard shortcuts.
\end{quote}

\subsubsection*{Dolphin SuperNova Users}\label{docs__pandoc__latex__src__cli_complete_curriculum.md__dolphin-supernova-users}

Dolphin SuperNova is available from \url{https://yourdolphin.com/supernova/}

\begin{enumerate}
\tightlist
\item
  \textbf{After running a command}, press \textbf{CapsLock+L} to read the current line.
\item
  \textbf{To read all output}, press \textbf{CapsLock+Numpad Plus} (say all from current position).
\item
  \textbf{To move through output}, press \textbf{CapsLock+Numpad 8} (up) and \textbf{CapsLock+Numpad 2} (down) to move line by line through the review buffer.
\item
  \textbf{For long outputs}, redirect to a file: \texttt{command\ \textgreater{}\ output.txt}, then \texttt{notepad.exe\ output.txt}.
\item
  \textbf{SuperNova settings}: press \textbf{CapsLock+SpaceBar} to open the SuperNova Control Panel.
\end{enumerate}

\begin{quote}
\textbf{Dolphin key note:} The SuperNova modifier key is \textbf{CapsLock} by default. This can be changed in the SuperNova Control Panel \textgreater{} Keyboard.
\end{quote}

\begin{center}\rule{0.5\linewidth}{0.5pt}\end{center}

\subsection*{Common Issue: "I Can\textquotesingle t Hear the Output"}\label{docs__pandoc__latex__src__cli_complete_curriculum.md__common-issue-i-cant-hear-the-output-1}

This happens to everyone starting out. Here are the solutions:

\begin{enumerate}
\tightlist
\item
  \textbf{Make sure your screen reader was running before you opened Command Prompt.} If not, close Command Prompt, confirm your screen reader is running, then reopen it.
\item
  \textbf{Press End} to make sure your cursor is at the prompt.
\item
  \textbf{Redirect to a file}: \texttt{command\ \textgreater{}\ output.txt}, then \texttt{notepad.exe\ output.txt}. This always works with all screen readers.
\item
  \textbf{Use the screen reader\textquotesingle s review commands} listed above to scan back through previous output.
\end{enumerate}

\begin{center}\rule{0.5\linewidth}{0.5pt}\end{center}

\subsection*{Practice Exercises}\label{docs__pandoc__latex__src__cli_complete_curriculum.md__practice-exercises-2}

Complete these in order. Take your time.

\subsubsection*{Exercise 1: Basic Navigation}\label{docs__pandoc__latex__src__cli_complete_curriculum.md__exercise-1-basic-navigation-1}

\begin{enumerate}
\tightlist
\item
  Open Command Prompt.
\item
  Run \texttt{cd} and note your location.
\item
  Run \texttt{dir\ /B} and listen to the listing.
\item
  Type \texttt{cd\ Documents} and press Enter.
\item
  Run \texttt{cd} to confirm your new location.
\item
  Run \texttt{dir\ /B} in this new location.
\end{enumerate}

\textbf{Goal:} You know where you are and what is around you.

\subsubsection*{Exercise 2: Using Tab Completion}\label{docs__pandoc__latex__src__cli_complete_curriculum.md__exercise-2-using-tab-completion-1}

\begin{enumerate}
\tightlist
\item
  In your home directory, type \texttt{cd\ D} (just those two characters).
\item
  Press \textbf{Tab}.
\item
  Command Prompt should auto-complete to a folder starting with D.
\item
  Repeat with other folder names.
\item
  Try \texttt{cd\ Down} and Tab to \texttt{Downloads}.
\end{enumerate}

\textbf{Goal:} Tab completion feels natural.

\subsubsection*{Exercise 3: Creating and Viewing Files}\label{docs__pandoc__latex__src__cli_complete_curriculum.md__exercise-3-creating-and-viewing-files-1}

\begin{enumerate}
\tightlist
\item
  Create a file: \texttt{echo\ Test\ content\ \textgreater{}\ test.txt}
\item
  View it: \texttt{type\ test.txt}
\item
  Create another: \texttt{echo\ Line\ 2\ \textgreater{}\ another.txt}
\item
  List both: \texttt{dir\ /B\ *.txt}
\end{enumerate}

\textbf{Goal:} You understand create, view, and list operations.

\subsubsection*{Exercise 4: Going Up Levels}\label{docs__pandoc__latex__src__cli_complete_curriculum.md__exercise-4-going-up-levels-1}

\begin{enumerate}
\tightlist
\item
  Navigate into several folders: \texttt{cd\ Documents}, then \texttt{cd\ folder1}, etc.
\item
  From inside, run \texttt{cd\ ..} multiple times to come back up.
\item
  After each \texttt{cd\ ..}, run \texttt{cd} to confirm your location.
\end{enumerate}

\textbf{Goal:} You are comfortable with relative navigation using \texttt{..}

\subsubsection*{Exercise 5: Redirecting Output}\label{docs__pandoc__latex__src__cli_complete_curriculum.md__exercise-5-redirecting-output-1}

\begin{enumerate}
\tightlist
\item
  Create a listing: \texttt{dir\ /B\ \textgreater{}\ directorylist.txt}
\item
  Open it: \texttt{notepad.exe\ directorylist.txt}
\item
  Read it with your screen reader.
\item
  Close Notepad.
\end{enumerate}

\textbf{Goal:} You can save long outputs to files for easier reading.

\begin{center}\rule{0.5\linewidth}{0.5pt}\end{center}

\subsection*{Checkpoint Questions}\label{docs__pandoc__latex__src__cli_complete_curriculum.md__checkpoint-questions-1}

After this lesson, you should be able to answer all ten with confidence:

\begin{enumerate}
\tightlist
\item
  What does \texttt{cd} do with no arguments?
\item
  What does \texttt{dir\ /B} do?
\item
  Why do we use \texttt{/B} with \texttt{dir}?
\item
  What is your current path right now?
\item
  How do you navigate to a new folder?
\item
  How do you go up one level?
\item
  What is the Tab key for?
\item
  What does \texttt{echo\ text\ \textgreater{}\ file.txt} do?
\item
  How do you read a file back in the terminal?
\item
  How do you stop a command that is running?
\end{enumerate}

\textbf{Answer all 10 before moving to Command Line Lesson 0.}

\begin{center}\rule{0.5\linewidth}{0.5pt}\end{center}

\subsection*{Common Questions}\label{docs__pandoc__latex__src__cli_complete_curriculum.md__common-questions-1}

\textbf{Q: Should I use Command Prompt or PowerShell?}\\
A: Command Prompt is simpler and a good starting point. PowerShell is more powerful. Both work well with all four screen readers covered in this curriculum. Start with whichever your curriculum pathway specifies.

\textbf{Q: Why is my screen reader not reading the output?}\\
A: This is common. Use \texttt{command\ \textgreater{}\ file.txt} to save output to a file, then open it with \texttt{notepad.exe\ file.txt} for reliable reading.

\textbf{Q: What if I type something wrong?}\\
A: Press \textbf{Enter} and you will see an error message. Type the correct command on the next line. No harm done.

\textbf{Q: How do I get help with a command?}\\
A: Type \texttt{help\ command-name} (for example, \texttt{help\ cd}). Or type \texttt{command\ /?} for built-in help (for example, \texttt{dir\ /?}).

\begin{center}\rule{0.5\linewidth}{0.5pt}\end{center}

\subsection*{Next Steps}\label{docs__pandoc__latex__src__cli_complete_curriculum.md__next-steps-1}

Once comfortable with these basics, move to \textbf{Command Line Lesson 0: Getting Started} for deeper path understanding, then continue through Command Line Lesson 1 to Command Line Lesson 5 for full terminal mastery.

\begin{center}\rule{0.5\linewidth}{0.5pt}\end{center}

\protect\phantomsection\label{docs__pandoc__latex__src__cli_complete_curriculum.md__ps-troubleshooting}{}

\subsection*{Troubleshooting}\label{docs__pandoc__latex__src__cli_complete_curriculum.md__troubleshooting-3}

{\def\LTcaptype{none} % do not increment counter
\begin{longtable}[]{@{}
  >{\raggedright\arraybackslash}p{(\linewidth - 2\tabcolsep) * \real{0.2991}}
  >{\raggedright\arraybackslash}p{(\linewidth - 2\tabcolsep) * \real{0.7009}}@{}}
\toprule\noalign{}
\begin{minipage}[b]{\linewidth}\raggedright
Issue
\end{minipage} & \begin{minipage}[b]{\linewidth}\raggedright
Solution
\end{minipage} \\
\midrule\noalign{}
\endhead
\bottomrule\noalign{}
\endlastfoot
Command Prompt won\textquotesingle t open &
Press Windows+R, type \texttt{cmd}, press Enter \\
Cannot hear the output &
Redirect to a file: \texttt{command\ \textgreater{}\ output.txt}, then \texttt{notepad.exe\ output.txt} \\
Tab completion not working &
Type at least one character before pressing Tab \\
Command not found &
Check spelling; try \texttt{help} for a list of available commands \\
Stuck in a command & Press \textbf{Ctrl+C} to stop it \\
JAWS stops working mid-session &
If using JAWS demo: restart computer, demo sessions last \textasciitilde{}40 minutes \\
\end{longtable}
}

\begin{center}\rule{0.5\linewidth}{0.5pt}\end{center}

\subsection*{Resources}\label{docs__pandoc__latex__src__cli_complete_curriculum.md__resources-1}

\begin{itemize}
\tightlist
\item
  \textbf{NVDA:} \url{https://www.nvaccess.org/}
\item
  \textbf{JAWS:} \url{https://www.freedomscientific.com/products/software/jaws/}
\item
  \textbf{Windows Narrator guide:} \url{https://support.microsoft.com/narrator}
\item
  \textbf{Dolphin SuperNova:} \url{https://yourdolphin.com/supernova/}
\item
  \textbf{Windows CMD Reference:} \url{https://example.com}
\end{itemize}

\begin{center}\rule{0.5\linewidth}{0.5pt}\end{center}

\protect\phantomsection\label{docs__pandoc__latex__src__cli_complete_curriculum.md__cmd-lesson-0}{}

\section*{Lesson 0: Getting Started --- Layout, Paths, and the Shell}\label{docs__pandoc__latex__src__cli_complete_curriculum.md__lesson-0-getting-started--layout-paths-and-the-shell-1}

\begin{center}\rule{0.5\linewidth}{0.5pt}\end{center}

Estimated time: 20-30 minutes

\subsection*{Learning Objectives}\label{docs__pandoc__latex__src__cli_complete_curriculum.md__learning-objectives-3}

\begin{itemize}
\tightlist
\item
  Launch Command Prompt and locate the prompt
\item
  Understand path notation and shortcuts (\texttt{C:\textbackslash{}}, \texttt{..}, \texttt{.})
\item
  Use tab completion to navigate quickly
\end{itemize}

\subsection*{Materials}\label{docs__pandoc__latex__src__cli_complete_curriculum.md__materials-1}

\begin{itemize}
\tightlist
\item
  Computer with Windows
\item
  Editor (Notepad/VS Code)
\end{itemize}

\subsection*{Step-by-step Tasks}\label{docs__pandoc__latex__src__cli_complete_curriculum.md__step-by-step-tasks-2}

\begin{enumerate}
\tightlist
\item
  Open Command Prompt and note the prompt (it includes the current path).
\item
  Run \texttt{cd} and say or note the printed path.
\item
  Use \texttt{dir\ /B} to list names in your home directory.
\item
  Practice \texttt{cd\ Documents}, \texttt{cd\ ..} and \texttt{cd\ \textbackslash{}} until comfortable.
\item
  Try tab-completion: type \texttt{cd\ D} and press Tab.
\end{enumerate}

\subsection*{Checkpoints}\label{docs__pandoc__latex__src__cli_complete_curriculum.md__checkpoints-3}

\begin{itemize}
\tightlist
\item
  Confirm you can state your current path and move to \texttt{Documents}.
\end{itemize}

\subsection*{3DMake Quick-Start: Your First Real Project}\label{docs__pandoc__latex__src__cli_complete_curriculum.md__3dmake-quick-start-your-first-real-project-1}

You now know enough CMD to set up a real 3DMake project. Do this now --- it is the fastest path to a printed object.

\textbf{Step 1 --- Download and run the 3DMake setup wizard:}

\begin{lstlisting}[style=Alabaster, language=cmd]
REM Go to your Downloads folder (where you extracted 3dmake_windows.zip)
cd %USERPROFILE%\Downloads\3dmake

REM Run the guided setup (answer the questions about your printer)
3dm setup

\end{lstlisting}

\textbf{Step 2 --- Create your first project:}

\begin{lstlisting}[style=Alabaster, language=cmd]
REM Create a folder for all your 3D projects
mkdir %USERPROFILE%\Documents\3d_projects

REM Create your first project folder and initialize it
mkdir %USERPROFILE%\Documents\3d_projects\my_cube
cd %USERPROFILE%\Documents\3d_projects\my_cube
3dm new

\end{lstlisting}

\textbf{Step 3 --- Confirm the project structure:}

\begin{lstlisting}[style=Alabaster, language=cmd]
REM You should see three items: src, build, and 3dmake.toml
dir /B

\end{lstlisting}

You just used \texttt{mkdir}, \texttt{cd}, and \texttt{dir\ /B} on a real project. You are ready to print.

\subsection*{Checkpoints}\label{docs__pandoc__latex__src__cli_complete_curriculum.md__checkpoints-4}

\begin{itemize}
\tightlist
\item
  Confirm you can state your current path and move to \texttt{Documents}.
\end{itemize}

\subsection*{Quiz - Lesson CMD.0}\label{docs__pandoc__latex__src__cli_complete_curriculum.md__quiz---lesson-cmd0}

\begin{enumerate}
\tightlist
\item
  What is a path?
\item
  What does \texttt{..} mean?
\item
  How do you autocomplete a path?
\item
  How do you go up one directory?
\item
  What command lists only names (\texttt{dir} flag)?
\item
  True or False: On Windows, CMD uses backslashes (\texttt{\textbackslash{}}) in paths, but forward slashes (\texttt{/}) are also accepted in some contexts.
\item
  Explain the difference between an absolute path and a relative path.
\item
  If you are in \texttt{C:\textbackslash{}Users\textbackslash{}YourName\textbackslash{}Documents} and you type \texttt{cd\ ..}, where do you end up?
\item
  What happens when you press Tab while typing a folder name in Command Prompt?
\item
  Describe a practical reason why understanding paths is important for a 3D printing workflow.
\item
  What does \texttt{C:\textbackslash{}} mean in a path, and when would you use it?
\item
  If a folder path contains spaces (e.g., \texttt{Program\ Files}), how do you navigate to it with \texttt{cd}?
\item
  Explain what the prompt \texttt{C:\textbackslash{}Users\textbackslash{}YourName\textgreater{}} tells you about your current state.
\item
  How would you navigate to your home directory from any location using a single command?
\item
  What is the advantage of using relative paths (like \texttt{..}) versus absolute paths in automation scripts?
\end{enumerate}

\subsection*{Extension Problems}\label{docs__pandoc__latex__src__cli_complete_curriculum.md__extension-problems-7}

\begin{enumerate}
\tightlist
\item
  Create a nested folder and practice \texttt{cd} into it by typing partial names and using Tab.
\item
  Use \texttt{dir\ /B\ /A:F} to list only files in a folder.
\item
  Save \texttt{cd} output (the path) to a file and open it in Notepad.
\item
  Try \texttt{cd} into a folder whose name contains spaces; observe how quotes are handled.
\item
  Create a short note file and open it from Command Prompt.
\item
  Build a folder structure that mirrors your project organization; navigate to each level and document the path.
\item
  Use \texttt{echo\ \%cd\%} to print your current path and save it to a file.
\item
  Investigate special paths (e.g., \texttt{\%USERPROFILE\%}, \texttt{\%TEMP\%}); write down what each contains and when you\textquotesingle d use them.
\item
  Compare absolute vs. relative paths by navigating to the same folder using each method; explain which is easier for automation.
\item
  Create a batch file that changes to a frequently-used folder and lists its contents in one command; test it from different starting locations.
\item
  Navigate to three different locations and at each one note the prompt, the path from \texttt{cd}, and verify you understand what each shows.
\item
  Create a complex folder tree (at least 5 levels deep) and navigate it using only relative paths; verify your location at each step.
\item
  Document all path shortcuts you know (\texttt{C:\textbackslash{}}, \texttt{..}, \texttt{.}) and demonstrate each one works as expected.
\item
  Write a guide for a peer on how to understand the Command Prompt and path notation without using GUI file explorer.
\item
  Create a troubleshooting flowchart: if someone says "I don\textquotesingle t know where I am," what commands do you give them to find out?
\end{enumerate}

\subsection*{References}\label{docs__pandoc__latex__src__cli_complete_curriculum.md__references-7}

\begin{itemize}
\tightlist
\item
  Microsoft. (2024). \emph{Command line reference}. \url{https://learn.microsoft.com/en-us/windows-server/administration/windows-commands/}
\item
  Microsoft. (2024). \emph{Using the cd command}. \url{https://learn.microsoft.com/en-us/windows-server/administration/windows-commands/cd_1}
\end{itemize}

\subsection*{Helpful Resources}\label{docs__pandoc__latex__src__cli_complete_curriculum.md__helpful-resources-6}

\begin{itemize}
\tightlist
\item
  \href{https://docs.microsoft.com/en-us/windows-server/administration/windows-commands/windows-commands-glossary}{Windows Command Reference}
\item
  \href{https://docs.microsoft.com/en-us/windows-server/administration/windows-commands/cd}{CD Command Guide}
\item
  \href{https://docs.microsoft.com/en-us/windows-server/administration/windows-commands/dir}{DIR Command Reference}
\item
  \href{https://superuser.com/questions/167461/what-are-path-separators-on-windows}{Understanding Windows Paths}
\end{itemize}

\begin{center}\rule{0.5\linewidth}{0.5pt}\end{center}

\protect\phantomsection\label{docs__pandoc__latex__src__cli_complete_curriculum.md__cmd-lesson-1}{}

\section*{Lesson 1: Navigation --- Moving Around Your File System}\label{docs__pandoc__latex__src__cli_complete_curriculum.md__lesson-1-navigation--moving-around-your-file-system-1}

\begin{center}\rule{0.5\linewidth}{0.5pt}\end{center}

\textbf{Duration:} 1.5-2 hours (for screen reader users)\\
\textbf{Prerequisites:} Command Line Lesson Pre and Command Line Lesson 0\\
\textbf{Learning Objectives:}

\begin{itemize}
\tightlist
\item
  Navigate confidently to any folder location
\item
  Understand and use relative versus absolute paths
\item
  Use shortcuts to move between common folders
\item
  Organize your file system logically
\end{itemize}

\begin{center}\rule{0.5\linewidth}{0.5pt}\end{center}

\subsection*{Core Navigation Commands}\label{docs__pandoc__latex__src__cli_complete_curriculum.md__core-navigation-commands}

\subsubsection*{\texorpdfstring{Command: \texttt{cd} (Change Directory)}{Command: cd (Change Directory)}}\label{docs__pandoc__latex__src__cli_complete_curriculum.md__command-cd-change-directory}

\begin{lstlisting}[style=Alabaster, language=cmd]
cd FolderName          :: Go into a folder
cd ..                  :: Go up one level
cd \                   :: Go to root of current drive (e.g., C:\)
cd %USERPROFILE%       :: Go to home directory

\end{lstlisting}

\subsubsection*{\texorpdfstring{Command: \texttt{dir\ /B} (List Contents)}{Command: dir /B (List Contents)}}\label{docs__pandoc__latex__src__cli_complete_curriculum.md__command-dir-b-list-contents}

\begin{lstlisting}[style=Alabaster, language=cmd]
dir /B                 :: List all names (one per line, screen-reader friendly)
dir /B /A:D            :: List only directories/folders
dir /B /A:F            :: List only files

\end{lstlisting}

\subsubsection*{Command: Show Current Location}\label{docs__pandoc__latex__src__cli_complete_curriculum.md__command-show-current-location}

\begin{lstlisting}[style=Alabaster, language=cmd]
cd                     :: Running cd with no arguments shows current path

\end{lstlisting}

\begin{center}\rule{0.5\linewidth}{0.5pt}\end{center}

\subsection*{Understanding Paths}\label{docs__pandoc__latex__src__cli_complete_curriculum.md__understanding-paths}

\subsubsection*{Absolute Paths (Full Address)}\label{docs__pandoc__latex__src__cli_complete_curriculum.md__absolute-paths-full-address}

An absolute path starts from the root and shows the complete location:

\begin{lstlisting}[style=Alabaster, language=cmd]
C:\Users\YourName\Documents\3DPrinting

\end{lstlisting}

\begin{itemize}
\tightlist
\item
  \texttt{C:\textbackslash{}} --- The C drive root
\item
  \texttt{Users} --- First folder
\item
  \texttt{YourName} --- Your user folder
\item
  \texttt{Documents} --- Documents folder
\item
  \texttt{3DPrinting} --- Your 3D printing folder
\end{itemize}

\textbf{Navigate there directly:}

\begin{lstlisting}[style=Alabaster, language=cmd]
cd C:\Users\YourName\Documents\3DPrinting

\end{lstlisting}

\subsubsection*{Relative Paths (Directions From Here)}\label{docs__pandoc__latex__src__cli_complete_curriculum.md__relative-paths-directions-from-here}

\begin{lstlisting}[style=Alabaster, language=cmd]
cd Documents\3DPrinting    :: Go into Documents, then into 3DPrinting
cd ..                      :: Go up one level
cd ..\..                   :: Go up two levels

\end{lstlisting}

\textbf{Shortcuts:}

\begin{itemize}
\tightlist
\item
  \texttt{.} = Current folder
\item
  \texttt{..} = Parent folder (up one level)
\item
  \texttt{\%USERPROFILE\%} = Your home folder
\end{itemize}

\begin{center}\rule{0.5\linewidth}{0.5pt}\end{center}

\subsection*{Common Navigation Patterns}\label{docs__pandoc__latex__src__cli_complete_curriculum.md__common-navigation-patterns}

\subsubsection*{Pattern 1: Navigating Down}\label{docs__pandoc__latex__src__cli_complete_curriculum.md__pattern-1-navigating-down}

\begin{lstlisting}[style=Alabaster, language=cmd]
C:\>cd Users
C:\Users>cd YourName
C:\Users\YourName>cd Documents
C:\Users\YourName\Documents>

\end{lstlisting}

\subsubsection*{Pattern 2: Navigating Up}\label{docs__pandoc__latex__src__cli_complete_curriculum.md__pattern-2-navigating-up}

\begin{lstlisting}[style=Alabaster, language=cmd]
C:\Users\YourName\Documents>cd ..
C:\Users\YourName>cd ..
C:\Users>cd ..
C:\>

\end{lstlisting}

\subsubsection*{Pattern 3: Jump to a Known Path}\label{docs__pandoc__latex__src__cli_complete_curriculum.md__pattern-3-jump-to-a-known-path}

\begin{lstlisting}[style=Alabaster, language=cmd]
C:\Deep\Folder\Somewhere>cd C:\Users\YourName
C:\Users\YourName>

\end{lstlisting}

\subsubsection*{Pattern 4: Going Home}\label{docs__pandoc__latex__src__cli_complete_curriculum.md__pattern-4-going-home}

\begin{lstlisting}[style=Alabaster, language=cmd]
cd %USERPROFILE%

\end{lstlisting}

\begin{center}\rule{0.5\linewidth}{0.5pt}\end{center}

\subsection*{Special Folders and Shortcuts}\label{docs__pandoc__latex__src__cli_complete_curriculum.md__special-folders-and-shortcuts}

\begin{lstlisting}[style=Alabaster, language=cmd]
%USERPROFILE%              :: Your home folder
%USERPROFILE%\Desktop      :: Your Desktop
%USERPROFILE%\Documents    :: Your Documents
%USERPROFILE%\Downloads    :: Your Downloads
%ProgramFiles%             :: Program Files folder
C:\                        :: Root of C drive

\end{lstlisting}

\begin{center}\rule{0.5\linewidth}{0.5pt}\end{center}

\subsection*{Tab Completion (Essential Skill)}\label{docs__pandoc__latex__src__cli_complete_curriculum.md__tab-completion-essential-skill}

\begin{enumerate}
\tightlist
\item
  Type the first few characters of a folder name.
\item
  Press \textbf{Tab}.
\item
  Command Prompt completes it.
\end{enumerate}

\textbf{Example:}

\begin{lstlisting}[style=Alabaster, language=cmd]
C:\Users\YourName>cd Doc  [press Tab]
C:\Users\YourName>cd Documents

\end{lstlisting}

\begin{itemize}
\tightlist
\item
  Your screen reader announces the completed name immediately.
\item
  Press \textbf{Tab} multiple times to cycle through all matches.
\end{itemize}

\begin{center}\rule{0.5\linewidth}{0.5pt}\end{center}

\subsection*{Screen Reader Tips for Navigation}\label{docs__pandoc__latex__src__cli_complete_curriculum.md__screen-reader-tips-for-navigation}

\subsubsection*{NVDA Users}\label{docs__pandoc__latex__src__cli_complete_curriculum.md__nvda-users-2}

\begin{enumerate}
\tightlist
\item
  After each \texttt{cd} command, run \texttt{cd} alone to confirm location. Press \textbf{NVDA+Home} (Insert+Home) to read the current line.
\item
  Press \textbf{NVDA+Down Arrow} to read through \texttt{dir\ /B} output.
\item
  For long listings: \texttt{dir\ /B\ \textgreater{}\ listing.txt}, then \texttt{notepad.exe\ listing.txt}.
\end{enumerate}

\subsubsection*{JAWS Users}\label{docs__pandoc__latex__src__cli_complete_curriculum.md__jaws-users-2}

\begin{enumerate}
\tightlist
\item
  After each \texttt{cd} command, press \textbf{Insert+Up Arrow} to re-read the current line.
\item
  Press \textbf{Insert+Down Arrow} repeatedly to read output line by line.
\item
  For long listings: \texttt{dir\ /B\ \textgreater{}\ listing.txt}, then \texttt{notepad.exe\ listing.txt}.
\end{enumerate}

\subsubsection*{Windows Narrator Users}\label{docs__pandoc__latex__src__cli_complete_curriculum.md__windows-narrator-users-1}

\begin{enumerate}
\tightlist
\item
  After each \texttt{cd} command, press \textbf{Narrator+D} (CapsLock+D) to read the current line.
\item
  Press \textbf{Narrator+R} to read output from the current position.
\item
  For long listings, always redirect to a file and open in Notepad --- Notepad is more comfortable for Narrator users than reading directly from the terminal buffer.
\end{enumerate}

\subsubsection*{Dolphin SuperNova Users}\label{docs__pandoc__latex__src__cli_complete_curriculum.md__dolphin-supernova-users-1}

\begin{enumerate}
\tightlist
\item
  After each \texttt{cd} command, press \textbf{CapsLock+L} to read the current line.
\item
  Press \textbf{CapsLock+Numpad Plus} (say all) to read all output from the current position.
\item
  For long listings: \texttt{dir\ /B\ \textgreater{}\ listing.txt}, then \texttt{notepad.exe\ listing.txt}.
\end{enumerate}

\subsubsection*{Best Practice for All Screen Readers}\label{docs__pandoc__latex__src__cli_complete_curriculum.md__best-practice-for-all-screen-readers}

Always confirm your location after moving:

\begin{lstlisting}[style=Alabaster, language=cmd]
cd FolderName
cd                        :: Verify your new location

\end{lstlisting}

\begin{center}\rule{0.5\linewidth}{0.5pt}\end{center}

\subsection*{Practice Exercises}\label{docs__pandoc__latex__src__cli_complete_curriculum.md__practice-exercises-3}

\subsubsection*{Exercise 1: Basic Down Navigation}\label{docs__pandoc__latex__src__cli_complete_curriculum.md__exercise-1-basic-down-navigation}

\begin{enumerate}
\tightlist
\item
  \texttt{cd\ \%USERPROFILE\%}
\item
  Run \texttt{cd} to see location.
\item
  \texttt{dir\ /B}
\item
  \texttt{cd\ Documents}
\item
  Run \texttt{cd} to confirm.
\item
  \texttt{dir\ /B}
\end{enumerate}

\subsubsection*{Exercise 2: Using Tab Completion}\label{docs__pandoc__latex__src__cli_complete_curriculum.md__exercise-2-using-tab-completion-2}

\begin{enumerate}
\tightlist
\item
  \texttt{cd\ \%USERPROFILE\%}
\item
  Type \texttt{cd\ Doc} then press \textbf{Tab}.
\item
  Press \textbf{Enter}.
\item
  Run \texttt{cd} to confirm.
\end{enumerate}

\subsubsection*{Exercise 3: Navigating Up}\label{docs__pandoc__latex__src__cli_complete_curriculum.md__exercise-3-navigating-up}

\begin{enumerate}
\tightlist
\item
  Navigate several levels deep.
\item
  Run \texttt{cd} to see full path.
\item
  Use \texttt{cd\ ..} to go up one level at a time.
\item
  Run \texttt{cd} after each step.
\end{enumerate}

\subsubsection*{Exercise 4: Absolute Path Navigation}\label{docs__pandoc__latex__src__cli_complete_curriculum.md__exercise-4-absolute-path-navigation}

\begin{enumerate}
\tightlist
\item
  Navigate deep (3+ levels).
\item
  Run \texttt{cd\ \%USERPROFILE\%} to jump home.
\item
  Run \texttt{cd} to confirm.
\end{enumerate}

\subsubsection*{Exercise 5: Creating and Navigating a Structure}\label{docs__pandoc__latex__src__cli_complete_curriculum.md__exercise-5-creating-and-navigating-a-structure}

\begin{lstlisting}[style=Alabaster, language=cmd]
mkdir 3DPractice
cd 3DPractice
mkdir Models
cd Models
mkdir OpenSCAD
cd OpenSCAD

\end{lstlisting}

Run \texttt{cd} at each level. Then navigate back up using only \texttt{cd\ ..}.

\begin{center}\rule{0.5\linewidth}{0.5pt}\end{center}

\subsection*{Checkpoint Questions}\label{docs__pandoc__latex__src__cli_complete_curriculum.md__checkpoint-questions-2}

Answer all 10 before moving to Command Line Lesson 2:

\begin{enumerate}
\tightlist
\item
  What is the difference between \texttt{cd\ Documents} and \texttt{cd\ \textbackslash{}}?
\item
  What does \texttt{cd\ ..} do?
\item
  How do you go to your home folder from anywhere?
\item
  What is an absolute path? Give an example.
\item
  What is a relative path? Give an example.
\item
  How do you confirm your current location?
\item
  What does Tab completion do?
\item
  How would you navigate 3 levels deep, then back home?
\item
  What is the difference between \texttt{.} and \texttt{..}?
\item
  If you are lost, what command should you run first?
\end{enumerate}

\begin{center}\rule{0.5\linewidth}{0.5pt}\end{center}

\subsection*{Extension Problems}\label{docs__pandoc__latex__src__cli_complete_curriculum.md__extension-problems-8}

\begin{enumerate}
\tightlist
\item
  Create a folder with spaces in the name (e.g., \texttt{My\ Projects}) and navigate to it using quotes.
\item
  Use Tab completion to navigate 5+ levels deep without typing full names.
\item
  Create a script that saves \texttt{cd} output at each level to a log file.
\item
  Navigate to the same destination using both absolute and relative paths from different starting locations.
\item
  Build a complex folder tree (5+ levels) and navigate using only relative paths.
\item
  Document the full paths to your 5 most-used folders.
\item
  Create folders named \texttt{01Folder}, \texttt{02Folder}, etc. and practice Tab completion through them.
\item
  Navigate to a folder, save the path with \texttt{cd\ \textgreater{}\ mypath.txt}, navigate away, then return using the saved path.
\item
  Challenge: navigate to a destination using only relative paths without running \texttt{cd} to check at each step.
\item
  Create a batch file that navigates to a project folder and lists its contents in one step.
\end{enumerate}

\begin{center}\rule{0.5\linewidth}{0.5pt}\end{center}

\subsection*{Common Issues}\label{docs__pandoc__latex__src__cli_complete_curriculum.md__common-issues}

\subsubsection*{"The system cannot find the path specified"}\label{docs__pandoc__latex__src__cli_complete_curriculum.md__the-system-cannot-find-the-path-specified}

\begin{itemize}
\tightlist
\item
  Check spelling with \texttt{dir\ /B} to see correct folder names.
\item
  Use Tab completion to avoid typos.
\item
  Use the absolute path: \texttt{cd\ C:\textbackslash{}Users\textbackslash{}YourName\textbackslash{}Documents}.
\end{itemize}

\subsubsection*{"I\textquotesingle m lost"}\label{docs__pandoc__latex__src__cli_complete_curriculum.md__im-lost}

\begin{lstlisting}[style=Alabaster, language=cmd]
cd

\end{lstlisting}

This always shows your current location.

\subsubsection*{Tab Completion Not Working}\label{docs__pandoc__latex__src__cli_complete_curriculum.md__tab-completion-not-working}

\begin{itemize}
\tightlist
\item
  Type at least one character before pressing Tab.
\item
  Confirm the folder exists with \texttt{dir\ /B}.
\item
  Press Tab again to cycle through multiple matches.
\end{itemize}

\begin{center}\rule{0.5\linewidth}{0.5pt}\end{center}

\subsection*{Quick Reference}\label{docs__pandoc__latex__src__cli_complete_curriculum.md__quick-reference}

\begin{lstlisting}[style=Alabaster, language=cmd]
cd                         :: Show current location
cd FolderName              :: Go into a folder
cd ..                      :: Go up one level
cd ..\..                   :: Go up two levels
cd \                       :: Go to drive root
cd %USERPROFILE%           :: Go to home folder
dir /B                     :: List folder contents
dir /B /A:D                :: List only folders
dir /B /A:F                :: List only files

\end{lstlisting}

\begin{center}\rule{0.5\linewidth}{0.5pt}\end{center}

\subsection*{Next Steps}\label{docs__pandoc__latex__src__cli_complete_curriculum.md__next-steps-2}

Complete all exercises, pass the checkpoint questions, then move to \textbf{Command Line Lesson 2: File \& Folder Manipulation}.

\begin{center}\rule{0.5\linewidth}{0.5pt}\end{center}

\protect\phantomsection\label{docs__pandoc__latex__src__cli_complete_curriculum.md__cmd-lesson-2}{}

\section*{Lesson 2: File and Folder Manipulation}\label{docs__pandoc__latex__src__cli_complete_curriculum.md__lesson-2-file-and-folder-manipulation-1}

\begin{center}\rule{0.5\linewidth}{0.5pt}\end{center}

\textbf{Duration:} 2-2.5 hours (for screen reader users)\\
\textbf{Prerequisites:} Command Line Lesson Pre, Command Line Lesson 0, Command Line Lesson 1

\textbf{Learning Objectives:}

\begin{itemize}
\tightlist
\item
  Create, copy, move, and delete files and folders safely
\item
  Use wildcards to operate on multiple files
\item
  Understand dangerous operations and how to stay safe
\item
  Rename files and understand file extensions
\end{itemize}

\begin{center}\rule{0.5\linewidth}{0.5pt}\end{center}

\subsection*{Core File Manipulation Commands}\label{docs__pandoc__latex__src__cli_complete_curriculum.md__core-file-manipulation-commands}

\subsubsection*{\texorpdfstring{Create Folders: \texttt{mkdir}}{Create Folders: mkdir}}\label{docs__pandoc__latex__src__cli_complete_curriculum.md__create-folders-mkdir}

\begin{lstlisting}[style=Alabaster, language=cmd]
mkdir FolderName                     :: Create a single folder
mkdir Folder1 Folder2 Folder3        :: Create multiple folders at once
mkdir A\B\C                          :: Create nested folders (A, then B inside A, then C inside B)

\end{lstlisting}

\subsubsection*{\texorpdfstring{Create Files: \texttt{echo} with Redirection}{Create Files: echo with Redirection}}\label{docs__pandoc__latex__src__cli_complete_curriculum.md__create-files-echo-with-redirection}

\begin{lstlisting}[style=Alabaster, language=cmd]
echo This is some content > filename.txt    :: Create a file with text content
echo. > emptyfile.txt                       :: Create an empty file

\end{lstlisting}

Note: \texttt{echo.} (echo followed immediately by a period, no space) creates an empty file. \texttt{echo} with a space would put a space character in the file.

\subsubsection*{\texorpdfstring{Copy Files: \texttt{copy}}{Copy Files: copy}}\label{docs__pandoc__latex__src__cli_complete_curriculum.md__copy-files-copy}

\begin{lstlisting}[style=Alabaster, language=cmd]
copy source.txt destination.txt          :: Copy a file to a new name
copy source.txt backup\                  :: Copy a file into a folder
copy *.txt backup\                       :: Copy all .txt files to backup folder

\end{lstlisting}

\subsubsection*{\texorpdfstring{Copy Folders: \texttt{xcopy}}{Copy Folders: xcopy}}\label{docs__pandoc__latex__src__cli_complete_curriculum.md__copy-folders-xcopy}

\begin{lstlisting}[style=Alabaster, language=cmd]
xcopy sourcefolder destinationfolder /E /I    :: Copy a folder and all its contents

\end{lstlisting}

\begin{itemize}
\tightlist
\item
  \texttt{/E} copies all subdirectories including empty ones.
\item
  \texttt{/I} treats the destination as a folder (not a file) if it doesn\textquotesingle t exist.
\end{itemize}

\subsubsection*{\texorpdfstring{Move and Rename: \texttt{move}}{Move and Rename: move}}\label{docs__pandoc__latex__src__cli_complete_curriculum.md__move-and-rename-move}

\begin{lstlisting}[style=Alabaster, language=cmd]
move oldname.txt newname.txt             :: Rename a file
move file.txt folder\                    :: Move a file into a folder
move *.txt archive\                      :: Move all .txt files to archive folder

\end{lstlisting}

\texttt{move} renames when both source and destination are in the same folder. It moves when the destination is a different folder.

\subsubsection*{\texorpdfstring{Delete Files: \texttt{del}}{Delete Files: del}}\label{docs__pandoc__latex__src__cli_complete_curriculum.md__delete-files-del}

Caution

The \texttt{/Q} flag suppresses the confirmation prompt --- \texttt{del\ /Q} deletes immediately and silently. Only use \texttt{/Q} in scripts where you have already verified the targets; never use it interactively when learning.

\begin{lstlisting}[style=Alabaster, language=cmd]
del filename.txt                         :: Delete one specific file
del *.txt                                :: Delete all .txt files in current folder
del /Q *.tmp                             :: Delete quietly (no confirmation prompt)

\end{lstlisting}

Warning

\texttt{del} does \textbf{not} send files to the Recycle Bin. Deleted files are permanently gone with no undo. Always verify the target with \texttt{dir\ /B} before deleting.

\subsubsection*{\texorpdfstring{Delete Folders: \texttt{rmdir}}{Delete Folders: rmdir}}\label{docs__pandoc__latex__src__cli_complete_curriculum.md__delete-folders-rmdir}

\begin{lstlisting}[style=Alabaster, language=cmd]
rmdir foldername                         :: Delete an empty folder
rmdir /S /Q foldername                   :: Delete a folder and ALL its contents (no undo)

\end{lstlisting}

Warning

\texttt{rmdir\ /S\ /Q} permanently deletes a folder and \textbf{all its contents} with no confirmation and no undo. Always verify the exact folder name with \texttt{dir\ /B} first --- one wrong character can delete the wrong folder entirely.

\begin{center}\rule{0.5\linewidth}{0.5pt}\end{center}

\subsection*{Safe File Operations}\label{docs__pandoc__latex__src__cli_complete_curriculum.md__safe-file-operations}

Caution

\textbf{Check, then delete.} Never run \texttt{del} or \texttt{rmdir} without first running \texttt{dir\ /B} to confirm exactly what you\textquotesingle re about to remove. Wildcards like \texttt{del\ *.txt} will delete every matching file instantly and permanently.

\subsubsection*{Rule 1: Always Check Before Deleting}\label{docs__pandoc__latex__src__cli_complete_curriculum.md__rule-1-always-check-before-deleting}

\begin{lstlisting}[style=Alabaster, language=cmd]
dir /B                   :: See exactly what is here
dir /B *.txt             :: See exactly which .txt files exist
del *.txt                :: Only then delete

\end{lstlisting}

\subsubsection*{Rule 2: Make Backups First}\label{docs__pandoc__latex__src__cli_complete_curriculum.md__rule-2-make-backups-first}

\begin{lstlisting}[style=Alabaster, language=cmd]
mkdir backup
copy *.txt backup\       :: Copy all .txt files to a backup folder first
del *.txt                :: Now safe to delete originals

\end{lstlisting}

\subsubsection*{Rule 3: Test on One File First}\label{docs__pandoc__latex__src__cli_complete_curriculum.md__rule-3-test-on-one-file-first}

\begin{lstlisting}[style=Alabaster, language=cmd]
del test-one-file.txt    :: Test deletion on a single file
dir /B                   :: Confirm it is gone

\end{lstlisting}

Only then proceed with bulk operations.

\begin{center}\rule{0.5\linewidth}{0.5pt}\end{center}

\subsection*{Using Wildcards}\label{docs__pandoc__latex__src__cli_complete_curriculum.md__using-wildcards}

Wildcards let you operate on multiple files matching a pattern.

\subsubsection*{\texorpdfstring{\texttt{*} Wildcard (Match Any Number of Characters)}{* Wildcard (Match Any Number of Characters)}}\label{docs__pandoc__latex__src__cli_complete_curriculum.md__-wildcard-match-any-number-of-characters}

\begin{lstlisting}[style=Alabaster, language=cmd]
dir /B *.txt             :: List all files ending in .txt
copy *.scad backup\      :: Copy all .scad files to backup
del *.tmp                :: Delete all .tmp files

\end{lstlisting}

\subsubsection*{\texorpdfstring{\texttt{?} Wildcard (Match Exactly One Character)}{? Wildcard (Match Exactly One Character)}}\label{docs__pandoc__latex__src__cli_complete_curriculum.md__-wildcard-match-exactly-one-character}

\begin{lstlisting}[style=Alabaster, language=cmd]
dir /B file?.txt         :: Matches file1.txt, file2.txt, filea.txt, etc.
copy model?.scad models\ :: Copy model1.scad, model2.scad, etc.

\end{lstlisting}

\begin{center}\rule{0.5\linewidth}{0.5pt}\end{center}

\subsection*{Practical Examples}\label{docs__pandoc__latex__src__cli_complete_curriculum.md__practical-examples}

\subsection*{3DMake in Practice: Navigating a Project}\label{docs__pandoc__latex__src__cli_complete_curriculum.md__3dmake-in-practice-navigating-a-project-1}

Navigation commands are exactly what you need to explore and manage 3DMake projects:

\begin{lstlisting}[style=Alabaster, language=cmd]
REM Go to your project root
cd %USERPROFILE%\Documents\3d_projects\my_cube

REM See the project layout: src, build, 3dmake.toml
dir /B

REM Enter the source folder where your .scad model files live
cd src
dir /B

REM Open the main model in Notepad
3dm edit-model

REM Go back to project root, then check build output
cd ..
cd build
dir /B *.stl

REM Return to project root
cd ..

REM Use wildcards to find all projects with "cube" in the name
dir /B %USERPROFILE%\Documents\3d_projects\*cube*

\end{lstlisting}

Tab completion makes navigation fast: type \texttt{cd\ \%USERPROFILE\%\textbackslash{}Doc} and press Tab.

\begin{center}\rule{0.5\linewidth}{0.5pt}\end{center}

\subsection*{Practical Examples}\label{docs__pandoc__latex__src__cli_complete_curriculum.md__practical-examples-1}

\subsubsection*{Example 1: Create a Project Structure}\label{docs__pandoc__latex__src__cli_complete_curriculum.md__example-1-create-a-project-structure}

\begin{lstlisting}[style=Alabaster, language=cmd]
mkdir 3DProjects
cd 3DProjects
mkdir Models Prints Documentation Backups
dir /B /A:D

\end{lstlisting}

\subsubsection*{Example 2: Backup Before Editing}\label{docs__pandoc__latex__src__cli_complete_curriculum.md__example-2-backup-before-editing}

\begin{lstlisting}[style=Alabaster, language=cmd]
copy project.scad project-backup.scad    :: Create a backup copy
:: (edit project.scad)
copy project.scad project-v2.scad        :: Save a new versioned copy

\end{lstlisting}

\subsubsection*{Example 3: Organize Files by Type}\label{docs__pandoc__latex__src__cli_complete_curriculum.md__example-3-organize-files-by-type}

\begin{lstlisting}[style=Alabaster, language=cmd]
mkdir scad-files
mkdir text-files
move *.scad scad-files\
move *.txt text-files\
dir /B /A:D

\end{lstlisting}

\begin{center}\rule{0.5\linewidth}{0.5pt}\end{center}

\subsection*{Practice Exercises}\label{docs__pandoc__latex__src__cli_complete_curriculum.md__practice-exercises-4}

\subsubsection*{Exercise 1: Create a Folder Structure}\label{docs__pandoc__latex__src__cli_complete_curriculum.md__exercise-1-create-a-folder-structure}

\begin{enumerate}
\tightlist
\item
  \texttt{mkdir\ practice-session}
\item
  \texttt{cd\ practice-session}
\item
  \texttt{mkdir\ files\ documents\ models}
\item
  \texttt{dir\ /B\ /A:D}
\end{enumerate}

\textbf{Goal:} Create organized folder structures confidently.

\subsubsection*{Exercise 2: Create and Copy Files}\label{docs__pandoc__latex__src__cli_complete_curriculum.md__exercise-2-create-and-copy-files}

\begin{enumerate}
\tightlist
\item
  \texttt{echo\ Hello\ World\ \textgreater{}\ test.txt}
\item
  \texttt{type\ test.txt}
\item
  \texttt{copy\ test.txt\ test-backup.txt}
\item
  \texttt{dir\ /B\ *.txt}
\item
  \texttt{type\ test-backup.txt}
\end{enumerate}

\textbf{Goal:} Create and copy files without errors.

\subsubsection*{Exercise 3: Safe Deletion Practice}\label{docs__pandoc__latex__src__cli_complete_curriculum.md__exercise-3-safe-deletion-practice}

\begin{enumerate}
\tightlist
\item
  \texttt{echo\ content\ \textgreater{}\ file1.txt}
\item
  \texttt{echo\ content\ \textgreater{}\ file2.txt}
\item
  \texttt{echo\ content\ \textgreater{}\ file3.txt}
\item
  \texttt{dir\ /B\ *.txt} (confirm all three exist)
\item
  \texttt{del\ file1.txt} (delete just one)
\item
  \texttt{dir\ /B\ *.txt} (confirm only two remain)
\end{enumerate}

\textbf{Goal:} Practice the "check, then delete" pattern.

\subsubsection*{Exercise 4: Wildcard Operations}\label{docs__pandoc__latex__src__cli_complete_curriculum.md__exercise-4-wildcard-operations}

\subsection*{3DMake in Practice: Organizing Your Projects}\label{docs__pandoc__latex__src__cli_complete_curriculum.md__3dmake-in-practice-organizing-your-projects-1}

File management commands are essential for 3DMake project organization:

\begin{lstlisting}[style=Alabaster, language=cmd]
REM Create a new project from scratch
mkdir %USERPROFILE%\Documents\3d_projects\bottle_cap
cd %USERPROFILE%\Documents\3d_projects\bottle_cap
3dm new

REM Back up a working project before making big changes
cd %USERPROFILE%\Documents\3d_projects
xcopy /E /I bottle_cap bottle_cap_backup

REM Rename the main model file to something more specific
cd bottle_cap\src
rename main.scad bottle_cap.scad

REM List only .scad files to confirm the rename
dir /B *.scad

REM Delete a draft STL you no longer need
del %USERPROFILE%\Documents\3d_projects\bottle_cap\build\draft.stl

REM Build the renamed model (using the -m flag to specify the name)
cd %USERPROFILE%\Documents\3d_projects\bottle_cap
3dm build -m bottle_cap

REM See the fresh output in the build folder
dir /B build

\end{lstlisting}

\subsubsection*{Exercise 4: Wildcard Operations}\label{docs__pandoc__latex__src__cli_complete_curriculum.md__exercise-4-wildcard-operations-1}

\begin{enumerate}
\tightlist
\item
  Create three files: \texttt{echo\ a\ \textgreater{}\ doc1.txt}, \texttt{echo\ b\ \textgreater{}\ doc2.txt}, \texttt{echo\ c\ \textgreater{}\ doc3.txt}
\item
  \texttt{dir\ /B\ *.txt}
\item
  \texttt{mkdir\ archive}
\item
  \texttt{copy\ *.txt\ archive\textbackslash{}}
\item
  \texttt{cd\ archive}
\item
  \texttt{dir\ /B\ *.txt}
\end{enumerate}

\textbf{Goal:} Comfortable using wildcards for bulk operations.

\begin{center}\rule{0.5\linewidth}{0.5pt}\end{center}

\subsection*{Screen Reader Tips}\label{docs__pandoc__latex__src__cli_complete_curriculum.md__screen-reader-tips}

\subsubsection*{NVDA Users}\label{docs__pandoc__latex__src__cli_complete_curriculum.md__nvda-users-3}

After file operations, run \texttt{dir\ /B} and press \textbf{NVDA+Down Arrow} to verify results. Redirect output to a file for complex verification: \texttt{dir\ /B\ \textgreater{}\ check.txt}, then \texttt{notepad.exe\ check.txt}.

\subsubsection*{JAWS Users}\label{docs__pandoc__latex__src__cli_complete_curriculum.md__jaws-users-3}

After file operations, run \texttt{dir\ /B} and press \textbf{Insert+Down Arrow} to read results. For complex output: \texttt{dir\ /B\ \textgreater{}\ check.txt}, then \texttt{notepad.exe\ check.txt}.

\subsubsection*{Windows Narrator Users}\label{docs__pandoc__latex__src__cli_complete_curriculum.md__windows-narrator-users-2}

After file operations, run \texttt{dir\ /B} and press \textbf{Narrator+R} to read from the current position. Redirect to a file and open in Notepad for long output.

\subsubsection*{Dolphin SuperNova Users}\label{docs__pandoc__latex__src__cli_complete_curriculum.md__dolphin-supernova-users-2}

After file operations, run \texttt{dir\ /B} and press \textbf{CapsLock+Numpad Plus} to read results. For long output: \texttt{dir\ /B\ \textgreater{}\ check.txt}, then \texttt{notepad.exe\ check.txt}.

Important

After any \texttt{copy}, \texttt{move}, or \texttt{del} operation, always verify with \texttt{dir\ /B} to confirm the expected result. Never assume an operation succeeded without checking --- silent failures do occur.

\subsubsection*{General Tip for All Screen Readers}\label{docs__pandoc__latex__src__cli_complete_curriculum.md__general-tip-for-all-screen-readers}

After any \texttt{copy}, \texttt{move}, or \texttt{del} operation, always verify with \texttt{dir\ /B} to confirm the expected result. Never assume an operation worked without checking.

\begin{center}\rule{0.5\linewidth}{0.5pt}\end{center}

\subsection*{Checkpoint Questions}\label{docs__pandoc__latex__src__cli_complete_curriculum.md__checkpoint-questions-3}

\begin{enumerate}
\tightlist
\item
  How do you create a folder?
\item
  How do you copy a file?
\item
  How do you rename a file?
\item
  How do you delete a file safely (checking first)?
\item
  What does \texttt{*} match in a wildcard pattern?
\item
  What does \texttt{?} match in a wildcard pattern?
\item
  How would you copy all .scad files to a backup folder?
\item
  What should you always do before deleting?
\item
  How do you delete a folder and all its contents?
\item
  Why is it important to create a backup before editing?
\end{enumerate}

\begin{center}\rule{0.5\linewidth}{0.5pt}\end{center}

\subsection*{Extension Problems}\label{docs__pandoc__latex__src__cli_complete_curriculum.md__extension-problems-9}

\begin{enumerate}
\tightlist
\item
  Create a nested folder structure (5+ levels) using a single \texttt{mkdir} command with backslashes.
\item
  Create 10 files and organize them into subfolders using \texttt{move} with wildcards.
\item
  Create a date-based backup strategy: copy all files to a folder named with today\textquotesingle s date.
\item
  Use wildcards to select specific file types and copy them to separate organized folders.
\item
  Create a file-renaming system (file001.txt, file002.txt, etc.) and practice moving them between folders.
\end{enumerate}

\begin{center}\rule{0.5\linewidth}{0.5pt}\end{center}

\subsection*{Next Steps}\label{docs__pandoc__latex__src__cli_complete_curriculum.md__next-steps-3}

Complete all exercises, pass the checkpoint questions, then move to \textbf{Command Line Lesson 3: Input, Output \& Redirection}.

\begin{center}\rule{0.5\linewidth}{0.5pt}\end{center}

\protect\phantomsection\label{docs__pandoc__latex__src__cli_complete_curriculum.md__cmd-lesson-3}{}

\section*{Lesson 3: Input, Output, and Piping}\label{docs__pandoc__latex__src__cli_complete_curriculum.md__lesson-3-input-output-and-piping-1}

\begin{center}\rule{0.5\linewidth}{0.5pt}\end{center}

\textbf{Duration:} 1 class period
\textbf{Prerequisite:} Command Line Lesson 2 (File and Folder Manipulation)

\begin{center}\rule{0.5\linewidth}{0.5pt}\end{center}

\subsection*{Learning Objectives}\label{docs__pandoc__latex__src__cli_complete_curriculum.md__learning-objectives-4}

By the end of this lesson, you will be able to:

\begin{itemize}
\tightlist
\item
  Use \texttt{echo} to print text to the screen
\item
  Use \texttt{type} to read file contents
\item
  Use \texttt{\textgreater{}} to redirect output into a file
\item
  Use \texttt{\textbar{}} (pipe) to send one command\textquotesingle s output to another
\item
  Copy output to the clipboard with \texttt{clip}
\item
  Open files with a text editor from the command line
\end{itemize}

\begin{center}\rule{0.5\linewidth}{0.5pt}\end{center}

\subsection*{Commands Covered}\label{docs__pandoc__latex__src__cli_complete_curriculum.md__commands-covered-1}

{\def\LTcaptype{none} % do not increment counter
\begin{longtable}[]{@{}ll@{}}
\toprule\noalign{}
Command & What It Does \\
\midrule\noalign{}
\endhead
\bottomrule\noalign{}
\endlastfoot
\texttt{echo\ text} & Print text to the screen \\
\texttt{type\ filename} & Print the contents of a file \\
\texttt{\textgreater{}\ filename} &
Redirect output into a file (overwrites) \\
\texttt{\textgreater{}\textgreater{}\ filename} &
Append output to a file (adds to end) \\
\texttt{\textbar{}} & Pipe - send output from one command to the next \\
\texttt{clip} & Copy piped input to the Windows clipboard \\
\texttt{notepad\ filename} & Open a file in Notepad \\
\end{longtable}
}

\begin{center}\rule{0.5\linewidth}{0.5pt}\end{center}

\subsection*{\texorpdfstring{\texttt{echo} - Printing Text}{echo - Printing Text}}\label{docs__pandoc__latex__src__cli_complete_curriculum.md__echo---printing-text-1}

\texttt{echo} prints text to the screen. It is useful for testing, for writing text into files, and for understanding how piping works.

\begin{lstlisting}[style=Alabaster, language=cmd]
echo Hello, World
echo This is a test

\end{lstlisting}

\begin{center}\rule{0.5\linewidth}{0.5pt}\end{center}

\subsection*{\texorpdfstring{\texttt{type} - Reading Files}{type - Reading Files}}\label{docs__pandoc__latex__src__cli_complete_curriculum.md__type---reading-files}

\texttt{type} prints the contents of a file to the screen.

\begin{lstlisting}[style=Alabaster, language=cmd]
:: Read a text file
type %USERPROFILE%\Documents\notes.txt

:: Read an OpenSCAD file
type %USERPROFILE%\Documents\OpenSCAD_Projects\project0.scad

\end{lstlisting}

With a long file, use \texttt{type\ filename\ \textbar{}\ more} to read it page by page (press \texttt{Space} to advance, \texttt{Q} to quit).

\begin{center}\rule{0.5\linewidth}{0.5pt}\end{center}

\subsection*{\texorpdfstring{\texttt{\textgreater{}} - Redirecting Output to a File}{\textgreater{} - Redirecting Output to a File}}\label{docs__pandoc__latex__src__cli_complete_curriculum.md__---redirecting-output-to-a-file-1}

The \texttt{\textgreater{}} symbol redirects output from the screen into a file instead.

\begin{lstlisting}[style=Alabaster, language=cmd]
:: Create a file with a single line
echo Author: Your Name > header.txt

:: Confirm the file was created and has content
type header.txt

\end{lstlisting}

\textbf{Warning:} \texttt{\textgreater{}} overwrites the file if it already exists. Use \texttt{\textgreater{}\textgreater{}} to append instead:

\begin{lstlisting}[style=Alabaster, language=cmd]
echo Date: 2025 >> header.txt
echo Project: Floor Marker >> header.txt
type header.txt

\end{lstlisting}

\begin{center}\rule{0.5\linewidth}{0.5pt}\end{center}

\subsection*{\texorpdfstring{\texttt{\textbar{}} - Piping}{\textbar{} - Piping}}\label{docs__pandoc__latex__src__cli_complete_curriculum.md__---piping-1}

The pipe symbol \texttt{\textbar{}} sends the output of one command to the input of the next. This lets you chain commands together.

\begin{lstlisting}[style=Alabaster, language=cmd]
:: List files and send the list to clip (copies to clipboard)
dir /B | clip
## 3DMake in Practice: Capturing Build Output

Piping and redirection are essential for working with 3DMake's output accessibly:

```cmd
REM Build and save the full log to a file so you can read it in Notepad
3dm build > build_log.txt
notepad build_log.txt

REM Build and slice, keeping only lines with errors or warnings
3dm build slice 2>&1 | findstr /I "error warning"

REM List all STL files produced so far
dir /B build\*.stl

REM Get a list of available printer profiles and save for reference
3dm list-profiles > profiles.txt
notepad profiles.txt

REM Get a list of material overlays (PLA, PETG, supports, etc.)
3dm list-overlays > overlays.txt
notepad overlays.txt

REM Get an AI description of your model and save it
3dm info > model_description.txt
notepad model_description.txt

\end{lstlisting}

\begin{lstlisting}[style=Alabaster, language=cmd]
:: List files and send the list to clip (copies to clipboard)
dir /B | clip

:: Now paste with Ctrl + V anywhere

\end{lstlisting}

\begin{lstlisting}[style=Alabaster, language=cmd]
:: Search within a file's contents using find
type project0.scad | find "cube"

\end{lstlisting}

\begin{center}\rule{0.5\linewidth}{0.5pt}\end{center}

\subsection*{\texorpdfstring{\texttt{clip} - Copying to Clipboard}{clip - Copying to Clipboard}}\label{docs__pandoc__latex__src__cli_complete_curriculum.md__clip---copying-to-clipboard-1}

\texttt{clip} takes whatever is piped to it and puts it on the Windows clipboard.

\begin{lstlisting}[style=Alabaster, language=cmd]
:: Copy your current directory path to the clipboard
cd | clip

:: Copy a file listing to clipboard
dir /B | clip

:: Copy the contents of a file to clipboard
type notes.txt | clip

\end{lstlisting}

After any of these, press \texttt{Ctrl\ +\ V} in any application to paste.

\begin{center}\rule{0.5\linewidth}{0.5pt}\end{center}

\subsection*{Opening Files in Notepad}\label{docs__pandoc__latex__src__cli_complete_curriculum.md__opening-files-in-notepad-1}

\begin{lstlisting}[style=Alabaster, language=cmd]
:: Open a file in Notepad
notepad %USERPROFILE%\Documents\notes.txt

:: Open a .scad file
notepad %USERPROFILE%\Documents\OpenSCAD_Projects\project0.scad

:: Create a new file and open it
echo. > new_notes.txt
notepad new_notes.txt

\end{lstlisting}

\begin{center}\rule{0.5\linewidth}{0.5pt}\end{center}

\subsection*{Step-by-step Tasks}\label{docs__pandoc__latex__src__cli_complete_curriculum.md__step-by-step-tasks-3}

\begin{enumerate}
\tightlist
\item
  Create \texttt{practice.txt} with three lines using \texttt{echo} and \texttt{\textgreater{}}/\texttt{\textgreater{}\textgreater{}}.
\item
  Read the file with \texttt{type\ practice.txt}.
\item
  Pipe the file into \texttt{find} to search for a word.
\item
  Copy the file contents to clipboard with \texttt{type\ practice.txt\ \textbar{}\ clip}.
\item
  Redirect \texttt{dir\ /B} into \texttt{list.txt} and open it in Notepad.
\end{enumerate}

\subsection*{Checkpoints}\label{docs__pandoc__latex__src__cli_complete_curriculum.md__checkpoints-5}

\begin{itemize}
\tightlist
\item
  After step 3 you should be able to find a keyword using piping.
\end{itemize}

\subsection*{Quiz - Lesson CMD.3}\label{docs__pandoc__latex__src__cli_complete_curriculum.md__quiz---lesson-cmd3}

\begin{enumerate}
\tightlist
\item
  What is the difference between \texttt{\textgreater{}} and \texttt{\textgreater{}\textgreater{}}?
\item
  What does the pipe \texttt{\textbar{}} do?
\item
  How do you copy output to the clipboard?
\item
  How would you page through long output?
\item
  How do you suppress output (send it nowhere)?
\item
  True or False: The pipe operator \texttt{\textbar{}} connects the output of one command to the input of another.
\item
  Explain why redirecting output to a file is useful for screen reader users.
\item
  Write a command that would search for the word "sphere" in all \texttt{.scad} files in a directory.
\item
  How would you count the number of lines in a file using CMD piping?
\item
  Describe a practical scenario in 3D printing where you would pipe or redirect command output.
\item
  What would be the difference in output between \texttt{echo\ test\ \textgreater{}\ file.txt} (run twice) vs \texttt{echo\ test\ \textgreater{}\textgreater{}\ file.txt} (run twice)? Show the expected file contents.
\item
  Design a three-step piping chain: read a file, filter for specific content, and save the results; explain what each pipe does.
\item
  You have a 500-line \texttt{.scad} file and need to find all instances of \texttt{sphere()} and count them. Write the command.
\item
  Explain how \texttt{clip} is particularly valuable for screen reader users when working with file paths or long output strings.
\item
  Describe how you would use pipes and redirection to create a timestamped backup report of all \texttt{.stl} files in a 3D printing project folder.
\end{enumerate}

\subsection*{Extension Problems}\label{docs__pandoc__latex__src__cli_complete_curriculum.md__extension-problems-10}

\begin{enumerate}
\tightlist
\item
  Use piping to count lines in a file (hint: \texttt{type\ file.txt\ \textbar{}\ find\ /C\ /V\ ""}).
\item
  Save a long \texttt{dir\ /B} output and search it with \texttt{find}.
\item
  Chain multiple pipes to filter and then save results.
\item
  Practice copying different command outputs to clipboard and pasting.
\item
  Create a small batch script that generates a report (counts of files by extension).
\item
  Build a data processing pipeline: read a text file, filter rows, and export results; document each step.
\item
  Write a batch script that pipes directory listing to count occurrences of each file extension; create a summary report.
\item
  Create a log analysis command: read a log file, filter for errors, and save matching lines to a separate error log.
\item
  Design a piping workflow for 3D printing file management: find \texttt{.stl} files, extract their names, and generate a report.
\item
  Develop a reusable piping pattern library: create batch scripts for common filtering, sorting, and reporting patterns; test each with different inputs.
\item
  Build a complex filter pipeline: read a \texttt{.scad} file, extract lines containing specific geometry commands, count each type, and output a summary to both screen and file.
\item
  Create an interactive filtering tool: build a batch script that accepts a search term, pipes through multiple filters, and displays paginated results.
\item
  Develop a performance analysis tool: use piping to combine file listing, metadata extraction, and statistical reporting; export results to a dated report file.
\item
  Implement a comprehensive error-handling pipeline: read output, catch errors, log them separately, and generate a summary of successes vs failures.
\item
  Design and execute a real-world project backup workflow: use piping to verify file existence, count files by type, generate a backup manifest, and create audit logs --- all in one integrated command pipeline.
\end{enumerate}

\subsection*{References}\label{docs__pandoc__latex__src__cli_complete_curriculum.md__references-8}

\begin{itemize}
\tightlist
\item
  Microsoft. (2024). \emph{Using redirection operators in CMD}. \url{https://example.com}
\item
  Microsoft. (2024). \emph{FIND command reference}. \url{https://example.com}
\item
  Microsoft. (2024). \emph{CMD pipeline concepts}. \url{https://example.com}
\end{itemize}

\subsection*{Helpful Resources}\label{docs__pandoc__latex__src__cli_complete_curriculum.md__helpful-resources-7}

\begin{itemize}
\tightlist
\item
  \href{https://learn.microsoft.com/en-us/windows-server/administration/windows-commands/redirection}{Using Redirection in CMD}
\item
  \href{https://learn.microsoft.com/en-us/windows-server/administration/windows-commands/find}{Piping and FIND}
\item
  \href{https://learn.microsoft.com/en-us/windows-server/administration/windows-commands/type}{TYPE Command Reference}
\item
  \href{https://learn.microsoft.com/en-us/windows-server/administration/windows-commands/findstr}{FINDSTR for Advanced Searching}
\item
  \href{https://learn.microsoft.com/en-us/windows-server/administration/windows-commands/}{CMD Pipeline Concepts}
\end{itemize}

\begin{center}\rule{0.5\linewidth}{0.5pt}\end{center}

\protect\phantomsection\label{docs__pandoc__latex__src__cli_complete_curriculum.md__cmd-lesson-4}{}

\section*{Lesson 4: Environment Variables, PATH, and Aliases}\label{docs__pandoc__latex__src__cli_complete_curriculum.md__lesson-4-environment-variables-path-and-aliases-1}

\begin{center}\rule{0.5\linewidth}{0.5pt}\end{center}

Estimated time: 30-45 minutes

\subsection*{Learning Objectives}\label{docs__pandoc__latex__src__cli_complete_curriculum.md__learning-objectives-5}

\begin{itemize}
\tightlist
\item
  Read environment variables with \texttt{\%VARNAME\%}
\item
  Inspect and verify programs in the \texttt{PATH}
\item
  Create temporary aliases with \texttt{doskey} and understand making them persistent via a startup script
\end{itemize}

\subsection*{Materials}\label{docs__pandoc__latex__src__cli_complete_curriculum.md__materials-2}

\begin{itemize}
\tightlist
\item
  Command Prompt
\end{itemize}

\subsection*{Step-by-step Tasks}\label{docs__pandoc__latex__src__cli_complete_curriculum.md__step-by-step-tasks-4}

\begin{enumerate}
\tightlist
\item
  Show your username and home path with \texttt{echo\ \%USERNAME\%} and \texttt{echo\ \%USERPROFILE\%}.
\item
  Inspect \texttt{echo\ \%PATH\%} and identify whether \texttt{openscad} or \texttt{code} would be found.
\item
  Run \texttt{where\ openscad} and note the result.
\item
  Create a temporary alias: \texttt{doskey\ preview=openscad\ \$*} and test \texttt{preview\ myfile.scad}.
\item
  Open your startup script (\texttt{notepad.exe\ autorun.bat}) and add the doskey line to make it persistent (advanced).
\end{enumerate}

\subsection*{Checkpoints}\label{docs__pandoc__latex__src__cli_complete_curriculum.md__checkpoints-6}

\begin{itemize}
\tightlist
\item
  After step 3 you can determine whether a program will be found by \texttt{PATH}.
\end{itemize}

\subsection*{Quiz - Lesson CMD.4}\label{docs__pandoc__latex__src__cli_complete_curriculum.md__quiz---lesson-cmd4}

\subsection*{3DMake in Practice: DOSKEY Macros for 3DMake}\label{docs__pandoc__latex__src__cli_complete_curriculum.md__3dmake-in-practice-doskey-macros-for-3dmake}

DOSKEY macros give you fast shortcuts for your most-used 3DMake commands:

\begin{lstlisting}[style=Alabaster, language=cmd]
REM Quick macro to jump to your 3D projects folder
doskey 3dp=cd %USERPROFILE%\Documents\3d_projects

REM Quick macro to run the full build-slice-print pipeline
doskey 3dprint=3dm build slice print

REM Quick macro to build then describe the model
doskey 3dinfo=3dm build $T 3dm info

REM Test the macros in the current session
3dp
3dprint

REM To make these permanent, add them to a startup batch file
notepad %USERPROFILE%\cmd_startup.cmd
REM Add the doskey lines, save, and then register via Windows Registry:
REM HKCU\Software\Microsoft\Command Processor\AutoRun = %USERPROFILE%\cmd_startup.cmd

\end{lstlisting}

\subsection*{Checkpoints}\label{docs__pandoc__latex__src__cli_complete_curriculum.md__checkpoints-7}

\begin{itemize}
\tightlist
\item
  After step 3 you can determine whether a program will be found by \texttt{PATH}.
\end{itemize}

\subsection*{Quiz - Lesson CMD.4}\label{docs__pandoc__latex__src__cli_complete_curriculum.md__quiz---lesson-cmd4-1}

\begin{enumerate}
\tightlist
\item
  How do you print an environment variable?
\item
  What is the purpose of \texttt{PATH}?
\item
  How do you check whether \texttt{openscad} is available?
\item
  How do you create a temporary alias in CMD?
\item
  Where would you make an alias permanent?
\item
  True or False: Environment variable names in CMD are case-sensitive.
\item
  Explain why having a program in your PATH is useful compared to always using its full file path.
\item
  Write a command that would create a doskey alias called \texttt{slicer} for the OpenSCAD executable.
\item
  What file or technique would you use to make a doskey alias persist across CMD sessions?
\item
  Describe a practical benefit of using the \texttt{\%TEMP\%} directory for temporary files in a 3D printing workflow.
\item
  You have a custom batch script at \texttt{C:\textbackslash{}Scripts\textbackslash{}backup\_models.bat} that you want to run from anywhere as \texttt{backup-now}. What steps would you take to make this work?
\item
  Explain the difference between setting an environment variable in the current session with \texttt{set} vs. using \texttt{setx} for permanence.
\item
  Design a strategy for managing multiple 3D printing projects, each with different tool paths and directories; show how to structure environment variables for each.
\item
  If a program is not found by \texttt{where}, what are the possible reasons, and how would you troubleshoot?
\item
  Describe how you would verify that your CMD autorun script is loading correctly and how to debug issues if aliases or environment variables don\textquotesingle t appear after opening a new CMD session.
\end{enumerate}

\subsection*{Extension Problems}\label{docs__pandoc__latex__src__cli_complete_curriculum.md__extension-problems-11}

\begin{enumerate}
\tightlist
\item
  Add a folder to PATH for a test program (describe steps; do not change system PATH without admin).
\item
  Create a short autorun snippet that sets two aliases and test re-opening CMD.
\item
  Use \texttt{where} to list the path for several common programs.
\item
  Explore \texttt{\%TEMP\%} and create a file there.
\item
  Save a copy of your current PATH to a text file and examine it in your editor.
\item
  Create a CMD autorun script that loads custom aliases and environment variables for your 3D printing workflow; test it in a new session.
\item
  Build a "project profile" batch script that sets environment variables for CAD, slicing, and print directories; switch between profiles for different projects.
\item
  Write a batch script that audits your current environment variables and creates a summary report of what\textquotesingle s set and why.
\item
  Design a custom alias system using doskey for common 3D printing commands; document the aliases and their purposes.
\item
  Create a profile migration guide: document how to export and import your CMD aliases and variables across machines for consistent workflows.
\item
  Implement a safe PATH modification script: create a utility that allows you to add/remove directories from PATH for the current session only; show how to make it permanent with \texttt{setx}.
\item
  Build a comprehensive autorun framework: create separate \texttt{.bat} files for aliases, environment variables, and helper macros; have your main autorun load all of them.
\item
  Develop an environment validation tool: write a batch script that checks whether all required programs (OpenSCAD, slicers, etc.) are accessible via PATH; report findings and suggest fixes.
\item
  Create a project-switching alias system: design a batch script that changes all environment variables and aliases based on the current project; test switching between multiple projects.
\item
  Build a troubleshooting guide: create a batch script that exports your current environment state (variables, aliases, PATH) to a timestamped file, allowing you to compare states before and after changes.
\end{enumerate}

\subsection*{References}\label{docs__pandoc__latex__src__cli_complete_curriculum.md__references-9}

\begin{itemize}
\tightlist
\item
  Microsoft. (2024). \emph{Environment variables in CMD}. \url{https://example.com}
\item
  Microsoft. (2024). \emph{DOSKEY command reference}. \url{https://example.com}
\item
  Microsoft. (2024). \emph{WHERE command reference}. \url{https://example.com}
\end{itemize}

\subsection*{Helpful Resources}\label{docs__pandoc__latex__src__cli_complete_curriculum.md__helpful-resources-8}

\begin{itemize}
\tightlist
\item
  \href{https://learn.microsoft.com/en-us/windows-server/administration/windows-commands/set}{Environment Variables in CMD}
\item
  \href{https://learn.microsoft.com/en-us/windows-server/administration/windows-commands/path}{Understanding the PATH Variable}
\item
  \href{https://learn.microsoft.com/en-us/windows-server/administration/windows-commands/doskey}{DOSKEY Alias Reference}
\item
  \href{https://learn.microsoft.com/en-us/windows-server/administration/windows-commands/where}{WHERE Command for Locating Programs}
\item
  \href{https://learn.microsoft.com/en-us/windows-server/administration/windows-commands/setx}{SETX for Permanent Variables}
\end{itemize}

\begin{center}\rule{0.5\linewidth}{0.5pt}\end{center}

\protect\phantomsection\label{docs__pandoc__latex__src__cli_complete_curriculum.md__cmd-lesson-5}{}

\section*{Lesson 5: Filling in the Gaps --- Startup Scripts, History, and Useful Tricks}\label{docs__pandoc__latex__src__cli_complete_curriculum.md__lesson-5-filling-in-the-gaps--startup-scripts-history-and-useful-tricks}

\begin{center}\rule{0.5\linewidth}{0.5pt}\end{center}

Estimated time: 30-45 minutes

\subsection*{Learning Objectives}\label{docs__pandoc__latex__src__cli_complete_curriculum.md__learning-objectives-6}

\begin{itemize}
\tightlist
\item
  Use history and abort commands (\texttt{doskey\ /history}, \texttt{F7}, \texttt{Ctrl+C})
\item
  Inspect and edit your CMD autorun script for persistent settings
\item
  Run programs by full path using the \texttt{start} or call operator
\item
  Handle common screen reader edge cases when using the terminal
\end{itemize}

\subsection*{Materials}\label{docs__pandoc__latex__src__cli_complete_curriculum.md__materials-3}

\begin{itemize}
\tightlist
\item
  Command Prompt and an editor (Notepad/VS Code)
\end{itemize}

\subsection*{Step-by-step Tasks}\label{docs__pandoc__latex__src__cli_complete_curriculum.md__step-by-step-tasks-5}

\begin{enumerate}
\tightlist
\item
  Run several simple commands (e.g., \texttt{cd}, \texttt{dir\ /B}, \texttt{echo\ hi}) then press \texttt{F7} or run \texttt{doskey\ /history} to view them.
\item
  Use \texttt{F8} to search back through history, or use the up-arrow to re-run a previous command.
\item
  Practice aborting a long-running command with \texttt{Ctrl\ +\ C} (for example, \texttt{ping\ 8.8.8.8}).
\item
  Open your autorun script: \texttt{notepad.exe\ autorun.bat}; if it doesn\textquotesingle t exist, create it with \texttt{echo.\ \textgreater{}\ autorun.bat}.
\item
  Add a persistent alias line to your autorun script (example: \texttt{doskey\ preview=openscad\ \$*}), save, and configure CMD to use it by registering it in the registry (advanced).
\end{enumerate}

\subsection*{Checkpoints}\label{docs__pandoc__latex__src__cli_complete_curriculum.md__checkpoints-8}

\begin{itemize}
\tightlist
\item
  After step 2 you can re-run a recent command from history.
\item
  After step 5 your alias should persist across CMD sessions.
\end{itemize}

\subsection*{Quiz - Lesson CMD.5}\label{docs__pandoc__latex__src__cli_complete_curriculum.md__quiz---lesson-cmd5}

\begin{enumerate}
\tightlist
\item
  How do you view the command history in CMD?
\item
  Which key combination aborts a running command?
\item
  What does \texttt{echo\ \%CMDCMDLINE\%} show?
\item
  How does the \texttt{start} command help run executables?
\item
  What is one strategy if terminal output stops being announced by your screen reader?
\item
  True or False: Using \texttt{Ctrl+C} permanently deletes any files created by the command you abort.
\item
  Explain the difference between pressing \texttt{F7} and running \texttt{doskey\ /history} in CMD.
\item
  If you add a doskey macro to your autorun script but it doesn\textquotesingle t take effect after opening a new CMD window, what should you verify?
\item
  Write a command that would run a program at the path \texttt{C:\textbackslash{}Program\ Files\textbackslash{}OpenSCAD\textbackslash{}openscad.exe} directly.
\item
  Describe a practical workflow scenario where having keyboard shortcuts (doskey macros) in your autorun script would save time.
\item
  Explain how to re-run the 5th command from your history using \texttt{F7} and selection, and what would happen if that command had file operations (creates/deletes).
\item
  Design an autorun initialization strategy that separates utilities for different projects; explain how you would switch between them.
\item
  Walk through a troubleshooting workflow: your screen reader stops announcing output after running a long command. What steps would you take to diagnose and resolve the issue?
\item
  Create a safety checkpoint system: before any destructive operation (mass delete, overwrite), how would you use autorun macros and history to verify the command is correct?
\item
  Develop a comprehensive capstone scenario: integrate everything from Command Line Lesson 0 through Command Line Lesson 5 (navigation, file operations, piping, environment setup, history) to design an automated 3D printing project workflow with error handling and logging.
\end{enumerate}

\subsection*{Extension Problems}\label{docs__pandoc__latex__src__cli_complete_curriculum.md__extension-problems-12}

\begin{enumerate}
\tightlist
\item
  Add a doskey macro and an environment variable change to your autorun script and document the behavior after reopening CMD.
\item
  Create a short batch script that automates creating a project folder and an initial \texttt{.scad} file.
\item
  Experiment with running OpenSCAD by full path using \texttt{start} and by placing it in PATH; compare results.
\item
  Practice redirecting \texttt{help} output to a file and reading it in Notepad for screen reader clarity.
\item
  Document three screen reader troubleshooting steps you used and when they helped.
\item
  Build a comprehensive CMD autorun script that includes aliases, environment variables, and helper macros for your 3D printing workflow.
\item
  Create a batch script that troubleshoots common CMD issues (missing commands, permission errors, command not found); test at least three scenarios.
\item
  Write a batch script that coordinates multiple tasks: creates a project folder, starts OpenSCAD, and opens a notes file.
\item
  Design a screen-reader accessibility guide for CMD: document commands, outputs, and accessible navigation patterns.
\item
  Develop an advanced CMD workflow: implement error handling, logging, and confirmation prompts for risky operations.
\item
  Implement an "undo" system using history: create a batch script that logs destructive commands (\texttt{del}, \texttt{move}, \texttt{copy\ /Y}) and allows you to review the last operation.
\item
  Build an autorun debugger: create a script that compares two CMD sessions\textquotesingle{} environment states (variables, aliases, macros) to identify what loaded/failed to load.
\item
  Develop a multi-project autorun manager: design a system where you can switch entire environments (paths, aliases, variables) for different 3D printing projects by running a single script.
\item
  Create a comprehensive accessibility analyzer: write a batch script that tests whether key CMD commands produce screen-reader-friendly output; document workarounds for commands that don\textquotesingle t.
\item
  Design a complete capstone project: build an integrated automation suite that manages a 3D printing workflow (project setup, file organization, CAD/slicing tool automation, output logging, error recovery, and audit trails) with full error handling and documentation.
\end{enumerate}

\subsection*{References}\label{docs__pandoc__latex__src__cli_complete_curriculum.md__references-10}

\begin{itemize}
\tightlist
\item
  Microsoft. (2024). \emph{CMD history and recall functionality}. \url{https://example.com}
\item
  Microsoft. (2024). \emph{Using CMD autorun scripts}. \url{https://example.com}
\item
  Microsoft. (2024). \emph{The START command for running executables}. \url{https://example.com}
\end{itemize}

\subsection*{Helpful Resources}\label{docs__pandoc__latex__src__cli_complete_curriculum.md__helpful-resources-9}

\begin{itemize}
\tightlist
\item
  \href{https://learn.microsoft.com/en-us/windows-server/administration/windows-commands/doskey}{DOSKEY History and Recall}
\item
  \href{https://learn.microsoft.com/en-us/windows-server/administration/windows-commands/cmd}{CMD Autorun Scripts}
\item
  \href{https://learn.microsoft.com/en-us/windows-server/administration/windows-commands/start}{START Command Reference}
\item
  \href{https://learn.microsoft.com/en-us/windows-server/administration/windows-commands/}{CMD History Navigation}
\item
  \href{https://learn.microsoft.com/en-us/windows-server/administration/windows-commands/windows-commands-glossary}{Screen Reader Tips for CMD}
\end{itemize}

\begin{center}\rule{0.5\linewidth}{0.5pt}\end{center}

\protect\phantomsection\label{docs__pandoc__latex__src__cli_complete_curriculum.md__cmd-lesson-6}{}

\section*{Lesson 6: Advanced Techniques --- Batch Scripts, Subroutines, and Professional Workflows}\label{docs__pandoc__latex__src__cli_complete_curriculum.md__lesson-6-advanced-techniques--batch-scripts-subroutines-and-professional-workflows}

\begin{center}\rule{0.5\linewidth}{0.5pt}\end{center}

\textbf{Duration:} 4-4.5 hours (for screen reader users)
\textbf{Prerequisites:} Complete Command Line Lesson 0 through Command Line Lesson 5
\textbf{Skill Level:} Advanced intermediate

\subsection*{3DMake in Practice: Automating Your Print Workflow}\label{docs__pandoc__latex__src__cli_complete_curriculum.md__3dmake-in-practice-automating-your-print-workflow-1}

Use a batch script to fully automate a 3DMake workflow across multiple projects:

\begin{lstlisting}[style=Alabaster, language=cmd]
@echo off
REM batch_print.bat - build and print every 3DMake project subfolder

set PROJECTS=%USERPROFILE%\Documents\3d_projects

for /D %%P in ("%PROJECTS%\*") do (
    echo Processing: %%~nxP
    cd "%%P"

    3dm build slice > build_log.txt 2>&1
    if %ERRORLEVEL% EQU 0 (
        echo   OK - sending to printer
        3dm print
    ) else (
        echo   FAILED - check build_log.txt
        notepad build_log.txt
    )
)

echo Batch complete.

\end{lstlisting}

Save as \texttt{batch\_print.bat} and run by typing \texttt{batch\_print}. This script uses loops, conditionals, and output redirection --- all from this lesson.

This lesson extends CMD skills to professional-level workflows. You\textquotesingle ll learn to automate complex tasks, write reusable batch scripts, and integrate tools for 3D printing workflows.

\begin{center}\rule{0.5\linewidth}{0.5pt}\end{center}

\subsection*{Learning Objectives}\label{docs__pandoc__latex__src__cli_complete_curriculum.md__learning-objectives-7}

By the end of this lesson, you will be able to:

\begin{itemize}
\tightlist
\item
  Create and run batch files (.bat files)
\item
  Write subroutines that accept parameters
\item
  Use loops to repeat tasks automatically
\item
  Automate batch processing of 3D models
\item
  Debug batch scripts when something goes wrong
\item
  Create professional workflows combining multiple tools
\end{itemize}

\begin{center}\rule{0.5\linewidth}{0.5pt}\end{center}

\subsection*{Batch Script Basics}\label{docs__pandoc__latex__src__cli_complete_curriculum.md__batch-script-basics}

\subsubsection*{What\textquotesingle s a Batch Script?}\label{docs__pandoc__latex__src__cli_complete_curriculum.md__whats-a-batch-script}

A batch file (\texttt{.bat}) contains multiple CMD commands that run in sequence. Instead of typing commands one by one, you put them in a file and run them all at once.

\textbf{Why use batch scripts?}

\begin{itemize}
\tightlist
\item
  Repeatability: Run the same task 100 times identically
\item
  Documentation: Commands are written down for reference
\item
  Complexity: Combine many commands logically
\item
  Automation: Schedule scripts to run automatically
\end{itemize}

\subsubsection*{Creating Your First Batch Script}\label{docs__pandoc__latex__src__cli_complete_curriculum.md__creating-your-first-batch-script}

\textbf{Step 1: Open a text editor}

\begin{lstlisting}[style=Alabaster, language=cmd]
notepad.exe my-first-script.bat

\end{lstlisting}

\textbf{Step 2: Type this script}

\begin{lstlisting}[style=Alabaster, language=cmd]
@echo off
:: This is a comment - screen readers will read it
echo Script is running!
cd
dir /B
echo Script is done!

\end{lstlisting}

\textbf{Step 3: Save the file}

\begin{itemize}
\tightlist
\item
  In Notepad: Ctrl+S
\item
  Make sure filename ends in \texttt{.bat}
\item
  Save in an easy-to-find location (like Documents)
\end{itemize}

\textbf{Step 4: Run the script}

\begin{lstlisting}[style=Alabaster, language=cmd]
my-first-script.bat

\end{lstlisting}

\textbf{What happens:}
CMD runs each command in sequence and shows output.

\subsubsection*{\texorpdfstring{Important: \texttt{@echo\ off}}{Important: @echo off}}\label{docs__pandoc__latex__src__cli_complete_curriculum.md__important-echo-off}

On some scripts, each command is echoed before running. \texttt{@echo\ off} at the top suppresses this, showing only output --- not the commands themselves.

\begin{lstlisting}[style=Alabaster, language=cmd]
@echo off
:: Now only output is shown, not each command
echo Hello!

\end{lstlisting}

\begin{center}\rule{0.5\linewidth}{0.5pt}\end{center}

\subsection*{Variables and Parameters}\label{docs__pandoc__latex__src__cli_complete_curriculum.md__variables-and-parameters-1}

\subsubsection*{Using Variables}\label{docs__pandoc__latex__src__cli_complete_curriculum.md__using-variables-1}

Variables store values you want to use later.

\textbf{Example script:}

\begin{lstlisting}[style=Alabaster, language=cmd]
@echo off
set mypath=C:\Users\YourName\Documents
cd %mypath%
echo I am now in:
cd
dir /B

\end{lstlisting}

\textbf{Breaking it down:}

\begin{itemize}
\tightlist
\item
  \texttt{set\ mypath=...} assigns the variable
\item
  \texttt{\%mypath\%} uses the variable (wraps it in \texttt{\%})
\item
  Variables in CMD are always referenced with \texttt{\%VARNAME\%}
\end{itemize}

\subsubsection*{Subroutines with Parameters}\label{docs__pandoc__latex__src__cli_complete_curriculum.md__subroutines-with-parameters}

A subroutine is reusable code that you can call with different inputs using \texttt{:label} and \texttt{call}.

\textbf{Example: A subroutine that lists files in a folder}

\begin{lstlisting}[style=Alabaster, language=cmd]
@echo off
call :ListFolder "C:\Users\YourName\Documents"
call :ListFolder "C:\Users\YourName\Downloads"
goto :eof

:ListFolder
echo Contents of: %~1
cd /D %~1
dir /B
goto :eof

\end{lstlisting}

\textbf{What\textquotesingle s happening:}

\begin{itemize}
\tightlist
\item
  \texttt{:ListFolder} marks the start of the subroutine
\item
  \texttt{\%\textasciitilde{}1} is the first argument passed to the subroutine
\item
  \texttt{call\ :ListFolder\ "path"} calls the subroutine with a path
\item
  \texttt{goto\ :eof} exits the subroutine (or script)
\end{itemize}

\textbf{Screen reader tip:} When you call a subroutine, CMD will announce the results just like any command.

\begin{center}\rule{0.5\linewidth}{0.5pt}\end{center}

\subsection*{Loops - Repeating Tasks}\label{docs__pandoc__latex__src__cli_complete_curriculum.md__loops---repeating-tasks-1}

\subsubsection*{Loop Over Files}\label{docs__pandoc__latex__src__cli_complete_curriculum.md__loop-over-files-1}

Imagine you have 10 SCAD files and want to print their names. You could do it 10 times manually, or use a loop.

\textbf{Example: Print every .scad file in a folder}

\begin{lstlisting}[style=Alabaster, language=cmd]
@echo off
for %%f in (*.scad) do (
    echo === File: %%f ===
    type %%f
    echo.
)

\end{lstlisting}

\textbf{What\textquotesingle s happening:}

\begin{itemize}
\tightlist
\item
  \texttt{for\ \%\%f\ in\ (*.scad)\ do} iterates over each \texttt{.scad} file
\item
  \texttt{\%\%f} is the loop variable (use \texttt{\%f} in interactive CMD, \texttt{\%\%f} in batch files)
\item
  Inside the parentheses, do something with each file
\end{itemize}

\subsubsection*{Loop with a Counter}\label{docs__pandoc__latex__src__cli_complete_curriculum.md__loop-with-a-counter-1}

\textbf{Example: Do something 5 times}

\begin{lstlisting}[style=Alabaster, language=cmd]
@echo off
for /L %%i in (1,1,5) do (
    echo This is iteration number %%i
)

\end{lstlisting}

\textbf{What\textquotesingle s happening:}

\begin{itemize}
\tightlist
\item
  \texttt{for\ /L\ \%\%i\ in\ (start,step,end)} counts from start to end
\item
  \texttt{\%\%i} is the counter variable
\end{itemize}

\begin{center}\rule{0.5\linewidth}{0.5pt}\end{center}

\subsection*{Real-World Example - Batch Processing SCAD Files}\label{docs__pandoc__latex__src__cli_complete_curriculum.md__real-world-example---batch-processing-scad-files-1}

\subsubsection*{Scenario}\label{docs__pandoc__latex__src__cli_complete_curriculum.md__scenario-1}

You have 10 OpenSCAD (.scad) files in a folder. You want to:

\begin{enumerate}
\tightlist
\item
  List them all
\item
  Check how many there are
\item
  For each one, verify it exists
\end{enumerate}

\subsubsection*{The Script}\label{docs__pandoc__latex__src__cli_complete_curriculum.md__the-script-1}

\begin{lstlisting}[style=Alabaster, language=cmd]
@echo off
set scadFolder=C:\Users\YourName\Documents\3D_Projects
set count=0

echo Processing SCAD files in: %scadFolder%
echo.

for %%f in (%scadFolder%\*.scad) do (
    if exist "%%f" (
        echo  Found: %%~nxf
        set /A count+=1
    ) else (
        echo  Missing: %%~nxf
    )
)

echo.
echo Total files found: %count%
echo Batch processing complete!

\end{lstlisting}

\textbf{Breaking it down:}

\begin{itemize}
\tightlist
\item
  \texttt{set\ scadFolder=...} = where to look
\item
  \texttt{for\ \%\%f\ in\ (...\textbackslash{}*.scad)} = find all .scad files
\item
  \texttt{if\ exist\ "\%\%f"} = check if file exists
\item
  \texttt{\%\%\textasciitilde{}nxf} = just the filename with extension (not the full path)
\item
  \texttt{set\ /A\ count+=1} = increment the counter
\end{itemize}

\subsubsection*{Running the Script}\label{docs__pandoc__latex__src__cli_complete_curriculum.md__running-the-script-1}

\begin{enumerate}
\item
  Save as \texttt{batch-process.bat}
\item
  Edit \texttt{scadFolder} to match your real folder
\item
  Run it:

  \begin{lstlisting}[style=Alabaster, language=cmd]
  batch-process.bat

  \end{lstlisting}
\end{enumerate}

\textbf{Screen reader output:}

\begin{lstlisting}[style=Alabaster]
Processing SCAD files in: C:\Users\YourName\Documents\3D_Projects

 Found: model1.scad
 Found: model2.scad
 Found: model3.scad
[... more files ...]

Total files found: 10
Batch processing complete!

\end{lstlisting}

\begin{center}\rule{0.5\linewidth}{0.5pt}\end{center}

\subsection*{Error Handling}\label{docs__pandoc__latex__src__cli_complete_curriculum.md__error-handling-1}

\subsubsection*{Checking Error Levels}\label{docs__pandoc__latex__src__cli_complete_curriculum.md__checking-error-levels}

What if something goes wrong? Use \texttt{errorlevel} checks:

\textbf{Example:}

\begin{lstlisting}[style=Alabaster, language=cmd]
@echo off
type C:\nonexistent\path\file.txt
if %errorlevel% neq 0 (
    echo Error: Could not read file
    echo Errorlevel: %errorlevel%
)

\end{lstlisting}

\textbf{What\textquotesingle s happening:}

\begin{itemize}
\tightlist
\item
  After any command, \texttt{\%errorlevel\%} holds \texttt{0} for success or non-zero for failure
\item
  \texttt{if\ \%errorlevel\%\ neq\ 0} checks for failure
\item
  Handle the error gracefully
\end{itemize}

\textbf{Screen reader advantage:} Errors are announced clearly instead of crashing silently.

\subsubsection*{Validating Input}\label{docs__pandoc__latex__src__cli_complete_curriculum.md__validating-input-1}

\textbf{Example: Make sure a folder exists before processing}

\begin{lstlisting}[style=Alabaster, language=cmd]
@echo off
set folderPath=C:\Users\YourName\Documents

if not exist "%folderPath%\" (
    echo Error: Folder does not exist: %folderPath%
    goto :eof
)

echo Processing folder: %folderPath%
dir /B "%folderPath%"

\end{lstlisting}

\textbf{What\textquotesingle s happening:}

\begin{itemize}
\tightlist
\item
  \texttt{if\ not\ exist\ "\%folderPath\%\textbackslash{}"} checks if folder exists (trailing \texttt{\textbackslash{}} ensures it\textquotesingle s a folder)
\item
  \texttt{goto\ :eof} exits the script early if error
\end{itemize}

\begin{center}\rule{0.5\linewidth}{0.5pt}\end{center}

\subsection*{Debugging Batch Scripts}\label{docs__pandoc__latex__src__cli_complete_curriculum.md__debugging-batch-scripts}

\subsubsection*{Common Errors and Solutions}\label{docs__pandoc__latex__src__cli_complete_curriculum.md__common-errors-and-solutions-1}

\paragraph*{Error 1: "Command not found"}\label{docs__pandoc__latex__src__cli_complete_curriculum.md__error-1-command-not-found-1}

\textbf{Cause:} Typo in command name or program not in PATH

\textbf{Fix:} Check spelling and verify PATH

\begin{lstlisting}[style=Alabaster, language=cmd]
:: Wrong:
dri /B

:: Correct:
dir /B

\end{lstlisting}

\paragraph*{Error 2: "Variable is empty"}\label{docs__pandoc__latex__src__cli_complete_curriculum.md__error-2-variable-is-empty}

\textbf{Cause:} Variable was never set or has a typo

\textbf{Fix:} Make sure variable is set before using it

\begin{lstlisting}[style=Alabaster, language=cmd]
set myvar=hello
echo %myvar%

\end{lstlisting}

\paragraph*{Error 3: "The system cannot find the path specified"}\label{docs__pandoc__latex__src__cli_complete_curriculum.md__error-3-the-system-cannot-find-the-path-specified}

\textbf{Cause:} Wrong folder path

\textbf{Fix:} Verify path exists

\begin{lstlisting}[style=Alabaster, language=cmd]
:: Check if path exists:
if exist "C:\Users\YourName\Documents\" echo Path exists

\end{lstlisting}

\paragraph*{Error 4: "Access denied"}\label{docs__pandoc__latex__src__cli_complete_curriculum.md__error-4-access-denied-1}

\textbf{Cause:} Don\textquotesingle t have permission

\textbf{Fix:} Run CMD as administrator

\begin{itemize}
\tightlist
\item
  Right-click Command Prompt -\textgreater{} Run as administrator
\end{itemize}

\subsubsection*{Debugging Technique: Trace Output}\label{docs__pandoc__latex__src__cli_complete_curriculum.md__debugging-technique-trace-output-1}

Add \texttt{echo} statements to track what\textquotesingle s happening:

\begin{lstlisting}[style=Alabaster, language=cmd]
@echo off
set path_var=C:\Users\YourName\Documents
echo Starting script. Path is: %path_var%

for %%f in (%path_var%\*) do (
    echo Processing: %%f
    :: Do something with %%f
    echo Done with: %%f
)

echo Script complete

\end{lstlisting}

Your screen reader will announce each step, so you know where errors happen.

\begin{center}\rule{0.5\linewidth}{0.5pt}\end{center}

\subsection*{Creating Professional Workflows}\label{docs__pandoc__latex__src__cli_complete_curriculum.md__creating-professional-workflows-1}

\subsubsection*{Example 1: Automated Project Setup}\label{docs__pandoc__latex__src__cli_complete_curriculum.md__example-1-automated-project-setup-1}

\textbf{Scenario:} You start a new 3D printing project regularly. Instead of creating folders manually:

\begin{lstlisting}[style=Alabaster, language=cmd]
@echo off
set /p projectName=Enter project name: 
set baseFolder=C:\Users\YourName\Documents\3D_Projects
set projectFolder=%baseFolder%\%projectName%

:: Create folder structure
mkdir "%projectFolder%"
mkdir "%projectFolder%\designs"
mkdir "%projectFolder%\output"
mkdir "%projectFolder%\notes"

:: Create a README
echo # %projectName% > "%projectFolder%\README.txt"
echo. >> "%projectFolder%\README.txt"
echo Created: %date% >> "%projectFolder%\README.txt"
echo. >> "%projectFolder%\README.txt"
echo ## Designs >> "%projectFolder%\README.txt"
echo All .scad files go here. >> "%projectFolder%\README.txt"
echo. >> "%projectFolder%\README.txt"
echo ## Output >> "%projectFolder%\README.txt"
echo STL and other exports go here. >> "%projectFolder%\README.txt"

echo Project setup complete: %projectFolder%

\end{lstlisting}

\textbf{What it does:}

\begin{itemize}
\tightlist
\item
  Prompts for a project name
\item
  Creates folder structure for a new project
\item
  Sets up subfolders for designs, output, notes
\item
  Creates a README file automatically
\end{itemize}

\subsubsection*{Example 2: Batch File Verification}\label{docs__pandoc__latex__src__cli_complete_curriculum.md__example-2-batch-file-verification-1}

\textbf{Scenario:} Before processing, verify all required files exist:

\begin{lstlisting}[style=Alabaster, language=cmd]
@echo off
set projectFolder=C:\Users\YourName\Documents\3D_Projects\MyKeychain
set allGood=1

for %%i in (README.txt designs output notes) do (
    if exist "%projectFolder%\%%i" (
        echo  Found: %%i
    ) else (
        echo  Missing: %%i
        set allGood=0
    )
)

if %allGood%==1 (
    echo All checks passed!
) else (
    echo Some files are missing!
)

\end{lstlisting}

\begin{center}\rule{0.5\linewidth}{0.5pt}\end{center}

\subsection*{Screen Reader Tips for Batch Scripts}\label{docs__pandoc__latex__src__cli_complete_curriculum.md__screen-reader-tips-for-batch-scripts}

\subsubsection*{Making Script Output Readable}\label{docs__pandoc__latex__src__cli_complete_curriculum.md__making-script-output-readable-1}

\textbf{Problem:} Script runs but output scrolls too fast or is hard to follow

\textbf{Solution 1: Save to file}

\begin{lstlisting}[style=Alabaster, language=cmd]
my-script.bat > output.txt
notepad.exe output.txt

\end{lstlisting}

\textbf{Solution 2: Use echo with clear sections}

\begin{lstlisting}[style=Alabaster, language=cmd]
echo ========== STARTING ==========
echo.
:: ... script ...
echo.
echo ========== COMPLETE ==========

\end{lstlisting}

\textbf{Solution 3: Pause between major sections}

\begin{lstlisting}[style=Alabaster, language=cmd]
echo Pausing... Press any key to continue
pause > nul

\end{lstlisting}

Your screen reader will announce the pause, give you time to read output.

\subsubsection*{Announcing Progress}\label{docs__pandoc__latex__src__cli_complete_curriculum.md__announcing-progress-1}

\textbf{For long-running scripts:}

\begin{lstlisting}[style=Alabaster, language=cmd]
@echo off
set count=0
for %%f in (*.scad) do (
    set /A count+=1
    echo Processing file %%f
)
echo All %count% files processed!

\end{lstlisting}

\begin{center}\rule{0.5\linewidth}{0.5pt}\end{center}

\subsection*{Practice Exercises}\label{docs__pandoc__latex__src__cli_complete_curriculum.md__practice-exercises-5}

\subsubsection*{Exercise 1: Your First Batch Script}\label{docs__pandoc__latex__src__cli_complete_curriculum.md__exercise-1-your-first-batch-script}

\textbf{Goal:} Create and run a simple batch script

\textbf{Steps:}

\begin{enumerate}
\item
  Create file: \texttt{notepad.exe\ hello.bat}
\item
  Type:

  \begin{lstlisting}[style=Alabaster, language=cmd]
  @echo off
  echo Hello from my first CMD batch script!
  cd
  dir /B

  \end{lstlisting}
\item
  Save and run: \texttt{hello.bat}
\end{enumerate}

\textbf{Checkpoint:} You should see output for each command.

\subsubsection*{Exercise 2: Script with a Variable}\label{docs__pandoc__latex__src__cli_complete_curriculum.md__exercise-2-script-with-a-variable-1}

\textbf{Goal:} Use a variable to make the script flexible

\textbf{Steps:}

\begin{enumerate}
\item
  Create file: \texttt{notepad.exe\ smart-listing.bat}
\item
  Type:

  \begin{lstlisting}[style=Alabaster, language=cmd]
  @echo off
  set targetFolder=C:\Users\YourName\Documents
  echo Listing contents of: %targetFolder%
  dir /B "%targetFolder%"

  \end{lstlisting}
\item
  Edit \texttt{targetFolder} to a real folder on your computer
\item
  Run: \texttt{smart-listing.bat}
\end{enumerate}

\textbf{Checkpoint:} You should see listing of that specific folder.

\subsubsection*{Exercise 3: Subroutine}\label{docs__pandoc__latex__src__cli_complete_curriculum.md__exercise-3-subroutine}

\textbf{Goal:} Create a reusable subroutine

\textbf{Steps:}

\begin{enumerate}
\item
  Create file: \texttt{notepad.exe\ navigate.bat}
\item
  Type:

  \begin{lstlisting}[style=Alabaster, language=cmd]
  @echo off
  call :GoTo "C:\Users\YourName\Documents"
  call :GoTo "C:\Users\YourName\Downloads"
  goto :eof

  :GoTo
  if exist "%~1\" (
      cd /D "%~1"
      echo Now in:
      cd
      echo Contents:
      dir /B
  ) else (
      echo Path does not exist: %~1
  )
  goto :eof

  \end{lstlisting}
\item
  Run: \texttt{navigate.bat}
\end{enumerate}

\textbf{Checkpoint:} Both subroutine calls should work, showing contents of each folder.

\subsubsection*{Exercise 4: Loop}\label{docs__pandoc__latex__src__cli_complete_curriculum.md__exercise-4-loop-1}

\textbf{Goal:} Use a loop to repeat an action

\textbf{Steps:}

\begin{enumerate}
\item
  Create file: \texttt{notepad.exe\ repeat.bat}
\item
  Type:

  \begin{lstlisting}[style=Alabaster, language=cmd]
  @echo off
  echo Demonstrating a loop:
  for /L %%i in (1,1,5) do (
      echo Iteration %%i: Hello!
  )
  echo Loop complete!

  \end{lstlisting}
\item
  Run: \texttt{repeat.bat}
\end{enumerate}

\textbf{Checkpoint:} Should print "Iteration 1" through "Iteration 5".

\subsubsection*{Exercise 5: Real-World Script}\label{docs__pandoc__latex__src__cli_complete_curriculum.md__exercise-5-real-world-script-1}

\textbf{Goal:} Create a useful script for a real task

\textbf{Steps:}

\begin{enumerate}
\item
  Create a folder: \texttt{mkdir\ C:\textbackslash{}Users\textbackslash{}YourName\textbackslash{}Documents\textbackslash{}TestFiles}
\item
  Create some test files:

  \begin{lstlisting}[style=Alabaster, language=cmd]
  echo test > C:\Users\YourName\Documents\TestFiles\file1.txt
  echo test > C:\Users\YourName\Documents\TestFiles\file2.txt
  echo test > C:\Users\YourName\Documents\TestFiles\file3.txt

  \end{lstlisting}
\item
  Create script: \texttt{notepad.exe\ report.bat}
\item
  Type:

  \begin{lstlisting}[style=Alabaster, language=cmd]
  @echo off
  set folder=C:\Users\YourName\Documents\TestFiles
  set count=0

  echo === FILE REPORT ===
  echo Folder: %folder%
  echo.
  echo Files:
  for %%f in ("%folder%\*") do (
      echo   - %%~nxf
      set /A count+=1
  )
  echo.
  echo Total: %count% files
  echo === END REPORT ===

  \end{lstlisting}
\item
  Run: \texttt{report.bat}
\end{enumerate}

\textbf{Checkpoint:} Should show report of all files in the test folder.

\begin{center}\rule{0.5\linewidth}{0.5pt}\end{center}

\subsection*{Quiz - Lesson CMD.6}\label{docs__pandoc__latex__src__cli_complete_curriculum.md__quiz---lesson-cmd6}

\begin{enumerate}
\tightlist
\item
  What is a CMD batch script?
\item
  What file extension do CMD batch scripts use?
\item
  What is a variable in a batch script and how do you create one?
\item
  What is a subroutine (\texttt{:label}) and why would you use one?
\item
  How do you run a batch script?
\item
  What is a \texttt{for} loop and what does \texttt{for\ \%\%f\ in\ (*.scad)\ do} do?
\item
  What does \texttt{if\ exist} do?
\item
  How do you handle errors in a batch script?
\item
  When would you use \texttt{if\ \%errorlevel\%\ neq\ 0}?
\item
  What technique makes batch script output readable for screen readers?
\end{enumerate}

\begin{center}\rule{0.5\linewidth}{0.5pt}\end{center}

\subsection*{Extension Problems}\label{docs__pandoc__latex__src__cli_complete_curriculum.md__extension-problems-13}

\begin{enumerate}
\tightlist
\item
  \textbf{Auto-Backup Script:} Create a batch script that copies all files from one folder to another, announcing progress
\item
  \textbf{File Counter:} Write a subroutine that counts files by extension (.txt, .scad, .stl, etc.)
\item
  \textbf{Folder Cleaner:} Batch script that deletes files older than 30 days (with user confirmation)
\item
  \textbf{Project Template:} Subroutine that creates a complete project folder structure with all needed files
\item
  \textbf{Batch Rename:} Script that renames all files in a folder according to a pattern
\item
  \textbf{Log Generator:} Create a script that records what it does to a log file for later review
\item
  \textbf{Scheduled Task:} Set up a batch script to run automatically every day at a specific time
\item
  \textbf{File Verifier:} Check that all SCAD files in a folder have corresponding STL exports
\item
  \textbf{Report Generator:} Create a summary report of all projects in a folder
\item
  \textbf{Error Tracker:} Script that lists all commands that had errors and logs them with timestamps
\end{enumerate}

\begin{center}\rule{0.5\linewidth}{0.5pt}\end{center}

\subsection*{Important Notes}\label{docs__pandoc__latex__src__cli_complete_curriculum.md__important-notes-1}

\begin{itemize}
\tightlist
\item
  \textbf{Always test scripts on small sets of files first} before running them on important data
\item
  \textbf{Save your work regularly} --- use version naming if possible
\item
  \textbf{Test error handling} --- make sure errors don\textquotesingle t crash silently
\item
  \textbf{Document your scripts} --- use \texttt{::} comments so you remember what each part does
\item
  \textbf{Backup before batch operations} --- if something goes wrong, you have the original
\end{itemize}

\begin{center}\rule{0.5\linewidth}{0.5pt}\end{center}

\subsection*{References}\label{docs__pandoc__latex__src__cli_complete_curriculum.md__references-11}

\begin{itemize}
\tightlist
\item
  \textbf{Microsoft CMD Batch Scripting Guide:} \url{https://example.com}
\item
  \textbf{FOR Loop Documentation:} \url{https://example.com}
\item
  \textbf{IF Statement Reference:} \url{https://example.com}
\item
  \textbf{SET Variable Reference:} \url{https://example.com}
\end{itemize}

\begin{center}\rule{0.5\linewidth}{0.5pt}\end{center}

\textbf{Next Steps:} After mastering this lesson, explore advanced batch scripting, scheduled tasks, and 3D printing integration in the main curriculum.

\begin{center}\rule{0.5\linewidth}{0.5pt}\end{center}

\protect\phantomsection\label{docs__pandoc__latex__src__cli_complete_curriculum.md__cmd-unit-exam}{}

\section*{Unit Exam --- Comprehensive Assessment}\label{docs__pandoc__latex__src__cli_complete_curriculum.md__unit-exam--comprehensive-assessment-1}

\begin{center}\rule{0.5\linewidth}{0.5pt}\end{center}

Estimated time: 60-90 minutes

\subsection*{Key Learning Outcomes Assessed}\label{docs__pandoc__latex__src__cli_complete_curriculum.md__key-learning-outcomes-assessed-1}

By completing this unit test, you will demonstrate:

\begin{enumerate}
\tightlist
\item
  Understanding of file system navigation and path concepts
\item
  Proficiency with file and folder manipulation commands
\item
  Ability to redirect and pipe command output
\item
  Knowledge of environment variables and aliases
\item
  Screen-reader accessibility best practices in terminal environments
\item
  Problem-solving and command chaining skills
\end{enumerate}

\subsection*{Target Audience}\label{docs__pandoc__latex__src__cli_complete_curriculum.md__target-audience-1}

Users who have completed Command Line Lesson 0 through Command Line Lesson 6 and need to demonstrate mastery of Command Prompt fundamentals.

\subsection*{Instructions}\label{docs__pandoc__latex__src__cli_complete_curriculum.md__instructions-1}

Complete all sections below. For multiple choice, select the best answer. For short answers, write one to two sentences. For hands-on tasks, capture evidence (screenshots or output files) and submit alongside your answers.

\begin{center}\rule{0.5\linewidth}{0.5pt}\end{center}

\subsection*{Part A: Multiple Choice Questions (20 questions)}\label{docs__pandoc__latex__src__cli_complete_curriculum.md__part-a-multiple-choice-questions-20-questions-1}

Select the best answer for each question. Each question is worth 1 point.

\begin{enumerate}
\item
  What is the primary purpose of the \texttt{PATH} environment variable?

  \begin{itemize}
  \tightlist
  \item
    A) Store your home directory location
  \item
    B) Tell the shell where to find executable programs
  \item
    C) Configure the visual appearance of the terminal
  \item
    D) Store the current working directory name
  \end{itemize}
\item
  Which command shows your current working directory in CMD?

  \begin{itemize}
  \tightlist
  \item
    A) \texttt{dir\ /B}
  \item
    B) \texttt{cd} (with no arguments)
  \item
    C) \texttt{pwd}
  \item
    D) \texttt{whoami}
  \end{itemize}
\item
  What does \texttt{\%USERPROFILE\%} represent in CMD?

  \begin{itemize}
  \tightlist
  \item
    A) The root directory
  \item
    B) The current directory
  \item
    C) The parent directory
  \item
    D) Your home directory
  \end{itemize}
\item
  How do you list only file names (not full details) in a screen-reader-friendly way?

  \begin{itemize}
  \tightlist
  \item
    A) \texttt{dir}
  \item
    B) \texttt{dir\ /B}
  \item
    C) \texttt{dir\ /L}
  \item
    D) \texttt{type\ /B}
  \end{itemize}
\item
  Which command creates a new empty file in CMD?

  \begin{itemize}
  \tightlist
  \item
    A) \texttt{mkdir\ filename}
  \item
    B) \texttt{echo.\ \textgreater{}\ filename}
  \item
    C) \texttt{touch\ filename}
  \item
    D) \texttt{new\ filename}
  \end{itemize}
\item
  What is the difference between \texttt{\textgreater{}} and \texttt{\textgreater{}\textgreater{}}?

  \begin{itemize}
  \tightlist
  \item
    A) \texttt{\textgreater{}} redirects to file, \texttt{\textgreater{}\textgreater{}} displays on screen
  \item
    B) \texttt{\textgreater{}} overwrites a file, \texttt{\textgreater{}\textgreater{}} appends to a file
  \item
    C) They do the same thing
  \item
    D) \texttt{\textgreater{}} is for text, \texttt{\textgreater{}\textgreater{}} is for binary
  \end{itemize}
\item
  What does the pipe operator \texttt{\textbar{}} do?

  \begin{itemize}
  \tightlist
  \item
    A) Creates a folder
  \item
    B) Sends the output of one command to the input of another
  \item
    C) Deletes files matching a pattern
  \item
    D) Lists all processes
  \end{itemize}
\item
  Which command copies a file in CMD?

  \begin{itemize}
  \tightlist
  \item
    A) \texttt{move}
  \item
    B) \texttt{del}
  \item
    C) \texttt{copy}
  \item
    D) \texttt{cd}
  \end{itemize}
\item
  How do you rename a file from \texttt{oldname.txt} to \texttt{newname.txt} in CMD?

  \begin{itemize}
  \tightlist
  \item
    A) \texttt{copy\ oldname.txt\ newname.txt}
  \item
    B) \texttt{move\ oldname.txt\ newname.txt}
  \item
    C) \texttt{rename\ oldname.txt\ newname.txt}
  \item
    D) Both B and C are correct
  \end{itemize}
\item
  What is the purpose of \texttt{find} in CMD piping?

  \begin{itemize}
  \tightlist
  \item
    A) Find files in a directory
  \item
    B) Search for text patterns within output
  \item
    C) Find a string to copy to clipboard
  \item
    D) Find which shell to use
  \end{itemize}
\item
  Which key allows you to autocomplete a path in CMD?

  \begin{itemize}
  \tightlist
  \item
    A) \texttt{Ctrl\ +\ A}
  \item
    B) \texttt{Ctrl\ +\ E}
  \item
    C) \texttt{Tab}
  \item
    D) \texttt{Space}
  \end{itemize}
\item
  How do you copy text to the Windows clipboard from CMD?

  \begin{itemize}
  \tightlist
  \item
    A) \texttt{type\ filename\ \textgreater{}\ clipboard}
  \item
    B) \texttt{type\ filename\ \textbar{}\ clip}
  \item
    C) \texttt{copy\ filename}
  \item
    D) \texttt{type\ filename\ \textbar{}\ paste}
  \end{itemize}
\item
  What does \texttt{where\ openscad} do?

  \begin{itemize}
  \tightlist
  \item
    A) Opens the OpenSCAD application
  \item
    B) Gets help about the openscad command
  \item
    C) Locates the full path of the openscad executable
  \item
    D) Lists all available commands
  \end{itemize}
\item
  Which wildcard matches any single character?

  \begin{itemize}
  \tightlist
  \item
    A) \texttt{*}
  \item
    B) \texttt{?}
  \item
    C) \texttt{\%}
  \item
    D) \texttt{\#}
  \end{itemize}
\item
  What is the purpose of the \texttt{start} command?

  \begin{itemize}
  \tightlist
  \item
    A) Run a script or executable, optionally in a new window
  \item
    B) Execute all commands in parallel
  \item
    C) Combine multiple commands
  \item
    D) Create an alias
  \end{itemize}
\item
  How do you create a temporary alias (macro) in CMD?

  \begin{itemize}
  \tightlist
  \item
    A) \texttt{alias\ preview=\textquotesingle{}openscad\textquotesingle{}}
  \item
    B) \texttt{doskey\ preview=openscad\ \$*}
  \item
    C) \texttt{set\ preview=openscad}
  \item
    D) \texttt{macro\ preview\ openscad}
  \end{itemize}
\item
  How can doskey macros be made to persist across CMD sessions?

  \begin{itemize}
  \tightlist
  \item
    A) They are automatically saved
  \item
    B) By adding them to a startup batch script registered as an autorun
  \item
    C) By using the Windows Control Panel
  \item
    D) Doskey macros cannot be made permanent
  \end{itemize}
\item
  How do you abort a long-running command in CMD?

  \begin{itemize}
  \tightlist
  \item
    A) Press \texttt{Escape}
  \item
    B) Press \texttt{Ctrl\ +\ X}
  \item
    C) Press \texttt{Ctrl\ +\ C}
  \item
    D) Press \texttt{Alt\ +\ F4}
  \end{itemize}
\item
  What command shows the history of previously run commands in CMD?

  \begin{itemize}
  \tightlist
  \item
    A) \texttt{history}
  \item
    B) \texttt{doskey\ /history}
  \item
    C) \texttt{F7} (opens history popup)
  \item
    D) Both B and C are correct
  \end{itemize}
\item
  How do you make an environment variable permanent in CMD (for all future sessions)?

  \begin{itemize}
  \tightlist
  \item
    A) Use \texttt{set} in the terminal every time
  \item
    B) Use \texttt{setx} to write it to the registry
  \item
    C) Use the Windows Control Panel only
  \item
    D) Environment variables cannot be made permanent in CMD
  \end{itemize}
\end{enumerate}

\begin{center}\rule{0.5\linewidth}{0.5pt}\end{center}

\subsection*{Part B: Short Answer Questions (10 questions)}\label{docs__pandoc__latex__src__cli_complete_curriculum.md__part-b-short-answer-questions-10-questions-1}

Answer each question in one to two sentences. Each question is worth 2 points.

\begin{enumerate}
\item
  Explain the difference between absolute and relative paths. Give one example of each.
\item
  Why is \texttt{dir\ /B} preferred over \texttt{dir} for screen reader users? Describe what flag you would add to list only files.
\item
  What is the purpose of redirecting output to a file, and give an example of when you would use \texttt{\textgreater{}} instead of \texttt{\textgreater{}\textgreater{}}?
\end{enumerate}

Warning

\texttt{rmdir\ /S\ /Q\ C:\textbackslash{}Documents\textbackslash{}my\_folder} permanently deletes the folder and \textbf{all its contents} with no confirmation prompt and no undo. Always run \texttt{dir\ /B\ C:\textbackslash{}Documents\textbackslash{}my\_folder} first to confirm what you\textquotesingle re about to destroy.

\begin{enumerate}
\setcounter{enumi}{3}
\item
  Describe what would happen if you ran \texttt{rmdir\ /S\ /Q\ C:\textbackslash{}Documents\textbackslash{}my\_folder} and why this command should be used carefully.
\item
  How would you search for all files with a \texttt{.scad} extension in your current directory? Write the command.
\item
  Explain what happens when you pipe the output of \texttt{dir\ /B} into \texttt{clip}. What would you do next?
\item
  What is an environment variable, and give one example of how you might use it in CMD.
\item
  If a program is not in your \texttt{PATH}, what two methods could you use to run it from CMD?
\item
  Describe how you would open a file in Notepad and also add a line to it from CMD.
\item
  What is one strategy you would use if your screen reader stops announcing terminal output while using CMD?
\end{enumerate}

\begin{center}\rule{0.5\linewidth}{0.5pt}\end{center}

\subsection*{Part C: Hands-On Tasks (10 tasks)}\label{docs__pandoc__latex__src__cli_complete_curriculum.md__part-c-hands-on-tasks-10-tasks-1}

Complete each task and capture evidence (screenshots, output files, or command transcripts). Each task is worth 3 points.

\subsubsection*{Tasks 1-5: File System and Navigation}\label{docs__pandoc__latex__src__cli_complete_curriculum.md__tasks-1-5-file-system-and-navigation-1}

\begin{enumerate}
\item
  Create a folder structure \texttt{\%USERPROFILE\%\textbackslash{}Documents\textbackslash{}CMD\_Assessment\textbackslash{}Projects} using a single command. Capture the \texttt{dir\ /B} output showing the creation.
\item
  Create five files named \texttt{project\_1.scad}, \texttt{project\_2.scad}, \texttt{project\_3.txt}, \texttt{notes\_1.txt}, and \texttt{notes\_2.txt} inside the \texttt{Projects} folder. Use wildcards to list only \texttt{.scad} files, then capture the output.
\item
  Copy the entire \texttt{Projects} folder to \texttt{Projects\_Backup} using \texttt{xcopy\ /E\ /I}. Capture the \texttt{dir\ /B} output showing both folders exist.
\item
  Move (rename) \texttt{project\_1.scad} to \texttt{project\_1\_final.scad}. Capture the \texttt{dir\ /B} output showing the renamed file.
\item
  Delete \texttt{notes\_1.txt} and \texttt{notes\_2.txt} using a single \texttt{del} command with wildcards. Capture the final \texttt{dir\ /B} output.
\end{enumerate}

\subsubsection*{Tasks 6-10: Advanced Operations and Scripting}\label{docs__pandoc__latex__src__cli_complete_curriculum.md__tasks-6-10-advanced-operations-and-scripting-1}

\begin{enumerate}
\item
  Create a file called \texttt{my\_data.txt} with at least four lines using \texttt{echo} and \texttt{\textgreater{}\textgreater{}}. Then read it with \texttt{type\ my\_data.txt} and capture the output.
\item
  Use \texttt{find} to search for a keyword (e.g., "project") in \texttt{my\_data.txt} and pipe the results to \texttt{clip}. Paste the results into Notepad and capture a screenshot.
\item
  List all files in the \texttt{Projects} folder and redirect the output to \texttt{projects\_list.txt}. Open it in Notepad and capture a screenshot of the file.
\item
  Create a temporary \texttt{doskey} alias called \texttt{mydir} that runs \texttt{dir\ /B}, test it, and capture the output. Then explain what would be required to make it persistent.
\item
  Run \texttt{help\ dir} and redirect the output to \texttt{help\_output.txt}. Open the file in Notepad and capture a screenshot showing at least the first page of help content.
\end{enumerate}

\begin{center}\rule{0.5\linewidth}{0.5pt}\end{center}

\subsection*{Grading Rubric}\label{docs__pandoc__latex__src__cli_complete_curriculum.md__grading-rubric-1}

{\def\LTcaptype{none} % do not increment counter
\begin{longtable}[]{@{}llll@{}}
\toprule\noalign{}
Section & Questions & Points Each & Total \\
\midrule\noalign{}
\endhead
\bottomrule\noalign{}
\endlastfoot
Multiple Choice & 20 & 1 & 20 \\
Short Answer & 10 & 2 & 20 \\
Hands-On Tasks & 10 & 3 & 30 \\
\textbf{Total} & \textbf{40} & - & \textbf{70} \\
\end{longtable}
}

\textbf{Passing Score:} 49 points (70\%)

\begin{center}\rule{0.5\linewidth}{0.5pt}\end{center}

\subsection*{Helpful Resources for Review}\label{docs__pandoc__latex__src__cli_complete_curriculum.md__helpful-resources-for-review-1}

\begin{itemize}
\tightlist
\item
  \href{https://learn.microsoft.com/en-us/windows-server/administration/windows-commands/windows-commands-glossary}{CMD Command Reference}
\item
  \href{https://learn.microsoft.com/en-us/windows-server/administration/windows-commands/cd}{Navigation and File System}
\item
  \href{https://learn.microsoft.com/en-us/windows-server/administration/windows-commands/find}{Using Pipes and Filtering}
\item
  \href{https://learn.microsoft.com/en-us/windows-server/administration/windows-commands/doskey}{DOSKEY and Aliases}
\item
  \href{https://learn.microsoft.com/en-us/windows-server/administration/windows-commands/}{Screen Reader Accessibility Tips}
\end{itemize}

\begin{center}\rule{0.5\linewidth}{0.5pt}\end{center}

\subsection*{Submission Checklist}\label{docs__pandoc__latex__src__cli_complete_curriculum.md__submission-checklist-1}

\begin{itemize}
\tightlist
\item[$\square$]
  All 20 multiple choice questions answered
\item[$\square$]
  All 10 short answer questions answered (1-2 sentences each)
\item[$\square$]
  All 10 hands-on tasks completed with evidence captured
\item[$\square$]
  Files/screenshots organized and labeled clearly
\item[$\square$]
  Submission includes this checklist
\end{itemize}

\begin{center}\rule{0.5\linewidth}{0.5pt}\end{center}

\protect\phantomsection\label{docs__pandoc__latex__src__cli_complete_curriculum.md__part-3-gitbash}{}

\section*{Part 3 --- Git Bash}\label{docs__pandoc__latex__src__cli_complete_curriculum.md__part-3--git-bash-1}

\begin{center}\rule{0.5\linewidth}{0.5pt}\end{center}

\textbf{Last Updated:} February 2026\\
\textbf{Total Duration:} 30--45 hours of instruction and practice (for screen reader users)\\
\textbf{Target Users:} Anyone using a screen reader --- NVDA, JAWS, Windows Narrator, and Dolphin SuperNova are all covered throughout this curriculum.\\
\textbf{Platform:} Git Bash on Windows (commands transfer directly to macOS Terminal and Linux)\\
\textbf{Alternate to:} PowerShell curriculum and CMD curriculum (same concepts, Bash/Unix syntax)

\emph{This single file contains the complete Git Bash Foundation curriculum in reading order. All lessons, quizzes, extension problems, and the unit exam are included. No additional files are needed.}

\begin{center}\rule{0.5\linewidth}{0.5pt}\end{center}

\subsection*{Table of Contents}\label{docs__pandoc__latex__src__cli_complete_curriculum.md__table-of-contents-4}

{\def\LTcaptype{none} % do not increment counter
\begin{longtable}[]{@{}
  >{\raggedright\arraybackslash}p{(\linewidth - 4\tabcolsep) * \real{0.3333}}
  >{\raggedright\arraybackslash}p{(\linewidth - 4\tabcolsep) * \real{0.3333}}
  >{\raggedright\arraybackslash}p{(\linewidth - 4\tabcolsep) * \real{0.3333}}@{}}
\toprule\noalign{}
\begin{minipage}[b]{\linewidth}\raggedright
Section
\end{minipage} & \begin{minipage}[b]{\linewidth}\raggedright
Title
\end{minipage} & \begin{minipage}[b]{\linewidth}\raggedright
Estimated Time
\end{minipage} \\
\midrule\noalign{}
\endhead
\bottomrule\noalign{}
\endlastfoot
\hyperref[docs__pandoc__latex__src__cli_complete_curriculum.md__gb-accessibility-guide]{Screen Reader Accessibility Guide}
& NVDA, JAWS, Narrator, and SuperNova reference for Git Bash &
1--1.5 hours \\
\hyperref[docs__pandoc__latex__src__cli_complete_curriculum.md__gb-tutorial]{Tutorial: Git Bash Quick Start}
& Prompt orientation, core commands, path concepts & 30--45 minutes \\
\hyperref[docs__pandoc__latex__src__cli_complete_curriculum.md__gb-lesson-pre]{Lesson Pre: Your First Terminal}
& Installing Git Bash, opening it, first commands & 1.5--2 hours \\
\hyperref[docs__pandoc__latex__src__cli_complete_curriculum.md__gb-lesson-0]{Lesson 0: Getting Started}
&
Paths, \texttt{\textasciitilde{}}, \texttt{./}, \texttt{../}, tab completion
& 20--30 minutes \\
\hyperref[docs__pandoc__latex__src__cli_complete_curriculum.md__gb-lesson-1]{Lesson 1: Navigation}
& \texttt{pwd}, \texttt{cd}, \texttt{ls} and flags, wildcards &
1 class period \\
\hyperref[docs__pandoc__latex__src__cli_complete_curriculum.md__gb-lesson-2]{Lesson 2: File and Folder Manipulation}
&
\texttt{touch}, \texttt{mkdir}, \texttt{cp}, \texttt{mv}, \texttt{rm}, \texttt{rmdir}
& 30--45 minutes \\
\hyperref[docs__pandoc__latex__src__cli_complete_curriculum.md__gb-lesson-3]{Lesson 3: Input, Output, and Piping}
&
\texttt{echo}, \texttt{cat}, \texttt{\textgreater{}}, \texttt{\textgreater{}\textgreater{}}, `
& \texttt{,\ }clip\texttt{,\ }grep` \\
\hyperref[docs__pandoc__latex__src__cli_complete_curriculum.md__gb-lesson-4]{Lesson 4: Environment Variables and Aliases}
&
\texttt{\$VARNAME}, \texttt{PATH}, \texttt{which}, \texttt{alias}, \texttt{.bashrc}
& 30--45 minutes \\
\hyperref[docs__pandoc__latex__src__cli_complete_curriculum.md__gb-lesson-5]{Lesson 5: Filling in the Gaps}
&
\texttt{history}, \texttt{Ctrl+R}, \texttt{Ctrl+C}, \texttt{.bashrc}, \texttt{./}
& 30--45 minutes \\
\hyperref[docs__pandoc__latex__src__cli_complete_curriculum.md__gb-lesson-6]{Lesson 6: Advanced Techniques}
& Shell scripts, functions, loops, professional workflows &
4--4.5 hours \\
\hyperref[docs__pandoc__latex__src__cli_complete_curriculum.md__gb-unit-exam]{Unit Exam}
& Comprehensive assessment --- Lessons 0 through 6 & 60--90 minutes \\
\end{longtable}
}

\textbf{Total:} approximately 30--45 hours including practice exercises

\begin{center}\rule{0.5\linewidth}{0.5pt}\end{center}

\subsection*{How to Use This File}\label{docs__pandoc__latex__src__cli_complete_curriculum.md__how-to-use-this-file-2}

\textbf{If you have never used a terminal before:} start with the Accessibility Guide, then the Tutorial, then Lesson Pre. Read in order from the top --- each lesson builds on the previous one.

\textbf{If you have used a terminal (CMD or PowerShell) but not Git Bash:} read the Tutorial section for the key differences (Unix paths, forward slashes, \texttt{ls} vs \texttt{dir}, \texttt{cat} vs \texttt{type}). Then start at Lesson 0.

\textbf{If you are experienced with Bash on macOS or Linux:} jump to the Accessibility Guide for screen-reader-specific techniques, then skip to whatever lesson covers your gap.

\textbf{Reference while working:} the Accessibility Guide and Tutorial contain material you will return to throughout the curriculum. Bookmark those sections.

\textbf{CMD or PowerShell users:} this curriculum covers the same concepts as the CMD and PowerShell curricula but uses Bash/Unix syntax. Commands learned here transfer directly to macOS Terminal and Linux. Choose one pathway and stay consistent.

\begin{center}\rule{0.5\linewidth}{0.5pt}\end{center}

\subsection*{Curriculum Rules (Read These First)}\label{docs__pandoc__latex__src__cli_complete_curriculum.md__curriculum-rules-read-these-first-2}

Important

\textbf{Rule 1 --- Always know where you are.} Start every session with \texttt{pwd}. If you don\textquotesingle t know your current path, you risk operating on the wrong files or folders.

\textbf{Rule 1 --- Always know where you are.}\\
Every session, first thing: \texttt{pwd}. If you do not know your path, you will get lost.

Warning

\textbf{There is no Recycle Bin in the terminal.} Files deleted with \texttt{rm} are permanently gone immediately --- no undo, no recovery. Always run \texttt{ls\ -1} to confirm exactly what you\textquotesingle re about to delete before running any \texttt{rm} command.

\textbf{Rule 2 --- Check before you delete.}\\
Before any \texttt{rm} or \texttt{rm\ -r} command: \texttt{ls\ -1} to confirm exactly what you are about to remove. Deleted files are gone immediately --- there is no Recycle Bin in the terminal.

Tip

Always use \texttt{ls\ -1} (the number \textbf{1}, not the letter L) --- never bare \texttt{ls}. The \texttt{-1} flag outputs one name per line, which screen readers can follow reliably. Bare \texttt{ls} outputs a multi-column grid that is difficult to navigate audibly.

\textbf{Rule 3 --- Use \texttt{ls\ -1} for listings.}\\
Always \texttt{ls\ -1} --- never bare \texttt{ls}. The \texttt{-1} flag gives one name per line, which all screen readers can follow reliably. (That is the number one, not the letter L.)

Tip

When output is long or confusing, redirect it to a file: \texttt{command\ \textgreater{}\ output.txt}, then open with \texttt{notepad.exe\ output.txt}. This is always clearer than reading raw terminal output and works with all four screen readers.

\textbf{Rule 4 --- When output is confusing, redirect to a file.}\\
\texttt{command\ \textgreater{}\ output.txt} then \texttt{notepad.exe\ output.txt}. This works with all four screen readers and is always clearer than reading raw terminal output.

\textbf{Rule 5 --- Save everything you create.}\\
Keep all practice files in a dedicated learning folder: \texttt{mkdir\ -p\ \textasciitilde{}/Documents/GitBash\_Learning}

\begin{center}\rule{0.5\linewidth}{0.5pt}\end{center}

\begin{center}\rule{0.5\linewidth}{0.5pt}\end{center}

\protect\phantomsection\label{docs__pandoc__latex__src__cli_complete_curriculum.md__gb-accessibility-guide}{}

\section*{Screen Reader Accessibility Guide for Git Bash}\label{docs__pandoc__latex__src__cli_complete_curriculum.md__screen-reader-accessibility-guide-for-git-bash}

\begin{center}\rule{0.5\linewidth}{0.5pt}\end{center}

Target Users: NVDA, JAWS, and other screen reader users\\
Last Updated: 2026

This guide is used throughout the Git Bash Foundation curriculum to help screen reader users navigate and work efficiently with the terminal.

\subsection*{Table of Contents}\label{docs__pandoc__latex__src__cli_complete_curriculum.md__table-of-contents-5}

\begin{enumerate}
\tightlist
\item
  \hyperref[docs__pandoc__latex__src__cli_complete_curriculum.md__gb-getting-started]{Getting Started with Screen Readers}
\item
  \hyperref[docs__pandoc__latex__src__cli_complete_curriculum.md__gb-nvda-tips]{NVDA-Specific Tips}
\item
  \hyperref[docs__pandoc__latex__src__cli_complete_curriculum.md__gb-jaws-tips]{JAWS-Specific Tips}
\item
  \hyperref[docs__pandoc__latex__src__cli_complete_curriculum.md__gb-general-terminal]{General Terminal Accessibility}
\item
  \hyperref[docs__pandoc__latex__src__cli_complete_curriculum.md__gb-long-output]{Working with Long Output}
\item
  \hyperref[docs__pandoc__latex__src__cli_complete_curriculum.md__gb-shortcuts]{Keyboard Shortcuts Reference}
\item
  \hyperref[docs__pandoc__latex__src__cli_complete_curriculum.md__gb-troubleshooting]{Troubleshooting}
\end{enumerate}

\protect\phantomsection\label{docs__pandoc__latex__src__cli_complete_curriculum.md__gb-getting-started}{}

\subsection*{Getting Started with Screen Readers}\label{docs__pandoc__latex__src__cli_complete_curriculum.md__getting-started-with-screen-readers-2}

\subsubsection*{Which Screen Reader Should I Use?}\label{docs__pandoc__latex__src__cli_complete_curriculum.md__which-screen-reader-should-i-use-2}

Both NVDA and JAWS work well with Git Bash, but they have different strengths:

{\def\LTcaptype{none} % do not increment counter
\begin{longtable}[]{@{}lll@{}}
\toprule\noalign{}
Feature & NVDA & JAWS \\
\midrule\noalign{}
\endhead
\bottomrule\noalign{}
\endlastfoot
Cost & Free & Commercial (paid) \\
Installation & Simple & Complex but thorough \\
Git Bash Support & Excellent & Excellent \\
Learning Curve & Gentle & Steeper \\
Customization & Good & Extensive \\
\end{longtable}
}

Recommendation: Start with NVDA if you\textquotesingle re new to screen readers. Both will work for this curriculum.

Additional options: Dolphin SuperNova is a commercial screen reader/ magnifier popular in some education settings; Windows Narrator is built into Windows and useful for quick access without installing additional software. Choose based on availability and personal preference.

\subsubsection*{Before You Start}\label{docs__pandoc__latex__src__cli_complete_curriculum.md__before-you-start-2}

\begin{enumerate}
\tightlist
\item
  Make sure your screen reader is running before opening Git Bash
\item
  Open Git Bash and let your screen reader read the window title and prompt
\item
  If you don\textquotesingle t hear anything, press Alt+Tab to cycle windows and find Git Bash
\item
  Use your screen reader\textquotesingle s screen review features to understand the layout
\end{enumerate}

\subsubsection*{What is Git Bash?}\label{docs__pandoc__latex__src__cli_complete_curriculum.md__what-is-git-bash}

Git Bash is a terminal application for Windows that provides a Unix-style command-line experience (Bash shell). It is installed as part of Git for Windows, which is free software available at \url{https://git-scm.com/}. When you open Git Bash, you get the same \texttt{bash}, \texttt{ls}, \texttt{grep}, \texttt{cat}, and other Unix tools used on Linux and macOS - but running on your Windows computer.

\protect\phantomsection\label{docs__pandoc__latex__src__cli_complete_curriculum.md__gb-nvda-tips}{}

\subsection*{NVDA-Specific Tips}\label{docs__pandoc__latex__src__cli_complete_curriculum.md__nvda-specific-tips-2}

NVDA is free and available from \url{https://www.nvaccess.org/}

\subsubsection*{Key Commands for Git Bash}\label{docs__pandoc__latex__src__cli_complete_curriculum.md__key-commands-for-git-bash}

{\def\LTcaptype{none} % do not increment counter
\begin{longtable}[]{@{}
  >{\raggedright\arraybackslash}p{(\linewidth - 2\tabcolsep) * \real{0.3038}}
  >{\raggedright\arraybackslash}p{(\linewidth - 2\tabcolsep) * \real{0.6962}}@{}}
\toprule\noalign{}
\begin{minipage}[b]{\linewidth}\raggedright
Command
\end{minipage} & \begin{minipage}[b]{\linewidth}\raggedright
What It Does
\end{minipage} \\
\midrule\noalign{}
\endhead
\bottomrule\noalign{}
\endlastfoot
NVDA+Home & Read the current line (your command or output) \\
NVDA+Down Arrow & Read from cursor to end of screen \\
NVDA+Up Arrow & Read from top to cursor \\
NVDA+Page Down & Read next page \\
NVDA+Page Up & Read previous page \\
NVDA+F7 & Open the Review Mode viewer (can scroll through text) \\
NVDA+Shift+Right Arrow & Read next word \\
NVDA+Shift+Down Arrow & Read entire screen \\
NVDA+End & Jump to end of line \\
NVDA+Home & Jump to start of line \\
\end{longtable}
}

\subsubsection*{Example: Reading Long Output}\label{docs__pandoc__latex__src__cli_complete_curriculum.md__example-reading-long-output-4}

Scenario: You ran \texttt{ls} and it listed 50 files.

Solution with NVDA:

\begin{enumerate}
\tightlist
\item
  After the command finishes, press NVDA+Home to read the current line
\item
  Press NVDA+Down Arrow repeatedly to read all output
\item
  Or press NVDA+F7 to open Review Mode and use arrow keys to scroll
\end{enumerate}

\subsubsection*{\texorpdfstring{Tip: Use \texttt{ls\ -1} for Screen Reader Friendly Listings}{Tip: Use ls -1 for Screen Reader Friendly Listings}}\label{docs__pandoc__latex__src__cli_complete_curriculum.md__tip-use-ls--1-for-screen-reader-friendly-listings}

By default, \texttt{ls} in Git Bash may display files in columns. Use \texttt{ls\ -1} (the number one, not the letter L) to display one file per line - much easier to follow with a screen reader.

\protect\phantomsection\label{docs__pandoc__latex__src__cli_complete_curriculum.md__gb-jaws-tips}{}

\subsection*{JAWS-Specific Tips}\label{docs__pandoc__latex__src__cli_complete_curriculum.md__jaws-specific-tips-2}

JAWS is a commercial screen reader available from \url{https://www.freedomscientific.com/}

\subsubsection*{Dolphin SuperNova}\label{docs__pandoc__latex__src__cli_complete_curriculum.md__dolphin-supernova-2}

Dolphin SuperNova (commercial): \url{https://example.com} --- provides speech and braille support and tightly integrated magnification. Refer to vendor docs for keyboard mappings; many users rely on review-mode features similar to NVDA/JAWS.

\subsubsection*{Windows Narrator}\label{docs__pandoc__latex__src__cli_complete_curriculum.md__windows-narrator-2}

Windows Narrator (built-in): \url{https://example.com} --- good for quick checks and lightweight usage; Narrator command keys differ by Windows version (use Narrator key + arrow keys to read). Add Narrator to your workflow when a full third-party screen reader isn\textquotesingle t available.

\subsubsection*{Key Commands for Git Bash}\label{docs__pandoc__latex__src__cli_complete_curriculum.md__key-commands-for-git-bash-1}

{\def\LTcaptype{none} % do not increment counter
\begin{longtable}[]{@{}ll@{}}
\toprule\noalign{}
Command & What It Does \\
\midrule\noalign{}
\endhead
\bottomrule\noalign{}
\endlastfoot
Insert+Down Arrow & Read line by line downward \\
Insert+Up Arrow & Read line by line upward \\
Insert+Page Down & Read next page of text \\
Insert+Page Up & Read previous page of text \\
Insert+End & Jump to end of text on screen \\
Insert+Home & Jump to start of text on screen \\
Insert+Ctrl+Down & Read to end of screen \\
Insert+Ctrl+Up & Read to beginning of screen \\
Insert+F3 & Open JAWS menu \\
\end{longtable}
}

\subsubsection*{Example: Reading Long Output}\label{docs__pandoc__latex__src__cli_complete_curriculum.md__example-reading-long-output-5}

Scenario: You ran \texttt{ls\ \textgreater{}\ list.txt} and saved output to a file.

Solution with JAWS:

\begin{enumerate}
\tightlist
\item
  Open the file: \texttt{notepad\ list.txt}
\item
  In Notepad, press Insert+Ctrl+Down to hear all content
\item
  Use Insert+Down Arrow to read line by line at your own pace
\end{enumerate}

\protect\phantomsection\label{docs__pandoc__latex__src__cli_complete_curriculum.md__gb-general-terminal}{}

\subsection*{General Terminal Accessibility}\label{docs__pandoc__latex__src__cli_complete_curriculum.md__general-terminal-accessibility-2}

\subsubsection*{Understanding the Git Bash Layout}\label{docs__pandoc__latex__src__cli_complete_curriculum.md__understanding-the-git-bash-layout}

The Git Bash window contains:

\begin{enumerate}
\tightlist
\item
  Title bar: Window name (e.g., "MINGW64:/c/Users/YourName")
\item
  Content area: Command history and output
\item
  Prompt: Where you type (e.g., \texttt{YourName@COMPUTER\ MINGW64\ \textasciitilde{}\$})
\end{enumerate}

Your screen reader reads from top to bottom, but focus is at the prompt (bottom).

\subsubsection*{The Git Bash Prompt}\label{docs__pandoc__latex__src__cli_complete_curriculum.md__the-git-bash-prompt}

The default Git Bash prompt looks like:

\begin{lstlisting}[style=Alabaster, language=bash]
YourName@COMPUTERNAME MINGW64 ~/Documents
$

\end{lstlisting}

\begin{itemize}
\tightlist
\item
  \texttt{YourName} = your Windows username
\item
  \texttt{COMPUTERNAME} = your computer\textquotesingle s name
\item
  \texttt{MINGW64} = the environment type (64-bit)
\item
  \texttt{\textasciitilde{}/Documents} = your current location
\item
  \texttt{\$} = ready for input
\end{itemize}

Note: Paths in Git Bash use forward slashes (\texttt{/}) and your Windows \texttt{C:\textbackslash{}Users\textbackslash{}YourName} folder appears as \texttt{/c/Users/YourName} or simply \texttt{\textasciitilde{}}.

\subsubsection*{Navigation Sequence}\label{docs__pandoc__latex__src__cli_complete_curriculum.md__navigation-sequence-2}

When you open Git Bash:

\begin{enumerate}
\tightlist
\item
  Your screen reader announces the window title
\item
  Then it announces the prompt line
\item
  Anything before the prompt is previous output
\item
  Anything after the prompt is where new output will appear
\end{enumerate}

\subsubsection*{Reading Output Effectively}\label{docs__pandoc__latex__src__cli_complete_curriculum.md__reading-output-effectively}

Strategy 1: Immediate Output (Small Amount)

\begin{itemize}
\tightlist
\item
  Run a command
\item
  Your screen reader announces output immediately
\item
  This works well for short outputs (a few lines)
\end{itemize}

Strategy 2: Large Output (Many Lines)

\begin{itemize}
\tightlist
\item
  Redirect to a file: \texttt{command\ \textgreater{}\ output.txt}
\item
  Open the file: \texttt{notepad\ output.txt}
\item
  Read in Notepad (easier for long text)
\end{itemize}

Strategy 3: Searching Output

\begin{itemize}
\item
  Use \texttt{grep} to filter:

  \begin{lstlisting}[style=Alabaster, language=bash]
  ls | grep "pattern"

  \end{lstlisting}
\item
  Only results matching "pattern" are shown
\end{itemize}

\protect\phantomsection\label{docs__pandoc__latex__src__cli_complete_curriculum.md__gb-long-output}{}

\subsection*{Working with Long Output}\label{docs__pandoc__latex__src__cli_complete_curriculum.md__working-with-long-output-2}

\subsubsection*{Solution 1: Redirect to a File}\label{docs__pandoc__latex__src__cli_complete_curriculum.md__solution-1-redirect-to-a-file-2}

\begin{lstlisting}[style=Alabaster, language=bash]
ls -1 > list.txt
notepad list.txt

\end{lstlisting}

\subsubsection*{Solution 2: Use Pagination}\label{docs__pandoc__latex__src__cli_complete_curriculum.md__solution-2-use-pagination-2}

\begin{lstlisting}[style=Alabaster, language=bash]
ls | less

\end{lstlisting}

\begin{itemize}
\tightlist
\item
  Press Space to see next page
\item
  Press Q to quit
\item
  Note: Some screen readers struggle with \texttt{less}, so Solution 1 is preferred
\end{itemize}

\subsubsection*{Solution 3: Filter Output}\label{docs__pandoc__latex__src__cli_complete_curriculum.md__solution-3-filter-output-2}

\begin{lstlisting}[style=Alabaster, language=bash]
ls | grep "\.scad"

\end{lstlisting}

\subsubsection*{Solution 4: Count Before Displaying}\label{docs__pandoc__latex__src__cli_complete_curriculum.md__solution-4-count-before-displaying-2}

\begin{lstlisting}[style=Alabaster, language=bash]
ls | wc -l

\end{lstlisting}

Tells you how many items there are. If the count is over 20, use the file method.

\protect\phantomsection\label{docs__pandoc__latex__src__cli_complete_curriculum.md__gb-shortcuts}{}

\subsection*{Keyboard Shortcuts Reference}\label{docs__pandoc__latex__src__cli_complete_curriculum.md__keyboard-shortcuts-reference-2}

\subsubsection*{All Users (Works in Git Bash Regardless of Screen Reader)}\label{docs__pandoc__latex__src__cli_complete_curriculum.md__all-users-works-in-git-bash-regardless-of-screen-reader}

{\def\LTcaptype{none} % do not increment counter
\begin{longtable}[]{@{}ll@{}}
\toprule\noalign{}
Key & Action \\
\midrule\noalign{}
\endhead
\bottomrule\noalign{}
\endlastfoot
Up Arrow & Show previous command \\
Down Arrow & Show next command \\
Tab & Auto-complete file/folder names \\
Shift+Tab & Cycle backward through completions \\
Home & Jump to start of line \\
End & Jump to end of line \\
Ctrl+A & Jump to start of line (alternate) \\
Ctrl+E & Jump to end of line (alternate) \\
Ctrl+C & Stop command \\
Ctrl+L & Clear screen \\
Enter & Run command \\
Ctrl+R & Search command history interactively \\
\end{longtable}
}

\protect\phantomsection\label{docs__pandoc__latex__src__cli_complete_curriculum.md__gb-troubleshooting}{}

\subsection*{Troubleshooting}\label{docs__pandoc__latex__src__cli_complete_curriculum.md__troubleshooting-4}

\subsubsection*{Problem 1: "I Can\textquotesingle t Hear the Output After Running a Command"}\label{docs__pandoc__latex__src__cli_complete_curriculum.md__problem-1-i-cant-hear-the-output-after-running-a-command}

\begin{enumerate}
\tightlist
\item
  Redirect to file: \texttt{command\ \textgreater{}\ output.txt} then \texttt{notepad\ output.txt}
\item
  Press End or Ctrl+End to go to end of text
\item
  Use Up Arrow to review previous command
\end{enumerate}

\subsubsection*{Problem 2: "Tab Completion Isn\textquotesingle t Working"}\label{docs__pandoc__latex__src__cli_complete_curriculum.md__problem-2-tab-completion-isnt-working}

\begin{enumerate}
\tightlist
\item
  Need at least one character - type \texttt{cd\ D} then Tab (not just \texttt{cd} then Tab)
\item
  Check if item exists - use \texttt{ls} first to see available items
\item
  Multiple matches - press Tab twice to list all options
\end{enumerate}

\subsubsection*{Problem 3: "\textquotesingle Command Not Found\textquotesingle{} Error"}\label{docs__pandoc__latex__src__cli_complete_curriculum.md__problem-3-command-not-found-error}

\begin{enumerate}
\tightlist
\item
  Check spelling carefully
\item
  Git Bash uses Unix commands: use \texttt{ls} not \texttt{dir}, use \texttt{cat} not \texttt{type}
\item
  Verify the program is installed and in your PATH: \texttt{which\ openscad}
\end{enumerate}

\subsubsection*{Problem 4: "Paths Look Weird"}\label{docs__pandoc__latex__src__cli_complete_curriculum.md__problem-4-paths-look-weird}

Git Bash converts Windows paths to Unix style:

\begin{itemize}
\tightlist
\item
  Windows: \texttt{C:\textbackslash{}Users\textbackslash{}YourName\textbackslash{}Documents}
\item
  Git Bash: \texttt{/c/Users/YourName/Documents} or \texttt{\textasciitilde{}/Documents}
\end{itemize}

Both refer to the same location. When opening files in Windows apps (like Notepad), you may need to use the Windows path format.

\subsubsection*{Problem 5: "Command Runs Forever and Won\textquotesingle t Stop"}\label{docs__pandoc__latex__src__cli_complete_curriculum.md__problem-5-command-runs-forever-and-wont-stop}

Press Ctrl+C

\subsubsection*{Problem 6: "I Need to Edit My Last Command"}\label{docs__pandoc__latex__src__cli_complete_curriculum.md__problem-6-i-need-to-edit-my-last-command}

\begin{enumerate}
\tightlist
\item
  Press Up Arrow to show previous command
\item
  Use arrow keys to move through it
\item
  Edit the command
\item
  Press Enter to run the modified version
\end{enumerate}

\subsection*{Pro Tips for Efficiency}\label{docs__pandoc__latex__src__cli_complete_curriculum.md__pro-tips-for-efficiency}

\subsubsection*{\texorpdfstring{1. Use \texttt{ls\ -1} for Screen Reader Friendly Listings}{1. Use ls -1 for Screen Reader Friendly Listings}}\label{docs__pandoc__latex__src__cli_complete_curriculum.md__1-use-ls--1-for-screen-reader-friendly-listings}

\begin{lstlisting}[style=Alabaster, language=bash]
ls -1

\end{lstlisting}

One file per line - much easier to follow with a screen reader.

\subsubsection*{2. Create Aliases for Frequently Used Commands}\label{docs__pandoc__latex__src__cli_complete_curriculum.md__2-create-aliases-for-frequently-used-commands}

\begin{lstlisting}[style=Alabaster, language=bash]
alias la='ls -1a'

\end{lstlisting}

Add this to your \texttt{\textasciitilde{}/.bashrc} file to make it permanent.

\subsubsection*{3. Use Command History Effectively}\label{docs__pandoc__latex__src__cli_complete_curriculum.md__3-use-command-history-effectively}

\begin{lstlisting}[style=Alabaster, language=bash]
history

\end{lstlisting}

Run a previous command by number:

\begin{lstlisting}[style=Alabaster, language=bash]
!5

\end{lstlisting}

(Runs the 5th command in history)

Or search history interactively:

\begin{lstlisting}[style=Alabaster, language=bash]
# Press Ctrl+R then type part of a previous command

\end{lstlisting}

\subsubsection*{4. Redirect Everything to Files for Accessibility}\label{docs__pandoc__latex__src__cli_complete_curriculum.md__4-redirect-everything-to-files-for-accessibility}

\begin{lstlisting}[style=Alabaster, language=bash]
command-name > results.txt
notepad results.txt

\end{lstlisting}

\subsubsection*{5. Create a README for Yourself}\label{docs__pandoc__latex__src__cli_complete_curriculum.md__5-create-a-readme-for-yourself}

\begin{lstlisting}[style=Alabaster, language=bash]
echo "ls -1 means list files one per line (screen reader friendly)" > my-notes.txt
echo "cd means change directory" >> my-notes.txt
notepad my-notes.txt

\end{lstlisting}

\subsection*{Recommended Workflow}\label{docs__pandoc__latex__src__cli_complete_curriculum.md__recommended-workflow-2}

For every new task:

\begin{enumerate}
\tightlist
\item
  Know where you are: \texttt{pwd}
\item
  See what\textquotesingle s around: \texttt{ls\ -1}
\item
  Plan your next step: Think before typing
\item
  Run the command: Type and press Enter
\item
  Check the output: Use screen reader or redirect to file
\item
  Move forward: Next command or \texttt{cd} to next folder
\end{enumerate}

\subsection*{Quick Reference Card}\label{docs__pandoc__latex__src__cli_complete_curriculum.md__quick-reference-card}

\begin{lstlisting}[style=Alabaster]
EVERY COMMAND STARTS WITH:
1. pwd (where am I?)
2. ls -1 (what's here?)
3. cd path (go there)

LONG OUTPUT?
-> command > file.txt
-> notepad file.txt

STUCK?
-> Ctrl+C

WANT TO REPEAT?
-> Up Arrow
-> history

NEED HELP?
-> man command-name  (or: command --help)

\end{lstlisting}

\subsection*{Additional Resources}\label{docs__pandoc__latex__src__cli_complete_curriculum.md__additional-resources-2}

\begin{itemize}
\tightlist
\item
  NVDA Documentation: \url{https://www.nvaccess.org/documentation/}
\item
  JAWS Documentation: \url{https://www.freedomscientific.com/support/}
\item
  Git for Windows (includes Git Bash): \url{https://git-scm.com/}
\item
  Bash Manual: \url{https://www.gnu.org/software/bash/manual/}
\end{itemize}

\begin{center}\rule{0.5\linewidth}{0.5pt}\end{center}

\protect\phantomsection\label{docs__pandoc__latex__src__cli_complete_curriculum.md__gb-tutorial}{}

\section*{Tutorial: Git Bash Quick Start}\label{docs__pandoc__latex__src__cli_complete_curriculum.md__tutorial-git-bash-quick-start}

\begin{center}\rule{0.5\linewidth}{0.5pt}\end{center}

Estimated time: 30-45 minutes

Learning Objectives

\begin{itemize}
\tightlist
\item
  Launch Git Bash and identify the prompt
\item
  Understand and use basic path notation (\texttt{\textasciitilde{}}, \texttt{./}, \texttt{../})
\item
  Use \texttt{pwd}, \texttt{ls\ -1}, and \texttt{cd} to navigate the filesystem
\item
  Open files in an external editor and run simple commands
\end{itemize}

Materials

\begin{itemize}
\tightlist
\item
  A computer with Git Bash installed (install via \url{https://git-scm.com/})
\item
  Access to a text editor (Notepad, VS Code)
\end{itemize}

\subsection*{Pre-Requisite Knowledge}\label{docs__pandoc__latex__src__cli_complete_curriculum.md__pre-requisite-knowledge}

\begin{itemize}
\tightlist
\item
  Typing and basic text-editing skills
\item
  Familiarity with file/folder concepts and basic OS navigation
\item
  Basic screen-reading familiarity (if applicable)
\end{itemize}

\subsection*{What is Git Bash?}\label{docs__pandoc__latex__src__cli_complete_curriculum.md__what-is-git-bash-1}

Git Bash is a free terminal application for Windows that provides a Bash (Unix shell) environment. It is installed as part of Git for Windows. When you open Git Bash, you can use Unix commands like \texttt{ls}, \texttt{cat}, \texttt{grep}, and \texttt{echo} - the same commands used on Linux and macOS - right on your Windows computer.

Git Bash is useful for:

\begin{itemize}
\tightlist
\item
  Running CLI programs (like OpenSCAD, slicers, or 3DMake)
\item
  Navigating the filesystem with keyboard-only commands
\item
  Automating repetitive tasks with scripts
\item
  Accessibility: works cleanly with NVDA, JAWS, and other screen readers
\end{itemize}

\subsection*{What We\textquotesingle ll Do and Why}\label{docs__pandoc__latex__src__cli_complete_curriculum.md__what-well-do-and-why}

You\textquotesingle ll use Git Bash to run CLI programs, navigate the filesystem, and manipulate files - tasks that are especially efficient when using a keyboard or a screen reader.

\subsection*{Quick Tutorial \& Core Concepts}\label{docs__pandoc__latex__src__cli_complete_curriculum.md__quick-tutorial--core-concepts}

\subsubsection*{Paths and Navigation}\label{docs__pandoc__latex__src__cli_complete_curriculum.md__paths-and-navigation}

\begin{itemize}
\tightlist
\item
  \texttt{\textasciitilde{}} - home directory (e.g., \texttt{/c/Users/YourName})
\item
  \texttt{.} - current directory
\item
  \texttt{..} - parent directory
\item
  \texttt{./} - current directory shortcut (used to run scripts: \texttt{./script.sh})
\item
  \texttt{../} - parent directory shortcut
\item
  Use Tab to autocomplete files and folders
\end{itemize}

Git Bash uses forward slashes (\texttt{/}) for paths, just like Linux. Your Windows \texttt{C:\textbackslash{}Users\textbackslash{}YourName} folder is accessible as \texttt{/c/Users/YourName} or simply \texttt{\textasciitilde{}}.

\subsubsection*{Useful Commands (Examples)}\label{docs__pandoc__latex__src__cli_complete_curriculum.md__useful-commands-examples}

\begin{lstlisting}[style=Alabaster, language=bash]
pwd                # show current directory
ls -1              # list files, one per line (screen-reader friendly)
cd path/to/dir     # change directory
whoami             # current user

\end{lstlisting}

\subsubsection*{Wildcards}\label{docs__pandoc__latex__src__cli_complete_curriculum.md__wildcards}

\begin{itemize}
\tightlist
\item
  \texttt{*} matches zero or more characters
\item
  \texttt{?} matches a single character
\end{itemize}

Use \texttt{ls\ *.scad} to filter by extension, for example.

\subsection*{Common Operations}\label{docs__pandoc__latex__src__cli_complete_curriculum.md__common-operations}

\subsubsection*{File and Folder Manipulation}\label{docs__pandoc__latex__src__cli_complete_curriculum.md__file-and-folder-manipulation}

\begin{lstlisting}[style=Alabaster, language=bash]
mkdir my-folder          # create folder
mkdir -p a/b/c           # create nested folders at once
cp -r src dest           # copy (use -r for directories)
mv oldname newname       # rename or move
rm file                  # remove file
rm -r folder             # remove folder and contents
touch filename.txt       # create new empty file

\end{lstlisting}

\subsubsection*{Input, Output, and Piping}\label{docs__pandoc__latex__src__cli_complete_curriculum.md__input-output-and-piping}

\begin{lstlisting}[style=Alabaster, language=bash]
echo 'hello' | clip      # copy to clipboard
command > output.txt     # redirect output to file
command >> output.txt    # append output to file
command > /dev/null      # discard/suppress output

\end{lstlisting}

Use \texttt{\textbar{}} to pipe output and \texttt{\textgreater{}} to redirect into files.

\subsubsection*{Editing and Running Programs}\label{docs__pandoc__latex__src__cli_complete_curriculum.md__editing-and-running-programs}

\begin{lstlisting}[style=Alabaster, language=bash]
notepad file.txt         # open in Notepad (Windows)
code file.txt            # open in VS Code (if in PATH)
./script.sh              # run a script in the current directory
chmod +x script.sh       # make a script executable

\end{lstlisting}

\subsection*{Screen-Reader Friendly Tips}\label{docs__pandoc__latex__src__cli_complete_curriculum.md__screen-reader-friendly-tips}

\begin{itemize}
\tightlist
\item
  Prefer \texttt{ls\ -1} for name-only, one-per-line output.
\item
  Filter lists with \texttt{grep}: \texttt{ls\ -1\ \textasciitilde{}/Documents\ \textbar{}\ grep\ "\textbackslash{}.scad\$"} (files only ending in .scad).
\item
  Redirect very long outputs to a file and open it in an editor.
\end{itemize}

\subsection*{Error Handling and Control}\label{docs__pandoc__latex__src__cli_complete_curriculum.md__error-handling-and-control}

\begin{itemize}
\tightlist
\item
  Abort a running command: \texttt{Ctrl+C}
\item
  View history: Up/Down arrows or \texttt{history}
\item
  Clear screen: \texttt{clear} (or \texttt{Ctrl+L})
\item
  If an error is long, read the first few lines for the gist and copy short snippets into a file to examine.
\end{itemize}

\subsection*{Environment Variables \& PATH}\label{docs__pandoc__latex__src__cli_complete_curriculum.md__environment-variables--path}

Environment variables configure your session. \texttt{PATH} tells the shell where to find executables.

\begin{lstlisting}[style=Alabaster, language=bash]
echo $PATH
echo $HOME
echo $USER

\end{lstlisting}

To add a directory to PATH for the current session:

\begin{lstlisting}[style=Alabaster, language=bash]
export PATH="$PATH:/path/to/my/tools"

\end{lstlisting}

To make it permanent, add that line to \texttt{\textasciitilde{}/.bashrc}.

\subsection*{Running CLI Applications and Archives}\label{docs__pandoc__latex__src__cli_complete_curriculum.md__running-cli-applications-and-archives}

To extract ZIP archives in Git Bash:

\begin{lstlisting}[style=Alabaster, language=bash]
# Unzip a file
unzip file.zip -d destinationfolder

# Or use Windows' built-in tool (from Git Bash):
powershell.exe Expand-Archive -Path file.zip -DestinationPath folder

\end{lstlisting}

\subsection*{Aliases and Cross-Platform Notes}\label{docs__pandoc__latex__src__cli_complete_curriculum.md__aliases-and-cross-platform-notes}

Git Bash provides Unix commands that work identically on Linux and macOS. This means your Git Bash skills transfer directly if you ever work on a Mac or Linux system.

Useful aliases to add to \texttt{\textasciitilde{}/.bashrc}:

\begin{lstlisting}[style=Alabaster, language=bash]
alias ll='ls -la'
alias la='ls -1a'
alias ..='cd ..'
alias ...='cd ../..'

\end{lstlisting}

\subsection*{Step-by-step Tasks (Hands-On)}\label{docs__pandoc__latex__src__cli_complete_curriculum.md__step-by-step-tasks-hands-on}

\begin{enumerate}
\tightlist
\item
  Open Git Bash; listen for the prompt and current path.
\item
  Run \texttt{pwd} to confirm your location.
\item
  Run \texttt{ls\ -1} in your home directory and note the output.
\item
  Practice \texttt{cd\ Documents}, \texttt{cd\ ../}, and \texttt{cd\ \textasciitilde{}} to move between folders.
\item
  Create and open a file: \texttt{touch\ example.txt\ \&\&\ notepad\ example.txt} (or \texttt{code\ example.txt}).
\end{enumerate}

\subsection*{Checkpoints}\label{docs__pandoc__latex__src__cli_complete_curriculum.md__checkpoints-9}

\begin{itemize}
\tightlist
\item
  After step 3 you should be able to state your current directory.
\item
  After step 5 you should be able to create and open a text file from Git Bash.
\end{itemize}

\subsection*{Quick Quiz (10 questions)}\label{docs__pandoc__latex__src__cli_complete_curriculum.md__quick-quiz-10-questions}

\begin{enumerate}
\tightlist
\item
  What command prints your current directory?
\item
  What does \texttt{\textasciitilde{}} represent?
\item
  How do you list only names, one per line?
\item
  How do you go up one directory level?
\item
  How would you open \texttt{notes.txt} in Notepad from Git Bash?
\item
  True or False: The pipe operator \texttt{\textbar{}} sends output to a file.
\item
  Explain why running a script requires \texttt{chmod\ +x} first.
\item
  What is the difference between running a script with \texttt{./script.sh} versus \texttt{bash\ script.sh}?
\item
  Describe how you would handle a very long command output when using a screen reader.
\item
  What does the \texttt{PATH} environment variable do, and why is it important when running programs like OpenSCAD?
\end{enumerate}

\subsection*{Extension Problems}\label{docs__pandoc__latex__src__cli_complete_curriculum.md__extension-problems-14}

\begin{enumerate}
\tightlist
\item
  Create a folder \texttt{OpenSCADProjects} in Documents and verify its contents.
\item
  Create three files named \texttt{a.scad}, \texttt{b.scad}, \texttt{c.scad} and list them with a wildcard.
\item
  Save \texttt{ls\ -1\ \textasciitilde{}/Documents} output to \texttt{doclist.txt} and open it.
\item
  Try tab-completion in a deeply nested folder and note behavior.
\item
  Capture \texttt{pwd} output into a file and open it: \texttt{pwd\ \textgreater{}\ cwd.txt\ \&\&\ notepad\ cwd.txt}.
\item
  Build an automated setup script that creates a complete project directory structure, initializes placeholder files, and generates a README.
\item
  Create a Git Bash cheat sheet for your most-used commands; organize by category (navigation, files, scripting, troubleshooting).
\item
  Write a non-visual tutorial for Git Bash basics; use audio descriptions and keyboard-only navigation as the primary learning method.
\item
  Develop a workflow automation script: combines multiple Git Bash concepts (folders, aliases, piping) to solve a real 3D printing task.
\item
  Create a Git Bash proficiency self-assessment: list all concepts covered, provide test commands for each, and reflect on what you learned.
\end{enumerate}

\subsection*{References}\label{docs__pandoc__latex__src__cli_complete_curriculum.md__references-12}

\begin{itemize}
\tightlist
\item
  Git for Windows. (2024). \emph{Git Bash}. \url{https://gitforwindows.org/}
\item
  GNU. (2024). \emph{Bash reference manual}. \url{https://www.gnu.org/software/bash/manual/bash.html}
\item
  The Linux Documentation Project. (2024). \emph{Bash Beginners Guide}. \url{https://tldp.org/LDP/Bash-Beginners-Guide/html/}
\end{itemize}

\subsection*{Helpful Resources}\label{docs__pandoc__latex__src__cli_complete_curriculum.md__helpful-resources-10}

\begin{itemize}
\tightlist
\item
  \href{https://gitforwindows.org/}{Git for Windows}
\item
  \href{https://www.gnu.org/software/bash/manual/bash.html}{Bash Reference Manual (GNU)}
\item
  \href{https://tldp.org/LDP/Bash-Beginners-Guide/html/}{Bash Beginners Guide}
\item
  \href{https://www.gnu.org/software/coreutils/manual/}{GNU Coreutils}
\end{itemize}

Other Screen Readers

Dolphin SuperNova (commercial) and Windows Narrator (built-in) are also supported; the workflows and recommendations in this document apply to them. See \url{https://example.com} and \url{https://example.com} for vendor documentation.

\begin{center}\rule{0.5\linewidth}{0.5pt}\end{center}

\protect\phantomsection\label{docs__pandoc__latex__src__cli_complete_curriculum.md__gb-lesson-pre}{}

\section*{Lesson Pre: Your First Terminal --- Screen Reader Navigation Fundamentals}\label{docs__pandoc__latex__src__cli_complete_curriculum.md__lesson-pre-your-first-terminal--screen-reader-navigation-fundamentals-1}

\begin{center}\rule{0.5\linewidth}{0.5pt}\end{center}

Duration: 1.5-2 hours (for screen reader users)\\
Prerequisites: None - this is the starting point\\
Accessibility Note: This lesson is designed specifically for screen reader users (NVDA, JAWS)

Important

Git Bash uses Unix/bash commands, different from Windows CMD/PowerShell

\subsection*{What is Git Bash?}\label{docs__pandoc__latex__src__cli_complete_curriculum.md__what-is-git-bash-2}

Git Bash is a Unix/Linux shell running on Windows. It gives you bash commands (same as macOS and Linux) while staying on your Windows computer.

Why learn this?

\begin{itemize}
\tightlist
\item
  Faster and more precise work (especially for 3D printing scripts and automation)
\item
  Essential for programming and using tools like OpenSCAD
\item
  Accessibility: Git Bash works perfectly with screen readers
\item
  Transferable: Same commands work on macOS and Linux
\item
  Industry-standard: Professional developers use this everywhere
\end{itemize}

\subsection*{Installing and Opening Git Bash}\label{docs__pandoc__latex__src__cli_complete_curriculum.md__installing-and-opening-git-bash}

\subsubsection*{3DMake: Your First Real-World Goal}\label{docs__pandoc__latex__src__cli_complete_curriculum.md__3dmake-your-first-real-world-goal-2}

\textbf{3DMake} (\href{https://github.com/tdeck/3dmake}{github.com/tdeck/3dmake}) is a free, accessibility-first command-line tool for designing and 3D-printing models entirely from the terminal --- no mouse required. Every Git Bash skill in this curriculum directly supports a real 3DMake workflow.

Here is the end goal --- one command that builds your OpenSCAD model, slices it, and sends it to the printer:

\begin{lstlisting}[style=Alabaster, language=bash]
# Navigate to your project folder
cd ~/Documents/3d_projects/my_cube

# Build, slice, and print in a single command
3dm build slice print

\end{lstlisting}

Keep this goal in mind as you work through each lesson.

\subsection*{Installing and Opening Git Bash}\label{docs__pandoc__latex__src__cli_complete_curriculum.md__installing-and-opening-git-bash-1}

\subsubsection*{Installation (If Not Already Installed)}\label{docs__pandoc__latex__src__cli_complete_curriculum.md__installation-if-not-already-installed}

\begin{enumerate}
\tightlist
\item
  Download from: \url{https://git-scm.com/download/win}
\item
  Run the installer
\item
  Accept defaults (all options work fine)
\item
  Installation takes \textasciitilde{}5 minutes
\end{enumerate}

\subsubsection*{Opening Git Bash}\label{docs__pandoc__latex__src__cli_complete_curriculum.md__opening-git-bash}

Method 1: Search (Easiest)

\begin{enumerate}
\tightlist
\item
  Press the Windows key alone
\item
  You should hear "Search"
\item
  Type: \texttt{Git\ Bash}
\item
  You\textquotesingle ll hear search results appear
\item
  Press Enter to open
\end{enumerate}

Method 2: From File Explorer

\begin{enumerate}
\tightlist
\item
  Open File Explorer
\item
  Navigate to any folder
\item
  Right-click on the folder
\item
  Look for "Git Bash Here"
\item
  Press Enter
\end{enumerate}

\subsubsection*{First Connection: Understanding the Prompt}\label{docs__pandoc__latex__src__cli_complete_curriculum.md__first-connection-understanding-the-prompt-2}

When Git Bash opens, your screen reader will announce the window title and then the prompt.

What you\textquotesingle ll hear (typical example):

\begin{lstlisting}[style=Alabaster]
user@computer MINGW64 ~
$

\end{lstlisting}

What this means:

\begin{itemize}
\tightlist
\item
  \texttt{user} = Your username
\item
  \texttt{@computer} = Your computer name
\item
  \texttt{MINGW64} = Version of Git Bash
\item
  \texttt{\textasciitilde{}} = Your current location (home directory)
\item
  \texttt{\$} = The prompt is ready for your input
\end{itemize}

Important

Your cursor is blinking right after the \texttt{\$}. This is where you type.

\subsection*{Your First Commands (Screen Reader Edition)}\label{docs__pandoc__latex__src__cli_complete_curriculum.md__your-first-commands-screen-reader-edition-2}

\subsubsection*{\texorpdfstring{Command 1: "Where Am I?" - \texttt{pwd}}{Command 1: "Where Am I?" - pwd}}\label{docs__pandoc__latex__src__cli_complete_curriculum.md__command-1-where-am-i---pwd-1}

What it does: Tells you your current location

Type this:

\begin{lstlisting}[style=Alabaster, language=bash]
pwd

\end{lstlisting}

Press Enter

What you\textquotesingle ll hear:
Your screen reader will announce the current path, something like:

\begin{lstlisting}[style=Alabaster]
/c/Users/YourName

\end{lstlisting}

Understanding paths in Git Bash:

\begin{itemize}
\tightlist
\item
  Paths use forward slashes: \texttt{/c/Users/YourName/Documents}
\item
  \texttt{/c/} means the C: drive (Windows drive)
\item
  This is Unix-style, same as macOS and Linux
\item
  Think of it like folders inside folders: \texttt{/c/} (main drive) -\textgreater{} \texttt{Users} -\textgreater{} \texttt{YourName} -\textgreater{} \texttt{Documents}
\end{itemize}

\subsubsection*{\texorpdfstring{Command 2: "What\textquotesingle s Here?" - \texttt{ls}}{Command 2: "What\textquotesingle s Here?" - ls}}\label{docs__pandoc__latex__src__cli_complete_curriculum.md__command-2-whats-here---ls}

What it does: Lists all files and folders in your current location

Type this:

\begin{lstlisting}[style=Alabaster, language=bash]
ls

\end{lstlisting}

Press Enter

What you\textquotesingle ll hear:
Your screen reader will announce a list of file and folder names, one per line. Perfect for screen readers!

\subsubsection*{\texorpdfstring{Command 3: "Go There" - \texttt{cd\ Documents}}{Command 3: "Go There" - cd Documents}}\label{docs__pandoc__latex__src__cli_complete_curriculum.md__command-3-go-there---cd-documents-1}

What it does: Changes your location (navigates to a folder)

Type this:

\begin{lstlisting}[style=Alabaster, language=bash]
cd Documents

\end{lstlisting}

Press Enter

What you\textquotesingle ll hear:
The prompt changes to show your new location. You might hear something like:

\begin{lstlisting}[style=Alabaster]
user@computer MINGW64 ~/Documents
$

\end{lstlisting}

Practice navigation:

\begin{enumerate}
\tightlist
\item
  Run \texttt{pwd} to confirm you\textquotesingle re in Documents
\item
  Run \texttt{ls} to see what files are in Documents
\item
  Try going back: \texttt{cd\ ..} (the \texttt{..} means "go up one level")
\item
  Run \texttt{pwd} again to confirm
\item
  Go back to Documents: \texttt{cd\ Documents}
\end{enumerate}

\subsection*{Creating and Viewing Files}\label{docs__pandoc__latex__src__cli_complete_curriculum.md__creating-and-viewing-files-1}

\subsubsection*{Create a Simple File}\label{docs__pandoc__latex__src__cli_complete_curriculum.md__create-a-simple-file-2}

Type this:

\begin{lstlisting}[style=Alabaster, language=bash]
echo "Hello, Git Bash!" > hello.txt

\end{lstlisting}

What this does:

\begin{itemize}
\tightlist
\item
  \texttt{echo} sends text to the screen (or file)
\item
  \texttt{"Hello,\ Git\ Bash!"} is the text
\item
  \texttt{\textgreater{}} redirects it to a file called \texttt{hello.txt}
\end{itemize}

\subsubsection*{Read the File Back}\label{docs__pandoc__latex__src__cli_complete_curriculum.md__read-the-file-back-2}

Type this:

\begin{lstlisting}[style=Alabaster, language=bash]
cat hello.txt

\end{lstlisting}

What you\textquotesingle ll hear: Your screen reader announces:

\begin{lstlisting}[style=Alabaster]
Hello, Git Bash!

\end{lstlisting}

\subsubsection*{Open and Edit the File}\label{docs__pandoc__latex__src__cli_complete_curriculum.md__open-and-edit-the-file-2}

Type this:

\begin{lstlisting}[style=Alabaster, language=bash]
notepad.exe hello.txt

\end{lstlisting}

This opens the file in Notepad where you can edit it with your screen reader.

\subsection*{Essential Keyboard Shortcuts}\label{docs__pandoc__latex__src__cli_complete_curriculum.md__essential-keyboard-shortcuts-2}

These work in Git Bash and are crucial for screen reader users:

{\def\LTcaptype{none} % do not increment counter
\begin{longtable}[]{@{}
  >{\raggedright\arraybackslash}p{(\linewidth - 2\tabcolsep) * \real{0.2152}}
  >{\raggedright\arraybackslash}p{(\linewidth - 2\tabcolsep) * \real{0.7848}}@{}}
\toprule\noalign{}
\begin{minipage}[b]{\linewidth}\raggedright
Key Combination
\end{minipage} & \begin{minipage}[b]{\linewidth}\raggedright
What It Does
\end{minipage} \\
\midrule\noalign{}
\endhead
\bottomrule\noalign{}
\endlastfoot
Up Arrow &
Shows your previous command (press again to go further back) \\
Down Arrow & Shows your next command (if you went back) \\
Tab & Auto-completes folder/file names \\
Ctrl+C & Stops a running command \\
Ctrl+L & Clears the screen \\
Ctrl+A & Go to beginning of line \\
Ctrl+E & Go to end of line \\
Enter & Runs the command \\
\end{longtable}
}

Screen reader tip: These all work perfectly with your screen reader. Try them!

\subsection*{Screen Reader-Specific Tips}\label{docs__pandoc__latex__src__cli_complete_curriculum.md__screen-reader-specific-tips-2}

\subsubsection*{NVDA Users}\label{docs__pandoc__latex__src__cli_complete_curriculum.md__nvda-users-4}

\begin{enumerate}
\item
  Reading Command Output:

  \begin{itemize}
  \tightlist
  \item
    Use NVDA+Home to read the current line
  \item
    Use NVDA+Down Arrow to read to the end of the screen
  \item
    Use NVDA+Page Down to read the next page
  \end{itemize}
\item
  Reviewing Text:

  \begin{itemize}
  \tightlist
  \item
    Use NVDA+Shift+Page Up to review text above
  \end{itemize}
\end{enumerate}

\subsubsection*{JAWS Users}\label{docs__pandoc__latex__src__cli_complete_curriculum.md__jaws-users-4}

\begin{enumerate}
\item
  Reading Output:

  \begin{itemize}
  \tightlist
  \item
    Use Insert+Down Arrow to read line-by-line
  \item
    Use Insert+Page Down to read by page
  \item
    Use Insert+End to jump to the end of text
  \end{itemize}
\item
  Reading All Text:

  \begin{itemize}
  \tightlist
  \item
    Use Insert+Down Arrow repeatedly
  \item
    Or use Insert+Ctrl+Down to read to the end
  \end{itemize}
\end{enumerate}

\subsubsection*{Common Issue: "I Can\textquotesingle t Hear the Output"}\label{docs__pandoc__latex__src__cli_complete_curriculum.md__common-issue-i-cant-hear-the-output-2}

Problem: You run a command but don\textquotesingle t hear the output

Solutions:

\begin{enumerate}
\tightlist
\item
  Make sure your cursor is at the prompt (try pressing End or Ctrl+E)
\item
  Use Up Arrow to go back to your previous command and review it
\item
  Try redirecting to a file: \texttt{command\ \textgreater{}\ output.txt} then open the file
\item
  In NVDA: Try pressing NVDA+F7 to open the Review Mode viewer
\end{enumerate}

\subsection*{Practice Exercises}\label{docs__pandoc__latex__src__cli_complete_curriculum.md__practice-exercises-6}

Complete these in order. Take your time with each one:

\subsubsection*{Exercise 1: Basic Navigation}\label{docs__pandoc__latex__src__cli_complete_curriculum.md__exercise-1-basic-navigation-2}

\begin{enumerate}
\tightlist
\item
  Open Git Bash
\item
  Run \texttt{pwd} and note your location
\item
  Run \texttt{ls} and listen to what\textquotesingle s there
\item
  Try \texttt{cd\ Documents} or another folder
\item
  Run \texttt{pwd} to confirm your new location
\item
  Run \texttt{ls} in this new location
\end{enumerate}

Goal: You should be comfortable knowing where you are and what\textquotesingle s around you

\subsubsection*{Exercise 2: Using Tab Completion}\label{docs__pandoc__latex__src__cli_complete_curriculum.md__exercise-2-using-tab-completion-3}

\begin{enumerate}
\tightlist
\item
  In your home directory, type \texttt{cd\ D} (just the letter D)
\item
  Press Tab
\item
  Git Bash should auto-complete to a folder starting with D
\item
  Repeat with other folder names
\item
  Try typing a longer name: \texttt{cd\ Down} and Tab
\end{enumerate}

Goal: Tab completion should feel natural

\subsubsection*{Exercise 3: Creating and Viewing Files}\label{docs__pandoc__latex__src__cli_complete_curriculum.md__exercise-3-creating-and-viewing-files-2}

\begin{enumerate}
\tightlist
\item
  Create a file: \texttt{echo\ "Test\ content"\ \textgreater{}\ test.txt}
\item
  View it: \texttt{cat\ test.txt}
\item
  Create another: \texttt{echo\ "Line\ 2"\ \textgreater{}\ another.txt}
\item
  List both: \texttt{ls\ *.txt}
\end{enumerate}

Goal: You understand create, view, and list operations

\subsubsection*{Exercise 4: Going Up Levels}\label{docs__pandoc__latex__src__cli_complete_curriculum.md__exercise-4-going-up-levels-2}

\begin{enumerate}
\tightlist
\item
  Navigate into several folders: \texttt{cd\ Documents}, then \texttt{cd\ folder1}, etc.
\item
  From deep inside, use \texttt{cd\ ..} multiple times to go back up
\item
  After each \texttt{cd\ ..}, run \texttt{pwd} to confirm your location
\end{enumerate}

Goal: You understand relative navigation with \texttt{..}

\subsubsection*{Exercise 5: Redirecting Output}\label{docs__pandoc__latex__src__cli_complete_curriculum.md__exercise-5-redirecting-output-2}

\begin{enumerate}
\tightlist
\item
  Create a list: \texttt{ls\ \textgreater{}\ directorylist.txt}
\item
  Open it: \texttt{notepad.exe\ directorylist.txt}
\item
  Read it with your screen reader
\item
  Close Notepad
\item
  Verify the file exists: \texttt{ls\ \textbar{}\ grep\ directory}
\end{enumerate}

Goal: You can save long outputs to files for easier reading

\subsection*{Checkpoint Questions}\label{docs__pandoc__latex__src__cli_complete_curriculum.md__checkpoint-questions-4}

After completing this lesson, you should be able to answer:

\begin{enumerate}
\tightlist
\item
  What does \texttt{pwd} do?
\item
  What does \texttt{ls} do?
\item
  Why do we use \texttt{ls} for listings?
\item
  What path are you in right now?
\item
  How do you navigate to a new folder?
\item
  How do you go up one level?
\item
  What\textquotesingle s the Tab key for?
\item
  What does \texttt{echo\ "text"\ \textgreater{}\ file.txt} do?
\item
  How do you read a file back?
\item
  How do you stop a command that\textquotesingle s running?
\end{enumerate}

You should be able to answer all 10 with confidence before moving to Bash Lesson 0.

\subsection*{Common Questions}\label{docs__pandoc__latex__src__cli_complete_curriculum.md__common-questions-2}

Q: Is Git Bash the same as Command Prompt or PowerShell?
A: No. Git Bash uses Unix/bash commands. CMD and PowerShell use Windows commands. This curriculum teaches bash (Unix) style.

Q: Why is my screen reader not reading the output?
A: This is common. Use \texttt{command\ \textgreater{}\ file.txt} to save output to a file, then open it with Notepad for reliable reading.

Q: What if I type something wrong?
A: Just press Enter and you\textquotesingle ll see an error message. Type the correct command on the next line. No harm done!

Q: How do I get help with a command?
A: Type \texttt{man\ command-name} (we\textquotesingle ll cover this in Bash Lesson 0)

Q: Why do paths look different in Git Bash?
A: Git Bash uses Unix-style paths with forward slashes, not Windows backslashes. Don\textquotesingle t worry - you\textquotesingle ll get used to it quickly.

\subsection*{Path Comparison: Windows vs Git Bash}\label{docs__pandoc__latex__src__cli_complete_curriculum.md__path-comparison-windows-vs-git-bash}

{\def\LTcaptype{none} % do not increment counter
\begin{longtable}[]{@{}
  >{\raggedright\arraybackslash}p{(\linewidth - 4\tabcolsep) * \real{0.1389}}
  >{\raggedright\arraybackslash}p{(\linewidth - 4\tabcolsep) * \real{0.3750}}
  >{\raggedright\arraybackslash}p{(\linewidth - 4\tabcolsep) * \real{0.4861}}@{}}
\toprule\noalign{}
\begin{minipage}[b]{\linewidth}\raggedright
Style
\end{minipage} & \begin{minipage}[b]{\linewidth}\raggedright
Example
\end{minipage} & \begin{minipage}[b]{\linewidth}\raggedright
Where Used
\end{minipage} \\
\midrule\noalign{}
\endhead
\bottomrule\noalign{}
\endlastfoot
Windows &
\texttt{C:\textbackslash{}Users\textbackslash{}Name\textbackslash{}Documents}
& CMD, PowerShell, Windows Explorer \\
Git Bash & \texttt{/c/Users/Name/Documents} &
Git Bash, macOS Terminal, Linux \\
\end{longtable}
}

In Git Bash:

\begin{itemize}
\tightlist
\item
  \texttt{C:\textbackslash{}} becomes \texttt{/c/}
\item
  Backslashes \texttt{\textbackslash{}} become forward slashes \texttt{/}
\item
  Otherwise identical
\end{itemize}

\subsection*{Next Steps}\label{docs__pandoc__latex__src__cli_complete_curriculum.md__next-steps-4}

Once you\textquotesingle re comfortable with these basics:

\begin{itemize}
\tightlist
\item
  Move to Bash Lesson 0: Getting Started for deeper path understanding
\item
  Then continue through Bash Lesson 1 through Bash Lesson 5 for full terminal mastery
\end{itemize}

\subsection*{Resources}\label{docs__pandoc__latex__src__cli_complete_curriculum.md__resources-2}

\begin{itemize}
\tightlist
\item
  Git Bash Installation Guide: \url{https://git-scm.com/book/en/v2/Getting-Started-Installing-Git}
\item
  Bash Basics: \url{https://www.gnu.org/software/bash/manual/bash.html\#Basic-Shell-Features}
\item
  NVDA Screen Reader: \url{https://www.nvaccess.org/}
\item
  JAWS Screen Reader: \url{https://www.freedomscientific.com/products/software/jaws/}
\end{itemize}

\protect\phantomsection\label{docs__pandoc__latex__src__cli_complete_curriculum.md__ps-troubleshooting}{}

\subsection*{Troubleshooting}\label{docs__pandoc__latex__src__cli_complete_curriculum.md__troubleshooting-5}

{\def\LTcaptype{none} % do not increment counter
\begin{longtable}[]{@{}
  >{\raggedright\arraybackslash}p{(\linewidth - 2\tabcolsep) * \real{0.2857}}
  >{\raggedright\arraybackslash}p{(\linewidth - 2\tabcolsep) * \real{0.7143}}@{}}
\toprule\noalign{}
\begin{minipage}[b]{\linewidth}\raggedright
Issue
\end{minipage} & \begin{minipage}[b]{\linewidth}\raggedright
Solution
\end{minipage} \\
\midrule\noalign{}
\endhead
\bottomrule\noalign{}
\endlastfoot
Git Bash won\textquotesingle t open &
Make sure it\textquotesingle s installed; search for "Git Bash" in Start menu \\
Can\textquotesingle t hear the output &
Try redirecting to a file: \texttt{command\ \textgreater{}\ output.txt} \\
Tab completion not working &
Make sure you typed at least one character before pressing Tab \\
Command not found &
Make sure you spelled it correctly; try \texttt{man} for available commands \\
Stuck in a command & Press Ctrl+C to stop it \\
\end{longtable}
}

Still stuck? The checkpoint questions and exercises are your best teacher. Work through them multiple times until comfortable.

Other Screen Readers

Dolphin SuperNova (commercial) and Windows Narrator (built-in) are also supported; the workflows and recommendations in this document apply to them. See \url{https://yourdolphin.com/supernova/} and \url{https://support.microsoft.com/narrator} for vendor documentation.

\begin{center}\rule{0.5\linewidth}{0.5pt}\end{center}

\protect\phantomsection\label{docs__pandoc__latex__src__cli_complete_curriculum.md__gb-lesson-0}{}

\section*{Lesson 0: Getting Started --- Layout, Paths, and the Shell}\label{docs__pandoc__latex__src__cli_complete_curriculum.md__lesson-0-getting-started--layout-paths-and-the-shell-2}

\begin{center}\rule{0.5\linewidth}{0.5pt}\end{center}

Estimated time: 20-30 minutes

\subsection*{Learning Objectives}\label{docs__pandoc__latex__src__cli_complete_curriculum.md__learning-objectives-8}

\begin{itemize}
\tightlist
\item
  Launch Git Bash and locate the prompt
\item
  Understand Unix-style path notation and shortcuts (\texttt{\textasciitilde{}}, \texttt{./}, \texttt{../})
\item
  Use tab completion to navigate quickly
\end{itemize}

\subsection*{Materials}\label{docs__pandoc__latex__src__cli_complete_curriculum.md__materials-4}

\begin{itemize}
\tightlist
\item
  Computer with Git Bash installed
\item
  Editor (Notepad/VS Code)
\end{itemize}

\subsection*{Step-by-step Tasks}\label{docs__pandoc__latex__src__cli_complete_curriculum.md__step-by-step-tasks-6}

\begin{enumerate}
\tightlist
\item
  Open Git Bash and note the prompt (it includes the current path).
\item
  Run \texttt{pwd} and say or note the printed path.
\item
  Use \texttt{ls} to list names in your home directory.
\item
  Practice \texttt{cd\ Documents}, \texttt{cd\ ../} and \texttt{cd\ \textasciitilde{}} until comfortable.
\item
  Try tab-completion: type \texttt{cd\ \textasciitilde{}/D} and press Tab.
\end{enumerate}

\subsection*{Checkpoints}\label{docs__pandoc__latex__src__cli_complete_curriculum.md__checkpoints-10}

\begin{itemize}
\tightlist
\item
  Confirm you can state your current path and move to \texttt{Documents}.
\end{itemize}

\subsection*{Quiz - Lesson GitBash.0}\label{docs__pandoc__latex__src__cli_complete_curriculum.md__quiz---lesson-gitbash0}

\subsection*{3DMake Quick-Start: Your First Real Project}\label{docs__pandoc__latex__src__cli_complete_curriculum.md__3dmake-quick-start-your-first-real-project-2}

You now know enough Git Bash to set up a real 3DMake project. Do this now --- it is the fastest path to a printed object.

\textbf{Step 1 --- Download and run the 3DMake setup wizard:}

\begin{lstlisting}[style=Alabaster, language=bash]
# Go to your Downloads folder (where you extracted 3dmake)
cd ~/Downloads/3dmake

# Run the guided setup (answer the questions about your printer)
./3dm setup

\end{lstlisting}

\textbf{Step 2 --- Create your first project:}

\begin{lstlisting}[style=Alabaster, language=bash]
# Create a folder for all your 3D projects and initialize the first one
mkdir -p ~/Documents/3d_projects/my_cube
cd ~/Documents/3d_projects/my_cube
3dm new

\end{lstlisting}

\textbf{Step 3 --- Confirm the project structure:}

\begin{lstlisting}[style=Alabaster, language=bash]
# You should see three items: src/, build/, and 3dmake.toml
ls -1

\end{lstlisting}

You just used \texttt{mkdir}, \texttt{cd}, and \texttt{ls\ -1} on a real project. You are ready to print.

\subsection*{Checkpoints}\label{docs__pandoc__latex__src__cli_complete_curriculum.md__checkpoints-11}

\begin{itemize}
\tightlist
\item
  Confirm you can state your current path and move to \texttt{Documents}.
\end{itemize}

\subsection*{Quiz - Lesson GitBash.0}\label{docs__pandoc__latex__src__cli_complete_curriculum.md__quiz---lesson-gitbash0-1}

\begin{enumerate}
\tightlist
\item
  What is a path?
\item
  What does \texttt{\textasciitilde{}} mean?
\item
  How do you autocomplete a path?
\item
  How do you go up one directory?
\item
  What command lists file names?
\end{enumerate}

Note

Git Bash uses \textbf{forward slashes} (\texttt{/}) in paths, not Windows-style backslashes (\texttt{\textbackslash{}}). For example: \texttt{/c/Users/YourName/Documents}. This is the Unix convention and transfers directly to macOS and Linux.

\begin{enumerate}
\setcounter{enumi}{5}
\tightlist
\item
  True or False: Git Bash uses backslashes (\texttt{\textbackslash{}}) in paths like Windows CMD.
\item
  Explain the difference between an absolute path and a relative path.
\item
  If you are in \texttt{/c/Users/YourName/Documents} and you type \texttt{cd\ ../}, where do you end up?
\item
  What happens when you press Tab while typing a folder name in Git Bash?
\item
  Describe a practical reason why understanding paths is important for a 3D printing workflow.
\item
  What does \texttt{./} mean in a path, and when would you use it?
\end{enumerate}

Note

When a folder name contains spaces, always wrap the full path in \textbf{double quotes}: \texttt{cd\ "My\ Projects"}. Without quotes, the shell interprets the space as a separator between two separate arguments.

\begin{enumerate}
\setcounter{enumi}{11}
\tightlist
\item
  If a folder path contains spaces (e.g., \texttt{My\ Projects}), how do you navigate to it with \texttt{cd}?
\item
  Explain what the prompt \texttt{YourName@COMPUTER\ MINGW64\ \textasciitilde{}/Documents\ \$} tells you about your current state.
\item
  How would you navigate to your home directory from any location using a single command?
\item
  What is the advantage of using relative paths (like \texttt{../}) versus absolute paths in automation scripts?
\end{enumerate}

\subsection*{Extension Problems}\label{docs__pandoc__latex__src__cli_complete_curriculum.md__extension-problems-15}

\begin{enumerate}
\tightlist
\item
  Create a nested folder and practice \texttt{cd} into it by typing partial names and using Tab.
\item
  Use \texttt{ls\ -la} to list all files including hidden ones in a folder.
\item
  Save \texttt{pwd} output to a file and open it in Notepad.
\item
  Try \texttt{cd} into a folder whose name contains spaces; observe how quotes are handled.
\item
  Create a short note file and open it from Git Bash.
\item
  Build a folder structure that mirrors your project organization; navigate to each level and document the path.
\item
  Create a script that prints your current path and the total number of files in it; run it from different locations.
\item
  Investigate special paths (e.g., \texttt{\$HOME}, \texttt{\$USER}); write down what each contains and when you\textquotesingle d use them.
\item
  Compare absolute vs. relative paths by navigating to the same folder using each method; explain which is easier for automation.
\item
  Create a bash function that changes to a frequently-used folder and lists its contents in one command; test it from different starting locations.
\item
  Navigate to three different locations and at each one note the prompt, the path from \texttt{pwd}, and verify you understand what each shows.
\item
  Create a complex folder tree (at least 5 levels deep) and navigate it using only relative paths; verify your location at each step.
\item
  Document all shortcuts you know (\texttt{\textasciitilde{}}, \texttt{./}, \texttt{../}, \texttt{\$HOME}) and demonstrate each one works as expected.
\item
  Write a guide for a peer on how to understand the Git Bash prompt and path notation without using GUI file explorer.
\item
  Create a troubleshooting flowchart: if someone says "I don\textquotesingle t know where I am," what commands do you give them to find out?
\end{enumerate}

\subsection*{References}\label{docs__pandoc__latex__src__cli_complete_curriculum.md__references-13}

\begin{itemize}
\tightlist
\item
  GNU. (2024). \emph{Bash Manual}. \url{https://example.com}
\item
  Git SCM. (2024). \emph{Git Bash documentation}. \url{https://example.com}
\item
  Linux Foundation. (2024). \emph{The Linux Command Line}. \url{https://example.com}
\end{itemize}

\subsection*{Helpful Resources}\label{docs__pandoc__latex__src__cli_complete_curriculum.md__helpful-resources-11}

\begin{itemize}
\tightlist
\item
  \href{https://www.gnu.org/software/bash/manual/htmlnode/Bourne-Shell-Builtins.html}{Bash Manual - Navigation}
\item
  \href{https://git-scm.com/book/en/v2/Git-Basics-Getting-a-Git-Repository}{Git Bash Basics}
\item
  \href{https://linuxcommand.org/lc3lts0020.php}{Linux Path Guide}
\item
  \href{https://www.gnu.org/software/bash/manual/htmlnode/Controlling-the-Prompt.html}{Understanding Bash Prompts}
\item
  \href{https://www.gnu.org/software/bash/manual/htmlnode/Commands-For-Completion.html}{Tab Completion in Bash}
\end{itemize}

\begin{center}\rule{0.5\linewidth}{0.5pt}\end{center}

\protect\phantomsection\label{docs__pandoc__latex__src__cli_complete_curriculum.md__gb-lesson-1}{}

\section*{Lesson 1: Navigation --- Moving Around Your File System}\label{docs__pandoc__latex__src__cli_complete_curriculum.md__lesson-1-navigation--moving-around-your-file-system-2}

\begin{center}\rule{0.5\linewidth}{0.5pt}\end{center}

\textbf{Duration:} 1 class period
\textbf{Prerequisite:} Bash Lesson 0 (Getting Started)

\subsection*{Learning Objectives}\label{docs__pandoc__latex__src__cli_complete_curriculum.md__learning-objectives-9}

By the end of this lesson, you will be able to:

\begin{itemize}
\tightlist
\item
  Use \texttt{pwd} to print your current location
\item
  Use \texttt{cd} to move between directories
\item
  Use \texttt{ls} (and its flags) to list files and folders
\item
  Use wildcards \texttt{*} and \texttt{?} to filter listings
\item
  Navigate relative vs. absolute paths
\item
  Search for files by name and extension
\end{itemize}

\subsection*{Materials}\label{docs__pandoc__latex__src__cli_complete_curriculum.md__materials-5}

\begin{itemize}
\tightlist
\item
  Git Bash
\item
  Text editor (Notepad or VS Code)
\end{itemize}

\begin{center}\rule{0.5\linewidth}{0.5pt}\end{center}

\subsection*{Commands Covered in This Lesson}\label{docs__pandoc__latex__src__cli_complete_curriculum.md__commands-covered-in-this-lesson-1}

{\def\LTcaptype{none} % do not increment counter
\begin{longtable}[]{@{}
  >{\raggedright\arraybackslash}p{(\linewidth - 2\tabcolsep) * \real{0.2432}}
  >{\raggedright\arraybackslash}p{(\linewidth - 2\tabcolsep) * \real{0.7568}}@{}}
\toprule\noalign{}
\begin{minipage}[b]{\linewidth}\raggedright
Command
\end{minipage} & \begin{minipage}[b]{\linewidth}\raggedright
What It Does
\end{minipage} \\
\midrule\noalign{}
\endhead
\bottomrule\noalign{}
\endlastfoot
\texttt{pwd} & Print Working Directory - shows where you are \\
\texttt{cd\ path} & Change Directory - move to a new location \\
\texttt{ls} & List - shows files and folders in current location \\
\texttt{ls\ -1} &
List names only, one per line (screen reader friendly) \\
\texttt{ls\ -1\ -F} &
List names with type indicators (/ for directories) \\
\texttt{ls\ *.extension} & List files matching a pattern \\
\end{longtable}
}

\begin{center}\rule{0.5\linewidth}{0.5pt}\end{center}

\subsection*{\texorpdfstring{\texttt{pwd} - Where Am I?}{pwd - Where Am I?}}\label{docs__pandoc__latex__src__cli_complete_curriculum.md__pwd---where-am-i-1}

Type \texttt{pwd} and press \texttt{Enter}. Git Bash prints the full path to your current location.

\begin{lstlisting}[style=Alabaster, language=bash]
pwd
# Output: /c/Users/YourName

\end{lstlisting}

\textbf{When to use:} Always run this if you\textquotesingle re unsure of your current location.

\begin{center}\rule{0.5\linewidth}{0.5pt}\end{center}

\subsection*{\texorpdfstring{\texttt{cd} - Changing Directories}{cd - Changing Directories}}\label{docs__pandoc__latex__src__cli_complete_curriculum.md__cd---changing-directories-1}

\texttt{cd} stands for "change directory."

\begin{lstlisting}[style=Alabaster, language=bash]
# Go to Documents
cd Documents

# Go up one level to parent directory
cd ..

# Go to root of file system
cd /

# Go to home directory
cd ~

# Go to a specific path
cd /c/Users/YourName/Documents/3D_Projects

# Go to previous directory
cd -

\end{lstlisting}

\begin{center}\rule{0.5\linewidth}{0.5pt}\end{center}

\subsection*{\texorpdfstring{\texttt{ls} - Listing Files and Folders}{ls - Listing Files and Folders}}\label{docs__pandoc__latex__src__cli_complete_curriculum.md__ls---listing-files-and-folders-1}

Use \texttt{ls\ -1} for screen reader compatibility --- names only, one per line.

\begin{lstlisting}[style=Alabaster, language=bash]
# List all files and folders (names only, one per line)
ls -1

# List only files (no hidden, no directories)
ls -1 -p | grep -v /

# List only directories
ls -1 -d */

\end{lstlisting}

\begin{center}\rule{0.5\linewidth}{0.5pt}\end{center}

\subsection*{Wildcards - Finding Files by Pattern}\label{docs__pandoc__latex__src__cli_complete_curriculum.md__wildcards---finding-files-by-pattern-1}

Wildcards help you find files without typing the full name.

\textbf{\texttt{*} (asterisk)} matches any number of characters:

\begin{lstlisting}[style=Alabaster, language=bash]
# List all .scad files
ls *.scad

# List all files starting with "part"
ls part*

# List all files ending with "_final"
ls *_final*

\end{lstlisting}

\textbf{\texttt{?} (question mark)} matches exactly one character:

\begin{lstlisting}[style=Alabaster, language=bash]
# Find files like model1.scad, model2.scad (but not model12.scad)
ls model?.scad

\end{lstlisting}

\begin{center}\rule{0.5\linewidth}{0.5pt}\end{center}

\subsection*{Step-by-step Practice}\label{docs__pandoc__latex__src__cli_complete_curriculum.md__step-by-step-practice-2}

\begin{enumerate}
\tightlist
\item
  Run \texttt{pwd} and confirm your location
\item
  Move to \texttt{Documents}: \texttt{cd\ Documents}
\end{enumerate}

\subsection*{3DMake in Practice: Navigating a Project}\label{docs__pandoc__latex__src__cli_complete_curriculum.md__3dmake-in-practice-navigating-a-project-2}

Navigation commands are exactly what you need to explore and manage 3DMake projects:

\begin{lstlisting}[style=Alabaster, language=bash]
# Go to your project root
cd ~/Documents/3d_projects/my_cube

# See the project layout: src/, build/, 3dmake.toml
ls -1

# Enter the source folder where your .scad model files live
cd src
ls -1

# Open the main model in your text editor
3dm edit-model

# Go back to project root, then check build output
cd ..
cd build
ls -1 *.stl

# Return to project root
cd ..

# Use wildcards to find all projects with "cube" in the name
ls -1 ~/Documents/3d_projects/*cube*

\end{lstlisting}

Tab completion makes navigation fast: type \texttt{cd\ \textasciitilde{}/Doc} and press Tab.

\begin{center}\rule{0.5\linewidth}{0.5pt}\end{center}

\subsection*{Step-by-step Practice}\label{docs__pandoc__latex__src__cli_complete_curriculum.md__step-by-step-practice-3}

\begin{enumerate}
\tightlist
\item
  Run \texttt{pwd} and confirm your location
\item
  Move to \texttt{Documents}: \texttt{cd\ Documents}
\item
  Confirm you moved: \texttt{pwd}
\item
  List files and folders: \texttt{ls\ -1}
\item
  List only files: \texttt{ls\ -1\ -p\ \textbar{}\ grep\ -v\ /}
\item
  Go back up: \texttt{cd\ ..}
\item
  Search for files: \texttt{ls\ *.txt}
\end{enumerate}

\begin{center}\rule{0.5\linewidth}{0.5pt}\end{center}

\subsection*{Checkpoints}\label{docs__pandoc__latex__src__cli_complete_curriculum.md__checkpoints-12}

After this lesson, you should be able to:

\begin{itemize}
\tightlist
\item[$\square$]
  Navigate to any folder using \texttt{cd}
\item[$\square$]
  Confirm your location with \texttt{pwd}
\item[$\square$]
  List files and folders with \texttt{ls\ -1}
\item[$\square$]
  Use wildcards to find files by pattern
\item[$\square$]
  Move between absolute and relative paths confidently
\end{itemize}

\subsection*{Quiz - Lesson GitBash.1}\label{docs__pandoc__latex__src__cli_complete_curriculum.md__quiz---lesson-gitbash1}

\begin{enumerate}
\tightlist
\item
  What does \texttt{pwd} show?
\item
  How do you list directories only with \texttt{ls}?
\item
  What wildcard matches any number of characters?
\item
  How do you list files with the \texttt{.scad} extension?
\item
  Give an example of an absolute path and a relative path.
\item
  True or False: The \texttt{*} wildcard matches exactly one character.
\item
  Explain the difference between \texttt{ls\ -1} and \texttt{ls\ -1\ -d\ */}.
\item
  Write a command that would list all \texttt{.txt} files in your Documents folder using a wildcard.
\item
  How would you search for files containing "part" in their name across multiple files?
\item
  Describe a practical scenario where using wildcards saves time in a 3D printing workflow.
\item
  What happens when you use \texttt{ls\ model?.scad} versus \texttt{ls\ model*.scad}?
\item
  How would you navigate to a folder whose name contains both spaces and special characters?
\item
  If you\textquotesingle re in \texttt{\textasciitilde{}/Documents/Projects/3D} and you want to go to \texttt{\textasciitilde{}/Documents/Resources}, what command would you use?
\item
  Write a command sequence that navigates to the Downloads folder, lists only files, then returns to home.
\item
  Explain the purpose of using \texttt{ls\ -1} specifically in a screen reader context.
\end{enumerate}

\subsection*{Extension Problems}\label{docs__pandoc__latex__src__cli_complete_curriculum.md__extension-problems-16}

\begin{enumerate}
\tightlist
\item
  Write a one-line command that lists \texttt{.scad} files and saves to \texttt{scad\_list.txt}.
\item
  Use \texttt{ls\ -1\ \textasciitilde{}/Documents\ \textbar{}\ less} to page through long listings.
\item
  Combine \texttt{ls} with \texttt{grep} to search for a filename pattern.
\item
  Create a shortcut alias in the session for a long path and test it.
\item
  Practice tab-completion in a directory with many similarly named files.
\item
  Build a bash script that recursively lists all \texttt{.scad} and \texttt{.stl} files in a directory tree; save the results to a file.
\item
  Compare the output of \texttt{ls}, \texttt{ls\ -1}, \texttt{ls\ -la}, and \texttt{ls\ -1\ -d\ */} to understand the flags; document what each command does.
\item
  Create a filtering command that displays only files modified in the last 7 days; test it on your documents folder.
\item
  Write a non-visual guide to Git Bash navigation; include descriptions of common patterns and how to verify directory contents audibly.
\item
  Develop a navigation workflow for a typical 3D printing project: move between CAD, slicing, and print-log folders efficiently; document the commands.
\item
  Create a complex wildcard search: find all files in a folder and subfolders that match multiple patterns (e.g., \texttt{*\_v1.*} or \texttt{*\_final.*}).
\item
  Build a script that navigates through a folder tree, counts files at each level, and reports the structure.
\item
  Document the output differences between \texttt{ls}, \texttt{ls\ -1}, \texttt{ls\ -la}, and \texttt{ls\ -1\ -d\ */}; explain when to use each.
\item
  Create a navigation "cheat sheet" as a bash script that prints common paths and how to navigate to them.
\item
  Design a project folder structure on your computer, document each path, then create a script that validates all folders exist.
\end{enumerate}

\subsection*{References}\label{docs__pandoc__latex__src__cli_complete_curriculum.md__references-14}

\begin{itemize}
\tightlist
\item
  GNU. (2024). \emph{ls command reference}. \url{https://example.com}
\item
  GNU. (2024). \emph{Bash wildcards and globbing}. \url{https://example.com}
\item
  GNU. (2024). \emph{Navigation best practices in Bash}. \url{https://example.com}
\end{itemize}

\subsection*{Helpful Resources}\label{docs__pandoc__latex__src__cli_complete_curriculum.md__helpful-resources-12}

\begin{itemize}
\tightlist
\item
  \href{https://www.gnu.org/software/coreutils/manual/html_node/ls-invocation.html}{ls Command Reference}
\item
  \href{https://www.gnu.org/software/bash/manual/htmlnode/Pattern-Matching.html}{Bash Wildcards and Globbing}
\item
  \href{https://linuxcommand.org/lc3lts0020.php}{Navigation Best Practices}
\item
  \href{https://linuxize.com/post/linux-cd-command/}{Relative and Absolute Paths in Bash}
\item
  \href{https://www.nvaccess.org/documentation/}{Screen Reader Tips for Git Bash}
\end{itemize}

\begin{center}\rule{0.5\linewidth}{0.5pt}\end{center}

\protect\phantomsection\label{docs__pandoc__latex__src__cli_complete_curriculum.md__gb-lesson-2}{}

\section*{Lesson 2: File and Folder Manipulation}\label{docs__pandoc__latex__src__cli_complete_curriculum.md__lesson-2-file-and-folder-manipulation-2}

\begin{center}\rule{0.5\linewidth}{0.5pt}\end{center}

Estimated time: 30-45 minutes

\subsection*{Learning Objectives}\label{docs__pandoc__latex__src__cli_complete_curriculum.md__learning-objectives-10}

\begin{itemize}
\tightlist
\item
  Create, copy, move, and delete files and folders from Git Bash
\item
  Use \texttt{touch}, \texttt{mkdir}, \texttt{cp}, \texttt{mv}, \texttt{rm}, and \texttt{rmdir} safely
\item
  Understand when operations are permanent and how to confirm results
\end{itemize}

\subsection*{Materials}\label{docs__pandoc__latex__src__cli_complete_curriculum.md__materials-6}

\begin{itemize}
\tightlist
\item
  Git Bash
\item
  Small practice folder for exercises
\end{itemize}

\subsection*{Step-by-step Tasks}\label{docs__pandoc__latex__src__cli_complete_curriculum.md__step-by-step-tasks-7}

\begin{enumerate}
\tightlist
\item
  Create a practice directory: \texttt{mkdir\ \textasciitilde{}/Documents/GitBash\_Practice} and \texttt{cd} into it.
\item
  Create two files: \texttt{touch\ file1.txt} and \texttt{touch\ file2.txt}.
\item
  Copy \texttt{file1.txt} to \texttt{file1\_backup.txt} with \texttt{cp} and confirm with \texttt{ls\ -1}.
\item
  Rename \texttt{file2.txt} to \texttt{notes.txt} using \texttt{mv} and confirm.
\item
  Delete \texttt{file1.txt} with \texttt{rm} and verify the backup remains.
\end{enumerate}

\subsection*{Checkpoints}\label{docs__pandoc__latex__src__cli_complete_curriculum.md__checkpoints-13}

\begin{itemize}
\tightlist
\item
  After step 3 you should see both the original and the backup file.
\end{itemize}

\subsection*{3DMake in Practice: Organizing Your Projects}\label{docs__pandoc__latex__src__cli_complete_curriculum.md__3dmake-in-practice-organizing-your-projects-2}

File management commands are essential for 3DMake project organization:

\begin{lstlisting}[style=Alabaster, language=bash]
# Create a new project from scratch
mkdir -p ~/Documents/3d_projects/bottle_cap
cd ~/Documents/3d_projects/bottle_cap
3dm new

# Back up a working project before making big changes
cd ~/Documents/3d_projects
cp -r bottle_cap bottle_cap_backup

# Rename the main model file to something more specific
cd bottle_cap/src
mv main.scad bottle_cap.scad

# List only .scad files to confirm the rename
ls -1 *.scad

# Delete a draft STL you no longer need
rm ~/Documents/3d_projects/bottle_cap/build/draft.stl

# Build the renamed model (using the -m flag to specify the name)
cd ~/Documents/3d_projects/bottle_cap
3dm build -m bottle_cap

# See the fresh output in the build folder
ls -1 build/

\end{lstlisting}

\subsection*{Checkpoints}\label{docs__pandoc__latex__src__cli_complete_curriculum.md__checkpoints-14}

\begin{itemize}
\tightlist
\item
  After step 3 you should see both the original and the backup file.
\end{itemize}

\subsection*{Quiz - Lesson GitBash.2}\label{docs__pandoc__latex__src__cli_complete_curriculum.md__quiz---lesson-gitbash2}

\begin{enumerate}
\tightlist
\item
  How do you create an empty file from Git Bash?
\item
  What command copies a file?
\item
  How do you rename a file?
\item
  What does \texttt{rm\ -r} do?
\end{enumerate}

Warning

\texttt{rm} permanently deletes files with no undo and no Recycle Bin. There is no recovery from Git Bash. Always run \texttt{ls\ -1} first to confirm exactly what you\textquotesingle re targeting, especially when using wildcards.

\begin{enumerate}
\setcounter{enumi}{4}
\tightlist
\item
  Why is \texttt{rm} potentially dangerous?
\item
  True or False: \texttt{cp} requires the \texttt{-r} flag to copy both files and folders.
\item
  Explain the difference between \texttt{rm} and \texttt{rmdir}.
\item
  If you delete a file with \texttt{rm}, can you recover it from Git Bash?
\item
  Write a command that would copy an entire folder and all its contents to a new location.
\item
  Describe a practical safety check you would perform before running \texttt{rm\ -r} on a folder.
\item
  What happens if you \texttt{cp} a file to a destination where a file with the same name already exists? How would you handle this safely?
\item
  Compare \texttt{mv\ old\_name.txt\ new\_name.txt} vs \texttt{mv\ old\_name.txt\ \textasciitilde{}/Documents/new\_name.txt}. What is the key difference?
\item
  Design a workflow to safely delete 50 files matching the pattern \texttt{*.bak} from a folder containing 500 files. What commands and verifications would you use?
\item
  Explain how you could back up all \texttt{.scad} files from a project folder into a timestamped backup folder in one command.
\item
  When organizing a 3D printing project, you need to move completed designs to an archive folder and delete failed prototypes. How would you structure this as a safe, auditable process?
\end{enumerate}

\subsection*{Extension Problems}\label{docs__pandoc__latex__src__cli_complete_curriculum.md__extension-problems-17}

\begin{enumerate}
\tightlist
\item
  Create a folder tree and copy it to a new location with \texttt{cp\ -r}.
\item
  Write a one-line command that creates three files named \texttt{a}, \texttt{b}, \texttt{c} and lists them.
\item
  Move a file into a new folder and confirm the move.
\item
  Use wildcards to delete files matching a pattern in a safe test folder.
\item
  Export a listing of the practice folder to \texttt{practice\_listing.txt}.
\item
  Create a backup shell script that copies all \texttt{.scad} files from your project folder to a backup folder with timestamp naming.
\item
  Build a safe deletion workflow: list files matching a pattern, verify count, then delete with confirmation; document the steps.
\item
  Write a bash script that organizes files by extension into subfolders; test it on a sample folder tree.
\item
  Create a file operation audit trail: log all copy, move, and delete operations to a text file for review.
\item
  Develop a project template generator: a bash script that creates a standard folder structure for a new 3D printing project with essential subfolders.
\item
  Implement a file conflict handler: write a bash script that handles cases where \texttt{cp} would overwrite an existing file by renaming the existing file with a timestamp before copying.
\item
  Create a batch rename operation: use a script to rename all files in a folder from \texttt{old\_prefix\_*} to \texttt{new\_prefix\_*}; test with actual files and verify the results.
\item
  Build a folder comparison tool: list all files in two folders and identify which files exist in one but not the other; output to a report.
\item
  Write a destructive operation validator: before executing \texttt{rm\ -r}, create a script that lists exactly what will be deleted, shows file counts by type, and requires explicit user confirmation to proceed.
\item
  Design a complete project lifecycle workflow: create folders for active projects, completed designs, and archive; include move operations between folders, backup steps, and verification that all files arrive intact.
\end{enumerate}

\subsection*{References}\label{docs__pandoc__latex__src__cli_complete_curriculum.md__references-15}

\begin{itemize}
\tightlist
\item
  GNU. (2024). \emph{touch command reference}. \url{https://example.com}
\item
  GNU. (2024). \emph{cp and mv commands}. \url{https://example.com}
\item
  GNU. (2024). \emph{File system operations guide}. \url{https://example.com}
\end{itemize}

\subsection*{Helpful Resources}\label{docs__pandoc__latex__src__cli_complete_curriculum.md__helpful-resources-13}

\begin{itemize}
\tightlist
\item
  \href{https://www.gnu.org/software/coreutils/manual/html_node/touch-invocation.html}{touch Command Reference}
\item
  \href{https://www.gnu.org/software/coreutils/manual/html_node/cp-invocation.html}{cp Command Reference}
\item
  \href{https://www.gnu.org/software/coreutils/manual/html_node/mv-invocation.html}{mv Command Reference}
\item
  \href{https://www.gnu.org/software/coreutils/manual/html_node/rm-invocation.html}{rm Command Reference}
\item
  \href{https://linuxcommand.org/lc3lts0050.php}{Safe Deletion Practices}
\end{itemize}

\begin{center}\rule{0.5\linewidth}{0.5pt}\end{center}

\protect\phantomsection\label{docs__pandoc__latex__src__cli_complete_curriculum.md__gb-lesson-3}{}

\section*{Lesson 3: Input, Output, and Piping}\label{docs__pandoc__latex__src__cli_complete_curriculum.md__lesson-3-input-output-and-piping-2}

\begin{center}\rule{0.5\linewidth}{0.5pt}\end{center}

\textbf{Duration:} 1 class period
\textbf{Prerequisite:} Bash Lesson 2 (File and Folder Manipulation)

\begin{center}\rule{0.5\linewidth}{0.5pt}\end{center}

\subsection*{Learning Objectives}\label{docs__pandoc__latex__src__cli_complete_curriculum.md__learning-objectives-11}

By the end of this lesson, you will be able to:

\begin{itemize}
\tightlist
\item
  Use \texttt{echo} to print text to the screen
\item
  Use \texttt{cat} to read file contents
\item
  Use \texttt{\textgreater{}} to redirect output into a file
\item
  Use \texttt{\textbar{}} (pipe) to send one command\textquotesingle s output to another
\item
  Copy output to the clipboard with \texttt{clip}
\item
  Open files with a text editor from the command line
\end{itemize}

\begin{center}\rule{0.5\linewidth}{0.5pt}\end{center}

\subsection*{Commands Covered}\label{docs__pandoc__latex__src__cli_complete_curriculum.md__commands-covered-2}

{\def\LTcaptype{none} % do not increment counter
\begin{longtable}[]{@{}
  >{\raggedright\arraybackslash}p{(\linewidth - 2\tabcolsep) * \real{0.3077}}
  >{\raggedright\arraybackslash}p{(\linewidth - 2\tabcolsep) * \real{0.6923}}@{}}
\toprule\noalign{}
\begin{minipage}[b]{\linewidth}\raggedright
Command
\end{minipage} & \begin{minipage}[b]{\linewidth}\raggedright
What It Does
\end{minipage} \\
\midrule\noalign{}
\endhead
\bottomrule\noalign{}
\endlastfoot
\texttt{echo\ "text"} & Print text to the screen \\
\texttt{cat\ filename} & Print the contents of a file \\
\texttt{\textgreater{}\ filename} &
Redirect output into a file (overwrites) \\
\texttt{\textgreater{}\textgreater{}\ filename} &
Append output to a file (adds to end) \\
\texttt{\textbar{}} & Pipe - send output from one command to the next \\
\texttt{clip} & Copy piped input to the Windows clipboard (Git Bash) \\
\texttt{notepad.exe\ filename} & Open a file in Notepad \\
\end{longtable}
}

\begin{center}\rule{0.5\linewidth}{0.5pt}\end{center}

\subsection*{\texorpdfstring{\texttt{echo} - Printing Text}{echo - Printing Text}}\label{docs__pandoc__latex__src__cli_complete_curriculum.md__echo---printing-text-2}

\texttt{echo} prints text to the screen. It is useful for testing, for writing text into files, and for understanding how piping works.

\begin{lstlisting}[style=Alabaster, language=bash]
echo "Hello, World"
echo "This is a test"

\end{lstlisting}

\begin{center}\rule{0.5\linewidth}{0.5pt}\end{center}

\subsection*{\texorpdfstring{\texttt{cat} - Reading Files}{cat - Reading Files}}\label{docs__pandoc__latex__src__cli_complete_curriculum.md__cat---reading-files-1}

\texttt{cat} prints the contents of a file to the screen.

\begin{lstlisting}[style=Alabaster, language=bash]
# Read a text file
cat ~/Documents/notes.txt

# Read an OpenSCAD file
cat ~/Documents/OpenSCAD_Projects/project0.scad

\end{lstlisting}

With a long file, use \texttt{cat\ filename\ \textbar{}\ less} to read it page by page (press \texttt{Space} to advance, \texttt{Q} to quit).

\begin{center}\rule{0.5\linewidth}{0.5pt}\end{center}

\subsection*{\texorpdfstring{\texttt{\textgreater{}} - Redirecting Output to a File}{\textgreater{} - Redirecting Output to a File}}\label{docs__pandoc__latex__src__cli_complete_curriculum.md__---redirecting-output-to-a-file-2}

The \texttt{\textgreater{}} symbol redirects output from the screen into a file instead.

\begin{lstlisting}[style=Alabaster, language=bash]
# Create a file with a single line
echo "Author: Your Name" > header.txt

# Confirm the file was created and has content
cat header.txt

\end{lstlisting}

\textbf{Warning:} \texttt{\textgreater{}} overwrites the file if it already exists. Use \texttt{\textgreater{}\textgreater{}} to append instead:

\begin{lstlisting}[style=Alabaster, language=bash]
echo "Date: 2025" >> header.txt
echo "Project: Floor Marker" >> header.txt
cat header.txt

\end{lstlisting}

\begin{center}\rule{0.5\linewidth}{0.5pt}\end{center}

\subsection*{\texorpdfstring{\texttt{\textbar{}} - Piping}{\textbar{} - Piping}}\label{docs__pandoc__latex__src__cli_complete_curriculum.md__---piping-2}

The pipe symbol \texttt{\textbar{}} sends the output of one command to the input of the next. This lets you chain commands together.

\begin{lstlisting}[style=Alabaster, language=bash]
# List files and send the list to clip (copies to clipboard)
ls -1 | clip
## 3DMake in Practice: Capturing Build Output

Piping and redirection are essential for working with 3DMake's output accessibly:

```bash
# Build and save the full log to a file so you can read it
3dm build > build_log.txt
cat build_log.txt

# Build and slice, keeping only lines with errors or warnings
3dm build slice 2>&1 | grep -i "error\|warning"

# List all STL files produced so far
ls -1 build/*.stl

# Get a list of available printer profiles and save for reference
3dm list-profiles > profiles.txt
cat profiles.txt

# Get a list of material overlays (PLA, PETG, supports, etc.)
3dm list-overlays > overlays.txt
cat overlays.txt

# Get an AI description of your model, then copy to clipboard
3dm info > model_description.txt
cat model_description.txt | clip

\end{lstlisting}

\begin{lstlisting}[style=Alabaster, language=bash]
# List files and send the list to clip (copies to clipboard)
ls -1 | clip

# Now paste with Ctrl + V anywhere

\end{lstlisting}

\begin{lstlisting}[style=Alabaster, language=bash]
# Search within a file's contents using grep
cat project0.scad | grep "cube"

\end{lstlisting}

\begin{center}\rule{0.5\linewidth}{0.5pt}\end{center}

\subsection*{\texorpdfstring{\texttt{clip} - Copying to Clipboard}{clip - Copying to Clipboard}}\label{docs__pandoc__latex__src__cli_complete_curriculum.md__clip---copying-to-clipboard-2}

\texttt{clip} takes whatever is piped to it and puts it on the Windows clipboard.

\begin{lstlisting}[style=Alabaster, language=bash]
# Copy your current directory path to the clipboard
pwd | clip

# Copy a file listing to clipboard
ls -1 | clip

# Copy the contents of a file to clipboard
cat notes.txt | clip

\end{lstlisting}

After any of these, press \texttt{Ctrl\ +\ V} in any application to paste.

\begin{center}\rule{0.5\linewidth}{0.5pt}\end{center}

\subsection*{Opening Files in Notepad}\label{docs__pandoc__latex__src__cli_complete_curriculum.md__opening-files-in-notepad-2}

\begin{lstlisting}[style=Alabaster, language=bash]
# Open a file in Notepad
notepad.exe ~/Documents/notes.txt

# Open a .scad file
notepad.exe ~/Documents/OpenSCAD_Projects/project0.scad

# Create a new file and open it
touch new_notes.txt
notepad.exe new_notes.txt

\end{lstlisting}

\begin{center}\rule{0.5\linewidth}{0.5pt}\end{center}

\subsection*{Step-by-step Tasks}\label{docs__pandoc__latex__src__cli_complete_curriculum.md__step-by-step-tasks-8}

\begin{enumerate}
\tightlist
\item
  Create \texttt{practice.txt} with three lines using \texttt{echo} and \texttt{\textgreater{}}/\texttt{\textgreater{}\textgreater{}}.
\item
  Read the file with \texttt{cat\ practice.txt}.
\item
  Pipe the file into \texttt{grep} to search for a word.
\item
  Copy the file contents to clipboard with \texttt{cat\ practice.txt\ \textbar{}\ clip}.
\item
  Redirect \texttt{ls\ -1} into \texttt{list.txt} and open it in Notepad.
\end{enumerate}

\subsection*{Checkpoints}\label{docs__pandoc__latex__src__cli_complete_curriculum.md__checkpoints-15}

\begin{itemize}
\tightlist
\item
  After step 3 you should be able to find a keyword using piping.
\end{itemize}

\subsection*{Quiz - Lesson GitBash.3}\label{docs__pandoc__latex__src__cli_complete_curriculum.md__quiz---lesson-gitbash3}

\begin{enumerate}
\tightlist
\item
  What is the difference between \texttt{\textgreater{}} and \texttt{\textgreater{}\textgreater{}}?
\item
  What does the pipe \texttt{\textbar{}} do?
\item
  How do you copy output to the clipboard?
\item
  How would you page through long output?
\item
  How do you suppress output (send it to \texttt{/dev/null})?
\item
  True or False: The pipe operator \texttt{\textbar{}} connects the output of one command to the input of another.
\item
  Explain why redirecting output to a file is useful for screen reader users.
\item
  Write a command that would search for the word "sphere" in all \texttt{.scad} files in a directory.
\item
  How would you count the number of lines in a file using bash piping?
\item
  Describe a practical scenario in 3D printing where you would pipe or redirect command output.
\item
  What would be the difference in output between \texttt{echo\ "test"\ \textgreater{}\ file.txt} (run twice) vs \texttt{echo\ "test"\ \textgreater{}\textgreater{}\ file.txt} (run twice)? Show the expected file contents.
\item
  Design a three-step piping chain: read a file, filter for specific content, and save the results; explain what each pipe does.
\item
  You have a 500-line \texttt{.scad} file and need to find all instances of \texttt{sphere()} and count them. Write the command.
\item
  Explain how \texttt{clip} is particularly valuable for screen reader users when working with file paths or long output strings.
\item
  Describe how you would use pipes and redirection to create a timestamped backup report of all \texttt{.stl} files in a 3D printing project folder.
\end{enumerate}

\subsection*{Extension Problems}\label{docs__pandoc__latex__src__cli_complete_curriculum.md__extension-problems-18}

\begin{enumerate}
\tightlist
\item
  Use piping to count lines in a file (hint: \texttt{cat\ file.txt\ \textbar{}\ wc\ -l}).
\item
  Save a long \texttt{ls\ -1} output and search it with \texttt{grep}.
\item
  Chain multiple pipes to filter and then save results.
\item
  Practice copying different command outputs to clipboard and pasting.
\item
  Create a small bash script that generates a report (counts of files by extension).
\item
  Build a data processing pipeline: read a text file, filter rows, and export results; document each step.
\item
  Write a script that pipes directory listing to count occurrences of each file extension; create a summary report.
\item
  Create a log analysis command: read a log file, filter for errors, and save matching lines to a separate error log.
\item
  Design a piping workflow for 3D printing file management: find \texttt{.stl} files, extract their names, and generate a report.
\item
  Develop a reusable piping function library: create bash functions for common filtering, sorting, and reporting patterns; test each with different inputs.
\item
  Build a complex filter pipeline: read a \texttt{.scad} file, extract lines containing specific geometry commands, count each type, and output a summary to both screen and file.
\item
  Create an interactive filtering tool: build a bash script that accepts a search term, pipes through multiple filters, and displays paginated results.
\item
  Develop a performance analysis tool: use piping to combine file listing, metadata extraction, and statistical reporting; export results to a dated report file.
\item
  Implement a comprehensive error-handling pipeline: read output, catch errors, log them separately, and generate a summary of successes vs failures.
\item
  Design and execute a real-world project backup workflow: use piping to verify file existence, count files by type, generate a backup manifest, and create audit logs --- all in one integrated command pipeline.
\end{enumerate}

\subsection*{References}\label{docs__pandoc__latex__src__cli_complete_curriculum.md__references-16}

\begin{itemize}
\tightlist
\item
  GNU. (2024). \emph{Redirections in Bash}. \url{https://example.com}
\item
  GNU. (2024). \emph{grep command reference}. \url{https://example.com}
\item
  GNU. (2024). \emph{Bash pipeline concepts}. \url{https://example.com}
\end{itemize}

\subsection*{Helpful Resources}\label{docs__pandoc__latex__src__cli_complete_curriculum.md__helpful-resources-14}

\begin{itemize}
\tightlist
\item
  \href{https://www.gnu.org/software/bash/manual/htmlnode/Redirections.html}{Bash Redirections}
\item
  \href{https://www.gnu.org/software/grep/manual/grep.html}{Piping and grep}
\item
  \href{https://www.gnu.org/software/coreutils/manual/html_node/cat-invocation.html}{cat Command Reference}
\item
  \href{https://www.gnu.org/software/coreutils/manual/html_node/wc-invocation.html}{wc for Counting}
\item
  \href{https://www.gnu.org/software/bash/manual/htmlnode/Pipelines.html}{Bash Pipeline Concepts}
\end{itemize}

\begin{center}\rule{0.5\linewidth}{0.5pt}\end{center}

\protect\phantomsection\label{docs__pandoc__latex__src__cli_complete_curriculum.md__gb-lesson-4}{}

\section*{Lesson 4: Environment Variables, PATH, and Aliases}\label{docs__pandoc__latex__src__cli_complete_curriculum.md__lesson-4-environment-variables-path-and-aliases-2}

\begin{center}\rule{0.5\linewidth}{0.5pt}\end{center}

Estimated time: 30-45 minutes

\subsection*{Learning Objectives}\label{docs__pandoc__latex__src__cli_complete_curriculum.md__learning-objectives-12}

\begin{itemize}
\tightlist
\item
  Read environment variables with \texttt{\$VARNAME}
\item
  Inspect and verify programs in the \texttt{PATH}
\item
  Create temporary aliases and understand making them persistent via \texttt{.bashrc}
\end{itemize}

\subsection*{Materials}\label{docs__pandoc__latex__src__cli_complete_curriculum.md__materials-7}

\begin{itemize}
\tightlist
\item
  Git Bash (with rights to edit \texttt{.bashrc})
\end{itemize}

\subsection*{Step-by-step Tasks}\label{docs__pandoc__latex__src__cli_complete_curriculum.md__step-by-step-tasks-9}

\begin{enumerate}
\tightlist
\item
  Show your username and home path with \texttt{echo\ \$USER} and \texttt{echo\ \$HOME}.
\item
  Inspect \texttt{echo\ \$PATH} and identify whether \texttt{openscad} or \texttt{code} would be found.
\item
  Run \texttt{which\ openscad} and note the result.
\item
  Create a temporary alias: \texttt{alias\ preview=\textquotesingle{}openscad\textquotesingle{}} and test \texttt{preview\ myfile.scad}.
\item
  Open your profile (\texttt{notepad.exe\ \textasciitilde{}/.bashrc}) and add the alias line to make it persistent.
\end{enumerate}

\subsection*{Checkpoints}\label{docs__pandoc__latex__src__cli_complete_curriculum.md__checkpoints-16}

\begin{itemize}
\tightlist
\item
  After step 3 you can determine whether a program will be found by \texttt{PATH}.
\end{itemize}

\subsection*{Quiz - Lesson GitBash.4}\label{docs__pandoc__latex__src__cli_complete_curriculum.md__quiz---lesson-gitbash4}

\subsection*{3DMake in Practice: Aliases in .bashrc}\label{docs__pandoc__latex__src__cli_complete_curriculum.md__3dmake-in-practice-aliases-in-bashrc}

Bash aliases make your 3DMake workflow much faster:

\begin{lstlisting}[style=Alabaster, language=bash]
# Quick alias to jump to your 3D projects folder
alias 3dp='cd ~/Documents/3d_projects'

# Full build-slice-print pipeline in one word
alias 3dprint='3dm build slice print'

# Build and then get an AI description of the model
alias 3dinfo='3dm build && 3dm info'

# Make sure 3dm is in your PATH (add to .bashrc if needed)
echo 'export PATH="$PATH:$HOME/Downloads/3dmake"' >> ~/.bashrc

# Persist all three aliases in .bashrc
echo "alias 3dp='cd ~/Documents/3d_projects'" >> ~/.bashrc
echo "alias 3dprint='3dm build slice print'" >> ~/.bashrc
echo "alias 3dinfo='3dm build && 3dm info'" >> ~/.bashrc

# Reload .bashrc immediately without restarting
source ~/.bashrc

# Now use them:
3dp         # jump to your 3D projects folder
3dprint     # run the full pipeline
3dinfo      # build and describe the model

\end{lstlisting}

\subsection*{Checkpoints}\label{docs__pandoc__latex__src__cli_complete_curriculum.md__checkpoints-17}

\begin{itemize}
\tightlist
\item
  After step 3 you can determine whether a program will be found by \texttt{PATH}.
\end{itemize}

\subsection*{Quiz - Lesson GitBash.4}\label{docs__pandoc__latex__src__cli_complete_curriculum.md__quiz---lesson-gitbash4-1}

\begin{enumerate}
\tightlist
\item
  How do you print an environment variable?
\item
  What is the purpose of \texttt{PATH}?
\item
  How do you check whether \texttt{openscad} is available?
\item
  How do you create a temporary alias?
\item
  Where would you make an alias permanent?
\item
  True or False: Environment variables in bash are case-sensitive.
\item
  Explain why having a program in your PATH is useful compared to always using its full file path.
\item
  Write a command that would create an alias called \texttt{slicer} for the OpenSCAD executable.
\item
  What file would you edit to make an alias persist across Git Bash sessions?
\item
  Describe a practical benefit of using the \texttt{\$TMPDIR} or \texttt{/tmp} directory for temporary files in a 3D printing workflow.
\item
  You have a custom script at \texttt{\textasciitilde{}/scripts/backup\_models.sh} that you want to run from anywhere as \texttt{backup-now}. What steps would you take to make this work?
\item
  Explain the difference between setting an environment variable in the current session with \texttt{export} vs. adding it to \texttt{.bashrc} for permanence.
\item
  Design a \texttt{.bashrc} strategy for managing multiple 3D printing projects, each with different tool paths and directories; show how to structure environment variables for each.
\item
  If a program is not found by \texttt{which}, what are the possible reasons, and how would you troubleshoot?
\item
  Describe how you would verify that your \texttt{.bashrc} is loading correctly and how to debug issues if aliases or environment variables don\textquotesingle t appear after restarting Git Bash.
\end{enumerate}

\subsection*{Extension Problems}\label{docs__pandoc__latex__src__cli_complete_curriculum.md__extension-problems-19}

\begin{enumerate}
\tightlist
\item
  Add a folder to PATH for a test program (describe steps; do not change system PATH without admin).
\item
  Create a short \texttt{.bashrc} snippet that sets two aliases and test re-opening Git Bash.
\item
  Use \texttt{which} to list the path for several common programs.
\item
  Explore \texttt{\$TMPDIR} or \texttt{/tmp} and create a file there.
\item
  Save a copy of your current PATH to a text file and examine it in your editor.
\item
  Create a \texttt{.bashrc} script that loads custom aliases and environment variables for your 3D printing workflow; test it in a new session.
\item
  Build a "project profile" that sets environment variables for CAD, slicing, and print directories; switch between profiles for different projects.
\item
  Write a script that audits your current environment variables and creates a summary report of what\textquotesingle s set and why.
\item
  Design a custom alias system for common 3D printing commands; document the aliases and their purposes.
\item
  Create a profile migration guide: document how to export and import your \texttt{.bashrc} across machines for consistent workflows.
\item
  Implement a safe PATH modification script: create a utility that allows you to add/remove directories from PATH for the current session only; show how to make it permanent in \texttt{.bashrc}.
\item
  Build a comprehensive \texttt{.bashrc} framework with modular sourcing: create separate \texttt{.sh} files for aliases, environment variables, and functions; have your main \texttt{.bashrc} source all of them.
\item
  Develop an environment validation tool: write a bash script that checks whether all required programs (OpenSCAD, slicers, etc.) are accessible via PATH; report findings and suggest fixes.
\item
  Create a project-switching alias system: design a function that changes all environment variables and aliases based on the current project; test switching between multiple projects.
\item
  Build a \texttt{.bashrc} troubleshooting guide: create a script that exports your current environment state (variables, aliases, PATH) to a timestamped file, allowing you to compare states before and after changes.
\end{enumerate}

\subsection*{References}\label{docs__pandoc__latex__src__cli_complete_curriculum.md__references-17}

\begin{itemize}
\tightlist
\item
  GNU. (2024). \emph{Environment variables in Bash}. \url{https://example.com}
\item
  GNU. (2024). \emph{alias command reference}. \url{https://example.com}
\item
  GNU. (2024). \emph{Creating and using Bash profiles}. \url{https://example.com}
\end{itemize}

\subsection*{Helpful Resources}\label{docs__pandoc__latex__src__cli_complete_curriculum.md__helpful-resources-15}

\begin{itemize}
\tightlist
\item
  \href{https://www.gnu.org/software/bash/manual/htmlnode/Environment.html}{Environment Variables in Bash}
\item
  \href{https://linuxize.com/post/how-to-add-directory-to-path-in-linux/}{Understanding the PATH Variable}
\item
  \href{https://www.gnu.org/software/bash/manual/htmlnode/Aliases.html}{Bash alias Reference}
\item
  \href{https://www.gnu.org/software/bash/manual/htmlnode/Bash-Startup-Files.html}{Creating a Bash Profile (.bashrc)}
\item
  \href{https://linux.die.net/man/1/which}{which Command for Locating Programs}
\end{itemize}

\begin{center}\rule{0.5\linewidth}{0.5pt}\end{center}

\protect\phantomsection\label{docs__pandoc__latex__src__cli_complete_curriculum.md__gb-lesson-5}{}

\section*{Lesson 5: Filling in the Gaps --- Shell Profiles, History, and Useful Tricks}\label{docs__pandoc__latex__src__cli_complete_curriculum.md__lesson-5-filling-in-the-gaps--shell-profiles-history-and-useful-tricks}

\begin{center}\rule{0.5\linewidth}{0.5pt}\end{center}

Estimated time: 30-45 minutes

\subsection*{Learning Objectives}\label{docs__pandoc__latex__src__cli_complete_curriculum.md__learning-objectives-13}

\begin{itemize}
\tightlist
\item
  Use history and abort commands (\texttt{history}, \texttt{Ctrl+R}, \texttt{Ctrl+C})
\item
  Inspect and edit your \texttt{.bashrc} profile for persistent settings
\item
  Run programs by full path using \texttt{./} or absolute paths
\item
  Handle common screen reader edge cases when using the terminal
\end{itemize}

\subsection*{Materials}\label{docs__pandoc__latex__src__cli_complete_curriculum.md__materials-8}

\begin{itemize}
\tightlist
\item
  Git Bash and an editor (Notepad/VS Code)
\end{itemize}

\subsection*{Step-by-step Tasks}\label{docs__pandoc__latex__src__cli_complete_curriculum.md__step-by-step-tasks-10}

\begin{enumerate}
\tightlist
\item
  Run several simple commands (e.g., \texttt{pwd}, \texttt{ls\ -1}, \texttt{echo\ hi}) then run \texttt{history} to view them.
\item
  Use \texttt{!\textless{}n\textgreater{}} to re-run a previous command by its history number (replace \texttt{\textless{}n\textgreater{}} with a history number).
\item
  Practice aborting a long-running command with \texttt{Ctrl\ +\ C} (for example, \texttt{ping\ google.com}).
\item
  Open your profile: \texttt{notepad.exe\ \textasciitilde{}/.bashrc}; if it doesn\textquotesingle t exist, create it: \texttt{touch\ \textasciitilde{}/.bashrc}.
\item
  Add a persistent alias line to your profile (example: \texttt{alias\ preview=\textquotesingle{}openscad\textquotesingle{}}), save, and run \texttt{source\ \textasciitilde{}/.bashrc} or reopen Git Bash to verify.
\end{enumerate}

\subsection*{Checkpoints}\label{docs__pandoc__latex__src__cli_complete_curriculum.md__checkpoints-18}

\begin{itemize}
\tightlist
\item
  After step 2 you can re-run a recent command by history number.
\item
  After step 5 your alias should persist across sessions.
\end{itemize}

\subsection*{Quiz - Lesson GitBash.5}\label{docs__pandoc__latex__src__cli_complete_curriculum.md__quiz---lesson-gitbash5}

\begin{enumerate}
\tightlist
\item
  How do you view the command history in Git Bash?
\item
  Which key combination aborts a running command?
\item
  What does \texttt{echo\ \$BASH\_VERSION} show?
\item
  How does \texttt{./} help run scripts and executables in the current directory?
\item
  What is one strategy if terminal output stops being announced by your screen reader?
\item
  True or False: Using \texttt{Ctrl+C} permanently deletes any files created by the command you abort.
\item
  Explain the difference between \texttt{history} and \texttt{Ctrl+R} (reverse history search) in Git Bash.
\item
  If you add an alias to \texttt{.bashrc} but it doesn\textquotesingle t take effect after opening a new Git Bash window, what should you verify?
\item
  Write a command that would run a script at the path \texttt{\textasciitilde{}/scripts/openscad\_runner.sh} directly.
\item
  Describe a practical workflow scenario where having keyboard shortcuts (aliases) in your \texttt{.bashrc} would save time.
\item
  Explain how to re-run the 5th command from your history using \texttt{!5}, and what would happen if that command had file operations (creates/deletes).
\item
  Design a \texttt{.bashrc} initialization strategy that separates utilities for different projects; explain how you would switch between them.
\item
  Walk through a troubleshooting workflow: your screen reader stops announcing output after running a long command. What steps would you take to diagnose and resolve the issue?
\item
  Create a safety checkpoint system: before any destructive operation (mass delete, overwrite), how would you use \texttt{.bashrc} functions and history to verify the command is correct?
\item
  Develop a comprehensive capstone scenario: integrate everything from Bash Lesson 0 through Bash Lesson 5 (navigation, file operations, piping, environment setup, history) to design an automated 3D printing project workflow with error handling and logging.
\end{enumerate}

\subsection*{Extension Problems}\label{docs__pandoc__latex__src__cli_complete_curriculum.md__extension-problems-20}

\begin{enumerate}
\tightlist
\item
  Add an alias and an environment variable change to your \texttt{.bashrc} and document the behavior after reopening Git Bash.
\item
  Create a short bash script that automates creating a project folder and an initial \texttt{.scad} file.
\item
  Experiment with running OpenSCAD by full path using \texttt{./} and by placing it in PATH; compare results.
\item
  Practice redirecting \texttt{man\ ls} output to a file and reading it in Notepad for screen reader clarity.
\item
  Document three screen reader troubleshooting steps you used and when they helped.
\item
  Build a comprehensive \texttt{.bashrc} that includes aliases, environment variables, and helper functions for your 3D printing workflow.
\item
  Create a bash script that troubleshoots common Git Bash issues (module loading, permission errors, command not found); test at least three scenarios.
\item
  Write a bash function that coordinates multiple tasks: creates a project folder, starts OpenSCAD, and opens a notes file.
\item
  Design a screen-reader accessibility guide for Git Bash: document commands, outputs, and accessible navigation patterns.
\item
  Develop an advanced Git Bash workflow: implement error handling, logging, and confirmation prompts for risky operations.
\item
  Implement an "undo" system using history: create a function that logs destructive commands (\texttt{rm}, \texttt{mv}, \texttt{cp}) and allows you to review the last operation.
\item
  Build a \texttt{.bashrc} debugger: create a script that compares two Git Bash sessions\textquotesingle{} environment states (variables, aliases, functions) to identify what loaded/failed to load.
\item
  Develop a multi-project profile manager: design a system where you can switch entire environments (paths, aliases, variables) for different 3D printing projects by running a single command.
\item
  Create a comprehensive accessibility analyzer: write a bash script that tests whether key Git Bash commands produce screen-reader-friendly output; document workarounds for commands that don\textquotesingle t.
\item
  Design a complete capstone project: build an integrated automation suite that manages a 3D printing workflow (project setup, file organization, CAD/slicing tool automation, output logging, error recovery, and audit trails) with full error handling and documentation.
\end{enumerate}

\subsection*{References}\label{docs__pandoc__latex__src__cli_complete_curriculum.md__references-18}

\begin{itemize}
\tightlist
\item
  GNU. (2024). \emph{Bash history and recall functionality}. \url{https://example.com}
\item
  GNU. (2024). \emph{Understanding and creating Bash profiles}. \url{https://example.com}
\item
  GNU. (2024). \emph{Running scripts and executables in Bash}. \url{https://example.com}
\end{itemize}

\subsection*{Helpful Resources}\label{docs__pandoc__latex__src__cli_complete_curriculum.md__helpful-resources-16}

\begin{itemize}
\tightlist
\item
  \href{https://www.gnu.org/software/bash/manual/htmlnode/Bash-History-Facilities.html}{Bash History and Recall}
\item
  \href{https://www.gnu.org/software/bash/manual/htmlnode/Bash-Startup-Files.html}{Understanding .bashrc}
\item
  \href{https://www.gnu.org/software/bash/manual/htmlnode/Bash-History-Builtins.html}{history Command Reference}
\item
  \href{https://www.gnu.org/software/bash/manual/htmlnode/Executing-Scripts.html}{Running Scripts with ./}
\item
  \href{https://www.nvaccess.org/documentation/}{Screen Reader Tips and Tricks}
\end{itemize}

\begin{center}\rule{0.5\linewidth}{0.5pt}\end{center}

\protect\phantomsection\label{docs__pandoc__latex__src__cli_complete_curriculum.md__gb-lesson-6}{}

\section*{Lesson 6: Advanced Techniques --- Shell Scripts, Functions, and Professional Workflows}\label{docs__pandoc__latex__src__cli_complete_curriculum.md__lesson-6-advanced-techniques--shell-scripts-functions-and-professional-workflows}

\begin{center}\rule{0.5\linewidth}{0.5pt}\end{center}

\textbf{Duration:} 4-4.5 hours (for screen reader users)
\textbf{Prerequisites:} Complete Bash Lesson 0 through Bash Lesson 5
\textbf{Skill Level:} Advanced intermediate

\subsection*{3DMake in Practice: Automating Your Print Workflow}\label{docs__pandoc__latex__src__cli_complete_curriculum.md__3dmake-in-practice-automating-your-print-workflow-2}

Use a shell script to fully automate a 3DMake workflow across multiple projects:

\begin{lstlisting}[style=Alabaster, language=bash]
#!/bin/bash
# batch_print.sh - build and print every 3DMake project in a folder

PROJECTS_DIR="$HOME/Documents/3d_projects"

for project in "$PROJECTS_DIR"/*/; do
    name=$(basename "$project")
    echo "Processing: $name"
    cd "$project"

    if 3dm build slice > build_log.txt 2>&1; then
        echo "  OK - sending to printer"
        3dm print
    else
        echo "  FAILED - check build_log.txt"
        cat build_log.txt
    fi
done

echo "Batch complete."

\end{lstlisting}

Tip

Save as \texttt{batch\_print.sh}, run \texttt{chmod\ +x\ batch\_print.sh} once to make it executable, then run with \texttt{./batch\_print.sh}. If you forget \texttt{chmod\ +x} you\textquotesingle ll get "Permission denied" --- just run the \texttt{chmod} command and try again.

Save as \texttt{batch\_print.sh}, make executable with \texttt{chmod\ +x\ batch\_print.sh}, then run with \texttt{./batch\_print.sh}. This script uses loops, conditionals, and redirection --- all from this lesson.

This lesson extends Git Bash skills to professional-level workflows. You\textquotesingle ll learn to automate complex tasks, write reusable shell scripts, and integrate tools for 3D printing workflows.

\begin{center}\rule{0.5\linewidth}{0.5pt}\end{center}

\subsection*{Learning Objectives}\label{docs__pandoc__latex__src__cli_complete_curriculum.md__learning-objectives-14}

By the end of this lesson, you will be able to:

\begin{itemize}
\tightlist
\item
  Create and run shell scripts (.sh files)
\item
  Write functions that accept parameters
\item
  Use loops to repeat tasks automatically
\item
  Automate batch processing of 3D models
\item
  Debug scripts when something goes wrong
\item
  Create professional workflows combining multiple tools
\end{itemize}

\begin{center}\rule{0.5\linewidth}{0.5pt}\end{center}

\subsection*{Shell Script Basics}\label{docs__pandoc__latex__src__cli_complete_curriculum.md__shell-script-basics}

\subsubsection*{What\textquotesingle s a Shell Script?}\label{docs__pandoc__latex__src__cli_complete_curriculum.md__whats-a-shell-script}

A shell script (\texttt{.sh}) contains multiple bash commands that run in sequence. Instead of typing commands one by one, you put them in a file and run them all at once.

\textbf{Why use shell scripts?}

\begin{itemize}
\tightlist
\item
  Repeatability: Run the same task 100 times identically
\item
  Documentation: Commands are written down for reference
\item
  Complexity: Combine many commands logically
\item
  Automation: Schedule scripts to run automatically
\end{itemize}

\subsubsection*{Creating Your First Shell Script}\label{docs__pandoc__latex__src__cli_complete_curriculum.md__creating-your-first-shell-script}

\textbf{Step 1: Open a text editor}

\begin{lstlisting}[style=Alabaster, language=bash]
notepad.exe my-first-script.sh

\end{lstlisting}

\textbf{Step 2: Type this script}

\begin{lstlisting}[style=Alabaster, language=bash]
#!/bin/bash
# This is a comment - screen readers will read it
echo "Script is running!"
pwd
ls -1
echo "Script is done!"

\end{lstlisting}

\textbf{Step 3: Save the file}

\begin{itemize}
\tightlist
\item
  In Notepad: Ctrl+S
\item
  Make sure filename ends in \texttt{.sh}
\item
  Save in an easy-to-find location (like Documents)
\end{itemize}

Note

Before you can run a \texttt{.sh} script, you must make it executable with \texttt{chmod\ +x\ scriptname.sh}. This is a one-time step per script. Without it, running \texttt{./script.sh} will produce a "Permission denied" error even though the file exists.

\textbf{Step 4: Make it executable and run the script}

\begin{lstlisting}[style=Alabaster, language=bash]
chmod +x my-first-script.sh
./my-first-script.sh

\end{lstlisting}

\textbf{What happens:}
Bash runs each command in sequence and shows output.

Important

Always start shell scripts with \texttt{\#!/bin/bash} (the "shebang" line). Without it, the system may use a different shell to run your script, which can cause unexpected errors --- especially on systems where \texttt{/bin/sh} is not bash.

\subsubsection*{Important: The Shebang Line}\label{docs__pandoc__latex__src__cli_complete_curriculum.md__important-the-shebang-line}

The \texttt{\#!/bin/bash} at the top of your script (called a "shebang") tells the system which shell to use to run the script.

\begin{lstlisting}[style=Alabaster, language=bash]
#!/bin/bash
# Now bash runs every command below
echo "Hello!"

\end{lstlisting}

\begin{center}\rule{0.5\linewidth}{0.5pt}\end{center}

\subsection*{Variables and Parameters}\label{docs__pandoc__latex__src__cli_complete_curriculum.md__variables-and-parameters-2}

\subsubsection*{Using Variables}\label{docs__pandoc__latex__src__cli_complete_curriculum.md__using-variables-2}

Variables store values you want to use later.

\textbf{Example script:}

\begin{lstlisting}[style=Alabaster, language=bash]
#!/bin/bash
mypath="$HOME/Documents"
cd "$mypath"
echo "I am now in:"
pwd
ls -1

\end{lstlisting}

\textbf{Breaking it down:}

\begin{itemize}
\tightlist
\item
  \texttt{mypath="..."} assigns the variable (no spaces around \texttt{=})
\end{itemize}

Tip

Always quote variables that contain file paths: use \texttt{"\$mypath"} not \texttt{\$mypath}. Without quotes, spaces in folder names will break your script. This is one of the most common bash scripting mistakes.

\begin{itemize}
\tightlist
\item
  \texttt{"\$mypath"} uses the variable (quote it to handle spaces in paths)
\item
  Variables in bash are always referenced with \texttt{\$}
\end{itemize}

\subsubsection*{Functions with Parameters}\label{docs__pandoc__latex__src__cli_complete_curriculum.md__functions-with-parameters-1}

A function is reusable code that you can run with different inputs.

\textbf{Example: A function that lists files in a folder}

\begin{lstlisting}[style=Alabaster, language=bash]
#!/bin/bash
list_folder() {
    local path="$1"
    echo "Contents of: $path"
    cd "$path"
    ls -1
}

# Use the function:
list_folder "$HOME/Documents"
list_folder "$HOME/Downloads"

\end{lstlisting}

\textbf{What\textquotesingle s happening:}

\begin{itemize}
\tightlist
\item
  \texttt{list\_folder()} defines the function
\item
  \texttt{local\ path="\$1"} assigns the first argument to a local variable
\item
  \texttt{\$1} is the first argument passed to the function
\item
  Call the function with \texttt{list\_folder\ "path/to/folder"}
\end{itemize}

\textbf{Screen reader tip:} When you call a function, bash will announce the results just like any command.

\begin{center}\rule{0.5\linewidth}{0.5pt}\end{center}

\subsection*{Loops - Repeating Tasks}\label{docs__pandoc__latex__src__cli_complete_curriculum.md__loops---repeating-tasks-2}

\subsubsection*{Loop Over Files}\label{docs__pandoc__latex__src__cli_complete_curriculum.md__loop-over-files-2}

Imagine you have 10 SCAD files and want to print their contents. You could do it 10 times manually, or use a loop.

\textbf{Example: Print every .scad file in a folder}

\begin{lstlisting}[style=Alabaster, language=bash]
#!/bin/bash
for file in *.scad; do
    echo "=== File: $file ==="
    cat "$file"
    echo ""
done

\end{lstlisting}

\textbf{What\textquotesingle s happening:}

\begin{itemize}
\tightlist
\item
  \texttt{for\ file\ in\ *.scad;\ do} iterates over each \texttt{.scad} file
\item
  \texttt{\$file} holds the current filename
\item
  \texttt{done} ends the loop
\item
  Inside the loop, do something with each \texttt{\$file}
\end{itemize}

\subsubsection*{Loop with a Counter}\label{docs__pandoc__latex__src__cli_complete_curriculum.md__loop-with-a-counter-2}

\textbf{Example: Do something 5 times}

\begin{lstlisting}[style=Alabaster, language=bash]
#!/bin/bash
for i in $(seq 1 5); do
    echo "This is iteration number $i"
    # Do something here
done

\end{lstlisting}

\textbf{What\textquotesingle s happening:}

\begin{itemize}
\tightlist
\item
  \texttt{for\ i\ in\ \$(seq\ 1\ 5)} loops with \texttt{i} from 1 to 5
\item
  \texttt{\$i} is the counter variable
\end{itemize}

\begin{center}\rule{0.5\linewidth}{0.5pt}\end{center}

\subsection*{Real-World Example - Batch Processing SCAD Files}\label{docs__pandoc__latex__src__cli_complete_curriculum.md__real-world-example---batch-processing-scad-files-2}

\subsubsection*{Scenario}\label{docs__pandoc__latex__src__cli_complete_curriculum.md__scenario-2}

You have 10 OpenSCAD (.scad) files in a folder. You want to:

\begin{enumerate}
\tightlist
\item
  List them all
\item
  Check how many there are
\item
  For each one, verify it exists
\end{enumerate}

\subsubsection*{The Script}\label{docs__pandoc__latex__src__cli_complete_curriculum.md__the-script-2}

\begin{lstlisting}[style=Alabaster, language=bash]
#!/bin/bash
scad_folder="$HOME/Documents/3D_Projects"
count=0

echo "Processing SCAD files in: $scad_folder"
echo ""

for file in "$scad_folder"/*.scad; do
    if [ -f "$file" ]; then
        echo "  Found: $(basename "$file")"
        count=$((count + 1))
    else
        echo "  Missing: $(basename "$file")"
    fi
done

echo ""
echo "Total files found: $count"
echo "Batch processing complete!"

\end{lstlisting}

\textbf{Breaking it down:}

\begin{itemize}
\tightlist
\item
  \texttt{scad\_folder="..."} = where to look
\item
  \texttt{for\ file\ in\ "\$scad\_folder"/*.scad} = find all .scad files
\item
  \texttt{if\ {[}\ -f\ "\$file"\ {]}} = check if file exists and is a regular file
\item
  \texttt{basename\ "\$file"} = just the filename (not the full path)
\item
  \texttt{count=\$((count\ +\ 1))} = increment the counter
\end{itemize}

\subsubsection*{Running the Script}\label{docs__pandoc__latex__src__cli_complete_curriculum.md__running-the-script-2}

\begin{enumerate}
\tightlist
\item
  Save as \texttt{batch-process.sh}
\item
  Edit \texttt{scad\_folder} to match your real folder
\item
  Make it executable and run it:

  \begin{lstlisting}[style=Alabaster, language=bash]
  chmod +x batch-process.sh
  ./batch-process.sh

  \end{lstlisting}
\end{enumerate}

\textbf{Screen reader output:}

\begin{lstlisting}[style=Alabaster]
Processing SCAD files in: /c/Users/YourName/Documents/3D_Projects

  Found: model1.scad
  Found: model2.scad
  Found: model3.scad
[... more files ...]

Total files found: 10
Batch processing complete!

\end{lstlisting}

\begin{center}\rule{0.5\linewidth}{0.5pt}\end{center}

\subsection*{Error Handling}\label{docs__pandoc__latex__src__cli_complete_curriculum.md__error-handling-2}

\subsubsection*{Try-Style Checks with Exit Codes}\label{docs__pandoc__latex__src__cli_complete_curriculum.md__try-style-checks-with-exit-codes}

What if something goes wrong? Use exit code checks:

\textbf{Example:}

\begin{lstlisting}[style=Alabaster, language=bash]
#!/bin/bash
file="$HOME/nonexistent/path/file.txt"

if cat "$file" 2>/dev/null; then
    echo "File read successfully"
else
    echo "Error: Could not read file"
    echo "File path was: $file"
fi

\end{lstlisting}

\textbf{What\textquotesingle s happening:}

\begin{itemize}
\tightlist
\item
  \texttt{2\textgreater{}/dev/null} suppresses error messages from \texttt{cat}
\item
  \texttt{if\ cat\ ...;\ then} checks whether the command succeeded
\item
  \texttt{else} handles the failure gracefully
\end{itemize}

\textbf{Screen reader advantage:} Errors are announced clearly instead of crashing silently.

\subsubsection*{Validating Input}\label{docs__pandoc__latex__src__cli_complete_curriculum.md__validating-input-2}

\textbf{Example: Make sure a folder exists before processing}

\begin{lstlisting}[style=Alabaster, language=bash]
#!/bin/bash
process_folder() {
    local folder_path="$1"

    if [ ! -d "$folder_path" ]; then
        echo "Error: Folder does not exist: $folder_path"
        return 1
    fi

    echo "Processing folder: $folder_path"
    ls -1 "$folder_path"
}

process_folder "$HOME/Documents"

\end{lstlisting}

\textbf{What\textquotesingle s happening:}

\begin{itemize}
\tightlist
\item
  \texttt{{[}\ !\ -d\ "\$folder\_path"\ {]}} checks if the path is NOT a directory
\item
  \texttt{return\ 1} exits the function early with a non-zero (error) status
\end{itemize}

\begin{center}\rule{0.5\linewidth}{0.5pt}\end{center}

\subsection*{Debugging Shell Scripts}\label{docs__pandoc__latex__src__cli_complete_curriculum.md__debugging-shell-scripts}

\subsubsection*{Common Errors and Solutions}\label{docs__pandoc__latex__src__cli_complete_curriculum.md__common-errors-and-solutions-2}

\paragraph*{Error 1: "Command not found"}\label{docs__pandoc__latex__src__cli_complete_curriculum.md__error-1-command-not-found-2}

\textbf{Cause:} Typo in command name

\textbf{Fix:} Check spelling

\begin{lstlisting}[style=Alabaster, language=bash]
# Wrong:
ech "hello"

# Correct:
echo "hello"

\end{lstlisting}

\paragraph*{Error 2: "Variable is empty"}\label{docs__pandoc__latex__src__cli_complete_curriculum.md__error-2-variable-is-empty-1}

\textbf{Cause:} Variable was never assigned or has a typo

\textbf{Fix:} Make sure variable is set before using it

\begin{lstlisting}[style=Alabaster, language=bash]
myvar="hello"  # Set first
echo "$myvar"  # Then use

\end{lstlisting}

\paragraph*{Error 3: "No such file or directory"}\label{docs__pandoc__latex__src__cli_complete_curriculum.md__error-3-no-such-file-or-directory}

\textbf{Cause:} Wrong folder path

\textbf{Fix:} Verify path exists

\begin{lstlisting}[style=Alabaster, language=bash]
# Check if path exists:
if [ -d "$HOME/Documents" ]; then
    echo "Path exists"
fi

\end{lstlisting}

\paragraph*{Error 4: "Permission denied"}\label{docs__pandoc__latex__src__cli_complete_curriculum.md__error-4-permission-denied}

\textbf{Cause:} Script not executable, or no write permission

\textbf{Fix:} Make the script executable, or check file permissions

\begin{lstlisting}[style=Alabaster, language=bash]
chmod +x my-script.sh

\end{lstlisting}

\subsubsection*{\texorpdfstring{Debugging Technique: Trace Output with \texttt{set\ -x}}{Debugging Technique: Trace Output with set -x}}\label{docs__pandoc__latex__src__cli_complete_curriculum.md__debugging-technique-trace-output-with-set--x}

Add \texttt{set\ -x} at the top of your script to print each command before it runs:

\begin{lstlisting}[style=Alabaster, language=bash]
#!/bin/bash
set -x  # Enable trace mode
path_var="$HOME/Documents"
echo "Starting script. Path is: $path_var"

for file in "$path_var"/*; do
    echo "Processing: $file"
    # Do something with $file
    echo "Done with: $file"
done

echo "Script complete"

\end{lstlisting}

Your screen reader will announce each step, so you know where errors happen.

\begin{center}\rule{0.5\linewidth}{0.5pt}\end{center}

\subsection*{Creating Professional Workflows}\label{docs__pandoc__latex__src__cli_complete_curriculum.md__creating-professional-workflows-2}

\subsubsection*{Example 1: Automated Project Setup}\label{docs__pandoc__latex__src__cli_complete_curriculum.md__example-1-automated-project-setup-2}

\textbf{Scenario:} You start a new 3D printing project regularly. Instead of creating folders manually:

\begin{lstlisting}[style=Alabaster, language=bash]
#!/bin/bash
read -p "Enter project name: " project_name
base_folder="$HOME/Documents/3D_Projects"
project_folder="$base_folder/$project_name"

# Create folder structure
mkdir -p "$project_folder"
mkdir -p "$project_folder/designs"
mkdir -p "$project_folder/output"
mkdir -p "$project_folder/notes"

# Create a README
cat > "$project_folder/README.txt" << EOF
# $project_name

Created: $(date)

## Designs
All .scad files go here.

## Output
STL and other exports go here.

## Notes
Project notes and observations.
EOF

echo "Project setup complete: $project_folder"

\end{lstlisting}

\textbf{What it does:}

\begin{itemize}
\tightlist
\item
  Prompts for a project name
\item
  Creates folder structure for a new project
\item
  Sets up subfolders for designs, output, notes
\item
  Creates a README file automatically
\end{itemize}

\subsubsection*{Example 2: Batch File Verification}\label{docs__pandoc__latex__src__cli_complete_curriculum.md__example-2-batch-file-verification-2}

\textbf{Scenario:} Before processing, verify all required files exist:

\begin{lstlisting}[style=Alabaster, language=bash]
#!/bin/bash
verify_project() {
    local project_folder="$1"
    local required_items=("README.txt" "designs" "output" "notes")
    local all_good=true

    for item in "${required_items[@]}"; do
        local path="$project_folder/$item"
        if [ -e "$path" ]; then
            echo "  Found: $item"
        else
            echo "  Missing: $item"
            all_good=false
        fi
    done

    if $all_good; then
        echo "All checks passed!"
        return 0
    else
        echo "Some files are missing!"
        return 1
    fi
}

# Use it:
project="$HOME/Documents/3D_Projects/MyKeychain"
if verify_project "$project"; then
    echo "Safe to proceed with processing"
fi

\end{lstlisting}

\begin{center}\rule{0.5\linewidth}{0.5pt}\end{center}

\subsection*{Screen Reader Tips for Shell Scripts}\label{docs__pandoc__latex__src__cli_complete_curriculum.md__screen-reader-tips-for-shell-scripts}

\subsubsection*{Making Script Output Readable}\label{docs__pandoc__latex__src__cli_complete_curriculum.md__making-script-output-readable-2}

\textbf{Problem:} Script runs but output scrolls too fast or is hard to follow

\textbf{Solution 1: Save to file}

\begin{lstlisting}[style=Alabaster, language=bash]
./my-script.sh > output.txt
notepad.exe output.txt

\end{lstlisting}

\textbf{Solution 2: Use echo with clear sections}

\begin{lstlisting}[style=Alabaster, language=bash]
echo "========== STARTING =========="
echo ""
# ... script ...
echo ""
echo "========== COMPLETE =========="

\end{lstlisting}

\textbf{Solution 3: Pause between major sections}

\begin{lstlisting}[style=Alabaster, language=bash]
echo "Pausing... Press Enter to continue"
read

\end{lstlisting}

Your screen reader will announce the pause, give you time to read output.

\subsubsection*{Announcing Progress}\label{docs__pandoc__latex__src__cli_complete_curriculum.md__announcing-progress-2}

\textbf{For long-running scripts:}

\begin{lstlisting}[style=Alabaster, language=bash]
#!/bin/bash
count=0
total=$(ls *.scad | wc -l)

for file in *.scad; do
    count=$((count + 1))
    echo "Processing $count of $total: $file"
    # Do something with $file
done

echo "All $count files processed!"

\end{lstlisting}

\begin{center}\rule{0.5\linewidth}{0.5pt}\end{center}

\subsection*{Practice Exercises}\label{docs__pandoc__latex__src__cli_complete_curriculum.md__practice-exercises-7}

\subsubsection*{Exercise 1: Your First Shell Script}\label{docs__pandoc__latex__src__cli_complete_curriculum.md__exercise-1-your-first-shell-script}

\textbf{Goal:} Create and run a simple shell script

\textbf{Steps:}

\begin{enumerate}
\tightlist
\item
  Create file: \texttt{notepad.exe\ hello.sh}
\item
  Type:

  \begin{lstlisting}[style=Alabaster, language=bash]
  #!/bin/bash
  echo "Hello from my first Git Bash shell script!"
  pwd
  ls -1

  \end{lstlisting}
\item
  Save, make executable, and run:

  \begin{lstlisting}[style=Alabaster, language=bash]
  chmod +x hello.sh
  ./hello.sh

  \end{lstlisting}
\end{enumerate}

\textbf{Checkpoint:} You should see output for each command.

\subsubsection*{Exercise 2: Script with a Variable}\label{docs__pandoc__latex__src__cli_complete_curriculum.md__exercise-2-script-with-a-variable-2}

\textbf{Goal:} Use a variable to make the script flexible

\textbf{Steps:}

\begin{enumerate}
\tightlist
\item
  Create file: \texttt{notepad.exe\ smart-listing.sh}
\item
  Type:

  \begin{lstlisting}[style=Alabaster, language=bash]
  #!/bin/bash
  target_folder="$HOME/Documents"
  echo "Listing contents of: $target_folder"
  ls -1 "$target_folder"

  \end{lstlisting}
\item
  Edit \texttt{target\_folder} to a real folder on your computer
\item
  Run:

  \begin{lstlisting}[style=Alabaster, language=bash]
  chmod +x smart-listing.sh
  ./smart-listing.sh

  \end{lstlisting}
\end{enumerate}

\textbf{Checkpoint:} You should see listing of that specific folder.

\subsubsection*{Exercise 3: Function}\label{docs__pandoc__latex__src__cli_complete_curriculum.md__exercise-3-function-1}

\textbf{Goal:} Create a reusable function

\textbf{Steps:}

\begin{enumerate}
\tightlist
\item
  Create file: \texttt{notepad.exe\ navigate.sh}
\item
  Type:

  \begin{lstlisting}[style=Alabaster, language=bash]
  #!/bin/bash
  go_to() {
      local path="$1"
      if [ -d "$path" ]; then
          cd "$path"
          echo "Now in: $(pwd)"
          echo "Contents:"
          ls -1
      else
          echo "Path does not exist: $path"
      fi
  }

  # Test the function:
  go_to "$HOME/Documents"
  go_to "$HOME/Downloads"

  \end{lstlisting}
\item
  Run:

  \begin{lstlisting}[style=Alabaster, language=bash]
  chmod +x navigate.sh
  ./navigate.sh

  \end{lstlisting}
\end{enumerate}

\textbf{Checkpoint:} Both function calls should work, showing contents of each folder.

\subsubsection*{Exercise 4: Loop}\label{docs__pandoc__latex__src__cli_complete_curriculum.md__exercise-4-loop-2}

\textbf{Goal:} Use a loop to repeat an action

\textbf{Steps:}

\begin{enumerate}
\tightlist
\item
  Create file: \texttt{notepad.exe\ repeat.sh}
\item
  Type:

  \begin{lstlisting}[style=Alabaster, language=bash]
  #!/bin/bash
  echo "Demonstrating a loop:"

  for i in $(seq 1 5); do
      echo "Iteration $i: Hello!"
  done

  echo "Loop complete!"

  \end{lstlisting}
\item
  Run:

  \begin{lstlisting}[style=Alabaster, language=bash]
  chmod +x repeat.sh
  ./repeat.sh

  \end{lstlisting}
\end{enumerate}

\textbf{Checkpoint:} Should print "Iteration 1" through "Iteration 5".

\subsubsection*{Exercise 5: Real-World Script}\label{docs__pandoc__latex__src__cli_complete_curriculum.md__exercise-5-real-world-script-2}

\textbf{Goal:} Create a useful script for a real task

\textbf{Steps:}

\begin{enumerate}
\tightlist
\item
  Create a folder: \texttt{mkdir\ \textasciitilde{}/Documents/TestFiles}
\item
  Create some test files:

  \begin{lstlisting}[style=Alabaster, language=bash]
  echo "test" > ~/Documents/TestFiles/file1.txt
  echo "test" > ~/Documents/TestFiles/file2.txt
  echo "test" > ~/Documents/TestFiles/file3.txt

  \end{lstlisting}
\item
  Create script: \texttt{notepad.exe\ report.sh}
\item
  Type:

  \begin{lstlisting}[style=Alabaster, language=bash]
  #!/bin/bash
  folder="$HOME/Documents/TestFiles"
  count=0

  echo "=== FILE REPORT ==="
  echo "Folder: $folder"
  echo ""
  echo "Files:"
  for file in "$folder"/*; do
      echo "  - $(basename "$file")"
      count=$((count + 1))
  done
  echo ""
  echo "Total: $count files"
  echo "=== END REPORT ==="

  \end{lstlisting}
\item
  Run:

  \begin{lstlisting}[style=Alabaster, language=bash]
  chmod +x report.sh
  ./report.sh

  \end{lstlisting}
\end{enumerate}

\textbf{Checkpoint:} Should show report of all files in the test folder.

\begin{center}\rule{0.5\linewidth}{0.5pt}\end{center}

\subsection*{Quiz - Lesson GitBash.6}\label{docs__pandoc__latex__src__cli_complete_curriculum.md__quiz---lesson-gitbash6}

\begin{enumerate}
\tightlist
\item
  What is a shell script?
\item
  What file extension do bash shell scripts use?
\item
  What is a variable in bash and how do you create one?
\item
  What is a function and why would you use one?
\item
  How do you run a shell script?
\item
  What is a \texttt{for} loop and what does \texttt{for\ file\ in\ *.scad;\ do} do?
\item
  What does \texttt{{[}\ -f\ "\$file"\ {]}} check?
\item
  How do you handle errors in a bash script?
\item
  When would you use \texttt{if\ {[}\ !\ -d\ "\$path"\ {]};\ then}?
\item
  What technique makes shell script output readable for screen readers?
\end{enumerate}

\begin{center}\rule{0.5\linewidth}{0.5pt}\end{center}

\subsection*{Extension Problems}\label{docs__pandoc__latex__src__cli_complete_curriculum.md__extension-problems-21}

\begin{enumerate}
\tightlist
\item
  \textbf{Auto-Backup Script:} Create a bash script that copies all files from one folder to another, announcing progress
\item
  \textbf{File Counter:} Write a function that counts files by extension (.txt, .scad, .stl, etc.)
\item
  \textbf{Folder Cleaner:} Script that deletes files older than 30 days (with user confirmation)
\item
  \textbf{Project Template:} Function that creates a complete project folder structure with all needed files
\item
  \textbf{Batch Rename:} Script that renames all files in a folder according to a pattern
\item
  \textbf{Log Generator:} Create a script that records what it does to a log file for later review
\item
  \textbf{Scheduled Task:} Set up a script to run automatically using cron or Task Scheduler
\item
  \textbf{File Verifier:} Check that all SCAD files in a folder have corresponding STL exports
\item
  \textbf{Report Generator:} Create a summary report of all projects in a folder
\item
  \textbf{Error Tracker:} Script that lists all commands that had errors and logs them with timestamps
\end{enumerate}

\begin{center}\rule{0.5\linewidth}{0.5pt}\end{center}

\subsection*{Important Notes}\label{docs__pandoc__latex__src__cli_complete_curriculum.md__important-notes-2}

\begin{itemize}
\tightlist
\item
  \textbf{Always test scripts on small sets of files first} before running them on important data
\item
  \textbf{Save your work regularly} --- use version naming if possible
\item
  \textbf{Test error handling} --- make sure errors don\textquotesingle t crash silently
\item
  \textbf{Document your scripts} --- use \texttt{\#} comments so you remember what each part does
\item
  \textbf{Backup before batch operations} --- if something goes wrong, you have the original
\end{itemize}

\begin{center}\rule{0.5\linewidth}{0.5pt}\end{center}

\subsection*{References}\label{docs__pandoc__latex__src__cli_complete_curriculum.md__references-19}

\begin{itemize}
\tightlist
\item
  \textbf{GNU Bash Scripting Guide:} \url{https://example.com}
\item
  \textbf{Function Documentation:} \url{https://example.com}
\item
  \textbf{Error Handling:} \url{https://example.com}
\item
  \textbf{Loops:} \url{https://example.com}
\end{itemize}

\begin{center}\rule{0.5\linewidth}{0.5pt}\end{center}

\protect\phantomsection\label{docs__pandoc__latex__src__cli_complete_curriculum.md__gb-unit-exam}{}

\section*{Unit Exam --- Comprehensive Assessment}\label{docs__pandoc__latex__src__cli_complete_curriculum.md__unit-exam--comprehensive-assessment-2}

\begin{center}\rule{0.5\linewidth}{0.5pt}\end{center}

Estimated time: 60-90 minutes

\subsection*{Key Learning Outcomes Assessed}\label{docs__pandoc__latex__src__cli_complete_curriculum.md__key-learning-outcomes-assessed-2}

By completing this unit test, you will demonstrate:

\begin{enumerate}
\tightlist
\item
  Understanding of file system navigation and path concepts
\item
  Proficiency with file and folder manipulation commands
\item
  Ability to redirect and pipe command output
\item
  Knowledge of environment variables and aliases
\item
  Screen-reader accessibility best practices in terminal environments
\item
  Problem-solving and command chaining skills
\end{enumerate}

\subsection*{Target Audience}\label{docs__pandoc__latex__src__cli_complete_curriculum.md__target-audience-2}

Users who have completed Bash Lesson 0 through Bash Lesson 6 and need to demonstrate mastery of Git Bash fundamentals.

\subsection*{Instructions}\label{docs__pandoc__latex__src__cli_complete_curriculum.md__instructions-2}

Complete all sections below. For multiple choice, select the best answer. For short answers, write one to two sentences. For hands-on tasks, capture evidence (screenshots or output files) and submit alongside your answers.

\begin{center}\rule{0.5\linewidth}{0.5pt}\end{center}

\subsection*{Part A: Multiple Choice Questions (20 questions)}\label{docs__pandoc__latex__src__cli_complete_curriculum.md__part-a-multiple-choice-questions-20-questions-2}

Select the best answer for each question. Each question is worth 1 point.

\begin{enumerate}
\item
  What is the primary purpose of the \texttt{PATH} environment variable?

  \begin{itemize}
  \tightlist
  \item
    A) Store your home directory location
  \item
    B) Tell the shell where to find executable programs
  \item
    C) Configure the visual appearance of the terminal
  \item
    D) Store the current working directory name
  \end{itemize}
\item
  Which command prints your current working directory in Git Bash?

  \begin{itemize}
  \tightlist
  \item
    A) \texttt{ls\ -1}
  \item
    B) \texttt{cd}
  \item
    C) \texttt{pwd}
  \item
    D) \texttt{whoami}
  \end{itemize}
\item
  What does the \texttt{\textasciitilde{}} symbol represent in Git Bash paths?

  \begin{itemize}
  \tightlist
  \item
    A) The root directory
  \item
    B) The current directory
  \item
    C) The parent directory
  \item
    D) The home directory
  \end{itemize}
\item
  How do you list only file names (not full details) in a way that is screen-reader friendly?

  \begin{itemize}
  \tightlist
  \item
    A) \texttt{ls}
  \item
    B) \texttt{ls\ -1}
  \item
    C) \texttt{ls\ -l}
  \item
    D) \texttt{cat\ -1}
  \end{itemize}
\item
  Which command creates a new empty file in Git Bash?

  \begin{itemize}
  \tightlist
  \item
    A) \texttt{mkdir\ filename}
  \item
    B) \texttt{touch\ filename}
  \item
    C) \texttt{new\ filename}
  \item
    D) \texttt{echo\ filename}
  \end{itemize}
\item
  What is the difference between \texttt{\textgreater{}} and \texttt{\textgreater{}\textgreater{}}?

  \begin{itemize}
  \tightlist
  \item
    A) \texttt{\textgreater{}} redirects to file, \texttt{\textgreater{}\textgreater{}} displays on screen
  \item
    B) \texttt{\textgreater{}} overwrites a file, \texttt{\textgreater{}\textgreater{}} appends to a file
  \item
    C) They do the same thing
  \item
    D) \texttt{\textgreater{}} is for text, \texttt{\textgreater{}\textgreater{}} is for binary
  \end{itemize}
\item
  What does the pipe operator \texttt{\textbar{}} do?

  \begin{itemize}
  \tightlist
  \item
    A) Creates a folder
  \item
    B) Sends the output of one command to the input of another
  \item
    C) Deletes files matching a pattern
  \item
    D) Lists all processes
  \end{itemize}
\item
  Which command copies a file in Git Bash?

  \begin{itemize}
  \tightlist
  \item
    A) \texttt{mv}
  \item
    B) \texttt{rm}
  \item
    C) \texttt{cp}
  \item
    D) \texttt{cd}
  \end{itemize}
\item
  How do you rename a file from \texttt{oldname.txt} to \texttt{newname.txt} in Git Bash?

  \begin{itemize}
  \tightlist
  \item
    A) \texttt{cp\ oldname.txt\ newname.txt}
  \item
    B) \texttt{mv\ oldname.txt\ newname.txt}
  \item
    C) \texttt{rename\ oldname.txt\ newname.txt}
  \item
    D) \texttt{rn\ oldname.txt\ newname.txt}
  \end{itemize}
\item
  What is the purpose of \texttt{grep} in Git Bash piping?

  \begin{itemize}
  \tightlist
  \item
    A) Find files in a directory
  \item
    B) Search for text patterns within output or files
  \item
    C) Select a string to copy to clipboard
  \item
    D) Select which shell to use
  \end{itemize}
\item
  Which key allows you to autocomplete a path in Git Bash?

  \begin{itemize}
  \tightlist
  \item
    A) \texttt{Ctrl\ +\ A}
  \item
    B) \texttt{Ctrl\ +\ E}
  \item
    C) \texttt{Tab}
  \item
    D) \texttt{Space}
  \end{itemize}
\item
  How do you copy text to the Windows clipboard from Git Bash?

  \begin{itemize}
  \tightlist
  \item
    A) \texttt{cat\ filename\ \textgreater{}\ clipboard}
  \item
    B) \texttt{cat\ filename\ \textbar{}\ clip}
  \item
    C) \texttt{copy\ filename}
  \item
    D) \texttt{cat\ filename\ \textbar{}\ paste}
  \end{itemize}
\end{enumerate}

Tip

Use \texttt{which\ programname} to check whether a program is in your PATH and find its full path. If \texttt{which\ openscad} returns nothing, OpenSCAD is not in your PATH and you\textquotesingle ll need to add it or call it by its full path.

\begin{enumerate}
\setcounter{enumi}{12}
\item
  What does \texttt{which\ openscad} do?

  \begin{itemize}
  \tightlist
  \item
    A) Opens the OpenSCAD application
  \item
    B) Gets help about the openscad command
  \item
    C) Locates the full path of the openscad executable
  \item
    D) Lists all available commands
  \end{itemize}
\item
  Which wildcard matches any single character in Git Bash?

  \begin{itemize}
  \tightlist
  \item
    A) \texttt{*}
  \item
    B) \texttt{?}
  \item
    C) \texttt{\%}
  \item
    D) \texttt{\#}
  \end{itemize}
\item
  What is the purpose of \texttt{./} before a script name?

  \begin{itemize}
  \tightlist
  \item
    A) Run a script in the current directory
  \item
    B) Execute all commands in parallel
  \item
    C) Combine multiple commands
  \item
    D) Create an alias
  \end{itemize}
\item
  How do you create a temporary alias in Git Bash?

  \begin{itemize}
  \tightlist
  \item
    A) \texttt{set-alias\ preview=\textquotesingle{}openscad\textquotesingle{}}
  \item
    B) \texttt{alias\ preview=\textquotesingle{}openscad\textquotesingle{}}
  \item
    C) \texttt{new-alias\ preview\ openscad}
  \item
    D) \texttt{preview\ =\ openscad}
  \end{itemize}
\end{enumerate}

Important

Aliases saved only in the terminal session are lost when you close Git Bash. To make an alias permanent, add it to \texttt{\textasciitilde{}/.bashrc}. After editing \texttt{.bashrc}, run \texttt{source\ \textasciitilde{}/.bashrc} to apply changes without restarting.

\begin{enumerate}
\setcounter{enumi}{16}
\tightlist
\item
  Where is a Git Bash alias typically stored to persist across sessions?

  \begin{itemize}
  \tightlist
  \item
    A) \texttt{C:\textbackslash{}Program\ Files\textbackslash{}Git\textbackslash{}profile.sh}
  \item
    B) In the \texttt{\textasciitilde{}/.bashrc} file
  \item
    C) \texttt{\textasciitilde{}/bash\_profile}
  \item
    D) Aliases cannot be made permanent
  \end{itemize}
\end{enumerate}

Important

\textbf{Ctrl+C} is the universal "abort" key --- it stops any running command immediately and returns you to the prompt. This works in Git Bash, CMD, and PowerShell. Use it whenever a command hangs, loops unexpectedly, or produces unwanted output.

\begin{enumerate}
\setcounter{enumi}{17}
\item
  How do you abort a long-running command in Git Bash?

  \begin{itemize}
  \tightlist
  \item
    A) Press \texttt{Escape}
  \item
    B) Press \texttt{Ctrl\ +\ X}
  \item
    C) Press \texttt{Ctrl\ +\ C}
  \item
    D) Press \texttt{Alt\ +\ F4}
  \end{itemize}
\item
  What command shows the history of previously run commands in Git Bash?

  \begin{itemize}
  \tightlist
  \item
    A) \texttt{history}
  \item
    B) \texttt{get-history}
  \item
    C) \texttt{show-history}
  \item
    D) \texttt{bash-history}
  \end{itemize}
\end{enumerate}

Tip

After editing \texttt{\textasciitilde{}/.bashrc}, run \texttt{source\ \textasciitilde{}/.bashrc} to apply your changes immediately without restarting Git Bash. Without this step, new aliases and variables won\textquotesingle t take effect until the next time you open a terminal.

\begin{enumerate}
\setcounter{enumi}{19}
\tightlist
\item
  How do you reload your \texttt{.bashrc} without restarting Git Bash?

  \begin{itemize}
  \tightlist
  \item
    A) Use \texttt{reload\ \textasciitilde{}/.bashrc} in the terminal
  \item
    B) Use \texttt{source\ \textasciitilde{}/.bashrc}
  \item
    C) Use the Windows Control Panel
  \item
    D) \texttt{.bashrc} reloads automatically
  \end{itemize}
\end{enumerate}

\begin{center}\rule{0.5\linewidth}{0.5pt}\end{center}

\subsection*{Part B: Short Answer Questions (10 questions)}\label{docs__pandoc__latex__src__cli_complete_curriculum.md__part-b-short-answer-questions-10-questions-2}

Answer each question in one to two sentences. Each question is worth 2 points.

\begin{enumerate}
\item
  Explain the difference between absolute and relative paths. Give one example of each.
\item
  Why is \texttt{ls\ -1} preferred over \texttt{ls} for screen reader users? Describe what flag combination you would use to list only directories.
\item
  What is the purpose of redirecting output to a file, and give an example of when you would use \texttt{\textgreater{}} instead of \texttt{\textgreater{}\textgreater{}}?
\item
  Describe what would happen if you ran \texttt{rm\ -r\ \textasciitilde{}/Documents/my\_folder} and why this command should be used carefully.
\item
  How would you search for all files with a \texttt{.scad} extension in your current directory? Write the command.
\item
  Explain what happens when you pipe the output of \texttt{ls\ -1} into \texttt{clip}. What would you do next?
\item
  What is an environment variable, and give one example of how you might use it in Git Bash.
\item
  If a program is not in your \texttt{PATH}, what two methods could you use to run it from Git Bash?
\item
  Describe how you would open a file in Notepad and also add a line to it from Git Bash.
\item
  What is one strategy you would use if your screen reader stops announcing terminal output while using Git Bash?
\end{enumerate}

\begin{center}\rule{0.5\linewidth}{0.5pt}\end{center}

\subsection*{Part C: Hands-On Tasks (10 tasks)}\label{docs__pandoc__latex__src__cli_complete_curriculum.md__part-c-hands-on-tasks-10-tasks-2}

Complete each task and capture evidence (screenshots, output files, or command transcripts). Each task is worth 3 points.

\subsubsection*{Tasks 1-5: File System and Navigation}\label{docs__pandoc__latex__src__cli_complete_curriculum.md__tasks-1-5-file-system-and-navigation-2}

\begin{enumerate}
\item
  Create a folder structure \texttt{\textasciitilde{}/Documents/GitBash\_Assessment/Projects} using a single command. Capture the \texttt{ls\ -1} output showing the creation.
\item
  Create five files named \texttt{project\_1.scad}, \texttt{project\_2.scad}, \texttt{project\_3.txt}, \texttt{notes\_1.txt}, and \texttt{notes\_2.txt} inside the \texttt{Projects} folder. Use wildcards to list only \texttt{.scad} files, then capture the output.
\item
  Copy the entire \texttt{Projects} folder to \texttt{Projects\_Backup} using \texttt{cp\ -r}. Capture the \texttt{ls\ -1} output showing both folders exist.
\item
  Move (rename) \texttt{project\_1.scad} to \texttt{project\_1\_final.scad}. Capture the \texttt{ls\ -1} output showing the renamed file.
\item
  Delete \texttt{notes\_1.txt} and \texttt{notes\_2.txt} using a single \texttt{rm} command with wildcards. Capture the final \texttt{ls\ -1} output.
\end{enumerate}

\subsubsection*{Tasks 6-10: Advanced Operations and Scripting}\label{docs__pandoc__latex__src__cli_complete_curriculum.md__tasks-6-10-advanced-operations-and-scripting-2}

\begin{enumerate}
\setcounter{enumi}{5}
\item
  Create a file called \texttt{my\_data.txt} with at least four lines using \texttt{echo} and \texttt{\textgreater{}\textgreater{}}. Then read it with \texttt{cat\ my\_data.txt} and capture the output.
\item
  Use \texttt{grep} to search for a keyword (e.g., "project") in \texttt{my\_data.txt} and pipe the results to \texttt{clip}. Paste the results into Notepad and capture a screenshot.
\item
  List all files in the \texttt{Projects} folder and redirect the output to \texttt{projects\_list.txt}. Open it in Notepad and capture a screenshot of the file.
\item
  Create a temporary alias called \texttt{myls} that runs \texttt{ls\ -1}, test it, and capture the output. Then explain what would be required to make it persistent.
\item
  Run \texttt{man\ ls} (or \texttt{ls\ -\/-help}) and redirect the output to \texttt{help\_output.txt}. Open the file in Notepad and capture a screenshot showing at least the first page of help content.
\end{enumerate}

\begin{center}\rule{0.5\linewidth}{0.5pt}\end{center}

\subsection*{Grading Rubric}\label{docs__pandoc__latex__src__cli_complete_curriculum.md__grading-rubric-2}

{\def\LTcaptype{none} % do not increment counter
\begin{longtable}[]{@{}llll@{}}
\toprule\noalign{}
Section & Questions & Points Each & Total \\
\midrule\noalign{}
\endhead
\bottomrule\noalign{}
\endlastfoot
Multiple Choice & 20 & 1 & 20 \\
Short Answer & 10 & 2 & 20 \\
Hands-On Tasks & 10 & 3 & 30 \\
\textbf{Total} & \textbf{40} & - & \textbf{70} \\
\end{longtable}
}

\textbf{Passing Score:} 49 points (70\%)

\begin{center}\rule{0.5\linewidth}{0.5pt}\end{center}

\subsection*{Helpful Resources for Review}\label{docs__pandoc__latex__src__cli_complete_curriculum.md__helpful-resources-for-review-2}

\begin{itemize}
\tightlist
\item
  \href{https://git-scm.com/book/en/v2}{Git Bash Command Reference}
\item
  \href{https://linuxcommand.org/lc3lts0020.php}{Navigation and File System}
\item
  \href{https://www.gnu.org/software/grep/manual/grep.html}{Using Pipes and Filtering}
\item
  \href{https://www.gnu.org/software/bash/manual/htmlnode/Bash-Startup-Files.html}{Bash Profile and Aliases}
\item
  \href{https://www.nvaccess.org/documentation/}{Screen Reader Accessibility Tips}
\end{itemize}

\begin{center}\rule{0.5\linewidth}{0.5pt}\end{center}

\subsection*{Submission Checklist}\label{docs__pandoc__latex__src__cli_complete_curriculum.md__submission-checklist-2}

\begin{itemize}
\tightlist
\item[$\square$]
  All 20 multiple choice questions answered
\item[$\square$]
  All 10 short answer questions answered (1-2 sentences each)
\item[$\square$]
  All 10 hands-on tasks completed with evidence captured
\item[$\square$]
  Files/screenshots organized and labeled clearly
\item[$\square$]
  Submission includes this checklist
\end{itemize}

\newpage

\chapter{3dMake-Accessible Design with OpenSCAD}\label{docs__pandoc__latex__src__3dmake_foundation__part_2.md__3dmake-accessible-design-with-openscad}

This section covers the complete 3D design and printing workflow using OpenSCAD, 3DMake, and accessible tools. Students progress from basic primitives through advanced parametric design and stakeholder-centric projects.

Time commitment: 30-40 hours instruction + projects

Skills gained: 3D modeling, parametric design, automated workflows, tolerance testing, quality assurance

\section{3dMake Foundation \& Design --- Complete Course Compendium}\label{docs__pandoc__latex__src__3dmake_all_lessons.md__3dmake-foundation--design--complete-course-compendium}

\begin{quote}
\textbf{Course:} Accessible 3D Printing with OpenSCAD and 3DMake
\textbf{Track:} 3dMake Certification --- High School Distance Learning
\textbf{Contents:} Instructional Framework · Curriculum Guide · All 11 Lessons · Final Exam · Quick Reference Guide
\end{quote}

\begin{center}\rule{0.5\linewidth}{0.5pt}\end{center}

\section*{Part I --- Instructional Framework}\label{docs__pandoc__latex__src__3dmake_all_lessons.md__part-i--instructional-framework}

\emph{A pedagogical overview of the 3DMake and OpenSCAD ecosystem for secondary STEM education.}

\subsection*{The Architecture of Programmatic Design}\label{docs__pandoc__latex__src__3dmake_all_lessons.md__the-architecture-of-programmatic-design}

Programmatic CAD differs from standard industry tools like Fusion 360 or SolidWorks by utilizing a "code-as-model" philosophy. OpenSCAD, the primary engine supported by 3DMake, uses a functional programming language to define three-dimensional volumes.\footnote{Programming with OpenSCAD. https://programmingwithopenscad.github.io/} This approach is particularly robust for parametric design, where the dimensions of an object are defined as variables, allowing for instantaneous reconfiguration without manual rebuilding.\footnote{CadHub. \emph{OpenSCAD Review}. https://learn.cadhub.xyz/blog/openscad-review/} 3DMake enhances this by providing a unified command-line interface (CLI) to manage the entire process, from editing source files to triggering remote print jobs through interfaces like OctoPrint or Bambu Connect.\footnote{Deck, T. (2025). \emph{3DMake}. GitHub. https://github.com/tdeck/3dmake}

The educational value of this toolchain lies in its transparency. Students are not hidden behind a proprietary user interface; they interact directly with the file system, configuration files, and API integrations. This exposure fosters a deeper understanding of how modern software interacts with hardware --- a skill set increasingly vital in robotics and aerospace engineering.\footnote{Class Central. \emph{OpenSCAD Courses}. https://www.classcentral.com/subject/openscad} However, the shift to a terminal-based environment requires a structured instructional approach to overcome initial barriers to entry.

\subsection*{AI Integration and Model Descriptive Feedback}\label{docs__pandoc__latex__src__3dmake_all_lessons.md__ai-integration-and-model-descriptive-feedback}

The integration of artificial intelligence into the 3DMake workflow represents a significant advancement in democratizing CAD for students with varying levels of spatial reasoning skills. The \texttt{3dm\ info} command acts as a multimodal bridge between the deterministic world of OpenSCAD and the probabilistic world of LLMs.\footnote{Deck, T. (2025). \emph{3DMake}. GitHub. https://github.com/tdeck/3dmake}

When a student executes \texttt{3dm\ info}, the system initiates a rendering pipeline, generating multiple viewpoints of the current model, which are packaged as image data and sent to the configured AI (e.g., Gemini) along with a descriptive prompt. The AI\textquotesingle s response provides a summary that can include the model\textquotesingle s intended function, aesthetic qualities, and potential engineering failures.\footnote{Deck, T. (2025). \emph{3DMake}. GitHub. https://github.com/tdeck/3dmake}

\subsubsection*{The Limits of AI Spatial Reasoning}\label{docs__pandoc__latex__src__3dmake_all_lessons.md__the-limits-of-ai-spatial-reasoning}

While powerful, current LLMs do not possess a true "3D world model" --- they are trained primarily on 2D images and text.\footnote{ZenML. \emph{LLM-Powered 3D Model Generation}. https://www.zenml.io/llmops-database/llm-powered-3d-model-generation-for-3d-printing} Common failures include:

\begin{itemize}
\tightlist
\item
  \textbf{Detached Geometry:} The AI might describe a "table" even if the legs are hovering below the tabletop.\footnote{ZenML. \emph{LLM-Powered 3D Model Generation}. https://www.zenml.io/llmops-database/llm-powered-3d-model-generation-for-3d-printing}
\item
  \textbf{Scale Misinterpretation:} Without a reference object in the render, the AI may misjudge scale, leading to inappropriate feedback on wall thickness.\footnote{Deck, T. (2025). \emph{3DMake}. GitHub. https://github.com/tdeck/3dmake}
\end{itemize}

Caution

\textbf{Hallucination of Detail:} The AI may describe features it expects to see (like "engraved text") even if the student\textquotesingle s code failed to render them.\footnote{ZenML. \emph{LLM-Powered 3D Model Generation}. https://www.zenml.io/llmops-database/llm-powered-3d-model-generation-for-3d-printing}

Important

\textbf{AI is a verification assistant, not a source of truth.} The deterministic rendering of OpenSCAD remains the final authority on geometry. AI is most useful as a "sanity check" to catch obvious mistakes before wasting filament.\footnote{Deck, T. (2025). \emph{3DMake}. GitHub. https://github.com/tdeck/3dmake}

\subsection*{Occupational Health and Safety}\label{docs__pandoc__latex__src__3dmake_all_lessons.md__occupational-health-and-safety}

Safety in the 3DMake workflow extends to the physical act of printing, which involves thermal, chemical, and mechanical risks. The National Institute for Occupational Safety and Health (NIOSH) provides a clear framework for these risks.\footnote{Ohio State University EHS. https://ehs.osu.edu/kb/3d-printer-safety}

\subsubsection*{Chemical and Particulate Emissions}\label{docs__pandoc__latex__src__3dmake_all_lessons.md__chemical-and-particulate-emissions}

The melting of plastic filament is a thermal degradation process. ABS releases styrene, a known respiratory irritant and potential carcinogen.\footnote{Ohio State University EHS. https://ehs.osu.edu/kb/3d-printer-safety} Even PLA emits millions of ultrafine particles (UFPs) per minute during extrusion --- particles smaller than 100 nanometers that can penetrate deep into the lungs.\footnote{Ohio State University EHS. https://ehs.osu.edu/kb/3d-printer-safety}

{\def\LTcaptype{none} % do not increment counter
\begin{longtable}[]{@{}lll@{}}
\toprule\noalign{}
Emission Component & Primary Source Filaments & Mitigation Strategy \\
\midrule\noalign{}
\endhead
\bottomrule\noalign{}
\endlastfoot
\end{longtable}
}

Warning

ABS and ASA filaments release styrene --- a respiratory irritant and potential carcinogen. Always use an enclosed printer with carbon filtration for these materials.

\textbar{} Styrene \textbar{} ABS, ASA \textbar{} Enclosed printer with carbon filtration.\footnote{Ohio State University EHS. https://ehs.osu.edu/kb/3d-printer-safety} \textbar{}
\textbar{} Formaldehyde \textbar{} POM, Nylon \textbar{} High-efficiency external ventilation.\footnote{Ohio State University EHS. https://ehs.osu.edu/kb/3d-printer-safety} \textbar{}
\textbar{} Ultrafine Particles \textbar{} All filaments \textbar{} HEPA filtration and "20-minute" settling period.\footnote{CDC/NIOSH. https://www.cdc.gov/niosh/blogs/2024/safe-3d-printing.html} \textbar{}
\textbar{} Volatile Organic Compounds \textbar{} All filaments \textbar{} Minimum 6 air changes per hour in the room.\footnote{Washington State DoH. https://doh.wa.gov/community-and-environment/schools/3d-printers} \textbar{}

\subsubsection*{Physical and Mechanical Hazards}\label{docs__pandoc__latex__src__3dmake_all_lessons.md__physical-and-mechanical-hazards}

The extruder nozzle can reach 260°C and the heated build plate can reach 110°C.\footnote{RIT EHS. https://www.rit.edu/ehs/3-d-printer-safety} Standard operating procedures for a student lab must include:

\begin{itemize}
\tightlist
\item
  \textbf{Pre-use Inspection:} Check for frayed wires, loose belts, and a clear build surface.\footnote{Washington State DoH. https://doh.wa.gov/community-and-environment/schools/3d-printers}
\item
  \textbf{Environmental Controls:} Prohibit eating or drinking in the print area to avoid ingestion of contaminants.\footnote{RIT EHS. https://www.rit.edu/ehs/3-d-printer-safety}
\item
  \textbf{Emergency Response:} Locate the nearest ABC fire extinguisher (Class D for metal prints).\footnote{UTK EHS. https://ehs.utk.edu/index.php/table-of-policies-plans-procedures-guides/3d-printer-safety/}
\item
  \textbf{Post-Print Cooling:} Ensure the printer has cooled below 30°C before removing the model.\footnote{Washington State DoH. https://doh.wa.gov/community-and-environment/schools/3d-printers}
\end{itemize}

\subsection*{Challenges Inherent in the OpenSCAD Language}\label{docs__pandoc__latex__src__3dmake_all_lessons.md__challenges-inherent-in-the-openscad-language}

OpenSCAD\textquotesingle s declarative nature and unique rendering kernel present specific hurdles for students. Unlike imperative languages, OpenSCAD describes \emph{what} a shape is, not \emph{how} to build it step by step.\footnote{Gonzalez Avila et al. https://thomaspietrzak.com/bibliography/gonzalez24.pdf}

\subsubsection*{Immutable State}\label{docs__pandoc__latex__src__3dmake_all_lessons.md__immutable-state}

Note

Variables in OpenSCAD are more like constants within a scope. If a student writes \texttt{x\ =\ 5;\ x\ =\ 10;}, OpenSCAD will use the last assigned value (\texttt{x\ =\ 10}) for the entire script.\footnote{Reddit. https://www.reddit.com/r/openscad/comments/1fxj8xv/} This requires adopting a functional programming mindset where geometry is defined by its state rather than a sequence of movements.

\subsubsection*{Performance Bottlenecks}\label{docs__pandoc__latex__src__3dmake_all_lessons.md__performance-bottlenecks}

Note

OpenSCAD\textquotesingle s CGAL geometry kernel is highly accurate but slow for certain operations. \textbf{Minkowski sums} (used for rounding corners) can take minutes at high \texttt{\$fn} values. Keep \texttt{\$fn} low while developing and only raise it for final renders.\footnote{OpenSCAD Issue \#6297. https://github.com/openscad/openscad/issues/6297} OpenSCAD also has no native "fillet" command --- students must construct these manually using boolean subtractions or libraries like BOSL2.\footnote{CadHub. \emph{OpenSCAD Review}. https://learn.cadhub.xyz/blog/openscad-review/}

\subsubsection*{The Absolute Coordinate Barrier}\label{docs__pandoc__latex__src__3dmake_all_lessons.md__the-absolute-coordinate-barrier}

OpenSCAD has no concept of relative constraints --- every object is positioned in absolute (X, Y, Z) space.\footnote{Hacker News. https://news.ycombinator.com/item?id=46338565} If a student moves one part, they must manually calculate and update the translation of every related part. This necessitates heavy use of variables and mathematical offsets.\footnote{Hacker News. https://news.ycombinator.com/item?id=46338565}

\subsection*{Technical Limitations of 3DMake}\label{docs__pandoc__latex__src__3dmake_all_lessons.md__technical-limitations-of-3dmake}

{\def\LTcaptype{none} % do not increment counter
\begin{longtable}[]{@{}
  >{\raggedright\arraybackslash}p{(\linewidth - 4\tabcolsep) * \real{0.3333}}
  >{\raggedright\arraybackslash}p{(\linewidth - 4\tabcolsep) * \real{0.3333}}
  >{\raggedright\arraybackslash}p{(\linewidth - 4\tabcolsep) * \real{0.3333}}@{}}
\toprule\noalign{}
\begin{minipage}[b]{\linewidth}\raggedright
Limitation
\end{minipage} & \begin{minipage}[b]{\linewidth}\raggedright
Technical Root Cause
\end{minipage} & \begin{minipage}[b]{\linewidth}\raggedright
Educational Impact
\end{minipage} \\
\midrule\noalign{}
\endhead
\bottomrule\noalign{}
\endlastfoot
CLI Barrier & No graphical interface for configuration &
Steep learning curve for students with zero terminal experience.\footnote{Deck, T. (2025). \emph{3DMake}. GitHub. https://github.com/tdeck/3dmake} \\
\end{longtable}
}

Caution

AI features (\texttt{3dm\ info} descriptions, \texttt{3dm\ orient}) require an active internet connection and a configured API key. These features will fail silently on offline school networks --- plan accordingly.

\textbar{} API Dependency \textbar{} AI features require external internet and API keys \textbar{} Advanced features fail in offline school networks.\footnote{Deck, T. (2025). \emph{3DMake}. GitHub. https://github.com/tdeck/3dmake} \textbar{}
\textbar{} Slicer Lock-in \textbar{} Reliant on external templates for G-code generation \textbar{} Students may not learn the nuance of slicing settings.\footnote{Deck, T. (2025). \emph{3DMake}. GitHub. https://github.com/tdeck/3dmake} \textbar{}
\textbar{} Feedback Latency \textbar{} No real-time "live" preview in the editor \textbar{} The edit-compile-view cycle is slower than GUI-based CAD.\footnote{Deck, T. (2025). \emph{3DMake}. GitHub. https://github.com/tdeck/3dmake} \textbar{}

Caution

\textbf{Non-manifold geometry} (zero-thickness walls or improperly closed polyhedrons) will generate an STL without warning from \texttt{3dm\ build}, but the \texttt{slice} command may then fail or produce corrupt G-code. Always verify mesh integrity before printing.\footnote{OpenSCAD Issue \#6114. https://github.com/openscad/openscad/issues/6114} Students must learn mesh verification skills using tools like MeshLab or PrusaSlicer\textquotesingle s repair functions.\footnote{Tom\textquotesingle s Hardware. https://www.tomshardware.com/3d-printing/how-to-repair-stl-files-in-meshlab}

\subsection*{Local Resources --- Utah Makerspaces}\label{docs__pandoc__latex__src__3dmake_all_lessons.md__local-resources--utah-makerspaces}

Utah provides a broad network of makerspaces across public libraries, universities, and community nonprofits.

\subsubsection*{Public Library Makerspaces}\label{docs__pandoc__latex__src__3dmake_all_lessons.md__public-library-makerspaces}

\textbf{Salt Lake City Public Library -- Creative Lab} (Main + Marmalade, Glendale, Sprague branches): 3D printers, Cricut cutters, Adobe Creative Cloud, sound booth, and more.\footnote{https://services.slcpl.org/creativelab}

\textbf{Salt Lake County Library -- Create Spaces} (Daybreak, Granite, Holladay, Kearns, Magna, Sandy): Flashforge and LulzBot 3D printers, VR/robotics kits, audio/visual studios.\footnote{https://www.slcolibrary.org/what-we-have/create}

\textbf{Weber County Library Makerspaces:} Laser cutters, vinyl cutters, embroidery machines, and design software stations.\footnote{https://www.weberpl.lib.ut.us/use-your-library/makerspaces}

\textbf{Washington County Library -- Sid Atkin Makerspace (St. George):} Large-format 3D printers, CNC cutters, sewing machines, electronics kits, and woodworking tools.\footnote{https://library.washco.utah.gov/st-george/makerspace/}

\textbf{Logan Library -- Cytiva STEM Makerspace (Cache County):} Bambu 3D printers, laser engravers, heat presses, and digitization tools.\footnote{https://library.loganutah.org/about\_us/services/makerspace.php}

\subsubsection*{University \& Higher-Education Makerspaces}\label{docs__pandoc__latex__src__3dmake_all_lessons.md__university--higher-education-makerspaces}

\textbf{University of Utah -- Marriott Library ProtoSpace:} Prusa MK4 and Prusa XL printers, 3D scanners, open to students and staff.\footnote{https://www.lib.utah.edu/services/3d-printing.php}

\textbf{Utah Tech University -- Jr. Makerspace (St. George):} Supports K--12 and community visitors with training and STEM events.\footnote{https://stem.utah.gov/innovationhubnetwork/}

\textbf{Southern Utah University -- Thunderworks Makerspace (Cedar City):} CNC tools, 3D printers, electronics benches, and generous public hours.\footnote{https://stem.utah.gov/innovationhubnetwork/}

\subsubsection*{Regional STEM Hubs}\label{docs__pandoc__latex__src__3dmake_all_lessons.md__regional-stem-hubs}

\textbf{Make Salt Lake:} Nonprofit makerspace with woodshop, metal shop, laser cutters, 3D printer labs, and electronics benches.\footnote{https://stem.utah.gov/innovationhubnetwork/}

The \textbf{Utah STEM Action Center\textquotesingle s Innovation Hub Network} maintains a statewide directory across all counties: \href{https://stem.utah.gov/innovationhubnetwork/}{stem.utah.gov/innovationhubnetwork}.

\subsubsection*{National Makerspace Directories}\label{docs__pandoc__latex__src__3dmake_all_lessons.md__national-makerspace-directories}

\begin{itemize}
\tightlist
\item
  \href{https://www.nationofmakers.us/find-a-maker-organization}{Nation of Makers} --- U.S. crowdsourced map of makerspaces and hackerspaces
\item
  \href{https://makerspace.com/directory/}{Makerspace.com Global Directory} --- Worldwide listing with search and filters
\item
  \href{https://makerspacedir.com/}{Makerspace Directory (makerspacedir.com)} --- U.S.-focused, covers Fab Labs, youth programs, and coworking hybrids
\item
  \href{https://workmakecreate.com/}{WorkMakeCreate} --- Find free and public makerspaces by ZIP code
\end{itemize}

\subsection*{Supplemental Learning Resources}\label{docs__pandoc__latex__src__3dmake_all_lessons.md__supplemental-learning-resources}

\begin{itemize}
\tightlist
\item
  \href{docs\%5Cpandoc\%5Clatex\%5Csrc\%5C943ee2446365181c.epub}{Programming with OpenSCAD: A Beginner\textquotesingle s Guide to Coding 3D-Printable Objects} --- Complete reference covering syntax, geometry concepts, design patterns, and best practices
\item
  \href{docs\%5Cpandoc\%5Clatex\%5Csrc\%5Cee8ec7d1b3e2cc7c.epub}{Simplifying 3D Printing with OpenSCAD} --- Focused guide to practical workflows, optimization, and real-world printing
\item
  \href{https://github.com/ProgrammingWithOpenSCAD/CodeSolutions}{CodeSolutions Repository} --- Working OpenSCAD examples organized by topic and difficulty
\item
  \href{https://programmingwithopenscad.github.io/learning.html}{Practice Worksheets and Guides} --- Printable materials for visualization practice and assessment
\end{itemize}

\subsection*{References}\label{docs__pandoc__latex__src__3dmake_all_lessons.md__references}

Deck, T. (2025). \emph{3DMake: A command-line tool for 3D printing workflows}. GitHub. https://github.com/tdeck/3dmake
Gohde, J., \& Kintel, M. (2021). \emph{Programming with OpenSCAD: A beginner\textquotesingle s guide to coding 3D-printable objects}. No Starch Press.
Gonzalez Avila, J. F., Pietrzak, T., \& Casiez, G. (2024). \emph{Understanding the challenges of OpenSCAD users for 3D printing}. ACM UIST.
NIOSH. (2024). \emph{Approaches to safe 3D printing}. https://www.cdc.gov/niosh/blogs/2024/safe-3d-printing.html
Ohio State University EHS. (2026). \emph{3D printer safety concerns and ventilation}. https://ehs.osu.edu/kb/3d-printer-safety
Washington State Department of Health. (2026). \emph{3D printer and filament selection for safe school environments}. https://doh.wa.gov/community-and-environment/schools/3d-printers

\begin{itemize}
\tightlist
\item
  https://www.nationofmakers.us/find-a-maker-organization
\item
  https://makerspace.com/directory/
\item
  https://makerspacedir.com/
\item
  https://workmakecreate.com/
\end{itemize}

\begin{center}\rule{0.5\linewidth}{0.5pt}\end{center}

\section*{Part II --- Curriculum Guide}\label{docs__pandoc__latex__src__3dmake_all_lessons.md__part-ii--curriculum-guide}

\emph{Structure, learning paths, and reference material for the complete 11-lesson course.}

\subsection*{Overview}\label{docs__pandoc__latex__src__3dmake_all_lessons.md__overview}

This 11-lesson curriculum (plus 4 reference appendices) teaches non-visual 3D modeling, parametric design principles, and the complete 3D printing workflow through integrated projects. Each lesson builds on prior knowledge, introducing real-world examples and hands-on activities.

\textbf{Target Audience:} High school and early undergraduate students, makers, and anyone interested in programmatic CAD and 3D printing
\textbf{Estimated Total Time:} 30--40 hours (11 lessons + 4 appendices + projects)

\subsection*{Lesson Structure at a Glance}\label{docs__pandoc__latex__src__3dmake_all_lessons.md__lesson-structure-at-a-glance}

{\def\LTcaptype{none} % do not increment counter
\begin{longtable}[]{@{}
  >{\raggedright\arraybackslash}p{(\linewidth - 8\tabcolsep) * \real{0.0750}}
  >{\raggedright\arraybackslash}p{(\linewidth - 8\tabcolsep) * \real{0.1750}}
  >{\raggedright\arraybackslash}p{(\linewidth - 8\tabcolsep) * \real{0.2500}}
  >{\raggedright\arraybackslash}p{(\linewidth - 8\tabcolsep) * \real{0.1750}}
  >{\raggedright\arraybackslash}p{(\linewidth - 8\tabcolsep) * \real{0.3250}}@{}}
\toprule\noalign{}
\begin{minipage}[b]{\linewidth}\raggedright
\#
\end{minipage} & \begin{minipage}[b]{\linewidth}\raggedright
Title
\end{minipage} & \begin{minipage}[b]{\linewidth}\raggedright
Duration
\end{minipage} & \begin{minipage}[b]{\linewidth}\raggedright
Level
\end{minipage} & \begin{minipage}[b]{\linewidth}\raggedright
Key Project
\end{minipage} \\
\midrule\noalign{}
\endhead
\bottomrule\noalign{}
\endlastfoot
1 & Environmental Configuration & 2.5--3.5 hrs & Beginner &
Project scaffold \\
2 & Geometric Primitives \& CSG & 60 min & Beginner & 3 CSG examples \\
3 & Parametric Architecture & 2.5 hrs & Beginner+ &
Parametric bracket \\
4 & AI Verification & 2--2.5 hrs & Intermediate & AI-notes.md \\
5 & Safety \& Physical Interface & 2.5--3.5 hrs & Intermediate &
Pre-print checklist \\
6 & 3dm Commands \& Text Embossing & 2.5--3.5 hrs & Intermediate &
Keycap variants \\
7 & Parametric Transforms & 3--3.5 hrs & Intermediate+ & Phone Stand \\
8 & Advanced Parametric Design & 90--120 min & Advanced &
Stackable Bins \\
9 & Automation \& Workflows & 2.5--3.5 hrs & Advanced &
Batch build script \\
10 & Troubleshooting \& Mastery & 5--6 hrs & Advanced &
3 complete projects \\
11 & Stakeholder-Centric Design & 3.5--4.5 hrs & Advanced+ &
Jewelry Holder \\
\end{longtable}
}

\subsection*{Learning Paths}\label{docs__pandoc__latex__src__3dmake_all_lessons.md__learning-paths}

\textbf{Path 1 --- Complete Mastery (30--40 hours):} Lessons 1--11 + all appendices. Best for first-time learners wanting comprehensive understanding.

\textbf{Path 2 --- Design Focus (Fast Track):} Lessons 1 → 2 → 3 → 6 → 7 → 8 → 9. Best for experienced designers new to programmatic CAD.

\textbf{Path 3 --- Project-Based:} Lessons 1--3 → 6 (Keycap) → 7 (Phone Stand) → 8 (Bins) → 9 (Automation) → 10 (Troubleshooting) → 11 (Leadership). Best for learning by doing.

\textbf{Path 4 --- Safety \& Printing Focus:} Lessons 1, 2, 5, 6, 10 + Appendices A, B, C. Best for practical printing and quality focus.

\subsection*{Detailed Lesson Descriptions}\label{docs__pandoc__latex__src__3dmake_all_lessons.md__detailed-lesson-descriptions}

\subsubsection*{Lesson 1: Environmental Configuration and the Developer Workflow}\label{docs__pandoc__latex__src__3dmake_all_lessons.md__lesson-1-environmental-configuration-and-the-developer-workflow}

Install and configure 3dMake, create a project scaffold, and run your first build. Covers installation and tool verification, project structure (\texttt{src/}, \texttt{build/}, \texttt{3dmake.toml}), parametric design philosophy, and the \texttt{3dm\ build} command. \textbf{Checkpoint:} You can initialize a project, edit a parametric model, and generate an STL file.

\subsubsection*{Lesson 2: Geometric Primitives and Constructive Solid Geometry}\label{docs__pandoc__latex__src__3dmake_all_lessons.md__lesson-2-geometric-primitives-and-constructive-solid-geometry}

Combine basic shapes using boolean operations. Covers \texttt{cube()}, \texttt{sphere()}, \texttt{cylinder()}; CSG operations (\texttt{union()}, \texttt{difference()}, \texttt{intersection()}); the 0.001 offset rule for non-manifold geometry; and low-resolution renders for faster debugging. \textbf{Checkpoint:} You understand CSG operations and can diagnose common geometry issues.

\subsubsection*{Lesson 3: Parametric Architecture and Modular Libraries}\label{docs__pandoc__latex__src__3dmake_all_lessons.md__lesson-3-parametric-architecture-and-modular-libraries}

Create reusable modules and organize code into libraries using the DRY principle. Covers module definitions, derived calculations, library organization (\texttt{lib/} folder), BOSL2, and low-resolution testing with \texttt{\$fn}. \textbf{Checkpoint:} You can create parametric modules and organize code into libraries.

\subsubsection*{Lesson 4: AI-Enhanced Verification and Multimodal Feedback}\label{docs__pandoc__latex__src__3dmake_all_lessons.md__lesson-4-ai-enhanced-verification-and-multimodal-feedback}

Use \texttt{3dm\ info} to generate AI diagnostics and validate designs. Covers the \texttt{3dm\ info} command, AI-generated descriptions, AI limitations, privacy and governance, and prompt engineering basics. \textbf{Checkpoint:} You can use AI tools to supplement your design validation.

\subsubsection*{Lesson 5: Safety Protocols and the Physical Fabrication Interface}\label{docs__pandoc__latex__src__3dmake_all_lessons.md__lesson-5-safety-protocols-and-the-physical-fabrication-interface}

Transition from digital design to physical printing. Covers the Hierarchy of Controls, chemical and particulate emissions, pre-print environmental and equipment checks, post-print inspection, and spool metadata. \textbf{Checkpoint:} You understand safety procedures and can safely conduct supervised prints.

\subsubsection*{Lesson 6: Practical 3dm Commands and Text Embossing}\label{docs__pandoc__latex__src__3dmake_all_lessons.md__lesson-6-practical-3dm-commands-and-text-embossing}

Master key 3dm commands by building a customizable keycap. Covers \texttt{3dm\ info}, \texttt{3dm\ preview}, \texttt{3dm\ orient}, \texttt{3dm\ slice}, and text embossing with \texttt{linear\_extrude()} and \texttt{text()}. \textbf{Project:} Generate 3+ keycap variants (small, medium, large). \textbf{Checkpoint:} You can generate keycaps with embossed text and understand all major 3dm commands.

\subsubsection*{Lesson 7: Parametric Transforms and the Phone Stand Project}\label{docs__pandoc__latex__src__3dmake_all_lessons.md__lesson-7-parametric-transforms-and-the-phone-stand-project}

Apply transforms to build a multi-part assembly. Covers \texttt{translate()}, \texttt{rotate()}, \texttt{scale()}, Minkowski sum for edge rounding, multi-part assemblies, and parametric angle variations. \textbf{Project:} Phone stand with configurations for phones, tablets, and documents. \textbf{Checkpoint:} You can design multi-part assemblies with positioned components.

\subsubsection*{Lesson 8: Advanced Parametric Design and Interlocking Features}\label{docs__pandoc__latex__src__3dmake_all_lessons.md__lesson-8-advanced-parametric-design-and-interlocking-features}

Design tolerance-critical assemblies where parts snap together. Covers tolerance and clearance management, stack-up error, interlocking rims, snap-fit connectors, and chamfers. \textbf{Project:} Stackable storage bins --- test different \texttt{stack\_clear} values, create size variants, add optional dividers. \textbf{Checkpoint:} You understand tolerance management and can design stackable assemblies.

\subsubsection*{Lesson 9: Automation and 3dm Workflows}\label{docs__pandoc__latex__src__3dmake_all_lessons.md__lesson-9-automation-and-3dm-workflows}

Automate design workflows using shell scripts. Covers chaining 3dm commands with \texttt{\&\&}, shell script batch processing, library management (\texttt{3dm\ lib}), variant testing matrices, and production build workflows. \textbf{Checkpoint:} You can automate design workflows and manage parametric variants at scale.

\subsubsection*{Lesson 10: Hands-On Practice Exercises and Troubleshooting}\label{docs__pandoc__latex__src__3dmake_all_lessons.md__lesson-10-hands-on-practice-exercises-and-troubleshooting}

Capstone lesson synthesizing all prior learning across three exercise sets: \textbf{Set A} (guided projects: Phone Stand, Keycap Set, Storage System), \textbf{Set B} (problem diagnosis: non-manifold geometry, print failure prevention, dimensional accuracy), and \textbf{Set C} (validation and documentation templates). \textbf{Checkpoint:} You can complete real projects from concept to verified print.

\subsubsection*{Lesson 11: Stakeholder-Centric Design and the Beaded Jewelry Project}\label{docs__pandoc__latex__src__3dmake_all_lessons.md__lesson-11-stakeholder-centric-design-and-the-beaded-jewelry-project}

Design for real users, not just yourself. Covers stakeholder identification and interview techniques, extracting functional requirements, defining measurable acceptance criteria, design iteration based on feedback, and documentation for reproducibility. \textbf{Project:} Beaded jewelry bracelet holder --- interview a stakeholder, extract requirements, design parametrically, test with actual bracelets, iterate. \textbf{Checkpoint:} You can conduct a real interview, translate needs into measurable requirements, and iterate a design based on user feedback.

\subsection*{Project Completion Tracking}\label{docs__pandoc__latex__src__3dmake_all_lessons.md__project-completion-tracking}

{\def\LTcaptype{none} % do not increment counter
\begin{longtable}[]{@{}
  >{\raggedright\arraybackslash}p{(\linewidth - 6\tabcolsep) * \real{0.2500}}
  >{\raggedright\arraybackslash}p{(\linewidth - 6\tabcolsep) * \real{0.2500}}
  >{\raggedright\arraybackslash}p{(\linewidth - 6\tabcolsep) * \real{0.2500}}
  >{\raggedright\arraybackslash}p{(\linewidth - 6\tabcolsep) * \real{0.2500}}@{}}
\toprule\noalign{}
\begin{minipage}[b]{\linewidth}\raggedright
Project
\end{minipage} & \begin{minipage}[b]{\linewidth}\raggedright
Lessons
\end{minipage} & \begin{minipage}[b]{\linewidth}\raggedright
Output Files
\end{minipage} & \begin{minipage}[b]{\linewidth}\raggedright
Est. Print Time
\end{minipage} \\
\midrule\noalign{}
\endhead
\bottomrule\noalign{}
\endlastfoot
Keycap Set & 1--6 & 5+ SCAD files, 3+ STL files & 15--30 min total \\
Phone Stand & 1--7 & 1 SCAD with 3+ variants & 30--60 min total \\
Storage Bins & 1--8 & 1 parametric SCAD, 3 sizes & 60--90 min total \\
\end{longtable}
}

\subsection*{Assessment and Completion}\label{docs__pandoc__latex__src__3dmake_all_lessons.md__assessment-and-completion}

Each lesson includes learning objectives, step-by-step tasks, checkpoints, a 10-question self-assessment quiz, and 10 extension problems.

\textbf{Lesson completion criteria:} Watched/read entire lesson · Completed all step-by-step tasks · Reached all checkpoints · Answered quiz questions · Attempted at least 3 extension problems · Documented findings.

\textbf{Project completion criteria:} Code builds without errors · All parameters functional · STL generated and inspected · Measurements documented · Assembly tested (if multi-part) · README or documentation included.

\subsection*{Appendices Overview}\label{docs__pandoc__latex__src__3dmake_all_lessons.md__appendices-overview}

{\def\LTcaptype{none} % do not increment counter
\begin{longtable}[]{@{}
  >{\raggedright\arraybackslash}p{(\linewidth - 6\tabcolsep) * \real{0.3030}}
  >{\raggedright\arraybackslash}p{(\linewidth - 6\tabcolsep) * \real{0.2121}}
  >{\raggedright\arraybackslash}p{(\linewidth - 6\tabcolsep) * \real{0.1818}}
  >{\raggedright\arraybackslash}p{(\linewidth - 6\tabcolsep) * \real{0.3030}}@{}}
\toprule\noalign{}
\begin{minipage}[b]{\linewidth}\raggedright
Appendix
\end{minipage} & \begin{minipage}[b]{\linewidth}\raggedright
Focus
\end{minipage} & \begin{minipage}[b]{\linewidth}\raggedright
Size
\end{minipage} & \begin{minipage}[b]{\linewidth}\raggedright
Use When
\end{minipage} \\
\midrule\noalign{}
\endhead
\bottomrule\noalign{}
\endlastfoot
A: Comprehensive Slicing Guide &
PrusaSlicer, Bambu Studio, Cura, OrcaSlicer, and more & 1,500+ lines &
Slicing questions, slicer reference \\
B: Material Properties \& Selection &
PLA, PETG, ABS, TPU, Polycarbonate, Nylon --- properties, settings, decision tree
& 1,200+ lines & Choosing material, troubleshooting prints \\
C: Tolerance Testing \& QA Matrix &
Measurement procedures, functional testing, stack-up calculations &
1,200+ lines & Quality verification, measurement techniques \\
D: PowerShell Integration &
Batch processing, automation scripts, workflow integration &
1,100+ lines & Building automation, batch processing \\
E: Advanced OpenSCAD Concepts &
Gears, batch/statistical analysis, performance optimization, recursion &
Optional & Specialized/advanced applications \\
\end{longtable}
}

\subsection*{Accessibility Features}\label{docs__pandoc__latex__src__3dmake_all_lessons.md__accessibility-features}

This curriculum is designed for non-visual learners: text descriptions of all models via \texttt{3dm\ info}, tactile 2D previews via \texttt{3dm\ preview}, structured written documentation, command-line based (no graphical interface required), and measurement-based validation without visual inspection.

\subsection*{Next Steps After Completion}\label{docs__pandoc__latex__src__3dmake_all_lessons.md__next-steps-after-completion}

Upon completing this curriculum, students are ready for: advanced BOSL2 library features and parametric assemblies; multi-material printing (TPU, Nylon, etc.); mechanical assemblies with bearings, gears, and mechanisms; community contribution via shared designs and libraries; and professional applications in product design, prototyping, and manufacturing.

\subsection*{Curriculum Revision History}\label{docs__pandoc__latex__src__3dmake_all_lessons.md__curriculum-revision-history}

{\def\LTcaptype{none} % do not increment counter
\begin{longtable}[]{@{}
  >{\raggedright\arraybackslash}p{(\linewidth - 4\tabcolsep) * \real{0.3333}}
  >{\raggedright\arraybackslash}p{(\linewidth - 4\tabcolsep) * \real{0.3333}}
  >{\raggedright\arraybackslash}p{(\linewidth - 4\tabcolsep) * \real{0.3333}}@{}}
\toprule\noalign{}
\begin{minipage}[b]{\linewidth}\raggedright
Version
\end{minipage} & \begin{minipage}[b]{\linewidth}\raggedright
Date
\end{minipage} & \begin{minipage}[b]{\linewidth}\raggedright
Changes
\end{minipage} \\
\midrule\noalign{}
\endhead
\bottomrule\noalign{}
\endlastfoot
2.1 & Feb 2026 &
Added Appendix E (Advanced OpenSCAD); Enhanced Lessons 3, 6, 7, 8, 9 with advanced topics \\
2.0 & Feb 2026 &
Added Lesson 11 (Stakeholder Design) + 4 Appendices; Consolidated Units 0--3 content \\
1.0 & Feb 2026 &
Initial comprehensive curriculum with 10 lessons + 5 projects \\
\end{longtable}
}

\begin{center}\rule{0.5\linewidth}{0.5pt}\end{center}

\section*{Part III --- All 11 Lessons}\label{docs__pandoc__latex__src__3dmake_all_lessons.md__part-iii--all-11-lessons}

\section*{Lesson 1 --- Environmental Configuration and the Developer Workflow}\label{docs__pandoc__latex__src__3dmake_all_lessons.md__lesson-1--environmental-configuration-and-the-developer-workflow}

\textbf{Estimated Time:} 60--90 minutes\\
\textbf{Course:} 3dMake Certification --- High School Distance Learning Track

\begin{center}\rule{0.5\linewidth}{0.5pt}\end{center}

\subsection*{A Message to You Before We Begin}\label{docs__pandoc__latex__src__3dmake_all_lessons.md__a-message-to-you-before-we-begin}

Welcome to your first lesson. If you\textquotesingle ve never used a command line, never written code, or never sent a file to a 3D printer before, you are in exactly the right place. This lesson is designed so that someone who has done none of those things can follow every step.

If you \emph{have} done some of those things, you\textquotesingle ll still find this lesson useful --- because the way 3D printing tools fit together is different from most software you\textquotesingle ve used before, and understanding the full picture from the start will save you real trouble later.

By the time you finish this lesson, you will have your tools installed and working, you\textquotesingle ll understand the basic logic of the coordinate system your designs live in, and you\textquotesingle ll have run your first parametric 3D model through the complete build-and-verify workflow.

Take your time here. The concepts in this first lesson are the foundation everything else is built on.

\begin{center}\rule{0.5\linewidth}{0.5pt}\end{center}

\subsection*{Learning Objectives}\label{docs__pandoc__latex__src__3dmake_all_lessons.md__learning-objectives}

By the end of this lesson you will be able to:

\begin{itemize}
\tightlist
\item
  Install and verify the 3dMake toolchain on Windows, macOS, and Linux
\item
  Understand the OpenSCAD coordinate system and unit conventions
\item
  Write and build your first parametric \texttt{.scad} file
\item
  Read and interpret \texttt{3dm\ info} output without a visual display
\item
  Use the core developer workflow: edit → build → verify
\end{itemize}

\begin{center}\rule{0.5\linewidth}{0.5pt}\end{center}

\subsection*{Concept 1: What Is Parametric 3D Design?}\label{docs__pandoc__latex__src__3dmake_all_lessons.md__concept-1-what-is-parametric-3d-design}

\subsubsection*{The Two Ways to Design for 3D Printing}\label{docs__pandoc__latex__src__3dmake_all_lessons.md__the-two-ways-to-design-for-3d-printing}

There are two ways to design a 3D object for printing. The first way is to use a \textbf{visual editor} --- a program where you click, drag, and push shapes around the screen until they look right. Programs like Tinkercad work this way. They\textquotesingle re intuitive and fast for simple objects, but they have a limitation: if you want to change the size of something, you often have to go back and manually adjust every related dimension by hand.

The second way is to use \textbf{code}. You write instructions that describe your object mathematically. "Make a box that is 30 millimeters wide, 20 millimeters deep, and 10 millimeters tall." If you later decide that 30 mm is too narrow and 40 mm would be better, you change one number, and the model automatically updates everywhere that number is used. This is called \textbf{parametric design} --- the word "parametric" means the design is controlled by \emph{parameters}, which are the named values you can adjust.

Parametric design is how professional engineers create parts. An engineer designing a phone stand might create one file that works for phone models ranging from 60 to 90 mm wide just by changing a single variable. That single source file \emph{is the design concept}, and specific values turn it into specific objects.

\textbf{OpenSCAD} is the tool you\textquotesingle ll use for parametric design in this course. Instead of clicking and dragging, you write code that describes shapes and their relationships. The learning curve is steeper than a visual editor at the very beginning, but within a few lessons you\textquotesingle ll be creating things that would be extremely difficult or impossible in a drag-and-drop tool.

\subsubsection*{What 3dMake Does}\label{docs__pandoc__latex__src__3dmake_all_lessons.md__what-3dmake-does}

OpenSCAD by itself is powerful, but its workflow has a lot of friction. To go from code to a printable file, you\textquotesingle d normally need to open OpenSCAD, manually trigger a render, export an STL file, open a separate slicer program, check for geometry errors, and save the output. That\textquotesingle s a lot of steps just to see if your idea works.

\textbf{3dMake} is a tool that wraps around OpenSCAD and smooths out that friction. When you type \texttt{3dm\ build}, it runs the entire compile-and-export process for you automatically. When you type \texttt{3dm\ info}, it analyzes your model and tells you its dimensions, volume, and a description in plain language. The more time you spend designing rather than managing files, the faster you learn and the more you make.

Think of it this way: OpenSCAD is the engine, and 3dMake is the dashboard. You could use the engine directly, but the dashboard makes it dramatically easier to operate.

\begin{center}\rule{0.5\linewidth}{0.5pt}\end{center}

\subsection*{Concept 2: The Command Line}\label{docs__pandoc__latex__src__3dmake_all_lessons.md__concept-2-the-command-line}

\subsubsection*{What It Is and Why We Use It}\label{docs__pandoc__latex__src__3dmake_all_lessons.md__what-it-is-and-why-we-use-it}

You\textquotesingle re going to be working with the \textbf{command line} throughout this course. If you\textquotesingle ve mostly used computers by clicking icons and menus, the command line can feel strange at first --- but it\textquotesingle s really just a different way to talk to your computer.

Here is the basic idea: the command line is a text box where you type an instruction, press Enter, and the computer does something and prints a response. That\textquotesingle s the whole interaction model. Type a command. Press Enter. Read the result. Repeat.

For example, if you type \texttt{3dm\ build} and press Enter, the computer finds the 3dMake program, tells it to build your project, and then prints either a success message with your model\textquotesingle s dimensions, or an error message with a line number pointing to the problem. This feedback loop is tight, fast, and precise --- which is exactly what you want when you\textquotesingle re designing and testing.

The command line has another advantage: it\textquotesingle s repeatable. If you type a command that works, you can type it again. You can save it in a script and run it automatically. You can share it with someone else, and they can run the same command and get the same result. Clicking around menus is difficult to document and impossible to automate; text commands are both.

\subsubsection*{Which Terminal to Use}\label{docs__pandoc__latex__src__3dmake_all_lessons.md__which-terminal-to-use}

\textbf{Windows users:} You have two options. \textbf{PowerShell} is the modern choice and the one we recommend --- it\textquotesingle s more powerful and consistent with what you\textquotesingle ll find in professional environments. To open it, press the Windows key, type \texttt{PowerShell}, and press Enter. Alternatively, \textbf{Command Prompt (CMD)} is an older option that still works for basic tasks. This course provides code examples for both.

\textbf{macOS users:} Open \textbf{Terminal} by going to Applications → Utilities → Terminal. It runs a shell called \texttt{bash} (or \texttt{zsh} on newer Macs --- both work for this course).

\textbf{Linux users:} You almost certainly already know how to open a terminal. If not, look for "Terminal" in your applications menu.

\subsubsection*{A Few Things to Know Before We Start}\label{docs__pandoc__latex__src__3dmake_all_lessons.md__a-few-things-to-know-before-we-start}

\begin{itemize}
\tightlist
\item
  On macOS and Linux, \textbf{file paths use forward slashes}: \texttt{src/main.scad}
\item
  On Windows, \textbf{file paths use backslashes}: \texttt{src\textbackslash{}main.scad}
\item
  Commands are \textbf{case-sensitive on macOS and Linux} --- \texttt{3DM} and \texttt{3dm} are different things, and only the lowercase version will work
\item
  When a command says it "isn\textquotesingle t found," it almost never means the software isn\textquotesingle t installed --- it usually means the terminal opened before the installation updated its search path (we\textquotesingle ll cover how to fix this in Step 1)
\end{itemize}

\begin{center}\rule{0.5\linewidth}{0.5pt}\end{center}

\subsection*{Step 1 --- Install the 3dMake Toolchain}\label{docs__pandoc__latex__src__3dmake_all_lessons.md__step-1--install-the-3dmake-toolchain}

\subsubsection*{Understanding What You\textquotesingle re Installing}\label{docs__pandoc__latex__src__3dmake_all_lessons.md__understanding-what-youre-installing}

You\textquotesingle re installing two pieces of software that work together:

\textbf{OpenSCAD} is the underlying 3D geometry engine. It reads your \texttt{.scad} code files and turns them into 3D model files (specifically \texttt{.stl} files). You can think of it like a compiler for 3D geometry --- the same way a Python interpreter turns Python code into running behavior, OpenSCAD turns design code into physical geometry. You don\textquotesingle t need to interact with OpenSCAD directly very often; 3dMake does it for you.

\textbf{3dMake} is the project management layer that sits on top of OpenSCAD. It organizes your project files, runs OpenSCAD with the right settings, analyzes your model output, and provides the plain-English \texttt{3dm\ info} descriptions you\textquotesingle ll use to verify your work.

The installer sets up both. Run the appropriate command for your operating system below, let it finish completely, and then run the verification commands before doing anything else.

\subsubsection*{Installation Commands}\label{docs__pandoc__latex__src__3dmake_all_lessons.md__installation-commands}

\textbf{Linux and macOS (bash terminal):}

\begin{lstlisting}[style=Alabaster, language=bash]
# Download and run the official installer

curl -fsSL https://get.3dmake.dev | bash


# After it finishes, verify both tools are installed correctly.

# Both of these commands should print version numbers, not errors.

3dm --version

openscad --version


\end{lstlisting}

\textbf{Windows (PowerShell):}

\begin{lstlisting}[style=Alabaster, language=powershell]
# Install using Windows Package Manager

winget install 3dMake.3dMake


# If winget isn't available, download the installer directly:

# https://3dmake.dev/install

# Run the downloaded installer, then open a NEW PowerShell window.


# Verify the installation:

3dm --version

openscad --version


\end{lstlisting}

\textbf{Windows (Command Prompt / CMD):}

\begin{lstlisting}[style=Alabaster, language=cmd]
REM After installation, open a NEW CMD window and verify:

3dm --version

openscad --version


REM If you see "not recognized as an internal or external command",

REM see the troubleshooting section below.


\end{lstlisting}

\subsubsection*{What the Version Numbers Mean}\label{docs__pandoc__latex__src__3dmake_all_lessons.md__what-the-version-numbers-mean}

When you run \texttt{3dm\ -\/-version}, you should see output like \texttt{3dMake\ v1.4.2} (the exact number doesn\textquotesingle t matter as long as something prints). When you run \texttt{openscad\ -\/-version}, you should see \texttt{OpenSCAD\ version\ 2021.01} or later. If either command produces an error instead of a version number, something went wrong with the installation. Read the troubleshooting section before moving on.

\subsubsection*{Troubleshooting Installation}\label{docs__pandoc__latex__src__3dmake_all_lessons.md__troubleshooting-installation}

\textbf{"command not found" or "is not recognized as an internal or external command"}

This is the most common issue and it almost always has the same cause: your terminal session started before the installer ran. When the installer finishes, it adds the 3dMake folder to a system setting called the \textbf{PATH} --- a list of folders your computer searches when you type a command name. But a terminal that was already open when the installer ran doesn\textquotesingle t know about the change. The fix is simple: \textbf{close your terminal completely and open a new one}, then try the verification commands again.

\textbf{What is PATH?} When you type \texttt{3dm} at the command line, your computer doesn\textquotesingle t search your entire hard drive --- that would take forever. Instead, it searches only the folders listed in PATH. If 3dMake\textquotesingle s folder isn\textquotesingle t in that list, the computer genuinely can\textquotesingle t find the program even though it\textquotesingle s installed. Closing and reopening the terminal forces it to reload the PATH list.

If opening a new terminal doesn\textquotesingle t fix it, you may need to add the folder manually. On Windows, search the Start menu for "Edit the system environment variables," click "Environment Variables," and add the 3dMake installation folder to the \texttt{Path} variable. On macOS/Linux, add the folder path to your \texttt{\textasciitilde{}/.bashrc} or \texttt{\textasciitilde{}/.zshrc} file and restart the terminal.

\textbf{OpenSCAD version too old}

3dMake requires OpenSCAD 2021 or later. If you have an older version from a previous installation, download the latest release from openscad.org, uninstall the old version, and install the new one.

\begin{center}\rule{0.5\linewidth}{0.5pt}\end{center}

\subsection*{Concept 3: The Coordinate System --- Where Objects Live in 3D Space}\label{docs__pandoc__latex__src__3dmake_all_lessons.md__concept-3-the-coordinate-system--where-objects-live-in-3d-space}

\subsubsection*{Why You Need to Understand This Before Writing Code}\label{docs__pandoc__latex__src__3dmake_all_lessons.md__why-you-need-to-understand-this-before-writing-code}

Imagine you\textquotesingle re on the phone with a friend trying to tell them where to place a chair in an empty room. You\textquotesingle d need to agree on two things: where to measure from (a corner? the center of the room?), and which directions count as "forward," "sideways," and "up." Without that shared reference, your instructions are meaningless.

OpenSCAD works exactly the same way. Every object you create lives at a specific location in 3D space. That space has a fixed starting point called the \textbf{origin} and three directions called \textbf{axes}. If you don\textquotesingle t know these, you\textquotesingle ll place objects and wonder why they appear where they do.

\subsubsection*{The Three Axes}\label{docs__pandoc__latex__src__3dmake_all_lessons.md__the-three-axes}

OpenSCAD uses a \textbf{right-handed XYZ coordinate system} --- a standard that\textquotesingle s shared by engineering software worldwide, so learning it here transfers directly to other professional tools.

Picture a 3D printer\textquotesingle s build plate in front of you:

\begin{itemize}
\tightlist
\item
  \textbf{X axis} --- points to the \textbf{right} across the build plate. Positive X is rightward; negative X is leftward.
\item
  \textbf{Y axis} --- points \textbf{away from you}, toward the back of the build plate. Positive Y is away from you; negative Y is toward you.
\item
  \textbf{Z axis} --- points \textbf{straight up}, perpendicular to the build plate. Positive Z is upward; negative Z would be underground.
\end{itemize}

The \textbf{origin} --- the point (0, 0, 0) --- sits at the front-left corner of the build space, right at bed level. When you type \texttt{cube({[}20,\ 20,\ 10{]})} without any translation, it appears with its front-left-bottom corner at the origin, extending 20 mm to the right, 20 mm back, and 10 mm up.

\subsubsection*{Why This Matters for Printing}\label{docs__pandoc__latex__src__3dmake_all_lessons.md__why-this-matters-for-printing}

Because the Z axis points up and the origin is at bed level, any geometry with a negative Z coordinate would be "underground" --- below the build plate. The printer can\textquotesingle t print there. Part of your job when designing is to make sure your entire model lives in positive Z space (Z ≥ 0) so it sits on the bed correctly.

\subsubsection*{Units: Always Millimeters}\label{docs__pandoc__latex__src__3dmake_all_lessons.md__units-always-millimeters}

OpenSCAD has no built-in concept of units --- it just works with numbers. The convention in this course (and in the 3D printing community generally) is that \textbf{all dimensions are in millimeters}. This is not optional or adjustable --- it\textquotesingle s a universal agreement.

When you write \texttt{cube({[}30,\ 20,\ 10{]})}, that means 30 mm × 20 mm × 10 mm. If you accidentally think in centimeters and write \texttt{cube({[}3,\ 2,\ 1{]})} meaning "3 cm by 2 cm by 1 cm," OpenSCAD will treat those as millimeters and your print will come out ten times too small. Always, always, always think in millimeters.

\subsubsection*{\texorpdfstring{The \texttt{center} Parameter: Where Does the Object Anchor?}{The center Parameter: Where Does the Object Anchor?}}\label{docs__pandoc__latex__src__3dmake_all_lessons.md__the-center-parameter-where-does-the-object-anchor}

Every 3D primitive in OpenSCAD has a \texttt{center} parameter. This is one of the most important things to understand early, because it affects where objects appear and how easy your placement math is.

\textbf{\texttt{center=false}} (this is the default --- it applies if you don\textquotesingle t write anything): the object\textquotesingle s \textbf{bottom-left-front corner} sits at the position you specify. A \texttt{cube({[}20,\ 20,\ 10{]})} with no translation occupies the space from X=0 to X=20, Y=0 to Y=20, Z=0 to Z=10. The object lives entirely in positive space and sits neatly on the build plate.

\textbf{\texttt{center=true}}: the object\textquotesingle s \textbf{geometric center} sits at the position you specify. The same cube at the origin would extend from X=−10 to X=10, Y=−10 to Y=10, Z=−5 to Z=5. Half of it would be underground --- which means if you print it without moving it up first, you\textquotesingle ll only get the top half.

Neither option is wrong. You choose based on which makes your math simpler. For parts that should sit on the build plate, \texttt{center=false} is more natural. For parts you want to rotate around their center, \texttt{center=true} is much cleaner.

\begin{lstlisting}[style=Alabaster, language=openscad]
// Demonstrating center=false vs center=true

// Build this and use 3dm info to see the bounding box of each shape.


// This cube sits on the bed naturally — its bottom is at Z=0.

cube([20, 20, 10]);                    // corner at origin; occupies X:0–20, Y:0–20, Z:0–10


// This cube straddles the origin — half of it is underground!

// Move it sideways with translate() so it doesn't overlap the first one.

translate([35, 0, 0])

  cube([20, 20, 10], center=true);     // center at origin; occupies X:25–45, Y:-10–10, Z:-5–5


// A small red sphere marks where the origin is.

color("red") sphere(r=1.5, $fn=16);


// Colored cylinders mark the positive direction of each axis.

color("red")   translate([25, 0, 0]) cylinder(r=0.5, h=1, $fn=8);  // +X

color("green") translate([0, 25, 0]) cylinder(r=0.5, h=1, $fn=8);  // +Y

color("blue")  cylinder(r=0.5, h=25, $fn=8);                        // +Z


\end{lstlisting}

Note

When using \texttt{center=true}, notice how the bounding box may show negative coordinates --- that\textquotesingle s geometry underground. In a real project, always use \texttt{translate()} to lift centered objects so their bottom is at Z=0.

Build this code and run \texttt{3dm\ info} to read the bounding boxes.

\begin{center}\rule{0.5\linewidth}{0.5pt}\end{center}

\subsection*{Step 2 --- Create Your First Project}\label{docs__pandoc__latex__src__3dmake_all_lessons.md__step-2--create-your-first-project}

\subsubsection*{What a "Project" Is}\label{docs__pandoc__latex__src__3dmake_all_lessons.md__what-a-project-is}

3dMake organizes work into \textbf{projects} --- folders with a specific, standard structure. This structure keeps your source code separate from generated output, which makes it easy to clean, rebuild, back up, and share your work.

Here is what a project folder looks like:

\begin{lstlisting}[style=Alabaster]
my_project/

  src/          ← your .scad source files live here (YOU write and edit these)

  build/        ← generated STL files go here (3dMake creates these automatically)

  3dmake.toml   ← project configuration file (created when you initialize the project)


\end{lstlisting}

Important

Never manually edit files in the \texttt{build/} folder. Everything in there is generated automatically --- changes will be overwritten the next time you run \texttt{3dm\ build}. All your actual work goes in \texttt{src/}.

\subsubsection*{Creating the Project}\label{docs__pandoc__latex__src__3dmake_all_lessons.md__creating-the-project}

\textbf{Linux / macOS (bash):}

\begin{lstlisting}[style=Alabaster, language=bash]
# Create the project folder and its standard subdirectories

mkdir my_first_project

cd my_first_project

mkdir src build


# Create the main design file

cat > src/main.scad << 'EOF'

// My First Parametric Model

// ========================

// All dimensions in millimeters.

// To customize this design, change the values below and run: 3dm build


width  = 30;   // mm — how wide the box is (X dimension)

depth  = 20;   // mm — how deep the box is (Y dimension, front to back)

height = 10;   // mm — how tall the box is (Z dimension)


cube([width, depth, height]);

EOF


\end{lstlisting}

\textbf{Windows (PowerShell):}

\begin{lstlisting}[style=Alabaster, language=powershell]
# Create project folder and subdirectories

New-Item -ItemType Directory -Path "my_first_project\src","my_first_project\build"

Set-Location "my_first_project"


# Create the main design file

@'

// My First Parametric Model

// All dimensions in millimeters.

// Change values below and run: 3dm build


width  = 30;   // mm

depth  = 20;   // mm

height = 10;   // mm


cube([width, depth, height]);

'@ | Out-File -FilePath "src\main.scad" -Encoding UTF8


\end{lstlisting}

\textbf{Windows (CMD):}

\begin{lstlisting}[style=Alabaster, language=cmd]
mkdir my_first_project\src

mkdir my_first_project\build

cd my_first_project


echo // My First Parametric Model > src\main.scad

echo width  = 30; >> src\main.scad

echo depth  = 20; >> src\main.scad

echo height = 10; >> src\main.scad

echo cube([width, depth, height]); >> src\main.scad


\end{lstlisting}

\subsubsection*{Reading the Code Line by Line}\label{docs__pandoc__latex__src__3dmake_all_lessons.md__reading-the-code-line-by-line}

This small file introduces the core pattern you\textquotesingle ll use for the entire course. Let\textquotesingle s read every line carefully.

\texttt{//\ My\ First\ Parametric\ Model} --- this is a \textbf{comment}. In OpenSCAD, anything following \texttt{//} on a line is ignored by the computer --- it\textquotesingle s a note for human readers. Comments are essential. They explain what the code does, what the units are, and what you were thinking when you wrote it.

\texttt{width\ =\ 30;} --- this \textbf{declares a variable} named \texttt{width} and sets it equal to 30. The semicolon at the end is required --- it tells OpenSCAD "this statement is finished." The comment \texttt{//\ mm} reminds anyone reading the file that this value is in millimeters.

\texttt{depth\ =\ 20;} and \texttt{height\ =\ 10;} --- same pattern. Each variable represents one dimension of the box.

\texttt{cube({[}width,\ depth,\ height{]});} --- this creates a box. The brackets \texttt{{[}{]}} enclose a list of three values: the X size, the Y size, and the Z size. Instead of writing \texttt{cube({[}30,\ 20,\ 10{]})} with raw numbers, you use the variable names. This is the magic of parametric design: if you change \texttt{width\ =\ 30} to \texttt{width\ =\ 50} at the top, the cube automatically becomes 50 mm wide everywhere \texttt{width} is used --- which might be one place in a simple file, but could be dozens of places in a complex one.

This pattern --- declare named parameters at the top of the file, use them in geometry below --- is the single most important habit this course will teach you.

\begin{center}\rule{0.5\linewidth}{0.5pt}\end{center}

\subsection*{Step 3 --- Build and Verify}\label{docs__pandoc__latex__src__3dmake_all_lessons.md__step-3--build-and-verify}

\subsubsection*{Running Your First Build}\label{docs__pandoc__latex__src__3dmake_all_lessons.md__running-your-first-build}

With your terminal open inside the \texttt{my\_first\_project/} folder, run:

\begin{lstlisting}[style=Alabaster, language=bash]
3dm build


\end{lstlisting}

This single command does several things: it finds your \texttt{src/main.scad} file, runs OpenSCAD on it to turn the code into geometry, checks the result for basic validity, and saves the output as \texttt{build/main.stl}.

If the build succeeds, you\textquotesingle ll see output indicating the file was written. If there\textquotesingle s an error in your code, you\textquotesingle ll see an error message with a line number pointing to the problem. Error messages with line numbers are extremely helpful --- they tell you exactly where to look to fix the issue.

\subsubsection*{Understanding the Model Information}\label{docs__pandoc__latex__src__3dmake_all_lessons.md__understanding-the-model-information}

Now run:

\begin{lstlisting}[style=Alabaster, language=bash]
3dm info


\end{lstlisting}

This analyzes your built model and prints a report. You\textquotesingle ll see something like this:

\begin{lstlisting}[style=Alabaster]
File:         build/main.stl

Bounding Box: 30.0 x 20.0 x 10.0 mm

Volume:       6000.0 mm³

Triangles:    12

Description:  A small rectangular box measuring 30mm wide by 20mm deep by 10mm tall.


\end{lstlisting}

Let\textquotesingle s understand each piece.

\textbf{Bounding Box:} The bounding box is the smallest rectangular box that perfectly contains your entire model. \texttt{30.0\ x\ 20.0\ x\ 10.0\ mm} means your model is 30 mm in X, 20 mm in Y, and 10 mm in Z. This number is your primary verification tool --- if the bounding box doesn\textquotesingle t match what you designed, something went wrong. Check it every single time you build.

\textbf{Volume:} This is how much solid material fills the model. A solid box that is 30 × 20 × 10 mm contains exactly 30 × 20 × 10 = 6,000 cubic millimeters of material. If your model were hollow, the volume would be less. Volume directly predicts how much filament you\textquotesingle ll need.

\textbf{Triangles:} STL files represent surfaces as a mesh of triangles. More triangles means more geometric detail (and a larger file). A simple box needs only 12 triangles. A sphere with smooth curves needs thousands.

\textbf{Description:} The plain-English description comes from an AI analysis of the geometry. It\textquotesingle s useful for a quick sanity check, but it\textquotesingle s not a replacement for checking the bounding box numbers yourself.

Caution

If the bounding box shows \textbf{0 × 0 × 0}, the build failed and no geometry was produced. Check the output from \texttt{3dm\ build} --- error messages will point to the specific line in your code.

\subsubsection*{Estimating Material Cost from Volume}\label{docs__pandoc__latex__src__3dmake_all_lessons.md__estimating-material-cost-from-volume}

Before printing anything --- especially something that might take hours --- it\textquotesingle s professional practice to estimate the cost. Here is the formula:

\begin{lstlisting}[style=Alabaster]
mass (grams) = (volume_mm³ ÷ 1000) × filament_density_g_per_cm³

cost ($)     = mass_g × (spool_price_$ ÷ spool_weight_g)


\end{lstlisting}

For our 6,000 mm³ box using PLA filament:

\begin{lstlisting}[style=Alabaster]
mass = (6000 ÷ 1000) × 1.24 = 7.44 grams

cost = 7.44 × (20 ÷ 1000)   = $0.15


\end{lstlisting}

\textbf{Common filament densities:} PLA = 1.24 g/cm³, PETG = 1.27 g/cm³, ABS = 1.05 g/cm³, TPU = 1.20 g/cm³.

This calculation won\textquotesingle t be exact (slicer settings affect actual usage), but it gives you a reasonable estimate within about 20\% of the actual cost.

\begin{center}\rule{0.5\linewidth}{0.5pt}\end{center}

\subsection*{Concept 4: The FDM Printing Pipeline}\label{docs__pandoc__latex__src__3dmake_all_lessons.md__concept-4-the-fdm-printing-pipeline}

\subsubsection*{How Your Code Becomes a Physical Object}\label{docs__pandoc__latex__src__3dmake_all_lessons.md__how-your-code-becomes-a-physical-object}

Understanding how 3D printing actually works changes how you design. The printing process has physical constraints --- things it can and can\textquotesingle t do --- and every design decision you make either works with those constraints or fights against them. Here\textquotesingle s the complete journey from your code to a finished print.

\textbf{Stage 1 --- STL File}
When you run \texttt{3dm\ build}, OpenSCAD converts your \texttt{.scad} code into an STL file. STL stands for \textbf{Standard Triangle Language} (sometimes called Standard Tessellation Language). An STL file describes the \emph{surface} of your 3D model as a mesh of tiny triangles --- like wrapping a shape in a piece of paper made of triangles. It contains no color, no material information, no internal structure --- just the outer surface geometry.

\textbf{Stage 2 --- Slicing}
A program called a \textbf{slicer} (PrusaSlicer and Cura are popular options) reads the STL and slices the model into horizontal layers, like cutting through a loaf of bread to see the shape of each slice. The thickness of each slice is the \textbf{layer height} --- typically 0.15 to 0.30 mm.

\textbf{Stage 3 --- G-code}
The slicer converts those layer slices into a file of \textbf{G-code} --- a language of machine instructions telling the printer exactly what to do. G-code commands look like: "move the print head to X=50, Y=30; heat the nozzle to 215°C; extrude 2.3 mm of filament." The printer reads these instructions literally and executes them in sequence.

\textbf{Stage 4 --- Printing}
The printer reads the G-code and executes it. The nozzle heats to between 200°C and 240°C (depending on the material), melts a thin strand of plastic filament, and deposits it in the shape of each layer --- one layer on top of the last, slowly building up the 3D form.

\textbf{Stage 5 --- Cooling and Bonding}
As each layer is deposited, it needs to cool and bond to the layer below it. This inter-layer fusion is what holds the part together. However, it also means 3D-printed parts are slightly weaker in the vertical direction than in the horizontal direction --- because the bond between layers is never quite as strong as the plastic itself. Good design takes this into account: parts that will experience stress should be oriented so the layers are perpendicular to the direction of that stress.

\subsubsection*{The Three Key Slicer Settings to Know}\label{docs__pandoc__latex__src__3dmake_all_lessons.md__the-three-key-slicer-settings-to-know}

\textbf{Layer Height} --- how thick each printed layer is, typically 0.15 to 0.30 mm. Thinner layers produce smoother surfaces with more detail, but require more layers to reach the same total height, so prints take longer. For most classroom projects, 0.20 mm is a good starting point.

\textbf{Infill Percentage} --- the internal structure of a print. A 100\% infill print is completely solid. A 0\% print is just the outer walls (hollow). In practice, most parts work well at 15--20\% infill --- the slicer creates a grid-like internal structure that supports the walls without using much material. Structural parts that will bear loads or stress should use 40--50\% infill.

\textbf{Supports} --- if your design has overhanging geometry (parts that extend horizontally without anything underneath them), the slicer can add temporary scaffolding called supports. These get removed after printing. The problem is that supports can leave rough surfaces and add both print time and material cost. Good design minimizes overhangs, which reduces or eliminates the need for supports.

\subsubsection*{First Prints Are Tests, Not Final Products}\label{docs__pandoc__latex__src__3dmake_all_lessons.md__first-prints-are-tests-not-final-products}

Tip

You will almost never get a perfect-fitting part on the first print. Treat the first print of any new feature as a \textbf{test print}: a quick, small print specifically designed to tell you what needs to be adjusted. This saves time and material by testing deliberately rather than accidentally printing a large part only to find the hole is 0.3 mm too small.

\begin{center}\rule{0.5\linewidth}{0.5pt}\end{center}

\subsection*{Step 4 --- The Edit--Build--Verify Workflow}\label{docs__pandoc__latex__src__3dmake_all_lessons.md__step-4--the-editbuildverify-workflow}

\subsubsection*{The Fundamental Loop}\label{docs__pandoc__latex__src__3dmake_all_lessons.md__the-fundamental-loop}

Professional software development runs on a tight feedback cycle: write some code, run it, check the result, fix any problems, and repeat. The shorter each cycle is, the faster you learn and the fewer problems pile up. Parametric 3D design works exactly the same way:

\begin{lstlisting}[style=Alabaster]
EDIT → BUILD → VERIFY → repeat


Edit:   Change src/main.scad in any text editor. Save the file.

Build:  3dm build   → compiles your .scad code into build/main.stl

Verify: 3dm info    → confirms bounding box and volume match your design intent

        (optional)  → open build/main.stl in your slicer for a visual check


\end{lstlisting}

Tip

\textbf{Verify after every meaningful change}, not just when you think you\textquotesingle re done. Frequent verification means any problem is always close to the last thing you changed. Verifying only at the end of a long session means sifting through dozens of changes to find one small mistake.

\subsubsection*{Watch Mode: Automatic Rebuilds}\label{docs__pandoc__latex__src__3dmake_all_lessons.md__watch-mode-automatic-rebuilds}

If you\textquotesingle re iterating quickly on a design, you can tell 3dMake to automatically rebuild every time you save your source file:

\begin{lstlisting}[style=Alabaster, language=bash]
3dm watch


\end{lstlisting}

Run this in one terminal window, keep your text editor open in another. Every time you save \texttt{src/main.scad}, the terminal automatically runs \texttt{3dm\ build} and reports the result. This creates a nearly real-time feedback loop: you save the file and, a few seconds later, see whether it built successfully and what the new dimensions are.

\begin{center}\rule{0.5\linewidth}{0.5pt}\end{center}

\subsection*{Step 5 --- Building Parametric Variants}\label{docs__pandoc__latex__src__3dmake_all_lessons.md__step-5--building-parametric-variants}

\subsubsection*{Overriding Parameters at Build Time}\label{docs__pandoc__latex__src__3dmake_all_lessons.md__overriding-parameters-at-build-time}

Once your design uses variables for all its dimensions, you can generate different versions without editing the source file at all. The \texttt{-D} flag tells OpenSCAD to override a variable at build time:

\begin{lstlisting}[style=Alabaster, language=bash]
openscad -D "width=50" -o build/main_wide.stl src/main.scad


\end{lstlisting}

This builds a version with \texttt{width=50} instead of the \texttt{width=30} in the file --- without changing the file. You can combine multiple overrides:

\begin{lstlisting}[style=Alabaster, language=bash]
openscad -D "width=50" -D "height=15" -o build/main_large.stl src/main.scad


\end{lstlisting}

\subsubsection*{Building Multiple Variants Automatically}\label{docs__pandoc__latex__src__3dmake_all_lessons.md__building-multiple-variants-automatically}

You can use a shell loop to generate a whole family of variants automatically:

\textbf{Linux / macOS (bash):}

\begin{lstlisting}[style=Alabaster, language=bash]
>[!TIP]

>Use shell loops to generate multiple parametric variants automatically. Each run creates a separate STL file — a fast way to produce a tolerance test matrix without manually editing the source file.


```bash

# Build four width variants — generates four separate STL files

for w in 20 30 40 50; do

  openscad -D "width=$w" -o "build/main_w${w}.stl" src/main.scad

  echo "Built: width=${w}mm → build/main_w${w}.stl"

done


\end{lstlisting}

\begin{lstlisting}[style=Alabaster]

**Windows (PowerShell):**


```powershell

$widths = @(20, 30, 40, 50)

foreach ($w in $widths) {

  openscad -D "width=$w" -o "build\main_w$w.stl" "src\main.scad"

  Write-Host "Built: width=${w}mm → build\main_w$w.stl"

}


\end{lstlisting}

\textbf{Windows (CMD):}

\begin{lstlisting}[style=Alabaster, language=cmd]
FOR %%W IN (20 30 40 50) DO (

  openscad -D "width=%%W" -o build\main_w%%W.stl src\main.scad

  echo Built width=%%W

)


\end{lstlisting}

This is one of the most powerful things about parametric design: one source file, many specific outputs. The file represents the concept; the parameters turn it into instances.

\begin{center}\rule{0.5\linewidth}{0.5pt}\end{center}

\subsection*{Concept 5: Code Documentation Standards}\label{docs__pandoc__latex__src__3dmake_all_lessons.md__concept-5-code-documentation-standards}

\subsubsection*{Why Documented Code Is Part of Certification}\label{docs__pandoc__latex__src__3dmake_all_lessons.md__why-documented-code-is-part-of-certification}

When you return to a file weeks later, you will not remember what every variable was for. When a classmate, instructor, or future collaborator reads your code, they need to understand your design intent without asking you. When you submit work for this certification, well-documented code demonstrates professional-level practice --- not just that you got the geometry right, but that you thought about it clearly.

OpenSCAD documentation has three levels:

\textbf{File Header Comments} --- a block comment at the very top of the file explaining: what this part is, what it\textquotesingle s used for, print recommendations (layer height, infill, supports needed?), and which parameters are meant to be customized. This is the first thing anyone reads.

\textbf{Parameter Line Comments} --- every variable declaration should have a comment stating the unit and a reasonable range. This:

\begin{lstlisting}[style=Alabaster, language=openscad]
wall = 2;   // mm — wall thickness (1.5–3.0 for a 0.4mm nozzle)


\end{lstlisting}

is dramatically more useful than just \texttt{wall\ =\ 2;}.

\textbf{Module Comments} --- every module (a named, reusable chunk of geometry --- covered in depth in Lesson 3) should have a one-line description of what shape it creates.

\subsubsection*{A Well-Documented File Template}\label{docs__pandoc__latex__src__3dmake_all_lessons.md__a-well-documented-file-template}

\begin{lstlisting}[style=Alabaster, language=openscad]
// ==============================================================

// Parametric Box Lid

// ==============================================================

// Purpose:     Snap-fit lid for the storage box (box.scad)

// Print time:  ~8 minutes at 0.20 mm layers

// Material:    ~6 g PLA

//

// PARAMETERS TO CUSTOMIZE:

// - box_width:   width of the matching box (default: 60 mm)

// - box_depth:   depth of the matching box (default: 40 mm)

// - lip_height:  how deep the lid lip fits into the box (default: 3 mm)

//

// PRINT RECOMMENDATIONS:

// - Layer height: 0.20 mm

// - Infill:       15%

// - Supports:     not needed

// - Orientation:  print flat, face-down

// ==============================================================


box_width  = 60;   // mm — must match box.scad box_width

box_depth  = 40;   // mm — must match box.scad box_depth

lip_height = 3;    // mm — depth the lid lip extends downward (2–5 mm)

wall       = 2;    // mm — lid wall thickness (1.5–3.0 mm)


// (Geometry below this line — uses the parameters above)


\end{lstlisting}

This header comment takes about two minutes to write and saves many minutes of confusion later. Make it a habit starting with your very first project.

\begin{center}\rule{0.5\linewidth}{0.5pt}\end{center}

\subsection*{Summary Table}\label{docs__pandoc__latex__src__3dmake_all_lessons.md__summary-table}

{\def\LTcaptype{none} % do not increment counter
\begin{longtable}[]{@{}
  >{\raggedright\arraybackslash}p{(\linewidth - 2\tabcolsep) * \real{0.1848}}
  >{\raggedright\arraybackslash}p{(\linewidth - 2\tabcolsep) * \real{0.8152}}@{}}
\toprule\noalign{}
\begin{minipage}[b]{\linewidth}\raggedright
Concept
\end{minipage} & \begin{minipage}[b]{\linewidth}\raggedright
Key Point
\end{minipage} \\
\midrule\noalign{}
\endhead
\bottomrule\noalign{}
\endlastfoot
Parametric design &
Change one number, model updates everywhere it\textquotesingle s used \\
OpenSCAD &
Code-based 3D geometry engine; you describe shapes with text \\
3dMake &
Project manager layer; \texttt{3dm\ build} and \texttt{3dm\ info} are your main commands \\
Units & Always millimeters --- no exceptions \\
Origin (0,0,0) & Front-left corner of build space, at bed level \\
X axis & Points right \\
Y axis & Points away from you (back) \\
Z axis & Points straight up \\
\texttt{center=false} &
Default --- one corner of the object sits at the specified position \\
\texttt{center=true} &
Object straddles the specified position (center is at that point) \\
\texttt{3dm\ build} &
Compiles \texttt{src/main.scad} → \texttt{build/main.stl} \\
\texttt{3dm\ info} &
Reports bounding box, volume, triangle count, AI description \\
\texttt{3dm\ watch} &
Auto-rebuilds whenever you save a \texttt{.scad} file \\
Bounding box &
Smallest rectangle containing the whole model --- your \#1 verification metric \\
STL file & Triangle-mesh representation of surface geometry \\
FDM pipeline & STL → Slicer → G-code → Print → Cool \\
Layer height &
Thickness of each printed layer; 0.20 mm is a good default \\
Infill &
Internal density; 15--20\% for typical parts, 40--50\% for structural \\
Supports &
Temporary scaffolding for overhangs; minimize them in your design \\
\end{longtable}
}

\begin{center}\rule{0.5\linewidth}{0.5pt}\end{center}

\subsection*{Extension Resources}\label{docs__pandoc__latex__src__3dmake_all_lessons.md__extension-resources}

\subsubsection*{Lesson 1 Asset Folder}\label{docs__pandoc__latex__src__3dmake_all_lessons.md__lesson-1-asset-folder}

\begin{itemize}
\tightlist
\item
  \href{docs\%5Cpandoc\%5Clatex\%5Csrc\%5C3dMake_Foundation\%5CTemplates\%5CStudent\%5Cyour_first_print_student_template.md}{Your First Print guide}
\item
  \href{docs\%5Cpandoc\%5Clatex\%5Csrc\%5Cassets\%5C3dMake_Foundation\%5CLessons_3dMake_1\%5Cbasic_project_scaffold.scad}{Basic project scaffold template}
\end{itemize}

\subsubsection*{\texorpdfstring{Learning Series Sample Projects (\texttt{3dmake\_learning\_series/})}{Learning Series Sample Projects (3dmake\_learning\_series/)}}\label{docs__pandoc__latex__src__3dmake_all_lessons.md__learning-series-sample-projects-3dmake_learning_series}

\begin{itemize}
\tightlist
\item
  \texttt{01\_cube\_keycap} (Beginner) --- Text embossing basics
\item
  \texttt{02\_parametric\_phone\_stand} (Intermediate) --- Transforms and Minkowski fillets
\item
  \texttt{03\_stackable\_bins} (Advanced) --- Tolerance and assemblies
\end{itemize}

\subsubsection*{\texorpdfstring{Reference Materials (\texttt{Reference\_Materials/})}{Reference Materials (Reference\_Materials/)}}\label{docs__pandoc__latex__src__3dmake_all_lessons.md__reference-materials-reference_materials}

\begin{itemize}
\tightlist
\item
  \texttt{3dmake-setup-guide.md} --- Complete setup walkthrough
\item
  \texttt{openscad-cheat-sheet.md} --- Keyboard shortcuts and common functions
\item
  \texttt{filament-comparison-table.md} --- Material properties reference
\item
  \texttt{master-rubric.md} --- Assessment criteria
\end{itemize}

\begin{center}\rule{0.5\linewidth}{0.5pt}\end{center}

\subsection*{Quiz --- Lesson 1 (15 questions)}\label{docs__pandoc__latex__src__3dmake_all_lessons.md__quiz--lesson-1-15-questions}

\begin{enumerate}
\tightlist
\item
  What command initializes a 3dMake project?
\item
  What folder holds generated STL files?
\item
  How do you run 3dMake\textquotesingle s automatic rebuild mode?
\item
  Why is it useful to run \texttt{3dm\ build} frequently during development rather than only at the end?
\item
  Give one reason to prefer an external text editor over the built-in OpenSCAD editor for this workflow.
\item
  True or False: 3dMake requires a graphical user interface to use effectively.
\item
  Explain what the \texttt{3dmake.toml} file does in your project.
\item
  Describe what the \texttt{src/} and \texttt{build/} project folders are used for and which one you edit directly.
\item
  In your own words, explain the difference between a visual 3D editor (like Tinkercad) and a code-based parametric design tool (like OpenSCAD). What is the main advantage of the parametric approach?
\item
  What validation step should you always perform after running \texttt{3dm\ build} before sending a file to print?
\item
  If the bounding box reported by \texttt{3dm\ info} shows \texttt{0\ ×\ 0\ ×\ 0\ mm}, what does that tell you, and what should you do next?
\item
  Explain what \texttt{center=false} and \texttt{center=true} mean for a \texttt{cube()}. For each, where does the object appear relative to the position you specify?
\item
  What does FDM stand for, and what are the five stages of the FDM pipeline from STL file to cooled physical part?
\item
  A box has dimensions 50 × 40 × 25 mm and you\textquotesingle re printing it in PLA (density 1.24 g/cm³) at 100\% infill. Using a \$22 / 1 kg spool, what is the estimated material cost? Show your calculation.
\item
  You run \texttt{3dm\ build} and receive the error "No such file or directory: src/main.scad". List two likely causes and the fix for each.
\end{enumerate}

\begin{center}\rule{0.5\linewidth}{0.5pt}\end{center}

\subsection*{Extension Problems (15)}\label{docs__pandoc__latex__src__3dmake_all_lessons.md__extension-problems-15}

\begin{enumerate}
\tightlist
\item
  Add a README entry explaining your top-level parameters and expected units.
\item
  Create a parameter variant by changing \texttt{width} by 20\% and build both variants; compare dimensions using \texttt{3dm\ info} to confirm the change.
\item
  Script a command sequence that automates the entire new-project setup: create folder, create \texttt{src/main.scad}, and run \texttt{3dm\ build} --- in a single script file.
\item
  Intentionally introduce a syntax error in your \texttt{.scad} file (remove a semicolon), run \texttt{3dm\ build}, and document what the error message says and which line it points to. Then fix it.
\item
  Prepare a short instructor sign-off checklist describing the safety checks you should complete before sending a file to print.
\item
  Build a variant testing suite: create 5 different combinations of \texttt{width}, \texttt{depth}, and \texttt{height}, build all five, and record the bounding box and volume for each.
\item
  Create a 3dMake project template with best-practice structure, documentation header, and parameter comments. Write a README explaining how to use it.
\item
  Write a short accessibility guide explaining how a student using a screen reader could use \texttt{3dm\ info} output to understand a model\textquotesingle s geometry without a visual preview.
\item
  Design a simple parametric part library (three or more related parts) in a single project. Document all parameters with units and example values.
\item
  Write a comprehensive troubleshooting guide for common \texttt{3dm\ build} errors, with the cause and fix for each.
\item
  Add a \texttt{.gitignore} file to your project, commit it using Git, make a parameter change and create a second commit. In 2--3 sentences, explain why version control is valuable for parametric design.
\item
  Edit the global configuration file using \texttt{3dm\ edit-global-config} to change the default editor. Document the setting name and how you verified the change took effect.
\item
  Create a project with three \texttt{.scad} source files and build each using \texttt{3dm\ build\ -m\ \textless{}n\textgreater{}}. Record the output STL filenames.
\item
  Research the FDM layer height trade-off: slice the same model at 0.15 mm, 0.20 mm, and 0.30 mm. Record the estimated print time and filament use for each. What is the trade-off?
\item
  Write a one-page "new student onboarding guide" for 3dMake covering installation, project creation, first build, and first print. Write it as if your audience is a classmate who has never used the command line.
\end{enumerate}

\begin{center}\rule{0.5\linewidth}{0.5pt}\end{center}

\subsection*{References and Helpful Resources}\label{docs__pandoc__latex__src__3dmake_all_lessons.md__references-and-helpful-resources}

\begin{itemize}
\item
  3DMake GitHub Repository --- https://github.com/tdeck/3dmake
\item
  OpenSCAD Manual --- https://en.wikibooks.org/wiki/OpenSCAD\_User\_Manual
\item
  OpenSCAD Parametric Design --- https://en.wikibooks.org/wiki/OpenSCAD\_User\_Manual/The\_OpenSCAD\_Language\#Variables
\item
  OpenSCAD Review --- CadHub --- https://learn.cadhub.xyz/blog/openscad-review/
\item
  PrusaSlicer Validation Tools --- https://docs.prusa3d.com/en/guide/39012-validation-tools/
\item
  3dmake e2e\_test.py --- https://github.com/tdeck/3dmake/blob/main/e2e\_test.py
\item
  PrusaSlicer Documentation --- https://docs.prusa3d.com/en/
\item
  OpenSCAD Non-Manifold FAQ --- https://en.wikibooks.org/wiki/OpenSCAD\_User\_Manual/FAQ\#Why\_is\_my\_model\_not\_manifold.3F
\end{itemize}

\subsubsection*{Supplemental Resources}\label{docs__pandoc__latex__src__3dmake_all_lessons.md__supplemental-resources}

\begin{itemize}
\tightlist
\item
  Programming with OpenSCAD EPUB Textbook --- assets/Programming\_with\_OpenSCAD.epub
\item
  CodeSolutions Repository --- https://github.com/ProgrammingWithOpenSCAD/CodeSolutions
\item
  OpenSCAD Quick Reference --- https://programmingwithopenscad.github.io/quick-reference.html
\item
  All other reference documents are embedded as appendices below --- no separate files needed.
\end{itemize}

\begin{center}\rule{0.5\linewidth}{0.5pt}\end{center}

\section*{Part 2: Extension Project and Reference Appendices}\label{docs__pandoc__latex__src__3dmake_all_lessons.md__part-2-extension-project-and-reference-appendices}

\emph{All supporting documents for Lesson 1 are included below. You do not need to open any additional files.}

\begin{center}\rule{0.5\linewidth}{0.5pt}\end{center}

\subsection*{Appendix A: Your First Print --- Extension Project}\label{docs__pandoc__latex__src__3dmake_all_lessons.md__appendix-a-your-first-print--extension-project}

\textbf{Estimated time:} 2--4 hours (including setup and print monitoring)

\subsubsection*{Overview}\label{docs__pandoc__latex__src__3dmake_all_lessons.md__overview-1}

This extension project applies the workflow and concepts from Lesson 1 to a real physical print. Rather than designing from scratch, you\textquotesingle ll select a simple ready-made model, evaluate whether it\textquotesingle s appropriate for a beginner print, configure slicer settings, run the print safely, and document what happened. The goal is to complete one full cycle of the pipeline you just learned --- from file to printed object --- and to practice the observation and documentation habits you\textquotesingle ll use for every print in this course.

\subsubsection*{Learning Objectives}\label{docs__pandoc__latex__src__3dmake_all_lessons.md__learning-objectives-1}

By completing this project, you will:

\begin{itemize}
\tightlist
\item
  Select a simple ready-made model and evaluate its printability for a classroom printer
\item
  Configure slicer settings for a short-duration print and prepare the printer safely
\item
  Monitor a print in progress and recognize early warning signs of failure
\item
  Document print parameters and reflect on the physical outcome
\end{itemize}

\subsubsection*{Materials}\label{docs__pandoc__latex__src__3dmake_all_lessons.md__materials}

\begin{itemize}
\tightlist
\item
  Computer with slicer and access to an online model repository
\item
  Classroom-approved printer, filament spool
\end{itemize}

\subsubsection*{Tasks}\label{docs__pandoc__latex__src__3dmake_all_lessons.md__tasks}

\begin{enumerate}
\item
  \textbf{Choose a simple model} (under 2 hours estimated print time) from Thingiverse or Printables. Download the STL.
\item
  \textbf{Inspect the model} before loading it in the slicer. Look at it from all angles. Note any overhangs, thin features, or very small details. Write two short reasons why this model is appropriate for a first print --- specifically, why it is not risky.
\item
  \textbf{Load the model in your slicer}, select the classroom profile, and adjust settings only if necessary. Record the final print time estimate, layer height, infill \%, and filament weight estimate.
\item
  \textbf{Perform safety checks} using the checklist in Lesson 5 Appendix A (or the condensed checklist below). Start the print and \textbf{remain present and watching for the first 15 minutes}. Note what you observe about first-layer adhesion and extrusion quality.
\item
  \textbf{After the print cools}, measure three critical dimensions using calipers (or by comparing against the model\textquotesingle s stated dimensions). Record the measured values and any deviations.
\end{enumerate}

\subsubsection*{Condensed Pre-Print Safety Checklist}\label{docs__pandoc__latex__src__3dmake_all_lessons.md__condensed-pre-print-safety-checklist}

\begin{itemize}
\tightlist
\item[$\square$]
  Bed clean and level
\item[$\square$]
  Filament loaded and extruding cleanly (run a short purge)
\item[$\square$]
  Build area clear of obstructions
\item[$\square$]
  Print area free of flammable materials
\item[$\square$]
  You will remain present for the first 15 minutes
\end{itemize}

\subsubsection*{Probing Questions}\label{docs__pandoc__latex__src__3dmake_all_lessons.md__probing-questions}

\begin{enumerate}
\tightlist
\item
  Why did you select this model? What specific risks did you anticipate, and how does this model avoid or minimize them?
\item
  Which slicer setting most affects print time for this model, and why?
\item
  If a thin feature failed during printing, what is the minimum change you would make to give it a better chance on the next attempt?
\end{enumerate}

\subsubsection*{Starter Code}\label{docs__pandoc__latex__src__3dmake_all_lessons.md__starter-code}

If you want to build a simple printable object from scratch rather than downloading one, use this scaffold as your starting point:

\begin{lstlisting}[style=Alabaster, language=openscad]
// Basic Project Scaffold

// ============================================

// PROJECT CONFIGURATION

// ============================================

// Set your project dimensions here. All other geometry derives from these.


object_w  = 50;   // mm — overall width

object_d  = 50;   // mm — overall depth

object_h  = 20;   // mm — overall height

wall      = 2.0;  // mm — wall thickness (2–3mm recommended for FDM)

$fn       = 30;   // curve resolution (higher = smoother, slower to render)


// ============================================

// MAIN DESIGN

// ============================================


module base_shape() {

  // Replace this cube with your own geometry

  cube([object_w, object_d, object_h]);

}


module hollow_version() {

  // Creates a hollow version by removing an inner volume

  difference() {

    base_shape();

    translate([wall, wall, wall])

      cube([object_w - 2*wall,

            object_d - 2*wall,

            object_h]);  // open top — no overhang needed

  }

}


// Render your design:

hollow_version();


\end{lstlisting}

\subsubsection*{Deliverables}\label{docs__pandoc__latex__src__3dmake_all_lessons.md__deliverables}

\begin{itemize}
\tightlist
\item
  Short report: model selected, key slicer settings, measured dimensions and deviations, answers to probing questions
\item
  Photos of the final print (at least two angles)
\item
  Completed student documentation template (Appendix B)
\end{itemize}

\begin{center}\rule{0.5\linewidth}{0.5pt}\end{center}

\subsection*{Appendix B: Your First Print --- Student Documentation Template}\label{docs__pandoc__latex__src__3dmake_all_lessons.md__appendix-b-your-first-print--student-documentation-template}

\textbf{Author:} \_\_\_\_\_\_\_\_\_\_\_\_\_\_\_\_
\textbf{Date:} \_\_\_\_\_\_\_\_\_\_
\textbf{Description:} Select a simple 3D model, configure slicer settings, and complete your first independent print job.

\begin{center}\rule{0.5\linewidth}{0.5pt}\end{center}

\subsubsection*{Model Selection}\label{docs__pandoc__latex__src__3dmake_all_lessons.md__model-selection}

\begin{itemize}
\tightlist
\item
  Model name and source (Thingiverse, Printables, etc.):
\item
  Link or file name:
\item
  Why did you choose this model?
\item
  Why is this model appropriate for a first print (two specific reasons)?
\end{itemize}

\begin{center}\rule{0.5\linewidth}{0.5pt}\end{center}

\subsubsection*{Printability Assessment}\label{docs__pandoc__latex__src__3dmake_all_lessons.md__printability-assessment}

\begin{itemize}
\tightlist
\item
  Estimated print time (from slicer): \_\_\_\_ hours \_\_\_\_ min
\item
  Estimated filament: \_\_\_\_ g
\item
  Layer height selected: \_\_\_\_ mm
\item
  Infill \%: \_\_\_\_
\item
  Supports required? Yes / No --- if yes, describe where:
\item
  Any overhangs or thin features that concerned you?
\end{itemize}

\begin{center}\rule{0.5\linewidth}{0.5pt}\end{center}

\subsubsection*{Printer Setup Log}\label{docs__pandoc__latex__src__3dmake_all_lessons.md__printer-setup-log}

{\def\LTcaptype{none} % do not increment counter
\begin{longtable}[]{@{}ll@{}}
\toprule\noalign{}
Parameter & Value \\
\midrule\noalign{}
\endhead
\bottomrule\noalign{}
\endlastfoot
Printer model & \\
Nozzle diameter & mm \\
Nozzle temperature & °C \\
Bed temperature & °C \\
Layer height & mm \\
Infill \% & \\
Support settings & \\
Print speed & mm/s \\
\end{longtable}
}

\begin{center}\rule{0.5\linewidth}{0.5pt}\end{center}

\subsubsection*{Safety Checks}\label{docs__pandoc__latex__src__3dmake_all_lessons.md__safety-checks}

\begin{itemize}
\tightlist
\item[$\square$]
  Bed clean and level
\item[$\square$]
  Filament loaded and purged cleanly
\item[$\square$]
  Nozzle at correct first-layer height
\item[$\square$]
  Print area clear of obstructions
\item[$\square$]
  I remained present and watching for the first 15 minutes
\end{itemize}

\begin{center}\rule{0.5\linewidth}{0.5pt}\end{center}

\subsubsection*{Print Monitoring Notes}\label{docs__pandoc__latex__src__3dmake_all_lessons.md__print-monitoring-notes}

\begin{itemize}
\tightlist
\item
  What did you observe in the first layer?
\item
  Was first-layer adhesion good, marginal, or poor?
\item
  Any extrusion issues observed (clicking, skipping, thin lines)?
\item
  Print completed successfully? Yes / No / Partially
\item
  If partially or no --- what happened and at what layer?
\end{itemize}

\begin{center}\rule{0.5\linewidth}{0.5pt}\end{center}

\subsubsection*{Dimensional Verification (After Cooling)}\label{docs__pandoc__latex__src__3dmake_all_lessons.md__dimensional-verification-after-cooling}

Measure three features on the printed part and compare to the model\textquotesingle s stated dimensions:

{\def\LTcaptype{none} % do not increment counter
\begin{longtable}[]{@{}lcccl@{}}
\toprule\noalign{}
Feature & Model Spec (mm) & Measured (mm) & Deviation (mm) & Notes \\
\midrule\noalign{}
\endhead
\bottomrule\noalign{}
\endlastfoot
& & & & \\
& & & & \\
& & & & \\
\end{longtable}
}

\begin{center}\rule{0.5\linewidth}{0.5pt}\end{center}

\subsubsection*{Reflections}\label{docs__pandoc__latex__src__3dmake_all_lessons.md__reflections}

\textbf{What went well:}

\textbf{What didn\textquotesingle t go well:}

\textbf{What would you do differently next time:}

\textbf{If you printed this again, what settings would you change and why?}

\begin{center}\rule{0.5\linewidth}{0.5pt}\end{center}

\subsubsection*{Attachments Checklist}\label{docs__pandoc__latex__src__3dmake_all_lessons.md__attachments-checklist}

\begin{itemize}
\tightlist
\item[$\square$]
  STL file or link to source model
\item[$\square$]
  Slicer settings (screenshot or exported settings file)
\item[$\square$]
  Photo of final print --- at least two angles
\item[$\square$]
  Completed measurement table above
\item[$\square$]
  Answers to probing questions
\end{itemize}

\begin{center}\rule{0.5\linewidth}{0.5pt}\end{center}

\subsubsection*{Teacher Feedback}\label{docs__pandoc__latex__src__3dmake_all_lessons.md__teacher-feedback}

{\def\LTcaptype{none} % do not increment counter
\begin{longtable}[]{@{}
  >{\raggedright\arraybackslash}p{(\linewidth - 4\tabcolsep) * \real{0.8280}}
  >{\centering\arraybackslash}p{(\linewidth - 4\tabcolsep) * \real{0.1183}}
  >{\raggedright\arraybackslash}p{(\linewidth - 4\tabcolsep) * \real{0.0538}}@{}}
\toprule\noalign{}
\begin{minipage}[b]{\linewidth}\raggedright
Category
\end{minipage} & \begin{minipage}[b]{\linewidth}\centering
Score (0--3)
\end{minipage} & \begin{minipage}[b]{\linewidth}\raggedright
Notes
\end{minipage} \\
\midrule\noalign{}
\endhead
\bottomrule\noalign{}
\endlastfoot
Problem \& Solution --- appropriate model selected; print completed &
& \\
Setup \& Documentation --- printer setup log complete; safety checks done
& & \\
Reflection --- measurements recorded; reflection is specific and shows learning
& & \\
\textbf{Total (0--9)} & & \\
\end{longtable}
}

\textbf{Resubmission:} If resubmitting, attach a paragraph explaining what was changed and why. The resubmission score replaces the original.

\begin{center}\rule{0.5\linewidth}{0.5pt}\end{center}

\subsection*{Appendix C: Your First Print --- Instructor Reference}\label{docs__pandoc__latex__src__3dmake_all_lessons.md__appendix-c-your-first-print--instructor-reference}

\subsubsection*{Project Context}\label{docs__pandoc__latex__src__3dmake_all_lessons.md__project-context}

This is the first physical print extension project in the course. Its primary purpose is to establish safe printer operation habits and complete documentation practices --- not to produce a complex or impressive object. A student who selects an appropriate model, monitors the print safely, and documents the outcome accurately has succeeded, regardless of whether the print is visually impressive.

\subsubsection*{Key Learning Points to Reinforce}\label{docs__pandoc__latex__src__3dmake_all_lessons.md__key-learning-points-to-reinforce}

\textbf{Model selection discipline.} The most common mistake at this stage is selecting a model that is too complex, too large, or has features that will fail (overhangs without supports, very thin walls, very fine detail). Push students to justify their selection in terms of printability, not just aesthetics.

Important

\textbf{The first 15 minutes rule:} students must stay and watch the printer for the first 15 minutes of every print. A print that lifts off the bed in the first few layers becomes a tangled failure --- you must be present to catch it early.

\textbf{Documentation enables iteration.} A student who records their slicer settings has something to start from on the next print. A student who doesn\textquotesingle t has to guess. Reinforce this constantly.

\subsubsection*{Constraints (Enforce These)}\label{docs__pandoc__latex__src__3dmake_all_lessons.md__constraints-enforce-these}

\begin{itemize}
\tightlist
\item
  Model must be under 2 hours estimated print time
\item
  Minimal or no support requirements recommended
\item
  Student must record slicer settings, not just print the file
\item
  Dimensional measurements must be taken after the part has fully cooled
\end{itemize}

\subsubsection*{Assessment Notes}\label{docs__pandoc__latex__src__3dmake_all_lessons.md__assessment-notes}

\textbf{Strong submissions} show: clear model selection reasoning (not just "it looked cool"), a complete printer setup log, specific first-layer observations, measurements that reflect actual caliper readings (not round numbers), and a reflection that identifies something specific to do differently.

\textbf{Common weak areas:} Vague reflections ("it went well"), missing measurements, and skipped safety checks. Address all three in feedback.

\begin{center}\rule{0.5\linewidth}{0.5pt}\end{center}

\subsection*{Appendix D: 3dMake Setup and Workflow Guide}\label{docs__pandoc__latex__src__3dmake_all_lessons.md__appendix-d-3dmake-setup-and-workflow-guide}

\subsubsection*{Installing 3dMake}\label{docs__pandoc__latex__src__3dmake_all_lessons.md__installing-3dmake}

\paragraph*{Prerequisites}\label{docs__pandoc__latex__src__3dmake_all_lessons.md__prerequisites}

Before installing 3dMake, you need:

\begin{itemize}
\tightlist
\item
  \textbf{Node.js} version 18 or higher --- download from https://nodejs.org
\item
  \textbf{OpenSCAD} --- download from https://openscad.org/downloads.html
\end{itemize}

To verify both are installed, open a terminal (PowerShell on Windows, Terminal on macOS/Linux) and run:

\begin{lstlisting}[style=Alabaster, language=bash]
node --version    # should print v18.0.0 or higher

openscad --version  # should print a version number


\end{lstlisting}

If either command returns "not found," install that software first before continuing.

\paragraph*{Installing 3dMake}\label{docs__pandoc__latex__src__3dmake_all_lessons.md__installing-3dmake-1}

\begin{lstlisting}[style=Alabaster, language=bash]
npm install -g 3dmake


\end{lstlisting}

Verify the installation:

\begin{lstlisting}[style=Alabaster, language=bash]
3dm --version


\end{lstlisting}

You should see a version number. If you get "command not found," check that npm\textquotesingle s global bin directory is in your PATH (see Troubleshooting below).

\begin{center}\rule{0.5\linewidth}{0.5pt}\end{center}

\subsubsection*{Creating Your First Project}\label{docs__pandoc__latex__src__3dmake_all_lessons.md__creating-your-first-project}

\begin{lstlisting}[style=Alabaster, language=bash]
# Create a new folder for your project and enter it

mkdir my_project

cd my_project


# Initialize a 3dMake project

3dm init


\end{lstlisting}

\texttt{3dm\ init} creates this project structure:

\begin{lstlisting}[style=Alabaster]
my_project/

├── src/

│   └── main.scad      ← Your design code goes here

├── build/             ← Generated STL files appear here (auto-created)

└── 3dmake.toml        ← Project configuration


\end{lstlisting}

\textbf{The \texttt{src/} folder} is where you write your OpenSCAD code. You should never need to edit files in \texttt{build/} directly --- that folder is managed by 3dMake.

\textbf{\texttt{3dmake.toml}} stores your project settings: the slicer path, filament density, cost-per-gram, and other project-level configuration. Open it with any text editor to see the available settings.

\begin{center}\rule{0.5\linewidth}{0.5pt}\end{center}

\subsubsection*{The Core Workflow}\label{docs__pandoc__latex__src__3dmake_all_lessons.md__the-core-workflow}

Every design session follows this cycle:

\begin{lstlisting}[style=Alabaster]
edit src/main.scad  →  3dm build  →  3dm info  →  review / repeat


\end{lstlisting}

\textbf{\texttt{3dm\ build}} --- compiles your \texttt{.scad} source into an STL file in the \texttt{build/} folder.

\begin{lstlisting}[style=Alabaster, language=bash]
3dm build              # standard build

3dm build --clean      # delete build/ first, then rebuild

3dm build --watch      # auto-rebuild whenever you save src/main.scad


\end{lstlisting}

\textbf{\texttt{3dm\ info}} --- reports on the most recently built STL: bounding box, volume, triangle count, and an AI-generated description.

\begin{lstlisting}[style=Alabaster, language=bash]
3dm info

3dm info --view front

3dm info --view top,front,isometric


\end{lstlisting}

\textbf{\texttt{3dm\ preview}} --- renders an image of your model to \texttt{build/preview.png}.

\textbf{\texttt{3dm\ orient}} --- recommends a print orientation based on your model\textquotesingle s geometry.

\textbf{\texttt{3dm\ slice}} --- calls your configured slicer on the current STL.

\begin{center}\rule{0.5\linewidth}{0.5pt}\end{center}

\subsubsection*{Command Reference}\label{docs__pandoc__latex__src__3dmake_all_lessons.md__command-reference}

{\def\LTcaptype{none} % do not increment counter
\begin{longtable}[]{@{}
  >{\raggedright\arraybackslash}p{(\linewidth - 4\tabcolsep) * \real{0.2400}}
  >{\raggedright\arraybackslash}p{(\linewidth - 4\tabcolsep) * \real{0.3600}}
  >{\raggedright\arraybackslash}p{(\linewidth - 4\tabcolsep) * \real{0.4000}}@{}}
\toprule\noalign{}
\begin{minipage}[b]{\linewidth}\raggedright
Command
\end{minipage} & \begin{minipage}[b]{\linewidth}\raggedright
Purpose
\end{minipage} & \begin{minipage}[b]{\linewidth}\raggedright
Notes
\end{minipage} \\
\midrule\noalign{}
\endhead
\bottomrule\noalign{}
\endlastfoot
\texttt{3dm\ init} & Initialize a new project & Run once per project \\
\texttt{3dm\ build} & Compile source → STL &
Run after every code change \\
\texttt{3dm\ info} & Report dimensions and AI description &
Always run before sending to print \\
\texttt{3dm\ preview} & Render model image &
Use for quick visual check \\
\texttt{3dm\ orient} & AI print orientation suggestion &
Use as a starting point, verify yourself \\
\texttt{3dm\ slice} & Run slicer on current build &
Requires slicer path in \texttt{3dmake.toml} \\
\texttt{3dm\ edit-global-config} & Edit the global settings file &
For changing default editor, slicer path \\
\end{longtable}
}

\begin{center}\rule{0.5\linewidth}{0.5pt}\end{center}

\subsubsection*{Configuring Your Slicer}\label{docs__pandoc__latex__src__3dmake_all_lessons.md__configuring-your-slicer}

Edit \texttt{3dmake.toml} in your project folder and add your slicer path:

\textbf{Windows (PrusaSlicer):}

\begin{lstlisting}[style=Alabaster]
[slicer]

path = "C:/Program Files/Prusa3D/PrusaSlicer/prusa-slicer.exe"


\end{lstlisting}

\textbf{macOS (PrusaSlicer):}

\begin{lstlisting}[style=Alabaster]
[slicer]

path = "/Applications/PrusaSlicer.app/Contents/MacOS/PrusaSlicer"


\end{lstlisting}

\textbf{Linux:}

\begin{lstlisting}[style=Alabaster]
[slicer]

path = "/usr/bin/prusa-slicer"


\end{lstlisting}

\begin{center}\rule{0.5\linewidth}{0.5pt}\end{center}

\subsubsection*{Configuring Cost Estimation}\label{docs__pandoc__latex__src__3dmake_all_lessons.md__configuring-cost-estimation}

To use the cost estimation feature, set your filament density and spool price:

\begin{lstlisting}[style=Alabaster]
[filament]

density = 1.24        # g/cm³ — PLA default (see Lesson 5 Appendix C for other materials)

spool_weight = 1000   # grams per spool

spool_price = 22.00   # dollars per spool


\end{lstlisting}

3dMake will then include cost estimates in \texttt{3dm\ info} output.

\begin{center}\rule{0.5\linewidth}{0.5pt}\end{center}

\subsubsection*{Troubleshooting}\label{docs__pandoc__latex__src__3dmake_all_lessons.md__troubleshooting}

\textbf{\texttt{3dm:\ command\ not\ found} after installation}

npm\textquotesingle s global bin directory isn\textquotesingle t in your PATH. Find the directory:

\begin{lstlisting}[style=Alabaster, language=bash]
npm config get prefix


\end{lstlisting}

Add the \texttt{bin} subdirectory of that path to your PATH. On Windows, this is usually \texttt{C:\textbackslash{}Users\textbackslash{}YourName\textbackslash{}AppData\textbackslash{}Roaming\textbackslash{}npm}. On macOS/Linux, it\textquotesingle s usually \texttt{\textasciitilde{}/.npm-global/bin} or \texttt{/usr/local/bin}.

\textbf{\texttt{No\ such\ file\ or\ directory:\ src/main.scad}}

You either haven\textquotesingle t run \texttt{3dm\ init}, or you\textquotesingle re in the wrong folder. Run \texttt{pwd} to confirm your current location, and \texttt{ls} (or \texttt{dir} on Windows) to see if \texttt{src/} exists.

\textbf{\texttt{OpenSCAD:\ command\ not\ found} during build}

OpenSCAD isn\textquotesingle t in your PATH. On Windows, find where OpenSCAD is installed (usually \texttt{C:\textbackslash{}Program\ Files\textbackslash{}OpenSCAD\textbackslash{}}) and add it to your PATH in System Settings → Environment Variables.

\textbf{Build succeeds but \texttt{3dm\ info} shows \texttt{0\ ×\ 0\ ×\ 0\ mm}}

Your OpenSCAD code compiled without errors but produced no geometry. Check that your shapes are actually being rendered (not commented out), and that no \texttt{difference()} has subtracted more than it started with.

\textbf{\texttt{WARNING:\ Normalized\ tree\ is\ empty}}

Same as above from OpenSCAD\textquotesingle s perspective. The most common causes: a shape with zero volume, all geometry inside a \texttt{difference()} was subtracted away, or a module was defined but never called.

\begin{center}\rule{0.5\linewidth}{0.5pt}\end{center}

\subsection*{Appendix E: Navigating This Curriculum (mdBook Guide)}\label{docs__pandoc__latex__src__3dmake_all_lessons.md__appendix-e-navigating-this-curriculum-mdbook-guide}

\emph{This curriculum is published as a web book using mdBook. This guide explains how to find what you need, navigate between chapters, and use the book with a screen reader.}

\subsubsection*{Basic Navigation}\label{docs__pandoc__latex__src__3dmake_all_lessons.md__basic-navigation}

\textbf{Sidebar table of contents} --- the left side of the page lists all chapters by unit and lesson. Click or tap any title to jump to it. On a small screen, the sidebar may be hidden --- look for a hamburger menu icon (three horizontal lines) to show it.

\textbf{Arrow navigation} --- at the bottom of every page are Previous and Next links. You can also use keyboard arrow keys:

\begin{itemize}
\tightlist
\item
  \texttt{Left\ Arrow} --- go to the previous chapter
\item
  \texttt{Right\ Arrow} --- go to the next chapter
\end{itemize}

\textbf{Search} --- press \texttt{S} to open the search box. Type your search term; results appear as a dropdown. Click any result to jump to that chapter. Press \texttt{Escape} to close search.

{\def\LTcaptype{none} % do not increment counter
\begin{longtable}[]{@{}ll@{}}
\toprule\noalign{}
Key & Action \\
\midrule\noalign{}
\endhead
\bottomrule\noalign{}
\endlastfoot
\texttt{S} & Focus the search box \\
\texttt{Escape} & Close search results \\
\texttt{Left\ Arrow} & Previous chapter \\
\texttt{Right\ Arrow} & Next chapter \\
\texttt{T} & Toggle the table of contents sidebar \\
\end{longtable}
}

\begin{center}\rule{0.5\linewidth}{0.5pt}\end{center}

\subsubsection*{Screen Reader Navigation}\label{docs__pandoc__latex__src__3dmake_all_lessons.md__screen-reader-navigation}

\paragraph*{With NVDA (Chrome, Firefox, or Edge recommended)}\label{docs__pandoc__latex__src__3dmake_all_lessons.md__with-nvda-chrome-firefox-or-edge-recommended}

\begin{itemize}
\tightlist
\item
  \texttt{H} --- jump between headings; the fastest way to skim a long lesson
\item
  \texttt{D} --- move between landmark regions (use to reach the sidebar)
\item
  \texttt{NVDA\ +\ F7} --- open the Elements List; select "Landmarks" to navigate to the sidebar directly
\item
  Set punctuation level to Most or All before reading code blocks: \texttt{NVDA\ +\ P} to cycle
\item
  \texttt{Ctrl\ +\ F} --- browser Find, works alongside mdBook search
\end{itemize}

\paragraph*{With JAWS (Chrome or Edge recommended)}\label{docs__pandoc__latex__src__3dmake_all_lessons.md__with-jaws-chrome-or-edge-recommended}

\begin{itemize}
\tightlist
\item
  \texttt{H} --- jump between headings
\item
  \texttt{R} --- move between landmark regions to reach the sidebar
\item
  \texttt{JAWS\ Key\ +\ F6} --- list all headings on the current page
\item
  Set punctuation to All before reading code: \texttt{JAWS\ Key\ +\ Shift\ +\ 2}
\item
  \texttt{JAWS\ Key\ +\ F5} --- links list
\item
  \texttt{Ctrl\ +\ F} --- browser find
\end{itemize}

\paragraph*{With VoiceOver (Mac / iOS)}\label{docs__pandoc__latex__src__3dmake_all_lessons.md__with-voiceover-mac--ios}

\begin{itemize}
\tightlist
\item
  \texttt{VO\ +\ U} --- open the rotor; select Headings to navigate by heading
\item
  \texttt{H} (Quick Nav on) --- move between headings
\item
  iOS: use the rotor (\texttt{two-finger\ rotate}) to set navigation mode to Headings
\end{itemize}

\begin{center}\rule{0.5\linewidth}{0.5pt}\end{center}

\subsubsection*{Finding What You Need}\label{docs__pandoc__latex__src__3dmake_all_lessons.md__finding-what-you-need}

\textbf{If you know the unit or project:} open the table of contents and look for the unit name. Structure: Unit 0 (Foundation lessons) → Unit 1 (Guided projects) → Unit 2 (Intermediate skills) → Unit 3 (Open-ended projects) → Reference Materials.

\textbf{If you\textquotesingle re looking for a specific term or command:} use Search (\texttt{S}). Search for an OpenSCAD command (\texttt{difference}, \texttt{translate}), a vocabulary word (\texttt{infill}, \texttt{tolerance}), or a project name.

\textbf{While working:} keep a second browser tab open to the Reference Materials section. Useful pages to bookmark: OpenSCAD Cheat Sheet, Slicing Settings Quick Reference, Filament Comparison Table, Screen Reader Coding Tips.

\begin{center}\rule{0.5\linewidth}{0.5pt}\end{center}

\subsection*{Appendix F: Screen Reader Coding Tips (NVDA and JAWS)}\label{docs__pandoc__latex__src__3dmake_all_lessons.md__appendix-f-screen-reader-coding-tips-nvda-and-jaws}

\subsubsection*{General Principles}\label{docs__pandoc__latex__src__3dmake_all_lessons.md__general-principles}

\textbf{Turn punctuation up.} All OpenSCAD syntax --- semicolons, brackets, parentheses, commas --- is meaningful. If your screen reader skips punctuation, you will miss syntax errors. Set to Most or All before coding.

\textbf{Navigate by line.} Arrow up and down through code one line at a time. Use \texttt{Ctrl\ +\ Left/Right} to move word by word within a line.

\textbf{Use go-to-line.} OpenSCAD errors always give a line number. In VSCode: \texttt{Ctrl\ +\ G}, type the number, press Enter.

\textbf{Comment as you go.} A short comment after a block of code lets you navigate back to it by searching with \texttt{Ctrl\ +\ F}.

\begin{center}\rule{0.5\linewidth}{0.5pt}\end{center}

\subsubsection*{NVDA Quick Reference}\label{docs__pandoc__latex__src__3dmake_all_lessons.md__nvda-quick-reference}

\textbf{Setting punctuation level:} \texttt{NVDA\ +\ P} cycles through None → Some → Most → All. For code: Most or All.

{\def\LTcaptype{none} % do not increment counter
\begin{longtable}[]{@{}ll@{}}
\toprule\noalign{}
Action & Keys \\
\midrule\noalign{}
\endhead
\bottomrule\noalign{}
\endlastfoot
Read current line & \texttt{NVDA\ +\ Up\ Arrow} \\
Spell current line & \texttt{NVDA\ +\ Up\ Arrow} (twice quickly) \\
Read from cursor & \texttt{NVDA\ +\ Down\ Arrow} \\
Read current word & \texttt{NVDA\ +\ Numpad\ 5} \\
Spell current word & \texttt{NVDA\ +\ Numpad\ 5} (twice quickly) \\
Stop reading & \texttt{Ctrl} \\
\end{longtable}
}

\textbf{Navigation in VSCode with NVDA:}

{\def\LTcaptype{none} % do not increment counter
\begin{longtable}[]{@{}ll@{}}
\toprule\noalign{}
Action & Keys \\
\midrule\noalign{}
\endhead
\bottomrule\noalign{}
\endlastfoot
Move by character & \texttt{Left\ /\ Right\ Arrow} \\
Move by word & \texttt{Ctrl\ +\ Left\ /\ Right\ Arrow} \\
Move by line & \texttt{Up\ /\ Down\ Arrow} \\
Start / end of line & \texttt{Home\ /\ End} \\
Go to line number & \texttt{Ctrl\ +\ G}, type number, \texttt{Enter} \\
Find text & \texttt{Ctrl\ +\ F} \\
Toggle line comment & \texttt{Ctrl\ +\ /} \\
\end{longtable}
}

If NVDA stops reading the editor: press \texttt{NVDA\ +\ Space} to toggle between Browse and Application mode. For code editors, you want Application mode.

\begin{center}\rule{0.5\linewidth}{0.5pt}\end{center}

\subsubsection*{JAWS Quick Reference}\label{docs__pandoc__latex__src__3dmake_all_lessons.md__jaws-quick-reference}

\textbf{Setting punctuation level:} \texttt{JAWS\ Key\ +\ Shift\ +\ 2} cycles through levels. For code: All.

{\def\LTcaptype{none} % do not increment counter
\begin{longtable}[]{@{}ll@{}}
\toprule\noalign{}
Action & Keys \\
\midrule\noalign{}
\endhead
\bottomrule\noalign{}
\endlastfoot
Read current line & \texttt{JAWS\ Key\ +\ Up\ Arrow} \\
Spell current line & \texttt{JAWS\ Key\ +\ Up\ Arrow} (twice quickly) \\
Read from cursor & \texttt{JAWS\ Key\ +\ A} \\
Read current word & \texttt{JAWS\ Key\ +\ Numpad\ 5} \\
Stop reading & \texttt{Ctrl} \\
Increase speech rate & \texttt{Alt\ +\ Ctrl\ +\ Page\ Up} \\
Decrease speech rate & \texttt{Alt\ +\ Ctrl\ +\ Page\ Down} \\
\end{longtable}
}

\textbf{Navigation in VSCode with JAWS:}

{\def\LTcaptype{none} % do not increment counter
\begin{longtable}[]{@{}ll@{}}
\toprule\noalign{}
Action & Keys \\
\midrule\noalign{}
\endhead
\bottomrule\noalign{}
\endlastfoot
Move by character & \texttt{Left\ /\ Right\ Arrow} \\
Move by word & \texttt{Ctrl\ +\ Left\ /\ Right\ Arrow} \\
Move by line & \texttt{Up\ /\ Down\ Arrow} \\
Go to line number & \texttt{Ctrl\ +\ G} \\
Find text & \texttt{Ctrl\ +\ F} \\
Toggle comment & \texttt{Ctrl\ +\ /} \\
\end{longtable}
}

If JAWS stops reading the editor: press \texttt{JAWS\ Key\ +\ Z} to toggle Virtual mode off. You want Virtual mode OFF in VSCode.

\begin{center}\rule{0.5\linewidth}{0.5pt}\end{center}

\subsubsection*{OpenSCAD-Specific Tips}\label{docs__pandoc__latex__src__3dmake_all_lessons.md__openscad-specific-tips}

\textbf{Reading errors.} OpenSCAD errors always include a line number. Use \texttt{Ctrl\ +\ G} in VSCode to jump to it. Common errors:

{\def\LTcaptype{none} % do not increment counter
\begin{longtable}[]{@{}
  >{\raggedright\arraybackslash}p{(\linewidth - 2\tabcolsep) * \real{0.3646}}
  >{\raggedright\arraybackslash}p{(\linewidth - 2\tabcolsep) * \real{0.6354}}@{}}
\toprule\noalign{}
\begin{minipage}[b]{\linewidth}\raggedright
Error message
\end{minipage} & \begin{minipage}[b]{\linewidth}\raggedright
What it usually means
\end{minipage} \\
\midrule\noalign{}
\endhead
\bottomrule\noalign{}
\endlastfoot
\texttt{Expected\ \textquotesingle{};\textquotesingle{}\ ...} &
Missing semicolon at end of a statement \\
\texttt{Expected\ \textquotesingle{},\textquotesingle{}\ or\ \textquotesingle{})\textquotesingle{}\ ...}
& Missing comma between parameters, or unclosed parenthesis \\
\texttt{Identifier\ ...\ is\ undefined} &
Variable name typed wrong, or used before declaration \\
\texttt{WARNING:\ Normalized\ tree\ is\ empty} &
Shape has no geometry --- subtracted more than you started with \\
\end{longtable}
}

\textbf{Bracket matching.} Every \texttt{(} needs a \texttt{)}, every \texttt{{[}} needs a \texttt{{]}}, every \texttt{\{} needs a \texttt{\}}. In VSCode, cursor on a bracket highlights its match. Navigate to the matching bracket with \texttt{Ctrl\ +\ Shift\ +\ \textbackslash{}}.

\textbf{Commenting out code.} To test without deleting: \texttt{//} before a line, or \texttt{Ctrl\ +\ /} in VSCode to toggle comment on selected lines.

\begin{center}\rule{0.5\linewidth}{0.5pt}\end{center}

\subsubsection*{Caliper and OpenSCAD Workflow Tip}\label{docs__pandoc__latex__src__3dmake_all_lessons.md__caliper-and-openscad-workflow-tip}

When measuring and entering into OpenSCAD: say the measurement aloud before typing, then read it back after typing to confirm.

\begin{enumerate}
\tightlist
\item
  Measure → say "seventy point three millimeters"
\item
  Type \texttt{70.3} in OpenSCAD
\item
  Read back: "seven zero point three" --- confirm match
\end{enumerate}

\begin{center}\rule{0.5\linewidth}{0.5pt}\end{center}

\subsection*{Appendix G: VSCode Setup Guide for NVDA and JAWS}\label{docs__pandoc__latex__src__3dmake_all_lessons.md__appendix-g-vscode-setup-guide-for-nvda-and-jaws}

\subsubsection*{Why VSCode Instead of the OpenSCAD Editor}\label{docs__pandoc__latex__src__3dmake_all_lessons.md__why-vscode-instead-of-the-openscad-editor}

The built-in OpenSCAD editor has inconsistent behavior with screen readers --- focus can jump unexpectedly and the editor sometimes stops being read after certain actions. VSCode is a mainstream code editor with strong, well-tested accessibility support for both NVDA and JAWS.

You write code in VSCode. OpenSCAD runs in the background to render the preview. You never need to interact with the OpenSCAD editor itself.

\begin{center}\rule{0.5\linewidth}{0.5pt}\end{center}

\subsubsection*{Install Required Software}\label{docs__pandoc__latex__src__3dmake_all_lessons.md__install-required-software}

\textbf{Install VSCode} --- download from https://code.visualstudio.com/
During installation: check "Add to PATH" and "Register Code as an editor for supported file types."

\textbf{Install OpenSCAD} --- download from https://openscad.org/downloads.html
Use the installer version (not portable). After installing, verify it\textquotesingle s in your PATH:

\begin{lstlisting}[style=Alabaster, language=powershell]
openscad --version


\end{lstlisting}

\textbf{Optional: OpenSCAD VSCode Extension} --- press \texttt{Ctrl\ +\ Shift\ +\ X}, search for "OpenSCAD Language Support," install it. Adds syntax highlighting and keyword completion for \texttt{.scad} files.

\begin{center}\rule{0.5\linewidth}{0.5pt}\end{center}

\subsubsection*{Configure the Task Runner}\label{docs__pandoc__latex__src__3dmake_all_lessons.md__configure-the-task-runner}

The task runner lets you open your \texttt{.scad} file in OpenSCAD or export an STL with a single keypress.

\textbf{Create or open your workspace folder} in VSCode:

\begin{lstlisting}[style=Alabaster, language=powershell]
cd ~/Documents/OpenSCAD_Projects

code .


\end{lstlisting}

\textbf{Create tasks.json:}

\begin{enumerate}
\tightlist
\item
  Press \texttt{Ctrl\ +\ Shift\ +\ P} to open the Command Palette
\item
  Type \texttt{Tasks:\ Configure\ Task} and press Enter
\item
  Select "Create tasks.json file from template" → "Others"
\end{enumerate}

Replace the file\textquotesingle s entire contents with:

\begin{lstlisting}[style=Alabaster, language=json]
{

  "version": "2.0.0",

  "tasks": [

    {

      "label": "Preview in OpenSCAD",

      "type": "shell",

      "command": "openscad",

      "args": ["${file}"],

      "group": { "kind": "build", "isDefault": true },

      "presentation": { "reveal": "silent", "panel": "shared" },

      "problemMatcher": []

    },

    {

      "label": "Export STL",

      "type": "shell",

      "command": "openscad",

      "args": [

        "-o",

        "${fileDirname}/${fileBasenameNoExtension}.stl",

        "${file}"

      ],

      "group": "build",

      "presentation": { "reveal": "always", "panel": "shared" },

      "problemMatcher": []

    }

  ]

}


\end{lstlisting}

Save with \texttt{Ctrl\ +\ S}.

\textbf{Run the preview:} with a \texttt{.scad} file open, press \texttt{Ctrl\ +\ Shift\ +\ B}. OpenSCAD opens and displays your model. \texttt{Alt\ +\ Tab} to return to VSCode.

\textbf{Export STL:} \texttt{Ctrl\ +\ Shift\ +\ P} → "Tasks: Run Task" → "Export STL."

\begin{center}\rule{0.5\linewidth}{0.5pt}\end{center}

\subsubsection*{NVDA Settings for VSCode}\label{docs__pandoc__latex__src__3dmake_all_lessons.md__nvda-settings-for-vscode}

\begin{itemize}
\tightlist
\item
  Symbol level: Most (\texttt{NVDA\ +\ N} → Preferences → Settings → Speech → Symbol level: Most)
\item
  For code reading: \texttt{NVDA\ +\ P} to cycle to Most or All on the fly
\end{itemize}

\subsubsection*{JAWS Settings for VSCode}\label{docs__pandoc__latex__src__3dmake_all_lessons.md__jaws-settings-for-vscode}

\begin{itemize}
\tightlist
\item
  Toggle Virtual mode OFF: \texttt{JAWS\ Key\ +\ Z} (you want Virtual mode off in VSCode)
\item
  Punctuation level: All (\texttt{JAWS\ Key\ +\ Shift\ +\ 2})
\end{itemize}

\begin{center}\rule{0.5\linewidth}{0.5pt}\end{center}

\subsubsection*{Workflow Summary}\label{docs__pandoc__latex__src__3dmake_all_lessons.md__workflow-summary}

\begin{lstlisting}[style=Alabaster]
1. Open VSCode workspace          code ~/Documents/OpenSCAD_Projects

2. Open or create a .scad file    Ctrl + P → type filename

3. Write OpenSCAD code

4. Preview in OpenSCAD            Ctrl + Shift + B

5. Export STL                     Ctrl + Shift + P → Tasks: Run Task → Export STL

6. Open slicer, import STL, slice, export G-code

7. Print


\end{lstlisting}

\begin{center}\rule{0.5\linewidth}{0.5pt}\end{center}

\subsubsection*{Troubleshooting}\label{docs__pandoc__latex__src__3dmake_all_lessons.md__troubleshooting-1}

{\def\LTcaptype{none} % do not increment counter
\begin{longtable}[]{@{}
  >{\raggedright\arraybackslash}p{(\linewidth - 2\tabcolsep) * \real{0.2424}}
  >{\raggedright\arraybackslash}p{(\linewidth - 2\tabcolsep) * \real{0.7576}}@{}}
\toprule\noalign{}
\begin{minipage}[b]{\linewidth}\raggedright
Problem
\end{minipage} & \begin{minipage}[b]{\linewidth}\raggedright
Solution
\end{minipage} \\
\midrule\noalign{}
\endhead
\bottomrule\noalign{}
\endlastfoot
OpenSCAD doesn\textquotesingle t open &
Confirm \texttt{openscad\ -\/-version} works in PowerShell \\
VSCode not read by screen reader &
JAWS: press \texttt{JAWS\ Key\ +\ Z} to exit Virtual mode; NVDA: click inside the editor once to confirm focus \\
Can\textquotesingle t find error in code &
Error includes a line number --- use \texttt{Ctrl\ +\ G} to jump to it \\
\texttt{.scad} has an error &
\texttt{Alt\ +\ Tab} to OpenSCAD to read the console; arrow through the error message \\
-e & \\
\end{longtable}
}

\begin{center}\rule{0.5\linewidth}{0.5pt}\end{center}

\section*{Lesson 2 --- Geometric Primitives and Constructive Solid Geometry}\label{docs__pandoc__latex__src__3dmake_all_lessons.md__lesson-2--geometric-primitives-and-constructive-solid-geometry}

\textbf{Estimated Time:} 90--120 minutes\\
\textbf{Course:} 3dMake Certification --- High School Distance Learning Track

\begin{center}\rule{0.5\linewidth}{0.5pt}\end{center}

\subsection*{Before You Start}\label{docs__pandoc__latex__src__3dmake_all_lessons.md__before-you-start}

In Lesson 1, you made a box by calling \texttt{cube()}. You understood the coordinate system, installed the tools, and ran the complete build-verify workflow. That\textquotesingle s a real foundation --- but a single rectangular box isn\textquotesingle t going to take you very far as a designer.

This lesson introduces the complete vocabulary of 3D shapes OpenSCAD knows how to build, and --- more importantly --- the three operations that let you combine those shapes into virtually anything you can imagine. Those three operations are \textbf{union} (merge), \textbf{difference} (subtract), and \textbf{intersection} (keep only what overlaps). Together with the primitive shapes, they form a system called \textbf{Constructive Solid Geometry}, and nearly every 3D model ever built in OpenSCAD is made from these tools.

\begin{center}\rule{0.5\linewidth}{0.5pt}\end{center}

\subsection*{Learning Objectives}\label{docs__pandoc__latex__src__3dmake_all_lessons.md__learning-objectives-2}

By the end of this lesson you will be able to:

\begin{itemize}
\tightlist
\item
  Use all core OpenSCAD 3D primitives: \texttt{cube}, \texttt{sphere}, \texttt{cylinder}, \texttt{polygon}
\item
  Combine shapes with \texttt{union()}, \texttt{difference()}, and \texttt{intersection()}
\item
  Apply the 0.001 mm offset rule and understand exactly why it exists
\item
  Create 2D profiles with \texttt{polygon()} and extrude them into 3D geometry
\item
  Use \texttt{rotate\_extrude()} and \texttt{linear\_extrude()} to create complex rotational shapes
\item
  Use the four debugging modifier characters (\texttt{\#}, \texttt{\%}, \texttt{!}, \texttt{*})
\item
  Use \texttt{hull()} to create convex envelope shapes
\end{itemize}

\begin{center}\rule{0.5\linewidth}{0.5pt}\end{center}

\subsection*{Concept 1: What Are Primitives?}\label{docs__pandoc__latex__src__3dmake_all_lessons.md__concept-1-what-are-primitives}

\subsubsection*{The Raw Materials of 3D Design}\label{docs__pandoc__latex__src__3dmake_all_lessons.md__the-raw-materials-of-3d-design}

A \textbf{primitive} is the simplest, most basic shape a system can produce from scratch, with no prior shape as input. Everything more complex gets built by combining and modifying these primitives. Think of them the way a sculptor thinks about blocks of material: the primitive is the raw material, and the operations you\textquotesingle ll learn in this lesson are the tools you use to shape it.

OpenSCAD has four main 3D primitives. Every complex shape you\textquotesingle ll ever design in this course is built from some combination of these four, using the CSG operations to combine them.

\begin{center}\rule{0.5\linewidth}{0.5pt}\end{center}

\subsection*{Step 1 --- The Four Core 3D Primitives}\label{docs__pandoc__latex__src__3dmake_all_lessons.md__step-1--the-four-core-3d-primitives}

\subsubsection*{\texorpdfstring{\texttt{cube({[}width,\ depth,\ height{]})} --- The Rectangular Box}{cube({[}width, depth, height{]}) --- The Rectangular Box}}\label{docs__pandoc__latex__src__3dmake_all_lessons.md__cubewidth-depth-height--the-rectangular-box}

Despite the name \texttt{cube}, this primitive creates any rectangular box --- it doesn\textquotesingle t have to be a perfect cube. The three numbers you provide are the dimensions along the X, Y, and Z axes respectively. \texttt{cube({[}30,\ 20,\ 10{]})} creates a box that is 30 mm wide (X), 20 mm deep (Y), and 10 mm tall (Z).

The \texttt{center=true} option (which you learned about in Lesson 1) centers the box on the origin rather than placing one corner there. Remember: \texttt{center=false} is the default, and it\textquotesingle s usually what you want for parts that should sit on the build plate.

\subsubsection*{\texorpdfstring{\texttt{sphere(r=radius)} --- The Round Ball}{sphere(r=radius) --- The Round Ball}}\label{docs__pandoc__latex__src__3dmake_all_lessons.md__sphererradius--the-round-ball}

\texttt{sphere(r=10)} creates a sphere with a 10 mm radius (so 20 mm diameter). The sphere is centered on the origin by default --- it can\textquotesingle t be un-centered the way a cube can.

The key parameter to know here is \textbf{\texttt{\$fn}} --- short for "facet number." OpenSCAD doesn\textquotesingle t have a mathematically perfect curved surface; it approximates curves by using lots of flat faces. A sphere with \texttt{\$fn=8} looks like a faceted gem. A sphere with \texttt{\$fn=64} looks smooth and round. The catch is that higher \texttt{\$fn} values take longer to compute. During development, use \texttt{\$fn=16} for speed. Use \texttt{\$fn=64} or higher only when you\textquotesingle re generating your final print file. Many designers set \texttt{\$fn} globally at the top of their file.

\subsubsection*{\texorpdfstring{\texttt{cylinder(h=height,\ r=radius)} --- The Tube, Cone, or Pipe}{cylinder(h=height, r=radius) --- The Tube, Cone, or Pipe}}\label{docs__pandoc__latex__src__3dmake_all_lessons.md__cylinderhheight-rradius--the-tube-cone-or-pipe}

\texttt{cylinder(h=20,\ r=8)} creates a 20 mm tall cylinder with an 8 mm radius. Like spheres, \texttt{\$fn} controls how round it looks.

There\textquotesingle s an especially powerful variant: instead of a single \texttt{r} (which makes both the top and bottom the same radius), you can specify \textbf{\texttt{r1}} for the bottom radius and \textbf{\texttt{r2}} for the top radius. This creates tapered shapes:

\begin{itemize}
\tightlist
\item
  \texttt{r1=10,\ r2=5} creates a cone that\textquotesingle s wider at the bottom and narrower at the top
\item
  \texttt{r1=8,\ r2=0} creates a spike --- the top radius is zero, so the top comes to a point
\item
  \texttt{r1=0,\ r2=8} creates an upside-down spike --- a shape that tapers to a point at the bottom
\end{itemize}

This flexibility makes \texttt{cylinder} far more versatile than its name suggests.

\subsubsection*{\texorpdfstring{\texttt{polygon(points)} --- The Custom 2D Shape}{polygon(points) --- The Custom 2D Shape}}\label{docs__pandoc__latex__src__3dmake_all_lessons.md__polygonpoints--the-custom-2d-shape}

None of the above shapes can make an L-shape, a T-shape, a hexagon, or a star. That\textquotesingle s what \texttt{polygon()} is for.

\texttt{polygon()} accepts a list of (x, y) coordinate pairs and creates a flat 2D shape in the XY plane by connecting them in order. The shape automatically closes --- the last point connects back to the first. Points should be listed in \textbf{counterclockwise order} as seen from above --- this convention tells OpenSCAD which side is the "inside" of the polygon.

By itself, a \texttt{polygon()} is a flat 2D cutout. Combined with \texttt{linear\_extrude(height=N)}, it gets pushed straight up along the Z axis to become a 3D prism with that polygon as its cross-section. This is your escape hatch for any shape that can\textquotesingle t be described with boxes, spheres, and cylinders.

\begin{lstlisting}[style=Alabaster, language=openscad]
// All four primitives demonstrated:


// 1. A simple box (already familiar from Lesson 1)

cube([30, 20, 10]);


// 2. A sphere — note it's centered on origin by default

// Translate it sideways so it doesn't overlap the cube

translate([40, 0, 10])

  sphere(r=10, $fn=32);    // $fn=32 gives a smooth appearance


// 3. A standard cylinder

translate([70, 0, 0])

  cylinder(h=20, r=8, $fn=32);


// 4. A tapered cylinder (cone) — wider at the bottom

translate([100, 0, 0])

  cylinder(h=20, r1=10, r2=3, $fn=32);


// 5. A spike — r2=0 makes a pointed top

translate([130, 0, 0])

  cylinder(h=20, r1=8, r2=0, $fn=32);


// 6. A polygon extruded into a prism

// This L-shape profile gets pushed up 5 mm to become a 3D L-bracket

translate([0, 40, 0])

  linear_extrude(height=5)

    polygon([[0,0],[40,0],[40,8],[8,8],[8,30],[0,30]]);


\end{lstlisting}

After building this code, run \texttt{3dm\ info} to check the bounding box of each shape and make sure you understand the relationship between the numbers you typed and the dimensions reported.

\begin{center}\rule{0.5\linewidth}{0.5pt}\end{center}

\subsection*{Concept 2: Constructive Solid Geometry}\label{docs__pandoc__latex__src__3dmake_all_lessons.md__concept-2-constructive-solid-geometry}

\subsubsection*{The Logic Behind Combining Shapes}\label{docs__pandoc__latex__src__3dmake_all_lessons.md__the-logic-behind-combining-shapes}

\textbf{Constructive Solid Geometry (CSG)} is a technique for building complex 3D shapes by performing logical operations on simpler ones. The word "Boolean" comes from Boolean algebra --- the mathematics of true/false logic. In 3D space, a Boolean operation asks a question about every point in space: is this point inside shape A? Inside shape B? Inside both? Inside neither?

The three answers to those questions correspond to the three CSG operations.

\subsubsection*{Why This Approach Is So Powerful}\label{docs__pandoc__latex__src__3dmake_all_lessons.md__why-this-approach-is-so-powerful}

Think about how you\textquotesingle d make a box with a cylindrical hole drilled through it. You can\textquotesingle t do that with a single primitive. But you can:

\begin{enumerate}
\tightlist
\item
  Create a box
\item
  Create a cylinder the size of the hole
\item
  Subtract the cylinder from the box
\end{enumerate}

That\textquotesingle s \texttt{difference()}. Most real-world parts are designed exactly this way: start with a basic solid shape, then cut, drill, and hollow it using additional shapes as "tools."

Now think about a phone stand that has a rounded edge. You can\textquotesingle t describe that with a box. But you can:

\begin{enumerate}
\tightlist
\item
  Create the main body as a box
\item
  Create a sphere slightly larger than the corner radius
\item
  Use \texttt{hull()} to create a smooth, rounded version
\end{enumerate}

That\textquotesingle s the power of CSG: you\textquotesingle re building a vocabulary of operations that let you describe arbitrarily complex geometry by combining simple pieces.

\begin{center}\rule{0.5\linewidth}{0.5pt}\end{center}

\subsection*{Step 2 --- The Three CSG Operations}\label{docs__pandoc__latex__src__3dmake_all_lessons.md__step-2--the-three-csg-operations}

\subsubsection*{\texorpdfstring{\texttt{union()} --- Merge Shapes Together}{union() --- Merge Shapes Together}}\label{docs__pandoc__latex__src__3dmake_all_lessons.md__union--merge-shapes-together}

\texttt{union()} combines two or more shapes into a single solid. Any point that is inside \emph{any} of the shapes is inside the result.

An important subtlety: \texttt{union()} isn\textquotesingle t the same as just rendering two objects next to each other. If shapes overlap, \texttt{union()} properly merges them into a single watertight solid. Two overlapping objects without a CSG wrapper can cause problems for the slicer --- it may not know whether the overlapping region should be treated as solid once or twice. Always use \texttt{union()} when shapes share space.

\begin{lstlisting}[style=Alabaster, language=openscad]
// union() — combines a cube and a sphere into one solid piece

// The sphere sits on top of the cube, and the two are fused where they meet.

union() {

  cube([20, 20, 5]);

  translate([10, 10, 5])

    sphere(r=10, $fn=32);

}


\end{lstlisting}

Build this and notice: the cube and sphere are one object. You can\textquotesingle t select them separately.

\subsubsection*{\texorpdfstring{\texttt{difference()} --- Subtract Shapes}{difference() --- Subtract Shapes}}\label{docs__pandoc__latex__src__3dmake_all_lessons.md__difference--subtract-shapes}

\texttt{difference()} subtracts child shapes from the first child. The \textbf{first} shape listed inside \texttt{difference()} is the base --- the "workpiece." Every shape listed after it acts as a cutting tool, removing its volume from the base.

Think of a drill press: the block of material is the first child, and the drill bits are the shapes that follow. You position the drills where you want the holes, and \texttt{difference()} removes exactly the volume each drill bit occupies.

\begin{lstlisting}[style=Alabaster, language=openscad]
// difference() — a cube with a cylindrical hole drilled through it

// The cube is the base. The cylinder is the cutting tool.

difference() {

  cube([20, 20, 20]);

  translate([10, 10, -0.001])       // see Step 4 for why this -0.001 matters

    cylinder(h=20.002, r=5, $fn=32);

}


\end{lstlisting}

Caution

The most common mistake with \texttt{difference()} is accidentally listing the cutting tool first and the base second --- the result will be empty or backwards. Always check: is the first child in your \texttt{difference()} the piece you want to \textbf{keep}?

\subsubsection*{\texorpdfstring{\texttt{intersection()} --- Keep Only What Overlaps}{intersection() --- Keep Only What Overlaps}}\label{docs__pandoc__latex__src__3dmake_all_lessons.md__intersection--keep-only-what-overlaps}

\texttt{intersection()} keeps only the volume that exists inside \textbf{all} shapes simultaneously. If shape A and shape B overlap, \texttt{intersection()} gives you just the overlapping region --- everything that is inside both at the same time.

A practical mental model: \texttt{intersection()} is like a cookie cutter. The cookie cutter shape is one child, the dough is another. The result is the piece of dough within the cutter\textquotesingle s outline.

\begin{lstlisting}[style=Alabaster, language=openscad]
// intersection() — keeps only the region inside both shapes

// A cube intersected with a sphere = a cube with rounded corners

intersection() {

  cube([20, 20, 20], center=true);

  sphere(r=12, $fn=64);

}


\end{lstlisting}

Build this and observe: the straight edges of the cube are rounded off wherever the sphere\textquotesingle s surface cuts through them. The result is a shape that\textquotesingle s simultaneously box-like and rounded.

\subsubsection*{Visualizing Each Operation}\label{docs__pandoc__latex__src__3dmake_all_lessons.md__visualizing-each-operation}

The best way to internalize these operations is to build each example above separately and study the result before moving on. When you can predict what each operation will produce before you build it, you\textquotesingle re ready to combine them in more complex designs.

Here\textquotesingle s a quick mental checklist:

\begin{itemize}
\tightlist
\item
  \textbf{union()}: Think "merge everything together, fill in any gaps."
\item
  \textbf{difference()}: Think "start with the first thing, carve away everything else."
\item
  \textbf{intersection()}: Think "keep only the parts that are inside every single child."
\end{itemize}

\begin{center}\rule{0.5\linewidth}{0.5pt}\end{center}

\subsection*{\texorpdfstring{Step 3 --- \texttt{polygon()} for Custom 2D Profiles}{Step 3 --- polygon() for Custom 2D Profiles}}\label{docs__pandoc__latex__src__3dmake_all_lessons.md__step-3--polygon-for-custom-2d-profiles}

\subsubsection*{When Primitives Aren\textquotesingle t the Right Shape}\label{docs__pandoc__latex__src__3dmake_all_lessons.md__when-primitives-arent-the-right-shape}

Boxes, spheres, and cylinders get you surprisingly far. But eventually you\textquotesingle ll need a shape that doesn\textquotesingle t fit any of them --- an L-bracket, a triangular wedge, a hexagonal base, an irregular outline traced from a measurement. That\textquotesingle s where \texttt{polygon()} earns its place.

\texttt{polygon()} takes a list of (x, y) coordinate pairs and creates the 2D shape enclosed by connecting them in order. The shape is drawn in the XY plane. Because it\textquotesingle s 2D, you can\textquotesingle t print it directly --- but combined with \texttt{linear\_extrude(height=N)}, it becomes a 3D solid with that exact cross-section.

Here\textquotesingle s how to use it: sketch your cross-section shape on graph paper. Mark the coordinates of every corner. Transfer those coordinates into the \texttt{polygon()} list. Connect the last point back to the first (OpenSCAD does this automatically). List them \textbf{counterclockwise} when viewed from above.

\begin{lstlisting}[style=Alabaster, language=openscad]
// An L-shaped bracket: cross-section traced from corner coordinates

bracket_profile = [

  [0,  0],    // bottom-left

  [40, 0],    // bottom-right

  [40, 8],    // step up — beginning of horizontal arm's top edge

  [8,  8],    // inside corner of the L

  [8,  30],   // top of the vertical arm

  [0,  30]    // top-left

];


// Extrude 5 mm along Z — this produces a 3D L-bracket

linear_extrude(height=5)

  polygon(bracket_profile);


\end{lstlisting}

Walk through those coordinates: starting at the bottom-left, moving right along the bottom, up and left to form the horizontal arm, then up the vertical arm, and back to the start. The result is a shape you genuinely cannot make with \texttt{cube()} alone.

\textbf{The graph-paper method works for almost any traced shape.} Measure your object, draw it on grid paper, read off the corner coordinates, and type them in. This is a real engineering workflow.

\begin{center}\rule{0.5\linewidth}{0.5pt}\end{center}

\subsection*{Step 4 --- The 0.001 mm Offset Rule}\label{docs__pandoc__latex__src__3dmake_all_lessons.md__step-4--the-0001-mm-offset-rule}

\subsubsection*{The Co-Planar Face Problem}\label{docs__pandoc__latex__src__3dmake_all_lessons.md__the-co-planar-face-problem}

When you use \texttt{difference()} to cut one shape with another, there\textquotesingle s a subtle problem that occurs when the cutting shape shares an exact face with the base shape --- for example, when the bottom face of a cutting cylinder sits at exactly Z=0, the same as the bottom face of the box it\textquotesingle s cutting.

At that shared boundary, OpenSCAD faces an ambiguity: does this face belong to the outer solid or the inner (cutting) solid? The geometry engine can\textquotesingle t make that decision cleanly, and the result is \textbf{undefined behavior} --- sometimes a rendering artifact, sometimes a non-manifold edge that crashes the slicer, sometimes output that looks fine in preview but is broken in the exported STL.

This problem is called a \textbf{co-planar face problem}, and it\textquotesingle s one of the most common sources of mysterious errors in OpenSCAD files.

\subsubsection*{The Fix}\label{docs__pandoc__latex__src__3dmake_all_lessons.md__the-fix}

Important

\textbf{The 0.001 mm offset rule:} whenever subtracting a shape, extend the cutting shape \textbf{0.001 mm past each face} it needs to pass through (0.002 mm total extra length). This eliminates co-planar face ambiguity that causes rendering artifacts and slicer errors.

Specifically, move it \textbf{0.001 mm past the face on the entry side} and make it \textbf{0.002 mm longer total}.

Why 0.001 mm? It\textquotesingle s large enough to eliminate the ambiguity entirely --- the cutting shape now unambiguously passes through the surface. It\textquotesingle s also small enough that no printer in the world can detect it. The 0.001 mm offset is invisible in the physical print.

\begin{lstlisting}[style=Alabaster, language=openscad]
// WRONG — co-planar faces: the cutting shape is exactly as tall as the base

// This may produce rendering artifacts or slicer errors.

difference() {

  cube([20, 20, 10]);

  translate([1, 1, 0])      // cutting shape's bottom at Z=0, same as base

    cube([18, 18, 10]);     // cutting shape's top at Z=10, same as base's top

}


// CORRECT — the 0.001 offset rule applied

// The cutting shape starts 0.001 mm BELOW the base bottom

// and extends 0.001 mm ABOVE the base top.

difference() {

  cube([20, 20, 10]);

  translate([1, 1, -0.001])         // start 0.001 mm below the floor

    cube([18, 18, 10 + 0.002]);     // 0.001 below + 0.001 above = 0.002 extra height

}


\end{lstlisting}

\textbf{The pattern to memorize:} when cutting along the Z axis, use \texttt{translate({[}x,\ y,\ -0.001{]})} to start the cut 0.001 mm below the floor, and add 0.002 to the cutting shape\textquotesingle s height (0.001 below + 0.001 above = 0.002 total extra). For cuts along the X or Y axis, apply the same logic to those dimensions instead.

\begin{center}\rule{0.5\linewidth}{0.5pt}\end{center}

\subsection*{Step 5 --- Debugging Modifier Characters}\label{docs__pandoc__latex__src__3dmake_all_lessons.md__step-5--debugging-modifier-characters}

\subsubsection*{The Four Special Modifiers}\label{docs__pandoc__latex__src__3dmake_all_lessons.md__the-four-special-modifiers}

OpenSCAD provides four \textbf{modifier characters} that you can prefix onto any shape to temporarily change how it appears in the preview. These are debugging tools --- they don\textquotesingle t change the final exported geometry. Always remove them before exporting a file for printing.

\textbf{\texttt{\#} (hash / highlight)} --- renders the shape in translucent pink/red, \emph{and} keeps it in the model. Use this when you want to see exactly where a cutting shape is sitting relative to the base. Add \texttt{\#} to a cylinder inside a \texttt{difference()} and you\textquotesingle ll see the cutting tool highlighted.

\textbf{\texttt{\%} (percent / ghost)} --- renders the shape as a transparent ghost, and \textbf{excludes it from the model output}. Use this for reference geometry --- a ghost outline of a component you\textquotesingle re designing something to attach to --- that helps you position things but shouldn\textquotesingle t be in the final print.

\textbf{\texttt{!} (bang / isolate)} --- renders \textbf{only this shape}, ignoring everything else in the entire file. Use this when you want to check a single module or shape by itself, without every other object in the scene blocking the view.

\textbf{\texttt{*} (star / disable)} --- completely disables this shape; it\textquotesingle s not rendered at all and produces no output. Use this to temporarily turn off part of your design during debugging, without deleting the code.

\begin{lstlisting}[style=Alabaster, language=openscad]
// USING MODIFIERS FOR DEBUGGING:


difference() {

  cube([30, 30, 20]);


  // # highlights the cutting shape so you can see where it's cutting

  # translate([15, 15, -0.001])

    cylinder(h=20.002, r=8, $fn=32);

}


// % ghost reference — helps position but doesn't print

% translate([0, 0, 20])

  cube([30, 30, 5]);   // ghost of a lid that will sit on top


\end{lstlisting}

Important

\textbf{F5 vs F6:} F5 runs a fast approximate preview that is NOT manifold-safe --- good for quick visual checks only. F6 runs the full CGAL geometric kernel: slower, but accurate and manifold-safe. Always use \textbf{F6} (or \texttt{3dm\ build}) before sending anything to print.

\begin{center}\rule{0.5\linewidth}{0.5pt}\end{center}

\subsection*{\texorpdfstring{Step 6 --- \texttt{linear\_extrude()} and \texttt{rotate\_extrude()}}{Step 6 --- linear\_extrude() and rotate\_extrude()}}\label{docs__pandoc__latex__src__3dmake_all_lessons.md__step-6--linear_extrude-and-rotate_extrude}

\subsubsection*{Two Ways to Create 3D Shapes from 2D Profiles}\label{docs__pandoc__latex__src__3dmake_all_lessons.md__two-ways-to-create-3d-shapes-from-2d-profiles}

You\textquotesingle ve already used \texttt{linear\_extrude()} with \texttt{polygon()}. Now let\textquotesingle s understand it more fully, and introduce its sibling \texttt{rotate\_extrude()}.

\textbf{\texttt{linear\_extrude(height=N)}} takes a 2D profile (any 2D shape: \texttt{circle}, \texttt{square}, \texttt{polygon}, or a combination) and pushes it straight up along the Z axis, creating a prism. The cross-section of the prism matches the 2D profile exactly.

Two important optional parameters:

\begin{itemize}
\tightlist
\item
  \texttt{twist=angle} --- rotates the profile as it extrudes, creating a spiral or helix. A \texttt{twist=90} on a square profile creates a twisted square column.
\item
  \texttt{scale=factor} --- changes the size of the profile from bottom to top. \texttt{scale=0.5} means the top is half the size of the bottom, creating a tapered shape.
\end{itemize}

\textbf{\texttt{rotate\_extrude(angle=360)}} takes a 2D profile and rotates it around the \textbf{Z axis}, sweeping it in a circle to create a rotationally symmetric solid. Imagine a lathe: the profile is the outline of a piece, and the lathe spins it to create a solid of revolution.

The profile for \texttt{rotate\_extrude()} must be positioned in \textbf{positive X space} --- to the right of the Z axis. The X distance from the Z axis becomes the radius of the resulting shape. If you place a small circle at X=15 and rotate-extrude it, you get a torus (donut) with a 15 mm ring radius.

\begin{lstlisting}[style=Alabaster, language=openscad]
// linear_extrude: circle pushed straight up = a cylinder

linear_extrude(height=20) {

  circle(r=10, $fn=32);

}


// linear_extrude with twist: a twisted square column

translate([30, 0, 0])

linear_extrude(height=30, twist=90) {

  square([10, 10], center=true);

}


// linear_extrude with scale: a tapered shape (like a pyramid with a square base)

translate([60, 0, 0])

linear_extrude(height=30, scale=0.1) {

  square([15, 15], center=true);

}


// rotate_extrude: a circle swept around the Z axis = a torus (donut)

translate([0, 50, 0])

rotate_extrude(angle=360, $fn=64) {

  translate([15, 0, 0]) circle(r=5, $fn=32);

}


\end{lstlisting}

Any shape that is rotationally symmetric --- a bowl, a cup, a ring, a knob, a vase --- can and should be made with \texttt{rotate\_extrude()}. It\textquotesingle s dramatically more elegant than building the same shape from many overlapping cylinders.

\begin{center}\rule{0.5\linewidth}{0.5pt}\end{center}

\subsection*{\texorpdfstring{Step 7 --- \texttt{hull()} --- The Convex Envelope}{Step 7 --- hull() --- The Convex Envelope}}\label{docs__pandoc__latex__src__3dmake_all_lessons.md__step-7--hull--the-convex-envelope}

\subsubsection*{\texorpdfstring{What \texttt{hull()} Does}{What hull() Does}}\label{docs__pandoc__latex__src__3dmake_all_lessons.md__what-hull-does}

\texttt{hull()} computes the \textbf{convex hull} of all its children --- the smallest convex shape that completely contains all of them. The easiest mental model: imagine stretching a rubber band around a collection of objects laid on a table. The shape the rubber band makes when it snaps tight around all of them is the convex hull.

Notice the key word: \textbf{convex}. A convex shape has no indentations or concave regions --- every surface curves outward. \texttt{hull()} fills in any gaps or concavities between the shapes you give it.

\subsubsection*{\texorpdfstring{How \texttt{hull()} Differs from \texttt{union()}}{How hull() Differs from union()}}\label{docs__pandoc__latex__src__3dmake_all_lessons.md__how-hull-differs-from-union}

\texttt{union()} takes all the shapes and merges them, preserving their individual forms including any gaps or indentations between them.

\texttt{hull()} takes all the shapes and creates the smoothest possible outer envelope that touches or contains all of them, filling in any gaps.

This makes \texttt{hull()} excellent for:

\begin{itemize}
\tightlist
\item
  Creating smooth transitions between two shapes of different sizes
\item
  Making rounded boxes (place a small sphere at each corner, take the hull)
\item
  Creating organic, blob-like shapes
\item
  Generating ramps and smooth blends
\end{itemize}

\begin{lstlisting}[style=Alabaster, language=openscad]
// hull() of three spheres — creates a smooth rounded triangle

hull() {

  translate([0,  0,  0]) sphere(r=5, $fn=32);

  translate([30, 0,  0]) sphere(r=5, $fn=32);

  translate([15, 20, 0]) sphere(r=5, $fn=32);

}


// Practical use: a rounded box

// Place small spheres at the eight corners of a box

// The hull creates a box with perfectly rounded edges

translate([0, 40, 0])

hull() {

  for (x = [3, 27], y = [3, 17], z = [3, 7]) {

    translate([x, y, z]) sphere(r=3, $fn=16);

  }

}


\end{lstlisting}

The rounded box pattern is one of the most commonly used techniques in OpenSCAD. The sphere radius controls how round the edges are. The positions of the spheres control the overall box dimensions. Change the sphere radius to change how much rounding you get.

\begin{center}\rule{0.5\linewidth}{0.5pt}\end{center}

\subsection*{Advanced CSG: Putting It All Together}\label{docs__pandoc__latex__src__3dmake_all_lessons.md__advanced-csg-putting-it-all-together}

\subsubsection*{A Parametric Mounting Bracket}\label{docs__pandoc__latex__src__3dmake_all_lessons.md__a-parametric-mounting-bracket}

Here is a complete practical example that uses nested CSG operations to create a real functional part. Read through it carefully --- notice how the design intent is encoded as named variables at the top, and the geometry simply uses those variables.

\begin{lstlisting}[style=Alabaster, language=openscad]
// Parametric Mounting Bracket

// Change these values to resize the bracket

width  = 40;    // mm — total bracket width

height = 30;    // mm — total bracket height

depth  = 8;     // mm — bracket thickness (front to back)

hole_r = 4;     // mm — radius of mounting holes

slot_w = 6;     // mm — width of top slot

slot_h = 15;    // mm — height of top slot

wall   = 3;     // mm — material remaining above slot


module bracket() {

  difference() {

    // Base: a solid rectangular block — this is what everything gets cut from

    cube([width, depth, height]);


    // Two circular mounting holes drilled through the Y axis

    // The rotate([-90,0,0]) turns the cylinder to point in the Y direction

    translate([10, -0.001, 10])

      rotate([-90, 0, 0]) cylinder(r=hole_r, h=depth + 0.002, $fn=32);

    translate([30, -0.001, 10])

      rotate([-90, 0, 0]) cylinder(r=hole_r, h=depth + 0.002, $fn=32);


    // A slot cut through the top portion of the bracket

    translate([width/2 - slot_w/2, -0.001, height - slot_h - wall])

      cube([slot_w, depth + 0.002, slot_h]);

  }

}


bracket();


\end{lstlisting}

Walk through this mentally: start with a solid block. Drill two circular holes through it from front to back. Cut a slot down from the top. The result is a bracket with two mounting holes and an adjustment slot --- all from one parametric design.

Try changing \texttt{hole\_r}, \texttt{slot\_w}, or \texttt{width} and rebuilding. Watch how every feature adapts automatically.

\begin{center}\rule{0.5\linewidth}{0.5pt}\end{center}

\subsection*{Quiz --- Lesson 2 (15 questions)}\label{docs__pandoc__latex__src__3dmake_all_lessons.md__quiz--lesson-2-15-questions}

\begin{enumerate}
\tightlist
\item
  What are the four main 3D primitives in OpenSCAD?
\item
  What does \texttt{difference()} do? Which child is the base, and which children are the cutting tools?
\item
  Why do you add 0.001 mm to the cutting geometry in a \texttt{difference()} operation? What problem does it prevent?
\item
  What does the \texttt{\#} modifier character do, and when would you use it during debugging?
\item
  What is the difference between F5 and F6 in the OpenSCAD GUI? Which should you use before printing?
\item
  What does \texttt{hull()} produce? How is it different from \texttt{union()}?
\item
  What does the \texttt{\%} modifier do to a shape in OpenSCAD?
\item
  Describe what \texttt{rotate\_extrude()} does and give one example of a shape it could produce.
\item
  What does \texttt{intersection()} return when applied to a cube and a sphere that partially overlap?
\item
  True or False: the \texttt{*} modifier renders a shape as a transparent ghost for debugging purposes.
\item
  Describe what \texttt{linear\_extrude()} does and explain what the \texttt{twist} parameter changes.
\item
  What is a co-planar face, and what problem does it cause in a \texttt{difference()} operation?
\item
  If you want to subtract a cylinder from a cube and the cylinder is exactly as tall as the cube, what two adjustments do you need to make to ensure a clean cut?
\item
  When using \texttt{polygon()}, what order should points be listed in, and why?
\item
  What is the difference between combining two overlapping shapes with \texttt{union()} versus rendering two separate overlapping objects without any CSG operation?
\end{enumerate}

\begin{center}\rule{0.5\linewidth}{0.5pt}\end{center}

\subsection*{Extension Problems (15)}\label{docs__pandoc__latex__src__3dmake_all_lessons.md__extension-problems-15-1}

\begin{enumerate}
\tightlist
\item
  Build a hollow sphere with 1.5 mm walls using \texttt{difference()} and the 0.001 mm rule. Document the outer and inner radii you chose and verify the wall thickness using \texttt{3dm\ info}.
\item
  Create a vase shape using \texttt{rotate\_extrude()} and a custom 2D profile. The vase should be wider in the middle than at the top and bottom.
\item
  Design a bracket or clip using nested \texttt{difference()} and \texttt{union()}. Add a comment to every CSG step explaining what it does.
\item
  Use \texttt{hull()} to create a smooth, rounded box with 3 mm radius corners. Document the sphere placement strategy you used.
\item
  Use the \texttt{\%} modifier to create a ghost reference shape that helps you position a snap-fit clip around a cylinder. Screenshot the debugging view and explain it.
\item
  Create a parametric name badge: a flat rectangular base with your name text embossed (raised) on the top face, using \texttt{linear\_extrude()} combined with \texttt{text()}. Make the text height a parameter.
\item
  Build a compound hinge using two cylinders (the barrel), a \texttt{hull()} for the hinge leaf, and alignment holes through the barrel. Document the design intent.
\item
  Design a 10-sided (decagonal) prism using \texttt{cylinder()} with \texttt{\$fn=10} and a \texttt{difference()} to cut a hexagonal through-hole aligned with the prism\textquotesingle s center.
\item
  Create a single test print that exercises all three CSG operations (\texttt{union}, \texttt{difference}, \texttt{intersection}) in one part. Add header comments explaining what each operation does in the design.
\item
  Using only \texttt{cube()}, \texttt{sphere()}, \texttt{difference()}, and \texttt{hull()}, build a simple chess pawn. It should be recognizable as a pawn when printed.
\item
  Create a lattice cube: a solid cube with a grid of cylindrical holes drilled through it in all three axes. Calculate and document the relationship between hole spacing and remaining wall thickness.
\item
  Build a parametric ring using \texttt{rotate\_extrude()} where the cross-section shape (round, square, or custom polygon) is easily changeable by modifying a variable.
\item
  Research and document the \texttt{surface()} primitive: what input does it accept, what shape does it produce, and when would you use it instead of \texttt{polyhedron()}? Include a working example.
\item
  Design a tolerance test set: five pairs of pegs and holes with clearances from 0.0 mm to 0.4 mm in 0.1 mm increments. Print one set and document which clearance allows free movement on your specific printer.
\item
  Write a short guide explaining all four modifier characters (\texttt{\#}, \texttt{!}, \texttt{\%}, \texttt{*}) with one concrete use case for each, and a note on why you must always remove them before exporting the final file.
\end{enumerate}

\begin{center}\rule{0.5\linewidth}{0.5pt}\end{center}

\subsection*{References and Helpful Resources}\label{docs__pandoc__latex__src__3dmake_all_lessons.md__references-and-helpful-resources-1}

\begin{itemize}
\item
  OpenSCAD User Manual --- Primitive Solids and Boolean Operations --- https://en.wikibooks.org/wiki/OpenSCAD\_User\_Manual/Primitive\_Solids
\item
  OpenSCAD User Manual --- CSG Modelling --- https://en.wikibooks.org/wiki/OpenSCAD\_User\_Manual/CSG\_Modelling
\item
  OpenSCAD User Manual --- Transformations and Extrusions --- https://en.wikibooks.org/wiki/OpenSCAD\_User\_Manual/Using\_the\_2D\_Subsystem
\item
  OpenSCAD User Manual --- Modifier Characters --- https://en.wikibooks.org/wiki/OpenSCAD\_User\_Manual/Modifier\_Characters
\end{itemize}

\subsubsection*{Supplemental Resources}\label{docs__pandoc__latex__src__3dmake_all_lessons.md__supplemental-resources-1}

\begin{itemize}
\tightlist
\item
  Programming with OpenSCAD EPUB Textbook --- assets/Programming\_with\_OpenSCAD.epub --- Chapters on CSG and primitives
\item
  CodeSolutions Repository --- https://github.com/ProgrammingWithOpenSCAD/CodeSolutions --- Worked examples for CSG, hull, and extrusions
\item
  OpenSCAD Quick Reference --- https://programmingwithopenscad.github.io/quick-reference.html --- All primitive and CSG syntax at a glance
\item
  3DMake GitHub Repository --- https://github.com/tdeck/3dmake --- Build workflow reference
\item
  All other reference documents are embedded as appendices below --- no separate files needed.
\end{itemize}

\begin{center}\rule{0.5\linewidth}{0.5pt}\end{center}

\section*{Part 2: Extension Projects and Reference Appendices}\label{docs__pandoc__latex__src__3dmake_all_lessons.md__part-2-extension-projects-and-reference-appendices}

\emph{All supporting documents for Lesson 2 are included below. You do not need to open any additional files.}

\begin{center}\rule{0.5\linewidth}{0.5pt}\end{center}

\subsection*{Appendix A: OpenSCAD Quick Reference --- Cheat Sheet}\label{docs__pandoc__latex__src__3dmake_all_lessons.md__appendix-a-openscad-quick-reference--cheat-sheet}

\emph{Keep this handy during all OpenSCAD work. For full documentation: https://openscad.org/documentation.html}

\begin{center}\rule{0.5\linewidth}{0.5pt}\end{center}

\subsubsection*{Basic Shapes (Primitives)}\label{docs__pandoc__latex__src__3dmake_all_lessons.md__basic-shapes-primitives}

\begin{lstlisting}[style=Alabaster, language=openscad]
cube([length, width, height]);           // rectangular box at origin

cube([l, w, h], center=true);           // box centered at origin


sphere(r=radius);                        // sphere by radius

sphere(d=diameter);                      // sphere by diameter


cylinder(h=height, r=radius);            // cylinder

cylinder(h=height, d=diameter);          // by diameter

cylinder(h=height, r1=5, r2=2);         // cone (different top/bottom radii)


$fn = 50;  // curve resolution — higher = smoother, slower to render


\end{lstlisting}

\begin{center}\rule{0.5\linewidth}{0.5pt}\end{center}

\subsubsection*{Transformations}\label{docs__pandoc__latex__src__3dmake_all_lessons.md__transformations}

\begin{lstlisting}[style=Alabaster, language=openscad]
translate([x, y, z]) shape;             // move

rotate([x_deg, y_deg, z_deg]) shape;    // rotate around each axis

scale([x, y, z]) shape;                 // scale (1.5 = 150%)

mirror([1, 0, 0]) shape;               // mirror across YZ plane

                                        // use [0,1,0] for XZ, [0,0,1] for XY


\end{lstlisting}

\begin{center}\rule{0.5\linewidth}{0.5pt}\end{center}

\subsubsection*{Boolean Operations}\label{docs__pandoc__latex__src__3dmake_all_lessons.md__boolean-operations}

\begin{lstlisting}[style=Alabaster, language=openscad]
union() { shape1; shape2; }             // combine shapes

difference() { base; subtract; }        // remove one shape from another

intersection() { shape1; shape2; }      // keep only the overlapping region


\end{lstlisting}

\textbf{Critical rule for \texttt{difference()}:} always make the subtracting shape extend 0.001 mm past each shared face to avoid co-planar artifacts.

\begin{lstlisting}[style=Alabaster, language=openscad]
// Example — box with a through-hole:

difference() {

  cube([30, 30, 10]);

  translate([15, 15, -0.001])           // extends past bottom face

    cylinder(h=10.002, r=5, $fn=32);   // extends past top face

}


\end{lstlisting}

\begin{center}\rule{0.5\linewidth}{0.5pt}\end{center}

\subsubsection*{Variables and Parameters}\label{docs__pandoc__latex__src__3dmake_all_lessons.md__variables-and-parameters}

\begin{lstlisting}[style=Alabaster, language=openscad]
length = 70;      // define a variable

width  = 16;

height = 5;


cube([length, width, height]);          // use variables

cube([length * 2, width, height]);      // math in expressions


\end{lstlisting}

\begin{center}\rule{0.5\linewidth}{0.5pt}\end{center}

\subsubsection*{Modules (Reusable Functions)}\label{docs__pandoc__latex__src__3dmake_all_lessons.md__modules-reusable-functions}

\begin{lstlisting}[style=Alabaster, language=openscad]
// Define:

module my_box(l=20, w=15, h=10) {

  cube([l, w, h]);

}


// Call:

my_box();               // uses all defaults

my_box(30, 20, 8);     // positional arguments

my_box(l=50);           // named argument; others use defaults


\end{lstlisting}

\begin{center}\rule{0.5\linewidth}{0.5pt}\end{center}

\subsubsection*{Loops}\label{docs__pandoc__latex__src__3dmake_all_lessons.md__loops}

\begin{lstlisting}[style=Alabaster, language=openscad]
// Repeat a shape N times:

for (i = [0:4]) {                       // i goes 0, 1, 2, 3, 4

  translate([i * 20, 0, 0]) cube([15, 15, 5]);

}


// With step size:

for (i = [0:5:20]) {                    // i goes 0, 5, 10, 15, 20

  translate([i, 0, 0]) sphere(r=3);

}


\end{lstlisting}

\begin{center}\rule{0.5\linewidth}{0.5pt}\end{center}

\subsubsection*{2D Shapes (for Extrusion)}\label{docs__pandoc__latex__src__3dmake_all_lessons.md__2d-shapes-for-extrusion}

\begin{lstlisting}[style=Alabaster, language=openscad]
circle(r=10);                           // 2D circle

square([w, h]);                         // 2D rectangle

polygon([[0,0],[10,0],[5,10]]);         // arbitrary 2D shape


// Extrude into 3D:

linear_extrude(height=5) circle(r=10); // cylinder from 2D circle

linear_extrude(height=10, twist=90)    // twisted extrusion

  square([10, 5]);

rotate_extrude()                        // solid of revolution

  translate([15, 0]) circle(r=3);      // torus


\end{lstlisting}

\begin{center}\rule{0.5\linewidth}{0.5pt}\end{center}

\subsubsection*{Modifier Characters (Debugging --- Remove Before Final Export)}\label{docs__pandoc__latex__src__3dmake_all_lessons.md__modifier-characters-debugging--remove-before-final-export}

{\def\LTcaptype{none} % do not increment counter
\begin{longtable}[]{@{}
  >{\centering\arraybackslash}p{(\linewidth - 4\tabcolsep) * \real{0.0941}}
  >{\raggedright\arraybackslash}p{(\linewidth - 4\tabcolsep) * \real{0.3882}}
  >{\raggedright\arraybackslash}p{(\linewidth - 4\tabcolsep) * \real{0.5176}}@{}}
\toprule\noalign{}
\begin{minipage}[b]{\linewidth}\centering
Modifier
\end{minipage} & \begin{minipage}[b]{\linewidth}\raggedright
Effect
\end{minipage} & \begin{minipage}[b]{\linewidth}\raggedright
When to use
\end{minipage} \\
\midrule\noalign{}
\endhead
\bottomrule\noalign{}
\endlastfoot
\texttt{\#} & Highlight shape in red/orange &
Finding where a hole or subtractive shape is \\
\texttt{\%} & Render shape as transparent ghost &
Checking alignment of mating parts \\
\texttt{!} & Render ONLY this shape &
Isolating one component for inspection \\
\texttt{*} & Disable (ignore) this shape &
Temporarily removing geometry \\
\end{longtable}
}

Warning

Always remove all modifier characters (\texttt{\#}, \texttt{\%}, \texttt{!}, \texttt{*}) before exporting your final STL. Leaving them in can alter or omit geometry in the exported file.

\begin{center}\rule{0.5\linewidth}{0.5pt}\end{center}

\subsubsection*{Useful Functions}\label{docs__pandoc__latex__src__3dmake_all_lessons.md__useful-functions}

\begin{lstlisting}[style=Alabaster, language=openscad]
len([a, b, c])     // returns 3 — length of a vector

sqrt(25)           // returns 5

pow(2, 8)          // returns 256 (2^8)

abs(-5)            // returns 5

min(3, 5, 1)       // returns 1

max(3, 5, 1)       // returns 5


\end{lstlisting}

\begin{center}\rule{0.5\linewidth}{0.5pt}\end{center}

\subsubsection*{Keyboard Shortcuts}\label{docs__pandoc__latex__src__3dmake_all_lessons.md__keyboard-shortcuts}

{\def\LTcaptype{none} % do not increment counter
\begin{longtable}[]{@{}
  >{\raggedright\arraybackslash}p{(\linewidth - 2\tabcolsep) * \real{0.1389}}
  >{\raggedright\arraybackslash}p{(\linewidth - 2\tabcolsep) * \real{0.8611}}@{}}
\toprule\noalign{}
\begin{minipage}[b]{\linewidth}\raggedright
Key
\end{minipage} & \begin{minipage}[b]{\linewidth}\raggedright
Action
\end{minipage} \\
\midrule\noalign{}
\endhead
\bottomrule\noalign{}
\endlastfoot
\texttt{F5} &
Preview --- fast, approximate (use for checking while designing) \\
\texttt{F6} & Full render --- exact, required before export \\
\texttt{Ctrl\ +\ S} & Save \\
\texttt{Ctrl\ +\ Z} & Undo \\
\texttt{F3} & Reset camera view \\
\end{longtable}
}

\begin{center}\rule{0.5\linewidth}{0.5pt}\end{center}

\subsubsection*{Export Workflow}\label{docs__pandoc__latex__src__3dmake_all_lessons.md__export-workflow}

\begin{enumerate}
\tightlist
\item
  Press \texttt{F6} (full render --- wait for it to complete)
\item
  File → Export → Export as STL
\item
  Save with a descriptive filename: \texttt{projectname\_v2.stl}
\item
  Run \texttt{3dm\ info} on the exported file to verify dimensions before slicing
\end{enumerate}

\begin{center}\rule{0.5\linewidth}{0.5pt}\end{center}

\subsection*{Appendix B: Bonus Print --- Extension Project}\label{docs__pandoc__latex__src__3dmake_all_lessons.md__appendix-b-bonus-print--extension-project}

\textbf{Estimated time:} 3--5 hours

\subsubsection*{Overview}\label{docs__pandoc__latex__src__3dmake_all_lessons.md__overview-2}

This extension project builds on your first successful print by introducing purposeful modification. You\textquotesingle ll take an existing model from an online repository, decide on a specific change that improves or adapts it for a new purpose, apply that change in OpenSCAD or the slicer, and document the entire process. The goal is to experience the design iteration cycle --- make a change, understand why you made it, evaluate the result.

\subsubsection*{Learning Objectives}\label{docs__pandoc__latex__src__3dmake_all_lessons.md__learning-objectives-3}

By completing this project, you will:

\begin{itemize}
\tightlist
\item
  Apply scaling and simple modifications to an existing model
\item
  Understand how scaling affects tolerances and assembly fit
\item
  Document design changes with enough detail that someone else could reproduce your work
\item
  Practice logging print time, filament use, and outcome for reproducibility
\end{itemize}

\subsubsection*{Materials}\label{docs__pandoc__latex__src__3dmake_all_lessons.md__materials-1}

\begin{itemize}
\tightlist
\item
  Online repository access (Thingiverse, Printables), slicer, printer, filament
\end{itemize}

\subsubsection*{Tasks}\label{docs__pandoc__latex__src__3dmake_all_lessons.md__tasks-1}

\begin{enumerate}
\item
  \textbf{Choose a model} from Thingiverse or Printables. Note its original stated dimensions.
\item
  \textbf{Decide on a purposeful modification} --- scale it up or down, add a mounting hole, strengthen a thin feature, or combine it with another model. Write one paragraph explaining what the modification is and why you\textquotesingle re making it. Who benefits from this change?
\item
  \textbf{Apply the change} in OpenSCAD (preferred) or by scaling in the slicer. Record the new dimensions and what changed relative to the original.
\item
  \textbf{Slice and print} all parts in one session where possible. Log: print time, filament used (in grams), layer height, infill \%, and any slicer settings you adjusted.
\item
  \textbf{Create a short construction note}: describe how the modified print is used or assembled. Photograph the result.
\end{enumerate}

\subsubsection*{Probing Questions}\label{docs__pandoc__latex__src__3dmake_all_lessons.md__probing-questions-1}

\begin{enumerate}
\tightlist
\item
  What motivated the modification, and who benefits from it?
\item
  How did scaling affect tolerances or assembly fit? If the model connects to anything else, did the fit change?
\end{enumerate}

\subsubsection*{Starter Code}\label{docs__pandoc__latex__src__3dmake_all_lessons.md__starter-code-1}

If adapting a model using OpenSCAD, use this scaling scaffold as a starting point:

\begin{lstlisting}[style=Alabaster, language=openscad]
// Bonus Print Modification Scaffold

// Import an STL and apply modifications parametrically.


// ---- Parameters ----

scale_factor  = 1.5;    // 1.0 = original, 1.5 = 50% larger, 0.75 = 25% smaller

add_mount     = true;   // set false to skip the mount hole

mount_r       = 2.5;    // mm — mount hole radius (for M5 screw: 2.75mm)

mount_pos     = [25, 0, 0]; // mm — position of mount hole center


// ---- Base Model ----

// Replace "original.stl" with your actual file path

module base_model() {

  scale([scale_factor, scale_factor, scale_factor])

    import("original.stl", convexity=10);

}


// ---- Mount Hole (optional) ----

module mount_hole() {

  if (add_mount) {

    translate(mount_pos)

      cylinder(h=100, r=mount_r, center=true, $fn=32);

  }

}


// ---- Assembly ----

difference() {

  base_model();

  mount_hole();

}


\end{lstlisting}

\begin{center}\rule{0.5\linewidth}{0.5pt}\end{center}

\subsection*{Appendix C: Bonus Print --- Student Documentation Template}\label{docs__pandoc__latex__src__3dmake_all_lessons.md__appendix-c-bonus-print--student-documentation-template}

\textbf{Author:} \_\_\_\_\_\_\_\_\_\_\_\_\_\_\_\_
\textbf{Date:} \_\_\_\_\_\_\_\_\_\_
\textbf{Description:} Modify an existing 3D model and document the design changes and print outcomes.

\begin{center}\rule{0.5\linewidth}{0.5pt}\end{center}

\subsubsection*{Original Model}\label{docs__pandoc__latex__src__3dmake_all_lessons.md__original-model}

\begin{itemize}
\tightlist
\item
  Model name and source:
\item
  Link or file name:
\item
  Original stated dimensions:
\item
  Why did you choose this model for modification?
\end{itemize}

\begin{center}\rule{0.5\linewidth}{0.5pt}\end{center}

\subsubsection*{Modification Plan}\label{docs__pandoc__latex__src__3dmake_all_lessons.md__modification-plan}

\begin{itemize}
\tightlist
\item
  What change will you make, and why?
\item
  Who benefits from this modification?
\item
  Will this be done in OpenSCAD, or by scaling in the slicer? Why?
\end{itemize}

\begin{center}\rule{0.5\linewidth}{0.5pt}\end{center}

\subsubsection*{Variant Specifications}\label{docs__pandoc__latex__src__3dmake_all_lessons.md__variant-specifications}

{\def\LTcaptype{none} % do not increment counter
\begin{longtable}[]{@{}
  >{\centering\arraybackslash}p{(\linewidth - 8\tabcolsep) * \real{0.2651}}
  >{\raggedright\arraybackslash}p{(\linewidth - 8\tabcolsep) * \real{0.1807}}
  >{\centering\arraybackslash}p{(\linewidth - 8\tabcolsep) * \real{0.1807}}
  >{\centering\arraybackslash}p{(\linewidth - 8\tabcolsep) * \real{0.2048}}
  >{\raggedright\arraybackslash}p{(\linewidth - 8\tabcolsep) * \real{0.1687}}@{}}
\toprule\noalign{}
\begin{minipage}[b]{\linewidth}\centering
Variant
\end{minipage} & \begin{minipage}[b]{\linewidth}\raggedright
Scale or Change
\end{minipage} & \begin{minipage}[b]{\linewidth}\centering
Est. print time
\end{minipage} & \begin{minipage}[b]{\linewidth}\centering
Est. filament (g)
\end{minipage} & \begin{minipage}[b]{\linewidth}\raggedright
Key dimensions
\end{minipage} \\
\midrule\noalign{}
\endhead
\bottomrule\noalign{}
\endlastfoot
Original (reference) & --- & & & \\
V1 (your modification) & & & & \\
V2 (optional) & & & & \\
\end{longtable}
}

\begin{center}\rule{0.5\linewidth}{0.5pt}\end{center}

\subsubsection*{Print Results}\label{docs__pandoc__latex__src__3dmake_all_lessons.md__print-results}

{\def\LTcaptype{none} % do not increment counter
\begin{longtable}[]{@{}
  >{\centering\arraybackslash}p{(\linewidth - 8\tabcolsep) * \real{0.2000}}
  >{\centering\arraybackslash}p{(\linewidth - 8\tabcolsep) * \real{0.2833}}
  >{\centering\arraybackslash}p{(\linewidth - 8\tabcolsep) * \real{0.3167}}
  >{\raggedright\arraybackslash}p{(\linewidth - 8\tabcolsep) * \real{0.1167}}
  >{\raggedright\arraybackslash}p{(\linewidth - 8\tabcolsep) * \real{0.0833}}@{}}
\toprule\noalign{}
\begin{minipage}[b]{\linewidth}\centering
Variant
\end{minipage} & \begin{minipage}[b]{\linewidth}\centering
Actual print time
\end{minipage} & \begin{minipage}[b]{\linewidth}\centering
Actual filament (g)
\end{minipage} & \begin{minipage}[b]{\linewidth}\raggedright
Quality
\end{minipage} & \begin{minipage}[b]{\linewidth}\raggedright
Notes
\end{minipage} \\
\midrule\noalign{}
\endhead
\bottomrule\noalign{}
\endlastfoot
V1 & & & & \\
V2 (if made) & & & & \\
\end{longtable}
}

\begin{center}\rule{0.5\linewidth}{0.5pt}\end{center}

\subsubsection*{How Does Scaling Affect Fit?}\label{docs__pandoc__latex__src__3dmake_all_lessons.md__how-does-scaling-affect-fit}

If your model connects to anything (a wall mount, another printed part, a standard fastener), describe how the scaled dimensions affect the fit:

Original mating dimension: \_\_\_\_ mm → Scaled dimension: \_\_\_\_ mm → Fit result:

\begin{center}\rule{0.5\linewidth}{0.5pt}\end{center}

\subsubsection*{Reflections}\label{docs__pandoc__latex__src__3dmake_all_lessons.md__reflections-1}

\begin{itemize}
\tightlist
\item
  Which variant is most effective for its intended purpose, and why?
\item
  How did scaling or modification affect print time and material use?
\item
  What trade-offs did you observe (strength vs. speed, size vs. material cost)?
\item
  If you made a third variant, what would it be?
\end{itemize}

\begin{center}\rule{0.5\linewidth}{0.5pt}\end{center}

\subsubsection*{Attachments Checklist}\label{docs__pandoc__latex__src__3dmake_all_lessons.md__attachments-checklist-1}

\begin{itemize}
\tightlist
\item[$\square$]
  \texttt{.scad} file with modification code (or slicer project file if scaled in slicer)
\item[$\square$]
  Photos of original and modified prints side by side (if both available)
\item[$\square$]
  Print log for each variant
\item[$\square$]
  Completed variant specification table above
\end{itemize}

\begin{center}\rule{0.5\linewidth}{0.5pt}\end{center}

\subsubsection*{Teacher Feedback}\label{docs__pandoc__latex__src__3dmake_all_lessons.md__teacher-feedback-1}

{\def\LTcaptype{none} % do not increment counter
\begin{longtable}[]{@{}
  >{\raggedright\arraybackslash}p{(\linewidth - 4\tabcolsep) * \real{0.8367}}
  >{\centering\arraybackslash}p{(\linewidth - 4\tabcolsep) * \real{0.1122}}
  >{\raggedright\arraybackslash}p{(\linewidth - 4\tabcolsep) * \real{0.0510}}@{}}
\toprule\noalign{}
\begin{minipage}[b]{\linewidth}\raggedright
Category
\end{minipage} & \begin{minipage}[b]{\linewidth}\centering
Score (0--3)
\end{minipage} & \begin{minipage}[b]{\linewidth}\raggedright
Notes
\end{minipage} \\
\midrule\noalign{}
\endhead
\bottomrule\noalign{}
\endlastfoot
Problem \& Solution --- modification is purposeful; print(s) successful
& & \\
Design \& Code Quality --- modification is documented; parametric if done in OpenSCAD
& & \\
Documentation --- variant table complete; scaling effects addressed &
& \\
\textbf{Total (0--9)} & & \\
\end{longtable}
}

\begin{center}\rule{0.5\linewidth}{0.5pt}\end{center}

\subsection*{Appendix D: Real-Life Problem Solver --- Extension Project}\label{docs__pandoc__latex__src__3dmake_all_lessons.md__appendix-d-real-life-problem-solver--extension-project}

\textbf{Estimated time:} 4--8 hours (design, adapt, print, and document)

\subsubsection*{Overview}\label{docs__pandoc__latex__src__3dmake_all_lessons.md__overview-3}

This extension project asks you to do something harder than downloading and printing: identify a real problem in your life or environment, find or design a 3D-printed solution, adapt it to meet the actual constraints of that situation, and document the full process. This is the closest this course gets to professional design work at this stage --- it requires a problem statement, candidate research, a design decision, iteration, and evaluation.

\subsubsection*{Learning Objectives}\label{docs__pandoc__latex__src__3dmake_all_lessons.md__learning-objectives-4}

By completing this project, you will:

\begin{itemize}
\tightlist
\item
  Identify a user problem and write a problem statement that defines what success looks like
\item
  Evaluate candidate printed solutions and choose one to adapt
\item
  Make at least one documented design modification
\item
  Test and measure a prototype with documented results
\item
  Write a final report summarizing performance and next steps
\end{itemize}

\subsubsection*{Materials}\label{docs__pandoc__latex__src__3dmake_all_lessons.md__materials-2}

\begin{itemize}
\tightlist
\item
  Computer with slicer and access to repositories
\item
  Printer, filament, basic hand tools for post-processing
\end{itemize}

\subsubsection*{Tasks}\label{docs__pandoc__latex__src__3dmake_all_lessons.md__tasks-2}

\begin{enumerate}
\item
  \textbf{Write a problem statement.} Interview a potential user (yourself counts) and write one paragraph that describes: what the person is trying to do, what gets in their way, and what a successful solution would let them do. Be specific --- "it\textquotesingle s inconvenient" is not a problem statement; "my headphones fall off the desk and the cable gets tangled" is.
\item
  \textbf{Research candidate models.} Search Thingiverse or Printables for at least three options that could address your problem. List all three in your documentation, explain why each could or couldn\textquotesingle t work, and justify your final selection.
\item
  \textbf{Adapt the model.} Make at least one modification --- scale it, add or remove features, combine it with another model, or change the material or infill for durability. Document your changes in a short changelog with one entry per change.
\item
  \textbf{Print and test.} Run a supervised test of the prototype. Log any failures and what corrective actions you took. If the first print fails or doesn\textquotesingle t work as intended, print again with documented changes.
\item
  \textbf{Write a final report} summarizing: how well the prototype performed, measured deviations from specifications, any remaining limitations, and what you would change in the next iteration.
\end{enumerate}

\subsubsection*{Probing Questions}\label{docs__pandoc__latex__src__3dmake_all_lessons.md__probing-questions-2}

\begin{enumerate}
\tightlist
\item
  What assumptions did you make about the user\textquotesingle s context? How could you validate those assumptions with someone other than yourself?
\item
  Which modification had the biggest impact on whether the part worked, and why?
\end{enumerate}

\begin{center}\rule{0.5\linewidth}{0.5pt}\end{center}

\subsection*{Appendix E: Real-Life Problem Solver --- Student Documentation Template}\label{docs__pandoc__latex__src__3dmake_all_lessons.md__appendix-e-real-life-problem-solver--student-documentation-template}

\textbf{Author:} \_\_\_\_\_\_\_\_\_\_\_\_\_\_\_\_
\textbf{Date:} \_\_\_\_\_\_\_\_\_\_
\textbf{Description:} Adapt an existing 3D model to solve a real problem, with documented design changes and test results.

\begin{center}\rule{0.5\linewidth}{0.5pt}\end{center}

\subsubsection*{Problem Statement}\label{docs__pandoc__latex__src__3dmake_all_lessons.md__problem-statement}

Write one paragraph. Who is the user? What are they trying to do? What gets in their way? What would a successful solution let them do?

\begin{center}\rule{0.5\linewidth}{0.5pt}\end{center}

\subsubsection*{Candidate Models}\label{docs__pandoc__latex__src__3dmake_all_lessons.md__candidate-models}

{\def\LTcaptype{none} % do not increment counter
\begin{longtable}[]{@{}
  >{\centering\arraybackslash}p{(\linewidth - 6\tabcolsep) * \real{0.0484}}
  >{\raggedright\arraybackslash}p{(\linewidth - 6\tabcolsep) * \real{0.3387}}
  >{\raggedright\arraybackslash}p{(\linewidth - 6\tabcolsep) * \real{0.2742}}
  >{\raggedright\arraybackslash}p{(\linewidth - 6\tabcolsep) * \real{0.3387}}@{}}
\toprule\noalign{}
\begin{minipage}[b]{\linewidth}\centering
\#
\end{minipage} & \begin{minipage}[b]{\linewidth}\raggedright
Model name and source
\end{minipage} & \begin{minipage}[b]{\linewidth}\raggedright
Why it could work
\end{minipage} & \begin{minipage}[b]{\linewidth}\raggedright
Why it might not work
\end{minipage} \\
\midrule\noalign{}
\endhead
\bottomrule\noalign{}
\endlastfoot
1 & & & \\
2 & & & \\
3 & & & \\
\end{longtable}
}

\textbf{Selected model:} \_\_\_\_\_ \textbf{Reason for selection:}

\begin{center}\rule{0.5\linewidth}{0.5pt}\end{center}

\subsubsection*{Modification Changelog}\label{docs__pandoc__latex__src__3dmake_all_lessons.md__modification-changelog}

{\def\LTcaptype{none} % do not increment counter
\begin{longtable}[]{@{}clll@{}}
\toprule\noalign{}
\# & Change made & Why & Result \\
\midrule\noalign{}
\endhead
\bottomrule\noalign{}
\endlastfoot
1 & & & \\
2 & & & \\
3 & & & \\
\end{longtable}
}

\begin{center}\rule{0.5\linewidth}{0.5pt}\end{center}

\subsubsection*{Original vs. Adapted Specifications}\label{docs__pandoc__latex__src__3dmake_all_lessons.md__original-vs-adapted-specifications}

{\def\LTcaptype{none} % do not increment counter
\begin{longtable}[]{@{}lccl@{}}
\toprule\noalign{}
Parameter & Original & Adapted & Reason \\
\midrule\noalign{}
\endhead
\bottomrule\noalign{}
\endlastfoot
& & & \\
& & & \\
& & & \\
\end{longtable}
}

\begin{center}\rule{0.5\linewidth}{0.5pt}\end{center}

\subsubsection*{Design Iteration Cycle}\label{docs__pandoc__latex__src__3dmake_all_lessons.md__design-iteration-cycle}

\paragraph*{Print 1}\label{docs__pandoc__latex__src__3dmake_all_lessons.md__print-1}

\begin{itemize}
\tightlist
\item
  Date:
\item
  Settings (layer height, infill, supports):
\item
  Result:
\item
  Issues observed:
\item
  Next action:
\end{itemize}

\paragraph*{Print 2 (if made)}\label{docs__pandoc__latex__src__3dmake_all_lessons.md__print-2-if-made}

\begin{itemize}
\tightlist
\item
  Date:
\item
  What changed from Print 1:
\item
  Result:
\item
  Issues resolved / remaining:
\end{itemize}

\begin{center}\rule{0.5\linewidth}{0.5pt}\end{center}

\subsubsection*{Test Results}\label{docs__pandoc__latex__src__3dmake_all_lessons.md__test-results}

How did you test whether the prototype solves the problem? What did you observe?

\begin{itemize}
\tightlist
\item
  Test method:
\item
  Did the prototype succeed? Yes / Partially / No
\item
  Any deformation, failure, or fit issues?
\item
  User feedback (if you had a second person try it):
\end{itemize}

\begin{center}\rule{0.5\linewidth}{0.5pt}\end{center}

\subsubsection*{Reflections}\label{docs__pandoc__latex__src__3dmake_all_lessons.md__reflections-2}

\textbf{What worked:}

\textbf{What didn\textquotesingle t work:}

\textbf{What you would change in iteration 3:}

\textbf{What this taught you about the gap between a design on screen and a physical object that works in real life:}

\begin{center}\rule{0.5\linewidth}{0.5pt}\end{center}

\subsubsection*{Accessibility Considerations}\label{docs__pandoc__latex__src__3dmake_all_lessons.md__accessibility-considerations}

\begin{itemize}
\tightlist
\item
  How would you describe this modified design to someone who cannot see it?
\item
  What measurements or tactile features would help a non-visual user understand or use the part?
\end{itemize}

\begin{center}\rule{0.5\linewidth}{0.5pt}\end{center}

\subsubsection*{Attachments Checklist}\label{docs__pandoc__latex__src__3dmake_all_lessons.md__attachments-checklist-2}

\begin{itemize}
\tightlist
\item[$\square$]
  Problem statement (one paragraph, specific)
\item[$\square$]
  Three candidate models with analysis
\item[$\square$]
  Modification changelog
\item[$\square$]
  \texttt{.scad} or \texttt{.stl} files for each iteration
\item[$\square$]
  Photos of each printed iteration
\item[$\square$]
  Test results documentation
\item[$\square$]
  Final report (performance summary and next steps)
\end{itemize}

\begin{center}\rule{0.5\linewidth}{0.5pt}\end{center}

\subsubsection*{Teacher Feedback}\label{docs__pandoc__latex__src__3dmake_all_lessons.md__teacher-feedback-2}

{\def\LTcaptype{none} % do not increment counter
\begin{longtable}[]{@{}
  >{\raggedright\arraybackslash}p{(\linewidth - 4\tabcolsep) * \real{0.8298}}
  >{\centering\arraybackslash}p{(\linewidth - 4\tabcolsep) * \real{0.1170}}
  >{\raggedright\arraybackslash}p{(\linewidth - 4\tabcolsep) * \real{0.0532}}@{}}
\toprule\noalign{}
\begin{minipage}[b]{\linewidth}\raggedright
Category
\end{minipage} & \begin{minipage}[b]{\linewidth}\centering
Score (0--3)
\end{minipage} & \begin{minipage}[b]{\linewidth}\raggedright
Notes
\end{minipage} \\
\midrule\noalign{}
\endhead
\bottomrule\noalign{}
\endlastfoot
Problem \& Solution --- modification functional; problem addressed &
& \\
Design \& Code Quality --- changes documented; iteration evident & & \\
Documentation --- changelog complete; test results recorded; reflection specific
& & \\
\textbf{Total (0--9)} & & \\
\end{longtable}
}

\textbf{Resubmission:} attach a paragraph explaining what was changed and why. Resubmission score replaces original.

\begin{center}\rule{0.5\linewidth}{0.5pt}\end{center}

\subsection*{Appendix F: Real-Life Problem Solver --- Instructor Reference}\label{docs__pandoc__latex__src__3dmake_all_lessons.md__appendix-f-real-life-problem-solver--instructor-reference}

\subsubsection*{Project Context}\label{docs__pandoc__latex__src__3dmake_all_lessons.md__project-context-1}

This is the most demanding extension project in Lessons 1--2. It requires students to work through a full design loop: problem definition → research → selection → modification → testing → evaluation. The quality of the outcome depends heavily on the quality of the problem statement --- students who write a vague problem statement produce vague designs. Front-load the problem statement review.

\subsubsection*{Key Learning Points to Reinforce}\label{docs__pandoc__latex__src__3dmake_all_lessons.md__key-learning-points-to-reinforce-1}

\textbf{A problem statement is not a solution.} Students tend to jump to "I want to print a phone stand" without explaining why. The problem statement should describe a situation and a gap, not a predetermined output. "My phone falls off my desk and I can\textquotesingle t see notifications while working" is a problem statement. "I want a phone stand" is a solution.

\textbf{A changelog creates accountability.} Every design change should have a reason. Students who change things without documenting why have no basis for evaluating whether the change helped. Reinforce this by asking "why did you make this change?" at every iteration review.

\textbf{Testing must be observable.} A test that consists of "I tried it and it worked" is not documentation. Students should describe the test method (how many times, under what conditions, measured how), not just the result.

\subsubsection*{Constraints (Enforce These)}\label{docs__pandoc__latex__src__3dmake_all_lessons.md__constraints-enforce-these-1}

\begin{itemize}
\tightlist
\item
  Problem statement must be written before model selection
\item
  At least one print iteration is required --- first-attempt submissions are not acceptable
\item
  Changelog must have at least one entry per design change
\item
  Test results must include an observable measurement or observation, not just a pass/fail judgment
\end{itemize}

\subsubsection*{Assessment Notes}\label{docs__pandoc__latex__src__3dmake_all_lessons.md__assessment-notes-1}

\textbf{Strong submissions} show: a specific problem statement that identifies a real user, a justified model selection (not just "I found this one"), a changelog that explains the reasoning behind each change, and a test result that describes what was actually measured or observed.

\textbf{Common weak areas:} problem statements that are really solution statements, modification changelogs with a single vague entry, and test sections that say "it worked fine" without any specifics. Push back on all three in feedback.

\textbf{Extension for advanced students:} have them repeat the design loop with a second user who is different from themselves. Designing for yourself and discovering that a different person has different constraints is a powerful experience.
-e

\begin{center}\rule{0.5\linewidth}{0.5pt}\end{center}

\section*{Lesson 3 --- Parametric Architecture and Modular Libraries}\label{docs__pandoc__latex__src__3dmake_all_lessons.md__lesson-3--parametric-architecture-and-modular-libraries}

\textbf{Estimated Time:} 90--120 minutes\\
\textbf{Course:} 3dMake Certification --- High School Distance Learning Track

\begin{center}\rule{0.5\linewidth}{0.5pt}\end{center}

\subsection*{Before You Start}\label{docs__pandoc__latex__src__3dmake_all_lessons.md__before-you-start-1}

You can now create shapes and combine them using CSG. That\textquotesingle s the vocabulary. Now it\textquotesingle s time to learn the \emph{grammar} --- how to organize and generalize your code so it\textquotesingle s maintainable, reusable, and professional.

The difference between a beginner OpenSCAD file and a professional one isn\textquotesingle t usually the shapes --- it\textquotesingle s the structure. A beginner file has magic numbers scattered throughout, the same geometry copy-pasted in three places, and no indication of what any variable is for. A professional file has clearly named parameters at the top, organized modules, validation that catches errors before they become bad prints, and code you could hand to someone else and they\textquotesingle d immediately understand.

This lesson introduces the complete system for professional OpenSCAD structure: modules, functions, libraries, and the DRY principle. It also covers one of the most confusing aspects of OpenSCAD for programmers coming from other languages: how variables work here is completely different from Python, JavaScript, or Java --- and understanding that difference upfront will save you hours of debugging.

\begin{center}\rule{0.5\linewidth}{0.5pt}\end{center}

\subsection*{Learning Objectives}\label{docs__pandoc__latex__src__3dmake_all_lessons.md__learning-objectives-5}

By the end of this lesson you will be able to:

\begin{itemize}
\tightlist
\item
  Understand and work with OpenSCAD\textquotesingle s unique "last-wins" variable scoping rules
\item
  Write parametric modules with validated inputs using \texttt{assert()}
\item
  Organize reusable geometry into library files using \texttt{use} and \texttt{include}
\item
  Write recursive functions for mathematical and spatial computations
\item
  Apply list comprehensions and vector operations for advanced parametric design
\item
  Apply the DRY (Don\textquotesingle t Repeat Yourself) principle through modules and functions
\item
  Use \texttt{for} loops, conditionals, and list comprehensions to generate repetitive geometry
\item
  Write type-safe modules using \texttt{is\_number()}, \texttt{is\_list()}, and \texttt{is\_string()}
\end{itemize}

\begin{center}\rule{0.5\linewidth}{0.5pt}\end{center}

\subsection*{Concept 1: How Variables Work in OpenSCAD (It\textquotesingle s Different)}\label{docs__pandoc__latex__src__3dmake_all_lessons.md__concept-1-how-variables-work-in-openscad-its-different}

\subsubsection*{The Most Important Thing to Understand in This Lesson}\label{docs__pandoc__latex__src__3dmake_all_lessons.md__the-most-important-thing-to-understand-in-this-lesson}

If you\textquotesingle ve ever programmed in Python, JavaScript, Java, or almost any other language, you have a mental model of variables that works a certain way: you assign a value, the program runs, and that variable holds that value until you change it. Variables are like buckets --- you put a value in, you can read what\textquotesingle s in there, and you can put a new value in later.

\textbf{OpenSCAD does not work this way.} The difference is fundamental and will cause bugs if you misunderstand it. Let\textquotesingle s dig in.

\subsubsection*{How OpenSCAD Actually Works}\label{docs__pandoc__latex__src__3dmake_all_lessons.md__how-openscad-actually-works}

OpenSCAD processes your file in two distinct passes:

\textbf{Pass 1 --- Parse:} OpenSCAD reads the entire file from top to bottom, collects every variable assignment and module/function definition, and resolves what every variable equals. No geometry is created during this pass.

\textbf{Pass 2 --- Render:} OpenSCAD creates geometry using the values that were determined in Pass 1.

This two-pass system leads to three rules that govern all OpenSCAD code:

\textbf{Rule 1: There is no assignment during execution.} You cannot update a variable while the geometry is being computed. Variables are set during parse time and then read during render time. "Updating" a variable mid-computation isn\textquotesingle t a concept in OpenSCAD.

\textbf{Rule 2: The last assignment in a scope wins --- everywhere in that scope.} This is the most confusing rule. If you write \texttt{x\ =\ 5;} on line 3 and then \texttt{x\ =\ 10;} on line 20, the value of \texttt{x} is \texttt{10} \emph{everywhere} in that scope --- including on line 3, even though you hadn\textquotesingle t written the second assignment yet. It\textquotesingle s as if OpenSCAD reads the entire file first, decides what everything equals, and then uses those final values. Writing \texttt{x\ =\ 5;\ echo(x);\ x\ =\ 10;} will print \texttt{10}, not \texttt{5}.

\textbf{Rule 3: Variables inside modules are local.} A variable declared inside a module exists only within that module. The global variable with the same name is unchanged.

\begin{lstlisting}[style=Alabaster, language=openscad]
// COMMON MISTAKE — expecting x to change during execution

x = 5;

echo("x =", x);   // This prints 10, NOT 5 — because the last assignment wins

x = 10;


// CORRECT PATTERN — each value gets its own name

base_size    = 20;          // 20

padded_size  = base_size + 10;   // 30 — computed from base_size

doubled_size = base_size * 2;    // 40 — computed from base_size

echo(base_size, padded_size, doubled_size);  // 20, 30, 40


// MODULE SCOPE — variables inside a module are completely separate

wall_t = 3;   // global wall thickness


module demo_box(size) {

  wall_t = 1;           // This is a DIFFERENT variable — local to this module

  cube([size, size, wall_t]);   // uses the local wall_t = 1

}


demo_box(20);

echo("Global wall_t is still:", wall_t);  // Still 3 — unchanged by the module


// CHOOSING BETWEEN VALUES conditionally — use a ternary expression

mode = "large";

box_size = (mode == "large") ? 40 : 20;   // If mode is "large", 40; otherwise 20

cube([box_size, box_size, 5]);


\end{lstlisting}

\subsubsection*{What This Means in Practice}\label{docs__pandoc__latex__src__3dmake_all_lessons.md__what-this-means-in-practice}

Because you can\textquotesingle t update a variable during execution, you can\textquotesingle t accumulate a total like \texttt{total\ =\ total\ +\ item}. Instead, you describe \emph{relationships} between values mathematically. If the width of your inner wall depends on the outer wall and the wall thickness, you write that as a formula: \texttt{inner\_width\ =\ outer\_width\ -\ 2\ *\ wall\_t;}. The system computes the value once, and that value is used everywhere.

This is actually quite elegant once you adjust to it. You\textquotesingle re writing equations, not instructions. Each variable is defined once and means exactly one thing.

For cases where you genuinely need to choose between values or compute something iteratively, OpenSCAD provides \textbf{ternary expressions} (condition ? value\_if\_true : value\_if\_false), \textbf{recursive functions}, and \textbf{list comprehensions} --- all covered in this lesson.

\begin{center}\rule{0.5\linewidth}{0.5pt}\end{center}

\subsection*{Step 1 --- Modules: Named, Reusable Geometry}\label{docs__pandoc__latex__src__3dmake_all_lessons.md__step-1--modules-named-reusable-geometry}

\subsubsection*{What a Module Is}\label{docs__pandoc__latex__src__3dmake_all_lessons.md__what-a-module-is}

A \textbf{module} is a named, reusable piece of geometry. When you find yourself writing the same shape code in multiple places, that\textquotesingle s the signal to create a module. Modules accept parameters --- just like functions in other languages --- so you can write one module that creates many different variations of the same shape.

The syntax is:

\begin{lstlisting}[style=Alabaster, language=openscad]
module name(param1, param2=default_value) {

  // geometry code here

  // uses param1 and param2

}


\end{lstlisting}

Parameters with default values (\texttt{param2=default\_value}) are optional --- callers can omit them and the default will be used. Parameters without defaults are required.

You call a module just like a statement:

\begin{lstlisting}[style=Alabaster, language=openscad]
name(value1, value2);


\end{lstlisting}

\subsubsection*{The DRY Principle}\label{docs__pandoc__latex__src__3dmake_all_lessons.md__the-dry-principle}

Tip

\textbf{DRY --- Don\textquotesingle t Repeat Yourself.} If geometry is defined in one module, changing it once updates everywhere it\textquotesingle s used. Copy-pasting the same geometry multiple times means one missed update creates inconsistency across your design.

If you need four mounting holes in a plate, you could write the cylinder code four times --- once for each hole. But now if you need to change the hole size, you have to change it in four places. If you miss one, your design is inconsistent. If you need to add a fifth hole later, you have to remember the pattern and paste it again.

The DRY approach: write a module \texttt{mounting\_hole(r,\ depth)} once, then call it four times. Now the hole geometry is defined in exactly one place. Change it once, every hole updates. Add more holes as easily as one line each.

\begin{lstlisting}[style=Alabaster, language=openscad]
// WRONG — copy-pasted, hard to maintain

difference() {

  cube([60, 40, 5]);

  translate([5,  5,  -0.001]) cylinder(r=2, h=5.002, $fn=16);

  translate([55, 5,  -0.001]) cylinder(r=2, h=5.002, $fn=16);

  translate([5,  35, -0.001]) cylinder(r=2, h=5.002, $fn=16);

  translate([55, 35, -0.001]) cylinder(r=2, h=5.002, $fn=16);

}


// CORRECT — DRY: one module, used multiple times

module mounting_hole(r=2, depth=5) {

  // The 0.001 offsets prevent co-planar face problems (Lesson 2)

  translate([0, 0, -0.001])

    cylinder(r=r, h=depth + 0.002, $fn=16);

}


difference() {

  cube([60, 40, 5]);

  translate([5,  5,  0]) mounting_hole();

  translate([55, 5,  0]) mounting_hole();

  translate([5,  35, 0]) mounting_hole();

  translate([55, 35, 0]) mounting_hole();

}


\end{lstlisting}

The second version is easier to read, easier to change, and scales effortlessly. If you need a larger hole, change \texttt{r=2} to \texttt{r=3} in one place, and all four holes update.

\begin{center}\rule{0.5\linewidth}{0.5pt}\end{center}

\subsection*{\texorpdfstring{Step 2 --- \texttt{assert()} for Parameter Validation}{Step 2 --- assert() for Parameter Validation}}\label{docs__pandoc__latex__src__3dmake_all_lessons.md__step-2--assert-for-parameter-validation}

\subsubsection*{Why Your Modules Should Validate Their Inputs}\label{docs__pandoc__latex__src__3dmake_all_lessons.md__why-your-modules-should-validate-their-inputs}

In a parametric design, other people (and future-you) will change parameter values. What happens if someone sets \texttt{wall\_thickness\ =\ 0}? Or \texttt{cylinder\_radius\ =\ -5}? Depending on the code, it might silently produce garbage geometry, or it might produce a cryptic error that doesn\textquotesingle t explain what went wrong.

\texttt{assert()} is OpenSCAD\textquotesingle s built-in validation statement. It takes a condition and an optional message. If the condition is false, OpenSCAD \emph{stops rendering immediately} and displays your message in the console. This is vastly better than allowing invalid values to silently produce a broken model.

Think of \texttt{assert()} as making your design constraints machine-enforceable. Instead of just knowing that \texttt{wall\_t} must be at least 1.2 mm for a 0.4 mm nozzle, you write it as an assertion and the tool enforces it automatically.

\subsubsection*{\texorpdfstring{How \texttt{assert()} Works}{How assert() Works}}\label{docs__pandoc__latex__src__3dmake_all_lessons.md__how-assert-works}

\begin{lstlisting}[style=Alabaster, language=openscad]
assert(condition, "message if condition is false");


\end{lstlisting}

If \texttt{condition} evaluates to true, nothing happens --- the code continues normally. If \texttt{condition} is false, the render stops and the message is printed. The message should be descriptive enough that anyone reading it immediately understands what went wrong.

\begin{lstlisting}[style=Alabaster, language=openscad]
module safe_cylinder(r, h, fn=32) {

  // Validate all inputs before creating any geometry

  assert(is_num(r) && r > 0,

    str("safe_cylinder: radius must be a positive number, got: ", r));

  assert(is_num(h) && h > 0,

    str("safe_cylinder: height must be a positive number, got: ", h));

  assert(fn >= 3,

    str("safe_cylinder: fn must be at least 3, got: ", fn));

  assert(r < 500,

    str("safe_cylinder: radius ", r, "mm seems wrong — did you enter cm instead of mm?"));


  cylinder(r=r, h=h, $fn=fn);

}


safe_cylinder(5, 20);          // fine

// safe_cylinder(-3, 20);      // STOPS: "radius must be a positive number, got: -3"

// safe_cylinder(5, 20, fn=1); // STOPS: "fn must be at least 3, got: 1"


\end{lstlisting}

The \texttt{str()} function assembles a string from multiple values --- very useful for error messages that include the bad value so the user knows exactly what to fix.

Tip

Add \texttt{assert()} statements to every module at the very top of the module body, before any geometry. This turns bad-input problems into clear, descriptive messages instead of mysterious rendering failures.

\begin{center}\rule{0.5\linewidth}{0.5pt}\end{center}

\subsection*{\texorpdfstring{Step 3 --- Library Files: \texttt{use} vs. \texttt{include}}{Step 3 --- Library Files: use vs. include}}\label{docs__pandoc__latex__src__3dmake_all_lessons.md__step-3--library-files-use-vs-include}

\subsubsection*{Organizing Code Across Multiple Files}\label{docs__pandoc__latex__src__3dmake_all_lessons.md__organizing-code-across-multiple-files}

Once you have a collection of useful modules --- mounting holes, standoffs, snap clips, text labels --- you\textquotesingle ll want to reuse them across multiple projects without copy-pasting them every time. The solution is a \textbf{library file}: a separate \texttt{.scad} file containing only module and function definitions, which you can import into any project.

There are two ways to load a library:

\textbf{\texttt{use\ \textless{}filename.scad\textgreater{}}} --- imports only the module and function \textbf{definitions} from the file. Any top-level geometry statements in the file (bare \texttt{cube()}, \texttt{cylinder()}, etc.) are \emph{not} executed. This is almost always what you want: you want the tools, not any leftover test geometry the library file might contain.

\textbf{\texttt{include\ \textless{}filename.scad\textgreater{}}} --- executes the \textbf{entire file}, including any top-level geometry statements. If the library file has a \texttt{cube({[}10,10,10{]});} at the top level, that cube will appear in your model. Use this sparingly and intentionally.

Note

\texttt{use\ \textless{}file.scad\textgreater{}} imports only module/function \textbf{definitions} --- top-level geometry in the file is ignored. \texttt{include\ \textless{}file.scad\textgreater{}} executes the \textbf{entire file} including any top-level geometry. Prefer \texttt{use} for libraries to avoid unexpected geometry.

\begin{lstlisting}[style=Alabaster, language=openscad]
// === lib/hardware.scad ===

// A reusable hardware library. Load with: use <../lib/hardware.scad>


module m3_bolt_hole(depth=10, countersink=false) {

  // M3 clearance hole: 3.4mm diameter (ISO standard M3 clearance)

  assert(depth > 0, "m3_bolt_hole: depth must be positive");

  translate([0, 0, -0.001]) {

    cylinder(d=3.4, h=depth + 0.002, $fn=16);

    if (countersink) {

      // Add a countersunk cone for flat-head screws

      translate([0, 0, depth - 2.5])

        cylinder(d1=3.4, d2=7.0, h=3, $fn=32);

    }

  }

}


module m3_standoff(height=5, outer_d=6) {

  // Hollow cylindrical standoff with M3 hole through center

  assert(height > 2, "m3_standoff: height must be > 2mm");

  difference() {

    cylinder(d=outer_d, h=height, $fn=16);

    m3_bolt_hole(depth=height);  // Modules can call other modules from the library

  }

}


\end{lstlisting}

\begin{lstlisting}[style=Alabaster, language=openscad]
// === src/main.scad ===

use <../lib/hardware.scad>   // loads definitions; no top-level geometry executes


// Place M3 standoffs at all four corners of a mounting plate

cube([60, 40, 5]);

for (pos = [[5,5], [55,5], [5,35], [55,35]]) {

  translate([pos[0], pos[1], 0])

    m3_standoff(height=8);

}


\end{lstlisting}

\subsubsection*{Structuring a Good Library}\label{docs__pandoc__latex__src__3dmake_all_lessons.md__structuring-a-good-library}

A well-organized library groups related modules, documents each one, and validates inputs. Here\textquotesingle s a template:

\begin{lstlisting}[style=Alabaster, language=openscad]
// ============================================================

// fasteners.scad — Parametric fastener library

// Usage: use <fasteners.scad>

// All dimensions follow ISO metric standards.

// ============================================================


module m3_hole(depth=10) {

  // M3 clearance hole: 3.2mm diameter

  cylinder(r=1.6, h=depth + 0.001, $fn=16);

}


module m4_hole(depth=10) {

  // M4 clearance hole: 4.3mm diameter

  cylinder(r=2.15, h=depth + 0.001, $fn=16);

}


module m3_nut_trap(extra_depth=0) {

  // M3 hex nut trap: 6.4mm across-flats, 2.4mm thick

  cylinder(r=3.2, h=2.4 + extra_depth, $fn=6);  // $fn=6 makes a hexagon

}


\end{lstlisting}

\begin{center}\rule{0.5\linewidth}{0.5pt}\end{center}

\subsection*{\texorpdfstring{Step 4 --- \texttt{for} Loops and Conditionals}{Step 4 --- for Loops and Conditionals}}\label{docs__pandoc__latex__src__3dmake_all_lessons.md__step-4--for-loops-and-conditionals}

\subsubsection*{\texorpdfstring{Repeating Geometry with \texttt{for}}{Repeating Geometry with for}}\label{docs__pandoc__latex__src__3dmake_all_lessons.md__repeating-geometry-with-for}

OpenSCAD\textquotesingle s \texttt{for} loop iterates over a list or range and creates geometry for each iteration. Each iteration is \textbf{independent} --- you can\textquotesingle t pass values between iterations. But you can create geometry in each one.

Three forms:

\texttt{for\ (variable\ =\ list)} --- iterates over every element in a list.

\texttt{for\ (variable\ =\ {[}start\ :\ end{]})} --- iterates over a range with a step of 1. \texttt{{[}0:5{]}} gives 0, 1, 2, 3, 4, 5.

\texttt{for\ (variable\ =\ {[}start\ :\ step\ :\ end{]})} --- iterates over a range with a specified step. \texttt{{[}0:45:315{]}} gives 0, 45, 90, 135, 180, 225, 270, 315.

\begin{lstlisting}[style=Alabaster, language=openscad]
// Place pegs at four specific positions

peg_positions = [[0,0], [20,0], [40,0], [20,20]];


for (pos = peg_positions) {

  translate([pos[0], pos[1], 0]) cylinder(r=3, h=10, $fn=32);

}


// Place cylinders every 45° around a circle (8 cylinders total)

for (angle = [0 : 45 : 315]) {

  rotate([0, 0, angle])

    translate([20, 0, 0])

      cylinder(r=2, h=5, $fn=16);

}


// Conditional: if show_lid is true, draw a lid

show_lid = true;

if (show_lid) {

  translate([0, 0, 30]) cube([50, 30, 3]);

}


\end{lstlisting}

\textbf{Evenly spacing around a circle:} If you want N objects evenly spaced, the angular step is \texttt{360\ /\ N}. For 12 objects, that\textquotesingle s 30° each. Use \texttt{{[}0\ :\ 360/N\ :\ 360\ -\ 360/N{]}} to avoid duplicating the first position at 360°.

\subsubsection*{\texorpdfstring{Using \texttt{for} with Multiple Variables}{Using for with Multiple Variables}}\label{docs__pandoc__latex__src__3dmake_all_lessons.md__using-for-with-multiple-variables}

You can use multiple \texttt{for} clauses in a single statement to create grids:

\begin{lstlisting}[style=Alabaster, language=openscad]
// 4×4 grid of cylinders — nested for in a single statement

for (x = [0:3], y = [0:3]) {

  translate([x * 20, y * 20, 0]) cylinder(r=3, h=5, $fn=16);

}


\end{lstlisting}

This is equivalent to two nested \texttt{for} loops but more compact.

\begin{center}\rule{0.5\linewidth}{0.5pt}\end{center}

\subsection*{Step 5 --- Recursive Functions}\label{docs__pandoc__latex__src__3dmake_all_lessons.md__step-5--recursive-functions}

\subsubsection*{Why OpenSCAD Needs Recursion}\label{docs__pandoc__latex__src__3dmake_all_lessons.md__why-openscad-needs-recursion}

Most programming languages let you write a loop that accumulates a value: \texttt{total\ =\ 0;\ for\ item\ in\ list:\ total\ +=\ item}. OpenSCAD doesn\textquotesingle t allow this --- variables are fixed after parse time. So when you need to compute something iteratively, you use \textbf{recursive functions}.

A recursive function is one that calls itself with a smaller version of the problem. Every recursive function has two parts:

\textbf{The base case} --- the input is simple enough to handle directly, so return an answer without recursing.

\textbf{The recursive case} --- reduce the problem slightly and call yourself with the smaller version, then combine the result.

\subsubsection*{Functions vs. Modules --- A Critical Distinction}\label{docs__pandoc__latex__src__3dmake_all_lessons.md__functions-vs-modules--a-critical-distinction}

Before we look at examples, there is an important rule in OpenSCAD that trips up almost every programmer coming from other languages:

\textbf{Functions return values (numbers, strings, lists). They cannot create geometry.}

\textbf{Modules create geometry. They cannot return values.}

Important

This is a hard boundary in OpenSCAD: \textbf{Functions return values --- they cannot create geometry. Modules create geometry --- they cannot return values.} A function with a \texttt{cylinder()} call inside it won\textquotesingle t work, and a module cannot return a value.

\begin{lstlisting}[style=Alabaster, language=openscad]
// FUNCTION — computes and returns a value

// This is a function because it's defined with "function" and uses "="

function factorial(n) = (n <= 1) ? 1 : n * factorial(n - 1);


echo(factorial(5));   // prints 120

echo(factorial(0));   // prints 1


// Another function: Fibonacci sequence

function fib(n) = (n <= 1) ? n : fib(n-1) + fib(n-2);


// Use the function to space pegs with Fibonacci spacing

// fib(2)=1, fib(3)=2, fib(4)=3, fib(5)=5, fib(6)=8, fib(7)=13, fib(8)=21, fib(9)=34

for (i = [2:9]) {

  translate([fib(i) * 4, 0, 0])

    cylinder(r=1.5, h=5, $fn=16);

}


// A useful function: generate N evenly-spaced values between start and stop

function linspace(start, stop, n) =

  (n <= 1)

    ? [start]

    : concat([start], linspace(start + (stop - start) / (n - 1), stop, n - 1));


positions = linspace(0, 50, 6);   // evaluates to [0, 10, 20, 30, 40, 50]

echo("Positions:", positions);


for (p = positions) {

  translate([p, 0, 0]) sphere(r=2, $fn=16);

}


\end{lstlisting}

\subsubsection*{Reading a Recursive Function}\label{docs__pandoc__latex__src__3dmake_all_lessons.md__reading-a-recursive-function}

The \texttt{factorial} function works like this:

\begin{itemize}
\tightlist
\item
  \texttt{factorial(5)} → is \texttt{n\ \textless{}=\ 1}? No. → return \texttt{5\ *\ factorial(4)}
\item
  \texttt{factorial(4)} → is \texttt{n\ \textless{}=\ 1}? No. → return \texttt{4\ *\ factorial(3)}
\item
  ...and so on until...
\item
  \texttt{factorial(1)} → is \texttt{n\ \textless{}=\ 1}? Yes. → return \texttt{1}
\end{itemize}

Then the chain unwinds: \texttt{1\ *\ 2\ *\ 3\ *\ 4\ *\ 5\ =\ 120}.

The ternary expression \texttt{(condition)\ ?\ value\_if\_true\ :\ value\_if\_false} is the standard way to write the base-case branch in OpenSCAD functions.

\begin{center}\rule{0.5\linewidth}{0.5pt}\end{center}

\subsection*{Step 6 --- List Comprehensions}\label{docs__pandoc__latex__src__3dmake_all_lessons.md__step-6--list-comprehensions}

\subsubsection*{What a List Comprehension Is}\label{docs__pandoc__latex__src__3dmake_all_lessons.md__what-a-list-comprehension-is}

A \textbf{list comprehension} is a compact way to generate a list by applying an expression to every element of a range or existing list, optionally filtering elements. Instead of writing a loop that builds a list step by step, you describe the entire list in one expression.

The syntax:

\begin{lstlisting}[style=Alabaster]
[for (variable = range_or_list) expression]

[for (variable = range_or_list) if (condition) expression]


\end{lstlisting}

List comprehensions are OpenSCAD\textquotesingle s answer to the fact that you can\textquotesingle t use \texttt{for} loops to accumulate values. If you need a list of computed values, a comprehension gives you the entire list at once.

\begin{lstlisting}[style=Alabaster, language=openscad]
// Generate heights: [5, 10, 15, 20, 25]

heights = [for (i = [1:5]) i * 5];

echo(heights);   // [5, 10, 15, 20, 25]


// Generate 8 evenly-spaced positions along the X axis

positions = [for (i = [0:7]) [i * 12, 0, 0]];

for (pos = positions) {

  translate(pos)

    cylinder(h = 5 + pos[0] * 0.3, r=3, $fn=16);  // height increases with X position

}


// Filter: generate only even numbers from 0 to 18

evens = [for (i = [0:9]) if (i % 2 == 0) i * 2];

echo(evens);   // [0, 4, 8, 12, 16]


// 2D grid of points — nested comprehension generates a flat list of [x,y,0] vectors

grid = [for (x = [0:3]) for (y = [0:3]) [x * 20, y * 20, 0]];

for (pt = grid) {

  translate(pt) cylinder(r=2, h=8, $fn=16);

}


\end{lstlisting}

List comprehensions are the preferred way in OpenSCAD to generate lists of positions, sizes, or other computed values that drive the placement of geometry. Master them and you\textquotesingle ll be able to generate complex repetitive patterns with very compact code.

\begin{center}\rule{0.5\linewidth}{0.5pt}\end{center}

\subsection*{Step 7 --- Vector Math and Math Functions}\label{docs__pandoc__latex__src__3dmake_all_lessons.md__step-7--vector-math-and-math-functions}

\subsubsection*{Working with Vectors}\label{docs__pandoc__latex__src__3dmake_all_lessons.md__working-with-vectors}

Positions and dimensions in 3D space are naturally expressed as vectors --- lists of three numbers representing X, Y, and Z. OpenSCAD supports arithmetic directly on vectors, which makes geometric calculations much more readable.

\begin{lstlisting}[style=Alabaster, language=openscad]
v1 = [10, 20, 5];

v2 = [3, 4, 0];


sum    = v1 + v2;    // [13, 24, 5]  — add matching components

diff   = v1 - v2;    // [7,  16, 5]  — subtract matching components

scaled = v1 * 2;     // [20, 40, 10] — multiply each component by a scalar

dot    = v1 * v2;    // 110          — dot product: v1·v2 = 10*3 + 20*4 + 5*0

length = norm(v1);   // ≈ 22.9       — Euclidean length (√(10² + 20² + 5²))


\end{lstlisting}

\subsubsection*{Important Math Functions}\label{docs__pandoc__latex__src__3dmake_all_lessons.md__important-math-functions}

\begin{lstlisting}[style=Alabaster, language=openscad]
abs(-5)         // 5    — absolute value

ceil(3.2)       // 4    — round up to nearest integer

floor(3.8)      // 3    — round down to nearest integer

round(3.5)      // 4    — round to nearest integer

max(3, 7, 2)    // 7    — largest of all arguments

min(3, 7, 2)    // 2    — smallest of all arguments

pow(2, 8)       // 256  — 2 raised to the power 8

sqrt(144)       // 12   — square root


>[!IMPORTANT]

>OpenSCAD trigonometry uses **degrees, not radians**. `sin(90)` = 1 (correct). If you pass radians by mistake, your geometry will be wildly wrong.


// Trigonometry — CRITICAL: OpenSCAD uses DEGREES, not radians

sin(90)         // 1    — sine of 90 degrees

cos(0)          // 1    — cosine of 0 degrees

tan(45)         // 1    — tangent of 45 degrees

atan2(1, 1)     // 45   — arctangent of y/x, safe for all quadrants


\end{lstlisting}

The \texttt{atan2(y,\ x)} function deserves special attention. When you need to convert a (x, y) position to an angle, the naive approach \texttt{atan(y/x)} fails when \texttt{x} is zero or negative (wrong quadrant). \texttt{atan2(y,\ x)} handles all cases correctly by taking both components separately.

\begin{center}\rule{0.5\linewidth}{0.5pt}\end{center}

\subsection*{\texorpdfstring{Step 8 --- Controlling Curve Quality with \texttt{\$fn}, \texttt{\$fa}, \texttt{\$fs}}{Step 8 --- Controlling Curve Quality with \$fn, \$fa, \$fs}}\label{docs__pandoc__latex__src__3dmake_all_lessons.md__step-8--controlling-curve-quality-with-fn-fa-fs}

\subsubsection*{Why You Need to Think About This}\label{docs__pandoc__latex__src__3dmake_all_lessons.md__why-you-need-to-think-about-this}

Spheres, cylinders, and circles in OpenSCAD are approximations --- they\textquotesingle re actually polygons with many flat faces. The more faces, the smoother the curve looks, but the longer the render takes and the larger the file.

\textbf{\texttt{\$fn}} (facets number) --- sets the exact number of faces. \texttt{\$fn=8} gives a rough octagon. \texttt{\$fn=64} gives a very smooth circle. Setting it on a single shape overrides the global value for that shape only.

\textbf{\texttt{\$fa}} (minimum angle per face) --- the smaller this is, the more faces large curves get. Default: 12°.

\textbf{\texttt{\$fs}} (minimum face size in mm) --- very small features automatically use more faces. Default: 2 mm.

When all three are specified, OpenSCAD uses whichever produces the \emph{most} faces.

\begin{lstlisting}[style=Alabaster, language=openscad]
// During development: low $fn for fast rendering

sphere(r=10, $fn=8);     // rough but fast


// For final export: high $fn for smooth curves

translate([30, 0, 0])

  sphere(r=10, $fn=64);  // smooth for printing


// Set globally at the top of a file:

$fa = 5;    // max 5 degrees per face

$fs = 0.5;  // minimum 0.5mm face size (more faces on small features)


\end{lstlisting}

\textbf{Recommended practice:} set \texttt{\$fn=16} or \texttt{\$fn=24} during development for fast iteration. Use \texttt{\$fn=64} or higher for spheres and visible curves in the final exported file. Use \texttt{\$fn=32} for functional cylindrical features like bolt holes --- it\textquotesingle s smooth enough for function without being slow to render.

\begin{center}\rule{0.5\linewidth}{0.5pt}\end{center}

\subsection*{Exercises}\label{docs__pandoc__latex__src__3dmake_all_lessons.md__exercises}

\textbf{Exercise 3.1:} Write a module \texttt{safe\_box(w,\ d,\ h,\ wall=2)} that validates all parameters with \texttt{assert()} and creates a hollow box. Add a constraint that \texttt{wall} must be less than half of the smallest dimension.

\textbf{Exercise 3.2:} Create a library file \texttt{lib/fasteners.scad} with modules for M2, M3, and M4 bolt holes. Use \texttt{use} to load it in a main file and place holes at the four corners of a 60 × 40 mm plate.

\textbf{Exercise 3.3:} Write a recursive function \texttt{sum\_list(lst,\ i=0)} that returns the sum of all elements in a list. Use it to compute the total volume of a set of cylinders.

\textbf{Exercise 3.4 (Advanced):} Use a list comprehension to generate a 5×5 grid of cylinders where each cylinder\textquotesingle s height equals \texttt{row\_index\ +\ col\_index\ +\ 1}. Use \texttt{echo()} to verify the maximum and minimum height values.

\begin{center}\rule{0.5\linewidth}{0.5pt}\end{center}

\subsection*{Quiz --- Lesson 3 (15 questions)}\label{docs__pandoc__latex__src__3dmake_all_lessons.md__quiz--lesson-3-15-questions}

\begin{enumerate}
\tightlist
\item
  What is the DRY principle and why is it important in parametric OpenSCAD design?
\item
  What is the difference between a module and a function in OpenSCAD?
\item
  What does \texttt{include\ \textless{}library.scad\textgreater{}} do differently from \texttt{use\ \textless{}library.scad\textgreater{}}? When would you use each?
\item
  If you write \texttt{x\ =\ 5;} on line 3 and \texttt{x\ =\ 10;} on line 20 in the same file, what is the value of \texttt{x} everywhere in that file? Why?
\item
  What does \texttt{is\_number()} return, and give one example of when to use it in a module.
\item
  What does the list comprehension \texttt{{[}for\ (i={[}0:5:20{]})\ i*2{]}} evaluate to?
\item
  What does the \texttt{echo()} function do in OpenSCAD, and where does its output appear?
\item
  True or False: a function in OpenSCAD can contain a \texttt{cylinder()} call and create geometry.
\item
  How do you specify a default parameter value in an OpenSCAD module? Write an example.
\item
  What is the ternary operator in OpenSCAD? Write an example that picks a box size based on a \texttt{"small"} or \texttt{"large"} variable.
\item
  Write an OpenSCAD \texttt{for} loop that places 8 cylinders evenly spaced around a circle of radius 25.
\item
  What is the difference between \texttt{{[}0:10{]}} and \texttt{{[}0:2:10{]}} as OpenSCAD range expressions?
\item
  Explain why you might use \texttt{assert()} in a module rather than relying on the caller to pass correct values.
\item
  Describe how a recursive function works. What are the two required parts of any recursive function?
\item
  A library file called \texttt{shapes.scad} contains a \texttt{module\ rounded\_box(...)} and also a top-level \texttt{cube({[}5,5,5{]});}. What happens when you load it with \texttt{use}? What happens when you load it with \texttt{include}?
\end{enumerate}

\begin{center}\rule{0.5\linewidth}{0.5pt}\end{center}

\subsection*{Extension Problems (15)}\label{docs__pandoc__latex__src__3dmake_all_lessons.md__extension-problems-15-2}

\begin{enumerate}
\tightlist
\item
  Create a \texttt{bolts.scad} library with M3, M4, and M5 clearance holes and nut traps. Use it in a test plate with one of each type.
\item
  Build a parametric pegboard using a \texttt{for} loop. Parameters: \texttt{cols}, \texttt{rows}, \texttt{spacing}, \texttt{peg\_r}, \texttt{peg\_h}. Every parameter should have an \texttt{assert()} validation.
\item
  Implement a recursive function that generates a Fibonacci sequence up to \texttt{n} terms and use it to space objects with Fibonacci distances.
\item
  Build a module that validates all its inputs using \texttt{assert()} and logs informational output with \texttt{echo()}. Document when \texttt{assert()} halts rendering vs. what \texttt{echo()} does.
\item
  Create two library files --- one for mechanical fasteners and one for decorative shapes --- and use both in one combined project.
\item
  Design a parametric gear profile module with \texttt{teeth}, \texttt{module\_size}, and \texttt{thickness} parameters. Use \texttt{assert()} to validate that the number of teeth is at least 8.
\item
  Build a "shelf peg" system: a board with N parametrically positioned holes and matching pegs sized to press-fit into them.
\item
  Use a list comprehension to generate a spiral path of evenly-spaced points and place a small sphere at each point on the spiral.
\item
  Extend a fastener library with threaded insert pockets for M3, M4, and M5 heat-set inserts, with the insert dimensions as parameters.
\item
  Write a module that draws a text label with a surrounding rectangular border frame. Parameters: \texttt{text\_str}, \texttt{font\_size}, \texttt{padding}, \texttt{border\_thickness}.
\item
  Build a parametric snap-fit lid for a rectangular box using \texttt{for} loops to generate evenly-spaced snap clips around the perimeter.
\item
  Write a module that takes a list of 3D coordinate vectors as input and draws a connecting chain of cylinders between consecutive points.
\item
  Create a parametric honeycomb panel module. Parameters: \texttt{rows}, \texttt{cols}, \texttt{cell\_r}, \texttt{wall\_t}, \texttt{depth}. Each cell should be a hexagonal through-hole.
\item
  Research the BOSL2 library (https://github.com/BelfrySCAD/BOSL2). Find two modules that solve problems you\textquotesingle ve encountered in this course, and write a documented example using each.
\item
  Design a comprehensive test suite: build 10 small parametric shapes each specifically testing a different OpenSCAD feature from this lesson. Export them all from one main file that uses a \texttt{for} loop.
\end{enumerate}

\begin{center}\rule{0.5\linewidth}{0.5pt}\end{center}

\subsection*{References and Helpful Resources}\label{docs__pandoc__latex__src__3dmake_all_lessons.md__references-and-helpful-resources-2}

\begin{itemize}
\item
  OpenSCAD User Manual --- Modules --- https://en.wikibooks.org/wiki/OpenSCAD\_User\_Manual/User-Defined\_Functions\_and\_Modules
\item
  OpenSCAD User Manual --- For Loops and Conditionals --- https://en.wikibooks.org/wiki/OpenSCAD\_User\_Manual/The\_OpenSCAD\_Language\#for\_loop
\item
  BOSL2 Library Documentation --- https://github.com/BelfrySCAD/BOSL2/wiki
\item
  OpenSCAD User Manual --- Include and Use --- https://en.wikibooks.org/wiki/OpenSCAD\_User\_Manual/Include\_Statement
\end{itemize}

\subsubsection*{Supplemental Resources}\label{docs__pandoc__latex__src__3dmake_all_lessons.md__supplemental-resources-2}

\begin{itemize}
\tightlist
\item
  \href{docs\%5Cpandoc\%5Clatex\%5Csrc\%5C943ee2446365181c.epub}{Programming with OpenSCAD EPUB Textbook} --- Chapters on modules, functions, and libraries
\item
  \href{https://github.com/ProgrammingWithOpenSCAD/CodeSolutions}{CodeSolutions Repository} --- Worked examples for parametric architecture
\item
  \href{https://github.com/BelfrySCAD/BOSL2}{BOSL2 GitHub Repository} --- Comprehensive parametric library
\item
  \href{https://github.com/tdeck/3dmake}{3DMake GitHub Repository} --- Build workflow reference
  -e
\end{itemize}

\begin{center}\rule{0.5\linewidth}{0.5pt}\end{center}

\section*{Lesson 4 --- AI-Enhanced Verification and Multimodal Feedback}\label{docs__pandoc__latex__src__3dmake_all_lessons.md__lesson-4--ai-enhanced-verification-and-multimodal-feedback}

\textbf{Estimated Time:} 90--120 minutes\\
\textbf{Course:} 3dMake Certification --- High School Distance Learning Track

\begin{center}\rule{0.5\linewidth}{0.5pt}\end{center}

\subsection*{Before You Start}\label{docs__pandoc__latex__src__3dmake_all_lessons.md__before-you-start-2}

Every lesson so far has ended with "build and verify." But what does verification actually mean? And when can you trust the result? This lesson is about that question.

The workflow you\textquotesingle ve been using --- \texttt{3dm\ build} followed by \texttt{3dm\ info} --- involves two types of tools: deterministic tools (which always give the same answer for the same input) and AI tools (which interpret and describe, but can be vague or wrong). Knowing the difference, and knowing which to use when, is one of the most important skills in modern digital fabrication.

Many beginners treat AI output as authoritative. When \texttt{3dm\ info} says "a small peg-like protrusion," they assume everything is correct. But AI descriptions can be vague, can miss features, and can be flat-out wrong about specifics. This lesson gives you a structured, layered verification approach that combines AI descriptions, deterministic tools, and your own judgment --- and makes clear which tool to reach for in which situation.

\begin{center}\rule{0.5\linewidth}{0.5pt}\end{center}

\subsection*{Learning Objectives}\label{docs__pandoc__latex__src__3dmake_all_lessons.md__learning-objectives-6}

By the end of this lesson you will be able to:

\begin{itemize}
\tightlist
\item
  Use \texttt{3dm\ info} effectively and know its real capabilities and limitations
\item
  Verify AI-generated OpenSCAD code systematically before printing
\item
  Build a personal verification chain for any critical design feature
\item
  Use \texttt{3dm\ preview} and slicer layer view as verification methods
\item
  Apply AI description to accessibility workflows for visually impaired makers
\item
  Distinguish deterministic from AI-assisted validation
\item
  Apply prompt engineering techniques for better AI design feedback
\item
  Recognize when to trust, question, or reject AI suggestions
\item
  Understand data privacy considerations when using cloud AI tools
\end{itemize}

\begin{center}\rule{0.5\linewidth}{0.5pt}\end{center}

\subsection*{Concept 1: Deterministic vs. AI-Based Validation}\label{docs__pandoc__latex__src__3dmake_all_lessons.md__concept-1-deterministic-vs-ai-based-validation}

\subsubsection*{Two Fundamentally Different Types of Tools}\label{docs__pandoc__latex__src__3dmake_all_lessons.md__two-fundamentally-different-types-of-tools}

When you run \texttt{3dm\ info} after a build, you get two very different types of information in the same output. It\textquotesingle s critical to understand the difference.

\textbf{Deterministic output} comes from directly reading the geometry of the STL file --- no interpretation involved. The bounding box (\texttt{30.0\ ×\ 20.0\ ×\ 10.0\ mm}) is computed from the actual vertex coordinates in the mesh. The volume is computed from the triangle faces. These numbers are produced by a formula applied to the file, and they will be identical every time you run the command on the same file. There is no randomness, no ambiguity, no "interpretation." If the bounding box says 30.0 mm, your model is 30.0 mm wide --- full stop.

\textbf{AI-based output} is the natural-language description portion of \texttt{3dm\ info}. An AI vision model "looks at" a rendered image of your model and generates a text description. This is fundamentally different: it involves interpretation, and interpretations can vary. The same model might be described as "a small peg" one time and "a narrow cylindrical protrusion" another time. The AI might miss a feature that is too small to be clearly visible. It might describe a 3.2 mm hole as "approximately 3 mm" or "roughly 3--4 mm." The description is genuinely useful for a quick sanity check, but it is not a measurement.

\subsubsection*{The Rule That Follows from This}\label{docs__pandoc__latex__src__3dmake_all_lessons.md__the-rule-that-follows-from-this}

\textbf{Always verify critical dimensions using deterministic methods first.} Bounding box, volume, triangle count --- these are computed from the mesh and cannot lie if the build succeeded. Only after confirming the numbers should you look at the AI description, and then only for orientation and general-shape confirmation.

\begin{center}\rule{0.5\linewidth}{0.5pt}\end{center}

\subsection*{\texorpdfstring{Step 1 --- Using \texttt{3dm\ info} Effectively}{Step 1 --- Using 3dm info Effectively}}\label{docs__pandoc__latex__src__3dmake_all_lessons.md__step-1--using-3dm-info-effectively}

\subsubsection*{\texorpdfstring{What \texttt{3dm\ info} Does}{What 3dm info Does}}\label{docs__pandoc__latex__src__3dmake_all_lessons.md__what-3dm-info-does}

\texttt{3dm\ info} runs after \texttt{3dm\ build} and produces a report about your built STL file. It always includes deterministic data (bounding box, volume, triangle count) and optionally generates an AI description by rendering the model from a specific angle.

\begin{lstlisting}[style=Alabaster, language=bash]
# Standard usage — run after every build

3dm build

3dm info


# Request a description from a specific viewing angle

3dm info --view front

3dm info --view top

3dm info --view isometric


# Save the description to a log file

3dm info --output description.txt


# Windows PowerShell

3dm info | Out-File -FilePath description.txt -Encoding UTF8


# Windows CMD

3dm info > description.txt


\end{lstlisting}

\subsubsection*{\texorpdfstring{What \texttt{3dm\ info} Does Well}{What 3dm info Does Well}}\label{docs__pandoc__latex__src__3dmake_all_lessons.md__what-3dm-info-does-well}

\texttt{3dm\ info} is genuinely useful for:

\textbf{Overall shape identification} --- confirming a part is roughly the right shape. If you designed a bracket and the description says "an L-shaped bracket," that\textquotesingle s a useful confirmation that the basic geometry is correct.

\textbf{Relative positions} --- confirming which part is on top of which, or that a protrusion sticks out from the side rather than the front.

\textbf{Component counting} --- "a base plate with four cylindrical protrusions" tells you the four pegs you expected are there and visible.

\textbf{Accessibility} --- a student who cannot easily see a rendered 3D model can use \texttt{3dm\ info} to get a verbal description of what they\textquotesingle ve built.

\subsubsection*{\texorpdfstring{What \texttt{3dm\ info} Does Poorly}{What 3dm info Does Poorly}}\label{docs__pandoc__latex__src__3dmake_all_lessons.md__what-3dm-info-does-poorly}

\texttt{3dm\ info} is \textbf{not} reliable for:

\textbf{Precise measurements} --- it cannot replace checking the bounding box or using calipers. "Approximately 10 mm" is not the same as "exactly 9.8 mm."

\textbf{Small features} --- features smaller than about 1--2 mm may not be clearly visible in the rendered view and might be ignored or described vaguely.

\textbf{Internal geometry} --- walls, cavities, and internal structures that aren\textquotesingle t visible on the exterior surface won\textquotesingle t be described.

\textbf{Tolerance-critical features} --- a 0.2 mm clearance gap that determines whether two parts fit together is invisible to visual AI.

\begin{center}\rule{0.5\linewidth}{0.5pt}\end{center}

\subsection*{Step 2 --- The Verification Escalation Chain}\label{docs__pandoc__latex__src__3dmake_all_lessons.md__step-2--the-verification-escalation-chain}

\subsubsection*{A Layered Approach to Verification}\label{docs__pandoc__latex__src__3dmake_all_lessons.md__a-layered-approach-to-verification}

When you need to confirm that your design is correct, don\textquotesingle t jump straight to printing. Instead, follow a layered chain of verification methods from most reliable (and fastest) to most authoritative (and slowest). Start at the top and only go further down if you need more certainty.

\textbf{Level 1 --- Bounding Box and Volume (always accurate)}

The very first thing to check after every build. If these numbers don\textquotesingle t match your design intent, something is wrong and you stop here to fix it before going further.

\begin{lstlisting}[style=Alabaster, language=bash]
3dm info

# Look for: Bounding Box: [expected width] × [expected depth] × [expected height] mm

# Look for: Volume: [reasonable volume] mm³


\end{lstlisting}

\textbf{Level 2 --- Slicer Layer Preview}

Open \texttt{build/main.stl} in your slicer (PrusaSlicer, Cura, etc.) and step through layers one at a time. Each layer is a cross-section through your model at a specific height --- like slicing a loaf of bread and looking at each slice. This reveals:

\begin{itemize}
\tightlist
\item
  Wall thicknesses that the bounding box can\textquotesingle t show
\item
  Internal geometry and cavities
\item
  Whether holes cut all the way through or stop short
\item
  Overhanging geometry that may need supports
\end{itemize}

The slicer layer view is the best non-physical verification tool available. Use it for anything more complex than a simple box.

\textbf{Level 3 --- Isolate and Test a Single Module}

If your full model is complex and \texttt{3dm\ info} gives a vague or confusing description, temporarily modify your file to render only the specific module you\textquotesingle re trying to verify:

\begin{lstlisting}[style=Alabaster, language=openscad]
// Temporarily comment out everything else, render only one module

// full_assembly();     // commented out

// base_plate();        // commented out

snap_clip();            // only this renders — easier for 3dm info to describe


\end{lstlisting}

Build and run \texttt{3dm\ info} on just that isolated piece. A simpler model gives more reliable AI descriptions and easier-to-check bounding boxes.

\textbf{Level 4 --- Physical Measurement}

For anything tolerance-critical --- a hole size that determines fit, a slot width for a specific component, a snap-fit clearance --- print a test coupon and measure with digital calipers. Physical measurement is the ground truth. Nothing replaces it for precise fit verification.

Lesson 8 covers test coupons in detail, but the principle is simple: design a small, fast-printing test piece that includes only the critical feature you\textquotesingle re validating. Print it, measure it, and adjust before printing the full part.

\begin{center}\rule{0.5\linewidth}{0.5pt}\end{center}

\subsection*{Step 3 --- A Verification Log}\label{docs__pandoc__latex__src__3dmake_all_lessons.md__step-3--a-verification-log}

\subsubsection*{Creating an Audit Trail}\label{docs__pandoc__latex__src__3dmake_all_lessons.md__creating-an-audit-trail}

When you verify a design feature, record which method you used and what it confirmed. This creates a professional audit trail --- evidence that you thought carefully about whether your design was correct.

Here\textquotesingle s a simple verification log format:

\begin{lstlisting}[style=Alabaster]
DESIGN VERIFICATION LOG

Project: phone_stand_v2.scad

Date: 2026-03-10


Feature: USB-C charging slot

  Designed: 9.5mm wide × 3.5mm tall

  Level 1 — Bounding box of isolated slot_module(): 9.5 × 3.5mm ✓

  Level 2 — Slicer layer view: rectangular hole visible at Z=5–8.5mm ✓

  Level 3 (AI): "rectangular cutout, approximately 10mm wide" (too vague for precision)

  Level 4 — Calipers after first test print: 9.4mm wide (0.1mm under — acceptable) ✓


Feature: Rubber foot posts (4×)

  Designed: 3.0mm diameter × 2.0mm tall

  Level 1 — Bounding box of isolated foot_post(): 3.0 × 3.0 × 2.0mm ✓

  Level 2 — Slicer: four small protrusions visible on bottom face ✓

  Level 3 (AI): "four small round protrusions" ✓

  Level 4 — Not yet printed


\end{lstlisting}

This log takes a few minutes to fill in, and it demonstrates exactly the kind of professional thinking that certification assessments reward.

\begin{center}\rule{0.5\linewidth}{0.5pt}\end{center}

\subsection*{Step 4 --- AI-Generated OpenSCAD Code: How to Verify It}\label{docs__pandoc__latex__src__3dmake_all_lessons.md__step-4--ai-generated-openscad-code-how-to-verify-it}

\subsubsection*{AI Can Write Code, but You Must Check It}\label{docs__pandoc__latex__src__3dmake_all_lessons.md__ai-can-write-code-but-you-must-check-it}

AI tools like Claude or ChatGPT can write OpenSCAD code from a plain-English description. This is a genuinely useful way to get a first draft of a design quickly. But \textbf{AI-generated OpenSCAD code is frequently wrong in subtle ways.} The code will look plausible, it may even build without errors, and it can still produce geometry that doesn\textquotesingle t match your intent or that will fail in the slicer.

The most common errors AI makes in OpenSCAD:

\textbf{Forgetting the 0.001 co-planar rule.} The AI writes a \texttt{difference()} with touching surfaces and omits the critical offset. The code looks correct, but the cut may produce rendering artifacts or non-manifold geometry.

\textbf{Python-style variable reassignment.} The AI writes \texttt{x\ =\ 5;\ x\ =\ x\ +\ 1;} as if variables work like Python, which doesn\textquotesingle t work in OpenSCAD (remember from Lesson 3: last assignment wins at parse time).

\textbf{Wrong transform order.} The AI reverses the order of \texttt{translate()} and \texttt{rotate()}, producing geometry in the wrong location. In OpenSCAD, transforms are applied in reverse order --- \texttt{rotate()} then \texttt{translate()} is not the same as \texttt{translate()} then \texttt{rotate()}.

\textbf{Non-manifold geometry.} The code looks correct but produces a solid that isn\textquotesingle t a properly sealed, watertight mesh. Slicers will reject it or produce incorrect layers.

\textbf{Wrong clearance values.} The AI invents plausible-sounding tolerances (like "0.1 mm clearance") that don\textquotesingle t match your specific printer. Tolerances must be calibrated --- they can\textquotesingle t be looked up in a general reference.

\subsubsection*{The AI Code Verification Checklist}\label{docs__pandoc__latex__src__3dmake_all_lessons.md__the-ai-code-verification-checklist}

Apply these checks to every piece of AI-generated OpenSCAD code before treating it as correct:

\begin{lstlisting}[style=Alabaster]
AI CODE VERIFICATION CHECKLIST

Run each check in order. Do not proceed if a check fails.


[ ] 1. SYNTAX

      3dm build → does it compile without errors?


[ ] 2. DIMENSIONS

      3dm info → does the bounding box match what you asked for?


[ ] 3. BOOLEAN OPERATIONS

      Does every difference() use the 0.001 offset rule?

      Look for: translate([*, *, 0]) inside difference()

      Should be: translate([*, *, -0.001])


[ ] 4. PARAMETER BEHAVIOR

      Change each parameter to an extreme value (minimum, maximum, zero).

      Does the model change correctly? Are there any edge cases that break it?


[ ] 5. MANIFOLD CHECK

      Open build/main.stl in your slicer.

      Are there any geometry warnings or errors?


[ ] 6. PRINT FEASIBILITY

      Are all walls at least 1.2mm thick for your nozzle size?

      Are there any overhangs steeper than 45°?

      Are all features large enough to be printable?


\end{lstlisting}

\begin{center}\rule{0.5\linewidth}{0.5pt}\end{center}

\subsection*{Step 5 --- Prompt Engineering for Better AI Results}\label{docs__pandoc__latex__src__3dmake_all_lessons.md__step-5--prompt-engineering-for-better-ai-results}

\subsubsection*{Weak Prompts Produce Weak Results}\label{docs__pandoc__latex__src__3dmake_all_lessons.md__weak-prompts-produce-weak-results}

The quality of AI output depends heavily on the quality of your input prompt. A vague request produces a vague or incorrect response. A specific, structured request produces specific, actionable feedback.

\textbf{Weak prompt:}

\begin{lstlisting}[style=Alabaster]
Review my OpenSCAD code.


\end{lstlisting}

This gives the AI almost no information about what to look for. The response will be generic and may miss exactly the issues you care about.

\textbf{Strong prompt:}

\begin{lstlisting}[style=Alabaster]
I have an OpenSCAD model of a parametric mounting bracket with two M3 bolt holes.

Parameters: width=40, height=30, depth=8, hole_r=3.2.


Please review this code for:

1. Correct use of the 0.001 offset rule in all difference() operations

2. Manifold-safe geometry (no co-planar faces)

3. Whether hole_r=3.2mm is the correct M3 clearance hole diameter per ISO standards

4. Any parameter combinations that could produce invalid geometry


For each issue found, provide:

- Severity: critical / warning / note

- Line number (if applicable)

- Explanation of the problem

- Suggested fix


\end{lstlisting}

This prompt tells the AI exactly what to look for, what standard to check against, and how to format its response. The output is dramatically more useful.

\subsubsection*{Four Characteristics of a Strong Code Review Prompt}\label{docs__pandoc__latex__src__3dmake_all_lessons.md__four-characteristics-of-a-strong-code-review-prompt}

\textbf{Specificity} --- name the design, state the parameters, describe the intent. Don\textquotesingle t make the AI guess what you\textquotesingle re trying to build.

\textbf{Explicit criteria} --- list the specific things you want checked. If you\textquotesingle re worried about the 0.001 rule, say so. If you want to verify dimensional correctness, specify the intended dimensions.

\textbf{Standard references} --- mention the standards or rules you want applied ("ISO M3 clearance hole", "0.001 co-planar offset rule", "minimum 1.2mm wall for 0.4mm nozzle").

\textbf{Output format request} --- ask for a structured response (severity, line, explanation, fix) rather than a prose paragraph.

\begin{center}\rule{0.5\linewidth}{0.5pt}\end{center}

\subsection*{\texorpdfstring{Step 6 --- Design Intent Guards with \texttt{assert()}}{Step 6 --- Design Intent Guards with assert()}}\label{docs__pandoc__latex__src__3dmake_all_lessons.md__step-6--design-intent-guards-with-assert}

\subsubsection*{Making Constraints Automatic}\label{docs__pandoc__latex__src__3dmake_all_lessons.md__making-constraints-automatic}

In Lesson 3, you learned how to use \texttt{assert()} inside modules to validate parameters. You can also use \texttt{assert()} at the top level of your file to enforce the same constraints you\textquotesingle d put in your verification log. This turns manual checks into automatic ones.

\begin{lstlisting}[style=Alabaster, language=openscad]
// Top-level design constraints

// These run before any geometry is computed.

width  = 30;

height = 10;

wall_t = 2;


assert(width  >= 10, str("width must be >= 10mm for structural integrity, got: ", width));

assert(height >= 5,  str("height must be >= 5mm to fit standard hardware, got: ", height));

assert(wall_t >= 1.5, str("wall_t must be >= 1.5mm to print reliably with 0.4mm nozzle, got: ", wall_t));

assert(wall_t < width / 2,

  str("wall_t (", wall_t, ") must be less than half of width (", width/2, ")"));


// Geometry — guaranteed to be within valid bounds by the asserts above

difference() {

  cube([width, width, height]);

  translate([wall_t, wall_t, -0.001])

    cube([width - 2*wall_t, width - 2*wall_t, height + 0.002]);

}


\end{lstlisting}

Every constraint you can express as an \texttt{assert()} is a constraint that no longer requires manual verification. If someone changes \texttt{wall\_t\ =\ 0.5} (too thin to print), the build immediately stops with a clear message instead of silently producing a model that fails on the printer.

\begin{center}\rule{0.5\linewidth}{0.5pt}\end{center}

\subsection*{Concept 2: Privacy and AI Tools}\label{docs__pandoc__latex__src__3dmake_all_lessons.md__concept-2-privacy-and-ai-tools}

\subsubsection*{What Happens to Files You Upload}\label{docs__pandoc__latex__src__3dmake_all_lessons.md__what-happens-to-files-you-upload}

When you use cloud-based AI tools --- including the \texttt{3dm\ info} AI description feature when it connects to a cloud model, or an external AI chat like Claude or ChatGPT --- your files and prompts are processed on remote servers. Depending on the service and its privacy policy, that data may be:

\begin{itemize}
\tightlist
\item
  Used to improve the AI model (potentially meaning others could see examples from your session)
\item
  Stored for a period of time
\item
  Subject to the provider\textquotesingle s legal obligations in their jurisdiction
\end{itemize}

For classroom learning projects, this is generally a non-issue --- there\textquotesingle s nothing sensitive about learning how to make a parametric box. But if you\textquotesingle re ever working on a design that represents real intellectual property (a product idea, a client\textquotesingle s design, proprietary measurements), you need to understand what you\textquotesingle re sending and to whom.

\textbf{Best practices for privacy-sensitive work:}

\begin{itemize}
\tightlist
\item
  Check the AI service\textquotesingle s data retention policy before uploading
\item
  Simplify or abstract the design before sending for review (remove specific dimensions, replace proprietary shapes with generic approximations)
\item
  Use local/offline tools for final verification of sensitive designs --- the bounding box and slicer checks don\textquotesingle t require network access
\end{itemize}

For the purposes of this certification course, all your projects are learning work and standard cloud AI tools are appropriate. But developing the habit of thinking about this now is professional practice.

\begin{center}\rule{0.5\linewidth}{0.5pt}\end{center}

\subsection*{Validation Type Summary}\label{docs__pandoc__latex__src__3dmake_all_lessons.md__validation-type-summary}

{\def\LTcaptype{none} % do not increment counter
\begin{longtable}[]{@{}
  >{\raggedright\arraybackslash}p{(\linewidth - 4\tabcolsep) * \real{0.2832}}
  >{\raggedright\arraybackslash}p{(\linewidth - 4\tabcolsep) * \real{0.3894}}
  >{\raggedright\arraybackslash}p{(\linewidth - 4\tabcolsep) * \real{0.3274}}@{}}
\toprule\noalign{}
\begin{minipage}[b]{\linewidth}\raggedright
Validation Type
\end{minipage} & \begin{minipage}[b]{\linewidth}\raggedright
What It Checks
\end{minipage} & \begin{minipage}[b]{\linewidth}\raggedright
Reliability
\end{minipage} \\
\midrule\noalign{}
\endhead
\bottomrule\noalign{}
\endlastfoot
\texttt{3dm\ info} bounding box \& volume & Exact model dimensions &
Deterministic --- identical every run \\
\texttt{3dm\ info} AI description & Overall shape, component layout &
Non-deterministic --- may vary \\
Slicer layer view & Wall thickness, internal features, overhangs &
Deterministic --- reliable for geometry \\
Physical calipers & Printed part dimensions &
Ground truth --- most authoritative \\
\texttt{assert()} guards & Parameter validity before render &
Deterministic --- runs at build time \\
\end{longtable}
}

\begin{center}\rule{0.5\linewidth}{0.5pt}\end{center}

\subsection*{Quiz --- Lesson 4 (15 questions)}\label{docs__pandoc__latex__src__3dmake_all_lessons.md__quiz--lesson-4-15-questions}

\begin{enumerate}
\tightlist
\item
  What does \texttt{3dm\ info} do? What two fundamentally different types of information does it provide?
\item
  What is the key difference between deterministic validation and AI-based validation?
\item
  Give two specific examples of information that \texttt{3dm\ info}\textquotesingle s AI description cannot reliably provide.
\item
  In the verification escalation chain, which level provides "ground truth" for dimensional accuracy, and why?
\item
  List three common errors that AI tools make when generating OpenSCAD code.
\item
  True or False: the bounding box reported by \texttt{3dm\ info} might be slightly different each time you run it on the same file.
\item
  What data privacy concern should you consider before uploading a proprietary design to a cloud AI service?
\item
  What does the slicer layer view let you verify that the bounding box cannot show?
\item
  What two characteristics distinguish a strong AI code review prompt from a weak one?
\item
  If an AI describes your model as "a small cylindrical protrusion" and you designed it as 3.2 mm in diameter, why is this description insufficient for print verification?
\item
  What is "prompt engineering" and why does it matter when asking an AI to review OpenSCAD code?
\item
  Write an \texttt{assert()} statement that ensures a variable called \texttt{layer\_height} is between 0.05 mm and 0.4 mm.
\item
  What is the advantage of using \texttt{assert()} at the top of a file compared to manually checking parameter values before every build?
\item
  Describe one scenario where AI validation would be genuinely useful and one where you must rely on deterministic tools instead.
\item
  A classmate says "the AI said my model looked fine, so I sent it to print." What is wrong with this approach, and what should they have done instead?
\end{enumerate}

\begin{center}\rule{0.5\linewidth}{0.5pt}\end{center}

\subsection*{Extension Problems (15)}\label{docs__pandoc__latex__src__3dmake_all_lessons.md__extension-problems-15-3}

\begin{enumerate}
\tightlist
\item
  Design a five-item verification checklist for a new 3D printed part. Include items that are deterministic and items that benefit from AI assistance. Label each.
\item
  Build a "prompt library" for three common design review scenarios (mounting bracket, enclosure with lid, text embossing). Test each prompt against an AI tool and document the output quality.
\item
  Run \texttt{3dm\ info} on the same STL file five times. Document any variation in the AI description output. What does this variation tell you about using AI for precision work?
\item
  Create a structured accept/reject log for one full design cycle: design → verification checklist → AI review → physical test → decision to print or revise.
\item
  Research what "hallucination" means in the context of AI language models. Describe two specific ways an AI could hallucinate during a design review of OpenSCAD code.
\item
  Create a set of deliberately flawed OpenSCAD files --- one with a missing 0.001 offset, one with a Python-style variable reassignment, one with a thin wall below 1.2 mm. Test whether an AI code review catches each flaw.
\item
  Write a one-page data privacy policy for a classroom 3D printing lab that uses cloud AI tools. What should students be required to understand before uploading files?
\item
  Build a "verification script" that automatically runs \texttt{3dm\ build}, checks the exit code, runs \texttt{3dm\ info}, and extracts the bounding box for comparison against expected values.
\item
  Research the difference between "AI-assisted" and "AI-automated" workflows. At what point does removing a human from the verification loop become risky in a manufacturing context?
\item
  Design an experiment: ask three different AI tools to review the same OpenSCAD file. Document the similarities and differences in their feedback. What does this tell you about AI reliability?
\item
  Build a comparison table: for five design features (bolt holes, thin walls, snap-fit gaps, embossed text, threaded inserts), identify whether AI review, slicer validation, or physical calipers is most reliable.
\item
  Research and document what "manifold" means in geometry. What makes a mesh non-manifold? What are the three most common causes in OpenSCAD?
\item
  Write a one-page explanation of the verification escalation chain written for a student in the 8th grade who has never heard of 3D printing. Use no technical jargon without defining it.
\item
  Create an \texttt{assert()} suite for a parametric box design: validate width, depth, height, wall thickness, and the relationship between wall thickness and overall size. Test it with at least three invalid parameter sets.
\item
  Write a short guide titled "When to Trust AI and When Not To" specifically for 3D design workflows. Include at least five specific decision rules.
\end{enumerate}

\begin{center}\rule{0.5\linewidth}{0.5pt}\end{center}

\subsection*{References and Helpful Resources}\label{docs__pandoc__latex__src__3dmake_all_lessons.md__references-and-helpful-resources-3}

\begin{itemize}
\item
  OpenSCAD User Manual --- Manifold and Geometry Validation --- https://en.wikibooks.org/wiki/OpenSCAD\_User\_Manual/The\_OpenSCAD\_Language
\item
  Anthropic Prompt Engineering Guide --- https://docs.anthropic.com/en/docs/build-with-claude/prompt-engineering/overview
\item
  Data Privacy --- Review your AI provider\textquotesingle s terms of service and data retention policy before submitting proprietary design files.
\item
  NIST AI Risk Management Framework (AI RMF 1.0) --- https://www.nist.gov/system/files/documents/2023/01/26/AI\%20RMF\%201.0.pdf
\end{itemize}

\subsubsection*{Supplemental Resources}\label{docs__pandoc__latex__src__3dmake_all_lessons.md__supplemental-resources-3}

\begin{itemize}
\tightlist
\item
  3DMake GitHub Repository --- https://github.com/tdeck/3dmake
\item
  Anthropic Prompt Engineering Overview --- https://docs.anthropic.com/en/docs/build-with-claude/prompt-engineering/overview
\item
  NIST AI RMF --- https://airc.nist.gov/Home
\item
  PrusaSlicer Documentation --- https://docs.prusa3d.com/en/
\item
  All other reference documents are embedded as appendices below --- no separate files needed.
\end{itemize}

\begin{center}\rule{0.5\linewidth}{0.5pt}\end{center}

\section*{Part 2: Extension Project and Reference Appendices}\label{docs__pandoc__latex__src__3dmake_all_lessons.md__part-2-extension-project-and-reference-appendices-1}

\emph{All supporting documents for Lesson 4 are included below. You do not need to open any additional files.}

\begin{center}\rule{0.5\linewidth}{0.5pt}\end{center}

\subsection*{Appendix A: Dice, Dice, Dice --- Extension Project}\label{docs__pandoc__latex__src__3dmake_all_lessons.md__appendix-a-dice-dice-dice--extension-project}

\textbf{Estimated time:} 3--6 hours

\subsubsection*{Overview}\label{docs__pandoc__latex__src__3dmake_all_lessons.md__overview-4}

This extension project connects the verification and AI-assisted feedback skills from Lesson 4 to a physical object where design precision is directly measurable: a six-sided die. A fair die is one of the most demanding manufactured objects you can design --- it requires uniform geometry, controlled mass distribution, and a randomness that can be tested statistically. The design is straightforward. The verification is where the real learning happens.

This project also inverts the lesson\textquotesingle s AI-skepticism lesson in a productive way: it uses AI as part of the design ideation process, then requires \emph{you} to verify every claim with physical testing and statistical analysis. A die that the AI describes as "balanced and symmetric" should be tested with 100+ real rolls. If the data disagrees with the description, the data wins.

\subsubsection*{Learning Objectives}\label{docs__pandoc__latex__src__3dmake_all_lessons.md__learning-objectives-7}

By completing this project, you will:

\begin{itemize}
\tightlist
\item
  Design parametric dice in OpenSCAD with controlled geometry and mass distribution
\item
  Apply the verification chain from Lesson 4 to a physical fabrication outcome
\item
  Test durability and randomness for small, thrown objects
\item
  Conduct and document statistical fairness testing
\item
  Understand how infill, wall thickness, and pip geometry affect roll behavior
\end{itemize}

\subsubsection*{Materials}\label{docs__pandoc__latex__src__3dmake_all_lessons.md__materials-3}

\begin{itemize}
\tightlist
\item
  Computer with OpenSCAD and slicer
\item
  Printer and filament (PLA recommended --- easy to print and dimensionally consistent)
\item
  A small throwing surface (book, box lid, or low-sided tray)
\item
  Calipers if available (for dimensional verification)
\end{itemize}

\subsubsection*{Tasks}\label{docs__pandoc__latex__src__3dmake_all_lessons.md__tasks-3}

\begin{enumerate}
\item
  \textbf{Design three distinct die variants} in OpenSCAD using parameters (\texttt{size}, \texttt{wall}, \texttt{infill\_proxy}, \texttt{pip\_depth}, \texttt{pip\_r}). Each variant must differ in at least one meaningful way --- for example: different pip geometry (sunken vs. flat), different wall thickness, or different overall size. All three must be generated from the same parametric module by changing variables at the top of the file.
\item
  \textbf{Document your design decisions.} For each variant, write one short paragraph explaining which parameters you changed and why you expected that change to affect fairness or durability.
\item
  \textbf{Print one sample of each variant.} Use the same material and layer height for all three to isolate the design variable.
\item
  \textbf{Durability test:} five throws onto a soft surface for each die. Record whether any delamination, corner chipping, or pip damage is visible after testing.
\item
  \textbf{Fairness test:} roll each die at least 100 times on a consistent surface. Record every result in the roll data table (Appendix B). Compute the frequency of each face and compare to the expected 16.67\%.
\item
  \textbf{Statistical analysis:} calculate the deviation from expected for each face. A die within ±3\% per face on 100 rolls is performing within normal random variation. A face appearing more than 20\% of the time suggests a design or print bias that warrants investigation.
\end{enumerate}

\subsubsection*{Probing Questions}\label{docs__pandoc__latex__src__3dmake_all_lessons.md__probing-questions-3}

\begin{itemize}
\tightlist
\item
  How does internal infill or a hollow cavity affect mass distribution and roll randomness?
\item
  Which parameter change had the largest effect on fairness or durability, and why do you think that is?
\item
  If one face came up significantly more often than expected, what is the most likely physical cause?
\end{itemize}

\begin{center}\rule{0.5\linewidth}{0.5pt}\end{center}

\subsubsection*{Starter Code}\label{docs__pandoc__latex__src__3dmake_all_lessons.md__starter-code-2}

\begin{lstlisting}[style=Alabaster, language=openscad]
// Parametric Six-Sided Die

// All critical dimensions are top-level parameters.

// Change these to create your three variants.


// ---- Parameters ----

die_size   = 16;     // mm — overall cube side length

wall       = 2.0;    // mm — wall thickness (increase for more durability)

pip_r      = 1.0;    // mm — radius of pip hemispheres

pip_depth  = 0.6;    // mm — how deep pips are pressed into the face

corner_r   = 1.2;    // mm — chamfer/rounding radius on die edges

$fn        = 32;     // curve resolution


// Verify printability

assert(wall >= 1.2,

  str("wall (", wall, "mm) is below minimum printable thickness (1.2mm)"));

assert(pip_r > 0 && pip_r < die_size / 6,

  str("pip_r (", pip_r, "mm) is unreasonable for this die size"));

assert(pip_depth < wall,

  str("pip_depth (", pip_depth, "mm) must be less than wall (", wall, "mm)"));


// ---- Pip Positions (for each face, relative to face center) ----

// Each face is a list of [x, y] offsets from center, in pip-spacing units.

// spacing = die_size / 5 gives a comfortable pip layout.

pip_spacing = die_size / 5;


face_1_pips = [[0, 0]];

face_2_pips = [[-1, -1], [1, 1]];

face_3_pips = [[-1, -1], [0, 0], [1, 1]];

face_4_pips = [[-1, -1], [1, -1], [-1, 1], [1, 1]];

face_5_pips = [[-1, -1], [1, -1], [0, 0], [-1, 1], [1, 1]];

face_6_pips = [[-1, -1], [1, -1], [-1, 0], [1, 0], [-1, 1], [1, 1]];


// ---- Single Pip (hemisphere subtracted from a face) ----

module pip(depth=pip_depth, r=pip_r) {

  sphere(r=r, $fn=16);

}


// ---- Pips on a face (2D positions, extruded onto face surface) ----

// This module places pips at the given positions on the +Z face of the die.

module place_pips(positions, face_z, r=pip_r, depth=pip_depth, spacing=pip_spacing) {

  for (pos = positions) {

    translate([die_size/2 + pos[0]*spacing,

               die_size/2 + pos[1]*spacing,

               face_z + r - depth])

      pip(depth, r);

  }

}


// ---- Die Body ----

module die_body() {

  // Simple cube — replace with minkowski sphere for rounded edges

  cube([die_size, die_size, die_size]);

}


// ---- Full Die with Pips Subtracted ----

module die() {

  difference() {

    die_body();


    // Face +Z (top): 1 pip

    place_pips(face_1_pips, die_size);

    // Face -Z (bottom): 6 pips (opposite of 1)

    mirror([0, 0, 1])

      place_pips(face_6_pips, 0);

    // Face +X (right): 2 pips (opposite of 5)

    rotate([0, -90, 0]) translate([0, 0, -die_size])

      place_pips(face_2_pips, die_size);

    // Face -X (left): 5 pips

    rotate([0, 90, 0])

      place_pips(face_5_pips, die_size);

    // Face +Y (front): 3 pips (opposite of 4)

    rotate([90, 0, 0]) translate([0, 0, -die_size])

      place_pips(face_3_pips, die_size);

    // Face -Y (back): 4 pips

    rotate([-90, 0, 0])

      place_pips(face_4_pips, die_size);

  }

}


die();


\end{lstlisting}

\textbf{Notes on using this scaffold:}

\begin{itemize}
\tightlist
\item
  To create Variant 2, change \texttt{wall} and \texttt{pip\_depth}. To create Variant 3, change \texttt{die\_size}. Keep all other parameters the same between variants so you isolate the effect.
\item
  The \texttt{assert()} statements will catch unreasonable parameter combinations before rendering.
\item
  For rounded-corner dice (more realistic feel, less edge chipping), replace \texttt{die\_body()} with a \texttt{minkowski()} of a smaller cube and a \texttt{sphere(r=corner\_r)} --- see the AI code gallery in Appendix D for a worked example of this technique.
\end{itemize}

\subsubsection*{Quiz Questions (from workshop)}\label{docs__pandoc__latex__src__3dmake_all_lessons.md__quiz-questions-from-workshop}

\begin{enumerate}
\tightlist
\item
  What parameter controls wall thickness in your \texttt{.scad} file?
\item
  Name one method to increase a die\textquotesingle s impact resistance without changing its overall size.
\item
  How would you test a die for randomness? What is the minimum number of rolls for a meaningful result?
\item
  Why might a hollow die behave differently from a solid one of the same external dimensions?
\item
  What is a simple visual sign of print-layer delamination on a die?
\item
  True or False: A perfectly fair die must have uniform weight distribution throughout, including the internal structure.
\item
  Describe two design approaches for rendering the face numbers: one using \texttt{difference()} to create sunken pips, one using \texttt{linear\_extrude()} for raised numbers.
\item
  Your die lands on face 1 significantly more often than any other face in 100 rolls. Name two possible physical causes and describe how you would diagnose each.
\item
  If you increase infill from 15\% to 50\%, what is the primary effect? (A) Dramatically longer print time, same weight (B) Significantly longer print time and higher weight (C) Same weight, different surface finish. Answer: B
\item
  Explain why a fair die is an excellent test of manufacturing precision. How does this project connect design decisions to real-world measurable outcomes?
\end{enumerate}

\begin{center}\rule{0.5\linewidth}{0.5pt}\end{center}

\subsection*{Appendix B: Dice, Dice, Dice --- Student Documentation Template}\label{docs__pandoc__latex__src__3dmake_all_lessons.md__appendix-b-dice-dice-dice--student-documentation-template}

\textbf{Author:} \_\_\_\_\_\_\_\_\_\_\_\_\_\_\_\_
\textbf{Date:} \_\_\_\_\_\_\_\_\_\_
\textbf{Description:} Design parametric 3D-printed dice and test for fairness through physical testing and statistical analysis.

\begin{center}\rule{0.5\linewidth}{0.5pt}\end{center}

\subsubsection*{Design Concept}\label{docs__pandoc__latex__src__3dmake_all_lessons.md__design-concept}

\begin{itemize}
\tightlist
\item
  Overall die size and material:
\item
  How will you address uniform weight distribution in your design?
\item
  What design details will differentiate your three variants?
\end{itemize}

\begin{center}\rule{0.5\linewidth}{0.5pt}\end{center}

\subsubsection*{Design Specifications}\label{docs__pandoc__latex__src__3dmake_all_lessons.md__design-specifications}

{\def\LTcaptype{none} % do not increment counter
\begin{longtable}[]{@{}lccc@{}}
\toprule\noalign{}
Parameter & Variant 1 & Variant 2 & Variant 3 \\
\midrule\noalign{}
\endhead
\bottomrule\noalign{}
\endlastfoot
\texttt{die\_size} (mm) & & & \\
\texttt{wall} (mm) & & & \\
\texttt{pip\_r} (mm) & & & \\
\texttt{pip\_depth} (mm) & & & \\
\texttt{corner\_r} (mm) & & & \\
Infill \% (slicer) & & & \\
Estimated mass (g) & & & \\
Actual mass (g) & & & \\
\end{longtable}
}

\textbf{Variant 1 design rationale} (one paragraph --- what did you change and what outcome did you expect?):

\textbf{Variant 2 design rationale:}

\textbf{Variant 3 design rationale:}

\begin{center}\rule{0.5\linewidth}{0.5pt}\end{center}

\subsubsection*{Dimensional Verification}\label{docs__pandoc__latex__src__3dmake_all_lessons.md__dimensional-verification}

Run \texttt{3dm\ build} then \texttt{3dm\ info} on each variant. Record the deterministic output:

{\def\LTcaptype{none} % do not increment counter
\begin{longtable}[]{@{}lccc@{}}
\toprule\noalign{}
Measurement & Variant 1 & Variant 2 & Variant 3 \\
\midrule\noalign{}
\endhead
\bottomrule\noalign{}
\endlastfoot
Bounding box X (mm) & & & \\
Bounding box Y (mm) & & & \\
Bounding box Z (mm) & & & \\
Volume (mm³) & & & \\
Expected die\_size (mm) & & & \\
Deviation from expected (mm) & & & \\
\end{longtable}
}

If any dimension deviates by more than 0.5 mm from \texttt{die\_size}, investigate before printing.

\begin{center}\rule{0.5\linewidth}{0.5pt}\end{center}

\subsubsection*{Durability Test}\label{docs__pandoc__latex__src__3dmake_all_lessons.md__durability-test}

Five throws per variant onto a soft surface (book cover, box lid, fabric).

{\def\LTcaptype{none} % do not increment counter
\begin{longtable}[]{@{}
  >{\centering\arraybackslash}p{(\linewidth - 6\tabcolsep) * \real{0.0649}}
  >{\raggedright\arraybackslash}p{(\linewidth - 6\tabcolsep) * \real{0.3117}}
  >{\raggedright\arraybackslash}p{(\linewidth - 6\tabcolsep) * \real{0.3117}}
  >{\raggedright\arraybackslash}p{(\linewidth - 6\tabcolsep) * \real{0.3117}}@{}}
\toprule\noalign{}
\begin{minipage}[b]{\linewidth}\centering
Throw
\end{minipage} & \begin{minipage}[b]{\linewidth}\raggedright
Variant 1 --- Observations
\end{minipage} & \begin{minipage}[b]{\linewidth}\raggedright
Variant 2 --- Observations
\end{minipage} & \begin{minipage}[b]{\linewidth}\raggedright
Variant 3 --- Observations
\end{minipage} \\
\midrule\noalign{}
\endhead
\bottomrule\noalign{}
\endlastfoot
1 & & & \\
2 & & & \\
3 & & & \\
4 & & & \\
5 & & & \\
\end{longtable}
}

Post-durability inspection:

\begin{itemize}
\tightlist
\item
  Any layer delamination visible? (Variant 1 / 2 / 3):
\item
  Any corner chipping? (Variant 1 / 2 / 3):
\item
  Any pip damage? (Variant 1 / 2 / 3):
\end{itemize}

\begin{center}\rule{0.5\linewidth}{0.5pt}\end{center}

\subsubsection*{Fairness Testing --- Roll Results}\label{docs__pandoc__latex__src__3dmake_all_lessons.md__fairness-testing--roll-results}

Complete this table for your best-performing variant (or all three if time permits). Minimum 100 rolls per die tested.

\textbf{Variant tested:} \_\_\_\_ \textbf{Total rolls:} \_\_\_\_ \textbf{Date:} \_\_\_\_ \textbf{Surface used:} \_\_\_\_

{\def\LTcaptype{none} % do not increment counter
\begin{longtable}[]{@{}
  >{\centering\arraybackslash}p{(\linewidth - 10\tabcolsep) * \real{0.1184}}
  >{\raggedright\arraybackslash}p{(\linewidth - 10\tabcolsep) * \real{0.2895}}
  >{\centering\arraybackslash}p{(\linewidth - 10\tabcolsep) * \real{0.1184}}
  >{\centering\arraybackslash}p{(\linewidth - 10\tabcolsep) * \real{0.1316}}
  >{\centering\arraybackslash}p{(\linewidth - 10\tabcolsep) * \real{0.2237}}
  >{\centering\arraybackslash}p{(\linewidth - 10\tabcolsep) * \real{0.1184}}@{}}
\toprule\noalign{}
\begin{minipage}[b]{\linewidth}\centering
Face
\end{minipage} & \begin{minipage}[b]{\linewidth}\raggedright
Tally (mark each roll)
\end{minipage} & \begin{minipage}[b]{\linewidth}\centering
Frequency
\end{minipage} & \begin{minipage}[b]{\linewidth}\centering
Percentage
\end{minipage} & \begin{minipage}[b]{\linewidth}\centering
Expected (16.67\%)
\end{minipage} & \begin{minipage}[b]{\linewidth}\centering
Deviation
\end{minipage} \\
\midrule\noalign{}
\endhead
\bottomrule\noalign{}
\endlastfoot
1 & & & \% & 16.67\% & \\
2 & & & \% & 16.67\% & \\
3 & & & \% & 16.67\% & \\
4 & & & \% & 16.67\% & \\
5 & & & \% & 16.67\% & \\
6 & & & \% & 16.67\% & \\
\textbf{Total} & & \textbf{100} & \textbf{100\%} & \textbf{100\%} & \\
\end{longtable}
}

\begin{center}\rule{0.5\linewidth}{0.5pt}\end{center}

\subsubsection*{Statistical Analysis}\label{docs__pandoc__latex__src__3dmake_all_lessons.md__statistical-analysis}

\textbf{Maximum deviation from expected} (the highest "Deviation" value in your table): \_\_\_\_ \%

\textbf{Interpretation:}

\begin{itemize}
\tightlist
\item
  Deviation ≤ 3\% on all faces: Normal random variation --- die performs within acceptable range
\item
  Deviation 3--5\% on one face: Mild bias --- investigate print quality and pip geometry
\item
  Deviation \textgreater{} 5\% on any face: Significant bias --- design or print issue warrants redesign
\end{itemize}

\textbf{Your conclusion} (is this die fair within statistical tolerance?):

\textbf{Most frequent face:} \_\_\_\_ \textbf{Least frequent face:} \_\_\_\_

If the die shows bias, what physical cause do you suspect? (Pip depth on the heavy face? Bed adhesion producing a flatter bottom? Asymmetric corner rounding?):

\begin{center}\rule{0.5\linewidth}{0.5pt}\end{center}

\subsubsection*{Reflections}\label{docs__pandoc__latex__src__3dmake_all_lessons.md__reflections-3}

\begin{enumerate}
\item
  Were you surprised by the fairness test results? Did the die perform better or worse than you expected?
\item
  Which design parameter had the largest effect on the test outcome, and what does that tell you about the physics of the die?
\item
  What would you change in a fourth iteration to improve fairness?
\item
  What did this project teach you about the gap between a design that looks correct on screen and one that performs correctly in the physical world?
\end{enumerate}

\begin{center}\rule{0.5\linewidth}{0.5pt}\end{center}

\subsubsection*{Attachments Checklist}\label{docs__pandoc__latex__src__3dmake_all_lessons.md__attachments-checklist-3}

\begin{itemize}
\tightlist
\item[$\square$]
  \texttt{.scad} file with parametric die module (all three variants producible from one file)
\item[$\square$]
  Photo of all three printed die variants
\item[$\square$]
  Raw roll data tally sheet (the actual marks for each roll)
\item[$\square$]
  Completed frequency and statistical analysis table
\item[$\square$]
  Dimensional verification output from \texttt{3dm\ info} for each variant
\end{itemize}

\begin{center}\rule{0.5\linewidth}{0.5pt}\end{center}

\subsubsection*{Teacher Feedback}\label{docs__pandoc__latex__src__3dmake_all_lessons.md__teacher-feedback-3}

{\def\LTcaptype{none} % do not increment counter
\begin{longtable}[]{@{}
  >{\raggedright\arraybackslash}p{(\linewidth - 4\tabcolsep) * \real{0.8678}}
  >{\centering\arraybackslash}p{(\linewidth - 4\tabcolsep) * \real{0.0909}}
  >{\raggedright\arraybackslash}p{(\linewidth - 4\tabcolsep) * \real{0.0413}}@{}}
\toprule\noalign{}
\begin{minipage}[b]{\linewidth}\raggedright
Category
\end{minipage} & \begin{minipage}[b]{\linewidth}\centering
Score (0--3)
\end{minipage} & \begin{minipage}[b]{\linewidth}\raggedright
Notes
\end{minipage} \\
\midrule\noalign{}
\endhead
\bottomrule\noalign{}
\endlastfoot
Problem \& Solution --- dice print successfully; fairness testing conducted rigorously
& & \\
Design \& Code Quality --- parametric module produces all three variants; \texttt{assert()} used; code is commented
& & \\
Documentation --- roll data complete; statistical analysis accurate; reflections show understanding
& & \\
\textbf{Total (0--9)} & & \\
\end{longtable}
}

\begin{center}\rule{0.5\linewidth}{0.5pt}\end{center}

\subsection*{Appendix C: Dice, Dice, Dice --- Instructor Reference}\label{docs__pandoc__latex__src__3dmake_all_lessons.md__appendix-c-dice-dice-dice--instructor-reference}

\subsubsection*{Project Context}\label{docs__pandoc__latex__src__3dmake_all_lessons.md__project-context-2}

This extension project reinforces the lesson\textquotesingle s central theme --- trusting data over descriptions --- by making the verification stakes concrete and measurable. The fairness of a printed die is something a student can test with their own hands. A die that the AI describes as "perfectly symmetric" but lands on face 1 thirty percent of the time is an immediate, visceral demonstration of why physical testing matters.

The chi-squared test is mentioned in the teacher rubric but is not required for baseline completion. Students who complete it should receive recognition for the additional rigor; students who compute simple percentage deviation are satisfying the core requirement.

\subsubsection*{Key Learning Points to Reinforce}\label{docs__pandoc__latex__src__3dmake_all_lessons.md__key-learning-points-to-reinforce-2}

\textbf{Verification is empirical, not descriptive.} The workflow mirrors Lesson 4\textquotesingle s verification chain directly: deterministic check (\texttt{3dm\ info} bounding box), then visual check (slicer layer view), then physical test (the roll data). A student who prints without running \texttt{3dm\ info} first has skipped the deterministic layer.

\textbf{Parametric design means one module, multiple variants.} Students who create three separate \texttt{.scad} files (one per variant) miss the core parametric design goal. The correct approach is a single file where changing variables at the top produces each variant. This is directly assessable in the code review.

\textbf{Statistical fairness has a physical explanation.} When a face appears too frequently, the most common causes are: (1) pip depth asymmetry --- faces with fewer/shallower pips have more mass on that face, causing it to face down more often; (2) bed adhesion flattening the bottom face slightly during printing; (3) infill asymmetry if the model origin isn\textquotesingle t centered. Students should be guided toward diagnosing the \emph{cause}, not just noting the result.

\subsubsection*{Constraints (Enforce These)}\label{docs__pandoc__latex__src__3dmake_all_lessons.md__constraints-enforce-these-2}

\begin{itemize}
\tightlist
\item
  All three variants must be produced from a single parametric \texttt{.scad} file
\item
  \texttt{assert()} statements must validate key parameters (wall thickness minimum, pip depth less than wall)
\item
  Minimum 100 rolls for fairness testing --- 50 rolls is not statistically meaningful
\item
  Raw roll data (the actual tally) must be submitted, not just the summary table
\end{itemize}

\subsubsection*{Assessment Notes}\label{docs__pandoc__latex__src__3dmake_all_lessons.md__assessment-notes-2}

\textbf{Strong submissions} show: roll data with a tally sheet (not just a filled-in table), a physical explanation for any bias observed, and a genuine iteration --- the student changed something between variants because of a hypothesis, not randomly.

\textbf{Common weak areas:} Single-variant testing because "it looked fine," statistical tables with numbers that don\textquotesingle t add up to 100\%, and designs where pip depth equals wall thickness (causing holes through the die). The \texttt{assert()} in the starter code catches the last one at render time.

\textbf{Extension for advanced students:} A formal chi-squared goodness-of-fit test. With 100 rolls and 6 faces, the critical value at p=0.05 significance is χ² = 11.07 (5 degrees of freedom). If the computed chi-squared exceeds this value, the die has a statistically significant bias. This is a meaningful cross-curricular mathematics application.

\begin{center}\rule{0.5\linewidth}{0.5pt}\end{center}

\subsection*{Appendix D: AI Guidance --- Advantages, Disadvantages, and Flawed Code Gallery}\label{docs__pandoc__latex__src__3dmake_all_lessons.md__appendix-d-ai-guidance--advantages-disadvantages-and-flawed-code-gallery}

\emph{This appendix expands on the lesson\textquotesingle s AI discussion and provides concrete debuggable examples of typical AI mistakes. Use the flawed examples as debugging exercises --- fix each one, then run \texttt{3dm\ build} → \texttt{3dm\ info} → slicer layer view to confirm the fix.}

\begin{center}\rule{0.5\linewidth}{0.5pt}\end{center}

\subsubsection*{D1 --- Advantages of AI in OpenSCAD Development}\label{docs__pandoc__latex__src__3dmake_all_lessons.md__d1--advantages-of-ai-in-openscad-development}

\textbf{Rapid ideation and scaffolding.} AI can convert a natural-language description into a first-draft parametric scaffold quickly. This is useful for getting a starting structure, but the generated code must always be verified --- never sent directly to print.

\textbf{Accessibility support.} Natural-language descriptions via \texttt{3dm\ info} make model verification more accessible to visually impaired makers who cannot rely on visual inspection of the 3D preview. AI-assisted verbal description of model geometry can substitute for visual inspection when implemented carefully.

\textbf{Promptable code review.} A well-structured prompt asking an AI to review OpenSCAD code for boolean pitfalls, coplanar faces, and thin walls can catch issues that a visual scan might miss --- as long as you verify any issues it raises rather than trusting its "all clear."

\textbf{Iteration speed.} AI can propose parameterizations and printability tips rapidly, letting you evaluate more design directions before committing to print. The key is that you verify deterministically; AI accelerates the ideation, not the verification.

\begin{center}\rule{0.5\linewidth}{0.5pt}\end{center}

\subsubsection*{D2 --- Disadvantages and Risks}\label{docs__pandoc__latex__src__3dmake_all_lessons.md__d2--disadvantages-and-risks}

\textbf{Non-determinism.} The same prompt given twice may produce different answers. AI is not a source of truth for design specifications or dimensional requirements.

\textbf{No geometric guarantees.} AI generates descriptions based on inference from rendered images or code text --- it does not compute solid geometry. It can describe a model as correct that has non-manifold geometry or zero-thickness walls that will fail in the slicer.

\textbf{Standards and tolerances gaps.} AI may suggest clearances, wall thicknesses, or print settings that contradict established practice for your specific printer and material. Always check AI suggestions against the lesson\textquotesingle s reference values.

\textbf{Hallucinations.} AI tools sometimes generate code using functions that do not exist in OpenSCAD, such as \texttt{rounded\_cube()}, \texttt{fillet3d()}, or \texttt{bevel\_edge()}. These functions will cause a compile error if you attempt to use them. The AI is not aware it has invented them.

\textbf{Privacy and IP exposure.} Cloud AI tools may retain uploaded files or use them for training data, depending on their terms of service. Before uploading a design file to any cloud AI service, check the provider\textquotesingle s data retention and privacy policy. In a classroom context, no personally identifiable or commercially sensitive designs should be uploaded without explicit instructor approval.

\begin{center}\rule{0.5\linewidth}{0.5pt}\end{center}

\subsubsection*{D3 --- Flawed AI-Generated OpenSCAD --- Debugging Gallery}\label{docs__pandoc__latex__src__3dmake_all_lessons.md__d3--flawed-ai-generated-openscad--debugging-gallery}

Use each example below as a debugging exercise. Identify the flaw, fix it, and verify the fix with \texttt{3dm\ build} followed by \texttt{3dm\ info}.

\begin{center}\rule{0.5\linewidth}{0.5pt}\end{center}

\textbf{Example 1: Hallucinated Convenience Function}

The AI claims OpenSCAD has a \texttt{rounded\_cube()} built-in. It does not.

\begin{lstlisting}[style=Alabaster, language=openscad]
// ❌ FLAWED — rounded_cube() does not exist in OpenSCAD

module enclosure() {

  rounded_cube(size=[60, 40, 20], r=3);

  translate([10, 10, 0]) cylinder(h=20, r=2, center=false);

}

enclosure();


\end{lstlisting}

\textbf{Fixed version} --- use \texttt{minkowski()} with a \texttt{sphere()} instead:

\begin{lstlisting}[style=Alabaster, language=openscad]
// ✅ CORRECT — native minkowski approach for rounded corners

module rounded_box(x=60, y=40, z=20, r=3) {

  minkowski() {

    cube([x - 2*r, y - 2*r, z - 2*r]);

    sphere(r=r, $fn=32);

  }

}

rounded_box();


\end{lstlisting}

\emph{What to learn from this:} If an AI gives you OpenSCAD code that uses a function name you haven\textquotesingle t seen before, look it up in the OpenSCAD documentation before using it. A compile error at render time is the warning sign.

\begin{center}\rule{0.5\linewidth}{0.5pt}\end{center}

\textbf{Example 2: Coplanar Boolean (Missing 0.001 Offset)}

When the subtracted shape shares an exact face with the outer shape, some renderers produce Z-fighting artifacts and unreliable geometry.

\begin{lstlisting}[style=Alabaster, language=openscad]
// ❌ FLAWED — bottom faces are coplanar (Z=0 on both shapes)

difference() {

  cube([30, 30, 10]);

  translate([5, 5, 0]) cube([20, 20, 10]);

}


\end{lstlisting}

\textbf{Fixed version} --- extend the subtracted shape by 0.001 mm past each coplanar face:

\begin{lstlisting}[style=Alabaster, language=openscad]
// ✅ CORRECT — 0.001 overshoot on both Z faces

difference() {

  cube([30, 30, 10]);

  translate([5, 5, -0.001]) cube([20, 20, 10 + 0.002]);

}


\end{lstlisting}

\emph{What to learn from this:} Any time a \texttt{difference()} subtraction should go all the way through a face, add a 0.001 mm overshoot. This is a mechanical rule, not an approximation --- apply it consistently.

\begin{center}\rule{0.5\linewidth}{0.5pt}\end{center}

\textbf{Example 3: Wrong Transform Order (Misplaced Holes)}

Rotate-then-translate and translate-then-rotate produce different results. AI often gets this wrong.

\begin{lstlisting}[style=Alabaster, language=openscad]
// ❌ FLAWED — rotate applied before translate; hole ends up in wrong position

module bracket(width=40, height=20, th=3, hole_r=3) {

  difference() {

    cube([width, th, height]);

    translate([width/2, th/2, height/2])

      rotate([90, 0, 0])

        cylinder(h=th, r=hole_r, center=true);

  }

}

bracket();


\end{lstlisting}

\textbf{Fixed version} --- the translate sets the center, then the rotate orients correctly:

\begin{lstlisting}[style=Alabaster, language=openscad]
// ✅ CORRECT — translate to center, then rotate into the correct plane

module bracket(width=40, height=20, th=3, hole_r=3) {

  difference() {

    cube([width, th, height]);

    translate([width/2, th/2, height/2])

      rotate([90, 0, 0])

        cylinder(h=th + 0.002, r=hole_r, center=true);

  }

}

bracket();


\end{lstlisting}

\emph{What to learn from this:} Always check hole placement in the slicer layer view. The AI\textquotesingle s code may look syntactically correct but produce geometry that is correct only for a specific, unverified set of assumptions about how transforms compose.

\begin{center}\rule{0.5\linewidth}{0.5pt}\end{center}

\textbf{Example 4: Procedural Variable Reassignment}

OpenSCAD is a declarative language --- variables are not mutable. AI trained on Python, JavaScript, or other imperative languages sometimes generates code that assumes variables can be updated.

\begin{lstlisting}[style=Alabaster, language=openscad]
// ❌ FLAWED — OpenSCAD does not allow reassignment of variables

width = 60;

width = width + 10; // This silently uses the LAST assignment; the first is ignored

cube([width, 20, 10]);


\end{lstlisting}

\textbf{Fixed version} --- use distinct names for derived values:

\begin{lstlisting}[style=Alabaster, language=openscad]
// ✅ CORRECT — derive a new variable from the base value

base_w  = 60;

final_w = base_w + 10;

cube([final_w, 20, 10]);


\end{lstlisting}

\emph{What to learn from this:} In OpenSCAD, a variable that appears assigned twice will use the last assignment for the entire file --- including places that appear before the second assignment in the code. This is a source of subtle bugs. Always use a new name for derived values.

\begin{center}\rule{0.5\linewidth}{0.5pt}\end{center}

\textbf{Example 5: Unsafe Parameter Limits (Geometry Inversion)}

AI-generated modules often lack parameter validation. When given edge-case values, the geometry inverts silently.

\begin{lstlisting}[style=Alabaster, language=openscad]
// ❌ FLAWED — no validation; wall=12 on a w=20 box inverts the geometry

module box(w=30, d=30, h=20, wall=3) {

  difference() {

    cube([w, d, h]);

    translate([wall, wall, -0.001])

      cube([w - 2*wall, d - 2*wall, h + 0.002]);

  }

}

box(w=20, d=20, h=15, wall=12); // wall larger than half the width — silent failure


\end{lstlisting}

\textbf{Fixed version} --- add \texttt{assert()} to catch invalid combinations:

\begin{lstlisting}[style=Alabaster, language=openscad]
// ✅ CORRECT — assert() catches invalid parameters before geometry is computed

module box(w=30, d=30, h=20, wall=3) {

  assert(w > 0 && d > 0 && h > 0, "All dimensions must be positive");

  assert(wall > 0 && wall < min(w, d) / 2,

    str("wall (", wall, "mm) must be less than min(w,d)/2 = ", min(w,d)/2, "mm"));

  difference() {

    cube([w, d, h]);

    translate([wall, wall, -0.001])

      cube([w - 2*wall, d - 2*wall, h + 0.002]);

  }

}

box(w=20, d=20, h=15, wall=12); // will halt with a clear error message


\end{lstlisting}

\emph{What to learn from this:} \texttt{assert()} is your safety net for parametric designs. Any relationship between parameters that could produce invalid geometry should have an \texttt{assert()} checking it. This is especially important for designs you share with others.

\begin{center}\rule{0.5\linewidth}{0.5pt}\end{center}

\textbf{Example 6: Hallucinated Text Function + Misplaced Emboss}

Two compound errors: a function that doesn\textquotesingle t exist (\texttt{fillet\_text()}), and a text position that extends outside the model face.

\begin{lstlisting}[style=Alabaster, language=openscad]
// ❌ FLAWED — fillet_text() does not exist; second translate puts text outside the body

difference() {

  cube([80, 40, 6]);

  translate([0, 0, 0])

    fillet_text("RYAN", size=10, radius=1); // bogus function

  translate([70, 20, 3])

    linear_extrude(2) text("RYAN", size=10); // misplaced — extends off the edge

}


\end{lstlisting}

\textbf{Fixed version} --- use standard \texttt{text()} with \texttt{linear\_extrude()}, centered on the face:

\begin{lstlisting}[style=Alabaster, language=openscad]
// ✅ CORRECT — centered debossed text using standard OpenSCAD functions

module top_deboss(txt="RYAN", size=10, depth=1) {

  translate([40, 20, 6 + 0.001])  // position at center of top face

    linear_extrude(height=depth)

      text(txt, size=size, halign="center", valign="center");

}


difference() {

  cube([80, 40, 6]);

  top_deboss();

}


\end{lstlisting}

\emph{What to learn from this:} Text positioning requires knowing the face center, not just a rough location. Use \texttt{halign="center"} and \texttt{valign="center"} with a translate to the face center. The 0.001 mm overshoot on the Z translate ensures the subtracted text breaks through the face cleanly.

\begin{center}\rule{0.5\linewidth}{0.5pt}\end{center}

\textbf{Example 7: Zero Clearance on a Fitted Lid}

AI-generated box and lid designs often omit clearance, making the lid impossible to remove after printing.

\begin{lstlisting}[style=Alabaster, language=openscad]
// ❌ FLAWED — zero clearance means the lid walls will fuse to the box walls

module box_outer(sz=[60, 40, 20], wall=2) { cube(sz); }

module box_inner(sz=[60, 40, 20], wall=2) { cube([sz[0]-2*wall, sz[1]-2*wall, sz[2]]); }


difference() {

  box_outer();

  translate([2, 2, -0.001]) box_inner(); // clearance = 0; lid won't come off

}


\end{lstlisting}

\textbf{Fixed version} --- add a clearance parameter to the socket dimensions:

\begin{lstlisting}[style=Alabaster, language=openscad]
// ✅ CORRECT — clearance creates the physical gap needed for the lid to fit and move

module box_with_clearance(w=60, d=40, h=20, wall=2, clr=0.25) {

  assert(clr > 0 && clr < 1.0, "Clearance should be between 0 and 1mm for most FDM fits");

  difference() {

    cube([w, d, h]);

    translate([wall + clr, wall + clr, -0.001])

      cube([w - 2*(wall + clr), d - 2*(wall + clr), h + 0.002]);

  }

}

box_with_clearance();


\end{lstlisting}

\emph{What to learn from this:} FDM-printed parts that are designed to fit together always need clearance. A nominal 0mm gap at design time produces a stuck or impossible fit after printing. Typical clearance for a slip fit is 0.2--0.3 mm; for a press fit, 0.0--0.1 mm. This topic is covered in depth in Lesson 8.
-e

\begin{center}\rule{0.5\linewidth}{0.5pt}\end{center}

\section*{Safety Protocols and the Physical Fabrication Interface}\label{docs__pandoc__latex__src__3dmake_all_lessons.md__safety-protocols-and-the-physical-fabrication-interface}

\textbf{Estimated Time:} 90--120 minutes\\
\textbf{Course:} 3dMake Certification --- High School Distance Learning Track

\begin{center}\rule{0.5\linewidth}{0.5pt}\end{center}

\subsection*{Before You Start}\label{docs__pandoc__latex__src__3dmake_all_lessons.md__before-you-start-3}

Every lesson so far has been entirely digital: code on a screen, commands in a terminal, virtual geometry. Now we\textquotesingle re moving into the physical world. 3D printers are machines that get very hot, draw substantial electrical current, and produce airborne particles and vapors. This lesson exists because these machines are genuinely safe when you know what you\textquotesingle re doing, and genuinely dangerous when you don\textquotesingle t.

The good news is that 3D printing safety is not complicated. The hazards are well-understood, the mitigations are straightforward, and the incidents that do happen are almost always caused by ignoring simple rules. The goal of this lesson is to make sure you know the rules and understand \emph{why} they exist --- not just as things to memorize for a quiz, but as habits you can apply in any shop or makerspace you ever work in.

\begin{center}\rule{0.5\linewidth}{0.5pt}\end{center}

\subsection*{Learning Objectives}\label{docs__pandoc__latex__src__3dmake_all_lessons.md__learning-objectives-8}

By the end of this lesson you will be able to:

\begin{itemize}
\tightlist
\item
  Identify and apply the Hierarchy of Controls for 3D printing hazards
\item
  Apply fire, electrical, and fume safety protocols for FDM printing
\item
  Understand filament-specific hazards and safe material handling
\item
  Follow a safe printer startup, monitoring, and shutdown sequence
\item
  Select appropriate filament materials for given projects
\item
  Use pre-flight build scripts to estimate cost before committing to a print
\item
  Design parts that minimize safety risks during printing
\end{itemize}

\begin{center}\rule{0.5\linewidth}{0.5pt}\end{center}

\subsection*{Concept 1: The Hierarchy of Controls}\label{docs__pandoc__latex__src__3dmake_all_lessons.md__concept-1-the-hierarchy-of-controls}

\subsubsection*{Why We Don\textquotesingle t Just "Be Careful"}\label{docs__pandoc__latex__src__3dmake_all_lessons.md__why-we-dont-just-be-careful}

When it comes to workplace and lab safety, the common advice "just be careful" is the least effective approach available. It relies entirely on individual attention, and attention is inconsistent --- everyone gets tired, distracted, or hurried sometimes. Professional safety systems are designed to work even when individual attention lapses.

The \textbf{Hierarchy of Controls} is a framework developed by OSHA (the U.S. Occupational Safety and Health Administration) and adopted worldwide. It ranks safety strategies from most effective to least effective, and it says: apply the most effective strategy you can afford, and layer additional strategies below it.

The five levels, from most to least effective:

\textbf{1. Elimination} --- Remove the hazard entirely. If a filament requires a sealed enclosure with active filtration and you\textquotesingle re in an open classroom, the safest choice is simply not to use that filament. If a print requires supports made from difficult material, redesign the print to not need supports. Elimination always wins.

\textbf{2. Substitution} --- Replace the hazard with a less dangerous one. Use PLA instead of ABS. Use a PEI spring steel bed instead of a glass bed that can shatter. Substitution keeps the functionality while reducing the risk.

\textbf{3. Engineering Controls} --- Build protection into the environment. A fume enclosure with a HEPA + activated carbon filter. Thermal runaway protection in the printer firmware. A smoke detector mounted above the printer. These controls work automatically without relying on human behavior.

\textbf{4. Administrative Controls} --- Rules, procedures, and training. The startup checklist, the monitoring schedule, the training program you\textquotesingle re in right now. Administrative controls are less reliable than engineering ones because they depend on people following rules consistently.

\textbf{5. Personal Protective Equipment (PPE)} --- Gloves, goggles, respirators. PPE is the last resort --- it protects one person, it can fail or be used incorrectly, and it doesn\textquotesingle t actually remove the hazard. Always apply higher-level controls first and use PPE as a supplement.

\textbf{The key insight:} many beginners think that putting on gloves and goggles makes them safe. In the hierarchy, PPE is the weakest protection. Real safety comes from eliminating or substituting hazards, and then engineering the remaining risks away. PPE is what you use when those options aren\textquotesingle t available.

\begin{center}\rule{0.5\linewidth}{0.5pt}\end{center}

\subsection*{Step 1 --- Fire and Electrical Safety}\label{docs__pandoc__latex__src__3dmake_all_lessons.md__step-1--fire-and-electrical-safety}

\subsubsection*{The Golden Rule}\label{docs__pandoc__latex__src__3dmake_all_lessons.md__the-golden-rule}

\textbf{Never leave an FDM printer completely unattended during the first layer.} Print failures most often happen at the start --- the first layer either bonds to the bed or it doesn\textquotesingle t. If it doesn\textquotesingle t bond and nobody is watching, the printer may continue for hours, dragging filament across the bed in a mess called "spaghetti" that can tangle around the print head, cause jams, or --- in rare worst cases --- create a fire hazard.

Be present and watching for at least the first 15 minutes of every print. For long prints, check in every 30 minutes or set up remote monitoring via OctoPrint or a camera.

\subsubsection*{Fire Hazards}\label{docs__pandoc__latex__src__3dmake_all_lessons.md__fire-hazards}

Fire risk in FDM printing comes from two main sources. First, \textbf{electrical failure}: wiring that cracks from repeated movement, loose connections at terminal blocks, or a faulty temperature sensor that allows the heater to run without limit. Second, \textbf{filament accumulation}: plastic that melts and collects on the hot nozzle can sometimes ignite if the build-up becomes large enough.

Fire safety protocols:

\begin{itemize}
\tightlist
\item
  Keep a \textbf{Class ABC fire extinguisher} within arm\textquotesingle s reach of every printer
\item
  Never run the printer while away from the building
\item
  Keep the print area clear of flammable materials: paper, fabric, solvent containers
\item
  Mount a \textbf{smoke detector} directly above the printer so rising smoke reaches it quickly
\item
  Know where the power switch is \emph{before} you start --- cutting power is the first response to any emergency
\item
  If smoke appears: cut power immediately, do not try to continue the print, and report the incident
\end{itemize}

\subsubsection*{Electrical Safety}\label{docs__pandoc__latex__src__3dmake_all_lessons.md__electrical-safety}

\begin{itemize}
\tightlist
\item
  Use a \textbf{surge protector} to protect the printer\textquotesingle s control board from voltage spikes
\item
  Inspect all wiring periodically, especially the flexible cable bundle that runs to the heated bed --- this bundle flexes with every print and can crack the insulation over time
\end{itemize}

Warning

\textbf{Thermal runaway protection must be enabled.} This firmware feature monitors the temperature sensor and shuts the heater off if the temperature climbs unexpectedly or the sensor fails. A printer without thermal runaway protection can overheat catastrophically. Never disable this protection. If your printer doesn\textquotesingle t have it, do not use it.

\begin{itemize}
\tightlist
\item
  Never open the power supply enclosure or work on any electrical component while the printer is plugged in. The power supply contains mains voltage (120V or 240V) that can kill.
\end{itemize}

\begin{center}\rule{0.5\linewidth}{0.5pt}\end{center}

\subsection*{Step 2 --- Fume Safety}\label{docs__pandoc__latex__src__3dmake_all_lessons.md__step-2--fume-safety}

\subsubsection*{Why Fumes Matter}\label{docs__pandoc__latex__src__3dmake_all_lessons.md__why-fumes-matter}

All FDM printing produces some airborne emissions. These fall into two categories: \textbf{particulates} (tiny plastic particles, especially ultrafine particles below 0.1 microns) and \textbf{VOCs} (volatile organic compounds --- chemical vapors that off-gas from the hot plastic). Both can affect respiratory health with repeated exposure, and the risk varies significantly by material.

{\def\LTcaptype{none} % do not increment counter
\begin{longtable}[]{@{}
  >{\raggedright\arraybackslash}p{(\linewidth - 4\tabcolsep) * \real{0.1280}}
  >{\raggedright\arraybackslash}p{(\linewidth - 4\tabcolsep) * \real{0.0800}}
  >{\raggedright\arraybackslash}p{(\linewidth - 4\tabcolsep) * \real{0.7920}}@{}}
\toprule\noalign{}
\begin{minipage}[b]{\linewidth}\raggedright
Material
\end{minipage} & \begin{minipage}[b]{\linewidth}\raggedright
Fume Risk
\end{minipage} & \begin{minipage}[b]{\linewidth}\raggedright
Required Mitigation
\end{minipage} \\
\midrule\noalign{}
\endhead
\bottomrule\noalign{}
\endlastfoot
PLA & Low & Good room ventilation (open window or HVAC on) \\
PETG & Low--Medium &
Active ventilation; avoid enclosed spaces without air exchange \\
\end{longtable}
}

Warning

ABS and ASA are \textbf{NOT safe for open classrooms} without a dedicated fume enclosure with HEPA + activated carbon filtration. Do not print ABS or ASA in any space without proper exhaust infrastructure.

\textbar{} ABS              \textbar{} High       \textbar{} Dedicated fume enclosure with HEPA + activated carbon filter required; NOT safe for open classrooms \textbar{}
\textbar{} ASA              \textbar{} High       \textbar{} Same as ABS                                                                                         \textbar{}
\textbar{} TPU              \textbar{} Medium     \textbar{} Active ventilation; keep nozzle temperature at the minimum needed                                   \textbar{}
\textbar{} Resin (SLA/MSLA) \textbar{} Very High  \textbar{} Full enclosure; nitrile gloves; UV-blocking goggles; no open-workspace use                          \textbar{}

\subsubsection*{The Safe Choice: PLA}\label{docs__pandoc__latex__src__3dmake_all_lessons.md__the-safe-choice-pla}

Tip

\textbf{PLA is the safest classroom filament available.} Made from bio-derived sources, it prints at low temperatures (190--220°C) and produces the lowest emissions of any common FDM material. Default to PLA for any project where mechanical requirements don\textquotesingle t specifically demand something else.

\subsubsection*{What "Good Ventilation" Actually Means}\label{docs__pandoc__latex__src__3dmake_all_lessons.md__what-good-ventilation-actually-means}

"Good ventilation" means that the air in the room is actively exchanging with outside air. This is not the same as "a window that\textquotesingle s crackable." Active ventilation means:

\begin{itemize}
\tightlist
\item
  A window open with a fan pulling air through the room, or
\item
  An HVAC system that is running and moving air, not just recirculating internal air
\end{itemize}

If you\textquotesingle re printing ABS or ASA, room ventilation is not sufficient. You need an enclosure with a filtered exhaust system --- a box around the printer with an internal fan that pulls fumes through a HEPA filter (for particles) and an activated carbon filter (for VOCs) before exhausting outside or into a large ventilated space.

\begin{center}\rule{0.5\linewidth}{0.5pt}\end{center}

\subsection*{Step 3 --- Material Selection}\label{docs__pandoc__latex__src__3dmake_all_lessons.md__step-3--material-selection}

\subsubsection*{Choosing the Right Filament}\label{docs__pandoc__latex__src__3dmake_all_lessons.md__choosing-the-right-filament}

The right material for a project depends on what the part needs to do and where it will be used. This table gives you a practical guide:

{\def\LTcaptype{none} % do not increment counter
\begin{longtable}[]{@{}
  >{\raggedright\arraybackslash}p{(\linewidth - 10\tabcolsep) * \real{0.0530}}
  >{\raggedright\arraybackslash}p{(\linewidth - 10\tabcolsep) * \real{0.0728}}
  >{\raggedright\arraybackslash}p{(\linewidth - 10\tabcolsep) * \real{0.1126}}
  >{\raggedright\arraybackslash}p{(\linewidth - 10\tabcolsep) * \real{0.1589}}
  >{\raggedright\arraybackslash}p{(\linewidth - 10\tabcolsep) * \real{0.1060}}
  >{\raggedright\arraybackslash}p{(\linewidth - 10\tabcolsep) * \real{0.4967}}@{}}
\toprule\noalign{}
\begin{minipage}[b]{\linewidth}\raggedright
Filament
\end{minipage} & \begin{minipage}[b]{\linewidth}\raggedright
Strength
\end{minipage} & \begin{minipage}[b]{\linewidth}\raggedright
Flexibility
\end{minipage} & \begin{minipage}[b]{\linewidth}\raggedright
Heat Resistance
\end{minipage} & \begin{minipage}[b]{\linewidth}\raggedright
Print Difficulty
\end{minipage} & \begin{minipage}[b]{\linewidth}\raggedright
Best Uses
\end{minipage} \\
\midrule\noalign{}
\endhead
\bottomrule\noalign{}
\endlastfoot
PLA & Medium & Rigid & Low (softens at 50--60°C) & Easy &
Prototypes, display models, artwork, educational parts \\
PETG & Medium-High & Slightly flexible & Medium (75--80°C) &
Easy--Medium &
Functional parts, food-adjacent applications, water contact \\
ABS & High & Semi-rigid & High (85--100°C) & Hard &
Hot environments, automotive interiors, post-processing (acetone smoothing) \\
TPU & Medium & Very flexible & Medium & Medium &
Gaskets, grips, phone cases, flexible joints \\
Nylon & Very High & Flexible & High & Hard &
Mechanical parts, gears, load-bearing functional parts \\
ASA & High & Semi-rigid & Very High (90°C+) & Hard &
Outdoor use, UV-exposed parts \\
\end{longtable}
}

\subsubsection*{What "Hard" to Print Actually Means}\label{docs__pandoc__latex__src__3dmake_all_lessons.md__what-hard-to-print-actually-means}

When a material is listed as "hard," that means it requires one or more of the following: a higher nozzle temperature, a heated enclosure to prevent warping, special bed adhesion (glue stick, hairspray, specialized plates), slower print speeds, or more careful tuning of retraction settings. Starting with a difficult material before you have PLA dialed in is a recipe for frustration and wasted filament.

Tip

\textbf{The practical material rule:} use PLA until you have a specific functional reason not to. Need better heat resistance? Try PETG next. Reserve ABS and specialty materials for when you have the appropriate safety infrastructure.

\subsubsection*{PLA and Sustainability --- A Realistic Picture}\label{docs__pandoc__latex__src__3dmake_all_lessons.md__pla-and-sustainability--a-realistic-picture}

PLA is often marketed as "biodegradable" and "eco-friendly." These claims deserve some scrutiny. PLA \emph{is} bio-derived --- made from fermented plant starch rather than petroleum --- which reduces its carbon footprint somewhat compared to petroleum plastics. But PLA only biodegrades under very specific industrial composting conditions (temperatures above 58°C sustained for weeks). It does not break down in a home compost bin or a landfill.

The most sustainable approach to FDM printing is not any particular material --- it\textquotesingle s designing durable parts that don\textquotesingle t need to be reprinted, using the minimum infill that meets your mechanical requirements, and avoiding test prints you don\textquotesingle t need (which is exactly why this course emphasizes verification before printing).

\begin{center}\rule{0.5\linewidth}{0.5pt}\end{center}

\subsection*{Step 4 --- Estimating Print Cost Before You Print}\label{docs__pandoc__latex__src__3dmake_all_lessons.md__step-4--estimating-print-cost-before-you-print}

\subsubsection*{Always Estimate First}\label{docs__pandoc__latex__src__3dmake_all_lessons.md__always-estimate-first}

Tip

Always estimate filament cost and print time \textbf{before} sending a model to print. For small parts the cost is trivial, but large high-infill parts can cost several dollars and several hours --- time you can\textquotesingle t recover if the part needs revision.

The estimate formula from Lesson 1 applies here:

\begin{lstlisting}[style=Alabaster]
mass (grams)  = (volume_mm³ ÷ 1000) × filament_density_g_per_cm³

cost ($)      = mass × (spool_price_$ ÷ spool_weight_g)


\end{lstlisting}

The preflight scripts below automate this calculation and require you to confirm before proceeding.

\subsubsection*{Pre-Flight Cost Check Scripts}\label{docs__pandoc__latex__src__3dmake_all_lessons.md__pre-flight-cost-check-scripts}

\textbf{Linux / macOS (bash):}

\begin{lstlisting}[style=Alabaster, language=bash]
#!/bin/bash

set -e


3dm build

INFO=$(3dm info)

echo "$INFO"


# Extract volume from 3dm info output

VOLUME=$(echo "$INFO" | grep -oP '(?<=Volume:\s)\d+\.?\d*')

if [ -z "$VOLUME" ]; then

  echo "ERROR: Could not parse volume from 3dm info output"

  exit 1

fi


DENSITY=1.24    # g/cm³ — PLA density (change for other materials)

COST_PER_G=0.02 # $/g  — adjust for your spool price


MASS=$(echo "$VOLUME $DENSITY" | awk '{printf "%.1f", ($1/1000)*$2}')

COST=$(echo "$MASS $COST_PER_G" | awk '{printf "%.2f", $1*$2}')


echo ""

echo "=== PRE-FLIGHT SUMMARY ==="

echo "Volume:   ${VOLUME} mm³"

echo "Mass:     ${MASS} g  (PLA @ ${DENSITY} g/cm³)"

echo "Filament: \$${COST}  (@ \$${COST_PER_G}/g)"

echo ""

echo "Proceed? [y/N]"

read CONFIRM

if [ "$CONFIRM" != "y" ]; then

  echo "Cancelled."

  exit 0

fi

echo "Approved. Send to slicer."


\end{lstlisting}

\textbf{Windows (PowerShell):}

\begin{lstlisting}[style=Alabaster, language=powershell]
3dm build

$info = 3dm info

Write-Host $info


$volumeMatch = $info | Select-String -Pattern "Volume:\s+([\d.]+)"

if (-not $volumeMatch) { Write-Error "Could not parse volume"; exit 1 }

$volume = [double]$volumeMatch.Matches.Groups[1].Value


$density  = 1.24    # PLA

$costPerG = 0.02

$mass     = ($volume / 1000) * $density

$cost     = $mass * $costPerG


Write-Host ""

Write-Host "=== PRE-FLIGHT SUMMARY ==="

Write-Host ("Volume:   {0:F0} mm3" -f $volume)

Write-Host ("Mass:     {0:F1} g  (PLA)" -f $mass)

Write-Host ("`$Filament: {0:F2}" -f $cost)

Write-Host ""

$confirm = Read-Host "Proceed? [y/N]"

if ($confirm -ne "y") { Write-Host "Cancelled."; exit 0 }

Write-Host "Approved. Send to slicer."


\end{lstlisting}

\begin{center}\rule{0.5\linewidth}{0.5pt}\end{center}

\subsection*{Step 5 --- The Printer Startup Checklist}\label{docs__pandoc__latex__src__3dmake_all_lessons.md__step-5--the-printer-startup-checklist}

\subsubsection*{Before Every Print Session}\label{docs__pandoc__latex__src__3dmake_all_lessons.md__before-every-print-session}

Use this checklist before starting any print. It takes about 5 minutes and prevents the most common causes of print failure.

\begin{lstlisting}[style=Alabaster]
PRINTER STARTUP CHECKLIST

□ Bed level verified (run auto-level, or manually check at all four corners)

□ Build plate clean (wipe with IPA — remove finger oils and dust)

□ Nozzle cleared (cold pull if returning from more than a day idle)

□ Filament loaded and extruding cleanly (run 50mm manual extrude purge)

□ First layer height set correctly (0.2mm gap typical for 0.4mm nozzle)

□ Print fan operational (watch it spin up at the start)

□ Workspace clear of flammable materials

□ You will remain present for the first 15 minutes of the print

□ Power switch is accessible and you know where it is


\end{lstlisting}

\subsubsection*{Why Each Item Matters}\label{docs__pandoc__latex__src__3dmake_all_lessons.md__why-each-item-matters}

Important

\textbf{Bed leveling is the single most critical setting for successful prints.} A bed that is just 0.1 mm too high or too low can cause every print to fail. Always verify leveling before each print session. Most modern printers have automatic bed leveling (ABL), but you should still understand how it works and check that it\textquotesingle s running correctly.

\textbf{Clean build plate} prevents adhesion failures. Skin oils from touching the plate are invisible but dramatically reduce how well plastic sticks. An isopropyl alcohol (IPA) wipe before every print takes 30 seconds and prevents a common failure mode.

\textbf{Cold pull (nozzle clearing)} removes degraded filament from inside the nozzle. Heat the nozzle to print temperature, push some filament through, then let it cool to about 90°C and pull the filament out firmly. The end of the pulled filament should bring out any debris or degraded material from inside the hot end.

\textbf{Filament loaded correctly} means the extruder is feeding without clicking or skipping, the filament is straight in the Bowden tube (if present), and the purge line at the start of the print flows smoothly with no bubbles or gaps.

\begin{center}\rule{0.5\linewidth}{0.5pt}\end{center}

\subsection*{Step 6 --- Monitoring During Prints}\label{docs__pandoc__latex__src__3dmake_all_lessons.md__step-6--monitoring-during-prints}

\subsubsection*{What to Watch For}\label{docs__pandoc__latex__src__3dmake_all_lessons.md__what-to-watch-for}

Checking in periodically during a print takes only a few seconds each time and lets you catch problems early. A problem caught at layer 10 is a minor inconvenience. The same problem discovered at layer 100 means you\textquotesingle ve wasted 90 layers of material and print time.

Warning

\textbf{Spaghetti} is the most common print failure: the part lifts off the bed and the nozzle drags filament through the air. If you see this, \textbf{pause or cancel the print immediately} --- catching it at layer 10 is a minor inconvenience; catching it at layer 100 wastes 90 layers of material and time.

Caution

\textbf{Layer shifts} (horizontal displacements in the print) usually indicate a mechanical problem --- belt tension, loose pulley, or a motor skipping steps. Cancel the print and diagnose the cause before trying again.

\textbf{Nozzle clog} --- the extruder starts making a clicking or ticking sound as the gear slips because material isn\textquotesingle t flowing. Filament may start grinding. Pause the print and clear the clog before restarting.

\textbf{Stringing} --- thin wisps of filament connecting different parts of the print. Usually a retraction or temperature setting issue rather than an emergency, but it may require post-processing cleanup.

Warning

\textbf{Smoke is an immediate emergency.} Cut power immediately. Do not attempt to continue or diagnose while the printer is running. Notify your instructor. Do not restart the printer until the cause has been fully identified and corrected.

\begin{center}\rule{0.5\linewidth}{0.5pt}\end{center}

\subsection*{Step 7 --- Safe Shutdown and Part Removal}\label{docs__pandoc__latex__src__3dmake_all_lessons.md__step-7--safe-shutdown-and-part-removal}

\subsubsection*{The Right Way to End a Print}\label{docs__pandoc__latex__src__3dmake_all_lessons.md__the-right-way-to-end-a-print}

A print ending doesn\textquotesingle t mean the work is done. Follow these steps after every print:

Caution

\textbf{Wait for the nozzle to cool below 50°C} before reaching anywhere near the hot end. The nozzle can reach 200°C+ during printing and retains heat for several minutes after shutdown.

\begin{enumerate}
\item
  \textbf{Wait for the nozzle to cool below 50°C} before reaching anywhere near the hot end.
\item
  \textbf{Wait for the bed to cool below 30°C} before removing the part. Parts printed on a hot bed can warp slightly if removed while the bed is still warm. Many PEI-coated beds release parts naturally as they cool --- you may not even need tools.
\item
  \textbf{Use a spatula or palette knife} to gently separate the part from the build plate. Never force parts off with your fingers --- the edge of a part against a spring steel bed can cut skin. Apply the spatula at a low angle and use a gentle rocking motion.
\end{enumerate}

Note

\textbf{Store filament in a sealed bag with desiccant.} Filament absorbs moisture (hygroscopy), causing bubbling, popping sounds, and reduced mechanical properties during extrusion. This is especially critical for Nylon and PETG. A silica gel desiccant packet in a ziplock bag is sufficient for short-term storage.

\begin{center}\rule{0.5\linewidth}{0.5pt}\end{center}

\subsection*{Concept 2: Design for Printability and Safety}\label{docs__pandoc__latex__src__3dmake_all_lessons.md__concept-2-design-for-printability-and-safety}

\subsubsection*{Your Design Choices Affect Safety}\label{docs__pandoc__latex__src__3dmake_all_lessons.md__your-design-choices-affect-safety}

The decisions you make in OpenSCAD directly affect how safe and reliable the printing process will be. Good design doesn\textquotesingle t just mean a part that looks right --- it means a part that prints reliably without requiring dangerous materials, excessive supports, or heroic printer tuning.

\textbf{Minimize overhangs.} Any geometry that extends more than about 45° from vertical needs supports --- temporary scaffolding that the printer adds automatically. Supports increase print time, use more material, often leave rough marks where they\textquotesingle re removed, and increase the risk of print failure due to tangled support material. Good design reorients parts or modifies geometry to eliminate or reduce overhangs.

Caution

\textbf{Minimum wall thickness for a 0.4 mm nozzle:} walls below 0.8 mm (2× nozzle width) may not print reliably, and walls below \textasciitilde{}0.8 mm may be omitted entirely by the slicer. Use at least 1.2 mm (3× nozzle width) for any structural wall.

\textbf{Use the safest material that meets requirements.} If PLA is strong enough, use PLA. Don\textquotesingle t introduce PETG, ABS, or specialty materials unless you have a specific reason and the appropriate safety infrastructure.

\textbf{Encode printability constraints in your code:}

\begin{lstlisting}[style=Alabaster, language=openscad]
// Printability constants — these drive safe design decisions

wall_min     = 1.2;    // mm — minimum wall for 0.4mm nozzle (3× nozzle width)

hole_min_r   = 1.5;    // mm — minimum hole radius for reliable circular holes

clearance    = 0.2;    // mm — fitting clearance between mating faces

overhang_max = 45;     // degrees — maximum overhang without supports


module printable_box(w, d, h, wall=2) {

  // Validate against printability limits

  assert(wall >= wall_min, str("wall (", wall, "mm) must be >= ", wall_min, "mm"));

  assert(w > 2 * wall, "width must be greater than two wall thicknesses");

  assert(d > 2 * wall, "depth must be greater than two wall thicknesses");


  // Open-top box: no overhangs on the inside

  difference() {

    cube([w, d, h]);

    translate([wall, wall, wall])

      cube([w - 2*wall, d - 2*wall, h]);  // open at top — no overhang required

  }

}


printable_box(50, 40, 30);


\end{lstlisting}

\begin{center}\rule{0.5\linewidth}{0.5pt}\end{center}

\subsection*{Exercises}\label{docs__pandoc__latex__src__3dmake_all_lessons.md__exercises-1}

\textbf{Exercise 5.1:} Run \texttt{3dm\ info} on your model from Lesson 3. Calculate the material cost in PLA and in PETG (density 1.27 g/cm³). If PLA spools cost \$20/kg and PETG spools cost \$25/kg, what is the cost difference for your part?

\textbf{Exercise 5.2:} Write a pre-flight script (bash or PowerShell) that builds your model, reports the estimated cost, and requires a \texttt{y} confirmation before proceeding.

\textbf{Exercise 5.3 (Advanced):} Modify the pre-flight script to accept a \texttt{-\/-material} argument (options: \texttt{pla}, \texttt{petg}, \texttt{abs}, \texttt{tpu}) and calculate cost using the correct density for the selected material.

\begin{center}\rule{0.5\linewidth}{0.5pt}\end{center}

\subsection*{Quiz --- Lesson 5 (15 questions)}\label{docs__pandoc__latex__src__3dmake_all_lessons.md__quiz--lesson-5-15-questions}

\begin{enumerate}
\tightlist
\item
  What are the five levels of the Hierarchy of Controls, from most to least effective?
\item
  Why is Elimination considered the most effective safety control?
\item
  Which filament requires a dedicated fume enclosure with HEPA + activated carbon filtration and is NOT recommended for open classrooms?
\item
  At what temperature should you wait before removing a print from the bed?
\item
  What should you do immediately if you see smoke coming from your 3D printer?
\item
  Give three reasons why PLA is considered the safest classroom filament.
\item
  What is the minimum reliable wall thickness for a 0.4 mm nozzle, and why does this minimum exist?
\item
  What is the purpose of storing filament with desiccant?
\item
  True or False: It is safe to leave a 3D printer completely unattended for hours once the first layer has printed successfully.
\item
  What is "spaghetti" in the context of a 3D printing failure, and what causes it?
\item
  What is thermal runaway protection, and why is it critical that it remain enabled?
\item
  What are the two types of airborne emissions from FDM printing, and which filaments produce the most hazardous levels of each?
\item
  Describe one engineering control and one administrative control you could implement in a classroom 3D printing lab.
\item
  Explain why minimizing overhangs in your design is both a safety consideration and a quality consideration.
\item
  A student wants to print a car dashboard mount in ABS because it will be in a hot car. The classroom has no fume enclosure. Using the Hierarchy of Controls, what recommendations would you make?
\end{enumerate}

\begin{center}\rule{0.5\linewidth}{0.5pt}\end{center}

\subsection*{Extension Problems (15)}\label{docs__pandoc__latex__src__3dmake_all_lessons.md__extension-problems-15-4}

\begin{enumerate}
\tightlist
\item
  Create a material selection flowchart that guides a user from mechanical requirements to the safest filament that meets them.
\item
  Conduct a ventilation assessment of your classroom or makerspace. Document air exchange methods, window locations, and any existing filtration.
\item
  Write a Standard Operating Procedure (SOP) for ABS printing that includes all required safety equipment and ventilation infrastructure.
\item
  Design and describe a "first layer test tile" --- a 100 × 100 × 0.4 mm flat tile --- and explain how you\textquotesingle d use it to verify bed leveling.
\item
  Build a parametric filament storage clip in OpenSCAD: a clip that holds the free end of a filament spool. Parameters: filament diameter, clip width.
\item
  Research OSHA\textquotesingle s published guidance on VOC exposure in 3D printing environments. Summarize the key recommendations.
\item
  Compare the Safety Data Sheets (SDS) for PLA and ABS filament from two manufacturers. Document differences in recommended exposure limits and PPE.
\item
  Design a simple printer enclosure in OpenSCAD: four walls, a front door opening, and a top panel with a circular vent hole. Document the parametric variables.
\item
  Create a safety poster covering the five most important 3D printing safety rules, organized using the Hierarchy of Controls framework.
\item
  Write a one-page risk assessment for a new FDM printer being added to a classroom.
\item
  Design and describe a parametric filament moisture indicator holder: a small box that holds a humidity indicator card inside a resealable filament storage bag.
\item
  Build a "print monitoring log" template: a paper form with columns for time, nozzle temperature, bed temperature, layer number, visual observations, and action taken.
\item
  Research the difference between particle emissions and VOC emissions from FDM printing. Which is considered more hazardous at typical classroom distances from the printer?
\item
  Design a parametric build plate corner protector: a small clip that attaches to the edge of a glass or spring steel build plate to protect it from spatula damage.
\item
  Write a hypothetical "near-miss incident report" for a fictional 3D printing incident using a standard workplace incident report format (date, location, description, contributing factors, corrective actions).
\end{enumerate}

\begin{center}\rule{0.5\linewidth}{0.5pt}\end{center}

\subsection*{References and Helpful Resources}\label{docs__pandoc__latex__src__3dmake_all_lessons.md__references-and-helpful-resources-4}

\begin{itemize}
\item
  OSHA Hierarchy of Controls --- https://www.osha.gov/hierarchy-of-controls / NIOSH --- https://www.cdc.gov/niosh/topics/hierarchy/default.html
\item
  UL Research Institutes --- 3D Printing Emissions Study --- https://www.ul.com/news/ul-research-institutes-releases-3d-printing-emissions-study
\item
  NIOSH Science Blog --- Health and Safety Considerations for 3D Printing --- https://blogs.cdc.gov/niosh-science-blog/2020/05/14/3d-printing/
\item
  PrusaSlicer Documentation --- https://docs.prusa3d.com/en/
\item
  All3DP Filament Types --- https://all3dp.com/1/3d-printer-filament-types-3d-printing-3d-filament/
\item
  MatterHackers Filament Compare --- https://www.matterhackers.com/3d-printer-filament-compare
\item
  Prusa Research Materials Guide --- https://help.prusa3d.com/materials
\end{itemize}

\subsubsection*{Supplemental Resources}\label{docs__pandoc__latex__src__3dmake_all_lessons.md__supplemental-resources-4}

\begin{itemize}
\tightlist
\item
  3DMake GitHub Repository --- https://github.com/tdeck/3dmake
\item
  OpenSCAD User Manual --- https://en.wikibooks.org/wiki/OpenSCAD\_User\_Manual
\item
  All other reference documents are embedded as appendices below --- no separate files needed.
\end{itemize}

\begin{center}\rule{0.5\linewidth}{0.5pt}\end{center}

\section*{Part 2: Reference Appendices}\label{docs__pandoc__latex__src__3dmake_all_lessons.md__part-2-reference-appendices}

\emph{All supporting documents for Lesson 5 are included below. You do not need to open any additional files.}

\begin{center}\rule{0.5\linewidth}{0.5pt}\end{center}

\subsection*{Appendix A: Safety Checklist for 3D Printing}\label{docs__pandoc__latex__src__3dmake_all_lessons.md__appendix-a-safety-checklist-for-3d-printing}

\emph{Complete this checklist before and after each printing session. Print it out and keep a copy at your printer station.}

\begin{center}\rule{0.5\linewidth}{0.5pt}\end{center}

\subsubsection*{Pre-Print Setup}\label{docs__pandoc__latex__src__3dmake_all_lessons.md__pre-print-setup}

\textbf{Work Area}

\begin{itemize}
\tightlist
\item[$\square$]
  Desk or table cleared of unnecessary items
\item[$\square$]
  Tripping hazards removed from around the printer
\item[$\square$]
  Adequate ventilation confirmed (window open, fan running, or fume enclosure active)
\item[$\square$]
  Class ABC fire extinguisher within arm\textquotesingle s reach
\item[$\square$]
  Smoke detector mounted above printer is functional
\end{itemize}

\textbf{Printer Inspection}

\begin{itemize}
\tightlist
\item[$\square$]
  Nozzle is clean --- no old filament residue baked on
\item[$\square$]
  Build plate is level (re-level if the printer has been moved or bumped)
\item[$\square$]
  Heated bed temperature sensor is attached and secure
\item[$\square$]
  All cables are secure and none are frayed or cracked --- especially the flexible bed cable bundle
\item[$\square$]
  Thermal runaway protection is enabled (check firmware settings if unsure)
\end{itemize}

\textbf{Filament Preparation}

\begin{itemize}
\tightlist
\item[$\square$]
  Filament spool rotates freely on the holder --- no binding or tangling
\item[$\square$]
  Filament path is clear of obstructions from spool to extruder
\item[$\square$]
  Extruder drive gear is clean --- not clogged with ground plastic
\item[$\square$]
  Filament is correctly loaded and primed (a few mm of filament extruded before starting)
\end{itemize}

\textbf{Environmental Conditions}

\begin{itemize}
\tightlist
\item[$\square$]
  Room temperature is adequate (18--25°C ideal --- extreme cold slows bed heating; extreme heat can soften PLA)
\item[$\square$]
  No direct drafts from windows or air conditioning blowing on the printer (drafts cause warping and layer splitting)
\item[$\square$]
  Adequate lighting to clearly see the first layer as it prints
\end{itemize}

\begin{center}\rule{0.5\linewidth}{0.5pt}\end{center}

\subsubsection*{During Print}\label{docs__pandoc__latex__src__3dmake_all_lessons.md__during-print}

\textbf{First Layer Monitoring (Do Not Leave)}

\begin{itemize}
\tightlist
\item[$\square$]
  Watch the first 2--3 layers without interruption
\item[$\square$]
  Bed adhesion is correct --- filament sticks flat, not lifting at corners or curling
\item[$\square$]
  Nozzle temperature is stable at target value
\item[$\square$]
  No unusual sounds (grinding, clicking, popping)
\end{itemize}

\textbf{Regular Checks (Every 15--30 Minutes)}

\begin{itemize}
\tightlist
\item[$\square$]
  Print is building correctly --- no layer shifting, no spaghetti
\item[$\square$]
  Filament is feeding smoothly from the spool
\item[$\square$]
  No grinding or skipping sounds from the extruder
\item[$\square$]
  No warping visible at part edges
\end{itemize}

\textbf{Temperature Stability}

\begin{itemize}
\tightlist
\item[$\square$]
  Heated bed maintains consistent temperature throughout
\item[$\square$]
  Nozzle temperature does not fluctuate more than ±5°C
\item[$\square$]
  No thermal runaway warnings on the printer display
\end{itemize}

\begin{center}\rule{0.5\linewidth}{0.5pt}\end{center}

\subsubsection*{Stop the Print Immediately If}\label{docs__pandoc__latex__src__3dmake_all_lessons.md__stop-the-print-immediately-if}

\begin{itemize}
\tightlist
\item[$\square$]
  The nozzle jams or extrudes unevenly
\item[$\square$]
  Filament stops feeding entirely
\item[$\square$]
  Any burning smell is detected
\item[$\square$]
  Visible layer shifting occurs
\item[$\square$]
  Grinding or skipping sounds from the extruder
\item[$\square$]
  Any unusual chemical or burning odor
\item[$\square$]
  Smoke appears --- \textbf{cut power immediately}
\end{itemize}

\begin{center}\rule{0.5\linewidth}{0.5pt}\end{center}

\subsubsection*{Post-Print}\label{docs__pandoc__latex__src__3dmake_all_lessons.md__post-print}

\textbf{Cool Down}

\begin{itemize}
\tightlist
\item[$\square$]
  Allow the heated bed to cool naturally for 10--15 minutes before attempting print removal
\item[$\square$]
  Keep hands clear of the nozzle until it reads below 50°C
\item[$\square$]
  Confirm the printer is idle and all temperatures are dropping before touching
\end{itemize}

\textbf{Print Removal}

\begin{itemize}
\tightlist
\item[$\square$]
  Use a proper spatula or scraper --- not fingers or improvised tools
\item[$\square$]
  Remove the print only when the bed has fully cooled (PLA releases easily from PEI beds when cool)
\item[$\square$]
  Inspect the print for sharp edges, support stubs, or layer delamination
\item[$\square$]
  Sand or file sharp edges if the part will be handled
\end{itemize}

\textbf{Equipment Cleaning}

\begin{itemize}
\tightlist
\item[$\square$]
  Wipe the nozzle with a brass wire brush when cooled to 50°C (never cold --- residue must be slightly soft)
\item[$\square$]
  Remove any plastic debris and failed print material from the build plate
\item[$\square$]
  Check the extruder for embedded filament residue
\item[$\square$]
  Clear any visible dust from the electronics area
\end{itemize}

\textbf{Workspace Cleanup}

\begin{itemize}
\tightlist
\item[$\square$]
  Return all tools to their proper storage location
\item[$\square$]
  Dispose of support material and failed prints in a bin (not on the floor)
\item[$\square$]
  Store filament in a sealed bag with desiccant
\item[$\square$]
  Confirm the workspace is clear and safe for the next user
\end{itemize}

\begin{center}\rule{0.5\linewidth}{0.5pt}\end{center}

\subsubsection*{Hazard Reference Summary}\label{docs__pandoc__latex__src__3dmake_all_lessons.md__hazard-reference-summary}

{\def\LTcaptype{none} % do not increment counter
\begin{longtable}[]{@{}
  >{\raggedright\arraybackslash}p{(\linewidth - 4\tabcolsep) * \real{0.2841}}
  >{\raggedright\arraybackslash}p{(\linewidth - 4\tabcolsep) * \real{0.2614}}
  >{\raggedright\arraybackslash}p{(\linewidth - 4\tabcolsep) * \real{0.4545}}@{}}
\toprule\noalign{}
\begin{minipage}[b]{\linewidth}\raggedright
Hazard
\end{minipage} & \begin{minipage}[b]{\linewidth}\raggedright
Temperature / Condition
\end{minipage} & \begin{minipage}[b]{\linewidth}\raggedright
Response
\end{minipage} \\
\midrule\noalign{}
\endhead
\bottomrule\noalign{}
\endlastfoot
Hot nozzle & 190--260°C & Do not touch; wait for \textless{} 50°C \\
Hot bed & 50--110°C & Do not touch; allow full cool-down \\
Fumes (ABS/PETG) & At operating temp &
Ventilate; use enclosure with filtration \\
Electrical (power supply) & 120V / 240V mains &
Never open while plugged in \\
Tangled filament & Any & Clear before printing; never force-pull \\
\end{longtable}
}

\subsubsection*{Emergency Actions}\label{docs__pandoc__latex__src__3dmake_all_lessons.md__emergency-actions}

\begin{itemize}
\tightlist
\item
  \textbf{Fire:} Unplug immediately. Use a dry chemical (Class ABC) extinguisher --- NOT water on electrical fires. Evacuate and call emergency services.
\item
  \textbf{Burns:} Rinse with cool (not cold) running water for 15+ minutes. Seek medical attention for any burn larger than a coin.
\item
  \textbf{Fume inhalation:} Move to fresh air immediately. Seek medical attention if symptoms (headache, dizziness, nausea) persist.
\end{itemize}

\emph{Last Reviewed:} \_\_\_\_\_\_\_\_\_\_ \emph{Reviewed By:} \_\_\_\_\_\_\_\_\_\_

\begin{center}\rule{0.5\linewidth}{0.5pt}\end{center}

\subsection*{Appendix B: Filament Comparison Table}\label{docs__pandoc__latex__src__3dmake_all_lessons.md__appendix-b-filament-comparison-table}

\emph{Quick reference for choosing the right filament. All temperature values are typical ranges --- always check the label on your specific spool. The spool always wins over this table.}

\begin{center}\rule{0.5\linewidth}{0.5pt}\end{center}

\subsubsection*{Side-by-Side Comparison}\label{docs__pandoc__latex__src__3dmake_all_lessons.md__side-by-side-comparison}

{\def\LTcaptype{none} % do not increment counter
\begin{longtable}[]{@{}
  >{\raggedright\arraybackslash}p{(\linewidth - 8\tabcolsep) * \real{0.2830}}
  >{\raggedright\arraybackslash}p{(\linewidth - 8\tabcolsep) * \real{0.1226}}
  >{\raggedright\arraybackslash}p{(\linewidth - 8\tabcolsep) * \real{0.1981}}
  >{\raggedright\arraybackslash}p{(\linewidth - 8\tabcolsep) * \real{0.1604}}
  >{\raggedright\arraybackslash}p{(\linewidth - 8\tabcolsep) * \real{0.2358}}@{}}
\toprule\noalign{}
\begin{minipage}[b]{\linewidth}\raggedright
Property
\end{minipage} & \begin{minipage}[b]{\linewidth}\raggedright
PLA
\end{minipage} & \begin{minipage}[b]{\linewidth}\raggedright
PETG
\end{minipage} & \begin{minipage}[b]{\linewidth}\raggedright
TPU (Flex)
\end{minipage} & \begin{minipage}[b]{\linewidth}\raggedright
ABS
\end{minipage} \\
\midrule\noalign{}
\endhead
\bottomrule\noalign{}
\endlastfoot
\textbf{Difficulty} & Easiest & Moderate & Hard & Very Hard \\
\textbf{Nozzle temp} & 190--220°C & 230--250°C & 220--240°C &
230--250°C \\
\textbf{Bed temp} & 50--60°C & 70--85°C & 30--60°C & 90--110°C \\
\textbf{Enclosure needed?} & No & No & No & Yes \\
\textbf{Rigidity} & High & Medium & None --- flexible & High \\
\textbf{Impact resistance} & Low (brittle) & Medium & Very high &
Medium \\
\textbf{Heat resistance} & Low (\textasciitilde{}60°C) &
Medium (\textasciitilde{}80°C) & Medium &
High (\textasciitilde{}100°C) \\
\textbf{Flexibility} & None & Slight & Rubber-like & None \\
\textbf{Moisture sensitivity} & Low & Medium & High & Medium \\
\textbf{Fume concern} & Lowest & Low & Moderate & High \\
\textbf{Bed adhesion} & Easy & Use glue stick on PEI & Easy &
Requires enclosure + glue \\
\textbf{Print speed} & Normal & Normal & Slow (20--30 mm/s) & Normal \\
\textbf{Stringing tendency} & Low & Medium--High & High & Low \\
\textbf{Approximate cost} & \$15--25/kg & \$15--30/kg & \$20--40/kg &
\$15--25/kg \\
\textbf{Recommended for beginners?} & ✅ Yes & After PLA &
After experience & ❌ No \\
\end{longtable}
}

\begin{center}\rule{0.5\linewidth}{0.5pt}\end{center}

\subsubsection*{When to Use Each Material}\label{docs__pandoc__latex__src__3dmake_all_lessons.md__when-to-use-each-material}

\textbf{PLA --- Use for:}

\begin{itemize}
\tightlist
\item
  Prototypes and test prints
\item
  Classroom projects
\item
  Decorative objects and tactile models
\item
  Anything that won\textquotesingle t be exposed to heat, moisture, or heavy mechanical stress
\end{itemize}

\textbf{PLA --- Avoid for:}

\begin{itemize}
\tightlist
\item
  Objects left in a hot car or direct sun (softens above \textasciitilde{}60°C)
\item
  Parts that need to flex or bend without cracking
\item
  High-impact applications
\end{itemize}

\textbf{PETG --- Use for:}

\begin{itemize}
\tightlist
\item
  Functional parts that need to be tougher than PLA
\item
  Parts exposed to mild heat or moisture
\item
  Mechanical components: brackets, clips, mounts
\item
  Food-contact applications (check your specific brand for food-safe certification)
\end{itemize}

\textbf{PETG --- Avoid for:}

\begin{itemize}
\tightlist
\item
  Very fine surface detail (it strings more than PLA)
\item
  Projects where you need the easiest possible first-time print
\end{itemize}

\textbf{TPU / Flexible --- Use for:}

\begin{itemize}
\tightlist
\item
  Wearable objects (wristbands, phone cases, watch straps)
\item
  Bumpers and shock absorbers
\item
  Grips and handles
\item
  Objects that must deform under pressure and return to shape
\end{itemize}

\textbf{TPU --- Avoid for:}

\begin{itemize}
\tightlist
\item
  Fine surface detail
\item
  Printers with Bowden extruders (the flexible filament buckles in the tube --- direct drive only)
\item
  Your first few prints while still learning
\end{itemize}

\textbf{ABS --- Use for:}

\begin{itemize}
\tightlist
\item
  High-heat environments (car interiors, near motors)
\item
  Parts requiring post-processing (ABS sands, bonds with acetone vapor, and paints easily)
\item
  Professional or industrial contexts with proper dedicated ventilation
\end{itemize}

\textbf{ABS --- Avoid for:}

\begin{itemize}
\tightlist
\item
  Any classroom without a dedicated sealed fume enclosure and active filtration
\item
  Beginners
\item
  Any print where warping at the edges would be a problem and you have no enclosure
\end{itemize}

\begin{center}\rule{0.5\linewidth}{0.5pt}\end{center}

\subsubsection*{Slicer Profile Quick Reference}\label{docs__pandoc__latex__src__3dmake_all_lessons.md__slicer-profile-quick-reference}

{\def\LTcaptype{none} % do not increment counter
\begin{longtable}[]{@{}
  >{\raggedright\arraybackslash}p{(\linewidth - 4\tabcolsep) * \real{0.1026}}
  >{\raggedright\arraybackslash}p{(\linewidth - 4\tabcolsep) * \real{0.2436}}
  >{\raggedright\arraybackslash}p{(\linewidth - 4\tabcolsep) * \real{0.6538}}@{}}
\toprule\noalign{}
\begin{minipage}[b]{\linewidth}\raggedright
Filament
\end{minipage} & \begin{minipage}[b]{\linewidth}\raggedright
PrusaSlicer Profile
\end{minipage} & \begin{minipage}[b]{\linewidth}\raggedright
Key Changes from PLA Default
\end{minipage} \\
\midrule\noalign{}
\endhead
\bottomrule\noalign{}
\endlastfoot
PLA & Generic PLA & Baseline --- no changes needed \\
PETG & Generic PETG &
Higher temps; slower cooling; glue stick on PEI bed \\
TPU & Generic Flex & Reduce speed to 20--30 mm/s; reduce retraction \\
ABS & Generic ABS & Higher temps; disable cooling fan; use enclosure \\
\end{longtable}
}

\begin{center}\rule{0.5\linewidth}{0.5pt}\end{center}

\subsubsection*{Filament Storage}\label{docs__pandoc__latex__src__3dmake_all_lessons.md__filament-storage}

All filament absorbs moisture from the air over time. Wet filament causes: popping or crackling sounds during printing, bubbles in extruded plastic, excessive stringing, and weak or brittle finished parts.

{\def\LTcaptype{none} % do not increment counter
\begin{longtable}[]{@{}
  >{\raggedright\arraybackslash}p{(\linewidth - 4\tabcolsep) * \real{0.0777}}
  >{\raggedright\arraybackslash}p{(\linewidth - 4\tabcolsep) * \real{0.1942}}
  >{\raggedright\arraybackslash}p{(\linewidth - 4\tabcolsep) * \real{0.7282}}@{}}
\toprule\noalign{}
\begin{minipage}[b]{\linewidth}\raggedright
Filament
\end{minipage} & \begin{minipage}[b]{\linewidth}\raggedright
Moisture Sensitivity
\end{minipage} & \begin{minipage}[b]{\linewidth}\raggedright
Storage Recommendation
\end{minipage} \\
\midrule\noalign{}
\endhead
\bottomrule\noalign{}
\endlastfoot
PLA & Low & Sealed bag with desiccant when not in use \\
PETG & Medium &
Sealed bag with desiccant; dry before use if stored open for weeks \\
TPU & High &
Always store sealed; dry in oven at 65°C for 4--6 hours if moisture-absorbed \\
ABS & Medium & Sealed bag with desiccant \\
\end{longtable}
}

\emph{Sources: All3DP Filament Types Guide; MatterHackers Filament Compare; Prusa Research Materials Guide}

\begin{center}\rule{0.5\linewidth}{0.5pt}\end{center}

\subsection*{Appendix C: Material Properties Quick Reference}\label{docs__pandoc__latex__src__3dmake_all_lessons.md__appendix-c-material-properties-quick-reference}

\emph{Essential material data for FDM/FFF printing. Use alongside the Filament Comparison Table for material selection decisions.}

\begin{center}\rule{0.5\linewidth}{0.5pt}\end{center}

\subsubsection*{Properties at a Glance}\label{docs__pandoc__latex__src__3dmake_all_lessons.md__properties-at-a-glance}

{\def\LTcaptype{none} % do not increment counter
\begin{longtable}[]{@{}
  >{\raggedright\arraybackslash}p{(\linewidth - 8\tabcolsep) * \real{0.2879}}
  >{\raggedright\arraybackslash}p{(\linewidth - 8\tabcolsep) * \real{0.1515}}
  >{\raggedright\arraybackslash}p{(\linewidth - 8\tabcolsep) * \real{0.1515}}
  >{\raggedright\arraybackslash}p{(\linewidth - 8\tabcolsep) * \real{0.1970}}
  >{\raggedright\arraybackslash}p{(\linewidth - 8\tabcolsep) * \real{0.2121}}@{}}
\toprule\noalign{}
\begin{minipage}[b]{\linewidth}\raggedright
Property
\end{minipage} & \begin{minipage}[b]{\linewidth}\raggedright
PLA
\end{minipage} & \begin{minipage}[b]{\linewidth}\raggedright
PETG
\end{minipage} & \begin{minipage}[b]{\linewidth}\raggedright
ABS
\end{minipage} & \begin{minipage}[b]{\linewidth}\raggedright
TPU
\end{minipage} \\
\midrule\noalign{}
\endhead
\bottomrule\noalign{}
\endlastfoot
Nozzle temp (°C) & 200--220 & 235--250 & 240--260 & 220--240 \\
Bed temp (°C) & 20--60 & 70--90 & 100--110 & 20--60 \\
Strength & Moderate & High & Very High & Low (flexible) \\
Flexibility & None & Slight & Slight & High \\
Ease of printing & Very Easy & Easy & Hard & Moderate \\
Cost (\$/kg) & \$15--25 & \$20--30 & \$18--28 & \$30--50 \\
Print speed & 40--60 mm/s & 30--50 mm/s & 20--40 mm/s & 10--30 mm/s \\
Bed adhesion & Easy & Moderate & Requires prep & Easy \\
UV resistance & Low & Moderate & High & Low \\
Chemical resistance & No & Moderate & Yes & No \\
\end{longtable}
}

\begin{center}\rule{0.5\linewidth}{0.5pt}\end{center}

\subsubsection*{Heat Resistance After Printing}\label{docs__pandoc__latex__src__3dmake_all_lessons.md__heat-resistance-after-printing}

This is one of the most common reasons to choose a material other than PLA.

{\def\LTcaptype{none} % do not increment counter
\begin{longtable}[]{@{}
  >{\raggedright\arraybackslash}p{(\linewidth - 6\tabcolsep) * \real{0.0870}}
  >{\centering\arraybackslash}p{(\linewidth - 6\tabcolsep) * \real{0.1848}}
  >{\centering\arraybackslash}p{(\linewidth - 6\tabcolsep) * \real{0.1087}}
  >{\raggedright\arraybackslash}p{(\linewidth - 6\tabcolsep) * \real{0.6196}}@{}}
\toprule\noalign{}
\begin{minipage}[b]{\linewidth}\raggedright
Material
\end{minipage} & \begin{minipage}[b]{\linewidth}\centering
Safe to use up to
\end{minipage} & \begin{minipage}[b]{\linewidth}\centering
Softens at
\end{minipage} & \begin{minipage}[b]{\linewidth}\raggedright
Notes
\end{minipage} \\
\midrule\noalign{}
\endhead
\bottomrule\noalign{}
\endlastfoot
PLA & \textasciitilde{}50--60°C & \textasciitilde{}60--70°C &
Not suitable for car interiors or items near heat sources \\
PETG & \textasciitilde{}80--100°C & \textasciitilde{}100--110°C &
Handles warm environments well \\
ABS & \textasciitilde{}90--110°C & \textasciitilde{}110--120°C &
Good heat resistance; most demanding to print \\
TPU & \textasciitilde{}60--80°C & \textasciitilde{}80--100°C &
Limited heat resistance despite flex durability \\
\end{longtable}
}

\begin{center}\rule{0.5\linewidth}{0.5pt}\end{center}

\subsubsection*{Chemical Resistance}\label{docs__pandoc__latex__src__3dmake_all_lessons.md__chemical-resistance}

{\def\LTcaptype{none} % do not increment counter
\begin{longtable}[]{@{}lccccc@{}}
\toprule\noalign{}
Material & Water & Alcohol & Acetone & Oils & Acids \\
\midrule\noalign{}
\endhead
\bottomrule\noalign{}
\endlastfoot
PLA & Low & Low & Dissolves & Low & Low \\
PETG & High & Moderate & Moderate & High & Moderate \\
ABS & High & Moderate & Dissolves & High & Moderate \\
TPU & Moderate & Low & Low & Moderate & Low \\
\end{longtable}
}

\textbf{Acetone note:} Can be used to smooth ABS surfaces (vapor or liquid application). Also dissolves PLA --- keep acetone away from PLA prints and filament spools.

\begin{center}\rule{0.5\linewidth}{0.5pt}\end{center}

\subsubsection*{Strength After Printing --- Curing Time}\label{docs__pandoc__latex__src__3dmake_all_lessons.md__strength-after-printing--curing-time}

Freshly printed parts are not at full strength.

{\def\LTcaptype{none} % do not increment counter
\begin{longtable}[]{@{}ll@{}}
\toprule\noalign{}
Time After Printing & Approximate Strength \\
\midrule\noalign{}
\endhead
\bottomrule\noalign{}
\endlastfoot
0--24 hours & 70--80\% of final strength \\
24--48 hours & 90--95\% of final strength \\
48--72 hours & \textasciitilde{}99\% of final strength \\
7+ days & Maximum strength \\
\end{longtable}
}

Do not stress-test functional parts immediately after printing. Allow at least 24 hours before load testing.

\begin{center}\rule{0.5\linewidth}{0.5pt}\end{center}

\subsubsection*{Storage Conditions}\label{docs__pandoc__latex__src__3dmake_all_lessons.md__storage-conditions}

{\def\LTcaptype{none} % do not increment counter
\begin{longtable}[]{@{}ll@{}}
\toprule\noalign{}
Factor & Recommended Range \\
\midrule\noalign{}
\endhead
\bottomrule\noalign{}
\endlastfoot
Storage temperature & 15--25°C \\
Relative humidity & 30--50\% \\
Light exposure & Away from direct sunlight \\
Container & Sealed with desiccant packets \\
\end{longtable}
}

\begin{center}\rule{0.5\linewidth}{0.5pt}\end{center}

\subsubsection*{Quick Material Selection Guide}\label{docs__pandoc__latex__src__3dmake_all_lessons.md__quick-material-selection-guide}

{\def\LTcaptype{none} % do not increment counter
\begin{longtable}[]{@{}ll@{}}
\toprule\noalign{}
Question & Answer \\
\midrule\noalign{}
\endhead
\bottomrule\noalign{}
\endlastfoot
First time printing? & PLA \\
Need durability? & PETG or ABS \\
Need flexibility? & TPU \\
Need high strength? & ABS or PETG \\
Quick test / prototype? & PLA \\
Outdoor use? & PETG or ABS \\
Food contact? & Food-grade certified PETG (verify brand) \\
Heat resistant? & ABS \\
Easiest possible print? & PLA \\
\end{longtable}
}

\begin{center}\rule{0.5\linewidth}{0.5pt}\end{center}

\subsection*{Appendix D: Printer Maintenance Log}\label{docs__pandoc__latex__src__3dmake_all_lessons.md__appendix-d-printer-maintenance-log}

\emph{Copy this log and keep a printed or digital copy with your printer. Fill in the Printer Information section first, then log every maintenance event and issue as it occurs. A complete log is your best diagnostic tool when problems arise.}

\begin{center}\rule{0.5\linewidth}{0.5pt}\end{center}

\subsubsection*{Printer Information}\label{docs__pandoc__latex__src__3dmake_all_lessons.md__printer-information}

{\def\LTcaptype{none} % do not increment counter
\begin{longtable}[]{@{}ll@{}}
\toprule\noalign{}
Field & Value \\
\midrule\noalign{}
\endhead
\bottomrule\noalign{}
\endlastfoot
Printer model & \\
Serial number & \\
Purchase date & \\
Firmware version & \\
Last service date & \\
Warranty expiry & \\
Manufacturer support contact & \\
\end{longtable}
}

\begin{center}\rule{0.5\linewidth}{0.5pt}\end{center}

\subsubsection*{Maintenance Schedule}\label{docs__pandoc__latex__src__3dmake_all_lessons.md__maintenance-schedule}

\textbf{Before each use (Daily)}

\begin{itemize}
\tightlist
\item[$\square$]
  Visual inspection --- no visible damage, cables intact
\item[$\square$]
  Clean nozzle tip if residue is present
\item[$\square$]
  Verify build plate is level (re-level if printer was moved)
\item[$\square$]
  Confirm filament spool rotates freely with no tangles
\end{itemize}

\textbf{Weekly}

\begin{itemize}
\tightlist
\item[$\square$]
  Clean extruder drive gear (compressed air or brush out ground plastic)
\item[$\square$]
  Inspect and wipe Z-axis rails and lead screw
\item[$\square$]
  Check all cable connections at the control board
\item[$\square$]
  Test emergency stop / power cut
\end{itemize}

\textbf{Monthly}

\begin{itemize}
\tightlist
\item[$\square$]
  Full build plate leveling procedure (manual or automatic)
\item[$\square$]
  Clean interior of print chamber
\item[$\square$]
  Inspect heating elements and thermistor connections
\item[$\square$]
  Run a temperature calibration test print
\end{itemize}

\textbf{Quarterly (every 3 months)}

\begin{itemize}
\tightlist
\item[$\square$]
  Replace nozzle if inner bore shows wear or repeat clogging
\item[$\square$]
  Full mechanical inspection --- belt tension, eccentric nuts, frame screws
\item[$\square$]
  Check firmware for updates
\item[$\square$]
  Run a full calibration test print and document results
\end{itemize}

\begin{center}\rule{0.5\linewidth}{0.5pt}\end{center}

\subsubsection*{Maintenance Log}\label{docs__pandoc__latex__src__3dmake_all_lessons.md__maintenance-log}

{\def\LTcaptype{none} % do not increment counter
\begin{longtable}[]{@{}
  >{\raggedright\arraybackslash}p{(\linewidth - 12\tabcolsep) * \real{0.0741}}
  >{\raggedright\arraybackslash}p{(\linewidth - 12\tabcolsep) * \real{0.0741}}
  >{\raggedright\arraybackslash}p{(\linewidth - 12\tabcolsep) * \real{0.2037}}
  >{\raggedright\arraybackslash}p{(\linewidth - 12\tabcolsep) * \real{0.1852}}
  >{\raggedright\arraybackslash}p{(\linewidth - 12\tabcolsep) * \real{0.2222}}
  >{\raggedright\arraybackslash}p{(\linewidth - 12\tabcolsep) * \real{0.1852}}
  >{\raggedright\arraybackslash}p{(\linewidth - 12\tabcolsep) * \real{0.0556}}@{}}
\toprule\noalign{}
\begin{minipage}[b]{\linewidth}\raggedright
Date
\end{minipage} & \begin{minipage}[b]{\linewidth}\raggedright
Type
\end{minipage} & \begin{minipage}[b]{\linewidth}\raggedright
Description
\end{minipage} & \begin{minipage}[b]{\linewidth}\raggedright
Time Spent
\end{minipage} & \begin{minipage}[b]{\linewidth}\raggedright
Issues Found
\end{minipage} & \begin{minipage}[b]{\linewidth}\raggedright
Resolution
\end{minipage} & \begin{minipage}[b]{\linewidth}\raggedright
By
\end{minipage} \\
\midrule\noalign{}
\endhead
\bottomrule\noalign{}
\endlastfoot
& & & & & & \\
& & & & & & \\
& & & & & & \\
& & & & & & \\
& & & & & & \\
& & & & & & \\
& & & & & & \\
& & & & & & \\
\end{longtable}
}

\emph{Types: Routine / Repair / Calibration / Cleaning / Part replacement}

\begin{center}\rule{0.5\linewidth}{0.5pt}\end{center}

\subsubsection*{Issue Tracking Log}\label{docs__pandoc__latex__src__3dmake_all_lessons.md__issue-tracking-log}

{\def\LTcaptype{none} % do not increment counter
\begin{longtable}[]{@{}llllll@{}}
\toprule\noalign{}
Date & Symptom & Diagnosis & Action Taken & Status & Notes \\
\midrule\noalign{}
\endhead
\bottomrule\noalign{}
\endlastfoot
& & & & & \\
& & & & & \\
& & & & & \\
& & & & & \\
& & & & & \\
\end{longtable}
}

\emph{Statuses: Open / In progress / Resolved / Monitoring}

\begin{center}\rule{0.5\linewidth}{0.5pt}\end{center}

\subsubsection*{Parts Replacement Record}\label{docs__pandoc__latex__src__3dmake_all_lessons.md__parts-replacement-record}

{\def\LTcaptype{none} % do not increment counter
\begin{longtable}[]{@{}llllll@{}}
\toprule\noalign{}
Date & Part Replaced & Reason & Supplier & Cost & Notes \\
\midrule\noalign{}
\endhead
\bottomrule\noalign{}
\endlastfoot
& & & & & \\
& & & & & \\
& & & & & \\
\end{longtable}
}

\begin{center}\rule{0.5\linewidth}{0.5pt}\end{center}

\subsubsection*{Filament Compatibility Notes}\label{docs__pandoc__latex__src__3dmake_all_lessons.md__filament-compatibility-notes}

Record your tested settings for each material on this specific printer. Published temperature ranges are starting points --- your printer may run hotter or cooler than its display indicates.

{\def\LTcaptype{none} % do not increment counter
\begin{longtable}[]{@{}
  >{\raggedright\arraybackslash}p{(\linewidth - 12\tabcolsep) * \real{0.1481}}
  >{\raggedright\arraybackslash}p{(\linewidth - 12\tabcolsep) * \real{0.0926}}
  >{\centering\arraybackslash}p{(\linewidth - 12\tabcolsep) * \real{0.2037}}
  >{\centering\arraybackslash}p{(\linewidth - 12\tabcolsep) * \real{0.1481}}
  >{\centering\arraybackslash}p{(\linewidth - 12\tabcolsep) * \real{0.0926}}
  >{\centering\arraybackslash}p{(\linewidth - 12\tabcolsep) * \real{0.2222}}
  >{\raggedright\arraybackslash}p{(\linewidth - 12\tabcolsep) * \real{0.0926}}@{}}
\toprule\noalign{}
\begin{minipage}[b]{\linewidth}\raggedright
Filament
\end{minipage} & \begin{minipage}[b]{\linewidth}\raggedright
Brand
\end{minipage} & \begin{minipage}[b]{\linewidth}\centering
Nozzle Temp
\end{minipage} & \begin{minipage}[b]{\linewidth}\centering
Bed Temp
\end{minipage} & \begin{minipage}[b]{\linewidth}\centering
Speed
\end{minipage} & \begin{minipage}[b]{\linewidth}\centering
Success Rate
\end{minipage} & \begin{minipage}[b]{\linewidth}\raggedright
Notes
\end{minipage} \\
\midrule\noalign{}
\endhead
\bottomrule\noalign{}
\endlastfoot
PLA & & & & & & \\
PETG & & & & & & \\
ABS & & & & & & \\
TPU & & & & & & \\
Other & & & & & & \\
\end{longtable}
}

\begin{center}\rule{0.5\linewidth}{0.5pt}\end{center}

\subsubsection*{Troubleshooting Quick Reference}\label{docs__pandoc__latex__src__3dmake_all_lessons.md__troubleshooting-quick-reference}

\textbf{Nozzle Clog}

\begin{itemize}
\tightlist
\item
  First response: Cold pull (heat to printing temp, push filament through manually, cool to 90°C, pull sharply)
\item
  Second response: Acupuncture needle or nozzle cleaning filament
\item
  Last resort: Replace nozzle
\end{itemize}

\textbf{Bed Adhesion Problems}

\begin{itemize}
\tightlist
\item
  Check: Bed level, nozzle height from bed, bed surface condition (clean? worn?)
\item
  Fix by material: PLA → clean PEI with IPA; PETG → apply glue stick; ABS → glue stick + enclosure
\end{itemize}

\textbf{Layer Shifting}

\begin{itemize}
\tightlist
\item
  Common causes: Belt too loose, motor current too low, print speed too high, X/Y axis obstruction
\item
  Check belt tension first --- it\textquotesingle s the most common cause
\end{itemize}

\textbf{Extrusion Issues (Under-extrusion)}

\begin{itemize}
\tightlist
\item
  Common causes: Partial clog, worn drive gear, incorrect extruder steps/mm, wet filament
\item
  Run a flow rate test: extrude 100 mm and measure actual output
\end{itemize}

\begin{center}\rule{0.5\linewidth}{0.5pt}\end{center}

\subsubsection*{Monthly Performance Summary}\label{docs__pandoc__latex__src__3dmake_all_lessons.md__monthly-performance-summary}

{\def\LTcaptype{none} % do not increment counter
\begin{longtable}[]{@{}lcll@{}}
\toprule\noalign{}
Month & Print Success Rate & Common Issues & Overall Assessment \\
\midrule\noalign{}
\endhead
\bottomrule\noalign{}
\endlastfoot
& & & \\
& & & \\
& & & \\
& & & \\
\end{longtable}
}

\emph{Last log entry:} \_\_\_\_\_\_\_\_\_\_ \emph{Logged by:} \_\_\_\_\_\_\_\_\_\_
-e

\begin{center}\rule{0.5\linewidth}{0.5pt}\end{center}

\section*{Lesson 6 --- Practical 3dm Commands and Text Embossing}\label{docs__pandoc__latex__src__3dmake_all_lessons.md__lesson-6--practical-3dm-commands-and-text-embossing}

\textbf{Estimated Time:} 90--120 minutes\\
\textbf{Course:} 3dMake Certification --- High School Distance Learning Track

\begin{center}\rule{0.5\linewidth}{0.5pt}\end{center}

\subsection*{Before You Start}\label{docs__pandoc__latex__src__3dmake_all_lessons.md__before-you-start-4}

You\textquotesingle ve built solid 3D geometry, organized it into modules, learned to verify it rigorously, and understood the physical realities of printing it safely. Now we\textquotesingle re going to add two powerful new capabilities to your toolkit.

The first is \textbf{text}. Adding text to 3D parts is more useful than it might sound at first. Engineers put part numbers on manufactured components so they can be identified after assembly. Makers label storage bins so they can find what they need. Designers emboss names and logos on products. Accessibility designers create braille labels. OpenSCAD\textquotesingle s \texttt{text()} function makes all of this parametric --- you can generate a hundred labeled bins automatically, make the text content a parameter, and produce variants for different languages or quantities without touching the geometry code at all.

The second is the \textbf{full \texttt{3dm} command suite}. You\textquotesingle ve been using \texttt{3dm\ build} and \texttt{3dm\ info}. There are three more commands --- \texttt{preview}, \texttt{orient}, and \texttt{slice} --- that round out the workflow from design to print.

\begin{center}\rule{0.5\linewidth}{0.5pt}\end{center}

\subsection*{Learning Objectives}\label{docs__pandoc__latex__src__3dmake_all_lessons.md__learning-objectives-9}

By the end of this lesson you will be able to:

\begin{itemize}
\tightlist
\item
  Use \texttt{text()} to add embossed and debossed labels to 3D models
\item
  Find available fonts on your system and write portable font code
\item
  Size text to fit a specific container using the cap-height model
\item
  Use the full \texttt{3dm} command suite: \texttt{build}, \texttt{info}, \texttt{preview}, \texttt{orient}, \texttt{slice}
\item
  Use OpenSCAD string functions: \texttt{str()}, \texttt{len()}, \texttt{search()}, \texttt{substr()}
\item
  Apply \texttt{let()} for scoped variable declarations
\item
  Automate multi-variant text label generation in bash, PowerShell, and CMD
\end{itemize}

\begin{center}\rule{0.5\linewidth}{0.5pt}\end{center}

\subsection*{\texorpdfstring{Concept 1: The Full \texttt{3dm} Command Suite}{Concept 1: The Full 3dm Command Suite}}\label{docs__pandoc__latex__src__3dmake_all_lessons.md__concept-1-the-full-3dm-command-suite}

\subsubsection*{Five Commands, One Workflow}\label{docs__pandoc__latex__src__3dmake_all_lessons.md__five-commands-one-workflow}

You\textquotesingle ve been using two commands. Here\textquotesingle s the complete set:

{\def\LTcaptype{none} % do not increment counter
\begin{longtable}[]{@{}
  >{\raggedright\arraybackslash}p{(\linewidth - 4\tabcolsep) * \real{0.1605}}
  >{\raggedright\arraybackslash}p{(\linewidth - 4\tabcolsep) * \real{0.5062}}
  >{\raggedright\arraybackslash}p{(\linewidth - 4\tabcolsep) * \real{0.3333}}@{}}
\toprule\noalign{}
\begin{minipage}[b]{\linewidth}\raggedright
Command
\end{minipage} & \begin{minipage}[b]{\linewidth}\raggedright
What It Does
\end{minipage} & \begin{minipage}[b]{\linewidth}\raggedright
Output
\end{minipage} \\
\midrule\noalign{}
\endhead
\bottomrule\noalign{}
\endlastfoot
\texttt{3dm\ build} &
Compiles \texttt{.scad} source → \texttt{.stl} geometry &
\texttt{build/main.stl} \\
\texttt{3dm\ info} & Reports dimensions and AI description &
Console text \\
\texttt{3dm\ preview} & Renders model to an image file &
\texttt{build/preview.png} \\
\texttt{3dm\ orient} & AI-suggested print orientation &
Console text recommendation \\
\texttt{3dm\ slice} & Calls your slicer on the current STL &
G-code file \\
\end{longtable}
}

\begin{lstlisting}[style=Alabaster, language=bash]
# BUILD VARIANTS

3dm build                          # standard build

3dm build --clean                  # delete build/ folder first, then rebuild

3dm build --watch                  # auto-rebuild when src/ files change


# INFORMATION AND DESCRIPTION

3dm info                           # bounding box, volume, triangle count, AI description

3dm info --file path/to/file.stl   # query a specific STL (not just the current build)

3dm info --view front              # generate description from the front view

3dm info --view top,front,isometric  # generate from multiple views


# ORIENTATION, PREVIEW, SLICING

3dm preview                        # render model to build/preview.png

3dm orient                         # AI recommendation for best print orientation

3dm slice                          # call slicer on build/main.stl


\end{lstlisting}

\subsubsection*{\texorpdfstring{\texttt{3dm\ orient} --- What It Does and How to Use It}{3dm orient --- What It Does and How to Use It}}\label{docs__pandoc__latex__src__3dmake_all_lessons.md__3dm-orient--what-it-does-and-how-to-use-it}

\texttt{3dm\ orient} uses AI to analyze your model\textquotesingle s geometry and recommend which face should be placed down on the build plate for printing. This recommendation considers two things: minimizing the need for supports (overhangs), and maximizing the quality of visible surfaces (which sit better when they\textquotesingle re on the top or sides, not facing down toward the bed).

Use \texttt{3dm\ orient} as a starting suggestion, but verify it with your own judgment. For complex parts, you may have competing priorities --- perhaps the AI recommends an orientation that minimizes supports but puts a cosmetically important surface face-down against the bed, where it will look rough. Your own analysis matters.

\begin{center}\rule{0.5\linewidth}{0.5pt}\end{center}

\subsection*{Concept 2: Text on 3D Parts}\label{docs__pandoc__latex__src__3dmake_all_lessons.md__concept-2-text-on-3d-parts}

\subsubsection*{Why 3D Text Is Useful}\label{docs__pandoc__latex__src__3dmake_all_lessons.md__why-3d-text-is-useful}

Before diving into syntax, it\textquotesingle s worth thinking about why you\textquotesingle d want text on a 3D print:

\textbf{Identification} --- part numbers, serial numbers, batch codes. Manufacturing parts need to be traceable. A part with its number embossed on it can never be confused with another part.

\textbf{Directional labels} --- "FRONT," "TOP," "THIS SIDE UP." Assembly instructions encoded directly in the part.

\textbf{Storage and organization} --- labels for bins, drawers, shelves. Parametric text means you can batch-generate 30 labeled bins automatically.

\textbf{Personalization} --- names, logos, messages on gifts, art, and personal projects.

\textbf{Accessibility} --- braille labels for visually impaired users. OpenSCAD can generate braille if you encode the dot patterns.

\begin{center}\rule{0.5\linewidth}{0.5pt}\end{center}

\subsection*{\texorpdfstring{Step 1 --- The \texttt{text()} Function}{Step 1 --- The text() Function}}\label{docs__pandoc__latex__src__3dmake_all_lessons.md__step-1--the-text-function}

\subsubsection*{How Text Works in OpenSCAD}\label{docs__pandoc__latex__src__3dmake_all_lessons.md__how-text-works-in-openscad}

\texttt{text()} creates a \textbf{2D} shape from a text string. By itself, it\textquotesingle s flat --- it has no height. You add height by wrapping it in \texttt{linear\_extrude()}, which pushes the 2D letterforms straight up along the Z axis.

\textbf{Embossed text} is raised above the surface --- the letters stand up above the base. This is more visible and readable, especially at small sizes, because there\textquotesingle s contrast between the raised letters and the surrounding surface.

\textbf{Debossed text} (also called engraved or incised) is carved into the surface --- the letters are sunken below the surrounding surface. This reads well when the sunken areas catch shadow, or when filled with paint. It also doesn\textquotesingle t add any height to the overall part dimensions.

\begin{lstlisting}[style=Alabaster, language=openscad]
// Full text() parameter reference

linear_extrude(height=1.5)

  text(

    text      = "3dMake",           // the string to display

    size      = 10,                  // cap height in mm (capital letter height)

    font      = "Liberation Sans",   // font family name (more on this in Step 2)

    halign    = "center",            // horizontal alignment: "left", "center", "right"

    valign    = "baseline",          // vertical alignment: "top", "center", "baseline", "bottom"

    spacing   = 1.0,                 // letter spacing multiplier (1.0 = normal)

    direction = "ltr",               // text direction: "ltr" (left-to-right) or "rtl"

    language  = "en",

    script    = "latin"

  );


// EMBOSSED: raised letters sitting on top of the base

cube([80, 15, 5]);

translate([40, 7.5, 5])             // sit text on top face (Z=5)

  linear_extrude(height=1.5)        // letters are 1.5mm tall

    text("FRONT", size=6, font="Liberation Sans:style=Bold",

         halign="center", valign="center");


// DEBOSSED (ENGRAVED): letters carved into the surface

difference() {

  translate([0, 20, 0]) cube([80, 15, 5]);

  translate([40, 27.5, 5 - 0.8])   // start 0.8mm below the top face

    linear_extrude(height=1)        // 1mm cut: 0.8 into + 0.2 through (clean cut)

      text("BACK", size=6, font="Liberation Sans:style=Bold",

           halign="center", valign="center");

}


\end{lstlisting}

\textbf{Understanding the debossed depth calculation:} If the plate is 5 mm tall and you want 0.8 mm deep engraving, start the extrusion at \texttt{z\ =\ 5\ -\ 0.8\ =\ 4.2}. Extrude 1.0 mm (so it reaches z=5.2 --- 0.2 mm past the top). This uses the 0.001 rule logic from Lesson 2: the cutter goes slightly beyond the face for a clean cut. The visible result is letters engraved 0.8 mm deep.

\begin{center}\rule{0.5\linewidth}{0.5pt}\end{center}

\subsection*{Step 2 --- Finding and Using Fonts}\label{docs__pandoc__latex__src__3dmake_all_lessons.md__step-2--finding-and-using-fonts}

\subsubsection*{Why Font Availability Is a Portability Problem}\label{docs__pandoc__latex__src__3dmake_all_lessons.md__why-font-availability-is-a-portability-problem}

OpenSCAD uses fonts installed on your operating system. Different computers have different fonts. If you write code that uses "Arial" and another student or instructor tries to render it on a Linux machine that doesn\textquotesingle t have Arial, OpenSCAD will fall back to a substitute font --- and the text proportions may change, potentially causing the text to overflow its container or look wrong.

The solution has two parts: check which fonts are available on your machine before writing your code, and define your font choice as a single parameter at the top of the file so it\textquotesingle s easy to change.

\subsubsection*{Finding Available Fonts}\label{docs__pandoc__latex__src__3dmake_all_lessons.md__finding-available-fonts}

\textbf{Linux (bash):}

\begin{lstlisting}[style=Alabaster, language=bash]
# List all available fonts

fc-list | sort | less


# Find fonts with "liberation" in the name

fc-list | grep -i "liberation"


# Get the exact format needed for OpenSCAD's font= parameter

fc-list --format="%{family}:style=%{style}\n" | grep -i liberation | sort | head -20

# Output examples:

# Liberation Sans:style=Bold

# Liberation Sans:style=Regular

# Liberation Mono:style=Bold


\end{lstlisting}

\textbf{Windows (PowerShell):}

\begin{lstlisting}[style=Alabaster, language=powershell]
# List all installed fonts

[System.Reflection.Assembly]::LoadWithPartialName("System.Drawing") | Out-Null

$fonts = (New-Object System.Drawing.Text.InstalledFontCollection).Families

$fonts | Select-Object Name | Sort-Object Name


# Search for a specific font

$fonts | Where-Object { $_.Name -like "*Arial*" } | Select-Object Name


\end{lstlisting}

On Windows, the most convenient method is the \textbf{Help → Font List} menu inside the OpenSCAD application --- it shows every available font with its exact name string.

\textbf{macOS (bash):}

\begin{lstlisting}[style=Alabaster, language=bash]
fc-list | sort

# or

system_profiler SPFontsDataType | grep "Full Name" | sort


\end{lstlisting}

\subsubsection*{Classroom-Safe Fonts (Available on Most Linux Systems)}\label{docs__pandoc__latex__src__3dmake_all_lessons.md__classroom-safe-fonts-available-on-most-linux-systems}

These fonts are widely pre-installed and safe to use in code you\textquotesingle ll share with classmates:

\begin{itemize}
\tightlist
\item
  \texttt{Liberation\ Sans} --- clean, modern sans-serif, compatible proportionally with Arial
\item
  \texttt{Liberation\ Mono} --- monospaced, excellent for part numbers and code
\item
  \texttt{DejaVu\ Serif} --- serif font with good contrast for labels
\end{itemize}

Note

If OpenSCAD can\textquotesingle t find your specified font, it prints \texttt{WARNING:\ No\ match\ for\ requested\ font\ family} in the console. Font availability varies by OS --- if text looks wrong, check the console for this warning and substitute an available font.

\begin{center}\rule{0.5\linewidth}{0.5pt}\end{center}

\subsection*{Step 3 --- Font-Agnostic Code Pattern}\label{docs__pandoc__latex__src__3dmake_all_lessons.md__step-3--font-agnostic-code-pattern}

\subsubsection*{Define Your Font Once}\label{docs__pandoc__latex__src__3dmake_all_lessons.md__define-your-font-once}

Instead of hardcoding a font name at every text call, define it once as a constant at the top of your file. This way, switching fonts is a one-line change, and the font choice is clearly documented.

\begin{lstlisting}[style=Alabaster, language=openscad]
// Define font choice once at the top — comment/uncomment for your OS

LABEL_FONT = "Liberation Sans:style=Bold";   // Linux

// LABEL_FONT = "DejaVu Sans:style=Bold";    // Linux alternative

// LABEL_FONT = "Arial:style=Bold";          // Windows / macOS

// LABEL_FONT = "Helvetica:style=Bold";      // macOS


module label(txt, size=8) {

  text(txt, font=LABEL_FONT, size=size, halign="center", valign="center");

}


module embossed_plate(message, plate_w=80, plate_d=15, plate_h=5,

                      text_size=8, relief=1.5) {

  cube([plate_w, plate_d, plate_h]);

  translate([plate_w/2, plate_d/2, plate_h])

    linear_extrude(height=relief)

      label(message, size=text_size);

}


embossed_plate("FRONT");

translate([0, 20, 0]) embossed_plate("BACK");

translate([0, 40, 0]) embossed_plate("LEFT");


\end{lstlisting}

\begin{center}\rule{0.5\linewidth}{0.5pt}\end{center}

\subsection*{Step 4 --- Text Sizing: The Cap-Height Model}\label{docs__pandoc__latex__src__3dmake_all_lessons.md__step-4--text-sizing-the-cap-height-model}

\subsubsection*{\texorpdfstring{How the \texttt{size} Parameter Works}{How the size Parameter Works}}\label{docs__pandoc__latex__src__3dmake_all_lessons.md__how-the-size-parameter-works}

The \texttt{size} parameter in \texttt{text()} sets the \textbf{cap height} --- the height of a capital letter in millimeters. This is not the total height of the text block; lowercase letters with descenders (g, p, q, y) extend further below the baseline.

Useful rules of thumb:

\begin{itemize}
\tightlist
\item
  Total text block height ≈ \texttt{size\ ×\ 1.25} (to include descenders)
\item
  Width of a capital letter ≈ \texttt{size\ ×\ 0.65} to \texttt{size\ ×\ 0.75}
\item
  Total width of an all-caps string ≈ \texttt{size\ ×\ 0.65\ ×\ character\_count}
\end{itemize}

These are approximations --- the exact values depend on the specific font.

\subsubsection*{Auto-Fitting Text to a Container}\label{docs__pandoc__latex__src__3dmake_all_lessons.md__auto-fitting-text-to-a-container}

When you need text to fill a specific width and height automatically, this module computes the right size:

\begin{lstlisting}[style=Alabaster, language=openscad]
module fitted_text(txt, container_w, container_h,

                   font="Liberation Sans:style=Bold") {

  char_count  = len(txt);

  size_by_w   = container_w / (char_count * 0.68);  // size that fills width

  size_by_h   = container_h * 0.75;                  // size that fills height

  fitted_size = min(size_by_w, size_by_h);            // use whichever is smaller


  echo(str("Fitted text size for '", txt, "': ", fitted_size, "mm"));


  text(txt, size=fitted_size, font=font,

       halign="center", valign="center");

}


// These labels auto-size to fit their containers

difference() {

  cube([60, 12, 5]);

  translate([30, 6, 5 - 0.8])

    linear_extrude(height=1)

      fitted_text("HELLO", 60, 12);

}


translate([0, 20, 0]) {

  cube([60, 12, 5]);

  translate([30, 6, 5])

    linear_extrude(height=1.5)

      fitted_text("WORLD", 60, 12);

}


\end{lstlisting}

\begin{center}\rule{0.5\linewidth}{0.5pt}\end{center}

\subsection*{Step 5 --- String Functions}\label{docs__pandoc__latex__src__3dmake_all_lessons.md__step-5--string-functions}

\subsubsection*{Building Labels Dynamically}\label{docs__pandoc__latex__src__3dmake_all_lessons.md__building-labels-dynamically}

The \texttt{str()} function converts any value to its string representation and concatenates multiple values into a single string. This is how you create part numbers, serial numbers, and version stamps programmatically.

\begin{lstlisting}[style=Alabaster, language=openscad]
// str() — convert values to strings and concatenate

part_id = str("PART-", 2026, "-", 42);

echo(part_id);   // PART-2026-42


// len() — length of a string or list

name = "OpenSCAD";

echo(len(name));   // 8


// search() — find the position of a character in a string

// Returns a list; [[4]] means found at index 4

echo(search("S", "OpenSCAD"));   // [[4]]


// substr() — extract a substring

// substr(string, start_index, length)

full = "BATCH-001";

batch  = substr(full, 0, 5);   // "BATCH"

number = substr(full, 6, 3);   // "001"

echo(batch, number);


\end{lstlisting}

\subsubsection*{Auto-Generating Serial Numbers}\label{docs__pandoc__latex__src__3dmake_all_lessons.md__auto-generating-serial-numbers}

\begin{lstlisting}[style=Alabaster, language=openscad]
// Generate 5 parts with serial numbers SN-01 through SN-05

for (i = [1:5]) {

  serial = str("SN-", i < 10 ? "0" : "", i);   // zero-pad single digits

  translate([i * 35, 0, 0]) {

    cube([30, 12, 4]);

    translate([15, 6, 4])

      linear_extrude(1.5)

        text(serial, size=5, font="Liberation Sans:style=Bold",

             halign="center", valign="center");

  }

}


\end{lstlisting}

\begin{center}\rule{0.5\linewidth}{0.5pt}\end{center}

\subsection*{\texorpdfstring{Step 6 --- \texttt{let()} for Scoped Variables}{Step 6 --- let() for Scoped Variables}}\label{docs__pandoc__latex__src__3dmake_all_lessons.md__step-6--let-for-scoped-variables}

\subsubsection*{\texorpdfstring{What \texttt{let()} Does}{What let() Does}}\label{docs__pandoc__latex__src__3dmake_all_lessons.md__what-let-does}

In Lesson 3 you learned that OpenSCAD global variables are resolved at parse time, and that last-assignment wins. \texttt{let()} provides a way to declare variables that are scoped to a specific block --- they only exist inside the \texttt{let()} block and don\textquotesingle t affect global scope.

\texttt{let()} is useful when you need intermediate calculated values in one place and don\textquotesingle t want them polluting the global namespace, or when you want to make a series of dependent calculations explicit within one block.

\begin{lstlisting}[style=Alabaster, language=openscad]
// let() scopes variables to the block they're declared in

let (

  base_w          = 80,

  base_d          = 50,

  plate_height    = 5,

  text_center_x   = base_w / 2,

  text_center_y   = base_d / 2

) {

  cube([base_w, base_d, plate_height]);

  translate([text_center_x, text_center_y, plate_height])

    linear_extrude(2)

      text("v1.0", size=6, halign="center", valign="center");

}

// text_center_x and text_center_y are not accessible outside the let() block


\end{lstlisting}

\begin{center}\rule{0.5\linewidth}{0.5pt}\end{center}

\subsection*{Step 7 --- A Complete Parametric Label Plate}\label{docs__pandoc__latex__src__3dmake_all_lessons.md__step-7--a-complete-parametric-label-plate}

Here is a production-quality label module that uses every concept from this lesson:

\begin{lstlisting}[style=Alabaster, language=openscad]
// Parametric Label Plate

// Adjust these parameters and run: 3dm build

label_text   = "STORAGE BOX";

label_w      = 80;        // mm — plate width

label_h      = 18;        // mm — plate height

plate_depth  = 3;         // mm — plate thickness

text_depth   = 1.2;       // mm — engraving depth

font_size    = 7;         // mm — cap height of text

corner_r     = 2;         // mm — corner rounding radius


module label_plate(txt, w, h, t, td, fs, cr) {

  difference() {

    // Rounded rectangle base (minkowski creates rounded corners)

    minkowski() {

      cube([w - 2*cr, h - 2*cr, t], center=true);

      cylinder(r=cr, h=0.01, $fn=16);

    }

    // Engraved text, centered

    translate([0, 0, t/2 - td + 0.001])

      linear_extrude(td + 0.001)

        text(txt, size=fs, font="Liberation Sans:style=Bold",

             halign="center", valign="center", $fn=4);

    // Mounting hole for wall attachment

    translate([w/2 - 6, 0, -0.001])

      cylinder(r=1.5, h=t + 0.002, $fn=16);

  }

}


label_plate(label_text, label_w, label_h, plate_depth, text_depth, font_size, corner_r);


\end{lstlisting}

\begin{center}\rule{0.5\linewidth}{0.5pt}\end{center}

\subsection*{Step 8 --- Automating Multi-Label Generation}\label{docs__pandoc__latex__src__3dmake_all_lessons.md__step-8--automating-multi-label-generation}

Once your label design is parametric, you can generate a whole set of labels automatically --- one build command per label variant.

\textbf{Linux / macOS (bash):}

\begin{lstlisting}[style=Alabaster, language=bash]
#!/bin/bash

LABELS=("FRONT" "BACK" "LEFT" "RIGHT" "TOP" "BOTTOM")

mkdir -p build/labels


for label in "${LABELS[@]}"; do

  filename=$(echo "$label" | tr '[:upper:]' '[:lower:]')

  output="build/labels/${filename}.stl"

  echo "Generating: $label → $output"

  openscad -D "label_text=\"${label}\"" -o "$output" src/main.scad

  [ $? -eq 0 ] && echo "  OK" || echo "  FAILED"

done

echo "Done. Output: build/labels/"


\end{lstlisting}

\textbf{Windows (PowerShell):}

\begin{lstlisting}[style=Alabaster, language=powershell]
$labels = @("FRONT", "BACK", "LEFT", "RIGHT", "TOP", "BOTTOM")

New-Item -ItemType Directory -Force -Path "build\labels" | Out-Null


foreach ($label in $labels) {

  $filename = $label.ToLower()

  $output   = "build\labels\$filename.stl"

  Write-Host "Generating: $label → $output"

  openscad -D "label_text=`"$label`"" -o $output "src\main.scad"

  if ($LASTEXITCODE -eq 0) { Write-Host "  OK" } else { Write-Host "  FAILED" }

}

Write-Host "Done."


\end{lstlisting}

\textbf{Windows (CMD):}

\begin{lstlisting}[style=Alabaster, language=cmd]
@echo off

mkdir build\labels 2>nul

FOR %%L IN (FRONT BACK LEFT RIGHT TOP BOTTOM) DO (

  echo Generating: %%L

  openscad -D "label_text=\"%%L\"" -o "build\labels\%%L.stl" src\main.scad

  IF %ERRORLEVEL% EQU 0 (echo   OK) ELSE (echo   FAILED)

)

echo Done.


\end{lstlisting}

\begin{center}\rule{0.5\linewidth}{0.5pt}\end{center}

\subsection*{Exercises}\label{docs__pandoc__latex__src__3dmake_all_lessons.md__exercises-2}

\textbf{Exercise 6.1:} Find the fonts available on your computer using the appropriate command for your OS. List three font names in the exact format required by OpenSCAD\textquotesingle s \texttt{font=} parameter (e.g., \texttt{Liberation\ Sans:style=Bold}).

\textbf{Exercise 6.2:} Create a parametric name plate module where the text is auto-sized to fill 75\% of the plate width. Test it with strings of 3, 6, 9, and 12 characters --- verify that the text scales correctly and doesn\textquotesingle t overflow.

\textbf{Exercise 6.3 (Advanced):} Write a shell script that generates 10 identical parts with serial numbers SN-001 through SN-010 embossed on each. Display success or failure for each part and the total file count at the end.

\begin{center}\rule{0.5\linewidth}{0.5pt}\end{center}

\subsection*{Quiz --- Lesson 6 (15 questions)}\label{docs__pandoc__latex__src__3dmake_all_lessons.md__quiz--lesson-6-15-questions}

\begin{enumerate}
\tightlist
\item
  What does \texttt{3dm\ orient} output and when would you use it?
\item
  What is the difference between embossed text and debossed (engraved) text in 3D printing? How does the OpenSCAD code differ?
\item
  What does the \texttt{size} parameter in \texttt{text()} represent specifically? Is it the total height of the text block?
\item
  What parameter in \texttt{text()} controls horizontal alignment, and what are the three main options?
\item
  What does \texttt{str("PART-",\ 2024,\ "-",\ 1)} return?
\item
  What does \texttt{len("OpenSCAD")} return?
\item
  What is the purpose of \texttt{let()} in OpenSCAD, and how does it differ from declaring a global variable?
\item
  What does setting \texttt{\$fn=4} do to letters rendered with \texttt{text()}?
\item
  What does \texttt{3dm\ preview} produce and where does the output file go?
\item
  What is the format for specifying a font style (e.g., bold) in OpenSCAD\textquotesingle s \texttt{font=} parameter?
\item
  Explain the depth calculation for debossed text: if a plate is 6 mm tall and you want 1 mm engraving depth, where do you position the \texttt{linear\_extrude()} call and what height do you give it?
\item
  What does \texttt{substr("BATCH-001",\ 6,\ 3)} return?
\item
  What happens if you use \texttt{text()} without wrapping it in \texttt{linear\_extrude()}?
\item
  Why is defining the font as a constant at the top of your file (e.g., \texttt{LABEL\_FONT\ =\ "Liberation\ Sans:style=Bold"}) better than hardcoding it in every \texttt{text()} call?
\item
  Write the OpenSCAD code to create a 50×20×4mm plate with the text "LOT-42" engraved 1.5mm deep, centered on the top face.
\end{enumerate}

\begin{center}\rule{0.5\linewidth}{0.5pt}\end{center}

\subsection*{Extension Problems (15)}\label{docs__pandoc__latex__src__3dmake_all_lessons.md__extension-problems-15-5}

\begin{enumerate}
\tightlist
\item
  Build a parametric serial number generator: accept a prefix string and a starting number, and generate a set of embossed labels with incrementing serial numbers.
\item
  Create a dynamic version stamp: use \texttt{str()} to combine a product name, major version, and minor version, all as separate parameters.
\item
  Design a 4-up label sheet: four identical labels in a 2×2 grid, using \texttt{translate()} and a \texttt{for} loop with two variables.
\item
  Build a two-line label: stack two \texttt{text()} calls at different heights. Line 1 is the label text; line 2 is a smaller subtitle.
\item
  Create a label with a decorative border: use \texttt{difference()} to cut a framing slot into the top face of the label plate.
\item
  Use \texttt{3dm\ orient} on three different models (a flat slab, a tall cylinder, an L-bracket) and document whether you agree with the AI suggestion in each case.
\item
  Build a "batch tag" system: a module that generates 12 small tags in a row, each with a different serial number from a list.
\item
  Design a keychain tag: a rounded rectangle with a ring hole, parametric text, and a border around the text.
\item
  Use \texttt{search()} to find the position of a separator character (dash or underscore) in a part-number string. Explain a practical use case for this.
\item
  Build a "negative space" text plate: a plate with all letters cut completely through, creating a stencil or silhouette effect.
\item
  Design a parametric drawer label holder: a clip that attaches to the front edge of a drawer and holds a flat label insert.
\item
  Create a font comparison plate: render the same text string using Liberation Sans, Liberation Mono, and DejaVu Serif side by side on one base plate.
\item
  Write a screen-reader accessibility guide for the five \texttt{3dm} commands. For each command: what it does, expected output, and how to interpret that output without visual reference.
\item
  Write a module that auto-reduces font size to fit text within a fixed-width container. The module should reduce size in 0.5 mm steps until the estimated width fits.
\item
  Design a production part-marking system: a jig that holds 10 small labels in a row, each with an incremented part number, ready to print all at once as one STL file.
\end{enumerate}

\begin{center}\rule{0.5\linewidth}{0.5pt}\end{center}

\subsection*{References and Helpful Resources}\label{docs__pandoc__latex__src__3dmake_all_lessons.md__references-and-helpful-resources-5}

\begin{itemize}
\item
  3DMake GitHub Repository --- Command Reference --- https://github.com/tdeck/3dmake
\item
  OpenSCAD User Manual --- Text and Fonts --- https://en.wikibooks.org/wiki/OpenSCAD\_User\_Manual/Text
\end{itemize}

\subsubsection*{Supplemental Resources}\label{docs__pandoc__latex__src__3dmake_all_lessons.md__supplemental-resources-5}

\begin{itemize}
\tightlist
\item
  Programming with OpenSCAD EPUB Textbook --- assets/Programming\_with\_OpenSCAD.epub --- String functions and text embossing examples
\item
  OpenSCAD Quick Reference --- https://programmingwithopenscad.github.io/quick-reference.html --- All string function syntax
\item
  CodeSolutions Repository --- https://github.com/ProgrammingWithOpenSCAD/CodeSolutions --- Worked text embossing examples
\item
  All other reference documents are embedded as appendices below --- no separate files needed.
\end{itemize}

\begin{center}\rule{0.5\linewidth}{0.5pt}\end{center}

\section*{Part 2: Extension Project and Reference Appendices}\label{docs__pandoc__latex__src__3dmake_all_lessons.md__part-2-extension-project-and-reference-appendices-2}

\emph{All supporting documents for Lesson 6 are included below. You do not need to open any additional files.}

\begin{center}\rule{0.5\linewidth}{0.5pt}\end{center}

\subsection*{Appendix A: Parametric Keychain --- Extension Project}\label{docs__pandoc__latex__src__3dmake_all_lessons.md__appendix-a-parametric-keychain--extension-project}

\textbf{Estimated time:} 2--4 hours

\subsubsection*{Overview}\label{docs__pandoc__latex__src__3dmake_all_lessons.md__overview-5}

This extension project applies the text embossing and parametric design skills from Lesson 6 to build something immediately usable: a personalized keychain. The goal is not just to make one keychain --- it\textquotesingle s to build a \emph{system} that can produce any variant (different names, sizes, materials) by changing a small set of top-level variables. That\textquotesingle s what makes it a professional parametric design rather than a one-off model.

\subsubsection*{Learning Objectives}\label{docs__pandoc__latex__src__3dmake_all_lessons.md__learning-objectives-10}

By completing this project, you will:

\begin{itemize}
\tightlist
\item
  Create parametric OpenSCAD modules that accept user inputs as top-level variables
\item
  Implement embossed and debossed text using \texttt{linear\_extrude()} of 2D text shapes
\item
  Generate and test multiple design variants systematically
\item
  Document design parameters for reproducibility and user customization
\end{itemize}

\subsubsection*{Objective}\label{docs__pandoc__latex__src__3dmake_all_lessons.md__objective}

Create a \texttt{keychain.scad} file that adapts to custom text, dimensions, and attachment ring size --- all through top-level parameter variables. Anyone with your \texttt{.scad} file should be able to produce a custom keychain by changing only the variables at the top of the file.

\subsubsection*{Tasks}\label{docs__pandoc__latex__src__3dmake_all_lessons.md__tasks-4}

\begin{enumerate}
\item
  \textbf{Create \texttt{keychain.scad}} with these top-level parameters at minimum: \texttt{name}, \texttt{tag\_w}, \texttt{tag\_h}, \texttt{tag\_thickness}, \texttt{text\_size}, \texttt{ring\_d}. Add others as your design requires. Every dimension in the file should trace back to one of these variables.
\item
  \textbf{Implement text} using \texttt{linear\_extrude()} of a 2D \texttt{text()} shape. Add either embossed text (the letters rise above the surface) or debossed text (the letters are engraved into the surface). Embossed text is more durable; debossed text is easier to read when painted or filled with ink.
\item
  \textbf{Produce three size variants} --- small (25 mm wide), medium (35 mm wide), large (45 mm wide) --- by changing only \texttt{tag\_w} and \texttt{tag\_h} proportionally. Export an STL for each and record print time and filament weight.
\item
  \textbf{Test the key ring attachment} by inserting a standard key ring (typically 25--30 mm diameter wire) through the attachment loop. Document what ring diameter fits, and whether the attachment point survives a reasonable tug test.
\end{enumerate}

\subsubsection*{Deliverables}\label{docs__pandoc__latex__src__3dmake_all_lessons.md__deliverables-1}

\begin{itemize}
\tightlist
\item
  \texttt{keychain.scad} --- fully parametric source file with all key dimensions as named variables and comments explaining each
\item
  Three STL variant files (small, medium, large)
\item
  Completed student documentation template (Appendix B of this document)
\item
  Print settings log and fit-test report for key ring attachment
\end{itemize}

\begin{center}\rule{0.5\linewidth}{0.5pt}\end{center}

\subsubsection*{Starter Code}\label{docs__pandoc__latex__src__3dmake_all_lessons.md__starter-code-3}

Use the minimal keycap example as your starting structure for the shell + emboss pattern. The more advanced examples below show how to add circular text arrangements, multi-line compositions, and font variation.

\textbf{Minimal Starting Point:}

\begin{lstlisting}[style=Alabaster, language=openscad]
// Cube Keycap — Structural Reference

// This shows the shell + emboss pattern. Replace with keychain geometry.


keysize    = 18;    // mm — tag size

keyheight  = 12;    // mm — tag thickness

wall       = 1.2;   // mm — shell wall

letter     = "R";   // change to your text

lettersize = 10;    // mm — text size

letteraise = 0.8;   // mm — emboss height above surface


module shell() {

  difference() {

    cube([keysize, keysize, keyheight], center=false);

    translate([wall, wall, wall])

      cube([keysize - 2*wall, keysize - 2*wall, keyheight], center=false);

  }

}


module emboss() {

  translate([keysize/2, keysize/2, keyheight - 0.01])

    linear_extrude(height=letteraise)

      text(letter, size=lettersize, halign="center", valign="center");

}


union() {

  shell();

  emboss();

}


\end{lstlisting}

\begin{center}\rule{0.5\linewidth}{0.5pt}\end{center}

\textbf{Parametric Keychain Nameplate} --- use this as your primary starting scaffold:

\begin{lstlisting}[style=Alabaster, language=openscad]
// Customizable Keychain Nameplate

// Change the parameters below to create your own variant.


// ---- Top-Level Parameters (change these) ----

name          = "ALEX";    // text to emboss on the tag

tag_w         = 35;        // mm — tag width

tag_h         = 20;        // mm — tag height (vertical dimension)

tag_thickness = 5;         // mm — tag thickness (Z height)

wall          = 1.2;       // mm — shell wall thickness

text_size     = 12;        // mm — font cap-height

text_raise    = 0.8;       // mm — emboss height above top face

ring_d        = 6;         // mm — key ring hole inner diameter

ring_wall     = 2;         // mm — wall thickness around ring hole


// ---- Derived Values ----

ring_outer_d  = ring_d + 2 * ring_wall;


// Verify ring attachment won't be too thin to print reliably

assert(ring_wall >= 1.2,

  str("ring_wall too thin to print reliably: ", ring_wall, "mm. Minimum: 1.2mm"));


// ---- Shell Module ----

module tag_shell() {

  difference() {

    cube([tag_w, tag_h, tag_thickness]);

    translate([wall, wall, wall])

      cube([tag_w - 2*wall, tag_h - 2*wall, tag_thickness]);  // open top — hollow shell

  }

}


// ---- Text Module ----

module tag_text() {

  translate([tag_w/2, tag_h/2, tag_thickness - 0.01])

    linear_extrude(height=text_raise)

      text(name,

           size=text_size,

           font="Liberation Sans:style=Bold",

           halign="center",

           valign="center",

           $fn=4);   // $fn=4 speeds up preview; use $fn=32 for final render

}


// ---- Key Ring Attachment Loop ----

// A short cylinder projecting from the top edge center with a hole through it.

module ring_attachment() {

  translate([tag_w/2, tag_h + ring_outer_d/2 - 0.001, tag_thickness/2])

    rotate([90, 0, 0])

      difference() {

        cylinder(d=ring_outer_d, h=tag_thickness, center=true, $fn=32);

        cylinder(d=ring_d,       h=tag_thickness + 0.002, center=true, $fn=32);

      }

}


// ---- Assembly ----

union() {

  tag_shell();

  tag_text();

  ring_attachment();

}


\end{lstlisting}

\begin{center}\rule{0.5\linewidth}{0.5pt}\end{center}

\textbf{Advanced Technique: Circular Text Arrangement}

\begin{lstlisting}[style=Alabaster, language=openscad]
// Arrange text in a circle around a center point.

// Useful for medallion-style keychains, badge designs, or decorative borders.


module rotate_text(display_text,

    text_size    = 10,

    distance     = 20,

    rotation_span = 360,

    tilt         = 0)

{

  rotate([0, 0, tilt])

  for (i = [0 : len(display_text) - 1]) {

    rotate([0, 0, -i * rotation_span / len(display_text)])

    translate([0, distance, 0])

      text(display_text[i],

           font="Liberation Sans:style=Bold",

           size=text_size,

           halign="center");

  }

}


// Usage: extrude the circular text arrangement

linear_extrude(height=2)

  rotate_text("MAKER", text_size=12, distance=30, rotation_span=75, tilt=30);


\end{lstlisting}

\begin{center}\rule{0.5\linewidth}{0.5pt}\end{center}

\textbf{Advanced Technique: Multi-Line Text Composition}

\begin{lstlisting}[style=Alabaster, language=openscad]
// Combine multiple text elements at different sizes for a professional nameplate.


module nameplate(name, role, plate_w=120, plate_d=60, plate_h=3) {

  union() {

    // Base plate

    cube([plate_w, plate_d, plate_h], center=true);


    // Main name — large, centered near top

    translate([0, 15, plate_h/2 + 1.5])

      linear_extrude(height=2)

        text(name, size=24, font="Liberation Sans:style=Bold",

             halign="center", valign="center");


    // Role — medium

    translate([0, 0, plate_h/2 + 1.5])

      linear_extrude(height=2)

        text(role, size=14, font="Liberation Sans:style=Regular",

             halign="center", valign="center");

  }

}


nameplate("Alex Chen", "3D Design Engineer");


\end{lstlisting}

\begin{center}\rule{0.5\linewidth}{0.5pt}\end{center}

\textbf{Font Reference: Cross-Platform Options}

The fonts below are widely available. Always fall back gracefully --- define your font as a constant at the top of your file so you change it in one place.

\begin{lstlisting}[style=Alabaster, language=openscad]
// Define font as a constant — change here, applies everywhere

BODY_FONT  = "Liberation Sans:style=Bold";    // clean, reliable on all platforms

MONO_FONT  = "Liberation Mono:style=Regular"; // technical/code style

SERIF_FONT = "DejaVu Serif:style=Regular";    // classic, decorative


// Usage

linear_extrude(height=2)

  text("BOLD", size=20, font=BODY_FONT, halign="center");


\end{lstlisting}

\subsubsection*{Assessment Questions}\label{docs__pandoc__latex__src__3dmake_all_lessons.md__assessment-questions}

\begin{enumerate}
\tightlist
\item
  How did you use OpenSCAD parameters to enable customization without requiring users to edit code logic?
\item
  What were the key differences in print time and material usage between your three size variants?
\item
  How did you test the key ring attachment, and what adjustments did your testing reveal?
\end{enumerate}

\begin{center}\rule{0.5\linewidth}{0.5pt}\end{center}

\subsection*{Appendix B: Parametric Keychain --- Student Documentation Template}\label{docs__pandoc__latex__src__3dmake_all_lessons.md__appendix-b-parametric-keychain--student-documentation-template}

\textbf{Author:} \_\_\_\_\_\_\_\_\_\_\_\_\_\_\_\_
\textbf{Date:} \_\_\_\_\_\_\_\_\_\_
\textbf{Description:} Design a fully parametric keychain that supports personalization and customization through top-level variables.

\begin{center}\rule{0.5\linewidth}{0.5pt}\end{center}

\subsubsection*{Design Concept}\label{docs__pandoc__latex__src__3dmake_all_lessons.md__design-concept-1}

\begin{itemize}
\tightlist
\item
  Keychain theme or purpose (personal use, gift, label system, other):
\item
  Design elements (shape, attachment method, text or symbol, any decorative features):
\item
  Parametric strategy --- what will be variables, and what will be fixed geometry?
\end{itemize}

\begin{center}\rule{0.5\linewidth}{0.5pt}\end{center}

\subsubsection*{Parametric Variables}\label{docs__pandoc__latex__src__3dmake_all_lessons.md__parametric-variables}

List every named variable in your file and what it controls.

{\def\LTcaptype{none} % do not increment counter
\begin{longtable}[]{@{}lccl@{}}
\toprule\noalign{}
Variable & Default Value & Unit & Purpose / What it affects \\
\midrule\noalign{}
\endhead
\bottomrule\noalign{}
\endlastfoot
& & mm & \\
& & mm & \\
& & mm & \\
& & mm & \\
& & --- & \\
\end{longtable}
}

\begin{center}\rule{0.5\linewidth}{0.5pt}\end{center}

\subsubsection*{Variant Configurations}\label{docs__pandoc__latex__src__3dmake_all_lessons.md__variant-configurations}

Document the three variants you produced. Change only size-related variables between variants --- text, attachment, and structural relationships should scale automatically.

{\def\LTcaptype{none} % do not increment counter
\begin{longtable}[]{@{}
  >{\centering\arraybackslash}p{(\linewidth - 10\tabcolsep) * \real{0.0886}}
  >{\centering\arraybackslash}p{(\linewidth - 10\tabcolsep) * \real{0.1519}}
  >{\centering\arraybackslash}p{(\linewidth - 10\tabcolsep) * \real{0.1519}}
  >{\centering\arraybackslash}p{(\linewidth - 10\tabcolsep) * \real{0.2025}}
  >{\centering\arraybackslash}p{(\linewidth - 10\tabcolsep) * \real{0.1899}}
  >{\centering\arraybackslash}p{(\linewidth - 10\tabcolsep) * \real{0.2152}}@{}}
\toprule\noalign{}
\begin{minipage}[b]{\linewidth}\centering
Variant
\end{minipage} & \begin{minipage}[b]{\linewidth}\centering
\texttt{tag\_w}
\end{minipage} & \begin{minipage}[b]{\linewidth}\centering
 (mm)
\end{minipage} & \begin{minipage}[b]{\linewidth}\centering
\texttt{tag\_h}
\end{minipage} & \begin{minipage}[b]{\linewidth}\centering
 (mm)
\end{minipage} & \begin{minipage}[b]{\linewidth}\centering
\texttt{text\_size}
\end{minipage} \\
\midrule\noalign{}
\endhead
\bottomrule\noalign{}
\endlastfoot
Small & & & & & \\
Medium & & & & & \\
Large & & & & & \\
\end{longtable}
}

\begin{center}\rule{0.5\linewidth}{0.5pt}\end{center}

\subsubsection*{Print Settings Used}\label{docs__pandoc__latex__src__3dmake_all_lessons.md__print-settings-used}

{\def\LTcaptype{none} % do not increment counter
\begin{longtable}[]{@{}lll@{}}
\toprule\noalign{}
Setting & Value & Notes \\
\midrule\noalign{}
\endhead
\bottomrule\noalign{}
\endlastfoot
Layer height & & \\
Infill & & \\
Supports & Yes / No & \\
Nozzle temperature & °C & \\
Bed temperature & °C & \\
Material & & \\
\end{longtable}
}

\begin{center}\rule{0.5\linewidth}{0.5pt}\end{center}

\subsubsection*{Key Ring Attachment Test}\label{docs__pandoc__latex__src__3dmake_all_lessons.md__key-ring-attachment-test}

{\def\LTcaptype{none} % do not increment counter
\begin{longtable}[]{@{}ll@{}}
\toprule\noalign{}
Test & Result \\
\midrule\noalign{}
\endhead
\bottomrule\noalign{}
\endlastfoot
Ring diameter tested & mm \\
Did ring fit through attachment loop? & Yes / No / Tight \\
Does the loop survive a firm tug? & Yes / No \\
Wall thickness at attachment point & mm \\
Any deformation or cracking observed? & \\
\end{longtable}
}

Adjustments needed based on testing:

\begin{center}\rule{0.5\linewidth}{0.5pt}\end{center}

\subsubsection*{Print and Assembly Results}\label{docs__pandoc__latex__src__3dmake_all_lessons.md__print-and-assembly-results}

For each variant, describe the print result:

\textbf{Small variant:}

\begin{itemize}
\tightlist
\item
  Text legibility (clear / readable / hard to read):
\item
  Surface quality:
\item
  Any print issues:
\end{itemize}

\textbf{Medium variant:}

\begin{itemize}
\tightlist
\item
  Text legibility:
\item
  Surface quality:
\item
  Any print issues:
\end{itemize}

\textbf{Large variant:}

\begin{itemize}
\tightlist
\item
  Text legibility:
\item
  Surface quality:
\item
  Any print issues:
\end{itemize}

\begin{center}\rule{0.5\linewidth}{0.5pt}\end{center}

\subsubsection*{Customization Guide}\label{docs__pandoc__latex__src__3dmake_all_lessons.md__customization-guide}

Write this section for a future user who will download your \texttt{.scad} file and want to make their own version:

\begin{enumerate}
\tightlist
\item
  \textbf{To change the name:} set \texttt{name\ =\ "YOUR\ TEXT"} at line \_\_\_.
\item
  \textbf{To change the size:} adjust \texttt{tag\_w} and \texttt{tag\_h} proportionally. Recommended ratio: \texttt{tag\_h\ =\ tag\_w\ *\ 0.57}.
\item
  \textbf{To change the font:} update the \texttt{LABEL\_FONT} constant at line \_\_\_.
\item
  \textbf{Common customizations:}

  \begin{itemize}
  \tightlist
  \item
    Smaller text that doesn\textquotesingle t overflow: reduce \texttt{text\_size} until the preview looks right.
  \item
    Different ring size: change \texttt{ring\_d} to match your key ring\textquotesingle s wire diameter + 1 mm clearance.
  \end{itemize}
\end{enumerate}

\begin{center}\rule{0.5\linewidth}{0.5pt}\end{center}

\subsubsection*{Reflections}\label{docs__pandoc__latex__src__3dmake_all_lessons.md__reflections-4}

\begin{enumerate}
\item
  Which parameterization decision had the most impact on the usefulness of the final design?
\item
  How could someone who doesn\textquotesingle t know OpenSCAD still customize this design? What instructions would you write for them?
\item
  What would you add in a future iteration --- a second line of text, a logo, a different attachment style?
\end{enumerate}

\begin{center}\rule{0.5\linewidth}{0.5pt}\end{center}

\subsubsection*{Attachments Checklist}\label{docs__pandoc__latex__src__3dmake_all_lessons.md__attachments-checklist-4}

\begin{itemize}
\tightlist
\item[$\square$]
  \texttt{.scad} file with full parametric structure and comments on all variables
\item[$\square$]
  Photos of all three printed variants side by side
\item[$\square$]
  Photos of key ring fit test (ring inserted through attachment loop)
\item[$\square$]
  Completed variant configuration table
\item[$\square$]
  Customization guide written for a future user
\end{itemize}

\begin{center}\rule{0.5\linewidth}{0.5pt}\end{center}

\subsubsection*{Teacher Feedback}\label{docs__pandoc__latex__src__3dmake_all_lessons.md__teacher-feedback-4}

{\def\LTcaptype{none} % do not increment counter
\begin{longtable}[]{@{}
  >{\raggedright\arraybackslash}p{(\linewidth - 4\tabcolsep) * \real{0.8476}}
  >{\centering\arraybackslash}p{(\linewidth - 4\tabcolsep) * \real{0.1048}}
  >{\raggedright\arraybackslash}p{(\linewidth - 4\tabcolsep) * \real{0.0476}}@{}}
\toprule\noalign{}
\begin{minipage}[b]{\linewidth}\raggedright
Category
\end{minipage} & \begin{minipage}[b]{\linewidth}\centering
Score (0--3)
\end{minipage} & \begin{minipage}[b]{\linewidth}\raggedright
Notes
\end{minipage} \\
\midrule\noalign{}
\endhead
\bottomrule\noalign{}
\endlastfoot
Problem \& Solution --- keychains print and function; key ring fits &
& \\
Design \& Code Quality --- fully parametric; all dimensions are variables; code is commented
& & \\
Documentation --- variant table complete; customization guide usable by a stranger
& & \\
\textbf{Total (0--9)} & & \\
\end{longtable}
}

\begin{center}\rule{0.5\linewidth}{0.5pt}\end{center}

\subsection*{Appendix C: Parametric Keychain --- Instructor Reference}\label{docs__pandoc__latex__src__3dmake_all_lessons.md__appendix-c-parametric-keychain--instructor-reference}

\subsubsection*{Project Context}\label{docs__pandoc__latex__src__3dmake_all_lessons.md__project-context-3}

This extension project reinforces the core text embossing and parametric design content of Lesson 6. The key educational goal is \textbf{parameter-driven design}: every visible dimension should trace back to a named variable, so changing one number updates the entire design proportionally. Students who hardcode numbers (e.g., \texttt{translate({[}35/2,\ 20/2,\ 5{]})} instead of \texttt{translate({[}tag\_w/2,\ tag\_h/2,\ tag\_thickness{]})}) miss the point even if the print looks fine.

\subsubsection*{Key Learning Points to Reinforce}\label{docs__pandoc__latex__src__3dmake_all_lessons.md__key-learning-points-to-reinforce-3}

\textbf{All dimensions must be variables.} A fully parametric design produces a different variant for every combination of input values, without any code changes. If changing \texttt{tag\_w} from 25 to 45 produces a broken design (text overflow, attachment loop detaches, walls become invalid), the design is not truly parametric.

\textbf{The font constant pattern.} Students should define \texttt{LABEL\_FONT\ =\ "Liberation\ Sans:style=Bold"} at the top of the file and reference it in every \texttt{text()} call. This is a direct application of the DRY principle from Lesson 3. Hardcoding font strings in every text call is a common mistake.

\textbf{Text overflow is a design failure, not a rendering quirk.} OpenSCAD does not automatically fit text to a container --- if \texttt{text\_size} is too large, the text will extend beyond the tag boundary and produce an illegal geometry when combined with the shell. Students should either add an \texttt{assert()} that checks estimated text width, or test each variant in preview before exporting.

\textbf{The key ring attachment geometry is the most technically demanding part.} The cylinder that forms the attachment loop must: (1) be large enough for a real key ring wire to fit through, (2) have walls thick enough to print reliably (minimum \textasciitilde{}1.2 mm), and (3) connect cleanly to the main tag body. The 0.001 mm overlap in the starter code\textquotesingle s \texttt{translate()} prevents non-manifold geometry at the junction.

\subsubsection*{Constraints (Enforce These)}\label{docs__pandoc__latex__src__3dmake_all_lessons.md__constraints-enforce-these-3}

\begin{itemize}
\tightlist
\item
  All key dimensions (width, height, thickness, text size, ring diameter) must be named variables --- no hardcoded numbers in geometry
\item
  At least three variants must be printed, not just rendered
\item
  Key ring attachment must be physically tested, not just visually inspected
\end{itemize}

\subsubsection*{Assessment Notes}\label{docs__pandoc__latex__src__3dmake_all_lessons.md__assessment-notes-3}

\textbf{Strong submissions} show: text that scales proportionally with \texttt{text\_size} as a variable, a ring attachment that has been tested with a real key ring, and a customization guide that actually tells a non-programmer how to make their own version.

\textbf{Common weak areas:} Text overflow not detected until printing (add \texttt{assert()} or preview check to starter code reminder), attachment loop too thin to be durable (check \texttt{ring\_wall\ \textgreater{}=\ 1.2\ mm}), and documentation that describes the final design but doesn\textquotesingle t guide future customization.

\textbf{Extension opportunity:} Students who finish early can investigate font rendering differences on their platform (Linux vs. macOS vs. Windows font availability), or add a second line of text (a subtitle or date) as a separate variable.

\begin{center}\rule{0.5\linewidth}{0.5pt}\end{center}

\subsection*{Appendix D: Slicing Settings Quick Reference}\label{docs__pandoc__latex__src__3dmake_all_lessons.md__appendix-d-slicing-settings-quick-reference}

\emph{Use this reference when preparing your keychain (or any Lesson 6 design) for printing. Find your slicer and printer below. Recommended settings for most keychain prints: 0.20 mm layer height, 20\% infill, no supports.}

\begin{center}\rule{0.5\linewidth}{0.5pt}\end{center}

\subsubsection*{D1 --- PrusaSlicer (Prusa Printers)}\label{docs__pandoc__latex__src__3dmake_all_lessons.md__d1--prusaslicer-prusa-printers}

\paragraph*{Recommended Settings by Use Case}\label{docs__pandoc__latex__src__3dmake_all_lessons.md__recommended-settings-by-use-case}

{\def\LTcaptype{none} % do not increment counter
\begin{longtable}[]{@{}
  >{\raggedright\arraybackslash}p{(\linewidth - 8\tabcolsep) * \real{0.2785}}
  >{\centering\arraybackslash}p{(\linewidth - 8\tabcolsep) * \real{0.1519}}
  >{\centering\arraybackslash}p{(\linewidth - 8\tabcolsep) * \real{0.0759}}
  >{\centering\arraybackslash}p{(\linewidth - 8\tabcolsep) * \real{0.1139}}
  >{\raggedright\arraybackslash}p{(\linewidth - 8\tabcolsep) * \real{0.3797}}@{}}
\toprule\noalign{}
\begin{minipage}[b]{\linewidth}\raggedright
Use Case
\end{minipage} & \begin{minipage}[b]{\linewidth}\centering
Layer Height
\end{minipage} & \begin{minipage}[b]{\linewidth}\centering
Infill
\end{minipage} & \begin{minipage}[b]{\linewidth}\centering
Supports
\end{minipage} & \begin{minipage}[b]{\linewidth}\raggedright
Notes
\end{minipage} \\
\midrule\noalign{}
\endhead
\bottomrule\noalign{}
\endlastfoot
Quick test / prototype & 0.30 mm & 10\% & As needed &
Draft --- fastest, roughest \\
Standard project & 0.20 mm & 15--20\% & As needed &
Best all-around starting point \\
Functional part & 0.20 mm & 30--40\% & As needed &
Use for parts under stress \\
Fine detail / display & 0.15 mm & 15\% & As needed &
Smoother surface; slower \\
Solid reference part & 0.20 mm & 40--50\% & Rarely &
Rarely needed; long print \\
\end{longtable}
}

\paragraph*{Filament Temperature Settings}\label{docs__pandoc__latex__src__3dmake_all_lessons.md__filament-temperature-settings}

{\def\LTcaptype{none} % do not increment counter
\begin{longtable}[]{@{}
  >{\raggedright\arraybackslash}p{(\linewidth - 6\tabcolsep) * \real{0.1159}}
  >{\centering\arraybackslash}p{(\linewidth - 6\tabcolsep) * \real{0.1594}}
  >{\centering\arraybackslash}p{(\linewidth - 6\tabcolsep) * \real{0.1159}}
  >{\raggedright\arraybackslash}p{(\linewidth - 6\tabcolsep) * \real{0.6087}}@{}}
\toprule\noalign{}
\begin{minipage}[b]{\linewidth}\raggedright
Filament
\end{minipage} & \begin{minipage}[b]{\linewidth}\centering
Nozzle Temp
\end{minipage} & \begin{minipage}[b]{\linewidth}\centering
Bed Temp
\end{minipage} & \begin{minipage}[b]{\linewidth}\raggedright
Notes
\end{minipage} \\
\midrule\noalign{}
\endhead
\bottomrule\noalign{}
\endlastfoot
PLA & 200--215°C & 50--60°C & Easiest to print; default choice \\
PETG & 230--250°C & 70--85°C & Use glue stick on PEI bed \\
TPU & 220--240°C & 30--60°C &
Print at 20--30 mm/s max; direct drive only \\
ABS & 230--250°C & 90--110°C & Requires enclosure; ventilate the room \\
\end{longtable}
}

\emph{Always check the temperature range printed on your filament spool --- it varies by brand.}

\paragraph*{Support Settings Guide}\label{docs__pandoc__latex__src__3dmake_all_lessons.md__support-settings-guide}

{\def\LTcaptype{none} % do not increment counter
\begin{longtable}[]{@{}ccl@{}}
\toprule\noalign{}
Overhang Angle & Supports Needed? & Recommendation \\
\midrule\noalign{}
\endhead
\bottomrule\noalign{}
\endlastfoot
\textless{} 45° & No & None \\
45--60° & Maybe & Preview first; add if sagging \\
\textgreater{} 60° & Yes & Support on build plate only \\
Bridge \textless{} 20 mm & No & Bridges usually fine \\
Bridge \textgreater{} 20 mm & Maybe & Preview; consider reorienting \\
\end{longtable}
}

\paragraph*{Common Problems and Quick Fixes}\label{docs__pandoc__latex__src__3dmake_all_lessons.md__common-problems-and-quick-fixes}

{\def\LTcaptype{none} % do not increment counter
\begin{longtable}[]{@{}
  >{\raggedright\arraybackslash}p{(\linewidth - 4\tabcolsep) * \real{0.2593}}
  >{\raggedright\arraybackslash}p{(\linewidth - 4\tabcolsep) * \real{0.3611}}
  >{\raggedright\arraybackslash}p{(\linewidth - 4\tabcolsep) * \real{0.3796}}@{}}
\toprule\noalign{}
\begin{minipage}[b]{\linewidth}\raggedright
Problem
\end{minipage} & \begin{minipage}[b]{\linewidth}\raggedright
Likely Cause
\end{minipage} & \begin{minipage}[b]{\linewidth}\raggedright
Fix
\end{minipage} \\
\midrule\noalign{}
\endhead
\bottomrule\noalign{}
\endlastfoot
Print lifts off bed & Poor adhesion / warping &
Add brim; use glue stick; re-level bed \\
Stringing between parts & Temp too high / retraction &
Lower temp 5°C; check retraction settings \\
Layer lines very visible & Layer height too thick &
Use 0.15 mm or 0.20 mm \\
Print takes too long & Layer height too thin / infill too high &
Use 0.30 mm draft; reduce infill \\
Holes too small & FDM tolerance --- always undersized &
Add 0.2--0.3 mm to hole diameter \\
Part broke at layer boundary & Load perpendicular to layers &
Reorient so load is parallel to layers \\
First layer not sticking & Bed not level & Run bed leveling routine \\
Spaghetti / print failure & Overhang without supports &
Add supports or reorient \\
\end{longtable}
}

\paragraph*{G-code Export Checklist}\label{docs__pandoc__latex__src__3dmake_all_lessons.md__g-code-export-checklist}

Before exporting, confirm:

\begin{itemize}
\tightlist
\item[$\square$]
  Correct printer profile selected
\item[$\square$]
  Correct filament profile selected
\item[$\square$]
  Layer height appropriate for use case
\item[$\square$]
  Infill percentage set
\item[$\square$]
  Supports enabled if needed
\item[$\square$]
  Layer preview reviewed (no floating parts; supports where needed)
\item[$\square$]
  Print time and filament weight noted for your records
\end{itemize}

\emph{Sources: Prusa Research PrusaSlicer knowledge base (help.prusa3d.com); Hubs FDM guide (hubs.com)}

\begin{center}\rule{0.5\linewidth}{0.5pt}\end{center}

\subsubsection*{D2 --- Cura / Anycubic i3 Mega}\label{docs__pandoc__latex__src__3dmake_all_lessons.md__d2--cura--anycubic-i3-mega}

\paragraph*{Recommended Settings by Use Case}\label{docs__pandoc__latex__src__3dmake_all_lessons.md__recommended-settings-by-use-case-1}

{\def\LTcaptype{none} % do not increment counter
\begin{longtable}[]{@{}
  >{\raggedright\arraybackslash}p{(\linewidth - 8\tabcolsep) * \real{0.2785}}
  >{\centering\arraybackslash}p{(\linewidth - 8\tabcolsep) * \real{0.1519}}
  >{\centering\arraybackslash}p{(\linewidth - 8\tabcolsep) * \real{0.0759}}
  >{\centering\arraybackslash}p{(\linewidth - 8\tabcolsep) * \real{0.1139}}
  >{\raggedright\arraybackslash}p{(\linewidth - 8\tabcolsep) * \real{0.3797}}@{}}
\toprule\noalign{}
\begin{minipage}[b]{\linewidth}\raggedright
Use Case
\end{minipage} & \begin{minipage}[b]{\linewidth}\centering
Layer Height
\end{minipage} & \begin{minipage}[b]{\linewidth}\centering
Infill
\end{minipage} & \begin{minipage}[b]{\linewidth}\centering
Supports
\end{minipage} & \begin{minipage}[b]{\linewidth}\raggedright
Notes
\end{minipage} \\
\midrule\noalign{}
\endhead
\bottomrule\noalign{}
\endlastfoot
Quick test / prototype & 0.30 mm & 10\% & As needed &
Draft --- fastest, roughest \\
Standard project & 0.20 mm & 15--20\% & As needed &
Best all-around starting point \\
Functional part & 0.20 mm & 30--40\% & As needed &
Use for parts under stress \\
Fine detail / display & 0.15 mm & 15\% & As needed &
Smoother surface; slower \\
Solid reference part & 0.20 mm & 40--50\% & Rarely &
Rarely needed; long print \\
\end{longtable}
}

\paragraph*{Filament Temperature Settings}\label{docs__pandoc__latex__src__3dmake_all_lessons.md__filament-temperature-settings-1}

{\def\LTcaptype{none} % do not increment counter
\begin{longtable}[]{@{}
  >{\raggedright\arraybackslash}p{(\linewidth - 6\tabcolsep) * \real{0.1026}}
  >{\centering\arraybackslash}p{(\linewidth - 6\tabcolsep) * \real{0.1410}}
  >{\centering\arraybackslash}p{(\linewidth - 6\tabcolsep) * \real{0.1154}}
  >{\raggedright\arraybackslash}p{(\linewidth - 6\tabcolsep) * \real{0.6410}}@{}}
\toprule\noalign{}
\begin{minipage}[b]{\linewidth}\raggedright
Filament
\end{minipage} & \begin{minipage}[b]{\linewidth}\centering
Nozzle Temp
\end{minipage} & \begin{minipage}[b]{\linewidth}\centering
Bed Temp
\end{minipage} & \begin{minipage}[b]{\linewidth}\raggedright
Notes
\end{minipage} \\
\midrule\noalign{}
\endhead
\bottomrule\noalign{}
\endlastfoot
PLA & 200--210°C & 50--60°C & Easiest to print; Anycubic default \\
PETG & 230--240°C & 70--80°C &
Use painter\textquotesingle s tape on steel bed \\
TPU & 220--235°C & 30--50°C & Reduce speed to 25--30 mm/s \\
ABS & 240--250°C & 100--110°C &
Requires ventilation; enclosed chamber recommended \\
\end{longtable}
}

\paragraph*{Support Settings Guide}\label{docs__pandoc__latex__src__3dmake_all_lessons.md__support-settings-guide-1}

{\def\LTcaptype{none} % do not increment counter
\begin{longtable}[]{@{}
  >{\centering\arraybackslash}p{(\linewidth - 4\tabcolsep) * \real{0.2029}}
  >{\centering\arraybackslash}p{(\linewidth - 4\tabcolsep) * \real{0.2319}}
  >{\raggedright\arraybackslash}p{(\linewidth - 4\tabcolsep) * \real{0.5652}}@{}}
\toprule\noalign{}
\begin{minipage}[b]{\linewidth}\centering
Overhang Angle
\end{minipage} & \begin{minipage}[b]{\linewidth}\centering
Supports Needed?
\end{minipage} & \begin{minipage}[b]{\linewidth}\raggedright
Recommendation
\end{minipage} \\
\midrule\noalign{}
\endhead
\bottomrule\noalign{}
\endlastfoot
\textless{} 45° & No & None \\
45--60° & Maybe & Preview first; add if sagging \\
\textgreater{} 60° & Yes & Tree supports (Cura) --- easier to remove \\
Bridge \textless{} 20 mm & No & Bridges usually fine \\
Bridge \textgreater{} 20 mm & Maybe & Preview; consider reorienting \\
\end{longtable}
}

\paragraph*{Common Problems and Quick Fixes}\label{docs__pandoc__latex__src__3dmake_all_lessons.md__common-problems-and-quick-fixes-1}

{\def\LTcaptype{none} % do not increment counter
\begin{longtable}[]{@{}
  >{\raggedright\arraybackslash}p{(\linewidth - 4\tabcolsep) * \real{0.3010}}
  >{\raggedright\arraybackslash}p{(\linewidth - 4\tabcolsep) * \real{0.3301}}
  >{\raggedright\arraybackslash}p{(\linewidth - 4\tabcolsep) * \real{0.3689}}@{}}
\toprule\noalign{}
\begin{minipage}[b]{\linewidth}\raggedright
Problem
\end{minipage} & \begin{minipage}[b]{\linewidth}\raggedright
Likely Cause
\end{minipage} & \begin{minipage}[b]{\linewidth}\raggedright
Fix
\end{minipage} \\
\midrule\noalign{}
\endhead
\bottomrule\noalign{}
\endlastfoot
Print not adhering to steel bed & Poor leveling / worn tape &
Re-level; refresh painter\textquotesingle s tape \\
Filament oozing on travel & Temperature too high &
Lower temp 5°C; check Z-hop setting \\
Rough bottom layer & Bed too close & Adjust leveling knob \\
Holes too small & FDM tolerance --- always undersized &
Add 0.2--0.3 mm to hole diameter \\
Part broke at layer boundary & Load perpendicular to layers &
Reorient so load is parallel to layers \\
Clogs at nozzle & Moisture in filament / debris &
Dry filament; clear nozzle with needle \\
Spaghetti / complete failure & Bed not leveled or supports missed &
Re-level; add supports or reorient \\
\end{longtable}
}

\paragraph*{G-code Export Checklist (Cura)}\label{docs__pandoc__latex__src__3dmake_all_lessons.md__g-code-export-checklist-cura}

\begin{itemize}
\tightlist
\item[$\square$]
  Anycubic i3 Mega printer profile selected
\item[$\square$]
  Correct filament type and diameter (1.75 mm)
\item[$\square$]
  Layer height set appropriately
\item[$\square$]
  Infill percentage set
\item[$\square$]
  Supports enabled if needed (tree supports preferred)
\item[$\square$]
  Nozzle and bed temperatures verified
\item[$\square$]
  Print preview checked (no floating parts)
\item[$\square$]
  Estimated time and filament weight noted
\end{itemize}

\emph{Sources: Anycubic official documentation (anycubic.com); Ultimaker Cura guide (ultimaker.com)}

\begin{center}\rule{0.5\linewidth}{0.5pt}\end{center}

\subsubsection*{D3 --- Bambu Studio / X1 Series}\label{docs__pandoc__latex__src__3dmake_all_lessons.md__d3--bambu-studio--x1-series}

\paragraph*{Recommended Settings by Use Case}\label{docs__pandoc__latex__src__3dmake_all_lessons.md__recommended-settings-by-use-case-2}

{\def\LTcaptype{none} % do not increment counter
\begin{longtable}[]{@{}
  >{\raggedright\arraybackslash}p{(\linewidth - 8\tabcolsep) * \real{0.2472}}
  >{\centering\arraybackslash}p{(\linewidth - 8\tabcolsep) * \real{0.1348}}
  >{\centering\arraybackslash}p{(\linewidth - 8\tabcolsep) * \real{0.0674}}
  >{\centering\arraybackslash}p{(\linewidth - 8\tabcolsep) * \real{0.1011}}
  >{\raggedright\arraybackslash}p{(\linewidth - 8\tabcolsep) * \real{0.4494}}@{}}
\toprule\noalign{}
\begin{minipage}[b]{\linewidth}\raggedright
Use Case
\end{minipage} & \begin{minipage}[b]{\linewidth}\centering
Layer Height
\end{minipage} & \begin{minipage}[b]{\linewidth}\centering
Infill
\end{minipage} & \begin{minipage}[b]{\linewidth}\centering
Supports
\end{minipage} & \begin{minipage}[b]{\linewidth}\raggedright
Notes
\end{minipage} \\
\midrule\noalign{}
\endhead
\bottomrule\noalign{}
\endlastfoot
Quick test / prototype & 0.30 mm & 10\% & As needed &
Draft --- use Standard mode \\
Standard project & 0.20 mm & 15--20\% & As needed &
Best all-around; Bambu default-optimized \\
Functional part & 0.20 mm & 30--40\% & As needed &
Use for load-bearing parts \\
Fine detail / display & 0.15 mm & 15\% & As needed &
Bambu achieves excellent detail at speed \\
Solid reference part & 0.20 mm & 40--50\% & Rarely &
Modern infill usually sufficient \\
\end{longtable}
}

\paragraph*{Filament Temperature Settings}\label{docs__pandoc__latex__src__3dmake_all_lessons.md__filament-temperature-settings-2}

{\def\LTcaptype{none} % do not increment counter
\begin{longtable}[]{@{}
  >{\raggedright\arraybackslash}p{(\linewidth - 6\tabcolsep) * \real{0.1667}}
  >{\centering\arraybackslash}p{(\linewidth - 6\tabcolsep) * \real{0.1310}}
  >{\centering\arraybackslash}p{(\linewidth - 6\tabcolsep) * \real{0.1071}}
  >{\raggedright\arraybackslash}p{(\linewidth - 6\tabcolsep) * \real{0.5952}}@{}}
\toprule\noalign{}
\begin{minipage}[b]{\linewidth}\raggedright
Filament
\end{minipage} & \begin{minipage}[b]{\linewidth}\centering
Nozzle Temp
\end{minipage} & \begin{minipage}[b]{\linewidth}\centering
Bed Temp
\end{minipage} & \begin{minipage}[b]{\linewidth}\raggedright
Notes
\end{minipage} \\
\midrule\noalign{}
\endhead
\bottomrule\noalign{}
\endlastfoot
PLA (standard) & 200--210°C & 50--60°C &
Fastest print times; Bambu-optimized profiles \\
PETG & 230--250°C & 70--80°C &
Slight cooling fan reduction for layer adhesion \\
TPU & 220--240°C & 20--30°C &
Best quality with Bambu\textquotesingle s tuned retraction \\
ABS & 240--260°C & 100--110°C &
Requires chamber heating; enclosed AMS recommended \\
\end{longtable}
}

\emph{Check Bambu\textquotesingle s material database for current profile updates --- these are dynamically optimized.}

\paragraph*{Support Settings Guide}\label{docs__pandoc__latex__src__3dmake_all_lessons.md__support-settings-guide-2}

{\def\LTcaptype{none} % do not increment counter
\begin{longtable}[]{@{}
  >{\centering\arraybackslash}p{(\linewidth - 4\tabcolsep) * \real{0.1795}}
  >{\centering\arraybackslash}p{(\linewidth - 4\tabcolsep) * \real{0.2051}}
  >{\raggedright\arraybackslash}p{(\linewidth - 4\tabcolsep) * \real{0.6154}}@{}}
\toprule\noalign{}
\begin{minipage}[b]{\linewidth}\centering
Overhang Angle
\end{minipage} & \begin{minipage}[b]{\linewidth}\centering
Supports Needed?
\end{minipage} & \begin{minipage}[b]{\linewidth}\raggedright
Recommendation
\end{minipage} \\
\midrule\noalign{}
\endhead
\bottomrule\noalign{}
\endlastfoot
\textless{} 45° & No & None \\
45--60° & Maybe & Auto-supports may activate; accept defaults \\
\textgreater{} 60° & Yes &
Bambu auto-tree supports; minimal removal effort \\
Bridge \textless{} 20 mm & No &
Excellent bridging on Bambu --- no supports needed \\
Bridge \textgreater{} 20 mm & Maybe &
Often bridgeable; preview before confirming \\
\end{longtable}
}

\paragraph*{Common Problems and Quick Fixes}\label{docs__pandoc__latex__src__3dmake_all_lessons.md__common-problems-and-quick-fixes-2}

{\def\LTcaptype{none} % do not increment counter
\begin{longtable}[]{@{}
  >{\raggedright\arraybackslash}p{(\linewidth - 4\tabcolsep) * \real{0.2925}}
  >{\raggedright\arraybackslash}p{(\linewidth - 4\tabcolsep) * \real{0.3208}}
  >{\raggedright\arraybackslash}p{(\linewidth - 4\tabcolsep) * \real{0.3868}}@{}}
\toprule\noalign{}
\begin{minipage}[b]{\linewidth}\raggedright
Problem
\end{minipage} & \begin{minipage}[b]{\linewidth}\raggedright
Likely Cause
\end{minipage} & \begin{minipage}[b]{\linewidth}\raggedright
Fix
\end{minipage} \\
\midrule\noalign{}
\endhead
\bottomrule\noalign{}
\endlastfoot
Nozzle clogging mid-print & Thermal runaway / wet filament &
Dry filament in AMS; check nozzle temp \\
Layer splitting & Wrong material in AMS slot &
Verify correct material in Bambu Studio \\
Stringing / blobbing & Cooling fan too low &
Increase cooling; adjust material profile \\
Print lifts off bed & Bed leveling drift (rare on Bambu) &
Auto-calibrate via Bambu Studio \\
Holes undersized & FDM tolerance --- always undersized &
Add 0.3--0.4 mm to hole diameter \\
Part broke at layer boundary & Load perpendicular to layers &
Reorient so load is parallel to layers \\
Print aborts unexpectedly & AMS issue / filament jam &
Clear AMS; verify filament path \\
Multi-material print misaligned & Nozzle offset miscalibration &
Run auto-calibration in maintenance menu \\
\end{longtable}
}

\paragraph*{G-code Export Checklist (Bambu Studio)}\label{docs__pandoc__latex__src__3dmake_all_lessons.md__g-code-export-checklist-bambu-studio}

\begin{itemize}
\tightlist
\item[$\square$]
  Correct printer (X1 / X1E) and hotend selected
\item[$\square$]
  Filament material and color match actual AMS configuration
\item[$\square$]
  Layer height optimal for desired quality
\item[$\square$]
  Infill and supports set appropriately
\item[$\square$]
  Print preview shows correct color layering (if multi-material)
\item[$\square$]
  Estimated time and filament weight acceptable
\item[$\square$]
  Camera enabled for remote monitoring
\item[$\square$]
  Web upload or USB export ready
\end{itemize}

\emph{Sources: Bambu Lab documentation (bambulab.com); Bambu Lab Community forums (community.bambulab.com)}
-e

\begin{center}\rule{0.5\linewidth}{0.5pt}\end{center}

\section*{Lesson 7 --- Parametric Transforms and the Phone Stand Project}\label{docs__pandoc__latex__src__3dmake_all_lessons.md__lesson-7--parametric-transforms-and-the-phone-stand-project}

\textbf{Estimated Time:} 90--120 minutes\\
\textbf{Course:} 3dMake Certification --- High School Distance Learning Track

\begin{center}\rule{0.5\linewidth}{0.5pt}\end{center}

\subsection*{Before You Start}\label{docs__pandoc__latex__src__3dmake_all_lessons.md__before-you-start-5}

You can now create shapes, combine them with CSG, organize them into modules, and add text. This lesson is about the next layer of precision: \emph{moving, rotating, mirroring, and scaling} objects with mathematical exactness.

You\textquotesingle ve used \texttt{translate()} before to move shapes around. But there\textquotesingle s a subtlety about how OpenSCAD applies transforms --- specifically, the \emph{order} in which they execute --- that is one of the most common sources of confusing bugs in OpenSCAD designs. Understanding the order rule deeply is one of the most valuable things you\textquotesingle ll get from this lesson.

This lesson also builds your first complete, multi-module project: a parametric phone stand. It\textquotesingle s not just an exercise --- it\textquotesingle s a real, functional object. You\textquotesingle ll measure your actual phone, enter those measurements as parameters, and the math will generate the exact geometry needed to hold your phone at the angle you specify. Change the angle parameter and the entire stand recalculates.

\begin{center}\rule{0.5\linewidth}{0.5pt}\end{center}

\subsection*{Learning Objectives}\label{docs__pandoc__latex__src__3dmake_all_lessons.md__learning-objectives-11}

By the end of this lesson you will be able to:

\begin{itemize}
\tightlist
\item
  Apply \texttt{translate()}, \texttt{rotate()}, and \texttt{scale()} in the correct order for compound transforms
\item
  Use trigonometric functions (\texttt{sin()}, \texttt{cos()}, \texttt{atan2()}) to compute precise geometry positions
\item
  Use \texttt{minkowski()} to create rounded edges and \texttt{mirror()} to create perfectly symmetric parts
\item
  Build a complete parametric phone stand from three measured dimensions
\item
  Use \texttt{\$preview} to conditionally speed up rendering during development
\item
  Apply vector math functions: \texttt{norm()}, \texttt{cross()}, and the dot product
\end{itemize}

\begin{center}\rule{0.5\linewidth}{0.5pt}\end{center}

\subsection*{Concept 1: Transform Order --- The Most Important Rule in This Lesson}\label{docs__pandoc__latex__src__3dmake_all_lessons.md__concept-1-transform-order--the-most-important-rule-in-this-lesson}

\subsubsection*{How OpenSCAD Applies Multiple Transforms}\label{docs__pandoc__latex__src__3dmake_all_lessons.md__how-openscad-applies-multiple-transforms}

When you write multiple transforms on the same line, OpenSCAD applies them in \textbf{right-to-left order} --- the transform closest to the geometry object applies first.

Read this line:

\begin{lstlisting}[style=Alabaster, language=openscad]
translate([10, 0, 0]) rotate([0, 0, 45]) cube(5);


\end{lstlisting}

Your eye reads left to right: "translate, then rotate." But OpenSCAD executes right to left: first \texttt{rotate} (it\textquotesingle s closer to \texttt{cube}), then \texttt{translate}. So the cube is first rotated 45° around the origin, and \emph{then} moved 10 mm in X.

Now read this:

\begin{lstlisting}[style=Alabaster, language=openscad]
rotate([0, 0, 45]) translate([10, 0, 0]) cube(5);


\end{lstlisting}

This time: first \texttt{translate} (cube moves 10 mm right), then \texttt{rotate} (the cube, now positioned at X=10, rotates around the origin). The result is that the cube ends up in a completely different location --- it orbits around the origin rather than rotating in place.

\subsubsection*{Why This Matters}\label{docs__pandoc__latex__src__3dmake_all_lessons.md__why-this-matters}

This rule has a real consequence: \textbf{changing the order of transforms changes the position of your geometry}, even if the transforms themselves are identical. This is the source of a lot of frustrating bugs in OpenSCAD designs, especially when the code looks correct to you.

The way to internalize it: think of it as nested functions. In \texttt{translate(rotate(cube()))}, \texttt{cube()} runs first, then \texttt{rotate()} acts on its output, then \texttt{translate()} acts on that result. The same way \texttt{f(g(x))} in math means "apply g first, then f."

\begin{lstlisting}[style=Alabaster, language=openscad]
// Example A: cube rotates 45° in place, then moves 10mm right

// (rotate happens to the stationary cube, then translate moves it)

translate([10, 0, 0]) rotate([0, 0, 45]) cube(5);


// Example B: cube moves 10mm right first, then the rotated version orbits the origin

// (translate moves it first, then rotate spins it around origin — completely different result)

rotate([0, 0, 45]) translate([10, 0, 0]) cube(5);


\end{lstlisting}

\textbf{Build both of these side by side.} Look at where each cube ends up. The visual difference makes this rule concrete in a way that reading about it can\textquotesingle t.

\subsubsection*{The Rule in Plain Language}\label{docs__pandoc__latex__src__3dmake_all_lessons.md__the-rule-in-plain-language}

"The last transform written (closest to the shape) runs first." Repeat this until it\textquotesingle s automatic.

\begin{center}\rule{0.5\linewidth}{0.5pt}\end{center}

\subsection*{Step 1 --- Translate, Rotate, Scale}\label{docs__pandoc__latex__src__3dmake_all_lessons.md__step-1--translate-rotate-scale}

\subsubsection*{\texorpdfstring{\texttt{translate({[}x,\ y,\ z{]})} --- Move}{translate({[}x, y, z{]}) --- Move}}\label{docs__pandoc__latex__src__3dmake_all_lessons.md__translatex-y-z--move}

\texttt{translate()} moves its child object by the specified amount in each axis. \texttt{translate({[}10,\ 0,\ 0{]})} moves 10 mm to the right. \texttt{translate({[}0,\ 0,\ 5{]})} moves 5 mm up. Negative values move in the opposite direction.

You already know this one from previous lessons. The key reminder is the order rule above.

\subsubsection*{\texorpdfstring{\texttt{rotate({[}x\_deg,\ y\_deg,\ z\_deg{]})} --- Rotate}{rotate({[}x\_deg, y\_deg, z\_deg{]}) --- Rotate}}\label{docs__pandoc__latex__src__3dmake_all_lessons.md__rotatex_deg-y_deg-z_deg--rotate}

\texttt{rotate()} rotates its child around the X, Y, and/or Z axes by the specified angles in degrees. The rotation happens around the origin --- if the object isn\textquotesingle t at the origin, it will orbit around it (which is sometimes what you want, sometimes not).

\texttt{rotate({[}0,\ 0,\ 45{]})} rotates around the Z axis (spins in the XY plane).
\texttt{rotate({[}90,\ 0,\ 0{]})} rotates around the X axis (tips the object forward/backward).
\texttt{rotate({[}0,\ 90,\ 0{]})} rotates around the Y axis (tips the object left/right).

When you want to rotate an object around its own center (not the origin), translate it to the origin first, rotate it, then translate it to its final position:

\begin{lstlisting}[style=Alabaster, language=openscad]
// Rotate a cylinder around its own center, then position it

translate([20, 30, 0])        // 3. move to final position

  rotate([0, 90, 0])          // 2. rotate (runs before translate due to order rule)

    cylinder(h=20, r=5);      // 1. cylinder starts at origin


\end{lstlisting}

\subsubsection*{\texorpdfstring{\texttt{scale({[}x,\ y,\ z{]})} --- Stretch or Compress}{scale({[}x, y, z{]}) --- Stretch or Compress}}\label{docs__pandoc__latex__src__3dmake_all_lessons.md__scalex-y-z--stretch-or-compress}

\texttt{scale({[}1,\ 1,\ 2{]})} doubles the height of an object without changing its width or depth. \texttt{scale({[}0.5,\ 0.5,\ 0.5{]})} shrinks everything by half. \texttt{scale({[}1,\ 2,\ 1{]})} doubles depth only.

\texttt{scale()} is less common than the others because it changes proportions --- a sphere becomes an ellipsoid, a cube becomes a rectangular box. But it\textquotesingle s useful for quick dimension experiments or for stretching a symmetrical part in one direction.

\begin{lstlisting}[style=Alabaster, language=openscad]
// Stretch a sphere 2x in Z to create an egg shape

scale([1, 1, 2]) sphere(r=10, $fn=32);


\end{lstlisting}

\begin{center}\rule{0.5\linewidth}{0.5pt}\end{center}

\subsection*{Step 2 --- Trigonometry for Geometric Positioning}\label{docs__pandoc__latex__src__3dmake_all_lessons.md__step-2--trigonometry-for-geometric-positioning}

\subsubsection*{Why You Need Trig in Design}\label{docs__pandoc__latex__src__3dmake_all_lessons.md__why-you-need-trig-in-design}

Trigonometry shows up in 3D design whenever you need to:

\begin{itemize}
\tightlist
\item
  Place objects evenly around a circle
\item
  Compute an angle from a measured slope
\item
  Position something at a known distance and direction from a reference point
\item
  Calculate how deep a feature needs to be at a given angle
\end{itemize}

Fortunately, you don\textquotesingle t need to understand all of trigonometry deeply --- you need three functions and one rule.

\textbf{The rule: all angles in OpenSCAD are in degrees.} Not radians. This is different from most programming languages. \texttt{sin(90)} in OpenSCAD returns 1 (correct for 90°). If you accidentally use radians, your numbers will be wildly off.

\subsubsection*{\texorpdfstring{\texttt{sin()}, \texttt{cos()}, and Placing Things Around a Circle}{sin(), cos(), and Placing Things Around a Circle}}\label{docs__pandoc__latex__src__3dmake_all_lessons.md__sin-cos-and-placing-things-around-a-circle}

\texttt{sin(angle)} and \texttt{cos(angle)} return the Y and X components of a unit vector pointing at that angle. Multiplying by a radius gives you the actual coordinates.

For a circle of radius \texttt{r}, the (x, y) position at angle \texttt{a} is:

\begin{itemize}
\tightlist
\item
  \texttt{x\ =\ cos(a)\ *\ r}
\item
  \texttt{y\ =\ sin(a)\ *\ r}
\end{itemize}

\begin{lstlisting}[style=Alabaster, language=openscad]
// Place 6 cylinders evenly around a circle of radius 25mm

// 360° / 6 = 60° apart

for (i = [0:5]) {

  angle = i * 60;

  translate([cos(angle) * 25, sin(angle) * 25, 0])

    cylinder(r=3, h=5, $fn=16);

}


\end{lstlisting}

\subsubsection*{\texorpdfstring{\texttt{atan2(y,\ x)} --- The Safe Angle Calculator}{atan2(y, x) --- The Safe Angle Calculator}}\label{docs__pandoc__latex__src__3dmake_all_lessons.md__atan2y-x--the-safe-angle-calculator}

When you need to go the other direction --- from (x, y) coordinates back to an angle --- use \texttt{atan2(y,\ x)}.

\texttt{atan(y/x)} seems like it would work, but it fails in two quadrants. When x is negative (the left half of the circle), the division produces a value that places the angle in the wrong half of the circle. \texttt{atan2(y,\ x)} takes both components separately and always returns the correct angle for all four quadrants.

\begin{lstlisting}[style=Alabaster, language=openscad]
dx = -5;  dy = 5;

angle_correct = atan2(dy, dx);   // 135° — correct (upper-left quadrant)

angle_wrong   = atan(dy / dx);   // -45° — wrong (treats it as lower-right)

echo("atan2:", angle_correct, "  atan:", angle_wrong);


\end{lstlisting}

Always use \texttt{atan2()} when computing angles from coordinate pairs. Make this a habit.

\begin{center}\rule{0.5\linewidth}{0.5pt}\end{center}

\subsection*{\texorpdfstring{Step 3 --- \texttt{minkowski()} for Rounded Edges}{Step 3 --- minkowski() for Rounded Edges}}\label{docs__pandoc__latex__src__3dmake_all_lessons.md__step-3--minkowski-for-rounded-edges}

\subsubsection*{What the Minkowski Sum Does}\label{docs__pandoc__latex__src__3dmake_all_lessons.md__what-the-minkowski-sum-does}

\texttt{minkowski()} computes something called the \textbf{Minkowski sum} of its children. The geometric meaning: it "rolls" the second shape over every point on the surface of the first shape, and the result is all the space swept out by this rolling.

In practice, the most useful application is rounding the edges and corners of a box. If you take the Minkowski sum of a cube and a small sphere, you get a box with perfectly rounded corners and edges --- the rounding radius equals the sphere\textquotesingle s radius.

The variation used most often for printable parts is to use a \textbf{cylinder} instead of a sphere. A very flat cylinder (nearly zero height) rounds the vertical edges of a box but keeps the bottom face flat. This is ideal for parts that will sit on a surface, because a sphere-Minkowski box has a rounded bottom that makes it rock.

\begin{lstlisting}[style=Alabaster, language=openscad]
// Rounded box — sphere rounds all corners and edges including bottom

module rounded_cube(w, d, h, r=3) {

  minkowski() {

    cube([w - 2*r, d - 2*r, h - r], center=false);

    sphere(r=r, $fn=16);

  }

}


// Flat-base rounded box — cylinder rounds edges but keeps bottom flat

// The w/d dimensions shrink by 2r, and the height stays the same

module flat_rounded_cube(w, d, h, r=3) {

  minkowski() {

    cube([w - 2*r, d - 2*r, h], center=false);

    cylinder(r=r, h=0.01, $fn=24);   // nearly zero height: rounds edges, not top/bottom

  }

}


rounded_cube(40, 30, 20, r=4);

translate([50, 0, 0]) flat_rounded_cube(40, 30, 20, r=4);


\end{lstlisting}

\subsubsection*{Performance Warning}\label{docs__pandoc__latex__src__3dmake_all_lessons.md__performance-warning}

\texttt{minkowski()} is computationally expensive. With high \texttt{\$fn} values on the secondary shape, it can take several minutes to compute. During development, keep the secondary shape at \texttt{\$fn=12–16}. Only use \texttt{\$fn=32+} for final export. Even better: use the \texttt{\$preview} technique in Step 5 to skip Minkowski entirely during quick iterations.

\begin{center}\rule{0.5\linewidth}{0.5pt}\end{center}

\subsection*{\texorpdfstring{Step 4 --- \texttt{mirror()} for Symmetry}{Step 4 --- mirror() for Symmetry}}\label{docs__pandoc__latex__src__3dmake_all_lessons.md__step-4--mirror-for-symmetry}

\subsubsection*{Design One Half, Mirror the Other}\label{docs__pandoc__latex__src__3dmake_all_lessons.md__design-one-half-mirror-the-other}

Many designs have a line of symmetry --- two legs, two wings, two ears, two cable clips. Instead of designing both halves separately (which introduces the risk of slight misalignment), design one half and mirror it. The mirrored copy is guaranteed to be exactly symmetric.

\texttt{mirror({[}1,\ 0,\ 0{]})} reflects across the YZ plane --- negates X coordinates.
\texttt{mirror({[}0,\ 1,\ 0{]})} reflects across the XZ plane --- negates Y coordinates.
\texttt{mirror({[}0,\ 0,\ 1{]})} reflects across the XY plane --- negates Z coordinates.

\begin{lstlisting}[style=Alabaster, language=openscad]
// Design one ear, mirror it to get a symmetric pair

module ear_peg() {

  translate([20, 0, 0]) cylinder(r=4, h=8, $fn=32);

}


ear_peg();                     // right ear

mirror([1, 0, 0]) ear_peg();  // left ear — exact mirror image


// Mirror works for complex modules too

module mounting_clip() {

  difference() {

    cube([15, 8, 10]);

    translate([7.5, -0.001, 2])

      cylinder(r=3, h=8.002, $fn=32);  // cable channel

  }

}


translate([0, 20, 0]) {

  mounting_clip();

  translate([15, 0, 0]) mirror([1, 0, 0]) mounting_clip();  // symmetric partner

}


\end{lstlisting}

\begin{center}\rule{0.5\linewidth}{0.5pt}\end{center}

\subsection*{Step 5 --- Build the Complete Phone Stand}\label{docs__pandoc__latex__src__3dmake_all_lessons.md__step-5--build-the-complete-phone-stand}

\subsubsection*{Why Build This}\label{docs__pandoc__latex__src__3dmake_all_lessons.md__why-build-this}

The phone stand is your first complete multi-module project. It\textquotesingle s functional, parametric, and applies trigonometry to solve a real design problem. More importantly, it\textquotesingle s something you\textquotesingle ll actually use --- which is a good test of your own design work.

\subsubsection*{Measuring Your Phone}\label{docs__pandoc__latex__src__3dmake_all_lessons.md__measuring-your-phone}

Before writing code, use a ruler or calipers to measure:

\begin{itemize}
\tightlist
\item
  \textbf{Phone width}: the dimension across the face (X dimension in the cradle)
\item
  \textbf{Phone thickness}: how thick the phone is (determines cradle depth)
\item
  \textbf{Preferred angle}: how far you want the phone to lean back from vertical (65° from horizontal is a comfortable default for most people)
\end{itemize}

\subsubsection*{The Full Parametric Design}\label{docs__pandoc__latex__src__3dmake_all_lessons.md__the-full-parametric-design}

\begin{lstlisting}[style=Alabaster, language=openscad]
// ============================================================

// Parametric Phone Stand

// ============================================================

// Measure your phone and enter the three values below.

// All other dimensions are computed automatically.

//

// PARAMETERS TO CUSTOMIZE:

// - phone_w: width of your phone in mm

// - phone_d: thickness of your phone in mm

// - angle:   viewing angle from horizontal (60–80 degrees typically)

// ============================================================


phone_w = 75;   // mm — phone width (fits in cradle)

phone_d = 9;    // mm — phone thickness (sets cradle clearance)

angle   = 65;   // degrees — viewing angle from horizontal


// Internal dimensions — computed from the parameters above

lip_h       = 15;   // mm — depth of cradle lip

base_h      = 5;    // mm — base plate thickness

cradle_wall = 3;    // mm — wall thickness around phone

r_fillet    = 3;    // mm — rounded edge radius


// Trigonometry: compute base depth from the viewing angle

// The phone support must be at the correct lean angle

base_d = (phone_d + 10) / cos(90 - angle);


module flat_rounded_cube(w, d, h, r=3) {

  // Guard against degenerate dimensions after subtracting rounding radius

  assert(w > 2*r && d > 2*r, "Dimensions too small for the rounding radius");

  minkowski() {

    cube([w - 2*r, d - 2*r, h]);

    cylinder(r=r, h=0.01, $fn=24);

  }

}


module base_plate() {

  // Wide, flat base that anchors the stand

  flat_rounded_cube(phone_w + 20, base_d + 10, base_h, r_fillet);

}


module back_support() {

  // Angled wall that supports the back of the phone

  // rotate([angle-90, 0, 0]) tilts the wall to the viewing angle

  rotate([angle - 90, 0, 0])

    flat_rounded_cube(phone_w + 6, cradle_wall, 60, r_fillet);

}


module lip() {

  // Bottom ledge that the phone rests on

  rotate([angle - 90, 0, 0])

    translate([0, -lip_h, 0])

      flat_rounded_cube(phone_w + 6, lip_h + cradle_wall, cradle_wall, r_fillet);

}


// Final assembly: base + back support + lip

base_plate();

translate([0, 0, base_h]) {

  back_support();

  lip();

}


\end{lstlisting}

\subsubsection*{Understanding the Trigonometry}\label{docs__pandoc__latex__src__3dmake_all_lessons.md__understanding-the-trigonometry}

The rotation \texttt{rotate({[}angle\ -\ 90,\ 0,\ 0{]})} deserves an explanation.

In OpenSCAD, a rotation of 0° around the X axis means the object is vertical (standing straight up). A rotation of -90° tips it flat forward (horizontal). When we want the back support to lean at \texttt{angle} degrees from horizontal, we need a rotation of \texttt{angle\ -\ 90} degrees:

\begin{itemize}
\tightlist
\item
  If \texttt{angle\ =\ 90} (vertical phone): \texttt{rotate({[}0,\ 0,\ 0{]})} --- no rotation, support is vertical
\item
  If \texttt{angle\ =\ 65} (typical lean): \texttt{rotate({[}-25,\ 0,\ 0{]})} --- leans 25° back
\item
  If \texttt{angle\ =\ 45} (steeper lean): \texttt{rotate({[}-45,\ 0,\ 0{]})} --- leans 45° back
\end{itemize}

After building the stand, try changing \texttt{angle} from 45 to 80 in 5-degree steps and rebuild each time. Watch the support geometry recalculate automatically for every value.

\begin{center}\rule{0.5\linewidth}{0.5pt}\end{center}

\subsection*{\texorpdfstring{Step 6 --- Adaptive Quality with \texttt{\$preview}}{Step 6 --- Adaptive Quality with \$preview}}\label{docs__pandoc__latex__src__3dmake_all_lessons.md__step-6--adaptive-quality-with-preview}

\subsubsection*{Fast During Development, High-Quality for Export}\label{docs__pandoc__latex__src__3dmake_all_lessons.md__fast-during-development-high-quality-for-export}

\texttt{\$preview} is a special OpenSCAD variable that is automatically \texttt{true} during F5 preview (the fast OpenSCAD preview) and \texttt{false} during F6 full render and \texttt{3dm\ build}. You can use it to conditionally use lower-quality (faster) settings during development and automatically switch to high quality for final export.

\begin{lstlisting}[style=Alabaster, language=openscad]
// Low $fn during preview for fast renders; high $fn for final build

$fn = $preview ? 16 : 64;


sphere(r=20);  // uses $fn=16 in preview, $fn=64 during 3dm build


\end{lstlisting}

This is especially valuable with \texttt{minkowski()}, which can be very slow at high \texttt{\$fn}. During development, replace the Minkowski operation with a plain box placeholder:

\begin{lstlisting}[style=Alabaster, language=openscad]
module rounded_box_adaptive(w, d, h, r=4) {

  if ($preview) {

    // Fast placeholder during development

    color("orange", 0.6) cube([w, d, h]);

  } else {

    // Full, accurate Minkowski sum for final export

    minkowski() {

      cube([w - 2*r, d - 2*r, h - r]);

      sphere(r=r, $fn=12);

    }

  }

}


\end{lstlisting}

\begin{center}\rule{0.5\linewidth}{0.5pt}\end{center}

\subsection*{Concept 2: Vector Math Functions}\label{docs__pandoc__latex__src__3dmake_all_lessons.md__concept-2-vector-math-functions}

\subsubsection*{Three Functions Worth Knowing}\label{docs__pandoc__latex__src__3dmake_all_lessons.md__three-functions-worth-knowing}

\textbf{\texttt{norm(v)}} --- returns the length (magnitude) of a vector. \texttt{norm({[}3,\ 4,\ 0{]})} = 5. This is the Euclidean distance from the origin to the point (3, 4, 0), and it follows from the Pythagorean theorem: √(3² + 4² + 0²) = 5.

\textbf{Dot product (\texttt{a\ *\ b})} --- returns a single number that measures how "aligned" two vectors are. If both vectors point in exactly the same direction, the dot product equals the product of their lengths. If they\textquotesingle re perpendicular, it equals 0. If they point in opposite directions, it\textquotesingle s negative. Useful for checking angles between directions.

\textbf{\texttt{cross(a,\ b)}} --- returns a vector perpendicular to both input vectors. \texttt{cross({[}1,0,0{]},\ {[}0,1,0{]})} returns \texttt{{[}0,0,1{]}} --- the Z axis is perpendicular to both X and Y. Useful for computing surface normals or finding the axis of rotation between two vectors.

\begin{lstlisting}[style=Alabaster, language=openscad]
// norm: length of a vector

v = [3, 4, 0];

echo("Length:", norm(v));   // 5


// dot product: alignment measure

a = [1, 0, 0];   // X axis

b = [0.7, 0.7, 0];

echo("Dot product:", a * b);   // 0.7


// cross product: vector perpendicular to both

c = cross([1, 0, 0], [0, 1, 0]);

echo("Cross:", c);   // [0, 0, 1] — the Z axis


\end{lstlisting}

\begin{center}\rule{0.5\linewidth}{0.5pt}\end{center}

\subsection*{Quiz --- Lesson 7 (15 questions)}\label{docs__pandoc__latex__src__3dmake_all_lessons.md__quiz--lesson-7-15-questions}

\begin{enumerate}
\tightlist
\item
  In OpenSCAD, do transforms apply left-to-right or right-to-left (which transform runs first)?
\item
  What is the result of \texttt{cos(0)} and \texttt{sin(90)} in OpenSCAD (working in degrees)?
\item
  What is the key advantage of \texttt{atan2(y,\ x)} over \texttt{atan(y/x)} for computing angles?
\item
  What does \texttt{minkowski()} do geometrically? Describe it in plain language.
\item
  What does \texttt{mirror({[}1,\ 0,\ 0{]})} do? Which plane does it reflect across?
\item
  What does \texttt{scale({[}1,\ 1,\ 2{]})} do to a sphere?
\item
  Write OpenSCAD code to place 8 cylinders evenly spaced around a circle of radius 30 mm.
\item
  What does \texttt{norm({[}3,\ 4,\ 0{]})} return, and how is this computed?
\item
  True or False: \texttt{translate({[}10,0,0{]})\ rotate({[}0,0,45{]})\ cube(5)} rotates the cube first, then translates it.
\item
  What does setting \texttt{\$preview} to \texttt{true} mean in OpenSCAD, and when is it set automatically?
\item
  Why does using \texttt{minkowski()} with \texttt{cylinder(h=0.01)} instead of \texttt{sphere()} produce a flat bottom surface?
\item
  In the phone stand design, the back support uses \texttt{rotate({[}angle\ -\ 90,\ 0,\ 0{]})}. If \texttt{angle\ =\ 70}, what is the actual rotation applied? What does this mean physically?
\item
  What does \texttt{cross({[}1,0,0{]},\ {[}0,1,0{]})} return and what is the geometric meaning?
\item
  Explain why designing one half of a symmetric part and using \texttt{mirror()} is better than designing both halves separately.
\item
  A design requires placing 5 components at the vertices of a regular pentagon centered at the origin. Using \texttt{cos()} and \texttt{sin()}, what is the formula for the coordinates of each vertex, given a circumradius of 25 mm?
\end{enumerate}

\begin{center}\rule{0.5\linewidth}{0.5pt}\end{center}

\subsection*{Extension Problems (15)}\label{docs__pandoc__latex__src__3dmake_all_lessons.md__extension-problems-15-6}

\begin{enumerate}
\tightlist
\item
  Redesign your phone stand with a USB cable slot cut through the lip. Make the slot width and depth parameters.
\item
  Add rubber foot pockets to the base plate: four small rectangular cutouts on the bottom face, positioned at the corners.
\item
  Create a phone stand variant that holds the phone in landscape (horizontal) orientation. Document which parameters changed and why.
\item
  Build a radially symmetric decoration: 12 identical fins evenly spaced around a central cylinder using \texttt{for} loop and \texttt{rotate()}.
\item
  Use \texttt{atan2()} to compute the angle of a ramp and use that angle with \texttt{rotate()} to align a small object precisely on the ramp surface.
\item
  Design a phone stand for a large tablet (e.g., 180mm × 240mm). Update the parameters, rebuild, and verify the stand still holds at the correct angle.
\item
  Build a "compound arm" --- three rigid link segments connected end-to-end, each at a different angle. Display all three as an assembly with no gaps or overlaps.
\item
  Create an ergonomic keyboard wrist rest using \texttt{minkowski()} with a cylinder for smooth, flat-base contouring.
\item
  Build a visual transform order demo: two colored objects that differ only in the order of \texttt{translate()} and \texttt{rotate()}, placed side by side with labels.
\item
  Redesign the phone stand as a wall-mount: replace the base plate with a flat back plate that has two M3 mounting holes.
\item
  Build a mirror-symmetric earring holder: design one branch with pegs, then use \texttt{mirror()} to create the symmetric partner. Both branches should connect to a central hanging loop.
\item
  Use \texttt{norm()} to calculate the space diagonal of a rectangular prism (distance from one corner to the opposite corner) and use that value as a dimension in the design.
\item
  Research OpenSCAD\textquotesingle s \texttt{multmatrix()} function. Build a shear transformation example (slanting a box into a parallelogram-shaped prism) that can\textquotesingle t be achieved with standard transforms.
\item
  Add snap-fit clips to the phone stand cradle: two small flexible fingers that slightly grip the phone. Dimension the fingers using snap-fit principles from Lesson 8.
\item
  Design a dual phone stand: two cradles side by side at different angles, sharing a common base. Use a \texttt{for} loop and parametric spacing.
\end{enumerate}

\begin{center}\rule{0.5\linewidth}{0.5pt}\end{center}

\subsection*{References and Helpful Resources}\label{docs__pandoc__latex__src__3dmake_all_lessons.md__references-and-helpful-resources-6}

\begin{itemize}
\item
  OpenSCAD User Manual --- Transformations --- https://en.wikibooks.org/wiki/OpenSCAD\_User\_Manual/Transformations
\item
  OpenSCAD User Manual --- Mathematical Functions --- https://en.wikibooks.org/wiki/OpenSCAD\_User\_Manual/Mathematical\_Functions
\item
  OpenSCAD User Manual --- Minkowski and Hull --- https://en.wikibooks.org/wiki/OpenSCAD\_User\_Manual/Minkowski\_and\_Hull
\end{itemize}

\subsubsection*{Supplemental Resources}\label{docs__pandoc__latex__src__3dmake_all_lessons.md__supplemental-resources-6}

\begin{itemize}
\tightlist
\item
  \href{docs\%5Cpandoc\%5Clatex\%5Csrc\%5C943ee2446365181c.epub}{Programming with OpenSCAD EPUB Textbook} --- Transforms chapter
\item
  \href{https://github.com/ProgrammingWithOpenSCAD/CodeSolutions}{CodeSolutions Repository} --- Phone stand worked example
\item
  \href{https://programmingwithopenscad.github.io/quick-reference.html}{OpenSCAD Quick Reference} --- Math and transform functions
  -e
\end{itemize}

\begin{center}\rule{0.5\linewidth}{0.5pt}\end{center}

\section*{Lesson 8 --- Advanced Parametric Design and Interlocking Features}\label{docs__pandoc__latex__src__3dmake_all_lessons.md__lesson-8--advanced-parametric-design-and-interlocking-features}

\textbf{Estimated Time:} 90--120 minutes\\
\textbf{Course:} 3dMake Certification --- High School Distance Learning Track

\begin{center}\rule{0.5\linewidth}{0.5pt}\end{center}

\subsection*{Before You Start}\label{docs__pandoc__latex__src__3dmake_all_lessons.md__before-you-start-6}

Every design you\textquotesingle ve built so far has been a single-piece print. This lesson moves to \textbf{multi-part assemblies}: designs where two or more printed pieces must physically fit together. This is where the gap between "looks right on screen" and "works in real life" becomes most critical.

The central concept of this lesson is \textbf{tolerance}: the intentional gap between mating parts that makes them fit together correctly. Too tight and the pieces won\textquotesingle t go together at all. Too loose and they wobble, rattle, or fall apart. The exact right amount of clearance depends on your specific printer, your specific filament, and your slicer settings --- and the only reliable way to know it is to measure it.

This lesson gives you a systematic, professional approach: print a tolerance test coupon, measure the actual results, encode those measurements as named constants in your code, and build designs that work the first time instead of after three reprints.

\begin{center}\rule{0.5\linewidth}{0.5pt}\end{center}

\subsection*{Learning Objectives}\label{docs__pandoc__latex__src__3dmake_all_lessons.md__learning-objectives-12}

By the end of this lesson you will be able to:

\begin{itemize}
\tightlist
\item
  Explain why FDM-printed holes are always undersized and how to compensate
\item
  Design press-fit and slip-fit joints using printer-specific tolerance values
\item
  Print and use a tolerance test coupon to measure your printer\textquotesingle s actual clearances
\item
  Design snap-fit cantilever arms sized correctly for different materials
\item
  Design dovetail slide joints and stackable assemblies
\item
  Use threaded insert pockets for secure metal-to-plastic fastening
\item
  Apply true chamfer geometry
\item
  Apply tolerance values as named constants throughout a design
\end{itemize}

\begin{center}\rule{0.5\linewidth}{0.5pt}\end{center}

\subsection*{Concept 1: The Physics of FDM Tolerances}\label{docs__pandoc__latex__src__3dmake_all_lessons.md__concept-1-the-physics-of-fdm-tolerances}

\subsubsection*{Why Printed Holes Are Always Too Small}\label{docs__pandoc__latex__src__3dmake_all_lessons.md__why-printed-holes-are-always-too-small}

This is one of the most important physical facts about FDM printing, and it surprises almost every beginner the first time they encounter it.

When an FDM printer extrudes a bead of filament, that bead spreads slightly as it deposits --- it squishes outward from the center of the extrusion path. This happens at every wall, but the \emph{direction} of the squishing is important:

\begin{itemize}
\tightlist
\item
  On \textbf{outer walls} (the outside of a part), the extrusion squishes \emph{outward} from the wall\textquotesingle s nominal surface. The part comes out slightly larger than designed.
\item
  On \textbf{inner walls} (the inside of a hole), the extrusion squishes \emph{inward} --- toward the center of the hole. The hole comes out smaller than designed.
\end{itemize}

The result: a hole designed to be 10 mm in diameter prints as approximately 9.6--9.8 mm. A peg designed to be 10 mm prints as approximately 9.8--10.0 mm. Two parts with a 10 mm nominal dimension will not fit together --- the peg is too large for the hole.

This isn\textquotesingle t a flaw in your printer. It\textquotesingle s a fundamental consequence of how FDM extrusion works, and every designer working with FDM printing accounts for it.

\subsubsection*{What "Clearance" Means}\label{docs__pandoc__latex__src__3dmake_all_lessons.md__what-clearance-means}

\textbf{Clearance} is the intentional difference in size between a peg (or shaft, or tab) and its mating socket (or hole, or slot). You make the peg slightly smaller than the socket --- the amount of this size difference is the clearance.

Different clearances produce different fits:

{\def\LTcaptype{none} % do not increment counter
\begin{longtable}[]{@{}
  >{\raggedright\arraybackslash}p{(\linewidth - 4\tabcolsep) * \real{0.1810}}
  >{\raggedright\arraybackslash}p{(\linewidth - 4\tabcolsep) * \real{0.3190}}
  >{\raggedright\arraybackslash}p{(\linewidth - 4\tabcolsep) * \real{0.5000}}@{}}
\toprule\noalign{}
\begin{minipage}[b]{\linewidth}\raggedright
Fit Type
\end{minipage} & \begin{minipage}[b]{\linewidth}\raggedright
Clearance (typical PLA, 0.4mm nozzle)
\end{minipage} & \begin{minipage}[b]{\linewidth}\raggedright
Behavior
\end{minipage} \\
\midrule\noalign{}
\endhead
\bottomrule\noalign{}
\endlastfoot
Press-fit (permanent) & 0.0 -- 0.1 mm &
Requires force to assemble; doesn\textquotesingle t come apart easily \\
Slip-fit (removable) & 0.2 -- 0.3 mm &
Slides smoothly; can be assembled and disassembled by hand \\
Loose/clearance fit & 0.4 -- 0.5 mm &
Some play and wobble; appropriate for hinges and pivots \\
\end{longtable}
}

These are starting-point estimates. Your actual values will differ based on your printer and materials. That\textquotesingle s what the test coupon is for.

\begin{center}\rule{0.5\linewidth}{0.5pt}\end{center}

\subsection*{Step 1 --- The Tolerance Test Coupon}\label{docs__pandoc__latex__src__3dmake_all_lessons.md__step-1--the-tolerance-test-coupon}

\subsubsection*{Print This Before Every New Printer or Material}\label{docs__pandoc__latex__src__3dmake_all_lessons.md__print-this-before-every-new-printer-or-material}

The tolerance test coupon is a small, quick print that gives you empirical data about your specific printer\textquotesingle s actual clearances. Print it once with your standard material and settings. Test each peg-socket pair. Record the results. These numbers become the constants in your design files.

You\textquotesingle ll need to do this test if you:

\begin{itemize}
\tightlist
\item
  Use a new printer you haven\textquotesingle t printed on before
\item
  Switch to a significantly different filament (e.g., from PLA to PETG)
\item
  Change your nozzle size
\item
  Change your slicer settings significantly (especially layer height)
\end{itemize}

\begin{lstlisting}[style=Alabaster, language=openscad]
// Tolerance Test Coupon

// Prints a row of pegs and a row of matching sockets at five clearance values.

// After printing: test each peg in each socket. Record which clearances give

// press-fit, slip-fit, and loose fit.


nominal        = 10;     // mm — nominal peg diameter

test_h         = 8;      // mm — peg height

plate_t        = 3;      // mm — socket plate thickness

socket_d_outer = nominal + 8;   // mm — outer diameter of socket annulus


clearances = [0.0, 0.1, 0.2, 0.3, 0.4];   // mm — the five clearances to test


module test_peg(clearance) {

  // Peg diameter = nominal - clearance

  color("SteelBlue")

    cylinder(d=nominal - clearance, h=test_h, $fn=32);

}


module test_socket() {

  // All sockets have the same nominal hole; pegs vary

  color("Coral")

  difference() {

    cylinder(d=socket_d_outer, h=plate_t, $fn=32);

    translate([0, 0, -0.001])

      cylinder(d=nominal, h=plate_t + 0.002, $fn=32);

  }

}


module clearance_label(clearance) {

  // Small label showing the clearance value

  translate([0, -(nominal/2 + 5), 0])

    linear_extrude(1)

      text(str(clearance, "mm"), size=2.5, font="Liberation Sans",

           halign="center", $fn=4);

}


// Row 1: Pegs — each has a different clearance reduction applied

for (i = [0 : len(clearances) - 1]) {

  translate([i * (socket_d_outer + 4), 0, 0]) {

    test_peg(clearances[i]);

    clearance_label(clearances[i]);

  }

}


// Row 2: Sockets — all have the same nominal hole

translate([0, -(socket_d_outer + 8), 0])

  for (i = [0 : len(clearances) - 1]) {

    translate([i * (socket_d_outer + 4), 0, 0])

      test_socket();

  }


\end{lstlisting}

\subsubsection*{How to Read the Results}\label{docs__pandoc__latex__src__3dmake_all_lessons.md__how-to-read-the-results}

After printing, test each peg in each socket. Record:

{\def\LTcaptype{none} % do not increment counter
\begin{longtable}[]{@{}
  >{\raggedright\arraybackslash}p{(\linewidth - 4\tabcolsep) * \real{0.1444}}
  >{\raggedright\arraybackslash}p{(\linewidth - 4\tabcolsep) * \real{0.6333}}
  >{\raggedright\arraybackslash}p{(\linewidth - 4\tabcolsep) * \real{0.2222}}@{}}
\toprule\noalign{}
\begin{minipage}[b]{\linewidth}\raggedright
Fit Type
\end{minipage} & \begin{minipage}[b]{\linewidth}\raggedright
Test Method
\end{minipage} & \begin{minipage}[b]{\linewidth}\raggedright
What to Record
\end{minipage} \\
\midrule\noalign{}
\endhead
\bottomrule\noalign{}
\endlastfoot
Press-fit & Tightest peg that enters socket with firm thumb pressure &
Your press clearance \\
Slip-fit &
Loosest peg that doesn\textquotesingle t rattle in the socket when shaken
& Your slip clearance \\
Bearing/pivot & Loosest peg that rotates smoothly without visible wobble
& Slip + 0.05--0.1 mm \\
\end{longtable}
}

These measured values replace the rough estimates in the table above. Record them in a note or file, and use them for every design from now on.

\begin{center}\rule{0.5\linewidth}{0.5pt}\end{center}

\subsection*{Step 2 --- Encoding Tolerance as Named Constants}\label{docs__pandoc__latex__src__3dmake_all_lessons.md__step-2--encoding-tolerance-as-named-constants}

\subsubsection*{Never Hardcode Tolerance Values}\label{docs__pandoc__latex__src__3dmake_all_lessons.md__never-hardcode-tolerance-values}

Tolerance clearance values should appear in your design file exactly once --- at the top, as named constants. When you measure different clearances (because you\textquotesingle re on a different printer, or using PETG instead of PLA), you change the constant in one place and the entire design updates.

\begin{lstlisting}[style=Alabaster, language=openscad]
// ============================================================

// TOLERANCE CONSTANTS — measured from your test coupon print

// Run the test coupon, measure, and update these values.

// ============================================================

PRESS_FIT_C = 0.2;    // mm — tightest peg that enters with thumb pressure

SLIP_FIT_C  = 0.3;    // mm — loosest peg that doesn't rattle

SNAP_C      = PRESS_FIT_C;   // snap-fits use press-fit clearance


// Use these constants throughout the file:

// hole_r = peg_r - SLIP_FIT_C    (for a slip-fit hole)

// hole_r = peg_r - PRESS_FIT_C   (for a press-fit hole)


\end{lstlisting}

This is exactly the same DRY principle from Lesson 3 applied to physical measurements. The constants are derived from reality (your test print); the geometry uses those constants; change one measurement and everything that depends on it updates automatically.

\begin{center}\rule{0.5\linewidth}{0.5pt}\end{center}

\subsection*{Step 3 --- Snap-Fit Cantilever Arms}\label{docs__pandoc__latex__src__3dmake_all_lessons.md__step-3--snap-fit-cantilever-arms}

\subsubsection*{What a Snap-Fit Is}\label{docs__pandoc__latex__src__3dmake_all_lessons.md__what-a-snap-fit-is}

A snap-fit is a self-locking fastener made entirely from the printed plastic itself. A thin, flexible beam (the cantilever arm) deflects as the mating part pushes past it, then springs back to its original position, locking the parts together. Think of the plastic tab that holds a battery cover in place, or the clip that holds a phone case together.

Snap-fits are one of the most elegant features you can design. When they work well, they\textquotesingle re intuitive, require no separate hardware, and can be assembled and disassembled dozens of times without wearing out. When they\textquotesingle re designed wrong, they either crack (too stiff) or don\textquotesingle t hold (too flexible).

\subsubsection*{The Flexibility Rule}\label{docs__pandoc__latex__src__3dmake_all_lessons.md__the-flexibility-rule}

The key design parameter is the \textbf{beam flexibility ratio}: \texttt{beam\_thickness\ /\ beam\_length}. For PLA, this ratio must be \textbf{0.15 or less} --- meaning the beam\textquotesingle s thickness must be at most 15\% of its length. A beam that\textquotesingle s too thick relative to its length won\textquotesingle t flex and will crack instead of deflecting.

Different materials allow different ratios because they have different stiffness (elastic modulus):

{\def\LTcaptype{none} % do not increment counter
\begin{longtable}[]{@{}
  >{\raggedright\arraybackslash}p{(\linewidth - 4\tabcolsep) * \real{0.0988}}
  >{\raggedright\arraybackslash}p{(\linewidth - 4\tabcolsep) * \real{0.2840}}
  >{\raggedright\arraybackslash}p{(\linewidth - 4\tabcolsep) * \real{0.6173}}@{}}
\toprule\noalign{}
\begin{minipage}[b]{\linewidth}\raggedright
Material
\end{minipage} & \begin{minipage}[b]{\linewidth}\raggedright
Maximum beam\_t / beam\_l
\end{minipage} & \begin{minipage}[b]{\linewidth}\raggedright
Why
\end{minipage} \\
\midrule\noalign{}
\endhead
\bottomrule\noalign{}
\endlastfoot
PLA & 0.15 & Relatively stiff; snaps or cracks if over-flexed \\
PETG & 0.18 & Slightly more flexible than PLA \\
Nylon & 0.20 & Good flexibility and fatigue resistance \\
TPU & 0.25 & Very flexible; excellent for repeated flex \\
ABS & 0.12 & More brittle than PLA; needs shorter/thinner beams \\
\end{longtable}
}

\begin{lstlisting}[style=Alabaster, language=openscad]
module snap_arm(beam_l=20, beam_t=1.8, beam_w=8, hook_h=1.2, hook_angle=30,

                material="PLA") {


  // Maximum flexibility ratio by material

  max_ratio = (material == "TPU")  ? 0.25 :

              (material == "PETG") ? 0.18 :

              (material == "Nylon")? 0.20 :

              (material == "ABS")  ? 0.12 :

              0.15;  // default: PLA


  ratio = beam_t / beam_l;


  // Validate before any geometry is created

  assert(ratio <= max_ratio,

    str("snap_arm: beam_t/beam_l ratio (", ratio,

        ") exceeds maximum for ", material, " (", max_ratio,

        "). Reduce beam_t to ≤ ", floor(beam_l * max_ratio * 10) / 10, "mm."));


  union() {

    // The flexible beam itself

    cube([beam_w, beam_l, beam_t]);

    // The hook at the end of the beam that catches the mating part

    translate([0, beam_l, 0]) union() {

      rotate([0, -hook_angle, 0])

        cube([beam_w, 2, hook_h]);

      translate([0, 0, beam_t])

        cube([beam_w, 2, hook_h]);

    }

  }

}


// PLA snap arm — standard sizing

snap_arm(beam_l=20, beam_t=1.8, beam_w=8, hook_h=1.5, material="PLA");

// For a 20mm PLA beam: max thickness = 20 × 0.15 = 3.0mm. This design uses 1.8mm — safe.


// PETG snap arm — can be slightly thicker, more resilient

translate([15, 0, 0])

  snap_arm(beam_l=20, beam_t=2.2, beam_w=8, hook_h=1.5, material="PETG");


\end{lstlisting}

\subsubsection*{Orienting the Snap-Fit for Printing}\label{docs__pandoc__latex__src__3dmake_all_lessons.md__orienting-the-snap-fit-for-printing}

Snap-fit beams need to flex in one specific direction. For maximum strength and flexibility, the beam should be printed so that the flex direction is perpendicular to the layer lines. This means orienting the print so the beam runs horizontally and flexes up-down (in the Z direction), not sideways (where layer interfaces would create a weakness).

\begin{center}\rule{0.5\linewidth}{0.5pt}\end{center}

\subsection*{Step 4 --- Dovetail Slide Joint}\label{docs__pandoc__latex__src__3dmake_all_lessons.md__step-4--dovetail-slide-joint}

\subsubsection*{What a Dovetail Is and Why It\textquotesingle s Useful}\label{docs__pandoc__latex__src__3dmake_all_lessons.md__what-a-dovetail-is-and-why-its-useful}

A dovetail joint has a trapezoidal cross-section --- wider at one end than the other, like the tail of a dove. This shape means the joint can only slide in one direction. You can slide two dovetail pieces apart, but you cannot pull them perpendicular to the slide direction --- the trapezoidal shape prevents it.

This makes dovetails excellent for: drawer rails, camera slider tracks, adjustable phone mounts, and any assembly where you want parts that can be separated easily but don\textquotesingle t fall apart during normal use.

\begin{lstlisting}[style=Alabaster, language=openscad]
dt_l     = 50;    // mm — length of the slide

dt_w_top = 8;     // mm — narrower dimension (top of trapezoid)

dt_w_bot = 12;    // mm — wider dimension (base of trapezoid)

dt_h     = 5;     // mm — depth of the dovetail profile

clearance = 0.2;  // mm — your measured slip-fit clearance


module dovetail_male(l=dt_l, w_top=dt_w_top, w_bot=dt_w_bot, h=dt_h) {

  // The trapezoidal profile, extruded along the Y axis

  profile = [

    [-w_bot/2, 0], [ w_bot/2, 0],   // wide base

    [ w_top/2, h], [-w_top/2, h]    // narrow top

  ];

  rotate([90, 0, 0])

    linear_extrude(l)

      polygon(profile);

}


module dovetail_female(l=dt_l, w_top=dt_w_top, w_bot=dt_w_bot, h=dt_h, c=clearance) {

  // Female slot: same shape but all dimensions enlarged by clearance

  profile = [

    [-(w_bot/2 + c), -0.001], [ w_bot/2 + c, -0.001],

    [ w_top/2 + c,   h + c ], [-(w_top/2 + c), h + c]

  ];

  rotate([90, 0, 0])

    linear_extrude(l + 0.002)

      polygon(profile);

}


// Part A: a plate with a female dovetail slot

difference() {

  cube([30, dt_l, 15]);

  translate([15, -0.001, 5])

    dovetail_female();

}


// Part B: the sliding piece with a male dovetail

translate([40, 0, 0]) {

  cube([30, dt_l, 10]);

  translate([55 - 40, 0, 5])

    dovetail_male();

}


\end{lstlisting}

\begin{center}\rule{0.5\linewidth}{0.5pt}\end{center}

\subsection*{Step 5 --- Stackable Bin System}\label{docs__pandoc__latex__src__3dmake_all_lessons.md__step-5--stackable-bin-system}

\subsubsection*{Designing Bins That Stack}\label{docs__pandoc__latex__src__3dmake_all_lessons.md__designing-bins-that-stack}

For parts that stack vertically, one piece needs a raised lip on top that fits inside the walls of the next piece. The outer dimensions of this lip must be slightly smaller than the inner dimensions of the wall --- by your slip-fit clearance --- so the stacked bin sits snugly but can still be separated by hand.

\begin{lstlisting}[style=Alabaster, language=openscad]
bin_w     = 80;    // mm — bin width

bin_d     = 60;    // mm — bin depth

bin_h     = 40;    // mm — bin height

wall      = 2.5;   // mm — wall thickness

lip_h     = 5;     // mm — height of the stacking lip

lip_clear = 0.25;  // mm — your measured slip-fit clearance


module bin_body() {

  difference() {

    cube([bin_w, bin_d, bin_h]);

    translate([wall, wall, wall])

      cube([bin_w - 2*wall, bin_d - 2*wall, bin_h]);  // hollow interior, open top

  }

}


module stacking_lip() {

  // Lip outer dimensions = bin inner dimensions - clearance

  // This makes the lip fit inside the next bin's walls

  translate([lip_clear, lip_clear, bin_h])

    difference() {

      cube([bin_w - 2*lip_clear, bin_d - 2*lip_clear, lip_h]);

      translate([wall - lip_clear, wall - lip_clear, wall])

        cube([bin_w - 2*wall, bin_d - 2*wall, lip_h]);

    }

}


bin_body();

stacking_lip();


\end{lstlisting}

After building this, change \texttt{lip\_clear} and note how it affects the stacking fit. If you print two bins and they don\textquotesingle t stack or are too loose, adjust \texttt{lip\_clear} by 0.05--0.10 mm and reprint the test piece.

\begin{center}\rule{0.5\linewidth}{0.5pt}\end{center}

\subsection*{Step 6 --- Heat-Set Threaded Inserts}\label{docs__pandoc__latex__src__3dmake_all_lessons.md__step-6--heat-set-threaded-inserts}

\subsubsection*{The Problem with Plastic Threads}\label{docs__pandoc__latex__src__3dmake_all_lessons.md__the-problem-with-plastic-threads}

If you tap a thread directly into printed plastic and screw a bolt into it repeatedly, the plastic threads will wear down and eventually strip. This makes plastic-threaded holes unsuitable for any assembly that will be disassembled and reassembled multiple times.

\textbf{Heat-set threaded inserts} solve this problem. They are small brass cylinders with internal threads and external ridges. You heat them with a soldering iron and press them into a printed hole --- the heat melts the plastic around the insert, the ridges embed into the material, and when it cools, you have a permanent metal thread that won\textquotesingle t strip even after many assembly cycles.

Heat-set inserts are commonly available in M2, M3, and M4 sizes and are standard hardware in professional product design.

\begin{lstlisting}[style=Alabaster, language=openscad]
// M3 brass heat-set insert dimensions:

// Outer diameter: ~4.5mm; Length: ~5.7mm

module m3_insert_pocket(depth=6) {

  // Slightly undersized hole — insert is heated and pressed in with force

  cylinder(r=2.35, h=depth + 0.001, $fn=16);  // (4.5 + 0.2) / 2

}


// M4 brass heat-set insert:

// Outer diameter: ~5.6mm

module m4_insert_pocket(depth=8) {

  cylinder(r=2.9, h=depth + 0.001, $fn=16);   // (5.6 + 0.2) / 2

}


// Example: a plate with two M3 insert pockets

difference() {

  cube([40, 30, 10]);

  translate([10, 15, -0.001]) m3_insert_pocket(depth=10.002);

  translate([30, 15, -0.001]) m3_insert_pocket(depth=10.002);

}


\end{lstlisting}

\textbf{Note:} Heat-set insert hole dimensions vary between manufacturers. Check the datasheet for your specific inserts. The values above are typical for common brass inserts.

\begin{center}\rule{0.5\linewidth}{0.5pt}\end{center}

\subsection*{Step 7 --- True Chamfers}\label{docs__pandoc__latex__src__3dmake_all_lessons.md__step-7--true-chamfers}

\subsubsection*{What a Chamfer Is}\label{docs__pandoc__latex__src__3dmake_all_lessons.md__what-a-chamfer-is}

A chamfer is a 45° beveled edge where two surfaces meet at a sharp corner. Chamfers serve two purposes in printed parts: they look more professional than a sharp edge, and they can act as lead-ins that help parts align during assembly.

A common mistake is creating a chamfer by simply cropping the top of a part at an angle --- this produces a slanted surface, but it\textquotesingle s not a true chamfer created from the model\textquotesingle s geometry. A true chamfer is created using a \texttt{difference()} boolean cut.

\begin{lstlisting}[style=Alabaster, language=openscad]
module chamfered_box(w, d, h, c=2) {

  // c = chamfer size in mm (both horizontal and vertical dimension of the bevel)

  difference() {

    cube([w, d, h]);

    // Remove the top c millimeters of the edges

    // by cutting a hollow frame shape from the top

    translate([0, 0, h - c])

      difference() {

        cube([w, d, c + 0.001]);                      // top layer of the box

        translate([c, c, 0]) cube([w - 2*c, d - 2*c, c + 0.001]);  // inset rectangle

      }

  }

}


chamfered_box(50, 40, 20, c=3);


\end{lstlisting}

The chamfer logic: take the top \texttt{c} mm of the box as a separate region. Subtract an inset rectangle from it, leaving only the border frame. Remove that border frame from the original box. The result is a clean 45° bevel around the top perimeter.

\begin{center}\rule{0.5\linewidth}{0.5pt}\end{center}

\subsection*{Step 8 --- Dowel Pins for Alignment}\label{docs__pandoc__latex__src__3dmake_all_lessons.md__step-8--dowel-pins-for-alignment}

\subsubsection*{Why Alignment Pins Matter}\label{docs__pandoc__latex__src__3dmake_all_lessons.md__why-alignment-pins-matter}

When two parts are assembled by hand, they may not align perfectly --- especially if the mating surfaces are flat. Dowel pins are small cylindrical protrusions on one part that fit into matching pockets on the other. They guarantee precise alignment regardless of how the parts were held during assembly.

One part has the pins (protruding pegs). The other has the pockets (matching blind holes). The pins enter the pockets with a press-fit or slip-fit clearance, locking the parts into perfect alignment.

\begin{lstlisting}[style=Alabaster, language=openscad]
pin_r     = 3;      // mm — pin radius

pin_h     = 6;      // mm — pin height

pin_clear = 0.15;   // mm — slightly tighter than slip-fit for precise alignment


module alignment_pin() {

  cylinder(r=pin_r, h=pin_h, $fn=32);

}


module alignment_pocket() {

  // Slightly larger radius and slightly deeper — ensures pin seats fully

  cylinder(r=pin_r + pin_clear, h=pin_h + 0.5, $fn=32);

}


// Part A — has two alignment pins protruding from top face

cube([50, 40, 5]);

translate([10, 10, 5]) alignment_pin();

translate([40, 30, 5]) alignment_pin();


// Part B — has matching pockets in bottom face

translate([60, 0, 0])

  difference() {

    cube([50, 40, 5]);

    translate([10, 10, -0.001]) alignment_pocket();

    translate([40, 30, -0.001]) alignment_pocket();

  }


\end{lstlisting}

\begin{center}\rule{0.5\linewidth}{0.5pt}\end{center}

\subsection*{Exercises}\label{docs__pandoc__latex__src__3dmake_all_lessons.md__exercises-3}

\textbf{Exercise 8.1:} Print the tolerance test coupon. Test each peg-socket pair. Record your press-fit and slip-fit clearance values. Update \texttt{PRESS\_FIT\_C} and \texttt{SLIP\_FIT\_C} with your measured values.

\textbf{Exercise 8.2:} Design a snap-fit battery cover for a 60 × 40 × 20 mm space using your measured press-fit clearance and the \texttt{snap\_arm()} module. Verify the beam flexibility ratio before printing.

\textbf{Exercise 8.3 (Advanced):} Create a parametric box set where a \texttt{size\_scale} parameter (values 0.5, 1.0, 2.0) generates small, medium, and large boxes that all use the same lid design and tolerance values. Generate all three sizes with a bash/PowerShell loop and verify each with \texttt{3dm\ info}.

\begin{center}\rule{0.5\linewidth}{0.5pt}\end{center}

\subsection*{Quiz --- Lesson 8 (15 questions)}\label{docs__pandoc__latex__src__3dmake_all_lessons.md__quiz--lesson-8-15-questions}

\begin{enumerate}
\tightlist
\item
  Why do FDM-printed holes consistently come out smaller than designed? Explain the physical mechanism.
\item
  What clearance range (in mm) produces a slip-fit for a typical FDM printer with a 0.4mm nozzle and PLA?
\item
  What is a press-fit, and what application would you use one for?
\item
  What is the beam\_t / beam\_l flexibility rule for PLA snap-fit cantilever beams?
\item
  Why must you print a tolerance test coupon rather than using published clearance values from a chart?
\item
  What is a heat-set threaded insert and why is it preferred over a plastic-tapped hole for assemblies that are disassembled repeatedly?
\item
  What is a dovetail joint and in which direction does it lock?
\item
  What does a stacking lip do in a stackable bin design, and how does clearance factor into its dimensions?
\item
  True or False: a thicker snap-fit beam is easier to deflect than a thinner one of the same length.
\item
  What happens if you design a stacking lip with clearance = 0mm?
\item
  Explain the concept of "tolerance stack-up." If three parts each have ±0.2mm dimensional variation, what is the worst-case total gap variation in an assembled chain of three?
\item
  Why should tolerance clearance values be defined as named constants at the top of a file rather than hardcoded at each use?
\item
  Describe how a chamfer is created using \texttt{difference()}. What shape is subtracted from the original?
\item
  For a snap-fit clip made of PLA with \texttt{beam\_l\ =\ 18mm}, what is the maximum \texttt{beam\_t} allowed by the flexibility rule? Show your calculation.
\item
  Describe how you would design and conduct a complete tolerance test for a new printer. What prints would you make, what would you measure, and how would you decide on your final clearance values?
\end{enumerate}

\begin{center}\rule{0.5\linewidth}{0.5pt}\end{center}

\subsection*{Extension Problems (15)}\label{docs__pandoc__latex__src__3dmake_all_lessons.md__extension-problems-15-7}

\begin{enumerate}
\tightlist
\item
  Design a complete stackable bin set (small, medium, large) where each size nests inside the next. Use your measured tolerance values throughout.
\item
  Build a snap-fit box lid that attaches with four clips. Make the number of clips a parameter and validate the beam dimensions with \texttt{assert()}.
\item
  Design a press-fit peg and socket set with five clearances from 0.0 to 0.4 mm. Print one set, test all pairs, and create a table showing which clearance produced which fit type on your printer.
\item
  Create a wall-mount organizer using snap-fit clips that attach to a standard pegboard (hole pitch = 25.4mm, hole diameter = 6.35mm).
\item
  Build a two-part enclosure with heat-set insert pockets for M3 screws and matching through-holes for the bolts. Test with real hardware.
\item
  Design a snap-together dice tray: the tray snaps into a carrying case lid using four snap clips.
\item
  Create an alignment jig using two dowel pins and matching pockets. Use it to verify your printer\textquotesingle s positional accuracy by measuring pin spacing with calipers.
\item
  Build a cable management clip with a snap-fit that can be opened and closed repeatedly for easy cable insertion and removal.
\item
  Design a telescoping tube: an inner cylinder that slides freely inside an outer cylinder using your measured slip-fit clearance.
\item
  Create a bayonet mount: a part that locks with a 90° twist.
\item
  Design a "living hinge" --- a thin flexible connection between two rigid sections that can be folded. Research the appropriate wall thickness and document which materials are suitable.
\item
  Build a parametric drawer divider system: interlocking slotted panels that can be arranged in any grid configuration.
\item
  Design a self-locking dovetail joint that slides together in one direction and resists being pulled apart in any perpendicular direction.
\item
  Create a complete assembly drawing in OpenSCAD: an exploded view with each component translated away from the assembly center, with clear labels and a tolerance reference table as a comment block.
\item
  Conduct a failure mode analysis for a snap-fit clip: list five ways it could fail in use, rate each by likelihood and severity, and propose a design modification to reduce the top two risks.
\end{enumerate}

\begin{center}\rule{0.5\linewidth}{0.5pt}\end{center}

\subsection*{References and Helpful Resources}\label{docs__pandoc__latex__src__3dmake_all_lessons.md__references-and-helpful-resources-7}

\begin{itemize}
\item
  Snap-Fit Design Principles --- Bayer MaterialScience Snap-Fit Design Manual
\item
  OpenSCAD User Manual --- Difference and Boolean Operations --- https://en.wikibooks.org/wiki/OpenSCAD\_User\_Manual/CSG\_Modelling
\item
  Tolerance and Fit in FDM Printing --- All3DP Guide --- https://all3dp.com/2/fdm-3d-printing-tolerances/
\end{itemize}

\subsubsection*{Supplemental Resources}\label{docs__pandoc__latex__src__3dmake_all_lessons.md__supplemental-resources-7}

\begin{itemize}
\tightlist
\item
  Programming with OpenSCAD EPUB Textbook --- assets/Programming\_with\_OpenSCAD.epub --- Assemblies and interlocking features
\item
  3DMake GitHub Repository --- https://github.com/tdeck/3dmake
\item
  All other reference documents are embedded as appendices below --- no separate files needed.
\end{itemize}

\begin{center}\rule{0.5\linewidth}{0.5pt}\end{center}

\section*{Part 2: Extension Project and Reference Appendices}\label{docs__pandoc__latex__src__3dmake_all_lessons.md__part-2-extension-project-and-reference-appendices-3}

\emph{All supporting documents for Lesson 8 are included below. You do not need to open any additional files.}

\begin{center}\rule{0.5\linewidth}{0.5pt}\end{center}

\subsection*{Appendix A: Snap-Fit Clip --- Extension Project}\label{docs__pandoc__latex__src__3dmake_all_lessons.md__appendix-a-snap-fit-clip--extension-project}

\textbf{Estimated time:} 3--6 hours

\subsubsection*{Overview}\label{docs__pandoc__latex__src__3dmake_all_lessons.md__overview-6}

This extension project takes the snap-fit concepts from Lesson 8 and turns them into a systematic engineering investigation. Rather than building a snap-fit and hoping it works, you will design a parametric test structure, print it at three different tolerance values, measure what happens, and let the data drive your final design.

This is how snap-fit components are actually designed in professional practice: tolerance ranges are established empirically for a specific material-printer combination, then codified as constants for future designs.

\subsubsection*{Learning Objectives}\label{docs__pandoc__latex__src__3dmake_all_lessons.md__learning-objectives-13}

By completing this project, you will:

\begin{itemize}
\tightlist
\item
  Design flexible components that must bend without failing
\item
  Conduct tolerance and material testing systematically
\item
  Understand the relationship between geometry, material properties, and functional performance
\item
  Document design iterations and trade-offs in a format that would transfer to a colleague
\end{itemize}

\subsubsection*{Objective}\label{docs__pandoc__latex__src__3dmake_all_lessons.md__objective-1}

Design a parametric snap-fit clip with adjustable clip thickness and spacing. Use it to determine optimal tolerances for your specific printer and material combination.

\subsubsection*{Tasks}\label{docs__pandoc__latex__src__3dmake_all_lessons.md__tasks-5}

\begin{enumerate}
\item
  \textbf{Design} a parametric snap-fit test piece with adjustable clip thickness and spacing. All key dimensions must be variables. The design must include at least one cantilever beam using the beam\_t / beam\_l flexibility rule from Lesson 8.
\item
  \textbf{Print test pieces} for three tolerance values. For each variant, change \emph{only} the tolerance variable --- keep all other parameters fixed. Label each print (embossed text, or a note taped to the print).
\item
  \textbf{Test each clip} for snap-fit behavior. For each variant: assemble and disassemble at least 5 times. Record the effort required, any deformation observed, and whether the clip still functions after 5 cycles.
\item
  \textbf{Refine dimensions} based on results. Choose the tolerance value that produced the most reliable behavior and create a final clip design. Write a one-paragraph justification of your choice.
\end{enumerate}

\subsubsection*{Deliverables}\label{docs__pandoc__latex__src__3dmake_all_lessons.md__deliverables-2}

\begin{itemize}
\tightlist
\item
  Source \texttt{.scad} file with parametric clip modules (all key dimensions as named variables)
\item
  Test print results table (see student documentation template in Appendix B)
\item
  Final clip design with recommended print settings documented
\item
  Completed student documentation template
\end{itemize}

\begin{center}\rule{0.5\linewidth}{0.5pt}\end{center}

\subsubsection*{Starter Code}\label{docs__pandoc__latex__src__3dmake_all_lessons.md__starter-code-4}

The stackable bins example from Lesson 8 demonstrates the stacking rim and clearance pattern that snap-fit features build on. Use it as a reference for parametric structure, then design your own snap-fit geometry.

\begin{lstlisting}[style=Alabaster, language=openscad]
// Stackable Storage Bins — Structural Reference

// This example shows the clearance and rim pattern.

// Your snap-fit clip will use similar parametric constants.


bin_w        = 80;    // mm — bin width

bin_d        = 120;   // mm — bin depth

bin_h        = 60;    // mm — bin height

wall         = 2;     // mm — wall thickness

rim          = 3;     // mm — interlock rim height

chamfer      = 2;     // mm — top edge chamfer

stack_clear  = 0.6;   // mm — clearance between stacked bins (measure yours)


module outer() {

  cube([bin_w, bin_d, bin_h]);

}


module inner() {

  translate([wall, wall, wall])

    cube([bin_w - 2*wall, bin_d - 2*wall, bin_h - 2*wall]);

}


module body() {

  difference() { outer(); inner(); }

}


module rim_outer() {

  translate([0, 0, bin_h])

    cube([bin_w, bin_d, rim]);

}


module rim_inner() {

  translate([wall + stack_clear, wall + stack_clear, bin_h])

    cube([bin_w - 2*(wall + stack_clear),

          bin_d - 2*(wall + stack_clear), rim]);

}


module chamfer_top() {

  difference() {

    children();

    translate([-1, -1, bin_h - chamfer])

      cube([bin_w + 2, bin_d + 2, chamfer + 2]);

  }

}


union() {

  chamfer_top() { body(); }

  rim_outer();

  difference() { rim_outer(); rim_inner(); }

}


\end{lstlisting}

\textbf{Snap-fit beam starting scaffold} --- adapt this for your test piece:

\begin{lstlisting}[style=Alabaster, language=openscad]
// Snap-Fit Cantilever Beam — Parametric Test Scaffold

// Replace the base_body and target geometry with your own design.


// ---- Key Parameters ----

beam_l      = 15;    // mm — beam length (longer = more flex)

beam_t      = 1.5;   // mm — beam thickness (thinner = more flex)

beam_w      = 6;     // mm — beam width

hook_h      = 1.2;   // mm — hook catch height (the lip that locks)

tolerance   = 0.2;   // mm — clearance for the socket opening

base_w      = 30;    // mm — base block width

base_d      = 20;    // mm — base block depth

base_h      = 4;     // mm — base block height


// Flexibility check: beam_t / beam_l should be < 0.1 for PLA

echo("Flexibility ratio:", beam_t / beam_l,

     "(target < 0.1 for reliable snap-fit)");

assert(beam_t / beam_l < 0.12,

  str("Beam too stiff: beam_t/beam_l = ", beam_t/beam_l,

      ". Reduce beam_t or increase beam_l."));


// ---- Snap-Fit Arm ----

module snap_arm() {

  union() {

    // Main cantilever beam

    cube([beam_w, beam_l, beam_t]);

    // Hook at tip

    translate([0, beam_l - hook_h, 0])

      cube([beam_w, hook_h, beam_t + hook_h]);

  }

}


// ---- Socket (receiver) ----

// The socket opening is wider than the arm by 2×tolerance

module snap_socket() {

  socket_opening = beam_w + 2 * tolerance;

  difference() {

    cube([base_w, base_d, base_h]);

    // Socket slot

    translate([(base_w - socket_opening) / 2, 0, -0.001])

      cube([socket_opening, beam_l + hook_h + tolerance, base_h + 0.002]);

  }

}


// ---- Test Assembly (render both for sizing check) ----

snap_socket();

translate([(base_w - beam_w) / 2, 0, base_h])

  snap_arm();


\end{lstlisting}

\subsubsection*{Assessment Questions}\label{docs__pandoc__latex__src__3dmake_all_lessons.md__assessment-questions-1}

\begin{enumerate}
\tightlist
\item
  What clip thickness and spacing values produced a reliable snap-fit for your specific printer and material?
\item
  How did you balance flexibility (easy to snap and unsnap) with durability (no premature failure)?
\item
  What would you change in the design if this clip needed to survive 1,000+ assembly cycles?
\end{enumerate}

\begin{center}\rule{0.5\linewidth}{0.5pt}\end{center}

\subsection*{Appendix B: Snap-Fit Clip --- Student Documentation Template}\label{docs__pandoc__latex__src__3dmake_all_lessons.md__appendix-b-snap-fit-clip--student-documentation-template}

\textbf{Author:} \_\_\_\_\_\_\_\_\_\_\_\_\_\_\_\_
\textbf{Date:} \_\_\_\_\_\_\_\_\_\_
\textbf{Description:} Design a snap-fit connector that joins two 3D-printed parts without external fasteners.

\begin{center}\rule{0.5\linewidth}{0.5pt}\end{center}

\subsubsection*{Design Concept}\label{docs__pandoc__latex__src__3dmake_all_lessons.md__design-concept-2}

\begin{itemize}
\tightlist
\item
  What two parts will your snap-fit connect?
\item
  Design approach (cantilever arm, torsion box, other):
\item
  Material selected and expected flex limits:
\item
  How did you verify the beam flexibility ratio (beam\_t / beam\_l)?
\end{itemize}

\begin{center}\rule{0.5\linewidth}{0.5pt}\end{center}

\subsubsection*{Tolerance Specifications}\label{docs__pandoc__latex__src__3dmake_all_lessons.md__tolerance-specifications}

{\def\LTcaptype{none} % do not increment counter
\begin{longtable}[]{@{}
  >{\raggedright\arraybackslash}p{(\linewidth - 6\tabcolsep) * \real{0.3288}}
  >{\centering\arraybackslash}p{(\linewidth - 6\tabcolsep) * \real{0.2466}}
  >{\centering\arraybackslash}p{(\linewidth - 6\tabcolsep) * \real{0.3014}}
  >{\raggedright\arraybackslash}p{(\linewidth - 6\tabcolsep) * \real{0.1233}}@{}}
\toprule\noalign{}
\begin{minipage}[b]{\linewidth}\raggedright
Parameter
\end{minipage} & \begin{minipage}[b]{\linewidth}\centering
Nominal Value (mm)
\end{minipage} & \begin{minipage}[b]{\linewidth}\centering
Tolerance Applied (mm)
\end{minipage} & \begin{minipage}[b]{\linewidth}\raggedright
Rationale
\end{minipage} \\
\midrule\noalign{}
\endhead
\bottomrule\noalign{}
\endlastfoot
Beam thickness (beam\_t) & & & \\
Beam length (beam\_l) & & & \\
Hook height & & & \\
Socket opening clearance & & & \\
Contact area width & & & \\
\end{longtable}
}

\textbf{Flexibility ratio check:}

\begin{lstlisting}[style=Alabaster]
beam_t / beam_l = ___ / ___ = ___

Target: < 0.1 for PLA, < 0.08 for PETG

Result: [ ] Pass   [ ] Fail — adjusted to: ___


\end{lstlisting}

\begin{center}\rule{0.5\linewidth}{0.5pt}\end{center}

\subsubsection*{Test Print Results}\label{docs__pandoc__latex__src__3dmake_all_lessons.md__test-print-results}

Print three variants with different tolerance values. Change \textbf{only} the tolerance variable between variants.

{\def\LTcaptype{none} % do not increment counter
\begin{longtable}[]{@{}
  >{\centering\arraybackslash}p{(\linewidth - 10\tabcolsep) * \real{0.0769}}
  >{\centering\arraybackslash}p{(\linewidth - 10\tabcolsep) * \real{0.1538}}
  >{\centering\arraybackslash}p{(\linewidth - 10\tabcolsep) * \real{0.1648}}
  >{\centering\arraybackslash}p{(\linewidth - 10\tabcolsep) * \real{0.1978}}
  >{\centering\arraybackslash}p{(\linewidth - 10\tabcolsep) * \real{0.2857}}
  >{\centering\arraybackslash}p{(\linewidth - 10\tabcolsep) * \real{0.1209}}@{}}
\toprule\noalign{}
\begin{minipage}[b]{\linewidth}\centering
Variant
\end{minipage} & \begin{minipage}[b]{\linewidth}\centering
Tolerance (mm)
\end{minipage} & \begin{minipage}[b]{\linewidth}\centering
Assembly effort
\end{minipage} & \begin{minipage}[b]{\linewidth}\centering
Disassembly effort
\end{minipage} & \begin{minipage}[b]{\linewidth}\centering
Deformation after 5 cycles
\end{minipage} & \begin{minipage}[b]{\linewidth}\centering
Pass / Fail
\end{minipage} \\
\midrule\noalign{}
\endhead
\bottomrule\noalign{}
\endlastfoot
V1 & & & & & \\
V2 & & & & & \\
V3 & & & & & \\
\end{longtable}
}

\textbf{Selected tolerance value:} \_\_\_\_\_ mm
\textbf{Justification} (one paragraph --- reference your test data, not intuition):

\begin{center}\rule{0.5\linewidth}{0.5pt}\end{center}

\subsubsection*{Assembly Testing Log}\label{docs__pandoc__latex__src__3dmake_all_lessons.md__assembly-testing-log}

\paragraph*{Cycle 1 --- Initial Assembly}\label{docs__pandoc__latex__src__3dmake_all_lessons.md__cycle-1--initial-assembly}

\begin{itemize}
\tightlist
\item
  Date:
\item
  Effort required to snap together (easy / moderate / hard / wouldn\textquotesingle t snap):
\item
  Snap strength --- resistance to pulling apart (weak / moderate / strong / won\textquotesingle t release):
\item
  Visual inspection --- any deformation, whitening at bend, or cracking:
\end{itemize}

\paragraph*{Cycle 2}\label{docs__pandoc__latex__src__3dmake_all_lessons.md__cycle-2}

\begin{itemize}
\tightlist
\item
  Date:
\item
  Observations:
\end{itemize}

\paragraph*{Cycle 3}\label{docs__pandoc__latex__src__3dmake_all_lessons.md__cycle-3}

\begin{itemize}
\tightlist
\item
  Date:
\item
  Observations:
\end{itemize}

\paragraph*{Cycle 4}\label{docs__pandoc__latex__src__3dmake_all_lessons.md__cycle-4}

\begin{itemize}
\tightlist
\item
  Date:
\item
  Observations:
\end{itemize}

\paragraph*{Cycle 5}\label{docs__pandoc__latex__src__3dmake_all_lessons.md__cycle-5}

\begin{itemize}
\tightlist
\item
  Date:
\item
  Observations:
\end{itemize}

\paragraph*{Durability Assessment}\label{docs__pandoc__latex__src__3dmake_all_lessons.md__durability-assessment}

\begin{itemize}
\tightlist
\item
  After 5 cycles, is the snap-fit still functional? Yes / No / Partially
\item
  Any permanent deformation observed? Describe:
\item
  Estimated service life based on observed degradation rate:
\end{itemize}

\begin{center}\rule{0.5\linewidth}{0.5pt}\end{center}

\subsubsection*{Mechanical Analysis}\label{docs__pandoc__latex__src__3dmake_all_lessons.md__mechanical-analysis}

Answer in complete sentences:

\begin{enumerate}
\item
  Describe the snap mechanism --- explain \emph{why} it works mechanically, not just that it does.
\item
  What material properties enable this design? What would happen if you printed the same geometry in a rigid material with no flex (like resin)?
\item
  Where are the stress concentration points in your design? How did you address them?
\end{enumerate}

\begin{center}\rule{0.5\linewidth}{0.5pt}\end{center}

\subsubsection*{Reflections}\label{docs__pandoc__latex__src__3dmake_all_lessons.md__reflections-5}

\begin{enumerate}
\item
  Did the snap-fit work as expected on the first print? If not, what surprised you?
\item
  What tolerance adjustment had the most significant effect, and by how much?
\item
  How would you modify this design for a clip that needs to work at outdoor temperatures (hot car in summer, cold winter)?
\item
  What materials would fail with this exact geometry? Why?
\end{enumerate}

\begin{center}\rule{0.5\linewidth}{0.5pt}\end{center}

\subsubsection*{Attachments Checklist}\label{docs__pandoc__latex__src__3dmake_all_lessons.md__attachments-checklist-5}

\begin{itemize}
\tightlist
\item[$\square$]
  \texttt{.scad} file with parametric snap-fit module (all key dimensions as named variables)
\item[$\square$]
  Photos of test variants before assembly (labeled V1, V2, V3)
\item[$\square$]
  Photos of assembled connections for each variant
\item[$\square$]
  Photos showing any deformation after 5 cycles
\item[$\square$]
  Completed tolerance specifications table
\item[$\square$]
  Completed testing log with dates and observations
\end{itemize}

\begin{center}\rule{0.5\linewidth}{0.5pt}\end{center}

\subsubsection*{Teacher Feedback}\label{docs__pandoc__latex__src__3dmake_all_lessons.md__teacher-feedback-5}

{\def\LTcaptype{none} % do not increment counter
\begin{longtable}[]{@{}
  >{\raggedright\arraybackslash}p{(\linewidth - 4\tabcolsep) * \real{0.8242}}
  >{\centering\arraybackslash}p{(\linewidth - 4\tabcolsep) * \real{0.1209}}
  >{\raggedright\arraybackslash}p{(\linewidth - 4\tabcolsep) * \real{0.0549}}@{}}
\toprule\noalign{}
\begin{minipage}[b]{\linewidth}\raggedright
Category
\end{minipage} & \begin{minipage}[b]{\linewidth}\centering
Score (0--3)
\end{minipage} & \begin{minipage}[b]{\linewidth}\raggedright
Notes
\end{minipage} \\
\midrule\noalign{}
\endhead
\bottomrule\noalign{}
\endlastfoot
Problem \& Solution --- snap-fit works reliably; meets functional requirements
& & \\
Design \& Code Quality --- parametric, mechanically sound, beam rule applied
& & \\
Documentation --- testing log complete; tolerance rationale explained &
& \\
\textbf{Total (0--9)} & & \\
\end{longtable}
}

\begin{center}\rule{0.5\linewidth}{0.5pt}\end{center}

\subsection*{Appendix C: Snap-Fit Clip --- Instructor Reference}\label{docs__pandoc__latex__src__3dmake_all_lessons.md__appendix-c-snap-fit-clip--instructor-reference}

\subsubsection*{Project Context}\label{docs__pandoc__latex__src__3dmake_all_lessons.md__project-context-4}

This extension project reinforces the core tolerance and snap-fit content of Lesson 8 through hands-on empirical testing. The key distinction from a standard assignment is the \textbf{systematic testing approach}: students must let printed test data drive their final design decisions, not intuition or first-attempt results.

\subsubsection*{Key Learning Points to Reinforce}\label{docs__pandoc__latex__src__3dmake_all_lessons.md__key-learning-points-to-reinforce-4}

\textbf{Tolerance engineering is empirical.} Published clearance values are starting points. The only reliable way to know what works on a specific printer with a specific material is to print and measure. Students who skip the test coupon step and jump straight to a final design will encounter this lesson the hard way.

\textbf{The beam flexibility ratio is a design rule, not a guideline.} Beam\_t / beam\_l \textless{} 0.1 (PLA) is derived from material yield strain. A beam that violates this ratio will either not snap (too stiff) or crack on first use (overstressed). The \texttt{assert()} in the starter code enforces this at render time.

\textbf{Snap-fits degrade with cycles.} PLA creep under repeated stress means a PLA snap-fit that works perfectly at cycle 1 may be noticeably looser at cycle 20. This is not a design failure --- it\textquotesingle s an expected material behavior that students should document and account for.

\subsubsection*{Constraints (Enforce These)}\label{docs__pandoc__latex__src__3dmake_all_lessons.md__constraints-enforce-these-4}

\begin{itemize}
\tightlist
\item
  Snap-fit must join two printed parts without any external fasteners
\item
  All key dimensions must be variables (parametric code required)
\item
  At least three variants must be printed and tested --- no single-attempt submissions
\item
  Assembly must be tested over at least 5 cycles with observations recorded for each
\end{itemize}

\subsubsection*{Functional Requirements (Evaluation Criteria)}\label{docs__pandoc__latex__src__3dmake_all_lessons.md__functional-requirements-evaluation-criteria}

\begin{itemize}
\tightlist
\item
  Parts snap together securely without over-constraint or binding
\item
  Connection resists separation under intended use without permanent deformation at cycle 1
\item
  Snap mechanism is reusable for at least 5 assembly cycles
\item
  Design reflects understanding of tolerance and material behavior (visible in code comments and documentation)
\end{itemize}

\subsubsection*{Assessment Notes}\label{docs__pandoc__latex__src__3dmake_all_lessons.md__assessment-notes-4}

\textbf{Strong submissions} show: tolerance test data driving the final design decision (not guessing), multiple assembly cycles documented with specific observations, and a reflection that discusses material limits with reference to measured results.

\textbf{Common weak areas:} Single-variant testing ("it worked so I stopped"), vague reflection ("I would make it thicker"), and tolerance values with no rationale. Push back on these in feedback.

\textbf{Extension opportunity:} Students who finish early can investigate material property comparisons --- print the same geometry in PLA and PETG (if available) and compare flex behavior and cycle durability. This connects directly to the material selection content in Lesson 5.
-e

\begin{center}\rule{0.5\linewidth}{0.5pt}\end{center}

\section*{Lesson 9 --- Automation and 3dm Workflows}\label{docs__pandoc__latex__src__3dmake_all_lessons.md__lesson-9--automation-and-3dm-workflows}

\textbf{Estimated time:} 90--120 minutes

\begin{center}\rule{0.5\linewidth}{0.5pt}\end{center}

\subsection*{Before You Start}\label{docs__pandoc__latex__src__3dmake_all_lessons.md__before-you-start-7}

In Lessons 1--8, you built real, functional 3D models --- phone stands, keycaps, stackable bins --- and learned how parametric design lets a single \texttt{.scad} file produce many different variants just by changing numbers. You\textquotesingle ve been running builds manually: open a terminal, type \texttt{3dm\ build}, check the output, repeat. That\textquotesingle s fine when you\textquotesingle re developing one part. But what happens when you need to build 20 variants of a part for a batch test? Or when you need to \emph{prove} that every commit to your project builds correctly before you send files to a printer?

This lesson is about the next step: \textbf{automation}. You\textquotesingle ll learn to write scripts that build, verify, archive, and iterate without you sitting at the keyboard. These are real software engineering skills --- the same techniques used in companies that ship physical products with embedded software. By the end of this lesson, you\textquotesingle ll understand how professional fabrication workflows work and be able to apply those patterns to your own projects.

\begin{center}\rule{0.5\linewidth}{0.5pt}\end{center}

\subsection*{Learning Objectives}\label{docs__pandoc__latex__src__3dmake_all_lessons.md__learning-objectives-14}

By the end of this lesson you will be able to:

\begin{itemize}
\tightlist
\item
  Write robust automation scripts in bash (Linux/macOS), PowerShell (Windows), and CMD (Windows legacy)
\item
  Use proper error handling, logging, and trap/exit patterns for unattended builds
\item
  Pass string and numeric parameters correctly to \texttt{openscad\ -D} from all three shells
\item
  Automate batch builds, archiving, and variant generation
\item
  Implement a watch mode that rebuilds automatically when source files change
\item
  Integrate 3dMake automation into CI/CD pipelines with GitHub Actions
\end{itemize}

\begin{center}\rule{0.5\linewidth}{0.5pt}\end{center}

\subsection*{Concept: Why Automation Matters}\label{docs__pandoc__latex__src__3dmake_all_lessons.md__concept-why-automation-matters}

Before writing any code, let\textquotesingle s understand the problem automation solves.

\subsubsection*{The Manual Build Problem}\label{docs__pandoc__latex__src__3dmake_all_lessons.md__the-manual-build-problem}

Suppose you\textquotesingle re building a parametric shelf bracket that needs to be tested at three different wall thicknesses (1.2 mm, 1.6 mm, 2.0 mm) across four different bracket widths (60 mm, 80 mm, 100 mm, 120 mm). That\textquotesingle s 12 combinations. Running each build manually means:

\begin{enumerate}
\tightlist
\item
  Open terminal
\item
  Type the command with the right parameter values
\item
  Wait for OpenSCAD to render
\item
  Check that the output exists
\item
  Rename the file so it doesn\textquotesingle t overwrite the last one
\item
  Repeat 11 more times
\end{enumerate}

This takes time, introduces human error (typos in parameter values, forgetting to rename files), and produces results that are hard to reproduce later --- "which settings did I use for the 80mm bracket again?"

\subsubsection*{Automation Solves Four Problems}\label{docs__pandoc__latex__src__3dmake_all_lessons.md__automation-solves-four-problems}

\textbf{Repeatability:} A script runs exactly the same way every time. Parameter values, file names, output locations --- all defined once in code, not typed by hand each time.

\textbf{Error detection:} A good script notices when something goes wrong and stops immediately, rather than silently producing broken output and continuing.

\textbf{Documentation:} A build script \emph{is} documentation. Reading a well-written script tells you exactly how the project is supposed to be built --- what parameters matter, what the expected output looks like, how things are organized.

\textbf{Scalability:} Whether you need 12 builds or 1,200, the automation works the same way.

\subsubsection*{Automation Is a Safety Net}\label{docs__pandoc__latex__src__3dmake_all_lessons.md__automation-is-a-safety-net}

Think of build automation like the safety rails on a mountain trail. The trail was perfectly safe without them --- but the rails mean you can focus on the view instead of watching every step. Automation removes the tedious, error-prone parts of your workflow so you can focus on design.

\begin{center}\rule{0.5\linewidth}{0.5pt}\end{center}

\subsection*{\texorpdfstring{Step 1 --- The OpenSCAD \texttt{-D} Flag and Parameter Passing}{Step 1 --- The OpenSCAD -D Flag and Parameter Passing}}\label{docs__pandoc__latex__src__3dmake_all_lessons.md__step-1--the-openscad--d-flag-and-parameter-passing}

\subsubsection*{\texorpdfstring{What the \texttt{-D} Flag Does}{What the -D Flag Does}}\label{docs__pandoc__latex__src__3dmake_all_lessons.md__what-the--d-flag-does}

OpenSCAD\textquotesingle s \texttt{-D} flag overrides any variable in your \texttt{.scad} file at the command line. This is the bridge between your automation scripts and your parametric models. Instead of editing the \texttt{.scad} file to change a value, you pass the new value directly when you run the build.

Think of it like a recipe where the recipe card has the default quantities written in pencil. The \texttt{-D} flag lets you write different numbers in the margin for this particular batch --- without erasing the original recipe.

\begin{lstlisting}[style=Alabaster, language=bash]
# The .scad file has: width = 40;

# The -D flag overrides it to 60 for this build

openscad -D "width=60" -o output.stl src/main.scad


\end{lstlisting}

\subsubsection*{Numeric Parameters: Straightforward}\label{docs__pandoc__latex__src__3dmake_all_lessons.md__numeric-parameters-straightforward}

For numeric values, the syntax is identical on all platforms:

\begin{lstlisting}[style=Alabaster, language=bash]
# Works the same on Linux/macOS/Windows

openscad -D "width=50" -D "height=15" -o output.stl src/main.scad

openscad -D "width=50" -D "height=15" -D "wall=2.5" -o output.stl src/main.scad


\end{lstlisting}

You can pass as many \texttt{-D} flags as you need. Each one overrides one variable.

\subsubsection*{String Parameters: The Hard Part}\label{docs__pandoc__latex__src__3dmake_all_lessons.md__string-parameters-the-hard-part}

String values are much trickier, and this is where most beginners get tripped up. In OpenSCAD, a string value must be surrounded by double quotes: \texttt{label="FRONT"}. But you\textquotesingle re passing this string \emph{inside another string} on the command line, which creates a quoting conflict.

\textbf{This is the \#1 source of silent failures in automation scripts.} If you get the quoting wrong, OpenSCAD may accept the command without error but produce a model where the label is empty or incorrect. No error message. No warning. Just wrong output.

\paragraph*{Why Quoting Is Different on Each Platform}\label{docs__pandoc__latex__src__3dmake_all_lessons.md__why-quoting-is-different-on-each-platform}

Each shell --- bash (Linux/macOS), PowerShell (Windows), and CMD (Windows legacy) --- has its own rules for how quotes work inside quotes. There\textquotesingle s no single universal syntax. You need to learn the rule for each environment you use.

The underlying goal is always the same: you need OpenSCAD to receive the string \texttt{label="FRONT"} with the double quotes intact. Each shell provides a different way to achieve that.

\paragraph*{Linux / macOS (bash)}\label{docs__pandoc__latex__src__3dmake_all_lessons.md__linux--macos-bash}

\textbf{The rule:} Use single quotes around the entire \texttt{-D} value to protect inner double quotes. Single-quoted strings in bash are completely literal --- nothing inside them is interpreted.

\begin{lstlisting}[style=Alabaster, language=bash]
# CORRECT: single quotes around the entire -D value protect inner double quotes

openscad -D 'label="FRONT"' -o out.stl src/main.scad


# CORRECT: backslash-escaped double quotes inside a double-quoted string

openscad -D "label=\"FRONT\"" -o out.stl src/main.scad


# CORRECT: variable with proper escaping

LABEL="FRONT PANEL"

openscad -D "label=\"${LABEL}\"" -o out.stl src/main.scad


# DEBUGGING TIP: echo the command before running to verify quoting

echo "openscad -D 'label=\"${LABEL}\"' -o out.stl src/main.scad"


# WRONG — string value will be malformed; OpenSCAD may use an empty string

openscad -D "label=FRONT" -o out.stl src/main.scad    # missing inner quotes


\end{lstlisting}

\textbf{Debugging trick:} Before running a command, \texttt{echo} it first to see exactly what your shell will send to OpenSCAD. If the echoed command looks wrong, it \emph{is} wrong --- fix the quoting before running it for real.

\paragraph*{Windows (PowerShell)}\label{docs__pandoc__latex__src__3dmake_all_lessons.md__windows-powershell}

\textbf{The rule:} Single quotes work the same way as in bash --- they protect inner double quotes. PowerShell also supports backtick (\texttt{\textasciigrave{}}) as an escape character inside double-quoted strings.

\begin{lstlisting}[style=Alabaster, language=powershell]
# CORRECT: single quotes protect inner double quotes

openscad -D 'label="FRONT"' -o out.stl src\main.scad


# CORRECT: backtick escapes for double quotes inside a double-quoted string

$label = "FRONT PANEL"

openscad -D "label=`"$label`"" -o out.stl src\main.scad


# CORRECT: build the argument string explicitly

$dArg = 'label="' + $label + '"'

openscad -D $dArg -o out.stl src\main.scad


# DEBUGGING TIP: display the argument before using it

Write-Host "D argument: $dArg"


\end{lstlisting}

\paragraph*{Windows (CMD)}\label{docs__pandoc__latex__src__3dmake_all_lessons.md__windows-cmd}

\textbf{The rule:} CMD uses doubled double-quotes (\texttt{""}) to produce a single literal double quote inside a string. This is the most visually confusing notation --- \texttt{""FRONT""} means: open double-quote, then the word FRONT, then close double-quote.

\begin{lstlisting}[style=Alabaster, language=cmd]
REM CMD uses doubled double-quotes to produce a literal double-quote inside a string

openscad -D "label=""FRONT""" -o out.stl src\main.scad


REM With a variable

SET LABEL=FRONT

openscad -D "label=""%LABEL%""" -o out.stl src\main.scad


REM DEBUGGING: echo first to verify

echo label=""%LABEL%""


\end{lstlisting}

\subsubsection*{Common Mistake: Assuming Silent Success}\label{docs__pandoc__latex__src__3dmake_all_lessons.md__common-mistake-assuming-silent-success}

OpenSCAD will often not report an error when string quoting is wrong. The variable just gets an empty string or \texttt{undef} as its value. If your label text disappears from the rendered model, check your quoting before anything else.

\begin{center}\rule{0.5\linewidth}{0.5pt}\end{center}

\subsection*{Step 2 --- Robust Error Handling}\label{docs__pandoc__latex__src__3dmake_all_lessons.md__step-2--robust-error-handling}

\subsubsection*{What "Robust" Means}\label{docs__pandoc__latex__src__3dmake_all_lessons.md__what-robust-means}

A script is \emph{robust} if it handles failure gracefully: it detects when something went wrong, reports the failure clearly, and stops before making things worse. A \emph{fragile} script ignores errors and continues --- which often means you end up with a folder full of broken or empty output files.

\subsubsection*{The Three Failure Modes to Guard Against}\label{docs__pandoc__latex__src__3dmake_all_lessons.md__the-three-failure-modes-to-guard-against}

\textbf{Silent failure:} The command runs, returns an error code, but the script doesn\textquotesingle t check and keeps going. You discover the problem later when the output files don\textquotesingle t work.

\textbf{Partial output:} The build starts, produces a partially-written file, then crashes. The file exists but is corrupted. Your script needs to check that the output is complete and has a reasonable size.

\textbf{Cascade failure:} One failure causes subsequent steps to fail too, making the error messages confusing. Stop at the first failure --- don\textquotesingle t keep running.

\subsubsection*{Error Handling in bash}\label{docs__pandoc__latex__src__3dmake_all_lessons.md__error-handling-in-bash}

Bash has three critical safety settings that should be at the top of every automation script:

\begin{itemize}
\tightlist
\item
  \texttt{set\ -e} --- exit immediately if any command returns a non-zero exit code (non-zero = failure)
\item
  \texttt{set\ -u} --- error if you reference a variable that was never defined (catches typos in variable names)
\item
  \texttt{set\ -o\ pipefail} --- catch errors inside pipelines, not just the last command in the pipe
\end{itemize}

Think of these as your script\textquotesingle s "safety mode." Without them, bash will happily continue running after errors. With them, the first problem causes an immediate stop.

\begin{lstlisting}[style=Alabaster, language=bash]
#!/bin/bash

# build.sh — robust build script with logging and error handling


set -e           # exit immediately if any command returns a non-zero exit code

set -u           # error on use of any unset variable

set -o pipefail  # catch errors inside pipelines (not just the last command)


LOG_DIR="build/logs"

LOG_FILE="$LOG_DIR/build_$(date +%Y%m%d_%H%M%S).log"

mkdir -p "$LOG_DIR" build


# Logging function: prints with timestamp and appends to log file

log() {

  local level="$1"

  shift

  echo "[$(date +%H:%M:%S)] [$level] $*" | tee -a "$LOG_FILE"

}


# Trap: runs when any error occurs, reports which line failed

on_error() {

  log "ERROR" "Build failed at line $1. Check: $LOG_FILE"

  exit 1

}

trap 'on_error $LINENO' ERR


log "INFO" "=== Build started ==="

log "INFO" "Source: src/main.scad"


# Build — capture stderr, show on screen, and append to log

if ! openscad -o build/main.stl src/main.scad 2>&1 | tee -a "$LOG_FILE"; then

  log "ERROR" "OpenSCAD exited with error"

  exit 1

fi


# Verify the output file was actually created and has reasonable size

STL_PATH="build/main.stl"

if [ ! -f "$STL_PATH" ]; then

  log "ERROR" "STL file not created: $STL_PATH"

  exit 1

fi


STL_SIZE=$(stat -c%s "$STL_PATH")

if [ "$STL_SIZE" -lt 1000 ]; then

  log "WARN" "STL file suspiciously small: ${STL_SIZE} bytes — may be empty geometry"

fi


log "INFO" "Build complete: $STL_PATH (${STL_SIZE} bytes)"

3dm info 2>&1 | tee -a "$LOG_FILE"

log "INFO" "=== Build finished ==="


\end{lstlisting}

\subsubsection*{Error Handling in PowerShell}\label{docs__pandoc__latex__src__3dmake_all_lessons.md__error-handling-in-powershell}

PowerShell uses try/catch blocks and checks the \texttt{\$LASTEXITCODE} variable (which holds the exit code of the last external command) to detect failures. The key difference from bash: PowerShell doesn\textquotesingle t automatically stop on error --- you must check \texttt{\$LASTEXITCODE} yourself after every external command.

\begin{lstlisting}[style=Alabaster, language=powershell]
# build.ps1 — robust build script for Windows


param(

    [string]$Source    = "src\main.scad",

    [string]$Output    = "build\main.stl",

    [string]$LogDir    = "build\logs"

)


$timestamp = Get-Date -Format "yyyyMMdd_HHmmss"

New-Item -ItemType Directory -Force -Path $LogDir, "build" | Out-Null

$logFile = Join-Path $LogDir "build_$timestamp.log"


function Write-Log {

    param([string]$Level, [string]$Message)

    $entry = "[$(Get-Date -Format 'HH:mm:ss')] [$Level] $Message"

    Write-Host $entry

    Add-Content -Path $logFile -Value $entry

}


Write-Log "INFO" "=== Build started ==="

Write-Log "INFO" "Source: $Source"


try {

    $buildOutput = & openscad -o $Output $Source 2>&1

    $buildOutput | Tee-Object -FilePath $logFile -Append | Write-Host

    if ($LASTEXITCODE -ne 0) { throw "OpenSCAD returned exit code $LASTEXITCODE" }

} catch {

    Write-Log "ERROR" "Build failed: $_"

    exit 1

}


if (-not (Test-Path $Output)) {

    Write-Log "ERROR" "STL file not created: $Output"

    exit 1

}


$stlSize = (Get-Item $Output).Length

if ($stlSize -lt 1000) {

    Write-Log "WARN" "STL suspiciously small: $stlSize bytes"

}


Write-Log "INFO" "Build complete: $Output ($stlSize bytes)"

3dm info 2>&1 | Tee-Object -FilePath $logFile -Append

Write-Log "INFO" "=== Build finished ==="


\end{lstlisting}

\begin{center}\rule{0.5\linewidth}{0.5pt}\end{center}

\subsection*{Step 3 --- Batch Builds and Variant Generation}\label{docs__pandoc__latex__src__3dmake_all_lessons.md__step-3--batch-builds-and-variant-generation}

\subsubsection*{The Concept: Parametric Variants at Scale}\label{docs__pandoc__latex__src__3dmake_all_lessons.md__the-concept-parametric-variants-at-scale}

Now that you can build reliably with error handling, you can automate the generation of multiple variants. The pattern is simple: loop over a list of parameter combinations, run a build for each, name the output file to reflect the parameters used.

Good output naming is critical. A file called \texttt{main.stl} tells you nothing. A file called \texttt{bracket\_w80\_h60\_wall2.stl} tells you exactly what parameters produced it --- you can reconstruct the build command just from the filename.

\subsubsection*{Batch Build in bash}\label{docs__pandoc__latex__src__3dmake_all_lessons.md__batch-build-in-bash}

\begin{lstlisting}[style=Alabaster, language=bash]
#!/bin/bash

# batch_variants.sh


set -e

set -u

set -o pipefail


OUTPUT_DIR="build/variants"

mkdir -p "$OUTPUT_DIR"


SUCCESS=0

FAIL=0

FAILED_BUILDS=()


for WIDTH in 60 80 100 120; do

  for WALL in 1.2 1.6 2.0; do

    NAME="bracket_w${WIDTH}_wall${WALL/./_}"

    OUT="$OUTPUT_DIR/${NAME}.stl"


    echo "Building: $NAME ..."

    if openscad -D "width=${WIDTH}" -D "wall=${WALL}" \

                -o "$OUT" src/main.scad 2>/dev/null; then

      echo "  OK: $OUT"

      SUCCESS=$((SUCCESS + 1))

    else

      echo "  FAILED: $NAME"

      FAIL=$((FAIL + 1))

      FAILED_BUILDS+=("$NAME")

    fi

  done

done


echo ""

echo "=== Batch complete: $SUCCESS succeeded, $FAIL failed ==="

if [ ${#FAILED_BUILDS[@]} -gt 0 ]; then

  echo "Failed builds:"

  for b in "${FAILED_BUILDS[@]}"; do

    echo "  - $b"

  done

  exit 1

fi


\end{lstlisting}

\subsubsection*{Batch Build in PowerShell}\label{docs__pandoc__latex__src__3dmake_all_lessons.md__batch-build-in-powershell}

\begin{lstlisting}[style=Alabaster, language=powershell]
# batch_variants.ps1


$outputDir = "build\variants"

New-Item -ItemType Directory -Force -Path $outputDir | Out-Null


$success = 0

$fail    = 0

$failed  = @()


foreach ($width in @(60, 80, 100, 120)) {

  foreach ($wall in @(1.2, 1.6, 2.0)) {

    $wallStr = "$wall" -replace '\.', '_'

    $name    = "bracket_w${width}_wall${wallStr}"

    $out     = "$outputDir\${name}.stl"


    Write-Host "Building: $name ..."

    & openscad -D "width=$width" -D "wall=$wall" -o $out src\main.scad 2>$null

    if ($LASTEXITCODE -eq 0 -and (Test-Path $out)) {

      Write-Host "  OK: $out"

      $success++

    } else {

      Write-Host "  FAILED: $name"

      $fail++

      $failed += $name

    }

  }

}


Write-Host ""

Write-Host "=== Batch complete: $success succeeded, $fail failed ==="

if ($failed.Count -gt 0) {

  Write-Host "Failed builds:"

  $failed | ForEach-Object { Write-Host "  - $_" }

  exit 1

}


\end{lstlisting}

\begin{center}\rule{0.5\linewidth}{0.5pt}\end{center}

\subsection*{Step 4 --- Timestamped Archiving}\label{docs__pandoc__latex__src__3dmake_all_lessons.md__step-4--timestamped-archiving}

\subsubsection*{Why Archive Builds?}\label{docs__pandoc__latex__src__3dmake_all_lessons.md__why-archive-builds}

Imagine you printed a part two weeks ago that fit perfectly. Now you\textquotesingle ve made some changes to the model, and the new version doesn\textquotesingle t fit. You need to go back to the version that worked --- but you only have one \texttt{build/main.stl} file, which has been overwritten.

Timestamped archiving solves this. Every successful build is saved with a timestamp in its filename. You can always go back to any previous version. Over time, you build a history of exactly when each change was made.

\subsubsection*{Archive Script}\label{docs__pandoc__latex__src__3dmake_all_lessons.md__archive-script}

\begin{lstlisting}[style=Alabaster, language=bash]
#!/bin/bash

# archive_build.sh


set -e

set -u


ARCHIVE_DIR="build/archive"

mkdir -p "$ARCHIVE_DIR"


TIMESTAMP=$(date +%Y%m%d_%H%M%S)

ARCHIVE_NAME="main_${TIMESTAMP}.stl"


# Build

3dm build


# Archive the successful result

cp build/main.stl "$ARCHIVE_DIR/$ARCHIVE_NAME"

echo "Archived: $ARCHIVE_DIR/$ARCHIVE_NAME"


# Optional: remove archives older than 30 days

find "$ARCHIVE_DIR" -name "*.stl" -mtime +30 -delete

echo "Old archives cleaned."


\end{lstlisting}

\begin{center}\rule{0.5\linewidth}{0.5pt}\end{center}

\subsection*{Step 5 --- Watch Mode}\label{docs__pandoc__latex__src__3dmake_all_lessons.md__step-5--watch-mode}

\subsubsection*{The Concept: Automatic Rebuild on Save}\label{docs__pandoc__latex__src__3dmake_all_lessons.md__the-concept-automatic-rebuild-on-save}

Watch mode is a script that monitors your source files and automatically runs a build whenever a file changes. Instead of switching to your terminal and running \texttt{3dm\ build} every time you save a change in your editor, the script does it for you.

This creates a tight feedback loop: save your \texttt{.scad} file, glance at the terminal, and the new build result is already there. It\textquotesingle s like having an assistant who runs the build the instant you save.

\subsubsection*{How Watch Mode Works}\label{docs__pandoc__latex__src__3dmake_all_lessons.md__how-watch-mode-works}

There are two approaches:

\textbf{Poll-based:} Check the file\textquotesingle s modification timestamp every few seconds. If it changed since the last check, trigger a rebuild. Simple to implement; works everywhere. Slightly wasteful --- you\textquotesingle re checking constantly even when nothing changes.

\textbf{Event-based:} The operating system notifies you the instant a file changes (using \texttt{inotify} on Linux, \texttt{FSEvents} on macOS, \texttt{FileSystemWatcher} on Windows). More efficient; responds immediately. More complex to set up.

For development workflows, the poll-based approach is usually good enough.

\begin{lstlisting}[style=Alabaster, language=bash]
#!/bin/bash

# watch_and_build.sh — poll-based watch mode


set -u


WATCH_FILE="src/main.scad"

POLL_INTERVAL=2   # seconds between checks


echo "Watching: $WATCH_FILE (Ctrl+C to stop)"


LAST_MOD=""


while true; do

  CURRENT_MOD=$(stat -c%Y "$WATCH_FILE" 2>/dev/null || echo "0")


  if [ "$CURRENT_MOD" != "$LAST_MOD" ]; then

    LAST_MOD="$CURRENT_MOD"

    if [ -n "$LAST_MOD" ] && [ "$LAST_MOD" != "0" ]; then

      echo ""

      echo "[$(date +%H:%M:%S)] Change detected — rebuilding..."

      3dm build && echo "  Build OK" || echo "  Build FAILED"

    fi

  fi


  sleep "$POLL_INTERVAL"

done


\end{lstlisting}

\begin{lstlisting}[style=Alabaster, language=powershell]
# watch_and_build.ps1 — poll-based watch mode for Windows


$watchFile    = "src\main.scad"

$pollInterval = 2  # seconds


Write-Host "Watching: $watchFile (Ctrl+C to stop)"


$lastMod = $null


while ($true) {

    if (Test-Path $watchFile) {

        $currentMod = (Get-Item $watchFile).LastWriteTime


        if ($currentMod -ne $lastMod) {

            $lastMod = $currentMod

            if ($null -ne $lastMod) {

                Write-Host ""

                Write-Host "[$(Get-Date -Format 'HH:mm:ss')] Change detected — rebuilding..."

                & 3dm build

                if ($LASTEXITCODE -eq 0) {

                    Write-Host "  Build OK"

                } else {

                    Write-Host "  Build FAILED"

                }

            }

        }

    }

    Start-Sleep -Seconds $pollInterval

}


\end{lstlisting}

\begin{center}\rule{0.5\linewidth}{0.5pt}\end{center}

\subsection*{Step 6 --- CI/CD Integration}\label{docs__pandoc__latex__src__3dmake_all_lessons.md__step-6--cicd-integration}

\subsubsection*{What CI/CD Means}\label{docs__pandoc__latex__src__3dmake_all_lessons.md__what-cicd-means}

CI/CD stands for \textbf{Continuous Integration / Continuous Deployment}. In software development, CI/CD means that every time you push code to a shared repository, an automated system builds and tests your project in a clean environment --- not on your computer, but on a server in the cloud.

For a 3D printing project, CI/CD means:

\begin{itemize}
\tightlist
\item
  Every time you commit a change to your \texttt{.scad} files, GitHub automatically builds your STL
\item
  If the build fails (syntax error, empty geometry, etc.), you get an email notification immediately
\item
  The built STL is saved as a downloadable artifact --- you can prove to anyone that this exact commit produced this exact file
\item
  You can track exactly when changes were made and whether they broke anything
\end{itemize}

This sounds like overkill for a hobby project --- but even for personal projects, CI/CD is a useful safety net. It catches problems you might miss locally, and it creates a reliable record of your design history.

\subsubsection*{How GitHub Actions Works}\label{docs__pandoc__latex__src__3dmake_all_lessons.md__how-github-actions-works}

GitHub Actions is a CI/CD system built into GitHub. You define a "workflow" in a YAML file inside your repository. The workflow runs automatically on events you specify --- such as every push to the \texttt{main} branch.

The workflow runs on a fresh Ubuntu virtual machine in GitHub\textquotesingle s cloud. It has no memory of previous builds. Every run starts from scratch: clone the repository, install dependencies, run the build, save the output.

\begin{lstlisting}[style=Alabaster]
# .github/workflows/build.yml

# This file tells GitHub Actions what to do on every push.


name: Build STL


on:

  push:

    paths:

      - "src/**" # only trigger when .scad files change


jobs:

  build:

    runs-on: ubuntu-latest

    steps:

      - uses: actions/checkout@v4


      - name: Install OpenSCAD

        run: |

          sudo apt-get update

          sudo apt-get install -y openscad


      - name: Install 3dMake

        run: curl -fsSL https://get.3dmake.dev | bash


      - name: Build

        run: |

          3dm build

          3dm info


      - name: Verify STL exists and has reasonable size

        run: |

          test -f build/main.stl

          SIZE=$(stat -c%s build/main.stl)

          echo "STL size: $SIZE bytes"

          test "$SIZE" -gt 1000


      - name: Upload artifact

        uses: actions/upload-artifact@v4

        with:

          name: stl-output-${{ github.sha }}

          path: build/main.stl


\end{lstlisting}

When this is set up, every push to the repository triggers a fresh build in GitHub\textquotesingle s cloud. If the build fails, you get an email notification. The built STL is saved as an artifact you can download directly from the GitHub interface.

\subsubsection*{What Each Section Does}\label{docs__pandoc__latex__src__3dmake_all_lessons.md__what-each-section-does}

\textbf{\texttt{on:\ push:\ paths:}} --- Only run the workflow when files in the \texttt{src/} directory change. This prevents the workflow from running when you update the README, for example.

\textbf{\texttt{runs-on:\ ubuntu-latest}} --- Use GitHub\textquotesingle s latest Ubuntu virtual machine. This is a standard Linux environment with nothing pre-installed.

\textbf{\texttt{actions/checkout@v4}} --- Clone your repository into the virtual machine. Without this, the machine has no idea what your project looks like.

\textbf{\texttt{Install\ OpenSCAD} / \texttt{Install\ 3dMake}} --- Install your build tools. Every run starts fresh, so tools must be installed every time.

\textbf{\texttt{Build}} --- Run your build. If this fails (non-zero exit code), the whole workflow fails.

\textbf{\texttt{Verify\ STL\ exists\ and\ has\ reasonable\ size}} --- An explicit sanity check. \texttt{test\ -f} verifies the file exists. \texttt{test\ "\$SIZE"\ -gt\ 1000} verifies it\textquotesingle s not empty.

\textbf{\texttt{Upload\ artifact}} --- Save the built STL so you can download it from GitHub. The \texttt{\$\{\{\ github.sha\ \}\}} part includes the commit hash in the artifact name, so you always know which commit produced which STL.

\begin{center}\rule{0.5\linewidth}{0.5pt}\end{center}

\subsection*{Exercises}\label{docs__pandoc__latex__src__3dmake_all_lessons.md__exercises-4}

\textbf{Exercise 9.1:} Adapt the batch\_variants script to your platform. Generate your phone stand at 5 different \texttt{stand\_angle} values (30°, 45°, 60°, 70°, 80°). Log the bounding box for each to a file.

\textbf{Exercise 9.2:} Write a build script in your preferred shell that: (1) builds the model, (2) verifies the STL size is \textgreater{} 10,000 bytes, (3) archives the result with a timestamp if the build succeeds, (4) prints a failure message to the console if it fails.

\textbf{Exercise 9.3 (Advanced):} Create a parametric \texttt{Makefile} (Linux/macOS) or a PowerShell \texttt{Invoke-Build} task file (Windows) with separate targets: \texttt{build}, \texttt{clean}, \texttt{variants}, \texttt{archive}, and \texttt{test}. Running \texttt{make\ all} should run all targets in the correct order.

\begin{center}\rule{0.5\linewidth}{0.5pt}\end{center}

\subsection*{Quiz --- Lesson 3dMake.9 (15 questions)}\label{docs__pandoc__latex__src__3dmake_all_lessons.md__quiz--lesson-3dmake9-15-questions}

\begin{enumerate}
\tightlist
\item
  What does \texttt{set\ -e} do in a bash script?
\item
  Why is passing string parameters with \texttt{-D} more complex than passing numeric parameters?
\item
  What does \texttt{trap\ \textquotesingle{}on\_error\ \$LINENO\textquotesingle{}\ ERR} do in a bash script?
\item
  What is the correct way to pass \texttt{label="FRONT"} as an OpenSCAD \texttt{-D} argument in bash?
\item
  What does \texttt{set\ -u} protect against?
\item
  How would you verify that an output STL was actually created in a build script?
\item
  What is a "silent failure" in the context of automation, and how do you prevent one?
\item
  In PowerShell, what character is used to escape a double quote inside a double-quoted string?
\item
  What does \texttt{\$LASTEXITCODE} contain in a PowerShell script after running an external command?
\item
  Describe the purpose of a timestamped build archive and when you would use it.
\item
  What is CI/CD, and how does it apply to a 3D printing project workflow?
\item
  True or False: In CMD, doubled double-quotes inside a quoted string produce a single literal double-quote.
\item
  Explain the difference between a poll-based file watcher and an event-based file watcher. What are the trade-offs?
\item
  In a GitHub Actions workflow, what does \texttt{on:\ push:\ paths:\ -\ \textquotesingle{}src/**\textquotesingle{}} specify?
\item
  If a batch build script runs 10 builds and each may succeed or fail independently, describe the logic for reporting: total built, total failed, and which specific builds failed.
\end{enumerate}

\begin{center}\rule{0.5\linewidth}{0.5pt}\end{center}

\subsection*{Extension Problems (15)}\label{docs__pandoc__latex__src__3dmake_all_lessons.md__extension-problems-15-8}

\begin{enumerate}
\tightlist
\item
  Write a batch build script that tests 10 parameter combinations and generates a summary report with pass/fail status for each.
\item
  Create a "golden master" verification script: build your reference model, save the bounding box, then verify that future builds match those dimensions within 0.1 mm.
\item
  Implement the timestamped archive system and add a rotation policy: automatically delete archives older than 30 days.
\item
  Build a Git pre-commit hook that runs \texttt{3dm\ build} and fails the commit if the build doesn\textquotesingle t succeed.
\item
  Create a build dashboard: a script that runs all your projects\textquotesingle{} builds, collects their \texttt{3dm\ info} outputs, and formats them into a summary report.
\item
  Implement parallel builds: run multiple OpenSCAD builds simultaneously using background processes or PowerShell jobs. Measure the speedup.
\item
  Create a "changelog" generator: compare the bounding boxes of all archived builds and generate a list of which builds had dimension changes.
\item
  Write a CI/CD pipeline configuration for GitLab CI or Bitbucket Pipelines (instead of GitHub Actions).
\item
  Build a build notification system: send a Slack message or email when a build fails.
\item
  Create a "test suite" of deliberately bad \texttt{.scad} files (empty geometry, non-manifold shapes, syntax errors) and verify your build script catches all of them.
\item
  Implement semantic versioning in your build system: automatically increment a patch version number on each successful build.
\item
  Write a multi-platform build script that detects the current OS and uses the appropriate commands (bash on Linux/macOS, PowerShell on Windows).
\item
  Create a build matrix: for each of 5 designs × 4 materials, build the design with material-specific parameters and verify each output.
\item
  Implement a "regression test": store the bounding box of each successful build, and flag any build where the bounding box changes unexpectedly.
\item
  Research and implement "build caching": skip rebuilds when the source file hasn\textquotesingle t changed since the last successful build.
\end{enumerate}

\begin{center}\rule{0.5\linewidth}{0.5pt}\end{center}

\subsection*{References and Helpful Resources}\label{docs__pandoc__latex__src__3dmake_all_lessons.md__references-and-helpful-resources-8}

\begin{itemize}
\item
  Bash Reference Manual --- https://www.gnu.org/software/bash/manual/bash.html
\item
  PowerShell Documentation --- https://docs.microsoft.com/en-us/powershell/
\item
  GitHub Actions Documentation --- https://docs.github.com/en/actions
\item
  3DMake GitHub Repository --- https://github.com/tdeck/3dmake
\end{itemize}

\subsubsection*{Supplemental Resources}\label{docs__pandoc__latex__src__3dmake_all_lessons.md__supplemental-resources-8}

\begin{itemize}
\tightlist
\item
  Programming with OpenSCAD EPUB Textbook --- assets/Programming\_with\_OpenSCAD.epub
\item
  3DMake GitHub Repository --- https://github.com/tdeck/3dmake
\item
  OpenSCAD User Manual --- https://en.wikibooks.org/wiki/OpenSCAD\_User\_Manual
  -e
\end{itemize}

\begin{center}\rule{0.5\linewidth}{0.5pt}\end{center}

\section*{Lesson 10 --- Hands-On Practice Exercises and Troubleshooting}\label{docs__pandoc__latex__src__3dmake_all_lessons.md__lesson-10--hands-on-practice-exercises-and-troubleshooting}

\textbf{Estimated time:} 120--150 minutes

\begin{center}\rule{0.5\linewidth}{0.5pt}\end{center}

\subsection*{Before You Start}\label{docs__pandoc__latex__src__3dmake_all_lessons.md__before-you-start-8}

You\textquotesingle ve built a substantial toolkit over the first nine lessons: primitives, CSG operations, parametric modules, transforms, tolerances, interlocking features, automation scripts. This lesson is where those skills converge. Rather than introducing a single new concept, Lesson 10 is about \emph{depth} --- applying advanced OpenSCAD features to real design problems, understanding when to use each tool, and developing the troubleshooting instincts that separate experienced designers from beginners.

The advanced features in this lesson --- \texttt{hull()}, \texttt{minkowski()}, \texttt{projection()}, \texttt{surface()}, \texttt{children()}, and \texttt{search()} --- are not obscure corners of OpenSCAD. They\textquotesingle re the tools that unlock the difference between "looks like a 3D-printed object" and "looks like a professionally designed product."

\begin{center}\rule{0.5\linewidth}{0.5pt}\end{center}

\subsection*{Learning Objectives}\label{docs__pandoc__latex__src__3dmake_all_lessons.md__learning-objectives-15}

By the end of this lesson you will be able to:

\begin{itemize}
\tightlist
\item
  Apply \texttt{hull()} and \texttt{minkowski()} with awareness of their performance cost
\item
  Use \texttt{\$preview} to conditionally disable slow operations during development
\item
  Create 2D templates with \texttt{projection()} for drilling guides and laser cutting
\item
  Diagnose and fix all four common types of non-manifold geometry
\item
  Apply \texttt{surface()} for height-map terrain and embossed artwork
\item
  Write advanced parametric assemblies using \texttt{children()} and \texttt{\$children}
\item
  Use \texttt{search()} and \texttt{lookup()} for data-driven parametric design
\end{itemize}

\begin{center}\rule{0.5\linewidth}{0.5pt}\end{center}

\subsection*{Concept: From Functional to Refined}\label{docs__pandoc__latex__src__3dmake_all_lessons.md__concept-from-functional-to-refined}

Up to now, most designs have used right-angle geometry --- boxes, cylinders, straight edges. This produces functional parts, but real products rarely look like that. Consider a smartphone: it has rounded corners, smooth transitions between surfaces, organic curves. None of those were added by accident --- each one was a deliberate design choice that improves how the product feels in your hand and looks on a shelf.

The tools in this lesson let you add that kind of refinement to your OpenSCAD designs without abandoning the parametric approach. You\textquotesingle ll still define everything with numbers and variables --- but the shapes those numbers produce will be much more sophisticated.

\begin{center}\rule{0.5\linewidth}{0.5pt}\end{center}

\subsection*{\texorpdfstring{Step 1 --- \texttt{hull()} for Organic Shapes}{Step 1 --- hull() for Organic Shapes}}\label{docs__pandoc__latex__src__3dmake_all_lessons.md__step-1--hull-for-organic-shapes}

\subsubsection*{\texorpdfstring{What \texttt{hull()} Does}{What hull() Does}}\label{docs__pandoc__latex__src__3dmake_all_lessons.md__what-hull-does-1}

Imagine stretching a rubber band around several objects placed on a table. The rubber band forms the tightest possible loop that touches all of them, following the outside curve of each object. This is exactly what \texttt{hull()} does in 3D: it computes the \textbf{convex hull} --- the smallest convex shape that contains all its children.

"Convex" means the shape never curves inward. A sphere is convex. A donut is not (it has a hole --- a concave feature). The convex hull of two spheres is like a capsule --- the two spheres connected by a smooth, outward-curving surface.

\subsubsection*{\texorpdfstring{When to Use \texttt{hull()}}{When to Use hull()}}\label{docs__pandoc__latex__src__3dmake_all_lessons.md__when-to-use-hull}

\texttt{hull()} is excellent for:

\begin{itemize}
\tightlist
\item
  \textbf{Smooth transitions} between two different shapes at different heights (like the neck of a bottle smoothly widening into the body)
\item
  \textbf{Rounded boxes} --- place a small sphere at each corner, take the hull, and every corner and edge becomes rounded
\item
  \textbf{Organic profiles} --- aerofoil cross-sections, ergonomic handles, teardrop shapes
\item
  \textbf{Anywhere} you want to avoid sharp 90-degree corners without adding complex geometry
\end{itemize}

\texttt{hull()} is \textbf{much faster} than \texttt{minkowski()} (covered in Step 2) and should be your first choice for organic shapes.

\subsubsection*{Important Limitation}\label{docs__pandoc__latex__src__3dmake_all_lessons.md__important-limitation}

\texttt{hull()} only produces \textbf{convex} shapes. If you need a shape that curves inward --- a saddle, a cup, a grip with a waist --- \texttt{hull()} alone can\textquotesingle t do it. You\textquotesingle d need to use \texttt{hull()} as a building block and then use \texttt{difference()} to carve in the concave features.

\subsubsection*{Example 1: Smooth Transition Between a Hexagonal Cylinder and a Sphere}\label{docs__pandoc__latex__src__3dmake_all_lessons.md__example-1-smooth-transition-between-a-hexagonal-cylinder-and-a-sphere}

\begin{lstlisting}[style=Alabaster, language=openscad]
// hull() between two different shapes at different heights

// Creates a smooth transition (like a funnel or transition piece)

hull() {

  cylinder(r=20, h=1, $fn=6);       // hexagon at bottom

  translate([0, 0, 30])

    sphere(r=8, $fn=32);             // sphere at top

}


\end{lstlisting}

This produces a smooth funnel-like shape --- hexagonal at the base, tapering to a sphere at the top. Without \texttt{hull()}, you\textquotesingle d need complex math to define the intermediate surface.

\subsubsection*{\texorpdfstring{Example 2: Rounded Box Using \texttt{hull()} of 8 Spheres}{Example 2: Rounded Box Using hull() of 8 Spheres}}\label{docs__pandoc__latex__src__3dmake_all_lessons.md__example-2-rounded-box-using-hull-of-8-spheres}

Place one sphere at each of the 8 corners; \texttt{hull()} fills the space between them, rounding every corner and edge by the sphere radius.

\begin{lstlisting}[style=Alabaster, language=openscad]
r = 5;  // corner radius

w = 40; d = 30; h = 20;


hull() {

  for (x = [r, w-r]) {

    for (y = [r, d-r]) {

      translate([x, y, 0]) sphere(r=r, $fn=24);

      translate([x, y, h]) sphere(r=r, $fn=24);

    }

  }

}


\end{lstlisting}

The sphere radius \texttt{r} controls how rounded the corners are. Setting \texttt{r=1} gives subtle rounding; \texttt{r=10} gives very rounded, pillow-like corners.

\subsubsection*{Example 3: Wing / Aerofoil Profile}\label{docs__pandoc__latex__src__3dmake_all_lessons.md__example-3-wing--aerofoil-profile}

\begin{lstlisting}[style=Alabaster, language=openscad]
// hull() of two cylinders at different positions and radii

// produces a smooth aerofoil cross-section

translate([0, 50, 0])

hull() {

  cylinder(r=6, h=2, $fn=32);            // leading edge (round)

  translate([40, 0, 0])

    cylinder(r=1, h=2, $fn=16);          // trailing edge (sharp)

}


\end{lstlisting}

The larger cylinder becomes the rounded leading edge; the small cylinder becomes the sharp trailing edge. \texttt{hull()} constructs the smooth surface connecting them automatically.

\begin{center}\rule{0.5\linewidth}{0.5pt}\end{center}

\subsection*{\texorpdfstring{Step 2 --- \texttt{minkowski()} --- Performance Warning and Pattern}{Step 2 --- minkowski() --- Performance Warning and Pattern}}\label{docs__pandoc__latex__src__3dmake_all_lessons.md__step-2--minkowski--performance-warning-and-pattern}

\subsubsection*{What Minkowski Sum Means}\label{docs__pandoc__latex__src__3dmake_all_lessons.md__what-minkowski-sum-means}

\texttt{minkowski()} computes the \textbf{Minkowski sum}: it places a copy of the secondary shape at every point on the surface of the primary shape, then takes the union of all those copies. The result is as if you "rolled" the secondary shape across every surface, corner, and edge of the primary shape.

The practical effect: if your primary shape is a cube and your secondary shape is a small sphere, the result is a cube with perfectly rounded corners and edges --- where the rounding radius equals the sphere radius.

This is mathematically precise rounding that handles any geometry, including complex shapes with many faces. It\textquotesingle s also how you create "offset" versions of shapes --- the Minkowski sum with a sphere produces a shape that is uniformly larger in all directions by the sphere\textquotesingle s radius.

\subsubsection*{The Performance Warning}\label{docs__pandoc__latex__src__3dmake_all_lessons.md__the-performance-warning}

Here\textquotesingle s the problem: the face count of the Minkowski result is roughly \textbf{the product} of the face counts of the two input shapes. A cube has 6 faces. A sphere with \texttt{\$fn=32} has about 1,000 faces. The Minkowski sum has roughly 6,000 faces to compute. Now substitute a complex curved primary shape with 500 faces, and you\textquotesingle re computing 500,000 faces. This can take minutes on a modern computer.

\begin{quote}
⚠ A complex primary shape combined with a sphere at \texttt{\$fn=64} can take many minutes to render. Always use \texttt{\$preview} to skip the \texttt{minkowski()} call during development.
\end{quote}

\subsubsection*{\texorpdfstring{\texttt{\$preview}: Fast Development, Accurate Rendering}{\$preview: Fast Development, Accurate Rendering}}\label{docs__pandoc__latex__src__3dmake_all_lessons.md__preview-fast-development-accurate-rendering}

OpenSCAD has two rendering modes:

\begin{itemize}
\tightlist
\item
  \textbf{F5 (Preview):} Fast, approximate. Uses OpenCSG rendering. \texttt{\$preview} is \texttt{true}.
\item
  \textbf{F6 (Full Render):} Exact, uses CGAL library. Much slower but produces accurate manifold geometry for export. \texttt{\$preview} is \texttt{false}.
\end{itemize}

Use \texttt{\$preview} as a conditional: show a fast placeholder during development (F5), and run the expensive \texttt{minkowski()} only during final render (F6).

\begin{quote}
\textbf{Version note:} \texttt{\$preview} is available in \textbf{OpenSCAD 2019.05 and later}. In earlier versions it is undefined (treated as \texttt{false}). If you are using an older version, omit the \texttt{\$preview} guard and only run F6 when you need the final result.
\end{quote}

\begin{lstlisting}[style=Alabaster, language=openscad]
// $preview is TRUE during F5 (fast preview) and FALSE during F6 (full render).

// Use it to swap between a fast placeholder and the full Minkowski computation.


module rounded_box(w, d, h, r=4) {

  if ($preview) {

    // Fast preview mode: plain cube — renders instantly.

    // Color it orange as a reminder this is NOT the final shape.

    color("orange", 0.6) cube([w, d, h]);

  } else {

    // Full render mode: proper Minkowski sum.

    // Keep the secondary sphere LOW-POLY to limit face-count explosion.

    minkowski() {

      cube([w - 2*r, d - 2*r, h - r]);

      sphere(r=r, $fn=12);  // keep $fn LOW for minkowski — 12 is usually enough

    }

  }

}


rounded_box(50, 40, 20, r=5);


\end{lstlisting}

\subsubsection*{\texorpdfstring{\texttt{hull()} as a Faster \texttt{minkowski()} Alternative}{hull() as a Faster minkowski() Alternative}}\label{docs__pandoc__latex__src__3dmake_all_lessons.md__hull-as-a-faster-minkowski-alternative}

\texttt{hull()} only works for convex shapes, but it is far faster than \texttt{minkowski()}. For a simple rounded box you can get an equivalent result by placing a sphere at each corner and taking the hull:

\begin{lstlisting}[style=Alabaster, language=openscad]
translate([0, 60, 0])

hull() {

  // Equivalent to minkowski(cube, sphere) — much faster for convex shapes

  r = 4;

  w = 50; d = 40; h = 16;

  for (x = [r, w-r]) for (y = [r, d-r]) for (z = [r, h-r]) {

    translate([x, y, z]) sphere(r=r, $fn=12);

  }

}


\end{lstlisting}

\textbf{Performance guidelines for \texttt{minkowski()}:}

{\def\LTcaptype{none} % do not increment counter
\begin{longtable}[]{@{}
  >{\raggedright\arraybackslash}p{(\linewidth - 4\tabcolsep) * \real{0.3276}}
  >{\raggedright\arraybackslash}p{(\linewidth - 4\tabcolsep) * \real{0.2586}}
  >{\raggedright\arraybackslash}p{(\linewidth - 4\tabcolsep) * \real{0.4138}}@{}}
\toprule\noalign{}
\begin{minipage}[b]{\linewidth}\raggedright
Primary shape faces
\end{minipage} & \begin{minipage}[b]{\linewidth}\raggedright
Secondary 
\end{minipage} & \begin{minipage}[b]{\linewidth}\raggedright
\texttt{\$fn}
\end{minipage} \\
\midrule\noalign{}
\endhead
\bottomrule\noalign{}
\endlastfoot
\textless{} 100 & 8--16 & \textless{} 1 second \\
100--500 & 8--16 & 1--10 seconds \\
500--2000 & 8 & 10--60 seconds \\
\textgreater{} 2000 & any & Minutes to hours --- avoid \\
\end{longtable}
}

\textbf{Rule of thumb:} If the primary shape is a simple primitive (cube, cylinder, basic extrusion), \texttt{minkowski()} is manageable. If the primary shape is itself a complex CSG result, use \texttt{hull()} instead.

\begin{center}\rule{0.5\linewidth}{0.5pt}\end{center}

\subsection*{\texorpdfstring{Step 3 --- \texttt{projection()} for 2D Templates}{Step 3 --- projection() for 2D Templates}}\label{docs__pandoc__latex__src__3dmake_all_lessons.md__step-3--projection-for-2d-templates}

\subsubsection*{The Concept: Flattening 3D to 2D}\label{docs__pandoc__latex__src__3dmake_all_lessons.md__the-concept-flattening-3d-to-2d}

\texttt{projection()} takes a 3D shape and flattens it onto the XY plane, producing a 2D outline. Think of shining a light straight down onto a 3D object: the shadow it casts on the floor is the projection.

This is useful in fabrication contexts because real-world assembly often requires 2D templates:

\begin{itemize}
\tightlist
\item
  A \textbf{drilling template} --- print a flat sheet that shows exactly where to drill holes in a wall, shelf, or PCB
\item
  A \textbf{laser-cut profile} --- export the 2D shape of a part as a DXF for a laser cutter to follow
\item
  A \textbf{footprint check} --- verify that your design will actually fit in the space you have before printing
\item
  A \textbf{CNC machining path} --- export a cross-section as a DXF for a CNC router
\end{itemize}

\texttt{projection()} has two modes:

\begin{itemize}
\tightlist
\item
  \texttt{cut=false} --- projects the full silhouette (shadow from above): every part of the shape, regardless of height
\item
  \texttt{cut=true} --- projects a cross-section \emph{at z=0}: only the parts of the shape at exactly that height plane
\end{itemize}

\begin{lstlisting}[style=Alabaster, language=openscad]
// projection(cut=false) — complete silhouette (shadow from above)

// projection(cut=true)  — cross-section at z=0 (slice through the model)


module phone_stand_example() {

  cube([90, 70, 5]);

  translate([15, 20, 5]) cylinder(r=5, h=35, $fn=24);  // mounting post A

  translate([75, 20, 5]) cylinder(r=5, h=35, $fn=24);  // mounting post B

}


// 3D model (for printing)

phone_stand_example();


// 2D silhouette (footprint) — placed beside the 3D model for comparison

color("LightGreen")

translate([110, 0, 0])

  linear_extrude(1)   // 1mm thick so it is visible in 3D preview

    projection(cut=false)

      phone_stand_example();


// Cross-section at z=3 — shows internal structure at that height

color("LightCoral")

translate([220, 0, 0])

  linear_extrude(1)

    projection(cut=true)

      translate([0, 0, -3])  // shift model so the cut falls at z=3 of the part

        phone_stand_example();


\end{lstlisting}

\subsubsection*{Exporting DXF for Laser Cutting or CNC}\label{docs__pandoc__latex__src__3dmake_all_lessons.md__exporting-dxf-for-laser-cutting-or-cnc}

\begin{quote}
\textbf{Note:} \texttt{3dm} is designed for building STL output. Use the \texttt{openscad} CLI directly for DXF export --- do \textbf{not} use \texttt{3dm\ build} for this step.
\end{quote}

\begin{lstlisting}[style=Alabaster, language=bash]
# Linux / macOS

openscad -o template.dxf src/template_projection.scad


# Windows (PowerShell or CMD)

openscad -o template.dxf src\\template_projection.scad


\end{lstlisting}

\textbf{Best practice: keep projection files separate.} Create a dedicated \texttt{.scad} file for DXF export that imports your main model and wraps it in \texttt{projection()}. This keeps your main model file clean.

\begin{lstlisting}[style=Alabaster, language=openscad]
// src/template_projection.scad

// This file is used exclusively for DXF export.

// It should contain only projection() — no 3D geometry.


use <../src/main.scad>


// Project the full silhouette

projection(cut=false)

  main_assembly();


// Add mounting-hole markers as circles.

// DXF circles are standard and ready for CNC drilling.

hole_positions = [[10, 10], [80, 10], [10, 60], [80, 60]];

for (pos = hole_positions) {

  translate(pos)

    circle(r=1.75, $fn=16);  // 3.5 mm diameter holes

}


\end{lstlisting}

\begin{center}\rule{0.5\linewidth}{0.5pt}\end{center}

\subsection*{Step 4 --- Diagnosing Non-Manifold Geometry}\label{docs__pandoc__latex__src__3dmake_all_lessons.md__step-4--diagnosing-non-manifold-geometry}

\subsubsection*{What "Manifold" Means and Why It Matters}\label{docs__pandoc__latex__src__3dmake_all_lessons.md__what-manifold-means-and-why-it-matters}

A mesh is \textbf{manifold} (also called "watertight") when it follows two strict rules:

\begin{enumerate}
\tightlist
\item
  Every edge is shared by exactly two faces --- no more, no fewer
\item
  The surface has no holes, zero-thickness walls, or self-intersections
\end{enumerate}

Think of a manifold mesh like a sealed plastic bag. You can fill it with water and it won\textquotesingle t leak anywhere. Every point on the surface is clearly either "inside" or "outside." Slicers need this to determine which areas to fill with material.

A non-manifold mesh is like a bag with a hole in it, or two bags fused at a single edge with no material connecting them. Slicers can\textquotesingle t slice it correctly --- they might produce missing layers, extra walls, or simply refuse to import the file.

\subsubsection*{How OpenSCAD Produces Non-Manifold Geometry}\label{docs__pandoc__latex__src__3dmake_all_lessons.md__how-openscad-produces-non-manifold-geometry}

OpenSCAD\textquotesingle s CSG operations are mathematically defined, but floating-point arithmetic and co-planar face interactions can create subtle mesh problems. The four most common causes:

\subsubsection*{\texorpdfstring{Problem 1 --- Co-Planar Faces in \texttt{difference()}}{Problem 1 --- Co-Planar Faces in difference()}}\label{docs__pandoc__latex__src__3dmake_all_lessons.md__problem-1--co-planar-faces-in-difference}

When the cutter exactly coincides with the face of the base shape, OpenSCAD cannot determine which side of the face to keep. The result is an ambiguous, possibly non-manifold edge. Fix: extend the cutter by 0.001 mm on each side so it clearly passes through the material.

\begin{lstlisting}[style=Alabaster, language=openscad]
difference() {

  cube([20, 20, 20]);

  translate([5, 5, -0.001])           // extend 0.001 below base

    cylinder(r=5, h=20.002, $fn=24);  // extend 0.001 above top

}


\end{lstlisting}

This is the same 0.001 mm rule introduced in Lesson 2. It\textquotesingle s worth repeating because forgetting it is the most common cause of non-manifold geometry in student work.

\subsubsection*{Problem 2 --- Zero-Thickness or Negative-Thickness Walls}\label{docs__pandoc__latex__src__3dmake_all_lessons.md__problem-2--zero-thickness-or-negative-thickness-walls}

If your wall thickness calculation ever produces zero or a negative number, the resulting geometry has no material --- creating a degenerate mesh.

\begin{lstlisting}[style=Alabaster, language=openscad]
module thin_wall_box(w, d, h, wall=2) {

  assert(wall > 0 && wall < min(w,d)/2,

    str("wall must be 0 < wall < ", min(w,d)/2, ", got: ", wall));

  difference() {

    cube([w, d, h]);

    translate([wall, wall, wall])

      cube([w - 2*wall, d - 2*wall, h]);  // open top — valid manifold

  }

}


\end{lstlisting}

The \texttt{assert()} guard here catches the problem before it reaches the renderer. This is better than discovering it in the slicer.

\subsubsection*{Problem 3 --- Shapes That Only Touch (Not Overlap)}\label{docs__pandoc__latex__src__3dmake_all_lessons.md__problem-3--shapes-that-only-touch-not-overlap}

Two shapes that share a face but do not overlap create a non-manifold edge --- an edge that belongs to faces from both shapes but with no clear "inside." Fix: overlap by 0.001 mm and wrap in \texttt{union()}.

\begin{lstlisting}[style=Alabaster, language=openscad]
module two_boxes() {

  union() {

    cube([10, 10, 10]);

    translate([10 - 0.001, 0, 0])   // overlap by 0.001 — not just touch

      cube([10, 10, 10]);

  }

}


\end{lstlisting}

\subsubsection*{\texorpdfstring{Problem 4 --- Inverted Normals from Negative \texttt{scale()}}{Problem 4 --- Inverted Normals from Negative scale()}}\label{docs__pandoc__latex__src__3dmake_all_lessons.md__problem-4--inverted-normals-from-negative-scale}

\texttt{scale()} with a negative value flips face normals (the direction each face "points"), creating inside-out geometry. The shape looks correct in the preview but has every face pointing the wrong direction --- which slicers interpret as solid material where there should be empty space.

\begin{lstlisting}[style=Alabaster, language=openscad]
// WRONG — negative scale flips normals:

// scale([-1, 1, 1]) some_module();


// CORRECT — mirror() preserves normals:

mirror([1, 0, 0]) some_module();


\end{lstlisting}

\subsubsection*{Debugging Tip: The Ghost Modifier}\label{docs__pandoc__latex__src__3dmake_all_lessons.md__debugging-tip-the-ghost-modifier}

\begin{lstlisting}[style=Alabaster, language=openscad]
// Use the % modifier to render a shape as a transparent ghost

// so you can see inside the model and verify cutter placement.

difference() {

  cube([20, 20, 20]);

  % cylinder(r=5, h=25, $fn=24);  // ghost the cutter

}


\end{lstlisting}

The \texttt{\%} modifier makes the cutter semi-transparent in preview mode, letting you visually verify it\textquotesingle s positioned and sized correctly before committing to the render.

\subsubsection*{Manifold Check Workflow}\label{docs__pandoc__latex__src__3dmake_all_lessons.md__manifold-check-workflow}

\begin{lstlisting}[style=Alabaster, language=bash]
# Build the STL

3dm build


# OpenSCAD prints manifold warnings to the console, e.g.:

#   WARNING: Object may not be a valid 2-manifold ...

# Check for this message after every build.


# For more detail, use admesh:

# Linux:   apt install admesh

# macOS:   brew install admesh

admesh --analyze build/main.stl


# Or open the STL in PrusaSlicer and enable "Check for geometry errors"

# to see orange highlights on problem edges.


\end{lstlisting}

\begin{quote}
⚠ The correct OpenSCAD manifold warning is \textbf{"WARNING: Object may not be a valid 2-manifold"}. Searching for any other string will cause you to miss real errors.
\end{quote}

\begin{center}\rule{0.5\linewidth}{0.5pt}\end{center}

\subsection*{\texorpdfstring{Step 5 --- \texttt{surface()} for Height-Map Geometry}{Step 5 --- surface() for Height-Map Geometry}}\label{docs__pandoc__latex__src__3dmake_all_lessons.md__step-5--surface-for-height-map-geometry}

\subsubsection*{The Concept: Images as 3D Terrain}\label{docs__pandoc__latex__src__3dmake_all_lessons.md__the-concept-images-as-3d-terrain}

\texttt{surface()} reads a grayscale image or a \texttt{.dat} data file and creates a 3D surface where pixel brightness maps to Z height. Black pixels (value 0) produce no height; white pixels (value 255) produce maximum height. Every shade of gray in between produces an intermediate height, interpolated smoothly between pixels.

Think of a topographic map where elevation is encoded in color --- \texttt{surface()} works exactly like that. You can create terrain landscapes from satellite elevation data, embossed artwork from any image, textured panels from grayscale photographs, or precise surface profiles from engineering data.

\subsubsection*{\texorpdfstring{Why \texttt{surface()} Is Powerful}{Why surface() Is Powerful}}\label{docs__pandoc__latex__src__3dmake_all_lessons.md__why-surface-is-powerful}

Traditional CSG modeling works from the outside in: you define shapes geometrically and combine them. \texttt{surface()} lets you work from \emph{data} --- the shape emerges from a dataset rather than from geometric primitives. This is how architectural visualization software creates terrain, how prosthetic designers create textured grip surfaces, and how engineers create test panels with complex surface profiles.

\subsubsection*{\texorpdfstring{Example 1: Minimal Self-Contained \texttt{surface()} from a \texttt{.dat} File}{Example 1: Minimal Self-Contained surface() from a .dat File}}\label{docs__pandoc__latex__src__3dmake_all_lessons.md__example-1-minimal-self-contained-surface-from-a-dat-file}

Save the matrix below as \texttt{terrain.dat} in the same folder as your \texttt{.scad} file, then run the \texttt{surface()} call. No external image editor required.

\begin{lstlisting}[style=Alabaster]
 0  10  20  10   0

10  40  60  40  10

20  60 100  60  20

10  40  60  40  10

 0  10  20  10   0


\end{lstlisting}

\begin{lstlisting}[style=Alabaster, language=openscad]
// In your .scad file:

scale([5, 5, 0.3])   // each cell = 5 mm wide; height scaled from 0–100

  surface(file="terrain.dat", center=true);


\end{lstlisting}

This produces a small mountain shape --- highest at the center, tapering outward. The \texttt{scale()} controls the physical dimensions: X/Y scale determines how wide each data cell becomes, and Z scale determines the vertical exaggeration.

\subsubsection*{Example 2: Embossed Logo Panel}\label{docs__pandoc__latex__src__3dmake_all_lessons.md__example-2-embossed-logo-panel}

Create a grayscale PNG (logo = white, background = black), save it as \texttt{src/logo\_heightmap.png}, then use the module below.

\begin{lstlisting}[style=Alabaster, language=openscad]
module logo_panel(panel_w=60, panel_d=40, panel_h=3, relief=1.5) {

  // Base plate

  cube([panel_w, panel_d, panel_h]);


  // Embossed surface on top.

  // Assumes logo_heightmap.png is 100×100 pixels; scale to fit the panel.

  translate([0, 0, panel_h])

    scale([panel_w/100, panel_d/100, relief/255])

      surface(file="src/logo_heightmap.png", center=false, invert=false);

}


logo_panel();


\end{lstlisting}

\begin{quote}
✓ \texttt{surface()} supports PNG (grayscale or colour --- uses luminance) and \texttt{.dat} (whitespace-separated matrix). Always check that the image dimensions match the aspect ratio you want before scaling.
\end{quote}

\begin{center}\rule{0.5\linewidth}{0.5pt}\end{center}

\subsection*{\texorpdfstring{Step 6 --- Advanced: \texttt{children()} and Custom Container Modules}{Step 6 --- Advanced: children() and Custom Container Modules}}\label{docs__pandoc__latex__src__3dmake_all_lessons.md__step-6--advanced-children-and-custom-container-modules}

\subsubsection*{The Concept: Modules That Accept Geometry as Input}\label{docs__pandoc__latex__src__3dmake_all_lessons.md__the-concept-modules-that-accept-geometry-as-input}

So far, every module you\textquotesingle ve written produces geometry from parameters. \texttt{children()} lets you write modules that \emph{wrap} geometry --- they accept shapes as input and do something with them. This is like writing a picture frame that works with any picture, rather than one specific picture built into the frame.

The pattern is called a \textbf{container module}: a module that defines a structural context (a base plate, a grid, a mounting system) and places its children in that context. The caller provides the geometry; the module provides the infrastructure.

\subsubsection*{\texorpdfstring{\texttt{\$children}: How Many Children Were Passed}{\$children: How Many Children Were Passed}}\label{docs__pandoc__latex__src__3dmake_all_lessons.md__children-how-many-children-were-passed}

\texttt{\$children} is an automatic variable that contains the count of child objects passed to a module. You can use it to loop over children with \texttt{children(i)} to access each one individually by index (0-based).

\subsubsection*{\texorpdfstring{Example 1: \texttt{mounted\_on\_plate} --- Place Any Geometry on a Base with Corner Posts}{Example 1: mounted\_on\_plate --- Place Any Geometry on a Base with Corner Posts}}\label{docs__pandoc__latex__src__3dmake_all_lessons.md__example-1-mounted_on_plate--place-any-geometry-on-a-base-with-corner-posts}

\begin{lstlisting}[style=Alabaster, language=openscad]
module mounted_on_plate(plate_w, plate_d, plate_h=3, post_h=5) {

  // Base plate

  color("LightGray")

    cube([plate_w, plate_d, plate_h]);


  // Corner mounting posts

  for (x = [5, plate_w - 5]) {

    for (y = [5, plate_d - 5]) {

      translate([x, y, plate_h])

        cylinder(r=3, h=post_h, $fn=16);

    }

  }


  // Child geometry sits on top of the posts

  translate([0, 0, plate_h + post_h])

    children();

}


// Usage: any geometry wrapped in mounted_on_plate()

mounted_on_plate(80, 60, plate_h=3, post_h=8) {

  cylinder(r=15, h=20, $fn=32);

}


\end{lstlisting}

\texttt{children()} here is called without an index, which passes \emph{all} children as a group. The cylinder in the usage example will be placed on top of the posts --- without the module knowing or caring what shape it is.

\subsubsection*{\texorpdfstring{Example 2: \texttt{repeat\_children} --- Tile Multiple Named Children in a Grid}{Example 2: repeat\_children --- Tile Multiple Named Children in a Grid}}\label{docs__pandoc__latex__src__3dmake_all_lessons.md__example-2-repeat_children--tile-multiple-named-children-in-a-grid}

\texttt{\$children} gives you the count of children passed. Use \texttt{children(i)} to render the i-th child individually.

\begin{lstlisting}[style=Alabaster, language=openscad]
module repeat_children(n=3, spacing=25) {

  for (i = [0 : $children - 1]) {

    for (j = [0 : n - 1]) {

      translate([j * spacing, i * spacing, 0])

        children(i);  // render the i-th child

    }

  }

}


translate([0, 100, 0])

repeat_children(n=4, spacing=20) {

  sphere(r=5, $fn=16);           // child 0

  cylinder(r=4, h=10, $fn=16);   // child 1

  cube([8, 8, 6]);                // child 2

}


\end{lstlisting}

This tiles each child shape in its own row --- spheres in the first row, cylinders in the second, cubes in the third --- repeated 4 times horizontally. The module doesn\textquotesingle t know what shapes it\textquotesingle s tiling; it just repeats whatever is passed to it.

\begin{quote}
✓ OpenSCAD does \textbf{not} expose bounding-box queries at runtime. Always pass dimensions explicitly to container modules rather than trying to compute them from children.
\end{quote}

\begin{center}\rule{0.5\linewidth}{0.5pt}\end{center}

\subsection*{\texorpdfstring{Step 7 --- Advanced: Lookup Tables and \texttt{search()}}{Step 7 --- Advanced: Lookup Tables and search()}}\label{docs__pandoc__latex__src__3dmake_all_lessons.md__step-7--advanced-lookup-tables-and-search}

\subsubsection*{The Concept: Data-Driven Parametric Design}\label{docs__pandoc__latex__src__3dmake_all_lessons.md__the-concept-data-driven-parametric-design}

So far, parametric design meant: change a number, the shape changes. This step takes it further: \textbf{data-driven design} means the shape is determined by structured data --- tables, lists, databases of values.

Instead of hardcoding \texttt{density\ =\ 1.24} for PLA and remembering to change it manually when you switch materials, you define a material table and let the code look up the right values automatically. The design reacts to the material choice; you just change one variable.

This is how professional parametric CAD systems work. A single master file can describe a family of parts at different sizes, materials, and configurations --- all driven by lookup tables.

\subsubsection*{\texorpdfstring{\texttt{lookup()}: Smooth Interpolation from a Table}{lookup(): Smooth Interpolation from a Table}}\label{docs__pandoc__latex__src__3dmake_all_lessons.md__lookup-smooth-interpolation-from-a-table}

\texttt{lookup()} finds an output value by interpolating between entries in a \texttt{{[}input,\ output{]}} table. If the input falls between two entries, the output is linearly interpolated between them.

\begin{lstlisting}[style=Alabaster, language=openscad]
// Temperature-to-fan-speed lookup table

speed_table = [

  [0,   0],   //  0°C -> fan off

  [40,  0],   // 40°C -> fan off

  [50, 50],   // 50°C -> 50% speed

  [60, 80],   // 60°C -> 80% speed

  [70, 100]   // 70°C -> full speed

];


temp = 55;

fan_speed = lookup(temp, speed_table);

echo("Fan speed at", temp, "°C:", fan_speed, "%");  // -> 65%


\end{lstlisting}

\texttt{lookup()} is especially useful for smoothly varying geometric properties --- rounding radii that increase with part size, wall thicknesses that scale with load, fin counts that increase with diameter.

\subsubsection*{\texorpdfstring{\texttt{search()}: Exact Match in a List}{search(): Exact Match in a List}}\label{docs__pandoc__latex__src__3dmake_all_lessons.md__search-exact-match-in-a-list}

\texttt{search()} finds the index of a value in a list, enabling exact-match lookup for categorical data like material names.

\begin{quote}
⚠ \texttt{search()} returns a \textbf{list of lists}. Check \texttt{len(matches{[}0{]})\ \textgreater{}\ 0} for a valid result --- \texttt{idx\ !=\ {[}{]}} does \textbf{not} work correctly in OpenSCAD.
\end{quote}

\begin{lstlisting}[style=Alabaster, language=openscad]
materials = ["PLA",  "PETG", "ABS",  "TPU",  "Nylon"];

densities = [ 1.24,   1.27,   1.05,   1.20,   1.14];


selected = "PETG";

idx = search([selected], materials)[0];

density = densities[idx];

echo(str(selected, " density: ", density, " g/cm³"));


\end{lstlisting}

\subsubsection*{\texorpdfstring{Example 3: Print Cost Estimator with \texttt{search()} and \texttt{assert()}}{Example 3: Print Cost Estimator with search() and assert()}}\label{docs__pandoc__latex__src__3dmake_all_lessons.md__example-3-print-cost-estimator-with-search-and-assert}

This module demonstrates the full pattern: a lookup table for material properties, \texttt{search()} to find the right row, \texttt{assert()} to catch invalid material names, and computed geometry that reacts to all of it.

\begin{lstlisting}[style=Alabaster, language=openscad]
module print_cost_label(volume_mm3, material="PLA") {

  materials = ["PLA",  "PETG", "ABS",  "TPU",  "Nylon"];

  densities = [ 1.24,   1.27,   1.05,   1.20,   1.14];

  costs_g   = [ 0.020,  0.025,  0.018,  0.035,  0.045];


  // search() returns a list of lists.

  // matches[0] is the first result: an index on success, [] on failure.

  matches = search([material], materials);

  assert(len(matches[0]) > 0,

    str("Unknown material: \"", material, "\". Valid: ", materials));

  idx = matches[0];


  mass_g = (volume_mm3 / 1000) * densities[idx];

  cost   = mass_g * costs_g[idx];

  label  = str(material, ": ", round(mass_g * 10) / 10, "g / $",

               round(cost * 100) / 100);


  echo("Cost estimate:", label);

  linear_extrude(1)

    text(label, size=4, font="Liberation Sans", halign="center");

}


print_cost_label(12000, "PETG");


\end{lstlisting}

\begin{center}\rule{0.5\linewidth}{0.5pt}\end{center}

\subsection*{Exercise Set A: Phone Stand}\label{docs__pandoc__latex__src__3dmake_all_lessons.md__exercise-set-a-phone-stand}

\subsubsection*{A1 --- Hull Transition}\label{docs__pandoc__latex__src__3dmake_all_lessons.md__a1--hull-transition}

Add a smooth \texttt{hull()} transition between the base plate and the mounting post:

\begin{lstlisting}[style=Alabaster, language=openscad]
hull() {

  translate([15, 20, 0]) cylinder(r=8, h=1, $fn=32);   // post footprint on base

  translate([15, 20, 5]) cylinder(r=5, h=35, $fn=24);  // actual post

}


\end{lstlisting}

\subsubsection*{A2 --- Tolerance Stack-Up}\label{docs__pandoc__latex__src__3dmake_all_lessons.md__a2--tolerance-stack-up}

Measure the three main structural parts of your phone stand and record them in a table. Calculate worst-case stack-up:

\textbf{Worst case = sum of all individual tolerances}

Example for three parts each at ±0.15 mm: 0.15 + 0.15 + 0.15 = \textbf{±0.45 mm total}

Document your own measurements in a similar table.

\subsubsection*{A3 --- Add a Cable Slot}\label{docs__pandoc__latex__src__3dmake_all_lessons.md__a3--add-a-cable-slot}

Extend your phone stand design with a cable slot through the base:

\begin{lstlisting}[style=Alabaster, language=openscad]
cable_slot_w  = 12;   // mm

cable_slot_d  = 5;    // mm

cable_slot_z  = -0.001;


// Add inside your main difference() block:

translate([base_w/2 - cable_slot_w/2, 0, cable_slot_z])

  cube([cable_slot_w, cable_slot_d, base_h + 0.002]);


\end{lstlisting}

\begin{center}\rule{0.5\linewidth}{0.5pt}\end{center}

\subsection*{Exercise Set B: Keycap with Text}\label{docs__pandoc__latex__src__3dmake_all_lessons.md__exercise-set-b-keycap-with-text}

\subsubsection*{B1 --- Build a Mechanical Keyboard Keycap}\label{docs__pandoc__latex__src__3dmake_all_lessons.md__b1--build-a-mechanical-keyboard-keycap}

\begin{lstlisting}[style=Alabaster, language=openscad]
// Parametric keycap for Cherry MX switches

key_w      = 18;

key_d      = 18;

key_h      = 7;

stem_r     = 2.75;  // MX stem: 5.5 mm diameter

stem_h     = 3.8;

wall       = 1.5;

label_text = "A";


module keycap() {

  difference() {

    // Keycap body with slight top curve

    hull() {

      cube([key_w, key_d, key_h - 2], center=true);

      translate([0, 0, 1]) cube([key_w - 2, key_d - 2, key_h], center=true);

    }

    // Hollow inside

    translate([0, 0, -wall])

      cube([key_w - 2*wall, key_d - 2*wall, key_h], center=true);

    // MX stem hole

    translate([0, 0, -(key_h/2 + 0.001)])

      cylinder(r=stem_r + 0.1, h=stem_h + 0.001, $fn=16);

  }

}


module stem_mount() {

  translate([0, 0, -(key_h/2 + stem_h)])

    difference() {

      cylinder(r=stem_r + wall, h=stem_h, $fn=16);

      cylinder(r=stem_r, h=stem_h + 0.001, $fn=16);

    }

}


keycap();

stem_mount();


// Engrave label on top

translate([0, 0, key_h/2 - 0.8])

  linear_extrude(1.2)

    text(label_text, size=8, font="Liberation Sans:style=Bold",

         halign="center", valign="center", $fn=4);


\end{lstlisting}

\subsubsection*{\texorpdfstring{B2 --- Validate with \texttt{3dm\ info}\footnote{3DMake GitHub Repository --- Command reference including \texttt{3dm\ info}.
  https://github.com/tdeck/3dmake}}{B2 --- Validate with 3dm info}}\label{docs__pandoc__latex__src__3dmake_all_lessons.md__b2--validate-with-3dm-info}

\begin{lstlisting}[style=Alabaster, language=bash]
3dm info


\end{lstlisting}

Document what the model description reports and compare it to your design intent.

\subsubsection*{B3 --- Print and Test}\label{docs__pandoc__latex__src__3dmake_all_lessons.md__b3--print-and-test}

Print the keycap and test it on a Cherry MX switch (or compatible). If the stem is too tight, increase the clearance from \texttt{stem\_r\ +\ 0.1} to \texttt{stem\_r\ +\ 0.15}. If too loose, decrease to \texttt{stem\_r\ +\ 0.05}.

\begin{center}\rule{0.5\linewidth}{0.5pt}\end{center}

\subsection*{Exercise Set C: Stackable Bins}\label{docs__pandoc__latex__src__3dmake_all_lessons.md__exercise-set-c-stackable-bins}

\subsubsection*{C1 --- Build a Three-Size Bin Set}\label{docs__pandoc__latex__src__3dmake_all_lessons.md__c1--build-a-three-size-bin-set}

\begin{lstlisting}[style=Alabaster, language=openscad]
// Small bin

translate([0, 0, 0])

  bin_assembly(bin_w=60, bin_d=45, bin_h=30);


// Medium bin

translate([80, 0, 0])

  bin_assembly(bin_w=80, bin_d=60, bin_h=40);


// Large bin

translate([180, 0, 0])

  bin_assembly(bin_w=100, bin_d=80, bin_h=50);


\end{lstlisting}

\subsubsection*{C2 --- Diagnose and Fix Non-Manifold Geometry}\label{docs__pandoc__latex__src__3dmake_all_lessons.md__c2--diagnose-and-fix-non-manifold-geometry}

Non-manifold geometry occurs when faces share edges inconsistently (T-junctions, missing faces, zero-thickness walls). A common cause is two shapes that touch at exactly one face:

\begin{lstlisting}[style=Alabaster, language=openscad]
// PROBLEM: two cubes share a face exactly — may produce a non-manifold edge

cube([20, 20, 10]);

translate([20, 0, 0]) cube([20, 20, 10]);  // touching at x=20


// FIX 1: use union()

union() {

  cube([20, 20, 10]);

  translate([20, 0, 0]) cube([20, 20, 10]);

}


// FIX 2: overlap slightly

union() {

  cube([20.001, 20, 10]);

  translate([20, 0, 0]) cube([20, 20, 10]);

}


\end{lstlisting}

\begin{lstlisting}[style=Alabaster, language=bash]
# Diagnosis

3dm info  # AI will often flag non-manifold geometry

# Also open the STL in your slicer and enable "Check for geometry errors"


\end{lstlisting}

\subsubsection*{\texorpdfstring{C3 --- Advanced Geometry: \texttt{hull()} and \texttt{minkowski()}}{C3 --- Advanced Geometry: hull() and minkowski()}}\label{docs__pandoc__latex__src__3dmake_all_lessons.md__c3--advanced-geometry-hull-and-minkowski}

\begin{lstlisting}[style=Alabaster, language=openscad]
// hull() creates a convex envelope — useful for organic transitions

module smooth_transition() {

  hull() {

    translate([0, 0, 0])  cylinder(r=15, h=5, $fn=64);

    translate([0, 0, 30]) cylinder(r=5,  h=2, $fn=64);

  }

}

smooth_transition();


// minkowski() adds the secondary shape to every surface point

// Result: all edges rounded by 2 mm

module rounded_hull() {

  minkowski() {

    hull() {

      cylinder(r=10, h=3, $fn=8);           // octagonal prism

      translate([30, 0, 0]) sphere(r=8, $fn=32);

    }

    sphere(r=2, $fn=16);  // rounds all edges by 2 mm

  }

}

rounded_hull();


\end{lstlisting}

\begin{center}\rule{0.5\linewidth}{0.5pt}\end{center}

\subsection*{Quiz --- Lesson 3dMake.10 (15 questions)}\label{docs__pandoc__latex__src__3dmake_all_lessons.md__quiz--lesson-3dmake10-15-questions}

\begin{enumerate}
\tightlist
\item
  What tool do you use to measure printed part dimensions against the design specification?
\item
  What is tolerance stack-up, and why does it matter for multi-part assemblies?
\item
  What causes non-manifold geometry in OpenSCAD, and how do you detect it?
\item
  How does \texttt{hull()} differ from \texttt{union()}?
\item
  What does \texttt{3dm\ info} help you verify about your model?
\item
  What is the diameter of a Cherry MX stem, and what clearance would you add for a slip-fit keycap?
\item
  True or False: \texttt{find\ -newer} is an event-driven file change detection method.
\item
  If three parts each have ±0.15 mm tolerance, what is the worst-case total error for a three-part stack?
\item
  What does the \texttt{\$fn} parameter control in OpenSCAD?
\item
  Describe two methods for fixing non-manifold geometry caused by two touching (but not overlapping) shapes.
\item
  What is the difference between \texttt{hull()} and \texttt{minkowski()} for creating organic shapes? Give one use case for each.
\item
  What does \texttt{resize({[}50,\ 0,\ 0{]})} do, and why might \texttt{resize()} behave unexpectedly for non-uniform scaling?
\item
  When measuring a printed part with calipers\footnote{Digital Calipers Measurement Technique --- see Appendix B of this document.}, what is the difference between an inside measurement and an outside measurement, and when does that distinction matter for tolerance analysis?
\item
  Describe the iterative workflow for dialling in press-fit tolerances: what do you print, what do you measure, and how do you adjust?
\item
  If \texttt{3dm\ info} reports "the model appears non-manifold," what are three possible causes you would investigate in your OpenSCAD code?
\end{enumerate}

\begin{center}\rule{0.5\linewidth}{0.5pt}\end{center}

\subsection*{Extension Problems (15)}\label{docs__pandoc__latex__src__3dmake_all_lessons.md__extension-problems-15-9}

\begin{enumerate}
\tightlist
\item
  Create a tolerance sensitivity study: build 5 keycaps with stem clearance from 0.05--0.25 mm in 0.05 mm increments, print them, and record which values fit your switches.
\item
  Design a go/no-go gauge for a 10 mm nominal hole: a part with a "go" pin sized for slip-fit and a "no-go" pin sized for interference fit.
\item
  Write a printer calibration SOP: bed levelling, first-layer calibration, and dimension verification. Include a measurement checklist.
\item
  Build a three-tier stackable storage system for art supplies. Each tier has a different inner grid.
\item
  Conduct a tolerance stack-up analysis for your stackable bin system. Calculate worst-case misalignment.
\item
  Build a parametric test coupon that tests four wall thicknesses (0.8, 1.2, 1.6, 2.0 mm) in a single print.
\item
  Design a caliper stand: a holder that keeps your digital calipers at a comfortable angle for one-handed operation.
\item
  Build a non-manifold error catalogue: intentionally create 5 different types of non-manifold geometry, document how each was created, and how to fix it.
\item
  Use \texttt{hull()} to design a smooth ergonomic tool handle and compare it to a simple cylinder handle.
\item
  Create a printability checklist for new designs: overhangs, wall thickness, minimum feature size, support requirements. Apply it to your keycap and bin designs.
\item
  Research the \texttt{resize()} function. Build an example showing how it behaves differently from \texttt{scale()} for non-uniform resizing.
\item
  Design a multi-part assembly tutorial: a three-piece interlocking puzzle that teaches tolerance, alignment, and slip-fit.
\item
  Build a "measurement worksheet" template in OpenSCAD: render a flat sheet that lists all key dimensions of a part as text, for printing alongside the part.
\item
  Create a chi-squared goodness-of-fit test for your printer\textquotesingle s dimensional accuracy: measure 20 prints of the same part and determine whether the errors are normally distributed.
\item
  Write a comprehensive troubleshooting guide covering the 10 most common 3D printing failures you have encountered or researched, with causes, prevention, and fixes.
\end{enumerate}

\begin{center}\rule{0.5\linewidth}{0.5pt}\end{center}

\subsection*{References and Helpful Resources}\label{docs__pandoc__latex__src__3dmake_all_lessons.md__references-and-helpful-resources-9}

\begin{itemize}
\tightlist
\item
  OpenSCAD User Manual --- Hull and Minkowski.
  https://en.wikibooks.org/wiki/OpenSCAD\_User\_Manual/Minkowski\_and\_Hull
\end{itemize}

\subsubsection*{Supplemental Resources}\label{docs__pandoc__latex__src__3dmake_all_lessons.md__supplemental-resources-9}

\begin{itemize}
\tightlist
\item
  Programming with OpenSCAD EPUB Textbook --- assets/Programming\_with\_OpenSCAD.epub
\item
  CodeSolutions Repository: https://github.com/ProgrammingWithOpenSCAD/CodeSolutions
\item
  OpenSCAD Quick Reference: https://programmingwithopenscad.github.io/quick-reference.html
\item
  All other reference documents are embedded as appendices below --- no separate files needed.
\end{itemize}

\begin{center}\rule{0.5\linewidth}{0.5pt}\end{center}

\section*{Part 2: Reference Appendices}\label{docs__pandoc__latex__src__3dmake_all_lessons.md__part-2-reference-appendices-1}

\emph{All supporting documents for Lesson 10 are included below. You do not need to open any additional files.}

\begin{center}\rule{0.5\linewidth}{0.5pt}\end{center}

\subsection*{Appendix A: Measurement Worksheet}\label{docs__pandoc__latex__src__3dmake_all_lessons.md__appendix-a-measurement-worksheet}

\textbf{Use this for:} Measuring real-world objects before modeling. Record three measurements per feature, then use the average in your OpenSCAD code --- never a single measurement.

\textbf{Student Name:} \_\_\_\_\_\_\_\_\_\_\_\_\_\_\_\_\\
\textbf{Date:} \_\_\_\_\_\_\_\_\_\_\\
\textbf{Project / Assignment:} \_\_\_\_\_\_\_\_\_\_\_\_\_\_\_\_

\subsubsection*{How to Use This Worksheet}\label{docs__pandoc__latex__src__3dmake_all_lessons.md__how-to-use-this-worksheet}

\begin{enumerate}
\tightlist
\item
  Measure each feature three times and record all three values.
\item
  Calculate the average: add the three values and divide by 3.
\item
  Round to one decimal place (e.g., 23.4 mm).
\item
  Use the average in your OpenSCAD code --- not any single measurement.
\end{enumerate}

\emph{Units: All measurements in millimeters (mm) unless otherwise noted.}

\begin{center}\rule{0.5\linewidth}{0.5pt}\end{center}

\subsubsection*{Object 1}\label{docs__pandoc__latex__src__3dmake_all_lessons.md__object-1}

Object description: \_\_\_\_\_\_\_\_\_\_\_\_\_\_\_\_

{\def\LTcaptype{none} % do not increment counter
\begin{longtable}[]{@{}
  >{\raggedright\arraybackslash}p{(\linewidth - 10\tabcolsep) * \real{0.1562}}
  >{\raggedright\arraybackslash}p{(\linewidth - 10\tabcolsep) * \real{0.3281}}
  >{\centering\arraybackslash}p{(\linewidth - 10\tabcolsep) * \real{0.1094}}
  >{\centering\arraybackslash}p{(\linewidth - 10\tabcolsep) * \real{0.1094}}
  >{\centering\arraybackslash}p{(\linewidth - 10\tabcolsep) * \real{0.1094}}
  >{\centering\arraybackslash}p{(\linewidth - 10\tabcolsep) * \real{0.1875}}@{}}
\toprule\noalign{}
\begin{minipage}[b]{\linewidth}\raggedright
Feature
\end{minipage} & \begin{minipage}[b]{\linewidth}\raggedright
What You\textquotesingle re Measuring
\end{minipage} & \begin{minipage}[b]{\linewidth}\centering
M1 (mm)
\end{minipage} & \begin{minipage}[b]{\linewidth}\centering
M2 (mm)
\end{minipage} & \begin{minipage}[b]{\linewidth}\centering
M3 (mm)
\end{minipage} & \begin{minipage}[b]{\linewidth}\centering
Average (mm)
\end{minipage} \\
\midrule\noalign{}
\endhead
\bottomrule\noalign{}
\endlastfoot
Length (X) & & & & & \\
Width (Y) & & & & & \\
Height (Z) & & & & & \\
Feature 4 & & & & & \\
Feature 5 & & & & & \\
\end{longtable}
}

Notes / sketches:

\begin{lstlisting}[style=Alabaster]
(describe the object here — label which direction is X, Y, Z)


\end{lstlisting}

\begin{center}\rule{0.5\linewidth}{0.5pt}\end{center}

\subsubsection*{Object 2}\label{docs__pandoc__latex__src__3dmake_all_lessons.md__object-2}

Object description: \_\_\_\_\_\_\_\_\_\_\_\_\_\_\_\_

{\def\LTcaptype{none} % do not increment counter
\begin{longtable}[]{@{}
  >{\raggedright\arraybackslash}p{(\linewidth - 10\tabcolsep) * \real{0.1562}}
  >{\raggedright\arraybackslash}p{(\linewidth - 10\tabcolsep) * \real{0.3281}}
  >{\centering\arraybackslash}p{(\linewidth - 10\tabcolsep) * \real{0.1094}}
  >{\centering\arraybackslash}p{(\linewidth - 10\tabcolsep) * \real{0.1094}}
  >{\centering\arraybackslash}p{(\linewidth - 10\tabcolsep) * \real{0.1094}}
  >{\centering\arraybackslash}p{(\linewidth - 10\tabcolsep) * \real{0.1875}}@{}}
\toprule\noalign{}
\begin{minipage}[b]{\linewidth}\raggedright
Feature
\end{minipage} & \begin{minipage}[b]{\linewidth}\raggedright
What You\textquotesingle re Measuring
\end{minipage} & \begin{minipage}[b]{\linewidth}\centering
M1 (mm)
\end{minipage} & \begin{minipage}[b]{\linewidth}\centering
M2 (mm)
\end{minipage} & \begin{minipage}[b]{\linewidth}\centering
M3 (mm)
\end{minipage} & \begin{minipage}[b]{\linewidth}\centering
Average (mm)
\end{minipage} \\
\midrule\noalign{}
\endhead
\bottomrule\noalign{}
\endlastfoot
Length (X) & & & & & \\
Width (Y) & & & & & \\
Height (Z) & & & & & \\
Feature 4 & & & & & \\
Feature 5 & & & & & \\
\end{longtable}
}

Notes / sketches:

\begin{lstlisting}[style=Alabaster]
(describe the object here)


\end{lstlisting}

\begin{center}\rule{0.5\linewidth}{0.5pt}\end{center}

\subsubsection*{Object 3}\label{docs__pandoc__latex__src__3dmake_all_lessons.md__object-3}

Object description: \_\_\_\_\_\_\_\_\_\_\_\_\_\_\_\_

{\def\LTcaptype{none} % do not increment counter
\begin{longtable}[]{@{}
  >{\raggedright\arraybackslash}p{(\linewidth - 10\tabcolsep) * \real{0.1562}}
  >{\raggedright\arraybackslash}p{(\linewidth - 10\tabcolsep) * \real{0.3281}}
  >{\centering\arraybackslash}p{(\linewidth - 10\tabcolsep) * \real{0.1094}}
  >{\centering\arraybackslash}p{(\linewidth - 10\tabcolsep) * \real{0.1094}}
  >{\centering\arraybackslash}p{(\linewidth - 10\tabcolsep) * \real{0.1094}}
  >{\centering\arraybackslash}p{(\linewidth - 10\tabcolsep) * \real{0.1875}}@{}}
\toprule\noalign{}
\begin{minipage}[b]{\linewidth}\raggedright
Feature
\end{minipage} & \begin{minipage}[b]{\linewidth}\raggedright
What You\textquotesingle re Measuring
\end{minipage} & \begin{minipage}[b]{\linewidth}\centering
M1 (mm)
\end{minipage} & \begin{minipage}[b]{\linewidth}\centering
M2 (mm)
\end{minipage} & \begin{minipage}[b]{\linewidth}\centering
M3 (mm)
\end{minipage} & \begin{minipage}[b]{\linewidth}\centering
Average (mm)
\end{minipage} \\
\midrule\noalign{}
\endhead
\bottomrule\noalign{}
\endlastfoot
Length (X) & & & & & \\
Width (Y) & & & & & \\
Height (Z) & & & & & \\
Feature 4 & & & & & \\
Feature 5 & & & & & \\
\end{longtable}
}

\begin{center}\rule{0.5\linewidth}{0.5pt}\end{center}

\subsubsection*{Object 4}\label{docs__pandoc__latex__src__3dmake_all_lessons.md__object-4}

Object description: \_\_\_\_\_\_\_\_\_\_\_\_\_\_\_\_

{\def\LTcaptype{none} % do not increment counter
\begin{longtable}[]{@{}
  >{\raggedright\arraybackslash}p{(\linewidth - 10\tabcolsep) * \real{0.1562}}
  >{\raggedright\arraybackslash}p{(\linewidth - 10\tabcolsep) * \real{0.3281}}
  >{\centering\arraybackslash}p{(\linewidth - 10\tabcolsep) * \real{0.1094}}
  >{\centering\arraybackslash}p{(\linewidth - 10\tabcolsep) * \real{0.1094}}
  >{\centering\arraybackslash}p{(\linewidth - 10\tabcolsep) * \real{0.1094}}
  >{\centering\arraybackslash}p{(\linewidth - 10\tabcolsep) * \real{0.1875}}@{}}
\toprule\noalign{}
\begin{minipage}[b]{\linewidth}\raggedright
Feature
\end{minipage} & \begin{minipage}[b]{\linewidth}\raggedright
What You\textquotesingle re Measuring
\end{minipage} & \begin{minipage}[b]{\linewidth}\centering
M1 (mm)
\end{minipage} & \begin{minipage}[b]{\linewidth}\centering
M2 (mm)
\end{minipage} & \begin{minipage}[b]{\linewidth}\centering
M3 (mm)
\end{minipage} & \begin{minipage}[b]{\linewidth}\centering
Average (mm)
\end{minipage} \\
\midrule\noalign{}
\endhead
\bottomrule\noalign{}
\endlastfoot
Length (X) & & & & & \\
Width (Y) & & & & & \\
Height (Z) & & & & & \\
Feature 4 & & & & & \\
Feature 5 & & & & & \\
\end{longtable}
}

\begin{center}\rule{0.5\linewidth}{0.5pt}\end{center}

\subsubsection*{Object 5}\label{docs__pandoc__latex__src__3dmake_all_lessons.md__object-5}

Object description: \_\_\_\_\_\_\_\_\_\_\_\_\_\_\_\_

{\def\LTcaptype{none} % do not increment counter
\begin{longtable}[]{@{}
  >{\raggedright\arraybackslash}p{(\linewidth - 10\tabcolsep) * \real{0.1562}}
  >{\raggedright\arraybackslash}p{(\linewidth - 10\tabcolsep) * \real{0.3281}}
  >{\centering\arraybackslash}p{(\linewidth - 10\tabcolsep) * \real{0.1094}}
  >{\centering\arraybackslash}p{(\linewidth - 10\tabcolsep) * \real{0.1094}}
  >{\centering\arraybackslash}p{(\linewidth - 10\tabcolsep) * \real{0.1094}}
  >{\centering\arraybackslash}p{(\linewidth - 10\tabcolsep) * \real{0.1875}}@{}}
\toprule\noalign{}
\begin{minipage}[b]{\linewidth}\raggedright
Feature
\end{minipage} & \begin{minipage}[b]{\linewidth}\raggedright
What You\textquotesingle re Measuring
\end{minipage} & \begin{minipage}[b]{\linewidth}\centering
M1 (mm)
\end{minipage} & \begin{minipage}[b]{\linewidth}\centering
M2 (mm)
\end{minipage} & \begin{minipage}[b]{\linewidth}\centering
M3 (mm)
\end{minipage} & \begin{minipage}[b]{\linewidth}\centering
Average (mm)
\end{minipage} \\
\midrule\noalign{}
\endhead
\bottomrule\noalign{}
\endlastfoot
Length (X) & & & & & \\
Width (Y) & & & & & \\
Height (Z) & & & & & \\
Feature 4 & & & & & \\
Feature 5 & & & & & \\
\end{longtable}
}

\begin{center}\rule{0.5\linewidth}{0.5pt}\end{center}

\subsubsection*{Accuracy Check}\label{docs__pandoc__latex__src__3dmake_all_lessons.md__accuracy-check}

After completing all measurements, compare your averages with a partner who measured the same objects.

{\def\LTcaptype{none} % do not increment counter
\begin{longtable}[]{@{}
  >{\centering\arraybackslash}p{(\linewidth - 8\tabcolsep) * \real{0.0857}}
  >{\centering\arraybackslash}p{(\linewidth - 8\tabcolsep) * \real{0.2143}}
  >{\centering\arraybackslash}p{(\linewidth - 8\tabcolsep) * \real{0.3143}}
  >{\centering\arraybackslash}p{(\linewidth - 8\tabcolsep) * \real{0.2143}}
  >{\centering\arraybackslash}p{(\linewidth - 8\tabcolsep) * \real{0.1714}}@{}}
\toprule\noalign{}
\begin{minipage}[b]{\linewidth}\centering
Object
\end{minipage} & \begin{minipage}[b]{\linewidth}\centering
My Average (mm)
\end{minipage} & \begin{minipage}[b]{\linewidth}\centering
Partner\textquotesingle s Average (mm)
\end{minipage} & \begin{minipage}[b]{\linewidth}\centering
Difference (mm)
\end{minipage} & \begin{minipage}[b]{\linewidth}\centering
Within 1 mm?
\end{minipage} \\
\midrule\noalign{}
\endhead
\bottomrule\noalign{}
\endlastfoot
1 & & & & \\
2 & & & & \\
3 & & & & \\
4 & & & & \\
5 & & & & \\
\end{longtable}
}

If any difference is greater than 1 mm, remeasure together and find the source of the discrepancy.

\subsubsection*{Percent Error (Optional / Extension)}\label{docs__pandoc__latex__src__3dmake_all_lessons.md__percent-error-optional--extension}

If your instructor provides the "true" dimension of an object, calculate your percent error:

\texttt{\%\ error\ =\ (\textbar{}your\ average\ −\ true\ value\textbar{}\ /\ true\ value)\ ×\ 100}

{\def\LTcaptype{none} % do not increment counter
\begin{longtable}[]{@{}ccccc@{}}
\toprule\noalign{}
Object & Your Average & True Value & Difference & \% Error \\
\midrule\noalign{}
\endhead
\bottomrule\noalign{}
\endlastfoot
& & & & \\
& & & & \\
\end{longtable}
}

A percent error under 2\% is excellent for caliper work at this level.

\subsubsection*{Reflection}\label{docs__pandoc__latex__src__3dmake_all_lessons.md__reflection}

\emph{Answer in complete sentences.}

\begin{enumerate}
\item
  Which measurement was most difficult to take and why?
\item
  Did your three measurements for any feature vary significantly? What might cause that variation?
\item
  If you were designing an object in OpenSCAD that needed to fit \emph{over} one of these objects, which measurement would you use --- your smallest, your largest, or your average? Why?
\end{enumerate}

\begin{center}\rule{0.5\linewidth}{0.5pt}\end{center}

\subsection*{Appendix B: Measurement Calibration Guide}\label{docs__pandoc__latex__src__3dmake_all_lessons.md__appendix-b-measurement-calibration-guide}

\emph{Ensure accurate dimensions in your printed parts through proper calibration and verification.}

\subsubsection*{Why Calibration Matters}\label{docs__pandoc__latex__src__3dmake_all_lessons.md__why-calibration-matters}

Printed dimensions often deviate from designed dimensions due to nozzle width variations, inconsistent extrusion flow rate, plastic shrinkage during cooling, slicer interpretation differences, and material-specific shrinkage (ABS can shrink 0.3--1\%).

\begin{itemize}
\tightlist
\item
  Typical tolerances \textbf{without} calibration: ±0.3--0.5 mm
\item
  Typical tolerances \textbf{with} calibration: ±0.1--0.2 mm
\end{itemize}

\begin{center}\rule{0.5\linewidth}{0.5pt}\end{center}

\subsubsection*{Pre-Print Calibration}\label{docs__pandoc__latex__src__3dmake_all_lessons.md__pre-print-calibration}

\paragraph*{1. Nozzle Diameter Verification}\label{docs__pandoc__latex__src__3dmake_all_lessons.md__1-nozzle-diameter-verification}

Print a simple 10 mm × 10 mm × 5 mm solid box at default 0.4 mm line width with 100\% flow rate. Measure the result. If it deviates significantly from 10 mm, the nozzle may be partially clogged (clean it) or the wrong size (replace it).

\paragraph*{2. E-Steps Calibration (Extrusion Rate)}\label{docs__pandoc__latex__src__3dmake_all_lessons.md__2-e-steps-calibration-extrusion-rate}

E-steps control how many stepper-motor steps produce 1 mm of filament movement. If this is wrong, every print will be over- or under-extruded.

\textbf{Method:}

\begin{enumerate}
\tightlist
\item
  Heat extruder to printing temperature.
\item
  Mark filament 100 mm from the extruder entrance with a marker.
\item
  Command extrusion of exactly 100 mm (\texttt{G1\ E100\ F100}).
\item
  Measure the actual distance the mark moved.
\end{enumerate}

\textbf{Formula:}

\begin{lstlisting}[style=Alabaster]
New E-steps = Current E-steps × (100 mm / Actual distance moved)


Example:

  Current: 93 steps/mm

  Commanded: 100 mm

  Actual moved: 92 mm

  New: 93 × (100 / 92) = 101 steps/mm


\end{lstlisting}

Apply with \texttt{M92\ E101} then \texttt{M500} to save (Marlin), or update configuration file (Klipper).

\paragraph*{3. First Layer Height Verification}\label{docs__pandoc__latex__src__3dmake_all_lessons.md__3-first-layer-height-verification}

Print a single-layer square (base layer only). The first layer should be 0.2--0.25 mm thick --- roughly the thickness of a standard piece of printer paper. Adjust Z-offset by ±0.05 mm increments until correct.

\begin{center}\rule{0.5\linewidth}{0.5pt}\end{center}

\subsubsection*{XY Dimension Calibration}\label{docs__pandoc__latex__src__3dmake_all_lessons.md__xy-dimension-calibration}

\textbf{Test file (OpenSCAD):}

\begin{lstlisting}[style=Alabaster, language=openscad]
cube_size      = 20;

wall_thickness = 2;


difference() {

  cube([cube_size, cube_size, cube_size], center=true);

  cube([cube_size - 2*wall_thickness,

        cube_size - 2*wall_thickness,

        cube_size + 1], center=true);  // open top

}


\end{lstlisting}

Print at 100\% flow rate. Allow complete cooling (2+ hours). Measure the \emph{internal} dimension at 3 locations per axis. Expected internal width: \texttt{20\ -\ 2×2\ =\ 16\ mm}.

\textbf{Calibration formula:}

\begin{lstlisting}[style=Alabaster]
Flow rate adjustment = (Expected internal / Actual internal) × 100%


Example:

  Expected: 16.00 mm

  Actual: 15.75 mm

  Adjustment: (16.00 / 15.75) × 100% = 101.6%  →  set flow to 101.6% in slicer


\end{lstlisting}

\begin{center}\rule{0.5\linewidth}{0.5pt}\end{center}

\subsubsection*{Shrinkage Compensation by Material}\label{docs__pandoc__latex__src__3dmake_all_lessons.md__shrinkage-compensation-by-material}

{\def\LTcaptype{none} % do not increment counter
\begin{longtable}[]{@{}
  >{\raggedright\arraybackslash}p{(\linewidth - 4\tabcolsep) * \real{0.1127}}
  >{\centering\arraybackslash}p{(\linewidth - 4\tabcolsep) * \real{0.2394}}
  >{\raggedright\arraybackslash}p{(\linewidth - 4\tabcolsep) * \real{0.6479}}@{}}
\toprule\noalign{}
\begin{minipage}[b]{\linewidth}\raggedright
Material
\end{minipage} & \begin{minipage}[b]{\linewidth}\centering
Typical Shrinkage
\end{minipage} & \begin{minipage}[b]{\linewidth}\raggedright
Recommended Action
\end{minipage} \\
\midrule\noalign{}
\endhead
\bottomrule\noalign{}
\endlastfoot
PLA & 0.3--0.5\% & Usually acceptable; no action needed \\
PETG & 0.5--1\% & Scale design up 0.5--1\% for critical dimensions \\
ABS & 0.8--1.5\% & Scale up 1\% minimum \\
TPU & 1--2\% & Scale up 1--2\% for critical dimensions \\
\end{longtable}
}

Apply in OpenSCAD:

\begin{lstlisting}[style=Alabaster, language=openscad]
final_size        = 20;

material_shrink   = 1.01;  // 1% for ABS

designed_size     = final_size * material_shrink;


\end{lstlisting}

\begin{center}\rule{0.5\linewidth}{0.5pt}\end{center}

\subsubsection*{Critical Measurements to Track}\label{docs__pandoc__latex__src__3dmake_all_lessons.md__critical-measurements-to-track}

{\def\LTcaptype{none} % do not increment counter
\begin{longtable}[]{@{}
  >{\raggedright\arraybackslash}p{(\linewidth - 6\tabcolsep) * \real{0.2687}}
  >{\raggedright\arraybackslash}p{(\linewidth - 6\tabcolsep) * \real{0.2687}}
  >{\centering\arraybackslash}p{(\linewidth - 6\tabcolsep) * \real{0.2985}}
  >{\raggedright\arraybackslash}p{(\linewidth - 6\tabcolsep) * \real{0.1642}}@{}}
\toprule\noalign{}
\begin{minipage}[b]{\linewidth}\raggedright
Measurement
\end{minipage} & \begin{minipage}[b]{\linewidth}\raggedright
Method
\end{minipage} & \begin{minipage}[b]{\linewidth}\centering
Acceptable Tolerance
\end{minipage} & \begin{minipage}[b]{\linewidth}\raggedright
Frequency
\end{minipage} \\
\midrule\noalign{}
\endhead
\bottomrule\noalign{}
\endlastfoot
Wall thickness & Calipers (3 spots) & ±0.1 mm & Every print \\
Hole diameter & Calipers or gauge & ±0.1--0.2 mm & Every print \\
Overall dimensions & Calipers & ±0.2 mm & Monthly \\
Layer height & Stack on calipers & ±0.02 mm & Monthly \\
Vertical height & Measure side face & ±0.1 mm & Every print \\
\end{longtable}
}

\begin{center}\rule{0.5\linewidth}{0.5pt}\end{center}

\subsubsection*{Quick Calibration Checklist}\label{docs__pandoc__latex__src__3dmake_all_lessons.md__quick-calibration-checklist}

\textbf{Before first print with new settings:}

\begin{itemize}
\tightlist
\item[$\square$]
  E-steps calibration complete
\item[$\square$]
  First layer height verified
\item[$\square$]
  Nozzle diameter confirmed
\item[$\square$]
  Test print completed and measured
\end{itemize}

\textbf{Monthly maintenance:}

\begin{itemize}
\tightlist
\item[$\square$]
  Calibration cube printed and measured
\item[$\square$]
  Flow rate adjusted if needed
\item[$\square$]
  Layer height verified
\item[$\square$]
  Temperature consistency checked
\end{itemize}

\textbf{After any hardware change (nozzle swap, bed relevel, material change):}

\begin{itemize}
\tightlist
\item[$\square$]
  Re-verify nozzle diameter
\item[$\square$]
  Re-verify first layer height
\item[$\square$]
  Test print and measure before production prints
\end{itemize}

\begin{center}\rule{0.5\linewidth}{0.5pt}\end{center}

\subsubsection*{Recommended Measurement Tools}\label{docs__pandoc__latex__src__3dmake_all_lessons.md__recommended-measurement-tools}

{\def\LTcaptype{none} % do not increment counter
\begin{longtable}[]{@{}
  >{\raggedright\arraybackslash}p{(\linewidth - 6\tabcolsep) * \real{0.2388}}
  >{\centering\arraybackslash}p{(\linewidth - 6\tabcolsep) * \real{0.1791}}
  >{\centering\arraybackslash}p{(\linewidth - 6\tabcolsep) * \real{0.1194}}
  >{\raggedright\arraybackslash}p{(\linewidth - 6\tabcolsep) * \real{0.4627}}@{}}
\toprule\noalign{}
\begin{minipage}[b]{\linewidth}\raggedright
Tool
\end{minipage} & \begin{minipage}[b]{\linewidth}\centering
Typical Cost
\end{minipage} & \begin{minipage}[b]{\linewidth}\centering
Accuracy
\end{minipage} & \begin{minipage}[b]{\linewidth}\raggedright
Best Use
\end{minipage} \\
\midrule\noalign{}
\endhead
\bottomrule\noalign{}
\endlastfoot
Digital calipers & \$5--15 & ±0.05 mm & Primary tool for all work \\
Steel ruler & \$3--10 & ±1 mm & Quick rough checks \\
Vernier calipers & \$10--30 & ±0.05 mm & Precision work \\
Micrometer & \$20--50 & ±0.01 mm & Critical tolerance verification \\
\end{longtable}
}

\textbf{Recommendation:} Digital calipers are the most versatile and affordable starting point. Buy a pair before you print your first functional part.

\begin{center}\rule{0.5\linewidth}{0.5pt}\end{center}

\subsubsection*{Calibration Record}\label{docs__pandoc__latex__src__3dmake_all_lessons.md__calibration-record}

{\def\LTcaptype{none} % do not increment counter
\begin{longtable}[]{@{}ll@{}}
\toprule\noalign{}
Field & Value \\
\midrule\noalign{}
\endhead
\bottomrule\noalign{}
\endlastfoot
Last Calibration Date & \\
Printer Model & \\
Current E-Steps & \\
Current Flow Rate & \\
Materials Calibrated For & \\
\end{longtable}
}

\begin{center}\rule{0.5\linewidth}{0.5pt}\end{center}

\subsection*{Appendix C: Diagnostic Checklist}\label{docs__pandoc__latex__src__3dmake_all_lessons.md__appendix-c-diagnostic-checklist}

\emph{Use this checklist to systematically diagnose and troubleshoot printing issues before consulting the solutions guide.}

\subsubsection*{Quick Diagnosis Flowchart}\label{docs__pandoc__latex__src__3dmake_all_lessons.md__quick-diagnosis-flowchart}

\begin{lstlisting}[style=Alabaster]
Print Problem?

|

+---- [Before Print] Issues?

|  +---- Filament won't load  ->  CHECK: Temperature, drive gear, nozzle

|  +---- Printer won't heat   ->  CHECK: Power, temp sensor, firmware

|  +---- Bed not level        ->  CHECK: Leveling routine, bed surface

|

+---- [First Layer] Issues?

|  +---- Won't stick          ->  CHECK: Bed temp, cleanliness, nozzle height

|  +---- Too squished         ->  CHECK: Nozzle height, bed temp

|  +---- Gaps / uneven        ->  CHECK: Bed levelness, hot end alignment

|

+---- [Mid-Print] Issues?

|  +---- Stops extruding      ->  CHECK: Clog, temp drop, jam

|  +---- Layers shift         ->  CHECK: Loose belts, mechanical bind

|  +---- Print wobbles        ->  CHECK: Build plate, print stability

|

+---- [Print Quality] Issues?

   +---- Weak / brittle       ->  CHECK: Temperature, flow rate, layer adhesion

   +---- Rough surface        ->  CHECK: Flow rate, speed, temperature

   +---- Warped               ->  CHECK: Bed temperature, cooling, material


\end{lstlisting}

\begin{center}\rule{0.5\linewidth}{0.5pt}\end{center}

\subsubsection*{Pre-Print Diagnostics}\label{docs__pandoc__latex__src__3dmake_all_lessons.md__pre-print-diagnostics}

\paragraph*{A --- Power and Connectivity}\label{docs__pandoc__latex__src__3dmake_all_lessons.md__a--power-and-connectivity}

\begin{itemize}
\tightlist
\item[$\square$]
  Printer powered on; LED indicators normal
\item[$\square$]
  No error codes displayed; interface responding
\item[$\square$]
  USB connected (if applicable)
\end{itemize}

If failed: check power outlet, cable, power supply specs; power-cycle (off 30 sec, back on).

\paragraph*{B --- Temperature System}\label{docs__pandoc__latex__src__3dmake_all_lessons.md__b--temperature-system}

Heat hot end to 200°C; time to reach temperature should be 2--4 minutes with a steady rise (no plateau). Heat bed to 60°C; should reach target in 5--10 minutes.

\textbf{Stability test:} After reaching target, record temperature every minute for 10 minutes.

{\def\LTcaptype{none} % do not increment counter
\begin{longtable}[]{@{}cl@{}}
\toprule\noalign{}
Range & Assessment \\
\midrule\noalign{}
\endhead
\bottomrule\noalign{}
\endlastfoot
±2°C & Excellent \\
±5°C & Acceptable \\
±10°C & Marginal --- investigate thermal sensor \\
\textgreater±10°C & Problem --- do not print until resolved \\
\end{longtable}
}

\paragraph*{C --- Mechanical Systems}\label{docs__pandoc__latex__src__3dmake_all_lessons.md__c--mechanical-systems}

Disable motors (if firmware allows) and manually move each axis. Expected: smooth, no grinding.

{\def\LTcaptype{none} % do not increment counter
\begin{longtable}[]{@{}ll@{}}
\toprule\noalign{}
Axis & Result \\
\midrule\noalign{}
\endhead
\bottomrule\noalign{}
\endlastfoot
X-axis & smooth / rough / stuck \\
Y-axis & smooth / rough / stuck \\
Z-axis & smooth / rough / stuck \\
\end{longtable}
}

If rough or stuck: check for visible obstructions, inspect rails, verify pulleys turn freely, check belt tension, lubricate dry joints.

\paragraph*{D --- Build Plate Leveling (Paper Method)}\label{docs__pandoc__latex__src__3dmake_all_lessons.md__d--build-plate-leveling-paper-method}

\begin{enumerate}
\tightlist
\item
  Heat bed and nozzle to printing temperature.
\item
  Move nozzle to each corner; adjust leveling screw until a standard sheet of paper drags slightly under the nozzle.
\item
  Repeat for all 4 corners, then check center.
\item
  Target: consistent slight paper drag at all points; nozzle does not contact bed anywhere.
\end{enumerate}

\paragraph*{E --- Filament and Extruder}\label{docs__pandoc__latex__src__3dmake_all_lessons.md__e--filament-and-extruder}

\begin{itemize}
\tightlist
\item[$\square$]
  Filament diameter visually consistent over \textasciitilde{}50 cm
\item[$\square$]
  No visible cracks, kinks, or damage
\item[$\square$]
  Spool rotates freely without binding
\item[$\square$]
  Filament path clear of obstructions
\end{itemize}

\textbf{Manual extrusion test:} Heat to printing temperature. Push 10--20 mm of filament through manually. Expected: smooth push, consistent flow at nozzle tip.

\begin{center}\rule{0.5\linewidth}{0.5pt}\end{center}

\subsubsection*{First Layer Appearance Guide}\label{docs__pandoc__latex__src__3dmake_all_lessons.md__first-layer-appearance-guide}

{\def\LTcaptype{none} % do not increment counter
\begin{longtable}[]{@{}
  >{\raggedright\arraybackslash}p{(\linewidth - 4\tabcolsep) * \real{0.2697}}
  >{\raggedright\arraybackslash}p{(\linewidth - 4\tabcolsep) * \real{0.3371}}
  >{\raggedright\arraybackslash}p{(\linewidth - 4\tabcolsep) * \real{0.3933}}@{}}
\toprule\noalign{}
\begin{minipage}[b]{\linewidth}\raggedright
Appearance
\end{minipage} & \begin{minipage}[b]{\linewidth}\raggedright
Cause
\end{minipage} & \begin{minipage}[b]{\linewidth}\raggedright
Fix
\end{minipage} \\
\midrule\noalign{}
\endhead
\bottomrule\noalign{}
\endlastfoot
Wavy / embossed lines & Nozzle too close to bed & Raise Z-offset \\
Gaps between lines & Nozzle too far from bed & Lower Z-offset \\
Completely squished flat & Nozzle too low or bed too high &
Raise Z-offset or re-level \\
Partial adhesion only & Bed too cool or dirty &
Clean bed; increase bed temperature \\
Consistent, even lines & Correct & Continue print \\
\end{longtable}
}

\begin{center}\rule{0.5\linewidth}{0.5pt}\end{center}

\subsubsection*{Mid-Print Extrusion Failure}\label{docs__pandoc__latex__src__3dmake_all_lessons.md__mid-print-extrusion-failure}

When extrusion stops during a print:

\begin{enumerate}
\tightlist
\item
  \textbf{Pause} (do not cancel) immediately.
\item
  Listen for extruder grinding sounds (indicates jam).
\item
  Observe whether filament is still feeding into the extruder.
\end{enumerate}

\begin{lstlisting}[style=Alabaster]
Is filament stuck in the extruder?

+---- YES  -->  Nozzle clog likely

|               Try: cold pull, retract, then clean nozzle

|               See Appendix D (Common Issues) for full procedure

|

+---- NO   -->  Loading/path issue

               Is spool binding?     -->  Fix spool rotation

               Is path blocked?      -->  Clear obstruction

               Is drive gear slipping? -->  Clean and re-tension drive gear


\end{lstlisting}

\begin{center}\rule{0.5\linewidth}{0.5pt}\end{center}

\subsubsection*{Dimensional Accuracy Check (After Print Cools)}\label{docs__pandoc__latex__src__3dmake_all_lessons.md__dimensional-accuracy-check-after-print-cools}

Let the print cool for at least 2 hours before measuring.

\textbf{Measurement protocol:}

\begin{enumerate}
\tightlist
\item
  Measure each dimension 3 times at different locations.
\item
  Calculate average.
\item
  Compare to the design dimension.
\item
  Calculate deviation: \texttt{Deviation\ \%\ =\ ((Measured\ −\ Design)\ /\ Design)\ ×\ 100\%}
\end{enumerate}

{\def\LTcaptype{none} % do not increment counter
\begin{longtable}[]{@{}ccl@{}}
\toprule\noalign{}
Tolerance & Status & Action \\
\midrule\noalign{}
\endhead
\bottomrule\noalign{}
\endlastfoot
±0.5 mm or better & PASS & No adjustment needed \\
±0.5--1 mm & MARGINAL & Document and monitor \\
\textgreater±1 mm & FAIL & Adjust flow rate or recalibrate \\
\end{longtable}
}

\begin{center}\rule{0.5\linewidth}{0.5pt}\end{center}

\subsubsection*{Diagnostic Report Template}\label{docs__pandoc__latex__src__3dmake_all_lessons.md__diagnostic-report-template}

\emph{Fill this out before seeking help --- it ensures you\textquotesingle ve systematically checked the most common causes.}

\begin{lstlisting}[style=Alabaster]
DIAGNOSTIC REPORT

=================

Printer Model:

Problem Description:

When it occurs: (always / sometimes / first 5 layers / etc.)

Recent changes (nozzle swap, new filament, new settings):


DIAGNOSTICS PERFORMED:

[ ] Power / connectivity verified

[ ] Temperatures verified and stable

[ ] Mechanical movement tested on all axes

[ ] Bed leveling checked

[ ] Filament loading tested

[ ] First layer inspected and evaluated

[ ] Print quality evaluated mid-print


Key Findings:

1.

2.

3.


Attempted Solutions:

1.

2.


Result: (Solved / Partial / Ongoing)


\end{lstlisting}

\begin{center}\rule{0.5\linewidth}{0.5pt}\end{center}

\subsection*{Appendix D: Common Issues and Solutions}\label{docs__pandoc__latex__src__3dmake_all_lessons.md__appendix-d-common-issues-and-solutions}

\emph{Comprehensive troubleshooting guide for the most frequent 3D printing problems. Cross-reference with the Diagnostic Checklist (Appendix C) to identify which section applies to your situation.}

\begin{center}\rule{0.5\linewidth}{0.5pt}\end{center}

\subsubsection*{Pre-Print Issues}\label{docs__pandoc__latex__src__3dmake_all_lessons.md__pre-print-issues}

\paragraph*{Problem: Filament Won\textquotesingle t Load}\label{docs__pandoc__latex__src__3dmake_all_lessons.md__problem-filament-wont-load}

\textbf{Symptoms:} Extruder clicks or grinds but no filament moves; nothing extrudes when manual extrude is commanded.

\textbf{Diagnosis checklist:}

\begin{itemize}
\tightlist
\item[$\square$]
  Nozzle temperature high enough for the material?
\item[$\square$]
  Filament path clear of obstructions?
\item[$\square$]
  Drive gear still has grip (not worn smooth)?
\item[$\square$]
  Extruder tension set correctly?
\end{itemize}

\textbf{Solutions (try in order):}

\begin{enumerate}
\tightlist
\item
  Increase nozzle temperature by 10°C --- too cold makes extrusion harder.
\item
  Clean the extruder drive gear with a wire brush to restore grip.
\item
  Check filament diameter --- should be 1.75 mm ±0.03 mm. A spool that measured correctly when new may have inconsistent sections.
\item
  Verify extruder tension --- should grip firmly but not crush the filament.
\item
  Inspect nozzle for partial clog (see "Nozzle Clog" below).
\end{enumerate}

\paragraph*{Problem: Warped or Damaged Filament}\label{docs__pandoc__latex__src__3dmake_all_lessons.md__problem-warped-or-damaged-filament}

\textbf{Symptoms:} Filament appears kinked on the spool; inconsistent diameter; melted spots.

\textbf{Solutions:}

\begin{enumerate}
\tightlist
\item
  Cut off \textasciitilde{}30 cm of the damaged section and discard.
\item
  If the filament seems brittle or stringy: dry it in a low oven at 50°C for 2 hours (PLA). Moisture absorption causes significantly worse print quality.
\item
  Store filament sealed with desiccant between sessions --- especially important for Nylon, PETG, and TPU, which are hygroscopic.
\end{enumerate}

\begin{center}\rule{0.5\linewidth}{0.5pt}\end{center}

\subsubsection*{Print Quality Issues}\label{docs__pandoc__latex__src__3dmake_all_lessons.md__print-quality-issues}

\paragraph*{Problem: Poor Bed Adhesion (First Layer Lifts)}\label{docs__pandoc__latex__src__3dmake_all_lessons.md__problem-poor-bed-adhesion-first-layer-lifts}

\textbf{Symptoms:} Print lifts off bed during printing; edges curl up; print comes free entirely.

\textbf{Solutions (most effective first):}

\begin{enumerate}
\tightlist
\item
  Re-level the build plate --- cold plate leveling is more accurate.
\item
  Clean the build plate --- wipe with isopropyl alcohol to remove oils and dust.
\item
  Lower the nozzle --- if it\textquotesingle s too high, reduce Z-offset by 0.05--0.1 mm.
\item
  Increase bed temperature by 5--10°C.
\item
  Material-specific adhesion aids:

  \begin{itemize}
  \tightlist
  \item
    PLA: Painter\textquotesingle s tape, glue stick, or clean bare PEI/glass
  \item
    PETG: Textured PEI or thin glue layer (prevents it from bonding \emph{too} strongly to smooth surfaces)
  \item
    ABS: Heated enclosure + blue tape + ABS slurry (acetone + scraps)
  \item
    TPU: Clean bed, slightly elevated temperature
  \end{itemize}
\end{enumerate}

\paragraph*{Problem: Nozzle Clog}\label{docs__pandoc__latex__src__3dmake_all_lessons.md__problem-nozzle-clog}

\textbf{Symptoms:} No extrusion after initial layers; consistent gap between nozzle and surface; clicking from extruder; pressure building in hot end.

\textbf{Solutions (increasing intensity --- try each before escalating):}

\begin{enumerate}
\item
  \textbf{Cold Pull:} Heat nozzle to 200°C, grip filament, pull firmly while the nozzle cools. Repeat 5--10 times. The extracted filament should have a "plug" shape that brought debris with it.
\item
  \textbf{Poke from top:} Heat to printing temperature. Carefully insert a 0.3--0.4 mm drill bit or acupuncture needle through the top of the nozzle (from the heatbreak side). Do not force --- the goal is to break up, not compact, the clog.
\item
  \textbf{Soak and clean:} Remove the nozzle while hot (use wrench). Soak in acetone for 20--30 minutes. Use an ultrasonic cleaner if available. Clear the passage with a thin needle.
\item
  \textbf{Replace:} If the clog persists, the nozzle may be damaged internally. Replacement nozzles cost \$3--10 --- keep spares on hand.
\end{enumerate}

\paragraph*{Problem: Under-Extrusion (Thin Walls, Weak Prints)}\label{docs__pandoc__latex__src__3dmake_all_lessons.md__problem-under-extrusion-thin-walls-weak-prints}

\textbf{Symptoms:} Walls thinner than expected; layers separate easily; visible horizontal gaps; sections that appear lighter in color.

\textbf{Solutions:}

\begin{enumerate}
\tightlist
\item
  Check for partial nozzle clog (a fully extruding nozzle means a partial clog is possible).
\item
  Increase flow rate by 5--10\% in the slicer.
\item
  Raise temperature by 5--10°C.
\item
  Reduce print speed by 10--20\%.
\item
  Calibrate E-steps (see Appendix B).
\item
  Clean drive gear to remove accumulated plastic dust.
\end{enumerate}

\paragraph*{Problem: Over-Extrusion (Blobs, Rough Texture)}\label{docs__pandoc__latex__src__3dmake_all_lessons.md__problem-over-extrusion-blobs-rough-texture}

\textbf{Symptoms:} Excess material squishing between layers; rough, bumpy surface; blobs or zits on walls.

\textbf{Solutions:}

\begin{enumerate}
\tightlist
\item
  Reduce flow rate to 95\% in the slicer.
\item
  Lower temperature by 5--10°C.
\item
  Verify line width --- should be approximately 1.2× the nozzle diameter (0.4 mm nozzle → 0.48 mm line width).
\item
  Check slicer settings for nozzle size --- ensure it matches the actual nozzle installed.
\end{enumerate}

\paragraph*{Problem: Layer Shifting}\label{docs__pandoc__latex__src__3dmake_all_lessons.md__problem-layer-shifting}

\textbf{Symptoms:} Layers offset horizontally mid-print; top portion of print is misaligned; usually occurs at a specific layer height.

\textbf{Solutions:}

\begin{enumerate}
\tightlist
\item
  Check mechanics: manually move X/Y (motors disabled) --- should feel smooth, no resistance. Tighten all visible screws. Inspect belts for fraying.
\item
  Reduce travel speed by 20--30\%. Reduce acceleration in firmware.
\item
  Check print orientation --- tall, thin prints are vulnerable to the nozzle catching an edge and shifting. Rotate or add a brim.
\item
  If recently updated firmware: verify acceleration settings (try 1000 mm/s²) and step/mm calibration.
\end{enumerate}

\paragraph*{Problem: Stringing (Fine Threads Between Parts)}\label{docs__pandoc__latex__src__3dmake_all_lessons.md__problem-stringing-fine-threads-between-parts}

\textbf{Cause:} Nozzle oozes molten filament while traveling between print sections.

\textbf{Solutions:}

\begin{enumerate}
\tightlist
\item
  Reduce temperature by 5--10°C.
\item
  Increase retraction distance by 0.5--1.5 mm.
\item
  Increase retraction speed.
\item
  Enable "wipe on retract" in slicer.
\item
  Check for worn nozzle --- worn nozzles ooze more than new ones.
\end{enumerate}

\paragraph*{Problem: Warping (Corners Curl Up)}\label{docs__pandoc__latex__src__3dmake_all_lessons.md__problem-warping-corners-curl-up}

\textbf{Cause:} Material cools unevenly and contracts at different rates across the print.

\textbf{By material:}

\begin{itemize}
\tightlist
\item
  ABS: Use a heated enclosure, slow cooling fan, higher bed temperature (100--110°C).
\item
  PETG: Reduce cooling fan speed; increase bed temperature.
\item
  PLA: Usually warps if bed temperature is \emph{too high} --- reduce it.
\end{itemize}

\begin{center}\rule{0.5\linewidth}{0.5pt}\end{center}

\subsubsection*{Maintenance to Prevent Issues}\label{docs__pandoc__latex__src__3dmake_all_lessons.md__maintenance-to-prevent-issues}

{\def\LTcaptype{none} % do not increment counter
\begin{longtable}[]{@{}
  >{\raggedright\arraybackslash}p{(\linewidth - 4\tabcolsep) * \real{0.2128}}
  >{\raggedright\arraybackslash}p{(\linewidth - 4\tabcolsep) * \real{0.6064}}
  >{\centering\arraybackslash}p{(\linewidth - 4\tabcolsep) * \real{0.1809}}@{}}
\toprule\noalign{}
\begin{minipage}[b]{\linewidth}\raggedright
Issue
\end{minipage} & \begin{minipage}[b]{\linewidth}\raggedright
Prevention
\end{minipage} & \begin{minipage}[b]{\linewidth}\centering
Frequency
\end{minipage} \\
\midrule\noalign{}
\endhead
\bottomrule\noalign{}
\endlastfoot
Clogs & Clean nozzle; avoid heat-creep by maintaining cooling fan &
Before each print \\
Bed adhesion & Level bed; clean plate with IPA & Before each print \\
Layer shift & Check belt tension and drive gear screws & Monthly \\
Inconsistent quality & Calibrate E-steps; standardize environment &
Quarterly \\
Worn nozzle & Monitor extrusion quality; replace on schedule &
Every 6 months \\
Moisture in filament & Store sealed with desiccant & Ongoing \\
\end{longtable}
}

\begin{center}\rule{0.5\linewidth}{0.5pt}\end{center}

\subsection*{Appendix E: Accessibility Audit --- Extension Project}\label{docs__pandoc__latex__src__3dmake_all_lessons.md__appendix-e-accessibility-audit--extension-project}

\emph{Estimated time: 2--4 hours. This is an optional extension project for students who want to go deeper into inclusive design.}

\subsubsection*{Overview}\label{docs__pandoc__latex__src__3dmake_all_lessons.md__overview-7}

Accessibility means designing tools, documents, and workflows so that people with different abilities --- including visual impairment, motor differences, and cognitive differences --- can use them fully. In a fabrication context, accessibility applies to both the physical products you design and the documentation you produce about them.

This extension project asks you to evaluate one lesson or project for non-visual accessibility barriers, implement concrete improvements, and document what you found.

\subsubsection*{Learning Objectives}\label{docs__pandoc__latex__src__3dmake_all_lessons.md__learning-objectives-16}

By completing this project, you will:

\begin{itemize}
\tightlist
\item
  Evaluate digital content for accessibility barriers related to vision, blindness, and low vision
\item
  Apply concrete accessibility best practices (alt-text, code comments, plain language)
\item
  Document improvements and estimate implementation effort
\item
  Strengthen your awareness of inclusive design as a professional responsibility
\end{itemize}

\subsubsection*{Tasks}\label{docs__pandoc__latex__src__3dmake_all_lessons.md__tasks-6}

\begin{enumerate}
\tightlist
\item
  Select one lesson or project (an OpenSCAD file plus its markdown instructions) and review it for accessibility shortcomings.
\item
  Create a checklist covering: image alt-text, code walkthroughs, tactile cues, plain-language explanations, and file naming conventions. Apply it to the chosen resource.
\item
  Implement at least three small improvements (add alt text, improve code comments, write a text-only step-by-step walkthrough) and document every change.
\item
  Produce a one-page audit report summarizing what you found, what you fixed, and what larger improvements would require instructor time.
\end{enumerate}

\subsubsection*{Deliverable}\label{docs__pandoc__latex__src__3dmake_all_lessons.md__deliverable}

A one-page audit report plus a documented set of proposed changes (a "patch" or diff) for the target lesson.

\subsubsection*{Accessibility Principles to Apply}\label{docs__pandoc__latex__src__3dmake_all_lessons.md__accessibility-principles-to-apply}

\textbf{Alt-text for images:} Every image in a document should have a written description that conveys the same information. For a diagram showing a cross-section, the alt-text describes what the cross-section shows, not just "diagram."

\textbf{Code walkthroughs:} OpenSCAD files should include a plain-English comment block at the top describing what the file does and how to navigate it. Each module should have a comment explaining its purpose and parameters in plain language.

\textbf{Tactile cues in designs:} When designing objects that will be used physically, consider how they work for people who navigate by touch. Distinct ridges, chamfered corners, and different textures at key locations all improve usability for visually impaired users.

\textbf{Plain language:} Avoid jargon where a simpler word works. When technical terms are necessary, define them on first use.

\begin{center}\rule{0.5\linewidth}{0.5pt}\end{center}

\subsubsection*{Student Audit Documentation Template}\label{docs__pandoc__latex__src__3dmake_all_lessons.md__student-audit-documentation-template}

\textbf{Author:} \_\_\_\_\_\_\_\_\\
\textbf{Date:} \_\_\_\_\_\_\_\_\\
\textbf{Lesson/Project Audited:} \_\_\_\_\_\_\_\_

\paragraph*{Audit Scope}\label{docs__pandoc__latex__src__3dmake_all_lessons.md__audit-scope}

Resources reviewed:

\begin{enumerate}
\tightlist
\item
\item
\item
\end{enumerate}

Testing methodology used (screen reader, keyboard-only navigation, checklist review, other):

\begin{center}\rule{0.5\linewidth}{0.5pt}\end{center}

\paragraph*{Detailed Findings}\label{docs__pandoc__latex__src__3dmake_all_lessons.md__detailed-findings}

For each resource reviewed, complete the following:

\textbf{Resource 1: \_\_\_\_\_\_\_\_}

Screen reader assessment:

\begin{itemize}
\tightlist
\item
  Navigation: clear / unclear / broken
\item
  Output readability: readable / partially readable / unreadable
\item
  Error messages: helpful / unclear / missing
\item
  Specific barriers found:
\end{itemize}

Keyboard navigation:

\begin{itemize}
\tightlist
\item
  All functions accessible via keyboard: Yes / No
\item
  Tab order logical: Yes / No
\item
  Specific barriers:
\end{itemize}

{\def\LTcaptype{none} % do not increment counter
\begin{longtable}[]{@{}
  >{\raggedright\arraybackslash}p{(\linewidth - 6\tabcolsep) * \real{0.1972}}
  >{\centering\arraybackslash}p{(\linewidth - 6\tabcolsep) * \real{0.2254}}
  >{\centering\arraybackslash}p{(\linewidth - 6\tabcolsep) * \real{0.3803}}
  >{\centering\arraybackslash}p{(\linewidth - 6\tabcolsep) * \real{0.1972}}@{}}
\toprule\noalign{}
\begin{minipage}[b]{\linewidth}\raggedright
Recommendation
\end{minipage} & \begin{minipage}[b]{\linewidth}\centering
Priority (H/M/L)
\end{minipage} & \begin{minipage}[b]{\linewidth}\centering
Feasibility (Easy/Med/Hard)
\end{minipage} & \begin{minipage}[b]{\linewidth}\centering
Impact (H/M/L)
\end{minipage} \\
\midrule\noalign{}
\endhead
\bottomrule\noalign{}
\endlastfoot
& & & \\
& & & \\
\end{longtable}
}

\emph{(Repeat for each resource)}

\begin{center}\rule{0.5\linewidth}{0.5pt}\end{center}

\paragraph*{Summary of Barriers Found}\label{docs__pandoc__latex__src__3dmake_all_lessons.md__summary-of-barriers-found}

{\def\LTcaptype{none} % do not increment counter
\begin{longtable}[]{@{}llll@{}}
\toprule\noalign{}
Barrier & Frequency & Severity & Resources Affected \\
\midrule\noalign{}
\endhead
\bottomrule\noalign{}
\endlastfoot
& & & \\
& & & \\
\end{longtable}
}

\begin{center}\rule{0.5\linewidth}{0.5pt}\end{center}

\paragraph*{Recommendations Matrix}\label{docs__pandoc__latex__src__3dmake_all_lessons.md__recommendations-matrix}

{\def\LTcaptype{none} % do not increment counter
\begin{longtable}[]{@{}
  >{\centering\arraybackslash}p{(\linewidth - 8\tabcolsep) * \real{0.1356}}
  >{\centering\arraybackslash}p{(\linewidth - 8\tabcolsep) * \real{0.1864}}
  >{\centering\arraybackslash}p{(\linewidth - 8\tabcolsep) * \real{0.1017}}
  >{\raggedright\arraybackslash}p{(\linewidth - 8\tabcolsep) * \real{0.2373}}
  >{\raggedright\arraybackslash}p{(\linewidth - 8\tabcolsep) * \real{0.3390}}@{}}
\toprule\noalign{}
\begin{minipage}[b]{\linewidth}\centering
Priority
\end{minipage} & \begin{minipage}[b]{\linewidth}\centering
Feasibility
\end{minipage} & \begin{minipage}[b]{\linewidth}\centering
Impact
\end{minipage} & \begin{minipage}[b]{\linewidth}\raggedright
Recommendation
\end{minipage} & \begin{minipage}[b]{\linewidth}\raggedright
Implementation Steps
\end{minipage} \\
\midrule\noalign{}
\endhead
\bottomrule\noalign{}
\endlastfoot
High & Easy & High & & \\
High & Medium & High & & \\
Medium & Easy & Medium & & \\
\end{longtable}
}

\begin{center}\rule{0.5\linewidth}{0.5pt}\end{center}

\paragraph*{Reflection Questions}\label{docs__pandoc__latex__src__3dmake_all_lessons.md__reflection-questions}

Answer in complete sentences:

\begin{enumerate}
\tightlist
\item
  What surprised you about the accessibility barriers you found?
\item
  How would you prioritize improvements if you had limited time?
\item
  What role should users with disabilities play in accessibility testing, rather than non-disabled designers guessing at barriers?
\item
  What would it take to make accessibility review a standard part of the lesson development process?
\end{enumerate}

\begin{center}\rule{0.5\linewidth}{0.5pt}\end{center}

\paragraph*{Action Plan}\label{docs__pandoc__latex__src__3dmake_all_lessons.md__action-plan}

\begin{itemize}
\tightlist
\item
  Which improvements will you tackle first, and why?
\item
  How will you verify that your improvements actually help?
\item
  Timeline for implementation:
\end{itemize}

\begin{center}\rule{0.5\linewidth}{0.5pt}\end{center}

\subsubsection*{Assessment (if graded)}\label{docs__pandoc__latex__src__3dmake_all_lessons.md__assessment-if-graded}

{\def\LTcaptype{none} % do not increment counter
\begin{longtable}[]{@{}
  >{\raggedright\arraybackslash}p{(\linewidth - 4\tabcolsep) * \real{0.8161}}
  >{\centering\arraybackslash}p{(\linewidth - 4\tabcolsep) * \real{0.1264}}
  >{\raggedright\arraybackslash}p{(\linewidth - 4\tabcolsep) * \real{0.0575}}@{}}
\toprule\noalign{}
\begin{minipage}[b]{\linewidth}\raggedright
Category
\end{minipage} & \begin{minipage}[b]{\linewidth}\centering
Score (0--3)
\end{minipage} & \begin{minipage}[b]{\linewidth}\raggedright
Notes
\end{minipage} \\
\midrule\noalign{}
\endhead
\bottomrule\noalign{}
\endlastfoot
Problem \& Solution --- barriers identified, improvements actionable &
& \\
Design \& Code Quality --- testing methodology rigorous, findings specific
& & \\
Documentation --- template complete, recommendations prioritized & & \\
\textbf{Total (0--9)} & & \\
-e & & \\
\end{longtable}
}

\begin{center}\rule{0.5\linewidth}{0.5pt}\end{center}

\section*{Lesson 11 --- Stakeholder-Centric Design and the Final Project}\label{docs__pandoc__latex__src__3dmake_all_lessons.md__lesson-11--stakeholder-centric-design-and-the-final-project}

\textbf{Course:} 3dMake Certification --- Leadership Level
\textbf{Estimated Time:} 3--6 weeks (iterative; worked in phases)

\begin{center}\rule{0.5\linewidth}{0.5pt}\end{center}

\subsection*{A Message to You Before We Begin}\label{docs__pandoc__latex__src__3dmake_all_lessons.md__a-message-to-you-before-we-begin-1}

You have spent this course mastering the \emph{how} of digital fabrication: coding geometric primitives, managing tolerances, automating builds. This final lesson focuses on the \emph{why}. A professional designer does not simply build what they want --- they build what a stakeholder needs.

This lesson bridges the gap between being a maker and being a design professional. It introduces Design Thinking and Requirement Engineering, then puts those tools to work in a sustained final project you will carry through from a single conversation with a real person all the way to a finished, documented, physical artifact.

The goal is not perfection on the first try. The goal is to experience the full design cycle --- the same cycle that professional product designers, engineers, and makers go through when creating physical products. You will make decisions, encounter problems, revise your approach, and end up with something better than you could have imagined at the start.

\begin{center}\rule{0.5\linewidth}{0.5pt}\end{center}

\subsection*{Learning Objectives}\label{docs__pandoc__latex__src__3dmake_all_lessons.md__learning-objectives-17}

By the end of this lesson and project, you will be able to:

\begin{itemize}
\tightlist
\item
  Conduct a structured stakeholder interview using open-ended questioning strategies
\item
  Extract Functional Requirements and define measurable Acceptance Criteria from user data
\item
  Translate human needs into a Technical Design Specification
\item
  Build a fully parametric assembly in OpenSCAD that solves a specific user problem
\item
  Iterate a design based on systematic feedback collection
\item
  Demonstrate mastery of all prior curriculum concepts in a single integrated project
\end{itemize}

\begin{center}\rule{0.5\linewidth}{0.5pt}\end{center}

\subsection*{Part 1: Concepts --- Stakeholder-Centric Design}\label{docs__pandoc__latex__src__3dmake_all_lessons.md__part-1-concepts--stakeholder-centric-design}

\subsubsection*{What "Stakeholder-Centric" Means}\label{docs__pandoc__latex__src__3dmake_all_lessons.md__what-stakeholder-centric-means}

In engineering, a \textbf{stakeholder} is anyone affected by your design. This includes the person who asked for it, the person who will use it, the person who maintains it, and anyone who lives near it or depends on it. Professional designers do not start with shapes --- they start with \emph{people}.

The shift from "maker" to "design professional" is exactly this: a maker builds what they find interesting. A design professional builds what the stakeholder actually needs, which is often different.

This is harder than it sounds. People are not always able to articulate what they need. They describe symptoms, not root causes. They say "it needs to be bigger" when what they mean is "it\textquotesingle s too hard to grab with one hand." Your job is to ask the right questions, listen carefully, and translate what you hear into design constraints.

\subsubsection*{The Design Loop for Professionals}\label{docs__pandoc__latex__src__3dmake_all_lessons.md__the-design-loop-for-professionals}

Design Thinking is a structured method for solving problems that begins with understanding humans, not geometry. The five-stage loop is:

\begin{enumerate}
\tightlist
\item
  \textbf{Empathize} --- Interview the user to understand their "pain points." What do they struggle with? What workarounds do they currently use?
\item
  \textbf{Define} --- Create a list of requirements the part must satisfy. Make them specific and testable.
\item
  \textbf{Ideate} --- Use your OpenSCAD skills to create a parametric solution that addresses the requirements.
\item
  \textbf{Prototype \& Test} --- Build the part and have the stakeholder test it with their own hands.
\item
  \textbf{Iterate} --- Refine the parameters based on their feedback. Go back to step 1 if the problem changed.
\end{enumerate}

An important point: this is \textbf{not} a linear process. You will frequently loop back. You might build a prototype only to discover during testing that you defined the wrong problem. That is not a failure --- that is the process working correctly. The goal is to discover that early, while revisions are cheap.

\begin{center}\rule{0.5\linewidth}{0.5pt}\end{center}

\subsubsection*{Concept 2: The Stakeholder Interview}\label{docs__pandoc__latex__src__3dmake_all_lessons.md__concept-2-the-stakeholder-interview}

You cannot design for someone you don\textquotesingle t understand. A stakeholder interview is a structured conversation, not a quick survey. Done well, it reveals needs the user didn\textquotesingle t know how to express on their own.

\paragraph*{Strategy: Open-Ended Questions}\label{docs__pandoc__latex__src__3dmake_all_lessons.md__strategy-open-ended-questions}

The single most important skill in an interview is asking questions that cannot be answered with "yes" or "no."

Compare these two approaches:

{\def\LTcaptype{none} % do not increment counter
\begin{longtable}[]{@{}
  >{\raggedright\arraybackslash}p{(\linewidth - 2\tabcolsep) * \real{0.3434}}
  >{\raggedright\arraybackslash}p{(\linewidth - 2\tabcolsep) * \real{0.6566}}@{}}
\toprule\noalign{}
\begin{minipage}[b]{\linewidth}\raggedright
Closed Question
\end{minipage} & \begin{minipage}[b]{\linewidth}\raggedright
Open-Ended Alternative
\end{minipage} \\
\midrule\noalign{}
\endhead
\bottomrule\noalign{}
\endlastfoot
"Do you want it to be 10 cm wide?" &
"How would you use this in your daily space?" \\
"Is it too heavy?" &
"Tell me about a time you had trouble using something like this." \\
"Should it be red?" &
"What does the space where this will live look like?" \\
\end{longtable}
}

The open-ended version invites the user to tell you things you didn\textquotesingle t know to ask about. The closed version only confirms or denies what you already assumed.

\paragraph*{What to Listen For}\label{docs__pandoc__latex__src__3dmake_all_lessons.md__what-to-listen-for}

As the user talks, you\textquotesingle re listening for three categories of constraint:

\begin{itemize}
\tightlist
\item
  \textbf{Physical constraints:} "It needs to fit on my nightstand." "I only have 15 cm of clearance on the shelf."
\item
  \textbf{Accessibility constraints:} "I need to find the beads by touch alone." "My hands shake, so I need large controls."
\item
  \textbf{Durability/use constraints:} "I tend to drop things --- it can\textquotesingle t be fragile." "My kids will try to pull it apart."
\end{itemize}

Write down exact quotes. The user\textquotesingle s own words often contain the most precise requirements.

\begin{center}\rule{0.5\linewidth}{0.5pt}\end{center}

\subsubsection*{Concept 3: Engineering Functional Requirements}\label{docs__pandoc__latex__src__3dmake_all_lessons.md__concept-3-engineering-functional-requirements}

After the interview, you have a set of notes full of observations, stories, and opinions. Before writing any code, you must \textbf{translate} that material into Functional Requirements (FRs) --- specific, testable statements of what the design must do.

\paragraph*{The "User Story" Pattern}\label{docs__pandoc__latex__src__3dmake_all_lessons.md__the-user-story-pattern}

Use this format for every requirement:

\begin{quote}
"As a \textbf{{[}user type{]}}, I want to \textbf{{[}action{]}} so that \textbf{{[}benefit{]}}."
\end{quote}

This structure forces you to keep the user\textquotesingle s perspective central. Examples:

\begin{itemize}
\tightlist
\item
  \emph{"As a jewelry collector, I want the holder to have 2 mm clearance so that my bracelets slide on easily but don\textquotesingle t fall off."}
\item
  \emph{"As an elderly user, I want the clasp to require only one hand to operate so that I can put it on independently."}
\item
  \emph{"As a person with low vision, I want distinct tactile ridges on each section so that I can identify compartments without looking."}
\end{itemize}

\paragraph*{Acceptance Criteria}\label{docs__pandoc__latex__src__3dmake_all_lessons.md__acceptance-criteria}

Once you have a requirement, you need a \textbf{test} for it --- a specific condition that, when met, proves the requirement is satisfied. Without acceptance criteria, you can never know if you\textquotesingle re done.

{\def\LTcaptype{none} % do not increment counter
\begin{longtable}[]{@{}
  >{\raggedright\arraybackslash}p{(\linewidth - 2\tabcolsep) * \real{0.2713}}
  >{\raggedright\arraybackslash}p{(\linewidth - 2\tabcolsep) * \real{0.7287}}@{}}
\toprule\noalign{}
\begin{minipage}[b]{\linewidth}\raggedright
Functional Requirement
\end{minipage} & \begin{minipage}[b]{\linewidth}\raggedright
Acceptance Criterion
\end{minipage} \\
\midrule\noalign{}
\endhead
\bottomrule\noalign{}
\endlastfoot
Must support weight without tipping &
Holds 200 g with center of gravity at the rim; does not tip on a 5° incline \\
Bead hole accommodates 1.5 mm cord &
A 1.5 mm test cord passes through with finger pressure; does not fall out under its own weight \\
No sharp edges &
A flat palm pressed across all surfaces for 5 seconds produces no discomfort or marks \\
\end{longtable}
}

\begin{center}\rule{0.5\linewidth}{0.5pt}\end{center}

\subsubsection*{Concept 4: Technical Design Specification}\label{docs__pandoc__latex__src__3dmake_all_lessons.md__concept-4-technical-design-specification}

A Technical Design Specification is the bridge between requirements and code. It maps each FR to a named OpenSCAD parameter, defines the calculation behind each dimension, and records the reasoning.

This document is not bureaucracy. It is a record of \emph{decisions}. When you revisit a design three weeks later and wonder "why did I make the base this wide?", the specification has the answer.

\paragraph*{Code Pattern: Requirement Mapping}\label{docs__pandoc__latex__src__3dmake_all_lessons.md__code-pattern-requirement-mapping}

Variable names in your code should reflect the requirement they satisfy. Anyone reading the code should be able to trace a variable back to the FR that motivated it.

\begin{lstlisting}[style=Alabaster, language=openscad]
// FR1: Must fit standard bracelet diameters (typically 50–70 mm)

bracelet_max_d = 70;


// FR2: Must be stable without external support.

// Calculation: base width = 1.2 × total height (stability ratio)

base_w = 1.2 * total_height;


// FR3: Accessibility. Add tactile ridges for orientation by touch.

module tactile_marker() {

  sphere(r=2, $fn=16);

}


\end{lstlisting}

When a reviewer reads this code, they see not just \emph{what} you built, but \emph{why}. That\textquotesingle s the mark of professional-level documentation.

\begin{center}\rule{0.5\linewidth}{0.5pt}\end{center}

\subsubsection*{Concept 5: Feedback and Iteration}\label{docs__pandoc__latex__src__3dmake_all_lessons.md__concept-5-feedback-and-iteration}

After your first print, the work is not done --- it has just shifted. You move from designing to \emph{testing}, and the user\textquotesingle s hands are your most important instrument.

\paragraph*{Structured Feedback Collection}\label{docs__pandoc__latex__src__3dmake_all_lessons.md__structured-feedback-collection}

Don\textquotesingle t just ask "what do you think?" That produces vague answers. Instead, observe the user \emph{using} the object before asking any questions. Watch for:

\begin{itemize}
\tightlist
\item
  Hesitation (something isn\textquotesingle t intuitive)
\item
  Workarounds (they\textquotesingle re compensating for a flaw)
\item
  Surprise (positive or negative --- both are signals)
\end{itemize}

Then ask specific questions: "Was it easy to open with one hand?" "Did any part feel rough or uncomfortable?" "Is there anything you\textquotesingle d want to be different?"

\paragraph*{The Refinement Rule}\label{docs__pandoc__latex__src__3dmake_all_lessons.md__the-refinement-rule}

\textbf{Never guess at a fix based on vague feedback.} If the user says "it\textquotesingle s too tight," measure the actual fit with calipers before changing anything. Calculate the required clearance. Update your \texttt{TOLERANCE} constant. Print a test coupon before printing the full part.

This is the difference between engineering and guessing. Engineering produces predictable results. Guessing produces random results that occasionally work.

\begin{center}\rule{0.5\linewidth}{0.5pt}\end{center}

\subsection*{Part 2: The Project}\label{docs__pandoc__latex__src__3dmake_all_lessons.md__part-2-the-project}

You have two options for your final project. Both use the Design Thinking framework above. Your instructor will specify which option applies to your cohort.

\begin{center}\rule{0.5\linewidth}{0.5pt}\end{center}

\subsubsection*{Option A: The Beaded Jewelry Project}\label{docs__pandoc__latex__src__3dmake_all_lessons.md__option-a-the-beaded-jewelry-project}

\textbf{Unit:} 3 --- Open-Ended Projects
\textbf{Estimated Duration:} 1 week

\paragraph*{Project Brief}\label{docs__pandoc__latex__src__3dmake_all_lessons.md__project-brief}

Design and produce a wearable beaded jewelry piece that includes at least eight 3D-printed beads generated in OpenSCAD. Your design must use two distinct parametric bead modules and combine them into a completed, wearable piece (necklace, bracelet, or similar).

The jewelry project is an accessible entry point into stakeholder-centric design: the "stakeholder" is a real person (yourself, a friend, or a family member) with specific aesthetic preferences and physical requirements, and the "product" must be safe to wear, comfortable, and actually assembled.

\paragraph*{Constraints (must follow)}\label{docs__pandoc__latex__src__3dmake_all_lessons.md__constraints-must-follow}

\begin{itemize}
\tightlist
\item
  Your prototype must include a 3D-printed component designed in OpenSCAD.
\item
  Code your project in a single \texttt{.scad} file that parameterizes bead shapes for repeatability.
\item
  Use at least two different bead shapes in the final assembly.
\item
  The final piece must be wearable and assembled --- not just a set of loose beads.
\end{itemize}

\paragraph*{Learning Objectives}\label{docs__pandoc__latex__src__3dmake_all_lessons.md__learning-objectives-18}

\begin{itemize}
\tightlist
\item
  Create parametric OpenSCAD modules for repeated geometry
\item
  Combine modular parts into a coherent assembled object
\item
  Document design decisions and printing notes for reproducibility
\item
  Evaluate designs against measurable functional requirements
\end{itemize}

\paragraph*{Step-by-Step Milestones}\label{docs__pandoc__latex__src__3dmake_all_lessons.md__step-by-step-milestones}

\textbf{Day 1 --- Project Setup}

\begin{itemize}
\tightlist
\item
  Read this briefing and complete your Stakeholder Interview (use the template in Part 3 of this document).
\item
  Create a project folder. Initialize your Design Notes document.
\item
  Define at least 3 Functional Requirements using the FR template in Part 3.
\end{itemize}

\textbf{Days 1--2 --- Bead Module Development}

\begin{itemize}
\tightlist
\item
  Implement \texttt{bead\_A(size,\ detail)} and \texttt{bead\_B(size,\ detail)} parametric modules.
\item
  Your modules must accept at least two parameters: a size parameter and a detail/resolution parameter.
\item
  Test-print a single bead from each module. Record optimal print temperature and bed settings.
\end{itemize}

Here is a minimal starting scaffold --- replace the geometry with your own design:

\begin{lstlisting}[style=Alabaster, language=openscad]
// ============================================================

// Beaded Jewelry — Main Source File

// Author: [Your Name]

// Date: [YYYY-MM-DD]

// ============================================================


// ---- Key Parameters (change these to adapt the design) ----

bead_size    = 12;    // mm — controls overall bead scale

detail       = 32;    // $fn equivalent — increase for smoother curves

hole_d       = 2.8;   // mm — cord hole diameter (measure your cord + 0.3 mm clearance)

clearance    = 0.3;   // mm — added to hole for reliable cord passage


// ---- Bead A: Round/Sphere Variant ----

module bead_A(size=bead_size, fn=detail) {

  difference() {

    sphere(r=size/2, $fn=fn);

    // Cord hole through center

    translate([0, 0, -(size/2 + 0.001)])

      cylinder(r=(hole_d + clearance)/2, h=size + 0.002, $fn=16);

  }

}


// ---- Bead B: Bicone/Diamond Variant ----

module bead_B(size=bead_size, fn=detail) {

  difference() {

    // Bicone: hull of two cones meeting at midpoint

    hull() {

      translate([0, 0, size/4])  cylinder(r=size/2, h=0.1, $fn=fn, center=true);

      translate([0, 0, -size/2]) cylinder(r=0.1,    h=0.1, $fn=fn, center=true);

      translate([0, 0,  size/2]) cylinder(r=0.1,    h=0.1, $fn=fn, center=true);

    }

    // Cord hole

    translate([0, 0, -(size/2 + 0.001)])

      cylinder(r=(hole_d + clearance)/2, h=size + 0.002, $fn=16);

  }

}


// ---- Assembly: arrange 8+ beads in a line for spacing test ----

// Comment/uncomment individual modules to render one at a time for export


spacing = bead_size * 1.1;  // 10% gap between beads


for (i = [0:7]) {

  translate([i * spacing, 0, 0])

    if (i % 2 == 0) bead_A();

    else            bead_B();

}


\end{lstlisting}

\textbf{Days 3--4 --- Assembly and Iteration}

\begin{itemize}
\tightlist
\item
  Create an assembly that arranges at least 8 beads and tests fit and tolerances.
\item
  Print and string the beads on your chosen cord.
\item
  Iterate hole diameter and spacing until beads slide freely but do not fall off.
\item
  Record every iteration in your revision log.
\end{itemize}

\textbf{Day 5 --- Final Prototype and Documentation}

\begin{itemize}
\tightlist
\item
  Print the final bead set. Assemble the complete piece.
\item
  Photograph the result (multiple views, at least one worn).
\item
  Complete all deliverables listed below.
\end{itemize}

\paragraph*{Deliverables}\label{docs__pandoc__latex__src__3dmake_all_lessons.md__deliverables-3}

Submit both digitally and physically as instructed:

\begin{itemize}
\tightlist
\item
  \texttt{.scad} file containing parametric bead modules and the assembly script
\item
  \texttt{.stl} files (exported as needed) for printed parts
\item
  Technical documentation including:

  \begin{itemize}
  \tightlist
  \item
    Design Notes (ideas, measurements, parameter values, revision log)
  \item
    Construction / 3D printing notes (temperatures, speeds, supports, any failures)
  \item
    Completed Stakeholder Interview (Part 3 of this document)
  \item
    Completed Functional Requirements (Part 3)
  \item
    Photos of final prototype (multiple views)
  \end{itemize}
\item
  Physical turn-in: one assembled piece (if required by instructor)
\end{itemize}

\paragraph*{Functional Requirements (examples --- adapt to your design)}\label{docs__pandoc__latex__src__3dmake_all_lessons.md__functional-requirements-examples--adapt-to-your-design}

\begin{itemize}
\tightlist
\item
  The bead module must allow a hole diameter adjustable in 0.1 mm increments between 2.0--4.0 mm.
\item
  The assembly must include at least 8 beads and remain wearable (fits comfortably on intended body part).
\item
  The OpenSCAD file must be parameterized such that changing a single \texttt{scale} parameter adjusts bead size consistently.
\item
  The assembled piece must not have sharp edges that would injure skin under normal use.
\end{itemize}

\paragraph*{Accessibility Note}\label{docs__pandoc__latex__src__3dmake_all_lessons.md__accessibility-note}

Provide alt-text for all photos in your documentation. Include a short written walkthrough of how your \texttt{.scad} file generates the bead shapes so screen-reader users can understand the sequence and parameters.

\paragraph*{Grading Rubric --- Option A (0--9 scale)}\label{docs__pandoc__latex__src__3dmake_all_lessons.md__grading-rubric--option-a-09-scale}

See the Master Rubric in Part 3 of this document for the full scoring framework. For Option A:

{\def\LTcaptype{none} % do not increment counter
\begin{longtable}[]{@{}
  >{\raggedright\arraybackslash}p{(\linewidth - 4\tabcolsep) * \real{0.2857}}
  >{\centering\arraybackslash}p{(\linewidth - 4\tabcolsep) * \real{0.0893}}
  >{\raggedright\arraybackslash}p{(\linewidth - 4\tabcolsep) * \real{0.6250}}@{}}
\toprule\noalign{}
\begin{minipage}[b]{\linewidth}\raggedright
Category
\end{minipage} & \begin{minipage}[b]{\linewidth}\centering
Max Points
\end{minipage} & \begin{minipage}[b]{\linewidth}\raggedright
What Is Evaluated
\end{minipage} \\
\midrule\noalign{}
\endhead
\bottomrule\noalign{}
\endlastfoot
Implementation \& parametric code & 3 &
Modules are parametric, well-named, commented; assembly uses variables \\
Functionality \& wearability & 3 &
Piece is wearable; holes fit cord; no sharp edges; FRs met \\
Documentation \& print notes & 2 &
Design notes, revision log, FR document, print settings recorded \\
Presentation \& photos & 1 &
Multiple views; worn photo; alt-text provided \\
\end{longtable}
}

\paragraph*{Quiz --- Option A: Beaded Jewelry (15 questions)}\label{docs__pandoc__latex__src__3dmake_all_lessons.md__quiz--option-a-beaded-jewelry-15-questions}

\begin{enumerate}
\tightlist
\item
  Why use parametric modules for repeated parts rather than copying geometry? (one sentence)
\item
  What is a reliable way to test a bead hole size for a given cord diameter before printing the full set?
\item
  Name two OpenSCAD features or techniques useful for repeating geometry across a pattern.
\item
  What print setting is most likely to affect hole diameter accuracy in FDM printing?
\item
  How would you document an iteration that changed hole diameter from 2.5 mm to 2.7 mm?
\item
  True or False: Once your first bead prints successfully, all subsequent iterations will print exactly the same.
\item
  Explain what "wearability" means for a jewelry design. Name two specific checks you would perform to confirm a necklace is comfortable.
\item
  Your bead design uses a 3 mm hole. When strung on cord, it\textquotesingle s too tight to slide. What are two OpenSCAD-level solutions?
\item
  When documenting your design for reproducibility, should you record (A) final bead dimensions only, (B) all iterations including failed attempts, or (C) just the most recent version?
\item
  Why does the Stakeholder Interview happen \emph{before} you write any code?
\item
  Write a User Story for the following requirement: the cord hole must accommodate a 1.5 mm nylon cord.
\item
  What is an Acceptance Criterion, and how does it differ from a Functional Requirement?
\item
  You receive feedback that the beads "feel scratchy." Describe the engineering process for investigating and fixing this --- do not just say "sand them."
\item
  What is the \texttt{hull()} function useful for in bead design? Give one specific example.
\item
  Describe how the Design Loop (Empathize → Define → Ideate → Prototype → Iterate) applied to your specific jewelry project. Reference at least one concrete decision you made at each stage.
\end{enumerate}

\paragraph*{Extension Problems --- Option A (15)}\label{docs__pandoc__latex__src__3dmake_all_lessons.md__extension-problems--option-a-15}

\begin{enumerate}
\tightlist
\item
  Design an interlocking bead (snap-fit connection between beads) and document the clearances required.
\item
  Create a parametric clasp module that integrates with your bead string and includes a documented pass/fail test.
\item
  Modify one bead to include decorative text or a pattern generated procedurally in OpenSCAD.
\item
  Convert your single \texttt{.scad} project into a small library: separate bead modules into \texttt{use}-able include files and demonstrate reuse.
\item
  Propose a modification to make the piece weather-resistant for outdoor wear (materials, coatings, or geometry changes). Justify each choice.
\item
  Design and print a complete jewelry set: matching beads, clasp, and string connector with a consistent design language.
\item
  Create a variant generator that produces 10+ different bead designs from single-parameter changes. Document aesthetic and functional differences.
\item
  Build a design system document for your jewelry: modular bead library, material requirements, assembly instructions, and care guide.
\item
  Investigate material effects: print the same bead in two or more materials. Compare durability, aesthetics, and wearability with measurements.
\item
  Develop a parametric customization guide: enable another person to modify size, spacing, and aesthetics through top-level variables alone.
\item
  Conduct a tolerance sensitivity study: print 5 versions of a bead hole from 2.5 mm to 3.5 mm in 0.25 mm increments. Record which passes, which fails, and why.
\item
  Design a "tolerance coupon" --- a single flat print that tests 4 different hole diameters simultaneously --- so you can calibrate your printer for bead work in one print.
\item
  Use \texttt{surface()} or \texttt{minkowski()} to add a texture to your bead surface. Document the performance impact and the \texttt{\$preview} guard you used.
\item
  Write a batch build script that exports STLs for 8 bead variants (different sizes or shapes) automatically, with timestamped output names.
\item
  Conduct a formal stakeholder feedback session using the Feedback Collection Template (Part 3). Document the changes you made as a result.
\end{enumerate}

\begin{center}\rule{0.5\linewidth}{0.5pt}\end{center}

\subsubsection*{Option B: Open-Ended Final Project}\label{docs__pandoc__latex__src__3dmake_all_lessons.md__option-b-open-ended-final-project}

\textbf{Estimated Duration:} 3--6 weeks

\paragraph*{Project Brief}\label{docs__pandoc__latex__src__3dmake_all_lessons.md__project-brief-1}

Design, build, test, iterate, and document a real functional 3D-printed product of your choosing. This is not a tutorial --- there are no step-by-step instructions, no starter code, no single correct answer. You choose the problem; you design the solution; you prove it works.

\paragraph*{What Makes a Good Project?}\label{docs__pandoc__latex__src__3dmake_all_lessons.md__what-makes-a-good-project}

\textbf{Functional:} It does something real. Not a demonstration object --- an actual thing that solves a problem or serves a purpose.

\textbf{Multi-part:} At least two printed components that must fit together. This forces you to engage with tolerances, alignment, and interlocking features.

\textbf{Parametric:} Defined by variables, not hardcoded geometry. Someone else should be able to adapt your design to a different size or use case by changing a few numbers.

\textbf{Printable:} All components can be printed on a standard FDM printer without exotic supports or materials.

\textbf{Original:} Designed by you from scratch. You may be inspired by existing designs, but the code must be yours.

\paragraph*{Suggested Project Categories}\label{docs__pandoc__latex__src__3dmake_all_lessons.md__suggested-project-categories}

\textbf{Category A: Organizer / Storage} --- A system for organizing a specific collection of physical objects. Examples: parametric desktop organizer with interlocking cells; tool wall panel with snap-on holders; cable management clips sized for real cable diameters. \emph{Why it works:} real measurements to take, natural parametric axis, tolerance requirements throughout.

\textbf{Category B: Mounting / Enclosure} --- A mount, bracket, or enclosure for a specific real-world object. Examples: Raspberry Pi enclosure with lid and vent pattern; custom camera mount for a specific model; weatherproof sensor housing. \emph{Why it works:} must fit real hardware, heat-set insert practice, printability constraints.

\textbf{Category C: Ergonomic Tool or Accessory} --- Something designed to be held, worn, or operated. Examples: custom stylus grip, thumb ring for archery, ergonomic handle with grip texture. \emph{Why it works:} forces \texttt{hull()} for organic shapes, material selection for comfort, iterative refinement by feel.

\textbf{Category D: Mechanical Mechanism} --- A mechanism with moving parts. Examples: living hinge box, adjustable ratcheting stand, pen plotter stage. \emph{Why it works:} most demanding tolerances, deep engagement with clearance and snap-fit mechanics.

\textbf{Category E: Educational / Demonstrative Object} --- A physical object that teaches a concept. Examples: parametric gear train, manifold geometry demonstration set, printer calibration kit. \emph{Why it works:} strong parametric requirements, mathematical derivation in code, natural documentation need.

\begin{center}\rule{0.5\linewidth}{0.5pt}\end{center}

\subsection*{Project Phases (Option B)}\label{docs__pandoc__latex__src__3dmake_all_lessons.md__project-phases-option-b}

The project is organized into four phases. Work through them sequentially --- do not skip ahead.

\begin{center}\rule{0.5\linewidth}{0.5pt}\end{center}

\subsubsection*{Phase 1: Define and Plan (Week 1)}\label{docs__pandoc__latex__src__3dmake_all_lessons.md__phase-1-define-and-plan-week-1}

\textbf{Goal:} Understand the problem before writing any code.

The most common mistake in design projects is rushing to CAD. Phase 1 forces you to think before you model.

\paragraph*{Phase 1 Deliverables}\label{docs__pandoc__latex__src__3dmake_all_lessons.md__phase-1-deliverables}

\textbf{1.1 --- Stakeholder Interview}
Conduct a 45-minute interview using the Stakeholder Interview Template (Part 3). Your stakeholder is a real person who will use or benefit from your design.

\textbf{1.2 --- Problem Statement (200--400 words)}
Answer in prose: What specific problem does your design solve? Who will use it, and when? What must it do to be considered successful? What are the key constraints?

Do not describe the geometry yet. Describe the \emph{need}.

\textbf{1.3 --- Measurement Worksheet}
Before modeling, measure everything relevant in the real world with digital calipers.

{\def\LTcaptype{none} % do not increment counter
\begin{longtable}[]{@{}llccc@{}}
\toprule\noalign{}
Dimension & Description & Nominal (mm) & Measured (mm) & Tolerance \\
\midrule\noalign{}
\endhead
\bottomrule\noalign{}
\endlastfoot
... & ... & ... & ... & ±... \\
\end{longtable}
}

Measure at least 10 dimensions. For enclosures, measure the enclosed object completely. For mechanisms, measure all hardware that must interface with printed parts.

\textbf{1.4 --- Functional Requirements}
Use the FR Template (Part 3) to write at least 3 Functional Requirements, each with a User Story and at least one Acceptance Criterion.

\textbf{1.5 --- Design Specification}
Use the Design Specification Template (Part 3) to map your FRs to named OpenSCAD parameters.

\textbf{1.6 --- Sketch / Block Diagram}
Draw your design by hand (or any drawing tool). Label all key dimensions. Identify every interface between parts.

\textbf{1.7 --- Preliminary Bill of Materials}

{\def\LTcaptype{none} % do not increment counter
\begin{longtable}[]{@{}llccc@{}}
\toprule\noalign{}
Part & Material & Est. print time & Est. filament (g) & Qty \\
\midrule\noalign{}
\endhead
\bottomrule\noalign{}
\endlastfoot
\end{longtable}
}

Include both printed and non-printed components.

\textbf{1.8 --- Risk Assessment}
Identify three design challenges you anticipate. For each: what could go wrong, which tolerances are critical, and how will you address the risk.

\begin{center}\rule{0.5\linewidth}{0.5pt}\end{center}

\subsubsection*{Phase 2: Prototype and Test (Weeks 2--3)}\label{docs__pandoc__latex__src__3dmake_all_lessons.md__phase-2-prototype-and-test-weeks-23}

\textbf{Goal:} Build a working first version. Find the problems.

A prototype is a physical question: "Does this work?" Print fast, test aggressively, find the failure modes.

\paragraph*{Phase 2 Deliverables}\label{docs__pandoc__latex__src__3dmake_all_lessons.md__phase-2-deliverables}

\textbf{2.1 --- Parametric Source File}

\begin{lstlisting}[style=Alabaster, language=openscad]
// ============================================================

// [Project Name] — Main Source File

// Author: [Your Name]

// Date: [YYYY-MM-DD]

// Version: 0.1 (prototype)

// ============================================================


// ---- Key Parameters (change these to adapt the design) ----

// wall          = 2.0;   // mm — minimum print wall thickness

// clearance     = 0.2;   // mm — slip-fit clearance for interlocking parts

// corner_r      = 3.0;   // mm — exterior corner rounding radius


// ---- Derived Values (computed from parameters — do not edit directly) ----


// ---- Module Definitions ----


// ---- Assembly (uncomment the component to render/export) ----


\end{lstlisting}

Requirements for the source file:

\begin{itemize}
\tightlist
\item
  All key dimensions as named variables at the top
\item
  At least one module per major component
\item
  \texttt{assert()} guards on critical parameters
\item
  Comments explaining every non-obvious design decision
\end{itemize}

\textbf{2.2 --- First Print and Test Log}

{\def\LTcaptype{none} % do not increment counter
\begin{longtable}[]{@{}
  >{\raggedright\arraybackslash}p{(\linewidth - 10\tabcolsep) * \real{0.1429}}
  >{\raggedright\arraybackslash}p{(\linewidth - 10\tabcolsep) * \real{0.1111}}
  >{\centering\arraybackslash}p{(\linewidth - 10\tabcolsep) * \real{0.1587}}
  >{\centering\arraybackslash}p{(\linewidth - 10\tabcolsep) * \real{0.2540}}
  >{\centering\arraybackslash}p{(\linewidth - 10\tabcolsep) * \real{0.1429}}
  >{\raggedright\arraybackslash}p{(\linewidth - 10\tabcolsep) * \real{0.1905}}@{}}
\toprule\noalign{}
\begin{minipage}[b]{\linewidth}\raggedright
Component
\end{minipage} & \begin{minipage}[b]{\linewidth}\raggedright
Version
\end{minipage} & \begin{minipage}[b]{\linewidth}\centering
Print date
\end{minipage} & \begin{minipage}[b]{\linewidth}\centering
Key measurements
\end{minipage} & \begin{minipage}[b]{\linewidth}\centering
Pass/Fail
\end{minipage} & \begin{minipage}[b]{\linewidth}\raggedright
Issues found
\end{minipage} \\
\midrule\noalign{}
\endhead
\bottomrule\noalign{}
\endlastfoot
... & v0.1 & ... & ... & ... & ... \\
\end{longtable}
}

\textbf{2.3 --- Tolerance Test Results}
Print a tolerance test coupon for every critical interface before printing the full part. Document: nominal dimension, measured printed dimension, correction factor applied, whether the next print passed.

\textbf{2.4 --- Revision Log (first entries)}

\begin{lstlisting}[style=Alabaster]
v0.1 → v0.2 (YYYY-MM-DD)

  - Changed wall thickness from 1.5 to 2.0 mm: v0.1 cracked during snap-fit test

  - Increased lid clearance from 0.1 to 0.25 mm: lid was too tight to open one-handed

  - Added 2 mm chamfer to all interior edges: test part showed stress fracture at sharp corners


\end{lstlisting}

\begin{center}\rule{0.5\linewidth}{0.5pt}\end{center}

\subsubsection*{Phase 3: Refine and Iterate (Weeks 3--5)}\label{docs__pandoc__latex__src__3dmake_all_lessons.md__phase-3-refine-and-iterate-weeks-35}

\textbf{Goal:} Make it good. Systematic improvement based on test results.

Every decision in Phase 3 must be driven by test data, not intuition. "I think it will be stronger" is not enough --- print it, test it, measure it, decide.

\paragraph*{Required Iterations}\label{docs__pandoc__latex__src__3dmake_all_lessons.md__required-iterations}

You must complete at least three complete design-print-test-revise cycles. For each cycle, document:

\begin{enumerate}
\tightlist
\item
  What you changed and why (reference a specific failure mode)
\item
  What you printed (component name, version, parameters used)
\item
  How you tested it (specific test --- not just "I tried it")
\item
  What you measured (calipers, fit test, load test)
\item
  What you concluded (did the change fix the problem? Create new ones?)
\item
  What you\textquotesingle ll change next
\end{enumerate}

\paragraph*{Phase 3 Deliverables}\label{docs__pandoc__latex__src__3dmake_all_lessons.md__phase-3-deliverables}

\textbf{3.1 --- Iteration Cycles (3 minimum)} as described above.

\textbf{3.2 --- Feedback Collection}
Conduct a structured feedback session with your stakeholder using the Feedback Collection Template (Part 3). Document the changes you made as a result.

\textbf{3.3 --- Build Automation Script}
A script (bash or PowerShell) that builds all printable components, verifies each STL, logs results with timestamps, and archives all STLs with version tags.

\begin{lstlisting}[style=Alabaster, language=bash]
#!/bin/bash

# build_all.sh — builds all project components

set -euo pipefail


VERSION="v1.0"

OUTPUT_DIR="build/$VERSION"

mkdir -p "$OUTPUT_DIR"

echo "Building $VERSION..."


components=("lid" "base" "clip" "tray")  # replace with your component names


for comp in "${components[@]}"; do

  echo "  Building: $comp"

  openscad -D "component=\"${comp}\"" -o "$OUTPUT_DIR/${comp}.stl" src/main.scad

  SIZE=$(stat -c%s "$OUTPUT_DIR/${comp}.stl")

  echo "  OK: ${comp}.stl ($SIZE bytes)"

done


echo ""

echo "All components built: $OUTPUT_DIR"


\end{lstlisting}

\textbf{3.4 --- Advanced Feature Implementation}
Implement at least one advanced feature from Lesson 10:

\begin{itemize}
\tightlist
\item
  \texttt{hull()} or \texttt{minkowski()} for organic shaping
\item
  \texttt{surface()} for textured or embossed panels
\item
  \texttt{children()} container module
\item
  \texttt{search()} or \texttt{lookup()} table for material or size variants
\end{itemize}

Document which feature you chose, where it appears, and why.

\textbf{3.5 --- \texttt{projection()} Template (if applicable)}
If your design requires drilling or alignment of non-printed parts, produce a DXF template using \texttt{projection()}.

\begin{center}\rule{0.5\linewidth}{0.5pt}\end{center}

\subsubsection*{Phase 4: Document and Ship (Weeks 5--6)}\label{docs__pandoc__latex__src__3dmake_all_lessons.md__phase-4-document-and-ship-weeks-56}

\textbf{Goal:} Make the design reproducible by someone who wasn\textquotesingle t there when you built it.

A design that only works because you know all the tricks isn\textquotesingle t finished.

\paragraph*{Phase 4 Deliverables}\label{docs__pandoc__latex__src__3dmake_all_lessons.md__phase-4-deliverables}

\textbf{4.1 --- README.md}

\begin{lstlisting}[style=Alabaster]
# [Project Name]


[One-sentence description]


## What You Need


### Printed Parts


- [Part name] × [qty] — [material] — [est. print time]

  - Print settings: [layer height, infill, supports?]


### Hardware


- [Item] × [qty] — [spec / source]


### Tools


- Digital calipers

- [Other tools]


## Printing Instructions


[Orientation? Supports? Infill for structural parts? Temperature?]


## Assembly Instructions


[Step-by-step numbered. Note tolerance-sensitive steps.]


## Customization


[Key parameters and what they do. How to adapt for a different use case.]


## Known Limitations


[What doesn't work perfectly. What you'd improve with more time.]


## Version History


[Summary of major revisions]


## License


[Your chosen license]


\end{lstlisting}

\textbf{4.2 --- Final Bill of Materials}

{\def\LTcaptype{none} % do not increment counter
\begin{longtable}[]{@{}
  >{\raggedright\arraybackslash}p{(\linewidth - 12\tabcolsep) * \real{0.0606}}
  >{\raggedright\arraybackslash}p{(\linewidth - 12\tabcolsep) * \real{0.1970}}
  >{\raggedright\arraybackslash}p{(\linewidth - 12\tabcolsep) * \real{0.1212}}
  >{\raggedright\arraybackslash}p{(\linewidth - 12\tabcolsep) * \real{0.2121}}
  >{\centering\arraybackslash}p{(\linewidth - 12\tabcolsep) * \real{0.1364}}
  >{\centering\arraybackslash}p{(\linewidth - 12\tabcolsep) * \real{0.2273}}
  >{\centering\arraybackslash}p{(\linewidth - 12\tabcolsep) * \real{0.0455}}@{}}
\toprule\noalign{}
\begin{minipage}[b]{\linewidth}\raggedright
Part
\end{minipage} & \begin{minipage}[b]{\linewidth}\raggedright
File / Source
\end{minipage} & \begin{minipage}[b]{\linewidth}\raggedright
Material
\end{minipage} & \begin{minipage}[b]{\linewidth}\raggedright
Print Settings
\end{minipage} & \begin{minipage}[b]{\linewidth}\centering
Est. Time
\end{minipage} & \begin{minipage}[b]{\linewidth}\centering
Est. Weight (g)
\end{minipage} & \begin{minipage}[b]{\linewidth}\centering
Qty
\end{minipage} \\
\midrule\noalign{}
\endhead
\bottomrule\noalign{}
\endlastfoot
\end{longtable}
}

\textbf{4.3 --- Final STL Files}
Export one STL per printable component. Name them: \texttt{{[}project{]}-{[}component{]}-v{[}X.Y{]}.stl}

\textbf{4.4 --- Reflection Essay (400--600 words)}
Answer in essay form (not a list): What was the hardest design challenge, and how did you solve it? What would you do differently if you started over? What did you learn that surprised you? How has this project changed how you think about the physical objects you use every day?

\begin{center}\rule{0.5\linewidth}{0.5pt}\end{center}

\subsection*{Assessment Rubric (Option B)}\label{docs__pandoc__latex__src__3dmake_all_lessons.md__assessment-rubric-option-b}

Each criterion is scored 1--4.

{\def\LTcaptype{none} % do not increment counter
\begin{longtable}[]{@{}
  >{\raggedright\arraybackslash}p{(\linewidth - 8\tabcolsep) * \real{0.1082}}
  >{\raggedright\arraybackslash}p{(\linewidth - 8\tabcolsep) * \real{0.1861}}
  >{\raggedright\arraybackslash}p{(\linewidth - 8\tabcolsep) * \real{0.1645}}
  >{\raggedright\arraybackslash}p{(\linewidth - 8\tabcolsep) * \real{0.2121}}
  >{\raggedright\arraybackslash}p{(\linewidth - 8\tabcolsep) * \real{0.3290}}@{}}
\toprule\noalign{}
\begin{minipage}[b]{\linewidth}\raggedright
Criterion
\end{minipage} & \begin{minipage}[b]{\linewidth}\raggedright
1 --- Beginning
\end{minipage} & \begin{minipage}[b]{\linewidth}\raggedright
2 --- Developing
\end{minipage} & \begin{minipage}[b]{\linewidth}\raggedright
3 --- Proficient
\end{minipage} & \begin{minipage}[b]{\linewidth}\raggedright
4 --- Exemplary
\end{minipage} \\
\midrule\noalign{}
\endhead
\bottomrule\noalign{}
\endlastfoot
\textbf{Parametric Design} & Hardcoded dimensions; no variables &
Some variables; not all key dimensions &
All key dimensions parametric; named clearly &
Fully parametric; changing one variable produces a valid design at any scale \\
\textbf{CSG Modeling} & Primitives only; no CSG &
CSG used but with co-planar artifacts &
CSG correct; 0.001 offsets applied &
No non-manifold geometry; CSG used for all complex shapes \\
\textbf{Tolerance \& Fit} &
No tolerance consideration; parts don\textquotesingle t fit &
Tolerances guessed; some parts fit &
Tolerances measured and calibrated; all parts fit &
Test coupon printed; stack-up calculated; all fits confirmed \\
\textbf{Interlocking Features} & None & One feature, poorly fit &
One feature, well-designed and functional &
Two or more features; each type chosen appropriately \\
\textbf{Build Automation} & No build script &
Script builds but no error handling &
Script builds, verifies, and archives &
Handles all platforms; batch variants supported \\
\textbf{Documentation} & README missing or minimal &
README present but incomplete & README covers all required sections &
README enables a stranger to reproduce the design from scratch \\
\textbf{Iteration} & No documented revisions &
1--2 revisions, partially documented &
3 revisions, fully documented with rationale &
4+ revisions; each driven by measured test data \\
\textbf{Advanced Features} & None & One feature present &
One feature well-integrated &
Two or more; each justified in documentation \\
\textbf{Reflection} & Superficial or missing & Present but generic &
Genuine reflection on specific challenges &
Insightful; demonstrates growth and changed thinking \\
\end{longtable}
}

\textbf{Scoring:} 33--36 = Distinction / 27--32 = Merit / 18--26 = Pass / Below 18 = Needs revision

\begin{center}\rule{0.5\linewidth}{0.5pt}\end{center}

\subsubsection*{Quiz --- Option B: Final Project (15 questions)}\label{docs__pandoc__latex__src__3dmake_all_lessons.md__quiz--option-b-final-project-15-questions}

\begin{enumerate}
\tightlist
\item
  What is a stakeholder, and why is identifying stakeholders the first step in a design project?
\item
  What is the difference between a Maker and a Design Professional, as described in this lesson?
\item
  Why are open-ended interview questions preferred over yes/no questions?
\item
  Write a complete User Story for this requirement: "The lid must be openable with one hand."
\item
  What are Acceptance Criteria, and how do they make a Functional Requirement testable?
\item
  How does a Design Specification bridge requirements and OpenSCAD code? Give an example.
\item
  True or False: Design Thinking is a linear process that ends after the first prototype.
\item
  Why should you never adjust a design based on vague feedback like "it feels too tight" without first taking measurements?
\item
  How can \texttt{3dm\ info} involve a stakeholder who cannot see the 3D model visually?
\item
  Your stakeholder says the holder "feels wobbly." Describe the full engineering process from that comment to a revised design --- what do you measure, what do you change, how do you verify?
\item
  What is tolerance stack-up, and why does it become critical in Phase 1 rather than something to fix later?
\item
  Describe the purpose of printing a tolerance test coupon before printing a full assembly.
\item
  What should a revision log entry contain? Write an example entry for a wall thickness change.
\item
  Your build automation script reports that one STL is 340 bytes. What does this likely indicate, and what do you do next?
\item
  In your own words, explain why the Reflection Essay is a required deliverable --- what is it meant to capture that a build log cannot?
\end{enumerate}

\subsubsection*{Extension Problems --- Option B (15)}\label{docs__pandoc__latex__src__3dmake_all_lessons.md__extension-problems--option-b-15}

\begin{enumerate}
\tightlist
\item
  Conduct a full tolerance stack-up analysis for your assembly. Calculate worst-case misalignment across all interfaces.
\item
  Design a go/no-go gauge for your most critical fit --- a test part that passes only when the tolerance is correct.
\item
  Write a multi-platform build script that detects the OS and uses the appropriate shell commands automatically.
\item
  Implement semantic versioning in your build system: automatically increment a patch version on each successful build.
\item
  Build a Git pre-commit hook that runs your build script and refuses the commit if the build fails.
\item
  Create a "golden master" verification script: save the bounding box of a known-good build, then flag any future build that deviates by more than 0.1 mm.
\item
  Design a three-part interlocking system where all parts can be assembled in only one correct order. Document the design intent.
\item
  Write a GitHub Actions CI/CD workflow that builds all components on every push to the \texttt{main} branch and uploads STL artifacts.
\item
  Implement \texttt{lookup()} tables for material-specific parameters (density, clearance, recommended wall thickness). Make switching materials a one-variable change.
\item
  Produce a full DXF drilling template using \texttt{projection()} for any part of your design that requires holes in a non-printed substrate.
\item
  Build a "measurement worksheet" that renders in OpenSCAD as a flat printed sheet listing all key dimensions --- print it alongside your prototype.
\item
  Conduct a formal failure mode analysis: list every way your design could fail in real use, assess probability and impact, and describe your mitigation for each.
\item
  Create a parametric variant generator: a script that builds your design at 5 different scale factors and logs the bounding box and estimated weight for each.
\item
  Redesign your most complex module to use \texttt{children()} as a container pattern, then demonstrate it with three different child geometries.
\item
  Write a comprehensive retrospective document (600--1000 words) covering: what you planned vs. what you built, every significant decision point, what you would do differently, and what skills from earlier lessons proved most valuable.
\end{enumerate}

\begin{center}\rule{0.5\linewidth}{0.5pt}\end{center}

\subsection*{Milestone Check-Ins (Option B)}\label{docs__pandoc__latex__src__3dmake_all_lessons.md__milestone-check-ins-option-b}

{\def\LTcaptype{none} % do not increment counter
\begin{longtable}[]{@{}
  >{\raggedright\arraybackslash}p{(\linewidth - 4\tabcolsep) * \real{0.0949}}
  >{\raggedright\arraybackslash}p{(\linewidth - 4\tabcolsep) * \real{0.1387}}
  >{\raggedright\arraybackslash}p{(\linewidth - 4\tabcolsep) * \real{0.7664}}@{}}
\toprule\noalign{}
\begin{minipage}[b]{\linewidth}\raggedright
Week
\end{minipage} & \begin{minipage}[b]{\linewidth}\raggedright
Checkpoint
\end{minipage} & \begin{minipage}[b]{\linewidth}\raggedright
What to Show
\end{minipage} \\
\midrule\noalign{}
\endhead
\bottomrule\noalign{}
\endlastfoot
End of Week 1 & Phase 1 complete &
Interview notes, problem statement, measurement worksheet, FRs, design spec, sketch, BOM, risk assessment \\
End of Week 2 & First print &
At least one component printed and tested; Phase 2 test log started \\
End of Week 3 & Three revisions &
Phase 3 cycles 1--3 documented; feedback session complete; build script working \\
End of Week 4 & Feature complete &
All components printing and fitting; advanced feature implemented \\
End of Week 5 & Documentation draft &
README draft; STL files exported; reflection drafted \\
End of Week 6 & Final submission & All Phase 4 deliverables complete \\
\end{longtable}
}

\begin{center}\rule{0.5\linewidth}{0.5pt}\end{center}

\subsection*{Working Independently: Tips for Distance Learners}\label{docs__pandoc__latex__src__3dmake_all_lessons.md__working-independently-tips-for-distance-learners}

\textbf{Work in sessions, not marathons.} Two 90-minute sessions per week with clear goals produces better work than one 6-hour session. At the end of every session, write one sentence in your revision log about where you stopped and what you\textquotesingle ll do next.

\textbf{Print early, print often.} The instinct is to keep refining in CAD until the design is "perfect" before printing. Resist this. A bad print teaches you something a perfect render doesn\textquotesingle t. Print at 30\% infill and coarse layer height for test prints --- save fine settings for final parts.

\textbf{Describe your problems in writing before asking for help.} "My lid doesn\textquotesingle t close" is not a problem description. "My lid doesn\textquotesingle t close; the measured inner width of the base is 60.4 mm but the outer width of the lid is 60.2 mm, giving 0.2 mm per side clearance, but it still binds --- I suspect the corners because the lid snaps in at a tilt" is a problem description.

\textbf{Use your revision log as a thinking tool.} Don\textquotesingle t just record what changed. Record what you expected, what happened, and what you think that means.

\textbf{The design is done when it works, not when it\textquotesingle s perfect.}

\begin{center}\rule{0.5\linewidth}{0.5pt}\end{center}

\subsection*{Features Checklist (both options)}\label{docs__pandoc__latex__src__3dmake_all_lessons.md__features-checklist-both-options}

Use this checklist to verify you\textquotesingle ve applied required curriculum concepts:

\begin{itemize}
\tightlist
\item[$\square$]
  All key dimensions are named variables at the top of the source file
\item[$\square$]
  \texttt{assert()} guards protect critical parameters
\item[$\square$]
  At least two modules, each with a clear single responsibility
\item[$\square$]
  CSG operations use 0.001 mm offsets on co-planar faces
\item[$\square$]
  At least one interlocking feature (snap-fit, dovetail, thread, or friction fit)
\item[$\square$]
  Tolerance test coupon designed, printed, and measured
\item[$\square$]
  \texttt{clearance} constant used in at least one interface
\item[$\square$]
  Build automation script builds all components and verifies output
\item[$\square$]
  At least one advanced feature from Lesson 10
\item[$\square$]
  All STL files are manifold (no slicer warnings)
\item[$\square$]
  Stakeholder interview completed and documented
\item[$\square$]
  Functional requirements with acceptance criteria written
\item[$\square$]
  README complete (all required sections)
\item[$\square$]
  Bill of materials complete and accurate
\item[$\square$]
  At least three documented revision cycles
\item[$\square$]
  Feedback collection session completed
\item[$\square$]
  Reflection essay complete
\end{itemize}

\begin{center}\rule{0.5\linewidth}{0.5pt}\end{center}

\subsection*{Supplemental Resources}\label{docs__pandoc__latex__src__3dmake_all_lessons.md__supplemental-resources-10}

\begin{itemize}
\tightlist
\item
  Master Rubric: Part 3 of this document / \texttt{assets/3dMake\_Foundation/master-rubric.md}
\item
  Programming with OpenSCAD EPUB Textbook: \texttt{assets/Programming\_with\_OpenSCAD.epub}
\item
  OpenSCAD Quick Reference: https://programmingwithopenscad.github.io/quick-reference.html
\item
  3DMake GitHub Repository: https://github.com/tdeck/3dmake
\item
  CodeSolutions Repository: https://github.com/ProgrammingWithOpenSCAD/CodeSolutions
\end{itemize}

\begin{center}\rule{0.5\linewidth}{0.5pt}\end{center}

\subsection*{A Final Note}\label{docs__pandoc__latex__src__3dmake_all_lessons.md__a-final-note}

Every professional product you interact with went through a version of this process. The chair you\textquotesingle re sitting in was measured, prototyped, printed or machined, tested, revised, and documented before it reached you. The phone in your pocket has revision histories stretching back years, with engineers who made the exact kinds of decisions you\textquotesingle re making now --- just faster, with bigger teams, and with more sophisticated tools.

What you\textquotesingle re learning in this project is not just how to use OpenSCAD. You\textquotesingle re learning how to think through a physical design problem: how to start from a need, move through measurement and modeling, test against reality, iterate until it works, and communicate your process to others.

That thinking process transfers to any design tool, any material, and any scale. Take it seriously.

Good luck.

\begin{center}\rule{0.5\linewidth}{0.5pt}\end{center}

\section*{Part 3: Project Templates}\label{docs__pandoc__latex__src__3dmake_all_lessons.md__part-3-project-templates}

\emph{These templates are referenced throughout Lesson 11. Complete them as part of your project documentation. They are included here so the entire lesson is self-contained --- you do not need a separate file.}

\begin{center}\rule{0.5\linewidth}{0.5pt}\end{center}

\subsection*{Template 1: Stakeholder Interview}\label{docs__pandoc__latex__src__3dmake_all_lessons.md__template-1-stakeholder-interview}

\textbf{Use this for:} Phase 1 of Option B; Day 1 of Option A.
\textbf{Estimated duration:} 45--60 minutes.

\begin{center}\rule{0.5\linewidth}{0.5pt}\end{center}

\subsubsection*{Pre-Interview Preparation}\label{docs__pandoc__latex__src__3dmake_all_lessons.md__pre-interview-preparation}

{\def\LTcaptype{none} % do not increment counter
\begin{longtable}[]{@{}ll@{}}
\toprule\noalign{}
Field & Your Entry \\
\midrule\noalign{}
\endhead
\bottomrule\noalign{}
\endlastfoot
Date & \\
Stakeholder Name & \\
Their Role / Relationship to Project & \\
Interview Location & \\
Interview Objective (what do you need to learn?) & \\
\end{longtable}
}

\begin{center}\rule{0.5\linewidth}{0.5pt}\end{center}

\subsubsection*{Section 1: Context and Background (5 minutes)}\label{docs__pandoc__latex__src__3dmake_all_lessons.md__section-1-context-and-background-5-minutes}

\begin{enumerate}
\item
  "Tell me about your role and how this design would fit into your life."

  \begin{itemize}
  \tightlist
  \item
    \emph{Notes:}
  \end{itemize}
\item
  "What problem are we trying to solve together?"

  \begin{itemize}
  \tightlist
  \item
    \emph{Notes:}
  \end{itemize}
\item
  "Who else might be affected by this design?"

  \begin{itemize}
  \tightlist
  \item
    \emph{Notes:}
  \end{itemize}
\end{enumerate}

\begin{center}\rule{0.5\linewidth}{0.5pt}\end{center}

\subsubsection*{Section 2: Current Situation (10 minutes)}\label{docs__pandoc__latex__src__3dmake_all_lessons.md__section-2-current-situation-10-minutes}

\begin{enumerate}
\setcounter{enumi}{3}
\item
  "How do you currently handle this problem --- what do you do today?"

  \begin{itemize}
  \tightlist
  \item
    \emph{Notes:}
  \end{itemize}
\item
  "What works well about what you do now?"

  \begin{itemize}
  \tightlist
  \item
    \emph{Notes:}
  \end{itemize}
\item
  "What frustrates you most about the current situation?"

  \begin{itemize}
  \tightlist
  \item
    \emph{Notes (listen for exact words --- these often become requirements):}
  \end{itemize}
\end{enumerate}

\textbf{Pain Points Summary:}

\begin{itemize}
\tightlist
\item
  Primary:
\item
  Secondary:
\end{itemize}

\begin{center}\rule{0.5\linewidth}{0.5pt}\end{center}

\subsubsection*{Section 3: Ideal Solution (15 minutes)}\label{docs__pandoc__latex__src__3dmake_all_lessons.md__section-3-ideal-solution-15-minutes}

\emph{Use open-ended questions only. Do not suggest dimensions or shapes.}

\begin{enumerate}
\setcounter{enumi}{6}
\item
  "Ideally, what would a perfect solution do for you?"

  \begin{itemize}
  \tightlist
  \item
    \emph{Notes:}
  \end{itemize}
\item
  "Walk me through how you\textquotesingle d use it in a typical situation."

  \begin{itemize}
  \tightlist
  \item
    \emph{Notes:}
  \end{itemize}
\item
  "Where would it live --- what space does it need to fit into?"

  \begin{itemize}
  \tightlist
  \item
    \emph{Notes:}
  \end{itemize}
\end{enumerate}

\textbf{Constraints heard (fill in as they come up):}

{\def\LTcaptype{none} % do not increment counter
\begin{longtable}[]{@{}lll@{}}
\toprule\noalign{}
Type & What they said & Design implication \\
\midrule\noalign{}
\endhead
\bottomrule\noalign{}
\endlastfoot
Physical size & & \\
Weight & & \\
Temperature / environment & & \\
Durability & & \\
Budget & & \\
Timeline & & \\
\end{longtable}
}

\begin{center}\rule{0.5\linewidth}{0.5pt}\end{center}

\subsubsection*{Section 4: Usage Patterns (10 minutes)}\label{docs__pandoc__latex__src__3dmake_all_lessons.md__section-4-usage-patterns-10-minutes}

\begin{enumerate}
\setcounter{enumi}{9}
\item
  "How often would you use this?"

  \begin{itemize}
  \tightlist
  \item[$\square$]
    Daily {[} {]} Weekly {[} {]} Monthly {[} {]} Occasionally
  \item
    \emph{Notes:}
  \end{itemize}
\item
  "Who else might use it? Are there any accessibility considerations I should know about?"

  \begin{itemize}
  \tightlist
  \item
    \emph{Notes:}
  \end{itemize}
\end{enumerate}

\begin{center}\rule{0.5\linewidth}{0.5pt}\end{center}

\subsubsection*{Section 5: Success Criteria (10 minutes)}\label{docs__pandoc__latex__src__3dmake_all_lessons.md__section-5-success-criteria-10-minutes}

\begin{enumerate}
\setcounter{enumi}{11}
\item
  "If this design worked perfectly, how would you know? What would be different about your day?"

  \begin{itemize}
  \tightlist
  \item
    \emph{Notes:}
  \end{itemize}
\item
  "What would make you say this design failed?"

  \begin{itemize}
  \tightlist
  \item
    \emph{Notes:}
  \end{itemize}
\end{enumerate}

\textbf{Success Criteria (draft):}

\begin{enumerate}
\tightlist
\item
\item
\item
\end{enumerate}

\begin{center}\rule{0.5\linewidth}{0.5pt}\end{center}

\subsubsection*{Section 6: Priorities (10 minutes)}\label{docs__pandoc__latex__src__3dmake_all_lessons.md__section-6-priorities-10-minutes}

\begin{enumerate}
\setcounter{enumi}{13}
\tightlist
\item
  "If you could only have three features, which three would matter most to you?"

  \begin{enumerate}
  \tightlist
  \item
    \emph{(Priority: /10)}
  \item
    \emph{(Priority: /10)}
  \item
    \emph{(Priority: /10)}
  \end{enumerate}
\end{enumerate}

\begin{center}\rule{0.5\linewidth}{0.5pt}\end{center}

\subsubsection*{Post-Interview Analysis}\label{docs__pandoc__latex__src__3dmake_all_lessons.md__post-interview-analysis}

\textbf{Key Takeaways (your summary, not a transcript):}

\begin{enumerate}
\tightlist
\item
\item
\item
\end{enumerate}

\textbf{Surprising things you didn\textquotesingle t expect to hear:}

\textbf{Open questions to follow up on:}

\begin{enumerate}
\tightlist
\item
\item
\end{enumerate}

\textbf{How did this change your initial design idea?}

\emph{Conducted by:} \_\_\_\_\_\_\_\_ \emph{Date documented:} \_\_\_\_\_\_\_\_

\begin{center}\rule{0.5\linewidth}{0.5pt}\end{center}

\subsection*{Template 2: Functional Requirements}\label{docs__pandoc__latex__src__3dmake_all_lessons.md__template-2-functional-requirements}

\textbf{Use this for:} Phase 1 (Options A and B). Complete one FR block per requirement.

\begin{center}\rule{0.5\linewidth}{0.5pt}\end{center}

\subsubsection*{Project Information}\label{docs__pandoc__latex__src__3dmake_all_lessons.md__project-information}

{\def\LTcaptype{none} % do not increment counter
\begin{longtable}[]{@{}ll@{}}
\toprule\noalign{}
Field & Your Entry \\
\midrule\noalign{}
\endhead
\bottomrule\noalign{}
\endlastfoot
Project Name & \\
Version & \\
Date & \\
Prepared By & \\
\end{longtable}
}

\begin{center}\rule{0.5\linewidth}{0.5pt}\end{center}

\subsubsection*{FR1: {[}Name this requirement{]}}\label{docs__pandoc__latex__src__3dmake_all_lessons.md__fr1-name-this-requirement}

\textbf{Priority:} {[} {]} Critical {[} {]} High {[} {]} Medium {[} {]} Low

\textbf{Description} --- what the design must do:

\textbf{User Story:}

\begin{quote}
"As a \_\_\_\_\_\_, I want to \_\_\_\_\_\_ so that \_\_\_\_\_\_."
\end{quote}

\textbf{Acceptance Criteria} (each must be independently testable):

\begin{itemize}
\tightlist
\item[$\square$]
\item[$\square$]
\item[$\square$]
\end{itemize}

\textbf{Source} (which part of the interview led to this requirement?):

\begin{center}\rule{0.5\linewidth}{0.5pt}\end{center}

\subsubsection*{FR2: {[}Name this requirement{]}}\label{docs__pandoc__latex__src__3dmake_all_lessons.md__fr2-name-this-requirement}

\textbf{Priority:} {[} {]} Critical {[} {]} High {[} {]} Medium {[} {]} Low

\textbf{Description:}

\textbf{User Story:}

\begin{quote}
"As a \_\_\_\_\_\_, I want to \_\_\_\_\_\_ so that \_\_\_\_\_\_."
\end{quote}

\textbf{Acceptance Criteria:}

\begin{itemize}
\tightlist
\item[$\square$]
\item[$\square$]
\end{itemize}

\textbf{Source:}

\begin{center}\rule{0.5\linewidth}{0.5pt}\end{center}

\subsubsection*{FR3: {[}Name this requirement{]}}\label{docs__pandoc__latex__src__3dmake_all_lessons.md__fr3-name-this-requirement}

\textbf{Priority:} {[} {]} Critical {[} {]} High {[} {]} Medium {[} {]} Low

\textbf{Description:}

\textbf{User Story:}

\begin{quote}
"As a \_\_\_\_\_\_, I want to \_\_\_\_\_\_ so that \_\_\_\_\_\_."
\end{quote}

\textbf{Acceptance Criteria:}

\begin{itemize}
\tightlist
\item[$\square$]
\item[$\square$]
\end{itemize}

\textbf{Source:}

\begin{center}\rule{0.5\linewidth}{0.5pt}\end{center}

\emph{Add FR4, FR5, etc. by copying the block above.}

\begin{center}\rule{0.5\linewidth}{0.5pt}\end{center}

\subsubsection*{Scope}\label{docs__pandoc__latex__src__3dmake_all_lessons.md__scope}

\subsection*{\texorpdfstring{\textbf{In Scope (this design will address):}}{In Scope (this design will address):}}\label{docs__pandoc__latex__src__3dmake_all_lessons.md__in-scope-this-design-will-address}

\begin{itemize}
\item
\item
  \textbf{Out of Scope (explicitly excluded):}
\item
\item
  \begin{center}\rule{0.5\linewidth}{0.5pt}\end{center}
\end{itemize}

\subsubsection*{Verification Plan}\label{docs__pandoc__latex__src__3dmake_all_lessons.md__verification-plan}

{\def\LTcaptype{none} % do not increment counter
\begin{longtable}[]{@{}llll@{}}
\toprule\noalign{}
Requirement & Test Method & Pass Condition & Status \\
\midrule\noalign{}
\endhead
\bottomrule\noalign{}
\endlastfoot
FR1 & & & \\
FR2 & & & \\
FR3 & & & \\
\end{longtable}
}

\begin{center}\rule{0.5\linewidth}{0.5pt}\end{center}

\subsection*{Template 3: Design Specification}\label{docs__pandoc__latex__src__3dmake_all_lessons.md__template-3-design-specification}

\textbf{Use this for:} Mapping FRs to OpenSCAD parameters before writing code.

\begin{center}\rule{0.5\linewidth}{0.5pt}\end{center}

\subsubsection*{Project Information}\label{docs__pandoc__latex__src__3dmake_all_lessons.md__project-information-1}

{\def\LTcaptype{none} % do not increment counter
\begin{longtable}[]{@{}ll@{}}
\toprule\noalign{}
Field & Your Entry \\
\midrule\noalign{}
\endhead
\bottomrule\noalign{}
\endlastfoot
Project & \\
Version & \\
Date & \\
Designed By & \\
\end{longtable}
}

\begin{center}\rule{0.5\linewidth}{0.5pt}\end{center}

\subsubsection*{Design Overview}\label{docs__pandoc__latex__src__3dmake_all_lessons.md__design-overview}

\textbf{Problem Statement} (one paragraph):

\textbf{Solution Approach} (high-level):

\textbf{Key Design Decisions:}

\begin{enumerate}
\tightlist
\item
\item
\item
\end{enumerate}

\begin{center}\rule{0.5\linewidth}{0.5pt}\end{center}

\subsubsection*{Requirement → Parameter Mapping}\label{docs__pandoc__latex__src__3dmake_all_lessons.md__requirement--parameter-mapping}

For each Functional Requirement, identify the OpenSCAD variable(s) that implement it.

{\def\LTcaptype{none} % do not increment counter
\begin{longtable}[]{@{}llll@{}}
\toprule\noalign{}
FR & OpenSCAD Variable & Value / Formula & Rationale \\
\midrule\noalign{}
\endhead
\bottomrule\noalign{}
\endlastfoot
FR1 & & & \\
FR2 & & & \\
FR3 & & & \\
\end{longtable}
}

\textbf{Code pattern --- paste your parameter block here after writing it:}

\begin{lstlisting}[style=Alabaster, language=openscad]
// ---- Parameters derived from requirements ----


// FR1: [Requirement name]

// [variable] = [value];  // mm — [explanation]


// FR2: [Requirement name]

// [variable] = [value];  // [explanation]


\end{lstlisting}

\begin{center}\rule{0.5\linewidth}{0.5pt}\end{center}

\subsubsection*{Component Breakdown}\label{docs__pandoc__latex__src__3dmake_all_lessons.md__component-breakdown}

For each printed component:

\textbf{Component 1: {[}Name{]}}

\begin{itemize}
\tightlist
\item
  Purpose:
\item
  Dimensions (L × W × H):
\item
  Material:
\item
  Interfaces with: \emph{(which other parts?)}
\item
  Print orientation:
\item
  Supports needed: {[} {]} Yes {[} {]} No
\item
  Estimated print time:
\end{itemize}

\textbf{Component 2: {[}Name{]}} \emph{(copy block above)}

\begin{center}\rule{0.5\linewidth}{0.5pt}\end{center}

\subsubsection*{Tolerance and Fit Analysis}\label{docs__pandoc__latex__src__3dmake_all_lessons.md__tolerance-and-fit-analysis}

{\def\LTcaptype{none} % do not increment counter
\begin{longtable}[]{@{}
  >{\raggedright\arraybackslash}p{(\linewidth - 8\tabcolsep) * \real{0.2143}}
  >{\raggedright\arraybackslash}p{(\linewidth - 8\tabcolsep) * \real{0.2857}}
  >{\centering\arraybackslash}p{(\linewidth - 8\tabcolsep) * \real{0.1714}}
  >{\centering\arraybackslash}p{(\linewidth - 8\tabcolsep) * \real{0.2000}}
  >{\raggedright\arraybackslash}p{(\linewidth - 8\tabcolsep) * \real{0.1286}}@{}}
\toprule\noalign{}
\begin{minipage}[b]{\linewidth}\raggedright
Interface
\end{minipage} & \begin{minipage}[b]{\linewidth}\raggedright
Type of fit
\end{minipage} & \begin{minipage}[b]{\linewidth}\centering
Nominal dim.
\end{minipage} & \begin{minipage}[b]{\linewidth}\centering
Clearance used
\end{minipage} & \begin{minipage}[b]{\linewidth}\raggedright
Rationale
\end{minipage} \\
\midrule\noalign{}
\endhead
\bottomrule\noalign{}
\endlastfoot
Part A → Part B & Press / Slip / Loose & & & \\
\end{longtable}
}

\begin{center}\rule{0.5\linewidth}{0.5pt}\end{center}

\subsubsection*{Material Selection}\label{docs__pandoc__latex__src__3dmake_all_lessons.md__material-selection}

\textbf{Primary material:}

{\def\LTcaptype{none} % do not increment counter
\begin{longtable}[]{@{}
  >{\raggedright\arraybackslash}p{(\linewidth - 4\tabcolsep) * \real{0.3492}}
  >{\raggedright\arraybackslash}p{(\linewidth - 4\tabcolsep) * \real{0.1746}}
  >{\raggedright\arraybackslash}p{(\linewidth - 4\tabcolsep) * \real{0.4762}}@{}}
\toprule\noalign{}
\begin{minipage}[b]{\linewidth}\raggedright
Property
\end{minipage} & \begin{minipage}[b]{\linewidth}\raggedright
Requirement
\end{minipage} & \begin{minipage}[b]{\linewidth}\raggedright
Why this material satisfies it
\end{minipage} \\
\midrule\noalign{}
\endhead
\bottomrule\noalign{}
\endlastfoot
Strength & & \\
Flexibility & & \\
Temperature resistance & & \\
Cost & & \\
\end{longtable}
}

\begin{center}\rule{0.5\linewidth}{0.5pt}\end{center}

\subsubsection*{Risk Register}\label{docs__pandoc__latex__src__3dmake_all_lessons.md__risk-register}

{\def\LTcaptype{none} % do not increment counter
\begin{longtable}[]{@{}lccl@{}}
\toprule\noalign{}
Risk & Probability & Impact & Mitigation \\
\midrule\noalign{}
\endhead
\bottomrule\noalign{}
\endlastfoot
& H/M/L & H/M/L & \\
& H/M/L & H/M/L & \\
\end{longtable}
}

\begin{center}\rule{0.5\linewidth}{0.5pt}\end{center}

\subsubsection*{Design Change Log}\label{docs__pandoc__latex__src__3dmake_all_lessons.md__design-change-log}

{\def\LTcaptype{none} % do not increment counter
\begin{longtable}[]{@{}llll@{}}
\toprule\noalign{}
Version & Date & Change & Reason \\
\midrule\noalign{}
\endhead
\bottomrule\noalign{}
\endlastfoot
0.1 & & Initial design & \\
\end{longtable}
}

\begin{center}\rule{0.5\linewidth}{0.5pt}\end{center}

\subsection*{Template 4: Feedback Collection}\label{docs__pandoc__latex__src__3dmake_all_lessons.md__template-4-feedback-collection}

\textbf{Use this for:} Phase 3 --- structured feedback session with your stakeholder.

\begin{center}\rule{0.5\linewidth}{0.5pt}\end{center}

\subsubsection*{Session Information}\label{docs__pandoc__latex__src__3dmake_all_lessons.md__session-information}

{\def\LTcaptype{none} % do not increment counter
\begin{longtable}[]{@{}ll@{}}
\toprule\noalign{}
Field & Your Entry \\
\midrule\noalign{}
\endhead
\bottomrule\noalign{}
\endlastfoot
Date & \\
Stakeholder Name & \\
Prototype Version & \\
Session Duration & \\
\end{longtable}
}

\begin{center}\rule{0.5\linewidth}{0.5pt}\end{center}

\subsubsection*{Section 1: First Impressions (5 minutes)}\label{docs__pandoc__latex__src__3dmake_all_lessons.md__section-1-first-impressions-5-minutes}

\emph{Hand the prototype to the stakeholder without explanation. Observe for 60 seconds before asking anything.}

\textbf{Observations} (what did they do first?):

\begin{enumerate}
\item
  "What is your first reaction to this?"

  \begin{itemize}
  \tightlist
  \item
    \emph{Response:}
  \end{itemize}
\item
  "Is it what you expected?"

  \begin{itemize}
  \tightlist
  \item[$\square$]
    More than expected {[} {]} As expected {[} {]} Less than expected
  \item
    \emph{Why:}
  \end{itemize}
\item
  "Would you use this?"

  \begin{itemize}
  \tightlist
  \item[$\square$]
    Definitely {[} {]} Probably {[} {]} Maybe {[} {]} No
  \item
    \emph{Why:}
  \end{itemize}
\end{enumerate}

\begin{center}\rule{0.5\linewidth}{0.5pt}\end{center}

\subsubsection*{Section 2: Usability (15 minutes)}\label{docs__pandoc__latex__src__3dmake_all_lessons.md__section-2-usability-15-minutes}

\emph{Ask them to demonstrate using it as they normally would. Do not explain how it works --- observe.}

\textbf{Observations} (hesitations, workarounds, confusion):

\begin{enumerate}
\setcounter{enumi}{3}
\item
  "Was it easy to figure out how to use?"

  \begin{itemize}
  \tightlist
  \item[$\square$]
    Very easy {[} {]} Easy {[} {]} Neutral {[} {]} Difficult {[} {]} Very difficult
  \item
    \emph{What was confusing:}
  \end{itemize}
\item
  "Are there any parts that don\textquotesingle t work well together?"

  \begin{itemize}
  \tightlist
  \item
    \emph{Response:}
  \end{itemize}
\end{enumerate}

\textbf{Usability Issues Observed:}

{\def\LTcaptype{none} % do not increment counter
\begin{longtable}[]{@{}lcl@{}}
\toprule\noalign{}
Issue & Severity (H/M/L) & Stakeholder\textquotesingle s suggestion \\
\midrule\noalign{}
\endhead
\bottomrule\noalign{}
\endlastfoot
& & \\
& & \\
\end{longtable}
}

\begin{center}\rule{0.5\linewidth}{0.5pt}\end{center}

\subsubsection*{Section 3: Functional Assessment (10 minutes)}\label{docs__pandoc__latex__src__3dmake_all_lessons.md__section-3-functional-assessment-10-minutes}

Test each Functional Requirement explicitly.

{\def\LTcaptype{none} % do not increment counter
\begin{longtable}[]{@{}lll@{}}
\toprule\noalign{}
Requirement & Result & Notes \\
\midrule\noalign{}
\endhead
\bottomrule\noalign{}
\endlastfoot
FR1: & {[} {]} Pass {[} {]} Partial {[} {]} Fail & \\
FR2: & {[} {]} Pass {[} {]} Partial {[} {]} Fail & \\
FR3: & {[} {]} Pass {[} {]} Partial {[} {]} Fail & \\
\end{longtable}
}

\begin{enumerate}
\setcounter{enumi}{5}
\tightlist
\item
  "What\textquotesingle s missing that you expected to be there?"

  \begin{itemize}
  \tightlist
  \item
    \emph{Response:}
  \end{itemize}
\end{enumerate}

\begin{center}\rule{0.5\linewidth}{0.5pt}\end{center}

\subsubsection*{Section 4: Physical Assessment (10 minutes)}\label{docs__pandoc__latex__src__3dmake_all_lessons.md__section-4-physical-assessment-10-minutes}

\begin{enumerate}
\setcounter{enumi}{6}
\item
  Size: {[} {]} Too large {[} {]} Just right {[} {]} Too small

  \begin{itemize}
  \tightlist
  \item
    \emph{Notes:}
  \end{itemize}
\item
  Weight: {[} {]} Too heavy {[} {]} Just right {[} {]} Too light

  \begin{itemize}
  \tightlist
  \item
    \emph{Notes:}
  \end{itemize}
\item
  "Does the material / finish feel appropriate?"

  \begin{itemize}
  \tightlist
  \item
    \emph{Response:}
  \end{itemize}
\item
  "Any parts that feel fragile or rough?"

  \begin{itemize}
  \tightlist
  \item
    \emph{Response:}
  \end{itemize}
\end{enumerate}

\begin{center}\rule{0.5\linewidth}{0.5pt}\end{center}

\subsubsection*{Section 5: Overall Satisfaction (5 minutes)}\label{docs__pandoc__latex__src__3dmake_all_lessons.md__section-5-overall-satisfaction-5-minutes}

{\def\LTcaptype{none} % do not increment counter
\begin{longtable}[]{@{}lcl@{}}
\toprule\noalign{}
Criterion & Score (1--5) & Comments \\
\midrule\noalign{}
\endhead
\bottomrule\noalign{}
\endlastfoot
Overall satisfaction & & \\
Meets my needs & & \\
Design quality & & \\
Would recommend to someone else & & \\
\end{longtable}
}

\emph{(1 = Poor, 5 = Excellent)}

\begin{center}\rule{0.5\linewidth}{0.5pt}\end{center}

\subsubsection*{Section 6: Suggestions (10 minutes)}\label{docs__pandoc__latex__src__3dmake_all_lessons.md__section-6-suggestions-10-minutes}

\begin{enumerate}
\setcounter{enumi}{10}
\item
  "What one change would have the biggest positive impact?"

  \begin{itemize}
  \tightlist
  \item
    \emph{Response:}
  \end{itemize}
\item
  "What should we keep exactly as-is?"

  \begin{itemize}
  \tightlist
  \item
    \emph{Response:}
  \end{itemize}
\end{enumerate}

\begin{center}\rule{0.5\linewidth}{0.5pt}\end{center}

\subsubsection*{Post-Session Analysis}\label{docs__pandoc__latex__src__3dmake_all_lessons.md__post-session-analysis}

\textbf{Key Insights:}

\begin{enumerate}
\tightlist
\item
\item
\end{enumerate}

\subsection*{\texorpdfstring{\textbf{Critical Issues} (block use):}{Critical Issues (block use):}}\label{docs__pandoc__latex__src__3dmake_all_lessons.md__critical-issues-block-use}

\subsection*{\texorpdfstring{\textbf{Should-fix Issues} (reduce satisfaction):}{Should-fix Issues (reduce satisfaction):}}\label{docs__pandoc__latex__src__3dmake_all_lessons.md__should-fix-issues-reduce-satisfaction}

\textbf{Design Changes Resulting from This Session:}

{\def\LTcaptype{none} % do not increment counter
\begin{longtable}[]{@{}llll@{}}
\toprule\noalign{}
Original design & Feedback received & Change made & Rationale \\
\midrule\noalign{}
\endhead
\bottomrule\noalign{}
\endlastfoot
& & & \\
\end{longtable}
}

\emph{Session conducted by:} \_\_\_\_\_\_\_\_ \emph{Date:} \_\_\_\_\_\_\_\_

\begin{center}\rule{0.5\linewidth}{0.5pt}\end{center}

\subsection*{Template 5: Master Project Rubric}\label{docs__pandoc__latex__src__3dmake_all_lessons.md__template-5-master-project-rubric}

\emph{This rubric applies to all Unit 1, 2, and 3 projects unless the project briefing specifies otherwise. Project 0 is complete/incomplete only.}

\begin{center}\rule{0.5\linewidth}{0.5pt}\end{center}

\subsubsection*{Rubric Overview}\label{docs__pandoc__latex__src__3dmake_all_lessons.md__rubric-overview}

Projects are scored 0--9 across three categories, 3 points each:

{\def\LTcaptype{none} % do not increment counter
\begin{longtable}[]{@{}
  >{\raggedright\arraybackslash}p{(\linewidth - 4\tabcolsep) * \real{0.1837}}
  >{\centering\arraybackslash}p{(\linewidth - 4\tabcolsep) * \real{0.1020}}
  >{\raggedright\arraybackslash}p{(\linewidth - 4\tabcolsep) * \real{0.7143}}@{}}
\toprule\noalign{}
\begin{minipage}[b]{\linewidth}\raggedright
Category
\end{minipage} & \begin{minipage}[b]{\linewidth}\centering
Max Points
\end{minipage} & \begin{minipage}[b]{\linewidth}\raggedright
What Is Evaluated
\end{minipage} \\
\midrule\noalign{}
\endhead
\bottomrule\noalign{}
\endlastfoot
Problem \& Solution & 3 &
How well the prototype meets the identified problem or FRs \\
Design Quality & 3 &
Complexity, originality, technical quality; evidence of iteration \\
Documentation & 3 &
Completeness, accuracy, and thoughtfulness of all written deliverables \\
\end{longtable}
}

\begin{center}\rule{0.5\linewidth}{0.5pt}\end{center}

\subsubsection*{Category 1: Problem \& Solution (0--3 points)}\label{docs__pandoc__latex__src__3dmake_all_lessons.md__category-1-problem--solution-03-points}

{\def\LTcaptype{none} % do not increment counter
\begin{longtable}[]{@{}
  >{\centering\arraybackslash}p{(\linewidth - 2\tabcolsep) * \real{0.0403}}
  >{\raggedright\arraybackslash}p{(\linewidth - 2\tabcolsep) * \real{0.9597}}@{}}
\toprule\noalign{}
\begin{minipage}[b]{\linewidth}\centering
Score
\end{minipage} & \begin{minipage}[b]{\linewidth}\raggedright
Description
\end{minipage} \\
\midrule\noalign{}
\endhead
\bottomrule\noalign{}
\endlastfoot
3 &
Prototype clearly and effectively solves the stated problem. All FRs are met. Evidence of testing against requirements. \\
2 &
Prototype mostly meets the problem. Most FRs met. Minor gaps between design and requirements. \\
1 &
Prototype partially addresses the problem. Several FRs unmet or not clearly tested. \\
0 &
Prototype does not address the stated problem, or no FRs were established. \\
\end{longtable}
}

\begin{center}\rule{0.5\linewidth}{0.5pt}\end{center}

\subsubsection*{Category 2: Design Quality (0--3 points)}\label{docs__pandoc__latex__src__3dmake_all_lessons.md__category-2-design-quality-03-points}

{\def\LTcaptype{none} % do not increment counter
\begin{longtable}[]{@{}
  >{\centering\arraybackslash}p{(\linewidth - 2\tabcolsep) * \real{0.0244}}
  >{\raggedright\arraybackslash}p{(\linewidth - 2\tabcolsep) * \real{0.9756}}@{}}
\toprule\noalign{}
\begin{minipage}[b]{\linewidth}\centering
Score
\end{minipage} & \begin{minipage}[b]{\linewidth}\raggedright
Description
\end{minipage} \\
\midrule\noalign{}
\endhead
\bottomrule\noalign{}
\endlastfoot
3 &
Design is original, well-considered, technically executed. Code is clean, uses variables and modules appropriately, and is well-commented. Print is clean. Evidence of at least one meaningful revision. \\
2 &
Design is functional and shows thought. Code works but may lack structure (few comments, raw numbers). Print quality acceptable. Some iteration evident. \\
1 &
Design is basic or borrowed without modification. Code has issues. Print quality has unaddressed defects. \\
0 &
No meaningful original design. Print is non-functional or not completed. \\
\end{longtable}
}

\begin{center}\rule{0.5\linewidth}{0.5pt}\end{center}

\subsubsection*{Category 3: Documentation (0--3 points)}\label{docs__pandoc__latex__src__3dmake_all_lessons.md__category-3-documentation-03-points}

{\def\LTcaptype{none} % do not increment counter
\begin{longtable}[]{@{}
  >{\centering\arraybackslash}p{(\linewidth - 2\tabcolsep) * \real{0.0272}}
  >{\raggedright\arraybackslash}p{(\linewidth - 2\tabcolsep) * \real{0.9728}}@{}}
\toprule\noalign{}
\begin{minipage}[b]{\linewidth}\centering
Score
\end{minipage} & \begin{minipage}[b]{\linewidth}\raggedright
Description
\end{minipage} \\
\midrule\noalign{}
\endhead
\bottomrule\noalign{}
\endlastfoot
3 &
All required sections present, complete, and specific. Reflections are thoughtful and reference specific decisions, problems, and learning. Photos included. Measurements recorded. \\
2 &
Most sections present. Some sections vague or missing detail. Reflections show thought but are brief or generic. \\
1 &
Documentation incomplete. Major sections missing or one-line responses. Reflections minimal. \\
0 & Documentation not submitted or essentially empty. \\
\end{longtable}
}

\begin{center}\rule{0.5\linewidth}{0.5pt}\end{center}

\subsubsection*{Score Interpretation}\label{docs__pandoc__latex__src__3dmake_all_lessons.md__score-interpretation}

{\def\LTcaptype{none} % do not increment counter
\begin{longtable}[]{@{}
  >{\centering\arraybackslash}p{(\linewidth - 4\tabcolsep) * \real{0.1100}}
  >{\raggedright\arraybackslash}p{(\linewidth - 4\tabcolsep) * \real{0.3500}}
  >{\raggedright\arraybackslash}p{(\linewidth - 4\tabcolsep) * \real{0.5400}}@{}}
\toprule\noalign{}
\begin{minipage}[b]{\linewidth}\centering
Total Score
\end{minipage} & \begin{minipage}[b]{\linewidth}\raggedright
Interpretation
\end{minipage} & \begin{minipage}[b]{\linewidth}\raggedright
Next Step
\end{minipage} \\
\midrule\noalign{}
\endhead
\bottomrule\noalign{}
\endlastfoot
8--9 & Excellent work & Move to next project \\
6--7 & Good work with room for improvement &
Move on; instructor may suggest revisiting one element \\
4--5 & Meets basic expectations &
Resubmission of specific weak areas recommended \\
2--3 & Does not meet expectations & Resubmission required \\
0--1 & Missing major deliverables &
Meet with instructor; create a completion plan \\
\end{longtable}
}

\begin{center}\rule{0.5\linewidth}{0.5pt}\end{center}

\subsubsection*{Resubmission Policy}\label{docs__pandoc__latex__src__3dmake_all_lessons.md__resubmission-policy}

Students may resubmit any project as many times as needed to improve their score. Resubmissions must include a one-paragraph explanation of what was changed and why. The resubmission score replaces the original score.
-e

\begin{center}\rule{0.5\linewidth}{0.5pt}\end{center}

-e

\begin{center}\rule{0.5\linewidth}{0.5pt}\end{center}

\begin{center}\rule{0.5\linewidth}{0.5pt}\end{center}

\section*{Part IV --- Final Exam}\label{docs__pandoc__latex__src__3dmake_all_lessons.md__part-iv--final-exam}

\section*{3dMake Foundation --- Final Exam}\label{docs__pandoc__latex__src__3dmake_all_lessons.md__3dmake-foundation--final-exam}

Name: ~~~~~~~~~~~~~~~~~~~~~~~~~~~~~~~~~~~~~ Date:

\textbf{Total Points: 100 (4 points per problem)}

\textbf{Instructions:}

\begin{itemize}
\tightlist
\item
  Answer all 25 questions
\item
  For code errors, identify the specific problem and explain why it\textquotesingle s wrong
\item
  For behavioral questions, show your reasoning
\item
  For design problems, explain your approach and why it solves the stated challenge
\item
  You may reference the OpenSCAD documentation and appendices
\item
  All work must be your own
\end{itemize}

\begin{center}\rule{0.5\linewidth}{0.5pt}\end{center}

\subsection*{Section 1: Error Detection \& Code Analysis (Problems 1--10)}\label{docs__pandoc__latex__src__3dmake_all_lessons.md__section-1-error-detection--code-analysis-problems-110}

For each code block, identify any errors. If there is an error, explain what\textquotesingle s wrong and how to fix it. If the code is correct, state "No error" and explain what the code does.

\begin{center}\rule{0.5\linewidth}{0.5pt}\end{center}

\subsubsection*{Problem 1: Primitive Definition Error}\label{docs__pandoc__latex__src__3dmake_all_lessons.md__problem-1-primitive-definition-error}

\begin{lstlisting}[style=Alabaster, language=openscad]
cube([10, 10, 10], center=true);

sphere(r=5, $fn=16);

cylinder(h=20, r=8);


\end{lstlisting}

Is there an error in this code? If yes, identify it. If no, explain what this code renders.

\textbf{Answer:}

\begin{center}\rule{0.5\linewidth}{0.5pt}\end{center}

\subsubsection*{Problem 2: Transform Syntax Error}\label{docs__pandoc__latex__src__3dmake_all_lessons.md__problem-2-transform-syntax-error}

\begin{lstlisting}[style=Alabaster, language=openscad]
translate([5, 5, 0])

  rotate([0, 45, 0])

    cube([10, 10, 10], center=true);


\end{lstlisting}

Does this code have a syntax error? Explain what this code does.

\textbf{Answer:}

\begin{center}\rule{0.5\linewidth}{0.5pt}\end{center}

\subsubsection*{Problem 3: CSG Operation Error}\label{docs__pandoc__latex__src__3dmake_all_lessons.md__problem-3-csg-operation-error}

\begin{lstlisting}[style=Alabaster, language=openscad]
cube([20, 20, 20], center=true);

difference() {

  sphere(r=12);

  cylinder(h=30, r=5, center=true);

}


\end{lstlisting}

What is wrong with this CSG operation? Explain the fix.

\textbf{Answer:}

\begin{center}\rule{0.5\linewidth}{0.5pt}\end{center}

\subsubsection*{Problem 4: Module Definition Error}\label{docs__pandoc__latex__src__3dmake_all_lessons.md__problem-4-module-definition-error}

\begin{lstlisting}[style=Alabaster, language=openscad]
module bracket(width, height, depth) {

  cube([width, height, depth], center=true);

  translate([width/2 + 2, 0, 0])

    cube([4, height, depth/2], center=true);

}

bracket(20, 15, 10);


\end{lstlisting}

Is there an error in this module definition or call? Why or why not?

\textbf{Answer:}

\begin{center}\rule{0.5\linewidth}{0.5pt}\end{center}

\subsubsection*{Problem 5: Parameter \& Variable Scoping Error}\label{docs__pandoc__latex__src__3dmake_all_lessons.md__problem-5-parameter--variable-scoping-error}

\begin{lstlisting}[style=Alabaster, language=openscad]
wallthickness = 2;

module hollowcube(size) {

  difference() {

    cube([size, size, size], center=true);

    cube([size - wallthickness, size - wallthickness, size - wallthickness], center=true);

  }

}

hollowcube(20);


\end{lstlisting}

Will this code work correctly? If not, what is the problem and how would you fix it?

\textbf{Answer:}

\begin{center}\rule{0.5\linewidth}{0.5pt}\end{center}

\subsubsection*{Problem 6: Loop \& Iteration Error}\label{docs__pandoc__latex__src__3dmake_all_lessons.md__problem-6-loop--iteration-error}

\begin{lstlisting}[style=Alabaster, language=openscad]
for (i = [0:5:20]) {

  translate([i, 0, 0])

    cube([4, 4, 4]);

}


\end{lstlisting}

Will this code produce 5 cubes? Show the positions and explain why or why not.

\textbf{Answer:}

\begin{center}\rule{0.5\linewidth}{0.5pt}\end{center}

\subsubsection*{Problem 7: Center Parameter Misunderstanding}\label{docs__pandoc__latex__src__3dmake_all_lessons.md__problem-7-center-parameter-misunderstanding}

\begin{lstlisting}[style=Alabaster, language=openscad]
cube([10, 10, 20], center=false);

sphere(r=5);


\end{lstlisting}

What is the relationship between these two shapes? Where would the sphere appear relative to the cube?

\textbf{Answer:}

\begin{center}\rule{0.5\linewidth}{0.5pt}\end{center}

\subsubsection*{Problem 8: Intersection Error}\label{docs__pandoc__latex__src__3dmake_all_lessons.md__problem-8-intersection-error}

\begin{lstlisting}[style=Alabaster, language=openscad]
intersection() {

  cube([20, 20, 20], center=true);

  sphere(r=8, $fn=32);

}


\end{lstlisting}

Is there an error? What will this code render?

\textbf{Answer:}

\begin{center}\rule{0.5\linewidth}{0.5pt}\end{center}

\subsubsection*{Problem 9: Nested Transform Error}\label{docs__pandoc__latex__src__3dmake_all_lessons.md__problem-9-nested-transform-error}

\begin{lstlisting}[style=Alabaster, language=openscad]
translate([10, 0, 0])

  rotate([0, 0, 45])

    translate([5, 0, 0])

      cube([5, 5, 5], center=true);


\end{lstlisting}

Are the transforms applied in the correct order? Trace the final position of the cube.

\textbf{Answer:}

\begin{center}\rule{0.5\linewidth}{0.5pt}\end{center}

\subsubsection*{Problem 10: Resolution Parameter Error}\label{docs__pandoc__latex__src__3dmake_all_lessons.md__problem-10-resolution-parameter-error}

\begin{lstlisting}[style=Alabaster, language=openscad]
sphere(r=10, $fn=4);

cylinder(h=20, r=8, $fn=3);

cube([10, 10, 10]);


\end{lstlisting}

Identify the problem(s) with resolution in this code. What will happen when rendered?

\textbf{Answer:}

\begin{center}\rule{0.5\linewidth}{0.5pt}\end{center}

\subsection*{Section 2: Code Behavior \& Theory (Problems 11--17)}\label{docs__pandoc__latex__src__3dmake_all_lessons.md__section-2-code-behavior--theory-problems-1117}

For each question, show your reasoning. You may draw diagrams if helpful.

\begin{center}\rule{0.5\linewidth}{0.5pt}\end{center}

\subsubsection*{Problem 11: Vertex Coordinates}\label{docs__pandoc__latex__src__3dmake_all_lessons.md__problem-11-vertex-coordinates}

A cube is defined as \texttt{cube({[}10,\ 10,\ 20{]},\ center=false)}.

\textbf{a)} List the XYZ coordinates of all 8 vertices.

\textbf{Answer:}

\begin{itemize}
\tightlist
\item
  Vertex 1:
\item
  Vertex 2:
\item
  Vertex 3:
\item
  Vertex 4:
\item
  Vertex 5:
\item
  Vertex 6:
\item
  Vertex 7:
\item
  Vertex 8:
\end{itemize}

\textbf{b)} Now define the SAME cube with \texttt{center=true}. List the NEW coordinates of all 8 vertices.

\textbf{Answer:}

\begin{itemize}
\tightlist
\item
  Vertex 1:
\item
  Vertex 2:
\item
  Vertex 3:
\item
  Vertex 4:
\item
  Vertex 5:
\item
  Vertex 6:
\item
  Vertex 7:
\item
  Vertex 8:
\end{itemize}

\begin{center}\rule{0.5\linewidth}{0.5pt}\end{center}

\subsubsection*{Problem 12: Sphere Geometry}\label{docs__pandoc__latex__src__3dmake_all_lessons.md__problem-12-sphere-geometry}

Explain the difference between \texttt{sphere(r=10,\ \$fn=8)} and \texttt{sphere(r=10,\ \$fn=128)}.

Which would you use for a prototype and which for final printing? Why?

\textbf{Answer:}

\begin{center}\rule{0.5\linewidth}{0.5pt}\end{center}

\subsubsection*{Problem 13: Transform Order}\label{docs__pandoc__latex__src__3dmake_all_lessons.md__problem-13-transform-order}

Given this code:

\begin{lstlisting}[style=Alabaster, language=openscad]
translate([10, 0, 0])

  rotate([0, 0, 45])

    cube([5, 5, 5], center=true);


\end{lstlisting}

Does the order matter? What if you swap translate and rotate? Show both final positions.

\textbf{Answer:}

\begin{center}\rule{0.5\linewidth}{0.5pt}\end{center}

\subsubsection*{Problem 14: Boolean Operation Behavior}\label{docs__pandoc__latex__src__3dmake_all_lessons.md__problem-14-boolean-operation-behavior}

You have a solid cube and you want to create a hole through it. Which CSG operation would you use: \texttt{union()}, \texttt{difference()}, or \texttt{intersection()}?

Explain your choice and write pseudocode showing how you\textquotesingle d accomplish this.

\textbf{Answer:}

\begin{center}\rule{0.5\linewidth}{0.5pt}\end{center}

\subsubsection*{Problem 15: Parametric Design Advantage}\label{docs__pandoc__latex__src__3dmake_all_lessons.md__problem-15-parametric-design-advantage}

Compare these two approaches:

\textbf{Approach A:} Hard-coded cube with fixed dimensions

\begin{lstlisting}[style=Alabaster, language=openscad]
cube([10, 10, 20]);


\end{lstlisting}

\textbf{Approach B:} Parametric cube

\begin{lstlisting}[style=Alabaster, language=openscad]
module parametricbox(width, height, depth) {

  cube([width, height, depth], center=true);

}

parametricbox(10, 10, 20);


\end{lstlisting}

Why is Approach B better for design iteration? Give an example of how you\textquotesingle d use it.

\textbf{Answer:}

\begin{center}\rule{0.5\linewidth}{0.5pt}\end{center}

\subsubsection*{Problem 16: Scale Transform Behavior}\label{docs__pandoc__latex__src__3dmake_all_lessons.md__problem-16-scale-transform-behavior}

If you apply \texttt{scale({[}2,\ 1,\ 0.5{]})} to a \texttt{cube({[}10,\ 10,\ 10{]},\ center=true)}, what are the NEW dimensions of the cube?

\textbf{Answer --- New dimensions:}

Show your calculation:

\begin{center}\rule{0.5\linewidth}{0.5pt}\end{center}

\subsubsection*{Problem 17: Library Organization}\label{docs__pandoc__latex__src__3dmake_all_lessons.md__problem-17-library-organization}

You\textquotesingle ve created three useful modules:

\begin{itemize}
\tightlist
\item
  \texttt{bracket(width,\ height,\ depth)}
\item
  \texttt{hollowcube(size,\ wallthickness)}
\item
  \texttt{connectorpin(diameter,\ height)}
\end{itemize}

How would you organize these into a reusable library? What file structure would you create and why?

\textbf{Answer:}

\begin{center}\rule{0.5\linewidth}{0.5pt}\end{center}

\subsection*{Section 3: Design \& Problem-Solving (Problems 18--25)}\label{docs__pandoc__latex__src__3dmake_all_lessons.md__section-3-design--problem-solving-problems-1825}

These problems test your ability to design solutions, debug real-world issues, and think creatively.

\begin{center}\rule{0.5\linewidth}{0.5pt}\end{center}

\subsubsection*{Problem 18: Tolerance Design Challenge}\label{docs__pandoc__latex__src__3dmake_all_lessons.md__problem-18-tolerance-design-challenge}

You\textquotesingle re designing a snap-fit connector. The male part has a thickness of 2mm. The female slot needs to accommodate this part with enough flexibility to snap but not fall out.

Should the slot be:

\begin{itemize}
\tightlist
\item
  a) Exactly 2mm wide
\item
  b) 2.1mm wide
\item
  c) 1.9mm wide
\item
  d) 2.5mm wide
\end{itemize}

Explain your choice and the design thinking behind it.

\textbf{Answer:}

\begin{center}\rule{0.5\linewidth}{0.5pt}\end{center}

\subsubsection*{Problem 19: Design Iteration Problem}\label{docs__pandoc__latex__src__3dmake_all_lessons.md__problem-19-design-iteration-problem}

You print a keycap with \texttt{keysize=12} and the text embossing is too shallow to feel. Your code uses:

\begin{lstlisting}[style=Alabaster, language=openscad]
linear_extrude(height=1)

  text("A", size=8);


\end{lstlisting}

What parameter(s) would you adjust to make the embossing deeper? Show your new code.

\textbf{Answer:}

\begin{center}\rule{0.5\linewidth}{0.5pt}\end{center}

\subsubsection*{Problem 20: Error Diagnosis}\label{docs__pandoc__latex__src__3dmake_all_lessons.md__problem-20-error-diagnosis}

Your 3dMake build fails with this error: "Geometry is non-manifold." You have this code:

\begin{lstlisting}[style=Alabaster, language=openscad]
difference() {

  cube([20, 20, 20], center=true);

  cylinder(h=30, r=4, center=true);

}


\end{lstlisting}

Why might this fail? What\textquotesingle s the common fix for non-manifold geometry?

\textbf{Answer:}

\begin{center}\rule{0.5\linewidth}{0.5pt}\end{center}

\subsubsection*{Problem 21: Multi-Part Assembly}\label{docs__pandoc__latex__src__3dmake_all_lessons.md__problem-21-multi-part-assembly}

You\textquotesingle re designing a two-part box (lid + base). The base has dimensions {[}50, 30, 20{]}. The lid should sit on top of the base.

Write parametric modules for both parts and show how you\textquotesingle d position them together. Include appropriate positioning logic.

\textbf{Answer:}

\begin{lstlisting}[style=Alabaster, language=openscad]
module base(length, width, height) {

  // Your code here

}


module lid(length, width, height) {

  // Your code here

}


// Positioning code here:


\end{lstlisting}

\begin{center}\rule{0.5\linewidth}{0.5pt}\end{center}

\subsubsection*{Problem 22: Optimization Challenge}\label{docs__pandoc__latex__src__3dmake_all_lessons.md__problem-22-optimization-challenge}

You have a design that takes 5 minutes to render. You notice you have:

\begin{lstlisting}[style=Alabaster, language=openscad]
sphere(r=10, $fn=256);

cylinder(h=20, r=8, $fn=256);

cube([20, 20, 20]);


\end{lstlisting}

Which parameter(s) would you reduce to speed up rendering while maintaining acceptable quality for a prototype? Explain your choices.

\textbf{Answer:}

\begin{center}\rule{0.5\linewidth}{0.5pt}\end{center}

\subsubsection*{Problem 23: Real-World Constraint Problem}\label{docs__pandoc__latex__src__3dmake_all_lessons.md__problem-23-real-world-constraint-problem}

A stakeholder requests a custom handle for a tool. They specify:

\begin{itemize}
\tightlist
\item
  Must fit a hand (approximately 80mm long)
\item
  Must accommodate fingers 60mm long inside
\item
  Wall thickness must be at least 3mm for durability
\item
  Should be ergonomic (slightly curved)
\end{itemize}

Sketch or describe a parametric design for this handle. What parameters would you expose to allow customization?

\textbf{Answer:}

\begin{center}\rule{0.5\linewidth}{0.5pt}\end{center}

\subsubsection*{Problem 24: Code Reusability Challenge}\label{docs__pandoc__latex__src__3dmake_all_lessons.md__problem-24-code-reusability-challenge}

You\textquotesingle ve created a single keycap module. Now you need to create a keyboard with 5 keys arranged in a row, spaced 15mm apart.

Write code using a loop that creates 5 keycaps with letters A--E, properly spaced.

\begin{quote}
\textbf{Hint:} In OpenSCAD, strings are indexed like lists. \texttt{"ABCDE"{[}0{]}} returns \texttt{"A"}, \texttt{"ABCDE"{[}1{]}} returns \texttt{"B"}, and so on. The \texttt{len()} function returns the number of characters in a string.
\end{quote}

\textbf{Answer:}

\begin{lstlisting}[style=Alabaster, language=openscad]
module keycap(letter, keysize=10) {

  // Keycap code here (you can assume this exists)

}


// Your loop code here:


\end{lstlisting}

\begin{center}\rule{0.5\linewidth}{0.5pt}\end{center}

\subsubsection*{Problem 25: Design Thinking \& Iteration}\label{docs__pandoc__latex__src__3dmake_all_lessons.md__problem-25-design-thinking--iteration}

You\textquotesingle ve printed Iteration 1 of a product and measured the results. The wall thickness is 3mm but feels too fragile. In Iteration 2, you increased it to 5mm, and now it feels too rigid and won\textquotesingle t flex as intended.

For Iteration 3, what thickness would you try and why? How would you make this decision more scientific/data-driven?

\textbf{Answer:}

\begin{center}\rule{0.5\linewidth}{0.5pt}\end{center}

\subsection*{Bonus Challenge (Optional, +5 points)}\label{docs__pandoc__latex__src__3dmake_all_lessons.md__bonus-challenge-optional-5-points}

Design a parametric model for a custom assistive technology device (e.g., a tactile measuring tool, a custom gripper, an adapted eating utensil, etc.).

\begin{itemize}
\tightlist
\item
  Identify the user\textquotesingle s specific need
\item
  Specify the key dimensions and parameters
\item
  Write at least one module with realistic dimensions
\item
  Explain how the design would be tested and iterated
\end{itemize}

\textbf{Answer:}

\begin{lstlisting}[style=Alabaster, language=openscad]
// Your design here:


\end{lstlisting}

\textbf{User Need:}

\textbf{Parameters:}

\textbf{Testing Plan:}

\begin{center}\rule{0.5\linewidth}{0.5pt}\end{center}

\subsection*{Scoring Rubric}\label{docs__pandoc__latex__src__3dmake_all_lessons.md__scoring-rubric}

{\def\LTcaptype{none} % do not increment counter
\begin{longtable}[]{@{}
  >{\raggedright\arraybackslash}p{(\linewidth - 2\tabcolsep) * \real{0.0698}}
  >{\raggedright\arraybackslash}p{(\linewidth - 2\tabcolsep) * \real{0.9302}}@{}}
\toprule\noalign{}
\begin{minipage}[b]{\linewidth}\raggedright
Points
\end{minipage} & \begin{minipage}[b]{\linewidth}\raggedright
Criteria
\end{minipage} \\
\midrule\noalign{}
\endhead
\bottomrule\noalign{}
\endlastfoot
4 &
Correct answer with clear, complete explanation; demonstrates deep understanding \\
3 & Mostly correct answer; minor gaps in explanation or reasoning \\
2 &
Partially correct; shows some understanding but has significant gaps \\
1 & Minimal effort; shows limited understanding \\
0 & No answer or completely incorrect \\
\end{longtable}
}

\textbf{Total Possible: 100 + 5 bonus = 105 points}

\begin{center}\rule{0.5\linewidth}{0.5pt}\end{center}

\subsection*{References You May Use}\label{docs__pandoc__latex__src__3dmake_all_lessons.md__references-you-may-use}

\begin{itemize}
\tightlist
\item
  \hyperref[docs__pandoc__latex__src__3dmake_appendices_complete.md__appendix-a]{3dMake Quick Reference}
\item
  \hyperref[docs__pandoc__latex__src__3dmake_appendices_complete.md__appendix-a]{OpenSCAD Cheat Sheet}
\item
  \hyperref[docs__pandoc__latex__src__3dmake_appendices_complete.md__appendix-b]{Appendix A: Comprehensive Slicing Guide}
\item
  \hyperref[docs__pandoc__latex__src__3dmake_appendices_complete.md__appendix-d]{Appendix C: Tolerance Testing \& QA Matrix}
\item
  \hyperref[docs__pandoc__latex__src__cli_appendices_complete.md__appendix-b]{Appendix D: PowerShell Integration}
\end{itemize}

\begin{center}\rule{0.5\linewidth}{0.5pt}\end{center}

\emph{End of Final Exam}

\textbf{Submission Instructions:}

\begin{enumerate}
\tightlist
\item
  Answer all 25 questions completely
\item
  Show your work for calculations and reasoning
\item
  Include code samples where requested
\item
  Submit as a PDF with your name and date
\item
  Scoring will be based on correctness, clarity, and depth of understanding
\end{enumerate}

\emph{Good luck!}

\begin{center}\rule{0.5\linewidth}{0.5pt}\end{center}

\begin{center}\rule{0.5\linewidth}{0.5pt}\end{center}

\begin{center}\rule{0.5\linewidth}{0.5pt}\end{center}

\section*{INSTRUCTOR APPENDIX}\label{docs__pandoc__latex__src__3dmake_all_lessons.md__instructor-appendix}

\subsubsection*{Not for student distribution}\label{docs__pandoc__latex__src__3dmake_all_lessons.md__not-for-student-distribution}

\begin{center}\rule{0.5\linewidth}{0.5pt}\end{center}

\subsection*{Exam Administration Notes}\label{docs__pandoc__latex__src__3dmake_all_lessons.md__exam-administration-notes}

\subsubsection*{Sequencing --- Problem 20}\label{docs__pandoc__latex__src__3dmake_all_lessons.md__sequencing--problem-20}

Problem 20 tests non-manifold geometry diagnosis, which is formally introduced in \textbf{Lesson 10}.

\begin{itemize}
\tightlist
\item
  \textbf{If administered before Lesson 10 is complete:} Problem 20 is not counted in the passing threshold for Lessons 1--9. Consider adding a notice directing students to skip and return to it after Lesson 10.
\item
  \textbf{If administered after all 10 lessons:} No changes needed --- the problem is fair game.
\end{itemize}

\begin{center}\rule{0.5\linewidth}{0.5pt}\end{center}

\subsection*{Answer Keys \& Rubrics}\label{docs__pandoc__latex__src__3dmake_all_lessons.md__answer-keys--rubrics}

\subsubsection*{Problem 20 --- Non-Manifold Geometry (Extended Model Answer)}\label{docs__pandoc__latex__src__3dmake_all_lessons.md__problem-20--non-manifold-geometry-extended-model-answer}

The exam version tests the core concept with a simple case. The extended variant below can be used as an enrichment rubric or for a more advanced cohort. It contains \textbf{three non-manifold issues:}

\begin{lstlisting}[style=Alabaster, language=openscad]
// PROBLEM CODE — find and fix all issues

difference() {

  cube([40, 40, 25]);

  translate([10, 10, 0])        // NON-MANIFOLD: co-planar bottom face

    cylinder(r=8, h=25, $fn=24);

  translate([20, 20, 25])       // NON-MANIFOLD: co-planar top face

    cylinder(r=5, h=10, $fn=24);

}


// NON-MANIFOLD: these two cubes share a face (just touching, not overlapping)

translate([40, 0, 0])  cube([15, 40, 25]);

translate([55, 0, 0])  cube([15, 40, 25]);


\end{lstlisting}

\begin{enumerate}
\tightlist
\item
  \textbf{Co-planar bottom face} --- the first cylinder\textquotesingle s base is flush with the cube\textquotesingle s base (\texttt{z=0}).
\item
  \textbf{Co-planar top face} --- the second cylinder starts exactly at the cube\textquotesingle s top (\texttt{z=25}).
\item
  \textbf{Touching faces} --- the two side cubes share an exact face (just touching, not overlapping).
\end{enumerate}

\textbf{Corrected version:}

\begin{lstlisting}[style=Alabaster, language=openscad]
difference() {

  cube([40, 40, 25]);


  // FIX 1: extend cylinder 0.001 below base AND above top

  translate([10, 10, -0.001])

    cylinder(r=8, h=25.002, $fn=24);


  // FIX 2: extend the top-entry cylinder 0.001 into the solid

  translate([20, 20, 25 - 9.999])

    cylinder(r=5, h=10.001, $fn=24);

}


// FIX 3: overlap the two touching cubes by 0.001 and use explicit union()

union() {

  translate([40, 0, 0])  cube([15.001, 40, 25]);

  translate([55, 0, 0])  cube([15, 40, 25]);

}


\end{lstlisting}

\textbf{Extended rubric (10 points):}

\begin{itemize}
\tightlist
\item
  Identifies all 3 issues: 3 pts
\item
  Fixes issue 1 (co-planar bottom) correctly: 2 pts
\item
  Fixes issue 2 (co-planar top) correctly: 2 pts
\item
  Fixes issue 3 (touching faces) correctly: 2 pts
\item
  Model builds cleanly with no CGAL warnings: 1 pt
\end{itemize}

\begin{center}\rule{0.5\linewidth}{0.5pt}\end{center}

\subsubsection*{Problem 24 --- Keycap Loop (Model Answer)}\label{docs__pandoc__latex__src__3dmake_all_lessons.md__problem-24--keycap-loop-model-answer}

\begin{lstlisting}[style=Alabaster, language=openscad]
letters  = "ABCDE";

keysize  = 10;

key_gap  = 15;


for (i = [0 : len(letters) - 1]) {

  translate([i * (keysize + key_gap), 0, 0])

    keycap(letters[i], keysize);

}


\end{lstlisting}

\textbf{Expected output:} 5 keycaps labeled A, B, C, D, E, evenly spaced in a row.

\textbf{Rubric (4 points):}

\begin{itemize}
\tightlist
\item
  Correct \texttt{for} with range syntax: 1 pt
\item
  Correct string indexing \texttt{letters{[}i{]}}: 1 pt
\item
  Correct \texttt{translate()} to space keys apart: 1 pt
\item
  All 5 keycaps render without errors: 1 pt
\end{itemize}

\begin{center}\rule{0.5\linewidth}{0.5pt}\end{center}

\subsection*{Extension Projects (Optional / Post-Exam)}\label{docs__pandoc__latex__src__3dmake_all_lessons.md__extension-projects-optional--post-exam}

These projects are not part of the graded exam. Assign as enrichment after students complete all 10 lessons.

\begin{center}\rule{0.5\linewidth}{0.5pt}\end{center}

\subsubsection*{Extension A --- Recursive Fractal Geometry}\label{docs__pandoc__latex__src__3dmake_all_lessons.md__extension-a--recursive-fractal-geometry}

Recommended after Lesson 3. Demonstrates recursive OpenSCAD for visual and mathematical interest.

\begin{quote}
\textbf{Warning:} Level 4+ will take minutes to render --- stay at Level ≤ 3.
\end{quote}

\begin{lstlisting}[style=Alabaster, language=openscad]
module menger(size=27, level=2) {

  if (level == 0) {

    cube(size, center=true);

  } else {

    s3 = size / 3;

    difference() {

      cube(size, center=true);

      for (x=[-1,0,1]) for (y=[-1,0,1]) for (z=[-1,0,1]) {

        if ((x==0 ? 1 : 0) + (y==0 ? 1 : 0) + (z==0 ? 1 : 0) >= 2) {

          translate([x*s3, y*s3, z*s3])

            cube(s3 + 0.002, center=true);

        }

      }

    }

    for (x=[-1,0,1]) for (y=[-1,0,1]) for (z=[-1,0,1]) {

      if (!((x==0?1:0)+(y==0?1:0)+(z==0?1:0) >= 2)) {

        translate([x*s3, y*s3, z*s3])

          menger(s3, level-1);

      }

    }

  }

}


menger(size=27, level=2);  // change level to 0, 1, or 2


\end{lstlisting}

\begin{center}\rule{0.5\linewidth}{0.5pt}\end{center}

\subsubsection*{Extension B --- Generative Pattern Arrays}\label{docs__pandoc__latex__src__3dmake_all_lessons.md__extension-b--generative-pattern-arrays}

Demonstrates list comprehensions and math functions for artistic/functional array generation.

\begin{lstlisting}[style=Alabaster, language=openscad]
plate_w = 120;

plate_d = 80;

plate_h = 3;


hole_d    = 4;

grid_s    = 8;

wave_amp  = 3;

wave_freq = 0.1;


module perforated_panel() {

  difference() {

    cube([plate_w, plate_d, plate_h]);

    cols = floor(plate_w / grid_s);

    rows = floor(plate_d / grid_s);

    for (cx = [1 : cols - 1]) {

      for (ry = [1 : rows - 1]) {

        x = cx * grid_s;

        y = ry * grid_s;

        wave_val = sin(x * wave_freq * 360) * sin(y * wave_freq * 360);

        d = hole_d + wave_amp * wave_val;

        if (d > 0.5) {

          translate([x, y, -0.001])

            cylinder(d=d, h=plate_h + 0.002, $fn=16);

        }

      }

    }

  }

}


perforated_panel();


\end{lstlisting}

\begin{center}\rule{0.5\linewidth}{0.5pt}\end{center}

\subsubsection*{Extension C --- Multi-Platform Build Pipeline}\label{docs__pandoc__latex__src__3dmake_all_lessons.md__extension-c--multi-platform-build-pipeline}

Full build pipeline scripts for all three shell environments.

\paragraph*{Makefile (Linux/macOS)}\label{docs__pandoc__latex__src__3dmake_all_lessons.md__makefile-linuxmacos}

\begin{lstlisting}[style=Alabaster]
.PHONY: all build clean variants archive test


SRC     := src/main.scad

OUTPUT  := build/main.stl

WIDTHS  := 20 30 40 50 80


all: build test archive


build: $(OUTPUT)


$(OUTPUT): $(SRC) | build/

  openscad -o $@ $<

  3dm info


build/:

  mkdir -p build build/variants build/logs


clean:

  rm -rf build/


variants: build/

  @for w in $(WIDTHS); do \

    echo "Building width=$$w..."; \

    openscad -D "width=$$w" -o "build/variants/main_w$$w.stl" $(SRC); \

  done


archive: build/

  $(eval TS := $(shell date +%Y%m%d_%H%M%S))

  mkdir -p "archives/$(TS)"

  cp $(OUTPUT) "archives/$(TS)/"

  cp $(SRC)    "archives/$(TS)/"

  3dm info > "archives/$(TS)/BUILD_INFO.txt"


test: build/

  @for w in 10 30 80; do \

    openscad -D "width=$$w" -o "/tmp/test_w$$w.stl" $(SRC) && \

    echo "  width=$$w: OK" || echo "  width=$$w: FAILED"; \

  done


\end{lstlisting}

\paragraph*{PowerShell (Windows)}\label{docs__pandoc__latex__src__3dmake_all_lessons.md__powershell-windows}

\begin{lstlisting}[style=Alabaster, language=powershell]
param($Width = 30)


$src    = "src\main.scad"

$output = "build\main.stl"

$widths = @(20, 30, 40, 50, 80)


task init {

    New-Item -Force -ItemType Directory -Path "build","build\variants","build\logs" | Out-Null

}


task build init, {

    exec { openscad -D "width=$Width" -o $output $src }

    exec { 3dm info }

}


task clean {

    Remove-Item -Recurse -Force "build" -ErrorAction SilentlyContinue

}


task variants init, {

    foreach ($w in $widths) {

        exec { openscad -D "width=$w" -o "build\variants\main_w$w.stl" $src }

    }

}


task archive build, {

    $ts  = Get-Date -Format "yyyyMMdd_HHmmss"

    $dir = "archives\$ts"

    New-Item -Force -ItemType Directory -Path $dir | Out-Null

    Copy-Item $output $dir

    Copy-Item $src    $dir

    & 3dm info | Out-File "$dir\BUILD_INFO.txt" -Encoding UTF8

}


task test init, {

    foreach ($w in @(10, 30, 80)) {

        try {

            exec { openscad -D "width=$w" -o "build\test_w$w.stl" $src }

            Write-Host "  width=$w: OK"

        } catch { Write-Warning "  width=$w: FAILED" }

    }

}


task all build, test, archive


\end{lstlisting}

\paragraph*{CMD Batch (Windows Legacy)}\label{docs__pandoc__latex__src__3dmake_all_lessons.md__cmd-batch-windows-legacy}

\begin{lstlisting}[style=Alabaster, language=cmd]
@echo off

SET ACTION=%1

IF "%ACTION%"=="" SET ACTION=all


IF "%ACTION%"=="build"    GOTO :BUILD

IF "%ACTION%"=="clean"    GOTO :CLEAN

IF "%ACTION%"=="variants" GOTO :VARIANTS

IF "%ACTION%"=="archive"  GOTO :ARCHIVE

IF "%ACTION%"=="all"      GOTO :ALL


:ALL

CALL :BUILD_STEP

IF %ERRORLEVEL% NEQ 0 EXIT /b 1

CALL :VARIANTS_STEP

CALL :ARCHIVE_STEP

GOTO :EOF


:BUILD_STEP

mkdir build 2>nul

openscad -o build\main.stl src\main.scad

IF %ERRORLEVEL% NEQ 0 (echo [BUILD] FAILED & EXIT /b 1)

3dm info

EXIT /b 0


:CLEAN

rmdir /s /q build 2>nul

GOTO :EOF


:VARIANTS_STEP

mkdir build\variants 2>nul

FOR %%W IN (20 30 40 50 80) DO (

    openscad -D "width=%%W" -o "build\variants\main_w%%W.stl" src\main.scad

)

EXIT /b 0


:ARCHIVE_STEP

FOR /f "tokens=2 delims==" %%D IN ('wmic os get localdatetime /value') DO SET DT=%%D

SET TS=%DT:~0,8%_%DT:~8,6%

mkdir "archives\%TS%" 2>nul

copy build\main.stl "archives\%TS%\" >nul

copy src\main.scad  "archives\%TS%\" >nul

3dm info >> "archives\%TS%\BUILD_INFO.txt" 2>&1

EXIT /b 0


\end{lstlisting}

\begin{center}\rule{0.5\linewidth}{0.5pt}\end{center}

\subsection*{Quick Reference --- Shell Equivalents}\label{docs__pandoc__latex__src__3dmake_all_lessons.md__quick-reference--shell-equivalents}

{\def\LTcaptype{none} % do not increment counter
\begin{longtable}[]{@{}
  >{\raggedright\arraybackslash}p{(\linewidth - 6\tabcolsep) * \real{0.1374}}
  >{\raggedright\arraybackslash}p{(\linewidth - 6\tabcolsep) * \real{0.2519}}
  >{\raggedright\arraybackslash}p{(\linewidth - 6\tabcolsep) * \real{0.3130}}
  >{\raggedright\arraybackslash}p{(\linewidth - 6\tabcolsep) * \real{0.2977}}@{}}
\toprule\noalign{}
\begin{minipage}[b]{\linewidth}\raggedright
Task
\end{minipage} & \begin{minipage}[b]{\linewidth}\raggedright
bash (Linux/macOS)
\end{minipage} & \begin{minipage}[b]{\linewidth}\raggedright
PowerShell (Windows)
\end{minipage} & \begin{minipage}[b]{\linewidth}\raggedright
CMD (Windows)
\end{minipage} \\
\midrule\noalign{}
\endhead
\bottomrule\noalign{}
\endlastfoot
Build & \texttt{3dm\ build} & \texttt{3dm\ build} &
\texttt{3dm\ build} \\
Info & \texttt{3dm\ info} & \texttt{3dm\ info} & \texttt{3dm\ info} \\
Parameter (number) & \texttt{-D\ "width=50"} & \texttt{-D\ "width=50"} &
\texttt{-D\ "width=50"} \\
Parameter (string) &
\texttt{-D\ \textquotesingle{}label="A"\textquotesingle{}} &
\texttt{-D\ \textquotesingle{}label="A"\textquotesingle{}} &
\texttt{-D\ "label=""A"""} \\
Loop over list & \texttt{for\ x\ in\ a\ b\ c;\ do\ ...\ done} &
\texttt{foreach\ (\$x\ in\ @(...))\ \{...\}} &
\texttt{FOR\ \%\%X\ IN\ (a\ b\ c)\ DO\ ...} \\
Make directory & \texttt{mkdir\ -p\ dir} &
\texttt{New-Item\ -Force\ -ItemType\ Directory\ dir} &
\texttt{mkdir\ dir\ 2\textgreater{}nul} \\
Copy file & \texttt{cp\ src\ dst} & \texttt{Copy-Item\ src\ dst} &
\texttt{copy\ src\ dst} \\
Check exit code & \texttt{if\ {[}\ \$?\ -ne\ 0\ {]}} &
\texttt{if\ (\$LASTEXITCODE\ -ne\ 0)} &
\texttt{IF\ \%ERRORLEVEL\%\ NEQ\ 0} \\
Write to file & \texttt{echo\ text\ \textgreater{}\ file} &
\texttt{"text"\ \textbar{}\ Out-File\ file} &
\texttt{echo\ text\ \textgreater{}\ file} \\
Append to file & \texttt{echo\ text\ \textgreater{}\textgreater{}\ file}
& \texttt{"text"\ \textbar{}\ Add-Content\ file} &
\texttt{echo\ text\ \textgreater{}\textgreater{}\ file} \\
Date/time & \texttt{date\ +\%Y\%m\%d\_\%H\%M\%S} &
\texttt{Get-Date\ -Format\ "yyyyMMdd\_HHmmss"} &
\texttt{wmic\ os\ get\ localdatetime} \\
File size (bytes) & \texttt{stat\ -c\%s\ file} &
\texttt{(Get-Item\ file).Length} &
\texttt{FOR\ /F\ "tokens=4"\ \%\%S\ IN\ (\textquotesingle{}dir\ file\textquotesingle{})} \\
Watch for changes &
\texttt{ls\ src/*.scad\ \textbar{}\ entr\ 3dm\ build} &
(see Lesson 9 watcher) & (use PowerShell) \\
-e & & & \\
\end{longtable}
}

\begin{center}\rule{0.5\linewidth}{0.5pt}\end{center}

\begin{center}\rule{0.5\linewidth}{0.5pt}\end{center}

\section*{Part V --- Quick Reference Guide}\label{docs__pandoc__latex__src__3dmake_all_lessons.md__part-v--quick-reference-guide}

\section*{3dMake Foundation Quick Reference Guide}\label{docs__pandoc__latex__src__3dmake_all_lessons.md__3dmake_foundation-3dmake_quick_reference}

Fast lookup for lessons, projects, resources, and common tasks

\subsection*{Lesson Quick Reference}\label{docs__pandoc__latex__src__3dmake_all_lessons.md__lesson-quick-reference}

\subsubsection*{All 11 Lessons at a Glance}\label{docs__pandoc__latex__src__3dmake_all_lessons.md__all-11-lessons-at-a-glance}

{\def\LTcaptype{none} % do not increment counter
\begin{longtable}[]{@{}
  >{\raggedright\arraybackslash}p{(\linewidth - 10\tabcolsep) * \real{0.0308}}
  >{\raggedright\arraybackslash}p{(\linewidth - 10\tabcolsep) * \real{0.2231}}
  >{\raggedright\arraybackslash}p{(\linewidth - 10\tabcolsep) * \real{0.0769}}
  >{\raggedright\arraybackslash}p{(\linewidth - 10\tabcolsep) * \real{0.1154}}
  >{\raggedright\arraybackslash}p{(\linewidth - 10\tabcolsep) * \real{0.3692}}
  >{\raggedright\arraybackslash}p{(\linewidth - 10\tabcolsep) * \real{0.1846}}@{}}
\toprule\noalign{}
\begin{minipage}[b]{\linewidth}\raggedright
\#
\end{minipage} & \begin{minipage}[b]{\linewidth}\raggedright
Title
\end{minipage} & \begin{minipage}[b]{\linewidth}\raggedright
Duration
\end{minipage} & \begin{minipage}[b]{\linewidth}\raggedright
Level
\end{minipage} & \begin{minipage}[b]{\linewidth}\raggedright
Main Topics
\end{minipage} & \begin{minipage}[b]{\linewidth}\raggedright
Key Project
\end{minipage} \\
\midrule\noalign{}
\endhead
\bottomrule\noalign{}
\endlastfoot
1 & Environmental Configuration & 60-90m & Beginner &
Setup, project structure, \texttt{3dm\ build} & None \\
2 & Geometric Primitives \& CSG & 60m & Beginner &
Primitives, CSG operations, debugging & None \\
3 & Parametric Architecture & 60m & Beginner+ &
Modules, libraries, parameters & None \\
4 & AI Verification & 45-60m & Intermediate &
\texttt{3dm\ info}, validation, design documentation & None \\
5 & Safety \& Physical Interface & 60-90m & Intermediate &
Safety protocols, materials, pre-print checks & None \\
6 & 3dm Commands \& Text & 60-90m & Intermediate &
\texttt{3dm\ info/preview/orient/slice}, embossing & Keycap \\
7 & Parametric Transforms & 75-90m & Intermediate+ &
Transforms, multi-part design, assembly & Phone Stand \\
8 & Advanced Parametric Design & 90-120m & Advanced &
Tolerance, interlocking features, snap-fits & Stackable Bins \\
9 & Automation \& Workflows & 60-90m & Advanced &
PowerShell scripting, batch processing, CI/CD &
{[}key{]} Batch Automation \\
10 & Troubleshooting \& Mastery & 120-150m & Advanced &
Measurement, QA testing, diagnostics & {[}dice{]} QA +  Audit \\
11 & Stakeholder-Centric Design & 90-120m & Advanced+ &
Design thinking, user research, iteration &
{[}beads{]} Jewelry Holder \\
\end{longtable}
}

\subsection*{4 Reference Appendices}\label{docs__pandoc__latex__src__3dmake_all_lessons.md__4-reference-appendices}

Quick links to comprehensive reference materials:

{\def\LTcaptype{none} % do not increment counter
\begin{longtable}[]{@{}
  >{\raggedright\arraybackslash}p{(\linewidth - 6\tabcolsep) * \real{0.5307}}
  >{\raggedright\arraybackslash}p{(\linewidth - 6\tabcolsep) * \real{0.1341}}
  >{\raggedright\arraybackslash}p{(\linewidth - 6\tabcolsep) * \real{0.0782}}
  >{\raggedright\arraybackslash}p{(\linewidth - 6\tabcolsep) * \real{0.2570}}@{}}
\toprule\noalign{}
\begin{minipage}[b]{\linewidth}\raggedright
Appendix
\end{minipage} & \begin{minipage}[b]{\linewidth}\raggedright
Focus
\end{minipage} & \begin{minipage}[b]{\linewidth}\raggedright
Size
\end{minipage} & \begin{minipage}[b]{\linewidth}\raggedright
Use When
\end{minipage} \\
\midrule\noalign{}
\endhead
\bottomrule\noalign{}
\endlastfoot
\hyperref[docs__pandoc__latex__src__3dmake_appendices_complete.md__appendix-b]{A: Comprehensive Slicing Guide}
& PrusaSlicer, Bambu Studio, Cura, OrcaSlicer configuration &
1,500+ lines & Slicing questions, slicer reference \\
\hyperref[docs__pandoc__latex__src__3dmake_appendices_complete.md__appendix-c]{B: Material Properties \& Selection Guide}
& Shrinkage data, print settings, material properties & 1,200+ lines &
Choosing material, troubleshooting prints \\
\hyperref[docs__pandoc__latex__src__3dmake_appendices_complete.md__appendix-d]{C: Tolerance Testing \& Quality Assurance Matrix}
& QA procedures, tolerance validation methods & 1,200+ lines &
Quality verification, measurement techniques \\
\hyperref[docs__pandoc__latex__src__cli_appendices_complete.md__appendix-b]{D: PowerShell Integration for SCAD Workflows}
& Batch processing, automation scripts, workflow integration &
1,100+ lines & Building automation, batch processing \\
\end{longtable}
}

\subsection*{Learning Paths}\label{docs__pandoc__latex__src__3dmake_all_lessons.md__learning-paths-1}

\subsubsection*{Path 1: Complete Mastery (18-22 hours)}\label{docs__pandoc__latex__src__3dmake_all_lessons.md__path-1-complete-mastery-18-22-hours}

-\textgreater{} Lessons 1-11 + All Appendices

Best for: Complete skill development, comprehensive understanding

\subsubsection*{Path 2: Design Focus (12-15 hours)}\label{docs__pandoc__latex__src__3dmake_all_lessons.md__path-2-design-focus-12-15-hours}

-\textgreater{} Lessons 1-3, 6-8, 11 + Appendices A, B, C

Best for: Experienced makers new to programmatic CAD

\subsubsection*{Path 3: Project-Based (14-18 hours)}\label{docs__pandoc__latex__src__3dmake_all_lessons.md__path-3-project-based-14-18-hours}

-\textgreater{} Lessons 1-5 (Foundations) -\textgreater{} 6 (Keycap) -\textgreater{} 7 (Stand) -\textgreater{} 8 (Bins) -\textgreater{} 9 (Automation) -\textgreater{} 10 (Troubleshooting) -\textgreater{} 11 (Leadership)

Best for: Learning through building

\subsubsection*{Path 4: Safety \& Printing (10-12 hours)}\label{docs__pandoc__latex__src__3dmake_all_lessons.md__path-4-safety--printing-10-12-hours}

-\textgreater{} Lessons 1, 2, 5, 6, 10 + Appendices A, B, C

Best for: Focus on practical printing and quality

\subsection*{3dm Command Reference}\label{docs__pandoc__latex__src__3dmake_all_lessons.md__3dm-command-reference}

\subsubsection*{Essential Commands}\label{docs__pandoc__latex__src__3dmake_all_lessons.md__essential-commands}

\begin{lstlisting}[style=Alabaster, language=bash]
# Setup

./3dm setup                 # Initial configuration


# Development

3dm edit-model file.scad    # Open in editor

3dm build src/main.scad     # Generate STL from SCAD


# Inspection

3dm info file.scad      # Text analysis (AI if configured)

3dm preview file.scad       # Generate 2D tactile preview


# Optimization

3dm orient file.scad        # Suggest print orientation


# Production

3dm slice file.scad         # Generate G-code

3dm send build/main.gcode   # Send to printer


# Libraries

3dm lib list                # Show available libraries

3dm lib install BOSL2       # Install a library


\end{lstlisting}

\subsubsection*{Command Chaining}\label{docs__pandoc__latex__src__3dmake_all_lessons.md__command-chaining}

\begin{lstlisting}[style=Alabaster, language=bash]
# Sequential with error handling

3dm build src/main.scad && 3dm slice src/main.scad && echo "Ready to print"


# Loop through files

for f in src/*.scad; do 3dm build "$f" && 3dm slice "$f"; done


\end{lstlisting}

\subsection*{OpenSCAD Quick Reference}\label{docs__pandoc__latex__src__3dmake_all_lessons.md__openscad-quick-reference}

\subsubsection*{Primitives}\label{docs__pandoc__latex__src__3dmake_all_lessons.md__primitives}

\begin{lstlisting}[style=Alabaster, language=openscad]
cube([width, height, depth], center=false);

sphere(r=radius, $fn=32);

cylinder(r=radius, h=height, $fn=32);


\end{lstlisting}

\subsubsection*{Transforms}\label{docs__pandoc__latex__src__3dmake_all_lessons.md__transforms}

\begin{lstlisting}[style=Alabaster, language=openscad]
translate([x, y, z]) { ... }

rotate([x_deg, y_deg, z_deg]) { ... }

scale([x, y, z]) { ... }


\end{lstlisting}

\subsubsection*{Boolean Operations}\label{docs__pandoc__latex__src__3dmake_all_lessons.md__boolean-operations-1}

\begin{lstlisting}[style=Alabaster, language=openscad]
union() { shape1; shape2; }         // Combine

difference() { shape1; shape2; }    // Subtract

intersection() { shape1; shape2; }  // Keep overlap


\end{lstlisting}

\subsubsection*{Modules}\label{docs__pandoc__latex__src__3dmake_all_lessons.md__modules}

\begin{lstlisting}[style=Alabaster, language=openscad]
module my_shape(size) {

  cube([size, size, size]);

}

my_shape(20);   // Call module


\end{lstlisting}

\subsubsection*{Parameters}\label{docs__pandoc__latex__src__3dmake_all_lessons.md__parameters}

\begin{lstlisting}[style=Alabaster, language=openscad]
width = 50;     // mm

height = 30;    // mm

inner = width - 2*wall;


\end{lstlisting}

\subsection*{Projects Reference}\label{docs__pandoc__latex__src__3dmake_all_lessons.md__projects-reference}

\subsubsection*{Project 1: Parametric Keycap (Lesson 6)}\label{docs__pandoc__latex__src__3dmake_all_lessons.md__project-1-parametric-keycap-lesson-6}

Key Parameters:

\begin{lstlisting}[style=Alabaster, language=openscad]
key_size = 18;      // mm

key_height = 12;    // mm

wall = 1.2;         // mm

letter = "A";       // Character


\end{lstlisting}

Variants to Try:

\begin{itemize}
\tightlist
\item
  Small: 12mm, 10mm height
\item
  Medium: 18mm, 12mm height
\item
  Large: 24mm, 14mm height
\end{itemize}

Files:

\begin{itemize}
\tightlist
\item
  Code: Lesson 6 (Keycap section)
\item
  Output: keycap\_X.scad, keycap\_X.stl
\end{itemize}

\subsubsection*{Project 2: Phone Stand (Lesson 7)}\label{docs__pandoc__latex__src__3dmake_all_lessons.md__project-2-phone-stand-lesson-7}

Key Parameters:

\begin{lstlisting}[style=Alabaster, language=openscad]
phone_width = 75;   // mm

base_width = 85;    // mm

angle = 60;         // degrees

lip_height = 15;    // mm


\end{lstlisting}

Configurations:

{\def\LTcaptype{none} % do not increment counter
\begin{longtable}[]{@{}llll@{}}
\toprule\noalign{}
Phone & Width & Angle & Result \\
\midrule\noalign{}
\endhead
\bottomrule\noalign{}
\endlastfoot
iPhone & 60mm & 60 & Portrait viewing \\
iPad & 100mm & 40 & Landscape viewing \\
Tablet & 150mm & 35 & Document viewing \\
\end{longtable}
}

Files:

\begin{itemize}
\tightlist
\item
  Code: Lesson 7 (Phone Stand section)
\item
  Output: stand\_X.stl, stand\_X.gcode
\end{itemize}

\subsubsection*{Project 3: Stackable Bins (Lesson 8)}\label{docs__pandoc__latex__src__3dmake_all_lessons.md__project-3-stackable-bins-lesson-8}

Key Parameters:

\begin{lstlisting}[style=Alabaster, language=openscad]
bin_w = 80;         // width (mm)

bin_d = 120;        // depth (mm)

bin_h = 60;         // height (mm)

wall = 2;           // thickness (mm)

stack_clear = 0.6;  // tolerance (mm) - CRITICAL


\end{lstlisting}

Tolerance Testing:

\begin{lstlisting}[style=Alabaster]
stack_clear = 0.4mm  -> Too tight (hard to stack)

stack_clear = 0.6mm  -> Ideal (smooth fit)

stack_clear = 0.8mm  -> Too loose (unstable)


\end{lstlisting}

Files:

\begin{itemize}
\tightlist
\item
  Code: Lesson 8 (Bins section)
\item
  Output: bin\_*.stl, tolerance\_matrix.md
\end{itemize}

\subsection*{Code Template Library}\label{docs__pandoc__latex__src__3dmake_all_lessons.md__code-template-library}

\subsubsection*{Generic Parametric Part Template}\label{docs__pandoc__latex__src__3dmake_all_lessons.md__generic-parametric-part-template}

\begin{lstlisting}[style=Alabaster, language=openscad]
// ====== PARAMETERS (customize here) ======

param1 = 50;        // mm

param2 = 30;        // mm

param3 = 5;         // mm

$fn = 32;           // Resolution (lower = faster)

// ====== CALCULATED PARAMETERS ======

derived_param = param1 - 2*param3;

// ====== MODULES ======

module my_part() {

  cube([param1, param2, param3]);

}

// ====== MAIN ======

my_part();


\end{lstlisting}

\subsubsection*{Hollow Box Template}\label{docs__pandoc__latex__src__3dmake_all_lessons.md__hollow-box-template}

\begin{lstlisting}[style=Alabaster, language=openscad]
outer_size = 50;

inner_size = 40;

wall = 5;

difference() {

  cube([outer_size, outer_size, outer_size]);

  translate([wall, wall, wall])

    cube([inner_size, inner_size, inner_size]);

}


\end{lstlisting}

\subsubsection*{Batch Build Script Template}\label{docs__pandoc__latex__src__3dmake_all_lessons.md__batch-build-script-template}

\begin{lstlisting}[style=Alabaster, language=bash]
#!/bin/bash

# batch_build.sh


for scad in src/*.scad; do

  name=$(basename "$scad" .scad)

  echo "Building: $name"

  3dm build "$scad" || continue

  cp "build/main.stl" "build/${name}.stl"

done


\end{lstlisting}

\subsection*{Troubleshooting Quick Fixes}\label{docs__pandoc__latex__src__3dmake_all_lessons.md__troubleshooting-quick-fixes}

\subsubsection*{Problem: Model won\textquotesingle t build}\label{docs__pandoc__latex__src__3dmake_all_lessons.md__problem-model-wont-build}

Diagnosis:

\begin{lstlisting}[style=Alabaster, language=bash]
3dm info file.scad

# Look for error messages


\end{lstlisting}

Common Fixes:

\begin{itemize}
\tightlist
\item
  Check syntax (missing semicolons, parentheses)
\item
  Look for non-manifold geometry
\item
  Use \texttt{\$fn=12} for faster test renders
\end{itemize}

\subsubsection*{Problem: Parts don\textquotesingle t fit together}\label{docs__pandoc__latex__src__3dmake_all_lessons.md__problem-parts-dont-fit-together}

Diagnosis:

\begin{itemize}
\tightlist
\item
  Print and test fit
\item
  Measure with calipers
\end{itemize}

Solution:

\begin{itemize}
\tightlist
\item
  Adjust \texttt{stack\_clear} (smaller = tighter)
\item
  Increase \texttt{wall} thickness
\item
  Test with tolerance matrix
\end{itemize}

\subsubsection*{Problem: Embossed text looks bad}\label{docs__pandoc__latex__src__3dmake_all_lessons.md__problem-embossed-text-looks-bad}

Diagnosis:

\begin{itemize}
\tightlist
\item
  Check preview in slicer
\item
  Use \texttt{3dm\ preview} for tactile version
\end{itemize}

Solution:

\begin{itemize}
\tightlist
\item
  Increase \texttt{letter\_raise} (deeper emboss)
\item
  Use larger \texttt{\$fn} in \texttt{text()}
\item
  Simplify character or use different size
\end{itemize}

\subsubsection*{Problem: Print fails}\label{docs__pandoc__latex__src__3dmake_all_lessons.md__problem-print-fails}

Diagnosis:

\begin{itemize}
\tightlist
\item
  Check slicer layer preview
\item
  Verify bed adhesion and temperature
\end{itemize}

Solution:

\begin{itemize}
\tightlist
\item
  Check pre-print checklist (Lesson 5)
\item
  Adjust print temperature
\item
  Verify bed is level and clean
\end{itemize}

\subsection*{Assessment Checklist}\label{docs__pandoc__latex__src__3dmake_all_lessons.md__assessment-checklist}

\subsubsection*{Lesson Completion Criteria}\label{docs__pandoc__latex__src__3dmake_all_lessons.md__lesson-completion-criteria}

\begin{itemize}
\tightlist
\item[$\square$]
  Watched/read entire lesson
\item[$\square$]
  Completed all step-by-step tasks
\item[$\square$]
  Reached all checkpoints
\item[$\square$]
  Answered all quiz questions (self-assessed)
\item[$\square$]
  Attempted at least 3 extension problems
\item[$\square$]
  Documented findings
\end{itemize}

\subsubsection*{Project Completion Criteria}\label{docs__pandoc__latex__src__3dmake_all_lessons.md__project-completion-criteria}

\begin{itemize}
\tightlist
\item[$\square$]
  Code builds without errors
\item[$\square$]
  All parameters functional
\item[$\square$]
  STL generated and inspected
\item[$\square$]
  Measurements documented
\item[$\square$]
  Assembly tested (if multi-part)
\item[$\square$]
  README or documentation included
\end{itemize}

\subsubsection*{Quality Standards}\label{docs__pandoc__latex__src__3dmake_all_lessons.md__quality-standards}

\begin{itemize}
\tightlist
\item[$\square$]
  Code is well-commented
\item[$\square$]
  Parameters have clear names and units
\item[$\square$]
  Modules are reusable
\item[$\square$]
  Design follows DRY principle
\item[$\square$]
  Documentation is complete
\end{itemize}

\subsection*{Resources \& Links}\label{docs__pandoc__latex__src__3dmake_all_lessons.md__resources--links}

\subsubsection*{Official Docs}\label{docs__pandoc__latex__src__3dmake_all_lessons.md__official-docs}

\begin{itemize}
\tightlist
\item
  \href{https://en.wikibooks.org/wiki/OpenSCAD_User_Manual}{OpenSCAD Manual}
\item
  \href{https://github.com/tdeck/3dmake}{3DMake GitHub}
\item
  \href{https://github.com/revarbat/BOSL2}{BOSL2 Library}
\end{itemize}

\subsubsection*{Tutorials}\label{docs__pandoc__latex__src__3dmake_all_lessons.md__tutorials}

\begin{itemize}
\tightlist
\item
  \href{https://en.wikibooks.org/wiki/OpenSCAD_User_Manual/The_OpenSCAD_Language}{OpenSCAD Basics}
\item
  \href{https://www.prusa3d.com/support/}{3D Printing Guide}
\end{itemize}

\subsubsection*{Community}\label{docs__pandoc__latex__src__3dmake_all_lessons.md__community}

\begin{itemize}
\tightlist
\item
  \href{https://forum.openscad.org/}{OpenSCAD Forums}
\item
  \href{https://www.reddit.com/r/3Dprinting/}{r/3Dprinting}
\end{itemize}

\subsection*{Tips \& Tricks}\label{docs__pandoc__latex__src__3dmake_all_lessons.md__tips--tricks}

\subsubsection*{Debugging}\label{docs__pandoc__latex__src__3dmake_all_lessons.md__debugging}

Tip

Lower \texttt{\$fn} to \textbf{8--12 during development} for fast renders, then raise it to 32--64 for final export. High \texttt{\$fn} values (especially with Minkowski) can make renders take minutes.

\begin{itemize}
\tightlist
\item
  Lower \texttt{\$fn} to 8-12 for faster renders during development
\item
  Use \texttt{3dm\ info} frequently to catch issues early
\item
  Test components individually before assembling
\item
  Generate \texttt{3dm\ preview} for 2D tactile verification
\end{itemize}

\subsubsection*{Design}\label{docs__pandoc__latex__src__3dmake_all_lessons.md__design}

\begin{itemize}
\tightlist
\item
  Keep parameters at the top of file for easy modification
\item
  Use descriptive names (not \texttt{w}, use \texttt{width})
\item
  Include units in comments
\item
  Document parameter ranges (e.g., \texttt{//\ 0-100\ mm})
\end{itemize}

\subsubsection*{Organization}\label{docs__pandoc__latex__src__3dmake_all_lessons.md__organization}

\begin{itemize}
\tightlist
\item
  Use \texttt{src/} for SCAD files, \texttt{lib/} for modules, \texttt{build/} for outputs
\item
  Create variants by copying files and renaming
\item
  Use bash scripts for batch operations
\item
  Archive successful builds with timestamps
\end{itemize}

\subsubsection*{Accessibility}\label{docs__pandoc__latex__src__3dmake_all_lessons.md__accessibility}

Important

For accessible design, always use \texttt{3dm\ info} to verify non-visual usability and \texttt{3dm\ preview} to generate a tactile inspection image. Document all measurements clearly and test assembly without visual guidance.

\begin{itemize}
\tightlist
\item
  Always use \texttt{3dm\ info} to verify non-visual usability
\item
  Generate \texttt{3dm\ preview} for tactile inspection
\item
  Document measurements clearly
\item
  Test assembly without visual guidance
\end{itemize}

\subsection*{Glossary}\label{docs__pandoc__latex__src__3dmake_all_lessons.md__glossary}

{\def\LTcaptype{none} % do not increment counter
\begin{longtable}[]{@{}
  >{\raggedright\arraybackslash}p{(\linewidth - 2\tabcolsep) * \real{0.1446}}
  >{\raggedright\arraybackslash}p{(\linewidth - 2\tabcolsep) * \real{0.8554}}@{}}
\toprule\noalign{}
\begin{minipage}[b]{\linewidth}\raggedright
Term
\end{minipage} & \begin{minipage}[b]{\linewidth}\raggedright
Definition
\end{minipage} \\
\midrule\noalign{}
\endhead
\bottomrule\noalign{}
\endlastfoot
CSG &
Constructive Solid Geometry - combining shapes using union/difference \\
Manifold & Water-tight geometry with clear inside/outside \\
Parametric & Driven by variables; changing parameters updates design \\
Tolerance & Acceptable variation in dimensions \\
Stack-up & Cumulative error from multiple tolerances \\
Module & Reusable code block in OpenSCAD \\
\$fn & Resolution parameter (higher = more detail but slower) \\
G-code & Machine instructions for 3D printer \\
STL & 3D model file format for printing \\
\end{longtable}
}

\subsection*{Quick Answers}\label{docs__pandoc__latex__src__3dmake_all_lessons.md__quick-answers}

Q: Where do I put my SCAD files?\\
A: In the \texttt{src/} folder of your 3dMake project

Q: How do I test if my design will fit?\\
A: Use the tolerance testing matrix; print variants with different parameters

Q: What should I measure after printing?\\
A: Critical dimensions and compare to design specifications

Tip

\textbf{Non-manifold fix:} Use the 0.001 offset rule --- \texttt{translate({[}0,\ 0,\ -0.001{]})} before subtracting and add 0.002 to the cutting shape\textquotesingle s height.

Q: How do I fix non-manifold geometry?\\
A: Use the 0.001 offset rule: \texttt{translate({[}0,\ 0,\ 0.001{]})} before subtracting

Q: Can I combine multiple SCAD files?\\
A: Yes, use \texttt{include\ \textless{}path/to/file.scad\textgreater{}} or \texttt{use\ \textless{}path/to/file.scad\textgreater{}}

Q: How do I make designs accessible?\\
A: Use \texttt{3dm\ info} and \texttt{3dm\ preview} and include written measurements/descriptions

\newpage

\chapter{BackMatter}\label{docs__pandoc__latex__src__backmatter__backmatter.md__backmatter}

This BackMatter section collects supplementary materials, appendices, and reference content for the book.

\begin{itemize}
\tightlist
\item
  \textbf{Advanced Problem Sets:} \hyperref[docs__pandoc__latex__src__advanced_study_guide.md__advanced-study-guide]{Advanced Problem Sets} --- Extended practice and challenge exercises for instructors and students.
\item
  \textbf{Command Line Appendices:} \hyperref[docs__pandoc__latex__src__cli_appendices_complete.md__command-line-appendices--complete-reference]{Command Line Appendices} --- Reference material and in-depth command-line guides.
\item
  \textbf{3D Make \& OpenSCAD Appendices:} \hyperref[docs__pandoc__latex__src__3dmake_appendices_complete.md__3dmake-and-openscad-appendices--complete-reference]{3D Make and OpenSCAD Appendices} --- Additional tutorials, patterns, and design references.
\item
  \textbf{Technical References \& Troubleshooting:} \hyperref[docs__pandoc__latex__src__backmatter_complete.md__back-matter--complete-reference]{Technical References and Troubleshooting} --- Troubleshooting notes, technical tips, and reference tables.
\item
  \textbf{About the Author:} \hyperref[docs__pandoc__latex__src__contributors.md__contributors]{About the Author} --- Contributor and author information.
\end{itemize}

\section{Advanced Study Guide}\label{docs__pandoc__latex__src__advanced_study_guide.md__advanced-study-guide}

\subsection*{CLI Foundations \& 3dMake Certification}\label{docs__pandoc__latex__src__advanced_study_guide.md__cli-foundations--3dmake-certification}

\textbf{Format:} Each lesson contains 10 questions.

\begin{itemize}
\tightlist
\item
  \textbf{Questions 1--5 (Fix the Code):} Each block contains one or more deliberate errors. Identify every error and write a corrected version.
\item
  \textbf{Questions 6--10 (Comment the Code):} Each block is correct but uncommented. Add a meaningful inline or block comment to every significant line or section.
\end{itemize}

\begin{center}\rule{0.5\linewidth}{0.5pt}\end{center}

\section*{PART 1 --- PowerShell}\label{docs__pandoc__latex__src__advanced_study_guide.md__part-1--powershell}

\begin{center}\rule{0.5\linewidth}{0.5pt}\end{center}

\subsection*{PowerShell --- Lesson Pre: Your First Terminal}\label{docs__pandoc__latex__src__advanced_study_guide.md__powershell--lesson-pre-your-first-terminal}

\subsubsection*{Fix the Code (Questions 1--5)}\label{docs__pandoc__latex__src__advanced_study_guide.md__fix-the-code-questions-15}

\textbf{Q1.} This script is supposed to show the current directory and list files, but it uses the wrong commands.

\begin{lstlisting}[style=Alabaster, language=powershell]
location
list -names

\end{lstlisting}

\textbf{Q2.} The user wants to create a file with content, then read it back. Two commands are wrong.

\begin{lstlisting}[style=Alabaster, language=powershell]
print "Hello PowerShell" > notes.txt
read notes.txt

\end{lstlisting}

\textbf{Q3.} The user wants to open a file in Notepad. The path separator and command name are wrong.

\begin{lstlisting}[style=Alabaster, language=powershell]
notepad notes/txt

\end{lstlisting}

\textbf{Q4.} This snippet is supposed to stop a running command, clear the screen, and go home. One keyboard shortcut description is wrong, and one command is wrong.

\begin{lstlisting}[style=Alabaster, language=powershell]
# Press Ctrl+X to stop a running command
# Press Ctrl+L to clear the screen
# Type the following and press Enter to go home:
home

\end{lstlisting}

\textbf{Q5.} This script should save the file listing to a file, then open it. The redirection operator and the open command are both wrong.

\begin{lstlisting}[style=Alabaster, language=powershell]
ls -n >> directory_list.txt
open directory_list.txt

\end{lstlisting}

\begin{center}\rule{0.5\linewidth}{0.5pt}\end{center}

\subsubsection*{Comment the Code (Questions 6--10)}\label{docs__pandoc__latex__src__advanced_study_guide.md__comment-the-code-questions-610}

\textbf{Q6.} Add a comment to every line explaining what it does.

\begin{lstlisting}[style=Alabaster, language=powershell]
pwd
ls -n
cd Documents
pwd
ls -n

\end{lstlisting}

\textbf{Q7.} Add comments explaining what each part of this file-creation workflow does.

\begin{lstlisting}[style=Alabaster, language=powershell]
echo "Test content" > test.txt
cat test.txt
echo "Line 2" >> test.txt
cat test.txt

\end{lstlisting}

\textbf{Q8.} Add comments explaining each step of this tab-completion demonstration script.

\begin{lstlisting}[style=Alabaster, language=powershell]
cd ~
ls -n
cd Doc

\end{lstlisting}

\emph{(Note: the third line is intentionally incomplete --- the user would press Tab here.)}

\textbf{Q9.} Add comments to this output-redirection workflow explaining why each step is needed.

\begin{lstlisting}[style=Alabaster, language=powershell]
ls -n > big_list.txt
notepad.exe big_list.txt
ls -n | findstr "scad"

\end{lstlisting}

\textbf{Q10.} Add comments explaining each keyboard shortcut and what it accomplishes for screen reader users.

\begin{lstlisting}[style=Alabaster, language=powershell]
# Up Arrow
# Down Arrow
# Tab
# Ctrl+C
# Ctrl+L

\end{lstlisting}

\begin{center}\rule{0.5\linewidth}{0.5pt}\end{center}

\subsection*{PowerShell --- Lesson 0: Getting Started --- Paths and the Shell}\label{docs__pandoc__latex__src__advanced_study_guide.md__powershell--lesson-0-getting-started--paths-and-the-shell}

\subsubsection*{Fix the Code (Questions 1--5)}\label{docs__pandoc__latex__src__advanced_study_guide.md__fix-the-code-questions-15-1}

\textbf{Q1.} This script navigates to a nested folder but uses wrong path notation for Windows PowerShell.

\begin{lstlisting}[style=Alabaster, language=powershell]
cd ~/Documents/Projects
cd ../
pwd

\end{lstlisting}

\emph{(Hint: the second line uses a Linux-style relative path --- valid, but the first line also has an issue on some systems. Focus on the separator convention and the go-home shortcut used.)}

\textbf{Q2.} The user wants to navigate to a folder whose name contains a space. The command is missing required syntax.

\begin{lstlisting}[style=Alabaster, language=powershell]
cd My Projects

\end{lstlisting}

\textbf{Q3.} This script checks the current path, then navigates using an absolute path that is wrongly constructed.

\begin{lstlisting}[style=Alabaster, language=powershell]
pwd
cd C:/Users/Student/Documents

\end{lstlisting}

\textbf{Q4.} The script is supposed to list only file names (not full details). The flag is wrong.

\begin{lstlisting}[style=Alabaster, language=powershell]
ls -l

\end{lstlisting}

\textbf{Q5.} The 3DMake project setup sequence has a command error --- the wrong operator is used to initialize the project.

\begin{lstlisting}[style=Alabaster, language=powershell]
mkdir ~/Documents/3d_projects/my_cube
cd ~/Documents/3d_projects/my_cube
3dm create

\end{lstlisting}

\begin{center}\rule{0.5\linewidth}{0.5pt}\end{center}

\subsubsection*{Comment the Code (Questions 6--10)}\label{docs__pandoc__latex__src__advanced_study_guide.md__comment-the-code-questions-610-1}

\textbf{Q6.} Add comments explaining what each path shortcut means and when you would use it.

\begin{lstlisting}[style=Alabaster, language=powershell]
cd ~
cd ./subfolder
cd ../
cd $HOME

\end{lstlisting}

\textbf{Q7.} Add comments explaining the purpose of each line in this path-oriented script.

\begin{lstlisting}[style=Alabaster, language=powershell]
pwd
ls -n
cd Documents
pwd
cd ~
pwd

\end{lstlisting}

\textbf{Q8.} Add comments explaining what tab completion is doing at each step.

\begin{lstlisting}[style=Alabaster, language=powershell]
cd ~/D
cd ~/Documents/3d_p
ls -n

\end{lstlisting}

\emph{(Note: first two lines show partial paths where Tab would be pressed.)}

\textbf{Q9.} Add comments to this 3DMake quick-start sequence explaining each command\textquotesingle s role.

\begin{lstlisting}[style=Alabaster, language=powershell]
cd ~/Downloads/3dmake
./3dm setup
mkdir ~/Documents/3d_projects
mkdir ~/Documents/3d_projects/my_cube
cd ~/Documents/3d_projects/my_cube
3dm new
ls -n

\end{lstlisting}

\textbf{Q10.} Add comments explaining each line of this prompt-anatomy reference.

\begin{lstlisting}[style=Alabaster, language=powershell]
# PS C:\Users\Student\Documents>
pwd
ls -n
cd ..
pwd

\end{lstlisting}

\begin{center}\rule{0.5\linewidth}{0.5pt}\end{center}

\subsection*{PowerShell --- Lesson 1: Navigation}\label{docs__pandoc__latex__src__advanced_study_guide.md__powershell--lesson-1-navigation}

\subsubsection*{Fix the Code (Questions 1--5)}\label{docs__pandoc__latex__src__advanced_study_guide.md__fix-the-code-questions-15-2}

\textbf{Q1.} This script tries to list only \texttt{.scad} files but uses the wrong filter syntax for PowerShell.

\begin{lstlisting}[style=Alabaster, language=powershell]
ls -n *.scad
ls -n | findstr scad

\end{lstlisting}

\emph{(The first line is correct --- identify the error in the second line.)}

\textbf{Q2.} The script should navigate deep into a project structure and then return home in one command. The return command is wrong.

\begin{lstlisting}[style=Alabaster, language=powershell]
cd ~/Documents/3d_projects/my_cube/src
cd /home

\end{lstlisting}

\textbf{Q3.} This loop is supposed to list all files but it uses a nonexistent PowerShell command.

\begin{lstlisting}[style=Alabaster, language=powershell]
foreach ($f in ls -n *.txt) {
    print $f
}

\end{lstlisting}

\textbf{Q4.} The user wants to find all \texttt{.stl} files and save the results. The wildcard and redirection are both wrong.

\begin{lstlisting}[style=Alabaster, language=powershell]
ls -n build/%.stl > results.txt

\end{lstlisting}

\textbf{Q5.} This script navigates to a project and verifies the build output, but the path separator and verify command are wrong.

\begin{lstlisting}[style=Alabaster, language=powershell]
cd ~/Documents\3d_projects\my_cube
ls -n build/
Test-Exist build/main.stl

\end{lstlisting}

\begin{center}\rule{0.5\linewidth}{0.5pt}\end{center}

\subsubsection*{Comment the Code (Questions 6--10)}\label{docs__pandoc__latex__src__advanced_study_guide.md__comment-the-code-questions-610-2}

\textbf{Q6.} Add comments explaining each \texttt{ls} flag and what it produces.

\begin{lstlisting}[style=Alabaster, language=powershell]
ls -n
ls -n *.scad
ls -n | findstr ".stl"
ls -n -af

\end{lstlisting}

\textbf{Q7.} Add comments to this multi-level navigation script.

\begin{lstlisting}[style=Alabaster, language=powershell]
cd ~/Documents/3d_projects
pwd
cd my_cube
pwd
cd src
pwd
cd ../../..
pwd

\end{lstlisting}

\textbf{Q8.} Add comments explaining why each command exists in this verification workflow.

\begin{lstlisting}[style=Alabaster, language=powershell]
cd ~/Documents/3d_projects/my_cube
pwd
ls -n
ls -n src/
ls -n build/

\end{lstlisting}

\textbf{Q9.} Add comments explaining the wildcard patterns used in each line.

\begin{lstlisting}[style=Alabaster, language=powershell]
ls -n *.scad
ls -n part?.stl
ls -n *_final*
ls -n project_[1-3].scad

\end{lstlisting}

\textbf{Q10.} Add comments to this screen-reader-friendly output-saving workflow.

\begin{lstlisting}[style=Alabaster, language=powershell]
ls -n > file_list.txt
notepad.exe file_list.txt
ls -n | findstr "scad"
ls -n | findstr "scad" > scad_list.txt
notepad.exe scad_list.txt

\end{lstlisting}

\begin{center}\rule{0.5\linewidth}{0.5pt}\end{center}

\subsection*{PowerShell --- Lesson 2: File and Folder Manipulation}\label{docs__pandoc__latex__src__advanced_study_guide.md__powershell--lesson-2-file-and-folder-manipulation}

\subsubsection*{Fix the Code (Questions 1--5)}\label{docs__pandoc__latex__src__advanced_study_guide.md__fix-the-code-questions-15-3}

\textbf{Q1.} This script tries to copy a file but uses the wrong PowerShell cmdlet name.

\begin{lstlisting}[style=Alabaster, language=powershell]
cp source.txt destination.txt

\end{lstlisting}

\textbf{Q2.} The user wants to delete a file but mistakenly uses the Linux command.

\begin{lstlisting}[style=Alabaster, language=powershell]
rm old_draft.scad

\end{lstlisting}

\textbf{Q3.} This script creates a backup of an entire project folder, but the flag for recursive copy is missing.

\begin{lstlisting}[style=Alabaster, language=powershell]
Copy-Item my_cube backup_cube

\end{lstlisting}

\textbf{Q4.} The folder-creation command is right but the subsequent verification uses the wrong flag for listing directories only.

\begin{lstlisting}[style=Alabaster, language=powershell]
mkdir Models
mkdir Prints
mkdir Docs
ls -n -files

\end{lstlisting}

\textbf{Q5.} This safe-deletion sequence is out of order --- verification happens after deletion instead of before.

\begin{lstlisting}[style=Alabaster, language=powershell]
Remove-Item old_test.scad
ls -n *.scad

\end{lstlisting}

\begin{center}\rule{0.5\linewidth}{0.5pt}\end{center}

\subsubsection*{Comment the Code (Questions 6--10)}\label{docs__pandoc__latex__src__advanced_study_guide.md__comment-the-code-questions-610-3}

\textbf{Q6.} Add comments explaining what each file operation does and what to verify afterwards.

\begin{lstlisting}[style=Alabaster, language=powershell]
New-Item -ItemType File -Name "design.scad"
Copy-Item design.scad design_backup.scad
Move-Item design.scad archive/design_v1.scad
ls -n

\end{lstlisting}

\textbf{Q7.} Add comments to this folder-structure creation script.

\begin{lstlisting}[style=Alabaster, language=powershell]
mkdir 3DProjects
cd 3DProjects
mkdir Models
mkdir Prints
mkdir Docs
mkdir Backups
ls -n

\end{lstlisting}

\textbf{Q8.} Add comments explaining the safe backup-before-edit pattern.

\begin{lstlisting}[style=Alabaster, language=powershell]
Copy-Item project.scad project_backup.scad
ls -n *.scad
notepad.exe project.scad
Copy-Item project.scad project_v2.scad

\end{lstlisting}

\textbf{Q9.} Add comments to this wildcard file-management sequence.

\begin{lstlisting}[style=Alabaster, language=powershell]
ls -n *.tmp
Remove-Item *.tmp
ls -n *.scad
Copy-Item *.scad backup/
ls -n backup/

\end{lstlisting}

\textbf{Q10.} Add comments explaining what each step of this project-archiving script does.

\begin{lstlisting}[style=Alabaster, language=powershell]
$date = Get-Date -Format "yyyyMMdd"
$archiveName = "project_backup_$date"
Copy-Item -Recurse my_cube $archiveName
ls -n

\end{lstlisting}

\begin{center}\rule{0.5\linewidth}{0.5pt}\end{center}

\subsection*{PowerShell --- Lesson 3: Input, Output, and Piping}\label{docs__pandoc__latex__src__advanced_study_guide.md__powershell--lesson-3-input-output-and-piping}

\subsubsection*{Fix the Code (Questions 1--5)}\label{docs__pandoc__latex__src__advanced_study_guide.md__fix-the-code-questions-15-4}

\textbf{Q1.} The user wants to append a second line to a file without destroying the first. The operator is wrong.

\begin{lstlisting}[style=Alabaster, language=powershell]
echo "Line 1" > notes.txt
echo "Line 2" > notes.txt

\end{lstlisting}

\textbf{Q2.} This pipe should filter to only lines containing "error" but uses the wrong PowerShell command.

\begin{lstlisting}[style=Alabaster, language=powershell]
cat build_log.txt | grep "error"

\end{lstlisting}

\textbf{Q3.} The script is supposed to copy file contents to the Windows clipboard but uses the wrong pipe target.

\begin{lstlisting}[style=Alabaster, language=powershell]
cat notes.txt | clipboard

\end{lstlisting}

\textbf{Q4.} This pipe chain should count the number of \texttt{.scad} files, but the counting command is wrong.

\begin{lstlisting}[style=Alabaster, language=powershell]
ls -n *.scad | Measure-Lines

\end{lstlisting}

\textbf{Q5.} The output redirection is supposed to save \texttt{3dm\ info} output for review, but the file path uses a wrong separator and the command name is incorrect.

\begin{lstlisting}[style=Alabaster, language=powershell]
3dm information > build\info_report.txt
notepad.exe build\info_report.txt

\end{lstlisting}

\begin{center}\rule{0.5\linewidth}{0.5pt}\end{center}

\subsubsection*{Comment the Code (Questions 6--10)}\label{docs__pandoc__latex__src__advanced_study_guide.md__comment-the-code-questions-610-4}

\textbf{Q6.} Add comments explaining the difference between \texttt{\textgreater{}} and \texttt{\textgreater{}\textgreater{}} in each line.

\begin{lstlisting}[style=Alabaster, language=powershell]
echo "Project: my_cube" > log.txt
echo "Build started" >> log.txt
echo "Build complete" >> log.txt
cat log.txt

\end{lstlisting}

\textbf{Q7.} Add comments to this pipe chain explaining each stage of the pipeline.

\begin{lstlisting}[style=Alabaster, language=powershell]
ls -n | findstr ".scad" | Out-File scad_files.txt
cat scad_files.txt

\end{lstlisting}

\textbf{Q8.} Add comments explaining why each step of this screen-reader output workflow is used.

\begin{lstlisting}[style=Alabaster, language=powershell]
3dm build
3dm info > model_report.txt
notepad.exe model_report.txt

\end{lstlisting}

\textbf{Q9.} Add comments to this pipe-and-filter sequence used to search build logs.

\begin{lstlisting}[style=Alabaster, language=powershell]
cat build_log.txt | findstr "ERROR"
cat build_log.txt | findstr "ERROR" > errors_only.txt
cat build_log.txt | findstr "ERROR" | clip

\end{lstlisting}

\textbf{Q10.} Add comments to this combined input/output script.

\begin{lstlisting}[style=Alabaster, language=powershell]
echo "width=50" > params.txt
echo "height=30" >> params.txt
echo "wall=2.5" >> params.txt
cat params.txt
cat params.txt | clip

\end{lstlisting}

\begin{center}\rule{0.5\linewidth}{0.5pt}\end{center}

\subsection*{PowerShell --- Lesson 4: Environment Variables, PATH, and Aliases}\label{docs__pandoc__latex__src__advanced_study_guide.md__powershell--lesson-4-environment-variables-path-and-aliases}

\subsubsection*{Fix the Code (Questions 1--5)}\label{docs__pandoc__latex__src__advanced_study_guide.md__fix-the-code-questions-15-5}

\textbf{Q1.} The user tries to create a permanent alias but uses the wrong scope.

\begin{lstlisting}[style=Alabaster, language=powershell]
Set-Alias -Name preview -Value openscad -Scope Global

\end{lstlisting}

\emph{(Hint: for persistence across sessions, the alias must be stored somewhere specific.)}

\textbf{Q2.} This script sets an environment variable but uses the wrong syntax.

\begin{lstlisting}[style=Alabaster, language=powershell]
$env:PROJECT_DIR = ~/Documents/3d_projects

\end{lstlisting}

\textbf{Q3.} The script checks whether \texttt{openscad} is on the PATH, but uses the wrong command.

\begin{lstlisting}[style=Alabaster, language=powershell]
locate openscad

\end{lstlisting}

\textbf{Q4.} The user tries to add a directory to PATH but the concatenation operator is wrong.

\begin{lstlisting}[style=Alabaster, language=powershell]
$env:PATH = $env:PATH : "C:\Tools\openscad"

\end{lstlisting}

\textbf{Q5.} This alias definition has the wrong cmdlet name for creating aliases.

\begin{lstlisting}[style=Alabaster, language=powershell]
New-Variable -Name "build" -Value "3dm build"

\end{lstlisting}

\begin{center}\rule{0.5\linewidth}{0.5pt}\end{center}

\subsubsection*{Comment the Code (Questions 6--10)}\label{docs__pandoc__latex__src__advanced_study_guide.md__comment-the-code-questions-610-5}

\textbf{Q6.} Add comments explaining what each environment variable does and when you would set it.

\begin{lstlisting}[style=Alabaster, language=powershell]
$env:PROJECT_DIR = "$HOME\Documents\3d_projects"
$env:OPENSCAD_PATH = "C:\Program Files\OpenSCAD"
$env:PATH += ";$env:OPENSCAD_PATH"

\end{lstlisting}

\textbf{Q7.} Add comments to this alias-creation and verification sequence.

\begin{lstlisting}[style=Alabaster, language=powershell]
Set-Alias -Name build -Value "3dm"
Set-Alias -Name preview -Value openscad
Get-Alias build
Get-Alias preview

\end{lstlisting}

\textbf{Q8.} Add comments explaining each line of this PATH diagnostic script.

\begin{lstlisting}[style=Alabaster, language=powershell]
$env:PATH -split ";"
Get-Command openscad
Get-Command 3dm
(Get-Command openscad).Source

\end{lstlisting}

\textbf{Q9.} Add comments to this profile-editing sequence.

\begin{lstlisting}[style=Alabaster, language=powershell]
$PROFILE
Test-Path $PROFILE
New-Item -Path $PROFILE -Force
notepad.exe $PROFILE

\end{lstlisting}

\textbf{Q10.} Add comments explaining what each line of this environment-setup block (meant to go in a profile) does.

\begin{lstlisting}[style=Alabaster, language=powershell]
Set-Alias build "3dm build"
$env:PROJECT_DIR = "$HOME\Documents\3d_projects"
function go-projects { Set-Location $env:PROJECT_DIR }
Write-Host "3D printing environment loaded."

\end{lstlisting}

\begin{center}\rule{0.5\linewidth}{0.5pt}\end{center}

\subsection*{PowerShell --- Lesson 5: Profiles and History}\label{docs__pandoc__latex__src__advanced_study_guide.md__powershell--lesson-5-profiles-and-history}

\subsubsection*{Fix the Code (Questions 1--5)}\label{docs__pandoc__latex__src__advanced_study_guide.md__fix-the-code-questions-15-6}

\textbf{Q1.} The command to recall the 5th entry from history is wrong.

\begin{lstlisting}[style=Alabaster, language=powershell]
Get-History | Select -Index 5

\end{lstlisting}

\textbf{Q2.} The user wants to re-execute the last command from history. The command is wrong.

\begin{lstlisting}[style=Alabaster, language=powershell]
Invoke-Command -History -1

\end{lstlisting}

\textbf{Q3.} This script is supposed to search history for all commands containing "build" but the search syntax is wrong.

\begin{lstlisting}[style=Alabaster, language=powershell]
Get-History | Where-Object { $_.Command = "build" }

\end{lstlisting}

\textbf{Q4.} The profile path check uses the wrong variable name.

\begin{lstlisting}[style=Alabaster, language=powershell]
Test-Path $MYPROFILE
notepad.exe $MYPROFILE

\end{lstlisting}

\textbf{Q5.} This history-export script saves to a file but uses the wrong cmdlet for writing text.

\begin{lstlisting}[style=Alabaster, language=powershell]
Get-History | ForEach-Object { $_.CommandLine } | Write-File history_log.txt

\end{lstlisting}

\begin{center}\rule{0.5\linewidth}{0.5pt}\end{center}

\subsubsection*{Comment the Code (Questions 6--10)}\label{docs__pandoc__latex__src__advanced_study_guide.md__comment-the-code-questions-610-6}

\textbf{Q6.} Add comments to this history-recall workflow.

\begin{lstlisting}[style=Alabaster, language=powershell]
Get-History
Get-History | Select-Object -Last 10
Invoke-History 15
Get-History | Where-Object { $_.CommandLine -like "*build*" }

\end{lstlisting}

\textbf{Q7.} Add comments to this profile-loading sequence.

\begin{lstlisting}[style=Alabaster, language=powershell]
$PROFILE
Test-Path $PROFILE
. $PROFILE
Get-Alias build

\end{lstlisting}

\textbf{Q8.} Add comments explaining this screen-reader troubleshooting sequence.

\begin{lstlisting}[style=Alabaster, language=powershell]
3dm build > build_output.txt 2>&1
notepad.exe build_output.txt
3dm info > info_output.txt
notepad.exe info_output.txt

\end{lstlisting}

\textbf{Q9.} Add comments to this combined history-and-profile management script.

\begin{lstlisting}[style=Alabaster, language=powershell]
Get-History | Export-Csv -Path history_backup.csv -NoTypeInformation
notepad.exe $PROFILE
. $PROFILE
Get-History | Select-Object -Last 5

\end{lstlisting}

\textbf{Q10.} Add comments explaining each function in this profile snippet.

\begin{lstlisting}[style=Alabaster, language=powershell]
function build-project {
    Set-Location $env:PROJECT_DIR
    3dm build
    3dm info > build_report.txt
    notepad.exe build_report.txt
}

function new-project($name) {
    mkdir "$env:PROJECT_DIR\$name"
    Set-Location "$env:PROJECT_DIR\$name"
    3dm new
}

\end{lstlisting}

\begin{center}\rule{0.5\linewidth}{0.5pt}\end{center}

\subsection*{PowerShell --- Lesson 6: Scripts, Functions, and Professional Workflows}\label{docs__pandoc__latex__src__advanced_study_guide.md__powershell--lesson-6-scripts-functions-and-professional-workflows}

\subsubsection*{Fix the Code (Questions 1--5)}\label{docs__pandoc__latex__src__advanced_study_guide.md__fix-the-code-questions-15-7}

\textbf{Q1.} This script should loop over all \texttt{.scad} files and print their names, but the loop variable and iteration syntax are wrong.

\begin{lstlisting}[style=Alabaster, language=powershell]
$files = ls -n *.scad
for $file in $files {
    Write-Output $file
}

\end{lstlisting}

\textbf{Q2.} The function definition is missing the \texttt{param} block syntax, making parameters inaccessible.

\begin{lstlisting}[style=Alabaster, language=powershell]
function Process-Project {
    $projectPath = $args[0]
    Set-Location $projectPath
    3dm build
}

\end{lstlisting}

\emph{(Hint: \texttt{\$args} works but is fragile --- the function should declare a named parameter.)}

\textbf{Q3.} This try/catch block will never catch anything because the error action preference is wrong.

\begin{lstlisting}[style=Alabaster, language=powershell]
$ErrorActionPreference = "Continue"
try {
    3dm build
    if ($LASTEXITCODE -ne 0) { throw "Build failed" }
} catch {
    Write-Output "Error: $_"
}

\end{lstlisting}

\textbf{Q4.} The batch processing script creates output filenames but uses the wrong string interpolation.

\begin{lstlisting}[style=Alabaster, language=powershell]
$widths = @(60, 80, 100)
foreach ($w in $widths) {
    $outFile = "build/bracket_w" + w + ".stl"
    openscad -D "width=$w" -o $outFile src/main.scad
}

\end{lstlisting}

\textbf{Q5.} The execution policy override is set at the wrong scope --- it changes the machine-wide policy instead of just the current process.

\begin{lstlisting}[style=Alabaster, language=powershell]
Set-ExecutionPolicy -ExecutionPolicy Bypass -Scope LocalMachine
.\build-all.ps1

\end{lstlisting}

\begin{center}\rule{0.5\linewidth}{0.5pt}\end{center}

\subsubsection*{Comment the Code (Questions 6--10)}\label{docs__pandoc__latex__src__advanced_study_guide.md__comment-the-code-questions-610-7}

\textbf{Q6.} Add comments to this complete batch-build script explaining every significant block.

\begin{lstlisting}[style=Alabaster, language=powershell]
param(
    [string]$ProjectsDir = "$HOME\Documents\3d_projects"
)

$timestamp = Get-Date -Format "yyyyMMdd_HHmmss"
$logFile = "build_$timestamp.log"

Get-ChildItem -Path $ProjectsDir -Directory | ForEach-Object {
    $project = $_.FullName
    $name = $_.Name
    Write-Host "Processing: $name"
    Set-Location $project
    $output = & 3dm build slice 2>&1
    if ($LASTEXITCODE -eq 0) {
        Write-Host "  OK"
        & 3dm print
    } else {
        Write-Host "  FAILED"
        $output | Out-File "build_error.txt"
    }
}
Write-Host "Batch complete."

\end{lstlisting}

\textbf{Q7.} Add comments to this parameterized function.

\begin{lstlisting}[style=Alabaster, language=powershell]
function Build-Variant {
    param(
        [string]$ScadFile,
        [int]$Width,
        [double]$Wall,
        [string]$OutputDir = "build/variants"
    )
    $name = "model_w${Width}_wall${Wall}"
    $out  = "$OutputDir\${name}.stl"
    openscad -D "width=$Width" -D "wall=$Wall" -o $out $ScadFile
    if ($LASTEXITCODE -eq 0) {
        Write-Host "Built: $out"
    } else {
        Write-Host "FAILED: $name"
    }
}

\end{lstlisting}

\textbf{Q8.} Add comments to this error-handling and logging block.

\begin{lstlisting}[style=Alabaster, language=powershell]
$ErrorActionPreference = "Stop"
try {
    3dm build
    if ($LASTEXITCODE -ne 0) { throw "OpenSCAD returned $LASTEXITCODE" }
    $stlSize = (Get-Item "build\main.stl").Length
    if ($stlSize -lt 1000) { Write-Warning "STL suspiciously small: $stlSize bytes" }
    Write-Host "Build OK: $stlSize bytes"
} catch {
    Write-Host "Build failed: $_"
    exit 1
}

\end{lstlisting}

\textbf{Q9.} Add comments to this nested loop variant generator.

\begin{lstlisting}[style=Alabaster, language=powershell]
$widths = @(60, 80, 100, 120)
$walls  = @(1.2, 1.6, 2.0)
$ok = 0; $fail = 0

foreach ($w in $widths) {
    foreach ($wt in $walls) {
        $name = "bracket_w${w}_wall${wt}"
        $out  = "build\variants\${name}.stl"
        openscad -D "width=$w" -D "wall=$wt" -o $out src\main.scad 2>$null
        if ($LASTEXITCODE -eq 0) { $ok++ } else { $fail++ }
    }
}
Write-Host "Done: $ok succeeded, $fail failed"

\end{lstlisting}

\textbf{Q10.} Add comments to this complete project-setup function.

\begin{lstlisting}[style=Alabaster, language=powershell]
function New-PrintProject {
    param([string]$Name)
    $base = "$env:PROJECT_DIR\$Name"
    if (Test-Path $base) {
        Write-Warning "Project already exists: $base"
        return
    }
    New-Item -ItemType Directory -Path $base | Out-Null
    Set-Location $base
    3dm new
    Write-Host "Project created: $base"
    ls -n
}

\end{lstlisting}

\begin{center}\rule{0.5\linewidth}{0.5pt}\end{center}

\section*{PART 2 --- Windows Command Prompt (CMD)}\label{docs__pandoc__latex__src__advanced_study_guide.md__part-2--windows-command-prompt-cmd}

\begin{center}\rule{0.5\linewidth}{0.5pt}\end{center}

\subsection*{CMD --- Lesson Pre: Your First Terminal}\label{docs__pandoc__latex__src__advanced_study_guide.md__cmd--lesson-pre-your-first-terminal}

\subsubsection*{Fix the Code (Questions 1--5)}\label{docs__pandoc__latex__src__advanced_study_guide.md__fix-the-code-questions-15-8}

\textbf{Q1.} These two commands should show the current location and list files, but both commands are wrong for CMD.

\begin{lstlisting}[style=Alabaster, language=cmd]
pwd
ls

\end{lstlisting}

\textbf{Q2.} The user wants to create a file with content and then read it back. Both commands are wrong for CMD.

\begin{lstlisting}[style=Alabaster, language=cmd]
echo "Hello CMD" >> hello.txt
cat hello.txt

\end{lstlisting}

\textbf{Q3.} The user tries to open a file in Notepad but uses the wrong separator.

\begin{lstlisting}[style=Alabaster, language=cmd]
notepad hello/txt

\end{lstlisting}

\textbf{Q4.} The user wants to go up one level and confirm the new location. The go-up sequence is correct but the confirmation command is wrong.

\begin{lstlisting}[style=Alabaster, language=cmd]
cd ..
pwd

\end{lstlisting}

\textbf{Q5.} The long-output management command is wrong --- \texttt{less} is not available in CMD.

\begin{lstlisting}[style=Alabaster, language=cmd]
dir /B | less

\end{lstlisting}

\begin{center}\rule{0.5\linewidth}{0.5pt}\end{center}

\subsubsection*{Comment the Code (Questions 6--10)}\label{docs__pandoc__latex__src__advanced_study_guide.md__comment-the-code-questions-610-8}

\textbf{Q6.} Add comments to every line explaining what it does in CMD.

\begin{lstlisting}[style=Alabaster, language=cmd]
cd
dir /B
cd Documents
cd
dir /B

\end{lstlisting}

\textbf{Q7.} Add comments to this file-creation and reading workflow.

\begin{lstlisting}[style=Alabaster, language=cmd]
echo Hello, Command Prompt! > hello.txt
type hello.txt
notepad.exe hello.txt

\end{lstlisting}

\textbf{Q8.} Add comments explaining each step of this output-redirection workflow.

\begin{lstlisting}[style=Alabaster, language=cmd]
dir /B > list.txt
notepad.exe list.txt
dir /B | findstr "scad"

\end{lstlisting}

\textbf{Q9.} Add comments explaining what each keyboard shortcut does for screen reader users.

\begin{lstlisting}[style=Alabaster, language=cmd]
REM Up Arrow
REM Down Arrow
REM Tab
REM Ctrl+C
REM cls (type and press Enter)

\end{lstlisting}

\textbf{Q10.} Add comments to this tab-completion demonstration.

\begin{lstlisting}[style=Alabaster, language=cmd]
cd Doc
dir /B
cd ..
cd Down

\end{lstlisting}

\emph{(Note: lines 1 and 4 show partial paths where Tab would be pressed.)}

\begin{center}\rule{0.5\linewidth}{0.5pt}\end{center}

\subsection*{CMD --- Lesson 0: Getting Started}\label{docs__pandoc__latex__src__advanced_study_guide.md__cmd--lesson-0-getting-started}

\subsubsection*{Fix the Code (Questions 1--5)}\label{docs__pandoc__latex__src__advanced_study_guide.md__fix-the-code-questions-15-9}

\textbf{Q1.} The path to navigate to Documents uses the wrong separator for CMD.

\begin{lstlisting}[style=Alabaster, language=cmd]
cd C:/Users/Student/Documents

\end{lstlisting}

\textbf{Q2.} The user tries to list files and filter for \texttt{.scad} files using a Linux command.

\begin{lstlisting}[style=Alabaster, language=cmd]
dir /B | grep "scad"

\end{lstlisting}

\textbf{Q3.} The script tries to return to the home directory using the wrong variable.

\begin{lstlisting}[style=Alabaster, language=cmd]
cd $HOME

\end{lstlisting}

\textbf{Q4.} The 3DMake project initialization uses the wrong command name.

\begin{lstlisting}[style=Alabaster, language=cmd]
mkdir %USERPROFILE%\Documents\3d_projects\my_cube
cd %USERPROFILE%\Documents\3d_projects\my_cube
3dm init

\end{lstlisting}

\textbf{Q5.} The user tries to list only directory names using a flag that doesn\textquotesingle t exist in CMD.

\begin{lstlisting}[style=Alabaster, language=cmd]
dir /B /dirs

\end{lstlisting}

\begin{center}\rule{0.5\linewidth}{0.5pt}\end{center}

\subsubsection*{Comment the Code (Questions 6--10)}\label{docs__pandoc__latex__src__advanced_study_guide.md__comment-the-code-questions-610-9}

\textbf{Q6.} Add comments explaining the CMD path conventions shown in each line.

\begin{lstlisting}[style=Alabaster, language=cmd]
cd %USERPROFILE%
cd Documents
cd ..\Downloads
cd C:\Users\Student\Documents
cd

\end{lstlisting}

\textbf{Q7.} Add comments to this 3DMake quick-start sequence for CMD.

\begin{lstlisting}[style=Alabaster, language=cmd]
cd %USERPROFILE%\Downloads\3dmake
3dm setup
mkdir %USERPROFILE%\Documents\3d_projects
mkdir %USERPROFILE%\Documents\3d_projects\my_cube
cd %USERPROFILE%\Documents\3d_projects\my_cube
3dm new
dir /B

\end{lstlisting}

\textbf{Q8.} Add comments explaining each filter pattern used in these \texttt{dir} commands.

\begin{lstlisting}[style=Alabaster, language=cmd]
dir /B *.scad
dir /B *.stl
dir /B project?.scad
dir /B *_final*

\end{lstlisting}

\textbf{Q9.} Add comments to this prompt-orientation reference.

\begin{lstlisting}[style=Alabaster, language=cmd]
REM C:\Users\Student>
cd
dir /B
cd ..
cd

\end{lstlisting}

\textbf{Q10.} Add comments explaining every line of this screen-reader-first file listing workflow.

\begin{lstlisting}[style=Alabaster, language=cmd]
dir /B > list.txt
notepad.exe list.txt
dir /B | findstr /R "\.scad$"
dir /B | findstr /R "\.scad$" > scad_list.txt
notepad.exe scad_list.txt

\end{lstlisting}

\begin{center}\rule{0.5\linewidth}{0.5pt}\end{center}

\subsection*{CMD --- Lesson 1: Navigation}\label{docs__pandoc__latex__src__advanced_study_guide.md__cmd--lesson-1-navigation}

\subsubsection*{Fix the Code (Questions 1--5)}\label{docs__pandoc__latex__src__advanced_study_guide.md__fix-the-code-questions-15-10}

\textbf{Q1.} The script should navigate into a nested folder and verify the result, but the verification command is wrong.

\begin{lstlisting}[style=Alabaster, language=cmd]
cd Documents\3d_projects\my_cube
pwd

\end{lstlisting}

\textbf{Q2.} The user tries to navigate to a folder with a space in its name but the command will fail.

\begin{lstlisting}[style=Alabaster, language=cmd]
cd My Projects

\end{lstlisting}

\textbf{Q3.} The wildcard syntax for listing \texttt{.stl} files is wrong for CMD.

\begin{lstlisting}[style=Alabaster, language=cmd]
dir /B /.stl

\end{lstlisting}

\textbf{Q4.} The script tries to search all subdirectories for \texttt{.scad} files but uses the wrong flag.

\begin{lstlisting}[style=Alabaster, language=cmd]
dir /B /R *.scad

\end{lstlisting}

\textbf{Q5.} The \texttt{cd} command at the end should return to the home directory, but the syntax is wrong for CMD.

\begin{lstlisting}[style=Alabaster, language=cmd]
cd Documents
dir /B
cd ~

\end{lstlisting}

\begin{center}\rule{0.5\linewidth}{0.5pt}\end{center}

\subsubsection*{Comment the Code (Questions 6--10)}\label{docs__pandoc__latex__src__advanced_study_guide.md__comment-the-code-questions-610-10}

\textbf{Q6.} Add comments explaining each navigation step and what to expect.

\begin{lstlisting}[style=Alabaster, language=cmd]
cd %USERPROFILE%
cd Documents
cd 3d_projects
dir /B /A:D
cd my_cube
dir /B

\end{lstlisting}

\textbf{Q7.} Add comments to this deep-navigation and backtrack sequence.

\begin{lstlisting}[style=Alabaster, language=cmd]
cd Documents\3d_projects\my_cube\src
cd
dir /B
cd ..\..
cd

\end{lstlisting}

\textbf{Q8.} Add comments explaining what each line does in this project-exploration script.

\begin{lstlisting}[style=Alabaster, language=cmd]
cd %USERPROFILE%\Documents\3d_projects\my_cube
dir /B
dir /B src
dir /B build
dir /B build\*.stl

\end{lstlisting}

\textbf{Q9.} Add comments explaining the wildcard patterns used in each listing.

\begin{lstlisting}[style=Alabaster, language=cmd]
dir /B *.scad
dir /B model?.scad
dir /B *_v[0-9].*
dir /B *backup*

\end{lstlisting}

\textbf{Q10.} Add comments to this screen-reader navigation audit trail script.

\begin{lstlisting}[style=Alabaster, language=cmd]
cd
dir /B > location_check.txt
cd Documents
cd >> location_check.txt
dir /B >> location_check.txt
notepad.exe location_check.txt

\end{lstlisting}

\begin{center}\rule{0.5\linewidth}{0.5pt}\end{center}

\subsection*{CMD --- Lesson 2: File and Folder Manipulation}\label{docs__pandoc__latex__src__advanced_study_guide.md__cmd--lesson-2-file-and-folder-manipulation}

\subsubsection*{Fix the Code (Questions 1--5)}\label{docs__pandoc__latex__src__advanced_study_guide.md__fix-the-code-questions-15-11}

\textbf{Q1.} The user wants to copy an entire folder tree, but the \texttt{/E} and \texttt{/I} flags are missing.

\begin{lstlisting}[style=Alabaster, language=cmd]
xcopy source_folder destination_folder

\end{lstlisting}

\textbf{Q2.} The file deletion command will not ask for confirmation, which is dangerous here. The \texttt{/Q} flag should not be present.

\begin{lstlisting}[style=Alabaster, language=cmd]
del /Q *.scad

\end{lstlisting}

\textbf{Q3.} The user tries to rename a file using a command that doesn\textquotesingle t exist in CMD.

\begin{lstlisting}[style=Alabaster, language=cmd]
rename old_name.scad new_name.scad

\end{lstlisting}

\emph{(Hint: check the exact CMD command name.)}

\textbf{Q4.} This safe-deletion sequence is in the wrong order.

\begin{lstlisting}[style=Alabaster, language=cmd]
del *.tmp
dir /B *.tmp

\end{lstlisting}

\textbf{Q5.} The command to remove a folder and all its contents is missing the required flags.

\begin{lstlisting}[style=Alabaster, language=cmd]
rmdir old_project

\end{lstlisting}

\begin{center}\rule{0.5\linewidth}{0.5pt}\end{center}

\subsubsection*{Comment the Code (Questions 6--10)}\label{docs__pandoc__latex__src__advanced_study_guide.md__comment-the-code-questions-610-11}

\textbf{Q6.} Add comments explaining each file-operation command.

\begin{lstlisting}[style=Alabaster, language=cmd]
echo content > file1.txt
copy file1.txt file1_backup.txt
move file1.txt archive\file1.txt
dir /B

\end{lstlisting}

\textbf{Q7.} Add comments to this project-backup workflow.

\begin{lstlisting}[style=Alabaster, language=cmd]
mkdir backup
copy *.scad backup\
dir /B backup\
dir /B *.scad

\end{lstlisting}

\textbf{Q8.} Add comments explaining the safe-deletion pattern.

\begin{lstlisting}[style=Alabaster, language=cmd]
dir /B *.tmp
dir /B *.tmp | find /c ""
del *.tmp
dir /B *.tmp

\end{lstlisting}

\textbf{Q9.} Add comments to this folder-structure creation script.

\begin{lstlisting}[style=Alabaster, language=cmd]
mkdir 3DProjects
cd 3DProjects
mkdir Models
mkdir Prints
mkdir Docs
mkdir Backups
dir /B /A:D

\end{lstlisting}

\textbf{Q10.} Add comments to this project-archiving script.

\begin{lstlisting}[style=Alabaster, language=cmd]
for /f "tokens=1-3 delims=/ " %%a in ("%date%") do set TODAY=%%c%%a%%b
mkdir archive_%TODAY%
xcopy /E /I my_cube archive_%TODAY%\my_cube
dir /B

\end{lstlisting}

\begin{center}\rule{0.5\linewidth}{0.5pt}\end{center}

\subsection*{CMD --- Lesson 3: Input, Output, and Piping}\label{docs__pandoc__latex__src__advanced_study_guide.md__cmd--lesson-3-input-output-and-piping}

\subsubsection*{Fix the Code (Questions 1--5)}\label{docs__pandoc__latex__src__advanced_study_guide.md__fix-the-code-questions-15-12}

\textbf{Q1.} The user wants to append to a log file, but the first line overwrites the file.

\begin{lstlisting}[style=Alabaster, language=cmd]
echo Build started > build_log.txt
echo Build complete > build_log.txt
type build_log.txt

\end{lstlisting}

\textbf{Q2.} The pipe destination for copying to the Windows clipboard is wrong.

\begin{lstlisting}[style=Alabaster, language=cmd]
type notes.txt | copy

\end{lstlisting}

\textbf{Q3.} The \texttt{findstr} command is supposed to be case-insensitive but is missing the correct flag.

\begin{lstlisting}[style=Alabaster, language=cmd]
type build_log.txt | findstr "error"

\end{lstlisting}

\textbf{Q4.} This chain should count the number of \texttt{.scad} files but uses the wrong counting syntax.

\begin{lstlisting}[style=Alabaster, language=cmd]
dir /B *.scad | find /v ""

\end{lstlisting}

\emph{(Hint: check the \texttt{find} flags for counting lines.)}

\textbf{Q5.} The script should search a file for a keyword and save the results, but the command for reading a file in CMD is wrong.

\begin{lstlisting}[style=Alabaster, language=cmd]
cat build_log.txt | findstr "FAILED" > errors.txt

\end{lstlisting}

\begin{center}\rule{0.5\linewidth}{0.5pt}\end{center}

\subsubsection*{Comment the Code (Questions 6--10)}\label{docs__pandoc__latex__src__advanced_study_guide.md__comment-the-code-questions-610-12}

\textbf{Q6.} Add comments explaining the difference between \texttt{\textgreater{}} and \texttt{\textgreater{}\textgreater{}} in each line.

\begin{lstlisting}[style=Alabaster, language=cmd]
echo Project: my_cube > log.txt
echo Build started >> log.txt
echo Build complete >> log.txt
type log.txt

\end{lstlisting}

\textbf{Q7.} Add comments to this pipe chain.

\begin{lstlisting}[style=Alabaster, language=cmd]
dir /B *.scad | findstr "main"
dir /B *.scad | find /c "" 
dir /B | findstr /R "\.scad$" > scad_list.txt

\end{lstlisting}

\textbf{Q8.} Add comments explaining each step of this build-and-verify workflow.

\begin{lstlisting}[style=Alabaster, language=cmd]
3dm build > build_output.txt 2>&1
type build_output.txt
3dm info > info_report.txt
notepad.exe info_report.txt

\end{lstlisting}

\textbf{Q9.} Add comments explaining this clipboard workflow for screen readers.

\begin{lstlisting}[style=Alabaster, language=cmd]
dir /B | clip
type notes.txt | clip
dir /B *.scad | clip

\end{lstlisting}

\textbf{Q10.} Add comments to this multi-step log-building script.

\begin{lstlisting}[style=Alabaster, language=cmd]
echo === Build Log === > full_log.txt
echo Date: %date% >> full_log.txt
echo Time: %time% >> full_log.txt
3dm build >> full_log.txt 2>&1
3dm info >> full_log.txt
notepad.exe full_log.txt

\end{lstlisting}

\begin{center}\rule{0.5\linewidth}{0.5pt}\end{center}

\subsection*{CMD --- Lesson 4: Environment Variables, PATH, and Aliases (DOSKEY)}\label{docs__pandoc__latex__src__advanced_study_guide.md__cmd--lesson-4-environment-variables-path-and-aliases-doskey}

\subsubsection*{Fix the Code (Questions 1--5)}\label{docs__pandoc__latex__src__advanced_study_guide.md__fix-the-code-questions-15-13}

\textbf{Q1.} The user tries to set a persistent environment variable using the wrong command.

\begin{lstlisting}[style=Alabaster, language=cmd]
set PROJECT_DIR=C:\Users\Student\Documents\3d_projects

\end{lstlisting}

\emph{(Hint: \texttt{set} only persists for the current session.)}

\textbf{Q2.} The DOSKEY alias definition has a wrong syntax for passing arguments.

\begin{lstlisting}[style=Alabaster, language=cmd]
doskey build=3dm build $1

\end{lstlisting}

\textbf{Q3.} The script tries to print an environment variable but uses the wrong syntax.

\begin{lstlisting}[style=Alabaster, language=cmd]
echo $PROJECT_DIR

\end{lstlisting}

\textbf{Q4.} The user wants to check where \texttt{openscad} is installed, but uses a Linux command.

\begin{lstlisting}[style=Alabaster, language=cmd]
which openscad

\end{lstlisting}

\textbf{Q5.} The PATH modification syntax is wrong --- the semicolon and concatenation are missing.

\begin{lstlisting}[style=Alabaster, language=cmd]
set PATH=%PATH% C:\Tools\OpenSCAD

\end{lstlisting}

\begin{center}\rule{0.5\linewidth}{0.5pt}\end{center}

\subsubsection*{Comment the Code (Questions 6--10)}\label{docs__pandoc__latex__src__advanced_study_guide.md__comment-the-code-questions-610-13}

\textbf{Q6.} Add comments explaining each variable-setting command and its scope.

\begin{lstlisting}[style=Alabaster, language=cmd]
set MY_VAR=hello
setx MY_VAR hello
set MY_VAR
echo %MY_VAR%

\end{lstlisting}

\textbf{Q7.} Add comments to this DOSKEY alias setup.

\begin{lstlisting}[style=Alabaster, language=cmd]
doskey ls=dir /B $*
doskey cat=type $1
doskey build=3dm build $*
doskey info=3dm info $*

\end{lstlisting}

\textbf{Q8.} Add comments explaining this PATH diagnostic sequence.

\begin{lstlisting}[style=Alabaster, language=cmd]
echo %PATH%
where openscad
where 3dm
openscad --version

\end{lstlisting}

\textbf{Q9.} Add comments to this environment-variable usage script.

\begin{lstlisting}[style=Alabaster, language=cmd]
set PROJECT_DIR=%USERPROFILE%\Documents\3d_projects
cd %PROJECT_DIR%
dir /B
echo %PROJECT_DIR%

\end{lstlisting}

\textbf{Q10.} Add comments explaining the purpose of each block in this persistent-environment-setup script.

\begin{lstlisting}[style=Alabaster, language=cmd]
setx PROJECT_DIR "%USERPROFILE%\Documents\3d_projects"
setx OPENSCAD_PATH "C:\Program Files\OpenSCAD"
set PATH=%PATH%;%OPENSCAD_PATH%
where openscad
echo Setup complete.

\end{lstlisting}

\begin{center}\rule{0.5\linewidth}{0.5pt}\end{center}

\subsection*{CMD --- Lesson 5: Startup Scripts, History, and Useful Tricks}\label{docs__pandoc__latex__src__advanced_study_guide.md__cmd--lesson-5-startup-scripts-history-and-useful-tricks}

\subsubsection*{Fix the Code (Questions 1--5)}\label{docs__pandoc__latex__src__advanced_study_guide.md__fix-the-code-questions-15-14}

\textbf{Q1.} The history popup shortcut is wrong --- the user pressed the wrong function key.

\begin{lstlisting}[style=Alabaster, language=cmd]
REM Press F6 to open the command history popup

\end{lstlisting}

\textbf{Q2.} The DOSKEY macro file is loaded with the wrong flag.

\begin{lstlisting}[style=Alabaster, language=cmd]
doskey /load macros.txt

\end{lstlisting}

\textbf{Q3.} The script tries to re-run the previous command automatically, but the syntax is wrong.

\begin{lstlisting}[style=Alabaster, language=cmd]
doskey /history | tail -1 | cmd

\end{lstlisting}

\textbf{Q4.} The AUTORUN registry key path is misspelled, so the startup script will not load.

\begin{lstlisting}[style=Alabaster, language=cmd]
reg add "HKCU\Software\Microsoft\Command Processer" /v AutoRun /t REG_SZ /d "%USERPROFILE%\cmd_startup.bat"

\end{lstlisting}

\textbf{Q5.} This startup script tries to set an alias at launch but uses the wrong command name.

\begin{lstlisting}[style=Alabaster, language=cmd]
@echo off
alias build=3dm build
echo CMD environment loaded.

\end{lstlisting}

\begin{center}\rule{0.5\linewidth}{0.5pt}\end{center}

\subsubsection*{Comment the Code (Questions 6--10)}\label{docs__pandoc__latex__src__advanced_study_guide.md__comment-the-code-questions-610-14}

\textbf{Q6.} Add comments to this history-recall reference.

\begin{lstlisting}[style=Alabaster, language=cmd]
REM F7 key
REM F8 key
REM doskey /history
doskey /history > history_log.txt
notepad.exe history_log.txt

\end{lstlisting}

\textbf{Q7.} Add comments explaining each line of this startup script (\texttt{cmd\_startup.bat}).

\begin{lstlisting}[style=Alabaster, language=cmd]
@echo off
doskey ls=dir /B $*
doskey cat=type $1
doskey build=3dm build $*
set PROJECT_DIR=%USERPROFILE%\Documents\3d_projects
echo CMD environment ready. Project dir: %PROJECT_DIR%

\end{lstlisting}

\textbf{Q8.} Add comments to this screen-reader output-management sequence.

\begin{lstlisting}[style=Alabaster, language=cmd]
3dm build > build_out.txt 2>&1
if %errorlevel% neq 0 (
    echo BUILD FAILED >> build_out.txt
)
notepad.exe build_out.txt

\end{lstlisting}

\textbf{Q9.} Add comments explaining each step of this macro-persistence workflow.

\begin{lstlisting}[style=Alabaster, language=cmd]
doskey build=3dm build $*
doskey /macros > macros.txt
notepad.exe macros.txt
doskey /macrofile=macros.txt

\end{lstlisting}

\textbf{Q10.} Add comments to this end-of-session cleanup and log script.

\begin{lstlisting}[style=Alabaster, language=cmd]
3dm info > session_report.txt
echo Session ended: %date% %time% >> session_report.txt
doskey /history >> session_report.txt
notepad.exe session_report.txt

\end{lstlisting}

\begin{center}\rule{0.5\linewidth}{0.5pt}\end{center}

\subsection*{CMD --- Lesson 6: Batch Scripts, Subroutines, and Professional Workflows}\label{docs__pandoc__latex__src__advanced_study_guide.md__cmd--lesson-6-batch-scripts-subroutines-and-professional-workflows}

\subsubsection*{Fix the Code (Questions 1--5)}\label{docs__pandoc__latex__src__advanced_study_guide.md__fix-the-code-questions-15-15}

\textbf{Q1.} This batch script should loop over widths and build variants, but the loop variable syntax is wrong.

\begin{lstlisting}[style=Alabaster, language=cmd]
@echo off
for $WIDTH in (60 80 100 120) do (
    echo Building width: %WIDTH%
    openscad -D "width=%WIDTH%" -o build\bracket_%WIDTH%.stl src\main.scad
)

\end{lstlisting}

\textbf{Q2.} The error-check after running openscad uses the wrong variable.

\begin{lstlisting}[style=Alabaster, language=cmd]
openscad -o build\main.stl src\main.scad
if %exitcode% neq 0 echo BUILD FAILED

\end{lstlisting}

\textbf{Q3.} The subroutine call is missing the \texttt{call} keyword, so the script will exit after the subroutine.

\begin{lstlisting}[style=Alabaster, language=cmd]
@echo off
:main
echo Starting build...
:build_project
    3dm build
    goto :eof

\end{lstlisting}

\emph{(Hint: \texttt{:build\_project} needs to be invoked correctly.)}

\textbf{Q4.} The \texttt{setlocal} / \texttt{endlocal} block is in the wrong order, so local variables leak out.

\begin{lstlisting}[style=Alabaster, language=cmd]
@echo off
set TEMP_VAR=hello
setlocal
echo %TEMP_VAR%
endlocal

\end{lstlisting}

\textbf{Q5.} This script saves the ERRORLEVEL but uses the wrong syntax to capture it.

\begin{lstlisting}[style=Alabaster, language=cmd]
3dm build
set BUILD_RESULT = %errorlevel%
if %BUILD_RESULT% neq 0 echo Build failed with code %BUILD_RESULT%

\end{lstlisting}

\begin{center}\rule{0.5\linewidth}{0.5pt}\end{center}

\subsubsection*{Comment the Code (Questions 6--10)}\label{docs__pandoc__latex__src__advanced_study_guide.md__comment-the-code-questions-610-15}

\textbf{Q6.} Add comments to every significant line of this batch build script.

\begin{lstlisting}[style=Alabaster, language=cmd]
@echo off
setlocal enabledelayedexpansion
set OUTPUT_DIR=build\variants
if not exist %OUTPUT_DIR% mkdir %OUTPUT_DIR%
set SUCCESS=0
set FAIL=0
for %%W in (60 80 100 120) do (
    for %%T in (1.2 1.6 2.0) do (
        set NAME=bracket_w%%W_wall%%T
        openscad -D "width=%%W" -D "wall=%%T" -o "%OUTPUT_DIR%\!NAME!.stl" src\main.scad
        if !errorlevel! equ 0 (
            set /a SUCCESS+=1
        ) else (
            set /a FAIL+=1
        )
    )
)
echo Done: %SUCCESS% succeeded, %FAIL% failed
endlocal

\end{lstlisting}

\textbf{Q7.} Add comments to this subroutine-based build script.

\begin{lstlisting}[style=Alabaster, language=cmd]
@echo off
call :log "Build started"
call :build_step
call :verify_step
call :log "Build finished"
goto :eof

:log
echo [%time%] %~1 >> build_log.txt
goto :eof

:build_step
3dm build >> build_log.txt 2>&1
if %errorlevel% neq 0 ( echo FAILED >> build_log.txt & exit /b 1 )
goto :eof

:verify_step
3dm info >> build_log.txt
goto :eof

\end{lstlisting}

\textbf{Q8.} Add comments explaining each block of this error-handling script.

\begin{lstlisting}[style=Alabaster, language=cmd]
@echo off
3dm build > build_out.txt 2>&1
set BUILD_CODE=%errorlevel%
if %BUILD_CODE% neq 0 (
    echo Build failed with code %BUILD_CODE%
    type build_out.txt
    exit /b %BUILD_CODE%
)
for %%F in (build\main.stl) do set STL_SIZE=%%~zF
if %STL_SIZE% lss 1000 echo WARNING: STL file very small: %STL_SIZE% bytes
echo Build OK: %STL_SIZE% bytes

\end{lstlisting}

\textbf{Q9.} Add comments to this parameterized batch-project setup script.

\begin{lstlisting}[style=Alabaster, language=cmd]
@echo off
set NAME=%1
if "%NAME%"=="" ( echo Usage: new_project.bat project_name & exit /b 1 )
set BASE=%USERPROFILE%\Documents\3d_projects\%NAME%
if exist "%BASE%" ( echo Already exists: %BASE% & exit /b 1 )
mkdir "%BASE%"
cd /d "%BASE%"
3dm new
echo Project created: %BASE%
dir /B

\end{lstlisting}

\textbf{Q10.} Add comments explaining every section of this variant-archive script.

\begin{lstlisting}[style=Alabaster, language=cmd]
@echo off
for /f "tokens=1-3 delims=/ " %%a in ("%date%") do (
    set TODAY=%%c%%a%%b
)
set ARCHIVE=build\archive_%TODAY%
mkdir %ARCHIVE%
for %%F in (build\variants\*.stl) do (
    copy "%%F" "%ARCHIVE%\%%~nxF" > nul
)
echo Archived to: %ARCHIVE%
dir /B %ARCHIVE%

\end{lstlisting}

\begin{center}\rule{0.5\linewidth}{0.5pt}\end{center}

\section*{PART 3 --- Git Bash}\label{docs__pandoc__latex__src__advanced_study_guide.md__part-3--git-bash}

\begin{center}\rule{0.5\linewidth}{0.5pt}\end{center}

\subsection*{Git Bash --- Lesson Pre: Your First Terminal}\label{docs__pandoc__latex__src__advanced_study_guide.md__git-bash--lesson-pre-your-first-terminal}

\subsubsection*{Fix the Code (Questions 1--5)}\label{docs__pandoc__latex__src__advanced_study_guide.md__fix-the-code-questions-15-16}

\textbf{Q1.} The user tries to print the current directory and list files, but one command is wrong for Git Bash.

\begin{lstlisting}[style=Alabaster, language=bash]
pwd
dir

\end{lstlisting}

\textbf{Q2.} The screen-reader--friendly listing command uses the wrong flag (letter L instead of number 1).

\begin{lstlisting}[style=Alabaster, language=bash]
ls -l

\end{lstlisting}

\emph{(Note: this produces detail output, not the one-per-line names-only format.)}

\textbf{Q3.} The user wants to create a file and read it back, but uses the CMD-style read command.

\begin{lstlisting}[style=Alabaster, language=bash]
echo "Hello Git Bash" > hello.txt
type hello.txt

\end{lstlisting}

\textbf{Q4.} The user wants to navigate home and verify, but uses the wrong home-navigation command.

\begin{lstlisting}[style=Alabaster, language=bash]
cd /home
pwd

\end{lstlisting}

\textbf{Q5.} The long-output pagination command is wrong for Git Bash.

\begin{lstlisting}[style=Alabaster, language=bash]
ls -1 | more

\end{lstlisting}

\emph{(Hint: Git Bash uses a different pager than CMD.)}

\begin{center}\rule{0.5\linewidth}{0.5pt}\end{center}

\subsubsection*{Comment the Code (Questions 6--10)}\label{docs__pandoc__latex__src__advanced_study_guide.md__comment-the-code-questions-610-16}

\textbf{Q6.} Add comments explaining every command in this first-session orientation.

\begin{lstlisting}[style=Alabaster, language=bash]
pwd
ls -1
cd Documents
pwd
ls -1
cd ~
pwd

\end{lstlisting}

\textbf{Q7.} Add comments to this file-creation and reading workflow.

\begin{lstlisting}[style=Alabaster, language=bash]
echo "Hello, Git Bash!" > hello.txt
cat hello.txt
echo "Line 2" >> hello.txt
cat hello.txt

\end{lstlisting}

\textbf{Q8.} Add comments explaining the output-redirection strategy for screen reader users.

\begin{lstlisting}[style=Alabaster, language=bash]
ls -1 > file_list.txt
notepad.exe file_list.txt
ls -1 | grep "scad"

\end{lstlisting}

\textbf{Q9.} Add comments to this keyboard-shortcut reference for Git Bash.

\begin{lstlisting}[style=Alabaster, language=bash]
# Up Arrow
# Down Arrow
# Tab
# Ctrl+C
# Ctrl+R

\end{lstlisting}

\textbf{Q10.} Add comments to this tab-completion demonstration.

\begin{lstlisting}[style=Alabaster, language=bash]
cd ~/Doc
ls -1
cd ~
cd ~/Down

\end{lstlisting}

\emph{(Note: lines 1 and 4 show partial paths where Tab would be pressed.)}

\begin{center}\rule{0.5\linewidth}{0.5pt}\end{center}

\subsection*{Git Bash --- Lesson 0: Getting Started}\label{docs__pandoc__latex__src__advanced_study_guide.md__git-bash--lesson-0-getting-started}

\subsubsection*{Fix the Code (Questions 1--5)}\label{docs__pandoc__latex__src__advanced_study_guide.md__fix-the-code-questions-15-17}

\textbf{Q1.} Git Bash uses Unix-style paths, but this navigation uses a Windows-style path.

\begin{lstlisting}[style=Alabaster, language=bash]
cd C:\Users\Student\Documents

\end{lstlisting}

\textbf{Q2.} The listing command uses the wrong flag --- letter L instead of number 1.

\begin{lstlisting}[style=Alabaster, language=bash]
ls -l

\end{lstlisting}

\textbf{Q3.} The script should go up one level and then back home, but both navigation commands are wrong.

\begin{lstlisting}[style=Alabaster, language=bash]
cd ../
cd /home/student

\end{lstlisting}

\textbf{Q4.} The 3DMake project creation sequence uses the wrong path expansion variable.

\begin{lstlisting}[style=Alabaster, language=bash]
mkdir %USERPROFILE%/Documents/3d_projects/my_cube
cd %USERPROFILE%/Documents/3d_projects/my_cube
3dm new

\end{lstlisting}

\textbf{Q5.} The user wants to list hidden files along with regular files, but uses the wrong flag.

\begin{lstlisting}[style=Alabaster, language=bash]
ls -1 -h

\end{lstlisting}

\begin{center}\rule{0.5\linewidth}{0.5pt}\end{center}

\subsubsection*{Comment the Code (Questions 6--10)}\label{docs__pandoc__latex__src__advanced_study_guide.md__comment-the-code-questions-610-17}

\textbf{Q6.} Add comments explaining each path shortcut.

\begin{lstlisting}[style=Alabaster, language=bash]
cd ~
cd ./subfolder
cd ../
cd -
cd /c/Users/Student/Documents

\end{lstlisting}

\textbf{Q7.} Add comments to this 3DMake quick-start sequence.

\begin{lstlisting}[style=Alabaster, language=bash]
cd ~/Downloads/3dmake
./3dm setup
mkdir -p ~/Documents/3d_projects/my_cube
cd ~/Documents/3d_projects/my_cube
3dm new
ls -1

\end{lstlisting}

\textbf{Q8.} Add comments explaining what each line of the prompt anatomy reference means.

\begin{lstlisting}[style=Alabaster, language=bash]
# student@LAPTOP MINGW64 ~/Documents/3d_projects $
pwd
ls -1
cd ..
pwd

\end{lstlisting}

\textbf{Q9.} Add comments to this path-comparison demonstration.

\begin{lstlisting}[style=Alabaster, language=bash]
pwd
cd /c/Users/Student/Documents/3d_projects
pwd
cd ~
pwd
cd Documents/3d_projects
pwd

\end{lstlisting}

\textbf{Q10.} Add comments explaining every line of this screen-reader-first listing workflow.

\begin{lstlisting}[style=Alabaster, language=bash]
ls -1 > list.txt
notepad.exe list.txt
ls -1 | grep "scad"
ls -1 | grep "scad" > scad_list.txt
notepad.exe scad_list.txt

\end{lstlisting}

\begin{center}\rule{0.5\linewidth}{0.5pt}\end{center}

\subsection*{Git Bash --- Lesson 1: Navigation}\label{docs__pandoc__latex__src__advanced_study_guide.md__git-bash--lesson-1-navigation}

\subsubsection*{Fix the Code (Questions 1--5)}\label{docs__pandoc__latex__src__advanced_study_guide.md__fix-the-code-questions-15-18}

\textbf{Q1.} The wildcard pattern for listing \texttt{.scad} files is wrong for Git Bash.

\begin{lstlisting}[style=Alabaster, language=bash]
ls -1 *.scad
ls -1 | grep /.scad/

\end{lstlisting}

\emph{(The first line is correct --- identify the error in the second line.)}

\textbf{Q2.} The script should list only directories. The flag used is wrong for Git Bash.

\begin{lstlisting}[style=Alabaster, language=bash]
ls -1 /A:D

\end{lstlisting}

\textbf{Q3.} The \texttt{find} command should search recursively for \texttt{.scad} files but uses wrong syntax.

\begin{lstlisting}[style=Alabaster, language=bash]
find . -name "*.scad" -type d

\end{lstlisting}

\textbf{Q4.} The script saves a listing and opens it, but uses the wrong open command.

\begin{lstlisting}[style=Alabaster, language=bash]
ls -1 > list.txt
open list.txt

\end{lstlisting}

\textbf{Q5.} The recursive search for all STL files uses the wrong wildcard.

\begin{lstlisting}[style=Alabaster, language=bash]
ls -1 build/%.stl

\end{lstlisting}

\begin{center}\rule{0.5\linewidth}{0.5pt}\end{center}

\subsubsection*{Comment the Code (Questions 6--10)}\label{docs__pandoc__latex__src__advanced_study_guide.md__comment-the-code-questions-610-18}

\textbf{Q6.} Add comments to this navigation and listing sequence.

\begin{lstlisting}[style=Alabaster, language=bash]
pwd
cd ~/Documents/3d_projects
ls -1
cd my_cube
ls -1
cd src
ls -1 *.scad
cd ~

\end{lstlisting}

\textbf{Q7.} Add comments explaining each wildcard pattern.

\begin{lstlisting}[style=Alabaster, language=bash]
ls -1 *.scad
ls -1 model?.scad
ls -1 *_v[0-9].*
ls -1 *_final*

\end{lstlisting}

\textbf{Q8.} Add comments to this project-exploration script.

\begin{lstlisting}[style=Alabaster, language=bash]
cd ~/Documents/3d_projects/my_cube
pwd
ls -1
ls -1 src/
ls -1 build/
ls -1 build/*.stl

\end{lstlisting}

\textbf{Q9.} Add comments explaining this screen-reader listing workflow.

\begin{lstlisting}[style=Alabaster, language=bash]
ls -1 > all_files.txt
ls -1 *.scad > scad_files.txt
ls -1 -d */ > dirs_only.txt
notepad.exe all_files.txt

\end{lstlisting}

\textbf{Q10.} Add comments explaining what \texttt{find} does in each variation.

\begin{lstlisting}[style=Alabaster, language=bash]
find . -name "*.scad"
find . -name "*.scad" -type f
find ~/Documents/3d_projects -name "*.stl"
find . -name "*.scad" > scad_paths.txt

\end{lstlisting}

\begin{center}\rule{0.5\linewidth}{0.5pt}\end{center}

\subsection*{Git Bash --- Lesson 2: File and Folder Manipulation}\label{docs__pandoc__latex__src__advanced_study_guide.md__git-bash--lesson-2-file-and-folder-manipulation}

\subsubsection*{Fix the Code (Questions 1--5)}\label{docs__pandoc__latex__src__advanced_study_guide.md__fix-the-code-questions-15-19}

\textbf{Q1.} The command to recursively copy a folder is missing the required flag.

\begin{lstlisting}[style=Alabaster, language=bash]
cp my_cube my_cube_backup

\end{lstlisting}

\textbf{Q2.} The user deletes a folder but uses the wrong command for Git Bash.

\begin{lstlisting}[style=Alabaster, language=bash]
rmdir /S /Q old_project

\end{lstlisting}

\textbf{Q3.} The script creates a file, but uses the CMD command instead of the Git Bash equivalent.

\begin{lstlisting}[style=Alabaster, language=bash]
echo. > empty_file.txt

\end{lstlisting}

\textbf{Q4.} The safe-deletion sequence is out of order --- files are deleted before verification.

\begin{lstlisting}[style=Alabaster, language=bash]
rm *.tmp
ls -1 *.tmp

\end{lstlisting}

\textbf{Q5.} The move command renames a file but uses the wrong command for Git Bash.

\begin{lstlisting}[style=Alabaster, language=bash]
ren old_name.scad new_name.scad

\end{lstlisting}

\begin{center}\rule{0.5\linewidth}{0.5pt}\end{center}

\subsubsection*{Comment the Code (Questions 6--10)}\label{docs__pandoc__latex__src__advanced_study_guide.md__comment-the-code-questions-610-19}

\textbf{Q6.} Add comments to this file-creation and manipulation sequence.

\begin{lstlisting}[style=Alabaster, language=bash]
touch design.scad
cp design.scad design_backup.scad
mv design.scad archive/design_v1.scad
ls -1

\end{lstlisting}

\textbf{Q7.} Add comments to this project-backup workflow.

\begin{lstlisting}[style=Alabaster, language=bash]
cd ~/Documents/3d_projects
cp -r my_cube my_cube_backup
ls -1
cd my_cube_backup
ls -1

\end{lstlisting}

\textbf{Q8.} Add comments explaining the safe-deletion pattern.

\begin{lstlisting}[style=Alabaster, language=bash]
ls -1 *.tmp
rm *.tmp
ls -1 *.tmp

\end{lstlisting}

\textbf{Q9.} Add comments to this folder-structure creation script.

\begin{lstlisting}[style=Alabaster, language=bash]
mkdir -p 3DProjects/{Models,Prints,Docs,Backups}
ls -1 3DProjects/
ls -1 -d 3DProjects/*/

\end{lstlisting}

\textbf{Q10.} Add comments explaining every operation in this 3DMake project-management script.

\begin{lstlisting}[style=Alabaster, language=bash]
mkdir -p ~/Documents/3d_projects/bottle_cap
cd ~/Documents/3d_projects/bottle_cap
3dm new
cd ~/Documents/3d_projects
cp -r bottle_cap bottle_cap_backup
cd bottle_cap/src
mv main.scad bottle_cap.scad
ls -1 *.scad

\end{lstlisting}

\begin{center}\rule{0.5\linewidth}{0.5pt}\end{center}

\subsection*{Git Bash --- Lesson 3: Input, Output, and Piping}\label{docs__pandoc__latex__src__advanced_study_guide.md__git-bash--lesson-3-input-output-and-piping}

\subsubsection*{Fix the Code (Questions 1--5)}\label{docs__pandoc__latex__src__advanced_study_guide.md__fix-the-code-questions-15-20}

\textbf{Q1.} The second line overwrites the log instead of appending to it.

\begin{lstlisting}[style=Alabaster, language=bash]
echo "Build started" > build_log.txt
echo "Build complete" > build_log.txt
cat build_log.txt

\end{lstlisting}

\textbf{Q2.} The pipe destination for the Windows clipboard is wrong for Git Bash on Windows.

\begin{lstlisting}[style=Alabaster, language=bash]
cat notes.txt | clipboard

\end{lstlisting}

\textbf{Q3.} This \texttt{grep} command should be case-insensitive but is missing the flag.

\begin{lstlisting}[style=Alabaster, language=bash]
cat build_log.txt | grep "error"

\end{lstlisting}

\textbf{Q4.} The command counts lines but uses the wrong utility for Git Bash.

\begin{lstlisting}[style=Alabaster, language=bash]
ls -1 *.scad | find /c ""

\end{lstlisting}

\textbf{Q5.} This script reads a file and searches it, but uses the CMD-style read command.

\begin{lstlisting}[style=Alabaster, language=bash]
type build_log.txt | grep "FAILED"

\end{lstlisting}

\begin{center}\rule{0.5\linewidth}{0.5pt}\end{center}

\subsubsection*{Comment the Code (Questions 6--10)}\label{docs__pandoc__latex__src__advanced_study_guide.md__comment-the-code-questions-610-20}

\textbf{Q6.} Add comments explaining the difference between \texttt{\textgreater{}} and \texttt{\textgreater{}\textgreater{}}.

\begin{lstlisting}[style=Alabaster, language=bash]
echo "Project: my_cube" > log.txt
echo "Build started" >> log.txt
echo "Build complete" >> log.txt
cat log.txt

\end{lstlisting}

\textbf{Q7.} Add comments to this pipe chain.

\begin{lstlisting}[style=Alabaster, language=bash]
ls -1 *.scad | grep "main"
ls -1 *.scad | wc -l
ls -1 | grep ".stl" > stl_list.txt

\end{lstlisting}

\textbf{Q8.} Add comments to this build-and-verify workflow.

\begin{lstlisting}[style=Alabaster, language=bash]
3dm build
3dm info > model_report.txt
notepad.exe model_report.txt
cat model_report.txt | grep "Bounding"

\end{lstlisting}

\textbf{Q9.} Add comments explaining this clipboard workflow.

\begin{lstlisting}[style=Alabaster, language=bash]
ls -1 | clip
cat notes.txt | clip
ls -1 *.scad | clip

\end{lstlisting}

\textbf{Q10.} Add comments to this multi-step log-building script.

\begin{lstlisting}[style=Alabaster, language=bash]
echo "=== Build Log ===" > full_log.txt
echo "Date: $(date)" >> full_log.txt
3dm build >> full_log.txt 2>&1
3dm info >> full_log.txt
notepad.exe full_log.txt

\end{lstlisting}

\begin{center}\rule{0.5\linewidth}{0.5pt}\end{center}

\subsection*{Git Bash --- Lesson 4: Environment Variables, PATH, and Aliases}\label{docs__pandoc__latex__src__advanced_study_guide.md__git-bash--lesson-4-environment-variables-path-and-aliases}

\subsubsection*{Fix the Code (Questions 1--5)}\label{docs__pandoc__latex__src__advanced_study_guide.md__fix-the-code-questions-15-21}

\textbf{Q1.} The alias is created, but changes will be lost when the terminal is closed. Where should this line go to persist?

\begin{lstlisting}[style=Alabaster, language=bash]
alias build="3dm build"

\end{lstlisting}

\emph{(Hint: identify where this line must be saved for persistence --- this is a conceptual fix; write the corrected workflow, not just the line.)}

\textbf{Q2.} The variable is set but the export is missing, so child processes won\textquotesingle t see it.

\begin{lstlisting}[style=Alabaster, language=bash]
PROJECT_DIR="$HOME/Documents/3d_projects"

\end{lstlisting}

\textbf{Q3.} The PATH modification uses the wrong separator.

\begin{lstlisting}[style=Alabaster, language=bash]
export PATH="$PATH;/c/Program Files/OpenSCAD"

\end{lstlisting}

\textbf{Q4.} The \texttt{which} command is used correctly but \texttt{source} is used wrong to reload \texttt{.bashrc}.

\begin{lstlisting}[style=Alabaster, language=bash]
which openscad
reload ~/.bashrc

\end{lstlisting}

\textbf{Q5.} The variable reference uses wrong shell syntax.

\begin{lstlisting}[style=Alabaster, language=bash]
export LABEL_FONT="Liberation Sans"
echo %LABEL_FONT%

\end{lstlisting}

\begin{center}\rule{0.5\linewidth}{0.5pt}\end{center}

\subsubsection*{Comment the Code (Questions 6--10)}\label{docs__pandoc__latex__src__advanced_study_guide.md__comment-the-code-questions-610-21}

\textbf{Q6.} Add comments explaining each line of this environment-setup block.

\begin{lstlisting}[style=Alabaster, language=bash]
export PROJECT_DIR="$HOME/Documents/3d_projects"
export OPENSCAD_PATH="/c/Program Files/OpenSCAD"
export PATH="$PATH:$OPENSCAD_PATH"

\end{lstlisting}

\textbf{Q7.} Add comments to this alias-creation and verification sequence.

\begin{lstlisting}[style=Alabaster, language=bash]
alias build="3dm build"
alias info="3dm info"
alias preview="3dm preview"
alias ls1="ls -1"

\end{lstlisting}

\textbf{Q8.} Add comments explaining this PATH diagnostic script.

\begin{lstlisting}[style=Alabaster, language=bash]
echo $PATH | tr ':' '\n'
which openscad
which 3dm
openscad --version

\end{lstlisting}

\textbf{Q9.} Add comments to this \texttt{.bashrc}-editing sequence.

\begin{lstlisting}[style=Alabaster, language=bash]
cat ~/.bashrc
echo 'alias build="3dm build"' >> ~/.bashrc
echo 'export PROJECT_DIR="$HOME/Documents/3d_projects"' >> ~/.bashrc
source ~/.bashrc
alias build

\end{lstlisting}

\textbf{Q10.} Add comments explaining each function and variable in this environment block (intended for \texttt{.bashrc}).

\begin{lstlisting}[style=Alabaster, language=bash]
export PROJECT_DIR="$HOME/Documents/3d_projects"
alias build="3dm build"
alias info="3dm info"
function go-projects() { cd "$PROJECT_DIR"; }
function new-project() { mkdir -p "$PROJECT_DIR/$1" && cd "$PROJECT_DIR/$1" && 3dm new; }
echo "3D printing environment loaded."

\end{lstlisting}

\begin{center}\rule{0.5\linewidth}{0.5pt}\end{center}

\subsection*{Git Bash --- Lesson 5: Shell Profiles, History, and Useful Tricks}\label{docs__pandoc__latex__src__advanced_study_guide.md__git-bash--lesson-5-shell-profiles-history-and-useful-tricks}

\subsubsection*{Fix the Code (Questions 1--5)}\label{docs__pandoc__latex__src__advanced_study_guide.md__fix-the-code-questions-15-22}

\textbf{Q1.} The history search shortcut description is wrong.

\begin{lstlisting}[style=Alabaster, language=bash]
# Press Ctrl+F to search command history interactively

\end{lstlisting}

\textbf{Q2.} The alias is added to the wrong profile file for Git Bash on Windows.

\begin{lstlisting}[style=Alabaster, language=bash]
echo 'alias build="3dm build"' >> ~/.bash_profile

\end{lstlisting}

\emph{(Hint: on most Git Bash setups, interactive-shell configuration goes in the other file.)}

\textbf{Q3.} The script re-runs the last command from history using the wrong syntax.

\begin{lstlisting}[style=Alabaster, language=bash]
history | tail -1 | bash

\end{lstlisting}

\textbf{Q4.} The \texttt{source} command is used but the wrong file is sourced.

\begin{lstlisting}[style=Alabaster, language=bash]
source ~/.bash_history

\end{lstlisting}

\textbf{Q5.} The history-export pipeline uses the wrong file-write command for bash.

\begin{lstlisting}[style=Alabaster, language=bash]
history | Write-File history_log.txt

\end{lstlisting}

\begin{center}\rule{0.5\linewidth}{0.5pt}\end{center}

\subsubsection*{Comment the Code (Questions 6--10)}\label{docs__pandoc__latex__src__advanced_study_guide.md__comment-the-code-questions-610-22}

\textbf{Q6.} Add comments to this history-recall workflow.

\begin{lstlisting}[style=Alabaster, language=bash]
history
history | tail -20
history | grep "build"
!build

\end{lstlisting}

\textbf{Q7.} Add comments explaining each step of this \texttt{.bashrc}-management workflow.

\begin{lstlisting}[style=Alabaster, language=bash]
cat ~/.bashrc
notepad.exe ~/.bashrc
source ~/.bashrc
alias build
echo $PROJECT_DIR

\end{lstlisting}

\textbf{Q8.} Add comments to this screen-reader output-management sequence.

\begin{lstlisting}[style=Alabaster, language=bash]
3dm build > build_out.txt 2>&1
cat build_out.txt | grep -i "error"
3dm info > info_out.txt
notepad.exe info_out.txt

\end{lstlisting}

\textbf{Q9.} Add comments explaining each control sequence.

\begin{lstlisting}[style=Alabaster, language=bash]
# Ctrl+R — type to search
# Ctrl+C
# Ctrl+A
# Ctrl+E
# Ctrl+L

\end{lstlisting}

\textbf{Q10.} Add comments to this full environment-check script.

\begin{lstlisting}[style=Alabaster, language=bash]
echo "Shell: $SHELL"
echo "Home: $HOME"
echo "Project dir: $PROJECT_DIR"
which openscad
which 3dm
alias | grep build
history | wc -l

\end{lstlisting}

\begin{center}\rule{0.5\linewidth}{0.5pt}\end{center}

\subsection*{Git Bash --- Lesson 6: Shell Scripts, Functions, Loops}\label{docs__pandoc__latex__src__advanced_study_guide.md__git-bash--lesson-6-shell-scripts-functions-loops}

\subsubsection*{Fix the Code (Questions 1--5)}\label{docs__pandoc__latex__src__advanced_study_guide.md__fix-the-code-questions-15-23}

\textbf{Q1.} This shell script loops over widths and builds variants, but the loop syntax is wrong.

\begin{lstlisting}[style=Alabaster, language=bash]
#!/bin/bash
for WIDTH in (60 80 100 120); do
    openscad -D "width=$WIDTH" -o "build/bracket_$WIDTH.stl" src/main.scad
    echo "Built width: $WIDTH"
done

\end{lstlisting}

\textbf{Q2.} The script uses \texttt{set\ -e} but the OpenSCAD error will be silently suppressed by the pipe.

\begin{lstlisting}[style=Alabaster, language=bash]
#!/bin/bash
set -e
openscad -o build/main.stl src/main.scad | tee build_log.txt
echo "Build done"

\end{lstlisting}

\emph{(Hint: add the missing safety flag.)}

\textbf{Q3.} The function is defined but the parameter is read using the wrong syntax.

\begin{lstlisting}[style=Alabaster, language=bash]
#!/bin/bash
build_variant() {
    local width = $1
    openscad -D "width=$width" -o "build/variant_$width.stl" src/main.scad
}
build_variant 80

\end{lstlisting}

\textbf{Q4.} The script makes a file executable but uses the wrong command name.

\begin{lstlisting}[style=Alabaster, language=bash]
#!/bin/bash
chown +x batch_build.sh
./batch_build.sh

\end{lstlisting}

\textbf{Q5.} The error-handling trap is syntactically wrong.

\begin{lstlisting}[style=Alabaster, language=bash]
#!/bin/bash
set -e
trap "echo Build failed at line $LINENO" ERROR
3dm build
3dm info

\end{lstlisting}

\begin{center}\rule{0.5\linewidth}{0.5pt}\end{center}

\subsubsection*{Comment the Code (Questions 6--10)}\label{docs__pandoc__latex__src__advanced_study_guide.md__comment-the-code-questions-610-23}

\textbf{Q6.} Add comments to this complete batch-build script.

\begin{lstlisting}[style=Alabaster, language=bash]
#!/bin/bash
set -e
set -u
set -o pipefail

OUTPUT_DIR="build/variants"
mkdir -p "$OUTPUT_DIR"

SUCCESS=0
FAIL=0

for WIDTH in 60 80 100 120; do
    for WALL in 1.2 1.6 2.0; do
        NAME="bracket_w${WIDTH}_wall${WALL/./_}"
        OUT="$OUTPUT_DIR/${NAME}.stl"
        if openscad -D "width=${WIDTH}" -D "wall=${WALL}" \
                    -o "$OUT" src/main.scad 2>/dev/null; then
            SUCCESS=$((SUCCESS + 1))
        else
            FAIL=$((FAIL + 1))
        fi
    done
done

echo "Done: $SUCCESS succeeded, $FAIL failed"

\end{lstlisting}

\textbf{Q7.} Add comments to this function-based build script.

\begin{lstlisting}[style=Alabaster, language=bash]
#!/bin/bash
log() {
    local level="$1"; shift
    echo "[$(date +%H:%M:%S)] [$level] $*" | tee -a build_log.txt
}

build_project() {
    local project_dir="$1"
    cd "$project_dir"
    log "INFO" "Building: $(basename $project_dir)"
    if 3dm build > /dev/null 2>&1; then
        log "INFO" "OK"
        return 0
    else
        log "ERROR" "FAILED"
        return 1
    fi
}

for project in ~/Documents/3d_projects/*/; do
    build_project "$project"
done

\end{lstlisting}

\textbf{Q8.} Add comments to this error-handling script.

\begin{lstlisting}[style=Alabaster, language=bash]
#!/bin/bash
set -e
set -u
set -o pipefail

on_error() {
    echo "ERROR at line $1"
    exit 1
}
trap 'on_error $LINENO' ERR

3dm build

STL="build/main.stl"
[ -f "$STL" ] || { echo "STL not created"; exit 1; }

SIZE=$(stat -c%s "$STL")
[ "$SIZE" -gt 1000 ] || echo "WARNING: small STL ($SIZE bytes)"

echo "Build OK: $SIZE bytes"
3dm info

\end{lstlisting}

\textbf{Q9.} Add comments to this string-parameter quoting demonstration.

\begin{lstlisting}[style=Alabaster, language=bash]
#!/bin/bash
LABEL="FRONT PANEL"

openscad -D 'label="FRONT"' -o out1.stl src/main.scad
openscad -D "label=\"${LABEL}\"" -o out2.stl src/main.scad
echo "openscad -D 'label=\"${LABEL}\"' -o out3.stl src/main.scad"

\end{lstlisting}

\textbf{Q10.} Add comments to this complete project-setup script.

\begin{lstlisting}[style=Alabaster, language=bash]
#!/bin/bash
set -e
NAME="$1"
BASE="$HOME/Documents/3d_projects/$NAME"

[ -n "$NAME" ] || { echo "Usage: $0 project_name"; exit 1; }
[ ! -d "$BASE" ] || { echo "Exists: $BASE"; exit 1; }

mkdir -p "$BASE"
cd "$BASE"
3dm new
echo "Created: $BASE"
ls -1

\end{lstlisting}

\begin{center}\rule{0.5\linewidth}{0.5pt}\end{center}

\section*{PART 4 --- 3dMake / OpenSCAD Curriculum}\label{docs__pandoc__latex__src__advanced_study_guide.md__part-4--3dmake--openscad-curriculum}

\begin{center}\rule{0.5\linewidth}{0.5pt}\end{center}

\subsection*{3dMake --- Lesson 1: Environmental Configuration and the Developer Workflow}\label{docs__pandoc__latex__src__advanced_study_guide.md__3dmake--lesson-1-environmental-configuration-and-the-developer-workflow}

\subsubsection*{Fix the Code (Questions 1--5)}\label{docs__pandoc__latex__src__advanced_study_guide.md__fix-the-code-questions-15-24}

\textbf{Q1.} This \texttt{.scad} file uses a raw number instead of a named parameter. Rewrite it so it is fully parametric.

\begin{lstlisting}[style=Alabaster, language=openscad]
cube([30, 20, 10]);

\end{lstlisting}

\textbf{Q2.} The verification command is wrong --- the user types the wrong 3dm subcommand to get model dimensions.

\begin{lstlisting}[style=Alabaster, language=bash]
3dm check

\end{lstlisting}

\textbf{Q3.} This OpenSCAD snippet violates the coordinate system: the cube is underground (extends below Z=0 without intent).

\begin{lstlisting}[style=Alabaster, language=openscad]
width = 30;
depth = 20;
height = 10;
cube([width, depth, height], center=true);

\end{lstlisting}

\textbf{Q4.} The cost-estimation formula has wrong unit conversions.

\begin{lstlisting}[style=Alabaster, language=Python]
volume_mm3 = 6000
density = 1.24        # g/cm³
spool_price = 20      # dollars
spool_weight = 1000   # grams

mass = volume_mm3 * density         # ERROR: wrong unit conversion
cost = mass * (spool_price / spool_weight)
print(f"Mass: {mass:.2f} g, Cost: ${cost:.2f}")

\end{lstlisting}

\textbf{Q5.} The project initialization sequence uses the wrong directory for source files.

\begin{lstlisting}[style=Alabaster, language=bash]
mkdir my_project
cd my_project
3dm new
# User's SCAD file should go here:
notepad.exe main.scad
3dm build

\end{lstlisting}

\begin{center}\rule{0.5\linewidth}{0.5pt}\end{center}

\subsubsection*{Comment the Code (Questions 6--10)}\label{docs__pandoc__latex__src__advanced_study_guide.md__comment-the-code-questions-610-24}

\textbf{Q6.} Add comments explaining every line of this parametric cube file.

\begin{lstlisting}[style=Alabaster, language=openscad]
width  = 30;
depth  = 20;
height = 10;

cube([width, depth, height]);

\end{lstlisting}

\textbf{Q7.} Add comments to this complete build-and-verify workflow script.

\begin{lstlisting}[style=Alabaster, language=bash]
cd ~/Documents/3d_projects/my_cube
3dm build
3dm info
3dm info > model_report.txt
notepad.exe model_report.txt

\end{lstlisting}

\textbf{Q8.} Add comments explaining what each section of this installation script does (Linux/macOS).

\begin{lstlisting}[style=Alabaster, language=bash]
curl -fsSL https://get.3dmake.dev | bash
3dm --version
openscad --version
mkdir -p ~/Documents/3d_projects/my_cube
cd ~/Documents/3d_projects/my_cube
3dm new
ls -1

\end{lstlisting}

\textbf{Q9.} Add comments explaining the FDM pipeline stages referenced in this parameter file.

\begin{lstlisting}[style=Alabaster, language=openscad]
// Stage 1 — parametric source
width  = 40;
depth  = 30;
height = 15;
infill = 20;        // percent
layer_h = 0.20;     // mm

// Stage 2 — geometry (STL)
cube([width, depth, height]);

\end{lstlisting}

\textbf{Q10.} Add comments to this cost-estimation calculation.

\begin{lstlisting}[style=Alabaster, language=bash]
VOLUME=6000
DENSITY=1.24
SPOOL_PRICE=20
SPOOL_WEIGHT=1000

MASS=$(echo "scale=4; ($VOLUME / 1000) * $DENSITY" | bc)
COST=$(echo "scale=4; $MASS * ($SPOOL_PRICE / $SPOOL_WEIGHT)" | bc)
echo "Mass: ${MASS} g"
echo "Cost: \$${COST}"

\end{lstlisting}

\begin{center}\rule{0.5\linewidth}{0.5pt}\end{center}

\subsection*{3dMake --- Lesson 2: Geometric Primitives and CSG}\label{docs__pandoc__latex__src__advanced_study_guide.md__3dmake--lesson-2-geometric-primitives-and-csg}

\subsubsection*{Fix the Code (Questions 1--5)}\label{docs__pandoc__latex__src__advanced_study_guide.md__fix-the-code-questions-15-25}

\textbf{Q1.} The \texttt{difference()} operation is missing the 0.001 mm offset rule. Add the correct offsets.

\begin{lstlisting}[style=Alabaster, language=openscad]
difference() {
  cube([20, 20, 10]);
  translate([10, 10, 0])
    cylinder(h=10, r=5, $fn=32);
}

\end{lstlisting}

\textbf{Q2.} The children are in the wrong order for \texttt{difference()} --- the cutting tool is listed first.

\begin{lstlisting}[style=Alabaster, language=openscad]
difference() {
  cylinder(h=20, r=5, $fn=32);   // intended cutter
  cube([20, 20, 20]);             // intended base
}

\end{lstlisting}

\textbf{Q3.} The sphere is given an extremely low \texttt{\$fn} value that will produce a faceted, unusable mesh for printing.

\begin{lstlisting}[style=Alabaster, language=openscad]
sphere(r=15, $fn=3);

\end{lstlisting}

\textbf{Q4.} The \texttt{polygon()} points are listed in clockwise order (will produce an inverted or empty shape). Correct the winding order.

\begin{lstlisting}[style=Alabaster, language=openscad]
linear_extrude(height=5)
  polygon([[0,30],[0,0],[40,0],[40,8],[8,8],[8,30]]);

\end{lstlisting}

\textbf{Q5.} The \texttt{\#} debug modifier is left in the code before export. Identify and remove it.

\begin{lstlisting}[style=Alabaster, language=openscad]
difference() {
  cube([30, 30, 20]);
  # translate([15, 15, -0.001])
    cylinder(h=20.002, r=8, $fn=32);
}

\end{lstlisting}

\begin{center}\rule{0.5\linewidth}{0.5pt}\end{center}

\subsubsection*{Comment the Code (Questions 6--10)}\label{docs__pandoc__latex__src__advanced_study_guide.md__comment-the-code-questions-610-25}

\textbf{Q6.} Add comments explaining each primitive and its key parameters.

\begin{lstlisting}[style=Alabaster, language=openscad]
cube([30, 20, 10]);

translate([50, 0, 10])
  sphere(r=10, $fn=32);

translate([80, 0, 0])
  cylinder(h=20, r=8, $fn=32);

translate([110, 0, 0])
  cylinder(h=20, r1=10, r2=3, $fn=32);

\end{lstlisting}

\textbf{Q7.} Add comments to this \texttt{difference()} operation explaining each line\textquotesingle s purpose.

\begin{lstlisting}[style=Alabaster, language=openscad]
difference() {
  cube([20, 20, 20]);
  translate([10, 10, -0.001])
    cylinder(h=20.002, r=5, $fn=32);
}

\end{lstlisting}

\textbf{Q8.} Add comments explaining why the 0.001 offset is used and what each value means.

\begin{lstlisting}[style=Alabaster, language=openscad]
difference() {
  cube([20, 20, 10]);
  translate([1, 1, -0.001])
    cube([18, 18, 10.002]);
}

\end{lstlisting}

\textbf{Q9.} Add comments to this \texttt{polygon()} L-bracket extrusion.

\begin{lstlisting}[style=Alabaster, language=openscad]
bracket_profile = [
  [0,  0],
  [40, 0],
  [40, 8],
  [8,  8],
  [8,  30],
  [0,  30]
];

linear_extrude(height=5)
  polygon(bracket_profile);

\end{lstlisting}

\textbf{Q10.} Add comments explaining each debug modifier and when to use it.

\begin{lstlisting}[style=Alabaster, language=openscad]
difference() {
  cube([30, 30, 20]);

  # translate([15, 15, -0.001])
    cylinder(h=20.002, r=8, $fn=32);
}

% translate([0, 0, 20])
  cube([30, 30, 5]);

! sphere(r=10, $fn=32);

* cube([5, 5, 5]);

\end{lstlisting}

\begin{center}\rule{0.5\linewidth}{0.5pt}\end{center}

\subsection*{3dMake --- Lesson 3: Parametric Architecture and Modular Libraries}\label{docs__pandoc__latex__src__advanced_study_guide.md__3dmake--lesson-3-parametric-architecture-and-modular-libraries}

\subsubsection*{Fix the Code (Questions 1--5)}\label{docs__pandoc__latex__src__advanced_study_guide.md__fix-the-code-questions-15-26}

\textbf{Q1.} This OpenSCAD code attempts to use Python-style variable reassignment, which doesn\textquotesingle t work in OpenSCAD. Rewrite it correctly.

\begin{lstlisting}[style=Alabaster, language=openscad]
x = 5;
x = x + 3;
cube([x, x, 5]);

\end{lstlisting}

\textbf{Q2.} The \texttt{assert()} call uses the wrong comparison --- it accepts negative values.

\begin{lstlisting}[style=Alabaster, language=openscad]
module safe_box(w, d, h) {
  assert(w >= 0, "width must be non-negative");
  cube([w, d, h]);
}
safe_box(-10, 20, 5);

\end{lstlisting}

\textbf{Q3.} The module uses \texttt{include} to load a library when it should use \texttt{use} (the library file contains top-level test geometry that will pollute the scene).

\begin{lstlisting}[style=Alabaster, language=openscad]
include <../lib/hardware.scad>
mounting_hole(r=2, depth=5);

\end{lstlisting}

\textbf{Q4.} The DRY violation: the mounting hole is copy-pasted instead of using a module. Refactor it.

\begin{lstlisting}[style=Alabaster, language=openscad]
difference() {
  cube([60, 40, 5]);
  translate([5,  5,  -0.001]) cylinder(r=2, h=5.002, $fn=16);
  translate([55, 5,  -0.001]) cylinder(r=2, h=5.002, $fn=16);
  translate([5,  35, -0.001]) cylinder(r=2, h=5.002, $fn=16);
  translate([55, 35, -0.001]) cylinder(r=2, h=5.002, $fn=16);
}

\end{lstlisting}

\textbf{Q5.} The \texttt{for} loop uses wrong syntax for generating geometry in OpenSCAD.

\begin{lstlisting}[style=Alabaster, language=openscad]
for i in [0:3] {
  translate([i * 15, 0, 0])
    cylinder(r=3, h=10, $fn=16);
}

\end{lstlisting}

\begin{center}\rule{0.5\linewidth}{0.5pt}\end{center}

\subsubsection*{Comment the Code (Questions 6--10)}\label{docs__pandoc__latex__src__advanced_study_guide.md__comment-the-code-questions-610-26}

\textbf{Q6.} Add comments explaining OpenSCAD\textquotesingle s variable scoping rules as shown in this code.

\begin{lstlisting}[style=Alabaster, language=openscad]
x = 5;
echo("x =", x);
x = 10;

base_size   = 20;
padded_size = base_size + 10;
doubled     = base_size * 2;
echo(base_size, padded_size, doubled);

\end{lstlisting}

\textbf{Q7.} Add comments explaining every part of this validated module.

\begin{lstlisting}[style=Alabaster, language=openscad]
module safe_cylinder(r, h, fn=32) {
  assert(is_num(r) && r > 0,
    str("radius must be positive, got: ", r));
  assert(is_num(h) && h > 0,
    str("height must be positive, got: ", h));
  assert(fn >= 3,
    str("fn must be >= 3, got: ", fn));
  cylinder(r=r, h=h, $fn=fn);
}

\end{lstlisting}

\textbf{Q8.} Add comments to this DRY-compliant mounting hole module and its usage.

\begin{lstlisting}[style=Alabaster, language=openscad]
module mounting_hole(r=2, depth=5) {
  translate([0, 0, -0.001])
    cylinder(r=r, h=depth + 0.002, $fn=16);
}

difference() {
  cube([60, 40, 5]);
  translate([5,  5,  0]) mounting_hole();
  translate([55, 5,  0]) mounting_hole();
  translate([5,  35, 0]) mounting_hole();
  translate([55, 35, 0]) mounting_hole();
}

\end{lstlisting}

\textbf{Q9.} Add comments to this \texttt{use} vs \texttt{include} example showing correct library import.

\begin{lstlisting}[style=Alabaster, language=openscad]
use <../lib/hardware.scad>

mounting_hole(r=2, depth=5);
translate([20, 0, 0]) mounting_hole(r=3, depth=8);

\end{lstlisting}

\textbf{Q10.} Add comments explaining each line of this \texttt{for}-loop geometry generator.

\begin{lstlisting}[style=Alabaster, language=openscad]
for (i = [0:5]) {
  angle = i * 60;
  translate([cos(angle) * 25, sin(angle) * 25, 0])
    cylinder(r=3, h=5, $fn=16);
}

\end{lstlisting}

\begin{center}\rule{0.5\linewidth}{0.5pt}\end{center}

\subsection*{3dMake --- Lesson 4: AI-Enhanced Verification and Multimodal Feedback}\label{docs__pandoc__latex__src__advanced_study_guide.md__3dmake--lesson-4-ai-enhanced-verification-and-multimodal-feedback}

\subsubsection*{Fix the Code (Questions 1--5)}\label{docs__pandoc__latex__src__advanced_study_guide.md__fix-the-code-questions-15-27}

\textbf{Q1.} The verification sequence uses an AI description as the primary check instead of the deterministic bounding box. Rewrite it in the correct order.

\begin{lstlisting}[style=Alabaster, language=bash]
3dm build
3dm info --view front
# Accept design if description says "looks like a bracket"
# Skip bounding box check

\end{lstlisting}

\textbf{Q2.} The \texttt{3dm\ info} command is saved to a file using the wrong redirection for PowerShell.

\begin{lstlisting}[style=Alabaster, language=powershell]
3dm info > description.txt -Encoding UTF8

\end{lstlisting}

\textbf{Q3.} The AI code verification checklist is missing the 0.001 offset check. Add the missing item in the right position.

\begin{lstlisting}[style=Alabaster]
[ ] 1. SYNTAX — 3dm build compiles without error
[ ] 2. DIMENSIONS — bounding box matches design intent
[ ] 3. PARAMETER BEHAVIOR — test extreme values
[ ] 4. MANIFOLD CHECK — open in slicer, no errors
[ ] 5. PRINT FEASIBILITY — wall thickness >= 1.2mm

\end{lstlisting}

\textbf{Q4.} The verification log records an AI description as proof of precise dimensions. Correct the methodology.

\begin{lstlisting}[style=Alabaster]
Feature: USB-C slot
  Designed: 9.5mm wide × 3.5mm tall
  Verification: AI description says "rectangular cutout approximately 10mm wide" ✓
  Status: VERIFIED

\end{lstlisting}

\textbf{Q5.} The isolation technique for a complex model uses the wrong modifier to render only one module.

\begin{lstlisting}[style=Alabaster, language=openscad]
full_assembly();
% snap_clip();

\end{lstlisting}

\begin{center}\rule{0.5\linewidth}{0.5pt}\end{center}

\subsubsection*{Comment the Code (Questions 6--10)}\label{docs__pandoc__latex__src__advanced_study_guide.md__comment-the-code-questions-610-27}

\textbf{Q6.} Add comments explaining each \texttt{3dm\ info} variant and when to use it.

\begin{lstlisting}[style=Alabaster, language=bash]
3dm build
3dm info
3dm info --view front
3dm info --view top
3dm info --view isometric
3dm info --output description.txt

\end{lstlisting}

\textbf{Q7.} Add comments to this verification escalation chain.

\begin{lstlisting}[style=Alabaster, language=bash]
3dm info
# Level 1 check passed — proceed to Level 2

# Open slicer manually on build/main.stl
# Step through layers looking for issues

# Level 2 check passed — proceed to Level 3

# Level 3: isolate module

\end{lstlisting}

\textbf{Q8.} Add comments explaining each step of this isolation-and-verify workflow.

\begin{lstlisting}[style=Alabaster, language=openscad]
// full_assembly();
// base_plate();
snap_clip();

\end{lstlisting}

\begin{lstlisting}[style=Alabaster, language=bash]
3dm build
3dm info --view isometric
3dm info --view front

\end{lstlisting}

\textbf{Q9.} Add comments explaining the two types of output visible in this \texttt{3dm\ info} result.

\begin{lstlisting}[style=Alabaster]
File:         build/main.stl
Bounding Box: 30.0 x 20.0 x 10.0 mm
Volume:       6000.0 mm³
Triangles:    12
Description:  A small rectangular box measuring 30mm wide by 20mm deep by 10mm tall.

\end{lstlisting}

\textbf{Q10.} Add comments to this AI code verification checklist explaining why each check exists.

\begin{lstlisting}[style=Alabaster]
[ ] 1. SYNTAX
[ ] 2. DIMENSIONS
[ ] 3. BOOLEAN OPERATIONS — check every difference() for 0.001 offset
[ ] 4. PARAMETER BEHAVIOR — test extreme values
[ ] 5. MANIFOLD CHECK
[ ] 6. PRINT FEASIBILITY

\end{lstlisting}

\begin{center}\rule{0.5\linewidth}{0.5pt}\end{center}

\subsection*{3dMake --- Lesson 5: Safety, Filaments, and Print Environment}\label{docs__pandoc__latex__src__advanced_study_guide.md__3dmake--lesson-5-safety-filaments-and-print-environment}

\subsubsection*{Fix the Code (Questions 1--5)}\label{docs__pandoc__latex__src__advanced_study_guide.md__fix-the-code-questions-15-28}

\textbf{Q1.} The cost-estimation formula uses the wrong density for PLA.

\begin{lstlisting}[style=Alabaster, language=openscad]
// Cost estimate for PLA
volume_mm3 = 8500;
density    = 1.05;   // g/cm³ — WRONG material
spool_cost = 20;
spool_g    = 1000;

mass = (volume_mm3 / 1000) * density;
cost = mass * (spool_cost / spool_g);
echo("Est. cost: $", cost);

\end{lstlisting}

\textbf{Q2.} The safety checklist incorrectly marks ABS as safe to print in an open classroom without a fume enclosure.

\begin{lstlisting}[style=Alabaster]
Pre-print safety check:
[x] Filament loaded: ABS
[x] Ventilation: window open
[x] Ready to print in open classroom

\end{lstlisting}

\textbf{Q3.} The thermal runaway note is wrong --- this protection must be enabled, not optional.

\begin{lstlisting}[style=Alabaster]
Printer setup checklist:
[ ] Bed level confirmed
[ ] Nozzle clean
[ ] Thermal runaway: optional, enable if desired

\end{lstlisting}

\textbf{Q4.} The post-print removal instruction is dangerous --- it says to remove the part immediately at full temperature.

\begin{lstlisting}[style=Alabaster]
Post-print steps:
1. Print finishes.
2. Immediately remove part from heated bed with bare hands.
3. Store filament in a bag.

\end{lstlisting}

\textbf{Q5.} The filament density table has PLA and ABS swapped. Correct the table.

\begin{lstlisting}[style=Alabaster]
Material  | Density (g/cm³)
----------|----------------
PLA       | 1.05
ABS       | 1.24
PETG      | 1.27
TPU       | 1.20

\end{lstlisting}

\begin{center}\rule{0.5\linewidth}{0.5pt}\end{center}

\subsubsection*{Comment the Code (Questions 6--10)}\label{docs__pandoc__latex__src__advanced_study_guide.md__comment-the-code-questions-610-28}

\textbf{Q6.} Add comments explaining each line of this pre-print safety checklist script.

\begin{lstlisting}[style=Alabaster, language=bash]
echo "=== Pre-Print Safety Check ===" > safety_log.txt
echo "Date: $(date)" >> safety_log.txt
echo "Filament: PLA" >> safety_log.txt
echo "Ventilation confirmed: YES" >> safety_log.txt
echo "Fire extinguisher present: YES" >> safety_log.txt
echo "Thermal runaway enabled: YES" >> safety_log.txt
echo "First layer watch timer: 15 min" >> safety_log.txt

\end{lstlisting}

\textbf{Q7.} Add comments to this cost-estimation snippet explaining each formula step.

\begin{lstlisting}[style=Alabaster, language=openscad]
volume_mm3    = 8500;
density_pla   = 1.24;
spool_price   = 20;
spool_weight  = 1000;

mass = (volume_mm3 / 1000) * density_pla;
cost = mass * (spool_price / spool_weight);

echo("Volume (mm³):", volume_mm3);
echo("Mass (g):", mass);
echo("Est. cost ($):", cost);

\end{lstlisting}

\textbf{Q8.} Add comments explaining each row of this filament selection table embedded in OpenSCAD comments.

\begin{lstlisting}[style=Alabaster, language=openscad]
// FILAMENT REFERENCE
// PLA   — density 1.24 g/cm³ — easiest, classroom default
// PETG  — density 1.27 g/cm³ — tougher, mild heat, moderate stringing
// ABS   — density 1.05 g/cm³ — high heat, fumes — enclosure required
// TPU   — density 1.20 g/cm³ — flexible, direct-drive extruder only
material = "PLA";

\end{lstlisting}

\textbf{Q9.} Add comments to this filament-selection module.

\begin{lstlisting}[style=Alabaster, language=openscad]
function filament_density(material) =
  (material == "PLA")  ? 1.24 :
  (material == "PETG") ? 1.27 :
  (material == "ABS")  ? 1.05 :
  (material == "TPU")  ? 1.20 :
  undef;

volume   = 6000;
mat      = "PLA";
density  = filament_density(mat);
mass     = (volume / 1000) * density;
echo("Material:", mat, "| Mass:", mass, "g");

\end{lstlisting}

\textbf{Q10.} Add comments explaining each step of this first-layer monitoring checklist embedded in an automation log.

\begin{lstlisting}[style=Alabaster, language=bash]
echo "=== First Layer Monitor ===" >> print_log.txt
echo "Layer 1: bed adhesion — check in 2 min" >> print_log.txt
echo "Layer 2: no lifting at corners" >> print_log.txt
echo "Layer 3: temperature stable" >> print_log.txt
echo "Layer 3: no grinding sounds" >> print_log.txt
echo "First layer check complete — continue monitoring every 15 min" >> print_log.txt

\end{lstlisting}

\begin{center}\rule{0.5\linewidth}{0.5pt}\end{center}

\subsection*{3dMake --- Lesson 6: Practical 3dm Commands and Text Embossing}\label{docs__pandoc__latex__src__advanced_study_guide.md__3dmake--lesson-6-practical-3dm-commands-and-text-embossing}

\subsubsection*{Fix the Code (Questions 1--5)}\label{docs__pandoc__latex__src__advanced_study_guide.md__fix-the-code-questions-15-29}

\textbf{Q1.} The embossed text is positioned below the top face of the base plate instead of on top of it.

\begin{lstlisting}[style=Alabaster, language=openscad]
base_h = 5;
cube([80, 15, base_h]);
translate([40, 7.5, 0])
  linear_extrude(height=1.5)
    text("FRONT", size=6, font="Liberation Sans:style=Bold",
         halign="center", valign="center");

\end{lstlisting}

\textbf{Q2.} The debossed text cut is not deep enough --- the extrude height doesn\textquotesingle t reach past the top face.

\begin{lstlisting}[style=Alabaster, language=openscad]
difference() {
  cube([80, 15, 5]);
  translate([40, 7.5, 4.5])
    linear_extrude(height=0.4)
      text("BACK", size=6, font="Liberation Sans:style=Bold",
           halign="center", valign="center");
}

\end{lstlisting}

\textbf{Q3.} The font is hardcoded in three separate places instead of using the single-constant pattern. Refactor it.

\begin{lstlisting}[style=Alabaster, language=openscad]
cube([80, 15, 5]);
translate([40, 7.5, 5])
  linear_extrude(height=1.5)
    text("FRONT", size=6, font="Liberation Sans:style=Bold",
         halign="center", valign="center");

translate([0, 20, 0]) cube([80, 15, 5]);
translate([40, 27.5, 10])
  linear_extrude(height=1.5)
    text("BACK", size=6, font="Liberation Sans:style=Bold",
         halign="center", valign="center");

\end{lstlisting}

\textbf{Q4.} The \texttt{3dm\ slice} command is run before \texttt{3dm\ build}, meaning it will slice the previous (stale) STL file.

\begin{lstlisting}[style=Alabaster, language=bash]
3dm slice
3dm build
3dm info

\end{lstlisting}

\textbf{Q5.} The \texttt{let()} block uses wrong assignment syntax inside the expression.

\begin{lstlisting}[style=Alabaster, language=openscad]
module label_plate(msg, w=80, d=15, h=5) {
  let(text_size == h * 0.6) {
    cube([w, d, h]);
    translate([w/2, d/2, h])
      linear_extrude(height=1.5)
        text(msg, size=text_size, halign="center", valign="center");
  }
}

\end{lstlisting}

\begin{center}\rule{0.5\linewidth}{0.5pt}\end{center}

\subsubsection*{Comment the Code (Questions 6--10)}\label{docs__pandoc__latex__src__advanced_study_guide.md__comment-the-code-questions-610-29}

\textbf{Q6.} Add comments explaining every parameter of the \texttt{text()} function call.

\begin{lstlisting}[style=Alabaster, language=openscad]
linear_extrude(height=1.5)
  text(
    text      = "3dMake",
    size      = 10,
    font      = "Liberation Sans",
    halign    = "center",
    valign    = "baseline",
    spacing   = 1.0,
    direction = "ltr",
    language  = "en",
    script    = "latin"
  );

\end{lstlisting}

\textbf{Q7.} Add comments explaining each step of this embossed-plate module.

\begin{lstlisting}[style=Alabaster, language=openscad]
LABEL_FONT = "Liberation Sans:style=Bold";

module embossed_plate(message, w=80, d=15, h=5, text_size=8, relief=1.5) {
  cube([w, d, h]);
  translate([w/2, d/2, h])
    linear_extrude(height=relief)
      text(message, font=LABEL_FONT, size=text_size,
           halign="center", valign="center");
}

embossed_plate("FRONT");
translate([0, 20, 0]) embossed_plate("BACK");

\end{lstlisting}

\textbf{Q8.} Add comments explaining each command in the full \texttt{3dm} suite.

\begin{lstlisting}[style=Alabaster, language=bash]
3dm build
3dm build --clean
3dm build --watch
3dm info
3dm info --view front
3dm info --view top,front,isometric
3dm preview
3dm orient
3dm slice

\end{lstlisting}

\textbf{Q9.} Add comments to this debossed text implementation.

\begin{lstlisting}[style=Alabaster, language=openscad]
difference() {
  cube([80, 15, 5]);
  translate([40, 7.5, 5 - 0.8])
    linear_extrude(height=1.0)
      text("BACK", size=6, font="Liberation Sans:style=Bold",
           halign="center", valign="center");
}

\end{lstlisting}

\textbf{Q10.} Add comments to this variant-generation automation script.

\begin{lstlisting}[style=Alabaster, language=bash]
#!/bin/bash
LABELS=("FRONT" "BACK" "LEFT" "RIGHT" "TOP" "BOTTOM")
mkdir -p build/labels

for LABEL in "${LABELS[@]}"; do
    OUT="build/labels/plate_${LABEL}.stl"
    openscad -D "label=\"${LABEL}\"" -o "$OUT" src/label_plate.scad
    echo "Built: $OUT"
done
echo "All label variants built."

\end{lstlisting}

\begin{center}\rule{0.5\linewidth}{0.5pt}\end{center}

\subsection*{3dMake --- Lesson 7: Parametric Transforms and the Phone Stand}\label{docs__pandoc__latex__src__advanced_study_guide.md__3dmake--lesson-7-parametric-transforms-and-the-phone-stand}

\subsubsection*{Fix the Code (Questions 1--5)}\label{docs__pandoc__latex__src__advanced_study_guide.md__fix-the-code-questions-15-30}

\textbf{Q1.} The transform order is wrong --- the cube will orbit the origin instead of rotating in place.

\begin{lstlisting}[style=Alabaster, language=openscad]
rotate([0, 0, 45]) translate([20, 0, 0]) cube(5);

\end{lstlisting}

\emph{(Hint: to rotate in place then move, which should come first in reading order?)}

\textbf{Q2.} The \texttt{atan} function is used instead of \texttt{atan2}, which will give wrong results in some quadrants.

\begin{lstlisting}[style=Alabaster, language=openscad]
dx = -5;
dy = 5;
angle = atan(dy / dx);
echo("Angle:", angle);

\end{lstlisting}

\textbf{Q3.} The \texttt{minkowski()} call uses \texttt{\$fn=64} on the secondary sphere during development, making renders extremely slow.

\begin{lstlisting}[style=Alabaster, language=openscad]
module rounded_box(w, d, h, r=3) {
  minkowski() {
    cube([w - 2*r, d - 2*r, h - r]);
    sphere(r=r, $fn=64);
  }
}
rounded_box(40, 30, 20);

\end{lstlisting}

\textbf{Q4.} The \texttt{mirror()} call uses the wrong axis vector --- it reflects across the XZ plane instead of the YZ plane.

\begin{lstlisting}[style=Alabaster, language=openscad]
module ear_peg() {
  translate([20, 0, 0]) cylinder(r=4, h=8, $fn=32);
}
ear_peg();
mirror([0, 1, 0]) ear_peg();

\end{lstlisting}

\textbf{Q5.} The \texttt{\$preview} variable is used backwards --- the slow high-quality render runs during preview and the fast version runs during final build.

\begin{lstlisting}[style=Alabaster, language=openscad]
$fn = $preview ? 64 : 16;
sphere(r=20);

\end{lstlisting}

\begin{center}\rule{0.5\linewidth}{0.5pt}\end{center}

\subsubsection*{Comment the Code (Questions 6--10)}\label{docs__pandoc__latex__src__advanced_study_guide.md__comment-the-code-questions-610-30}

\textbf{Q6.} Add comments explaining the transform order rule demonstrated in these two examples.

\begin{lstlisting}[style=Alabaster, language=openscad]
translate([10, 0, 0]) rotate([0, 0, 45]) cube(5);

rotate([0, 0, 45]) translate([10, 0, 0]) cube(5);

\end{lstlisting}

\textbf{Q7.} Add comments explaining each trig function used in this circular placement loop.

\begin{lstlisting}[style=Alabaster, language=openscad]
for (i = [0:5]) {
  angle = i * 60;
  translate([cos(angle) * 25, sin(angle) * 25, 0])
    cylinder(r=3, h=5, $fn=16);
}

\end{lstlisting}

\textbf{Q8.} Add comments explaining each module and parameter in the phone stand assembly.

\begin{lstlisting}[style=Alabaster, language=openscad]
phone_w = 75;
phone_d = 9;
angle   = 65;

lip_h      = 15;
base_h     = 5;
cradle_wall= 3;
r_fillet   = 3;
base_d     = (phone_d + 10) / cos(90 - angle);

module base_plate() {
  flat_rounded_cube(phone_w + 20, base_d + 10, base_h, r_fillet);
}

module back_support() {
  rotate([angle - 90, 0, 0])
    flat_rounded_cube(phone_w + 6, cradle_wall, 60, r_fillet);
}

module lip() {
  rotate([angle - 90, 0, 0])
    translate([0, -lip_h, 0])
      flat_rounded_cube(phone_w + 6, lip_h + cradle_wall, cradle_wall, r_fillet);
}

base_plate();
translate([0, 0, base_h]) { back_support(); lip(); }

\end{lstlisting}

\textbf{Q9.} Add comments explaining how \texttt{\$preview} enables fast development.

\begin{lstlisting}[style=Alabaster, language=openscad]
$fn = $preview ? 16 : 64;

module flat_rounded_cube(w, d, h, r=3) {
  if ($preview) {
    color("orange", 0.6) cube([w, d, h]);
  } else {
    minkowski() {
      cube([w - 2*r, d - 2*r, h]);
      cylinder(r=r, h=0.01, $fn=24);
    }
  }
}

flat_rounded_cube(40, 30, 20, r=4);

\end{lstlisting}

\textbf{Q10.} Add comments to this \texttt{mirror()} symmetry example.

\begin{lstlisting}[style=Alabaster, language=openscad]
module mounting_clip() {
  difference() {
    cube([15, 8, 10]);
    translate([7.5, -0.001, 2])
      cylinder(r=3, h=8.002, $fn=32);
  }
}

translate([0, 20, 0]) {
  mounting_clip();
  translate([15, 0, 0]) mirror([1, 0, 0]) mounting_clip();
}

\end{lstlisting}

\begin{center}\rule{0.5\linewidth}{0.5pt}\end{center}

\subsection*{3dMake --- Lesson 8: Advanced Parametric Design and Interlocking Features}\label{docs__pandoc__latex__src__advanced_study_guide.md__3dmake--lesson-8-advanced-parametric-design-and-interlocking-features}

\subsubsection*{Fix the Code (Questions 1--5)}\label{docs__pandoc__latex__src__advanced_study_guide.md__fix-the-code-questions-15-31}

\textbf{Q1.} The tolerance constants are hardcoded inline instead of defined as named constants at the top. Refactor to use constants.

\begin{lstlisting}[style=Alabaster, language=openscad]
module peg(nominal=10, h=8) {
  cylinder(d=nominal - 0.3, h=h, $fn=32);
}
module socket(nominal=10, plate_t=3) {
  difference() {
    cylinder(d=nominal + 8, h=plate_t, $fn=32);
    translate([0, 0, -0.001])
      cylinder(d=nominal, h=plate_t + 0.002, $fn=32);
  }
}

\end{lstlisting}

\textbf{Q2.} The snap-fit beam parameters are specified for PLA but the comment says PETG. Fix the comment to match the geometry (do not change the numbers --- assume PLA values are correct).

\begin{lstlisting}[style=Alabaster, language=openscad]
// PETG snap arm — thicker beam for stiffer material
module snap_arm_pla(beam_l=18, beam_t=1.8, beam_w=8) {
  cube([beam_w, beam_l, beam_t]);
  translate([0, beam_l, 0])
    cube([beam_w, 3, 2.5]);
}

\end{lstlisting}

\textbf{Q3.} The stacking-bin lip is positioned at the wrong Z height --- it should start at \texttt{bin\_h}, not \texttt{bin\_h\ -\ lip\_h}.

\begin{lstlisting}[style=Alabaster, language=openscad]
bin_h  = 40;
lip_h  = 5;
bin_w  = 80;
bin_d  = 60;
lip_clear = 0.25;

module stacking_lip() {
  translate([lip_clear, lip_clear, bin_h - lip_h])
    cube([bin_w - 2*lip_clear, bin_d - 2*lip_clear, lip_h]);
}

\end{lstlisting}

\textbf{Q4.} The M3 insert pocket extends above the host material (depth is too large for a 10mm plate).

\begin{lstlisting}[style=Alabaster, language=openscad]
difference() {
  cube([40, 30, 10]);
  translate([10, 15, -0.001]) cylinder(r=2.35, h=20, $fn=16);
  translate([30, 15, -0.001]) cylinder(r=2.35, h=20, $fn=16);
}

\end{lstlisting}

\textbf{Q5.} The dovetail female profile is given a negative clearance, which would make the slot smaller than the male part.

\begin{lstlisting}[style=Alabaster, language=openscad]
clearance = -0.2;

module dovetail_female(w_top=8, w_bot=12, h=5, c=clearance) {
  profile = [
    [-(w_bot/2 + c), -0.001], [ w_bot/2 + c, -0.001],
    [ w_top/2 + c,   h + c ], [-(w_top/2 + c), h + c]
  ];
  rotate([90, 0, 0])
    linear_extrude(50)
      polygon(profile);
}

\end{lstlisting}

\begin{center}\rule{0.5\linewidth}{0.5pt}\end{center}

\subsubsection*{Comment the Code (Questions 6--10)}\label{docs__pandoc__latex__src__advanced_study_guide.md__comment-the-code-questions-610-31}

\textbf{Q6.} Add comments explaining each tolerance constant and how it was determined.

\begin{lstlisting}[style=Alabaster, language=openscad]
PRESS_FIT_C = 0.2;
SLIP_FIT_C  = 0.3;
LOOSE_FIT_C = 0.45;
SNAP_C      = PRESS_FIT_C;

\end{lstlisting}

\textbf{Q7.} Add comments explaining each step of the tolerance test coupon.

\begin{lstlisting}[style=Alabaster, language=openscad]
nominal     = 10;
test_h      = 8;
plate_t     = 3;
socket_d_out= nominal + 8;
clearances  = [0.0, 0.1, 0.2, 0.3, 0.4];

module test_peg(clearance) {
  color("SteelBlue")
    cylinder(d=nominal - clearance, h=test_h, $fn=32);
}

module test_socket() {
  color("Coral")
  difference() {
    cylinder(d=socket_d_out, h=plate_t, $fn=32);
    translate([0, 0, -0.001])
      cylinder(d=nominal, h=plate_t + 0.002, $fn=32);
  }
}

\end{lstlisting}

\textbf{Q8.} Add comments explaining the physics behind each snap-arm parameter.

\begin{lstlisting}[style=Alabaster, language=openscad]
module snap_arm(beam_l=18, beam_t=1.8, beam_w=8, hook_h=1.5) {
  cube([beam_w, beam_l, beam_t]);
  translate([0, beam_l, 0])
    cube([beam_w, 3, hook_h + beam_t]);
}

\end{lstlisting}

\textbf{Q9.} Add comments to this dovetail joint module explaining the geometry.

\begin{lstlisting}[style=Alabaster, language=openscad]
module dovetail_male(l=50, w_top=8, w_bot=12, h=5) {
  profile = [
    [-w_bot/2, 0], [ w_bot/2, 0],
    [ w_top/2, h], [-w_top/2, h]
  ];
  rotate([90, 0, 0])
    linear_extrude(l)
      polygon(profile);
}

\end{lstlisting}

\textbf{Q10.} Add comments to this chamfered-box module explaining the geometry operation.

\begin{lstlisting}[style=Alabaster, language=openscad]
module chamfered_box(w, d, h, c=2) {
  difference() {
    cube([w, d, h]);
    translate([0, 0, h - c])
      difference() {
        cube([w, d, c + 0.001]);
        translate([c, c, 0])
          cube([w - 2*c, d - 2*c, c + 0.001]);
      }
  }
}
chamfered_box(50, 40, 20, c=3);

\end{lstlisting}

\begin{center}\rule{0.5\linewidth}{0.5pt}\end{center}

\subsection*{3dMake --- Lesson 9: Automation and 3dm Workflows}\label{docs__pandoc__latex__src__advanced_study_guide.md__3dmake--lesson-9-automation-and-3dm-workflows}

\subsubsection*{Fix the Code (Questions 1--5)}\label{docs__pandoc__latex__src__advanced_study_guide.md__fix-the-code-questions-15-32}

\textbf{Q1.} The bash script is missing all three safety flags. Add them in the correct location.

\begin{lstlisting}[style=Alabaster, language=bash]
#!/bin/bash
LOG="build/build.log"
mkdir -p build
openscad -o build/main.stl src/main.scad 2>&1 | tee -a "$LOG"
echo "Done"

\end{lstlisting}

\textbf{Q2.} The string parameter quoting is wrong --- OpenSCAD will receive the label without inner double quotes.

\begin{lstlisting}[style=Alabaster, language=bash]
LABEL="FRONT PANEL"
openscad -D "label=$LABEL" -o out.stl src/main.scad

\end{lstlisting}

\textbf{Q3.} The PowerShell script checks \texttt{\$LASTEXITCODE} but doesn\textquotesingle t actually stop on failure.

\begin{lstlisting}[style=Alabaster, language=powershell]
$output = & openscad -o build\main.stl src\main.scad 2>&1
$output | Write-Host
if ($LASTEXITCODE -ne 0) { Write-Host "Build failed" }
Write-Host "Continuing..."

\end{lstlisting}

\textbf{Q4.} The batch script counter increment uses wrong CMD syntax.

\begin{lstlisting}[style=Alabaster, language=cmd]
@echo off
set SUCCESS=0
set FAIL=0
for %%W in (60 80 100) do (
    openscad -D "width=%%W" -o build\out_%%W.stl src\main.scad
    if %errorlevel% equ 0 (
        set SUCCESS = %SUCCESS% + 1
    ) else (
        set FAIL = %FAIL% + 1
    )
)
echo %SUCCESS% succeeded, %FAIL% failed

\end{lstlisting}

\textbf{Q5.} The file-size check uses the wrong comparison operator --- it warns when the file is large instead of when it\textquotesingle s suspiciously small.

\begin{lstlisting}[style=Alabaster, language=bash]
STL_SIZE=$(stat -c%s "build/main.stl")
if [ "$STL_SIZE" -gt 1000 ]; then
  echo "WARNING: STL suspiciously small"
fi

\end{lstlisting}

\begin{center}\rule{0.5\linewidth}{0.5pt}\end{center}

\subsubsection*{Comment the Code (Questions 6--10)}\label{docs__pandoc__latex__src__advanced_study_guide.md__comment-the-code-questions-610-32}

\textbf{Q6.} Add comments to every significant line of this robust bash build script.

\begin{lstlisting}[style=Alabaster, language=bash]
#!/bin/bash
set -e
set -u
set -o pipefail

LOG_DIR="build/logs"
LOG_FILE="$LOG_DIR/build_$(date +%Y%m%d_%H%M%S).log"
mkdir -p "$LOG_DIR" build

log() {
  local level="$1"; shift
  echo "[$(date +%H:%M:%S)] [$level] $*" | tee -a "$LOG_FILE"
}

on_error() { log "ERROR" "Failed at line $1"; exit 1; }
trap 'on_error $LINENO' ERR

log "INFO" "=== Build started ==="
openscad -o build/main.stl src/main.scad 2>&1 | tee -a "$LOG_FILE"
STL_SIZE=$(stat -c%s "build/main.stl")
[ "$STL_SIZE" -gt 1000 ] || log "WARN" "STL small: ${STL_SIZE} bytes"
log "INFO" "Build OK: ${STL_SIZE} bytes"
3dm info 2>&1 | tee -a "$LOG_FILE"

\end{lstlisting}

\textbf{Q7.} Add comments explaining each string-quoting technique.

\begin{lstlisting}[style=Alabaster, language=bash]
LABEL="FRONT PANEL"

openscad -D 'label="FRONT"' -o out1.stl src/main.scad
openscad -D "label=\"${LABEL}\"" -o out2.stl src/main.scad
echo "openscad -D 'label=\"${LABEL}\"' -o out3.stl src/main.scad"

\end{lstlisting}

\textbf{Q8.} Add comments to this PowerShell batch-build script.

\begin{lstlisting}[style=Alabaster, language=powershell]
$outputDir = "build\variants"
New-Item -ItemType Directory -Force -Path $outputDir | Out-Null

$widths = @(60, 80, 100, 120)
$walls  = @(1.2, 1.6, 2.0)
$ok = 0; $fail = 0

foreach ($w in $widths) {
    foreach ($wt in $walls) {
        $name = "bracket_w${w}_wall${wt}"
        $out  = "$outputDir\${name}.stl"
        & openscad -D "width=$w" -D "wall=$wt" -o $out src\main.scad 2>$null
        if ($LASTEXITCODE -eq 0) { $ok++ } else { $fail++ }
    }
}
Write-Host "Done: $ok succeeded, $fail failed"

\end{lstlisting}

\textbf{Q9.} Add comments to this CI/CD GitHub Actions workflow.

\begin{lstlisting}[style=Alabaster]
name: Build STL Artifacts
on: [push]
jobs:
  build:
    runs-on: ubuntu-latest
    steps:
      - uses: actions/checkout@v3
      - name: Install OpenSCAD
        run: sudo apt-get install -y openscad
      - name: Install 3dm
        run: curl -fsSL https://get.3dmake.dev | bash
      - name: Build
        run: 3dm build
      - name: Verify
        run: 3dm info
      - name: Upload STL
        uses: actions/upload-artifact@v3
        with:
          name: stl-output
          path: build/*.stl

\end{lstlisting}

\textbf{Q10.} Add comments to this watch-mode automation script.

\begin{lstlisting}[style=Alabaster, language=bash]
#!/bin/bash
set -e
SRC="src/main.scad"
LAST_MOD=""

echo "Watching: $SRC"
while true; do
    CURRENT_MOD=$(stat -c %Y "$SRC" 2>/dev/null || echo "0")
    if [ "$CURRENT_MOD" != "$LAST_MOD" ]; then
        LAST_MOD="$CURRENT_MOD"
        echo "[$(date +%H:%M:%S)] Change detected — rebuilding..."
        3dm build && 3dm info || echo "Build failed"
    fi
    sleep 2
done

\end{lstlisting}

\begin{center}\rule{0.5\linewidth}{0.5pt}\end{center}

\subsection*{\texorpdfstring{3dMake --- Lesson 10: Advanced Techniques --- \texttt{hull()}, \texttt{projection()}, Non-Manifold Geometry}{3dMake --- Lesson 10: Advanced Techniques --- hull(), projection(), Non-Manifold Geometry}}\label{docs__pandoc__latex__src__advanced_study_guide.md__3dmake--lesson-10-advanced-techniques--hull-projection-non-manifold-geometry}

\subsubsection*{Fix the Code (Questions 1--5)}\label{docs__pandoc__latex__src__advanced_study_guide.md__fix-the-code-questions-15-33}

\textbf{Q1.} The \texttt{hull()} rounded box uses corner spheres, but they are placed at the wrong positions --- at (0,0) instead of offset by the radius.

\begin{lstlisting}[style=Alabaster, language=openscad]
r = 5;
w = 40; d = 30; h = 20;

hull() {
  for (x = [0, w]) {
    for (y = [0, d]) {
      translate([x, y, 0]) sphere(r=r, $fn=24);
      translate([x, y, h]) sphere(r=r, $fn=24);
    }
  }
}

\end{lstlisting}

\textbf{Q2.} The \texttt{minkowski()} call uses \texttt{\$fn=64} on the secondary sphere during development instead of using \texttt{\$preview}.

\begin{lstlisting}[style=Alabaster, language=openscad]
module rounded_box(w, d, h, r=4) {
  minkowski() {
    cube([w - 2*r, d - 2*r, h - r]);
    sphere(r=r, $fn=64);
  }
}
rounded_box(50, 40, 20, r=5);

\end{lstlisting}

\textbf{Q3.} The \texttt{projection()} cross-section is taken at the wrong height --- the model is not shifted to put the cut plane at the intended Z level.

\begin{lstlisting}[style=Alabaster, language=openscad]
module phone_stand_example() {
  cube([90, 70, 5]);
  translate([15, 20, 5]) cylinder(r=5, h=35, $fn=24);
}

linear_extrude(1)
  projection(cut=true)
    phone_stand_example();

\end{lstlisting}

\emph{(The cut should be at z=10 of the model, not z=0.)}

\textbf{Q4.} The non-manifold geometry fix is wrong --- the zero-thickness wall is left in place.

\begin{lstlisting}[style=Alabaster, language=openscad]
// WRONG: zero-thickness internal wall (non-manifold)
difference() {
  cube([20, 20, 20]);
  translate([10, 0, 0]) cube([0, 20, 20]);  // zero-width slice — non-manifold
}

// ATTEMPTED FIX (still wrong):
difference() {
  cube([20, 20, 20]);
  translate([10, 0, 0]) cube([0, 20, 20]);  // still zero width
}

\end{lstlisting}

\emph{(Provide a correct fix using a non-zero wall thickness.)}

\textbf{Q5.} The \texttt{surface()} call uses an incorrect file path that will silently produce no geometry.

\begin{lstlisting}[style=Alabaster, language=openscad]
scale([1, 1, 0.1])
  surface(file="heightmap/terrain.png", center=true, invert=false);

\end{lstlisting}

\emph{(The file path issue: if the file doesn\textquotesingle t exist, OpenSCAD produces empty geometry with no error. Write the correct verification approach using \texttt{3dm\ info} output.)}

\begin{center}\rule{0.5\linewidth}{0.5pt}\end{center}

\subsubsection*{Comment the Code (Questions 6--10)}\label{docs__pandoc__latex__src__advanced_study_guide.md__comment-the-code-questions-610-33}

\textbf{Q6.} Add comments explaining what \texttt{hull()} produces in each example.

\begin{lstlisting}[style=Alabaster, language=openscad]
hull() {
  cylinder(r=20, h=1, $fn=6);
  translate([0, 0, 30])
    sphere(r=8, $fn=32);
}

translate([60, 0, 0])
hull() {
  cylinder(r=6, h=2, $fn=32);
  translate([40, 0, 0])
    cylinder(r=1, h=2, $fn=16);
}

\end{lstlisting}

\textbf{Q7.} Add comments to this \texttt{\$preview}-guarded \texttt{minkowski()} module.

\begin{lstlisting}[style=Alabaster, language=openscad]
module rounded_box(w, d, h, r=4) {
  if ($preview) {
    color("orange", 0.6) cube([w, d, h]);
  } else {
    minkowski() {
      cube([w - 2*r, d - 2*r, h - r]);
      sphere(r=r, $fn=12);
    }
  }
}

rounded_box(50, 40, 20, r=5);

\end{lstlisting}

\textbf{Q8.} Add comments explaining each \texttt{projection()} mode.

\begin{lstlisting}[style=Alabaster, language=openscad]
module phone_stand_example() {
  cube([90, 70, 5]);
  translate([15, 20, 5]) cylinder(r=5, h=35, $fn=24);
}

phone_stand_example();

color("LightGreen")
translate([110, 0, 0])
  linear_extrude(1)
    projection(cut=false)
      phone_stand_example();

color("LightCoral")
translate([220, 0, 0])
  linear_extrude(1)
    projection(cut=true)
      translate([0, 0, -3])
        phone_stand_example();

\end{lstlisting}

\textbf{Q9.} Add comments to this non-manifold geometry fix.

\begin{lstlisting}[style=Alabaster, language=openscad]
// PROBLEM: co-planar faces
difference() {
  cube([20, 20, 10]);
  translate([5, 5, 0])
    cube([10, 10, 10]);
}

// FIX: 0.001 offset rule applied
difference() {
  cube([20, 20, 10]);
  translate([5, 5, -0.001])
    cube([10, 10, 10.002]);
}

\end{lstlisting}

\textbf{Q10.} Add comments to this data-driven design using \texttt{search()} and \texttt{lookup()}.

\begin{lstlisting}[style=Alabaster, language=openscad]
SIZES = [
  ["small",  20, 15, 8],
  ["medium", 35, 25, 12],
  ["large",  50, 40, 18]
];

function get_size(name, col) =
  SIZES[search([name], SIZES, 1, 0)[0]][col];

variant = "medium";
w = get_size(variant, 1);
d = get_size(variant, 2);
h = get_size(variant, 3);
cube([w, d, h]);
echo("Building:", variant, w, "x", d, "x", h);

\end{lstlisting}

\begin{center}\rule{0.5\linewidth}{0.5pt}\end{center}

\subsection*{3dMake --- Lesson 11: Stakeholder-Centric Design and the Final Project}\label{docs__pandoc__latex__src__advanced_study_guide.md__3dmake--lesson-11-stakeholder-centric-design-and-the-final-project}

\subsubsection*{Fix the Code (Questions 1--5)}\label{docs__pandoc__latex__src__advanced_study_guide.md__fix-the-code-questions-15-34}

\textbf{Q1.} The requirement mapping uses a non-descriptive variable name --- rename the variable so it traces directly to the functional requirement.

\begin{lstlisting}[style=Alabaster, language=openscad]
// FR1: Must fit standard bracelet diameters (50–70 mm)
x = 70;
inner_r = x / 2;
cylinder(r=inner_r + 2, h=10, $fn=64);

\end{lstlisting}

\textbf{Q2.} The acceptance criterion test is based on AI description instead of a measurable, deterministic check. Rewrite it.

\begin{lstlisting}[style=Alabaster]
Feature: bead hole diameter
Acceptance Test: run 3dm info and confirm description mentions "small hole"
Result: PASS

\end{lstlisting}

\textbf{Q3.} The bead module is missing an \texttt{assert()} for the cord-hole diameter --- small holes will fail to print cleanly. Add the appropriate assertion.

\begin{lstlisting}[style=Alabaster, language=openscad]
module bead_A(size=15, detail=32) {
  difference() {
    sphere(r=size/2, $fn=detail);
    cylinder(r=1.0, h=size + 1, center=true, $fn=16);
  }
}

\end{lstlisting}

\textbf{Q4.} The Technical Design Specification variable is named \texttt{x} instead of a name that reflects its requirement origin.

\begin{lstlisting}[style=Alabaster, language=openscad]
// FR2: Base must be stable — width >= 1.2 × total height
total_height = 80;
x = 1.2 * total_height;
cube([x, x, 5]);

\end{lstlisting}

\textbf{Q5.} The iteration step changes a dimension by guessing without measuring first. Rewrite the comment to enforce the engineering approach.

\begin{lstlisting}[style=Alabaster, language=openscad]
// Iteration: user said "it's too tight"
// Changed slip_clearance from 0.25 to 0.35 — seems about right
slip_clearance = 0.35;

\end{lstlisting}

\begin{center}\rule{0.5\linewidth}{0.5pt}\end{center}

\subsubsection*{Comment the Code (Questions 6--10)}\label{docs__pandoc__latex__src__advanced_study_guide.md__comment-the-code-questions-610-34}

\textbf{Q6.} Add comments to this requirement-mapped parameter block.

\begin{lstlisting}[style=Alabaster, language=openscad]
// FR1
bracelet_max_d = 70;

// FR2
total_height = 80;
base_w = 1.2 * total_height;

// FR3
module tactile_marker() {
  sphere(r=2, $fn=16);
}

// FR4
SLIP_C = 0.3;

\end{lstlisting}

\textbf{Q7.} Add comments explaining each function in this final project\textquotesingle s modular structure.

\begin{lstlisting}[style=Alabaster, language=openscad]
module bead_A(size=15, detail=32, hole_r=1.0) {
  assert(size > 0, "size must be positive");
  assert(hole_r > 0 && hole_r < size/2, "hole_r out of range");
  difference() {
    sphere(r=size/2, $fn=detail);
    cylinder(r=hole_r, h=size + 1, center=true, $fn=16);
  }
}

module bead_B(size=12, detail=6, hole_r=1.0) {
  assert(size > 0, "size must be positive");
  difference() {
    cylinder(r=size/2, h=size, center=true, $fn=detail);
    cylinder(r=hole_r, h=size + 1, center=true, $fn=16);
  }
}

\end{lstlisting}

\textbf{Q8.} Add comments to this structured feedback log.

\begin{lstlisting}[style=Alabaster, language=bash]
echo "=== Stakeholder Feedback — Iteration 1 ===" > feedback.txt
echo "Date: $(date)" >> feedback.txt
echo "Observation: user hesitated to locate clasp" >> feedback.txt
echo "Quote: 'the clasp blends in with the rest'" >> feedback.txt
echo "Measurement: clasp width = 8.2mm (designed 9.0mm)" >> feedback.txt
echo "Action: increase clasp_w by 0.8mm; reprint test coupon first" >> feedback.txt

\end{lstlisting}

\textbf{Q9.} Add comments to this user story and acceptance-criteria block embedded in the SCAD file.

\begin{lstlisting}[style=Alabaster, language=openscad]
// USER STORY
// As a jewelry collector,
// I want the holder to have 2mm clearance
// so that my bracelets slide on easily but don't fall off.

SLIP_C = 0.3;  // measured from test coupon

bracelet_max_d = 70;
holder_inner_r = (bracelet_max_d / 2) + SLIP_C;
holder_outer_r = holder_inner_r + 3;
holder_h = 10;

difference() {
  cylinder(r=holder_outer_r, h=holder_h, $fn=64);
  translate([0, 0, -0.001])
    cylinder(r=holder_inner_r, h=holder_h + 0.002, $fn=64);
}

\end{lstlisting}

\textbf{Q10.} Add comments to this complete design-thinking workflow script.

\begin{lstlisting}[style=Alabaster, language=bash]
#!/bin/bash
PROJECT="$1"
BASE="$HOME/Documents/3d_projects/$PROJECT"

mkdir -p "$BASE"
cd "$BASE"
3dm new

cat > design_notes.md << EOF
# $PROJECT Design Notes
## Stakeholder Interview
## Functional Requirements
## Acceptance Criteria
## Technical Design Specification
## Iteration Log
EOF

echo "Project scaffolded: $BASE"
ls -1

\end{lstlisting}

\section{Command Line Appendices --- Complete Reference}\label{docs__pandoc__latex__src__cli_appendices_complete.md__command-line-appendices--complete-reference}

Supplemental automation references for integrating OpenSCAD and 3dMake workflows with the three command-line environments covered in this curriculum. Each appendix is self-contained and maps directly to the corresponding CLI foundation curriculum.

\begin{center}\rule{0.5\linewidth}{0.5pt}\end{center}

\subsection*{Contents}\label{docs__pandoc__latex__src__cli_appendices_complete.md__contents}

{\def\LTcaptype{none} % do not increment counter
\begin{longtable}[]{@{}
  >{\raggedright\arraybackslash}p{(\linewidth - 4\tabcolsep) * \real{0.3333}}
  >{\raggedright\arraybackslash}p{(\linewidth - 4\tabcolsep) * \real{0.3333}}
  >{\raggedright\arraybackslash}p{(\linewidth - 4\tabcolsep) * \real{0.3333}}@{}}
\toprule\noalign{}
\begin{minipage}[b]{\linewidth}\raggedright
Appendix
\end{minipage} & \begin{minipage}[b]{\linewidth}\raggedright
Shell
\end{minipage} & \begin{minipage}[b]{\linewidth}\raggedright
Title
\end{minipage} \\
\midrule\noalign{}
\endhead
\bottomrule\noalign{}
\endlastfoot
\hyperref[docs__pandoc__latex__src__cli_appendices_complete.md__appendix-a]{Appendix A}
& CMD / Batch &
Command Line (CMD/Batch) Integration for SCAD Workflows \\
\hyperref[docs__pandoc__latex__src__cli_appendices_complete.md__appendix-b]{Appendix B}
& PowerShell & PowerShell Integration for SCAD Workflows \\
\hyperref[docs__pandoc__latex__src__cli_appendices_complete.md__appendix-c]{Appendix C}
& Git Bash / POSIX &
Git Bash / POSIX Shell Integration for SCAD Workflows \\
\end{longtable}
}

All three appendices cover the same workflow stages --- single-file build, batch directory build, parametric sweeps, logging, USB transfer, and network printer monitoring --- using the idioms and tools native to each shell.

\begin{center}\rule{0.5\linewidth}{0.5pt}\end{center}

\begin{center}\rule{0.5\linewidth}{0.5pt}\end{center}

\protect\phantomsection\label{docs__pandoc__latex__src__cli_appendices_complete.md__appendix-a}{}

\section*{Appendix A: Command Line (CMD/Batch) Integration for SCAD Workflows}\label{docs__pandoc__latex__src__cli_appendices_complete.md__appendix-a-command-line-cmdbatch-integration-for-scad-workflows}

\begin{center}\rule{0.5\linewidth}{0.5pt}\end{center}

This appendix shows how traditional command-line (Windows CMD / batch) scripting automates SCAD workflows without requiring PowerShell. It mirrors the patterns in Appendix D but provides examples and idioms for \texttt{cmd} / batch files and systems that prefer native Windows command prompt scripts.

\subsection*{Overview: Why Automate with CMD/Batch?}\label{docs__pandoc__latex__src__cli_appendices_complete.md__overview-why-automate-with-cmdbatch}

\begin{itemize}
\tightlist
\item
  Minimal dependencies: works on basic Windows installations.
\item
  Easy to call from other tools and CI systems that expect \texttt{.bat} helpers.
\item
  Useful for environments where PowerShell policy restrictions exist.
\end{itemize}

\subsection*{Prerequisites \& Setup}\label{docs__pandoc__latex__src__cli_appendices_complete.md__prerequisites--setup}

\subsubsection*{Required Software}\label{docs__pandoc__latex__src__cli_appendices_complete.md__required-software}

\begin{lstlisting}[style=Alabaster, language=cmd]
where openscad      :: OpenSCAD (path in PATH or full path required)
where prusa-slicer  :: PrusaSlicer (or your slicer)
where python        :: Python (optional)

\end{lstlisting}

\subsubsection*{Directory Structure (Windows style)}\label{docs__pandoc__latex__src__cli_appendices_complete.md__directory-structure-windows-style}

\begin{lstlisting}[style=Alabaster, language=cmd]
C:\Projects\3dMake\
+------ src\
+------ stl\
+------ gcode\
+------ logs\
+------ scripts\
       +------ build.bat
       +------ batch_build.bat

\end{lstlisting}

\subsubsection*{Notes on CMD vs PowerShell}\label{docs__pandoc__latex__src__cli_appendices_complete.md__notes-on-cmd-vs-powershell}

\begin{itemize}
\tightlist
\item
  CMD/batch has more limited error handling and no structured objects; rely on return codes and file checks.
\item
  Use full executable paths to avoid PATH surprises.
\end{itemize}

\subsection*{Basic Workflow: Single-file build (batch)}\label{docs__pandoc__latex__src__cli_appendices_complete.md__basic-workflow-single-file-build-batch}

Create \texttt{build.bat} (simple example):

\begin{lstlisting}[style=Alabaster, language=cmd]
@echo off
rem build.bat - Convert SCAD -> STL -> G-code (minimal)
setlocal enabledelayedexpansion
set PROJECTROOT=%~dp0\..\
set SCADFILE=%1
if "%SCADFILE%"=="" (
  echo Usage: build.bat file.scad
  exit /b 1
)
set SCADPATH=%PROJECTROOT%src\%SCADFILE%
set STLPATH=%PROJECTROOT%stl\%~n1.stl
set GCODEPATH=%PROJECTROOT%gcode\%~n1.gcode

rem Export STL using OpenSCAD
"C:\Program Files\OpenSCAD\openscad.exe" -o "%STLPATH%" "%SCADPATH%"
if not exist "%STLPATH%" (
  echo ERROR: STL not created
  exit /b 1
)

rem Slice (example) - adjust path to slicer
"C:\Program Files\Prusa3D\PrusaSlicer\prusa-slicer.exe" --load-config-file "%~dp0\default.ini" -o "%GCODEPATH%" "%STLPATH%"
if not exist "%GCODEPATH%" (
  echo ERROR: G-code not created
  exit /b 1
)

echo BUILD COMPLETE: %GCODEPATH%
endlocal

\end{lstlisting}

\subsection*{Batch (directory) build}\label{docs__pandoc__latex__src__cli_appendices_complete.md__batch-directory-build}

\texttt{batch\_build.bat} can iterate files and call \texttt{build.bat}:

\begin{lstlisting}[style=Alabaster, language=cmd]
@echo off
set PROJECTROOT=%~dp0\..\
for /r "%PROJECTROOT%src" %%f in (*.scad) do (
  echo Processing %%~nxf
  call "%~dp0build.bat" "%%~nxf"
)

\end{lstlisting}

\subsection*{Parametric Sweeps (simple templating)}\label{docs__pandoc__latex__src__cli_appendices_complete.md__parametric-sweeps-simple-templating}

Because batch is limited for text templating, a common pattern is to use small helper scripts with \texttt{sed}/\texttt{python} or write a minimal templating helper in a language such as Python. Example invocation in batch:

\begin{lstlisting}[style=Alabaster, language=cmd]
rem Example: call a Python script to generate variants, then build
python "%~dp0/generate_variants.py" "%PROJECTROOT%src\bracelet_holder.scad"
call "%~dp0/batch_build.bat"

\end{lstlisting}

\subsection*{Logging and CSV output}\label{docs__pandoc__latex__src__cli_appendices_complete.md__logging-and-csv-output}

Batch scripts can append simple CSV lines using \texttt{echo} redirection:

\begin{lstlisting}[style=Alabaster, language=cmd]
echo %DATE% %TIME%,%~n1,Success >> "%PROJECTROOT%logs\build_history.csv"

\end{lstlisting}

\subsection*{Send to Printer (USB copy)}\label{docs__pandoc__latex__src__cli_appendices_complete.md__send-to-printer-usb-copy}

\begin{lstlisting}[style=Alabaster, language=cmd]
rem Copy gcode to removable drive (E: example)
copy "%GCODEPATH%" E:\

\end{lstlisting}

\subsection*{Monitoring (networked printer) - using curl or PowerShell helper}\label{docs__pandoc__latex__src__cli_appendices_complete.md__monitoring-networked-printer---using-curl-or-powershell-helper}

\texttt{curl} on Windows or a small PowerShell one-liner may be used to query APIs. For pure CMD, ship a small helper .exe (curl) or call \texttt{powershell\ -Command\ "Invoke-RestMethod\ ..."}.

\subsection*{Best Practices}\label{docs__pandoc__latex__src__cli_appendices_complete.md__best-practices}

\begin{itemize}
\tightlist
\item
  Use full paths for all tools.
\item
  Check return codes (\texttt{if\ errorlevel\ 1}) after each step.
\item
  Keep batch scripts small; delegate complex text processing to Python/PowerShell.
\item
  Log to simple CSVs for human and machine parsing.
\end{itemize}

\begin{center}\rule{0.5\linewidth}{0.5pt}\end{center}

\begin{center}\rule{0.5\linewidth}{0.5pt}\end{center}

\protect\phantomsection\label{docs__pandoc__latex__src__cli_appendices_complete.md__appendix-b}{}

\section*{Appendix B: PowerShell Integration for SCAD Workflows}\label{docs__pandoc__latex__src__cli_appendices_complete.md__appendix-b-powershell-integration-for-scad-workflows}

\begin{center}\rule{0.5\linewidth}{0.5pt}\end{center}

This appendix shows how PowerShell (command-line scripting) streamlines 3D design workflows, automating the repetitive tasks from design through printing. It bridges PowerShell\_Foundation concepts with 3dMake\_Foundation practical applications.

Referenced in: Lessons 9 (Automation), 10 (Mastery), and any lesson requiring batch operations

Prerequisites: Completion of PowerShell\_Foundation (Lessons 1-6)

\subsection*{Overview: Why Automate SCAD Workflows?}\label{docs__pandoc__latex__src__cli_appendices_complete.md__overview-why-automate-scad-workflows}

\subsubsection*{Manual Workflow (Time-Consuming)}\label{docs__pandoc__latex__src__cli_appendices_complete.md__manual-workflow-time-consuming}

\begin{enumerate}
\tightlist
\item
  Open SCAD manually -\textgreater{} Edit parameters -\textgreater{} Save -\textgreater{} Export STL
\item
  Open Slicer manually -\textgreater{} Load STL -\textgreater{} Adjust settings -\textgreater{} Slice -\textgreater{} Export G-code
\item
  Transfer G-code to printer via USB
\item
  Record results in notebook
\item
  Repeat for next variation (variant 2, variant 3, etc.)
\end{enumerate}

Time: \textasciitilde{}20-30 minutes per design iteration

\subsubsection*{Automated Workflow (Fast \& Reproducible)}\label{docs__pandoc__latex__src__cli_appendices_complete.md__automated-workflow-fast--reproducible}

\begin{enumerate}
\tightlist
\item
  Create PowerShell script with design parameters
\item
  Script automatically:

  \begin{itemize}
  \tightlist
  \item
    Generates SCAD code
  \item
    Exports to STL
  \item
    Slices to G-code
  \item
    Transfers to printer
  \item
    Logs results
  \end{itemize}
\item
  Run script once, walk away
\item
  Check results later
\end{enumerate}

Time: 2-3 minutes per iteration (plus print time)

\subsubsection*{Benefits of Automation}\label{docs__pandoc__latex__src__cli_appendices_complete.md__benefits-of-automation}

\begin{itemize}
\tightlist
\item
  {[}YES{]} Speed: 10x faster for batch operations
\item
  {[}YES{]} Consistency: No manual errors
\item
  {[}YES{]} Reproducibility: Same settings every time
\item
  {[}YES{]} Scalability: Test 5, 10, or 100 variations easily
\item
  {[}YES{]} Logging: Automatic documentation
\item
  {[}YES{]} Accessibility: Screen reader captures script output (not UI clicks)
\end{itemize}

\subsection*{Prerequisites \& Setup}\label{docs__pandoc__latex__src__cli_appendices_complete.md__prerequisites--setup-1}

\subsubsection*{Required Software}\label{docs__pandoc__latex__src__cli_appendices_complete.md__required-software-1}

\begin{lstlisting}[style=Alabaster, language=powershell]
# Check what you have installed
where openscad      # OpenSCAD
where prusa-slicer  # PrusaSlicer (or your slicer)
where python3       # Python (optional, for advanced scripts)

\end{lstlisting}

\subsubsection*{Directory Structure}\label{docs__pandoc__latex__src__cli_appendices_complete.md__directory-structure}

Create a project directory:

\begin{lstlisting}[style=Alabaster]
C:\Projects\3dMake\
+------ src/
|   +------ bracelet_holder.scad
|   +------ phone_stand.scad
|   +------ stackable_bins.scad
+------ stl/
|   +------ bracelet_holder.stl
|   +------ phone_stand.stl
|   +------ stackable_bins.stl
+------ gcode/
|   +------ bracelet_holder.gcode
|   +------ phone_stand.gcode
|   +------ stackable_bins.gcode
+------ logs/
|   +------ batch_2024-01-15.log
|   +------ print_history.csv
+------ scripts/
    +------ build.ps1      (main build script)
    +------ batch_test.ps1 (test variations)
    +------ monitor.ps1    (monitor printer)

\end{lstlisting}

\subsubsection*{PowerShell Execution Policy}\label{docs__pandoc__latex__src__cli_appendices_complete.md__powershell-execution-policy}

\begin{lstlisting}[style=Alabaster, language=powershell]
# Check current policy
Get-ExecutionPolicy
# If it's "Restricted", change it (requires admin)
Set-ExecutionPolicy -ExecutionPolicy RemoteSigned -Scope CurrentUser
# Verify change
Get-ExecutionPolicy

\end{lstlisting}

\subsection*{Basic Workflow Automation}\label{docs__pandoc__latex__src__cli_appendices_complete.md__basic-workflow-automation}

\subsubsection*{Script 1: Single-File Build}\label{docs__pandoc__latex__src__cli_appendices_complete.md__script-1-single-file-build}

This script takes a SCAD file and converts it through the entire workflow:

\begin{lstlisting}[style=Alabaster, language=powershell]
# build.ps1 - Convert SCAD -> STL -> G-code
param(
    [string]$ScadFile,          # Input: bracelet_holder.scad
    [string]$OutputDir = ".\",  # Output directory
    [string]$SlicerConfig = "default"
)
# Set up paths
$ProjectRoot = "C:\Projects\3dMake"
$ScadFile = Join-Path $ProjectRoot "src" $ScadFile
$StlFile = Join-Path $ProjectRoot "stl" ([System.IO.Path]::GetFileNameWithoutExtension($ScadFile) + ".stl")
$GcodeFile = Join-Path $ProjectRoot "gcode" ([System.IO.Path]::GetFileNameWithoutExtension($ScadFile) + ".gcode")
$LogFile = Join-Path $ProjectRoot "logs" "build_$(Get-Date -Format 'yyyy-MM-dd').log"
# Log function
function Write-Log {
    param([string]$Message)
    $timestamp = Get-Date -Format "yyyy-MM-dd HH:mm:ss"
    $logEntry = "[$timestamp] $Message"
    Write-Host $logEntry
    Add-Content -Path $LogFile -Value $logEntry
}
Write-Log "Starting build for $($ScadFile | Split-Path -Leaf)"
# Step 1: Export STL
Write-Log "Step 1: Exporting SCAD to STL..."
$startTime = Get-Date
& "C:\Program Files\OpenSCAD\openscad.exe" `
    -o "$StlFile" `
    "$ScadFile"
$duration = (Get-Date) - $startTime
Write-Log " STL exported in $($duration.TotalSeconds) seconds: $StlFile"
# Verify STL exists
if (-not (Test-Path $StlFile)) {
    Write-Log "[NO] ERROR: STL file not created"
    exit 1
}
# Step 2: Slice to G-code
Write-Log "Step 2: Slicing STL to G-code..."
$startTime = Get-Date
& "C:\Program Files\Prusa3D\PrusaSlicer\prusa-slicer.exe" `
    --load-config-file "$SlicerConfig" `
    --export-gcode "$GcodeFile" `
    "$StlFile"
$duration = (Get-Date) - $startTime
Write-Log " G-code generated in $($duration.TotalSeconds) seconds: $GcodeFile"
# Verify G-code exists
if (-not (Test-Path $GcodeFile)) {
    Write-Log "[NO] ERROR: G-code file not created"
    exit 1
}
# Step 3: Get file info
$stlSize = (Get-Item $StlFile).Length / 1MB
$gcodeSize = (Get-Item $GcodeFile).Length / 1MB
Write-Log "File sizes: STL=$([math]::Round($stlSize, 2))MB, G-code=$([math]::Round($gcodeSize, 2))MB"
Write-Log " BUILD COMPLETE"
Write-Log "Ready to print: $GcodeFile"

\end{lstlisting}

Usage:

\begin{lstlisting}[style=Alabaster, language=powershell]
# Run the script
.\build.ps1 -ScadFile "bracelet_holder.scad"
# Output example:
# [2024-01-15 14:30:22] Starting build for bracelet_holder.scad
# [2024-01-15 14:30:23] Step 1: Exporting SCAD to STL...
# [2024-01-15 14:30:25]  STL exported in 2.41 seconds: ...
# [2024-01-15 14:30:26] Step 2: Slicing STL to G-code...
# [2024-01-15 14:30:28]  G-code generated in 2.31 seconds: ...
# [2024-01-15 14:30:28]  BUILD COMPLETE

\end{lstlisting}

\subsubsection*{Script 2: Batch Build (Multiple Files)}\label{docs__pandoc__latex__src__cli_appendices_complete.md__script-2-batch-build-multiple-files}

\begin{lstlisting}[style=Alabaster, language=powershell]
# batch_build.ps1 - Build all SCAD files in a directory
param(
    [string]$SourceDir = "C:\Projects\3dMake\src",
    [string]$SlicerConfig = "default"
)
$ProjectRoot = Split-Path $SourceDir
$buildScript = Join-Path $ProjectRoot "scripts" "build.ps1"
# Find all SCAD files
$scadFiles = Get-ChildItem -Path $SourceDir -Filter "*.scad" -Recurse
Write-Host "Found $($scadFiles.Count) SCAD files"
if ($scadFiles.Count -eq 0) {
    Write-Host "No SCAD files found in $SourceDir"
    exit 1
}
# Process each file
$results = @()
foreach ($file in $scadFiles) {
    Write-Host "`n--- Processing: $($file.Name) ---"
    $startTime = Get-Date
    # Run build script for this file
    & $buildScript -ScadFile $file.Name -SlicerConfig $SlicerConfig
    $duration = (Get-Date) - $startTime
    $results += [PSCustomObject]@{
        FileName = $file.Name
        DurationSeconds = $duration.TotalSeconds
        Status = "Success"
        Timestamp = Get-Date
    }
}
# Summary report
Write-Host "`n=== BATCH BUILD SUMMARY ==="
Write-Host "Files processed: $($results.Count)"
Write-Host "Total time: $([math]::Round(($results | Measure-Object -Property DurationSeconds -Sum).Sum, 1)) seconds"
# Save results to CSV
$csvPath = Join-Path $ProjectRoot "logs" "batch_$(Get-Date -Format 'yyyy-MM-dd-HHmmss').csv"
$results | Export-Csv -Path $csvPath -NoTypeInformation
Write-Host "Results saved: $csvPath"

\end{lstlisting}

Usage:

\begin{lstlisting}[style=Alabaster, language=powershell]
.\batch_build.ps1 -SourceDir "C:\Projects\3dMake\src"
# Processes all .scad files automatically

\end{lstlisting}

\subsection*{Parametric Design Variation Testing}\label{docs__pandoc__latex__src__cli_appendices_complete.md__parametric-design-variation-testing}

\subsubsection*{The Problem: Testing Multiple Parameter Values}\label{docs__pandoc__latex__src__cli_appendices_complete.md__the-problem-testing-multiple-parameter-values}

You want to test the same design with different parameters (e.g., peg diameter: 5mm, 6mm, 7mm):

Manually:

\begin{enumerate}
\tightlist
\item
  Open bracelet\_holder.scad
\item
  Change peg\_diameter = 5
\item
  Save, export, slice, print
\item
  Change peg\_diameter = 6
\item
  Save, export, slice, print
\item
  Change peg\_diameter = 7
\item
  Save, export, slice, print Time: 30 minutes
\end{enumerate}

\subsubsection*{Automated Solution: Parametric Sweep}\label{docs__pandoc__latex__src__cli_appendices_complete.md__automated-solution-parametric-sweep}

\begin{lstlisting}[style=Alabaster, language=powershell]
# parametric_sweep.ps1 - Test multiple parameter values
param(
    [string]$TemplateScad = "C:\Projects\3dMake\src\bracelet_holder.scad",
    [hashtable]$ParameterRanges = @{
        "peg_diameter" = @(5, 5.5, 6, 6.5, 7)
        "holder_width" = @(120, 127, 135)
    }
)
$ProjectRoot = "C:\Projects\3dMake"
$variantsDir = Join-Path $ProjectRoot "variants"
$buildScript = Join-Path $ProjectRoot "scripts" "build.ps1"
$logFile = Join-Path $ProjectRoot "logs" "parametric_sweep_$(Get-Date -Format 'yyyy-MM-dd-HHmmss').csv"
# Create variants directory
New-Item -ItemType Directory -Path $variantsDir -Force | Out-Null
# Read template
$template = Get-Content $TemplateScad -Raw
# Generate variants
$results = @()
$variantNum = 0
# For each diameter value
foreach ($dia in $ParameterRanges["peg_diameter"]) {
    # For each width value
    foreach ($width in $ParameterRanges["holder_width"]) {
        $variantNum++
        $variantName = "variant_dia${dia}_width${width}"
        $variantScad = Join-Path $variantsDir "$variantName.scad"
        # Create variant SCAD file with parameters
        $content = $template `
            -replace 'peg_diameter = [\d.]+', "peg_diameter = $dia" `
            -replace 'holder_width = [\d.]+', "holder_width = $width"
        $content | Set-Content $variantScad
        Write-Host "Generated: $variantName"
        # Build variant (STL + G-code)
        $startTime = Get-Date
        # Export to STL
        & "C:\Program Files\OpenSCAD\openscad.exe" `
            -o (Join-Path $ProjectRoot "stl" "$variantName.stl") `
            $variantScad
        $duration = (Get-Date) - $startTime
        # Record result
        $results += [PSCustomObject]@{
            Variant = $variantName
            PegDiameter = $dia
            HolderWidth = $width
            BuildTimeSeconds = $duration.TotalSeconds
            Status = "Success"
            Timestamp = Get-Date
        }
        Write-Host "   Built in $($duration.TotalSeconds) seconds"
    }
}
# Save results
$results | Export-Csv -Path $logFile -NoTypeInformation
Write-Host "`nParametric sweep complete: $($results.Count) variants generated"
Write-Host "Results: $logFile"
# Summary statistics
Write-Host "`n=== PARAMETER RANGES TESTED ==="
Write-Host "Peg diameters: $($ParameterRanges['peg_diameter'] -join ', ')"
Write-Host "Holder widths: $($ParameterRanges['holder_width'] -join ', ')"
Write-Host "Total combinations: $($results.Count)"

\end{lstlisting}

Usage:

\begin{lstlisting}[style=Alabaster, language=powershell]
$ranges = @{
    "peg_diameter" = @(5, 5.5, 6, 6.5, 7)
    "holder_width" = @(120, 127, 135)
}
.\parametric_sweep.ps1 -TemplateScad "bracelet_holder.scad" -ParameterRanges $ranges
# Generates 15 variants automatically (5 x 3)

\end{lstlisting}

Output CSV:

\begin{lstlisting}[style=Alabaster]
Variant,PegDiameter,HolderWidth,BuildTimeSeconds,Status,Timestamp
variant_dia5_width120,5,120,2.4,Success,2024-01-15 14:35:22
variant_dia5_width127,5,127,2.5,Success,2024-01-15 14:35:25
variant_dia5_width135,5,135,2.3,Success,2024-01-15 14:35:28
...

\end{lstlisting}

\subsection*{Automated Print Documentation}\label{docs__pandoc__latex__src__cli_appendices_complete.md__automated-print-documentation}

\subsubsection*{Script: Print Logging \& Quality Tracking}\label{docs__pandoc__latex__src__cli_appendices_complete.md__script-print-logging--quality-tracking}

\begin{lstlisting}[style=Alabaster, language=powershell]
# log_print.ps1 - Document each print with metadata
param(
    [string]$GcodeFile,
    [string]$ProjectName,
    [string]$Material = "PLA",
    [string]$Notes = ""
)
$ProjectRoot = "C:\Projects\3dMake"
$printLogCsv = Join-Path $ProjectRoot "logs" "print_history.csv"
# Get G-code file info
$gcodeInfo = Get-Item $GcodeFile
$gcodeSize = $gcodeInfo.Length / 1MB
$estimatedTime = EstimateTime $GcodeFile  # See function below
# Calculate estimated cost
$weightG = EstimateWeight $GcodeFile
$materialCostPerKg = 20  # PLA at $20/kg
$estimatedCost = ($weightG / 1000) * $materialCostPerKg
# Create log entry
$entry = [PSCustomObject]@{
    PrintID = "PRINT_$(Get-Date -Format 'yyyyMMdd_HHmmss')"
    ProjectName = $ProjectName
    Date = Get-Date
    Material = $Material
    GcodeFile = $gcodeFile
    GcodeSizeMB = [math]::Round($gcodeSize, 2)
    EstimatedTimeHours = [math]::Round($estimatedTime / 3600, 2)
    EstimatedWeightG = [math]::Round($weightG, 1)
    EstimatedCostUSD = [math]::Round($estimatedCost, 2)
    Notes = $Notes
}
# Append to CSV
if (Test-Path $printLogCsv) {
    $entry | Export-Csv -Path $printLogCsv -NoTypeInformation -Append
} else {
    $entry | Export-Csv -Path $printLogCsv -NoTypeInformation
}
Write-Host "Print logged:"
Write-Host "  Project: $ProjectName"
Write-Host "  File: $($gcodeInfo.Name)"
Write-Host "  Estimated time: $([math]::Round($estimatedTime / 3600, 2)) hours"
Write-Host "  Estimated weight: $([math]::Round($weightG, 1))g"
Write-Host "  Estimated cost: $$([math]::Round($estimatedCost, 2))"
# Function to estimate print time from G-code (simplified)
function EstimateTime {
    param([string]$GcodeFile)
    $lines = @(Get-Content $GcodeFile | Select-String "^G1" | Measure-Object).Count
    # Rough estimate: 10 lines per second
    return $lines / 10
}
# Function to estimate weight from G-code (simplified)
function EstimateWeight {
    param([string]$GcodeFile)
    # G-code filament estimate (if supported by your slicer)
    $content = Get-Content $GcodeFile -Raw
    if ($content -match "filament used = ([\d.]+)") {
        return [double]$matches[1]
    } else {
        return 0  # Fallback
    }
}

\end{lstlisting}

Usage:

\begin{lstlisting}[style=Alabaster, language=powershell]
.\log_print.ps1 `
    -GcodeFile "C:\Projects\3dMake\gcode\bracelet_holder.gcode" `
    -ProjectName "Bracelet Holder" `
    -Material "PLA" `
    -Notes "Final design, tested with actual bracelets"

\end{lstlisting}

\subsection*{Printer Communication \& Monitoring}\label{docs__pandoc__latex__src__cli_appendices_complete.md__printer-communication--monitoring}

\subsubsection*{Script: Send G-code to Printer (USB)}\label{docs__pandoc__latex__src__cli_appendices_complete.md__script-send-g-code-to-printer-usb}

\begin{lstlisting}[style=Alabaster, language=powershell]
# send_to_printer.ps1 - Copy G-code to printer USB
param(
    [string]$GcodeFile,
    [string]$PrinterUSBLetter = "E"  # E:, F:, G:, etc.
)
$printerDrive = "${PrinterUSBLetter}:\"
# Check if USB drive connected
if (-not (Test-Path $printerDrive)) {
    Write-Host "[NO] ERROR: Printer USB not found at $PrinterUSBLetter"
    Write-Host "Available drives:"
    Get-PSDrive -PSProvider FileSystem | Where-Object Name -Like "[D-Z]" | Select-Object Name
    exit 1
}
# Copy file
$filename = Split-Path $GcodeFile -Leaf
$destination = Join-Path $printerDrive $filename
Copy-Item -Path $GcodeFile -Destination $destination -Force
Write-Host " G-code copied to printer USB:"
Write-Host "  From: $GcodeFile"
Write-Host "  To: $destination"
Write-Host "`nNext steps:"
Write-Host "  1. Eject USB safely from computer"
Write-Host "  2. Insert USB into printer"
Write-Host "  3. Select file on printer screen"
Write-Host "  4. Press print"

\end{lstlisting}

Usage:

\begin{lstlisting}[style=Alabaster, language=powershell]
.\send_to_printer.ps1 -GcodeFile "C:\Projects\3dMake\gcode\bracelet_holder.gcode" -PrinterUSBLetter "E"

\end{lstlisting}

\subsubsection*{Script: Monitor Printer Status (Network Printers)}\label{docs__pandoc__latex__src__cli_appendices_complete.md__script-monitor-printer-status-network-printers}

\begin{lstlisting}[style=Alabaster, language=powershell]
# monitor_printer.ps1 - Check printer status via API
param(
    [string]$PrinterIP,
    [int]$CheckInterval = 30,    # seconds
    [int]$MaxChecks = 240        # 2 hours max
)
$printerUrl = "[https://example.com](https://example.com)
$checksPerformed = 0
Write-Host "Monitoring printer at $PrinterIP"
Write-Host "Check interval: $CheckInterval seconds"
Write-Host "Ctrl+C to stop`n"
while ($checksPerformed -lt $MaxChecks) {
    try {
        # Get printer status
        $response = Invoke-WebRequest -Uri $printerUrl -UseBasicParsing -ErrorAction Stop
        $status = $response.Content | ConvertFrom-Json
        $state = $status.state.text
        $progress = $status.progress.completion
        $timeRemaining = $status.progress.printTimeLeft
        Write-Host "[$((Get-Date).ToString("HH:mm:ss"))] State: $state | Progress: $progress% | Time remaining: $timeRemaining seconds"
        # If print completed, exit loop
        if ($state -eq "Operational") {
            Write-Host "`n Print complete!"
            break
        }
    } catch {
        Write-Host "Connection error: $_"
    }
    Start-Sleep -Seconds $CheckInterval
    $checksPerformed++
}
if ($checksPerformed -ge $MaxChecks) {
    Write-Host "Max monitoring time reached. Stopping."
}

\end{lstlisting}

Usage:

\begin{lstlisting}[style=Alabaster, language=powershell]
.\monitor_printer.ps1 -PrinterIP "192.168.1.100" -CheckInterval 30

\end{lstlisting}

\subsection*{Complete Workflow Integration}\label{docs__pandoc__latex__src__cli_appendices_complete.md__complete-workflow-integration}

\subsubsection*{Master Script: Design -\textgreater{} Print -\textgreater{} Log}\label{docs__pandoc__latex__src__cli_appendices_complete.md__master-script-design---print---log}

\begin{lstlisting}[style=Alabaster, language=powershell]
# full_workflow.ps1 - Complete automation from design to printing
param(
    [string]$ProjectName,
    [string]$ScadFile,
    [string]$Material = "PLA",
    [string]$Notes = "",
    [switch]$SendToPrinter,
    [string]$PrinterUSBLetter = "E"
)
$ProjectRoot = "C:\Projects\3dMake"
Write-Host "=== 3dMake Full Workflow Automation ==="
Write-Host "Project: $ProjectName"
Write-Host ""
# Step 1: Build (SCAD -> STL -> G-code)
Write-Host "Step 1: Building..."
& "$ProjectRoot\scripts\build.ps1" -ScadFile $ScadFile
if ($LASTEXITCODE -ne 0) { exit 1 }
# Step 2: Log the print
Write-Host "`nStep 2: Logging print metadata..."
$gcodeFile = Join-Path $ProjectRoot "gcode" ([System.IO.Path]::GetFileNameWithoutExtension($ScadFile) + ".gcode")
& "$ProjectRoot\scripts\log_print.ps1" `
    -GcodeFile $gcodeFile `
    -ProjectName $ProjectName `
    -Material $Material `
    -Notes $Notes
# Step 3: Send to printer (optional)
if ($SendToPrinter) {
    Write-Host "`nStep 3: Sending to printer..."
    & "$ProjectRoot\scripts\send_to_printer.ps1" `
        -GcodeFile $gcodeFile `
        -PrinterUSBLetter $PrinterUSBLetter
}
Write-Host "`n Workflow complete!"

\end{lstlisting}

Usage:

\begin{lstlisting}[style=Alabaster, language=powershell]
# Full workflow with automatic printer transfer
.\full_workflow.ps1 `
    -ProjectName "Bracelet Holder" `
    -ScadFile "bracelet_holder.scad" `
    -Material "PLA" `
    -Notes "Final design, v2 with improved peg strength" `
    -SendToPrinter `
    -PrinterUSBLetter "E"

\end{lstlisting}

\subsection*{PowerShell Skills Applied to SCAD}\label{docs__pandoc__latex__src__cli_appendices_complete.md__powershell-skills-applied-to-scad}

\subsubsection*{Lesson Mapping: PowerShell -\textgreater{} SCAD Workflows}\label{docs__pandoc__latex__src__cli_appendices_complete.md__lesson-mapping-powershell---scad-workflows}

{\def\LTcaptype{none} % do not increment counter
\begin{longtable}[]{@{}
  >{\raggedright\arraybackslash}p{(\linewidth - 4\tabcolsep) * \real{0.2432}}
  >{\raggedright\arraybackslash}p{(\linewidth - 4\tabcolsep) * \real{0.3063}}
  >{\raggedright\arraybackslash}p{(\linewidth - 4\tabcolsep) * \real{0.4505}}@{}}
\toprule\noalign{}
\begin{minipage}[b]{\linewidth}\raggedright
PowerShell Lesson
\end{minipage} & \begin{minipage}[b]{\linewidth}\raggedright
SCAD Application
\end{minipage} & \begin{minipage}[b]{\linewidth}\raggedright
Example
\end{minipage} \\
\midrule\noalign{}
\endhead
\bottomrule\noalign{}
\endlastfoot
PS 1: Navigation & Working with project directories &
Using \texttt{cd} to navigate src/ -\textgreater{} stl/ -\textgreater{} gcode/ \\
PS 2: File Manipulation & Copying, organizing SCAD files &
Copy design files, organize variants by date \\
PS 3: Piping \& Objects & Pass data between scripts &
\texttt{\$results\ \textbar{}\ Export-Csv} exports analysis \\
PS 4: Variables \& Aliases & Parameterize SCAD workflows &
Variables store file paths, material types, etc. \\
PS 5: Functions \& Modules & Reusable automation blocks &
Build, slice, log = functions in one script \\
PS Unit Test & Verify workflow correctness &
Test that STL file exists before slicing \\
\end{longtable}
}

\subsubsection*{Example: File Navigation with SCAD Projects}\label{docs__pandoc__latex__src__cli_appendices_complete.md__example-file-navigation-with-scad-projects}

\begin{lstlisting}[style=Alabaster, language=powershell]
# Navigate between SCAD project folders
$ProjectRoot = "C:\Projects\3dMake"
# Go to source directory
cd "$ProjectRoot\src"
Get-ChildItem -Filter "*.scad"  # List all SCAD files
# Process each file
foreach ($file in Get-ChildItem -Filter "*.scad") {
    Write-Host "Processing: $($file.Name)"
    # Do something with each file
}
# Go to output directory and check results
cd "$ProjectRoot\stl"
Get-ChildItem -Filter "*.stl" | Measure-Object -Sum

\end{lstlisting}

\subsubsection*{Example: Working with Piping \& Objects}\label{docs__pandoc__latex__src__cli_appendices_complete.md__example-working-with-piping--objects}

\begin{lstlisting}[style=Alabaster, language=powershell]
# Build all files, then get statistics
Get-ChildItem -Path "$ProjectRoot\src" -Filter "*.scad" |
    ForEach-Object {
        # Build each one
        & .\build.ps1 -ScadFile $_.Name
    } |
    # Analyze results
    Group-Object -Property Status |
    Select-Object Name, Count

\end{lstlisting}

\subsection*{Accessibility Considerations}\label{docs__pandoc__latex__src__cli_appendices_complete.md__accessibility-considerations}

\subsubsection*{Why PowerShell for SCAD?}\label{docs__pandoc__latex__src__cli_appendices_complete.md__why-powershell-for-scad}

\begin{enumerate}
\tightlist
\item
  Screen Reader Friendly: Console output = text (not UI clicks)
\item
  Scriptable: Can run overnight, results logged
\item
  Auditable: Every step written to log file
\item
  Shareable: Scripts document the process for others
\item
  Testable: Can verify each step independently
\end{enumerate}

\subsubsection*{Example: Accessible Build Output}\label{docs__pandoc__latex__src__cli_appendices_complete.md__example-accessible-build-output}

\begin{lstlisting}[style=Alabaster]
[2024-01-15 14:30:22] Starting build for bracelet_holder.scad
[2024-01-15 14:30:23] Step 1: Exporting SCAD to STL...
[2024-01-15 14:30:25]  STL exported in 2.41 seconds
[2024-01-15 14:30:26] Step 2: Slicing STL to G-code...
[2024-01-15 14:30:28]  G-code generated in 2.31 seconds
[2024-01-15 14:30:28] File sizes: STL=2.15MB, G-code=4.32MB
[2024-01-15 14:30:28]  BUILD COMPLETE

Ready to print: C:\Projects\3dMake\gcode\bracelet_holder.gcode

\end{lstlisting}

All information is text-based and sequential-perfectly accessible to screen readers.

\subsection*{Best Practices \& Tips}\label{docs__pandoc__latex__src__cli_appendices_complete.md__best-practices--tips}

\subsubsection*{1. Always Log Everything}\label{docs__pandoc__latex__src__cli_appendices_complete.md__1-always-log-everything}

\begin{lstlisting}[style=Alabaster, language=powershell]
function Write-Log {
    param(
        [string]$Message,
        [string]$LogFile
    )
    $timestamp = Get-Date -Format "yyyy-MM-dd HH:mm:ss"
    $entry = "[$timestamp] $Message"
    Write-Host $entry
    Add-Content -Path $LogFile -Value $entry
}

\end{lstlisting}

\subsubsection*{2. Verify Steps Succeed Before Continuing}\label{docs__pandoc__latex__src__cli_appendices_complete.md__2-verify-steps-succeed-before-continuing}

\begin{lstlisting}[style=Alabaster, language=powershell]
# Don't just run commands-check they worked
& "C:\Program Files\OpenSCAD\openscad.exe" -o "$StlFile" "$ScadFile"
if (-not (Test-Path $StlFile)) {
    Write-Log "ERROR: STL creation failed"
    exit 1
}
Write-Log " STL created successfully"

\end{lstlisting}

\subsubsection*{3. Make Scripts Reusable}\label{docs__pandoc__latex__src__cli_appendices_complete.md__3-make-scripts-reusable}

\begin{lstlisting}[style=Alabaster, language=powershell]
# BAD: Hardcoded paths
$scadFile = "C:\Users\John\Desktop\bracelet_holder.scad"
# GOOD: Parameters with defaults
param(
    [string]$ScadFile = "bracelet_holder.scad",
    [string]$ProjectRoot = "C:\Projects\3dMake"
)
$fullPath = Join-Path $ProjectRoot "src" $ScadFile

\end{lstlisting}

\subsubsection*{4. Use Configuration Files}\label{docs__pandoc__latex__src__cli_appendices_complete.md__4-use-configuration-files}

\begin{lstlisting}[style=Alabaster, language=powershell]
# config.ps1 - Centralized settings
$Config = @{
    ProjectRoot = "C:\Projects\3dMake"
    OpenSCADPath = "C:\Program Files\OpenSCAD\openscad.exe"
    SlicerPath = "C:\Program Files\Prusa3D\PrusaSlicer\prusa-slicer.exe"
    PrinterUSB = "E"
    DefaultMaterial = "PLA"
    DefaultInfill = 20
}
# Use in scripts:
# . .\config.ps1
# & $Config.OpenSCADPath ...

\end{lstlisting}

\subsection*{Troubleshooting Automated Workflows}\label{docs__pandoc__latex__src__cli_appendices_complete.md__troubleshooting-automated-workflows}

{\def\LTcaptype{none} % do not increment counter
\begin{longtable}[]{@{}
  >{\raggedright\arraybackslash}p{(\linewidth - 4\tabcolsep) * \real{0.1695}}
  >{\raggedright\arraybackslash}p{(\linewidth - 4\tabcolsep) * \real{0.2458}}
  >{\raggedright\arraybackslash}p{(\linewidth - 4\tabcolsep) * \real{0.5847}}@{}}
\toprule\noalign{}
\begin{minipage}[b]{\linewidth}\raggedright
Problem
\end{minipage} & \begin{minipage}[b]{\linewidth}\raggedright
Cause
\end{minipage} & \begin{minipage}[b]{\linewidth}\raggedright
Solution
\end{minipage} \\
\midrule\noalign{}
\endhead
\bottomrule\noalign{}
\endlastfoot
Scripts won\textquotesingle t run & Execution policy &
\texttt{Set-ExecutionPolicy\ -ExecutionPolicy\ RemoteSigned} \\
Command not found & Tool not in PATH &
Specify full path: \texttt{\&\ "C:\textbackslash{}Program\ Files\textbackslash{}OpenSCAD\textbackslash{}openscad.exe"} \\
Files not created & Wrong working directory &
Use absolute paths with \texttt{Join-Path} \\
Script hangs & Waiting for user input &
Disable UI input; use \texttt{-o} flag for batch operations \\
CSV data corrupted & Special characters in notes &
Properly quote strings in CSV export \\
No output/logging & Path doesn\textquotesingle t exist &
Create directory first: \texttt{New-Item\ -Type\ Directory\ -Path\ ...\ -Force} \\
\end{longtable}
}

\subsection*{Summary: PowerShell + SCAD Integration Benefits}\label{docs__pandoc__latex__src__cli_appendices_complete.md__summary-powershell--scad-integration-benefits}

\subsubsection*{Time Savings}\label{docs__pandoc__latex__src__cli_appendices_complete.md__time-savings}

{\def\LTcaptype{none} % do not increment counter
\begin{longtable}[]{@{}llll@{}}
\toprule\noalign{}
Task & Manual & Automated & Savings \\
\midrule\noalign{}
\endhead
\bottomrule\noalign{}
\endlastfoot
Single design build & 5 min & 2.5 min & 50\% \\
Test 10 parameter variations & 50 min & 5 min & 90\% \\
Batch processing 20 designs & 100 min & 10 min & 90\% \\
\end{longtable}
}

\subsubsection*{Consistency}\label{docs__pandoc__latex__src__cli_appendices_complete.md__consistency}

\begin{itemize}
\tightlist
\item
  {[}YES{]} Same process every time (no missed steps)
\item
  {[}YES{]} Reproducible results (exact same settings)
\item
  {[}YES{]} Comprehensive logging (documentation automatic)
\end{itemize}

\subsubsection*{Scalability}\label{docs__pandoc__latex__src__cli_appendices_complete.md__scalability}

\begin{itemize}
\tightlist
\item
  {[}YES{]} 1 design or 1,000 designs = same time commitment
\item
  {[}YES{]} Run overnight for batch operations
\item
  {[}YES{]} Easy to expand with new features
\end{itemize}

\subsubsection*{Accessibility}\label{docs__pandoc__latex__src__cli_appendices_complete.md__accessibility}

\begin{itemize}
\tightlist
\item
  {[}YES{]} All interaction is text-based (screen reader friendly)
\item
  {[}YES{]} Results logged in CSV (machine-readable, sortable)
\item
  {[}YES{]} Repeatable (can verify steps later)
\end{itemize}

\subsection*{Next Steps}\label{docs__pandoc__latex__src__cli_appendices_complete.md__next-steps}

\begin{enumerate}
\tightlist
\item
  Try the basic build script (Script 1) with one of your SCAD files
\item
  Add logging (Script 3) to track what you create
\item
  Test parametric sweeps (Script 2) with design variations
\item
  Integrate full workflow (Script 5) for complete automation
\item
  Customize config.ps1 for your specific setup
\end{enumerate}

Remember: Start simple, test each script individually, then combine them into larger workflows. PowerShell is most powerful when built incrementally!

\begin{center}\rule{0.5\linewidth}{0.5pt}\end{center}

\protect\phantomsection\label{docs__pandoc__latex__src__cli_appendices_complete.md__appendix-c}{}

\section*{Appendix C: Git Bash / POSIX Shell Integration for SCAD Workflows}\label{docs__pandoc__latex__src__cli_appendices_complete.md__appendix-c-git-bash--posix-shell-integration-for-scad-workflows}

\begin{center}\rule{0.5\linewidth}{0.5pt}\end{center}

This appendix mirrors Appendix D (PowerShell) but provides examples for Git Bash (mingw) or other POSIX-compatible shells on Windows, macOS, and Linux. Use these scripts when you prefer shell utilities (\texttt{bash}, \texttt{sed}, \texttt{awk}, \texttt{curl}) and Unix-like tooling.

\subsection*{Overview: Why Git Bash / POSIX Shell?}\label{docs__pandoc__latex__src__cli_appendices_complete.md__overview-why-git-bash--posix-shell}

\begin{itemize}
\tightlist
\item
  Leverage standard Unix tools for text processing and automation.
\item
  Cross-platform: same scripts often work on Linux/macOS and Windows via Git Bash.
\item
  Good fit for containerized CI like GitHub Actions.
\end{itemize}

\subsection*{Prerequisites \& Setup}\label{docs__pandoc__latex__src__cli_appendices_complete.md__prerequisites--setup-2}

\subsubsection*{Required Software}\label{docs__pandoc__latex__src__cli_appendices_complete.md__required-software-2}

\begin{lstlisting}[style=Alabaster, language=bash]
command -v openscad   # OpenSCAD
command -v prusa-slicer # PrusaSlicer (or use CLI slicer)
command -v python3     # Python (optional)
command -v curl        # curl for API calls

\end{lstlisting}

\subsubsection*{Directory Structure (POSIX-style)}\label{docs__pandoc__latex__src__cli_appendices_complete.md__directory-structure-posix-style}

\begin{lstlisting}[style=Alabaster]
~/projects/3dMake/
|- src/
|- stl/
|- gcode/
|- logs/
`- scripts/
  |- build.sh
  `- batch_build.sh

\end{lstlisting}

\subsection*{Basic Workflow: Single-file build (bash)}\label{docs__pandoc__latex__src__cli_appendices_complete.md__basic-workflow-single-file-build-bash}

\texttt{scripts/build.sh}:

\begin{lstlisting}[style=Alabaster, language=bash]
#!/usr/bin/env bash
set -euo pipefail
ROOT_DIR="$(cd "$(dirname "$0")/.." && pwd)"
SCAD_FILE="$1"
SCAD_PATH="$ROOT_DIR/src/$SCAD_FILE"
STL_PATH="$ROOT_DIR/stl/${SCAD_FILE%.*}.stl"
GCODE_PATH="$ROOT_DIR/gcode/${SCAD_FILE%.*}.gcode"

echo "Exporting SCAD -> STL: $SCAD_FILE"
"/c/Program Files/OpenSCAD/openscad.exe" -o "$STL_PATH" "$SCAD_PATH" || { echo "STL export failed"; exit 1; }

echo "Slicing STL -> G-code"
"/c/Program Files/Prusa3D/PrusaSlicer/prusa-slicer.exe" --load-config-file "$ROOT_DIR/scripts/default.ini" -o "$GCODE_PATH" "$STL_PATH" || { echo "Slicing failed"; exit 1; }

echo "BUILD COMPLETE: $GCODE_PATH"

\end{lstlisting}

Notes:

\begin{itemize}
\tightlist
\item
  Paths to GUI apps on Windows may need the \texttt{/c/Program\textbackslash{}\ Files/...} form under Git Bash, or call the native \texttt{.exe} with full quoted path.
\item
  On Linux/macOS adjust executable paths accordingly.
\end{itemize}

\subsection*{Batch Build (directory)}\label{docs__pandoc__latex__src__cli_appendices_complete.md__batch-build-directory}

\texttt{scripts/batch\_build.sh}:

\begin{lstlisting}[style=Alabaster, language=bash]
#!/usr/bin/env bash
set -euo pipefail
ROOT_DIR="$(cd "$(dirname "$0")/.." && pwd)"
find "$ROOT_DIR/src" -name "*.scad" -print0 | while IFS= read -r -d '' file; do
  fname="$(basename "$file")"
  echo "Building $fname"
  "$ROOT_DIR/scripts/build.sh" "$fname"
done

\end{lstlisting}

\subsection*{Parametric Sweeps}\label{docs__pandoc__latex__src__cli_appendices_complete.md__parametric-sweeps}

Use \texttt{sed} or \texttt{python} to template parameter changes, then call \texttt{build.sh} for each variant. Example (bash + python hybrid):

\begin{lstlisting}[style=Alabaster, language=bash]
python3 scripts/gen_variants.py "$ROOT_DIR/src/bracelet_holder.scad" --out variants/
for v in variants/*.scad; do
  ./scripts/build.sh "$(basename "$v")"
done

\end{lstlisting}

\subsection*{Logging \& CSV}\label{docs__pandoc__latex__src__cli_appendices_complete.md__logging--csv}

Append simple CSV lines using \texttt{printf}:

\begin{lstlisting}[style=Alabaster, language=bash]
printf "%s,%s,Success\n" "$(date +'%Y-%m-%d %H:%M:%S')" "${SCAD_FILE%.*}" >> "$ROOT_DIR/logs/build_history.csv"

\end{lstlisting}

\subsection*{Send to Printer (USB)}\label{docs__pandoc__latex__src__cli_appendices_complete.md__send-to-printer-usb}

\begin{lstlisting}[style=Alabaster, language=bash]
# Copy to mounted USB drive (example mount point /media/usb)
cp "$GCODE_PATH" "/media/usb/"

\end{lstlisting}

\subsection*{Monitoring (network printers)}\label{docs__pandoc__latex__src__cli_appendices_complete.md__monitoring-network-printers}

Use \texttt{curl} to call a printer API (e.g., OctoPrint):

\begin{lstlisting}[style=Alabaster, language=bash]
curl -s "[https://example.com](https://example.com) -H "X-Api-Key: $API_KEY" | jq .

\end{lstlisting}

\subsection*{Best Practices}\label{docs__pandoc__latex__src__cli_appendices_complete.md__best-practices-1}

\begin{itemize}
\tightlist
\item
  Use \texttt{set\ -euo\ pipefail} for safer scripts.
\item
  Use \texttt{mktemp} or a \texttt{variants/} directory for generated files.
\item
  Keep heavy text processing to \texttt{python} if logic becomes complex.
\item
  Make scripts executable (\texttt{chmod\ +x\ scripts/*.sh}).
\end{itemize}

\begin{center}\rule{0.5\linewidth}{0.5pt}\end{center}

\section{3dMake and OpenSCAD Appendices --- Complete Reference}\label{docs__pandoc__latex__src__3dmake_appendices_complete.md__3dmake-and-openscad-appendices--complete-reference}

Comprehensive reference materials for the 3dMake Foundation and OpenSCAD curriculum, including makerspace locators, tool references, slicing guides, material science, quality assurance procedures, and advanced design concepts.

\begin{center}\rule{0.5\linewidth}{0.5pt}\end{center}

\subsection*{Contents}\label{docs__pandoc__latex__src__3dmake_appendices_complete.md__contents}

{\def\LTcaptype{none} % do not increment counter
\begin{longtable}[]{@{}
  >{\raggedright\arraybackslash}p{(\linewidth - 2\tabcolsep) * \real{0.5000}}
  >{\raggedright\arraybackslash}p{(\linewidth - 2\tabcolsep) * \real{0.5000}}@{}}
\toprule\noalign{}
\begin{minipage}[b]{\linewidth}\raggedright
Appendix
\end{minipage} & \begin{minipage}[b]{\linewidth}\raggedright
Title
\end{minipage} \\
\midrule\noalign{}
\endhead
\bottomrule\noalign{}
\endlastfoot
\hyperref[docs__pandoc__latex__src__3dmake_appendices_complete.md__makerspace-utah]{Makerspace Index: Utah}
& Utah Makerspace Index (By County) \\
\hyperref[docs__pandoc__latex__src__3dmake_appendices_complete.md__makerspace-national]{Makerspace Index: United States}
& National Makerspace Guide --- State-by-State Index \\
\hyperref[docs__pandoc__latex__src__3dmake_appendices_complete.md__appendix-a]{Appendix A}
& 3DMake CLI Cheat Sheet \\
\hyperref[docs__pandoc__latex__src__3dmake_appendices_complete.md__appendix-b]{Appendix B}
& Comprehensive Slicing Guide --- All Major Slicers \\
\hyperref[docs__pandoc__latex__src__3dmake_appendices_complete.md__appendix-c]{Appendix C}
& Material Properties and Selection Guide \\
\hyperref[docs__pandoc__latex__src__3dmake_appendices_complete.md__appendix-d]{Appendix D}
& Tolerance Testing and Quality Assurance Matrix \\
\hyperref[docs__pandoc__latex__src__3dmake_appendices_complete.md__appendix-e]{Appendix E}
& Advanced OpenSCAD Concepts \\
\end{longtable}
}

\begin{center}\rule{0.5\linewidth}{0.5pt}\end{center}

\begin{center}\rule{0.5\linewidth}{0.5pt}\end{center}

\protect\phantomsection\label{docs__pandoc__latex__src__3dmake_appendices_complete.md__makerspace-utah}{}

\section*{Makerspace Index: Utah (By County)}\label{docs__pandoc__latex__src__3dmake_appendices_complete.md__makerspace-index-utah-by-county}

\begin{center}\rule{0.5\linewidth}{0.5pt}\end{center}

\subsection*{Beaver County}\label{docs__pandoc__latex__src__3dmake_appendices_complete.md__beaver-county}

\emph{(No makerspace info discovered.)}

\subsection*{Box Elder County}\label{docs__pandoc__latex__src__3dmake_appendices_complete.md__box-elder-county}

\emph{(No makerspace info discovered.)}

\subsection*{Cache County}\label{docs__pandoc__latex__src__3dmake_appendices_complete.md__cache-county}

\subsubsection*{Logan Library -- Cytiva STEM Makerspace}\label{docs__pandoc__latex__src__3dmake_appendices_complete.md__logan-library--cytiva-stem-makerspace}

Bambu 3D printers, mini printers, vinyl cutter, sublimation printer, laser engraver, heat presses, craft cutters, and media digitization tools.\footnote{Logan Library Cytiva STEM Makerspace.}

\subsection*{Carbon County}\label{docs__pandoc__latex__src__3dmake_appendices_complete.md__carbon-county}

\subsubsection*{Southeast Education Service Center (SESC)}\label{docs__pandoc__latex__src__3dmake_appendices_complete.md__southeast-education-service-center-sesc}

STEM kits with coaching, curriculum, and project materials for regional schools.\footnote{Utah STEM Action Center -- Innovation Hub Network.}

\subsection*{Iron County}\label{docs__pandoc__latex__src__3dmake_appendices_complete.md__iron-county}

\subsubsection*{Southern Utah University -- Thunderworks Makerspace}\label{docs__pandoc__latex__src__3dmake_appendices_complete.md__southern-utah-university--thunderworks-makerspace}

Campus/community fabrication lab with CNC tools, 3D printers, and electronics benches.\footnote{Utah STEM Action Center -- Innovation Hub Network.}

\subsection*{Salt Lake County}\label{docs__pandoc__latex__src__3dmake_appendices_complete.md__salt-lake-county}

\subsubsection*{Salt Lake City Public Library -- Creative Lab}\label{docs__pandoc__latex__src__3dmake_appendices_complete.md__salt-lake-city-public-library--creative-lab}

3D printers, Cricut cutters, photography/film tools, Adobe Creative Cloud, recording booth.\footnote{Salt Lake City Public Library Creative Lab.}

\subsubsection*{Salt Lake County Library -- Create Spaces}\label{docs__pandoc__latex__src__3dmake_appendices_complete.md__salt-lake-county-library--create-spaces}

3D printers, VR, robotics, A/V studios, DIY repair tools across six branches.\footnote{Salt Lake County Library Create Spaces.}

\subsubsection*{Make Salt Lake}\label{docs__pandoc__latex__src__3dmake_appendices_complete.md__make-salt-lake}

Community makerspace with woodshop, metal shop, laser cutters, electronics benches, and 3D printer lab.\footnote{Utah STEM Action Center -- Innovation Hub Network.}

\subsubsection*{University of Utah -- Marriott Library ProtoSpace}\label{docs__pandoc__latex__src__3dmake_appendices_complete.md__university-of-utah--marriott-library-protospace}

Prusa MK4/XL 3D printers, 3D scanning, prototyping tools.\footnote{University of Utah -- Marriott Library ProtoSpace \& 3D Printing.}

\subsection*{Washington County}\label{docs__pandoc__latex__src__3dmake_appendices_complete.md__washington-county}

\subsubsection*{Washington County Library -- Sid Atkin Makerspace}\label{docs__pandoc__latex__src__3dmake_appendices_complete.md__washington-county-library--sid-atkin-makerspace}

3D printers, CNC cutter, sewing machines, woodworking tools, electronics kits, VR.\footnote{Washington County Library -- Sid Atkin Makerspace.}

\subsubsection*{Utah Tech University -- Jr. Makerspace}\label{docs__pandoc__latex__src__3dmake_appendices_complete.md__utah-tech-university--jr-makerspace}

K--12 and community maker training and fabrication tools.\footnote{Utah STEM Action Center -- Innovation Hub Network.}

\subsection*{Weber County}\label{docs__pandoc__latex__src__3dmake_appendices_complete.md__weber-county}

\subsubsection*{Weber County Library Makerspaces}\label{docs__pandoc__latex__src__3dmake_appendices_complete.md__weber-county-library-makerspaces}

Laser cutters, vinyl cutters, laminators, embroidery machines, scanners, and design software stations.\footnote{Weber County Library Makerspaces.}

\begin{center}\rule{0.5\linewidth}{0.5pt}\end{center}

\section*{Endnotes}\label{docs__pandoc__latex__src__3dmake_appendices_complete.md__endnotes}

\begin{center}\rule{0.5\linewidth}{0.5pt}\end{center}

\protect\phantomsection\label{docs__pandoc__latex__src__3dmake_appendices_complete.md__makerspace-national}{}

\section*{National Makerspace Guide (United States) --- State-by-State Index}\label{docs__pandoc__latex__src__3dmake_appendices_complete.md__national-makerspace-guide-united-states--state-by-state-index}

\begin{center}\rule{0.5\linewidth}{0.5pt}\end{center}

A comprehensive state‑by‑state index of ways to locate makerspaces, hackerspaces, Fab Labs, library labs, and community workshops across the U.S., including national directories and state‑level aggregators.

\subsection*{National Directories (Best Starting Points)}\label{docs__pandoc__latex__src__3dmake_appendices_complete.md__national-directories-best-starting-points}

\subsubsection*{Nation of Makers -- U.S. Maker Organization Map}\label{docs__pandoc__latex__src__3dmake_appendices_complete.md__nation-of-makers--us-maker-organization-map}

Crowdsourced map of U.S.-based maker organizations including makerspaces, hackerspaces, DIY bio labs, Fab Labs, school labs, library spaces, and maker events.\footnote{https://www.nationofmakers.us/find-a-maker-organization}

\subsubsection*{Makerspace.com -- Global \& U.S. Makerspace Directory}\label{docs__pandoc__latex__src__3dmake_appendices_complete.md__makerspacecom--global--us-makerspace-directory}

Large, continuously updated worldwide directory with city/state filters.\footnote{https://makerspace.com/directory/}

\subsubsection*{Makerspacedir.com -- Makerspace Directory}\label{docs__pandoc__latex__src__3dmake_appendices_complete.md__makerspacedircom--makerspace-directory}

U.S.-focused listing of makerspaces, hackerspaces, Fab Labs, coworking/maker hybrids, and creative community spaces.\footnote{https://makerspacedir.com/}

\subsubsection*{POIData Makerspace Industry Report (2026)}\label{docs__pandoc__latex__src__3dmake_appendices_complete.md__poidata-makerspace-industry-report-2026}

Lists \textbf{1,046 makerspaces in the U.S.}, with counts by state (e.g., CA 84, FL 45, CO 37, OR 35). Useful for statistical and geographical reference.\footnote{https://www.poidata.io/report/makerspace/united-states}

\subsubsection*{WorkMakeCreate -- Public \& Free-to-Use Makerspace Finder}\label{docs__pandoc__latex__src__3dmake_appendices_complete.md__workmakecreate--public--free-to-use-makerspace-finder}

Locates community workshops, tool libraries, and free/low‑cost makerspaces by ZIP code.\footnote{https://workmakecreate.com/}

\subsubsection*{GlobalSpec Makerspace Provider List}\label{docs__pandoc__latex__src__3dmake_appendices_complete.md__globalspec-makerspace-provider-list}

Allows filtering of makerspaces by U.S. state with business profiles and categories.\footnote{https://www.globalspec.com/local/6383/C\_US}

\begin{center}\rule{0.5\linewidth}{0.5pt}\end{center}

\section*{State-by-State Makerspace Index (United States)}\label{docs__pandoc__latex__src__3dmake_appendices_complete.md__state-by-state-makerspace-index-united-states}

Below is a national omnibus index generated from directories that explicitly list makerspaces by state. States not listed did not appear in any verified directory results.

\subsection*{Alabama}\label{docs__pandoc__latex__src__3dmake_appendices_complete.md__alabama}

\begin{itemize}
\tightlist
\item
  2 makerspaces listed on GlobalSpec.\footnote{https://www.globalspec.com/local/6383/C\_US}
\end{itemize}

\subsection*{Arizona}\label{docs__pandoc__latex__src__3dmake_appendices_complete.md__arizona}

\begin{itemize}
\tightlist
\item
  Makerspaces listed on Makerspace Directory and GlobalSpec. Example: Coco‑Op (Flagstaff).\footnote{https://makerspace.com/}
\end{itemize}

\subsection*{California}\label{docs__pandoc__latex__src__3dmake_appendices_complete.md__california}

\begin{itemize}
\tightlist
\item
  84 makerspaces (highest nationwide).\footnote{https://www.poidata.io/report/makerspace/united-states}
\item
  Multiple listings in Makerspace.com, GlobalSpec, and WorkMakeCreate.
\end{itemize}

\subsection*{Colorado}\label{docs__pandoc__latex__src__3dmake_appendices_complete.md__colorado}

\begin{itemize}
\tightlist
\item
  37 makerspaces statewide.\footnote{https://www.poidata.io/report/makerspace/united-states}
\item
  Makerspaces include Denhac (Denver).\footnote{https://www.globalspec.com/local/6383/C\_US}
\end{itemize}

\subsection*{Connecticut}\label{docs__pandoc__latex__src__3dmake_appendices_complete.md__connecticut}

\begin{itemize}
\tightlist
\item
  2 makerspaces listed on GlobalSpec (e.g., CT Hackerspace).\footnote{https://www.globalspec.com/local/6383/C\_US}
\end{itemize}

\subsection*{District of Columbia}\label{docs__pandoc__latex__src__3dmake_appendices_complete.md__district-of-columbia}

\begin{itemize}
\tightlist
\item
  3 makerspaces statewide.\footnote{https://www.poidata.io/report/makerspace/united-states}
\item
  Example: HacDC.\footnote{https://www.globalspec.com/local/6383/C\_US}
\end{itemize}

\subsection*{Florida}\label{docs__pandoc__latex__src__3dmake_appendices_complete.md__florida}

\begin{itemize}
\tightlist
\item
  45 makerspaces statewide.\footnote{https://www.poidata.io/report/makerspace/united-states}
\item
  Example: Marion County Maker Collective.\footnote{https://makerspace.com/}
\end{itemize}

\subsection*{Georgia}\label{docs__pandoc__latex__src__3dmake_appendices_complete.md__georgia}

\begin{itemize}
\tightlist
\item
  19 makerspaces statewide.\footnote{https://www.poidata.io/report/makerspace/united-states}
\item
  Example: SparkMacon.\footnote{https://makerspace.com/}
\end{itemize}

\subsection*{Hawaii}\label{docs__pandoc__latex__src__3dmake_appendices_complete.md__hawaii}

\begin{itemize}
\tightlist
\item
  2 makerspaces listed on GlobalSpec.\footnote{https://www.globalspec.com/local/6383/C\_US}
\end{itemize}

\subsection*{Idaho}\label{docs__pandoc__latex__src__3dmake_appendices_complete.md__idaho}

\begin{itemize}
\tightlist
\item
  5 makerspaces statewide (from POIData).\footnote{https://www.poidata.io/report/makerspace/united-states}
\end{itemize}

\subsection*{Illinois}\label{docs__pandoc__latex__src__3dmake_appendices_complete.md__illinois}

\begin{itemize}
\tightlist
\item
  4 makerspaces listed on GlobalSpec (e.g., mHUB Chicago).\footnote{https://www.globalspec.com/local/6383/C\_US}
\end{itemize}

\subsection*{Indiana}\label{docs__pandoc__latex__src__3dmake_appendices_complete.md__indiana}

\begin{itemize}
\tightlist
\item
  3 makerspaces listed on GlobalSpec; 22 statewide per POIData.\footnote{https://www.poidata.io/report/makerspace/united-states}
\end{itemize}

\subsection*{Iowa}\label{docs__pandoc__latex__src__3dmake_appendices_complete.md__iowa}

\begin{itemize}
\tightlist
\item
  11 statewide makerspaces.\footnote{https://www.poidata.io/report/makerspace/united-states}
\item
  Example: Cedar Valley Makers.\footnote{https://www.globalspec.com/local/6383/C\_US}
\end{itemize}

\subsection*{Kansas}\label{docs__pandoc__latex__src__3dmake_appendices_complete.md__kansas}

\begin{itemize}
\tightlist
\item
  14 statewide makerspaces.\footnote{https://www.poidata.io/report/makerspace/united-states}
\end{itemize}

\subsection*{Kentucky}\label{docs__pandoc__latex__src__3dmake_appendices_complete.md__kentucky}

\begin{itemize}
\tightlist
\item
  6 statewide makerspaces.\footnote{https://www.poidata.io/report/makerspace/united-states}
\end{itemize}

\subsection*{Louisiana}\label{docs__pandoc__latex__src__3dmake_appendices_complete.md__louisiana}

\begin{itemize}
\tightlist
\item
  3 statewide makerspaces.\footnote{https://www.poidata.io/report/makerspace/united-states}
\end{itemize}

\subsection*{Maine}\label{docs__pandoc__latex__src__3dmake_appendices_complete.md__maine}

\begin{itemize}
\tightlist
\item
  11 statewide makerspaces.\footnote{https://www.poidata.io/report/makerspace/united-states}
\end{itemize}

\subsection*{Maryland}\label{docs__pandoc__latex__src__3dmake_appendices_complete.md__maryland}

\begin{itemize}
\tightlist
\item
  12 statewide makerspaces.\footnote{https://www.poidata.io/report/makerspace/united-states}
\item
  Example: Baltimore Hackerspace.\footnote{https://www.globalspec.com/local/6383/C\_US}
\end{itemize}

\subsection*{Massachusetts}\label{docs__pandoc__latex__src__3dmake_appendices_complete.md__massachusetts}

\begin{itemize}
\tightlist
\item
  24 statewide makerspaces.\footnote{https://www.poidata.io/report/makerspace/united-states}
\item
  Example: Framingham Makerspace.\footnote{https://www.makerdirectory.com/makerspace-directory/}
\end{itemize}

\subsection*{Michigan}\label{docs__pandoc__latex__src__3dmake_appendices_complete.md__michigan}

\begin{itemize}
\tightlist
\item
  18 statewide makerspaces.\footnote{https://www.poidata.io/report/makerspace/united-states}
\item
  Large concentration documented by Maker Works directory.\footnote{https://www.maker-works.com/local-makerspaces}
\end{itemize}

\subsection*{Minnesota}\label{docs__pandoc__latex__src__3dmake_appendices_complete.md__minnesota}

\begin{itemize}
\tightlist
\item
  27 statewide makerspaces.\footnote{https://www.poidata.io/report/makerspace/united-states}
\end{itemize}

\subsection*{Mississippi}\label{docs__pandoc__latex__src__3dmake_appendices_complete.md__mississippi}

\begin{itemize}
\tightlist
\item
  3 statewide makerspaces.\footnote{https://www.poidata.io/report/makerspace/united-states}
\end{itemize}

\subsection*{Missouri}\label{docs__pandoc__latex__src__3dmake_appendices_complete.md__missouri}

\begin{itemize}
\tightlist
\item
  16 statewide makerspaces.\footnote{https://www.poidata.io/report/makerspace/united-states}
\item
  Example: Arch Reactor (St. Louis).\footnote{https://www.globalspec.com/local/6383/C\_US}
\end{itemize}

\subsection*{Montana}\label{docs__pandoc__latex__src__3dmake_appendices_complete.md__montana}

\begin{itemize}
\tightlist
\item
  9 statewide makerspaces.\footnote{https://www.poidata.io/report/makerspace/united-states}
\end{itemize}

\subsection*{Nebraska}\label{docs__pandoc__latex__src__3dmake_appendices_complete.md__nebraska}

\begin{itemize}
\tightlist
\item
  7 statewide makerspaces.\footnote{https://www.poidata.io/report/makerspace/united-states}
\end{itemize}

\subsection*{Nevada}\label{docs__pandoc__latex__src__3dmake_appendices_complete.md__nevada}

\begin{itemize}
\tightlist
\item
  6 statewide makerspaces.\footnote{https://www.poidata.io/report/makerspace/united-states}
\end{itemize}

\subsection*{New Hampshire}\label{docs__pandoc__latex__src__3dmake_appendices_complete.md__new-hampshire}

\begin{itemize}
\tightlist
\item
  4 statewide makerspaces.\footnote{https://www.poidata.io/report/makerspace/united-states}
\end{itemize}

\subsection*{New Jersey}\label{docs__pandoc__latex__src__3dmake_appendices_complete.md__new-jersey}

\begin{itemize}
\tightlist
\item
  12 statewide makerspaces.\footnote{https://www.poidata.io/report/makerspace/united-states}
\end{itemize}

\subsection*{New Mexico}\label{docs__pandoc__latex__src__3dmake_appendices_complete.md__new-mexico}

\begin{itemize}
\tightlist
\item
  8 statewide makerspaces.\footnote{https://www.poidata.io/report/makerspace/united-states}
\end{itemize}

\subsection*{New York}\label{docs__pandoc__latex__src__3dmake_appendices_complete.md__new-york}

\begin{itemize}
\tightlist
\item
  28 statewide makerspaces.\footnote{https://www.poidata.io/report/makerspace/united-states}
\item
  Dozens indexed by WorkMakeCreate (NYC region).\footnote{https://workmakecreate.com/}
\end{itemize}

\subsection*{North Carolina}\label{docs__pandoc__latex__src__3dmake_appendices_complete.md__north-carolina}

\begin{itemize}
\tightlist
\item
  28 statewide makerspaces.\footnote{https://www.poidata.io/report/makerspace/united-states}
\end{itemize}

\subsection*{North Dakota}\label{docs__pandoc__latex__src__3dmake_appendices_complete.md__north-dakota}

\begin{itemize}
\tightlist
\item
  \emph{(No verified listings in sources)}
\end{itemize}

\subsection*{Ohio}\label{docs__pandoc__latex__src__3dmake_appendices_complete.md__ohio}

\begin{itemize}
\tightlist
\item
  22 statewide makerspaces.\footnote{https://www.poidata.io/report/makerspace/united-states}
\item
  Example: Akron-Summit Public Library makerspace.\footnote{https://www.poidata.io/report/makerspace/united-states}
\end{itemize}

\subsection*{Oklahoma}\label{docs__pandoc__latex__src__3dmake_appendices_complete.md__oklahoma}

\begin{itemize}
\tightlist
\item
  8 statewide makerspaces.\footnote{https://www.poidata.io/report/makerspace/united-states}
\end{itemize}

\subsection*{Oregon}\label{docs__pandoc__latex__src__3dmake_appendices_complete.md__oregon}

\begin{itemize}
\tightlist
\item
  35 statewide makerspaces.\footnote{https://www.poidata.io/report/makerspace/united-states}
\end{itemize}

\subsection*{Pennsylvania}\label{docs__pandoc__latex__src__3dmake_appendices_complete.md__pennsylvania}

\begin{itemize}
\tightlist
\item
  15 statewide makerspaces.\footnote{https://www.poidata.io/report/makerspace/united-states}
\end{itemize}

\subsection*{Rhode Island}\label{docs__pandoc__latex__src__3dmake_appendices_complete.md__rhode-island}

\begin{itemize}
\tightlist
\item
  3 statewide makerspaces.\footnote{https://www.poidata.io/report/makerspace/united-states}
\item
  Example: Tinker\textbar Bristol.\footnote{https://makerspace.com/}
\end{itemize}

\subsection*{South Carolina}\label{docs__pandoc__latex__src__3dmake_appendices_complete.md__south-carolina}

\begin{itemize}
\tightlist
\item
  15 statewide makerspaces.\footnote{https://www.poidata.io/report/makerspace/united-states}
\end{itemize}

\subsection*{South Dakota}\label{docs__pandoc__latex__src__3dmake_appendices_complete.md__south-dakota}

\begin{itemize}
\tightlist
\item
  2 statewide makerspaces.\footnote{https://www.poidata.io/report/makerspace/united-states}
\end{itemize}

\subsection*{Tennessee}\label{docs__pandoc__latex__src__3dmake_appendices_complete.md__tennessee}

\begin{itemize}
\tightlist
\item
  16 statewide makerspaces; GlobalSpec lists 3 (e.g., CBI Maker Space).\footnote{https://www.globalspec.com/local/6383/C\_US}
\end{itemize}

\subsection*{Texas}\label{docs__pandoc__latex__src__3dmake_appendices_complete.md__texas}

\begin{itemize}
\tightlist
\item
  33 statewide makerspaces (3rd highest).\footnote{https://www.poidata.io/report/makerspace/united-states}
\item
  Example: 10BitWorks (San Antonio).\footnote{https://makerspace.com/directory/}
\end{itemize}

\subsection*{Utah}\label{docs__pandoc__latex__src__3dmake_appendices_complete.md__utah}

\begin{itemize}
\tightlist
\item
  9 statewide makerspaces.\footnote{https://www.poidata.io/report/makerspace/united-states}
\end{itemize}

\subsection*{Vermont}\label{docs__pandoc__latex__src__3dmake_appendices_complete.md__vermont}

\begin{itemize}
\tightlist
\item
  6 statewide makerspaces.\footnote{https://www.poidata.io/report/makerspace/united-states}
\end{itemize}

\subsection*{Virginia}\label{docs__pandoc__latex__src__3dmake_appendices_complete.md__virginia}

\begin{itemize}
\tightlist
\item
  19 statewide makerspaces.\footnote{https://www.poidata.io/report/makerspace/united-states}
\item
  Example: 757 Makerspace (Norfolk).\footnote{https://makerspace.com/directory/}
\end{itemize}

\subsection*{Washington}\label{docs__pandoc__latex__src__3dmake_appendices_complete.md__washington}

\begin{itemize}
\tightlist
\item
  28 statewide makerspaces.\footnote{https://www.poidata.io/report/makerspace/united-states}
\item
  Example: Confluent Space (Richland).\footnote{https://www.globalspec.com/local/6383/C\_US}
\end{itemize}

\subsection*{West Virginia}\label{docs__pandoc__latex__src__3dmake_appendices_complete.md__west-virginia}

\begin{itemize}
\tightlist
\item
  2 statewide makerspaces.\footnote{https://www.poidata.io/report/makerspace/united-states}
\end{itemize}

\subsection*{Wisconsin}\label{docs__pandoc__latex__src__3dmake_appendices_complete.md__wisconsin}

\begin{itemize}
\tightlist
\item
  17 statewide makerspaces.\footnote{https://www.poidata.io/report/makerspace/united-states}
\end{itemize}

\subsection*{Wyoming}\label{docs__pandoc__latex__src__3dmake_appendices_complete.md__wyoming}

\begin{itemize}
\tightlist
\item
  6 statewide makerspaces.\footnote{https://www.poidata.io/report/makerspace/united-states}
\end{itemize}

\begin{center}\rule{0.5\linewidth}{0.5pt}\end{center}

\section*{Endnotes}\label{docs__pandoc__latex__src__3dmake_appendices_complete.md__endnotes-1}

\begin{center}\rule{0.5\linewidth}{0.5pt}\end{center}

\protect\phantomsection\label{docs__pandoc__latex__src__3dmake_appendices_complete.md__appendix-a}{}

\section*{Appendix A: 3DMake CLI Cheat Sheet}\label{docs__pandoc__latex__src__3dmake_appendices_complete.md__appendix-a-3dmake-cli-cheat-sheet}

\begin{center}\rule{0.5\linewidth}{0.5pt}\end{center}

\begin{quote}
A comprehensive reference for the \texttt{3dm} command-line tool used alongside OpenSCAD for 3D printing.
\end{quote}

\begin{center}\rule{0.5\linewidth}{0.5pt}\end{center}

\subsection*{Basic Syntax}\label{docs__pandoc__latex__src__3dmake_appendices_complete.md__basic-syntax}

\begin{lstlisting}[style=Alabaster]
3dm ACTIONS... [OPTIONS]... [INPUT_FILE]

\end{lstlisting}

You can combine multiple actions in a single command, and most single-letter options can be used with a single dash (e.g. \texttt{-s\ 1.5} instead of \texttt{-\/-scale\ 1.5}).

\begin{center}\rule{0.5\linewidth}{0.5pt}\end{center}

\subsection*{Quick Examples}\label{docs__pandoc__latex__src__3dmake_appendices_complete.md__quick-examples}

\begin{lstlisting}[style=Alabaster, language=bash]
3dm build                                    # Build the default model
3dm build orient slice                       # Build, auto-orient, then slice
3dm build orient slice -m cover -o supports  # With a specific model and overlay
3dm info alpaca.stl                          # Get info about an STL file
3dm preview alpaca.stl                       # Preview an STL file
3dm slice print alpaca.stl                   # Slice and send an STL to printer
3dm image --angle top --angle front          # Export rendered images

\end{lstlisting}

\begin{center}\rule{0.5\linewidth}{0.5pt}\end{center}

\subsection*{All Actions}\label{docs__pandoc__latex__src__3dmake_appendices_complete.md__all-actions}

\subsubsection*{Project Setup \& Management}\label{docs__pandoc__latex__src__3dmake_appendices_complete.md__project-setup--management}

{\def\LTcaptype{none} % do not increment counter
\begin{longtable}[]{@{}
  >{\raggedright\arraybackslash}p{(\linewidth - 2\tabcolsep) * \real{0.0522}}
  >{\raggedright\arraybackslash}p{(\linewidth - 2\tabcolsep) * \real{0.9478}}@{}}
\toprule\noalign{}
\begin{minipage}[b]{\linewidth}\raggedright
Action
\end{minipage} & \begin{minipage}[b]{\linewidth}\raggedright
Description
\end{minipage} \\
\midrule\noalign{}
\endhead
\bottomrule\noalign{}
\endlastfoot
\texttt{setup} &
Set up 3DMake for the first time, or overwrite/upgrade existing settings (printer profile, Gemini AI key, OctoPrint connection) \\
\texttt{new} &
Interactively create a new 3DMake project directory structure (\texttt{src/}, \texttt{build/}, \texttt{3dmake.toml}, starter \texttt{.scad} file) \\
\end{longtable}
}

\begin{lstlisting}[style=Alabaster, language=bash]
3dm setup       # First-time setup or re-configure
3dm new         # Create a new project (prompts for directory name)

\end{lstlisting}

\begin{center}\rule{0.5\linewidth}{0.5pt}\end{center}

\subsubsection*{Build Pipeline}\label{docs__pandoc__latex__src__3dmake_appendices_complete.md__build-pipeline}

These actions can be chained together in sequence:

{\def\LTcaptype{none} % do not increment counter
\begin{longtable}[]{@{}
  >{\raggedright\arraybackslash}p{(\linewidth - 2\tabcolsep) * \real{0.0930}}
  >{\raggedright\arraybackslash}p{(\linewidth - 2\tabcolsep) * \real{0.9070}}@{}}
\toprule\noalign{}
\begin{minipage}[b]{\linewidth}\raggedright
Action
\end{minipage} & \begin{minipage}[b]{\linewidth}\raggedright
Description
\end{minipage} \\
\midrule\noalign{}
\endhead
\bottomrule\noalign{}
\endlastfoot
\texttt{build} &
Compile the OpenSCAD \texttt{.scad} source file into a \texttt{.stl} model \\
\texttt{orient} &
Auto-orient the model to minimize the need for supports (uses Tweaker3) \\
\texttt{slice} &
Slice the model into printer-ready G-code using PrusaSlicer \\
\texttt{print} &
Send the sliced G-code to the printer (OctoPrint, Bambu LAN, or Bambu Connect) \\
\end{longtable}
}

\begin{lstlisting}[style=Alabaster, language=bash]
3dm build                        # Just compile OpenSCAD -> STL
3dm build orient                 # Compile then auto-orient
3dm build orient slice           # Compile, orient, then slice
3dm build orient slice print     # Full pipeline: compile -> orient -> slice -> print
3dm slice print                  # Slice an already-built model and print it
3dm print                        # Send an already-sliced file to the printer

\end{lstlisting}

\begin{quote}
\textbf{Note:} Actions automatically trigger their dependencies. For example, \texttt{slice} will imply \texttt{build} if needed.
\end{quote}

\begin{center}\rule{0.5\linewidth}{0.5pt}\end{center}

\subsubsection*{Inspection \& Visualization}\label{docs__pandoc__latex__src__3dmake_appendices_complete.md__inspection--visualization}

{\def\LTcaptype{none} % do not increment counter
\begin{longtable}[]{@{}
  >{\raggedright\arraybackslash}p{(\linewidth - 2\tabcolsep) * \real{0.0612}}
  >{\raggedright\arraybackslash}p{(\linewidth - 2\tabcolsep) * \real{0.9388}}@{}}
\toprule\noalign{}
\begin{minipage}[b]{\linewidth}\raggedright
Action
\end{minipage} & \begin{minipage}[b]{\linewidth}\raggedright
Description
\end{minipage} \\
\midrule\noalign{}
\endhead
\bottomrule\noalign{}
\endlastfoot
\texttt{info} &
Display mesh dimensions and center point; optionally get an AI description of the model (requires Gemini or OpenRouter API key configured) \\
\texttt{preview} &
Produce a 2D silhouette projection of the model (outputs a flat STL for quick reference) \\
\texttt{image} &
Export one or more rendered PNG images of the model from specified viewpoints \\
\end{longtable}
}

\begin{lstlisting}[style=Alabaster, language=bash]
3dm info                         # Show mesh size/center info
3dm info --interactive           # Info + start AI chat session about the model
3dm info model.stl               # Info on a specific STL file
3dm preview                      # Generate 3-silhouette projection (default view)
3dm preview --view topsil        # Generate top-only silhouette
3dm image --angle top            # Render image from top viewpoint
3dm image --angle top --angle front --angle left  # Multiple viewpoints
3dm image --colorscheme slicer_dark --image-size 1920x1080

\end{lstlisting}

\textbf{Available \texttt{-\/-view} options for \texttt{preview}:}

{\def\LTcaptype{none} % do not increment counter
\begin{longtable}[]{@{}ll@{}}
\toprule\noalign{}
View Name & Description \\
\midrule\noalign{}
\endhead
\bottomrule\noalign{}
\endlastfoot
\texttt{3sil} & Three-view silhouette (top, left, front) --- default \\
\texttt{topsil} & Top-down silhouette only \\
\texttt{leftsil} & Left-side silhouette only \\
\texttt{rightsil} & Right-side silhouette only \\
\texttt{frontsil} & Front silhouette only \\
\texttt{backsil} & Back silhouette only \\
\end{longtable}
}

\textbf{Available \texttt{-\/-colorscheme} options for \texttt{image}:}

{\def\LTcaptype{none} % do not increment counter
\begin{longtable}[]{@{}ll@{}}
\toprule\noalign{}
Scheme Name & Description \\
\midrule\noalign{}
\endhead
\bottomrule\noalign{}
\endlastfoot
\texttt{slicer\_light} & Orange model on light gray background \\
\texttt{slicer\_dark} & Orange model on dark background \\
\texttt{light\_on\_dark} & Floral white model on dark background \\
\end{longtable}
}

\textbf{Available \texttt{-\/-angle} options for \texttt{image}:}\\
\texttt{above\_front\_left}, \texttt{above\_front\_right}, \texttt{above\_back\_left}, \texttt{above\_back\_right}, \texttt{top}, \texttt{bottom}

\begin{center}\rule{0.5\linewidth}{0.5pt}\end{center}

\subsubsection*{Configuration Editors}\label{docs__pandoc__latex__src__3dmake_appendices_complete.md__configuration-editors}

These are all isolated actions (cannot be combined with other actions):

{\def\LTcaptype{none} % do not increment counter
\begin{longtable}[]{@{}
  >{\raggedright\arraybackslash}p{(\linewidth - 2\tabcolsep) * \real{0.2062}}
  >{\raggedright\arraybackslash}p{(\linewidth - 2\tabcolsep) * \real{0.7938}}@{}}
\toprule\noalign{}
\begin{minipage}[b]{\linewidth}\raggedright
Action
\end{minipage} & \begin{minipage}[b]{\linewidth}\raggedright
Description
\end{minipage} \\
\midrule\noalign{}
\endhead
\bottomrule\noalign{}
\endlastfoot
\texttt{edit-model} &
Open the current model\textquotesingle s \texttt{.scad} source file in your configured editor \\
\texttt{edit-profile} &
Open the current printer profile \texttt{.ini} file in your editor \\
\texttt{edit-overlay} &
Open (or create) a slicer settings overlay file in your editor \\
\texttt{edit-profile-gcode} &
Interactively select and edit a G-code script embedded in the printer profile \\
\texttt{edit-prompt} &
Edit the AI prompt used by the \texttt{info} command \\
\texttt{edit-global-config} &
Edit the global \texttt{defaults.toml} settings file (printer, API keys, etc.) \\
\end{longtable}
}

\begin{lstlisting}[style=Alabaster, language=bash]
3dm edit-model                    # Edit the default model's .scad file
3dm edit-model -m lid             # Edit a specific model (lid.scad)
3dm edit-profile                  # Edit the current printer profile
3dm edit-profile -p prusa_mk4     # Edit a specific profile
3dm edit-overlay                  # Prompted to choose/create an overlay
3dm edit-overlay -o supports      # Edit the "supports" overlay directly
3dm edit-profile-gcode            # Choose and edit a G-code script in the profile
3dm edit-prompt                   # Edit the AI description prompt
3dm edit-global-config            # Edit global settings (defaults.toml)

\end{lstlisting}

\begin{center}\rule{0.5\linewidth}{0.5pt}\end{center}

\subsubsection*{Library Management}\label{docs__pandoc__latex__src__3dmake_appendices_complete.md__library-management}

{\def\LTcaptype{none} % do not increment counter
\begin{longtable}[]{@{}
  >{\raggedright\arraybackslash}p{(\linewidth - 2\tabcolsep) * \real{0.1597}}
  >{\raggedright\arraybackslash}p{(\linewidth - 2\tabcolsep) * \real{0.8403}}@{}}
\toprule\noalign{}
\begin{minipage}[b]{\linewidth}\raggedright
Action
\end{minipage} & \begin{minipage}[b]{\linewidth}\raggedright
Description
\end{minipage} \\
\midrule\noalign{}
\endhead
\bottomrule\noalign{}
\endlastfoot
\texttt{list-libraries} &
List all available OpenSCAD libraries in the catalog, showing which are installed and their versions \\
\texttt{install-libraries} &
Download and install any libraries required by the current project (as specified in \texttt{3dmake.toml}) \\
\end{longtable}
}

\begin{lstlisting}[style=Alabaster, language=bash]
3dm list-libraries        # See available and installed libraries
3dm install-libraries     # Install missing project libraries

\end{lstlisting}

\begin{center}\rule{0.5\linewidth}{0.5pt}\end{center}

\subsubsection*{Listing Configuration}\label{docs__pandoc__latex__src__3dmake_appendices_complete.md__listing-configuration}

{\def\LTcaptype{none} % do not increment counter
\begin{longtable}[]{@{}ll@{}}
\toprule\noalign{}
Action & Description \\
\midrule\noalign{}
\endhead
\bottomrule\noalign{}
\endlastfoot
\texttt{list-profiles} & List all available printer profiles \\
\texttt{list-overlays} & List all available slicer setting overlays \\
\end{longtable}
}

\begin{lstlisting}[style=Alabaster, language=bash]
3dm list-profiles         # See all printer profiles
3dm list-overlays         # See all overlays (with profile restrictions if any)

\end{lstlisting}

\begin{center}\rule{0.5\linewidth}{0.5pt}\end{center}

\subsubsection*{Diagnostics \& Info}\label{docs__pandoc__latex__src__3dmake_appendices_complete.md__diagnostics--info}

{\def\LTcaptype{none} % do not increment counter
\begin{longtable}[]{@{}
  >{\raggedright\arraybackslash}p{(\linewidth - 2\tabcolsep) * \real{0.1573}}
  >{\raggedright\arraybackslash}p{(\linewidth - 2\tabcolsep) * \real{0.8427}}@{}}
\toprule\noalign{}
\begin{minipage}[b]{\linewidth}\raggedright
Action
\end{minipage} & \begin{minipage}[b]{\linewidth}\raggedright
Description
\end{minipage} \\
\midrule\noalign{}
\endhead
\bottomrule\noalign{}
\endlastfoot
\texttt{test-connect} &
Test the connection to the configured print server (OctoPrint or Bambu LAN) \\
\texttt{help} & Display the help message with all actions and options \\
\texttt{version} &
Print the 3DMake version, program location, and config directory \\
\end{longtable}
}

\begin{lstlisting}[style=Alabaster, language=bash]
3dm test-connect          # Test printer connection
3dm help                  # Show help
3dm version               # Show version info

\end{lstlisting}

\begin{center}\rule{0.5\linewidth}{0.5pt}\end{center}

\subsection*{All Options}\label{docs__pandoc__latex__src__3dmake_appendices_complete.md__all-options}

{\def\LTcaptype{none} % do not increment counter
\begin{longtable}[]{@{}
  >{\raggedright\arraybackslash}p{(\linewidth - 6\tabcolsep) * \real{0.1345}}
  >{\raggedright\arraybackslash}p{(\linewidth - 6\tabcolsep) * \real{0.0420}}
  >{\raggedright\arraybackslash}p{(\linewidth - 6\tabcolsep) * \real{0.1429}}
  >{\raggedright\arraybackslash}p{(\linewidth - 6\tabcolsep) * \real{0.6807}}@{}}
\toprule\noalign{}
\begin{minipage}[b]{\linewidth}\raggedright
Option
\end{minipage} & \begin{minipage}[b]{\linewidth}\raggedright
Short
\end{minipage} & \begin{minipage}[b]{\linewidth}\raggedright
Type
\end{minipage} & \begin{minipage}[b]{\linewidth}\raggedright
Description
\end{minipage} \\
\midrule\noalign{}
\endhead
\bottomrule\noalign{}
\endlastfoot
\texttt{-\/-scale} & \texttt{-s} & \texttt{float} or \texttt{auto} &
Scale the model by a decimal factor (e.g. \texttt{1.5} for 150\%, \texttt{0.5} for 50\%) \\
\texttt{-\/-model} & \texttt{-m} & \texttt{string} &
Select a specific model by name in a multi-model project (omit \texttt{.scad} extension) \\
\texttt{-\/-profile} & \texttt{-p} & \texttt{string} &
Select a printer profile by name \\
\texttt{-\/-overlay} & \texttt{-o} & \texttt{string} &
Apply a slicer settings overlay; can be used multiple times \\
\texttt{-\/-view} & \texttt{-v} & \texttt{string} &
Select the preview view type (see view options above) \\
\texttt{-\/-angle} & \texttt{-a} & \texttt{string} &
Set a viewpoint for image export; can be used multiple times \\
\texttt{-\/-copies} & \texttt{-c} & \texttt{int} &
Number of copies to slice/print (default: \texttt{1}) \\
\texttt{-\/-interactive} & \texttt{-i} & flag &
Start an interactive AI chat after the \texttt{info} command \\
\texttt{-\/-colorscheme} & --- & \texttt{string} &
Color scheme name for image export \\
\texttt{-\/-image-size} & --- & \texttt{WxH} &
Image dimensions in pixels (e.g. \texttt{1920x1080}, default: \texttt{1080x720}) \\
\texttt{-\/-debug} & --- & flag &
Enable debug output (shows full subprocess logs, raises exceptions) \\
\texttt{-\/-help} / \texttt{-\/-h} & \texttt{-h} & flag &
Display help and exit \\
\end{longtable}
}

\begin{center}\rule{0.5\linewidth}{0.5pt}\end{center}

\subsection*{Input Files}\label{docs__pandoc__latex__src__3dmake_appendices_complete.md__input-files}

You can pass an \texttt{.stl} or \texttt{.scad} file directly instead of running inside a project directory:

\begin{lstlisting}[style=Alabaster, language=bash]
3dm info model.stl                  # Inspect a standalone STL
3dm preview model.stl               # Preview a standalone STL
3dm slice print model.stl           # Slice and print a standalone STL
3dm build model.scad                # Build a standalone SCAD file -> STL
3dm image model.stl --angle top     # Render image from standalone STL

\end{lstlisting}

\begin{quote}
When passing a \texttt{.scad} file, \texttt{build} is automatically added to the action list.
\end{quote}

\begin{center}\rule{0.5\linewidth}{0.5pt}\end{center}

\subsection*{\texorpdfstring{Project File (\texttt{3dmake.toml}) \& Global Config (\texttt{defaults.toml})}{Project File (3dmake.toml) \& Global Config (defaults.toml)}}\label{docs__pandoc__latex__src__3dmake_appendices_complete.md__project-file-3dmaketoml--global-config-defaultstoml}

Both files use identical TOML keys. \texttt{defaults.toml} lives in the config directory and sets machine-wide defaults. \texttt{3dmake.toml} lives in your project root and \textbf{overrides} the global defaults for that project only.

\begin{quote}
Run \texttt{3dm\ edit-global-config} to open \texttt{defaults.toml} in your editor.\\
Settings set via CLI flags (e.g. \texttt{-\/-scale}, \texttt{-\/-profile}) override both files at runtime.
\end{quote}

\begin{center}\rule{0.5\linewidth}{0.5pt}\end{center}

\subsubsection*{General / Project Settings}\label{docs__pandoc__latex__src__3dmake_appendices_complete.md__general--project-settings}

{\def\LTcaptype{none} % do not increment counter
\begin{longtable}[]{@{}
  >{\raggedright\arraybackslash}p{(\linewidth - 6\tabcolsep) * \real{0.1471}}
  >{\raggedright\arraybackslash}p{(\linewidth - 6\tabcolsep) * \real{0.0441}}
  >{\raggedright\arraybackslash}p{(\linewidth - 6\tabcolsep) * \real{0.1103}}
  >{\raggedright\arraybackslash}p{(\linewidth - 6\tabcolsep) * \real{0.6985}}@{}}
\toprule\noalign{}
\begin{minipage}[b]{\linewidth}\raggedright
Key
\end{minipage} & \begin{minipage}[b]{\linewidth}\raggedright
Type
\end{minipage} & \begin{minipage}[b]{\linewidth}\raggedright
Default
\end{minipage} & \begin{minipage}[b]{\linewidth}\raggedright
Description
\end{minipage} \\
\midrule\noalign{}
\endhead
\bottomrule\noalign{}
\endlastfoot
\texttt{model\_name} & string & \texttt{"main"} &
The default model name. Maps to \texttt{src/\textless{}name\textgreater{}.scad} and \texttt{build/\textless{}name\textgreater{}.stl} \\
\texttt{project\_name} & string & \emph{(folder name)} &
Project name used as a prefix on output filenames. Auto-detected from directory name if not set \\
\texttt{scale} & float & \texttt{1.0} &
Scale factor applied during slicing. \texttt{1.0} = 100\%, \texttt{1.5} = 150\%, \texttt{0.5} = 50\% \\
\texttt{copies} & int & \texttt{1} &
Number of copies to slice/print in one run \\
\texttt{strict\_warnings} & bool & \texttt{false} &
Pass \texttt{-\/-hardwarnings} to OpenSCAD --- treat all warnings as errors during \texttt{build} \\
\texttt{min\_3dmake\_version} & string & \emph{(none)} &
Minimum required 3DMake version to run this project. Prevents running with outdated installs \\
\texttt{interactive} & bool & \texttt{false} &
Default to interactive AI chat mode after \texttt{info} \\
\texttt{debug} & bool & \texttt{false} &
Enable debug/verbose output globally \\
\end{longtable}
}

\begin{lstlisting}[style=Alabaster]
# General project settings
model_name = "main"
project_name = "my-widget"
scale = 1.0
copies = 1
strict_warnings = true
min_3dmake_version = "1.2.0"
interactive = false
debug = false

\end{lstlisting}

\begin{center}\rule{0.5\linewidth}{0.5pt}\end{center}

\subsubsection*{OpenSCAD / Build Settings}\label{docs__pandoc__latex__src__3dmake_appendices_complete.md__openscad--build-settings}

{\def\LTcaptype{none} % do not increment counter
\begin{longtable}[]{@{}
  >{\raggedright\arraybackslash}p{(\linewidth - 6\tabcolsep) * \real{0.0895}}
  >{\raggedright\arraybackslash}p{(\linewidth - 6\tabcolsep) * \real{0.0789}}
  >{\raggedright\arraybackslash}p{(\linewidth - 6\tabcolsep) * \real{0.0368}}
  >{\raggedright\arraybackslash}p{(\linewidth - 6\tabcolsep) * \real{0.7947}}@{}}
\toprule\noalign{}
\begin{minipage}[b]{\linewidth}\raggedright
Key
\end{minipage} & \begin{minipage}[b]{\linewidth}\raggedright
Type
\end{minipage} & \begin{minipage}[b]{\linewidth}\raggedright
Default
\end{minipage} & \begin{minipage}[b]{\linewidth}\raggedright
Description
\end{minipage} \\
\midrule\noalign{}
\endhead
\bottomrule\noalign{}
\endlastfoot
\texttt{libraries} & list of strings & \texttt{{[}{]}} &
OpenSCAD libraries required by this project. Names must be \textbf{lowercase} and match entries in the library catalog. Install with \texttt{3dm\ install-libraries} \\
\texttt{local\_libraries} & list of strings & \texttt{{[}{]}} &
Paths to local library directories (absolute or relative to project root) to add to \texttt{OPENSCADPATH} during build \\
\end{longtable}
}

\begin{lstlisting}[style=Alabaster]
# OpenSCAD library dependencies
libraries = ["openscad-threads", "bosl2"]

# Local/custom libraries not in the catalog
local_libraries = ["./lib/my-custom-parts", "/home/user/shared-scad-libs"]

\end{lstlisting}

\begin{center}\rule{0.5\linewidth}{0.5pt}\end{center}

\subsubsection*{Printer \& Slicer Settings}\label{docs__pandoc__latex__src__3dmake_appendices_complete.md__printer--slicer-settings}

{\def\LTcaptype{none} % do not increment counter
\begin{longtable}[]{@{}
  >{\raggedright\arraybackslash}p{(\linewidth - 6\tabcolsep) * \real{0.0960}}
  >{\raggedright\arraybackslash}p{(\linewidth - 6\tabcolsep) * \real{0.0847}}
  >{\raggedright\arraybackslash}p{(\linewidth - 6\tabcolsep) * \real{0.0904}}
  >{\raggedright\arraybackslash}p{(\linewidth - 6\tabcolsep) * \real{0.7288}}@{}}
\toprule\noalign{}
\begin{minipage}[b]{\linewidth}\raggedright
Key
\end{minipage} & \begin{minipage}[b]{\linewidth}\raggedright
Type
\end{minipage} & \begin{minipage}[b]{\linewidth}\raggedright
Default
\end{minipage} & \begin{minipage}[b]{\linewidth}\raggedright
Description
\end{minipage} \\
\midrule\noalign{}
\endhead
\bottomrule\noalign{}
\endlastfoot
\texttt{printer\_profile} & string & \emph{(required)} &
Name of the printer profile \texttt{.ini} file to use (without the \texttt{.ini} extension). Set during \texttt{3dm\ setup} \\
\texttt{overlays} & list of strings & \texttt{{[}{]}} &
Default slicer setting overlays to apply on every slice. Profile-specific versions are preferred over default ones when available \\
\texttt{view} & string & \texttt{"3sil"} &
Default preview view type used by the \texttt{preview} action \\
\texttt{image\_angles} & list of strings & \texttt{{[}{]}} &
Default viewpoint angles for the \texttt{image} action \\
\texttt{colorscheme} & string & \texttt{"slicer\_light"} &
Default color scheme for the \texttt{image} action \\
\texttt{image\_size} & string & \texttt{"1080x720"} &
Default image export dimensions in \texttt{WIDTHxHEIGHT} pixels \\
\end{longtable}
}

\begin{lstlisting}[style=Alabaster]
# Slicer configuration
printer_profile = "prusa_mk4"
overlays = ["fast-draft"]

# Preview and image defaults
view = "3sil"
image_angles = ["above_front_left", "top"]
colorscheme = "slicer_light"
image_size = "1920x1080"

\end{lstlisting}

\begin{center}\rule{0.5\linewidth}{0.5pt}\end{center}

\subsubsection*{Print Sending Settings}\label{docs__pandoc__latex__src__3dmake_appendices_complete.md__print-sending-settings}

{\def\LTcaptype{none} % do not increment counter
\begin{longtable}[]{@{}
  >{\raggedright\arraybackslash}p{(\linewidth - 6\tabcolsep) * \real{0.1397}}
  >{\raggedright\arraybackslash}p{(\linewidth - 6\tabcolsep) * \real{0.0441}}
  >{\raggedright\arraybackslash}p{(\linewidth - 6\tabcolsep) * \real{0.1544}}
  >{\raggedright\arraybackslash}p{(\linewidth - 6\tabcolsep) * \real{0.6618}}@{}}
\toprule\noalign{}
\begin{minipage}[b]{\linewidth}\raggedright
Key
\end{minipage} & \begin{minipage}[b]{\linewidth}\raggedright
Type
\end{minipage} & \begin{minipage}[b]{\linewidth}\raggedright
Default
\end{minipage} & \begin{minipage}[b]{\linewidth}\raggedright
Description
\end{minipage} \\
\midrule\noalign{}
\endhead
\bottomrule\noalign{}
\endlastfoot
\texttt{print\_mode} & string & \emph{(required to print)} &
How to send prints to the printer. One of: \texttt{"octoprint"}, \texttt{"bambu\_lan"}, \texttt{"bambu\_connect"} \\
\texttt{auto\_start\_prints} & bool & \texttt{true} &
When using OctoPrint, automatically start the print after uploading the file \\
\end{longtable}
}

\begin{lstlisting}[style=Alabaster]
# Print sending mode
print_mode = "octoprint"       # Options: "octoprint", "bambu_lan", "bambu_connect"
auto_start_prints = true

\end{lstlisting}

\begin{center}\rule{0.5\linewidth}{0.5pt}\end{center}

\subsubsection*{OctoPrint Connection}\label{docs__pandoc__latex__src__3dmake_appendices_complete.md__octoprint-connection}

Used when \texttt{print\_mode\ =\ "octoprint"}. Test with \texttt{3dm\ test-connect}.

{\def\LTcaptype{none} % do not increment counter
\begin{longtable}[]{@{}
  >{\raggedright\arraybackslash}p{(\linewidth - 6\tabcolsep) * \real{0.1135}}
  >{\raggedright\arraybackslash}p{(\linewidth - 6\tabcolsep) * \real{0.0426}}
  >{\raggedright\arraybackslash}p{(\linewidth - 6\tabcolsep) * \real{0.0567}}
  >{\raggedright\arraybackslash}p{(\linewidth - 6\tabcolsep) * \real{0.7872}}@{}}
\toprule\noalign{}
\begin{minipage}[b]{\linewidth}\raggedright
Key
\end{minipage} & \begin{minipage}[b]{\linewidth}\raggedright
Type
\end{minipage} & \begin{minipage}[b]{\linewidth}\raggedright
Default
\end{minipage} & \begin{minipage}[b]{\linewidth}\raggedright
Description
\end{minipage} \\
\midrule\noalign{}
\endhead
\bottomrule\noalign{}
\endlastfoot
\texttt{octoprint\_host} & string & \emph{(none)} &
Full URL to your OctoPrint server, including protocol (e.g. \texttt{"http://octopi.local"} or \texttt{"http://192.168.1.50"}) \\
\texttt{octoprint\_key} & string & \emph{(none)} &
OctoPrint Application API key. Generate in OctoPrint \textgreater{} Settings \textgreater{} Application Keys \\
\end{longtable}
}

\begin{lstlisting}[style=Alabaster]
# OctoPrint connection (set in defaults.toml, not project toml)
octoprint_host = "http://octopi.local"
octoprint_key = "ABC123DEF456..."

\end{lstlisting}

\begin{center}\rule{0.5\linewidth}{0.5pt}\end{center}

\subsubsection*{Bambu Labs Connection}\label{docs__pandoc__latex__src__3dmake_appendices_complete.md__bambu-labs-connection}

Used when \texttt{print\_mode\ =\ "bambu\_lan"}. Test with \texttt{3dm\ test-connect}.

{\def\LTcaptype{none} % do not increment counter
\begin{longtable}[]{@{}
  >{\raggedright\arraybackslash}p{(\linewidth - 6\tabcolsep) * \real{0.1858}}
  >{\raggedright\arraybackslash}p{(\linewidth - 6\tabcolsep) * \real{0.0531}}
  >{\raggedright\arraybackslash}p{(\linewidth - 6\tabcolsep) * \real{0.0708}}
  >{\raggedright\arraybackslash}p{(\linewidth - 6\tabcolsep) * \real{0.6903}}@{}}
\toprule\noalign{}
\begin{minipage}[b]{\linewidth}\raggedright
Key
\end{minipage} & \begin{minipage}[b]{\linewidth}\raggedright
Type
\end{minipage} & \begin{minipage}[b]{\linewidth}\raggedright
Default
\end{minipage} & \begin{minipage}[b]{\linewidth}\raggedright
Description
\end{minipage} \\
\midrule\noalign{}
\endhead
\bottomrule\noalign{}
\endlastfoot
\texttt{bambu\_host} & string & \emph{(none)} &
Local IP address of your Bambu printer on your network (e.g. \texttt{"192.168.1.42"}) \\
\texttt{bambu\_serial\_number} & string & \emph{(none)} &
Printer serial number (found in the printer\textquotesingle s settings screen or Bambu Studio) \\
\texttt{bambu\_access\_code} & string & \emph{(none)} &
Local network access code shown on the printer screen (Settings \textgreater{} WLAN) \\
\end{longtable}
}

\begin{lstlisting}[style=Alabaster]
# Bambu Labs LAN connection (set in defaults.toml, not project toml)
bambu_host = "192.168.1.42"
bambu_serial_number = "01S00C123456789"
bambu_access_code = "12345678"

\end{lstlisting}

\begin{center}\rule{0.5\linewidth}{0.5pt}\end{center}

\subsubsection*{AI Description Settings}\label{docs__pandoc__latex__src__3dmake_appendices_complete.md__ai-description-settings}

Used by the \texttt{info} action. Set up during \texttt{3dm\ setup} or via \texttt{3dm\ edit-global-config}.

{\def\LTcaptype{none} % do not increment counter
\begin{longtable}[]{@{}
  >{\raggedright\arraybackslash}p{(\linewidth - 6\tabcolsep) * \real{0.0766}}
  >{\raggedright\arraybackslash}p{(\linewidth - 6\tabcolsep) * \real{0.0287}}
  >{\raggedright\arraybackslash}p{(\linewidth - 6\tabcolsep) * \real{0.0957}}
  >{\raggedright\arraybackslash}p{(\linewidth - 6\tabcolsep) * \real{0.7990}}@{}}
\toprule\noalign{}
\begin{minipage}[b]{\linewidth}\raggedright
Key
\end{minipage} & \begin{minipage}[b]{\linewidth}\raggedright
Type
\end{minipage} & \begin{minipage}[b]{\linewidth}\raggedright
Default
\end{minipage} & \begin{minipage}[b]{\linewidth}\raggedright
Description
\end{minipage} \\
\midrule\noalign{}
\endhead
\bottomrule\noalign{}
\endlastfoot
\texttt{gemini\_key} & string & \emph{(none)} &
Google Gemini API key for AI model descriptions. Free tier allows \textasciitilde{}50 requests/day. Get one at \href{https://aistudio.google.com/app/apikey}{aistudio.google.com/app/apikey} \\
\texttt{openrouter\_key} & string & \emph{(none)} &
OpenRouter API key (alternative to Gemini for AI descriptions, supports many LLM providers) \\
\texttt{llm\_name} & string & \emph{(provider default)} &
Specific LLM model name to use with Gemini or OpenRouter (e.g. \texttt{"gemini-1.5-flash"}, \texttt{"google/gemini-pro-vision"}) \\
\end{longtable}
}

\begin{lstlisting}[style=Alabaster]
# AI description keys (set in defaults.toml only — keep out of project toml!)
gemini_key = "AIzaSy..."
openrouter_key = "sk-or-..."
llm_name = "gemini-1.5-flash"

\end{lstlisting}

\begin{center}\rule{0.5\linewidth}{0.5pt}\end{center}

\subsubsection*{Editor Settings}\label{docs__pandoc__latex__src__3dmake_appendices_complete.md__editor-settings}

{\def\LTcaptype{none} % do not increment counter
\begin{longtable}[]{@{}
  >{\raggedright\arraybackslash}p{(\linewidth - 6\tabcolsep) * \real{0.1250}}
  >{\raggedright\arraybackslash}p{(\linewidth - 6\tabcolsep) * \real{0.0375}}
  >{\raggedright\arraybackslash}p{(\linewidth - 6\tabcolsep) * \real{0.1125}}
  >{\raggedright\arraybackslash}p{(\linewidth - 6\tabcolsep) * \real{0.7250}}@{}}
\toprule\noalign{}
\begin{minipage}[b]{\linewidth}\raggedright
Key
\end{minipage} & \begin{minipage}[b]{\linewidth}\raggedright
Type
\end{minipage} & \begin{minipage}[b]{\linewidth}\raggedright
Default
\end{minipage} & \begin{minipage}[b]{\linewidth}\raggedright
Description
\end{minipage} \\
\midrule\noalign{}
\endhead
\bottomrule\noalign{}
\endlastfoot
\texttt{editor} & string & \emph{(system default)} &
Command used to open files for editing (e.g. \texttt{"code"}, \texttt{"vim"}, \texttt{"notepad"}, \texttt{"subl"}). Used by all \texttt{edit-*} actions \\
\texttt{edit\_in\_background} & bool & \texttt{false} &
Launch the editor as a background process (non-blocking). Set to \texttt{true} on Windows where GUI editors are common \\
\end{longtable}
}

\begin{lstlisting}[style=Alabaster]
# Editor configuration (typically set in defaults.toml)
editor = "code"                 # VS Code
edit_in_background = true       # Don't block the terminal while editing

# Other common editor values:
# editor = "vim"
# editor = "nano"
# editor = "subl"               # Sublime Text
# editor = "notepad"            # Windows Notepad

\end{lstlisting}

\begin{center}\rule{0.5\linewidth}{0.5pt}\end{center}

\subsubsection*{Complete Annotated Example}\label{docs__pandoc__latex__src__3dmake_appendices_complete.md__complete-annotated-example}

A fully commented \texttt{3dmake.toml} for a real project:

\begin{lstlisting}[style=Alabaster]
# ── Project identity ──────────────────────────────────────────────
model_name = "enclosure"           # src/enclosure.scad -> build/enclosure.stl
project_name = "pi-case"           # Output files named pi-case-enclosure.gcode etc.
min_3dmake_version = "1.2.0"       # Fail loudly if tool is out of date

# ── Build ─────────────────────────────────────────────────────────
strict_warnings = true             # Treat OpenSCAD warnings as errors
libraries = ["openscad-threads"]   # Needs thread library (run: 3dm install-libraries)
local_libraries = ["./lib"]        # Also include project-local ./lib/ directory

# ── Slicing ───────────────────────────────────────────────────────
printer_profile = "prusa_mk4"      # Use profiles/prusa_mk4.ini
scale = 1.0                        # No scaling
copies = 1
overlays = ["fast-draft"]          # Always slice with fast-draft settings

# ── Images & Preview ──────────────────────────────────────────────
view = "3sil"
image_angles = ["above_front_left", "top"]
colorscheme = "slicer_light"
image_size = "1920x1080"

# ── Printing ──────────────────────────────────────────────────────
# (Print connection settings like octoprint_host belong in defaults.toml,
#  not here, since they are machine-specific not project-specific)
print_mode = "octoprint"
auto_start_prints = true

\end{lstlisting}

A fully commented \texttt{defaults.toml} for a machine-wide config:

\begin{lstlisting}[style=Alabaster]
# ── Machine-wide defaults (in ~/.config/3dmake/defaults.toml) ────

# General defaults
model_name = "main"
view = "3sil"
auto_start_prints = true

# Editor
editor = "code"
edit_in_background = true          # For GUI editors that return immediately

# Printer
printer_profile = "prusa_mk4"

# OctoPrint connection
print_mode = "octoprint"
octoprint_host = "http://octopi.local"
octoprint_key = "ABC123DEF456GHI789..."

# AI descriptions (keep sensitive keys here, not in project toml)
gemini_key = "AIzaSyABC123..."
llm_name = "gemini-1.5-flash"

\end{lstlisting}

\begin{center}\rule{0.5\linewidth}{0.5pt}\end{center}

\subsection*{Typical Workflows}\label{docs__pandoc__latex__src__3dmake_appendices_complete.md__typical-workflows}

\subsubsection*{New Project from Scratch}\label{docs__pandoc__latex__src__3dmake_appendices_complete.md__new-project-from-scratch}

\begin{lstlisting}[style=Alabaster, language=bash]
3dm new                    # Create project structure
cd my-project
3dm edit-model             # Write your OpenSCAD model
3dm build                  # Verify it compiles
3dm info                   # Check dimensions
3dm build orient slice print  # Full pipeline

\end{lstlisting}

\subsubsection*{Iterate on an Existing Model}\label{docs__pandoc__latex__src__3dmake_appendices_complete.md__iterate-on-an-existing-model}

\begin{lstlisting}[style=Alabaster, language=bash]
3dm edit-model             # Edit the .scad source
3dm build                  # Recompile
3dm preview                # Quick visual check
3dm slice -o supports      # Slice with supports overlay
3dm print                  # Send to printer

\end{lstlisting}

\subsubsection*{Working with a Standalone STL}\label{docs__pandoc__latex__src__3dmake_appendices_complete.md__working-with-a-standalone-stl}

\begin{lstlisting}[style=Alabaster, language=bash]
3dm info my-model.stl
3dm preview my-model.stl
3dm image my-model.stl --angle above_front_left --angle top
3dm slice print my-model.stl

\end{lstlisting}

\subsubsection*{Export Images for Documentation}\label{docs__pandoc__latex__src__3dmake_appendices_complete.md__export-images-for-documentation}

\begin{lstlisting}[style=Alabaster, language=bash]
3dm image \
  --angle above_front_left \
  --angle above_front_right \
  --angle top \
  --colorscheme slicer_light \
  --image-size 1920x1080

\end{lstlisting}

\begin{center}\rule{0.5\linewidth}{0.5pt}\end{center}

\subsection*{Environment Variables}\label{docs__pandoc__latex__src__3dmake_appendices_complete.md__environment-variables}

{\def\LTcaptype{none} % do not increment counter
\begin{longtable}[]{@{}
  >{\raggedright\arraybackslash}p{(\linewidth - 2\tabcolsep) * \real{0.2794}}
  >{\raggedright\arraybackslash}p{(\linewidth - 2\tabcolsep) * \real{0.7206}}@{}}
\toprule\noalign{}
\begin{minipage}[b]{\linewidth}\raggedright
Variable
\end{minipage} & \begin{minipage}[b]{\linewidth}\raggedright
Description
\end{minipage} \\
\midrule\noalign{}
\endhead
\bottomrule\noalign{}
\endlastfoot
\texttt{3DMAKE\_CONFIG\_DIR} &
Override the default configuration directory path \\
\end{longtable}
}

\begin{center}\rule{0.5\linewidth}{0.5pt}\end{center}

\subsection*{Configuration Directory Structure}\label{docs__pandoc__latex__src__3dmake_appendices_complete.md__configuration-directory-structure}

\begin{lstlisting}[style=Alabaster]
~/.config/3dmake/         (or $3DMAKE_CONFIG_DIR)
├── defaults.toml          Global settings (printer, API keys, etc.)
├── profiles/
│   └── my_printer.ini     Printer profiles (slicer settings)
├── overlays/
│   ├── default/
│   │   └── supports.ini   Default overlays (apply to all printers)
│   └── my_printer/
│       └── supports.ini   Printer-specific overlay overrides
├── libraries/             Installed OpenSCAD libraries
└── templates/             Starter templates for new projects/overlays

\end{lstlisting}

\begin{center}\rule{0.5\linewidth}{0.5pt}\end{center}

\subsection*{Overlays}\label{docs__pandoc__latex__src__3dmake_appendices_complete.md__overlays}

Overlays are \textbf{partial PrusaSlicer \texttt{.ini} files} that override specific slicer settings without replacing the entire printer profile. They are the primary way to switch between print quality modes, enable supports, adjust strength, or any other per-print tuning --- without touching your base printer profile.

\subsubsection*{How Overlays Work}\label{docs__pandoc__latex__src__3dmake_appendices_complete.md__how-overlays-work}

\begin{enumerate}
\tightlist
\item
  When slicing, 3DMake loads your base printer profile first, then applies each overlay \textbf{in order}, with later overlays winning over earlier ones.
\item
  Overlay files only need to contain the keys you want to change --- all other profile settings are preserved.
\item
  3DMake looks for overlays in two places, preferring the printer-specific version if it exists:

  \begin{itemize}
  \tightlist
  \item
    \texttt{overlays/\textless{}printer\_profile\textgreater{}/\textless{}name\textgreater{}.ini} ← preferred
  \item
    \texttt{overlays/default/\textless{}name\textgreater{}.ini} ← fallback
  \end{itemize}
\end{enumerate}

\subsubsection*{Creating and Editing Overlays}\label{docs__pandoc__latex__src__3dmake_appendices_complete.md__creating-and-editing-overlays}

\begin{lstlisting}[style=Alabaster, language=bash]
3dm edit-overlay                  # Prompted to pick or create an overlay
3dm edit-overlay -o strong        # Open (or create) the "strong" overlay directly
3dm list-overlays                 # See all available overlays

\end{lstlisting}

When creating a new overlay via \texttt{3dm\ edit-overlay}, you can optionally lock it to a specific printer profile. A locked overlay lives under \texttt{overlays/\textless{}profile\textgreater{}/} and only applies when that profile is active.

\subsubsection*{Using Overlays}\label{docs__pandoc__latex__src__3dmake_appendices_complete.md__using-overlays}

\begin{lstlisting}[style=Alabaster, language=bash]
# Apply a single overlay
3dm slice -o strong

# Apply multiple overlays (stacked in order, last wins on conflicts)
3dm slice -o fine -o supports

# Set a default overlay for a project in 3dmake.toml
overlays = ["strong"]

# CLI overlay replaces (not adds to) the toml default
3dm slice -o draft              # Uses only "draft", ignores toml default

\end{lstlisting}

\subsubsection*{Overlay File Format}\label{docs__pandoc__latex__src__3dmake_appendices_complete.md__overlay-file-format}

Overlay \texttt{.ini} files use the same key=value format as PrusaSlicer profiles. You only need to include the keys you want to override.

\begin{lstlisting}[style=Alabaster]
# overlays/default/strong.ini
# Keys and values are identical to PrusaSlicer's config export format.
# Lines starting with # are comments and are ignored.

perimeters = 4
top_solid_layers = 6
bottom_solid_layers = 6
fill_density = 40%
fill_pattern = grid

\end{lstlisting}

\begin{center}\rule{0.5\linewidth}{0.5pt}\end{center}

\subsubsection*{Key PrusaSlicer INI Settings Reference}\label{docs__pandoc__latex__src__3dmake_appendices_complete.md__key-prusaslicer-ini-settings-reference}

The most useful keys for overlay tuning, with types and sample values:

\paragraph*{Shell / Wall Strength}\label{docs__pandoc__latex__src__3dmake_appendices_complete.md__shell--wall-strength}

{\def\LTcaptype{none} % do not increment counter
\begin{longtable}[]{@{}
  >{\raggedright\arraybackslash}p{(\linewidth - 4\tabcolsep) * \real{0.2074}}
  >{\raggedright\arraybackslash}p{(\linewidth - 4\tabcolsep) * \real{0.4889}}
  >{\raggedright\arraybackslash}p{(\linewidth - 4\tabcolsep) * \real{0.3037}}@{}}
\toprule\noalign{}
\begin{minipage}[b]{\linewidth}\raggedright
Key
\end{minipage} & \begin{minipage}[b]{\linewidth}\raggedright
Description
\end{minipage} & \begin{minipage}[b]{\linewidth}\raggedright
Sample Values
\end{minipage} \\
\midrule\noalign{}
\endhead
\bottomrule\noalign{}
\endlastfoot
\texttt{perimeters} & Number of outer wall loops &
\texttt{2} (draft) · \texttt{3} (normal) · \texttt{5} (strong) \\
\texttt{top\_solid\_layers} & Solid layers on the top surface &
\texttt{3} (draft) · \texttt{5} (normal) · \texttt{8} (strong) \\
\texttt{bottom\_solid\_layers} & Solid layers on the bottom surface &
\texttt{3} (draft) · \texttt{4} (normal) · \texttt{6} (strong) \\
\texttt{top\_solid\_min\_thickness} &
Minimum top skin thickness in mm (overrides layer count if larger) &
\texttt{0.5} · \texttt{0.8} · \texttt{1.0} \\
\texttt{bottom\_solid\_min\_thickness} &
Minimum bottom skin thickness in mm &
\texttt{0.5} · \texttt{0.8} · \texttt{1.0} \\
\texttt{external\_perimeters\_first} &
Print outer walls before inner walls (better surface quality) &
\texttt{0} · \texttt{1} \\
\texttt{thin\_walls} &
Detect and print thin walls that would otherwise be missed &
\texttt{0} · \texttt{1} \\
\end{longtable}
}

\paragraph*{Infill}\label{docs__pandoc__latex__src__3dmake_appendices_complete.md__infill}

{\def\LTcaptype{none} % do not increment counter
\begin{longtable}[]{@{}
  >{\raggedright\arraybackslash}p{(\linewidth - 4\tabcolsep) * \real{0.1552}}
  >{\raggedright\arraybackslash}p{(\linewidth - 4\tabcolsep) * \real{0.3563}}
  >{\raggedright\arraybackslash}p{(\linewidth - 4\tabcolsep) * \real{0.4885}}@{}}
\toprule\noalign{}
\begin{minipage}[b]{\linewidth}\raggedright
Key
\end{minipage} & \begin{minipage}[b]{\linewidth}\raggedright
Description
\end{minipage} & \begin{minipage}[b]{\linewidth}\raggedright
Sample Values
\end{minipage} \\
\midrule\noalign{}
\endhead
\bottomrule\noalign{}
\endlastfoot
\texttt{fill\_density} & Infill percentage &
\texttt{5\%} · \texttt{15\%} · \texttt{20\%} · \texttt{40\%} · \texttt{60\%} \\
\texttt{fill\_pattern} & Infill pattern &
\texttt{grid} · \texttt{gyroid} · \texttt{honeycomb} · \texttt{cubic} · \texttt{rectilinear} · \texttt{lightning} · \texttt{triangles} \\
\texttt{top\_fill\_pattern} & Pattern for top solid layers &
\texttt{rectilinear} · \texttt{monotonic} · \texttt{concentric} · \texttt{smooth} \\
\texttt{bottom\_fill\_pattern} & Pattern for bottom solid layers &
\texttt{rectilinear} · \texttt{monotonic} · \texttt{concentric} \\
\texttt{infill\_every\_layers} &
Only infill every N layers (speeds up print, reduces strength) &
\texttt{1} · \texttt{2} · \texttt{3} \\
\texttt{solid\_infill\_every\_layers} &
Force solid infill every N layers & \texttt{0} (off) · \texttt{5} \\
\texttt{infill\_overlap} & How much infill overlaps with perimeters (\%)
& \texttt{15\%} · \texttt{25\%} \\
\end{longtable}
}

\paragraph*{Layer Height \& Speed}\label{docs__pandoc__latex__src__3dmake_appendices_complete.md__layer-height--speed}

{\def\LTcaptype{none} % do not increment counter
\begin{longtable}[]{@{}
  >{\raggedright\arraybackslash}p{(\linewidth - 4\tabcolsep) * \real{0.2342}}
  >{\raggedright\arraybackslash}p{(\linewidth - 4\tabcolsep) * \real{0.4955}}
  >{\raggedright\arraybackslash}p{(\linewidth - 4\tabcolsep) * \real{0.2703}}@{}}
\toprule\noalign{}
\begin{minipage}[b]{\linewidth}\raggedright
Key
\end{minipage} & \begin{minipage}[b]{\linewidth}\raggedright
Description
\end{minipage} & \begin{minipage}[b]{\linewidth}\raggedright
Sample Values
\end{minipage} \\
\midrule\noalign{}
\endhead
\bottomrule\noalign{}
\endlastfoot
\texttt{layer\_height} & Standard layer height in mm &
\texttt{0.1} · \texttt{0.15} · \texttt{0.2} · \texttt{0.3} \\
\texttt{first\_layer\_height} &
Height of the first layer (absolute mm or \%) &
\texttt{0.2} · \texttt{0.25} · \texttt{75\%} \\
\texttt{perimeter\_speed} & Speed for perimeter walls (mm/s) &
\texttt{40} · \texttt{60} · \texttt{80} \\
\texttt{external\_perimeter\_speed} &
Speed for outermost wall only (slower = better quality) &
\texttt{25} · \texttt{35} · \texttt{50\%} \\
\texttt{infill\_speed} & Speed for infill (mm/s) &
\texttt{60} · \texttt{80} · \texttt{120} \\
\texttt{solid\_infill\_speed} & Speed for solid infill layers (mm/s) &
\texttt{50} · \texttt{70} \\
\texttt{top\_solid\_infill\_speed} &
Speed for top surface layers (slower = smoother) &
\texttt{30} · \texttt{40} · \texttt{50} \\
\texttt{travel\_speed} & Non-extrusion travel speed (mm/s) &
\texttt{130} · \texttt{180} · \texttt{250} \\
\texttt{first\_layer\_speed} & First layer speed (mm/s or \%) &
\texttt{20} · \texttt{30} · \texttt{50\%} \\
\end{longtable}
}

\paragraph*{Supports}\label{docs__pandoc__latex__src__3dmake_appendices_complete.md__supports}

{\def\LTcaptype{none} % do not increment counter
\begin{longtable}[]{@{}
  >{\raggedright\arraybackslash}p{(\linewidth - 4\tabcolsep) * \real{0.2672}}
  >{\raggedright\arraybackslash}p{(\linewidth - 4\tabcolsep) * \real{0.4198}}
  >{\raggedright\arraybackslash}p{(\linewidth - 4\tabcolsep) * \real{0.3130}}@{}}
\toprule\noalign{}
\begin{minipage}[b]{\linewidth}\raggedright
Key
\end{minipage} & \begin{minipage}[b]{\linewidth}\raggedright
Description
\end{minipage} & \begin{minipage}[b]{\linewidth}\raggedright
Sample Values
\end{minipage} \\
\midrule\noalign{}
\endhead
\bottomrule\noalign{}
\endlastfoot
\texttt{support\_material} & Enable support structures &
\texttt{0} · \texttt{1} \\
\texttt{support\_material\_auto} & Auto-detect where supports are needed
& \texttt{0} · \texttt{1} \\
\texttt{support\_material\_threshold} &
Overhang angle above which supports are added (degrees) &
\texttt{40} · \texttt{45} · \texttt{55} \\
\texttt{support\_material\_pattern} & Support infill pattern &
\texttt{rectilinear} · \texttt{honeycomb} · \texttt{lightning} \\
\texttt{support\_material\_spacing} &
Distance between support lines (mm) &
\texttt{2} · \texttt{2.5} · \texttt{3} \\
\texttt{support\_material\_contact\_distance} &
Z gap between support and part (mm) &
\texttt{0.1} · \texttt{0.2} · \texttt{0.25} \\
\texttt{support\_material\_interface\_layers} &
Dense interface layers between support and part &
\texttt{0} · \texttt{2} · \texttt{3} \\
\end{longtable}
}

\paragraph*{Cooling \& Bridging}\label{docs__pandoc__latex__src__3dmake_appendices_complete.md__cooling--bridging}

{\def\LTcaptype{none} % do not increment counter
\begin{longtable}[]{@{}
  >{\raggedright\arraybackslash}p{(\linewidth - 4\tabcolsep) * \real{0.2568}}
  >{\raggedright\arraybackslash}p{(\linewidth - 4\tabcolsep) * \real{0.4595}}
  >{\raggedright\arraybackslash}p{(\linewidth - 4\tabcolsep) * \real{0.2838}}@{}}
\toprule\noalign{}
\begin{minipage}[b]{\linewidth}\raggedright
Key
\end{minipage} & \begin{minipage}[b]{\linewidth}\raggedright
Description
\end{minipage} & \begin{minipage}[b]{\linewidth}\raggedright
Sample Values
\end{minipage} \\
\midrule\noalign{}
\endhead
\bottomrule\noalign{}
\endlastfoot
\texttt{cooling} & Enable fan cooling & \texttt{0} · \texttt{1} \\
\texttt{min\_fan\_speed} & Minimum fan speed \% &
\texttt{20} · \texttt{35} · \texttt{50} \\
\texttt{max\_fan\_speed} & Maximum fan speed \% &
\texttt{50} · \texttt{80} · \texttt{100} \\
\texttt{bridge\_fan\_speed} & Fan speed during bridges \% &
\texttt{80} · \texttt{100} \\
\texttt{bridge\_speed} & Extrusion speed for bridges (mm/s) &
\texttt{20} · \texttt{30} · \texttt{40} \\
\texttt{bridge\_flow\_ratio} & Flow rate multiplier for bridges &
\texttt{0.8} · \texttt{0.9} · \texttt{1.0} \\
\end{longtable}
}

\paragraph*{Skirt, Brim \& Raft}\label{docs__pandoc__latex__src__3dmake_appendices_complete.md__skirt-brim--raft}

{\def\LTcaptype{none} % do not increment counter
\begin{longtable}[]{@{}
  >{\raggedright\arraybackslash}p{(\linewidth - 4\tabcolsep) * \real{0.2133}}
  >{\raggedright\arraybackslash}p{(\linewidth - 4\tabcolsep) * \real{0.5067}}
  >{\raggedright\arraybackslash}p{(\linewidth - 4\tabcolsep) * \real{0.2800}}@{}}
\toprule\noalign{}
\begin{minipage}[b]{\linewidth}\raggedright
Key
\end{minipage} & \begin{minipage}[b]{\linewidth}\raggedright
Description
\end{minipage} & \begin{minipage}[b]{\linewidth}\raggedright
Sample Values
\end{minipage} \\
\midrule\noalign{}
\endhead
\bottomrule\noalign{}
\endlastfoot
\texttt{skirts} & Number of skirt loops &
\texttt{0} · \texttt{1} · \texttt{3} \\
\texttt{skirt\_distance} & Distance from part to skirt (mm) &
\texttt{3} · \texttt{5} · \texttt{6} \\
\texttt{brim\_width} & Brim width in mm (\texttt{0} = disabled) &
\texttt{0} · \texttt{3} · \texttt{5} · \texttt{8} \\
\texttt{raft\_layers} & Number of raft layers (\texttt{0} = disabled) &
\texttt{0} · \texttt{2} · \texttt{3} \\
\end{longtable}
}

\paragraph*{Temperature}\label{docs__pandoc__latex__src__3dmake_appendices_complete.md__temperature}

{\def\LTcaptype{none} % do not increment counter
\begin{longtable}[]{@{}
  >{\raggedright\arraybackslash}p{(\linewidth - 4\tabcolsep) * \real{0.3152}}
  >{\raggedright\arraybackslash}p{(\linewidth - 4\tabcolsep) * \real{0.4239}}
  >{\raggedright\arraybackslash}p{(\linewidth - 4\tabcolsep) * \real{0.2609}}@{}}
\toprule\noalign{}
\begin{minipage}[b]{\linewidth}\raggedright
Key
\end{minipage} & \begin{minipage}[b]{\linewidth}\raggedright
Description
\end{minipage} & \begin{minipage}[b]{\linewidth}\raggedright
Sample Values
\end{minipage} \\
\midrule\noalign{}
\endhead
\bottomrule\noalign{}
\endlastfoot
\texttt{temperature} & Hotend temperature for all layers (°C) &
\texttt{200} · \texttt{215} · \texttt{230} \\
\texttt{first\_layer\_temperature} &
Hotend temperature for first layer only &
\texttt{215} · \texttt{225} · \texttt{235} \\
\texttt{bed\_temperature} & Bed temperature for all layers (°C) &
\texttt{0} · \texttt{55} · \texttt{60} · \texttt{80} \\
\texttt{first\_layer\_bed\_temperature} &
Bed temperature for first layer only &
\texttt{60} · \texttt{65} · \texttt{85} \\
\end{longtable}
}

\begin{center}\rule{0.5\linewidth}{0.5pt}\end{center}

\subsubsection*{\texorpdfstring{Sample Overlay: \texttt{strong.ini}}{Sample Overlay: strong.ini}}\label{docs__pandoc__latex__src__3dmake_appendices_complete.md__sample-overlay-strongini}

For functional parts, mechanical components, or anything that needs to survive stress. Heavier walls, more solid layers, higher infill, slower outer perimeters.

\begin{lstlisting}[style=Alabaster]
# overlays/default/strong.ini
# Heavy-duty print — maximizes strength at the cost of time and filament.

# Shell: thick walls and solid layers
perimeters = 5
top_solid_layers = 7
bottom_solid_layers = 6
top_solid_min_thickness = 1.2
bottom_solid_min_thickness = 1.0
external_perimeters_first = 1

# Infill: dense cubic for near-isotropic strength
fill_density = 50%
fill_pattern = cubic
infill_overlap = 25%

# Speed: slower outer wall for better layer adhesion
external_perimeter_speed = 25
perimeter_speed = 45
top_solid_infill_speed = 30

# First layer: slower for bed adhesion
first_layer_speed = 20

# Cooling: slightly reduced to improve layer bonding on perimeters
max_fan_speed = 70

\end{lstlisting}

\begin{center}\rule{0.5\linewidth}{0.5pt}\end{center}

\subsubsection*{\texorpdfstring{Sample Overlay: \texttt{draft.ini}}{Sample Overlay: draft.ini}}\label{docs__pandoc__latex__src__3dmake_appendices_complete.md__sample-overlay-draftini}

For quick test fits, rough prototypes, or checking dimensions before committing to a full print. Faster speeds, minimal walls, low infill, thicker layers.

\begin{lstlisting}[style=Alabaster]
# overlays/default/draft.ini
# Fast draft print — sacrifices surface quality and strength for speed.

# Shell: minimum viable walls
perimeters = 2
top_solid_layers = 3
bottom_solid_layers = 3

# Infill: low density, fast pattern
fill_density = 10%
fill_pattern = grid

# Layer height: thick layers = fewer passes = faster print
# NOTE: layer_height must be <= 75% of nozzle diameter.
# For a 0.4mm nozzle, max is 0.3mm.
layer_height = 0.3

# Speed: push everything faster
perimeter_speed = 80
external_perimeter_speed = 60
infill_speed = 120
solid_infill_speed = 90
top_solid_infill_speed = 60
travel_speed = 200
first_layer_speed = 30

# Cooling: full fan to keep up with fast extrusion
min_fan_speed = 50
max_fan_speed = 100

\end{lstlisting}

\begin{center}\rule{0.5\linewidth}{0.5pt}\end{center}

\subsubsection*{\texorpdfstring{Sample Overlay: \texttt{fine.ini}}{Sample Overlay: fine.ini}}\label{docs__pandoc__latex__src__3dmake_appendices_complete.md__sample-overlay-fineini}

For display models, miniatures, parts with fine surface detail, or anything where appearance matters more than speed. Thin layers, many perimeters, slow outer walls, smooth top surfaces.

\begin{lstlisting}[style=Alabaster]
# overlays/default/fine.ini
# Fine detail print — maximizes surface quality and dimensional accuracy.

# Shell: many perimeters for crisp edges and thick solid skins
perimeters = 4
top_solid_layers = 8
bottom_solid_layers = 5
top_solid_min_thickness = 1.0
external_perimeters_first = 1
thin_walls = 1

# Infill: moderate density with monotonic top for smooth finish
fill_density = 20%
fill_pattern = gyroid
top_fill_pattern = monotonic
bottom_fill_pattern = monotonic

# Layer height: thin layers for fine Z resolution
# NOTE: must be >= 25% of nozzle diameter to ensure reliable extrusion.
# For a 0.4mm nozzle, minimum is ~0.08mm; 0.1mm is safe.
layer_height = 0.1
first_layer_height = 0.2

# Speed: slow everything down for precision
perimeter_speed = 40
external_perimeter_speed = 20
infill_speed = 60
solid_infill_speed = 40
top_solid_infill_speed = 25
bridge_speed = 20
first_layer_speed = 15

# Cooling: strong cooling for clean overhangs and bridges
min_fan_speed = 50
max_fan_speed = 100
bridge_fan_speed = 100
bridge_flow_ratio = 0.85

# Brim: small brim helps adhesion for tall fine parts
brim_width = 3

\end{lstlisting}

\begin{center}\rule{0.5\linewidth}{0.5pt}\end{center}

\subsubsection*{\texorpdfstring{Sample Overlay: \texttt{supports.ini}}{Sample Overlay: supports.ini}}\label{docs__pandoc__latex__src__3dmake_appendices_complete.md__sample-overlay-supportsini}

Enable automatic support structures for models with steep overhangs. Stack with other quality overlays.

\begin{lstlisting}[style=Alabaster]
# overlays/default/supports.ini
# Enable automatic supports for overhangs > 45 degrees.

support_material = 1
support_material_auto = 1
support_material_threshold = 45
support_material_pattern = lightning
support_material_spacing = 2.5
support_material_contact_distance = 0.2
support_material_interface_layers = 2

\end{lstlisting}

\subsubsection*{\texorpdfstring{Sample Overlay: \texttt{no-supports.ini}}{Sample Overlay: no-supports.ini}}\label{docs__pandoc__latex__src__3dmake_appendices_complete.md__sample-overlay-no-supportsini}

Explicitly disable supports --- useful when your base profile has them on by default.

\begin{lstlisting}[style=Alabaster]
# overlays/default/no-supports.ini
support_material = 0
support_material_auto = 0

\end{lstlisting}

\begin{center}\rule{0.5\linewidth}{0.5pt}\end{center}

\subsubsection*{Stacking Overlays}\label{docs__pandoc__latex__src__3dmake_appendices_complete.md__stacking-overlays}

Overlays are applied \textbf{left to right}, so later overlays win on any conflicting keys. This lets you compose print settings modularly:

\begin{lstlisting}[style=Alabaster, language=bash]
# Fine quality WITH supports
3dm slice -o fine -o supports

# Strong walls but fast draft speed (fine conflicts: perimeter_speed in draft wins)
3dm slice -o strong -o draft

# Strong structural settings + supports, no brim (add a tiny no-brim override)
3dm slice -o strong -o supports -o no-brim

\end{lstlisting}

You can also set permanent default overlays for a project in \texttt{3dmake.toml}, while adding extra ones at the CLI:

\begin{lstlisting}[style=Alabaster]
# 3dmake.toml — always use strong as the base for this project
overlays = ["strong"]

\end{lstlisting}

\begin{lstlisting}[style=Alabaster, language=bash]
# At CLI: strong (from toml) is REPLACED by the CLI overlay list
3dm slice -o strong -o supports   # Both applied
3dm slice -o fine                 # Only fine — replaces toml default entirely

\end{lstlisting}

\begin{quote}
\textbf{Tip:} Name your overlays by intent (\texttt{strong}, \texttt{draft}, \texttt{fine}, \texttt{supports}, \texttt{vase-mode}, \texttt{no-brim}) rather than by settings (\texttt{40pct-infill}). This keeps the CLI readable and lets you tune the overlay file without renaming references everywhere.
\end{quote}

\begin{center}\rule{0.5\linewidth}{0.5pt}\end{center}

\subsection*{Action Dependency \& Chain Rules}\label{docs__pandoc__latex__src__3dmake_appendices_complete.md__action-dependency--chain-rules}

\subsubsection*{Implied Action Dependencies}\label{docs__pandoc__latex__src__3dmake_appendices_complete.md__implied-action-dependencies}

Some actions automatically trigger prerequisite actions. You never need to spell these out manually --- 3DMake inserts them for you.

\begin{lstlisting}[style=Alabaster]
print
  └── slice (implied)
        └── build (implied, if scad source exists)

image
  └── load-mesh (internal, implied)

info
  └── measure-mesh (internal, implied)
        └── load-mesh (internal, implied)

preview
  └── measure-mesh (internal, implied)
        └── load-mesh (internal, implied)

orient
  (no implied actions — operates on an already-built STL)

\end{lstlisting}

So running \texttt{3dm\ print} on a project with OpenSCAD source will automatically run \texttt{build} -\textgreater{} \texttt{slice} -\textgreater{} \texttt{print} in order. You only need to be explicit when you want to \textbf{stop partway through} (e.g. \texttt{3dm\ build} to just check compilation without slicing).

\subsubsection*{\texorpdfstring{\texttt{last\_in\_chain} --- Actions That Must Come Last}{last\_in\_chain --- Actions That Must Come Last}}\label{docs__pandoc__latex__src__3dmake_appendices_complete.md__last_in_chain--actions-that-must-come-last}

Three actions are marked \texttt{last\_in\_chain}, meaning \textbf{no other pipeline actions can follow them} in the same command. Attempting to do so is an error.

{\def\LTcaptype{none} % do not increment counter
\begin{longtable}[]{@{}
  >{\raggedright\arraybackslash}p{(\linewidth - 2\tabcolsep) * \real{0.1286}}
  >{\raggedright\arraybackslash}p{(\linewidth - 2\tabcolsep) * \real{0.8714}}@{}}
\toprule\noalign{}
\begin{minipage}[b]{\linewidth}\raggedright
Action
\end{minipage} & \begin{minipage}[b]{\linewidth}\raggedright
Why it\textquotesingle s last
\end{minipage} \\
\midrule\noalign{}
\endhead
\bottomrule\noalign{}
\endlastfoot
\texttt{image} &
Produces image files; nothing meaningful can follow a render \\
\texttt{info} &
Terminal inspection command; no build artifact is produced \\
\texttt{preview} &
Produces a projection STL; slicing a silhouette is not useful \\
\end{longtable}
}

\begin{lstlisting}[style=Alabaster, language=bash]
# ✅ Valid — last_in_chain actions used alone or at the end
3dm build image
3dm build orient image
3dm info
3dm build preview

# ❌ Invalid — actions after a last_in_chain action
3dm image slice        # ERROR: 'image' cannot be followed by 'slice'
3dm info build         # ERROR: 'info' cannot be followed by 'build'
3dm preview print      # ERROR: 'preview' cannot be followed by 'print'

\end{lstlisting}

\subsubsection*{Isolated Actions}\label{docs__pandoc__latex__src__3dmake_appendices_complete.md__isolated-actions}

A separate category called \textbf{isolated} means the action cannot be combined with \emph{any} other action at all --- it must be the only verb on the command line. All \texttt{edit-*} actions, \texttt{setup}, \texttt{new}, \texttt{version}, \texttt{help}, \texttt{test-connect}, \texttt{list-profiles}, \texttt{list-overlays}, \texttt{list-libraries}, and \texttt{install-libraries} are isolated.

\begin{lstlisting}[style=Alabaster, language=bash]
# ✅ Valid
3dm edit-model
3dm setup
3dm list-profiles

# ❌ Invalid
3dm edit-model build    # ERROR: edit-model is isolated
3dm setup new           # ERROR: setup is isolated

\end{lstlisting}

\subsubsection*{Full Action Execution Order}\label{docs__pandoc__latex__src__3dmake_appendices_complete.md__full-action-execution-order}

When you request multiple actions, they always execute in this fixed internal order regardless of how you type them on the command line:

\begin{lstlisting}[style=Alabaster]
setup -> new -> build -> load-mesh* -> measure-mesh* -> info -> orient ->
image -> preview -> slice -> print -> list-profiles -> list-overlays ->
edit-model -> edit-overlay -> edit-profile -> edit-profile-gcode ->
edit-prompt -> edit-global-config -> list-libraries -> install-libraries ->
test-connect -> help -> version

\end{lstlisting}

\_(Actions marked \texttt{\_} are internal and cannot be requested directly.)*

This means \texttt{3dm\ slice\ build} and \texttt{3dm\ build\ slice} are identical --- order of verbs on the command line does not matter.

\begin{center}\rule{0.5\linewidth}{0.5pt}\end{center}

\subsection*{Multi-Model Projects}\label{docs__pandoc__latex__src__3dmake_appendices_complete.md__multi-model-projects}

A single 3DMake project can contain multiple OpenSCAD models. Each model is a separate \texttt{.scad} file in the \texttt{src/} directory.

\subsubsection*{Directory Layout}\label{docs__pandoc__latex__src__3dmake_appendices_complete.md__directory-layout}

\begin{lstlisting}[style=Alabaster]
my-project/
├── 3dmake.toml
├── src/
│   ├── main.scad        ← default model (model_name = "main")
│   ├── lid.scad
│   ├── bracket.scad
│   └── insert.scad
└── build/               ← generated STLs and GCode land here
    ├── main.stl
    ├── lid.stl
    └── my-project-main.gcode

\end{lstlisting}

\subsubsection*{Selecting a Model}\label{docs__pandoc__latex__src__3dmake_appendices_complete.md__selecting-a-model}

\begin{lstlisting}[style=Alabaster, language=bash]
# Uses model_name from 3dmake.toml (default: "main")
3dm build slice

# Override at the CLI with --model / -m
3dm build slice -m lid
3dm build slice -m bracket
3dm info -m insert

# edit-model also respects -m
3dm edit-model -m lid

\end{lstlisting}

\subsubsection*{\texorpdfstring{Setting a Per-Model Default in \texttt{3dmake.toml}}{Setting a Per-Model Default in 3dmake.toml}}\label{docs__pandoc__latex__src__3dmake_appendices_complete.md__setting-a-per-model-default-in-3dmaketoml}

\begin{lstlisting}[style=Alabaster]
# 3dmake.toml
model_name = "bracket"    # Everyone on this project defaults to the bracket model

\end{lstlisting}

\subsubsection*{Output File Naming}\label{docs__pandoc__latex__src__3dmake_appendices_complete.md__output-file-naming}

Output G-code files are named using the pattern:

\begin{lstlisting}[style=Alabaster]
<project_name>-<model_name>.gcode          # 1 copy
<project_name>-<model_name>-x3.gcode       # 3 copies (--copies 3)

\end{lstlisting}

For example, a project named \texttt{pi-case} building the \texttt{lid} model with 2 copies produces:

\begin{lstlisting}[style=Alabaster]
build/pi-case-lid-x2.gcode

\end{lstlisting}

If \texttt{project\_name} is not set in \texttt{3dmake.toml}, it is auto-detected from the directory name.

\subsubsection*{Tip: Shared Parameters Across Models}\label{docs__pandoc__latex__src__3dmake_appendices_complete.md__tip-shared-parameters-across-models}

A common OpenSCAD pattern is a shared \texttt{params.scad} file included by all models:

\begin{lstlisting}[style=Alabaster]
src/
├── params.scad      ← shared dimensions, tolerances, material choices
├── main.scad        ← include <params.scad>
├── lid.scad         ← include <params.scad>
└── bracket.scad     ← include <params.scad>

\end{lstlisting}

This way changing one dimension in \texttt{params.scad} propagates to every model, and you can rebuild all variants with a shell loop:

\begin{lstlisting}[style=Alabaster, language=bash]
for model in main lid bracket; do
    3dm build slice -m $model
done

\end{lstlisting}

\begin{center}\rule{0.5\linewidth}{0.5pt}\end{center}

\subsection*{Special INI Keys}\label{docs__pandoc__latex__src__3dmake_appendices_complete.md__special-ini-keys}

\subsubsection*{\texorpdfstring{\texttt{3dm\_bed\_center}}{3dm\_bed\_center}}\label{docs__pandoc__latex__src__3dmake_appendices_complete.md__3dm_bed_center}

This is a \textbf{3DMake-specific} key (not a standard PrusaSlicer key) that you can add to a printer profile or overlay. It works around a PrusaSlicer CLI bug where models are not correctly centered on non-rectangular or offset build volumes.

\begin{lstlisting}[style=Alabaster]
# In a printer profile .ini or overlay .ini
# Value is "X,Y" coordinates of the bed center in mm
3dm_bed_center = 117.5,117.5    # Prusa MK3/MK4 (235x235mm bed)
3dm_bed_center = 110,110        # 220x220mm bed
3dm_bed_center = 150,150        # 300x300mm bed
3dm_bed_center = 175,162.5      # Prusa XL (350x325mm bed)

\end{lstlisting}

When this key is present, 3DMake passes \texttt{-\/-center\ X,Y} to the PrusaSlicer CLI automatically. You should set this in your \textbf{printer profile} rather than an overlay, since it is specific to the physical printer\textquotesingle s bed geometry.

\begin{center}\rule{0.5\linewidth}{0.5pt}\end{center}

\subsection*{AI Prompt Customization}\label{docs__pandoc__latex__src__3dmake_appendices_complete.md__ai-prompt-customization}

The \texttt{info} action sends rendered images of your model to an AI (Gemini or OpenRouter) along with a text prompt. The default prompt asks for a general physical description of the object. You can customize this prompt to get more targeted or useful output.

\subsubsection*{Editing the Prompt}\label{docs__pandoc__latex__src__3dmake_appendices_complete.md__editing-the-prompt}

\begin{lstlisting}[style=Alabaster, language=bash]
3dm edit-prompt    # Opens the prompt file in your configured editor

\end{lstlisting}

The prompt file lives at:

\begin{lstlisting}[style=Alabaster]
~/.config/3dmake/prompt.txt    (or $3DMAKE_CONFIG_DIR/prompt.txt)

\end{lstlisting}

If the file does not exist when you run \texttt{3dm\ edit-prompt}, 3DMake creates it from a built-in default automatically.

\subsubsection*{What the Prompt Receives}\label{docs__pandoc__latex__src__3dmake_appendices_complete.md__what-the-prompt-receives}

The AI receives your prompt text \textbf{plus six rendered PNG images} of the model taken from these fixed viewpoints: \texttt{above\_front\_left}, \texttt{above\_front\_right}, \texttt{above\_back\_left}, \texttt{above\_back\_right}, \texttt{top}, and \texttt{bottom}. Each image is 768×768 pixels.

\subsubsection*{Sample Prompts}\label{docs__pandoc__latex__src__3dmake_appendices_complete.md__sample-prompts}

\textbf{Default-style --- general physical description:}

\begin{lstlisting}[style=Alabaster]
You are analyzing a 3D model for 3D printing. Describe the object's physical
shape, identifying any notable geometric features, symmetry, overhangs, thin
walls, or small details that might affect printability. Be concise and factual.

\end{lstlisting}

\textbf{Printability review --- focused on print problems:}

\begin{lstlisting}[style=Alabaster]
You are a 3D printing expert reviewing a model before printing. Examine these
views and identify:
1. Any overhangs greater than 45 degrees that will need supports
2. Thin walls or small features that may not print reliably at 0.4mm nozzle diameter
3. The best orientation for printing (which face should be on the bed)
4. Any other potential print quality issues

Be specific about locations (e.g. "the underside of the lip on the left side").

\end{lstlisting}

\textbf{Dimensional estimation --- for sanity-checking scale:}

\begin{lstlisting}[style=Alabaster]
Based on the shape and proportions of this 3D model, estimate what real-world
object it most resembles and what its approximate dimensions might be in
millimeters. Note any features that suggest a specific scale (screw holes,
finger clearances, standard hardware interfaces, etc.).

\end{lstlisting}

\textbf{Design feedback --- for iterating on a model:}

\begin{lstlisting}[style=Alabaster]
You are reviewing a 3D model for functional and aesthetic design quality.
Comment on: whether the design looks structurally sound, any places where
the geometry looks unintentional or messy, whether any features look like
they might be too fragile, and one specific improvement suggestion.

\end{lstlisting}

\textbf{Interactive mode} (\texttt{3dm\ info\ -\/-interactive}) lets you follow up with specific questions after the initial AI response, such as asking it to focus on a particular feature or compare proportions.

\begin{center}\rule{0.5\linewidth}{0.5pt}\end{center}

\subsection*{Troubleshooting}\label{docs__pandoc__latex__src__3dmake_appendices_complete.md__troubleshooting}

\subsubsection*{Setup \& Configuration}\label{docs__pandoc__latex__src__3dmake_appendices_complete.md__setup--configuration}

\textbf{\texttt{3DMake\ settings\ and\ print\ options\ have\ not\ been\ set\ up\ on\ this\ machine.}}
Run \texttt{3dm\ setup} before any other command. This creates the config directory, copies default profiles, and walks you through initial configuration.

\textbf{\texttt{Printer\ profile\ \textquotesingle{}X\textquotesingle{}\ does\ not\ exist.}}
Run \texttt{3dm\ list-profiles} to see valid profile names. Profile names are case-sensitive and must match the \texttt{.ini} filename exactly (without the extension). Set the correct one with \texttt{3dm\ setup} or by editing \texttt{defaults.toml} with \texttt{3dm\ edit-global-config}.

\textbf{\texttt{This\ project\ requires\ 3DMake\ version\ X\ or\ newer.}}
The \texttt{min\_3dmake\_version} key in \texttt{3dmake.toml} is higher than your installed version. Run \texttt{3dm\ version} to check your version, then update 3DMake. If you are the project author, lower \texttt{min\_3dmake\_version} if the feature requirement does not actually apply.

\begin{center}\rule{0.5\linewidth}{0.5pt}\end{center}

\subsubsection*{Build Errors}\label{docs__pandoc__latex__src__3dmake_appendices_complete.md__build-errors}

\textbf{\texttt{Source\ file\ src/main.scad\ does\ not\ exist.}}
Either you are not in a 3DMake project directory, or \texttt{model\_name} in your \texttt{3dmake.toml} does not match an actual \texttt{.scad} file in \texttt{src/}. Run \texttt{3dm\ edit-model} to create the file, or correct \texttt{model\_name}.

\textbf{\texttt{Some\ needed\ libraries\ are\ not\ installed:\ \textless{}name\textgreater{}}}
Run \texttt{3dm\ install-libraries} to download and install the missing libraries listed in your \texttt{3dmake.toml} \texttt{libraries} key. Run \texttt{3dm\ list-libraries} to verify the library name is spelled correctly (names are lowercase).

\textbf{OpenSCAD prints warnings but does not fail.}
Add \texttt{strict\_warnings\ =\ true} to your \texttt{3dmake.toml} to treat warnings as errors. This is recommended for all projects to catch issues like undefined variables and deprecated syntax early.

\textbf{Build succeeds but STL looks wrong / has missing geometry.}
Run with \texttt{-\/-debug} to see the full OpenSCAD output including all warnings:

\begin{lstlisting}[style=Alabaster, language=bash]
3dm build --debug

\end{lstlisting}

\begin{center}\rule{0.5\linewidth}{0.5pt}\end{center}

\subsubsection*{Slicing Errors}\label{docs__pandoc__latex__src__3dmake_appendices_complete.md__slicing-errors}

\textbf{\texttt{Could\ not\ find\ overlay\ \textquotesingle{}X\textquotesingle{}\ for\ profile\ \textquotesingle{}Y\textquotesingle{}.}}
The overlay file does not exist at \texttt{overlays/Y/X.ini} or \texttt{overlays/default/X.ini}. Run \texttt{3dm\ list-overlays} to see what is available, or \texttt{3dm\ edit-overlay\ -o\ X} to create it.

\textbf{\texttt{Could\ not\ fit\ the\ object\ outline\ on\ the\ build\ surface.}}
The model is larger than the printer\textquotesingle s bed, or is being placed off-center. Check your model dimensions with \texttt{3dm\ info}, verify your \texttt{printer\_profile} is correct for your actual printer, and confirm \texttt{3dm\_bed\_center} is set correctly in your profile if you have a non-standard bed.

\textbf{Slicer produces unexpected print warnings.}
Run with \texttt{-\/-debug} to see the full PrusaSlicer stdout output:

\begin{lstlisting}[style=Alabaster, language=bash]
3dm slice --debug

\end{lstlisting}

\begin{center}\rule{0.5\linewidth}{0.5pt}\end{center}

\subsubsection*{Print / Connection Errors}\label{docs__pandoc__latex__src__3dmake_appendices_complete.md__print--connection-errors}

\textbf{\texttt{Unknown\ print\ mode\ \textquotesingle{}X\textquotesingle{}.}}
\texttt{print\_mode} in your config must be exactly one of: \texttt{octoprint}, \texttt{bambu\_lan}, or \texttt{bambu\_connect}. Check spelling in \texttt{defaults.toml}.

\textbf{OctoPrint: \texttt{ERROR:\ Authentication\ failed\ (401)}}
Your \texttt{octoprint\_key} is wrong or expired. In OctoPrint, go to Settings \textgreater{} Application Keys and generate a new key, then update \texttt{defaults.toml} via \texttt{3dm\ edit-global-config}.

\textbf{OctoPrint: \texttt{ERROR:\ Could\ not\ connect\ to\ server}}
Check that \texttt{octoprint\_host} is correct and includes the protocol (\texttt{http://} or \texttt{https://}). Verify OctoPrint is running and the machine running 3DMake is on the same network. Run \texttt{3dm\ test-connect} for a detailed diagnosis.

\textbf{Bambu: FTP or MQTT connection fails.}
Verify \texttt{bambu\_host} is the printer\textquotesingle s current local IP (it may change if your router assigns IPs dynamically --- consider setting a static IP for the printer). Check \texttt{bambu\_access\_code} on the printer screen under Settings \textgreater{} WLAN. Run \texttt{3dm\ test-connect} for step-by-step diagnosis.

\textbf{Bambu Connect mode:} \texttt{3dm\ print} with \texttt{print\_mode\ =\ "bambu\_connect"} does \textbf{not} connect directly to the printer. It packages your G-code into a \texttt{.3mf} file and opens the Bambu Connect desktop application via a URI handoff. Bambu Connect must be installed and running on your machine. \texttt{3dm\ test-connect} will report that it cannot test this mode --- this is expected.

\begin{center}\rule{0.5\linewidth}{0.5pt}\end{center}

\subsubsection*{General}\label{docs__pandoc__latex__src__3dmake_appendices_complete.md__general}

\textbf{\texttt{Unknown\ action\ \textquotesingle{}X\textquotesingle{}}}
Action names are lowercase and use hyphens, not underscores (e.g. \texttt{edit-model}, not \texttt{edit\_model}). Run \texttt{3dm\ help} for the full list.

\textbf{\texttt{The\ action\ \textquotesingle{}X\textquotesingle{}\ can\ only\ be\ used\ on\ its\ own}}
You combined an isolated action with other actions. Isolated actions (\texttt{edit-*}, \texttt{setup}, \texttt{new}, \texttt{version}, \texttt{help}, \texttt{test-connect}, \texttt{list-*}, \texttt{install-libraries}) must be run alone.

\textbf{\texttt{Action\ \textquotesingle{}X\textquotesingle{}\ cannot\ be\ followed\ by\ other\ actions}}
You placed a \texttt{last\_in\_chain} action (\texttt{image}, \texttt{info}, \texttt{preview}) before other pipeline actions. These must be the final (or only) action in your command.

\textbf{\texttt{Cannot\ select\ a\ model\ name\ when\ using\ an\ input\ file}}
\texttt{-\/-model} / \texttt{-m} and a direct input file path are mutually exclusive. Use one or the other.

\textbf{\texttt{Multiple\ inputs\ not\ supported\ yet}}
Only one input file can be passed per invocation. To process multiple files, use a shell loop.

\begin{center}\rule{0.5\linewidth}{0.5pt}\end{center}

\subsection*{OpenSCAD Tips \& Conventions for 3DMake}\label{docs__pandoc__latex__src__3dmake_appendices_complete.md__openscad-tips--conventions-for-3dmake}

\subsubsection*{Project Conventions}\label{docs__pandoc__latex__src__3dmake_appendices_complete.md__project-conventions}

3DMake expects OpenSCAD source files in \texttt{src/\textless{}model\_name\textgreater{}.scad} and writes STL output to \texttt{build/\textless{}model\_name\textgreater{}.stl}. As long as you follow this layout (created automatically by \texttt{3dm\ new}), everything works without extra configuration.

\subsubsection*{Parameterized Models}\label{docs__pandoc__latex__src__3dmake_appendices_complete.md__parameterized-models}

The most powerful pattern for 3DMake projects is a fully parameterized model where all key dimensions are defined as variables at the top of the file. This makes it trivial to create variants without duplicating code.

\begin{lstlisting}[style=Alabaster, language=openscad]
// src/box.scad
// All tunable parameters at the top

length    = 80;     // mm
width     = 50;     // mm
height    = 30;     // mm
wall      = 2.5;    // mm — wall thickness
tolerance = 0.2;    // mm — fit clearance for mating parts
fillet_r  = 1.5;    // mm — corner radius

// Model below uses only the variables above
module box_shell() {
    difference() {
        rounded_rect([length, width, height], r=fillet_r);
        translate([wall, wall, wall])
            cube([length-wall*2, width-wall*2, height]);
    }
}

box_shell();

\end{lstlisting}

\subsubsection*{Multi-Model Shared Parameters}\label{docs__pandoc__latex__src__3dmake_appendices_complete.md__multi-model-shared-parameters}

For projects with multiple related models (e.g. a box and its lid), put shared dimensions in a separate include file:

\begin{lstlisting}[style=Alabaster, language=openscad]
// src/params.scad  — shared across all models
outer_length  = 80;
outer_width   = 50;
wall          = 2.5;
tolerance     = 0.2;
lid_height    = 8;
body_height   = 35;

\end{lstlisting}

\begin{lstlisting}[style=Alabaster, language=openscad]
// src/body.scad
include <params.scad>

difference() {
    cube([outer_length, outer_width, body_height]);
    translate([wall, wall, wall])
        cube([outer_length-wall*2, outer_width-wall*2, body_height]);
}

\end{lstlisting}

\begin{lstlisting}[style=Alabaster, language=openscad]
// src/lid.scad
include <params.scad>

cube([outer_length - tolerance*2, outer_width - tolerance*2, lid_height]);

\end{lstlisting}

\subsubsection*{Useful OpenSCAD Modules for 3D Printing}\label{docs__pandoc__latex__src__3dmake_appendices_complete.md__useful-openscad-modules-for-3d-printing}

\begin{lstlisting}[style=Alabaster, language=openscad]
// Rounded rectangle (better bed adhesion, less stress concentration)
module rounded_rect(size, r=1) {
    hull() {
        for (x=[r, size[0]-r], y=[r, size[1]-r])
            translate([x, y, 0]) cylinder(r=r, h=size[2], $fn=32);
    }
}

// Countersunk screw hole (M3 example)
module m3_hole(depth=10, countersink=true) {
    cylinder(d=3.4, h=depth, $fn=24);
    if (countersink)
        translate([0, 0, depth-3.2])
            cylinder(d1=3.4, d2=6.5, h=3.2, $fn=24);
}

// Snap-fit clip tab
module snap_tab(length=10, width=4, height=3, overhang=0.8) {
    union() {
        cube([length, width, height]);
        translate([0, 0, height])
            linear_extrude(overhang, scale=[1, 1, 0])
                square([length, width]);
    }
}

\end{lstlisting}

\subsubsection*{\texorpdfstring{\texttt{\$fn}, \texttt{\$fa}, \texttt{\$fs} --- Circle Resolution}{\$fn, \$fa, \$fs --- Circle Resolution}}\label{docs__pandoc__latex__src__3dmake_appendices_complete.md__fn-fa-fs--circle-resolution}

OpenSCAD renders circles and cylinders as polygons. Control resolution with these special variables:

{\def\LTcaptype{none} % do not increment counter
\begin{longtable}[]{@{}
  >{\raggedright\arraybackslash}p{(\linewidth - 4\tabcolsep) * \real{0.0964}}
  >{\raggedright\arraybackslash}p{(\linewidth - 4\tabcolsep) * \real{0.4337}}
  >{\raggedright\arraybackslash}p{(\linewidth - 4\tabcolsep) * \real{0.4699}}@{}}
\toprule\noalign{}
\begin{minipage}[b]{\linewidth}\raggedright
Variable
\end{minipage} & \begin{minipage}[b]{\linewidth}\raggedright
Meaning
\end{minipage} & \begin{minipage}[b]{\linewidth}\raggedright
Recommended
\end{minipage} \\
\midrule\noalign{}
\endhead
\bottomrule\noalign{}
\endlastfoot
\texttt{\$fn} & Fixed number of fragments &
\texttt{32} for small circles, \texttt{64}+ for large \\
\texttt{\$fa} & Minimum angle per fragment (degrees) &
\texttt{1} for smooth curves \\
\texttt{\$fs} & Minimum fragment size (mm) &
\texttt{0.5} for print resolution \\
\end{longtable}
}

\begin{lstlisting}[style=Alabaster, language=openscad]
// Set globally at top of file for consistent quality
$fa = 1;
$fs = 0.4;    // Match your nozzle diameter

// Or set per-object for performance (preview fast, export high quality)
cylinder(r=10, h=5, $fn=64);

\end{lstlisting}

\begin{quote}
\textbf{Tip:} Use \texttt{\$fn=16} or \texttt{\$fn=24} during development for fast preview renders, then raise it to \texttt{\$fn=64} or higher before your final \texttt{3dm\ build} for export.
\end{quote}

\subsubsection*{Designing for Printability}\label{docs__pandoc__latex__src__3dmake_appendices_complete.md__designing-for-printability}

\textbf{Avoid overhangs steeper than 45°} without supports. Use chamfers instead of horizontal overhangs where possible:

\begin{lstlisting}[style=Alabaster, language=openscad]
// ❌ Horizontal lip — needs supports
translate([0, 0, 10]) cube([20, 5, 2]);

// ✅ 45° chamfer — self-supporting
hull() {
    cube([20, 2, 10]);
    translate([0, 5, 8]) cube([20, 0.01, 2]);
}

\end{lstlisting}

\textbf{Minimum feature sizes} at 0.4mm nozzle:

{\def\LTcaptype{none} % do not increment counter
\begin{longtable}[]{@{}
  >{\raggedright\arraybackslash}p{(\linewidth - 2\tabcolsep) * \real{0.3867}}
  >{\raggedright\arraybackslash}p{(\linewidth - 2\tabcolsep) * \real{0.6133}}@{}}
\toprule\noalign{}
\begin{minipage}[b]{\linewidth}\raggedright
Feature
\end{minipage} & \begin{minipage}[b]{\linewidth}\raggedright
Minimum
\end{minipage} \\
\midrule\noalign{}
\endhead
\bottomrule\noalign{}
\endlastfoot
Wall thickness & 0.8mm (2× nozzle) --- ideally 1.2mm+ \\
Hole diameter & 2mm (smaller holes tend to close up) \\
Gap / clearance between parts &
0.2--0.3mm for sliding fit, 0.1mm for press fit \\
Text / embossed detail & 0.5mm depth, 4pt font minimum \\
Pin diameter & 2mm minimum for structural pins \\
\end{longtable}
}

\textbf{Orientation matters.} Layer lines are weakest in tension perpendicular to the print direction. Design critical stress paths to run parallel to layer lines, or use the \texttt{orient} action to auto-optimize placement.

\textbf{Tolerance for fits.} PLA typically needs 0.2--0.3mm clearance for a sliding fit between printed parts. PETG and ABS may need slightly more. Build a quick tolerance test into your workflow:

\begin{lstlisting}[style=Alabaster, language=openscad]
// Tolerance test — print this and measure which gap fits your printer
for (i=[0:4])
    translate([i*15, 0, 0])
        difference() {
            cube([12, 12, 5]);
            translate([6, 6, 0])
                cylinder(d=5 + i*0.1, h=5, $fn=32);
        }

\end{lstlisting}

\subsubsection*{Using OpenSCAD Libraries with 3DMake}\label{docs__pandoc__latex__src__3dmake_appendices_complete.md__using-openscad-libraries-with-3dmake}

Declare library dependencies in \texttt{3dmake.toml}, then install them:

\begin{lstlisting}[style=Alabaster]
# 3dmake.toml
libraries = ["openscad-threads", "bosl2"]

\end{lstlisting}

\begin{lstlisting}[style=Alabaster, language=bash]
3dm install-libraries    # Downloads and installs to config dir
3dm list-libraries       # Verify installation

\end{lstlisting}

Then use them in your \texttt{.scad} files normally:

\begin{lstlisting}[style=Alabaster, language=openscad]
// BOSL2 example
include <BOSL2/std.scad>

cuboid([20, 20, 10], rounding=2, edges="Z");

\end{lstlisting}

\begin{lstlisting}[style=Alabaster, language=openscad]
// openscad-threads example
use <threads.scad>

metric_thread(diameter=8, pitch=1.25, length=20);

\end{lstlisting}

\begin{center}\rule{0.5\linewidth}{0.5pt}\end{center}

\protect\phantomsection\label{docs__pandoc__latex__src__3dmake_appendices_complete.md__appendix-b}{}

\section*{Appendix B: Comprehensive Slicing Guide --- All Major Slicers}\label{docs__pandoc__latex__src__3dmake_appendices_complete.md__appendix-b-comprehensive-slicing-guide--all-major-slicers}

\begin{center}\rule{0.5\linewidth}{0.5pt}\end{center}

This appendix covers settings, workflows, and troubleshooting for 7 major 3D printer slicers. Each slicer is represented with:

\begin{itemize}
\tightlist
\item
  Recommended settings for beginners
\item
  Screen-reader-accessible settings explanation
\item
  Common troubleshooting
\item
  Accessibility features built-in
\item
  Command-line usage (for PowerShell integration)
\end{itemize}

Referenced in: Lessons 5 (Safety), 8 (Design), 10 (Verification)

\subsection*{Overview: What is Slicing?}\label{docs__pandoc__latex__src__3dmake_appendices_complete.md__overview-what-is-slicing}

Slicing converts a 3D model into printer instructions:

\begin{lstlisting}[style=Alabaster]
SCAD Design (bracelet_holder.scad)
   v
Export to STL (3D shape file)
   v
Load into Slicer
   v
Apply settings (temperature, speed, supports, etc.)
   v
Generate G-code (printer instructions)
   v
Send to Printer
   v
Physical part

\end{lstlisting}

\subsubsection*{Core Slicing Parameters (All Slicers Share These)}\label{docs__pandoc__latex__src__3dmake_appendices_complete.md__core-slicing-parameters-all-slicers-share-these}

{\def\LTcaptype{none} % do not increment counter
\begin{longtable}[]{@{}
  >{\raggedright\arraybackslash}p{(\linewidth - 6\tabcolsep) * \real{0.1429}}
  >{\raggedright\arraybackslash}p{(\linewidth - 6\tabcolsep) * \real{0.2449}}
  >{\raggedright\arraybackslash}p{(\linewidth - 6\tabcolsep) * \real{0.1735}}
  >{\raggedright\arraybackslash}p{(\linewidth - 6\tabcolsep) * \real{0.4388}}@{}}
\toprule\noalign{}
\begin{minipage}[b]{\linewidth}\raggedright
Parameter
\end{minipage} & \begin{minipage}[b]{\linewidth}\raggedright
What It Does
\end{minipage} & \begin{minipage}[b]{\linewidth}\raggedright
Typical Range
\end{minipage} & \begin{minipage}[b]{\linewidth}\raggedright
Impact
\end{minipage} \\
\midrule\noalign{}
\endhead
\bottomrule\noalign{}
\endlastfoot
Nozzle Temp & Filament melting heat & 200-250C &
Too cold -\textgreater{} weak; too hot -\textgreater{} oozing \\
Bed Temp & Build plate heat & 50-110C &
Helps adhesion; prevents warping \\
Layer Height & Z-axis precision & 0.1-0.4mm &
Finer = slower, better detail \\
Print Speed & Movement velocity & 30-150 mm/s &
Faster = worse quality; slower = stronger \\
Infill \% & Interior density & 10-100\% & Higher = stronger + heavier \\
Support & Temporary scaffolding & On/Off &
Required for overhangs \textgreater45 \\
Bed Adhesion & First layer stickiness & Brim/Raft/Skirt &
Prevents parts lifting mid-print \\
\end{longtable}
}

\subsection*{1. PrusaSlicer (Prusa)}\label{docs__pandoc__latex__src__3dmake_appendices_complete.md__1-prusaslicer-prusa}

\subsubsection*{Overview}\label{docs__pandoc__latex__src__3dmake_appendices_complete.md__overview}

\begin{itemize}
\tightlist
\item
  Developer: Prusa Research (Czech company)
\item
  Platforms: Windows, Mac, Linux (open-source)
\item
  Best For: Beginner-friendly, excellent support, strong community
\item
  Download: \url{https://www.prusa3d.com/page/prusaslicer_410/}
\item
  Accessibility: Good font sizes, text-based profiles
\end{itemize}

\subsubsection*{Quick Setup for Beginners}\label{docs__pandoc__latex__src__3dmake_appendices_complete.md__quick-setup-for-beginners}

Step 1: Install \& Select Your Printer

\begin{lstlisting}[style=Alabaster]
1. Open PrusaSlicer
2. Go to "Help" -> "Check for Updates"
3. When prompted, select your printer model
4. Choose default profile (matches printer exactly)

\end{lstlisting}

Step 2: Load Your STL

\begin{lstlisting}[style=Alabaster]
1. Click "File" -> "Open STL Model"
2. Select your bracelet_holder.stl
3. Model appears in 3D view

\end{lstlisting}

Step 3: Essential Settings

These are the most important adjustments:

{\def\LTcaptype{none} % do not increment counter
\begin{longtable}[]{@{}
  >{\raggedright\arraybackslash}p{(\linewidth - 6\tabcolsep) * \real{0.1772}}
  >{\raggedright\arraybackslash}p{(\linewidth - 6\tabcolsep) * \real{0.2405}}
  >{\raggedright\arraybackslash}p{(\linewidth - 6\tabcolsep) * \real{0.2152}}
  >{\raggedright\arraybackslash}p{(\linewidth - 6\tabcolsep) * \real{0.3671}}@{}}
\toprule\noalign{}
\begin{minipage}[b]{\linewidth}\raggedright
Setting
\end{minipage} & \begin{minipage}[b]{\linewidth}\raggedright
Location
\end{minipage} & \begin{minipage}[b]{\linewidth}\raggedright
Beginner Value
\end{minipage} & \begin{minipage}[b]{\linewidth}\raggedright
Why
\end{minipage} \\
\midrule\noalign{}
\endhead
\bottomrule\noalign{}
\endlastfoot
Layer Height & Print Settings & 0.15mm & Balance speed \& quality \\
Infill & Print Settings & 20\% & Enough strength; fast print \\
Support & Print Settings & Yes (if needed) &
For overhangs \textgreater45 \\
Nozzle Temp & Filament Settings & 210C & Default for PLA \\
Bed Temp & Filament Settings & 60C & Standard PLA adhesion \\
Print Speed & Print Settings & 150 mm/s & Balanced quality \\
\end{longtable}
}

Step 4: Preview \& Export

\begin{lstlisting}[style=Alabaster]
1. Click "Slice now" (or G-code icon)
2. Left panel shows preview of each layer
3. Look for issues:
   - Supports covering entire model? (OK)
   - Model floating in air? (Not OK-likely error)
   - Top surface quality acceptable?
4. If satisfied, click "Export G-code"
5. Save to USB or send to printer

\end{lstlisting}

\subsubsection*{Accessible Parameter Explanations}\label{docs__pandoc__latex__src__3dmake_appendices_complete.md__accessible-parameter-explanations}

When adjusting settings, use these descriptions to understand what each does:

Layer Height (0.1-0.4mm)

\begin{itemize}
\tightlist
\item
  Lower (0.1mm): Smoother surface, more layers, slower (best for detail)
\item
  Higher (0.4mm): Faster, rougher surface (best for speed)
\item
  Recommendation: 0.15mm for balanced quality
\end{itemize}

Infill Percentage (10-100\%)

\begin{itemize}
\tightlist
\item
  10\%: Fast, uses less filament, weaker (good for prototypes)
\item
  20\%: Good balance (recommended for most prints)
\item
  50\%: Stronger, slower, heavier
\item
  100\%: Solid interior, strongest, slowest (waste of plastic)
\end{itemize}

Support Type

\begin{itemize}
\tightlist
\item
  None: Fast, good surface finish, but risky if overhangs exist
\item
  Linear (Default): Good balance-easy to remove, provides support
\item
  Grid: Extra strong support, takes longer to remove
\end{itemize}

First Layer

\begin{itemize}
\tightlist
\item
  Brim: Adds border to help adhesion (recommended for beginners)
\item
  Raft: Sacrificial platform (good if bed isn\textquotesingle t level)
\item
  Skirt: Just outline, doesn\textquotesingle t help adhesion (fastest)
\end{itemize}

\subsubsection*{Command-Line Usage (PowerShell Integration)}\label{docs__pandoc__latex__src__3dmake_appendices_complete.md__command-line-usage-powershell-integration}

\begin{lstlisting}[style=Alabaster, language=powershell]
# Slice a model automatically with PrusaSlicer
$model = "C:\Models\bracelet_holder.stl"
$output = "C:\GCode\bracelet_holder.gcode"
$config = "default"  # Use default printer profile
# Run PrusaSlicer in batch mode
& "C:\Program Files\Prusa3D\PrusaSlicer\prusa-slicer.exe" `
    --load-config-file "$config" `
    --export-gcode "$output" `
    "$model"
Write-Host "Slicing complete: $output"

\end{lstlisting}

\subsubsection*{Troubleshooting}\label{docs__pandoc__latex__src__3dmake_appendices_complete.md__troubleshooting-1}

{\def\LTcaptype{none} % do not increment counter
\begin{longtable}[]{@{}
  >{\raggedright\arraybackslash}p{(\linewidth - 4\tabcolsep) * \real{0.3061}}
  >{\raggedright\arraybackslash}p{(\linewidth - 4\tabcolsep) * \real{0.2653}}
  >{\raggedright\arraybackslash}p{(\linewidth - 4\tabcolsep) * \real{0.4286}}@{}}
\toprule\noalign{}
\begin{minipage}[b]{\linewidth}\raggedright
Problem
\end{minipage} & \begin{minipage}[b]{\linewidth}\raggedright
Cause
\end{minipage} & \begin{minipage}[b]{\linewidth}\raggedright
Solution
\end{minipage} \\
\midrule\noalign{}
\endhead
\bottomrule\noalign{}
\endlastfoot
First layer not sticking & Bed not level &
Bed leveling procedure in printer manual \\
Supports everywhere & No support type selected &
Change to "Linear" or "Grid" \\
Nozzle drags through model & Z-offset too low & Raise Z-offset +0.1mm \\
Oozing strings between parts & Temp too high &
Lower nozzle temp 5-10C \\
Print breaks mid-way & Adhesion problem & Add brim; check bed level \\
\end{longtable}
}

\subsection*{2. Bambu Studio (Bambu Lab)}\label{docs__pandoc__latex__src__3dmake_appendices_complete.md__2-bambu-studio-bambu-lab}

\subsubsection*{Overview}\label{docs__pandoc__latex__src__3dmake_appendices_complete.md__overview-1}

\begin{itemize}
\tightlist
\item
  Developer: Bambu Lab (printer manufacturer)
\item
  Platforms: Windows, Mac, Linux
\item
  Best For: Modern X1-series printers, excellent speed, AMS support
\item
  Download: \url{https://bambulab.com/en/download/studio}
\item
  Accessibility: Good contrast, keyboard navigation
\end{itemize}

\subsubsection*{Quick Setup}\label{docs__pandoc__latex__src__3dmake_appendices_complete.md__quick-setup}

Step 1: Create Account \& Connect Printer

\begin{enumerate}
\tightlist
\item
  Launch Bambu Studio
\item
  Sign in with Bambu Lab account
\item
  Select printer from network
\item
  Studio auto-detects printer settings
\end{enumerate}

Step 2: Load Model \& Configure

\begin{enumerate}
\tightlist
\item
  Drag STL into workspace
\item
  Auto-arranges on bed plate
\item
  Default profile applied automatically
\end{enumerate}

Step 3: Key Settings (Bambu-Specific)

{\def\LTcaptype{none} % do not increment counter
\begin{longtable}[]{@{}
  >{\raggedright\arraybackslash}p{(\linewidth - 6\tabcolsep) * \real{0.2179}}
  >{\raggedright\arraybackslash}p{(\linewidth - 6\tabcolsep) * \real{0.1154}}
  >{\raggedright\arraybackslash}p{(\linewidth - 6\tabcolsep) * \real{0.2949}}
  >{\raggedright\arraybackslash}p{(\linewidth - 6\tabcolsep) * \real{0.3718}}@{}}
\toprule\noalign{}
\begin{minipage}[b]{\linewidth}\raggedright
Setting
\end{minipage} & \begin{minipage}[b]{\linewidth}\raggedright
Default
\end{minipage} & \begin{minipage}[b]{\linewidth}\raggedright
Adjustment
\end{minipage} & \begin{minipage}[b]{\linewidth}\raggedright
Why
\end{minipage} \\
\midrule\noalign{}
\endhead
\bottomrule\noalign{}
\endlastfoot
AMS Multi-Color & Off & On (if AMS attached) & Auto-switch filament \\
Auto-Leveling & Enabled & Keep On & Bambu feature-very reliable \\
Nozzle Temp & 220C & Keep unless specified & Bambu-optimized \\
Layer Height & 0.2mm & 0.15mm for detail & Balance speed/quality \\
Bed Temp & 65C & 60C for PLA & Standard adhesion \\
\end{longtable}
}

Step 4: Send to Printer

\begin{enumerate}
\tightlist
\item
  Click "Prepare" (bottom right)
\item
  Review preview
\item
  Click "Send to Device"
\item
  Printer receives over WiFi
\item
  Start print from printer screen
\end{enumerate}

\subsubsection*{Accessible Parameter Explanations}\label{docs__pandoc__latex__src__3dmake_appendices_complete.md__accessible-parameter-explanations-1}

Auto-Leveling

\begin{itemize}
\tightlist
\item
  Bambu printers automatically level the nozzle before every print
\item
  Explanation: Saves manual calibration; extremely reliable
\item
  For VI users: Provides confidence that bed is properly prepared
\end{itemize}

Filament Calibration

\begin{itemize}
\tightlist
\item
  Before first use of new filament color, run "Filament Calibration"
\item
  This optimizes temperature and speed for that specific filament
\item
  Explanation: Ensures consistent color and strength
\end{itemize}

Multi-Material (AMS)

\begin{itemize}
\tightlist
\item
  If you have Auto Material System (AMS):

  \begin{itemize}
  \tightlist
  \item
    Load up to 4 filament colors
  \item
    Studio auto-switches during print
  \item
    Explanation: Multi-color prints without manual intervention
  \end{itemize}
\end{itemize}

\subsubsection*{Command-Line Usage (PowerShell)}\label{docs__pandoc__latex__src__3dmake_appendices_complete.md__command-line-usage-powershell}

\begin{lstlisting}[style=Alabaster, language=powershell]
# Slice with Bambu Studio (command-line interface)
$model = "C:\Models\bracelet_holder.stl"
$output = "C:\GCode\bracelet_holder.3mf"  # Bambu uses .3mf format
# Bambu Studio CLI
& "C:\Program Files\BambuStudio\bambu-studio.exe" `
    --output "$output" `
    "$model"
Write-Host "Slice saved: $output"
# Send to printer directly
# (Requires API key-see Bambu documentation)

\end{lstlisting}

\subsubsection*{Troubleshooting}\label{docs__pandoc__latex__src__3dmake_appendices_complete.md__troubleshooting-2}

{\def\LTcaptype{none} % do not increment counter
\begin{longtable}[]{@{}
  >{\raggedright\arraybackslash}p{(\linewidth - 4\tabcolsep) * \real{0.3407}}
  >{\raggedright\arraybackslash}p{(\linewidth - 4\tabcolsep) * \real{0.2527}}
  >{\raggedright\arraybackslash}p{(\linewidth - 4\tabcolsep) * \real{0.4066}}@{}}
\toprule\noalign{}
\begin{minipage}[b]{\linewidth}\raggedright
Problem
\end{minipage} & \begin{minipage}[b]{\linewidth}\raggedright
Cause
\end{minipage} & \begin{minipage}[b]{\linewidth}\raggedright
Solution
\end{minipage} \\
\midrule\noalign{}
\endhead
\bottomrule\noalign{}
\endlastfoot
WiFi not connecting & Network issue &
Restart printer WiFi; check SSID \\
AMS not switching & Filament not detected &
Load filament into AMS; recalibrate \\
Print quality inconsistent & Wrong filament type &
Run filament calibration \\
Nozzle crashes on first layer & Auto-level failed &
Manually check nozzle height \\
\end{longtable}
}

\subsection*{3. Cura (Ultimaker)}\label{docs__pandoc__latex__src__3dmake_appendices_complete.md__3-cura-ultimaker}

\subsubsection*{Overview}\label{docs__pandoc__latex__src__3dmake_appendices_complete.md__overview-2}

\begin{itemize}
\tightlist
\item
  Developer: Ultimaker (Dutch company, open-source)
\item
  Platforms: Windows, Mac, Linux
\item
  Best For: Broad printer support, user-friendly, good documentation
\item
  Download: \url{https://ultimaker.com/software/ultimaker-cura}
\item
  Accessibility: Clear UI, good contrast
\end{itemize}

\subsubsection*{Quick Setup}\label{docs__pandoc__latex__src__3dmake_appendices_complete.md__quick-setup-1}

Step 1: Add Your Printer

\begin{enumerate}
\tightlist
\item
  Launch Cura
\item
  Go to "Settings" (top-right)
\item
  Click "Printers" -\textgreater{} "Add Printer"
\item
  Select your printer model from list
\item
  Confirm network connection
\end{enumerate}

Step 2: Load \& Prepare Model

\begin{enumerate}
\tightlist
\item
  Drag STL into workspace
\item
  Model auto-scales if needed (confirm size)
\item
  Right-click -\textgreater{} "Support" (if overhangs need support)
\end{enumerate}

Step 3: Recommended Settings

{\def\LTcaptype{none} % do not increment counter
\begin{longtable}[]{@{}
  >{\raggedright\arraybackslash}p{(\linewidth - 4\tabcolsep) * \real{0.3143}}
  >{\raggedright\arraybackslash}p{(\linewidth - 4\tabcolsep) * \real{0.1429}}
  >{\raggedright\arraybackslash}p{(\linewidth - 4\tabcolsep) * \real{0.5429}}@{}}
\toprule\noalign{}
\begin{minipage}[b]{\linewidth}\raggedright
Setting
\end{minipage} & \begin{minipage}[b]{\linewidth}\raggedright
Value
\end{minipage} & \begin{minipage}[b]{\linewidth}\raggedright
Notes
\end{minipage} \\
\midrule\noalign{}
\endhead
\bottomrule\noalign{}
\endlastfoot
Profile & Standard & Good balance for most prints \\
Layer Height & 0.2mm & Default; change to 0.15mm for detail \\
Infill & 20\% & 100\% waste for solid parts \\
Support Angle & 50 & Auto-generates support for overhangs \\
Build Plate Adhesion & Brim & Helps first layer stick \\
Nozzle Temp & 200C & Standard PLA \\
Bed Temp & 60C & Standard PLA \\
\end{longtable}
}

Step 4: Print

\begin{enumerate}
\tightlist
\item
  Click "Slice" (bottom right)
\item
  Review layer-by-layer preview
\item
  Click "Print Over Network" or "Print to File"
\item
  Model sends to printer or saves as .gcode
\end{enumerate}

\subsubsection*{Accessible Settings Explanation}\label{docs__pandoc__latex__src__3dmake_appendices_complete.md__accessible-settings-explanation}

Combing Mode

\begin{itemize}
\tightlist
\item
  Off: Nozzle retracts on every travel (quality)
\item
  All: Never retracts (faster, possible stringing)
\item
  Not in Skin: Smart compromise
\item
  For VI users: Retraction prevents nozzle oozing on visible surfaces
\end{itemize}

Z-Offset (Z Clearance)

\begin{itemize}
\tightlist
\item
  Adjusts first-layer distance
\item
  Too low -\textgreater{} nozzle scrapes bed (bad)
\item
  Too high -\textgreater{} filament doesn\textquotesingle t stick (bad)
\item
  Correct -\textgreater{} thin line of plastic sticks to bed
\end{itemize}

Gradual Infill

\begin{itemize}
\tightlist
\item
  Automatically reduces infill strength away from surface
\item
  Saves filament while maintaining strength
\item
  Explanation: The core doesn\textquotesingle t need to be solid
\end{itemize}

\subsubsection*{Command-Line Usage (PowerShell)}\label{docs__pandoc__latex__src__3dmake_appendices_complete.md__command-line-usage-powershell-1}

\begin{lstlisting}[style=Alabaster, language=powershell]
# Cura engine CLI (CuraEngine)
$model = "C:\Models\bracelet_holder.stl"
$output = "C:\GCode\bracelet_holder.gcode"
$config = "default.cfg"
# Requires Cura to be installed; command-line slicing
& "C:\Program Files\Ultimaker Cura\CuraEngine.exe" `
    -c "$config" `
    -o "$output" `
    "$model"
Write-Host "Sliced: $output"

\end{lstlisting}

\subsubsection*{Troubleshooting}\label{docs__pandoc__latex__src__3dmake_appendices_complete.md__troubleshooting-3}

{\def\LTcaptype{none} % do not increment counter
\begin{longtable}[]{@{}
  >{\raggedright\arraybackslash}p{(\linewidth - 4\tabcolsep) * \real{0.3765}}
  >{\raggedright\arraybackslash}p{(\linewidth - 4\tabcolsep) * \real{0.2471}}
  >{\raggedright\arraybackslash}p{(\linewidth - 4\tabcolsep) * \real{0.3765}}@{}}
\toprule\noalign{}
\begin{minipage}[b]{\linewidth}\raggedright
Problem
\end{minipage} & \begin{minipage}[b]{\linewidth}\raggedright
Cause
\end{minipage} & \begin{minipage}[b]{\linewidth}\raggedright
Solution
\end{minipage} \\
\midrule\noalign{}
\endhead
\bottomrule\noalign{}
\endlastfoot
Model appears too small on bed & Scale wrong &
Right-click -\textgreater{} Scale to fit \\
Stringing between parts & Retraction disabled &
Enable retraction in settings \\
Support doesn\textquotesingle t generate & Auto-support off &
Enable "Generate Support" \\
Printer not found & Network/USB issue &
Check connection; restart Cura \\
\end{longtable}
}

\subsection*{4. SuperSlicer (Modification of Prusa)}\label{docs__pandoc__latex__src__3dmake_appendices_complete.md__4-superslicer-modification-of-prusa}

\subsubsection*{Overview}\label{docs__pandoc__latex__src__3dmake_appendices_complete.md__overview-3}

\begin{itemize}
\tightlist
\item
  Developer: Community fork of PrusaSlicer
\item
  Platforms: Windows, Mac, Linux
\item
  Best For: Advanced users wanting more control than Prusa offers
\item
  Download: \url{https://github.com/supermerill/SuperSlicer}
\item
  Accessibility: Similar to Prusa, more advanced options
\end{itemize}

\subsubsection*{Key Differences from Prusa}\label{docs__pandoc__latex__src__3dmake_appendices_complete.md__key-differences-from-prusa}

{\def\LTcaptype{none} % do not increment counter
\begin{longtable}[]{@{}lll@{}}
\toprule\noalign{}
Feature & PrusaSlicer & SuperSlicer \\
\midrule\noalign{}
\endhead
\bottomrule\noalign{}
\endlastfoot
Arachne Engine & No & Yes-better edges \\
Seam Position & Limited options & Full control \\
Pressure Equalization & No & Yes-better bridging \\
Stealth Mode & No & Yes-quieter/higher quality \\
\end{longtable}
}

\subsubsection*{When to Use SuperSlicer}\label{docs__pandoc__latex__src__3dmake_appendices_complete.md__when-to-use-superslicer}

\begin{itemize}
\tightlist
\item
  Printing challenging geometries with tight tolerances
\item
  Need advanced surface finish control
\item
  Familiar with PrusaSlicer already and want more power
\end{itemize}

\subsubsection*{Quick Start}\label{docs__pandoc__latex__src__3dmake_appendices_complete.md__quick-start}

\begin{enumerate}
\tightlist
\item
  Download SuperSlicer from GitHub
\item
  Export profile from PrusaSlicer (if you have it)
\item
  Import into SuperSlicer
\item
  All settings are compatible with PrusaSlicer
\end{enumerate}

\subsubsection*{Advanced Settings for SuperSlicer}\label{docs__pandoc__latex__src__3dmake_appendices_complete.md__advanced-settings-for-superslicer}

Arachne Engine

\begin{itemize}
\tightlist
\item
  Enables finer edges on walls
\item
  Results in cleaner, more accurate prints
\item
  Takes slightly longer to slice but worth it
\end{itemize}

Seam Positioning

\begin{itemize}
\tightlist
\item
  Random: Hides seams (good for aesthetic)
\item
  Aligned: Consistent location (good for debugging)
\item
  Rear: Always at back (recommended)
\end{itemize}

Pressure Equalization

\begin{itemize}
\tightlist
\item
  Helps with bridging (printing across gaps)
\item
  Reduces sagging on overhangs
\item
  Recommended: Enable for complex designs
\end{itemize}

\subsection*{5. OrcaSlicer (Modern Bamboo Alternative)}\label{docs__pandoc__latex__src__3dmake_appendices_complete.md__5-orcaslicer-modern-bamboo-alternative}

\subsubsection*{Overview}\label{docs__pandoc__latex__src__3dmake_appendices_complete.md__overview-4}

\begin{itemize}
\tightlist
\item
  Developer: Community (independent open-source)
\item
  Platforms: Windows, Mac, Linux
\item
  Best For: Users who want Bambu features without Bambu printer
\item
  Download: \url{https://github.com/SoftFever/OrcaSlicer}
\item
  Accessibility: Modern UI, good keyboard support
\end{itemize}

\subsubsection*{Why Orca?}\label{docs__pandoc__latex__src__3dmake_appendices_complete.md__why-orca}

OrcaSlicer brings Bambu Studio\textquotesingle s best features to any printer:

\begin{itemize}
\tightlist
\item
  Excellent defaults
\item
  Fast slicing
\item
  Good preview
\item
  Modern interface
\end{itemize}

\subsubsection*{Quick Setup}\label{docs__pandoc__latex__src__3dmake_appendices_complete.md__quick-setup-2}

\begin{enumerate}
\tightlist
\item
  Download OrcaSlicer
\item
  Select your printer (not just Bambu models)
\item
  Load STL
\item
  Uses good defaults-usually ready to print
\end{enumerate}

\subsubsection*{Key Settings}\label{docs__pandoc__latex__src__3dmake_appendices_complete.md__key-settings}

{\def\LTcaptype{none} % do not increment counter
\begin{longtable}[]{@{}lll@{}}
\toprule\noalign{}
Setting & Value & Why \\
\midrule\noalign{}
\endhead
\bottomrule\noalign{}
\endlastfoot
Wall Loops & 2 & Strong walls, visible detail \\
Internal Solid Layer & 4 & Strength for brackets/connectors \\
Infill Pattern & Grid & Balanced strength \\
Line Width & Auto & Matches nozzle diameter \\
Speed & 80 mm/s & Balanced quality/speed \\
\end{longtable}
}

\subsubsection*{Accessible Features}\label{docs__pandoc__latex__src__3dmake_appendices_complete.md__accessible-features}

Filament Manager

\begin{itemize}
\tightlist
\item
  Track filament type, color, weight used
\item
  Explanation: Know when to buy new filament
\end{itemize}

Print Time Estimation

\begin{itemize}
\tightlist
\item
  Accurate prediction of print duration
\item
  Updates during slicing
\end{itemize}

Material Presets

\begin{itemize}
\tightlist
\item
  Pre-configured settings for common filaments
\item
  Just select material type; settings auto-apply
\end{itemize}

\subsection*{6. IdeaMaker (Raise3D)}\label{docs__pandoc__latex__src__3dmake_appendices_complete.md__6-ideamaker-raise3d}

\subsubsection*{Overview}\label{docs__pandoc__latex__src__3dmake_appendices_complete.md__overview-5}

\begin{itemize}
\tightlist
\item
  Developer: Raise3D
\item
  Platforms: Windows, Mac, Linux
\item
  Best For: Raise3D printer users; advanced features
\item
  Download: \url{https://www.raise3d.com/ideamaker}
\item
  Accessibility: Professional UI, detailed settings
\end{itemize}

\subsubsection*{When to Use}\label{docs__pandoc__latex__src__3dmake_appendices_complete.md__when-to-use}

\begin{itemize}
\tightlist
\item
  Using a Raise3D printer (excellent multi-nozzle support)
\item
  Need industrial-strength slicing
\item
  Want dual-extrusion (2-color) printing
\end{itemize}

\subsubsection*{Quick Start}\label{docs__pandoc__latex__src__3dmake_appendices_complete.md__quick-start-1}

\begin{enumerate}
\tightlist
\item
  Launch IdeaMaker
\item
  Add printer (Raise3D models have built-in profiles)
\item
  Load STL
\item
  Adjust layer height and infill
\item
  Slice and send to printer
\end{enumerate}

\subsubsection*{Key Features}\label{docs__pandoc__latex__src__3dmake_appendices_complete.md__key-features}

Dual Extrusion

\begin{itemize}
\tightlist
\item
  Two nozzles = two colors in one print
\item
  Useful for: Bracelets (core + colored rim)
\item
  Requires: Coordinating two materials
\end{itemize}

Advanced Support

\begin{itemize}
\tightlist
\item
  Tree support (uses less material)
\item
  Grid support (stronger for larger parts)
\end{itemize}

Print Balancing

\begin{itemize}
\tightlist
\item
  Optimizes nozzle movement for efficiency
\item
  Reduces print time without sacrificing quality
\end{itemize}

\subsection*{7. Fusion 360 (CAD + Slicer)}\label{docs__pandoc__latex__src__3dmake_appendices_complete.md__7-fusion-360-cad--slicer}

\subsubsection*{Overview}\label{docs__pandoc__latex__src__3dmake_appendices_complete.md__overview-6}

\begin{itemize}
\tightlist
\item
  Developer: Autodesk
\item
  Platforms: Windows, Mac
\item
  Best For: If you already use Fusion 360 for CAD
\item
  Download: \url{https://www.autodesk.com/products/fusion-360}
\item
  Accessibility: Integrated environment; good for learning CAD+Slicing together
\end{itemize}

\subsubsection*{Why Integrate CAD + Slicing?}\label{docs__pandoc__latex__src__3dmake_appendices_complete.md__why-integrate-cad--slicing}

Traditional workflow:

\begin{lstlisting}[style=Alabaster]
Design in CAD -> Export to STL -> Open in Slicer -> Slice -> Print

\end{lstlisting}

Fusion 360 workflow:

\begin{lstlisting}[style=Alabaster]
Design in Fusion -> Run Print Preparation -> Slice -> Print

\end{lstlisting}

Advantage: Stay in one program; no format conversion

\subsubsection*{Quick Start}\label{docs__pandoc__latex__src__3dmake_appendices_complete.md__quick-start-2}

\begin{enumerate}
\tightlist
\item
  Design in Fusion 360 (or import STL)
\item
  Select part
\item
  Go to "3D Print" tab
\item
  Click "Print Preparation"
\item
  View auto-generated support and alignment
\item
  Adjust layer height, infill, etc.
\item
  Slice and export G-code
\end{enumerate}

\subsubsection*{Integration with 3dMake}\label{docs__pandoc__latex__src__3dmake_appendices_complete.md__integration-with-3dmake}

If using Fusion 360 for CAD:

\begin{enumerate}
\tightlist
\item
  Export SCAD design to STL
\item
  3dm export-stl src/main.scad output/main.stl
\item
  Then import into Fusion 360
\item
  For more complex designs that need refinement
\end{enumerate}

\subsection*{Universal Troubleshooting Guide}\label{docs__pandoc__latex__src__3dmake_appendices_complete.md__universal-troubleshooting-guide}

{\def\LTcaptype{none} % do not increment counter
\begin{longtable}[]{@{}
  >{\raggedright\arraybackslash}p{(\linewidth - 4\tabcolsep) * \real{0.3652}}
  >{\raggedright\arraybackslash}p{(\linewidth - 4\tabcolsep) * \real{0.2522}}
  >{\raggedright\arraybackslash}p{(\linewidth - 4\tabcolsep) * \real{0.3826}}@{}}
\toprule\noalign{}
\begin{minipage}[b]{\linewidth}\raggedright
Problem
\end{minipage} & \begin{minipage}[b]{\linewidth}\raggedright
Diagnosis
\end{minipage} & \begin{minipage}[b]{\linewidth}\raggedright
Solution
\end{minipage} \\
\midrule\noalign{}
\endhead
\bottomrule\noalign{}
\endlastfoot
Print won\textquotesingle t stick to bed & Bed temperature too low &
Raise bed temp +5C; check bed level \\
& Bed not level & Manual leveling procedure (printer manual) \\
& Build plate dirty & Clean with isopropyl alcohol \\
Nozzle hits model mid-print & Z-offset wrong &
Adjust Z-offset; re-level bed \\
& Model placed too low on bed & Use "Arrange on Bed" tool in slicer \\
Supports won\textquotesingle t remove & Too much support generated &
Reduce support density or angle \\
& Support too strong & Reduce support material percentage \\
Stringy/Oozing & Nozzle too hot & Reduce temp by 5-10C \\
& Retraction disabled & Enable retraction in settings \\
& Travel speed too fast & Reduce travel speed \\
Layer shifting (X/Y) & Belt tension off &
Check belt tension (printer manual) \\
& Stepper motor power issue & Firmware issue-check printer logs \\
Model prints poorly but slices look good & Filament quality issue &
Try different filament batch \\
& Nozzle clogged &
Unclog nozzle (heat -\textgreater{} purge -\textgreater{} clean) \\
\end{longtable}
}

\subsection*{Recommended Slicer for Each Situation}\label{docs__pandoc__latex__src__3dmake_appendices_complete.md__recommended-slicer-for-each-situation}

{\def\LTcaptype{none} % do not increment counter
\begin{longtable}[]{@{}
  >{\raggedright\arraybackslash}p{(\linewidth - 4\tabcolsep) * \real{0.3111}}
  >{\raggedright\arraybackslash}p{(\linewidth - 4\tabcolsep) * \real{0.2556}}
  >{\raggedright\arraybackslash}p{(\linewidth - 4\tabcolsep) * \real{0.4333}}@{}}
\toprule\noalign{}
\begin{minipage}[b]{\linewidth}\raggedright
Scenario
\end{minipage} & \begin{minipage}[b]{\linewidth}\raggedright
Best Slicer
\end{minipage} & \begin{minipage}[b]{\linewidth}\raggedright
Why
\end{minipage} \\
\midrule\noalign{}
\endhead
\bottomrule\noalign{}
\endlastfoot
I\textquotesingle m a beginner & PrusaSlicer &
Clear defaults, excellent UI \\
I have a Bambu printer & Bambu Studio & Native, best features \\
I have an Ultimaker & Cura & Official support, broad compatibility \\
I have a Raise3D & IdeaMaker & Official, dual-extrusion support \\
I want advanced control & SuperSlicer & Maximum customization \\
I want modern/fast slicing & OrcaSlicer & Great defaults, any printer \\
I use Fusion 360 for CAD & Fusion 360 Print Prep &
Integrated workflow \\
\end{longtable}
}

\subsection*{PowerShell Integration: Batch Slicing}\label{docs__pandoc__latex__src__3dmake_appendices_complete.md__powershell-integration-batch-slicing}

\subsubsection*{Slice Multiple Files Automatically}\label{docs__pandoc__latex__src__3dmake_appendices_complete.md__slice-multiple-files-automatically}

\begin{lstlisting}[style=Alabaster, language=powershell]
# Batch slice all SCAD designs in a project
$scadDir = "C:\Projects\3dMake\src"
$slicerConfig = "default.cfg"
$outputDir = "C:\Projects\3dMake\gcode"
# Find all SCAD files
$scadFiles = Get-ChildItem -Path $scadDir -Filter "*.scad" -Recurse
foreach ($scadFile in $scadFiles) {
    $stlFile = $scadFile.FullName -replace ".scad$", ".stl"
    $gcodeFile = Join-Path $outputDir ($scadFile.BaseName + ".gcode")
    # Export SCAD to STL
    Write-Host "Exporting: $($scadFile.Name)"
    & "C:\Program Files\OpenSCAD\openscad.exe" -o "$stlFile" "$($scadFile.FullName)"
    # Slice STL
    Write-Host "Slicing: $($scadFile.Name)"
    & "C:\Program Files\Prusa3D\PrusaSlicer\prusa-slicer.exe" `
        --load-config-file "$slicerConfig" `
        --export-gcode "$gcodeFile" `
        "$stlFile"
    Write-Host "Complete: $gcodeFile"
}
Write-Host "Batch slicing finished."

\end{lstlisting}

\subsubsection*{Monitor Printing Progress}\label{docs__pandoc__latex__src__3dmake_appendices_complete.md__monitor-printing-progress}

\begin{lstlisting}[style=Alabaster, language=powershell]
# Connect to printer API and monitor print status
# (Requires printer to support API-check documentation)
$printerIP = "192.168.1.100"  # Your printer's IP
$printerPort = 8080            # Typical API port
# Get current print status
$status = Invoke-WebRequest -Uri "https://yourdolphin.com/supernova/" `
    -UseBasicParsing | ConvertFrom-Json
Write-Host "State: $($status.state.text)"
Write-Host "Progress: $($status.progress.completion)%"
Write-Host "Time remaining: $($status.progress.printTimeLeft) seconds"

\end{lstlisting}

\subsection*{Accessibility Best Practices for Slicing}\label{docs__pandoc__latex__src__3dmake_appendices_complete.md__accessibility-best-practices-for-slicing}

\subsubsection*{Describing Sliced Models}\label{docs__pandoc__latex__src__3dmake_appendices_complete.md__describing-sliced-models}

\begin{lstlisting}[style=Alabaster, language=powershell]
# Use 3dm info to understand the slicer's output
# (Best practice: describe after slicing to verify settings)
3dm info output/bracelet_holder.stl
# Output includes:
# - Number of parts
# - Dimensions
# - Surface area
# - Volume

\end{lstlisting}

\subsubsection*{Testing Sliced Parts}\label{docs__pandoc__latex__src__3dmake_appendices_complete.md__testing-sliced-parts}

\begin{enumerate}
\item
  Weight Check:

  \begin{itemize}
  \tightlist
  \item
    Calculate expected weight from volume + material density
  \item
    Compare to actual printed weight
  \item
    Indicates if infill is correct
  \end{itemize}
\item
  Dimensional Check:

  \begin{itemize}
  \tightlist
  \item
    Use calipers to verify critical dimensions
  \item
    Compare to SCAD parameters
  \item
    Check tolerance stack-up
  \end{itemize}
\item
  Functional Test:

  \begin{itemize}
  \tightlist
  \item
    Assemble with other parts
  \item
    Test strength by loading with known weight
  \item
    Verify support removal didn\textquotesingle t damage part
  \end{itemize}
\end{enumerate}

\subsection*{Summary}\label{docs__pandoc__latex__src__3dmake_appendices_complete.md__summary}

This appendix provides:

\begin{itemize}
\tightlist
\item
  {[}YES{]} 7 slicer workflows, settings, and troubleshooting
\item
  {[}YES{]} Accessible parameter explanations
\item
  {[}YES{]} Command-line integration for PowerShell
\item
  {[}YES{]} Comparison table for choosing a slicer
\item
  {[}YES{]} Batch processing automation examples
\item
  {[}YES{]} Accessibility best practices
\end{itemize}

Use this guide whenever you:

\begin{itemize}
\tightlist
\item
  Start a new print
\item
  Switch slicers
\item
  Encounter quality issues
\item
  Want to automate slicing workflows
\item
  Need troubleshooting help
\end{itemize}

\begin{center}\rule{0.5\linewidth}{0.5pt}\end{center}

\protect\phantomsection\label{docs__pandoc__latex__src__3dmake_appendices_complete.md__appendix-c}{}

\section*{Appendix C: Material Properties and Selection Guide}\label{docs__pandoc__latex__src__3dmake_appendices_complete.md__appendix-c-material-properties-and-selection-guide}

\begin{center}\rule{0.5\linewidth}{0.5pt}\end{center}

This appendix covers material properties, characteristics, and selection criteria for 3D printing filaments. Each material includes:

\begin{itemize}
\tightlist
\item
  Physical properties (strength, flexibility, temperature tolerance)
\item
  Printing parameters (nozzle temp, bed temp, speed)
\item
  Suitability for different projects
\item
  Accessibility considerations (measurement-based verification)
\item
  Cost/availability comparison
\end{itemize}

Referenced in: Lessons 5 (Safety), 6-10 (Projects), 11 (Customer Requirements)

\subsection*{Overview: Why Material Matters}\label{docs__pandoc__latex__src__3dmake_appendices_complete.md__overview-why-material-matters}

The material you choose affects:

\begin{itemize}
\tightlist
\item
  Strength: Can the part hold weight?
\item
  Flexibility: Will it bend or break?
\item
  Temperature: Can it withstand heat?
\item
  Durability: Will it last months or years?
\item
  Cost: Budget constraints?
\item
  Printability: Ease of use for beginners?
\item
  Appearance: Texture, color, finish?
\end{itemize}

\subsubsection*{The Material Selection Flowchart}\label{docs__pandoc__latex__src__3dmake_appendices_complete.md__the-material-selection-flowchart}

\begin{lstlisting}[style=Alabaster]
What's the primary requirement?
+-- Strength & Detail?
    +-- PLA (best for beginners, detail)
        or PETG (stronger, tougher)
+-- Flexibility?
    +-- TPU/TPE (rubber-like)
+-- Heat Resistance?
    +-- ABS or Polycarbonate
+-- Transparency?
    +-- PETG or Polycarbonate
+-- Food Contact?
    +-- FDA-approved PETG or PLA
+-- Cost-Conscious?
    +-- PLA (cheapest, easiest)

\end{lstlisting}

\subsection*{1. PLA (Polylactic Acid)}\label{docs__pandoc__latex__src__3dmake_appendices_complete.md__1-pla-polylactic-acid}

\subsubsection*{Properties at a Glance}\label{docs__pandoc__latex__src__3dmake_appendices_complete.md__properties-at-a-glance}

{\def\LTcaptype{none} % do not increment counter
\begin{longtable}[]{@{}
  >{\raggedright\arraybackslash}p{(\linewidth - 4\tabcolsep) * \real{0.3038}}
  >{\raggedright\arraybackslash}p{(\linewidth - 4\tabcolsep) * \real{0.1013}}
  >{\raggedright\arraybackslash}p{(\linewidth - 4\tabcolsep) * \real{0.5949}}@{}}
\toprule\noalign{}
\begin{minipage}[b]{\linewidth}\raggedright
Property
\end{minipage} & \begin{minipage}[b]{\linewidth}\raggedright
Rating
\end{minipage} & \begin{minipage}[b]{\linewidth}\raggedright
Notes
\end{minipage} \\
\midrule\noalign{}
\endhead
\bottomrule\noalign{}
\endlastfoot
Strength & {[}3/5{]} & Good for most projects; not for dynamic loads \\
Flexibility & {[}1/5{]} & Brittle; will snap under stress \\
Temperature Resistance & {[}2/5{]} &
Softens around 60C; bad for hot environments \\
Ease of Printing & {[}5/5{]} & Most forgiving; best for beginners \\
Cost & {[}5/5{]} & Cheapest option (\textasciitilde{}\$20/kg) \\
Appearance & {[}5/5{]} & Excellent surface finish; many colors \\
Availability & {[}5/5{]} & Available everywhere \\
\end{longtable}
}

\subsubsection*{Ideal Projects}\label{docs__pandoc__latex__src__3dmake_appendices_complete.md__ideal-projects}

\begin{itemize}
\tightlist
\item
  {[}YES{]} Decorative pieces (jewelry, miniatures)
\item
  {[}YES{]} Enclosures/shells (light loads)
\item
  {[}YES{]} Prototype/mockups
\item
  {[}YES{]} Low-stress connector clips
\item
  {[}YES{]} Educational demonstrations
\item
  {[}NO{]} Load-bearing brackets
\item
  {[}NO{]} Hinges or flexing parts
\item
  {[}NO{]} Heat-resistant applications
\end{itemize}

\subsubsection*{Printing Parameters}\label{docs__pandoc__latex__src__3dmake_appendices_complete.md__printing-parameters}

{\def\LTcaptype{none} % do not increment counter
\begin{longtable}[]{@{}lll@{}}
\toprule\noalign{}
Parameter & Value & Tolerance \\
\midrule\noalign{}
\endhead
\bottomrule\noalign{}
\endlastfoot
Nozzle Temp & 200C & 190-210C (varies by brand) \\
Bed Temp & 60C & 50-65C \\
Print Speed & 50 mm/s & 40-60 mm/s \\
Retraction & 5mm @ 40 mm/s & Yes, prevents stringing \\
Cooling Fan & 100\% & High cooling improves quality \\
First Layer & Slower (25 mm/s) & Ensures adhesion \\
Layer Height & 0.2mm & 0.1-0.3mm depending on detail \\
\end{longtable}
}

\subsubsection*{PLA Variants}\label{docs__pandoc__latex__src__3dmake_appendices_complete.md__pla-variants}

Standard PLA

\begin{itemize}
\tightlist
\item
  Most common, reliable
\item
  Best for beginners
\end{itemize}

PLA Pro / Enhanced PLA

\begin{itemize}
\tightlist
\item
  Slightly stronger than standard PLA
\item
  Same temperature parameters
\item
  \textasciitilde{}10\% cost premium
\end{itemize}

Silk PLA

\begin{itemize}
\tightlist
\item
  Glossy finish instead of matte
\item
  Same strength as standard PLA
\item
  Slightly slower to print
\end{itemize}

Marble or Color-Changing PLA

\begin{itemize}
\tightlist
\item
  Visual effects
\item
  Same printing parameters
\item
  Mostly for aesthetics
\end{itemize}

\subsubsection*{Common Issues \& Solutions}\label{docs__pandoc__latex__src__3dmake_appendices_complete.md__common-issues--solutions}

{\def\LTcaptype{none} % do not increment counter
\begin{longtable}[]{@{}
  >{\raggedright\arraybackslash}p{(\linewidth - 4\tabcolsep) * \real{0.2373}}
  >{\raggedright\arraybackslash}p{(\linewidth - 4\tabcolsep) * \real{0.2797}}
  >{\raggedright\arraybackslash}p{(\linewidth - 4\tabcolsep) * \real{0.4831}}@{}}
\toprule\noalign{}
\begin{minipage}[b]{\linewidth}\raggedright
Problem
\end{minipage} & \begin{minipage}[b]{\linewidth}\raggedright
Cause
\end{minipage} & \begin{minipage}[b]{\linewidth}\raggedright
Solution
\end{minipage} \\
\midrule\noalign{}
\endhead
\bottomrule\noalign{}
\endlastfoot
Warping on corners & Bed too hot or cooling too fast &
Reduce bed temp to 50C; disable cooling for first layer \\
Stringing & Temp too high & Lower nozzle temp 5C \\
Poor layer adhesion & Nozzle too high & Lower Z-offset 0.1mm \\
Brittleness after print & Normal for PLA &
Not a problem; expected behavior \\
Nozzle clogs on retraction & Temperature inconsistency &
Ensure stable nozzle temp +/- 5C \\
\end{longtable}
}

\subsubsection*{Why PLA for This Course}\label{docs__pandoc__latex__src__3dmake_appendices_complete.md__why-pla-for-this-course}

PLA is recommended for all Lessons 1-11 because:

\begin{enumerate}
\tightlist
\item
  Beginner-friendly: Most forgiving material
\item
  Predictable: Consistent across different printers
\item
  Cost-effective: Maximize printing volume on student budget
\item
  Accessibility: Easier to troubleshoot for new users
\item
  Safe: Non-toxic, low fume emission
\item
  Available: Found at any 3D printing supplier
\end{enumerate}

\subsection*{2. PETG (Polyethylene Terephthalate Glycol)}\label{docs__pandoc__latex__src__3dmake_appendices_complete.md__2-petg-polyethylene-terephthalate-glycol}

\subsubsection*{Properties at a Glance}\label{docs__pandoc__latex__src__3dmake_appendices_complete.md__properties-at-a-glance-1}

{\def\LTcaptype{none} % do not increment counter
\begin{longtable}[]{@{}
  >{\raggedright\arraybackslash}p{(\linewidth - 4\tabcolsep) * \real{0.2963}}
  >{\raggedright\arraybackslash}p{(\linewidth - 4\tabcolsep) * \real{0.0988}}
  >{\raggedright\arraybackslash}p{(\linewidth - 4\tabcolsep) * \real{0.6049}}@{}}
\toprule\noalign{}
\begin{minipage}[b]{\linewidth}\raggedright
Property
\end{minipage} & \begin{minipage}[b]{\linewidth}\raggedright
Rating
\end{minipage} & \begin{minipage}[b]{\linewidth}\raggedright
Notes
\end{minipage} \\
\midrule\noalign{}
\endhead
\bottomrule\noalign{}
\endlastfoot
Strength & {[}4/5{]} & Tougher than PLA; better for load-bearing \\
Flexibility & {[}2/5{]} & Better than PLA but still primarily rigid \\
Temperature Resistance & {[}3/5{]} &
Softens around 85C; better than PLA \\
Ease of Printing & {[}4/5{]} & Slightly more challenging than PLA \\
Cost & {[}4/5{]} & \textasciitilde{}\$25/kg (slightly more than PLA) \\
Appearance & {[}4/5{]} &
Good surface finish; slightly glossier than PLA \\
Availability & {[}5/5{]} & Widely available \\
\end{longtable}
}

\subsubsection*{When to Use PETG Instead of PLA}\label{docs__pandoc__latex__src__3dmake_appendices_complete.md__when-to-use-petg-instead-of-pla}

\begin{itemize}
\tightlist
\item
  Functional brackets or mounts (where strength matters)
\item
  Parts that may be load-bearing (shelves, holders)
\item
  Outdoors/higher temperature environments
\item
  Parts requiring transparency (clear PETG available)
\item
  Better moisture resistance (vs. PLA)
\end{itemize}

\subsubsection*{Printing Parameters}\label{docs__pandoc__latex__src__3dmake_appendices_complete.md__printing-parameters-1}

{\def\LTcaptype{none} % do not increment counter
\begin{longtable}[]{@{}lll@{}}
\toprule\noalign{}
Parameter & Value & Tolerance \\
\midrule\noalign{}
\endhead
\bottomrule\noalign{}
\endlastfoot
Nozzle Temp & 235C & 230-245C \\
Bed Temp & 80C & 75-85C \\
Print Speed & 50 mm/s & 40-60 mm/s (same as PLA) \\
Retraction & 4mm @ 40 mm/s & Slightly less than PLA \\
Cooling Fan & 30-50\% & Less cooling than PLA \\
First Layer & Normal (50 mm/s) & Harder to adjust than PLA \\
Layer Height & 0.2mm & 0.1-0.3mm \\
\end{longtable}
}

\subsubsection*{Comparison: PETG vs. PLA}\label{docs__pandoc__latex__src__3dmake_appendices_complete.md__comparison-petg-vs-pla}

{\def\LTcaptype{none} % do not increment counter
\begin{longtable}[]{@{}lll@{}}
\toprule\noalign{}
Aspect & PLA & PETG \\
\midrule\noalign{}
\endhead
\bottomrule\noalign{}
\endlastfoot
Nozzle Temp & 200C & 235C \\
Bed Temp & 60C & 80C \\
Strength & Good & Better (20\% stronger) \\
Flexibility & Brittle & More resilient \\
Heat Tolerance & 60C & 85C \\
Ease & Very easy & Easy (needs tweaking) \\
Cost & \$20/kg & \$25/kg \\
Outdoor Use & Not ideal & Better \\
\end{longtable}
}

\subsubsection*{Projects for PETG}\label{docs__pandoc__latex__src__3dmake_appendices_complete.md__projects-for-petg}

\begin{itemize}
\tightlist
\item
  {[}YES{]} Phone stand (needs strength)
\item
  {[}YES{]} Bracket or shelf support
\item
  {[}YES{]} Flexible clip (needs resilience)
\item
  {[}YES{]} Outdoor item
\item
  {[}YES{]} Clear enclosure (if using clear PETG)
\item
  {[}NO{]} Fine detail work (slightly courser finish)
\item
  {[}NO{]} Food-contact items (not food-safe unless specified)
\end{itemize}

\subsection*{3. ABS (Acrylonitrile Butadiene Styrene)}\label{docs__pandoc__latex__src__3dmake_appendices_complete.md__3-abs-acrylonitrile-butadiene-styrene}

\subsubsection*{Properties at a Glance}\label{docs__pandoc__latex__src__3dmake_appendices_complete.md__properties-at-a-glance-2}

{\def\LTcaptype{none} % do not increment counter
\begin{longtable}[]{@{}
  >{\raggedright\arraybackslash}p{(\linewidth - 4\tabcolsep) * \real{0.3038}}
  >{\raggedright\arraybackslash}p{(\linewidth - 4\tabcolsep) * \real{0.1013}}
  >{\raggedright\arraybackslash}p{(\linewidth - 4\tabcolsep) * \real{0.5949}}@{}}
\toprule\noalign{}
\begin{minipage}[b]{\linewidth}\raggedright
Property
\end{minipage} & \begin{minipage}[b]{\linewidth}\raggedright
Rating
\end{minipage} & \begin{minipage}[b]{\linewidth}\raggedright
Notes
\end{minipage} \\
\midrule\noalign{}
\endhead
\bottomrule\noalign{}
\endlastfoot
Strength & {[}4/5{]} & Similar to PETG; good impact resistance \\
Flexibility & {[}3/5{]} & More flexible than PETG \\
Temperature Resistance & {[}5/5{]} &
Softens around 105C; best of common materials \\
Ease of Printing & {[}2/5{]} &
Requires enclosure/heated bed; challenging \\
Cost & {[}4/5{]} & \textasciitilde{}\$25-30/kg \\
Appearance & {[}3/5{]} & Rougher than PLA; requires post-processing \\
Availability & {[}4/5{]} & Good availability; not as universal as PLA \\
\end{longtable}
}

\subsubsection*{When to Use ABS}\label{docs__pandoc__latex__src__3dmake_appendices_complete.md__when-to-use-abs}

\begin{itemize}
\tightlist
\item
  High-temperature environments (near heat sources)
\item
  Mechanical parts (gears, bearings)
\item
  Durability (parts lasting years)
\item
  Post-processing (can be sanded, glued, vapor-smoothed)
\item
  Professional applications (not toys/decorative)
\end{itemize}

\subsubsection*{Why ABS is Challenging}\label{docs__pandoc__latex__src__3dmake_appendices_complete.md__why-abs-is-challenging}

ABS requires:

\begin{enumerate}
\tightlist
\item
  Enclosed environment: Minimize temperature fluctuations
\item
  Heated bed: 100C+ (much higher than PLA)
\item
  Controlled cooling: Too-fast cooling causes warping
\item
  Ventilation: ABS emits fumes (acetone-like smell)
\end{enumerate}

\subsubsection*{Printing Parameters}\label{docs__pandoc__latex__src__3dmake_appendices_complete.md__printing-parameters-2}

{\def\LTcaptype{none} % do not increment counter
\begin{longtable}[]{@{}lll@{}}
\toprule\noalign{}
Parameter & Value & Tolerance \\
\midrule\noalign{}
\endhead
\bottomrule\noalign{}
\endlastfoot
Nozzle Temp & 240C & 230-250C \\
Bed Temp & 100C & 95-105C \\
Enclosure & Required & Maintains heat (reduces warping) \\
Print Speed & 40 mm/s & Slower than PLA \\
Retraction & 3mm @ 30 mm/s & Very short \\
Cooling Fan & 0\% & OFF (causes warping) \\
First Layer & Slow (30 mm/s) & Critical for adhesion \\
\end{longtable}
}

\subsubsection*{Projects for ABS}\label{docs__pandoc__latex__src__3dmake_appendices_complete.md__projects-for-abs}

\begin{itemize}
\tightlist
\item
  {[}YES{]} Mechanical components
\item
  {[}YES{]} Heat-resistant enclosure
\item
  {[}YES{]} Durable outdoor item
\item
  {[}YES{]} Professional prototypes
\item
  {[}NO{]} Beginner projects (too challenging)
\item
  {[}NO{]} Decorative/aesthetic work (not recommended)
\item
  {[}NO{]} Flexible parts
\end{itemize}

\subsubsection*{Accessibility Note}\label{docs__pandoc__latex__src__3dmake_appendices_complete.md__accessibility-note}

ABS is not recommended for this course because:

\begin{enumerate}
\tightlist
\item
  Requires enclosed printer (expensive for beginners)
\item
  High failure rate for inexperienced users
\item
  Strong odor (ventilation concerns for some users)
\item
  Requires extra equipment (acetone for post-processing)
\end{enumerate}

\subsection*{4. TPU / TPE (Thermoplastic Polyurethane / Elastomer)}\label{docs__pandoc__latex__src__3dmake_appendices_complete.md__4-tpu--tpe-thermoplastic-polyurethane--elastomer}

\subsubsection*{Properties at a Glance}\label{docs__pandoc__latex__src__3dmake_appendices_complete.md__properties-at-a-glance-3}

{\def\LTcaptype{none} % do not increment counter
\begin{longtable}[]{@{}
  >{\raggedright\arraybackslash}p{(\linewidth - 4\tabcolsep) * \real{0.3038}}
  >{\raggedright\arraybackslash}p{(\linewidth - 4\tabcolsep) * \real{0.1013}}
  >{\raggedright\arraybackslash}p{(\linewidth - 4\tabcolsep) * \real{0.5949}}@{}}
\toprule\noalign{}
\begin{minipage}[b]{\linewidth}\raggedright
Property
\end{minipage} & \begin{minipage}[b]{\linewidth}\raggedright
Rating
\end{minipage} & \begin{minipage}[b]{\linewidth}\raggedright
Notes
\end{minipage} \\
\midrule\noalign{}
\endhead
\bottomrule\noalign{}
\endlastfoot
Strength & {[}3/5{]} & Good; impact-resistant \\
Flexibility & {[}5/5{]} & Very flexible; rubber-like \\
Temperature Resistance & {[}3/5{]} & Softens around 80C; moderate \\
Ease of Printing & {[}3/5{]} &
Needs tweaking; flexible materials are tricky \\
Cost & {[}3/5{]} & \textasciitilde{}\$30-40/kg (expensive) \\
Appearance & {[}4/5{]} & Smooth; feels good tactilely \\
Availability & {[}4/5{]} & Growing availability \\
\end{longtable}
}

\subsubsection*{What is TPU?}\label{docs__pandoc__latex__src__3dmake_appendices_complete.md__what-is-tpu}

TPU is a flexible rubber-like plastic that:

\begin{itemize}
\tightlist
\item
  Doesn\textquotesingle t crack when bent
\item
  Absorbs impact
\item
  Returns to original shape
\item
  Bridges the gap between plastic and rubber
\end{itemize}

\subsubsection*{Printing Parameters}\label{docs__pandoc__latex__src__3dmake_appendices_complete.md__printing-parameters-3}

{\def\LTcaptype{none} % do not increment counter
\begin{longtable}[]{@{}
  >{\raggedright\arraybackslash}p{(\linewidth - 4\tabcolsep) * \real{0.1688}}
  >{\raggedright\arraybackslash}p{(\linewidth - 4\tabcolsep) * \real{0.2078}}
  >{\raggedright\arraybackslash}p{(\linewidth - 4\tabcolsep) * \real{0.6234}}@{}}
\toprule\noalign{}
\begin{minipage}[b]{\linewidth}\raggedright
Parameter
\end{minipage} & \begin{minipage}[b]{\linewidth}\raggedright
Value
\end{minipage} & \begin{minipage}[b]{\linewidth}\raggedright
Tolerance
\end{minipage} \\
\midrule\noalign{}
\endhead
\bottomrule\noalign{}
\endlastfoot
Nozzle Temp & 215C & 210-225C \\
Bed Temp & 60C & 50-70C \\
Print Speed & 20-30 mm/s & VERY SLOW (flexibility needs time) \\
Retraction & Minimal or Off &
0-1mm (flexible material doesn\textquotesingle t retract well) \\
Cooling Fan & 0\% & Off (material needs heat) \\
Line Width & 0.5mm & Wider than normal (flexible material bridges) \\
\end{longtable}
}

\subsubsection*{Projects for TPU}\label{docs__pandoc__latex__src__3dmake_appendices_complete.md__projects-for-tpu}

\begin{itemize}
\tightlist
\item
  {[}YES{]} Phone case (needs flexibility + protection)
\item
  {[}YES{]} Flex joints / hinges
\item
  {[}YES{]} Gasket or seal
\item
  {[}YES{]} Shoe insert or orthotic
\item
  {[}YES{]} Tactile button (for accessibility)
\item
  {[}NO{]} Fine detail work (too stretchy)
\item
  {[}NO{]} Decorative items (usually not aesthetic)
\item
  {[}NO{]} Precision parts
\end{itemize}

\subsubsection*{Challenges with TPU}\label{docs__pandoc__latex__src__3dmake_appendices_complete.md__challenges-with-tpu}

\begin{enumerate}
\tightlist
\item
  Very slow printing: 5-10x slower than PLA
\item
  Stringing: Flexible material tends to ooze
\item
  Flexible bed needed: Standard beds may not work
\item
  Post-processing difficult: Hard to sand/glue
\end{enumerate}

\subsubsection*{Why TPU for Accessibility}\label{docs__pandoc__latex__src__3dmake_appendices_complete.md__why-tpu-for-accessibility}

TPU is excellent for accessibility because:

\begin{itemize}
\tightlist
\item
  Can create tactile buttons/indicators
\item
  Flexible grips for ergonomic handles
\item
  Gaskets that don\textquotesingle t damage delicate equipment
\item
  Accessible because: Achievable by all users if given proper guidance
\end{itemize}

\subsection*{5. Polycarbonate (PC)}\label{docs__pandoc__latex__src__3dmake_appendices_complete.md__5-polycarbonate-pc}

\subsubsection*{Properties at a Glance}\label{docs__pandoc__latex__src__3dmake_appendices_complete.md__properties-at-a-glance-4}

{\def\LTcaptype{none} % do not increment counter
\begin{longtable}[]{@{}lll@{}}
\toprule\noalign{}
Property & Rating & Notes \\
\midrule\noalign{}
\endhead
\bottomrule\noalign{}
\endlastfoot
Strength & {[}5/5{]} & Extremely strong; impact-resistant \\
Flexibility & {[}2/5{]} & Rigid; similar to PETG \\
Temperature Resistance & {[}5/5{]} & Best; softens around 130C \\
Ease of Printing & {[}2/5{]} & Difficult; prone to warping \\
Cost & {[}2/5{]} & Most expensive (\textasciitilde{}\$50+/kg) \\
Appearance & {[}4/5{]} & Transparent/translucent options \\
Availability & {[}2/5{]} & Limited; specialty supply \\
\end{longtable}
}

\subsubsection*{When to Use Polycarbonate}\label{docs__pandoc__latex__src__3dmake_appendices_complete.md__when-to-use-polycarbonate}

\begin{itemize}
\tightlist
\item
  Highest strength required
\item
  Transparent/bullet-proof enclosure needed
\item
  Extreme temperature environment
\item
  Professional/industrial applications
\end{itemize}

\subsubsection*{Why Polycarbonate is Not for This Course}\label{docs__pandoc__latex__src__3dmake_appendices_complete.md__why-polycarbonate-is-not-for-this-course}

\begin{enumerate}
\tightlist
\item
  Extreme difficulty (high failure rate)
\item
  Very expensive (3-5x cost of PLA)
\item
  Requires industrial-grade printer
\item
  Post-processing complex
\end{enumerate}

\subsection*{6. Nylon (PA)}\label{docs__pandoc__latex__src__3dmake_appendices_complete.md__6-nylon-pa}

\subsubsection*{Properties at a Glance}\label{docs__pandoc__latex__src__3dmake_appendices_complete.md__properties-at-a-glance-5}

{\def\LTcaptype{none} % do not increment counter
\begin{longtable}[]{@{}
  >{\raggedright\arraybackslash}p{(\linewidth - 4\tabcolsep) * \real{0.3380}}
  >{\raggedright\arraybackslash}p{(\linewidth - 4\tabcolsep) * \real{0.1127}}
  >{\raggedright\arraybackslash}p{(\linewidth - 4\tabcolsep) * \real{0.5493}}@{}}
\toprule\noalign{}
\begin{minipage}[b]{\linewidth}\raggedright
Property
\end{minipage} & \begin{minipage}[b]{\linewidth}\raggedright
Rating
\end{minipage} & \begin{minipage}[b]{\linewidth}\raggedright
Notes
\end{minipage} \\
\midrule\noalign{}
\endhead
\bottomrule\noalign{}
\endlastfoot
Strength & {[}5/5{]} & Very strong; can be flexible \\
Flexibility & {[}4/5{]} & More flexible than ABS \\
Temperature Resistance & {[}4/5{]} & Good; softens around 120C \\
Ease of Printing & {[}2/5{]} & Difficult; very temperature-sensitive \\
Cost & {[}3/5{]} & \textasciitilde{}\$30-40/kg \\
Appearance & {[}3/5{]} & Matte finish; less aesthetic than PLA \\
Availability & {[}3/5{]} & Growing but limited \\
\end{longtable}
}

\subsubsection*{Nylon Use Cases}\label{docs__pandoc__latex__src__3dmake_appendices_complete.md__nylon-use-cases}

\begin{itemize}
\tightlist
\item
  {[}YES{]} Mechanical parts (gears, hinges)
\item
  {[}YES{]} Flexible connectors
\item
  {[}YES{]} Threads/screws
\item
  {[}YES{]} High-stress applications
\item
  {[}NO{]} Not for beginners
\end{itemize}

\subsection*{Material Comparison Table}\label{docs__pandoc__latex__src__3dmake_appendices_complete.md__material-comparison-table}

{\def\LTcaptype{none} % do not increment counter
\begin{longtable}[]{@{}
  >{\raggedright\arraybackslash}p{(\linewidth - 16\tabcolsep) * \real{0.1124}}
  >{\raggedright\arraybackslash}p{(\linewidth - 16\tabcolsep) * \real{0.0899}}
  >{\raggedright\arraybackslash}p{(\linewidth - 16\tabcolsep) * \real{0.0674}}
  >{\raggedright\arraybackslash}p{(\linewidth - 16\tabcolsep) * \real{0.1124}}
  >{\raggedright\arraybackslash}p{(\linewidth - 16\tabcolsep) * \real{0.1461}}
  >{\raggedright\arraybackslash}p{(\linewidth - 16\tabcolsep) * \real{0.0787}}
  >{\raggedright\arraybackslash}p{(\linewidth - 16\tabcolsep) * \real{0.0787}}
  >{\raggedright\arraybackslash}p{(\linewidth - 16\tabcolsep) * \real{0.0674}}
  >{\raggedright\arraybackslash}p{(\linewidth - 16\tabcolsep) * \real{0.2472}}@{}}
\toprule\noalign{}
\begin{minipage}[b]{\linewidth}\raggedright
Material
\end{minipage} & \begin{minipage}[b]{\linewidth}\raggedright
Nozzle
\end{minipage} & \begin{minipage}[b]{\linewidth}\raggedright
Bed
\end{minipage} & \begin{minipage}[b]{\linewidth}\raggedright
Strength
\end{minipage} & \begin{minipage}[b]{\linewidth}\raggedright
Flexibility
\end{minipage} & \begin{minipage}[b]{\linewidth}\raggedright
Heat
\end{minipage} & \begin{minipage}[b]{\linewidth}\raggedright
Ease
\end{minipage} & \begin{minipage}[b]{\linewidth}\raggedright
Cost
\end{minipage} & \begin{minipage}[b]{\linewidth}\raggedright
Best For
\end{minipage} \\
\midrule\noalign{}
\endhead
\bottomrule\noalign{}
\endlastfoot
PLA & 200C & 60C & {[}3/5{]} & {[}1/5{]} & {[}2/5{]} & {[}5/5{]} & \$ &
Beginners, detail \\
PETG & 235C & 80C & {[}4/5{]} & {[}2/5{]} & {[}3/5{]} & {[}4/5{]} & \$\$
& Functional parts \\
ABS & 240C & 100C & {[}4/5{]} & {[}3/5{]} & {[}5/5{]} & {[}2/5{]} & \$\$
& High-temp/mechanical \\
TPU & 215C & 60C & {[}3/5{]} & {[}5/5{]} & {[}3/5{]} & {[}3/5{]} &
\$\$\$ & Flexibility/tactile \\
PC & 280C & 110C & {[}5/5{]} & {[}2/5{]} & {[}5/5{]} & {[}1/5{]} &
\$\$\$\$ & Extreme strength \\
Nylon & 250C & 85C & {[}5/5{]} & {[}4/5{]} & {[}4/5{]} & {[}2/5{]} &
\$\$\$ & Mechanical/flexible \\
\end{longtable}
}

\subsection*{Filament Quality Factors}\label{docs__pandoc__latex__src__3dmake_appendices_complete.md__filament-quality-factors}

\subsubsection*{Why Not All PLA is the Same}\label{docs__pandoc__latex__src__3dmake_appendices_complete.md__why-not-all-pla-is-the-same}

Diameter Tolerance

\begin{itemize}
\tightlist
\item
  Good filament: +/- 0.03mm
\item
  Poor filament: +/- 0.1mm or worse
\item
  Impact: Inconsistent extrusion, layer quality varies
\end{itemize}

Dryness

\begin{itemize}
\tightlist
\item
  PLA absorbs moisture from air
\item
  Wet filament = weak prints + bubbles
\item
  Solution: Store in sealed container with desiccant
\end{itemize}

Color Consistency

\begin{itemize}
\tightlist
\item
  Good brands: Same color throughout
\item
  Poor brands: Color varies spool-to-spool
\end{itemize}

Impurities

\begin{itemize}
\tightlist
\item
  Good: Minimal impurities
\item
  Poor: Visible specs/contaminants -\textgreater{} possible clogs
\end{itemize}

\subsubsection*{How to Check Filament Quality (Non-Visually)}\label{docs__pandoc__latex__src__3dmake_appendices_complete.md__how-to-check-filament-quality-non-visually}

\begin{enumerate}
\item
  Weight Check

  \begin{lstlisting}[style=Alabaster]
  Known: 1kg spool
  Weigh spool + filament
  Calculate remaining filament
  Should match spool markings

  \end{lstlisting}
\item
  Diameter Check

  \begin{lstlisting}[style=Alabaster]
  Use caliper to measure multiple points
  Should be consistent +/- 0.03mm
  If varying, likely lower quality

  \end{lstlisting}
\item
  Dryness Test

  \begin{lstlisting}[style=Alabaster]
  Feel texture: Should be smooth, not tacky
  Smell: Should be neutral (not musty)
  If wet, store in sealed container with desiccant

  \end{lstlisting}
\item
  Print Test

  \begin{lstlisting}[style=Alabaster]
  Print small cube (20mm x 20mm x 20mm)
  Inspect surface: Smooth or bubbly?
  Weight: Does it match expected weight?

  \end{lstlisting}
\end{enumerate}

\subsection*{Storage \& Maintenance}\label{docs__pandoc__latex__src__3dmake_appendices_complete.md__storage--maintenance}

\subsubsection*{Proper Filament Storage}\label{docs__pandoc__latex__src__3dmake_appendices_complete.md__proper-filament-storage}

Temperature

\begin{itemize}
\tightlist
\item
  Store between 15-25C
\item
  Avoid direct sunlight (fades color, degrades material)
\end{itemize}

Humidity

\begin{itemize}
\tightlist
\item
  Keep below 40\% humidity
\item
  Use desiccant packets in sealed containers
\item
  Change desiccant every 2-3 months
\end{itemize}

Organization

\begin{itemize}
\tightlist
\item
  Label spools with: Material, color, date opened, approx. remaining
\item
  Store vertically or on spindle (prevents kinking)
\end{itemize}

\subsubsection*{Filament Degradation Signs}\label{docs__pandoc__latex__src__3dmake_appendices_complete.md__filament-degradation-signs}

{\def\LTcaptype{none} % do not increment counter
\begin{longtable}[]{@{}
  >{\raggedright\arraybackslash}p{(\linewidth - 4\tabcolsep) * \real{0.3523}}
  >{\raggedright\arraybackslash}p{(\linewidth - 4\tabcolsep) * \real{0.3182}}
  >{\raggedright\arraybackslash}p{(\linewidth - 4\tabcolsep) * \real{0.3295}}@{}}
\toprule\noalign{}
\begin{minipage}[b]{\linewidth}\raggedright
Sign
\end{minipage} & \begin{minipage}[b]{\linewidth}\raggedright
Cause
\end{minipage} & \begin{minipage}[b]{\linewidth}\raggedright
Solution
\end{minipage} \\
\midrule\noalign{}
\endhead
\bottomrule\noalign{}
\endlastfoot
Weak prints / breaking easily & Filament aged or wet &
Replace with fresh filament \\
Discoloration or spots & Oxidation or contamination &
Not usable; discard safely \\
Brittle or crumbly & Overheated or UV damage & Not usable; discard \\
Slight fading (color) & UV exposure & Still usable; just faded \\
\end{longtable}
}

\subsection*{Cost Analysis}\label{docs__pandoc__latex__src__3dmake_appendices_complete.md__cost-analysis}

\subsubsection*{Cost per Project}\label{docs__pandoc__latex__src__3dmake_appendices_complete.md__cost-per-project}

\begin{lstlisting}[style=Alabaster]
Material cost = (Filament weight used) x (Cost per kg)

Example: Bracelet holder
- Weight: 25 grams
- Material: PLA at $20/kg
- Cost: (25g / 1000g) x $20 = $0.50

vs. PETG at $25/kg: $0.625
vs. ABS at $30/kg: $0.75

\end{lstlisting}

\subsubsection*{Budget Tips}\label{docs__pandoc__latex__src__3dmake_appendices_complete.md__budget-tips}

\begin{enumerate}
\tightlist
\item
  Buy bulk: 5kg spool is cheaper per gram than 1kg
\item
  Buy sales: 30-40\% discounts common during sales
\item
  Brand matters: Premium brands slightly more but more reliable
\item
  Generic brands: Often acceptable if reviews are good
\end{enumerate}

\subsubsection*{Recommended Brands (Ranked by Beginner-Friendliness)}\label{docs__pandoc__latex__src__3dmake_appendices_complete.md__recommended-brands-ranked-by-beginner-friendliness}

{\def\LTcaptype{none} % do not increment counter
\begin{longtable}[]{@{}
  >{\raggedright\arraybackslash}p{(\linewidth - 8\tabcolsep) * \real{0.0659}}
  >{\raggedright\arraybackslash}p{(\linewidth - 8\tabcolsep) * \real{0.1648}}
  >{\raggedright\arraybackslash}p{(\linewidth - 8\tabcolsep) * \real{0.3736}}
  >{\raggedright\arraybackslash}p{(\linewidth - 8\tabcolsep) * \real{0.0659}}
  >{\raggedright\arraybackslash}p{(\linewidth - 8\tabcolsep) * \real{0.3297}}@{}}
\toprule\noalign{}
\begin{minipage}[b]{\linewidth}\raggedright
Rank
\end{minipage} & \begin{minipage}[b]{\linewidth}\raggedright
Brand
\end{minipage} & \begin{minipage}[b]{\linewidth}\raggedright
Known For
\end{minipage} & \begin{minipage}[b]{\linewidth}\raggedright
Cost
\end{minipage} & \begin{minipage}[b]{\linewidth}\raggedright
Notes
\end{minipage} \\
\midrule\noalign{}
\endhead
\bottomrule\noalign{}
\endlastfoot
1 & Prusament & Reliability, Prusa compatibility & \$\$\$ &
Best; excellent support \\
2 & MatterHackers & Quality, variety & \$\$ &
Very good; educational focus \\
3 & Fillamentum & European quality & \$\$ & Excellent; eco-friendly \\
4 & Overture & Value, consistency & \$ & Good budget option \\
5 & eSUN & Variety, affordable & \$ & Decent; variable quality \\
\end{longtable}
}

\subsection*{Material Selection Decision Tree}\label{docs__pandoc__latex__src__3dmake_appendices_complete.md__material-selection-decision-tree}

\begin{lstlisting}[style=Alabaster]
START HERE: What's most important?

GOAL: Beginner success?
  YES - PLA (best choice)
  NO  - Next question

GOAL: Strength matters?
  YES - PETG or Nylon
  NO  - Next question

GOAL: Flexibility needed?
  YES - TPU (rubber-like)
  NO  - Next question

GOAL: High temperature?
  YES - ABS or Polycarbonate
  NO  - Next question

GOAL: Transparent?
  YES - Clear PETG or Polycarbonate
  NO  - Use PLA

FINAL CHOICE:
- If unsure, use PLA
- If needs strength, use PETG
- If needs flexibility, use TPU
- If needs heat, use ABS

\end{lstlisting}

\subsection*{PowerShell Integration: Track Material Usage}\label{docs__pandoc__latex__src__3dmake_appendices_complete.md__powershell-integration-track-material-usage}

\begin{lstlisting}[style=Alabaster, language=powershell]
# Track filament usage across all projects
$materialLog = @"
ProjectName,Material,ColorName,WeightUsed(g),DatePrinted,Notes
"@
# Example entries
$materialLog += "`nBraceletHolder,PLA,NaturalWhite,25.3,2024-01-15,Final design"
$materialLog += "`nPhoneStand,PETG,Black,47.2,2024-01-16,Needs strength"
$materialLog += "`nKeycap,PLA,Red,12.5,2024-01-17,Prototype"
# Save to CSV
$materialLog | Out-File "C:\Projects\material-log.csv"
# Analyze usage
$materials = Import-Csv "C:\Projects\material-log.csv"
$totalWeight = ($materials | Measure-Object -Property "WeightUsed(g)" -Sum).Sum
$avgWeight = ($materials | Measure-Object -Property "WeightUsed(g)" -Average).Average
Write-Host "Total weight used: $totalWeight grams"
Write-Host "Average per project: $avgWeight grams"
# Calculate cost (PLA at $20/kg)
$costPerKg = 20
$totalCost = ($totalWeight / 1000) * $costPerKm
Write-Host "Estimated cost: $$($totalCost.ToString("F2"))"

\end{lstlisting}

\subsection*{Summary}\label{docs__pandoc__latex__src__3dmake_appendices_complete.md__summary-1}

Key Takeaways:

\begin{enumerate}
\tightlist
\item
  Start with PLA: Best for learning; most forgiving
\item
  Understand the tradeoffs: Strength vs. ease, cost vs. quality
\item
  Match material to project: Decorative = PLA; functional = PETG
\item
  Store properly: Desiccant, cool, dark location
\item
  Track usage: Know what works for future projects
\item
  Test before committing: Print small test on new filament
\end{enumerate}

Recommended Progression:

\begin{itemize}
\tightlist
\item
  Lessons 6-7: PLA (simplest)
\item
  Lessons 8-9: PLA or PETG (functional parts)
\item
  Lesson 10: PETG or TPU (testing materials)
\item
  Lesson 11: Student choice (depends on stakeholder requirements)
\end{itemize}

\begin{center}\rule{0.5\linewidth}{0.5pt}\end{center}

\protect\phantomsection\label{docs__pandoc__latex__src__3dmake_appendices_complete.md__appendix-d}{}

\section*{Appendix D: Tolerance Testing and Quality Assurance Matrix}\label{docs__pandoc__latex__src__3dmake_appendices_complete.md__appendix-d-tolerance-testing-and-quality-assurance-matrix}

\begin{center}\rule{0.5\linewidth}{0.5pt}\end{center}

This appendix provides measurement-based testing methodology for verifying that 3D-printed parts meet design specifications. It\textquotesingle s designed to be used non-visually-focusing on calipers, scales, and functional tests rather than visual inspection.

Referenced in: Lessons 8-10 (Complex Design, Troubleshooting, Mastery)

\subsection*{Overview: What is Tolerance?}\label{docs__pandoc__latex__src__3dmake_appendices_complete.md__overview-what-is-tolerance}

Tolerance is the acceptable range of variation in dimensions:

\begin{lstlisting}[style=Alabaster]
Design spec:  Hole diameter = 6mm
Tolerance:    +/-0.5mm
Acceptable range:  5.5mm to 6.5mm
Actual print:  5.8mm [YES] (within tolerance)
or            7.2mm [NO] (exceeds tolerance)

\end{lstlisting}

\subsubsection*{Why Tolerance Matters}\label{docs__pandoc__latex__src__3dmake_appendices_complete.md__why-tolerance-matters}

\begin{enumerate}
\tightlist
\item
  Assembly: Parts must fit together
\item
  Function: Fit too tight = stuck; too loose = falls apart
\item
  Safety: Wrong tolerance = part failure
\item
  Cost: Tight tolerance = slower, more waste
\end{enumerate}

\subsection*{Essential Measurement Tools}\label{docs__pandoc__latex__src__3dmake_appendices_complete.md__essential-measurement-tools}

\subsubsection*{1. Digital Calipers}\label{docs__pandoc__latex__src__3dmake_appendices_complete.md__1-digital-calipers}

What they measure:

\begin{itemize}
\tightlist
\item
  Outside diameter (part width)
\item
  Inside diameter (hole width)
\item
  Depth
\end{itemize}

How to use non-visually:

\begin{enumerate}
\tightlist
\item
  Gently close calipers until they barely touch the part
\item
  Feel the resistance (should be light, not forced)
\item
  Read digital display with audio feedback or manually
\item
  Record three measurements at different locations
\item
  Average the three measurements
\end{enumerate}

Accuracy: +/-0.05mm (very precise for 3D printing)

\subsubsection*{2. Digital Scale (Kitchen Scale)}\label{docs__pandoc__latex__src__3dmake_appendices_complete.md__2-digital-scale-kitchen-scale}

What it measures:

\begin{itemize}
\tightlist
\item
  Part weight (indicates infill, material type)
\end{itemize}

How to use non-visually:

\begin{enumerate}
\tightlist
\item
  Place part on scale
\item
  Wait for reading to stabilize (1-2 seconds)
\item
  Read digital display (in grams)
\item
  Compare to expected weight
\end{enumerate}

Accuracy: +/-1g (good enough for verification)

Why it matters:

\begin{itemize}
\tightlist
\item
  Too light = infill too low or void inside part
\item
  Expected weight = indicates proper slicing
\end{itemize}

\subsubsection*{3. Test Jig / Go/No-Go Gauge}\label{docs__pandoc__latex__src__3dmake_appendices_complete.md__3-test-jig--gono-go-gauge}

What it measures:

\begin{itemize}
\tightlist
\item
  Pass/fail tolerance testing without calipers
\end{itemize}

How to make:

\begin{lstlisting}[style=Alabaster, language=openscad]
// Go/No-Go gauge for bracelet peg holes
// Tests if hole is within acceptable range
pegdiameter = 6;        // Design spec
tolerance = 0.5;         // Tolerance
godiameter = pegdiameter - tolerance;     // Min acceptable (5.5)
nogodiameter = pegdiameter + tolerance;  // Max acceptable (6.5)
module gogauge() {
  // Part should fit through this easily
  cylinder(h=10, r=godiameter/2);
}
module nogogauge() {
  // Part should NOT fit through this
  cylinder(h=10, r=nogodiameter/2);
}
// Test jig with both gauges
union() {
  translate([0, 0, 0]) gogauge();
  translate([0, 15, 0]) nogogauge();
}

\end{lstlisting}

Use non-visually:

\begin{itemize}
\tightlist
\item
  Try inserting peg into go-gauge -\textgreater{} Should slide easily
\item
  Try inserting peg into no-go-gauge -\textgreater{} Should NOT fit
\item
  If both tests pass = tolerance correct
\end{itemize}

\subsection*{Quality Assurance Testing Matrix}\label{docs__pandoc__latex__src__3dmake_appendices_complete.md__quality-assurance-testing-matrix}

\subsubsection*{Critical Dimensions to Test}\label{docs__pandoc__latex__src__3dmake_appendices_complete.md__critical-dimensions-to-test}

Create a Test Plan before printing:

{\def\LTcaptype{none} % do not increment counter
\begin{longtable}[]{@{}
  >{\raggedright\arraybackslash}p{(\linewidth - 8\tabcolsep) * \real{0.1809}}
  >{\raggedright\arraybackslash}p{(\linewidth - 8\tabcolsep) * \real{0.1915}}
  >{\raggedright\arraybackslash}p{(\linewidth - 8\tabcolsep) * \real{0.0745}}
  >{\raggedright\arraybackslash}p{(\linewidth - 8\tabcolsep) * \real{0.1170}}
  >{\raggedright\arraybackslash}p{(\linewidth - 8\tabcolsep) * \real{0.4362}}@{}}
\toprule\noalign{}
\begin{minipage}[b]{\linewidth}\raggedright
Part
\end{minipage} & \begin{minipage}[b]{\linewidth}\raggedright
Dimension
\end{minipage} & \begin{minipage}[b]{\linewidth}\raggedright
Spec
\end{minipage} & \begin{minipage}[b]{\linewidth}\raggedright
Tolerance
\end{minipage} & \begin{minipage}[b]{\linewidth}\raggedright
How to Test
\end{minipage} \\
\midrule\noalign{}
\endhead
\bottomrule\noalign{}
\endlastfoot
Bracelet Holder & Base width & 127mm & +/-2mm &
Measure with calipers (multiple points) \\
& Peg diameter & 6mm & +/-0.5mm & Test fit with go/no-go gauge \\
& Peg spacing & 8mm & +/-1mm & Measure distance between pegs \\
& Back wall height & 120mm & +/-2mm & Measure with calipers \\
Phone Stand & Slope angle & 20 & +/-3 &
Calculate from height/depth ratio \\
& Weight capacity & 200g & 150-250g & Load test (see below) \\
& Stability & N/A & Pass/fail & 1-hour load test without tipping \\
\end{longtable}
}

\subsubsection*{Pre-Print Planning}\label{docs__pandoc__latex__src__3dmake_appendices_complete.md__pre-print-planning}

Before slicing, define:

\begin{enumerate}
\tightlist
\item
  Critical dimensions: Which measurements matter most?
\item
  Acceptable range: What tolerance is realistic?
\item
  Test method: How will you verify?
\item
  Pass/fail criteria: What does "success" look like?
\end{enumerate}

Example Plan for Phone Stand:

\begin{lstlisting}[style=Alabaster]
# Phone Stand - Quality Assurance Plan

## Critical Dimensions
1. Slope angle: 20 +/-3
2. Base width: 80mm +/-1mm
3. Stand height: 60mm +/-2mm

## Functional Tests
1. Stability: Hold 200g for 1 hour without tipping
2. Grip: Phone doesn't slide during tilt
3. Assembly: Back brace attaches without force

## Pass/Fail Criteria
[YES] PASS if:
- All dimensions within tolerance
- Phone holds weight for 1 hour
- No cracks or layer separation

[NO] FAIL if:
- Any dimension >2mm off spec
- Phone slides or part tips
- Visible cracks

\end{lstlisting}

\subsection*{Measurement Procedures}\label{docs__pandoc__latex__src__3dmake_appendices_complete.md__measurement-procedures}

\subsubsection*{Procedure 1: Linear Dimension (Width, Height, Depth)}\label{docs__pandoc__latex__src__3dmake_appendices_complete.md__procedure-1-linear-dimension-width-height-depth}

Equipment needed: Digital calipers

Steps:

\begin{enumerate}
\tightlist
\item
  Place part on flat surface
\item
  Position caliper jaws perpendicular to surface
\item
  Gently close jaws until they just touch part
\item
  Feel for light resistance (not forced)
\item
  Read digital display
\item
  Record measurement
\item
  Repeat at 3 different locations
\item
  Average the three readings
\item
  Compare to design spec +/- tolerance
\end{enumerate}

Example:

\begin{lstlisting}[style=Alabaster]
Design spec: 127mm (bracelet holder width)
Tolerance:   +/-2mm (acceptable: 125-129mm)
Measurements:
  Location 1: 126.8mm
  Location 2: 127.1mm
  Location 3: 126.5mm
  Average: 126.8mm [YES] (within tolerance)

\end{lstlisting}

\subsubsection*{Procedure 2: Hole or Peg Diameter}\label{docs__pandoc__latex__src__3dmake_appendices_complete.md__procedure-2-hole-or-peg-diameter}

Equipment needed: Digital calipers, Go/No-Go gauges (optional)

Method A: Direct Measurement

\begin{enumerate}
\tightlist
\item
  Insert caliper jaws into hole/around peg
\item
  Adjust jaws to gently touch surfaces
\item
  Feel for light resistance on both sides
\item
  Read digital display (inside measurement mode)
\item
  Record three measurements
\item
  Average the readings
\end{enumerate}

Method B: Go/No-Go Gauge

\begin{enumerate}
\tightlist
\item
  Print test jigs (go \& no-go gauges)
\item
  Attempt to insert peg/hole into go-gauge
  -\textgreater{} Should slide through easily with light resistance
\item
  Attempt to insert peg/hole into no-go-gauge
  -\textgreater{} Should NOT fit or fit with visible resistance
\item
  If both pass -\textgreater{} dimension is acceptable
\end{enumerate}

Example: Peg Diameter

Design spec: 6.0mm Tolerance:   +/-0.5mm (acceptable: 5.5-6.5mm)
Method A (Direct): Measurement 1: 5.9mm Measurement 2: 6.0mm
Measurement 3: 5.8mm Average: 5.9mm {[}YES{]} Method B (Go/No-Go):
Slides through go-gauge (5.5mm) -\textgreater{} {[}YES{]}
Doesn\textquotesingle t fit no-go-gauge (6.5mm) -\textgreater{} {[}YES{]}
Result: PASS

\subsubsection*{Procedure 3: Surface Finish / Layer Quality}\label{docs__pandoc__latex__src__3dmake_appendices_complete.md__procedure-3-surface-finish--layer-quality}

Equipment needed: Caliper, ruler, touch/texture assessment

Step 1: Surface Texture

Run fingers/hand over surface:
{[}YES{]} Smooth -\textgreater{} Good quality
Slightly rough -\textgreater{} Acceptable
{[}NO{]} Very rough/bumpy -\textgreater{} Quality issue

Step 2: Layer Line Visibility

Feel horizontal ridges (layer lines):
{[}YES{]} Barely perceptible -\textgreater{} Good (0.2mm layers)
Noticeable but even -\textgreater{} OK (0.25mm layers)
{[}NO{]} Very pronounced -\textgreater{} Quality issue

Step 3: Dimensional Consistency

\begin{lstlisting}[style=Alabaster]
Measure thickness at multiple points:
  Design spec: 3mm wall thickness
  Measure at 5 locations
  Record all measurements
  All should be within +/-0.2mm of each other
  Example:
    Pt 1: 3.0mm
    Pt 2: 3.1mm
    Pt 3: 3.0mm
    Pt 4: 2.9mm
    Pt 5: 3.1mm
  Std Dev: 0.08mm [YES] (very consistent)

\end{lstlisting}

\subsubsection*{Procedure 4: Weight Verification (Confirms Infill)}\label{docs__pandoc__latex__src__3dmake_appendices_complete.md__procedure-4-weight-verification-confirms-infill}

Equipment needed: Digital scale

Step 1: Calculate Expected Weight

Expected weight = Volume x Density x Infill\% For bracelet holder (PLA):
Volume = 127mm x 80mm x 120mm = 1,219,200 mm = 1,219.2 cm
PLA density = 1.24 g/cm Infill = 20\%
Expected weight = 1,219.2 x 1.24 x 0.20 = 302.4g
Acceptable range: 290-315g (+/-5\%)

Step 2: Measure Actual Weight

\begin{enumerate}
\tightlist
\item
  Place part on digital scale
\item
  Wait for reading to stabilize (1-2 seconds)
\item
  Read display in grams
\item
  Compare to expected weight
\end{enumerate}

Step 3: Interpret Result

{\def\LTcaptype{none} % do not increment counter
\begin{longtable}[]{@{}ll@{}}
\toprule\noalign{}
Actual Weight & Interpretation \\
\midrule\noalign{}
\endhead
\bottomrule\noalign{}
\endlastfoot
290-315g & {[}YES{]} Correct infill, no internal voids \\
\textless{}280g & Infill too low or significant voids \\
\textgreater320g & Infill too high (was it supposed to be 20\%?) \\
\end{longtable}
}

\subsection*{Functional Testing}\label{docs__pandoc__latex__src__3dmake_appendices_complete.md__functional-testing}

\subsubsection*{Test 1: Load Testing (Strength)}\label{docs__pandoc__latex__src__3dmake_appendices_complete.md__test-1-load-testing-strength}

Purpose: Verify part can hold weight without failure

Equipment needed:

\begin{itemize}
\tightlist
\item
  Digital scale
\item
  Test weights (or books, water jugs)
\item
  Calipers (to check for deflection)
\end{itemize}

Procedure:

\begin{enumerate}
\tightlist
\item
  Measure baseline dimensions (part unloaded) Baseline height: 60mm
\item
  Place test weight on part Added weight: 200g
\item
  Wait 5 minutes (allow part to settle)
\item
  Measure dimensions again New height: 59.8mm
  Deflection: 0.2mm (acceptable)
\item
  Observe for cracks (visual or tactile) No cracks: {[}YES{]}
\item
  Remove weight and wait 5 minutes
\item
  Measure dimensions again (should return to baseline)
  Height after unload: 60.0mm Recovery: Complete {[}YES{]}
  Result: PASS (part handles load without permanent deformation)
\end{enumerate}

Acceptance Criteria for Phone Stand:

{[}YES{]} PASS if:

\begin{itemize}
\tightlist
\item
  Deflection \textless{}0.5mm under 200g load
\item
  No cracks visible/felt
\item
  Full recovery after load removed {[}NO{]} FAIL if:
\item
  Deflection \textgreater1mm
\item
  Visible cracks
\item
  Permanent deformation after unload
\end{itemize}

\subsubsection*{Test 2: Assembly Testing}\label{docs__pandoc__latex__src__3dmake_appendices_complete.md__test-2-assembly-testing}

Purpose: Verify parts fit together as designed

Equipment needed:

\begin{itemize}
\tightlist
\item
  Calipers
\item
  Go/No-Go gauges
\end{itemize}

Procedure for Multi-Part Assembly (e.g., Stackable Bins):

Part A (Bin body) dimensions:

\begin{itemize}
\tightlist
\item
  Top opening: 50mm x 50mm +/-1mm
\item
  Wall thickness: 2mm +/-0.2mm Part B (Stacking rim) dimensions:
\item
  Base diameter: 50mm +/-0.5mm
\item
  Should nest inside Part A Test procedure:
\end{itemize}

\begin{enumerate}
\tightlist
\item
  Measure Part A opening: 50.1mm {[}YES{]}
\item
  Measure Part B base: 49.8mm {[}YES{]}
\item
  Attempt to nest Part B into Part A
  -\textgreater{} Should fit with light resistance
  -\textgreater{} Feel for smooth insertion (no catching)
\item
  Check for rocking (part should sit stable)
\item
  Apply 500g load (simulate stacking)
  -\textgreater{} Should hold without slipping
\end{enumerate}

\subsubsection*{Test 3: Durability Testing (Repeated Use)}\label{docs__pandoc__latex__src__3dmake_appendices_complete.md__test-3-durability-testing-repeated-use}

Purpose: Verify part doesn\textquotesingle t fail after repeated cycles

Example: Phone Stand Tilt Cycles

\begin{enumerate}
\tightlist
\item
  Record baseline dimensions
\item
  Cycle 1: Place phone, tilt to max angle, remove
\item
  Inspect for cracks or damage
\item
  Repeat cycle 10 times
\item
  Measure dimensions again
\item
  Compare to baseline: Should be \textless{}0.1mm change
  If no damage after 10 cycles -\textgreater{} Durable {[}YES{]}
\end{enumerate}

\subsection*{Tolerance Stack-Up (Multi-Part Designs)}\label{docs__pandoc__latex__src__3dmake_appendices_complete.md__tolerance-stack-up-multi-part-designs}

When multiple parts are assembled, tolerances add:

Design: Part A opening: 50mm +/-0.5mm Part B base: 50mm +/-0.5mm
Worst-case fit: Part A minimum: 49.5mm Part B maximum: 50.5mm
Difference: 1.0mm (VERY TIGHT or won\textquotesingle t fit!)
Solution: Increase tolerance on Part B to +/-0.3mm
Part B minimum: 49.7mm Part B maximum: 50.3mm
Now fits inside Part A: 49.5-50.5mm {[}YES{]}

\subsubsection*{Tolerance Stack-Up Calculation}\label{docs__pandoc__latex__src__3dmake_appendices_complete.md__tolerance-stack-up-calculation}

For N parts in assembly:

\begin{lstlisting}[style=Alabaster]
Total tolerance = (tolerance + tolerance + ... + tolerance)
Example with 3 parts (each +/-0.5mm):
Total = (0.5 + 0.5 + 0.5)
      = 0.75
      = 0.87mm

\end{lstlisting}

\subsection*{Common Dimensional Problems \& Fixes}\label{docs__pandoc__latex__src__3dmake_appendices_complete.md__common-dimensional-problems--fixes}

{\def\LTcaptype{none} % do not increment counter
\begin{longtable}[]{@{}
  >{\raggedright\arraybackslash}p{(\linewidth - 6\tabcolsep) * \real{0.2450}}
  >{\raggedright\arraybackslash}p{(\linewidth - 6\tabcolsep) * \real{0.1788}}
  >{\raggedright\arraybackslash}p{(\linewidth - 6\tabcolsep) * \real{0.3311}}
  >{\raggedright\arraybackslash}p{(\linewidth - 6\tabcolsep) * \real{0.2450}}@{}}
\toprule\noalign{}
\begin{minipage}[b]{\linewidth}\raggedright
Problem
\end{minipage} & \begin{minipage}[b]{\linewidth}\raggedright
Typical Cause
\end{minipage} & \begin{minipage}[b]{\linewidth}\raggedright
How to Fix
\end{minipage} & \begin{minipage}[b]{\linewidth}\raggedright
Prevent Next Time
\end{minipage} \\
\midrule\noalign{}
\endhead
\bottomrule\noalign{}
\endlastfoot
Part too small (all dimensions off) & Scale wrong in slicer &
Scale STL up in CAD or slicer & Verify scale before slicing \\
Holes too small & Compensation shrinkage & Increase hole diameter +0.5mm
& Add shrinkage factor to design \\
Walls too thin & Layer squishing &
Check first layer height; may need recalibration &
Measure first layer thickness \\
Infill showing through (loose fill) & Infill too low &
Increase infill \% to 25-30\% & Recalculate for part type \\
Inconsistent across print & Nozzle clogging mid-print &
Clean nozzle; check filament quality &
Use quality filament; monitor print \\
Z-axis dimensions off & Z-axis uncalibrated &
Run Z-calibration procedure & Calibrate before critical prints \\
Dimensions change between prints & Thermal drift &
Print in stable temperature & Use enclosed printer if possible \\
\end{longtable}
}

\subsection*{Testing Checklist Template}\label{docs__pandoc__latex__src__3dmake_appendices_complete.md__testing-checklist-template}

Print this checklist before starting a new project:

\begin{lstlisting}[style=Alabaster]
# Quality Assurance Checklist: [Project Name]

## Design Specs
- [ ] Part A width:  mm +/-  mm
- [ ] Part B height:  mm +/-  mm
- [ ] Hole diameter:  mm +/-  mm
- [ ] Assembly fit tolerance:  mm

## Pre-Print
- [ ] STL exported correctly
- [ ] Scale verified
- [ ] Supports configured
- [ ] Slice file reviewed

## Post-Print (Dimensional)
- [ ] Part A width measured (3 points): mm, mm, mm
  - [ ] Within tolerance? YES/NO
- [ ] Part B height measured (3 points): mm, mm, mm
  - [ ] Within tolerance? YES/NO
- [ ] Hole diameter measured (go/no-go): PASS/FAIL
- [ ] Part weight measured: g (expected: g +/-g)
  - [ ] Within expected weight? YES/NO

## Post-Print (Functional)
- [ ] Load test (if applicable): PASS/FAIL
- [ ] Assembly test: PASS/FAIL
- [ ] Surface quality acceptable: YES/NO
- [ ] No cracks or voids: YES/NO

## Result
- [ ] ACCEPT (all tests pass)
- [ ] REWORK (minor issue, can fix)
- [ ] REJECT (major issue, reprint)

## Notes
[Space for observations]

\end{lstlisting}

\subsection*{PowerShell Integration: Track Quality Metrics}\label{docs__pandoc__latex__src__3dmake_appendices_complete.md__powershell-integration-track-quality-metrics}

\begin{lstlisting}[style=Alabaster, language=powershell]
# Track print quality across multiple projects
$qualityLog = @"
ProjectName,Date,Material,PartCount,DimensionsPass,FunctionalPass,Weight(g),Notes
"@
# Log entries
$qualityLog += "`nBraceletHolder,2024-01-15,PLA,1,PASS,PASS,302.1,Perfect fit"
$qualityLog += "`nPhoneStand,2024-01-16,PETG,2,PASS,PASS,487.3,Strong joints"
$qualityLog += "`nStackableBins,2024-01-17,PLA,3,FAIL,N/A,N/A,Holes too small reprint"
# Save to CSV
$qualityLog | Out-File "C:\Projects\quality-log.csv"
# Analyze pass/fail rate
$log = Import-Csv "C:\Projects\quality-log.csv"
$passCount = ($log | Where-Object { $.DimensionsPass -eq "PASS" }).Count
$totalCount = $log.Count
$passRate = ($passCount / $totalCount) * 100
Write-Host "Print success rate: $passRate%"
# Show recent failures
Write-Host "`nRecent issues:"
$log | Where-Object { $.DimensionsPass -eq "FAIL" } | Select-Object ProjectName, Notes

\end{lstlisting}

\subsection*{Accessibility-Focused QA Best Practices}\label{docs__pandoc__latex__src__3dmake_appendices_complete.md__accessibility-focused-qa-best-practices}

\subsubsection*{Why Measurement-Based Testing Matters}\label{docs__pandoc__latex__src__3dmake_appendices_complete.md__why-measurement-based-testing-matters}

Traditional visual QA (looking at surface finish, checking dimensions by eye) is inherently inaccessible. This appendix prioritizes:

\begin{enumerate}
\tightlist
\item
  Caliper measurements: Objective, quantifiable
\item
  Weight verification: Numerical result
\item
  Functional testing: Pass/fail criteria
\item
  Go/No-Go gauges: Tactile pass/fail
\item
  Test jigs: Can be shared/standardized
\end{enumerate}

\subsubsection*{Screen Reader Integration}\label{docs__pandoc__latex__src__3dmake_appendices_complete.md__screen-reader-integration}

When documenting QA results:

\begin{lstlisting}[style=Alabaster]
[YES] GOOD: "Bracket width measured 49.8mm (spec 50+/-1mm)"
[NO] AVOID: "Bracket looks good" (not measurable)
[YES] GOOD: "Part weighs 298g (spec 300+/-10g)"
[NO] AVOID: "Feels about right" (subjective)
[YES] GOOD: "Slides through go-gauge, blocks no-go-gauge"
[NO] AVOID: "Hole looks the right size" (visual only)

\end{lstlisting}

\subsubsection*{Documentation Template (Accessible)}\label{docs__pandoc__latex__src__3dmake_appendices_complete.md__documentation-template-accessible}

\begin{lstlisting}[style=Alabaster]
## Part: Bracelet Holder Peg (Test Date: Jan 15, 2024)

### Dimensional Verification
| Dimension    | Specification  | Measurement 1 | Measurement 2 | Measurement 3 | Average | Result     |
|--------------|----------------|---------------|---------------|---------------|---------|------------|
| Peg Diameter | 6.0 +/- 0.5mm  | 5.9mm         | 6.0mm         | 5.8mm         | 5.9mm   | [YES] PASS |
| Peg Length   | 25.0 +/- 0.5mm | 25.1mm        | 25.0mm        | 25.0mm        | 25.03mm | [YES] PASS |

### Functional Test
- [x] Slides into bracelet loop easily (no forcing)
- [x] No cracks visible or felt
- [x] Survives 1-hour load test with 20 bracelets
- [x] No permanent deformation after load removed

### Overall Result
[YES] PASS - All dimensions within tolerance; functionally sound

\end{lstlisting}

\subsection*{Summary}\label{docs__pandoc__latex__src__3dmake_appendices_complete.md__summary-2}

Key Principles:

\begin{enumerate}
\tightlist
\item
  Plan before printing: Define tolerances and test methods
\item
  Measure precisely: Use calibrated tools (digital calipers, scale)
\item
  Test functionally: Will parts assemble and work?
\item
  Document quantitatively: Use numbers, not adjectives
\item
  Iterate based on data: If test fails, adjust design or process
\end{enumerate}

When to Use Each Test:

{\def\LTcaptype{none} % do not increment counter
\begin{longtable}[]{@{}ll@{}}
\toprule\noalign{}
Situation & Recommended Test \\
\midrule\noalign{}
\endhead
\bottomrule\noalign{}
\endlastfoot
Prototype/mockup & Linear dimensions only \\
Assembly with other parts & Tolerance stack-up analysis \\
Load-bearing part & Functional load test \\
Repeated-use item & Durability cycling test \\
Multi-part design & Assembly fit test \\
\end{longtable}
}

Accessibility Reminder:

All QA testing in this appendix is measurement-based and non-visual, making it equally accessible to all users. The goal is objective, quantifiable verification-not subjective assessment.

\begin{center}\rule{0.5\linewidth}{0.5pt}\end{center}

\protect\phantomsection\label{docs__pandoc__latex__src__3dmake_appendices_complete.md__appendix-e}{}

\section*{Appendix E: Advanced OpenSCAD Concepts}\label{docs__pandoc__latex__src__3dmake_appendices_complete.md__appendix-e-advanced-openscad-concepts}

\begin{center}\rule{0.5\linewidth}{0.5pt}\end{center}

This appendix covers specialized topics for experienced users seeking to tackle complex parametric designs. These are optional topics not required for foundational mastery but valuable for professional application development.

\subsection*{Topic 1: Gears and Mechanical Components}\label{docs__pandoc__latex__src__3dmake_appendices_complete.md__topic-1-gears-and-mechanical-components}

Gears are one of the most challenging parametric designs, requiring careful calculation of tooth geometry. Understanding gear mathematics enables creation of mechanical systems-power transmission, speed reduction, or precise positioning.

\subsubsection*{Gear Terminology}\label{docs__pandoc__latex__src__3dmake_appendices_complete.md__gear-terminology}

\begin{itemize}
\tightlist
\item
  Pitch Diameter (PD): The reference diameter for meshing calculations
\item
  Module: PD divided by tooth count; determines tooth size
\item
  Pressure Angle: Typically 20deg or 14.5deg; affects tooth shape and strength
\item
  Clearance: Space between teeth to allow smooth meshing
\item
  Backlash: Intentional gap to prevent binding at tolerance extremes
\end{itemize}

\subsubsection*{Simple Involute Gear Algorithm}\label{docs__pandoc__latex__src__3dmake_appendices_complete.md__simple-involute-gear-algorithm}

\begin{lstlisting}[style=Alabaster, language=openscad]
// Simplified gear with circular teeth (adequate for most 3D prints)
function gear_module(pitch_diameter, teeth_count) = pitch_diameter / teeth_count;
module simple_gear(pitch_diameter, teeth_count, bore_diameter, thickness) {
  module_m = gear_module(pitch_diameter, teeth_count);
  outer_diameter = pitch_diameter + 2 * module_m;
  tooth_angle = 360 / teeth_count;
  difference() {
    // Main gear body
    cylinder(r=outer_diameter/2, h=thickness, $fn=teeth_count*4);
    // Center bore
    cylinder(r=bore_diameter/2, h=thickness + 2);
  }
  // Add teeth as small protrusions
  for (i = [0:teeth_count-1]) {
    angle = i * tooth_angle;
    translate([
      (pitch_diameter/2 + module_m/2) * cos(angle),
      (pitch_diameter/2 + module_m/2) * sin(angle),
      0
    ])
      rotate([0, 0, angle])
        cube([module_m * 0.8, module_m, thickness], center=true);
  }
}
// Create meshing gear pair
module gear_pair_demo() {
  // Gear 1: 20 teeth, 40mm pitch diameter
  simple_gear(40, 20, 8, 10);
  // Gear 2: 30 teeth, positioned to mesh
  center_distance = (40 + 60) / 4;  // Sum of radii
  translate([center_distance, 0, 0])
    simple_gear(60, 30, 8, 10);
}
gear_pair_demo();

\end{lstlisting}

\subsubsection*{Practical Application: Servo Gearbox}\label{docs__pandoc__latex__src__3dmake_appendices_complete.md__practical-application-servo-gearbox}

\begin{lstlisting}[style=Alabaster, language=openscad]
// Parametric servo gearbox with 3:1 reduction
module servo_gearbox(motor_torque, reduction_ratio, bore_size) {
  // Input gear (motor)
  input_teeth = 12;
  input_pd = 30;
  // Output gear (load)
  output_teeth = input_teeth * reduction_ratio;
  output_pd = input_pd * reduction_ratio;
  // Gearbox housing
  housing_size = output_pd + 40;
  difference() {
    // Main box
    cube([housing_size, housing_size, 30]);
    // Interior chamber
    translate([20, 20, 5])
      cube([housing_size - 40, housing_size - 40, 25]);
  }
  // Mount input gear
  translate([housing_size/4, housing_size/2, 15]) {
    simple_gear(input_pd, input_teeth, bore_size, 15);
    // Motor coupling
    cylinder(r=bore_size/2, h=5);
  }
  // Mount output gear
  translate([3*housing_size/4, housing_size/2, 15])
    simple_gear(output_pd, output_teeth, bore_size, 15);
}
servo_gearbox(10, 3, 5);

\end{lstlisting}

\subsubsection*{Belt and Pulley Systems}\label{docs__pandoc__latex__src__3dmake_appendices_complete.md__belt-and-pulley-systems}

\begin{lstlisting}[style=Alabaster, language=openscad]
// Timing pulley with tooth grooves
module timing_pulley(bore_diameter, pitch_diameter, teeth_count, width) {
  module_m = pitch_diameter / teeth_count;
  tooth_height = module_m * 0.3;
  difference() {
    // Main pulley body
    cylinder(r=pitch_diameter/2 + tooth_height/2, h=width);
    // Center bore
    cylinder(r=bore_diameter/2, h=width + 2);
    // Tooth grooves
    for (i = [0:teeth_count-1]) {
      angle = i * (360 / teeth_count);
      rotate([0, 0, angle])
        translate([pitch_diameter/2, -module_m*0.3, 0])
          cube([module_m * 0.6, module_m * 0.6, width], center=true);
    }
  }
}
// Belt drive demonstration
module belt_drive_system() {
  // Motor pulley (input)
  timing_pulley(5, 20, 16, 10);
  // Load pulley (output) - positioned for belt routing
  translate([80, 0, 0])
    timing_pulley(5, 40, 32, 10);
}
belt_drive_system();

\end{lstlisting}

\subsection*{Topic 2: Batch Processing and Statistical Analysis}\label{docs__pandoc__latex__src__3dmake_appendices_complete.md__topic-2-batch-processing-and-statistical-analysis}

Professional applications often generate many variants and need to analyze results. Batch processing enables systematic exploration of parameter spaces.

\subsubsection*{Parameter Sweep Generation}\label{docs__pandoc__latex__src__3dmake_appendices_complete.md__parameter-sweep-generation}

\begin{lstlisting}[style=Alabaster, language=openscad]
// Generate models with all combinations of parameters
// Define parameter ranges
wall_thicknesses = [1, 1.5, 2, 2.5, 3];
part_sizes = [20, 30, 40, 50];
materials = ["pla", "petg", "nylon"];
// Theory: Generate 5 x 4 x 3 = 60 variants automatically
// In practice, export each to separate STL for analysis

\end{lstlisting}

\subsubsection*{Statistical Summary Generation Script}\label{docs__pandoc__latex__src__3dmake_appendices_complete.md__statistical-summary-generation-script}

\begin{lstlisting}[style=Alabaster, language=cmd]
@echo off
rem analyze_batch.bat - Analyze batch results (CMD / Batch)
setlocal enableextensions enabledelayedexpansion
set "OUTPUT_DIR=batch_results"
set "REPORT_FILE=batch_analysis.csv"
echo Part,Wall,Size,File_Size_KB,Est_Print_Time_min,Est_Weight_g > "%REPORT_FILE%"

for %%w in (1 2 3) do (
  for %%s in (20 30 40 50) do (
    set "SCAD_FILE=part_%%s_wall%%w.scad"
    set "STL_FILE=%OUTPUT_DIR%\%%~nSCAD_FILE.stl"
    "C:\Program Files\OpenSCAD\openscad.exe" -D "part_size=%%s" -D "wall=%%w" -o "%OUTPUT_DIR%\part_%%s_wall%%w.stl" src\main.scad
    if exist "%OUTPUT_DIR%\part_%%s_wall%%w.stl" (
      for /f "usebackq" %%F in (`powershell -NoProfile -Command "(Get-Item -Path '%OUTPUT_DIR%\part_%%s_wall%%w.stl').Length"`) do set FILE_BYTES=%%F
      set /a FILE_KB=FILE_BYTES / 1024
      set /a EST_TIME=(FILE_KB / 1024) * 60
      set /a EST_WEIGHT=FILE_KB / 1024
    ) else (
      set FILE_KB=0
      set EST_TIME=0
      set EST_WEIGHT=0
    )
    echo part_%%s_wall%%w,%%w,%%s,!FILE_KB!,!EST_TIME!,!EST_WEIGHT!>>"%REPORT_FILE%"
  )
)
endlocal
echo Analysis complete. Results in %REPORT_FILE%

\end{lstlisting}

\begin{lstlisting}[style=Alabaster, language=powershell]
# analyze_batch.ps1 - Analyze batch results (PowerShell)
$OutputDir = "batch_results"
$ReportFile = "batch_analysis.csv"
"Part,Wall,Size,File_Size_KB,Est_Print_Time_min,Est_Weight_g" | Out-File -FilePath $ReportFile -Encoding utf8

foreach ($wall in @(1.0,1.5,2.0,2.5,3.0)) {
  foreach ($size in @(20,30,40,50)) {
    $scad = "part_${size}_wall${wall}.scad"
    $stl = Join-Path $OutputDir ("{0}.stl" -f ([System.IO.Path]::GetFileNameWithoutExtension($scad)))
    & "C:\Program Files\OpenSCAD\openscad.exe" -D ("part_size={0}" -f $size) -D ("wall={0}" -f $wall) -o $stl src\main.scad
    if (Test-Path $stl) {
      $bytes = (Get-Item $stl).Length
      $kb = [math]::Round($bytes / 1KB, 2)
      $est_time = [math]::Round(($kb / 1024) * 60, 2)
      $est_weight = [math]::Round($kb / 1024, 2)
    } else {
      $kb = 0; $est_time = 0; $est_weight = 0
    }
    "{0},{1},{2},{3},{4},{5}" -f ("part_{0}_wall{1}" -f $size,$wall), $wall, $size, $kb, $est_time, $est_weight | Out-File -FilePath $ReportFile -Append -Encoding utf8
  }
}
Write-Host "Analysis complete. Results in $ReportFile"

\end{lstlisting}

\begin{lstlisting}[style=Alabaster, language=bash]
#!/bin/bash
# analyze_batch.sh - Analyze batch results statistically

OUTPUT_DIR="batch_results"
REPORT_FILE="batch_analysis.csv"

echo "Part,Wall,Size,File_Size_KB,Est_Print_Time_min,Est_Weight_g" > "$REPORT_FILE"

for wall in 1.0 1.5 2.0 2.5 3.0; do
  for size in 20 30 40 50; do
    # Generate model
    SCAD_FILE="part_${size}_wall${wall}.scad"
    STL_FILE="$OUTPUT_DIR/${SCAD_FILE%.scad}.stl"
    
    # Export to STL
    openscad -D "part_size=$size" -D "wall=$wall" \
      -o "$STL_FILE" src/main.scad
    
    # Get file size
    FILE_SIZE=$(stat -c%s "$STL_FILE" 2>/dev/null | awk '{print $1/1024}')
    
    # Estimate print time (rough: 1 hour per MB)
    EST_TIME=$((FILE_SIZE / 1024 * 60))
    
    # Estimate weight (rough: 1g per MB for PLA)
    EST_WEIGHT=$((FILE_SIZE / 1024))
    
    # Log results
    echo "part_${size}_wall${wall},$wall,$size,$FILE_SIZE,$EST_TIME,$EST_WEIGHT" >> "$REPORT_FILE"
  done
done

echo "Analysis complete. Results in $REPORT_FILE"

\end{lstlisting}

\subsubsection*{Data-Driven Design Selection}\label{docs__pandoc__latex__src__3dmake_appendices_complete.md__data-driven-design-selection}

\begin{lstlisting}[style=Alabaster, language=openscad]
// Load results and select optimal configuration
// wall=2.0mm gives best balance of strength/weight for most sizes
function optimal_configuration(target_size, priority) =
  priority == "strength" ? 3.0 :
  priority == "speed" ? 1.0 :
  priority == "balanced" ? 2.0 :
  2.0;  // Default
module batch_optimized_part(size, priority) {
  wall = optimal_configuration(size, priority);
  echo(str("Generating optimized part: size=", size, " wall=", wall, " priority=", priority));
  // Create part with optimal parameters
  difference() {
    cube([size, size, size]);
    translate([wall, wall, wall])
      cube([size - 2*wall, size - 2*wall, size]);
  }
}
// Generate three variants with different priorities
batch_optimized_part(40, "strength");    // thickest walls
translate([60, 0, 0]) batch_optimized_part(40, "speed");    // thinnest walls
translate([120, 0, 0]) batch_optimized_part(40, "balanced"); // medium walls

\end{lstlisting}

\subsection*{Topic 3: Performance Optimization}\label{docs__pandoc__latex__src__3dmake_appendices_complete.md__topic-3-performance-optimization}

Complex models can take hours to render. Strategic optimization keeps iteration cycles fast.

\subsubsection*{Measuring Render Performance}\label{docs__pandoc__latex__src__3dmake_appendices_complete.md__measuring-render-performance}

\begin{lstlisting}[style=Alabaster, language=bash]
# Time how long rendering takes
time openscad -o output.stl -D "quality=32" src/main.scad
# Output: real 5m 32s, user 5m 28s, sys 0m 2s

\end{lstlisting}

\subsubsection*{Resolution Parameter Strategy}\label{docs__pandoc__latex__src__3dmake_appendices_complete.md__resolution-parameter-strategy}

\begin{lstlisting}[style=Alabaster, language=openscad]
// Use lower resolution during development, high during export
quality = $preview ? 16 : 64;  // 16 segments in preview, 64 in export
module efficient_sphere(radius) {
  sphere(r=radius, $fn=quality);
}
// In array operations, efficiency matters most
module circular_array(radius, count) {
  for (i = [0:count-1]) {
    angle = (360 / count) * i;
    x = radius * cos(angle);
    y = radius * sin(angle);
    translate([x, y, 0])
      efficient_sphere(5);  // Uses quality parameter
  }
}
circular_array(30, 100);  // 100 spheres: slow with quality=64!

\end{lstlisting}

\subsubsection*{Cache Complex Calculations}\label{docs__pandoc__latex__src__3dmake_appendices_complete.md__cache-complex-calculations}

\begin{lstlisting}[style=Alabaster, language=openscad]
// Pre-compute expensive calculations outside loops
// SLOW: Recalculates sin/cos for every iteration
module slow_spiral() {
  for (i = [0:200]) {
    angle = i * 2;
    x = 40 * cos(angle);
    y = 40 * sin(angle);
    z = i * 0.5;
    translate([x, y, z])
      sphere(r=1);
  }
}
// FAST: Pre-compute positions, then use them
spiral_positions = [for (i = [0:200]) [
  40 * cos(i * 2),
  40 * sin(i * 2),
  i * 0.5
]];
module fast_spiral() {
  for (pos = spiral_positions)
    translate(pos)
      sphere(r=1);
}
// Use the fast version
fast_spiral();

\end{lstlisting}

\subsubsection*{Simplification During Preview}\label{docs__pandoc__latex__src__3dmake_appendices_complete.md__simplification-during-preview}

\begin{lstlisting}[style=Alabaster, language=openscad]
// Complex model that simplifies in preview mode
module detailed_model() {
  if ($preview) {
    // Simplified version for fast preview
    cube([100, 100, 50]);
  } else {
    // Detailed version for export
    difference() {
      cube([100, 100, 50]);
      // Complex inner geometry
      for (x = [10:10:90])
        for (y = [10:10:90])
          translate([x, y, 20])
            cylinder(h=15, d=3, $fn=32);
    }
  }
}
detailed_model();

\end{lstlisting}

\subsubsection*{Profile-Driven Optimization}\label{docs__pandoc__latex__src__3dmake_appendices_complete.md__profile-driven-optimization}

\begin{lstlisting}[style=Alabaster, language=bash]
#!/bin/bash
# profile_render.sh - Identify slow rendering hotspots

echo "Profiling render performance..."

for fn_value in 8 16 32 64 128; do
  echo ""
  echo "Testing with \$fn=$fn_value..."
  
  /usr/bin/time -v openscad -D "fn=$fn_value" \
    -o test_$fn_value.stl src/main.scad 2>&1 | \
    grep "Maximum resident set size"
done

\end{lstlisting}

\subsection*{Topic 4: Print Orientation and Support Structure Algorithms}\label{docs__pandoc__latex__src__3dmake_appendices_complete.md__topic-4-print-orientation-and-support-structure-algorithms}

Strategic orientation dramatically affects support material, print time, and surface quality.

\subsubsection*{Strength Orientation Analysis}\label{docs__pandoc__latex__src__3dmake_appendices_complete.md__strength-orientation-analysis}

\begin{lstlisting}[style=Alabaster, language=openscad]
// Different orientations affect strength differently
// SCENARIO: Bracket with cantilever load
module bracket() {
  difference() {
    cube([50, 100, 10]);
    translate([25, 50, 5]) cylinder(r=5, h=10);
  }
}
// Orientation A: Lay flat (XY plane) - WEAK for cantilever
echo("Orientation A: Flat - Low strength for cantilever");
bracket();
// Orientation B: Stand tall (Z axis) - STRONG for cantilever
translate([70, 0, 0]) {
  echo("Orientation B: Tall - High strength for cantilever");
  rotate([90, 0, 0]) bracket();
}
// Orientation C: Diagonal - MODERATE strength
translate([0, 150, 0]) {
  echo("Orientation C: Diagonal - Moderate strength");
  rotate([45, 0, 0]) bracket();
}

\end{lstlisting}

\subsubsection*{Support Material Minimization}\label{docs__pandoc__latex__src__3dmake_appendices_complete.md__support-material-minimization}

\begin{lstlisting}[style=Alabaster, language=openscad]
// Design features to reduce supports needed
// POOR: Overhanging feature needs extensive supports
module poor_design() {
  cube([50, 50, 50]);
  translate([25, 25, 50])
    cube([30, 30, 10]);  // 30mm overhang!
}
// GOOD: Bridge angle keeps overhang under 45deg
module good_design() {
  cube([50, 50, 50]);
  // Add ramp instead of overhang
  translate([20, 20, 50])
    rotate([20, 0, 0])
      cube([30, 30, 10]);
}
// Compare
poor_design();
translate([80, 0, 0])
  good_design();

\end{lstlisting}

\subsubsection*{Bridge Span Calculation}\label{docs__pandoc__latex__src__3dmake_appendices_complete.md__bridge-span-calculation}

\begin{lstlisting}[style=Alabaster, language=openscad]
// Calculate maximum unsupported span for given material/settings
function max_bridge_span(material) =
  material == "pla" ? 10 :
  material == "petg" ? 12 :
  material == "nylon" ? 15 :
  material == "abs" ? 8 :
  10;  // Default
module bridged_connector(material, gap_width) {
  max_span = max_bridge_span(material);
  if (gap_width <= max_span) {
    echo(str("Bridge OK: gap ", gap_width, "mm fits within ", max_span, "mm limit"));
    // Create part with bridge
    union() {
      cube([50, 10, 10]);  // Side A
      translate([gap_width, 0, 0])
        cube([50, 10, 10]);  // Side B
      // Bridge connecting them
      translate([0, 5, 10])
        cube([gap_width + 100, 2, 1]);
    }
  } else {
    echo(str("ERROR: gap ", gap_width, "mm exceeds ", max_span, "mm max bridge span"));
  }
}
bridged_connector("pla", 8);      // OK: 8 < 10
// bridged_connector("pla", 15);  // ERROR: 15 > 10

\end{lstlisting}

\subsubsection*{Slicing Parameter Calculation}\label{docs__pandoc__latex__src__3dmake_appendices_complete.md__slicing-parameter-calculation}

\begin{lstlisting}[style=Alabaster, language=openscad]
// Calculate optimal slicing parameters based on geometry
function recommended_layer_height(nozzle_diameter, quality_priority) =
  quality_priority == "fast" ? nozzle_diameter * 0.75 :
  quality_priority == "balanced" ? nozzle_diameter * 0.5 :
  quality_priority == "detailed" ? nozzle_diameter * 0.25 :
  nozzle_diameter * 0.5;
function recommended_infill(structural_load, safety_factor) =
  structural_load < 1 ? 10 :    // Decorative: minimal fill
  structural_load < 5 ? 20 :    // Light duty: low fill
  structural_load < 20 ? 50 :   // Medium duty: standard fill
  100;                           // High duty: solid
module optimized_part(
  nozzle_0p4,
  quality_priority,
  structural_load,
  safety_factor
) {
  layer_h = recommended_layer_height(nozzle_0p4, quality_priority);
  infill_pct = recommended_infill(structural_load, safety_factor);
  echo(str("Recommended layer height: ", layer_h, "mm"));
  echo(str("Recommended infill: ", infill_pct, "%"));
  // Create part
  cube([50, 50, 30]);
}
optimized_part(0.4, "detailed", 15, 2);

\end{lstlisting}

\subsection*{Topic 5: Recursive Function Patterns}\label{docs__pandoc__latex__src__3dmake_appendices_complete.md__topic-5-recursive-function-patterns}

Recursion enables elegant solutions to problems with self-similar structure: trees, fractals, nested components.

\subsubsection*{Basic Recursive Pattern}\label{docs__pandoc__latex__src__3dmake_appendices_complete.md__basic-recursive-pattern}

\begin{lstlisting}[style=Alabaster, language=openscad]
// Tree structure: branches recursively smaller
function tree_depth(level) = level > 0 ? level + tree_depth(level - 1) : 0;
// Calculate: tree_depth(5) = 5 + 4 + 3 + 2 + 1 = 15
result = tree_depth(5);
echo(result);  // Prints: 15
// Recursive tree drawing
module tree_branch(length, angle, recursion_depth) {
  if (recursion_depth > 0) {
    // Draw this branch
    rotate([angle, 0, 0])
      cube([2, 2, length]);
    // Recursively draw sub-branches
    translate([0, 0, length])
      rotate([-angle/2, 0, 0])
        tree_branch(length * 0.7, angle, recursion_depth - 1);
    translate([0, 0, length])
      rotate([angle/2, 0, 0])
        tree_branch(length * 0.7, angle, recursion_depth - 1);
  }
}
tree_branch(30, 25, 4);  // Tree with 4 levels of recursion

\end{lstlisting}

\subsubsection*{Fractal Generation}\label{docs__pandoc__latex__src__3dmake_appendices_complete.md__fractal-generation}

\begin{lstlisting}[style=Alabaster, language=openscad]
// Sierpinski Triangle fractal
function sierpinski_triangle(size, depth) =
  depth == 0 ?
    [[0, 0], [size, 0], [size/2, size * sqrt(3)/2]] :
    // Recursive case: three smaller triangles at corners
    concat(
      sierpinski_triangle(size/2, depth - 1),
      sierpinski_triangle(size/2, depth - 1),
      sierpinski_triangle(size/2, depth - 1)
    );
// Create 3D fractal structure
module fractal_spiral(base_size, depth, height_per_level) {
  if (depth > 0) {
    // Create current level
    cube([base_size, base_size, 5], center=true);
    // Recursively create smaller level on top
    translate([0, 0, height_per_level])
      scale([0.6, 0.6, 1])
        fractal_spiral(base_size, depth - 1, height_per_level);
  }
}
fractal_spiral(40, 5, 10);  // 5-level spiral

\end{lstlisting}

\subsubsection*{Nested Component Assembly}\label{docs__pandoc__latex__src__3dmake_appendices_complete.md__nested-component-assembly}

\begin{lstlisting}[style=Alabaster, language=openscad]
// Russian doll boxes: each box contains a smaller copy
module russian_doll(size, wall, nesting_depth) {
  if (nesting_depth > 0) {
    // Current level: hollow box
    difference() {
      cube([size, size, size]);
      translate([wall, wall, wall])
        cube([size - 2*wall, size - 2*wall, size - wall]);
    }
    // Recursively place next smaller doll inside
    interior_size = size - 4*wall;
    translate([2*wall, 2*wall, wall])
      russian_doll(interior_size, wall, nesting_depth - 1);
  } else {
    // Smallest doll (solid)
    cube([size, size, size]);
  }
}
russian_doll(60, 2, 4);  // 4 nested boxes

\end{lstlisting}

\subsubsection*{Performance Considerations}\label{docs__pandoc__latex__src__3dmake_appendices_complete.md__performance-considerations}

\begin{lstlisting}[style=Alabaster, language=openscad]
// Recursion can be expensive - monitor depth
// EFFICIENT: Tail recursion (result computed immediately)
function sum_to(n, accumulator) =
  n == 0 ? accumulator :
  sum_to(n - 1, accumulator + n);
result = sum_to(100, 0);  // Computes: 5050
// INEFFICIENT: Deep recursion without optimization
function fibonacci(n) =
  n <= 1 ? n :
  fibonacci(n-1) + fibonacci(n-2);
// fibonacci(20) recalculates sub-problems thousands of times
// Don't use for depth > 20 without memoization
// BETTER: Use iteration where possible
function fibonacci_iter(n) =
  let(
    fib = [for (i = [0:n])
      i <= 1 ? i :
      let(prev = [for (j = [0:i-1]) i==j ? 1 : i==j+1 ? 0 : 0])
        1  // Placeholder
    ]
  ) fib[n];

\end{lstlisting}

\subsubsection*{Practical Example: Cable Management}\label{docs__pandoc__latex__src__3dmake_appendices_complete.md__practical-example-cable-management}

\begin{lstlisting}[style=Alabaster, language=openscad]
// Recursive cable tray with branching paths
module cable_tray(width, depth, height, levels, branch_angle) {
  if (levels > 0) {
    // Main tray section
    difference() {
      cube([width, depth, height]);
      translate([2, 2, 2])
        cube([width - 4, depth - 4, height - 2]);
    }
    // Left branch (recursively smaller)
    translate([-width/2, 0, height])
      rotate([0, 0, -branch_angle])
        cable_tray(width * 0.6, depth, height, levels - 1, branch_angle);
    // Right branch
    translate([width/2, 0, height])
      rotate([0, 0, branch_angle])
        cable_tray(width * 0.6, depth, height, levels - 1, branch_angle);
  }
}
cable_tray(40, 10, 5, 3, 30);  // 3-level branching cable tray

\end{lstlisting}

\subsection*{Summary: When to Use Each Advanced Technique}\label{docs__pandoc__latex__src__3dmake_appendices_complete.md__summary-when-to-use-each-advanced-technique}

{\def\LTcaptype{none} % do not increment counter
\begin{longtable}[]{@{}
  >{\raggedright\arraybackslash}p{(\linewidth - 6\tabcolsep) * \real{0.2989}}
  >{\raggedright\arraybackslash}p{(\linewidth - 6\tabcolsep) * \real{0.3563}}
  >{\raggedright\arraybackslash}p{(\linewidth - 6\tabcolsep) * \real{0.1379}}
  >{\raggedright\arraybackslash}p{(\linewidth - 6\tabcolsep) * \real{0.2069}}@{}}
\toprule\noalign{}
\begin{minipage}[b]{\linewidth}\raggedright
Technique
\end{minipage} & \begin{minipage}[b]{\linewidth}\raggedright
Use Case
\end{minipage} & \begin{minipage}[b]{\linewidth}\raggedright
Complexity
\end{minipage} & \begin{minipage}[b]{\linewidth}\raggedright
Performance
\end{minipage} \\
\midrule\noalign{}
\endhead
\bottomrule\noalign{}
\endlastfoot
Gears & Mechanical power transmission & High & Medium \\
Batch Processing & Design space exploration & Medium &
Depends on scope \\
Performance Optimization & Reducing render time & Medium &
High payoff \\
Orientation Analysis & Strength optimization & Medium &
Simulation cost \\
Recursion & Fractal/nested structures & High & Can be expensive \\
\end{longtable}
}

\subsection*{References and Further Learning}\label{docs__pandoc__latex__src__3dmake_appendices_complete.md__references-and-further-learning}

\begin{itemize}
\tightlist
\item
  OpenSCAD Documentation - Advanced Features: \url{https://en.wikibooks.org/wiki/OpenSCAD_User_Manual}
\item
  Gear Design Theory: \url{https://www.mekanizmalar.com/}
\item
  3D Print Orientation Optimization: \url{https://stratasys.com/}
\item
  Fractal Geometry: \url{https://en.wikipedia.org/wiki/Fractal}
\item
  Recursive Algorithms: \url{https://www.khanacademy.org/computing/computer-science/algorithms}
\end{itemize}

For Educators and Students

These advanced topics are intended for:

\begin{itemize}
\tightlist
\item
  Experienced users tackling professional applications
\item
  Specialized problem domains (robotics, mechanical design, etc.)
\item
  Research and optimization work
\item
  Custom library and framework development
\end{itemize}

For accessibility:

\begin{itemize}
\tightlist
\item
  All recursive examples include base cases clearly marked
\item
  Each section provides simplified versions before advanced variants
\item
  Code comments explain both "what" and "why"
\item
  Practical examples show real-world applications
\item
  Performance considerations documented to prevent frustration
\end{itemize}

\section{Back Matter --- Complete Reference}\label{docs__pandoc__latex__src__backmatter_complete.md__back-matter--complete-reference}

Supplementary reference material for \emph{Accessible 3D Printing with OpenSCAD and 3dMake}. These sections are designed for lookup rather than linear reading --- consult them as needed throughout the curriculum.

\begin{center}\rule{0.5\linewidth}{0.5pt}\end{center}

\subsection*{Contents}\label{docs__pandoc__latex__src__backmatter_complete.md__contents}

{\def\LTcaptype{none} % do not increment counter
\begin{longtable}[]{@{}
  >{\raggedright\arraybackslash}p{(\linewidth - 4\tabcolsep) * \real{0.3333}}
  >{\raggedright\arraybackslash}p{(\linewidth - 4\tabcolsep) * \real{0.3333}}
  >{\raggedright\arraybackslash}p{(\linewidth - 4\tabcolsep) * \real{0.3333}}@{}}
\toprule\noalign{}
\begin{minipage}[b]{\linewidth}\raggedright
Section
\end{minipage} & \begin{minipage}[b]{\linewidth}\raggedright
Title
\end{minipage} & \begin{minipage}[b]{\linewidth}\raggedright
Purpose
\end{minipage} \\
\midrule\noalign{}
\endhead
\bottomrule\noalign{}
\endlastfoot
\hyperref[docs__pandoc__latex__src__backmatter_complete.md__troubleshooting]{Technical Reference and Troubleshooting}
&
Systems architecture, CLI reference, and fault-by-fault troubleshooting compendium
&
Diagnose and resolve errors in OpenSCAD, 3dMake, CMD, PowerShell, Git Bash, slicers, and printers \\
\hyperref[docs__pandoc__latex__src__backmatter_complete.md__student-glossary]{Student Glossary}
& Terms, definitions, and code examples across all curriculum systems &
Look up any term encountered in lessons, with syntax examples for each shell \\
\hyperref[docs__pandoc__latex__src__backmatter_complete.md__teacher-glossary]{Teacher Glossary}
&
Pedagogical frameworks, instructional strategies, and assessment rubrics
&
Instructors: curriculum design, scaffolding, computational thinking, and cross-curricular connections \\
\hyperref[docs__pandoc__latex__src__backmatter_complete.md__further-reading]{Further Reading and References}
&
Scholarly works, official documentation, industry resources, and community links
& Continue learning; cite sources; find community support \\
\end{longtable}
}

\begin{center}\rule{0.5\linewidth}{0.5pt}\end{center}

\begin{center}\rule{0.5\linewidth}{0.5pt}\end{center}

\protect\phantomsection\label{docs__pandoc__latex__src__backmatter_complete.md__troubleshooting}{}

\section*{Technical Reference and Troubleshooting Guide}\label{docs__pandoc__latex__src__backmatter_complete.md__technical-reference-and-troubleshooting-guide}

\begin{center}\rule{0.5\linewidth}{0.5pt}\end{center}

\section*{PART I --- SYSTEMS ARCHITECTURE}\label{docs__pandoc__latex__src__backmatter_complete.md__part-i--systems-architecture}

\subsection*{1.1 Digital Fabrication Pipeline}\label{docs__pandoc__latex__src__backmatter_complete.md__11-digital-fabrication-pipeline}

\begin{lstlisting}[style=Alabaster]
Source (.scad)
    ↓
OpenSCAD Engine
    ↓
STL (Mesh Geometry)
    ↓
Slicer
    ↓
G-code
    ↓
FDM Printer Firmware
    ↓
Physical Object

\end{lstlisting}

Automation Overlay:

\begin{lstlisting}[style=Alabaster]
CLI -> Script -> Build Directory -> Version Control -> Reproducible Artifact

\end{lstlisting}

\begin{center}\rule{0.5\linewidth}{0.5pt}\end{center}

\section*{PART II --- OPENSCAD PROFESSIONAL MODELING REFERENCE}\label{docs__pandoc__latex__src__backmatter_complete.md__part-ii--openscad-professional-modeling-reference}

\subsection*{2.1 Language Architecture}\label{docs__pandoc__latex__src__backmatter_complete.md__21-language-architecture}

OpenSCAD is:

\begin{itemize}
\tightlist
\item
  Declarative
\item
  CSG-based
\item
  Deterministic
\item
  Single-pass evaluation
\item
  Non-interactive
\end{itemize}

Implication: Models must be architected intentionally.

\begin{center}\rule{0.5\linewidth}{0.5pt}\end{center}

\subsection*{2.2 Advanced Geometry Techniques}\label{docs__pandoc__latex__src__backmatter_complete.md__22-advanced-geometry-techniques}

\subsubsection*{2.2.1 Transform Stack Order}\label{docs__pandoc__latex__src__backmatter_complete.md__221-transform-stack-order}

Transformations apply from inside outward.

\begin{lstlisting}[style=Alabaster, language=openscad]
translate([10,0,0])
    rotate([0,0,45])
        cube(10);

\end{lstlisting}

Order matters.

\begin{center}\rule{0.5\linewidth}{0.5pt}\end{center}

\subsubsection*{2.2.2 Hull()}\label{docs__pandoc__latex__src__backmatter_complete.md__222-hull}

Creates convex hull of objects.

\begin{lstlisting}[style=Alabaster, language=openscad]
hull() {
    translate([0,0,0]) sphere(5);
    translate([20,0,0]) sphere(5);
}

\end{lstlisting}

\begin{center}\rule{0.5\linewidth}{0.5pt}\end{center}

\subsubsection*{2.2.3 Minkowski()}\label{docs__pandoc__latex__src__backmatter_complete.md__223-minkowski}

Expands object by shape kernel.

\begin{lstlisting}[style=Alabaster, language=openscad]
minkowski() {
    cube([10,10,10]);
    sphere(1);
}

\end{lstlisting}

Warning: Computationally expensive.

\begin{center}\rule{0.5\linewidth}{0.5pt}\end{center}

\subsubsection*{2.2.4 Linear and Rotational Extrusion}\label{docs__pandoc__latex__src__backmatter_complete.md__224-linear-and-rotational-extrusion}

\begin{lstlisting}[style=Alabaster, language=openscad]
linear_extrude(height=10)
    square(5);

rotate_extrude()
    translate([10,0,0])
        circle(5);

\end{lstlisting}

\begin{center}\rule{0.5\linewidth}{0.5pt}\end{center}

\subsubsection*{2.2.5 Projection}\label{docs__pandoc__latex__src__backmatter_complete.md__225-projection}

\begin{lstlisting}[style=Alabaster, language=openscad]
projection(cut=true)
    sphere(10);

\end{lstlisting}

\begin{center}\rule{0.5\linewidth}{0.5pt}\end{center}

\subsection*{2.3 CLI Operation of OpenSCAD}\label{docs__pandoc__latex__src__backmatter_complete.md__23-cli-operation-of-openscad}

\subsubsection*{Render STL (All shells)}\label{docs__pandoc__latex__src__backmatter_complete.md__render-stl-all-shells}

\begin{lstlisting}[style=Alabaster, language=bash]
openscad -o output.stl input.scad

\end{lstlisting}

\subsubsection*{Specify Variables via CLI}\label{docs__pandoc__latex__src__backmatter_complete.md__specify-variables-via-cli}

\begin{lstlisting}[style=Alabaster, language=bash]
openscad -D width=50 -o model.stl model.scad

\end{lstlisting}

\begin{center}\rule{0.5\linewidth}{0.5pt}\end{center}

\section*{PART III --- 3DMAKE PROFESSIONAL USAGE}\label{docs__pandoc__latex__src__backmatter_complete.md__part-iii--3dmake-professional-usage}

\subsection*{3.1 3dMake Overview}\label{docs__pandoc__latex__src__backmatter_complete.md__31-3dmake-overview}

3dMake automates model generation and batch workflows.

Core capabilities:

\begin{itemize}
\tightlist
\item
  Batch STL generation
\item
  File processing pipelines
\item
  Integration with OpenSCAD CLI
\item
  Script-driven builds
\end{itemize}

Repository: https://github.com/tdeck/3dmake

\begin{center}\rule{0.5\linewidth}{0.5pt}\end{center}

\subsection*{3.2 Example Automated Pipeline (PowerShell)}\label{docs__pandoc__latex__src__backmatter_complete.md__32-example-automated-pipeline-powershell}

\begin{lstlisting}[style=Alabaster, language=powershell]
New-Item -ItemType Directory -Force build
openscad -D width=40 -o build/model.stl src/model.scad
if ($LASTEXITCODE -ne 0) { exit 1 }

\end{lstlisting}

\begin{center}\rule{0.5\linewidth}{0.5pt}\end{center}

\subsection*{3.3 Batch Processing (cmd.exe)}\label{docs__pandoc__latex__src__backmatter_complete.md__33-batch-processing-cmdexe}

\begin{lstlisting}[style=Alabaster, language=cmd]
for %%f in (src\*.scad) do openscad -o build\%%~nf.stl %%f

\end{lstlisting}

\begin{center}\rule{0.5\linewidth}{0.5pt}\end{center}

\subsection*{3.4 Batch Processing (Git Bash)}\label{docs__pandoc__latex__src__backmatter_complete.md__34-batch-processing-git-bash}

\begin{lstlisting}[style=Alabaster, language=bash]
for f in src/*.scad; do
  openscad -o build/$(basename "$f" .scad).stl "$f"
done

\end{lstlisting}

\begin{center}\rule{0.5\linewidth}{0.5pt}\end{center}

\section*{PART IV --- WINDOWS COMMAND PROMPT (CMD.EXE) FULL REFERENCE}\label{docs__pandoc__latex__src__backmatter_complete.md__part-iv--windows-command-prompt-cmdexe-full-reference}

\subsection*{4.1 cmd Architecture}\label{docs__pandoc__latex__src__backmatter_complete.md__41-cmd-architecture}

\begin{itemize}
\tightlist
\item
  Text-based
\item
  Legacy DOS-compatible syntax
\item
  String-based (not object-based)
\end{itemize}

Executable:

\begin{lstlisting}[style=Alabaster]
C:\Windows\System32\cmd.exe

\end{lstlisting}

\begin{center}\rule{0.5\linewidth}{0.5pt}\end{center}

\subsection*{4.2 File System Operations}\label{docs__pandoc__latex__src__backmatter_complete.md__42-file-system-operations}

\subsubsection*{List Files}\label{docs__pandoc__latex__src__backmatter_complete.md__list-files}

\begin{lstlisting}[style=Alabaster, language=cmd]
dir

\end{lstlisting}

\subsubsection*{Change Directory}\label{docs__pandoc__latex__src__backmatter_complete.md__change-directory}

\begin{lstlisting}[style=Alabaster, language=cmd]
cd builds

\end{lstlisting}

\subsubsection*{Create Directory}\label{docs__pandoc__latex__src__backmatter_complete.md__create-directory}

\begin{lstlisting}[style=Alabaster, language=cmd]
mkdir build

\end{lstlisting}

\subsubsection*{Remove Directory}\label{docs__pandoc__latex__src__backmatter_complete.md__remove-directory}

\begin{lstlisting}[style=Alabaster, language=cmd]
rmdir /s /q build

\end{lstlisting}

\begin{center}\rule{0.5\linewidth}{0.5pt}\end{center}

\subsection*{4.3 Running Programs}\label{docs__pandoc__latex__src__backmatter_complete.md__43-running-programs}

\begin{lstlisting}[style=Alabaster, language=cmd]
openscad.exe -o model.stl model.scad

\end{lstlisting}

\begin{center}\rule{0.5\linewidth}{0.5pt}\end{center}

\subsection*{4.4 PATH Inspection}\label{docs__pandoc__latex__src__backmatter_complete.md__44-path-inspection}

\begin{lstlisting}[style=Alabaster, language=cmd]
echo %PATH%

\end{lstlisting}

\begin{center}\rule{0.5\linewidth}{0.5pt}\end{center}

\subsection*{4.5 Temporary Environment Variables}\label{docs__pandoc__latex__src__backmatter_complete.md__45-temporary-environment-variables}

\begin{lstlisting}[style=Alabaster, language=cmd]
set MYVAR=value

\end{lstlisting}

\begin{center}\rule{0.5\linewidth}{0.5pt}\end{center}

\subsection*{4.6 Exit Codes}\label{docs__pandoc__latex__src__backmatter_complete.md__46-exit-codes}

\begin{lstlisting}[style=Alabaster, language=cmd]
echo %ERRORLEVEL%

\end{lstlisting}

Conditional:

\begin{lstlisting}[style=Alabaster, language=cmd]
if %ERRORLEVEL% NEQ 0 echo Build failed

\end{lstlisting}

\begin{center}\rule{0.5\linewidth}{0.5pt}\end{center}

\subsection*{4.7 Batch Files (.bat)}\label{docs__pandoc__latex__src__backmatter_complete.md__47-batch-files-bat}

Example build.bat:

\begin{lstlisting}[style=Alabaster, language=cmd]
@echo off
mkdir build
openscad -o build\model.stl src\model.scad
if %ERRORLEVEL% NEQ 0 exit /b 1

\end{lstlisting}

\begin{center}\rule{0.5\linewidth}{0.5pt}\end{center}

\section*{PART V --- POWERSHELL PROFESSIONAL REFERENCE}\label{docs__pandoc__latex__src__backmatter_complete.md__part-v--powershell-professional-reference}

\subsection*{5.1 Architecture}\label{docs__pandoc__latex__src__backmatter_complete.md__51-architecture}

\begin{itemize}
\tightlist
\item
  Object-based pipeline
\item
  .NET-backed
\item
  Strong scripting support
\end{itemize}

\begin{center}\rule{0.5\linewidth}{0.5pt}\end{center}

\subsection*{5.2 Execution Policy}\label{docs__pandoc__latex__src__backmatter_complete.md__52-execution-policy}

\begin{lstlisting}[style=Alabaster, language=powershell]
Get-ExecutionPolicy
Set-ExecutionPolicy RemoteSigned -Scope CurrentUser

\end{lstlisting}

\begin{center}\rule{0.5\linewidth}{0.5pt}\end{center}

\subsection*{5.3 Structured Error Handling}\label{docs__pandoc__latex__src__backmatter_complete.md__53-structured-error-handling}

\begin{lstlisting}[style=Alabaster, language=powershell]
try {
    openscad -o model.stl model.scad
}
catch {
    Write-Host "Failure"
}

\end{lstlisting}

\begin{center}\rule{0.5\linewidth}{0.5pt}\end{center}

\section*{PART VI --- GIT BASH (UNIX-LIKE SHELL)}\label{docs__pandoc__latex__src__backmatter_complete.md__part-vi--git-bash-unix-like-shell}

\subsection*{6.1 POSIX-Compatible Commands}\label{docs__pandoc__latex__src__backmatter_complete.md__61-posix-compatible-commands}

\begin{lstlisting}[style=Alabaster, language=bash]
ls
pwd
mkdir
rm -rf build

\end{lstlisting}

\begin{center}\rule{0.5\linewidth}{0.5pt}\end{center}

\subsection*{6.2 Shell Variables}\label{docs__pandoc__latex__src__backmatter_complete.md__62-shell-variables}

\begin{lstlisting}[style=Alabaster, language=bash]
VAR=value
echo $VAR

\end{lstlisting}

\begin{center}\rule{0.5\linewidth}{0.5pt}\end{center}

\section*{PART VII --- ADVANCED TROUBLESHOOTING COMPENDIUM}\label{docs__pandoc__latex__src__backmatter_complete.md__part-vii--advanced-troubleshooting-compendium}

\begin{center}\rule{0.5\linewidth}{0.5pt}\end{center}

\section*{A. OpenSCAD Failures}\label{docs__pandoc__latex__src__backmatter_complete.md__a-openscad-failures}

\subsection*{A.1 Segmentation Fault}\label{docs__pandoc__latex__src__backmatter_complete.md__a1-segmentation-fault}

Cause:

\begin{itemize}
\tightlist
\item
  Deep Minkowski
\item
  Memory exhaustion
\end{itemize}

Fix:

\begin{itemize}
\tightlist
\item
  Reduce \$fn
\item
  Simplify geometry
\end{itemize}

\begin{center}\rule{0.5\linewidth}{0.5pt}\end{center}

\subsection*{A.2 Empty STL Output}\label{docs__pandoc__latex__src__backmatter_complete.md__a2-empty-stl-output}

Cause:

\begin{itemize}
\tightlist
\item
  Geometry evaluates to null
\end{itemize}

Diagnosis: Add temporary primitive to confirm render pipeline.

\begin{center}\rule{0.5\linewidth}{0.5pt}\end{center}

\subsection*{A.3 Boolean Artifacts}\label{docs__pandoc__latex__src__backmatter_complete.md__a3-boolean-artifacts}

Cause:

\begin{itemize}
\tightlist
\item
  Coplanar faces
\end{itemize}

Fix:

\begin{itemize}
\tightlist
\item
  Add small offsets (0.01mm)
\end{itemize}

\begin{center}\rule{0.5\linewidth}{0.5pt}\end{center}

\subsection*{A.4 Syntax Errors}\label{docs__pandoc__latex__src__backmatter_complete.md__a4-syntax-errors}

\subsubsection*{Missing Semicolons}\label{docs__pandoc__latex__src__backmatter_complete.md__missing-semicolons}

Cause:

\begin{lstlisting}[style=Alabaster, language=openscad]
cube([10,10,10])  // Missing semicolon

\end{lstlisting}

Fix:

\begin{lstlisting}[style=Alabaster, language=openscad]
cube([10,10,10]);

\end{lstlisting}

\begin{center}\rule{0.5\linewidth}{0.5pt}\end{center}

\subsubsection*{Unmatched Brackets}\label{docs__pandoc__latex__src__backmatter_complete.md__unmatched-brackets}

Common errors:

\begin{itemize}
\tightlist
\item
  \texttt{\{\ \}} curly brackets
\item
  \texttt{{[}\ {]}} square brackets
\item
  \texttt{(\ )} parentheses
\end{itemize}

Diagnosis: Check error line number in console.

\begin{center}\rule{0.5\linewidth}{0.5pt}\end{center}

\subsubsection*{Comment Syntax Errors}\label{docs__pandoc__latex__src__backmatter_complete.md__comment-syntax-errors}

Invalid multiline comment:

\begin{lstlisting}[style=Alabaster, language=openscad]
/* Comment
// Still in comment
*/

\end{lstlisting}

Single-line comments use \texttt{//}
Multiline comments use \texttt{/*\ */}

\begin{center}\rule{0.5\linewidth}{0.5pt}\end{center}

\subsection*{A.5 Logical Operator Errors}\label{docs__pandoc__latex__src__backmatter_complete.md__a5-logical-operator-errors}

\subsubsection*{\&\& (AND) Usage}\label{docs__pandoc__latex__src__backmatter_complete.md__-and-usage}

\begin{lstlisting}[style=Alabaster, language=openscad]
if (x > 5 && y < 10) {
    cube([10,10,10]);
}

\end{lstlisting}

\subsubsection*{\textbar\textbar{} (OR) Usage}\label{docs__pandoc__latex__src__backmatter_complete.md__-or-usage}

\begin{lstlisting}[style=Alabaster, language=openscad]
if (x > 5 || y < 10) {
    cube([10,10,10]);
}

\end{lstlisting}

Missing parentheses can cause evaluation errors.

\begin{center}\rule{0.5\linewidth}{0.5pt}\end{center}

\subsection*{A.6 \$fn Performance Issues}\label{docs__pandoc__latex__src__backmatter_complete.md__a6-fn-performance-issues}

Cause:

\begin{itemize}
\tightlist
\item
  Excessive fragment count
\end{itemize}

Example:

\begin{lstlisting}[style=Alabaster, language=openscad]
$fn = 500;  // Too high
sphere(5);

\end{lstlisting}

Fix:

\begin{lstlisting}[style=Alabaster, language=openscad]
$fn = 50;   // Sufficient for most cases
sphere(5);

\end{lstlisting}

\begin{center}\rule{0.5\linewidth}{0.5pt}\end{center}

\section*{B. 3dMake Failures}\label{docs__pandoc__latex__src__backmatter_complete.md__b-3dmake-failures}

\subsection*{B.1 Command Not Found}\label{docs__pandoc__latex__src__backmatter_complete.md__b1-command-not-found}

Verify:

\begin{lstlisting}[style=Alabaster, language=cmd]
where openscad

\end{lstlisting}

Or check 3dm:

\begin{lstlisting}[style=Alabaster, language=cmd]
where 3dm

\end{lstlisting}

\begin{center}\rule{0.5\linewidth}{0.5pt}\end{center}

\subsection*{B.2 Silent Failure}\label{docs__pandoc__latex__src__backmatter_complete.md__b2-silent-failure}

Check exit code:

\begin{lstlisting}[style=Alabaster, language=cmd]
echo %ERRORLEVEL%

\end{lstlisting}

\begin{center}\rule{0.5\linewidth}{0.5pt}\end{center}

\subsection*{B.3 3dm build Failures}\label{docs__pandoc__latex__src__backmatter_complete.md__b3-3dm-build-failures}

Common causes:

\begin{itemize}
\tightlist
\item
  Source file not found
\item
  Syntax errors in .scad
\item
  OpenSCAD not in PATH
\end{itemize}

Diagnosis:

\begin{lstlisting}[style=Alabaster, language=bash]
3dm build --verbose

\end{lstlisting}

\begin{center}\rule{0.5\linewidth}{0.5pt}\end{center}

\subsection*{B.4 3dm info Not Working}\label{docs__pandoc__latex__src__backmatter_complete.md__b4-3dm-info-not-working}

Requirements:

\begin{itemize}
\tightlist
\item
  Valid STL file
\item
  Network connection (for AI service)
\end{itemize}

Check:

\begin{lstlisting}[style=Alabaster, language=bash]
3dm info

\end{lstlisting}

If info works but describe doesn\textquotesingle t, check network.

\begin{center}\rule{0.5\linewidth}{0.5pt}\end{center}

\section*{C. Command Line Lesson Specific Failures}\label{docs__pandoc__latex__src__backmatter_complete.md__c-command-line-lesson-specific-failures}

\subsection*{C.1 "Access is denied"}\label{docs__pandoc__latex__src__backmatter_complete.md__c1-access-is-denied}

Cause:

\begin{itemize}
\tightlist
\item
  Insufficient privileges
\end{itemize}

Fix: Run as Administrator.

\begin{center}\rule{0.5\linewidth}{0.5pt}\end{center}

\subsection*{C.2 Incorrect Variable Expansion in Loop}\label{docs__pandoc__latex__src__backmatter_complete.md__c2-incorrect-variable-expansion-in-loop}

In batch files use:

\begin{lstlisting}[style=Alabaster]
%%f

\end{lstlisting}

In interactive shell use:

\begin{lstlisting}[style=Alabaster]
%f

\end{lstlisting}

\begin{center}\rule{0.5\linewidth}{0.5pt}\end{center}

\section*{D. PowerShell-Specific Failures}\label{docs__pandoc__latex__src__backmatter_complete.md__d-powershell-specific-failures}

\subsection*{D.1 Script Not Running}\label{docs__pandoc__latex__src__backmatter_complete.md__d1-script-not-running}

Fix:

\begin{lstlisting}[style=Alabaster, language=powershell]
Set-ExecutionPolicy RemoteSigned

\end{lstlisting}

\begin{center}\rule{0.5\linewidth}{0.5pt}\end{center}

\subsection*{D.2 Command Returns But File Not Created}\label{docs__pandoc__latex__src__backmatter_complete.md__d2-command-returns-but-file-not-created}

Check:

\begin{lstlisting}[style=Alabaster, language=powershell]
Test-Path build\model.stl

\end{lstlisting}

\begin{center}\rule{0.5\linewidth}{0.5pt}\end{center}

\section*{E. Git Bash Issues}\label{docs__pandoc__latex__src__backmatter_complete.md__e-git-bash-issues}

\subsection*{E.1 Windows Path Confusion}\label{docs__pandoc__latex__src__backmatter_complete.md__e1-windows-path-confusion}

Use forward slashes:

\begin{lstlisting}[style=Alabaster, language=bash]
openscad -o build/model.stl src/model.scad

\end{lstlisting}

\begin{center}\rule{0.5\linewidth}{0.5pt}\end{center}

\section*{F. 3D Printing Mechanical Failures}\label{docs__pandoc__latex__src__backmatter_complete.md__f-3d-printing-mechanical-failures}

\subsection*{F.1 Warping}\label{docs__pandoc__latex__src__backmatter_complete.md__f1-warping}

Solutions:

\begin{itemize}
\tightlist
\item
  Increase bed temp
\item
  Use enclosure
\item
  Add brim
\end{itemize}

\begin{center}\rule{0.5\linewidth}{0.5pt}\end{center}

\subsection*{F.2 Under-Extrusion}\label{docs__pandoc__latex__src__backmatter_complete.md__f2-under-extrusion}

Cause:

\begin{itemize}
\tightlist
\item
  Clogged nozzle
\item
  Incorrect flow rate
\end{itemize}

\begin{center}\rule{0.5\linewidth}{0.5pt}\end{center}

\subsection*{F.3 Dimensional Error}\label{docs__pandoc__latex__src__backmatter_complete.md__f3-dimensional-error}

Test with calibration cube.

\begin{center}\rule{0.5\linewidth}{0.5pt}\end{center}

\section*{G. Performance Diagnostics}\label{docs__pandoc__latex__src__backmatter_complete.md__g-performance-diagnostics}

\subsection*{G.1 Slow Render}\label{docs__pandoc__latex__src__backmatter_complete.md__g1-slow-render}

Reduce:

\begin{itemize}
\tightlist
\item
  \$fn
\item
  Minkowski usage
\item
  Nested difference()
\end{itemize}

\begin{center}\rule{0.5\linewidth}{0.5pt}\end{center}

\subsection*{G.2 High CPU Usage}\label{docs__pandoc__latex__src__backmatter_complete.md__g2-high-cpu-usage}

Expected during full render.

\begin{center}\rule{0.5\linewidth}{0.5pt}\end{center}

\subsection*{G.3 For Loop Performance}\label{docs__pandoc__latex__src__backmatter_complete.md__g3-for-loop-performance}

Inefficient:

\begin{lstlisting}[style=Alabaster, language=openscad]
for (i = [0:1000]) {
    translate([i, 0, 0])
        minkowski() { ... }  // Very slow
}

\end{lstlisting}

Better:

\begin{lstlisting}[style=Alabaster, language=openscad]
for (i = [0:100]) {  // Reduce iterations
    translate([i, 0, 0])
        hull() { ... }   // Faster than minkowski
}

\end{lstlisting}

\begin{center}\rule{0.5\linewidth}{0.5pt}\end{center}

\subsection*{G.4 Hull vs Minkowski Performance}\label{docs__pandoc__latex__src__backmatter_complete.md__g4-hull-vs-minkowski-performance}

\begin{itemize}
\tightlist
\item
  \texttt{hull()} --- Fast for convex shapes
\item
  \texttt{minkowski()} --- Slow but precise
\end{itemize}

Use hull() when possible.

\begin{center}\rule{0.5\linewidth}{0.5pt}\end{center}

\section*{H. SYSTEM DIAGNOSTIC CHECKLIST}\label{docs__pandoc__latex__src__backmatter_complete.md__h-system-diagnostic-checklist}

\begin{enumerate}
\tightlist
\item
  Is OpenSCAD in PATH?
\item
  Does CLI produce STL?
\item
  Is STL manifold?
\item
  Does slicer preview correctly?
\item
  Is G-code generated?
\item
  Does printer firmware accept G-code?
\item
  Is first layer adhering?
\item
  Are dimensions within tolerance?
\item
  Are all brackets matched?
\item
  Are comments properly closed?
\item
  Is \$fn set reasonably?
\item
  Are transform operations in correct order?
\end{enumerate}

\begin{center}\rule{0.5\linewidth}{0.5pt}\end{center}

\section*{J. EDITOR AND IDE TROUBLESHOOTING}\label{docs__pandoc__latex__src__backmatter_complete.md__j-editor-and-ide-troubleshooting}

\subsection*{J.1 VS Code OpenSCAD Extension Issues}\label{docs__pandoc__latex__src__backmatter_complete.md__j1-vs-code-openscad-extension-issues}

Symptom:

\begin{itemize}
\tightlist
\item
  No syntax highlighting
\end{itemize}

Fix:

\begin{enumerate}
\tightlist
\item
  Install "OpenSCAD" extension
\item
  Restart VS Code
\item
  Check file extension is \texttt{.scad}
\end{enumerate}

\begin{center}\rule{0.5\linewidth}{0.5pt}\end{center}

\subsection*{J.2 Notepad++ Syntax Highlighting}\label{docs__pandoc__latex__src__backmatter_complete.md__j2-notepad-syntax-highlighting}

Setup:

\begin{enumerate}
\tightlist
\item
  Settings \textgreater{} Style Configurator
\item
  Select Language: C or C++
\item
  Apply to \texttt{.scad} files
\end{enumerate}

\begin{center}\rule{0.5\linewidth}{0.5pt}\end{center}

\subsection*{J.3 Indentation Issues}\label{docs__pandoc__latex__src__backmatter_complete.md__j3-indentation-issues}

Mixed tabs/spaces:

Fix:

\begin{lstlisting}[style=Alabaster]
Convert tabs to spaces
Set indent to 4 spaces

\end{lstlisting}

\begin{center}\rule{0.5\linewidth}{0.5pt}\end{center}

\subsection*{J.4 Screen Reader + Editor Issues}\label{docs__pandoc__latex__src__backmatter_complete.md__j4-screen-reader--editor-issues}

\subsubsection*{NVDA with VS Code}\label{docs__pandoc__latex__src__backmatter_complete.md__nvda-with-vs-code}

Settings:

\begin{itemize}
\tightlist
\item
  Enable "Editor \textgreater{} Accessibility: Support"
\item
  Set "Editor \textgreater{} Render Whitespace": "all"
\end{itemize}

\subsubsection*{Indent Announcement}\label{docs__pandoc__latex__src__backmatter_complete.md__indent-announcement}

NVDA:

\begin{itemize}
\tightlist
\item
  Document Formatting \textgreater{} Report line indentation
\end{itemize}

JAWS:

\begin{itemize}
\tightlist
\item
  Settings Center \textgreater{} Reading \textgreater{} Report Indentation
\end{itemize}

\begin{center}\rule{0.5\linewidth}{0.5pt}\end{center}

\section*{K. SLICER-SPECIFIC TROUBLESHOOTING}\label{docs__pandoc__latex__src__backmatter_complete.md__k-slicer-specific-troubleshooting}

\subsection*{K.1 PrusaSlicer Import Failures}\label{docs__pandoc__latex__src__backmatter_complete.md__k1-prusaslicer-import-failures}

Symptom:

\begin{itemize}
\tightlist
\item
  STL won\textquotesingle t load
\end{itemize}

Causes:

\begin{itemize}
\tightlist
\item
  Non-manifold geometry
\item
  File corruption
\item
  Incorrect file path
\end{itemize}

Fix:

\begin{lstlisting}[style=Alabaster, language=bash]
openscad -o model.stl model.scad --check

\end{lstlisting}

\begin{center}\rule{0.5\linewidth}{0.5pt}\end{center}

\subsection*{K.2 Cura Slicing Errors}\label{docs__pandoc__latex__src__backmatter_complete.md__k2-cura-slicing-errors}

Common:

\begin{itemize}
\tightlist
\item
  Model outside build volume
\item
  Unsupported overhangs
\end{itemize}

Fix:

\begin{itemize}
\tightlist
\item
  Scale model down
\item
  Enable supports
\end{itemize}

\begin{center}\rule{0.5\linewidth}{0.5pt}\end{center}

\subsection*{K.3 Bambu Studio Issues}\label{docs__pandoc__latex__src__backmatter_complete.md__k3-bambu-studio-issues}

Plate adhesion:

\begin{itemize}
\tightlist
\item
  Clean build plate
\item
  Increase bed temperature
\item
  Add brim/raft
\end{itemize}

\begin{center}\rule{0.5\linewidth}{0.5pt}\end{center}

\section*{L. MATERIAL-SPECIFIC FAILURES}\label{docs__pandoc__latex__src__backmatter_complete.md__l-material-specific-failures}

\subsection*{L.1 PLA Printing Issues}\label{docs__pandoc__latex__src__backmatter_complete.md__l1-pla-printing-issues}

\subsubsection*{Stringing}\label{docs__pandoc__latex__src__backmatter_complete.md__stringing}

Fix:

\begin{itemize}
\tightlist
\item
  Enable retraction
\item
  Lower temperature by 5-10°C
\end{itemize}

\subsubsection*{Warping}\label{docs__pandoc__latex__src__backmatter_complete.md__warping}

\begin{itemize}
\tightlist
\item
  Rare with PLA
\item
  Add brim if occurs
\end{itemize}

\begin{center}\rule{0.5\linewidth}{0.5pt}\end{center}

\subsection*{L.2 PETG Challenges}\label{docs__pandoc__latex__src__backmatter_complete.md__l2-petg-challenges}

\subsubsection*{Stringing}\label{docs__pandoc__latex__src__backmatter_complete.md__stringing-1}

Common with PETG.

Fix:

\begin{itemize}
\tightlist
\item
  Increase retraction distance
\item
  Lower temperature
\end{itemize}

\subsubsection*{Bed Adhesion}\label{docs__pandoc__latex__src__backmatter_complete.md__bed-adhesion}

Too good --- may damage bed.

Use:

\begin{itemize}
\tightlist
\item
  Glue stick as release agent
\end{itemize}

\begin{center}\rule{0.5\linewidth}{0.5pt}\end{center}

\subsection*{L.3 TPU (Flexible) Problems}\label{docs__pandoc__latex__src__backmatter_complete.md__l3-tpu-flexible-problems}

\subsubsection*{Under-extrusion}\label{docs__pandoc__latex__src__backmatter_complete.md__under-extrusion}

Cause:

\begin{itemize}
\tightlist
\item
  Too fast print speed
\end{itemize}

Fix:

\begin{itemize}
\tightlist
\item
  Reduce to 20-30 mm/s
\end{itemize}

\subsubsection*{Stringing}\label{docs__pandoc__latex__src__backmatter_complete.md__stringing-2}

Minimal retraction:

\begin{itemize}
\tightlist
\item
  Set to 1-2mm only
\end{itemize}

\begin{center}\rule{0.5\linewidth}{0.5pt}\end{center}

\subsection*{L.4 ABS Warping}\label{docs__pandoc__latex__src__backmatter_complete.md__l4-abs-warping}

Major challenge.

Required:

\begin{itemize}
\tightlist
\item
  Heated enclosure
\item
  100°C+ bed temp
\item
  Large brim
\end{itemize}

\begin{center}\rule{0.5\linewidth}{0.5pt}\end{center}

\section*{M. MEASUREMENT AND QA FAILURES}\label{docs__pandoc__latex__src__backmatter_complete.md__m-measurement-and-qa-failures}

\subsection*{M.1 Caliper Measurement Errors}\label{docs__pandoc__latex__src__backmatter_complete.md__m1-caliper-measurement-errors}

\subsubsection*{Over-tightening}\label{docs__pandoc__latex__src__backmatter_complete.md__over-tightening}

Cause:

\begin{itemize}
\tightlist
\item
  Pressing too hard
\end{itemize}

Result:

\begin{itemize}
\tightlist
\item
  Part deforms
\item
  Inaccurate reading
\end{itemize}

Fix:

\begin{itemize}
\tightlist
\item
  Use gentle pressure
\end{itemize}

\begin{center}\rule{0.5\linewidth}{0.5pt}\end{center}

\subsection*{M.2 Tolerance Stack-up Issues}\label{docs__pandoc__latex__src__backmatter_complete.md__m2-tolerance-stack-up-issues}

Multiple parts don\textquotesingle t fit.

Analysis:

\begin{lstlisting}[style=Alabaster]
Part A: +0.15mm error
Part B: +0.12mm error
Total: +0.27mm

\end{lstlisting}

Fix:

\begin{itemize}
\tightlist
\item
  Adjust design tolerances
\item
  Print test pieces first
\end{itemize}

\begin{center}\rule{0.5\linewidth}{0.5pt}\end{center}

\subsection*{M.3 Dimensional Accuracy Check}\label{docs__pandoc__latex__src__backmatter_complete.md__m3-dimensional-accuracy-check}

Print calibration cube:

\begin{lstlisting}[style=Alabaster, language=openscad]
cube([20,20,20]);

\end{lstlisting}

Measure all dimensions. Expected: 20.00 ± 0.2mm

\begin{center}\rule{0.5\linewidth}{0.5pt}\end{center}

\section*{I. FULL BUILD VALIDATION SCRIPT (CMD)}\label{docs__pandoc__latex__src__backmatter_complete.md__i-full-build-validation-script-cmd}

\begin{lstlisting}[style=Alabaster, language=cmd]
@echo off
echo Validating environment...
where openscad >nul 2>nul
if %ERRORLEVEL% NEQ 0 (
    echo OpenSCAD not found.
    exit /b 1
)

mkdir build
openscad -o build\model.stl src\model.scad
if %ERRORLEVEL% NEQ 0 (
    echo Build failed.
    exit /b 1
)

echo Build succeeded.

\end{lstlisting}

\begin{center}\rule{0.5\linewidth}{0.5pt}\end{center}

\section*{N. TEXT EMBOSSING TROUBLESHOOTING}\label{docs__pandoc__latex__src__backmatter_complete.md__n-text-embossing-troubleshooting}

\subsection*{N.1 Font Not Found}\label{docs__pandoc__latex__src__backmatter_complete.md__n1-font-not-found}

Symptom:

\begin{lstlisting}[style=Alabaster]
WARNING: No match for font family

\end{lstlisting}

Cause:

\begin{itemize}
\tightlist
\item
  Font name incorrect
\item
  Font not installed
\end{itemize}

Fix:

Windows (PowerShell):

\begin{lstlisting}[style=Alabaster, language=powershell]
[System.Drawing.Text.InstalledFontCollection]::new().Families | Select-Object Name

\end{lstlisting}

Linux:

\begin{lstlisting}[style=Alabaster, language=bash]
fc-list | grep -i "font name"

\end{lstlisting}

\begin{center}\rule{0.5\linewidth}{0.5pt}\end{center}

\subsection*{N.2 Text Not Visible}\label{docs__pandoc__latex__src__backmatter_complete.md__n2-text-not-visible}

Cause:

\begin{itemize}
\tightlist
\item
  Forgot \texttt{linear\_extrude()}
\end{itemize}

Wrong:

\begin{lstlisting}[style=Alabaster, language=openscad]
text("HELLO");

\end{lstlisting}

Correct:

\begin{lstlisting}[style=Alabaster, language=openscad]
linear_extrude(height=2)
    text("HELLO");

\end{lstlisting}

\begin{center}\rule{0.5\linewidth}{0.5pt}\end{center}

\subsection*{N.3 Text Size Too Large/Small}\label{docs__pandoc__latex__src__backmatter_complete.md__n3-text-size-too-largesmall}

Cause:

\begin{itemize}
\tightlist
\item
  Size parameter is cap height in mm
\end{itemize}

Fix:

\begin{lstlisting}[style=Alabaster, language=openscad]
text("HELLO", size=8);  // 8mm tall letters

\end{lstlisting}

\begin{center}\rule{0.5\linewidth}{0.5pt}\end{center}

\section*{O. TRANSFORM OPERATION ERRORS}\label{docs__pandoc__latex__src__backmatter_complete.md__o-transform-operation-errors}

\subsection*{O.1 Transform Order Confusion}\label{docs__pandoc__latex__src__backmatter_complete.md__o1-transform-order-confusion}

Transforms apply \textbf{inside-out}.

Example:

\begin{lstlisting}[style=Alabaster, language=openscad]
translate([10,0,0])    // Applied SECOND
    rotate([0,0,45])    // Applied FIRST
        cube([10,10,10]);

\end{lstlisting}

Order matters!

\begin{center}\rule{0.5\linewidth}{0.5pt}\end{center}

\subsection*{O.2 Rotate Angles Wrong}\label{docs__pandoc__latex__src__backmatter_complete.md__o2-rotate-angles-wrong}

Angles in \textbf{degrees}, not radians.

\begin{lstlisting}[style=Alabaster, language=openscad]
rotate([0, 0, 45]);  // 45 degrees

\end{lstlisting}

\begin{center}\rule{0.5\linewidth}{0.5pt}\end{center}

\subsection*{O.3 Scale Issues}\label{docs__pandoc__latex__src__backmatter_complete.md__o3-scale-issues}

\begin{lstlisting}[style=Alabaster, language=openscad]
scale([2, 1, 1])  // 2x wider only
    cube([10,10,10]);

\end{lstlisting}

Result: 20×10×10 cube

\begin{center}\rule{0.5\linewidth}{0.5pt}\end{center}

\subsection*{O.4 Mirror Confusion}\label{docs__pandoc__latex__src__backmatter_complete.md__o4-mirror-confusion}

Mirror plane:

\begin{lstlisting}[style=Alabaster, language=openscad]
mirror([1, 0, 0])  // Mirror across YZ plane
    cube([10,10,10]);

\end{lstlisting}

\begin{center}\rule{0.5\linewidth}{0.5pt}\end{center}

\section*{P. IMPORT AND EXPORT ERRORS}\label{docs__pandoc__latex__src__backmatter_complete.md__p-import-and-export-errors}

\subsection*{P.1 Import STL Failed}\label{docs__pandoc__latex__src__backmatter_complete.md__p1-import-stl-failed}

Cause:

\begin{itemize}
\tightlist
\item
  File path incorrect
\item
  File corrupted
\end{itemize}

Fix:

\begin{lstlisting}[style=Alabaster, language=openscad]
import("path/to/file.stl");

\end{lstlisting}

Use forward slashes even on Windows.

\begin{center}\rule{0.5\linewidth}{0.5pt}\end{center}

\subsection*{P.2 Export Failed}\label{docs__pandoc__latex__src__backmatter_complete.md__p2-export-failed}

Check:

\begin{enumerate}
\tightlist
\item
  Write permissions
\item
  Sufficient disk space
\item
  Valid file path
\end{enumerate}

\begin{center}\rule{0.5\linewidth}{0.5pt}\end{center}

\section*{Q. RANDOMIZATION ISSUES}\label{docs__pandoc__latex__src__backmatter_complete.md__q-randomization-issues}

\subsection*{Q.1 rands() Not Random}\label{docs__pandoc__latex__src__backmatter_complete.md__q1-rands-not-random}

Cause:

\begin{itemize}
\tightlist
\item
  Same seed produces same output
\end{itemize}

\begin{lstlisting}[style=Alabaster, language=openscad]
vals = rands(0, 100, 10, 42);  // Seed = 42

\end{lstlisting}

Change seed for different results.

\begin{center}\rule{0.5\linewidth}{0.5pt}\end{center}

\subsection*{Q.2 Random Position Overlap}\label{docs__pandoc__latex__src__backmatter_complete.md__q2-random-position-overlap}

Objects placed randomly may overlap.

Need collision detection logic.

\begin{center}\rule{0.5\linewidth}{0.5pt}\end{center}

\begin{center}\rule{0.5\linewidth}{0.5pt}\end{center}

\protect\phantomsection\label{docs__pandoc__latex__src__backmatter_complete.md__student-glossary}{}

\section*{Student Glossary --- OpenSCAD • 3dMake • 3D Printing • PowerShell • Command Prompt • Git Bash}\label{docs__pandoc__latex__src__backmatter_complete.md__student-glossary--openscad--3dmake--3d-printing--powershell--command-prompt--git-bash}

\begin{center}\rule{0.5\linewidth}{0.5pt}\end{center}

\emph{OpenSCAD • 3dMake • 3D Printing • PowerShell • Command Prompt • Git Bash}

This glossary provides:

\begin{itemize}
\tightlist
\item
  Formal definitions
\item
  Operational context
\item
  Cross-system distinctions
\item
  CLI examples
\item
  Practical implications in 3D printing workflows
\end{itemize}

\begin{center}\rule{0.5\linewidth}{0.5pt}\end{center}

\section*{SECTION I --- Additive Manufacturing \& Fabrication}\label{docs__pandoc__latex__src__backmatter_complete.md__section-i--additive-manufacturing--fabrication}

\subsection*{Additive Manufacturing}\label{docs__pandoc__latex__src__backmatter_complete.md__additive-manufacturing}

A manufacturing methodology in which objects are fabricated layer-by-layer from digital models.\\
In desktop contexts, this typically refers to FDM (Fused Deposition Modeling).

\subsubsection*{Technical Characteristics}\label{docs__pandoc__latex__src__backmatter_complete.md__technical-characteristics}

\begin{itemize}
\tightlist
\item
  Layer-based deposition
\item
  Toolpath-generated geometry
\item
  STL-to-G-code workflow
\item
  Thermoplastic extrusion (in FDM)
\end{itemize}

\subsubsection*{Workflow Position}\label{docs__pandoc__latex__src__backmatter_complete.md__workflow-position}

OpenSCAD -\textgreater{} STL -\textgreater{} Slicer -\textgreater{} G-code -\textgreater{} Printer -\textgreater{} Physical object

\begin{center}\rule{0.5\linewidth}{0.5pt}\end{center}

\subsection*{FDM (Fused Deposition Modeling)}\label{docs__pandoc__latex__src__backmatter_complete.md__fdm-fused-deposition-modeling}

A thermoplastic extrusion process in which filament is melted and deposited in sequential layers.

\subsubsection*{Critical Variables}\label{docs__pandoc__latex__src__backmatter_complete.md__critical-variables}

\begin{itemize}
\tightlist
\item
  Nozzle temperature
\item
  Bed temperature
\item
  Layer height
\item
  Print speed
\item
  Cooling rate
\end{itemize}

\begin{center}\rule{0.5\linewidth}{0.5pt}\end{center}

\subsection*{Layer Height}\label{docs__pandoc__latex__src__backmatter_complete.md__layer-height}

The vertical thickness of each printed layer.

\begin{itemize}
\tightlist
\item
  Smaller values increase surface fidelity.
\item
  Larger values increase speed.
\end{itemize}

Example:

\begin{itemize}
\tightlist
\item
  0.2 mm = standard
\item
  0.1 mm = high resolution
\end{itemize}

\begin{center}\rule{0.5\linewidth}{0.5pt}\end{center}

\subsection*{Infill}\label{docs__pandoc__latex__src__backmatter_complete.md__infill}

Internal structural lattice inside a model.

Typical values:

\begin{itemize}
\tightlist
\item
  10--20\% for visual models
\item
  40\%+ for structural parts
\end{itemize}

\begin{center}\rule{0.5\linewidth}{0.5pt}\end{center}

\subsection*{Material Properties}\label{docs__pandoc__latex__src__backmatter_complete.md__material-properties}

\subsubsection*{PLA (Polylactic Acid)}\label{docs__pandoc__latex__src__backmatter_complete.md__pla-polylactic-acid}

\begin{itemize}
\tightlist
\item
  Easy to print
\item
  Low warping
\item
  Biodegradable
\item
  Not heat-resistant
\end{itemize}

\subsubsection*{PETG}\label{docs__pandoc__latex__src__backmatter_complete.md__petg}

\begin{itemize}
\tightlist
\item
  Stronger than PLA
\item
  Chemical resistant
\item
  Slight flexibility
\end{itemize}

\subsubsection*{TPU (Thermoplastic Polyurethane)}\label{docs__pandoc__latex__src__backmatter_complete.md__tpu-thermoplastic-polyurethane}

\begin{itemize}
\tightlist
\item
  Flexible
\item
  Rubber-like
\item
  Requires slow printing
\end{itemize}

\subsubsection*{ABS (Acrylonitrile Butadiene Styrene)}\label{docs__pandoc__latex__src__backmatter_complete.md__abs-acrylonitrile-butadiene-styrene}

\begin{itemize}
\tightlist
\item
  Strong and durable
\item
  Higher temperature
\item
  Requires enclosure
\item
  Prone to warping
\end{itemize}

\begin{center}\rule{0.5\linewidth}{0.5pt}\end{center}

\subsection*{Tolerance}\label{docs__pandoc__latex__src__backmatter_complete.md__tolerance}

Intentional dimensional offset to ensure mechanical fit.

Example: A 10mm peg may require a 10.2mm hole for clearance.

\begin{center}\rule{0.5\linewidth}{0.5pt}\end{center}

\subsection*{Manifold (Watertight Model)}\label{docs__pandoc__latex__src__backmatter_complete.md__manifold-watertight-model}

A printable model with:

\begin{itemize}
\tightlist
\item
  No self-intersections
\item
  No zero-thickness faces
\item
  No holes in mesh
\end{itemize}

Non-manifold geometry causes slicer failure.

\begin{center}\rule{0.5\linewidth}{0.5pt}\end{center}

\section*{SECTION II --- OpenSCAD (Parametric Modeling)}\label{docs__pandoc__latex__src__backmatter_complete.md__section-ii--openscad-parametric-modeling}

\subsection*{OpenSCAD}\label{docs__pandoc__latex__src__backmatter_complete.md__openscad}

A script-based solid modeling system using Constructive Solid Geometry (CSG).

Official site: https://openscad.org

\begin{center}\rule{0.5\linewidth}{0.5pt}\end{center}

\subsection*{Constructive Solid Geometry (CSG)}\label{docs__pandoc__latex__src__backmatter_complete.md__constructive-solid-geometry-csg}

Modeling technique combining primitives using Boolean operations.

\begin{center}\rule{0.5\linewidth}{0.5pt}\end{center}

\subsection*{Primitive}\label{docs__pandoc__latex__src__backmatter_complete.md__primitive}

Basic geometric object.

Examples:

\begin{lstlisting}[style=Alabaster, language=openscad]
cube([10,10,10]);
sphere(5);
cylinder(h=20, d=10);

\end{lstlisting}

\begin{center}\rule{0.5\linewidth}{0.5pt}\end{center}

\subsection*{Cone}\label{docs__pandoc__latex__src__backmatter_complete.md__cone}

A tapered cylinder created using different top and bottom diameters.

\begin{lstlisting}[style=Alabaster, language=openscad]
cylinder(h=20, d1=10, d2=0);

\end{lstlisting}

\begin{center}\rule{0.5\linewidth}{0.5pt}\end{center}

\subsection*{2D Shapes}\label{docs__pandoc__latex__src__backmatter_complete.md__2d-shapes}

Flat geometry that can be extruded to 3D.

Examples:

\begin{lstlisting}[style=Alabaster, language=openscad]
square([10,10]);
circle(5);
polygon([[0,0], [10,0], [5,10]]);

\end{lstlisting}

\begin{center}\rule{0.5\linewidth}{0.5pt}\end{center}

\subsection*{Polygon}\label{docs__pandoc__latex__src__backmatter_complete.md__polygon}

Custom 2D shape defined by vertices.

\begin{lstlisting}[style=Alabaster, language=openscad]
polygon(points=[[0,0], [10,0], [10,10], [0,10]]);

\end{lstlisting}

\begin{center}\rule{0.5\linewidth}{0.5pt}\end{center}

\subsection*{Text Commands}\label{docs__pandoc__latex__src__backmatter_complete.md__text-commands}

Create text geometry for embossing or debossing.

\begin{lstlisting}[style=Alabaster, language=openscad]
linear_extrude(height=2)
    text("HELLO", size=10, font="Liberation Sans");

\end{lstlisting}

\begin{center}\rule{0.5\linewidth}{0.5pt}\end{center}

\subsection*{Boolean Operations}\label{docs__pandoc__latex__src__backmatter_complete.md__boolean-operations}

\subsubsection*{union()}\label{docs__pandoc__latex__src__backmatter_complete.md__union}

Combines shapes.

\begin{lstlisting}[style=Alabaster, language=openscad]
union() {
    cube([10,10,10]);
    sphere(6);
}

\end{lstlisting}

\begin{center}\rule{0.5\linewidth}{0.5pt}\end{center}

\subsubsection*{difference()}\label{docs__pandoc__latex__src__backmatter_complete.md__difference}

Subtracts shapes.

\begin{lstlisting}[style=Alabaster, language=openscad]
difference() {
    cube([20,20,20]);
    cylinder(h=25, d=5);
}

\end{lstlisting}

\begin{center}\rule{0.5\linewidth}{0.5pt}\end{center}

\subsubsection*{intersection()}\label{docs__pandoc__latex__src__backmatter_complete.md__intersection}

Keeps overlapping geometry.

\begin{lstlisting}[style=Alabaster, language=openscad]
intersection() {
    sphere(10);
    cube([15,15,15], center=true);
}

\end{lstlisting}

\begin{center}\rule{0.5\linewidth}{0.5pt}\end{center}

\subsection*{Transform Operations}\label{docs__pandoc__latex__src__backmatter_complete.md__transform-operations}

\subsubsection*{Translate}\label{docs__pandoc__latex__src__backmatter_complete.md__translate}

Move geometry in 3D space.

\begin{lstlisting}[style=Alabaster, language=openscad]
translate([10, 20, 30])
    cube([10,10,10]);

\end{lstlisting}

\begin{center}\rule{0.5\linewidth}{0.5pt}\end{center}

\subsubsection*{Rotate}\label{docs__pandoc__latex__src__backmatter_complete.md__rotate}

Rotate geometry around axes.

\begin{lstlisting}[style=Alabaster, language=openscad]
rotate([0, 0, 45])
    cube([10,10,10]);

\end{lstlisting}

\begin{center}\rule{0.5\linewidth}{0.5pt}\end{center}

\subsubsection*{Scale}\label{docs__pandoc__latex__src__backmatter_complete.md__scale}

Resize geometry by multiplier.

\begin{lstlisting}[style=Alabaster, language=openscad]
scale([2, 1, 1])
    cube([10,10,10]);

\end{lstlisting}

\begin{center}\rule{0.5\linewidth}{0.5pt}\end{center}

\subsubsection*{Mirror}\label{docs__pandoc__latex__src__backmatter_complete.md__mirror}

Flip geometry across a plane.

\begin{lstlisting}[style=Alabaster, language=openscad]
mirror([1, 0, 0])
    cube([10,10,10]);

\end{lstlisting}

\begin{center}\rule{0.5\linewidth}{0.5pt}\end{center}

\subsection*{Advanced Operations}\label{docs__pandoc__latex__src__backmatter_complete.md__advanced-operations}

\subsubsection*{Hull}\label{docs__pandoc__latex__src__backmatter_complete.md__hull}

Creates convex hull of children objects.

\begin{lstlisting}[style=Alabaster, language=openscad]
hull() {
    sphere(5);
    translate([20,0,0]) sphere(5);
}

\end{lstlisting}

\begin{center}\rule{0.5\linewidth}{0.5pt}\end{center}

\subsubsection*{Minkowski Sum}\label{docs__pandoc__latex__src__backmatter_complete.md__minkowski-sum}

Expands geometry by rolling one shape around another.

\begin{lstlisting}[style=Alabaster, language=openscad]
minkowski() {
    cube([10,10,10]);
    sphere(2, $fn=12);
}

\end{lstlisting}

Warning: Computationally expensive.

\begin{center}\rule{0.5\linewidth}{0.5pt}\end{center}

\subsubsection*{Linear Extrude}\label{docs__pandoc__latex__src__backmatter_complete.md__linear-extrude}

Extends 2D shapes into 3D.

\begin{lstlisting}[style=Alabaster, language=openscad]
linear_extrude(height=10)
    square([10,10]);

\end{lstlisting}

\begin{center}\rule{0.5\linewidth}{0.5pt}\end{center}

\subsubsection*{Offset}\label{docs__pandoc__latex__src__backmatter_complete.md__offset}

Expands or contracts 2D shapes.

\begin{lstlisting}[style=Alabaster, language=openscad]
offset(r=2)
    square([10,10]);

\end{lstlisting}

\begin{center}\rule{0.5\linewidth}{0.5pt}\end{center}

\subsection*{Variable}\label{docs__pandoc__latex__src__backmatter_complete.md__variable}

A named parameter used to control geometry.

\begin{lstlisting}[style=Alabaster, language=openscad]
width = 30;
cube([width,10,5]);

\end{lstlisting}

\begin{center}\rule{0.5\linewidth}{0.5pt}\end{center}

\subsection*{Parametric Design}\label{docs__pandoc__latex__src__backmatter_complete.md__parametric-design}

Geometry controlled by variables.

Advantages:

\begin{itemize}
\tightlist
\item
  Rapid iteration
\item
  Reusability
\item
  Scalable models
\end{itemize}

\begin{center}\rule{0.5\linewidth}{0.5pt}\end{center}

\subsection*{Module}\label{docs__pandoc__latex__src__backmatter_complete.md__module}

Reusable geometry block.

\begin{lstlisting}[style=Alabaster, language=openscad]
module peg(d, h) {
    cylinder(d=d, h=h);
}

peg(5, 20);

\end{lstlisting}

\begin{center}\rule{0.5\linewidth}{0.5pt}\end{center}

\subsection*{Control Structures}\label{docs__pandoc__latex__src__backmatter_complete.md__control-structures}

\subsubsection*{For Loop}\label{docs__pandoc__latex__src__backmatter_complete.md__for-loop}

Repeat operations.

\begin{lstlisting}[style=Alabaster, language=openscad]
for (i = [0:5]) {
    translate([i*15, 0, 0])
        cube([10,10,10]);
}

\end{lstlisting}

\begin{center}\rule{0.5\linewidth}{0.5pt}\end{center}

\subsubsection*{If Statement}\label{docs__pandoc__latex__src__backmatter_complete.md__if-statement}

Conditional execution.

\begin{lstlisting}[style=Alabaster, language=openscad]
if (width > 20) {
    cube([width, 10, 10]);
}

\end{lstlisting}

\begin{center}\rule{0.5\linewidth}{0.5pt}\end{center}

\subsubsection*{Else and ElseIf}\label{docs__pandoc__latex__src__backmatter_complete.md__else-and-elseif}

Extended conditionals.

\begin{lstlisting}[style=Alabaster, language=openscad]
if (width > 20) {
    cube([width, 10, 10]);
} else if (width > 10) {
    cylinder(h=10, d=width);
} else {
    sphere(width/2);
}

\end{lstlisting}

\begin{center}\rule{0.5\linewidth}{0.5pt}\end{center}

\subsection*{Echo Function}\label{docs__pandoc__latex__src__backmatter_complete.md__echo-function}

Print debug information to console.

\begin{lstlisting}[style=Alabaster, language=openscad]
width = 30;
echo("Width is:", width);

\end{lstlisting}

\begin{center}\rule{0.5\linewidth}{0.5pt}\end{center}

\subsection*{Random Numbers}\label{docs__pandoc__latex__src__backmatter_complete.md__random-numbers}

Generate random values.

\begin{lstlisting}[style=Alabaster, language=openscad]
rands(min, max, count, seed);

// Example:
positions = rands(0, 100, 5);

\end{lstlisting}

\begin{center}\rule{0.5\linewidth}{0.5pt}\end{center}

\subsection*{Import}\label{docs__pandoc__latex__src__backmatter_complete.md__import}

Load external STL or other files.

\begin{lstlisting}[style=Alabaster, language=openscad]
import("model.stl");

\end{lstlisting}

\begin{center}\rule{0.5\linewidth}{0.5pt}\end{center}

\subsection*{\$fn (Fragment Number)}\label{docs__pandoc__latex__src__backmatter_complete.md__fn-fragment-number}

Controls circular resolution.

\begin{lstlisting}[style=Alabaster, language=openscad]
$fn = 100;
sphere(10);

\end{lstlisting}

Higher values increase smoothness and render time.

\begin{center}\rule{0.5\linewidth}{0.5pt}\end{center}

\subsection*{\# (Hash Mark - Debug Modifier)}\label{docs__pandoc__latex__src__backmatter_complete.md__-hash-mark---debug-modifier}

Highlights geometry in transparent red during preview.

\begin{lstlisting}[style=Alabaster, language=openscad]
#cube([10,10,10]);

\end{lstlisting}

Used for debugging geometry positions.

\begin{center}\rule{0.5\linewidth}{0.5pt}\end{center}

\subsection*{Comments}\label{docs__pandoc__latex__src__backmatter_complete.md__comments}

Non-executable text for documentation.

\subsubsection*{Single-line Comments}\label{docs__pandoc__latex__src__backmatter_complete.md__single-line-comments}

\begin{lstlisting}[style=Alabaster, language=openscad]
// This is a comment
cube([10,10,10]);

\end{lstlisting}

\subsubsection*{Multiline Comments}\label{docs__pandoc__latex__src__backmatter_complete.md__multiline-comments}

\begin{lstlisting}[style=Alabaster, language=openscad]
/*
  This is a
  multiline comment
*/
cube([10,10,10]);

\end{lstlisting}

\begin{center}\rule{0.5\linewidth}{0.5pt}\end{center}

\subsection*{Logical Operators}\label{docs__pandoc__latex__src__backmatter_complete.md__logical-operators}

\subsubsection*{\&\& (And)}\label{docs__pandoc__latex__src__backmatter_complete.md__-and}

Logical AND operator.

\begin{lstlisting}[style=Alabaster, language=openscad]
if (x > 5 && y < 10) {
    cube([10,10,10]);
}

\end{lstlisting}

\subsubsection*{\textbar\textbar{} (Or)}\label{docs__pandoc__latex__src__backmatter_complete.md__-or}

Logical OR operator.

\begin{lstlisting}[style=Alabaster, language=openscad]
if (x > 5 || y < 10) {
    cube([10,10,10]);
}

\end{lstlisting}

\begin{center}\rule{0.5\linewidth}{0.5pt}\end{center}

\subsection*{Brackets and Delimiters}\label{docs__pandoc__latex__src__backmatter_complete.md__brackets-and-delimiters}

\subsubsection*{{[} {]} (Square Brackets)}\label{docs__pandoc__latex__src__backmatter_complete.md__--square-brackets}

Used for vectors and arrays.

\begin{lstlisting}[style=Alabaster, language=openscad]
position = [10, 20, 30];
cube([10,10,10]);

\end{lstlisting}

\subsubsection*{\{ \} (Curly Brackets)}\label{docs__pandoc__latex__src__backmatter_complete.md__--curly-brackets}

Define code blocks.

\begin{lstlisting}[style=Alabaster, language=openscad]
if (x > 5) {
    cube([10,10,10]);
}

\end{lstlisting}

\subsubsection*{\textless{} \textgreater{} (Angle Brackets)}\label{docs__pandoc__latex__src__backmatter_complete.md__--angle-brackets}

Used in some operations and comparisons.

\begin{lstlisting}[style=Alabaster, language=openscad]
if (x < 10 && y > 5) { }

\end{lstlisting}

\begin{center}\rule{0.5\linewidth}{0.5pt}\end{center}

\subsection*{Preview vs Render}\label{docs__pandoc__latex__src__backmatter_complete.md__preview-vs-render}

\begin{itemize}
\tightlist
\item
  F5 = Preview (OpenGL approximation)
\item
  F6 = Render (full CSG evaluation)
\end{itemize}

CLI equivalent:

\begin{lstlisting}[style=Alabaster, language=bash]
openscad -o model.stl model.scad

\end{lstlisting}

\begin{center}\rule{0.5\linewidth}{0.5pt}\end{center}

\section*{SECTION III --- 3dMake (Automation Tool)}\label{docs__pandoc__latex__src__backmatter_complete.md__section-iii--3dmake-automation-tool}

\subsection*{3dMake}\label{docs__pandoc__latex__src__backmatter_complete.md__3dmake}

A command-line tool for automating 3D model generation and processing.

Repository: https://github.com/tdeck/3dmake

\begin{center}\rule{0.5\linewidth}{0.5pt}\end{center}

\subsection*{3dm build}\label{docs__pandoc__latex__src__backmatter_complete.md__3dm-build}

Command to generate STL from OpenSCAD source.

\begin{lstlisting}[style=Alabaster, language=bash]
3dm build

\end{lstlisting}

Outputs to \texttt{build/} directory.

\begin{center}\rule{0.5\linewidth}{0.5pt}\end{center}

\subsection*{3dm info}\label{docs__pandoc__latex__src__backmatter_complete.md__3dm-info}

AI-powered geometric description for non-visual validation.

\begin{lstlisting}[style=Alabaster, language=bash]
3dm info --view front

\end{lstlisting}

\begin{center}\rule{0.5\linewidth}{0.5pt}\end{center}

\subsection*{Headless Rendering}\label{docs__pandoc__latex__src__backmatter_complete.md__headless-rendering}

Running OpenSCAD without GUI.

\begin{lstlisting}[style=Alabaster, language=bash]
openscad -o output.stl input.scad

\end{lstlisting}

\begin{center}\rule{0.5\linewidth}{0.5pt}\end{center}

\subsection*{Automated Build Pipeline}\label{docs__pandoc__latex__src__backmatter_complete.md__automated-build-pipeline}

Example PowerShell pipeline:

\begin{lstlisting}[style=Alabaster, language=powershell]
openscad -o model.stl model.scad

\end{lstlisting}

Example Bash:

\begin{lstlisting}[style=Alabaster, language=bash]
openscad -o model.stl model.scad

\end{lstlisting}

\begin{center}\rule{0.5\linewidth}{0.5pt}\end{center}

\section*{SECTION IV --- Command-Line Systems}\label{docs__pandoc__latex__src__backmatter_complete.md__section-iv--command-line-systems}

\section*{CLI (Command-Line Interface)}\label{docs__pandoc__latex__src__backmatter_complete.md__cli-command-line-interface}

A text-based interface for interacting with the operating system.

\begin{center}\rule{0.5\linewidth}{0.5pt}\end{center}

\section*{Shell}\label{docs__pandoc__latex__src__backmatter_complete.md__shell}

A program that interprets commands.

Examples:

\begin{itemize}
\tightlist
\item
  PowerShell
\item
  Command Prompt
\item
  Git Bash
\end{itemize}

\begin{center}\rule{0.5\linewidth}{0.5pt}\end{center}

\section*{PowerShell}\label{docs__pandoc__latex__src__backmatter_complete.md__powershell}

Object-oriented Windows shell.

\subsubsection*{List Files}\label{docs__pandoc__latex__src__backmatter_complete.md__list-files-1}

\begin{lstlisting}[style=Alabaster, language=powershell]
Get-ChildItem

\end{lstlisting}

\subsubsection*{Change Directory}\label{docs__pandoc__latex__src__backmatter_complete.md__change-directory-1}

\begin{lstlisting}[style=Alabaster, language=powershell]
Set-Location Documents

\end{lstlisting}

\subsubsection*{Run Script}\label{docs__pandoc__latex__src__backmatter_complete.md__run-script}

\begin{lstlisting}[style=Alabaster, language=powershell]
.\build.ps1

\end{lstlisting}

\subsubsection*{Check Exit Code}\label{docs__pandoc__latex__src__backmatter_complete.md__check-exit-code}

\begin{lstlisting}[style=Alabaster, language=powershell]
$LASTEXITCODE

\end{lstlisting}

\begin{center}\rule{0.5\linewidth}{0.5pt}\end{center}

\section*{Command Prompt (cmd.exe)}\label{docs__pandoc__latex__src__backmatter_complete.md__command-prompt-cmdexe}

Traditional Windows shell.

\subsubsection*{List Files}\label{docs__pandoc__latex__src__backmatter_complete.md__list-files-2}

\begin{lstlisting}[style=Alabaster, language=cmd]
dir

\end{lstlisting}

\subsubsection*{Change Directory}\label{docs__pandoc__latex__src__backmatter_complete.md__change-directory-2}

\begin{lstlisting}[style=Alabaster, language=cmd]
cd Documents

\end{lstlisting}

\subsubsection*{Run Program}\label{docs__pandoc__latex__src__backmatter_complete.md__run-program}

\begin{lstlisting}[style=Alabaster, language=cmd]
openscad.exe -o model.stl model.scad

\end{lstlisting}

\subsubsection*{Check Exit Code}\label{docs__pandoc__latex__src__backmatter_complete.md__check-exit-code-1}

\begin{lstlisting}[style=Alabaster, language=cmd]
echo %ERRORLEVEL%

\end{lstlisting}

\begin{center}\rule{0.5\linewidth}{0.5pt}\end{center}

\section*{Git Bash}\label{docs__pandoc__latex__src__backmatter_complete.md__git-bash}

Unix-like shell for Windows.

\subsubsection*{List Files}\label{docs__pandoc__latex__src__backmatter_complete.md__list-files-3}

\begin{lstlisting}[style=Alabaster, language=bash]
ls

\end{lstlisting}

\subsubsection*{Change Directory}\label{docs__pandoc__latex__src__backmatter_complete.md__change-directory-3}

\begin{lstlisting}[style=Alabaster, language=bash]
cd Documents

\end{lstlisting}

\subsubsection*{Make Directory}\label{docs__pandoc__latex__src__backmatter_complete.md__make-directory}

\begin{lstlisting}[style=Alabaster, language=bash]
mkdir builds

\end{lstlisting}

\begin{center}\rule{0.5\linewidth}{0.5pt}\end{center}

\section*{Working Directory}\label{docs__pandoc__latex__src__backmatter_complete.md__working-directory}

The directory in which commands execute.

\subsubsection*{Print Working Directory}\label{docs__pandoc__latex__src__backmatter_complete.md__print-working-directory}

\begin{lstlisting}[style=Alabaster, language=bash]
pwd

\end{lstlisting}

\begin{lstlisting}[style=Alabaster, language=powershell]
Get-Location

\end{lstlisting}

\begin{center}\rule{0.5\linewidth}{0.5pt}\end{center}

\section*{PATH (Environment Variable)}\label{docs__pandoc__latex__src__backmatter_complete.md__path-environment-variable}

Tells system where executables are located.

\subsubsection*{View PATH (PowerShell)}\label{docs__pandoc__latex__src__backmatter_complete.md__view-path-powershell}

\begin{lstlisting}[style=Alabaster, language=powershell]
$env:PATH

\end{lstlisting}

\subsubsection*{View PATH (cmd)}\label{docs__pandoc__latex__src__backmatter_complete.md__view-path-cmd}

\begin{lstlisting}[style=Alabaster, language=cmd]
echo %PATH%

\end{lstlisting}

\begin{center}\rule{0.5\linewidth}{0.5pt}\end{center}

\section*{Environment Variable}\label{docs__pandoc__latex__src__backmatter_complete.md__environment-variable}

A system-level configuration variable.

\subsubsection*{Set Temporarily (PowerShell)}\label{docs__pandoc__latex__src__backmatter_complete.md__set-temporarily-powershell}

\begin{lstlisting}[style=Alabaster, language=powershell]
$env:TESTVAR="hello"

\end{lstlisting}

\subsubsection*{Set Temporarily (cmd)}\label{docs__pandoc__latex__src__backmatter_complete.md__set-temporarily-cmd}

\begin{lstlisting}[style=Alabaster, language=cmd]
set TESTVAR=hello

\end{lstlisting}

\begin{center}\rule{0.5\linewidth}{0.5pt}\end{center}

\section*{Exit Code}\label{docs__pandoc__latex__src__backmatter_complete.md__exit-code}

Numeric status returned by command.

\begin{itemize}
\tightlist
\item
  0 = Success
\item
  Non-zero = Error
\end{itemize}

Example:

\begin{lstlisting}[style=Alabaster, language=powershell]
if ($LASTEXITCODE -ne 0) { Write-Host "Build Failed" }

\end{lstlisting}

\begin{lstlisting}[style=Alabaster, language=bash]
if [ $? -ne 0 ]; then echo "Build Failed"; fi

\end{lstlisting}

\begin{center}\rule{0.5\linewidth}{0.5pt}\end{center}

\section*{Pipe}\label{docs__pandoc__latex__src__backmatter_complete.md__pipe}

Pass output from one command to another.

\begin{lstlisting}[style=Alabaster, language=powershell]
Get-ChildItem | Select-String ".scad"

\end{lstlisting}

\begin{lstlisting}[style=Alabaster, language=bash]
ls | grep ".scad"

\end{lstlisting}

\begin{center}\rule{0.5\linewidth}{0.5pt}\end{center}

\section*{SECTION V --- Git \& Version Control}\label{docs__pandoc__latex__src__backmatter_complete.md__section-v--git--version-control}

\subsection*{Git}\label{docs__pandoc__latex__src__backmatter_complete.md__git}

Distributed version control system.

\begin{center}\rule{0.5\linewidth}{0.5pt}\end{center}

\subsection*{Repository}\label{docs__pandoc__latex__src__backmatter_complete.md__repository}

Tracked project folder.

\begin{lstlisting}[style=Alabaster, language=bash]
git init

\end{lstlisting}

\begin{center}\rule{0.5\linewidth}{0.5pt}\end{center}

\subsection*{Commit}\label{docs__pandoc__latex__src__backmatter_complete.md__commit}

Snapshot of changes.

\begin{lstlisting}[style=Alabaster, language=bash]
git add .
git commit -m "Initial model"

\end{lstlisting}

\begin{center}\rule{0.5\linewidth}{0.5pt}\end{center}

\subsection*{Branch}\label{docs__pandoc__latex__src__backmatter_complete.md__branch}

Parallel development path.

\begin{lstlisting}[style=Alabaster, language=bash]
git branch feature-peg
git checkout feature-peg

\end{lstlisting}

\begin{center}\rule{0.5\linewidth}{0.5pt}\end{center}

\subsection*{Clone}\label{docs__pandoc__latex__src__backmatter_complete.md__clone}

Download remote repository.

\begin{lstlisting}[style=Alabaster, language=bash]
git clone https://github.com/user/project.git

\end{lstlisting}

\begin{center}\rule{0.5\linewidth}{0.5pt}\end{center}

\subsection*{Status}\label{docs__pandoc__latex__src__backmatter_complete.md__status}

Check changes.

\begin{lstlisting}[style=Alabaster, language=bash]
git status

\end{lstlisting}

\begin{center}\rule{0.5\linewidth}{0.5pt}\end{center}

\section*{SECTION VI --- File Formats}\label{docs__pandoc__latex__src__backmatter_complete.md__section-vi--file-formats}

\subsection*{STL}\label{docs__pandoc__latex__src__backmatter_complete.md__stl}

Triangle mesh file used in 3D printing.

Generated by:

\begin{lstlisting}[style=Alabaster, language=bash]
openscad -o model.stl model.scad

\end{lstlisting}

\begin{center}\rule{0.5\linewidth}{0.5pt}\end{center}

\subsection*{Slicer Software}\label{docs__pandoc__latex__src__backmatter_complete.md__slicer-software}

Converts STL to G-code.

\subsubsection*{PrusaSlicer}\label{docs__pandoc__latex__src__backmatter_complete.md__prusaslicer}

Open-source slicer with advanced features.

\subsubsection*{Cura}\label{docs__pandoc__latex__src__backmatter_complete.md__cura}

User-friendly slicer by Ultimaker.

\subsubsection*{Bambu Studio}\label{docs__pandoc__latex__src__backmatter_complete.md__bambu-studio}

Slicer for Bambu Lab printers.

\begin{center}\rule{0.5\linewidth}{0.5pt}\end{center}

\subsection*{G-code}\label{docs__pandoc__latex__src__backmatter_complete.md__g-code}

Machine instructions generated by slicer.

Example fragment:

\begin{lstlisting}[style=Alabaster]
G1 X10 Y10 Z0.2 F1500

\end{lstlisting}

\begin{center}\rule{0.5\linewidth}{0.5pt}\end{center}

\section*{SECTION VII --- Automation Concepts}\label{docs__pandoc__latex__src__backmatter_complete.md__section-vii--automation-concepts}

\subsection*{Script}\label{docs__pandoc__latex__src__backmatter_complete.md__script}

File containing executable instructions.

Examples:

\begin{itemize}
\tightlist
\item
  \texttt{.ps1}
\item
  \texttt{.sh}
\item
  \texttt{.scad}
\end{itemize}

\begin{center}\rule{0.5\linewidth}{0.5pt}\end{center}

\subsection*{Debugging}\label{docs__pandoc__latex__src__backmatter_complete.md__debugging}

Process of finding and fixing errors in code.

Techniques:

\begin{itemize}
\tightlist
\item
  Use \texttt{echo()} statements
\item
  Use \texttt{\#} modifier to highlight geometry
\item
  Test modules independently
\item
  Check console for warnings
\end{itemize}

\begin{center}\rule{0.5\linewidth}{0.5pt}\end{center}

\subsection*{Reproducibility}\label{docs__pandoc__latex__src__backmatter_complete.md__reproducibility}

Ability to regenerate identical outputs from identical inputs.

\begin{center}\rule{0.5\linewidth}{0.5pt}\end{center}

\subsection*{Automation}\label{docs__pandoc__latex__src__backmatter_complete.md__automation}

Using scripts instead of manual steps.

Example build script (PowerShell):

\begin{lstlisting}[style=Alabaster, language=powershell]
openscad -o build/model.stl src/model.scad

\end{lstlisting}

\begin{center}\rule{0.5\linewidth}{0.5pt}\end{center}

\section*{SECTION VIII --- Accessibility \& Documentation}\label{docs__pandoc__latex__src__backmatter_complete.md__section-viii--accessibility--documentation}

\subsection*{Markdown}\label{docs__pandoc__latex__src__backmatter_complete.md__markdown}

Lightweight markup language.

Example:

\begin{lstlisting}[style=Alabaster]
# Heading
- Bullet

\end{lstlisting}

\begin{center}\rule{0.5\linewidth}{0.5pt}\end{center}

\subsection*{Semantic Structure}\label{docs__pandoc__latex__src__backmatter_complete.md__semantic-structure}

Proper heading hierarchy:

\begin{lstlisting}[style=Alabaster]
# Title
## Section
### Subsection

\end{lstlisting}

Important for screen readers.

\begin{center}\rule{0.5\linewidth}{0.5pt}\end{center}

\subsection*{Editor Setup}\label{docs__pandoc__latex__src__backmatter_complete.md__editor-setup}

\subsubsection*{VS Code}\label{docs__pandoc__latex__src__backmatter_complete.md__vs-code}

Recommended extensions:

\begin{itemize}
\tightlist
\item
  OpenSCAD Language Support
\item
  Bracket Pair Colorizer
\end{itemize}

\begin{center}\rule{0.5\linewidth}{0.5pt}\end{center}

\subsubsection*{Notepad++}\label{docs__pandoc__latex__src__backmatter_complete.md__notepad}

Lightweight editor with syntax highlighting.

Configure:

\begin{itemize}
\tightlist
\item
  Settings \textgreater{} Preferences \textgreater{} Language
\item
  Select OpenSCAD or C-style syntax
\end{itemize}

\begin{center}\rule{0.5\linewidth}{0.5pt}\end{center}

\subsection*{Indentation}\label{docs__pandoc__latex__src__backmatter_complete.md__indentation}

Proper code formatting for readability.

\begin{lstlisting}[style=Alabaster, language=openscad]
module example() {
    if (condition) {
        cube([10,10,10]);
    }
}

\end{lstlisting}

Use consistent spacing (2 or 4 spaces per level).

\begin{center}\rule{0.5\linewidth}{0.5pt}\end{center}

\subsection*{Screen Reader}\label{docs__pandoc__latex__src__backmatter_complete.md__screen-reader}

Software that reads text aloud.

Examples:

\begin{itemize}
\tightlist
\item
  NVDA
\item
  JAWS
\item
  Orca
\item
  VoiceOver
\end{itemize}

\begin{center}\rule{0.5\linewidth}{0.5pt}\end{center}

\section*{SECTION IX --- Mechanical Design Concepts}\label{docs__pandoc__latex__src__backmatter_complete.md__section-ix--mechanical-design-concepts}

\subsection*{Clearance Fit}\label{docs__pandoc__latex__src__backmatter_complete.md__clearance-fit}

Loose fit allowing movement.

\subsection*{Press Fit}\label{docs__pandoc__latex__src__backmatter_complete.md__press-fit}

Tight fit requiring force.

\subsection*{Overhang}\label{docs__pandoc__latex__src__backmatter_complete.md__overhang}

Unsupported geometry exceeding safe angle (≈45° for FDM).

\begin{center}\rule{0.5\linewidth}{0.5pt}\end{center}

\subsection*{Tolerance}\label{docs__pandoc__latex__src__backmatter_complete.md__tolerance-1}

Intentional dimensional offset for mechanical fit.

Example:

\begin{itemize}
\tightlist
\item
  Designed hole: 10.2mm
\item
  Designed peg: 10.0mm
\item
  Clearance: 0.2mm
\end{itemize}

\begin{center}\rule{0.5\linewidth}{0.5pt}\end{center}

\subsection*{Calipers}\label{docs__pandoc__latex__src__backmatter_complete.md__calipers}

Precision measurement tool for verifying printed dimensions.

Types:

\begin{itemize}
\tightlist
\item
  Digital calipers (recommended)
\item
  Vernier calipers
\end{itemize}

Accuracy: ±0.01mm typical

\begin{center}\rule{0.5\linewidth}{0.5pt}\end{center}

\section*{SECTION X --- Advanced Workflow Concepts}\label{docs__pandoc__latex__src__backmatter_complete.md__section-x--advanced-workflow-concepts}

\subsection*{Headless CI Build}\label{docs__pandoc__latex__src__backmatter_complete.md__headless-ci-build}

Automated rendering using scripts or CI tools.

\begin{lstlisting}[style=Alabaster, language=bash]
openscad -o output.stl model.scad

\end{lstlisting}

\begin{center}\rule{0.5\linewidth}{0.5pt}\end{center}

\subsection*{Deterministic Modeling}\label{docs__pandoc__latex__src__backmatter_complete.md__deterministic-modeling}

Same input -\textgreater{} same output every time.

OpenSCAD is deterministic.

\begin{center}\rule{0.5\linewidth}{0.5pt}\end{center}

\subsection*{Dependency}\label{docs__pandoc__latex__src__backmatter_complete.md__dependency}

External program required for workflow.

Example:

\begin{itemize}
\tightlist
\item
  OpenSCAD must be in PATH.
\end{itemize}

\begin{center}\rule{0.5\linewidth}{0.5pt}\end{center}

\begin{center}\rule{0.5\linewidth}{0.5pt}\end{center}

\protect\phantomsection\label{docs__pandoc__latex__src__backmatter_complete.md__teacher-glossary}{}

\section*{Teacher Glossary --- Pedagogical Concepts for 3D Design and Computational Thinking Instruction}\label{docs__pandoc__latex__src__backmatter_complete.md__teacher-glossary--pedagogical-concepts-for-3d-design-and-computational-thinking-instruction}

\begin{center}\rule{0.5\linewidth}{0.5pt}\end{center}

\emph{Pedagogical Concepts for 3D Design \& Computational Thinking Instruction}

This glossary provides instructors with:

\begin{itemize}
\tightlist
\item
  Pedagogical frameworks
\item
  Teaching strategies
\item
  Assessment concepts
\item
  Accessibility guidance
\item
  Curriculum design principles
\end{itemize}

\begin{center}\rule{0.5\linewidth}{0.5pt}\end{center}

\section*{SECTION I --- Computational Thinking Framework}\label{docs__pandoc__latex__src__backmatter_complete.md__section-i--computational-thinking-framework}

\subsection*{Computational Thinking}\label{docs__pandoc__latex__src__backmatter_complete.md__computational-thinking}

Problem-solving approach using concepts fundamental to computer science.

Core pillars:

\begin{itemize}
\tightlist
\item
  \textbf{Decomposition} --- Breaking problems into smaller parts
\item
  \textbf{Pattern Recognition} --- Identifying similarities
\item
  \textbf{Abstraction} --- Focusing on relevant details
\item
  \textbf{Algorithms} --- Creating step-by-step solutions
\end{itemize}

Application in OpenSCAD:
Students decompose complex models into primitive shapes, recognize geometric patterns, abstract parameters, and write algorithmic module definitions.

\begin{center}\rule{0.5\linewidth}{0.5pt}\end{center}

\subsection*{Decomposition}\label{docs__pandoc__latex__src__backmatter_complete.md__decomposition}

Breaking a complex problem into manageable sub-problems.

\subsubsection*{In 3D Design Context}\label{docs__pandoc__latex__src__backmatter_complete.md__in-3d-design-context}

Example: A phone stand is decomposed into:

\begin{enumerate}
\tightlist
\item
  Base plate
\item
  Support angle
\item
  Lip to hold phone
\item
  Optional cable slot
\end{enumerate}

Each component becomes a separate module.

\begin{center}\rule{0.5\linewidth}{0.5pt}\end{center}

\subsection*{Abstraction}\label{docs__pandoc__latex__src__backmatter_complete.md__abstraction}

Removing unnecessary details to focus on essential features.

\subsubsection*{In Parametric Design}\label{docs__pandoc__latex__src__backmatter_complete.md__in-parametric-design}

\begin{lstlisting}[style=Alabaster, language=openscad]
module bolt(diameter, length) {
    // Abstract "bolt" concept
    // Hide implementation details
}

\end{lstlisting}

Students learn to identify which details matter and which can be parameterized.

\begin{center}\rule{0.5\linewidth}{0.5pt}\end{center}

\subsection*{Pattern Recognition}\label{docs__pandoc__latex__src__backmatter_complete.md__pattern-recognition}

Identifying repeating structures or concepts.

\subsubsection*{Examples in OpenSCAD}\label{docs__pandoc__latex__src__backmatter_complete.md__examples-in-openscad}

\begin{itemize}
\tightlist
\item
  Repeated geometric elements (array of holes)
\item
  Similar transformation sequences
\item
  Common parameter relationships
\end{itemize}

\begin{center}\rule{0.5\linewidth}{0.5pt}\end{center}

\subsection*{Algorithms}\label{docs__pandoc__latex__src__backmatter_complete.md__algorithms}

Step-by-step procedures for solving problems.

\subsubsection*{In Code}\label{docs__pandoc__latex__src__backmatter_complete.md__in-code}

\begin{lstlisting}[style=Alabaster, language=openscad]
// Algorithm for creating gear teeth
for (i = [0:teeth-1]) {
    rotate([0, 0, i * 360/teeth])
        translate([radius, 0, 0])
            tooth_profile();
}

\end{lstlisting}

\begin{center}\rule{0.5\linewidth}{0.5pt}\end{center}

\section*{SECTION II --- Instructional Design Strategies}\label{docs__pandoc__latex__src__backmatter_complete.md__section-ii--instructional-design-strategies}

\subsection*{Walking Skeleton Approach}\label{docs__pandoc__latex__src__backmatter_complete.md__walking-skeleton-approach}

Agile development practice: build minimal working version first, then iterate.

\subsubsection*{In 3D Design Curriculum}\label{docs__pandoc__latex__src__backmatter_complete.md__in-3d-design-curriculum}

\begin{enumerate}
\tightlist
\item
  Create basic cube
\item
  Add one parameter
\item
  Add one transformation
\item
  Build and verify
\item
  Repeat
\end{enumerate}

Prevents overwhelming students with complexity.

\begin{center}\rule{0.5\linewidth}{0.5pt}\end{center}

\subsection*{Design Cycle}\label{docs__pandoc__latex__src__backmatter_complete.md__design-cycle}

Iterative process:

\begin{lstlisting}[style=Alabaster]
Define -> Ideate -> Prototype -> Test -> Refine -> Repeat

\end{lstlisting}

\subsubsection*{In 3dMake Context}\label{docs__pandoc__latex__src__backmatter_complete.md__in-3dmake-context}

\begin{lstlisting}[style=Alabaster]
Edit .scad -> 3dm build -> 3dm info -> Measure print -> Adjust parameters

\end{lstlisting}

\begin{center}\rule{0.5\linewidth}{0.5pt}\end{center}

\subsection*{Scaffolding}\label{docs__pandoc__latex__src__backmatter_complete.md__scaffolding}

Temporary support structures that enable learning.

Examples:

\begin{itemize}
\tightlist
\item
  Starter code templates
\item
  Pre-defined modules
\item
  Graduated exercises
\item
  Checklists
\end{itemize}

Remove scaffolding as competency increases.

\begin{center}\rule{0.5\linewidth}{0.5pt}\end{center}

\subsection*{Physical Computing}\label{docs__pandoc__latex__src__backmatter_complete.md__physical-computing}

Learning that connects digital design to physical artifacts.

3D printing provides immediate tangible feedback:

\begin{itemize}
\tightlist
\item
  Code -\textgreater{} STL -\textgreater{} Physical object
\item
  Reinforces cause-and-effect
\item
  Engages kinesthetic learners
\end{itemize}

\begin{center}\rule{0.5\linewidth}{0.5pt}\end{center}

\subsection*{Stakeholder-Centric Design}\label{docs__pandoc__latex__src__backmatter_complete.md__stakeholder-centric-design}

Designing with end-user needs in focus.

\subsubsection*{Teaching Application}\label{docs__pandoc__latex__src__backmatter_complete.md__teaching-application}

Have students:

\begin{enumerate}
\tightlist
\item
  Interview stakeholders
\item
  Define requirements
\item
  Create design constraints
\item
  Iterate based on feedback
\end{enumerate}

Example: Design assistive device for specific user needs.

\begin{center}\rule{0.5\linewidth}{0.5pt}\end{center}

\section*{SECTION III --- Assessment \& Validation}\label{docs__pandoc__latex__src__backmatter_complete.md__section-iii--assessment--validation}

\subsection*{Code Review Criteria}\label{docs__pandoc__latex__src__backmatter_complete.md__code-review-criteria}

Assess student code on:

\begin{enumerate}
\tightlist
\item
  \textbf{Functionality} --- Does it work?
\item
  \textbf{Parametrization} --- Are values abstracted?
\item
  \textbf{Documentation} --- Are comments clear?
\item
  \textbf{Modularity} --- Is code reusable?
\item
  \textbf{Efficiency} --- Is approach optimal?
\end{enumerate}

\begin{center}\rule{0.5\linewidth}{0.5pt}\end{center}

\subsection*{Quality Assurance}\label{docs__pandoc__latex__src__backmatter_complete.md__quality-assurance}

Systematic verification process.

\subsubsection*{For 3D Prints}\label{docs__pandoc__latex__src__backmatter_complete.md__for-3d-prints}

\begin{itemize}
\tightlist
\item
  Measure with calipers
\item
  Compare to specification
\item
  Document deviations
\item
  Adjust parameters
\item
  Retest
\end{itemize}

Teaches scientific method and iterative refinement.

\begin{center}\rule{0.5\linewidth}{0.5pt}\end{center}

\subsection*{Safety Protocols}\label{docs__pandoc__latex__src__backmatter_complete.md__safety-protocols}

\subsubsection*{3D Printing Safety}\label{docs__pandoc__latex__src__backmatter_complete.md__3d-printing-safety}

Required instruction:

\begin{itemize}
\tightlist
\item
  Ventilation requirements
\item
  Hot surface warnings
\item
  Electrical safety
\item
  Material handling (MSDS)
\item
  Emergency shutdown procedures
\end{itemize}

\begin{center}\rule{0.5\linewidth}{0.5pt}\end{center}

\section*{SECTION IV --- Accessibility \& Universal Design}\label{docs__pandoc__latex__src__backmatter_complete.md__section-iv--accessibility--universal-design}

\subsection*{Screen Reader Accessibility Guide}\label{docs__pandoc__latex__src__backmatter_complete.md__screen-reader-accessibility-guide}

Ensure all students can access curriculum.

\subsubsection*{Key Principles}\label{docs__pandoc__latex__src__backmatter_complete.md__key-principles}

\begin{enumerate}
\tightlist
\item
  \textbf{Semantic HTML/Markdown} --- Proper heading structure
\item
  \textbf{Alt text for images} --- Describe geometric concepts
\item
  \textbf{CLI-first workflows} --- Reduce GUI dependence
\item
  \textbf{Descriptive output} --- \texttt{3dm\ info} provides text descriptions
\end{enumerate}

\begin{center}\rule{0.5\linewidth}{0.5pt}\end{center}

\subsection*{Braille Display Setup}\label{docs__pandoc__latex__src__backmatter_complete.md__braille-display-setup}

For students using refreshable braille displays:

\begin{enumerate}
\tightlist
\item
  Ensure code editors support screen readers
\item
  Use semantic indentation (not tabs)
\item
  Provide geometric descriptions in text
\item
  Test commands with NVDA/JAWS
\end{enumerate}

\begin{center}\rule{0.5\linewidth}{0.5pt}\end{center}

\subsection*{Indent Announcement Configuration}\label{docs__pandoc__latex__src__backmatter_complete.md__indent-announcement-configuration}

Configure screen readers to verbalize indentation level.

\subsubsection*{NVDA}\label{docs__pandoc__latex__src__backmatter_complete.md__nvda}

Settings \textgreater{} Document Formatting \textgreater{} Report line indentation

\subsubsection*{JAWS}\label{docs__pandoc__latex__src__backmatter_complete.md__jaws}

Settings Center \textgreater{} Reading \textgreater{} Report Indentation

Critical for understanding OpenSCAD block structure.

\begin{center}\rule{0.5\linewidth}{0.5pt}\end{center}

\subsection*{Non-Visual Toolchain Introduction}\label{docs__pandoc__latex__src__backmatter_complete.md__non-visual-toolchain-introduction}

Guide blind/low-vision students through:

\begin{enumerate}
\tightlist
\item
  CLI navigation (\texttt{pwd}, \texttt{ls}, \texttt{cd})
\item
  Text editor setup (VS Code with extensions)
\item
  \texttt{3dm\ build} command
\item
  \texttt{3dm\ info} for validation
\item
  Caliper measurement techniques (tactile)
\end{enumerate}

\begin{center}\rule{0.5\linewidth}{0.5pt}\end{center}

\subsection*{Editor Setup for Accessibility}\label{docs__pandoc__latex__src__backmatter_complete.md__editor-setup-for-accessibility}

\subsubsection*{VS Code}\label{docs__pandoc__latex__src__backmatter_complete.md__vs-code-1}

Recommended extensions:

\begin{itemize}
\tightlist
\item
  Screen Reader Optimized
\item
  Bracket Pair Colorizer (announces nesting)
\item
  OpenSCAD Language Support
\end{itemize}

\subsubsection*{Notepad++}\label{docs__pandoc__latex__src__backmatter_complete.md__notepad-1}

Lightweight alternative:

\begin{itemize}
\tightlist
\item
  Configure indentation announcement
\item
  Enable line number reading
\item
  Set up hotkeys for navigation
\end{itemize}

\begin{center}\rule{0.5\linewidth}{0.5pt}\end{center}

\section*{SECTION V --- Technical Teaching Concepts}\label{docs__pandoc__latex__src__backmatter_complete.md__section-v--technical-teaching-concepts}

\subsection*{Parametric Architecture}\label{docs__pandoc__latex__src__backmatter_complete.md__parametric-architecture}

Design philosophy where dimensions are variables, not fixed values.

\subsubsection*{Why Teach This}\label{docs__pandoc__latex__src__backmatter_complete.md__why-teach-this}

\begin{itemize}
\tightlist
\item
  Promotes computational thinking
\item
  Enables design reuse
\item
  Teaches variable scope
\item
  Reinforces algebra concepts
\end{itemize}

\begin{center}\rule{0.5\linewidth}{0.5pt}\end{center}

\subsection*{Geometric Primitives}\label{docs__pandoc__latex__src__backmatter_complete.md__geometric-primitives}

Basic 3D shapes that combine to form complex objects.

\subsubsection*{Teaching Sequence}\label{docs__pandoc__latex__src__backmatter_complete.md__teaching-sequence}

\begin{enumerate}
\tightlist
\item
  Cube
\item
  Sphere
\item
  Cylinder (including cones)
\item
  Boolean operations
\item
  Transformations
\item
  Compound models
\end{enumerate}

\begin{center}\rule{0.5\linewidth}{0.5pt}\end{center}

\subsection*{Tolerance Management}\label{docs__pandoc__latex__src__backmatter_complete.md__tolerance-management}

Critical mechanical engineering concept.

\subsubsection*{Teaching Approach}\label{docs__pandoc__latex__src__backmatter_complete.md__teaching-approach}

\begin{enumerate}
\tightlist
\item
  Print test parts
\item
  Measure with calipers
\item
  Calculate error
\item
  Adjust tolerance parameters
\item
  Reprint and verify
\end{enumerate}

Teaches:

\begin{itemize}
\tightlist
\item
  Precision vs accuracy
\item
  Measurement techniques
\item
  Iterative refinement
\end{itemize}

\begin{center}\rule{0.5\linewidth}{0.5pt}\end{center}

\subsection*{Interlocking Features}\label{docs__pandoc__latex__src__backmatter_complete.md__interlocking-features}

Mechanical joints and assemblies.

Examples:

\begin{itemize}
\tightlist
\item
  Snap fits
\item
  Press fits
\item
  Sliding joints
\item
  Threaded connections
\end{itemize}

Requires understanding of:

\begin{itemize}
\tightlist
\item
  Clearance tolerances
\item
  Material properties
\item
  Force analysis
\end{itemize}

\begin{center}\rule{0.5\linewidth}{0.5pt}\end{center}

\section*{SECTION VI --- Curriculum Design Elements}\label{docs__pandoc__latex__src__backmatter_complete.md__section-vi--curriculum-design-elements}

\subsection*{Design Mode vs Print Mode}\label{docs__pandoc__latex__src__backmatter_complete.md__design-mode-vs-print-mode}

OpenSCAD workflows:

\begin{itemize}
\tightlist
\item
  \textbf{Design Mode} (F5) --- Fast preview
\item
  \textbf{Print Mode} (F6) --- Full render for STL export
\end{itemize}

Teach distinction early to avoid confusion.

\begin{center}\rule{0.5\linewidth}{0.5pt}\end{center}

\subsection*{Editor Window vs Preview Window}\label{docs__pandoc__latex__src__backmatter_complete.md__editor-window-vs-preview-window}

OpenSCAD interface:

\begin{itemize}
\tightlist
\item
  Left: Code editor
\item
  Right: 3D preview
\end{itemize}

Non-visual students use editor + CLI tools instead of preview.

\begin{center}\rule{0.5\linewidth}{0.5pt}\end{center}

\subsection*{File Format Progression}\label{docs__pandoc__latex__src__backmatter_complete.md__file-format-progression}

Teach file types in order:

\begin{enumerate}
\tightlist
\item
  \texttt{.scad} --- Source code
\item
  \texttt{.stl} --- 3D mesh
\item
  \texttt{.gcode} --- Machine instructions
\end{enumerate}

Each represents different abstraction level.

\begin{center}\rule{0.5\linewidth}{0.5pt}\end{center}

\subsection*{Debugging Strategies}\label{docs__pandoc__latex__src__backmatter_complete.md__debugging-strategies}

Systematic approaches to finding errors.

\subsubsection*{For Students}\label{docs__pandoc__latex__src__backmatter_complete.md__for-students}

\begin{enumerate}
\tightlist
\item
  Read error messages carefully
\item
  Use \texttt{echo()} to print values
\item
  Test modules independently
\item
  Use \texttt{\#} modifier to highlight geometry
\item
  Simplify to minimal failing case
\end{enumerate}

\begin{center}\rule{0.5\linewidth}{0.5pt}\end{center}

\section*{SECTION VII --- Cross-Curricular Connections}\label{docs__pandoc__latex__src__backmatter_complete.md__section-vii--cross-curricular-connections}

\subsection*{Mathematics Integration}\label{docs__pandoc__latex__src__backmatter_complete.md__mathematics-integration}

3D design reinforces:

\begin{itemize}
\tightlist
\item
  Coordinate geometry (X, Y, Z axes)
\item
  Trigonometry (rotation angles)
\item
  Algebra (parametric equations)
\item
  Measurement \& units
\end{itemize}

\begin{center}\rule{0.5\linewidth}{0.5pt}\end{center}

\subsection*{Physics Applications}\label{docs__pandoc__latex__src__backmatter_complete.md__physics-applications}

\begin{itemize}
\tightlist
\item
  Center of mass calculations
\item
  Material properties (stress, strain)
\item
  Thermal expansion
\item
  Friction coefficients
\end{itemize}

\begin{center}\rule{0.5\linewidth}{0.5pt}\end{center}

\subsection*{Engineering Design Process}\label{docs__pandoc__latex__src__backmatter_complete.md__engineering-design-process}

Professional workflow:

\begin{enumerate}
\tightlist
\item
  Define problem
\item
  Research requirements
\item
  Brainstorm solutions
\item
  Build prototype
\item
  Test and evaluate
\item
  Refine design
\end{enumerate}

3D printing makes this tangible.

\begin{center}\rule{0.5\linewidth}{0.5pt}\end{center}

\section*{SECTION VIII --- Slicing \& Print Management}\label{docs__pandoc__latex__src__backmatter_complete.md__section-viii--slicing--print-management}

\subsection*{Slicing Guides for Instruction}\label{docs__pandoc__latex__src__backmatter_complete.md__slicing-guides-for-instruction}

\subsubsection*{PrusaSlicer}\label{docs__pandoc__latex__src__backmatter_complete.md__prusaslicer-1}

Recommended teaching slicer:

\begin{itemize}
\tightlist
\item
  Visual layer preview
\item
  Built-in supports
\item
  Print time estimates
\end{itemize}

\subsubsection*{Cura}\label{docs__pandoc__latex__src__backmatter_complete.md__cura-1}

Alternative:

\begin{itemize}
\tightlist
\item
  Simpler interface
\item
  Good for beginners
\end{itemize}

\subsubsection*{Bambu Studio}\label{docs__pandoc__latex__src__backmatter_complete.md__bambu-studio-1}

For Bambu Lab printers:

\begin{itemize}
\tightlist
\item
  Automated features
\item
  Multi-color support
\end{itemize}

\begin{center}\rule{0.5\linewidth}{0.5pt}\end{center}

\subsection*{Material Selection Pedagogy}\label{docs__pandoc__latex__src__backmatter_complete.md__material-selection-pedagogy}

\subsubsection*{PLA}\label{docs__pandoc__latex__src__backmatter_complete.md__pla}

Start here:

\begin{itemize}
\tightlist
\item
  Easiest to print
\item
  Lowest temperature
\item
  Minimal warping
\end{itemize}

\subsubsection*{PETG}\label{docs__pandoc__latex__src__backmatter_complete.md__petg-1}

Intermediate:

\begin{itemize}
\tightlist
\item
  Stronger
\item
  More forgiving than ABS
\end{itemize}

\subsubsection*{TPU}\label{docs__pandoc__latex__src__backmatter_complete.md__tpu}

Advanced:

\begin{itemize}
\tightlist
\item
  Flexible prints
\item
  Requires special techniques
\end{itemize}

\subsubsection*{ABS}\label{docs__pandoc__latex__src__backmatter_complete.md__abs}

Advanced:

\begin{itemize}
\tightlist
\item
  Professional applications
\item
  Requires enclosure
\item
  Warping challenges teach troubleshooting
\end{itemize}

\begin{center}\rule{0.5\linewidth}{0.5pt}\end{center}

\section*{SECTION IX --- Language \& Syntax Teaching}\label{docs__pandoc__latex__src__backmatter_complete.md__section-ix--language--syntax-teaching}

\subsection*{Code Statement}\label{docs__pandoc__latex__src__backmatter_complete.md__code-statement}

Single executable instruction.

\begin{lstlisting}[style=Alabaster, language=openscad]
cube([10,10,10]);  // One statement

\end{lstlisting}

\begin{center}\rule{0.5\linewidth}{0.5pt}\end{center}

\subsection*{Comments as Documentation}\label{docs__pandoc__latex__src__backmatter_complete.md__comments-as-documentation}

Teach commenting best practices:

\begin{enumerate}
\tightlist
\item
  Explain \textbf{why}, not \textbf{what}
\item
  Document parameters
\item
  Include units
\item
  Note constraints
\end{enumerate}

\begin{center}\rule{0.5\linewidth}{0.5pt}\end{center}

\subsection*{Arithmetic in Code}\label{docs__pandoc__latex__src__backmatter_complete.md__arithmetic-in-code}

Reinforce math concepts:

\begin{lstlisting}[style=Alabaster, language=openscad]
width = 10;
spacing = width * 1.5;  // Arithmetic

\end{lstlisting}

\begin{center}\rule{0.5\linewidth}{0.5pt}\end{center}

\section*{SECTION X --- Assessment Rubrics}\label{docs__pandoc__latex__src__backmatter_complete.md__section-x--assessment-rubrics}

\subsection*{Parametric Design Rubric}\label{docs__pandoc__latex__src__backmatter_complete.md__parametric-design-rubric}

{\def\LTcaptype{none} % do not increment counter
\begin{longtable}[]{@{}
  >{\raggedright\arraybackslash}p{(\linewidth - 8\tabcolsep) * \real{0.2075}}
  >{\raggedright\arraybackslash}p{(\linewidth - 8\tabcolsep) * \real{0.1509}}
  >{\raggedright\arraybackslash}p{(\linewidth - 8\tabcolsep) * \real{0.2264}}
  >{\raggedright\arraybackslash}p{(\linewidth - 8\tabcolsep) * \real{0.2264}}
  >{\raggedright\arraybackslash}p{(\linewidth - 8\tabcolsep) * \real{0.1887}}@{}}
\toprule\noalign{}
\begin{minipage}[b]{\linewidth}\raggedright
Criterion
\end{minipage} & \begin{minipage}[b]{\linewidth}\raggedright
Novice
\end{minipage} & \begin{minipage}[b]{\linewidth}\raggedright
Developing
\end{minipage} & \begin{minipage}[b]{\linewidth}\raggedright
Proficient
\end{minipage} & \begin{minipage}[b]{\linewidth}\raggedright
Advanced
\end{minipage} \\
\midrule\noalign{}
\endhead
\bottomrule\noalign{}
\endlastfoot
Parameterization & Hard-coded values & Some parameters &
Most values parameterized & Full parametric architecture \\
Documentation & No comments & Minimal comments & Clear parameter docs &
Comprehensive documentation \\
Modularity & No modules & Basic modules & Reusable modules &
Library-quality modules \\
Functionality & Doesn\textquotesingle t build & Builds with errors &
Meets requirements & Exceeds requirements \\
\end{longtable}
}

\begin{center}\rule{0.5\linewidth}{0.5pt}\end{center}

\subsection*{Manufacturing Readiness Rubric}\label{docs__pandoc__latex__src__backmatter_complete.md__manufacturing-readiness-rubric}

{\def\LTcaptype{none} % do not increment counter
\begin{longtable}[]{@{}ll@{}}
\toprule\noalign{}
Criterion & Assessment \\
\midrule\noalign{}
\endhead
\bottomrule\noalign{}
\endlastfoot
Manifold geometry & Pass/Fail \\
Tolerance specified & Pass/Fail \\
Documentation complete & Pass/Fail \\
Quality measurements recorded & Pass/Fail \\
\end{longtable}
}

\begin{center}\rule{0.5\linewidth}{0.5pt}\end{center}

\section*{SECTION XI --- Troubleshooting Pedagogy}\label{docs__pandoc__latex__src__backmatter_complete.md__section-xi--troubleshooting-pedagogy}

\subsection*{Teaching Debugging}\label{docs__pandoc__latex__src__backmatter_complete.md__teaching-debugging}

Framework:

\begin{enumerate}
\tightlist
\item
  \textbf{Reproduce} --- Can you make it fail consistently?
\item
  \textbf{Isolate} --- What\textquotesingle s the minimal failing case?
\item
  \textbf{Hypothesize} --- What might cause this?
\item
  \textbf{Test} --- How can you verify?
\item
  \textbf{Document} --- What did you learn?
\end{enumerate}

\begin{center}\rule{0.5\linewidth}{0.5pt}\end{center}

\subsection*{Common Student Errors}\label{docs__pandoc__latex__src__backmatter_complete.md__common-student-errors}

\subsubsection*{Syntax Errors}\label{docs__pandoc__latex__src__backmatter_complete.md__syntax-errors}

\begin{itemize}
\tightlist
\item
  Missing semicolons
\item
  Unmatched brackets
\item
  Typos in keywords
\end{itemize}

\subsubsection*{Logical Errors}\label{docs__pandoc__latex__src__backmatter_complete.md__logical-errors}

\begin{itemize}
\tightlist
\item
  Wrong parameter order
\item
  Incorrect units
\item
  Transform order confusion
\end{itemize}

\subsubsection*{Practical Errors}\label{docs__pandoc__latex__src__backmatter_complete.md__practical-errors}

\begin{itemize}
\tightlist
\item
  Model too large for print bed
\item
  Overhangs without support
\item
  Wall thickness too thin
\end{itemize}

\begin{center}\rule{0.5\linewidth}{0.5pt}\end{center}

\section*{SECTION XII --- Project-Based Learning}\label{docs__pandoc__latex__src__backmatter_complete.md__section-xii--project-based-learning}

\subsection*{Scaffolded Project Sequence}\label{docs__pandoc__latex__src__backmatter_complete.md__scaffolded-project-sequence}

\begin{enumerate}
\tightlist
\item
  \textbf{Week 1-2:} Primitives and transformations
\item
  \textbf{Week 3-4:} Modules and parameters
\item
  \textbf{Week 5-6:} Boolean operations
\item
  \textbf{Week 7-8:} Multi-part assemblies
\item
  \textbf{Week 9-10:} Student-designed project
\end{enumerate}

\begin{center}\rule{0.5\linewidth}{0.5pt}\end{center}

\subsection*{Portfolio Assessment}\label{docs__pandoc__latex__src__backmatter_complete.md__portfolio-assessment}

Students maintain portfolio showing:

\begin{itemize}
\tightlist
\item
  Design evolution
\item
  Failed prototypes (with analysis)
\item
  Final products
\item
  Measurements and QA data
\item
  Reflection on process
\end{itemize}

\begin{center}\rule{0.5\linewidth}{0.5pt}\end{center}

\begin{center}\rule{0.5\linewidth}{0.5pt}\end{center}

\protect\phantomsection\label{docs__pandoc__latex__src__backmatter_complete.md__further-reading}{}

\section*{Further Reading and References}\label{docs__pandoc__latex__src__backmatter_complete.md__further-reading-and-references}

\begin{center}\rule{0.5\linewidth}{0.5pt}\end{center}

\subsection*{Peer-Reviewed and Scholarly Works}\label{docs__pandoc__latex__src__backmatter_complete.md__peer-reviewed-and-scholarly-works}

Gonzalez Avila, J. F., Pietrzak, T., Girouard, A., \& Casiez, G. (2024). \emph{Understanding the challenges of OpenSCAD users for 3D printing.} In \emph{Proceedings of the CHI Conference on Human Factors in Computing Systems (CHI '24)} (Article 351, pp. 1--20). Association for Computing Machinery. https://doi.org/10.1145/3613904.3642566 :contentReference{[}oaicite:0{]}\{index=0\}

Li, G. (2024). \emph{Chinese dragon modeling based on OpenSCAD.} \emph{Applied and Computational Engineering, 67}, 1--12. https://doi.org/10.54254/2755-2721/67/20240585 :contentReference{[}oaicite:1{]}\{index=1\}

Kwon, N., Sun, T., Gao, Y., Zhao, L., Wang, X., Kim, J., \& Hong, S. R. (2024). \emph{3DPFIX: Improving remote novices' 3D printing troubleshooting through human-AI collaboration.} arXiv. https://doi.org/10.48550/arXiv.2401.15877 :contentReference{[}oaicite:2{]}\{index=2\}

\subsection*{Government \& Standards}\label{docs__pandoc__latex__src__backmatter_complete.md__government--standards}

Centers for Disease Control and Prevention, National Institute for Occupational Safety and Health. (2020, May 14). \emph{Health and safety considerations for 3D printing.} https://blogs.cdc.gov/niosh-science-blog/2020/05/14/3d-printing/

Centers for Disease Control and Prevention, National Institute for Occupational Safety and Health. (2024). \emph{Approaches to safe 3D printing.} https://www.cdc.gov/niosh/blogs/2024/safe-3d-printing.html

National Institute of Standards and Technology. (2023). \emph{Artificial Intelligence Risk Management Framework (AI RMF 1.0).} https://www.nist.gov/system/files/documents/2023/01/26/AI\%20RMF\%201.0.pdf

Occupational Safety and Health Administration. (n.d.). \emph{Hierarchy of controls.} https://www.osha.gov/hierarchy-of-controls

Washington State Department of Health. (n.d.). \emph{3D printers in schools.} https://doh.wa.gov/community-and-environment/schools/3d-printers

\subsection*{Industry \& Commercial Documentation}\label{docs__pandoc__latex__src__backmatter_complete.md__industry--commercial-documentation}

All3DP. (n.d.). \emph{3D printer filament types explained.} https://all3dp.com/1/3d-printer-filament-types-3d-printing-3d-filament/

All3DP. (n.d.). \emph{FDM 3D printing tolerances.} https://all3dp.com/2/fdm-3d-printing-tolerances/

Anycubic. (n.d.). \emph{Anycubic.} https://www.anycubic.com/

Autodesk. (n.d.). \emph{Fusion 360.} https://www.autodesk.com/products/fusion-360

Bambu Lab. (n.d.). \emph{Downloads.} https://bambulab.com/en/download

Bambu Lab. (n.d.). \emph{Bambu Studio.} https://bambulab.com/en/download/studio

Bambu Lab Wiki. (n.d.). \emph{Bambu Lab Wiki.} https://wiki.bambulab.com

Hubs. (n.d.). \emph{What is FDM 3D printing?} https://www.hubs.com/knowledge-base/what-is-fdm-3d-printing/

IDEA Maker (Raise3D). (n.d.). \emph{ideaMaker.} https://www.raise3d.com/ideamaker

MatterHackers. (n.d.). \emph{3D printer filament comparison.} https://www.matterhackers.com/3d-printer-filament-compare

Prusa Research. (n.d.). \emph{PrusaSlicer 4.2.4.} https://www.prusa3d.com/page/prusaslicer424/

Prusa Research. (n.d.). \emph{PrusaSlicer 4.1.0.} https://www.prusa3d.com/page/prusaslicer\_410/

Prusa Research. (n.d.). \emph{Support.} https://www.prusa3d.com/support/

Stratasys. (n.d.). \emph{Stratasys.} https://stratasys.com/

UL Research Institutes. (n.d.). \emph{3D printing emissions study.} https://www.ul.com/news/ul-research-institutes-releases-3d-printing-emissions-study

Ultimaker. (n.d.). \emph{Ultimaker Cura.} https://ultimaker.com/software/ultimaker-cura

ZenML. (n.d.). \emph{LLM-powered 3D model generation for 3D printing.} https://www.zenml.io/llmops-database/llm-powered-3d-model-generation-for-3d-printing

\subsection*{Open-Source Software \& Technical Documentation}\label{docs__pandoc__latex__src__backmatter_complete.md__open-source-software--technical-documentation}

Anthropic. (n.d.). \emph{Prompt engineering overview.} https://docs.anthropic.com/en/docs/build-with-claude/prompt-engineering/overview

Apache Software Foundation. (2004). \emph{Apache License, Version 2.0.} https://www.apache.org/licenses/LICENSE-2.0

BelfrySCAD. (n.d.). \emph{BOSL2.} https://github.com/BelfrySCAD/BOSL2

Deck, T. (2025). \emph{3DMake repository.} https://github.com/tdeck/3dmake

Free Software Foundation. (n.d.). \emph{GNU Bash manual.} https://www.gnu.org/software/bash/manual/

Free Software Foundation. (n.d.). \emph{GNU Coreutils manual.} https://www.gnu.org/software/coreutils/manual/

Free Software Foundation. (n.d.). \emph{GNU Grep manual.} https://www.gnu.org/software/grep/manual/grep.html

Git SCM. (n.d.). \emph{Git documentation.} https://git-scm.com/docs

Google Cloud. (2025). \emph{Get started with Gemini models.} https://docs.cloud.google.com/vertex-ai/generative-ai/docs/start/get-started-with-gemini-3

Microsoft. (n.d.). \emph{PowerShell documentation.} https://learn.microsoft.com/powershell/

Microsoft. (n.d.). \emph{Windows commands documentation.} https://learn.microsoft.com/en-us/windows-server/administration/windows-commands/

OpenSCAD. (n.d.). \emph{OpenSCAD documentation.} https://openscad.org/documentation.html

OpenSCAD User Manual Contributors. (n.d.). \emph{OpenSCAD User Manual.} https://en.wikibooks.org/wiki/OpenSCAD\_User\_Manual

Programming with OpenSCAD. (n.d.). \emph{Programming with OpenSCAD.} https://programmingwithopenscad.github.io/

Rust Project Developers. (n.d.). \emph{mdBook.} https://rust-lang.github.io/mdBook/

SuperSlicer. (n.d.). \emph{SuperSlicer repository.} https://github.com/supermerill/SuperSlicer

The Linux Documentation Project. (n.d.). \emph{Bash Beginner's Guide.} https://tldp.org/LDP/Bash-Beginners-Guide/html/

\subsection*{Educational \& Community Resources}\label{docs__pandoc__latex__src__backmatter_complete.md__educational--community-resources}

A11Y-101. (n.d.). \emph{Inclusive documentation design.} https://www.a11y-101.com/design/inclusive-documentation

Class Central. (n.d.). \emph{OpenSCAD courses.} https://www.classcentral.com/subject/openscad

Khan Academy. (n.d.). \emph{Algorithms.} https://www.khanacademy.org/computing/computer-science/algorithms

Salt Lake City Public Library. (n.d.). \emph{Creative Lab.} https://services.slcpl.org/creativelab

Salt Lake County Library. (n.d.). \emph{Create spaces.} https://www.slcolibrary.org/what-we-have/create

University of Utah. (n.d.). \emph{ProtoSpace.} https://lib.utah.edu/protospace.php

\subsubsection*{Accessibility \& Assistive Tech}\label{docs__pandoc__latex__src__backmatter_complete.md__accessibility--assistive-tech}

Freedom Scientific. (n.d.). \emph{JAWS screen reader.} https://www.freedomscientific.com/products/software/jaws/

Microsoft. (n.d.). \emph{Narrator guide.} https://support.microsoft.com/narrator

NV Access. (n.d.). \emph{NVDA screen reader.} https://www.nvaccess.org/

Your Dolphin. (n.d.). \emph{SuperNova.} https://yourdolphin.com/supernova/

\subsubsection*{Community \& Media}\label{docs__pandoc__latex__src__backmatter_complete.md__community--media}

Discord. (n.d.). \emph{3D printing community server.} https://discord.gg/F2Nx2VxTB7

Reddit. (n.d.). \emph{r/3Dprinting subreddit.} https://www.reddit.com/r/3Dprinting/

Reddit. (n.d.). \emph{r/OpenSCAD subreddit.} https://www.reddit.com/r/openscad/

YouTube. (n.d.). \emph{Luke's Lab channel.} https://www.youtube.com/@LukesBlab

YouTube. (n.d.). \emph{Teaching Tech channel.} https://www.youtube.com/@TeachingTech

\section{Contributors to this Curriculum}\label{docs__pandoc__latex__src__contributors.md__contributors}

Michael Ryan Hunsaker, M.Ed., Ph.D. (he/him)

\subsection*{About Dr. Hunsaker}\label{docs__pandoc__latex__src__contributors.md__about-dr-hunsaker}

Michael Ryan Hunsaker holds a Ph.D. in neuroscience (behavioral/translational neuroscience focus) from UC Davis and a Master\textquotesingle s degree in Special Education from the University of Utah with a credential to teach blind and visually impaired students. His career uniquely bridges scientific research, clinical experience, and classroom practice-positioning him to design accessible educational materials grounded in both neuroscience and practical pedagogy.

\subsubsection*{Research \& Clinical Background}\label{docs__pandoc__latex__src__contributors.md__research--clinical-background}

Following his doctoral work at UC Davis (2012), Dr. Hunsaker completed postdoctoral research in primate neural development at the UC Davis MIND Institute, where he developed sophisticated research methods to evaluate developmental disorders. During this period, he gained extensive clinical experience working directly with children and adults across a broad spectrum of diagnoses, including autism spectrum disorder, cerebral palsy, genetic conditions (fragile X syndrome, Williams syndrome, 22q11.2DS, Down syndrome), ADHD, and Tourette\textquotesingle s syndrome. This diverse exposure to different neurological profiles taught him a crucial lesson: each learner\textquotesingle s disability uniquely shapes their cognitive strengths, learning preferences, and educational needs.

\subsubsection*{Transition to Education}\label{docs__pandoc__latex__src__contributors.md__transition-to-education}

In 2013, recognizing the gap between academic research and classroom reality, Dr. Hunsaker made a deliberate career shift to direct education. Rather than pursue traditional academic pathways, he committed to applying his scientific training to help students in real classroom environments. As a special education teacher, he focused on understanding the individual relationship between disability and learning-developing instructional approaches that enable students to access their full potential rather than working around perceived limitations.

\subsubsection*{Expertise in Accessibility \& Inclusive Design}\label{docs__pandoc__latex__src__contributors.md__expertise-in-accessibility--inclusive-design}

His specialization in teaching blind and visually impaired students, particularly those with additional disabilities, has given Dr. Hunsaker deep expertise in accessible pedagogy and non-visual problem-solving. He grounds all curriculum design in evidence-based practice informed by neuroscience, accessibility research, and years of direct classroom experience with diverse learners. He understands that accessibility is not an afterthought or accommodation; it is the foundation upon which effective education for all students is built.

\subsubsection*{Creator of Accessible Educational Resources}\label{docs__pandoc__latex__src__contributors.md__creator-of-accessible-educational-resources}

Motivated by observing persistent gaps in affordable, practical, and culturally responsive accessible educational materials, Dr. Hunsaker created TVI Resources, a collection of free, openly licensed instructional materials designed to meet the real needs of educators and students. This commitment to accessible education informed his development of this curriculum: a comprehensive, text-based approach to 3D design and digital fabrication specifically engineered for blind and visually impaired learners, emphasizing independence, technical mastery, and real-world application.

\subsection*{Acknowledgements}\label{docs__pandoc__latex__src__contributors.md__acknowledgements}

This curriculum stands on the shoulders of innovators, educators, and open-source maintainers whose commitment to accessibility and open education made this work possible.

\subsubsection*{3DMake}\label{docs__pandoc__latex__src__contributors.md__3dmake}

Special thanks to Troy Deck, the creator and maintainer of \href{https://github.com/tdeck/3dmake}{3DMake}, an open-source command-line tool that automates the workflow from code-based 3D models to print-ready files. 3DMake is essential to this curriculum-it eliminates the visual GUI entirely and makes parametric 3D design accessible to screen reader users through text-based commands and configuration files. Without 3DMake, a truly non-visual pathway from design to fabrication would not exist.

\subsubsection*{Texas School for the Blind and Visually Impaired (TSBVI)}\label{docs__pandoc__latex__src__contributors.md__texas-school-for-the-blind-and-visually-impaired-tsbvi}

This curriculum was directly inspired by the innovative work at TSBVI, a pioneer in accessible education for blind and visually impaired students. The foundational idea of using text-driven tools and command-line workflows for technical education originated from  collaboration with TSBVI educators. Their commitment to demonstrating that blind and visually impaired students can excel in technical fields provided the vision for this project.

\subsubsection*{OpenSCAD Community}\label{docs__pandoc__latex__src__contributors.md__openscad-community}

Deep appreciation to the OpenSCAD project maintainers and contributors for creating and maintaining a code-driven, parametric 3D modeling system that is inherently more accessible than traditional GUI-based CAD tools. OpenSCAD\textquotesingle s text-based programming model and screen-reader-compatible approach to geometry design makes it uniquely suited for learners who use assistive technology. The OpenSCAD community\textquotesingle s commitment to open-source principles and accessibility embodies the philosophy that drives this curriculum.

\subsubsection*{Open Education \& Accessibility Movement}\label{docs__pandoc__latex__src__contributors.md__open-education--accessibility-movement}

Finally, this curriculum is part of a broader movement toward open educational resources that prioritize accessibility from the ground up. It reflects the conviction that education should be barrier-free, that students with disabilities bring unique perspectives and strengths, and that accessible design benefits all learners. The work of accessibility advocates, neurodivergent educators, and disability justice activists informed every decision in building this material.

\backmatter

% Print any remaining endnotes
\ifnum\value{endnote}>0
  \begingroup
  \parindent 0pt
  \parskip 2ex
  \def\enotesize{\normalsize}
  \theendnotes
  \endgroup
\fi





\cleardoublepage
\phantomsection
\end{document}